\documentclass[11pt,oneside]{book}

\usepackage[margin=1in]{geometry}
\usepackage{amsmath,amssymb,amsthm,mathtools}
\usepackage{booktabs,longtable,array}
\usepackage{enumitem}
\usepackage[colorlinks=true,linkcolor=blue,urlcolor=blue,citecolor=blue]{hyperref}
\usepackage[nameinlink,noabbrev]{cleveref}

\newtheorem{theorem}{Theorem}[chapter]
\newtheorem{lemma}[theorem]{Lemma}
\newtheorem{proposition}[theorem]{Proposition}
\newtheorem{corollary}[theorem]{Corollary}
\newtheorem{assessment}[theorem]{Assessment}
\newtheorem{firewall}[theorem]{Claim firewall}
\newtheorem{openproblem}[theorem]{Open problem}
\theoremstyle{definition}
\newtheorem{definition}[theorem]{Definition}
\theoremstyle{remark}
\newtheorem{remark}[theorem]{Remark}
\newcommand{\MGclosed}{\textsc{closed}}
\newcommand{\MGopen}{\textsc{open}}
\newcommand{\MGpartial}{\textsc{partial}}
\newcommand{\opnorm}[1]{\lVert#1\rVert_{\mathrm{op}}}
\newcommand{\trnorm}[1]{\lVert#1\rVert_{1}}
\newcommand{\Cov}{\operatorname{Cov}}
\newcommand{\Tr}{\operatorname{Tr}}
\newcommand{\htop}{\boldsymbol 1}

\title{Renewal Yang--Mills Mass-Gap Research Notes\\
Compact Series, Version V1}
\author{Research continuation after legacy archive f1}
\date{September 2026}

\begin{document}
\maketitle
\tableofcontents

\chapter{Context, legacy checkpoint, and proof architecture}
\label{ch:V1-context}

\section{What this notebook is doing}
\label{sec:V1-purpose}

The long-form Yang--Mills notebook reached legacy version V135 and became
too large to navigate reliably.  That complete source is frozen as
\begin{center}
\texttt{Renewal\_Yang\_Mills\_Mass\_Gap\_Research\_Notes\_V135\_f1.tex}.
\end{center}
This compact series restarts at V1.  It does not restart the mathematics:
the results below summarize the dependency-bearing part of the frozen
archive, record the remaining gaps, correct one proof shortcut in its last
chapter, and continue the programme at the first unresolved normalized
quartic carrier.

The intended route is
\begin{equation}
\begin{split}
\text{local Wilson cumulants}
&\Longrightarrow \text{marked-tree atlas (MTA)}
\Longrightarrow \text{weighted uniform fusion control}\\
&\Longrightarrow \text{connected Perron/source estimates}
\Longrightarrow \text{finite-edge and capacity bounds}\\
&\Longrightarrow \text{reflection-positive continuum construction}
\Longrightarrow \text{positive physical spectral gap}.
\end{split}
\label{eq:V1-route}
\end{equation}
Only the first arrow is presently under active construction.  Every later
arrow remains conditional on the preceding ones.

\begin{assessment}[Versioning and audit protocol]
\label{ass:V1-protocol}
The active compact file is replaced at every iteration by the next version,
while the frozen \(\texttt{\_f1}\) archive is retained.  At V5, V10, and
every subsequent multiple of five, the newest Standard-Model--Einstein and
Navier--Stokes notes in \texttt{latex\_notes} are inspected read-only before
the next proof direction is chosen.  When the compact series reaches V100,
it is frozen with suffix \(\texttt{\_f2}\) and a new compact V1 is started;
the same rule is repeated thereafter.
\end{assessment}

\section{Frozen-source identity}
\label{sec:V1-archive-identity}

The frozen f1 archive has \(88517\) newline-delimited source records,
\(3083467\) bytes, and SHA--256 digest
\begin{equation}
\texttt{82f6bbffeecdad3c5c63551660b19a0bdc7f3d6dffe4c68b4e9bba3a783a2b71}.
\label{eq:V1-archive-hash}
\end{equation}
All legacy theorem identifiers mentioned below refer to that exact source.

\section{Major mathematical results inherited from f1}
\label{sec:V1-legacy-summary}

\subsection{Exact Fourier normalization and coefficient mass}

For irreducible input representations \(R_1,\dots,R_m\), write
\[
d_{\mathrm{in}}:=\prod_i d_{R_i},\qquad
\bigotimes_iH_{R_i}
=\bigoplus_U W_U(H_U\otimes M_U),
\]
and, for \(A:=\bigotimes_iA_i\), put \(X_U:=W_U^*AW_U\).
Legacy V46 proves the exact product coefficient identity
\begin{equation}
\widehat h(U)
=\frac{d_{\mathrm{in}}}{d_U}
\Tr_{M_U}\!\left[X_U(I_{H_U}\otimes K_U)\right].
\label{eq:V1-exact-Fourier}
\end{equation}
Consequently the output BFA dimension \(d_U\) cancels, and complete raw
isotypic pinching gives
\begin{equation}
\sum_U\trnorm{X_U}\le\prod_i\trnorm{A_i}.
\label{eq:V1-raw-pinching}
\end{equation}
This corrects the earlier, false identification of arbitrary Fourier
coefficient mass with a representation-dimension formation law.  The
dual-singlet regression has raw mass \(d_R^{-1}\), not dimension-law mass.

\subsection{Marked link probabilities and the spatial compiler}

For a nontrivial exact product representation \(\boldsymbol R\), let
\[
\Xi(\boldsymbol R):=\sum_{e}A_{R_e},\qquad
\vartheta_e(\boldsymbol R):=\frac{A_{R_e}}{\Xi(\boldsymbol R)}.
\]
The marks form an exact probability budget over active coordinates.
Legacy V31 proves the binary MTA estimate without support loss and proves
that an all-order MTA estimate compiles into the required rooted spatial
fusion estimate.  It also gives a sharp obstruction to any unmarked
support-counting proof: one large support can meet arbitrarily many
singleton supports while all unmarked rooted norms remain one.

The uniform ternary row was subsequently repaired using the exact V46
normalization and the coefficient-sensitive physical partitions developed
through V50.  Thus order two and order three are not the present
bottleneck.

\subsection{Quartic topology and local operator control}

Legacy V40 proves that every minimal connected four-row coordinate
skeleton belongs to exactly one of
\[
[4],\qquad[3,2],\qquad[3,3],\qquad[2,2,2].
\]
Legacy V41 proves a finite-\(s\), representation-uniform local
four-junction tree decomposition, including all sixteen labeled
four-vertex trees, and multiplicity-uniform final-channel compression.
Legacy V53 proves complete normalized quartic MTA for the exclusive-pair
\([2,2,2]\) sector.  Legacy V56 proves it for the noncollision
\([3,2]\) sector.

\subsection{Collision capacity and global numerator aggregation}

The long V57--V129 analysis constructs the physical outer partition,
keeps noncommuting collision orientations separate, and develops positive
product-state capacity estimates.  Its final result closes both
orientations of every completely regrouped fixed collision monomial.
It does not by itself sum all collision monomials after normalization.

Legacy V130 replaces the coordinate-subset sum by an exact
Boolean--Möbius selector on the finite quartic partition alphabet.  Legacy
V131 obtains an exact sixteen-tree representation.  V132 proposes a
Bernoulli partition interpolation, but V133 proves that its local mixed
derivatives can collapse several nominal bridges to one quotient
covariance; this invalidates the proposed product bridge estimate.

V134 therefore returns upstream to exact connected coordinate words and
proves a support-uniform quartic numerator in the low-activity branch.
V135 combines a weak-cut/strong-component threshold decomposition with
that branch and proves
\begin{equation}
\opnorm{\mathbf K^{\mathrm{conn}}_{\{1,2,3,4\}}}
\le C\sum_{\tau\in\mathfrak T_4}
\prod_{ij\in E(\tau)}\Lambda_{ij},
\qquad
\Lambda_{ij}=s_\alpha H_{ij}.
\label{eq:V1-legacy-numerator}
\end{equation}
This is an unnormalized operator numerator.  It contains neither the
repeated input probability marks nor the final resolvent quotient.

\section{Current rigorous status}
\label{sec:V1-status}

\begin{longtable}{>{\raggedright\arraybackslash}p{0.31\textwidth}
                  >{\raggedright\arraybackslash}p{0.14\textwidth}
                  >{\raggedright\arraybackslash}p{0.47\textwidth}}
\toprule
gate & status & exact present boundary\\
\midrule
\endfirsthead
\toprule
gate & status & exact present boundary\\
\midrule
\endhead
Binary and ternary normalized MTA
& \MGclosed
& Correct Fourier normalization, marked support sums, raw pinching, and
  one-resolvent estimates are proved.\\[0.35em]
Exclusive-pair quartic \([2,2,2]\)
& \MGclosed
& Legacy V53 proves the complete arbitrary-coefficient sector.\\[0.35em]
Noncollision quartic \([3,2]\)
& \MGclosed
& Legacy V56 proves the complete assembled sector.\\[0.35em]
Collision quartic \([3,2]\)
& \MGpartial
& Every fixed regrouped monomial is closed by V129; normalized global
  aggregation is not.\\[0.35em]
Embedded \([4]\) and \([3,3]\)
& \MGopen
& Their local numerators are controlled, but no complete
  coefficient-sensitive normalized spectator theorem is proved.\\[0.35em]
Global quartic and all-order MTA
& \MGopen
& The V135 numerator cannot be marked after taking its operator norm.\\[0.35em]
Continuum Yang--Mills and mass gap
& \MGopen
& WUFC/CPM, finite-edge margins, capacity, a nontrivial
  reflection-positive continuum state, and the spectral-gap theorem all
  remain downstream.\\
\bottomrule
\end{longtable}

\begin{firewall}[No mass-gap claim]
\label{fire:V1-no-gap}
Neither the frozen archive nor this compact V1 constructs four-dimensional
continuum Yang--Mills theory or proves a positive mass gap.  The programme
contains conditional compilers and several closed finite-regulator
estimates, not a solution of the Clay problem.
\end{firewall}

\chapter{A corrected weak-cut proof}
\label{ch:V1-cut-correction}

\section{Setup}

Let \(F_1,\dots,F_4\) be bounded row operator functions with
\(\opnorm{F_i}\le1\).  Rows act on distinct tensor factors, so row products
commute.  For disjoint nonempty row sets \(C,D\), set
\[
F_C:=\prod_{i\in C}F_i,\qquad
\Lambda(C,D):=\sum_{\substack{i\in C\\j\in D}}\Lambda_{ij}.
\]
We use the legacy V50 covariance estimate in its spectator-stable form:
\begin{equation}
\opnorm{\Cov(F_C,F_D)}\le C\Lambda(C,D).
\label{eq:V1-composite-covariance}
\end{equation}

\begin{lemma}[Exact ternary product-cut identity]
\label{lem:V1-ternary-cut}
For three distinct rows \(i,j,k\),
\begin{equation}
\Cov(F_i,F_jF_k)
=\kappa_3(F_i,F_j,F_k)
+\mathbb EF_j\,\Cov(F_i,F_k)
+\mathbb EF_k\,\Cov(F_i,F_j).
\label{eq:V1-ternary-cut-identity}
\end{equation}
Consequently, for every nontrivial cut \(E\mid F\) of
\(\{i,j,k\}\),
\begin{equation}
\boxed{
\opnorm{\kappa_3(F_i,F_j,F_k)}
\le C_3\min\{1,\Lambda(E,F)\}.}
\label{eq:V1-ternary-cut-bound}
\end{equation}
\end{lemma}

\begin{proof}
Expanding the right side of
\eqref{eq:V1-ternary-cut-identity} gives
\[
\mathbb E(F_iF_jF_k)-\mathbb EF_i\,\mathbb E(F_jF_k),
\]
which is the left side.  For \(E=\{i\}\), equation
\eqref{eq:V1-composite-covariance} bounds the left side by
\(C(\Lambda_{ij}+\Lambda_{ik})\); ordinary covariance bounds give the
same estimate for the two correction terms, and the means are
contractions.  This proves the linear bound.  The moment--cumulant formula
has finitely many terms, each a contraction, and therefore gives a
universal order-three cap.  The other cuts follow by relabeling.  Taking
the smaller bound proves the display.
\end{proof}

\begin{theorem}[Corrected complete quartic weak-cut estimate]
\label{thm:V1-quartic-weak-cut}
For every nontrivial cut \(C\mid D\) of four rows,
\begin{equation}
\boxed{
\opnorm{\kappa_4(F_1,F_2,F_3,F_4)}
\le C_4\min\{1,\Lambda(C,D)\}.}
\label{eq:V1-quartic-weak-cut}
\end{equation}
\end{theorem}

\begin{proof}
Let \(\rho=\{C,D\}\).  Expanding moments into cumulants yields the exact
product-cumulant identity
\begin{equation}
\Cov(F_C,F_D)
=
\sum_{\substack{\pi\in\Pi_4\\\pi\vee\rho=\htop}}
\prod_{B\in\pi}\kappa(F_i:i\in B).
\label{eq:V1-product-cumulant}
\end{equation}
The top partition contributes the desired fourth cumulant once.  Every
other retained partition contains a block \(B\) meeting both \(C\) and
\(D\).  If \(|B|=2\), its cumulant is a cross-cut covariance.  If
\(|B|=3\), apply \Cref{lem:V1-ternary-cut} to
\((B\cap C)\mid(B\cap D)\).  All other block cumulants and singleton
means have universal order-at-most-three caps.  If two pair blocks cross,
use the cross-link estimate on one and cap the other.  Hence every proper
correction term in \eqref{eq:V1-product-cumulant} is
\(O(\Lambda(C,D))\).  The left side has the same bound by
\eqref{eq:V1-composite-covariance}.  Solving for the fourth cumulant gives
the linear estimate.  Its finite moment--cumulant formula supplies a
universal cap, proving the minimum.
\end{proof}

\begin{remark}[Correction to legacy V135]
\label{rem:V1-V135-correction}
This proof replaces the single V135 sentence that tried to cap the second
edge in a ternary tree estimate.  That shortcut was not justified for an
untruncated link.  The theorem itself remains valid by the exact cut
identity above.
\end{remark}

\chapter{Coefficient-sensitive normalized output control}
\label{ch:V1-capacity}

\section{Product-state capacity is the correct coefficient gate}
\label{sec:V1-product-capacity}

Let
\(\mathcal H=\bigotimes_{i=1}^4\mathcal H_i\).
For \(Q\ge0\) on \(\mathcal H\), define its four-row product-state
capacity by
\begin{equation}
\kappa_{\otimes}(Q)
:=
\sup_{\substack{\rho_i\ge0\\\Tr\rho_i=1}}
\Tr\!\left[Q(\rho_1\otimes\rho_2\otimes\rho_3\otimes\rho_4)\right].
\label{eq:V1-product-capacity}
\end{equation}
This retains all correlations inside each row while forbidding the
artificial replacement of arbitrary coefficients by a representation
dimension law.

\begin{lemma}[Positive decomposition of a trace-class input]
\label{lem:V1-positive-decomposition}
Every finite-dimensional trace-class operator \(A\) can be written
\[
A=B_1-B_2+iB_3-iB_4,\qquad B_r\ge0,
\]
with
\begin{equation}
\sum_{r=1}^4\Tr B_r\le2\trnorm A.
\label{eq:V1-positive-decomposition}
\end{equation}
\end{lemma}

\begin{proof}
Write \(A=\operatorname{Re}A+i\operatorname{Im}A\), and use the positive
and negative spectral parts of the two self-adjoint operators.  Their
total traces are
\(\trnorm{\operatorname{Re}A}+\trnorm{\operatorname{Im}A}\).
Both real and imaginary part maps are contractions in trace norm, so
this sum is at most \(2\trnorm A\).
\end{proof}

\begin{theorem}[Weighted raw pinching is controlled by product-state
capacity]
\label{thm:V1-weighted-pinching-capacity}
Let \((P_\nu)_{\nu\in\mathcal N}\) be mutually orthogonal projections on
\(\mathcal H\), and let \(c_\nu\ge0\).  Put
\begin{equation}
Q:=\sum_{\nu\in\mathcal N}c_\nu P_\nu.
\label{eq:V1-scalar-carrier}
\end{equation}
Then, for arbitrary trace-class row inputs \(A_i\),
\begin{equation}
\boxed{
\sum_{\nu\in\mathcal N}
c_\nu\trnorm{P_\nu(A_1\otimes\cdots\otimes A_4)P_\nu}
\le
16\,\kappa_\otimes(Q)\prod_{i=1}^4\trnorm{A_i}.}
\label{eq:V1-weighted-pinching-capacity}
\end{equation}
If all \(A_i\) are positive, the factor \(16\) is replaced by \(1\);
after trace normalization the left side is exactly the expectation of
\(Q\) in a product state.
\end{theorem}

\begin{proof}
First suppose \(A_i\ge0\).  Every compressed block
\(P_\nu(\bigotimes_iA_i)P_\nu\) is positive, so its trace norm equals its
trace.  Orthogonality is not even needed for the identity
\[
\sum_\nu c_\nu
\trnorm{P_\nu(\bigotimes_iA_i)P_\nu}
=
\Tr\!\left[Q\bigotimes_iA_i\right].
\]
If no \(A_i\) is zero, divide each by its trace and apply
\eqref{eq:V1-product-capacity}; the zero case is immediate.

For general \(A_i\), apply \Cref{lem:V1-positive-decomposition} in every
row and expand the tensor product.  Compression and the trace norm obey
the triangle inequality.  Apply the positive-input result to every
product of positive summands.  The sum of the resulting scalar masses is
\[
\prod_{i=1}^4\left(\sum_{r=1}^4\Tr B_{i,r}\right)
\le2^4\prod_{i=1}^4\trnorm{A_i},
\]
which proves \eqref{eq:V1-weighted-pinching-capacity}.
\end{proof}

\begin{remark}[Why ordinary operator norm is too coarse]
\label{rem:V1-capacity-vs-opnorm}
Since \(\kappa_\otimes(Q)\le\opnorm Q=\sup_\nu c_\nu\), a pointwise
final-channel bound is sufficient.  It is not necessary.  Product-state
capacity can exploit the fact that a large output block has very small
overlap with every row-factorized coefficient, exactly the mechanism
needed in the collision analysis.
\end{remark}

\section{Exact reduction of quartic MTA to a positive carrier}
\label{sec:V1-MTA-reduction}

Fix exact-support input representations \(\boldsymbol R_i\), put
\[
\Xi_i:=\Xi(\boldsymbol R_i),\qquad
d_{\mathrm{in}}:=\prod_i d_{\boldsymbol R_i},
\qquad
\vartheta_{i,e}:=\frac{A_{R_{i,e}}}{\Xi_i},
\]
and let \(\ell\) mark each edge of
\(\tau\in\mathfrak T_4\) by a coordinate in the corresponding support
intersection.  Define
\begin{equation}
\mathcal W_{\tau,\ell}(\boldsymbol R)
:=
\left(\prod_{i=1}^4\Xi_i\right)
\prod_{ij\in E(\tau)}
\vartheta_{i,\ell(ij)}\vartheta_{j,\ell(ij)}.
\label{eq:V1-marked-tree-weight}
\end{equation}
This is exactly the product of the four fixed-representation marked BFA
masses.

Resolve the physical final output before taking any norm.  Let
\(\nu=(\boldsymbol U,a)\) denote a final representation and its complete
physical multiplicity/outer label, let \(P_\nu\) be the corresponding
orthogonal projection in the fourfold input tensor product, and let
\(K_\nu\) be the associated central multiplier.  Define the positive
normalized carrier
\begin{equation}
\boxed{
\mathcal Q_{\mathrm{phys}}
:=
\sum_{\substack{\nu=(\boldsymbol U,a)\\
                 \boldsymbol U\ne\boldsymbol1}}
\frac{\Xi(\boldsymbol U)}
     {1-q_{\alpha,\boldsymbol U}}
\opnorm{K_\nu}\,P_\nu.}
\label{eq:V1-physical-carrier}
\end{equation}
The resolvent occurs exactly once in this definition.

\begin{theorem}[Capacity criterion for the fixed-input quartic MTA row]
\label{thm:V1-capacity-MTA}
If the complete physical carrier satisfies
\begin{equation}
\boxed{
\kappa_\otimes(\mathcal Q_{\mathrm{phys}})
\le
C_{\mathrm{cap}}
\sum_{\tau\in\mathfrak T_4}
\sum_{\ell\ {\rm marking}\ \tau}
\mathcal W_{\tau,\ell}(\boldsymbol R),}
\label{eq:V1-capacity-gate}
\end{equation}
then the fixed-input coefficient contribution obeys the fourth
marked-tree atlas inequality, uniformly in final representations,
multiplicities, and the number of physical outer blocks.
\end{theorem}

\begin{proof}
Let \(X=\bigotimes_iA_i^{\rm in}\) and
\(X_\nu=P_\nu XP_\nu\), after the harmless unitary identification of
the \(\nu\)-block with
\(H_{\boldsymbol U}\otimes M_\nu\).  The exact Fourier formula
\eqref{eq:V1-exact-Fourier}, partial-trace contractivity, and the ideal
property of the trace norm give
\begin{align}
&\sum_{\nu=(\boldsymbol U,a)}
e^{\rho\operatorname{ht}(\boldsymbol U)}
\Xi(\boldsymbol U)d_{\boldsymbol U}
\frac{\trnorm{\widehat{\mathbf K}_\nu(\boldsymbol U)}}
     {1-q_{\alpha,\boldsymbol U}}
\notag\\
&\quad\le
d_{\mathrm{in}}
\sum_\nu e^{\rho\operatorname{ht}(\boldsymbol U)}
\frac{\Xi(\boldsymbol U)}
     {1-q_{\alpha,\boldsymbol U}}
\opnorm{K_\nu}\trnorm{X_\nu}.
\label{eq:V1-pre-capacity}
\end{align}
Fusion height is subadditive:
\[
e^{\rho\operatorname{ht}(\boldsymbol U)}
\le\prod_i e^{\rho\operatorname{ht}(\boldsymbol R_i)}.
\]
Apply \Cref{thm:V1-weighted-pinching-capacity} to the remaining sum and
then use \eqref{eq:V1-capacity-gate}.  For each
\((\tau,\ell)\), the resulting factor is
\begin{align*}
&16C_{\mathrm{cap}}d_{\mathrm{in}}
\prod_i e^{\rho\operatorname{ht}(\boldsymbol R_i)}
\mathcal W_{\tau,\ell}(\boldsymbol R)
\prod_i\trnorm{A_i^{\rm in}}\\
&\quad=
16C_{\mathrm{cap}}
\prod_{i=1}^4
\left[
e^{\rho\operatorname{ht}(\boldsymbol R_i)}
\Xi_i d_{\boldsymbol R_i}\trnorm{A_i^{\rm in}}
\prod_{\substack{g\in E(\tau)\\g\ni i}}
\vartheta_{i,\ell(g)}
\right].
\end{align*}
This is the product of the four marked BFA summands.  Summing over the
sixteen trees and their actual coordinate markings is exactly the
fixed-input \(m=4\) MTA right side.  No output count appears, because all
output labels were retained in the single positive carrier before
\Cref{thm:V1-weighted-pinching-capacity} was applied.
\end{proof}

\begin{corollary}[Pointwise control is a stronger sufficient gate]
\label{cor:V1-pointwise-sufficient}
If
\[
\frac{\Xi(\boldsymbol U)}
     {1-q_{\alpha,\boldsymbol U}}\opnorm{K_\nu}
\le
C\sum_{\tau,\ell}\mathcal W_{\tau,\ell}(\boldsymbol R)
\]
for every physical block \(\nu\), then
\eqref{eq:V1-capacity-gate} holds with the same constant.
\end{corollary}

\begin{proof}
The hypothesis is the operator inequality
\[
\mathcal Q_{\mathrm{phys}}
\preccurlyeq
C\left(\sum_{\tau,\ell}\mathcal W_{\tau,\ell}\right)I.
\]
Take product-state expectations.
\end{proof}

\section{Regression checks and a newly revalidated spectator class}
\label{sec:V1-regressions}

\begin{proposition}[The product-state carrier passes the dual-singlet
regression]
\label{prop:V1-singlet-capacity}
Let \(R\) be irreducible and
\[
\Omega_R=d_R^{-1/2}\sum_{a=1}^{d_R}e_a\otimes e_a^*,
\qquad
P_{\boldsymbol1}=|\Omega_R\rangle\langle\Omega_R|.
\]
Then
\begin{equation}
\boxed{\kappa_\otimes(P_{\boldsymbol1})=\frac1{d_R}.}
\label{eq:V1-singlet-capacity}
\end{equation}
\end{proposition}

\begin{proof}
For unit vectors \(u,v\),
\[
|\langle\Omega_R,u\otimes v\rangle|^2
=\frac1{d_R}|\langle\overline v,u\rangle|^2
\le\frac1{d_R},
\]
with equality when \(u=\overline v\).  A product of mixed states is a
convex combination of pure product states, so it cannot increase the
supremum.  This proves the formula.  It agrees exactly with the raw
singlet coefficient mass and therefore avoids the false dimension-law
normalization withdrawn in legacy V45.
\end{proof}

\begin{lemma}[Row-private positive spectators do not increase capacity]
\label{lem:V1-private-spectators}
Suppose
\(\mathcal H_i=\mathcal C_i\otimes\mathcal S_i\),
\(Q_C\ge0\) acts on \(\bigotimes_i\mathcal C_i\), and
\(0\le S_i\le I_{\mathcal S_i}\).  Then
\begin{equation}
\boxed{
\kappa_\otimes\!\left(
Q_C\otimes\bigotimes_{i=1}^4S_i
\right)
\le\kappa_\otimes(Q_C).}
\label{eq:V1-private-spectator-capacity}
\end{equation}
Equality holds when every \(S_i=I\).
\end{lemma}

\begin{proof}
For a row density matrix \(\rho_i\) on
\(\mathcal C_i\otimes\mathcal S_i\), define the subnormalized positive
core state
\[
\sigma_i:=
\Tr_{\mathcal S_i}\!\left[
(I\otimes S_i^{1/2})\rho_i(I\otimes S_i^{1/2})
\right].
\]
Then \(0\le\Tr\sigma_i\le1\), and direct use of the defining property of
partial trace gives
\[
\Tr\!\left[
\left(Q_C\otimes\bigotimes_iS_i\right)\bigotimes_i\rho_i
\right]
=\Tr\!\left[Q_C\bigotimes_i\sigma_i\right].
\]
Normalize every nonzero \(\sigma_i\), apply
\eqref{eq:V1-product-capacity}, and multiply by
\(\prod_i\Tr\sigma_i\le1\).  Taking the supremum proves the inequality.
When \(S_i=I\), every product core state is obtained by tensoring with an
arbitrary private spectator state, giving equality.
\end{proof}

\begin{corollary}[Raw revalidation of the strictly private spectator
branch]
\label{cor:V1-private-branch}
Any local V41 four-junction carrier whose remaining spectator operator
is a tensor product of row-private positive contractions has the same
product-state capacity bound as its local core carrier.  This
amplification requires no representation-dimension formation law.
\end{corollary}

\begin{proof}
Apply \Cref{lem:V1-private-spectators} to the positive scalar carrier
obtained after physical block resolution.  The V41 local core estimate
is unchanged, and \Cref{thm:V1-capacity-MTA} supplies the coefficient
reduction whenever that local core estimate has the marked weight.
\end{proof}

\begin{proposition}[Private input mass forbids post-hoc marking of the
V135 numerator]
\label{prop:V1-no-posthoc}
The unnormalized star monomial
\(H_{12}H_{13}H_{14}\) cannot be bounded by a universal constant times
its marked counterpart
\[
\frac{H_{12}H_{13}H_{14}}{\Xi_1^2}
\]
using only \(0\le H_{ij}\le\Xi_i\Xi_j\).
\end{proposition}

\begin{proof}
Set
\[
H_{12}=H_{13}=H_{14}=1,\quad
\Xi_2=\Xi_3=\Xi_4=1,\quad
\Xi_1=L\ge1,
\]
and set the unused overlaps to zero.  All stated constraints hold.  The
unnormalized monomial equals one and the marked monomial equals
\(L^{-2}\).  A uniform comparison would force \(L^2\) to be bounded.
In the Wilson problem the missing decay must therefore be earned from a
retained multiplier, a coefficient capacity, or the final resolvent; it
is not contained in the V135 numerator norm.
\end{proof}

\section{The first exact unresolved carrier}
\label{sec:V1-first-open-carrier}

Let
\(\mathcal Q_{\tau,\ell}^{[4]}\)
denote \eqref{eq:V1-physical-carrier} restricted to one V41
four-junction tree and one coordinate marking, after removing the
strictly row-private spectator contractions covered by
\Cref{cor:V1-private-branch}.  The first remaining target is
\begin{equation}
\boxed{
\kappa_\otimes(\mathcal Q_{\tau,\ell}^{[4]})
\le C\,\mathcal W_{\tau,\ell}(\boldsymbol R).}
\label{eq:V1-first-open-carrier}
\end{equation}
Its unresolved spectator part has a bank shared by at least two rows.
The legacy V44 localization further separates a triple-shared-bank
branch from a local-output-dominant branch.  Those branches must be
reproved with raw coefficient capacity; the representation-dimension
formation argument formerly used on them is unavailable.

\begin{assessment}[Why this is upstream of collision aggregation]
\label{ass:V1-upstream-order}
The capacity theorem applies equally to \([4]\), \([3,3]\), and the
collision-word sum.  The \([4]\) family is attacked first because its
local core is already a single V41 tree carrier.  Proving
\eqref{eq:V1-first-open-carrier} tests the normalized output mechanism
without the additional Boolean collision selector.  Once the mechanism
is valid, it can be inserted before that selector and before any
triangle inequality over collision words.
\end{assessment}

\chapter{V1 ledger and V2 acquisition}
\label{ch:V1-ledger}

\begin{theorem}[Integrated compact-series V1 statement]
\label{thm:V1-integrated}
This restarted V1:
\begin{enumerate}[label=\textup{(V1.\arabic*)},leftmargin=6em]
\item records the exact frozen f1 source identity and the dependency-aware
      summary of its proved and open gates;
\item repairs the weak-cut proof by
      \eqref{eq:V1-ternary-cut-identity}--%
      \eqref{eq:V1-quartic-weak-cut};
\item proves the weighted raw-pinching theorem
      \eqref{eq:V1-weighted-pinching-capacity};
\item reduces normalized quartic MTA to the positive product-state
      capacity bound \eqref{eq:V1-capacity-gate};
\item proves the exact dual-singlet capacity
      \eqref{eq:V1-singlet-capacity};
\item proves stability under row-private positive spectator contractions;
      and
\item isolates the first unresolved embedded four-junction carrier
      \eqref{eq:V1-first-open-carrier}.
\end{enumerate}
\end{theorem}

\begin{proof}
The claims are respectively
\Cref{sec:V1-archive-identity,sec:V1-legacy-summary,
thm:V1-quartic-weak-cut,thm:V1-weighted-pinching-capacity,
thm:V1-capacity-MTA,prop:V1-singlet-capacity,
lem:V1-private-spectators,sec:V1-first-open-carrier}.
\end{proof}

\begin{openproblem}[V2 mandate: the shared-spectator four-junction
capacity]
\label{op:V2-mandate}
\begin{enumerate}
\item Keep one V41 tree core, the complete V98 physical outer block,
      and the single factor
      \(\Xi(\boldsymbol U)/(1-q_{\alpha,\boldsymbol U})\) inside the
      positive carrier.
\item Split the nonprivate spectator incidence into pair-shared and
      triple-shared banks before taking product-state capacity.
\item Prove a conditional-capacity contraction for a pair-shared
      spectator, or exhibit a two-row state showing why its endpoint
      Casimir tax is necessary.
\item On the triple-shared branch, test the raw carrier on the dual-singlet
      family of \Cref{prop:V1-singlet-capacity}; no dimension-law
      probability may be inserted.
\item On the local-output-dominant branch, combine the final local output
      block with the V41 tree multiplier before applying the sole
      resolvent.
\item Prove \eqref{eq:V1-first-open-carrier} on at least one complete
      physical branch, or isolate the exact missing operator moment.
\item No companion audit is due until compact-series V5.
\end{enumerate}
\end{openproblem}

\begin{firewall}[Claim boundary after compact V1]
\label{fire:V1-current-boundary}
The new capacity criterion is a rigorous reduction and closes strictly
row-private spectator amplification.  It does not prove
\eqref{eq:V1-first-open-carrier} for shared spectators, normalized
collision aggregation, the \([3,3]\) sector, global quartic or all-order
MTA, WUFC, CPM, continuum Yang--Mills existence, or a positive mass gap.
\end{firewall}

\section*{Compact-series V1 record}
\addcontentsline{toc}{section}{Compact-series V1 record}
This file begins the compact series after the exact legacy archive
\eqref{eq:V1-archive-hash}.  It summarizes rather than copies the frozen
135-version corpus and then continues the proof.  The next iteration must
copy this complete file, append new material before its unique active
document-closing command, save it as \texttt{\_V2.tex}, and remove only
the superseded active V1.  The archive \texttt{\_V135\_f1.tex} is retained.

% The compact V1 document terminator is intentionally disabled here:
% \end{document}

% =============================================================================
% BEGIN APPENDED COMPACT VERSION V2 MATERIAL
% =============================================================================

\chapter{V2: Shared-spectator capacity and exact resolvent allocation}
\label{ch:V2}

\section{Why shared positive spectators are not the obstruction}
\label{sec:V2-shared-capacity}

The V1 private-spectator lemma used a rowwise spectator factorization.
The following result removes that restriction.  The spectator may couple
several rows; the price is that only one of the two factors may use
product-state capacity at a time.

\begin{theorem}[One-sided tensor law for regrouped product-state
capacity]
\label{thm:V2-one-sided-tensor}
For each row let
\(\mathcal H_i=\mathcal C_i\otimes\mathcal S_i\).
Let \(Q_C\ge0\) act on \(\bigotimes_i\mathcal C_i\), and let
\(Q_S\ge0\) act on \(\bigotimes_i\mathcal S_i\).  Product-state capacity
on the left is taken across the regrouped row factors
\(\mathcal C_i\otimes\mathcal S_i\).  Then
\begin{equation}
\boxed{
\kappa_\otimes(Q_C\otimes Q_S)
\le
\min\!\left\{
\opnorm{Q_C}\kappa_\otimes(Q_S),
\kappa_\otimes(Q_C)\opnorm{Q_S}
\right\}.}
\label{eq:V2-one-sided-tensor}
\end{equation}
\end{theorem}

\begin{proof}
Fix row density matrices
\(\rho_i\) on \(\mathcal C_i\otimes\mathcal S_i\), and put
\(\rho=\bigotimes_i\rho_i\).  Define the subnormalized positive
spectator operator
\[
\sigma_S:=
\Tr_C\!\left[
(Q_C^{1/2}\otimes I)\rho(Q_C^{1/2}\otimes I)
\right].
\]
Then
\begin{align*}
\Tr[(Q_C\otimes Q_S)\rho]
&=\Tr(Q_S\sigma_S)
\le\opnorm{Q_S}\Tr\sigma_S,\\
\Tr\sigma_S
&=\Tr\!\left[Q_C\bigotimes_i\Tr_{\mathcal S_i}\rho_i\right]
\le\kappa_\otimes(Q_C).
\end{align*}
This proves the second bound.  Interchanging core and spectator proves
the first.  Take the supremum over the row states.
\end{proof}

\begin{corollary}[Arbitrary shared positive contractions are harmless]
\label{cor:V2-shared-contractions}
If \(0\le Q_S\le I\), then
\begin{equation}
\kappa_\otimes(Q_C\otimes Q_S)
\le\kappa_\otimes(Q_C).
\label{eq:V2-shared-contraction}
\end{equation}
Thus pair-shared, triple-shared, and four-shared spectator incidence does
not by itself increase capacity whenever the resolved positive carrier
actually factors.
\end{corollary}

\begin{proof}
Use the second term of
\eqref{eq:V2-one-sided-tensor} and
\(\opnorm{Q_S}\le1\).
\end{proof}

\begin{lemma}[Finite sums of factored positive carriers]
\label{lem:V2-factored-majorant}
If
\[
0\le Q\preccurlyeq\sum_{r=1}^N Q_{C,r}\otimes Q_{S,r},
\qquad
Q_{C,r},Q_{S,r}\ge0,
\]
then
\begin{equation}
\kappa_\otimes(Q)
\le
\sum_{r=1}^N
\min\!\left\{
\opnorm{Q_{C,r}}\kappa_\otimes(Q_{S,r}),
\kappa_\otimes(Q_{C,r})\opnorm{Q_{S,r}}
\right\}.
\label{eq:V2-factored-majorant}
\end{equation}
\end{lemma}

\begin{proof}
Take the expectation of the operator inequality in an arbitrary product
state.  Apply \Cref{thm:V2-one-sided-tensor} to every summand and then
take the supremum.
\end{proof}

\section{Entanglement swapping is a compulsory regression}
\label{sec:V2-swapping}

\begin{proposition}[Conditioning on a shared spectator need not preserve
product core states]
\label{prop:V2-entanglement-swapping}
For every \(d\ge2\), there are two row-product pure states and a
rank-one shared-spectator projection for which the induced normalized
two-core state is maximally entangled.
\end{proposition}

\begin{proof}
Let row one be
\(\mathcal C_1\otimes\mathcal S_1\), row two be
\(\mathcal C_2\otimes\mathcal S_2\), with every factor equal to
\(\mathbb C^d\).  Put
\[
\Omega_{CS}:=d^{-1/2}\sum_{a=1}^d e_a\otimes e_a,
\qquad
\rho:=
|\Omega_{C_1S_1}\rangle\langle\Omega_{C_1S_1}|
\otimes
|\Omega_{C_2S_2}\rangle\langle\Omega_{C_2S_2}|,
\]
which is a product across the two rows.  Let
\(P_{\Omega,S}\) project onto the maximally entangled vector of
\(\mathcal S_1\otimes\mathcal S_2\).  A direct contraction of the two
spectator indices gives
\begin{equation}
\Tr_{S_1S_2}\!\left[
(I_C\otimes P_{\Omega,S})\rho(I_C\otimes P_{\Omega,S})
\right]
=d^{-2}P_{\Omega,C}.
\label{eq:V2-entanglement-swap}
\end{equation}
The trace is \(d^{-2}\), and division by it leaves the maximally
entangled core projection \(P_{\Omega,C}\).
\end{proof}

\begin{remark}[What the regression rules out]
\label{rem:V2-swapping-scope}
Equation \eqref{eq:V2-entanglement-swap} does not disprove the one-sided
tensor law.  It proves that one may not condition on a shared output and
then apply product-state capacity to the conditional state on the other
factor.  The proof of \Cref{thm:V2-one-sided-tensor} correctly uses an
ordinary operator norm on that conditioned side.
\end{remark}

\section{The sole resolvent can be assigned to one output coordinate}
\label{sec:V2-resolvent-allocation}

\begin{lemma}[Scalar one-resolvent allocation]
\label{lem:V2-resolvent-allocation}
Let \(0\le q_e<1\) and \(x_e\ge0\) on a finite coordinate set \(X\).
Then
\begin{equation}
\boxed{
\frac{\sum_{e\in X}x_e}
     {1-\prod_{e\in X}q_e}
\le
\sum_{e\in X}\frac{x_e}{1-q_e}.}
\label{eq:V2-resolvent-allocation}
\end{equation}
\end{lemma}

\begin{proof}
For every \(e\),
\(\prod_fq_f\le q_e\), hence
\[
\frac{x_e}{1-\prod_fq_f}\le\frac{x_e}{1-q_e}.
\]
Sum over \(e\).  The left denominator remains the same in every summand,
so this is an allocation of one resolvent, not a product of resolvents.
\end{proof}

For a product final representation
\(\boldsymbol U=(U_e)_{e\in X}\), the Wilson multiplier and Casimir mass
have the product/sum form
\[
q_{\alpha,\boldsymbol U}=\prod_{e\in X}q_{\alpha,U_e},
\qquad
\Xi(\boldsymbol U)=\sum_{e\in X}A_{U_e}.
\]
Thus \Cref{lem:V2-resolvent-allocation} applies with
\(q_e=q_{\alpha,U_e}\) and \(x_e=A_{U_e}\).

\begin{corollary}[Blockwise positive-carrier allocation]
\label{cor:V2-carrier-allocation}
Let \((P_{\boldsymbol U,a})\) be an orthogonal physical output
resolution, and let \(k_{\boldsymbol U,a}\ge0\).  Then
\begin{align}
&\sum_{\boldsymbol U,a}
\frac{\Xi(\boldsymbol U)}
     {1-q_{\alpha,\boldsymbol U}}
k_{\boldsymbol U,a}P_{\boldsymbol U,a}
\notag\\
&\qquad\preccurlyeq
\sum_{e\in X}\sum_{\boldsymbol U,a}
\frac{A_{U_e}}{1-q_{\alpha,U_e}}
k_{\boldsymbol U,a}P_{\boldsymbol U,a}.
\label{eq:V2-carrier-allocation}
\end{align}
\end{corollary}

\begin{proof}
The scalar coefficient inequality
\eqref{eq:V2-resolvent-allocation} holds on every mutually orthogonal
block.  Summing those block inequalities gives the Loewner inequality.
\end{proof}

\section{Exact factorization for a unique-junction coordinate word}
\label{sec:V2-unique-junction-factorization}

Fix one four-row junction coordinate \(e_0\), one V41 tree fibre, and a
coordinatewise physical resolution
\[
P_{\boldsymbol U,\boldsymbol a}
=\bigotimes_{e\in X}P_{U_e,a_e}^{e}.
\]
Suppose the norm of the resolved coordinate word factors as
\begin{equation}
k_{\boldsymbol U,\boldsymbol a}
\le
k_{e_0}(U_{e_0},a_{e_0})
\prod_{f\ne e_0}m_f(U_f,a_f),
\qquad
0\le m_f\le1.
\label{eq:V2-word-factorization}
\end{equation}
This is the exact situation for a unique local cumulant multiplied by
complete positive spectator moment contractions; it is asserted here
only for words for which the displayed factorization has been verified.

Define the local junction carrier
\[
J_{e_0}:=
\sum_{U,a}k_{e_0}(U,a)P_{U,a}^{e_0},
\]
the spectator contraction
\[
M_f:=\sum_{U,a}m_f(U,a)P_{U,a}^{f},
\qquad 0\le M_f\le I,
\]
and the payer carriers
\begin{align}
J_{e_0}^{\rm pay}
&:=
\sum_{U,a}
\frac{A_U}{1-q_{\alpha,U}}
k_{e_0}(U,a)P_{U,a}^{e_0},
\label{eq:V2-junction-payer}\\
D_f^{\rm pay}
&:=
\sum_{U,a}
\frac{A_U}{1-q_{\alpha,U}}
m_f(U,a)P_{U,a}^{f}.
\label{eq:V2-spectator-payer}
\end{align}

\begin{theorem}[Payer-coordinate capacity reduction]
\label{thm:V2-payer-reduction}
Under \eqref{eq:V2-word-factorization}, the normalized physical carrier
of that coordinate word satisfies
\begin{align}
\kappa_\otimes(\mathcal Q_{\rm word})
\le{}&
\kappa_\otimes(J_{e_0}^{\rm pay})\notag\\
&+
\sum_{f\ne e_0}
\min\!\left\{
\opnorm{J_{e_0}}\kappa_\otimes(D_f^{\rm pay}),
\kappa_\otimes(J_{e_0})\opnorm{D_f^{\rm pay}}
\right\}.
\label{eq:V2-payer-reduction}
\end{align}
All capacities use the original row grouping at the indicated coordinate.
No coordinate count is hidden: the displayed sum is the exact output-mass
allocation, and it must later be summed by input marks or multiplier
moments.
\end{theorem}

\begin{proof}
Apply \Cref{cor:V2-carrier-allocation}.  The \(e_0\)-payer summand is
bounded, using \eqref{eq:V2-word-factorization}, by
\[
J_{e_0}^{\rm pay}\otimes\bigotimes_{h\ne e_0}M_h.
\]
Repeated use of \Cref{cor:V2-shared-contractions} bounds its capacity by
\(\kappa_\otimes(J_{e_0}^{\rm pay})\), even when an \(M_h\) couples
several rows.

For \(f\ne e_0\), the \(f\)-payer summand is bounded by
\[
J_{e_0}\otimes D_f^{\rm pay}
\otimes\bigotimes_{h\notin\{e_0,f\}}M_h.
\]
Discard the contraction factors by
\Cref{cor:V2-shared-contractions}, and apply
\Cref{thm:V2-one-sided-tensor} to the remaining two coordinate
carriers.  Finally use subadditivity of capacity for the positive sum
over payer coordinates.
\end{proof}

\begin{firewall}[Factorization must precede final recoupling]
\label{fire:V2-factorization-order}
The theorem does not assert
\eqref{eq:V2-word-factorization} for an arbitrary topology projection,
signed collision sum, or output basis which mixes coordinates.  It
applies to a complete coordinate word whose physical resolution is
coordinatewise and whose positive spectator multipliers remain intact.
If final recoupling mixes the proposed tensor factors, one must first
construct a common positive carrier or keep the larger unfactored
capacity.
\end{firewall}

\section{What the new reduction changes}
\label{sec:V2-consequence}

\begin{corollary}[Shared incidence is removed as a standalone
bottleneck]
\label{cor:V2-incidence-removed}
For a factorized unique-junction word, pair-, triple-, or four-row
incidence of a nonpayer positive spectator contraction creates no
capacity loss.  The only unpaid quantities in
\eqref{eq:V2-payer-reduction} are:
\begin{enumerate}[label=\textup{(\roman*)},leftmargin=4em]
\item the locally normalized four-junction capacity
      \(\kappa_\otimes(J_{e_0}^{\rm pay})\); and
\item for a spectator payer \(f\), one one-sided combination of the
      junction norm/capacity with the weighted spectator
      norm/capacity.
\end{enumerate}
\end{corollary}

\begin{proof}
Every nonpayer spectator is removed by
\Cref{cor:V2-shared-contractions}; the two listed quantities are exactly
the remaining terms of \eqref{eq:V2-payer-reduction}.
\end{proof}

\begin{assessment}[Upstream bottleneck after V2]
\label{ass:V2-upstream-bottleneck}
The legacy phrase \emph{triple-shared spectator obstruction} was too
coarse.  A shared positive contraction is harmless after exact
factorization.  The actual problem is the payer carrier
\(D_f^{\rm pay}\), whose factor
\(A_U/(1-q_{\alpha,U})\) is unbounded unless a retained spectator
multiplier supplies an output first moment, together with the
local-junction payer \(J_{e_0}^{\rm pay}\).  V3 should prove those two
coordinate-local moment bounds before returning to global collision
aggregation.
\end{assessment}

\section{V2 ledger and V3 acquisition}
\label{sec:V2-ledger}

\begin{theorem}[Integrated compact-series V2 statement]
\label{thm:V2-integrated}
Preserving compact V1, V2 proves:
\begin{enumerate}[label=\textup{(V2.\arabic*)},leftmargin=6em]
\item the one-sided regrouped tensor law
      \eqref{eq:V2-one-sided-tensor};
\item contraction of arbitrary shared positive spectators;
\item the entanglement-swapping regression
      \eqref{eq:V2-entanglement-swap};
\item the exact one-resolvent allocation
      \eqref{eq:V2-resolvent-allocation};
\item its blockwise positive-carrier form
      \eqref{eq:V2-carrier-allocation}; and
\item the unique-junction payer-coordinate reduction
      \eqref{eq:V2-payer-reduction}.
\end{enumerate}
\end{theorem}

\begin{proof}
These are
\Cref{thm:V2-one-sided-tensor,cor:V2-shared-contractions,
prop:V2-entanglement-swapping,lem:V2-resolvent-allocation,
cor:V2-carrier-allocation,thm:V2-payer-reduction}.
\end{proof}

\begin{openproblem}[V3 mandate: local payer moment estimates]
\label{op:V3-mandate}
\begin{enumerate}
\item For a positive Wilson spectator moment
      \(m_f(U,a)\), prove the sharp capacity and operator bounds for
      \(D_f^{\rm pay}\) in \eqref{eq:V2-spectator-payer} from the
      one-link first-moment debit; do not replace raw mass by a
      representation-dimension law.
\item For the V41 four-junction fibre, compute
      \(\kappa_\otimes(J_{e_0}^{\rm pay})\) with the local output block
      and resolvent kept together.
\item Test both estimates on high-weight coherent representations and
      on the dual-singlet family.
\item Insert the resulting bound into
      \eqref{eq:V2-payer-reduction} and identify which tree vertex earns
      each repeated global mark.
\item If the payer-coordinate sum introduces support loss, sum it by the
      exact mark normalization before taking the global input sums.
\item No companion audit is due until compact-series V5.
\end{enumerate}
\end{openproblem}

\begin{firewall}[Claim boundary after compact V2]
\label{fire:V2-current-boundary}
V2 removes shared positive contraction incidence as an independent
capacity obstruction and isolates exact coordinate-local payer carriers.
It does not yet bound those payer carriers at the full marked scale,
close the embedded \([4]\) row, aggregate normalized collisions, treat
\([3,3]\), prove global quartic or all-order MTA, construct continuum
Yang--Mills theory, or prove a positive mass gap.
\end{firewall}

\section{Append-only record for compact V2}
\label{sec:V2-record}

V2 was appended to the complete compact V1 file.  The exact V1 parent had
\(795\) newline-delimited source records, \(28782\) bytes, and SHA--256
digest
\[
\texttt{aea626ae295c5627c4ac8063f84aae4ebf21490810e81a17149f4d69d6a25a91}.
\]
No compact V1 mathematical content was changed.  The sole final
\verb|\end{document}| is restored below.

% =============================================================================
% END APPENDED COMPACT VERSION V2 MATERIAL
% =============================================================================

% The compact V2 document terminator is intentionally disabled here:
% \end{document}

% =============================================================================
% BEGIN APPENDED COMPACT VERSION V3 MATERIAL
% =============================================================================

\chapter{V3: Coordinate-payer moments and the high-output capacity gate}
\label{ch:V3}

\section{Standing one-link Wilson inequalities}
\label{sec:V3-Wilson-inputs}

Write \(q_{\alpha,U}\in[0,1]\) for a nontrivial one-coordinate Wilson
multiplier and \(A_U>0\) for its Casimir mass.  The established one-link
gap and first-moment estimates have the form
\begin{align}
1-q_{\alpha,U}
&\ge c_{\rm gap}\min\{1,s_\alpha A_U\},
\label{eq:V3-one-link-gap}\\
s_\alpha A_Uq_{\alpha,U}
&\le C_{\rm mom}.
\label{eq:V3-one-link-first-moment}
\end{align}
All constants below depend only on these fixed Wilson constants and the
compact group.

\begin{lemma}[Ordinary spectator payer bound]
\label{lem:V3-spectator-payer}
For every nontrivial one-link output,
\begin{equation}
\boxed{
\frac{A_Uq_{\alpha,U}}{1-q_{\alpha,U}}
\le\frac{C}{s_\alpha}.}
\label{eq:V3-spectator-payer}
\end{equation}
Consequently, if the spectator payer in
\eqref{eq:V2-spectator-payer} is an ordinary positive Wilson moment,
then
\begin{equation}
0\le D_f^{\rm pay}\preccurlyeq\frac{C}{s_\alpha}I,
\qquad
\opnorm{D_f^{\rm pay}},
\ \kappa_\otimes(D_f^{\rm pay})
\le\frac{C}{s_\alpha}.
\label{eq:V3-spectator-carrier-cap}
\end{equation}
\end{lemma}

\begin{proof}
If \(s_\alpha A_U\le1\), use
\eqref{eq:V3-one-link-gap}:
\[
\frac{A_Uq_{\alpha,U}}{1-q_{\alpha,U}}
\le\frac{q_{\alpha,U}}{c_{\rm gap}s_\alpha}
\le\frac1{c_{\rm gap}s_\alpha}.
\]
If \(s_\alpha A_U>1\), the gap is at least \(c_{\rm gap}\), while
\eqref{eq:V3-one-link-first-moment} bounds the numerator by
\(C_{\rm mom}/s_\alpha\).  This proves the scalar estimate.
The carrier is diagonal with these nonnegative coefficients, so its
operator norm has the same bound; product-state capacity is no larger
than operator norm.
\end{proof}

\begin{remark}[The \(s_\alpha^{-1}\) scale is sharp]
\label{rem:V3-payer-sharp}
When \(s_\alpha A_U\) is small,
\(1-q_{\alpha,U}\) has first order \(s_\alpha A_U\), and the quotient in
\eqref{eq:V3-spectator-payer} has order \(s_\alpha^{-1}\).  Therefore a
proof requiring an additional power of \(s_\alpha\) from this payer alone
cannot hold uniformly.  The missing powers must come from the connected
junction or from marked input mass.
\end{remark}

\section{A fourth-cumulant output first moment}
\label{sec:V3-four-output-moment}

At one coordinate, let \(K_{4,U}\) be the complete four-row cumulant
multiplier on a resolved final channel \(U\).  For a partition
\(\pi\in\Pi_4\), resolve the output \(S_B\) of every block
\(B\in\pi\).  The corresponding moment product has scalar Wilson factor
\(\prod_{B\in\pi}q_{\alpha,S_B}\) and contraction recouplings.  Fusion
geometry gives
\begin{equation}
A_U\le C_{\rm fus}\sum_{B\in\pi}A_{S_B}
\label{eq:V3-fusion-output}
\end{equation}
whenever \(U\) occurs in \(\bigotimes_{B\in\pi}S_B\).

\begin{theorem}[Uniform output first moment of the complete local fourth
cumulant]
\label{thm:V3-four-output-moment}
Under \eqref{eq:V3-one-link-first-moment} and
\eqref{eq:V3-fusion-output},
\begin{equation}
\boxed{
s_\alpha A_U\opnorm{K_{4,U}}\le C_4}
\label{eq:V3-four-output-moment}
\end{equation}
uniformly in the four input representations, the final representation,
and all intermediate and final multiplicities.
\end{theorem}

\begin{proof}
Use the finite moment--cumulant identity
\begin{equation}
K_{4,U}
=
\sum_{\pi\in\Pi_4}
(|\pi|-1)!(-1)^{|\pi|-1}M_{\pi,U},
\label{eq:V3-four-moment-cumulant}
\end{equation}
where \(M_{\pi,U}\) is the complete contraction of the block moment
operators onto final channel \(U\).  Fix one fully resolved intermediate
path in \(M_{\pi,U}\).  Recoupling maps and physical compressions are
contractions, hence
\[
\opnorm{M_{\pi,U}}
\le\prod_{B\in\pi}q_{\alpha,S_B}.
\]
Multiplying by \(s_\alpha A_U\), using
\eqref{eq:V3-fusion-output}, and assigning the sum one block at a time
gives
\begin{align*}
s_\alpha A_U\prod_{B\in\pi}q_{\alpha,S_B}
&\le
C_{\rm fus}\sum_{B\in\pi}
\left(s_\alpha A_{S_B}q_{\alpha,S_B}\right)
\prod_{\substack{D\in\pi\\D\ne B}}q_{\alpha,S_D}\\
&\le C_{\rm fus}C_{\rm mom}|\pi|.
\end{align*}
The estimate is an operator estimate on the complete resolved
multiplicity block; it does not sum matrix entries.  Summing the fifteen
partition terms with their fixed Moebius coefficients proves
\eqref{eq:V3-four-output-moment}.
\end{proof}

\begin{corollary}[Universal local-junction payer cap]
\label{cor:V3-junction-payer-cap}
Let
\[
J_4^{\rm pay}
:=
\sum_{U\ne\boldsymbol1}
\frac{A_U}{1-q_{\alpha,U}}
\opnorm{K_{4,U}}P_U.
\]
Then
\begin{equation}
\boxed{
\opnorm{J_4^{\rm pay}}\le\frac{C}{s_\alpha}.}
\label{eq:V3-junction-payer-cap}
\end{equation}
\end{corollary}

\begin{proof}
If \(s_\alpha A_U>1\), combine the constant lower gap in
\eqref{eq:V3-one-link-gap} with
\eqref{eq:V3-four-output-moment}.  If
\(s_\alpha A_U\le1\), then
\[
\frac{A_U}{1-q_{\alpha,U}}\opnorm{K_{4,U}}
\le\frac{C}{s_\alpha}\opnorm{K_{4,U}},
\]
and the universal cumulant cap bounds the last norm.  Taking the
supremum over the orthogonal blocks proves the display.
\end{proof}

\section{Tree-refined low-output improvement}
\label{sec:V3-tree-low-output}

For a V41 tree fibre \(\tau\), retain its local valence weight
\begin{equation}
P_\tau(\boldsymbol R;e)
:=
\prod_{i=1}^4 A_{R_{i,e}}^{\deg_\tau(i)/2}.
\label{eq:V3-local-valence}
\end{equation}
The representation-uniform V41 estimate is
\begin{equation}
\opnorm{K_{\tau,U}}
\le C s_\alpha^3P_\tau(\boldsymbol R;e).
\label{eq:V3-V41-tree}
\end{equation}

\begin{lemma}[Low-output payer retains the local tree numerator]
\label{lem:V3-low-output-tree}
On \(s_\alpha A_U\le1\),
\begin{equation}
\boxed{
\frac{A_U}{1-q_{\alpha,U}}
\opnorm{K_{\tau,U}}
\le
C s_\alpha^2P_\tau(\boldsymbol R;e).}
\label{eq:V3-low-output-tree}
\end{equation}
\end{lemma}

\begin{proof}
The low-output part of
\eqref{eq:V3-one-link-gap} gives
\[
\frac{A_U}{1-q_{\alpha,U}}\le\frac{C}{s_\alpha}.
\]
Multiply by \eqref{eq:V3-V41-tree}.
\end{proof}

\begin{corollary}[Exact cold/hot decomposition of the local payer]
\label{cor:V3-cold-hot-payer}
Writing \(P_{\rm cold}\) and \(P_{\rm hot}\) for the physical output
projections \(s_\alpha A_U\le1\) and \(s_\alpha A_U>1\),
\begin{align}
\opnorm{P_{\rm cold}J_\tau^{\rm pay}}
&\le Cs_\alpha^2P_\tau(\boldsymbol R;e),
\label{eq:V3-cold-payer}\\
\opnorm{P_{\rm hot}J_\tau^{\rm pay}}
&\le C/s_\alpha.
\label{eq:V3-hot-payer}
\end{align}
The first estimate preserves the complete tree valence.  The second is
uniform but too coarse to imply repeated global input marks.
\end{corollary}

\begin{proof}
The cold estimate is
\Cref{lem:V3-low-output-tree}.  The proof of
\Cref{cor:V3-junction-payer-cap}, restricted to hot blocks, gives the
second.
\end{proof}

\section{Insertion into the V2 payer reduction}
\label{sec:V3-insertion}

\begin{theorem}[What coordinate-local first moments do and do not close]
\label{thm:V3-payer-consequence}
For a factorized unique-junction word satisfying
\eqref{eq:V2-word-factorization}, all terms of
\eqref{eq:V2-payer-reduction} are finite and independent of final-channel
and multiplicity counts.  More precisely:
\begin{enumerate}[label=\textup{(\roman*)},leftmargin=4em]
\item every ordinary spectator payer has both norm and capacity
      \(O(s_\alpha^{-1})\);
\item the local four-junction payer has norm
      \(O(s_\alpha^{-1})\);
\item its cold-output tree fibre has the sharper bound
      \(Cs_\alpha^2P_\tau\).
\end{enumerate}
The remaining normalized \([4]\) problem is exactly the product-state
capacity of the hot-output local payer together with the conversion of
the \(s_\alpha^{-1}\) spectator payer into the repeated global marks.
\end{theorem}

\begin{proof}
Items \textup{(i)--(iii)} are
\Cref{lem:V3-spectator-payer,cor:V3-junction-payer-cap,
lem:V3-low-output-tree}, inserted into
\eqref{eq:V2-payer-reduction}.  These estimates remove output and
multiplicity counts, but \Cref{prop:V1-no-posthoc} shows that a uniform
operator cap cannot create the branching-row denominators.  Therefore
the two stated capacity/mark conversions are precisely what remains.
\end{proof}

\begin{assessment}[The next upstream object]
\label{ass:V3-next-object}
The hot local carrier must not be estimated by
\eqref{eq:V3-hot-payer} before the input coefficient is inserted.
The next object is the positive hot carrier
\[
\mathcal H_{\tau,e}
:=
\sum_{\substack{U,a\\s_\alpha A_U>1}}
\frac{A_U}{1-q_{\alpha,U}}
\opnorm{K_{\tau,U,a}}P_{U,a}.
\]
V4 must estimate
\(\kappa_\otimes(\mathcal H_{\tau,e})\), not
\(\opnorm{\mathcal H_{\tau,e}}\).  The dual-singlet calculation in
\eqref{eq:V1-singlet-capacity} shows the possible gain: a large block can
have small overlap with every product coefficient even when its operator
norm is order one.
\end{assessment}

\section{V3 ledger and V4 acquisition}
\label{sec:V3-ledger}

\begin{theorem}[Integrated compact-series V3 statement]
\label{thm:V3-integrated}
Preserving compact V1--V2, V3 proves:
\begin{enumerate}[label=\textup{(V3.\arabic*)},leftmargin=6em]
\item the sharp ordinary spectator payer estimate
      \eqref{eq:V3-spectator-payer};
\item the complete local fourth-cumulant output moment
      \eqref{eq:V3-four-output-moment};
\item the universal local-junction payer cap
      \eqref{eq:V3-junction-payer-cap};
\item the tree-refined cold-output bound
      \eqref{eq:V3-low-output-tree}; and
\item isolation of the hot product-state carrier
      \(\mathcal H_{\tau,e}\).
\end{enumerate}
\end{theorem}

\begin{proof}
These are
\Cref{lem:V3-spectator-payer,thm:V3-four-output-moment,
cor:V3-junction-payer-cap,lem:V3-low-output-tree,
ass:V3-next-object}.
\end{proof}

\begin{openproblem}[V4 mandate: hot local-junction product capacity]
\label{op:V4-mandate}
\begin{enumerate}
\item Resolve \(\mathcal H_{\tau,e}\) in the actual fourfold
      Clebsch--Gordan carrier, but retain the full multiplicity block.
\item Bound its expectation on arbitrary product row states before
      replacing any factor by an operator norm.
\item Use the high-output condition only through a proved fusion or
      Casimir relation; do not infer a large individual input without a
      disjoint priority split.
\item Partition by the largest local input mass and compare its location
      with the degree-two or degree-three vertices of \(\tau\).
\item Test the resulting estimate on the coherent high-weight family
      from legacy V40 and the dual-singlet family from compact V1.
\item Either prove the hot part of
      \eqref{eq:V1-first-open-carrier}, or exhibit the exact valence
      profile for which one global input mark is still absent.
\item Compact V5 must begin with the mandatory SM/NS companion audit.
\end{enumerate}
\end{openproblem}

\begin{firewall}[Claim boundary after compact V3]
\label{fire:V3-current-boundary}
V3 proves output first moments and removes all final-channel counting
from the coordinate payer carriers.  It does not prove the hot carrier
at the marked product-state scale, close the full embedded \([4]\) row,
aggregate normalized collisions, treat \([3,3]\), prove quartic or
all-order MTA, construct continuum Yang--Mills theory, or prove a positive
mass gap.
\end{firewall}

\section{Append-only record for compact V3}
\label{sec:V3-record}

V3 was appended to the complete compact V2 file.  The exact V2 parent had
\(1233\) newline-delimited source records, \(43133\) bytes, and SHA--256
digest
\[
\texttt{6a60062aa0af2539ad985a8cdce1cd1b07614f010cf0f6a0ebf9d1a17699c151}.
\]
No compact V2 mathematical content was changed.  The sole final
\verb|\end{document}| is restored below.

% =============================================================================
% END APPENDED COMPACT VERSION V3 MATERIAL
% =============================================================================

% The compact V3 document terminator is intentionally disabled here:
% \end{document}

% =============================================================================
% BEGIN APPENDED COMPACT VERSION V4 MATERIAL
% =============================================================================

\chapter{V4: Cartan-output saturation and the irreducible hot-junction
gate}
\label{ch:V4}

\section{Highest-weight output has full product-state capacity}
\label{sec:V4-Cartan-capacity}

Let \(G\) be the fixed compact connected gauge group.  For dominant
highest weights \(\lambda_i\), let \(R_{\lambda_i}\) be the corresponding
irreducible representation and choose unit highest-weight vectors
\(v_i\in H_{\lambda_i}\).

\begin{lemma}[Cartan component and its product highest-weight vector]
\label{lem:V4-Cartan-vector}
The tensor product
\(\bigotimes_{i=1}^mR_{\lambda_i}\) contains the irreducible Cartan
component \(R_{\lambda_1+\cdots+\lambda_m}\) with multiplicity one, and
\begin{equation}
v_1\otimes\cdots\otimes v_m
\in H_{\lambda_1+\cdots+\lambda_m}.
\label{eq:V4-product-highest-vector}
\end{equation}
\end{lemma}

\begin{proof}
The tensor product vector in
\eqref{eq:V4-product-highest-vector} has weight
\(\lambda_1+\cdots+\lambda_m\).  Every positive-root raising operator
acts through the coproduct as the sum of its actions on the tensor
factors and therefore annihilates the vector.  It generates an
irreducible highest-weight submodule with that highest weight.

The weight space of weight \(\lambda_1+\cdots+\lambda_m\) in the full
tensor product is one-dimensional: reaching the sum of the highest
weights forces every factor to lie in its own one-dimensional
highest-weight space.  Hence only one copy of the generated irreducible
can occur.
\end{proof}

\begin{theorem}[Cartan output projector saturates product-state capacity]
\label{thm:V4-Cartan-capacity}
Let \(P_{\rm Car}\) be the isotypic projection onto
\(R_{\lambda_1+\cdots+\lambda_m}\).  Across the original \(m\) input
rows,
\begin{equation}
\boxed{\kappa_\otimes(P_{\rm Car})=1.}
\label{eq:V4-Cartan-capacity}
\end{equation}
More generally, for \(c\ge0\),
\(\kappa_\otimes(cP_{\rm Car})=c\).
\end{theorem}

\begin{proof}
Every projection has operator norm one, so its product-state capacity is
at most one.  By \Cref{lem:V4-Cartan-vector},
\[
\langle v_1\otimes\cdots\otimes v_m,
P_{\rm Car}(v_1\otimes\cdots\otimes v_m)\rangle=1.
\]
The pure product state of the \(v_i\) therefore attains the upper bound.
Positive homogeneity gives the final assertion.
\end{proof}

\begin{corollary}[No uniform hot-output rarity principle]
\label{cor:V4-no-hot-rarity}
There is no function \(\varepsilon(A)\to0\) as \(A\to\infty\) such that
every isotypic output projection \(P_U\) with \(A_U\ge A\) satisfies
\(\kappa_\otimes(P_U)\le\varepsilon(A)\).
\end{corollary}

\begin{proof}
Choose a sequence of dominant input weights whose sum tends to infinity
and take the Cartan component.  Its Casimir tends to infinity, while
\eqref{eq:V4-Cartan-capacity} remains one.
\end{proof}

\section{Application to the V3 hot local carrier}
\label{sec:V4-hot-carrier}

Restrict the hot carrier from
\Cref{ass:V3-next-object} to the Cartan output:
\[
\mathcal H_{\tau,e}^{\rm Car}
:=
\frac{A_{\rm Car}}{1-q_{\alpha,{\rm Car}}}
\opnorm{K_{\tau,{\rm Car}}}P_{\rm Car},
\qquad
s_\alpha A_{\rm Car}>1.
\]

\begin{proposition}[Cartan capacity equals the unresolved scalar
multiplier]
\label{prop:V4-Cartan-hot-exact}
One has the exact identity
\begin{equation}
\boxed{
\kappa_\otimes(\mathcal H_{\tau,e}^{\rm Car})
=
\frac{A_{\rm Car}}{1-q_{\alpha,{\rm Car}}}
\opnorm{K_{\tau,{\rm Car}}}.}
\label{eq:V4-Cartan-hot-exact}
\end{equation}
\end{proposition}

\begin{proof}
The scalar coefficient is nonnegative and
\Cref{thm:V4-Cartan-capacity} applies.
\end{proof}

\begin{corollary}[Necessary Cartan regression for the capacity route]
\label{cor:V4-Cartan-necessary}
Any proof of the local hot capacity bound
\[
\kappa_\otimes(\mathcal H_{\tau,e})
\le C\mathcal W_{\tau,\ell}(\boldsymbol R)
\]
must, on the Cartan family, prove the pointwise estimate
\begin{equation}
\boxed{
\frac{A_{\rm Car}}{1-q_{\alpha,{\rm Car}}}
\opnorm{K_{\tau,{\rm Car}}}
\le C\mathcal W_{\tau,\ell}(\boldsymbol R).}
\label{eq:V4-Cartan-pointwise-gate}
\end{equation}
Raw pinching, output multiplicity compression, and product-state
capacity cannot supply an additional small factor on this family.
\end{corollary}

\begin{proof}
The Cartan carrier is a positive subcarrier of
\(\mathcal H_{\tau,e}\).  Capacity is monotone under Loewner order.
Apply \eqref{eq:V4-Cartan-hot-exact}.
\end{proof}

\section{The existing first moment is insufficient}
\label{sec:V4-first-moment-insufficient}

\begin{proposition}[Logical separation between the V3 first moment and
the marked Cartan gate]
\label{prop:V4-first-moment-insufficient}
The estimate
\[
s_\alpha A_{\rm Car}\opnorm{K_{\tau,{\rm Car}}}\le C
\]
does not, by itself, imply
\eqref{eq:V4-Cartan-pointwise-gate} uniformly in spectator input masses.
\end{proposition}

\begin{proof}
On the hot branch, the one-link gap only yields
\[
\frac{A_{\rm Car}}{1-q_{\alpha,{\rm Car}}}
\opnorm{K_{\tau,{\rm Car}}}
\le C/s_\alpha.
\]
The right side contains no global input masses.  For a star centered at
row one, the marked factor contributed by that tree is proportional to
\[
\frac{H_{12}H_{13}H_{14}}{\Xi_1^2}.
\]
The private-mass family in \Cref{prop:V1-no-posthoc} keeps the overlaps
fixed while letting \(\Xi_1\to\infty\), so this marked quantity tends to
zero while \(C/s_\alpha\) does not.  Hence the first-moment conclusion
alone cannot imply the marked gate.  This is a logical insufficiency of
the estimate, not a claim that the exact Wilson cumulant stays large on
that family.
\end{proof}

\begin{assessment}[Where the missing decay must occur]
\label{ass:V4-missing-decay}
The Cartan projection has full product capacity.  Therefore the missing
branching-row factors must already be visible in the exact Cartan
matrix element of the complete coordinate word.  They can arise from:
\begin{enumerate}[label=\textup{(\roman*)},leftmargin=4em]
\item spectator Wilson multipliers retained in the same row as a
      high-degree tree vertex;
\item cancellation between the fifteen moment-partition terms before
      the Cartan matrix element is absolutized; or
\item a positive endpoint tax formed jointly from the local junction and
      a selected spectator bank.
\end{enumerate}
Neither final-channel rarity nor an additional application of raw
pinching is available.
\end{assessment}

\section{A safe Cartan/non-Cartan split}
\label{sec:V4-output-split}

\begin{lemma}[Orthogonal capacity split]
\label{lem:V4-capacity-split}
Let
\[
\mathcal H_{\tau,e}
=\mathcal H_{\tau,e}^{\rm Car}
+\mathcal H_{\tau,e}^{\rm nCar}
\]
be the decomposition by the Cartan projector and its orthogonal
complement.  Then
\begin{equation}
\max\!\left\{
\kappa_\otimes(\mathcal H_{\tau,e}^{\rm Car}),
\kappa_\otimes(\mathcal H_{\tau,e}^{\rm nCar})
\right\}
\le
\kappa_\otimes(\mathcal H_{\tau,e})
\le
\kappa_\otimes(\mathcal H_{\tau,e}^{\rm Car})
+\kappa_\otimes(\mathcal H_{\tau,e}^{\rm nCar}).
\label{eq:V4-capacity-split}
\end{equation}
\end{lemma}

\begin{proof}
Both summands are positive.  The lower estimates follow from monotonicity.
For the upper estimate, take a product-state expectation, use additivity
of the expectation, bound the two terms by their separate suprema, and
then take the supremum.
\end{proof}

\begin{corollary}[Two genuinely different hot-output tasks]
\label{cor:V4-two-hot-tasks}
The hot local-junction problem splits without overlap into:
\begin{enumerate}[label=\textup{(\roman*)},leftmargin=4em]
\item the Cartan matrix-element gate
      \eqref{eq:V4-Cartan-pointwise-gate}; and
\item a non-Cartan product-capacity tail for
      \(\mathcal H_{\tau,e}^{\rm nCar}\).
\end{enumerate}
Closing both implies the hot part of the local carrier; neither task may
be substituted for the other.
\end{corollary}

\begin{proof}
Apply \Cref{lem:V4-capacity-split}.  The exact form of the first task is
\Cref{prop:V4-Cartan-hot-exact}.
\end{proof}

\section{V4 ledger and mandatory V5 audit}
\label{sec:V4-ledger}

\begin{theorem}[Integrated compact-series V4 statement]
\label{thm:V4-integrated}
Preserving compact V1--V3, V4 proves:
\begin{enumerate}[label=\textup{(V4.\arabic*)},leftmargin=6em]
\item multiplicity one of the Cartan component and its product
      highest-weight vector;
\item the exact capacity identity
      \eqref{eq:V4-Cartan-capacity};
\item failure of a universal hot-output rarity principle;
\item the exact Cartan hot-carrier value
      \eqref{eq:V4-Cartan-hot-exact};
\item necessity of the pointwise Cartan gate
      \eqref{eq:V4-Cartan-pointwise-gate}; and
\item the orthogonal Cartan/non-Cartan capacity split
      \eqref{eq:V4-capacity-split}.
\end{enumerate}
\end{theorem}

\begin{proof}
These are
\Cref{lem:V4-Cartan-vector,thm:V4-Cartan-capacity,
cor:V4-no-hot-rarity,prop:V4-Cartan-hot-exact,
cor:V4-Cartan-necessary,lem:V4-capacity-split}.
\end{proof}

\begin{openproblem}[V5 mandate: companion audit and exact Cartan
matrix element]
\label{op:V5-mandate}
\begin{enumerate}
\item Begin with the mandatory read-only audit of the newest
      Standard-Model--Einstein and Navier--Stokes notes; record their
      versions, sizes, hashes, and only type-correct structural lessons.
\item Write the exact V41 tree-fibre matrix element on
      \(v_1\otimes\cdots\otimes v_4\), with every spectator multiplier
      retained.
\item Keep the fifteen fourth-cumulant partition terms assembled until
      their Cartan scalar cancellation has been computed.
\item Split spectator mass at each degree-two or degree-three vertex
      into junction and off-junction parts.
\item Prove the required branching-row decay on one priority branch, or
      produce the first explicit high-weight/private-mass scaling of the
      exact Cartan matrix element.
\item Treat the non-Cartan capacity tail only after the Cartan regression
      passes.
\end{enumerate}
\end{openproblem}

\begin{firewall}[Claim boundary after compact V4]
\label{fire:V4-current-boundary}
V4 proves that product-state capacity offers no output-rarity gain on
the Cartan family and isolates the exact pointwise matrix-element gate.
It does not prove that gate, the non-Cartan tail, the full embedded
\([4]\) row, normalized collision aggregation, \([3,3]\), quartic or
all-order MTA, continuum Yang--Mills existence, or a positive mass gap.
\end{firewall}

\section{Append-only record for compact V4}
\label{sec:V4-record}

V4 was appended to the complete compact V3 file.  The exact V3 parent had
\(1590\) newline-delimited source records, \(55150\) bytes, and SHA--256
digest
\[
\texttt{333f0efcf0d455ed261258b5cdb06e767c2e0a0ee2ae0316f1d337194bec8d08}.
\]
No compact V3 mathematical content was changed.  The sole final
\verb|\end{document}| is restored below.

% =============================================================================
% END APPENDED COMPACT VERSION V4 MATERIAL
% =============================================================================

% The compact V4 document terminator is intentionally disabled here:
% \end{document}

% =============================================================================
% BEGIN APPENDED COMPACT VERSION V5 MATERIAL
% =============================================================================

\chapter{V5: Companion audit and the exact Cartan cumulant polynomial}
\label{ch:V5}

\section{Mandatory companion audit}
\label{sec:V5-audit}

\begin{assessment}[Read-only V5 audit]
\label{ass:V5-audit}
The newest companion sources at the start of V5 were:
\begin{center}
\begin{tabular}{llll}
\toprule
programme & version & records/bytes & SHA--256\\
\midrule
SM--Einstein & V157 f1 & \(98871/3623989\) &
\texttt{ea08ce7537c7327a80b5aa5a653111afaedae82cc18215d87534f17b563aad94}\\
Navier--Stokes & V31 & \(23052/887537\) &
\texttt{f1121b62238ad5491bb7cf03635ea9934942edf573ed8faaa9ee79e1d38a6696}\\
\bottomrule
\end{tabular}
\end{center}
SM V157 proves an algebraic subtraction of the complete invariant local
jet before testing irrelevant-shell summability; it explicitly warns that
componentwise subtraction can violate the Ward row.  NS V31 proves that a
flat probability mark does not remove the weighted first moment and that a
signed polarization hinge has a nontrivial balanced-sector nullspace.
\end{assessment}

\begin{firewall}[Audit transfer boundary]
\label{fire:V5-audit}
No SM renormalization estimate, Ward theorem, NS Duhamel bound, or
regularity assertion is imported.  Two proof-order lessons are retained:
form the complete symmetry-compatible object before taking norms, and do
not treat probability normalization as though it removes a required first
moment.  In the present YM problem this means: compute the complete Cartan
fourth-cumulant scalar before splitting its fifteen partition terms, and
keep the payer factor \(A_U/(1-q_{\alpha,U})\) visible.
\end{firewall}

\section{Deterministic Cartan outputs of every partition block}
\label{sec:V5-Cartan-blocks}

Fix dominant weights \(\lambda_1,\ldots,\lambda_4\).  For
\(\varnothing\ne B\subseteq I:=\{1,2,3,4\}\), put
\[
\lambda_B:=\sum_{i\in B}\lambda_i,
\qquad
q_B:=q_{\alpha,\lambda_B}.
\]
Let \(v_B:=\bigotimes_{i\in B}v_i\), in the fixed row order.

\begin{lemma}[Every block moment is scalar on its Cartan vector]
\label{lem:V5-block-Cartan-scalar}
Let \(M_B\) be the central one-coordinate Wilson moment operator on
\(\bigotimes_{i\in B}H_{\lambda_i}\).  Then
\begin{equation}
\boxed{M_Bv_B=q_Bv_B.}
\label{eq:V5-block-Cartan-scalar}
\end{equation}
For a set partition \(\pi\in\Pi_4\),
\begin{equation}
\left(\bigotimes_{B\in\pi}M_B\right)
(v_1\otimes v_2\otimes v_3\otimes v_4)
=
\left(\prod_{B\in\pi}q_B\right)
(v_1\otimes v_2\otimes v_3\otimes v_4).
\label{eq:V5-partition-Cartan-scalar}
\end{equation}
\end{lemma}

\begin{proof}
By \Cref{lem:V4-Cartan-vector}, \(v_B\) lies in the unique Cartan
irreducible of highest weight \(\lambda_B\).  The Wilson moment is central,
so Schur's lemma makes it scalar on that irreducible, with scalar
\(q_{\alpha,\lambda_B}=q_B\).  This proves the first identity.  Operators
belonging to distinct blocks of \(\pi\) act on disjoint tensor factors;
multiplication of the scalar identities proves the second.
\end{proof}

\section{The complete Cartan fourth-cumulant scalar}
\label{sec:V5-Cartan-polynomial}

\begin{theorem}[Exact Cartan Moebius polynomial]
\label{thm:V5-Cartan-polynomial}
Let \(K_{4,\rm Car}\) be the complete local fourth-cumulant multiplier
restricted to the Cartan component of
\(\bigotimes_{i=1}^4R_{\lambda_i}\).  Then it is scalar and
\begin{equation}
\boxed{
K_{4,\rm Car}
=
\sum_{\pi\in\Pi_4}
(|\pi|-1)!(-1)^{|\pi|-1}
\prod_{B\in\pi}q_B.}
\label{eq:V5-Cartan-polynomial}
\end{equation}
There is no intermediate multiplicity or recoupling coefficient in this
formula.
\end{theorem}

\begin{proof}
Insert the moment--cumulant formula
\eqref{eq:V3-four-moment-cumulant}.  Evaluate each partition moment on
the product highest-weight vector.  Equation
\eqref{eq:V5-partition-Cartan-scalar} supplies its scalar.  The product
vector spans the highest-weight line of the multiplicity-one final Cartan
component, so the central cumulant multiplier has the same scalar on the
whole irreducible by Schur's lemma.
\end{proof}

\begin{corollary}[Exact cancellation in the multiplicative model]
\label{cor:V5-multiplicative-cancellation}
If \(q_B=\prod_{i\in B}q_{\{i\}}\) for every nonempty \(B\subseteq I\),
then \(K_{4,\rm Car}=0\).
\end{corollary}

\begin{proof}
Every product in \eqref{eq:V5-Cartan-polynomial} then equals
\(\prod_{i=1}^4q_{\{i\}}\).  Its remaining coefficient is
\[
\sum_{\pi\in\Pi_4}(|\pi|-1)!(-1)^{|\pi|-1}=0,
\]
the total Moebius coefficient from the bottom to the top of the
nontrivial partition lattice.  Equivalently, all cumulants of order at
least two vanish for a multiplicative moment functional.
\end{proof}

\begin{assessment}[Source of the required Cartan decay]
\label{ass:V5-Cartan-decay-source}
The exact scalar is a defect of multiplicativity of
\(\lambda\mapsto q_{\alpha,\lambda}\).  The marked Cartan gate must
therefore be proved from quantitative mixed finite differences of this
multiplier, with spectator damping retained.  Bounding the fifteen
summands of \eqref{eq:V5-Cartan-polynomial} separately destroys the
cancellation exhibited by
\Cref{cor:V5-multiplicative-cancellation}.
\end{assessment}

\section{An exact cut recursion for the Cartan polynomial}
\label{sec:V5-Cartan-cut}

Define
\begin{align}
\delta_{A,B}&:=q_{A\cup B}-q_Aq_B
\quad(A\cap B=\varnothing),
\label{eq:V5-binary-defect}\\
\kappa_{A,B,C}
&:=q_{A\cup B\cup C}
-q_Aq_{B\cup C}-q_Bq_{A\cup C}-q_Cq_{A\cup B}
+2q_Aq_Bq_C.
\label{eq:V5-ternary-defect}
\end{align}

\begin{proposition}[Singleton-versus-triple Cartan cut identity]
\label{prop:V5-Cartan-cut}
The exact scalar in \eqref{eq:V5-Cartan-polynomial} satisfies
\begin{align}
\delta_{\{1\},\{2,3,4\}}
={}&K_{4,\rm Car}
+q_4\kappa_{\{1\},\{2\},\{3\}}
+q_3\kappa_{\{1\},\{2\},\{4\}}
+q_2\kappa_{\{1\},\{3\},\{4\}}
\notag\\
&+\delta_{\{1\},\{2\}}\delta_{\{3\},\{4\}}
+\delta_{\{1\},\{3\}}\delta_{\{2\},\{4\}}
+\delta_{\{1\},\{4\}}\delta_{\{2\},\{3\}}
\notag\\
&+\delta_{\{1\},\{2\}}q_3q_4
+\delta_{\{1\},\{3\}}q_2q_4
+\delta_{\{1\},\{4\}}q_2q_3.
\label{eq:V5-Cartan-cut}
\end{align}
The three other singleton cuts follow by relabeling.
\end{proposition}

\begin{proof}
Apply the exact product-cumulant identity
\eqref{eq:V1-product-cumulant} to the cut
\(\{1\}\mid\{2,3,4\}\), and evaluate every block cumulant on the Cartan
highest-weight vector using
\Cref{lem:V5-block-Cartan-scalar}.  Exactly ten partitions join this cut
to the top: the top block; three triple-plus-singleton partitions; three
two-pair partitions; and three crossing-pair-plus-two-singleton
partitions.  They give the ten terms on the right.
\end{proof}

\begin{corollary}[A cancellation-preserving acquisition route]
\label{cor:V5-Cartan-route}
The pointwise Cartan gate
\eqref{eq:V4-Cartan-pointwise-gate} can be attacked using one complete
binary defect across a singleton cut, three intact ternary defects, and
six pair-defect products.  All lower-order objects already have
coefficient-sensitive operator estimates in the frozen archive.  What is
not yet proved is that their Cartan specializations retain the same
branching-row marks simultaneously with the local output payer.
\end{corollary}

\begin{proof}
Solve \eqref{eq:V5-Cartan-cut} for \(K_{4,\rm Car}\).  The statement
about available lower-order estimates is the compact V1 legacy ledger.
No norm estimate is asserted here; the last sentence records the exact
remaining compatibility condition.
\end{proof}

\section{V5 ledger and V6 acquisition}
\label{sec:V5-ledger}

\begin{theorem}[Integrated compact-series V5 statement]
\label{thm:V5-integrated}
Preserving compact V1--V4, V5 proves:
\begin{enumerate}[label=\textup{(V5.\arabic*)},leftmargin=6em]
\item the mandatory, hashed SM V157 f1 and NS V31 audit with a strict
      type firewall;
\item deterministic Cartan output of every partition block,
      \eqref{eq:V5-block-Cartan-scalar};
\item the exact fourth-cumulant Moebius polynomial
      \eqref{eq:V5-Cartan-polynomial};
\item exact cancellation for a multiplicative multiplier law; and
\item the cancellation-preserving cut recursion
      \eqref{eq:V5-Cartan-cut}.
\end{enumerate}
\end{theorem}

\begin{proof}
These are
\Cref{ass:V5-audit,fire:V5-audit,
lem:V5-block-Cartan-scalar,thm:V5-Cartan-polynomial,
cor:V5-multiplicative-cancellation,prop:V5-Cartan-cut}.
\end{proof}

\begin{openproblem}[V6 mandate: quantitative Cartan mixed difference]
\label{op:V6-mandate}
\begin{enumerate}
\item Insert the exact Wilson multiplier formula into
      \(\delta_{A,B}\) and \(\kappa_{A,B,C}\); retain signs until the
      complete finite difference has been formed.
\item On each singleton cut, combine
      \(\delta_{\{i\},I\setminus\{i\}}\) with the local payer
      \(A_{\rm Car}/(1-q_{\alpha,\rm Car})\) before applying a norm.
\item Use the restored global ternary estimate only on the intact
      \(\kappa_{A,B,C}\) terms in \eqref{eq:V5-Cartan-cut}.
\item Partition by the degree-two or degree-three vertices of the fixed
      tree; do not charge a leaf with an unnecessary repeated mark.
\item Prove one complete Cartan star or path bound, or identify the exact
      mixed difference whose available Wilson derivative order is too
      low.
\item The next companion audit is compact V10.
\end{enumerate}
\end{openproblem}

\begin{firewall}[Claim boundary after compact V5]
\label{fire:V5-current-boundary}
V5 computes the exact Cartan fourth-cumulant scalar and its
cancellation-preserving lower-order recursion.  It does not yet prove the
quantitative marked Cartan inequality, the non-Cartan capacity tail, the
full embedded \([4]\) sector, collision aggregation, \([3,3]\), quartic
or all-order MTA, continuum Yang--Mills existence, or a positive mass gap.
\end{firewall}

\section{Append-only record for compact V5}
\label{sec:V5-record}

V5 was appended to the complete compact V4 file.  The exact V4 parent had
\(1909\) newline-delimited source records, \(66293\) bytes, and SHA--256
digest
\[
\texttt{ec4931854f06a772604a54216ba34e0ba96f8169bec3675992df054f29b23a1b}.
\]
No compact V4 mathematical content was changed.  The sole final
\verb|\end{document}| is restored below.

% =============================================================================
% END APPENDED COMPACT VERSION V5 MATERIAL
% =============================================================================

% The compact V5 document terminator is intentionally disabled here:
% \end{document}

% =============================================================================
% BEGIN APPENDED COMPACT VERSION V6 MATERIAL
% =============================================================================

\chapter{V6: Highest-weight random variables and the exact global
Cartan word}
\label{ch:V6}

\section{The actual Wilson probability space}
\label{sec:V6-Wilson-space}

Let \(\nu_\alpha\) be the normalized one-link Wilson probability measure
on \(G\),
\[
\mathrm d\nu_\alpha(U)
=Z_\alpha^{-1}
\exp\!\left\{\frac{\alpha}{3}\operatorname{ReTr}U\right\}\mathrm dU.
\]
For an irreducible representation of highest weight \(\lambda\), choose a
unit highest-weight vector \(v_\lambda\) and define
\begin{equation}
z_\lambda(U):=
\langle v_\lambda,D^\lambda(U)v_\lambda\rangle.
\label{eq:V6-highest-coefficient}
\end{equation}
Then \(|z_\lambda(U)|\le1\).

\begin{lemma}[Cartan multiplication of highest-weight coefficients]
\label{lem:V6-Cartan-multiplication}
For dominant weights \(\lambda,\mu\),
\begin{equation}
\boxed{z_{\lambda+\mu}(U)=z_\lambda(U)z_\mu(U).}
\label{eq:V6-Cartan-multiplication}
\end{equation}
More generally,
\[
z_{\sum_{i\in B}\lambda_i}(U)
=\prod_{i\in B}z_{\lambda_i}(U).
\]
\end{lemma}

\begin{proof}
The Cartan component of
\(R_\lambda\otimes R_\mu\) contains the highest-weight vector
\(v_\lambda\otimes v_\mu\), by
\Cref{lem:V4-Cartan-vector}.  Its matrix coefficient under the tensor
product action is
\[
\langle v_\lambda\otimes v_\mu,
(D^\lambda(U)\otimes D^\mu(U))
(v_\lambda\otimes v_\mu)\rangle
=z_\lambda(U)z_\mu(U).
\]
The restricted action is \(D^{\lambda+\mu}\), proving the first formula.
Iterate it for a finite set \(B\).
\end{proof}

\begin{lemma}[Wilson multiplier as a highest-weight expectation]
\label{lem:V6-multiplier-expectation}
For every dominant weight \(\lambda\),
\begin{equation}
\boxed{
q_{\alpha,\lambda}
=\mathbb E_{\nu_\alpha}z_\lambda.}
\label{eq:V6-multiplier-expectation}
\end{equation}
Consequently, in the notation of V5,
\begin{equation}
q_B
=
\mathbb E_{\nu_\alpha}
\prod_{i\in B}z_{\lambda_i}.
\label{eq:V6-block-moment}
\end{equation}
\end{lemma}

\begin{proof}
The density of \(\nu_\alpha\) is central.  Therefore
\[
\int_GD^\lambda(U)\,\mathrm d\nu_\alpha(U)
=q_{\alpha,\lambda}I
\]
by Schur's lemma and the definition of the central Wilson multiplier.
Take its matrix element in \(v_\lambda\).  The second identity follows
from \Cref{lem:V6-Cartan-multiplication}.
\end{proof}

\section{The Cartan scalar is one genuine fourth cumulant}
\label{sec:V6-genuine-cumulant}

Put \(Z_i:=z_{\lambda_i}(U)\) on the single probability space
\((G,\nu_\alpha)\).

\begin{theorem}[Probabilistic representation of the exact Cartan
cumulant]
\label{thm:V6-Cartan-cumulant}
The scalar computed in \eqref{eq:V5-Cartan-polynomial} is exactly
\begin{equation}
\boxed{
K_{4,\rm Car}
=\kappa_{\nu_\alpha}(Z_1,Z_2,Z_3,Z_4).}
\label{eq:V6-Cartan-cumulant}
\end{equation}
Equivalently, writing
\(\widetilde Z_i:=Z_i-\mathbb EZ_i\),
\begin{align}
K_{4,\rm Car}
={}&
\mathbb E\prod_{i=1}^4\widetilde Z_i
-\mathbb E(\widetilde Z_1\widetilde Z_2)
 \mathbb E(\widetilde Z_3\widetilde Z_4)
\notag\\
&-\mathbb E(\widetilde Z_1\widetilde Z_3)
 \mathbb E(\widetilde Z_2\widetilde Z_4)
-\mathbb E(\widetilde Z_1\widetilde Z_4)
 \mathbb E(\widetilde Z_2\widetilde Z_3).
\label{eq:V6-centered-fourth}
\end{align}
\end{theorem}

\begin{proof}
The joint moment of a block \(B\) is \eqref{eq:V6-block-moment}.
Substitution into the ordinary moment--cumulant formula gives precisely
\eqref{eq:V5-Cartan-polynomial}.  The centered identity is the standard
fourth joint-cumulant identity, obtained by expanding the four centered
variables and cancelling singleton partitions.
\end{proof}

\begin{corollary}[The Gaussian pairings are already subtracted]
\label{cor:V6-Gaussian-subtraction}
Any estimate of
\(\mathbb E\prod_i\widetilde Z_i\) which is followed by a triangle
inequality before subtracting the three covariance products generally
loses the cancellation responsible for the fourth connected scale.
The correct object is the complete right side of
\eqref{eq:V6-centered-fourth}.
\end{corollary}

\begin{proof}
For a centered jointly Gaussian vector the fourth moment is exactly the
sum of the three pairings, so the right side vanishes while the separate
absolute values need not.  This is an algebraic regression; it does not
assume that \(\nu_\alpha\) is Gaussian.
\end{proof}

\section{Independent-copy form of every binary Cartan defect}
\label{sec:V6-independent-copy}

Let \(U'\) be an independent copy of \(U\), and write
\[
Z_A(U):=\prod_{i\in A}z_{\lambda_i}(U).
\]

\begin{lemma}[Exact symmetrized defect]
\label{lem:V6-symmetrized-defect}
For disjoint nonempty \(A,B\),
\begin{equation}
\boxed{
\delta_{A,B}
=\frac12\mathbb E
\left[
(Z_A(U)-Z_A(U'))
(Z_B(U)-Z_B(U'))
\right].}
\label{eq:V6-symmetrized-defect}
\end{equation}
\end{lemma}

\begin{proof}
Expand the product.  The two diagonal expectations both equal
\(\mathbb E(Z_AZ_B)=q_{A\cup B}\), while independence makes each cross
expectation equal \(q_Aq_B\).  Division by two gives
\(\delta_{A,B}=q_{A\cup B}-q_Aq_B\).
\end{proof}

\begin{remark}[What a derivative estimate must act on]
\label{rem:V6-defect-derivative}
Equation \eqref{eq:V6-symmetrized-defect} shows that a binary bridge is a
product of two actual increments under the same resampling
\((U,U')\).  Iterating a normed covariance estimate would reuse that
resampling and can falsely count two independent bridges.  This is the
Cartan-line version of the quotient-rank obstruction in legacy V133.
\end{remark}

\section{Exact spectator factorization of a global Cartan word}
\label{sec:V6-global-word}

Let \(X\) be a finite coordinate set and fix one junction coordinate
\(e_0\).  At every \(f\ne e_0\), let \(\pi_f\) be the retained local
moment partition of the active rows.  Write
\[
\lambda_{B,f}:=\sum_{i\in B}\lambda_{i,f},
\qquad
q_{B,f}:=q_{\alpha,\lambda_{B,f}}.
\]
The global Cartan vector is the tensor product of all row-coordinate
highest-weight vectors.

\begin{theorem}[Complete scalar factorization before normalization]
\label{thm:V6-global-Cartan-word}
For the unique-\([4]\) coordinate word having the complete fourth
cumulant at \(e_0\) and the retained moment partition \(\pi_f\) at every
other coordinate,
\begin{equation}
\boxed{
K_{\rm word,Car}
=
\kappa_{\nu_\alpha}
\bigl(
z_{\lambda_{1,e_0}},
z_{\lambda_{2,e_0}},
z_{\lambda_{3,e_0}},
z_{\lambda_{4,e_0}}
\bigr)
\prod_{f\ne e_0}\prod_{B\in\pi_f}q_{B,f}.}
\label{eq:V6-global-Cartan-word}
\end{equation}
This equality is taken before any absolute value, output sum, or
resolvent estimate.
\end{theorem}

\begin{proof}
The local junction scalar is
\Cref{thm:V6-Cartan-cumulant}.  At a spectator coordinate \(f\), apply
\Cref{lem:V5-block-Cartan-scalar} to every block of \(\pi_f\); their
product is \(\prod_{B\in\pi_f}q_{B,f}\).  Different coordinates act on
tensor factors, so their scalars multiply.  Multiplicity-one of every
Cartan block removes all intermediate path sums.
\end{proof}

\begin{corollary}[Exact normalized Cartan word gate]
\label{cor:V6-normalized-word-gate}
Let \(\boldsymbol U_{\rm Car}\) be the coordinatewise global Cartan
output.  The contribution of the word in
\eqref{eq:V6-global-Cartan-word} to the normalized physical carrier is
the single nonnegative scalar
\begin{align}
\mathcal C_{\rm word}
:={}&
\frac{\Xi(\boldsymbol U_{\rm Car})}
     {1-q_{\alpha,\boldsymbol U_{\rm Car}}}
\left|
\kappa_{\nu_\alpha}
\bigl(
z_{\lambda_{1,e_0}},z_{\lambda_{2,e_0}},
z_{\lambda_{3,e_0}},z_{\lambda_{4,e_0}}
\bigr)
\right|
\notag\\
&\times
\prod_{f\ne e_0}\prod_{B\in\pi_f}|q_{B,f}|.
\label{eq:V6-normalized-Cartan-word}
\end{align}
Because the global Cartan projector has product-state capacity one,
the required marked estimate for this word is exactly
\begin{equation}
\mathcal C_{\rm word}
\le C\sum_{\tau,\ell}
\mathcal W_{\tau,\ell}(\boldsymbol R).
\label{eq:V6-scalar-Cartan-gate}
\end{equation}
\end{corollary}

\begin{proof}
Insert \eqref{eq:V6-global-Cartan-word} into the definition of the
physical carrier.  Apply \Cref{thm:V4-Cartan-capacity}.  No inequality
other than the final absolute value has been used.
\end{proof}

\begin{assessment}[Precise gain over V5]
\label{ass:V6-gain}
The unresolved Cartan problem is no longer an operator-valued
Clebsch--Gordan estimate.  It is the explicit scalar inequality
\eqref{eq:V6-scalar-Cartan-gate}.  All global input-mass decay must come
from the exact spectator product in
\eqref{eq:V6-normalized-Cartan-word}, while the junction scale must be
estimated on the complete centered combination
\eqref{eq:V6-centered-fourth}.  These two factors must be combined
before the payer-coordinate sum of V2.
\end{assessment}

\section{V6 ledger and V7 acquisition}
\label{sec:V6-ledger}

\begin{theorem}[Integrated compact-series V6 statement]
\label{thm:V6-integrated}
Preserving compact V1--V5, V6 proves:
\begin{enumerate}[label=\textup{(V6.\arabic*)},leftmargin=6em]
\item Cartan multiplication of highest-weight coefficients,
      \eqref{eq:V6-Cartan-multiplication};
\item the exact Wilson expectation
      \eqref{eq:V6-multiplier-expectation};
\item the genuine fourth-cumulant representation
      \eqref{eq:V6-Cartan-cumulant};
\item the centered Gaussian-pairing subtraction
      \eqref{eq:V6-centered-fourth};
\item the independent-copy binary defect
      \eqref{eq:V6-symmetrized-defect}; and
\item the complete global Cartan word and normalized scalar gate
      \eqref{eq:V6-global-Cartan-word}--%
      \eqref{eq:V6-scalar-Cartan-gate}.
\end{enumerate}
\end{theorem}

\begin{proof}
These are
\Cref{lem:V6-Cartan-multiplication,lem:V6-multiplier-expectation,
thm:V6-Cartan-cumulant,lem:V6-symmetrized-defect,
thm:V6-global-Cartan-word,cor:V6-normalized-word-gate}.
\end{proof}

\begin{openproblem}[V7 mandate: combine spectator damping with the
connected Cartan scale]
\label{op:V7-mandate}
\begin{enumerate}
\item Apply the established Wilson product-moment debits directly to
      the spectator product in
      \eqref{eq:V6-normalized-Cartan-word}, grouped by original row and
      retained block; do not drop its partition structure.
\item Estimate the complete centered expression
      \eqref{eq:V6-centered-fourth}, not its four terms separately.
\item For a fixed tree, use spectator moments only at vertices of degree
      at least two; leaves require no inverse global mass.
\item Split local versus spectator mass at each internal vertex and prove
      the star case before the path case.
\item Retain the sole output payer through
      \eqref{eq:V2-resolvent-allocation}.
\item Prove \eqref{eq:V6-scalar-Cartan-gate} for one complete
      private-spectator priority branch, or give an explicit scalar
      scaling which violates the proposed allocation.
\item The next companion audit is compact V10.
\end{enumerate}
\end{openproblem}

\begin{firewall}[Claim boundary after compact V6]
\label{fire:V6-current-boundary}
V6 removes final representation theory from the Cartan word and reduces
it to an exact scalar Wilson cumulant times exact spectator multipliers.
It does not yet prove the marked scalar inequality, the non-Cartan
capacity tail, the full embedded \([4]\) sector, normalized collision
aggregation, \([3,3]\), quartic or all-order MTA, continuum Yang--Mills
existence, or a positive mass gap.
\end{firewall}

\section{Append-only record for compact V6}
\label{sec:V6-record}

V6 was appended to the complete compact V5 file.  The exact V5 parent had
\(2191\) newline-delimited source records, \(76844\) bytes, and SHA--256
digest
\[
\texttt{fb047187fbc744ba95f76a65c5e0d0c671c22d2ebf9d1e659ab5f1c32643a164}.
\]
No compact V5 mathematical content was changed.  The sole final
\verb|\end{document}| is restored below.

% =============================================================================
% END APPENDED COMPACT VERSION V6 MATERIAL
% =============================================================================

% The compact V6 document terminator is intentionally disabled here:
% \end{document}

% =============================================================================
% BEGIN APPENDED COMPACT VERSION V7 MATERIAL
% =============================================================================

\chapter{V7: Post-correction closure of the private-dominant Cartan
four-junction sector}
\label{ch:V7}

\section{A finite Bernoulli compiler for Wilson product moments}
\label{sec:V7-product-moments}

The compact proof needs only the already-established one-link Wilson
debits through order three.  The following elementary lemma records
exactly how they combine without a support count.

\begin{lemma}[Product-moment compiler]
\label{lem:V7-product-moment-compiler}
Let \(F\) be finite, \(x_f\ge0\), and \(0\le r_f\le1\).  Fix
\(m\in\mathbb N\).  Suppose that for every \(1\le k\le m\),
\begin{equation}
x_f^k r_f\le C_k(1-r_f)
\quad(f\in F).
\label{eq:V7-one-site-debits}
\end{equation}
Then
\begin{equation}
\boxed{
\left(\sum_{f\in F}x_f\right)^m
\prod_{f\in F}r_f
\le C_m^{\rm prod},}
\label{eq:V7-product-moment}
\end{equation}
where \(C_m^{\rm prod}\) depends only on \(m,C_1,\ldots,C_m\), not on
\(|F|\).
\end{lemma}

\begin{proof}
Expand the \(m\)-th power and group a monomial by its distinct indices
\(f_1,\ldots,f_j\) and positive multiplicities
\(k_1+\cdots+k_j=m\).  After multiplication by
\(\prod_fr_f\), apply \eqref{eq:V7-one-site-debits} at each distinct
index:
\[
\prod_{\ell=1}^j x_{f_\ell}^{k_\ell}r_{f_\ell}
\prod_{f\notin\{f_1,\ldots,f_j\}}r_f
\le
\left(\prod_{\ell=1}^jC_{k_\ell}\right)
\prod_{\ell=1}^j(1-r_{f_\ell})
\prod_{f\notin\{f_1,\ldots,f_j\}}r_f.
\]
For fixed \(j\), the sum over unordered distinct indices is the
probability of exactly \(j\) failures in independent Bernoulli trials
with success probabilities \(r_f\), hence is at most one.  The number of
ordered multiplicity patterns and their multinomial coefficients depends
only on \(m\).  Summing \(1\le j\le m\) proves the claim.
\end{proof}

\begin{corollary}[Private Wilson products through cubic order]
\label{cor:V7-private-products}
Let \(P\) be any finite exact private coordinate bank, set
\[
Y:=\sum_{f\in P}(1+C_2(R_f)),
\qquad
q_P:=\prod_{f\in P}\eta_{R_f}(\alpha),
\]
and suppose \(Y>0\).  The legacy one-link linear, quadratic, and cubic
tail debits imply
\begin{equation}
\boxed{
(s_\alpha Y)^kq_P\le C_k
\qquad(k=1,2,3),}
\label{eq:V7-private-product-debits}
\end{equation}
uniformly in the size of \(P\), all representations, and \(\alpha\) in
the established weak-coupling range.
\end{corollary}

\begin{proof}
Put
\(x_f=s_\alpha(1+C_2(R_f))\) and \(r_f=\eta_{R_f}(\alpha)\).
The exact-support Casimir floor absorbs the harmless added one into the
legacy one-link debits.  Apply
\Cref{lem:V7-product-moment-compiler} for \(m=1,2,3\).
\end{proof}

\section{The exact unique-junction carrier}
\label{sec:V7-junction-carrier}

Let \(X_i\) be the four exact supports, and assume their only common
coordinate is \(e\):
\[
X_1\cap X_2\cap X_3\cap X_4=\{e\}.
\]
Put
\[
P_i:=X_i\setminus\bigcup_{j\ne i}X_j,
\qquad
P:=\mathop{\dot\bigcup}_{i=1}^4P_i,
\]
and use \(Y,q_P\) from
\Cref{cor:V7-private-products}.  For a tree
\(\tau\in\mathfrak T_4\), write
\[
d_i:=\deg_\tau(i),\quad
a_i:=1+C_2(R_{i,e}),\quad
\Xi_i:=\Xi(\boldsymbol R_i),
\]
and
\begin{equation}
w_{\tau,e}:=
s_\alpha^3\prod_{i=1}^4a_i^{d_i/2}.
\label{eq:V7-local-tree-weight}
\end{equation}
The degree identity is
\begin{equation}
\sum_{i=1}^4(d_i-1)=2.
\label{eq:V7-degree-excess}
\end{equation}

\begin{lemma}[Exact factorization and operator bound]
\label{lem:V7-exact-junction-factorization}
The complete coordinate-word operator with a four-row cumulant at \(e\)
and complete moment spectators elsewhere has the tree refinement
\begin{equation}
\mathbf J_e=\sum_{\tau\in\mathfrak T_4}\mathbf J_{\tau,e},
\qquad
\mathbf J_{\tau,e}
=q_P\mathbf K_{\tau,e}
\otimes
\bigotimes_{\substack{f\ne e\\|I_f|\ge2}}\mathbf Q_{I_f,f},
\label{eq:V7-junction-factorization}
\end{equation}
where each shared spectator moment is a contraction and
\begin{equation}
\boxed{
\opnorm{\mathbf J_{\tau,e}}
\le Cq_Pw_{\tau,e}.}
\label{eq:V7-junction-operator}
\end{equation}
\end{lemma}

\begin{proof}
At a private coordinate only one row is active, so its complete moment is
the scalar \(\eta_{R_f}\); their product is \(q_P\).  At a shared
spectator coordinate, the sum of all local cumulant partitions is the
complete moment
\[
\mathbf Q_{I_f,f}
=\mathbb E_{\nu_\alpha}
\bigotimes_{i\in I_f}D^{R_{i,f}}(U_f),
\]
which is an average of unitaries and hence a contraction.  Independence
of coordinates gives the tensor product.  Insert the V41 sixteen-tree
decomposition at \(e\) and its local tree estimate to obtain
\eqref{eq:V7-junction-operator}.
\end{proof}

\section{Private-dominant scalar payment}
\label{sec:V7-private-payment}

Let \(\boldsymbol T\) be a nontrivial resolved product output and write
\(\Xi_T:=\Xi(\boldsymbol T)\).  The established two-sided product gap
gives
\begin{equation}
\frac{\Xi_T}{1-q_{\alpha,\boldsymbol T}}
\le C
\begin{cases}
s_\alpha^{-1},&s_\alpha\Xi_T\le c_0,\\
\Xi_T,&s_\alpha\Xi_T>c_0.
\end{cases}
\label{eq:V7-output-resolvent-split}
\end{equation}

\begin{theorem}[Corrected private-dominant scalar theorem]
\label{thm:V7-private-scalar}
Fix \(D\ge1\).  Suppose \(Y>0\) and
\begin{equation}
\Xi_T\le DY,
\qquad
\prod_{i=1}^4\Xi_i^{d_i-1}\le DY^2.
\label{eq:V7-private-dominance}
\end{equation}
Then
\begin{equation}
\boxed{
\frac{\Xi_T}{1-q_{\alpha,\boldsymbol T}}
q_Pw_{\tau,e}
\le
C_D
\prod_{i=1}^4a_i^{d_i}\Xi_i^{1-d_i}.}
\label{eq:V7-private-scalar}
\end{equation}
\end{theorem}

\begin{proof}
On \(s_\alpha\Xi_T\le c_0\), use the first line of
\eqref{eq:V7-output-resolvent-split}, the quadratic product debit in
\eqref{eq:V7-private-product-debits}, and
\eqref{eq:V7-local-tree-weight}:
\[
\frac{\Xi_T}{1-q_{\alpha,\boldsymbol T}}
q_Pw_{\tau,e}
\le
Cq_Ps_\alpha^2\prod_i a_i^{d_i/2}
\le
\frac{C}{Y^2}\prod_i a_i^{d_i/2}.
\]
The second condition in \eqref{eq:V7-private-dominance} gives
\[
Y^{-2}
\le
D\prod_i\Xi_i^{1-d_i}.
\]
Since \(a_i\ge1\) and \(d_i\ge1\),
\(a_i^{d_i/2}\le a_i^{d_i}\).  This proves the desired bound on the
low-output branch.

On \(s_\alpha\Xi_T>c_0\), use the second line of
\eqref{eq:V7-output-resolvent-split}, the first condition in
\eqref{eq:V7-private-dominance}, and the cubic debit:
\begin{align*}
\frac{\Xi_T}{1-q_{\alpha,\boldsymbol T}}
q_Pw_{\tau,e}
&\le
CDYq_Ps_\alpha^3\prod_i a_i^{d_i/2}\\
&=
\frac{CD}{Y^2}(s_\alpha Y)^3q_P
\prod_i a_i^{d_i/2}\\
&\le
\frac{C_D}{Y^2}\prod_i a_i^{d_i/2}.
\end{align*}
Apply the same branching-mass and local-Casimir comparisons.
\end{proof}

\section{Coefficient-sensitive post-V45 closure}
\label{sec:V7-coefficient-closure}

\begin{theorem}[Private-dominant unique-\([4]\) MTA sector is
revalidated]
\label{thm:V7-private-MTA}
For every tree and junction coordinate, restrict the physical output
resolution to the blocks satisfying
\eqref{eq:V7-private-dominance}.  The resulting unique-junction
contribution obeys its complete \(m=4\) marked-tree atlas term for
arbitrary trace-class Fourier coefficients, uniformly in support,
volume, representations, output multiplicities, and \(\alpha\).
This theorem uses the corrected exact Fourier factor
\(d_{\rm in}/d_{\boldsymbol T}\) and raw pinching; it uses no
representation-dimension formation law.
\end{theorem}

\begin{proof}
On each retained orthogonal physical block,
\Cref{lem:V7-exact-junction-factorization} and
\Cref{thm:V7-private-scalar} give
\[
\frac{\Xi_T}{1-q_{\alpha,\boldsymbol T}}
\opnorm{\mathbf J_{\tau,e;\boldsymbol T}}
\le
C_D\prod_i a_i^{d_i}\Xi_i^{1-d_i}.
\]
The right side is exactly the fixed-representation marked tree weight:
\begin{align*}
\prod_i a_i^{d_i}\Xi_i^{1-d_i}
&=
\left(\prod_i\Xi_i\right)
\prod_{ij\in E(\tau)}
\frac{a_i}{\Xi_i}\frac{a_j}{\Xi_j}.
\end{align*}
Thus the pointwise sufficient hypothesis of
\Cref{cor:V1-pointwise-sufficient} holds on this restricted carrier.
That corollary uses \eqref{eq:V1-exact-Fourier}, cancels the BFA output
dimension against \(d_{\rm in}/d_{\boldsymbol T}\), and sums all raw
isotypic/multiplicity blocks by
\Cref{thm:V1-weighted-pinching-capacity}.  Restricting to an orthogonal
output family can only decrease the positive carrier.  Hence there is no
output or multiplicity count, and the marked-tree estimate follows.
\end{proof}

\begin{corollary}[Exact complement after the revalidated branch]
\label{cor:V7-private-complement}
An unpaid unique-junction physical block must satisfy at least one of
\begin{equation}
Y=0,\qquad
\Xi_T>DY,\qquad
\prod_i\Xi_i^{d_i-1}>DY^2.
\label{eq:V7-private-complement}
\end{equation}
\end{corollary}

\begin{proof}
This is the logical complement of the three hypotheses of
\Cref{thm:V7-private-scalar}.
\end{proof}

\begin{assessment}[Nonduplication and corrected status]
\label{ass:V7-nonduplication}
Legacy V42 contained the scalar mechanism, but its surrounding V42--V45
quartic conclusions were later placed behind a correction firewall when
the coefficient-formation law failed.  V7 rederives the scalar bound and
then supplies the missing post-correction coefficient step using the
exact V46 normalization recast as the compact V1 capacity theorem.
Accordingly, the private-dominant unique-\([4]\) sector is now a durable
closed row.  The complement \eqref{eq:V7-private-complement}, not the full
embedded sector, remains open.
\end{assessment}

\section{V7 ledger and V8 acquisition}
\label{sec:V7-ledger}

\begin{theorem}[Integrated compact-series V7 statement]
\label{thm:V7-integrated}
Preserving compact V1--V6, V7 proves:
\begin{enumerate}[label=\textup{(V7.\arabic*)},leftmargin=6em]
\item the finite Bernoulli product-moment compiler
      \eqref{eq:V7-product-moment};
\item support-uniform private product debits through order three;
\item the exact unique-junction factorization
      \eqref{eq:V7-junction-factorization};
\item the private-dominant scalar estimate
      \eqref{eq:V7-private-scalar};
\item post-V45 coefficient-sensitive closure of the corresponding
      normalized \([4]\) sector; and
\item the exact three-way complement
      \eqref{eq:V7-private-complement}.
\end{enumerate}
\end{theorem}

\begin{proof}
These are
\Cref{lem:V7-product-moment-compiler,cor:V7-private-products,
lem:V7-exact-junction-factorization,thm:V7-private-scalar,
thm:V7-private-MTA,cor:V7-private-complement}.
\end{proof}

\begin{openproblem}[V8 mandate: output-dominant and
shared-branching complements]
\label{op:V8-mandate}
\begin{enumerate}
\item On \(\Xi_T>DY\), allocate the sole resolvent to the coordinate
      carrying a fixed fraction of the output mass, using
      \eqref{eq:V2-carrier-allocation}.
\item If that coordinate is the junction, combine its payer with the
      exact Cartan cumulant \eqref{eq:V6-Cartan-cumulant} before norming.
\item If it is a shared spectator, use the one-sided capacity tensor law
      \eqref{eq:V2-one-sided-tensor}; do not condition the remaining core
      into a falsely product state.
\item On
      \(\prod_i\Xi_i^{d_i-1}>DY^2\), identify an internal tree vertex
      whose off-junction shared mass carries a fixed fraction of its
      required inverse mark.
\item Close one full output-dominant or shared-branching priority family,
      or exhibit the exact payer capacity still lacking.
\item The next mandatory companion audit is compact V10.
\end{enumerate}
\end{openproblem}

\begin{firewall}[Claim boundary after compact V7]
\label{fire:V7-current-boundary}
V7 rigorously and coefficient-sensitively closes the private-dominant
unique-junction \([4]\) sector.  It does not close its output-dominant or
shared-branching complements, the non-Cartan tail, multiple junctions,
normalized collision aggregation, \([3,3]\), global quartic or all-order
MTA, continuum Yang--Mills existence, or a positive mass gap.
\end{firewall}

\section{Append-only record for compact V7}
\label{sec:V7-record}

V7 was appended to the complete compact V6 file.  The exact V6 parent had
\(2554\) newline-delimited source records, \(88921\) bytes, and SHA--256
digest
\[
\texttt{a4364bd88a9949d45c9252fb305f155f15dde60bd68ca4235da1badddb0c5a60}.
\]
No compact V6 mathematical content was changed.  The sole final
\verb|\end{document}| is restored below.

% =============================================================================
% END APPENDED COMPACT VERSION V7 MATERIAL
% =============================================================================

% The compact V7 document terminator is intentionally disabled here:
% \end{document}

% =============================================================================
% BEGIN APPENDED COMPACT VERSION V8 MATERIAL
% =============================================================================

\chapter{V8: Exact localization of every residual unique-junction block}
\label{ch:V8}

\section{Internal-vertex localization from the tree degree identity}
\label{sec:V8-internal-localization}

Fix a labeled four-vertex tree \(\tau\), put
\[
d_i:=\deg_\tau(i),
\qquad
J_\tau:=\{i:d_i\ge2\},
\]
and recall
\[
\sum_{i=1}^4(d_i-1)=2.
\]
Only vertices in \(J_\tau\) require inverse input-mass powers.

\begin{lemma}[A failed private branching payment selects a large
internal row]
\label{lem:V8-large-internal-row}
Let \(Y>0\), \(D\ge1\), and suppose
\begin{equation}
\prod_{i=1}^4\Xi_i^{d_i-1}>DY^2.
\label{eq:V8-failed-branching}
\end{equation}
Then some \(i\in J_\tau\) satisfies
\begin{equation}
\boxed{\Xi_i>\sqrt D\,Y.}
\label{eq:V8-large-internal-row}
\end{equation}
\end{lemma}

\begin{proof}
If every internal vertex satisfied
\(\Xi_i\le\sqrt D\,Y\), then
\[
\prod_i\Xi_i^{d_i-1}
\le(\sqrt D\,Y)^{\sum_i(d_i-1)}
=DY^2,
\]
contradicting \eqref{eq:V8-failed-branching}.
\end{proof}

\begin{remark}[Star and path forms]
\label{rem:V8-star-path}
For a star, \(J_\tau\) is its center and
\eqref{eq:V8-failed-branching} is exactly
\(\Xi_c^2>DY^2\).  For a path, \(J_\tau\) consists of its two internal
vertices and the condition is
\(\Xi_u\Xi_v>DY^2\).  Thus the lemma introduces no leaf denominator.
\end{remark}

\section{Local, private, and shared input mass}
\label{sec:V8-input-split}

At the junction coordinate \(e\), write
\[
a_i:=1+C_2(R_{i,e}).
\]
For row \(i\), let
\[
p_i:=\sum_{f\in P_i}(1+C_2(R_{i,f})),
\qquad
h_i:=\sum_{\substack{f\in X_i\setminus\{e\}\\|I_f|\ge2}}
(1+C_2(R_{i,f})).
\]
Then
\begin{equation}
\Xi_i=a_i+p_i+h_i,
\qquad
0\le p_i\le Y.
\label{eq:V8-row-mass-split}
\end{equation}

\begin{lemma}[Large internal row is local- or shared-dominant]
\label{lem:V8-local-shared-dichotomy}
Assume \(D\ge16\) and let \(i\in J_\tau\) satisfy
\eqref{eq:V8-large-internal-row}.  Then at least one of
\begin{equation}
\boxed{
a_i\ge\frac14\Xi_i,
\qquad
h_i\ge\frac12\Xi_i}
\label{eq:V8-local-shared-dichotomy}
\end{equation}
holds.
\end{lemma}

\begin{proof}
Since
\(\Xi_i>\sqrt D\,Y\ge4Y\), one has
\(p_i\le Y<\Xi_i/4\).  If
\(a_i\ge\Xi_i/4\), the first alternative holds.  Otherwise
\[
h_i=\Xi_i-a_i-p_i
>\Xi_i-\frac14\Xi_i-\frac14\Xi_i
=\frac12\Xi_i.
\]
\end{proof}

\begin{corollary}[Complete branching complement]
\label{cor:V8-branching-complement}
For \(D\ge16\), every block satisfying the failed branching condition
\eqref{eq:V8-failed-branching} has an internal tree vertex which is
either junction-local-dominant or off-junction-shared-dominant in the
precise sense of \eqref{eq:V8-local-shared-dichotomy}.
\end{corollary}

\begin{proof}
Combine
\Cref{lem:V8-large-internal-row,lem:V8-local-shared-dichotomy}.
\end{proof}

\section{Private output mass cannot dominate the output complement}
\label{sec:V8-output-localization}

Decompose a resolved product output mass into junction, private, and
shared-coordinate parts:
\begin{equation}
\Xi_T=\Xi_{T,e}+\Xi_{T,P}+\Xi_{T,S}.
\label{eq:V8-output-split}
\end{equation}
At a private coordinate only one input row is active, so the output
representation equals that input representation.  Hence
\begin{equation}
\Xi_{T,P}=Y
\label{eq:V8-private-output-exact}
\end{equation}
with the same balanced-mass convention used in V7.

\begin{lemma}[Output-dominant blocks have a nonprivate output payer]
\label{lem:V8-nonprivate-output}
If \(D>1\) and
\begin{equation}
\Xi_T>DY,
\label{eq:V8-output-dominant}
\end{equation}
then
\begin{equation}
\boxed{
\Xi_{T,e}+\Xi_{T,S}
>
\left(1-\frac1D\right)\Xi_T,}
\label{eq:V8-nonprivate-output-fraction}
\end{equation}
and therefore
\begin{equation}
\frac{\Xi_T}{1-q_{\alpha,\boldsymbol T}}
\le
\frac{D}{D-1}
\frac{\Xi_{T,e}+\Xi_{T,S}}
     {1-q_{\alpha,\boldsymbol T}}.
\label{eq:V8-nonprivate-resolvent}
\end{equation}
\end{lemma}

\begin{proof}
Equations
\eqref{eq:V8-output-split}--%
\eqref{eq:V8-private-output-exact} give
\[
\Xi_{T,e}+\Xi_{T,S}=\Xi_T-Y.
\]
The hypothesis says \(Y<\Xi_T/D\), proving
\eqref{eq:V8-nonprivate-output-fraction}.  Rearrangement gives
\eqref{eq:V8-nonprivate-resolvent}; the denominator is positive on the
nontrivial output sector.
\end{proof}

\begin{corollary}[No coordinatewise private-payer sum is needed]
\label{cor:V8-no-private-payer-sum}
On the output-dominant complement
\eqref{eq:V8-output-dominant}, the sole resolvent may be allocated only
to the junction coordinate and shared spectator coordinates, at the
fixed cost \(D/(D-1)\).  In particular, no sum of
\(A_{T_f}/(1-q_{\alpha,T_f})\) over private coordinates is introduced.
\end{corollary}

\begin{proof}
Use \eqref{eq:V8-nonprivate-resolvent}, then apply the blockwise
one-resolvent allocation
\eqref{eq:V2-carrier-allocation} only to the two nonprivate terms.
\end{proof}

\begin{remark}[Why bankwise treatment matters]
\label{rem:V8-bankwise-private}
A coordinatewise private allocation would produce a sum of
\((1-q_f)^{-1}\) terms and can lose the size of the private bank.  The
exact private product instead obeys
\[
s_\alpha Yq_P\le C(1-q_P)
\]
by the one-link debit and the telescoping Bernoulli identity.  V8 avoids
the unnecessary coordinatewise loss altogether on
\eqref{eq:V8-output-dominant}.
\end{remark}

\section{Exhaustive residual atlas}
\label{sec:V8-residual-atlas}

\begin{theorem}[Every unpaid unique-junction block has a nonprivate
certificate]
\label{thm:V8-residual-atlas}
Fix \(D\ge16\).  After removing the private-dominant sector closed in
\Cref{thm:V7-private-MTA}, every remaining unique-junction tree block
belongs to at least one of the following physical families:
\begin{enumerate}[label=\textup{(\mathrm{O})},leftmargin=5em]
\item \(\Xi_T>DY\), with the output resolvent assigned to the junction
      or a shared spectator as in
      \eqref{eq:V8-nonprivate-resolvent};
\end{enumerate}
or
\begin{enumerate}[label=\textup{(\mathrm{L})},leftmargin=5em]
\item an internal vertex \(i\in J_\tau\) satisfies
      \(a_i\ge\Xi_i/4\);
\end{enumerate}
or
\begin{enumerate}[label=\textup{(\mathrm{S})},leftmargin=5em]
\item an internal vertex \(i\in J_\tau\) satisfies
      \(h_i\ge\Xi_i/2\).
\end{enumerate}
If \(Y=0\), the same conclusion holds with family
\textup{(O)} or with \textup{(L)}/\textup{(S)} after replacing the
vacuous comparison to \(Y\) by
\(\Xi_i=a_i+h_i\).
\end{theorem}

\begin{proof}
For \(Y>0\), the exact complement
\eqref{eq:V7-private-complement} contains the output-dominant condition,
the failed branching condition, or both.  The first gives
\textup{(O)} by \Cref{lem:V8-nonprivate-output}; the second gives
\textup{(L)} or \textup{(S)} by
\Cref{cor:V8-branching-complement}.

If \(Y=0\), there is no private mass.  If the output is assigned by the
output priority, it belongs to \textup{(O)}.  Otherwise choose any
internal vertex.  Since \(\Xi_i=a_i+h_i\), either
\(a_i\ge\Xi_i/2\), which implies \textup{(L)}, or
\(h_i\ge\Xi_i/2\), which is \textup{(S)}.
\end{proof}

\begin{firewall}[The atlas is a cover, not yet a disjoint payment]
\label{fire:V8-cover-not-payment}
Families \textup{(O)}, \textup{(L)}, and \textup{(S)} may overlap.
A fixed priority order makes them disjoint, but V8 does not claim their
analytic payment.  Its theorem proves that no residual can hide in a
diffuse private bank: every residual exposes either a nonprivate output
payer or a quantitatively dominant local/shared mass at an internal
tree vertex.
\end{firewall}

\section{V8 ledger and V9 acquisition}
\label{sec:V8-ledger}

\begin{theorem}[Integrated compact-series V8 statement]
\label{thm:V8-integrated}
Preserving compact V1--V7, V8 proves:
\begin{enumerate}[label=\textup{(V8.\arabic*)},leftmargin=6em]
\item the large-internal-row consequence
      \eqref{eq:V8-large-internal-row};
\item the quantitative local/shared dichotomy
      \eqref{eq:V8-local-shared-dichotomy};
\item exact separation of private output mass;
\item elimination of coordinatewise private-payer sums on the
      output-dominant branch; and
\item the exhaustive nonprivate residual atlas
      \Cref{thm:V8-residual-atlas}.
\end{enumerate}
\end{theorem}

\begin{proof}
These are
\Cref{lem:V8-large-internal-row,lem:V8-local-shared-dichotomy,
lem:V8-nonprivate-output,cor:V8-no-private-payer-sum,
thm:V8-residual-atlas}.
\end{proof}

\begin{openproblem}[V9 mandate: pay the local-dominant priority family]
\label{op:V9-mandate}
\begin{enumerate}
\item Give family \textup{(O)} first priority, then
      \textup{(L)}, then \textup{(S)}, so every physical block is used
      once.
\item On \textup{(L)}, use \(a_i\ge\Xi_i/4\) only at an internal tree
      vertex and retain the exact V41 valence of that vertex.
\item Separate star and path trees.  A star requires two repeated marks
      at its center; a path requires one at each internal vertex.
\item Keep the Cartan scalar in the exact form
      \eqref{eq:V6-centered-fourth} and the global spectator product in
      \eqref{eq:V6-normalized-Cartan-word}.
\item For non-Cartan blocks, use the product-state capacity theorem
      before an operator norm.
\item Close the complete local-dominant family, or isolate the exact
      second internal-vertex mark which is missing on a path.
\item Compact V10 must begin with the next mandatory companion audit.
\end{enumerate}
\end{openproblem}

\begin{firewall}[Claim boundary after compact V8]
\label{fire:V8-current-boundary}
V8 exhaustively localizes the complement of the closed
private-dominant unique-\([4]\) sector to nonprivate output, local, or
shared certificates.  It does not yet pay those residual families, close
the full embedded \([4]\) sector, treat multiple junctions, normalized
collision aggregation, \([3,3]\), quartic or all-order MTA, continuum
Yang--Mills existence, or a positive mass gap.
\end{firewall}

\section{Append-only record for compact V8}
\label{sec:V8-record}

V8 was appended to the complete compact V7 file.  The exact V7 parent had
\(2939\) newline-delimited source records, \(101455\) bytes, and
SHA--256 digest
\[
\texttt{ea1a9c761fe3e9e14603b29c0e19685ade2ef043b0a2b1b8acbee9d063b4dc2f}.
\]
No compact V7 mathematical content was changed.  The sole final
\verb|\end{document}| is restored below.

% =============================================================================
% END APPENDED COMPACT VERSION V8 MATERIAL
% =============================================================================

% The compact V8 document terminator is intentionally disabled here:
% \end{document}

% =============================================================================
% BEGIN APPENDED COMPACT VERSION V9 MATERIAL
% =============================================================================

\chapter{V9: The local-dominant complete Cartan sector}
\label{ch:V9}

\section{Extracting a pointwise Cartan bound from one-link MTA}
\label{sec:V9-one-link-extraction}

The exact one-link-support quartic MTA theorem in the frozen archive is
independent of the formation law withdrawn later: it assembles the
complete cumulant, all final Peter--Weyl channels, and the single
resolvent before taking the norm.  On the multiplicity-one Cartan output
it yields a scalar inequality.

\begin{lemma}[Rank-one extraction on the Cartan component]
\label{lem:V9-Cartan-extraction}
Fix four irreducible one-coordinate inputs of masses
\(a_i=1+C_2(R_i)\), and let \(U_{\rm Car}\) be their Cartan output.
For the complete local fourth-cumulant scalar
\(K_{4,\rm Car}\),
\begin{equation}
\boxed{
\frac{A_{U_{\rm Car}}}
     {1-q_{\alpha,U_{\rm Car}}}
|K_{4,\rm Car}|
\le C\prod_{i=1}^4a_i.}
\label{eq:V9-local-Cartan-MTA}
\end{equation}
\end{lemma}

\begin{proof}
In each input representation choose the rank-one Fourier coefficient
\[
B_i:=|v_i\rangle\langle v_i|,
\qquad
\trnorm{B_i}=1.
\]
The product highest-weight vector lies in the multiplicity-one Cartan
component by \Cref{lem:V4-Cartan-vector}.  The exact Fourier formula
\eqref{eq:V1-exact-Fourier} therefore gives, on this output,
\[
d_{U_{\rm Car}}\trnorm{\widehat K(U_{\rm Car})}
=d_{\rm in}|K_{4,\rm Car}|.
\]
The exponential height factor cancels because
\(\operatorname{ht}(U_{\rm Car})=\sum_i\operatorname{ht}(R_i)\).
On one-link exact supports every mark equals one, so the one-link quartic
MTA right side is a fixed constant times
\[
d_{\rm in}\prod_i a_i.
\]
Restricting its nonnegative output norm sum to \(U_{\rm Car}\) proves
\eqref{eq:V9-local-Cartan-MTA}.  Notice that this argument concerns the
complete fourth cumulant.  It makes no assignment of that cumulant to
an individual contraction tree.
\end{proof}

\begin{lemma}[The local Cartan cumulant has a uniform \(s^{-1}\)
normalized-free bound]
\label{lem:V9-local-Cartan-over-s}
The two-sided one-link character gap implies
\begin{equation}
\boxed{
\frac1{s_\alpha}|K_{4,\rm Car}|
\le C\prod_{i=1}^4a_i.}
\label{eq:V9-local-Cartan-over-s}
\end{equation}
\end{lemma}

\begin{proof}
Put \(A=A_{U_{\rm Car}}\).  If \(s_\alpha A\le1\), the upper character
gap gives
\[
1-q_{\alpha,U_{\rm Car}}\le Cs_\alpha A,
\qquad
\frac{A}{1-q_{\alpha,U_{\rm Car}}}\ge\frac{c}{s_\alpha}.
\]
If \(s_\alpha A>1\), then
\[
\frac{A}{1-q_{\alpha,U_{\rm Car}}}\ge A>\frac1{s_\alpha}.
\]
In either case, \eqref{eq:V9-local-Cartan-MTA} proves the result.
\end{proof}

\section{Global Cartan word below the output-dominant priority}
\label{sec:V9-global-Cartan-payment}

Retain the notation of V7--V8.  Give family \textup{(O)} of
\Cref{thm:V8-residual-atlas} first priority.  On its complement,
\begin{equation}
\Xi_T\le DY.
\label{eq:V9-not-output-dominant}
\end{equation}
In particular \(Y>0\), because every retained final output is nontrivial.

\begin{lemma}[Private product pays the global Cartan resolvent]
\label{lem:V9-private-pays-global}
For the complete Cartan word
\eqref{eq:V6-global-Cartan-word}, on
\eqref{eq:V9-not-output-dominant},
\begin{equation}
\boxed{
\frac{\Xi_T}{1-q_{\alpha,\boldsymbol T}}
q_P
\left|K_{4,\rm Car}\right|
\prod_{\substack{f\ne e\\f\notin P}}
\prod_{B\in\pi_f}|q_{B,f}|
\le
\frac{C_D}{s_\alpha}|K_{4,\rm Car}|.}
\label{eq:V9-private-global-payment}
\end{equation}
\end{lemma}

\begin{proof}
All nonprivate spectator multipliers have modulus at most one.  The
global product multiplier has the form
\[
q_{\alpha,\boldsymbol T}=q_Pq_{\rm rest},
\qquad 0\le q_{\rm rest}\le1,
\]
so
\[
1-q_{\alpha,\boldsymbol T}\ge1-q_P.
\]
The product version of the one-link tail debit, obtained by the
Bernoulli telescoping proof used in
\Cref{lem:V7-product-moment-compiler}, is
\[
s_\alpha Yq_P\le C(1-q_P).
\]
Use \(\Xi_T\le DY\) to obtain
\[
\frac{\Xi_Tq_P}{1-q_{\alpha,\boldsymbol T}}
\le
\frac{DYq_P}{1-q_P}
\le\frac{C_D}{s_\alpha}.
\]
Multiply by the local cumulant and the discarded contraction factors.
\end{proof}

\begin{corollary}[Unmarked local Cartan product bound]
\label{cor:V9-unmarked-local-product}
On the complement of output priority \textup{(O)},
\begin{equation}
\boxed{
\textup{normalized Cartan word}
\le C_D\prod_{i=1}^4a_i.}
\label{eq:V9-unmarked-local-product}
\end{equation}
\end{corollary}

\begin{proof}
Combine
\Cref{lem:V9-private-pays-global,lem:V9-local-Cartan-over-s}.
\end{proof}

\section{Mark recovery from the complete tree sum}
\label{sec:V9-mark-recovery}

For a unique marking at the junction coordinate \(e\), the marked tree
weight is
\begin{equation}
\mathcal W_{\tau,e}
=
\prod_{i=1}^4a_i^{d_i}\Xi_i^{1-d_i}
=
\left(\prod_i a_i\right)
\prod_{i\in J_\tau}
\left(\frac{a_i}{\Xi_i}\right)^{d_i-1}.
\label{eq:V9-marked-factorization}
\end{equation}

Write
\begin{equation}
\mathfrak W_e:=\sum_{\tau\in\mathfrak T_4}\mathcal W_{\tau,e},
\label{eq:V9-complete-tree-weight}
\end{equation}
where \(\mathfrak T_4\) is the set of the sixteen labelled quartic
trees.  This is the complete marked one-link contribution on the MTA
right side.  The sum, rather than an individual summand, is the correct
object to compare with the complete Cartan cumulant.

\begin{theorem}[The complete local-dominant Cartan sector is paid]
\label{thm:V9-Cartan-star}
Suppose that the residual priority block is not output-dominant and
that at least one input row \(c\) is local-dominant:
\begin{equation}
a_c\ge\frac14\Xi_c,
\label{eq:V9-star-local}
\end{equation}
Then the complete global Cartan word obeys
\begin{equation}
\boxed{
\textup{normalized Cartan word}
\le C_D\mathfrak W_e.}
\label{eq:V9-Cartan-star-payment}
\end{equation}
\end{theorem}

\begin{proof}
Among the sixteen labelled trees is the star \(\tau_c\) with center
\(c\).  Its three internal incidences give
\[
\mathcal W_{\tau_c,e}
=
\left(\prod_i a_i\right)
\left(\frac{a_c}{\Xi_c}\right)^2
\ge\frac1{16}\prod_i a_i.
\]
Consequently
\(\mathfrak W_e\ge\mathcal W_{\tau_c,e}\ge
16^{-1}\prod_i a_i\).  Apply
\eqref{eq:V9-unmarked-local-product}.  No decomposition of
\(K_{4,\rm Car}\) among the sixteen trees is used.
\end{proof}

\begin{corollary}[No Cartan residual remains in family \textup{(L)}]
\label{cor:V9-local-Cartan-residual}
After output priority \textup{(O)} and the private-dominant closure of
V7, every complete Cartan block in family \textup{(L)} is closed by
\Cref{thm:V9-Cartan-star}.  Thus a residual complete Cartan block has
\begin{equation}
a_i<\frac14\Xi_i
\qquad(1\le i\le4).
\label{eq:V9-no-local-row}
\end{equation}
\end{corollary}

\begin{proof}
Family \textup{(L)} supplies at least one row \(c\) satisfying
\eqref{eq:V9-star-local}.  The complete MTA tree sum contains the star
centered at that row, whether or not the tree used to discover \(c\)
was itself a star.  The theorem therefore closes the block.  Negating
the existence of such a row gives \eqref{eq:V9-no-local-row}.
\end{proof}

\begin{firewall}[Cartan closure is not full physical closure]
\label{fire:V9-Cartan-only}
\Cref{thm:V9-Cartan-star}
applies to the multiplicity-one Cartan output line.  It does not bound
non-Cartan physical blocks, where the local one-link MTA estimate cannot
automatically be converted to an operator-norm carrier estimate.
Those blocks still require the product-state capacity route of V1--V4.
\end{firewall}

\section{V9 ledger and mandatory V10 audit}
\label{sec:V9-ledger}

\begin{theorem}[Integrated compact-series V9 statement]
\label{thm:V9-integrated}
Preserving compact V1--V8, V9 proves:
\begin{enumerate}[label=\textup{(V9.\arabic*)},leftmargin=6em]
\item the pointwise one-link Cartan estimate
      \eqref{eq:V9-local-Cartan-MTA};
\item its \(s_\alpha^{-1}\)-weighted consequence
      \eqref{eq:V9-local-Cartan-over-s};
\item bankwise payment of the global resolvent by the exact private
      product;
\item recovery of the marked MTA weight from the complete sixteen-tree
      sum whenever any input row is local-dominant; and
\item closure of the entire complete-Cartan family \textup{(L)}, with
      no tree-fibre allocation assumption.
\end{enumerate}
\end{theorem}

\begin{proof}
These are
\Cref{lem:V9-Cartan-extraction,lem:V9-local-Cartan-over-s,
lem:V9-private-pays-global,thm:V9-Cartan-star,
cor:V9-local-Cartan-residual}.
\end{proof}

\begin{openproblem}[V10 mandate: audit and the all-nonlocal Cartan
residual]
\label{op:V10-mandate}
\begin{enumerate}
\item Begin with the mandatory read-only audit of the newest SM and NS
      compact or legacy notes.
\item Work under the exact residual condition
      \(a_i<\Xi_i/4\) for all four rows.
\item Reapply the V8 branching dichotomy to show that every failed
      private branching product forces either output priority
      \textup{(O)} or a genuinely shared-dominant internal row.
\item On the shared-dominant branch, construct a positive shared-bank
      payer without conditioning the remaining tensor coordinates into
      a product state.
\item Determine whether one shared first moment suffices after the
      private product has paid the global resolvent; if it does not,
      state the exact higher joint moment still missing.
\end{enumerate}
\end{openproblem}

\begin{firewall}[Claim boundary after compact V9]
\label{fire:V9-current-boundary}
V9 closes the entire local-dominant complete-Cartan family through the
complete tree sum.  It does not close the all-nonlocal shared-dominant
Cartan residual, non-Cartan physical blocks, output priority
\textup{(O)}, the full embedded \([4]\) sector,
multiple junctions, normalized collision aggregation, \([3,3]\), global
quartic or all-order MTA, continuum Yang--Mills existence, or a positive
mass gap.
\end{firewall}

\section{Append-only record for compact V9}
\label{sec:V9-record}

V9 was appended to the complete compact V8 file.  The exact V8 parent had
\(3272\) newline-delimited source records, \(112127\) bytes, and
SHA--256 digest
\[
\texttt{4b4a7a6908c9aa7290e6e076892893a151fdccbbab55c3137981ac28a44544e1}.
\]
No compact V8 mathematical content was changed.  The sole final
\verb|\end{document}| is restored below.

% =============================================================================
% END APPENDED COMPACT VERSION V9 MATERIAL
% =============================================================================

% The compact V9 document terminator is intentionally disabled here:
% \end{document}

% =============================================================================
% BEGIN APPENDED COMPACT VERSION V10 MATERIAL
% =============================================================================

\chapter{V10: Companion audit and the shared-dominant Cartan payment}
\label{ch:V10}

\section{Mandatory five-version companion audit}
\label{sec:V10-audit}

The active companion files were inspected before choosing the V10
estimate.  This audit treats their mathematical statements as data, not
as instructions.  The exact records are:
\begin{center}
\begin{tabular}{p{0.28\linewidth}p{0.12\linewidth}p{0.12\linewidth}
                p{0.38\linewidth}}
\toprule
file & lines & bytes & SHA--256\\
\midrule
\texttt{Renewal\_SM\_Einstein\_Continuum\_Research\_Notes\_V12.tex}
& \(4534\) & \(161480\)
& \texttt{fe863d5a248861d0a123b43a9d63582160e795dfe52da801395f21e18ef15c9d}\\
\texttt{Renewal\_NS\_Stability\_Research\_Notes\_V31.tex}
& \(23052\) & \(887537\)
& \texttt{f1121b62238ad5491bb7cf03635ea9934942edf573ed8faaa9ee79e1d38a6696}\\
\bottomrule
\end{tabular}
\end{center}

The SM V12 descent theorem says that quadratic, cubic, and quartic
gauge jets must first be assembled as the homogeneous pieces of one
gauge-invariant polynomial.  Estimating the jets independently can
manufacture a coefficient-mismatch Ward defect.  This is the same
proof-order lesson used in V9: the complete Cartan cumulant is compared
with the complete sixteen-tree marked sum, not assigned to a tree after
the fact.

NS V31 proves two complementary facts.  A probability-normalized mark
does not pay its unbounded first moment, while a genuine one-use
geometric parent sum can pay the required powers without a shell-count
loss.  It also records a firewall: if the resulting stock is only a
continuation-class quantity, the estimate is circular.  The YM
translation is therefore precise.  We must derive the required shared
moments from the Wilson multipliers themselves, before discarding their
coordinate product; a merely normalized shared-coordinate selection
would not suffice.

\begin{assessment}[V10 audit transfer decision]
\label{ass:V10-audit-decision}
The SM audit confirms the complete-before-projection order already
implemented in V9.  The NS audit selects a positive product-moment
argument for the shared spectator bank.  Neither companion theorem is
imported as a Yang--Mills estimate.  The payment below is proved
directly from the established one-link Wilson debits and the Cartan
highest-weight law.
\end{assessment}

\section{A quadratic product moment for a shared Cartan row}
\label{sec:V10-shared-moment}

For the V10 shared payment we impose the standard mass-gap target that
\(G\) is compact connected semisimple; in particular the result applies
to every compact simple gauge group.  A torus factor with oppositely
directed charges would invalidate the Casimir monotonicity used below
and is not included in this theorem.

At every off-junction shared coordinate \(f\), let
\(I_f\subset\{1,2,3,4\}\) be its active input rows.  On the Cartan
spectator output write
\begin{equation}
\lambda_f:=\sum_{j\in I_f}\lambda_{j,f},
\qquad
r_f:=q_{\alpha,\lambda_f}.
\label{eq:V10-shared-Cartan-output}
\end{equation}
For a fixed row \(i\), set
\begin{align}
F_i&:=\{f\ne e:i\in I_f,\ |I_f|\ge2\},\notag\\
a_{i,f}&:=1+C_2(\lambda_{i,f}),\qquad
h_i:=\sum_{f\in F_i}a_{i,f},\label{eq:V10-shared-row-mass}\\
q_{S,i}&:=\prod_{f\in F_i}r_f,\qquad
q_S:=\prod_{\substack{f\ne e\\|I_f|\ge2}}r_f.
\label{eq:V10-shared-products}
\end{align}
Thus \(0\le q_S\le q_{S,i}\le1\).

\begin{lemma}[Cartan output mass dominates every active input mass]
\label{lem:V10-Cartan-mass-monotonicity}
For \(f\in F_i\),
\begin{equation}
\boxed{a_{i,f}\le 1+C_2(\lambda_f).}
\label{eq:V10-Cartan-mass-monotonicity}
\end{equation}
\end{lemma}

\begin{proof}
For the invariant inner product used to define the Casimir,
\[
C_2(\lambda)=\langle\lambda,\lambda+2\rho\rangle.
\]
Put \(\mu=\sum_{j\in I_f\setminus\{i\}}\lambda_{j,f}\).  Both
\(\lambda_{i,f}\) and \(\mu\) are dominant.  Hence
\[
C_2(\lambda_{i,f}+\mu)-C_2(\lambda_{i,f})
=2\langle\lambda_{i,f},\mu\rangle
2\langle\mu,\rho\rangle+\langle\mu,\mu\rangle\ge0.
\]
The nonnegativity follows from the nonnegative inverse Cartan matrix on
the dominant chamber, component by component on the semisimple group.
Adding one proves the claim.
\end{proof}

\begin{theorem}[Shared-row Wilson moments through order two]
\label{thm:V10-shared-row-moments}
Uniformly in the number of shared coordinates and their
representations,
\begin{equation}
\boxed{
(s_\alpha h_i)^kq_S\le C_k,
\qquad k=1,2.}
\label{eq:V10-shared-row-moments}
\end{equation}
\end{theorem}

\begin{proof}
For \(f\in F_i\), put \(x_f=s_\alpha a_{i,f}\).  By
\Cref{lem:V10-Cartan-mass-monotonicity} and the established one-link
linear and quadratic Wilson debits,
\[
x_f^kr_f
\le
\bigl(s_\alpha(1+C_2(\lambda_f))\bigr)^kr_f
\le C_k(1-r_f),
\qquad k=1,2.
\]
Apply the Bernoulli product-moment compiler
\Cref{lem:V7-product-moment-compiler} to the family
\(\{(x_f,r_f):f\in F_i\}\).  It gives
\((s_\alpha h_i)^kq_{S,i}\le C_k\).
Since all multipliers are nonnegative and \(q_S\le q_{S,i}\), the same
bound holds with \(q_S\).  No shared coordinate is counted after the
product has been replaced by separate norms.
\end{proof}

\begin{firewall}[Why the first shared moment alone is insufficient]
\label{fire:V10-first-moment-insufficient}
At a star center the marked weight contains
\(\Xi_i^{-2}\).  A first-moment estimate supplies at most
\(h_i^{-1}\), so it cannot close that tree uniformly.  The second line
of \eqref{eq:V10-shared-row-moments} is exactly the required order.
This is the finite-product analogue of the first-moment warning in the
NS audit, but its proof here is entirely Wilson-theoretic.
\end{firewall}

\section{Selecting one shared row which pays both tree excesses}
\label{sec:V10-row-selection}

For a labelled quartic tree \(\tau\), define its branching denominator
\begin{equation}
B_\tau:=\prod_{j=1}^4\Xi_j^{d_j-1}.
\label{eq:V10-branching-denominator}
\end{equation}

\begin{lemma}[Maximal internal row dominates the branching denominator]
\label{lem:V10-maximal-shared-row}
Let \(D\ge16\), \(Y>0\), and suppose
\begin{equation}
B_\tau>DY^2,
\qquad
a_j<\frac14\Xi_j\quad(1\le j\le4).
\label{eq:V10-all-nonlocal-failed}
\end{equation}
Choose \(i\in J_\tau\) with
\(\Xi_i=\max_{j\in J_\tau}\Xi_j\).  Then
\begin{equation}
\boxed{
\Xi_i>\sqrt D\,Y,\qquad
h_i>\frac12\Xi_i,\qquad
B_\tau\le\Xi_i^2<4h_i^2.}
\label{eq:V10-maximal-shared-row}
\end{equation}
\end{lemma}

\begin{proof}
If \(\tau\) is a star, then
\(B_\tau=\Xi_i^2\), so the first and third claims are immediate.
If \(\tau\) is a path with internal vertices \(i,j\), then
\[
B_\tau=\Xi_i\Xi_j\le\Xi_i^2.
\]
The strict failed-branching inequality gives
\(\Xi_i>\sqrt D\,Y\).  In either case
\[
p_i\le Y<\frac1{\sqrt D}\Xi_i\le\frac14\Xi_i.
\]
Together with \(a_i<\Xi_i/4\) and
\(\Xi_i=a_i+p_i+h_i\), this yields
\(h_i>\Xi_i/2\).  Therefore
\(\Xi_i^2<4h_i^2\), completing the proof.
\end{proof}

\section{Payment of the shared-dominant Cartan tree sector}
\label{sec:V10-shared-payment}

Let \(K_{\tau,\rm Car}\) denote the Cartan scalar of the V41 tree fibre
at the junction, and put
\begin{equation}
L_\tau:=\prod_{j=1}^4a_j^{d_j/2},
\qquad
\mathcal W_{\tau,e}
=\frac{\prod_j a_j^{d_j}}{B_\tau}.
\label{eq:V10-local-and-marked-weights}
\end{equation}
The established local tree estimate is
\begin{equation}
|K_{\tau,\rm Car}|\le Cs_\alpha^3L_\tau.
\label{eq:V10-local-tree-estimate}
\end{equation}

\begin{theorem}[All-nonlocal shared-dominant Cartan payment]
\label{thm:V10-shared-Cartan-payment}
Fix \(D\ge16\).  Suppose \(Y>0\), the global product output is not in
priority family \textup{(O)}, and a tree \(\tau\) satisfies
\eqref{eq:V10-all-nonlocal-failed}.  Then its normalized Cartan
contribution obeys
\begin{equation}
\boxed{
\frac{\Xi_T}{1-q_{\alpha,\boldsymbol T}}\,
q_Pq_S|K_{\tau,\rm Car}|
\le C_D\mathcal W_{\tau,e}.}
\label{eq:V10-shared-Cartan-payment}
\end{equation}
\end{theorem}

\begin{proof}
Choose the maximal internal row \(i\) from
\Cref{lem:V10-maximal-shared-row}.  Since the block is not in
\textup{(O)}, one has \(\Xi_T\le DY\).

First suppose \(s_\alpha\Xi_T\le c_0\).  The low-output line of
\eqref{eq:V7-output-resolvent-split},
\eqref{eq:V10-local-tree-estimate}, and \(q_P\le1\) give
\[
\frac{\Xi_T}{1-q_{\alpha,\boldsymbol T}}\,
q_Pq_S|K_{\tau,\rm Car}|
\le Cs_\alpha^2q_SL_\tau.
\]
The quadratic case of \eqref{eq:V10-shared-row-moments} and
\eqref{eq:V10-maximal-shared-row} imply
\[
s_\alpha^2q_S
\le\frac{C}{h_i^2}
\le\frac{C}{B_\tau}.
\]
Since every \(a_j\ge1\),
\[
L_\tau\le\prod_j a_j^{d_j}.
\]
These inequalities prove \eqref{eq:V10-shared-Cartan-payment} on the
low-output branch.

Now suppose \(s_\alpha\Xi_T>c_0\).  The high-output line of
\eqref{eq:V7-output-resolvent-split} and
\(\Xi_T\le DY\) give
\[
\frac{\Xi_T}{1-q_{\alpha,\boldsymbol T}}\,
q_Pq_S|K_{\tau,\rm Car}|
\le C_DYs_\alpha^3q_Pq_SL_\tau.
\]
The private first-moment debit and the shared quadratic debit yield
\[
(s_\alpha Y)q_P\le C,
\qquad
(s_\alpha h_i)^2q_S\le C.
\]
Multiplying them gives
\[
Ys_\alpha^3q_Pq_S\le\frac{C}{h_i^2}
\le\frac{C}{B_\tau}.
\]
The same comparison \(L_\tau\le\prod_j a_j^{d_j}\) proves the result.
\end{proof}

\begin{remark}[No false tensor conditioning]
\label{rem:V10-no-conditioning}
The proof evaluates the genuine Cartan scalar at every shared
coordinate and retains the complete product \(q_S\) until its quadratic
moment is taken.  It neither conditions the remaining coordinates into
a product state nor converts a positive shared contraction into an
operator norm prematurely.
\end{remark}

\begin{corollary}[Unique-junction Cartan closure outside output priority]
\label{cor:V10-Cartan-outside-output}
For a unique-junction complete Cartan block, every contribution outside
priority family \textup{(O)} obeys the marked quartic tree atlas.
\end{corollary}

\begin{proof}
If \(Y=0\), a nontrivial block satisfies
\(\Xi_T>DY=0\) and is in \textup{(O)}, so there is nothing to prove.
Assume \(Y>0\).  Give first priority to the global event that some row
is local-dominant.  V9 closes the complete sixteen-tree Cartan sum on
that event.  On its complement, fix a tree.  If both inequalities in
\eqref{eq:V7-private-dominance} hold, V7 closes that tree.  Otherwise,
because the output inequality holds outside \textup{(O)}, the branching
inequality fails and
\Cref{thm:V10-shared-Cartan-payment} closes the tree.  Sum the finitely
many tree estimates.
\end{proof}

\begin{firewall}[Scope of the V10 closure]
\label{fire:V10-scope}
\Cref{cor:V10-Cartan-outside-output} concerns only the
multiplicity-one Cartan output at the junction and Cartan outputs at
shared spectators, for compact connected semisimple \(G\).  It does not
prove a capacity bound for non-Cartan physical blocks or for cancelling
abelian charges.  It also deliberately leaves output priority
\textup{(O)} open.
\end{firewall}

\section{V10 ledger and V11 acquisition}
\label{sec:V10-ledger}

\begin{theorem}[Integrated compact-series V10 statement]
\label{thm:V10-integrated}
Preserving compact V1--V9, V10 proves:
\begin{enumerate}[label=\textup{(V10.\arabic*)},leftmargin=6.5em]
\item the mandatory hashed audit of active SM V12 and NS V31;
\item Cartan shared-output mass monotonicity
      \eqref{eq:V10-Cartan-mass-monotonicity};
\item support-uniform shared-row Wilson moments through order two,
      \eqref{eq:V10-shared-row-moments};
\item maximal-row payment of both quartic branching excesses,
      \eqref{eq:V10-maximal-shared-row};
\item the all-nonlocal shared-dominant Cartan tree estimate
      \eqref{eq:V10-shared-Cartan-payment}; and
\item closure of the complete unique-junction Cartan sector outside
      output priority \textup{(O)}.
\end{enumerate}
\end{theorem}

\begin{proof}
These are
\Cref{ass:V10-audit-decision,lem:V10-Cartan-mass-monotonicity,
thm:V10-shared-row-moments,lem:V10-maximal-shared-row,
thm:V10-shared-Cartan-payment,cor:V10-Cartan-outside-output}.
\end{proof}

\begin{openproblem}[V11 mandate: the output-dominant Cartan sector]
\label{op:V11-mandate}
\begin{enumerate}
\item Work on \(\Xi_T>DY\) and retain the exact split
      \(\Xi_{T,e}+\Xi_{T,S}>(1-D^{-1})\Xi_T\).
\item Separate a junction-output payer from a shared-output payer before
      any coordinate sum is normed.
\item For a shared-output payer, prove the Cartan output analogue of
      \Cref{thm:V10-shared-row-moments} and preserve the remaining
      spectator damping.
\item For the junction-output payer, combine the output mass with the
      complete local fourth cumulant before resolving contraction
      trees.
\item Close the full Cartan output-priority family, or identify the
      exact joint output/input moment which remains missing.
\end{enumerate}
\end{openproblem}

\begin{firewall}[Claim boundary after compact V10]
\label{fire:V10-current-boundary}
V10 closes the unique-junction Cartan sector outside output priority.
It does not close output priority \textup{(O)}, non-Cartan physical
blocks, the full embedded \([4]\) sector, multiple junctions, normalized
collision aggregation, \([3,3]\), global quartic or all-order MTA,
continuum Yang--Mills existence, or a positive mass gap.
\end{firewall}

\section{Append-only record for compact V10}
\label{sec:V10-record}

V10 was appended to the complete compact V9 file.  The exact V9 parent
had \(3596\) newline-delimited source records, \(122715\) bytes, and
SHA--256 digest
\[
\texttt{9397407fe1dc5ae43d95f9a4907c1a803db27e2cb18dc17b2824b027975da07e}.
\]
No compact V9 mathematical content was changed.  The sole final
\verb|\end{document}| is restored below.

% =============================================================================
% END APPENDED COMPACT VERSION V10 MATERIAL
% =============================================================================

% The compact V10 document terminator is intentionally disabled here:
% \end{document}

% =============================================================================
% BEGIN APPENDED COMPACT VERSION V11 MATERIAL
% =============================================================================

\chapter{V11: Bankwise closure of the output-dominant Cartan sector}
\label{ch:V11}

\section{A failure-retaining product-moment compiler}
\label{sec:V11-reinforced-compiler}

The V7 compiler discarded the probability that at least one one-link
trial fails.  An output-bank resolvent requires that factor to be kept.

\begin{lemma}[Product moments with one retained failure]
\label{lem:V11-retained-failure}
Let \(F\) be finite, \(x_f\ge0\), and \(0\le r_f\le1\).  Suppose that
for \(1\le k\le m\),
\begin{equation}
x_f^kr_f\le C_k(1-r_f).
\label{eq:V11-site-debits}
\end{equation}
Then
\begin{equation}
\boxed{
\left(\sum_{f\in F}x_f\right)^m
\prod_{f\in F}r_f
\le C_m^\sharp
\left(1-\prod_{f\in F}r_f\right).}
\label{eq:V11-retained-failure}
\end{equation}
The constant depends only on \(m,C_1,\ldots,C_m\), not on
\(\lvert F\rvert\).
\end{lemma}

\begin{proof}
Expand the power and group a monomial by its \(j\) distinct indices
\(f_1,\ldots,f_j\) and positive multiplicities
\(k_1+\cdots+k_j=m\).  Applying
\eqref{eq:V11-site-debits} at those indices bounds the monomial times
\(\prod_fr_f\) by a fixed coefficient times
\[
\prod_{\ell=1}^j(1-r_{f_\ell})
\prod_{f\notin\{f_1,\ldots,f_j\}}r_f.
\]
For fixed \(j\), summing over unordered distinct indices gives exactly
the probability of \(j\) failures in independent Bernoulli trials.
The sum over \(1\le j\le m\) is bounded by the probability of at least
one failure,
\[
1-\prod_fr_f.
\]
Only finitely many ordered multiplicity patterns occur for fixed \(m\),
which proves the claim.
\end{proof}

\section{The complete shared-output bank}
\label{sec:V11-shared-output-bank}

Continue under the compact connected semisimple hypothesis of V10.
At a shared spectator coordinate \(f\), let
\[
b_f:=1+C_2(\lambda_f),
\qquad
r_f:=q_{\alpha,\lambda_f},
\]
where \(\lambda_f\) is its Cartan output.  Define
\begin{equation}
H_T:=\sum_{\substack{f\ne e\\|I_f|\ge2}}b_f,
\qquad
q_S:=\prod_{\substack{f\ne e\\|I_f|\ge2}}r_f.
\label{eq:V11-shared-output-bank}
\end{equation}
If the shared bank is empty, expressions containing
\(H_Tq_S/(1-q_S)\) are interpreted as zero.

\begin{theorem}[Resolvent-weighted shared-output moments]
\label{thm:V11-shared-output-moments}
For \(m=1,2,3\),
\begin{equation}
\boxed{
\frac{(s_\alpha H_T)^mq_S}{1-q_S}\le C_m.}
\label{eq:V11-shared-output-moments}
\end{equation}
\end{theorem}

\begin{proof}
Put \(x_f=s_\alpha b_f\).  The established one-link Wilson debits give
\[
x_f^kr_f\le C_k(1-r_f),
\qquad 1\le k\le3.
\]
Apply \Cref{lem:V11-retained-failure} and divide by
\(1-q_S\).  Every nonempty shared Cartan output is nontrivial, so the
denominator is positive.
\end{proof}

Let
\begin{equation}
H:=\sum_{i=1}^4h_i
=\sum_{\substack{f\ne e\\|I_f|\ge2}}
  \sum_{i\in I_f}a_{i,f}
\label{eq:V11-total-shared-input}
\end{equation}
be the shared input mass counted with row incidence.

\begin{lemma}[Input/output comparison on Cartan coordinates]
\label{lem:V11-input-output-comparison}
There is a group-uniform constant \(C_G\) such that
\begin{align}
H&\le4H_T,\label{eq:V11-shared-input-output}\\
A_{T,e}&\le C_G\sum_{i=1}^4a_i.
\label{eq:V11-junction-output-input}
\end{align}
If no row is local-dominant, so
\(a_i<\Xi_i/4\) for every \(i\), then
\begin{equation}
\boxed{A_{T,e}\le C_G'(Y+H_T).}
\label{eq:V11-junction-output-bank-comparison}
\end{equation}
\end{lemma}

\begin{proof}
At a shared coordinate,
\Cref{lem:V10-Cartan-mass-monotonicity} gives
\(a_{i,f}\le b_f\) for every active row.  There are at most four rows,
so summing gives \eqref{eq:V11-shared-input-output}.

For dominant weights \(\lambda_1,\ldots,\lambda_4\),
Cauchy--Schwarz and
\(C_2(\lambda)=\|\lambda\|^2+2\langle\rho,\lambda\rangle\)
give
\[
C_2\!\left(\sum_i\lambda_i\right)
\le4\sum_i\|\lambda_i\|^2
2\sum_i\langle\rho,\lambda_i\rangle
\le4\sum_i C_2(\lambda_i).
\]
This proves \eqref{eq:V11-junction-output-input}, after harmlessly
absorbing the added ones.

Finally, \(4a_i<\Xi_i=a_i+p_i+h_i\) implies
\[
3a_i<p_i+h_i\le Y+h_i.
\]
Sum over the four rows and use
\eqref{eq:V11-shared-input-output}.  Substitution into
\eqref{eq:V11-junction-output-input} proves
\eqref{eq:V11-junction-output-bank-comparison}.
\end{proof}

\begin{corollary}[Ordinary total shared-input moments]
\label{cor:V11-total-shared-moments}
For \(m=1,2,3\),
\begin{equation}
\boxed{(s_\alpha H)^mq_S\le C_m.}
\label{eq:V11-total-shared-moments}
\end{equation}
\end{corollary}

\begin{proof}
Use \(H\le4H_T\), then either
\Cref{thm:V11-shared-output-moments} and
\(1-q_S\le1\), or the V7 product-moment compiler directly.
\end{proof}

\section{Bankwise decomposition of output priority}
\label{sec:V11-output-decomposition}

For a Cartan product output,
\[
q_{\alpha,\boldsymbol T}
=q_{\alpha,T_e}q_Pq_S.
\]
On priority family \textup{(O)}, V8 gives
\[
\Xi_T\le\frac{D}{D-1}(A_{T,e}+H_T).
\]

\begin{lemma}[Only two nonprivate payer banks]
\label{lem:V11-two-payer-banks}
For any nonnegative junction coefficient \(K\),
\begin{align}
&\frac{\Xi_T}{1-q_{\alpha,\boldsymbol T}}\,
q_Pq_S|K|
\notag\\
&\quad\le
\frac{D}{D-1}
\left[
\frac{A_{T,e}}{1-q_{\alpha,T_e}}q_Pq_S|K|
+
\frac{H_Tq_S}{1-q_S}q_P|K|
\right].
\label{eq:V11-two-payer-banks}
\end{align}
\end{lemma}

\begin{proof}
The numerator comparison is the output-dominant estimate above.
Moreover
\[
1-q_{\alpha,\boldsymbol T}\ge1-q_{\alpha,T_e},
\qquad
1-q_{\alpha,\boldsymbol T}\ge1-q_S.
\]
Apply the first inequality to the junction mass and the second to the
shared-output mass.  No private-coordinate payer sum appears.
\end{proof}

\section{Local-dominant output-priority blocks}
\label{sec:V11-local-output}

\begin{theorem}[Output priority with a local-dominant row]
\label{thm:V11-local-output-payment}
Suppose \(\Xi_T>DY\) and
\(a_c\ge\Xi_c/4\) for at least one row \(c\).  Then the complete
unique-junction Cartan word obeys
\begin{equation}
\boxed{
\textup{normalized complete Cartan word}
\le C_D\mathfrak W_e.}
\label{eq:V11-local-output-payment}
\end{equation}
\end{theorem}

\begin{proof}
Apply \Cref{lem:V11-two-payer-banks} with the complete scalar
\(K_{4,\rm Car}\).  The junction-payer term is bounded by
\(C\prod_i a_i\) by \eqref{eq:V9-local-Cartan-MTA}.
The \(m=1\) case of
\eqref{eq:V11-shared-output-moments} bounds the shared-payer term by
\[
\frac{C}{s_\alpha}|K_{4,\rm Car}|
\le C\prod_i a_i
\]
using \eqref{eq:V9-local-Cartan-over-s}.  The star centered at \(c\)
belongs to the complete tree sum and satisfies
\[
\mathfrak W_e\ge\mathcal W_{\tau_c,e}
\ge\frac1{16}\prod_i a_i
\]
by the V9 calculation.  Combine these estimates.
\end{proof}

\section{All-nonlocal output-priority blocks}
\label{sec:V11-nonlocal-output}

Fix a tree \(\tau\), and retain
\[
B_\tau=\prod_i\Xi_i^{d_i-1},
\qquad
L_\tau=\prod_i a_i^{d_i/2}.
\]

\begin{lemma}[Tree-local junction payer]
\label{lem:V11-tree-junction-payer}
For the Cartan tree scalar,
\begin{equation}
\boxed{
\frac{A_{T,e}}{1-q_{\alpha,T_e}}
|K_{\tau,\rm Car}|
\le
C s_\alpha^2(1+s_\alpha A_{T,e})L_\tau.}
\label{eq:V11-tree-junction-payer}
\end{equation}
\end{lemma}

\begin{proof}
If \(s_\alpha A_{T,e}\le1\), this is the low-output payer
\eqref{eq:V3-low-output-tree}.  If
\(s_\alpha A_{T,e}>1\), the one-link gap is bounded below by a positive
constant.  Hence \eqref{eq:V10-local-tree-estimate} gives
\[
\frac{A_{T,e}}{1-q_{\alpha,T_e}}
|K_{\tau,\rm Car}|
\le C A_{T,e}s_\alpha^3L_\tau.
\]
The two cases give the displayed bound.
\end{proof}

\begin{lemma}[All-nonlocal branching alternatives]
\label{lem:V11-nonlocal-branching}
Assume \(a_i<\Xi_i/4\) for every row.
\begin{enumerate}[label=\textup{(\roman*)},leftmargin=4em]
\item If \(B_\tau\le DY^2\), then \(Y>0\) and
\begin{equation}
s_\alpha^2q_P\le\frac{C_D}{B_\tau}.
\label{eq:V11-private-two-power}
\end{equation}
\item If \(B_\tau>DY^2\), there is a maximal internal row \(i\) such
that
\begin{equation}
h_i>\frac12\Xi_i,\qquad
B_\tau<4h_i^2.
\label{eq:V11-shared-two-power}
\end{equation}
This remains true when \(Y=0\).
\end{enumerate}
\end{lemma}

\begin{proof}
In the first case \(B_\tau>0\) forces \(Y>0\).  The quadratic private
debit gives \(s_\alpha^2q_P\le C/Y^2\), and
\(B_\tau\le DY^2\) proves \eqref{eq:V11-private-two-power}.

For the second case choose an internal row of maximal \(\Xi_i\).
Exactly as in \Cref{lem:V10-maximal-shared-row},
\(B_\tau\le\Xi_i^2\) and
\(\Xi_i>\sqrt D\,Y\).  Thus \(p_i\le Y<\Xi_i/4\).
Together with \(a_i<\Xi_i/4\), this gives
\(h_i>\Xi_i/2\) and then \(B_\tau<4h_i^2\).
When \(Y=0\), the same proof has \(p_i=0\).
\end{proof}

\begin{theorem}[All-nonlocal output-dominant Cartan tree payment]
\label{thm:V11-nonlocal-output-payment}
Fix \(D\ge16\).  Suppose \(\Xi_T>DY\) and
\(a_i<\Xi_i/4\) for all rows.  Then every Cartan tree fibre satisfies
\begin{equation}
\boxed{
\frac{\Xi_T}{1-q_{\alpha,\boldsymbol T}}\,
q_Pq_S|K_{\tau,\rm Car}|
\le C_D\mathcal W_{\tau,e}.}
\label{eq:V11-nonlocal-output-payment}
\end{equation}
\end{theorem}

\begin{proof}
Apply \Cref{lem:V11-two-payer-banks} with
\(K=K_{\tau,\rm Car}\).

First suppose \(B_\tau\le DY^2\).  The shared-output payer is, by
\eqref{eq:V11-shared-output-moments} with \(m=1\) and the local tree
bound,
\[
\frac{H_Tq_S}{1-q_S}q_P|K_{\tau,\rm Car}|
\le Cs_\alpha^2q_PL_\tau
\le\frac{C_D}{B_\tau}L_\tau.
\]
For the junction payer, use
\eqref{eq:V11-tree-junction-payer}.  Its term without \(A_{T,e}\) is
paid by \eqref{eq:V11-private-two-power}.  For the remaining term,
\eqref{eq:V11-junction-output-bank-comparison} gives
\[
s_\alpha^3A_{T,e}q_Pq_S
\le C\left(
s_\alpha^3Yq_P
+s_\alpha^3H_Tq_Pq_S
\right).
\]
The cubic private debit bounds the first summand by \(C/Y^2\).  The
quadratic private debit and the ordinary first shared moment give
\[
s_\alpha^3H_Tq_Pq_S
=
\frac{(s_\alpha Y)^2q_P\,
      (s_\alpha H_T)q_S}{Y^2}
\le\frac{C}{Y^2}.
\]
Since \(Y^{-2}\le D/B_\tau\), the junction term is also bounded by
\(C_DL_\tau/B_\tau\).

Now suppose \(B_\tau>DY^2\), and choose \(i\) from
\Cref{lem:V11-nonlocal-branching}.  The shared-output term retains the
needed two powers because
\[
\frac{s_\alpha^3H_Th_i^2q_S}{1-q_S}
\le
\frac{s_\alpha^3(4H_T)^2H_Tq_S}{1-q_S}
\le C
\]
by \(h_i\le H\le4H_T\) and the \(m=3\) case of
\eqref{eq:V11-shared-output-moments}.  Hence
\[
\frac{H_Tq_S}{1-q_S}q_P|K_{\tau,\rm Car}|
\le\frac{C}{h_i^2}L_\tau
\le\frac{C}{B_\tau}L_\tau.
\]

For the junction term, its part without \(A_{T,e}\) is paid by the
ordinary quadratic shared-row moment
\eqref{eq:V10-shared-row-moments}.  For its remaining part use
\eqref{eq:V11-junction-output-bank-comparison}.  The \(Y\)-summand
obeys
\[
s_\alpha^3Yq_Pq_S
=
\frac{(s_\alpha Y)q_P
      (s_\alpha h_i)^2q_S}{h_i^2}
\le\frac{C}{h_i^2}.
\]
For the \(H_T\)-summand, the ordinary cubic shared-output moment gives
\[
s_\alpha^3H_Tq_S
\le\frac{C}{H_T^2}
\le\frac{16C}{h_i^2},
\]
because \(h_i\le4H_T\).  Thus the junction term is also at most
\(CL_\tau/B_\tau\).

Finally \(L_\tau\le\prod_i a_i^{d_i}\), since all \(a_i\ge1\).
Therefore \(L_\tau/B_\tau\le\mathcal W_{\tau,e}\), completing both
branches.
\end{proof}

\begin{corollary}[Full unique-junction Cartan closure]
\label{cor:V11-full-Cartan-closure}
For compact connected semisimple \(G\), every unique-junction complete
Cartan contribution obeys the quartic marked-tree atlas, uniformly in
support, representations, multiplicities, and the established
weak-coupling range.
\end{corollary}

\begin{proof}
Outside output priority use
\Cref{cor:V10-Cartan-outside-output}.  Inside output priority, use
\Cref{thm:V11-local-output-payment} when some row is local-dominant.
Otherwise apply
\Cref{thm:V11-nonlocal-output-payment} to each of the sixteen tree
fibres and sum.
\end{proof}

\begin{firewall}[Cartan is not the whole physical carrier]
\label{fire:V11-Cartan-only}
The corollary closes the exact worst-case full-product-capacity Cartan
line, but it is not a proof for non-Cartan output blocks.  On those
blocks one must control the positive carrier by product-state capacity
before taking an operator norm.  The theorem also does not address
groups with cancelling torus charges.
\end{firewall}

\section{V11 ledger and V12 acquisition}
\label{sec:V11-ledger}

\begin{theorem}[Integrated compact-series V11 statement]
\label{thm:V11-integrated}
Preserving compact V1--V10, V11 proves:
\begin{enumerate}[label=\textup{(V11.\arabic*)},leftmargin=6.5em]
\item the retained-failure product compiler
      \eqref{eq:V11-retained-failure};
\item bankwise shared-output moments through cubic order,
      \eqref{eq:V11-shared-output-moments};
\item exact reduction of output priority to one junction bank and one
      shared-output bank, with no private-coordinate sum;
\item payment of local-dominant output-priority Cartan blocks;
\item payment of all-nonlocal output-priority Cartan tree fibres; and
\item full unique-junction Cartan closure for compact connected
      semisimple gauge groups.
\end{enumerate}
\end{theorem}

\begin{proof}
These are
\Cref{lem:V11-retained-failure,thm:V11-shared-output-moments,
lem:V11-two-payer-banks,thm:V11-local-output-payment,
thm:V11-nonlocal-output-payment,cor:V11-full-Cartan-closure}.
\end{proof}

\begin{openproblem}[V12 mandate: non-Cartan capacity at one junction]
\label{op:V12-mandate}
\begin{enumerate}
\item Return to the positive local junction payer carrier
      \(J_e^{\rm pay}\) and shared-output payer carrier
      \(D_f^{\rm pay}\) of V2--V3.
\item Split off the Cartan isotypic component, now closed by V11.
\item For the orthogonal non-Cartan carrier, seek a quantitative
      product-state capacity gain tied to the Casimir drop from the
      Cartan highest weight.
\item Test the proposed gain on the V1 dual-singlet family, where
      capacity is exactly \(1/d_R\), and on near-Cartan weights, where
      no dimension-only exponential gain is possible.
\item Insert the capacity estimate before raw pinching and before
      summing output multiplicities.
\item Close one non-Cartan spectral band or state the exact
      representation-theoretic capacity profile still missing.
\end{enumerate}
\end{openproblem}

\begin{firewall}[Claim boundary after compact V11]
\label{fire:V11-current-boundary}
V11 closes the complete unique-junction Cartan sector.  It does not
close non-Cartan physical blocks, the full embedded \([4]\) sector,
multiple junctions, normalized collision aggregation, \([3,3]\), global
quartic or all-order MTA, continuum Yang--Mills existence, or a
positive mass gap.
\end{firewall}

\section{Append-only record for compact V11}
\label{sec:V11-record}

V11 was appended to the complete compact V10 file.  The exact V10
parent had \(3991\) newline-delimited source records, \(136747\) bytes,
and SHA--256 digest
\[
\texttt{6e2f86ae8f607f6d584f64c936ef84f20ed014b73e0fea06d29dda1db8341ea3}.
\]
No compact V10 mathematical content was changed.  The sole final
\verb|\end{document}| is restored below.

% =============================================================================
% END APPENDED COMPACT VERSION V11 MATERIAL
% =============================================================================

% The compact V11 document terminator is intentionally disabled here:
% \end{document}

% =============================================================================
% BEGIN APPENDED COMPACT VERSION V12 MATERIAL
% =============================================================================

\chapter{V12: A sharp non-Cartan capacity regression}
\label{ch:V12}

\section{Why Casimir drop is not a capacity parameter}
\label{sec:V12-purpose}

V11 closes the complete unique-junction Cartan sector.  The first V12
proposal was to control the orthogonal sector by a product-state
capacity gain depending on its Casimir drop below the Cartan highest
weight.  The proposal is false.  The obstruction already appears in
the elementary Clebsch--Gordan decomposition of
\(SU(2)\), and it persists in a four-row physical coupling path.

\section{Exact capacity in \(\boldsymbol{\frac12\otimes j}\)}
\label{sec:V12-two-factor}

Let \(V_{1/2}\cong\mathbb C^2\) and \(V_j\) be the spin-\(j\)
irreducible representation, \(j\in\frac12\mathbb N\), \(j\ge1/2\).
Write
\[
V_{1/2}\otimes V_j
=V_{j+1/2}\oplus V_{j-1/2},
\]
and denote the two orthogonal projections by \(P_+\) and \(P_-\).
Product-state capacity in this section is taken across these two
tensor factors.

\begin{theorem}[Exact lower-channel product-state capacity]
\label{thm:V12-lower-channel-capacity}
One has
\begin{equation}
\boxed{
\kappa_\otimes(P_-)=\frac{2j}{2j+1}.}
\label{eq:V12-lower-channel-capacity}
\end{equation}
In particular, this non-Cartan capacity tends to one as
\(j\to\infty\).
\end{theorem}

\begin{proof}
Let \(S=(S_1,S_2,S_3)\) and \(J=(J_1,J_2,J_3)\) be the spin operators
on \(V_{1/2}\) and \(V_j\), and put
\[
X:=2\sum_{a=1}^3S_a\otimes J_a.
\]
The total-Casimir identity shows that \(X\) has eigenvalue \(j\) on
\(V_{j+1/2}\) and eigenvalue \(-(j+1)\) on \(V_{j-1/2}\).  Therefore
\begin{equation}
P_-=\frac{jI-X}{2j+1}.
\label{eq:V12-projector-spin-formula}
\end{equation}

For a pure product vector \(u\otimes v\), let \(n\in S^2\) be the
Bloch vector of \(u\), so
\(\langle u,Su\rangle=n/2\).  Then
\[
\langle u\otimes v,X(u\otimes v)\rangle
=n\cdot\langle v,Jv\rangle.
\]
The operator \(n\cdot J\) has spectrum
\(-j,-j+1,\ldots,j\), hence this expectation is at least \(-j\).
Equation \eqref{eq:V12-projector-spin-formula} gives
\[
\langle P_-\rangle_{u\otimes v}
\le\frac{j-(-j)}{2j+1}.
\]
Equality is attained by taking the qubit spin along \(n\) and the
spin-\(j\) coherent lowest-weight vector in direction \(n\).
Mixed product states are convex combinations of pure product states,
so they cannot increase the supremum.
\end{proof}

\begin{corollary}[The exact two-channel weighted capacity]
\label{cor:V12-weighted-two-channel}
For \(c_+,c_-\ge0\),
\begin{equation}
\boxed{
\kappa_\otimes(c_+P_++c_-P_-)
=
\begin{cases}
c_+,&c_+\ge c_-,\\[3pt]
\displaystyle\frac{c_++2jc_-}{2j+1},&c_-\ge c_+.
\end{cases}}
\label{eq:V12-weighted-two-channel}
\end{equation}
\end{corollary}

\begin{proof}
Every product-state expectation is
\[
c_++(c_--c_+)p_-,
\qquad
p_-:=\Tr(P_-\rho_1\otimes\rho_2).
\]
The proof of
\Cref{thm:V12-lower-channel-capacity} gives
\(0\le p_-\le2j/(2j+1)\).  The value zero is attained by aligned
highest-weight coherent vectors, and the upper endpoint is attained by
oppositely aligned coherent vectors.  Maximize this affine function of
\(p_-\).
\end{proof}

\begin{assessment}[Capacity acts on the weighted carrier, not on
individual projectors]
\label{ass:V12-weighted-not-projector}
Equation \eqref{eq:V12-weighted-two-channel} is the finite-dimensional
model of the V1 capacity gate.  The Clebsch--Gordan probabilities and
the channel coefficients \(c_\pm\) must remain in the same expectation.
Replacing the weighted sum by separate estimates
\(c_\pm\kappa_\otimes(P_\pm)\) discards their normalization and can
create an output-count loss.
\end{assessment}

\section{A four-row non-Cartan physical block with capacity near one}
\label{sec:V12-four-row}

Fix positive spins \(k,l\).  Couple the first two factors through the
lower channel and then take the Cartan channel at the remaining two
couplings:
\begin{equation}
\left(V_{1/2}\otimes V_j\right)_{j-1/2}
\otimes V_k\otimes V_l
\longrightarrow V_{j-1/2+k+l}.
\label{eq:V12-four-row-path}
\end{equation}
Let \(P_{j;k,l}^{\rm path}\) be the orthogonal physical projection
onto this resolved coupling-path block in
\(V_{1/2}\otimes V_j\otimes V_k\otimes V_l\).

\begin{theorem}[Four nontrivial rows inherit the near-unit obstruction]
\label{thm:V12-four-row-obstruction}
Across the original four rows,
\begin{equation}
\boxed{
\kappa_\otimes(P_{j;k,l}^{\rm path})
\ge\frac{2j}{2j+1}.}
\label{eq:V12-four-row-obstruction}
\end{equation}
The final spin in \eqref{eq:V12-four-row-path} is below the Cartan final
spin by one, while their Casimir difference is
\begin{equation}
C_2(j+1/2+k+l)-C_2(j-1/2+k+l)
=2(j+k+l)+1.
\label{eq:V12-Casimir-drop}
\end{equation}
Thus the Casimir drop tends to infinity while the non-Cartan
product-state capacity tends to one.
\end{theorem}

\begin{proof}
Choose a common quantization axis.  In the first row take the
highest-weight qubit, and in the second row take the spin-\(j\)
lowest-weight vector.  The proof of
\Cref{thm:V12-lower-channel-capacity} shows that projection of this
product onto \(V_{j-1/2}\) has squared norm \(2j/(2j+1)\); the resulting
normalized vector is the lowest-weight vector of \(V_{j-1/2}\).
Choose lowest-weight vectors in \(V_k\) and \(V_l\).  Their tensor
product with that projected vector lies on the lowest-weight line of
the successive Cartan couplings, so the remaining path projections
preserve it.  Hence the four-row product input has the expectation
claimed in \eqref{eq:V12-four-row-obstruction}.

For \(SU(2)\), \(C_2(r)=r(r+1)\).  Put \(a=j+k+l\).  The difference in
\eqref{eq:V12-Casimir-drop} is
\[
(a+1/2)(a+3/2)-(a-1/2)(a+1/2)=2a+1.
\]
\end{proof}

\begin{corollary}[No non-Cartan rarity from dimension or Casimir drop]
\label{cor:V12-no-nonCartan-rarity}
There is no function \(\varepsilon(\Delta)\to0\) such that every
non-Cartan physical output projector with Cartan Casimir drop at least
\(\Delta\) has product-state capacity at most
\(\varepsilon(\Delta)\).  In particular, removing the Cartan projector
does not make the operator-norm bound asymptotically improvable by a
universal representation-dimension factor.
\end{corollary}

\begin{proof}
Use the family in
\Cref{thm:V12-four-row-obstruction}.  Its drop tends to infinity and
its capacity has lower bound tending to one.
\end{proof}

\begin{firewall}[Scope of the regression]
\label{fire:V12-regression-scope}
The example disproves a universal projector-rarity mechanism.  It does
not say that the actual coefficient-weighted Yang--Mills carrier is
large on this family.  Its local cumulant and payer coefficients may
vary between adjacent channels, and that variation is precisely the
information which a successful capacity proof must retain.
\end{firewall}

\section{The corrected non-Cartan target}
\label{sec:V12-corrected-target}

For fixed input representations, write the non-Cartan part of the
positive physical carrier as
\begin{equation}
\mathcal Q_{\rm nCar}
=\sum_{\nu\notin{\rm Car}}c_\nu P_\nu,
\qquad
c_\nu:=
\frac{\Xi(\boldsymbol U_\nu)}
     {1-q_{\alpha,\boldsymbol U_\nu}}
\|K_\nu\|.
\label{eq:V12-nonCartan-carrier}
\end{equation}
The required statement is still
\[
\sup_{\rho_1,\ldots,\rho_4}
\sum_{\nu\notin{\rm Car}}c_\nu
\Tr(P_\nu\rho_1\otimes\cdots\otimes\rho_4)
\le C\sum_{\tau,\ell}\mathcal W_{\tau,\ell},
\]
but V12 proves that it cannot be obtained by multiplying a pointwise
coefficient cap by a decaying capacity profile depending only on
Casimir drop.

\begin{proposition}[Adjacent-channel variation is the first upstream
quantity]
\label{prop:V12-adjacent-target}
On the \(SU(2)\) two-channel regression, a positive central carrier is
completely controlled on product states by
\begin{equation}
c_+,\qquad
\frac{c_++2jc_-}{2j+1}.
\label{eq:V12-adjacent-data}
\end{equation}
Consequently, any proposed non-Cartan capacity theorem must control the
actual adjacent coefficient \(c_-\), or a cancellation-preserving
combination of \(c_-\) and \(c_+\); projector capacity alone supplies
no small factor.
\end{proposition}

\begin{proof}
This is exactly
\Cref{cor:V12-weighted-two-channel}.  Since
\(2j/(2j+1)\to1\), no estimate of \(P_-\) alone can replace control of
its coefficient.
\end{proof}

\begin{assessment}[Upstream change of direction]
\label{ass:V12-direction}
The next object is not a universal capacity profile for
\(P_\nu\).  It is the coefficient-weighted Clebsch--Gordan expectation
of the complete local payer.  The first test should compute or bound
the adjacent \(SU(2)\) coefficients before summing channels.  If their
difference has a Wilson finite-difference representation, the exact
formula \eqref{eq:V12-weighted-two-channel} can retain it; if it does
not, then the hot non-Cartan carrier remains an explicit independent
input to the programme.
\end{assessment}

\section{V12 ledger and V13 acquisition}
\label{sec:V12-ledger}

\begin{theorem}[Integrated compact-series V12 statement]
\label{thm:V12-integrated}
Preserving compact V1--V11, V12 proves:
\begin{enumerate}[label=\textup{(V12.\arabic*)},leftmargin=6.5em]
\item the exact lower-channel capacity
      \(2j/(2j+1)\) in \(V_{1/2}\otimes V_j\);
\item the exact capacity of every positive two-channel central carrier;
\item a four-nontrivial-row physical path whose non-Cartan capacity
      tends to one;
\item divergence of its Cartan Casimir drop;
\item failure of every dimension- or Casimir-drop-only non-Cartan
      rarity principle; and
\item reduction of the next capacity attack to coefficient-weighted
      adjacent-channel data.
\end{enumerate}
\end{theorem}

\begin{proof}
These are
\Cref{thm:V12-lower-channel-capacity,
cor:V12-weighted-two-channel,thm:V12-four-row-obstruction,
cor:V12-no-nonCartan-rarity,prop:V12-adjacent-target}.
\end{proof}

\begin{openproblem}[V13 mandate: the adjacent weighted payer]
\label{op:V13-mandate}
\begin{enumerate}
\item Specialize the one-link local payer to
      \(V_{1/2}\otimes V_j\), retaining both output channels and the
      complete cumulant coefficient.
\item Express the two positive carrier coefficients \(c_+,c_-\) before
      taking any channel norm or triangle inequality.
\item Insert them into
      \eqref{eq:V12-weighted-two-channel}.
\item Test whether the weighted average has a finite-difference or
      covariance representation that retains the local tree powers.
\item If the adjacent calculation closes, lift it along the resolved
      four-row path \eqref{eq:V12-four-row-path}; otherwise isolate the
      exact coefficient estimate which is absent from the one-link
      Wilson hypotheses.
\end{enumerate}
\end{openproblem}

\begin{firewall}[Claim boundary after compact V12]
\label{fire:V12-current-boundary}
V12 is a rigorous no-go and target correction, not a non-Cartan
closure theorem.  The unique-junction Cartan sector remains closed by
V11.  Non-Cartan physical blocks, the full embedded \([4]\) sector,
multiple junctions, normalized collision aggregation, \([3,3]\), global
quartic or all-order MTA, continuum Yang--Mills existence, and a
positive mass gap remain open.
\end{firewall}

\section{Append-only record for compact V12}
\label{sec:V12-record}

V12 was appended to the complete compact V11 file.  The exact V11
parent had \(4496\) newline-delimited source records, \(152085\) bytes,
and SHA--256 digest
\[
\texttt{694d2d7c6d5ea752faed5bffa99d8cee63660fb0eba8bc9e1bf81b4d01164dcf}.
\]
No compact V11 mathematical content was changed.  The sole final
\verb|\end{document}| is restored below.

% =============================================================================
% END APPENDED COMPACT VERSION V12 MATERIAL
% =============================================================================

% The compact V12 document terminator is intentionally disabled here:
% \end{document}

% =============================================================================
% BEGIN APPENDED COMPACT VERSION V13 MATERIAL
% =============================================================================

\chapter{V13: Exact partial-trace capacity and the retention split}
\label{ch:V13}

\section{Invariant projections have scalar one-row marginals}
\label{sec:V13-partial-trace}

Let \(G\) be compact, let \(R_i\) be irreducible unitary
representations on \(H_i\), and put
\[
H:=\bigotimes_{i=1}^mH_i,
\qquad d_i:=\dim H_i.
\]
All product-state capacities below use this original row grouping.

\begin{lemma}[Scalar marginal of a diagonal-invariant operator]
\label{lem:V13-scalar-marginal}
If \(Q\ge0\) commutes with the diagonal representation
\(\bigotimes_iR_i(g)\) for every \(g\in G\), then
\begin{equation}
\boxed{
\Tr_{\widehat i}Q=\frac{\Tr Q}{d_i}I_{H_i}}
\label{eq:V13-scalar-marginal}
\end{equation}
for every row \(i\).  Here \(\Tr_{\widehat i}\) traces all rows except
\(i\).
\end{lemma}

\begin{proof}
Write \(R_{\widehat i}(g)=\bigotimes_{j\ne i}R_j(g)\).  Diagonal
invariance gives
\[
(R_i(g)\otimes I)Q(R_i(g)^*\otimes I)
=(I\otimes R_{\widehat i}(g)^*)Q
 (I\otimes R_{\widehat i}(g)).
\]
Partial trace is invariant under unitary conjugation on the traced
factor.  Thus \(\Tr_{\widehat i}Q\) commutes with \(R_i(g)\) for every
\(g\).  Schur's lemma makes it a scalar multiple of the identity.
Taking the trace identifies the scalar as \(\Tr Q/d_i\).
\end{proof}

\begin{theorem}[Exact trace-ratio upper bound for product capacity]
\label{thm:V13-trace-ratio-capacity}
Under the hypotheses of
\Cref{lem:V13-scalar-marginal},
\begin{equation}
\boxed{
\kappa_\otimes(Q)
\le
\min\left\{\|Q\|,
\frac{\Tr Q}{d_1},\ldots,\frac{\Tr Q}{d_m}\right\}.}
\label{eq:V13-trace-ratio-capacity}
\end{equation}
In particular, if \(P\) is a diagonal-invariant projection of rank
\(r\), then
\begin{equation}
\boxed{
\kappa_\otimes(P)
\le\min\left\{1,\frac r{d_1},\ldots,\frac r{d_m}\right\}.}
\label{eq:V13-projection-rank-capacity}
\end{equation}
\end{theorem}

\begin{proof}
The operator-norm bound follows directly from the definition.  For a
product density matrix
\(\rho=\rho_i\otimes\rho_{\widehat i}\), one has
\(0\le\rho_{\widehat i}\le I\).  Positivity of \(Q\) therefore gives
\begin{align*}
\Tr(Q\rho)
&\le\Tr\!\left[Q(\rho_i\otimes I)\right]\\
&=\Tr\!\left[\rho_i\Tr_{\widehat i}Q\right]
=\frac{\Tr Q}{d_i}.
\end{align*}
Take the supremum and then the minimum over \(i\).  For a projection,
\(\|P\|=1\) and \(\Tr P=r\).
\end{proof}

\begin{corollary}[One physical irreducible copy]
\label{cor:V13-one-copy}
Let \(P_{U,a}\) project onto one resolved copy of an irreducible output
\(U\) inside \(\bigotimes_iR_i\).  Then
\begin{equation}
\boxed{
\kappa_\otimes(P_{U,a})
\le
\min\left\{1,\frac{d_U}{d_{R_1}},\ldots,
\frac{d_U}{d_{R_m}}\right\}.}
\label{eq:V13-one-copy}
\end{equation}
If a physical block retains \(r_a\) copies, replace \(d_U\) by
\(r_ad_U\).
\end{corollary}

\begin{proof}
The range of \(P_{U,a}\) is diagonal-invariant and has rank \(d_U\).
Apply \Cref{thm:V13-trace-ratio-capacity}.  The multi-copy statement
uses its actual rank.
\end{proof}

\begin{assessment}[This is not the withdrawn dimension law]
\label{ass:V13-raw-not-formation}
Equation \eqref{eq:V13-one-copy} is a theorem about the overlap of a
specified invariant projection with arbitrary product density
matrices.  It neither replaces a Fourier coefficient by
\(d_R^{-1}\) nor asserts a representation-dimension formation law.
The exact input coefficients remain in the product-state expectation,
as required by V1.
\end{assessment}

\section{The two previous regressions are extremal cases}
\label{sec:V13-regressions}

\begin{corollary}[Dual singlet and near-Cartan lower channel]
\label{cor:V13-two-extremes}
\begin{enumerate}[label=\textup{(\roman*)},leftmargin=4em]
\item In \(R\otimes R^*\), the singlet rank is one, so
\(\kappa_\otimes(P_{\boldsymbol1})\le1/d_R\); the V1 product vector
shows equality.
\item In \(V_{1/2}\otimes V_j\), the lower channel has rank \(2j\), so
\[
\kappa_\otimes(P_-)\le
\min\left\{1,j,\frac{2j}{2j+1}\right\}
=\frac{2j}{2j+1};
\]
V12 shows equality.
\end{enumerate}
\end{corollary}

\begin{proof}
Apply \eqref{eq:V13-projection-rank-capacity} and the matching lower
bounds already proved in
\Cref{prop:V1-singlet-capacity,thm:V12-lower-channel-capacity}.
\end{proof}

\begin{remark}[Why near-unit non-Cartan capacity is forced]
\label{rem:V13-near-unit-forced}
The V12 lower channel fills all but one dimension relative to the
large input row.  Its capacity \(2j/(2j+1)\) is not an accident of a
loose variational estimate: it saturates the exact one-row marginal
bound.
\end{remark}

\section{A coefficient-weighted carrier can be bounded before blocks
are separated}
\label{sec:V13-weighted-carrier}

\begin{corollary}[Trace budget for a complete invariant carrier]
\label{cor:V13-weighted-trace-budget}
Let
\[
Q=\sum_\nu c_\nu P_\nu,\qquad c_\nu\ge0,
\]
where the \(P_\nu\) are mutually orthogonal diagonal-invariant physical
projections.  Then
\begin{equation}
\boxed{
\kappa_\otimes(Q)
\le
\min\left\{
\sup_\nu c_\nu,\,
\min_i\frac1{d_i}
\sum_\nu c_\nu\operatorname{rank}P_\nu
\right\}.}
\label{eq:V13-weighted-trace-budget}
\end{equation}
\end{corollary}

\begin{proof}
The operator norm of the orthogonal block sum is
\(\sup_\nu c_\nu\), and its trace is
\(\sum_\nu c_\nu\operatorname{rank}P_\nu\).
Apply \Cref{thm:V13-trace-ratio-capacity} directly to the complete
operator \(Q\), before estimating individual block capacities.
\end{proof}

\begin{firewall}[The trace budget exposes, rather than hides, output
multiplicity]
\label{fire:V13-trace-budget}
The second term of \eqref{eq:V13-weighted-trace-budget} contains the
actual weighted output rank.  It is useful only if that complete trace
sum has an independent Wilson or fusion estimate.  Bounding each
\(c_\nu\) uniformly and then counting the projections would reintroduce
the output multiplicity which the V1 raw-pinching theorem was designed
to avoid.
\end{firewall}

\section{The exact \(SU(2)\) retention/cancellation split}
\label{sec:V13-retention}

For \(SU(2)\), write
\[
A_j:=1+j(j+1).
\]
Then the dimension and Casimir satisfy the exact identity
\begin{equation}
(2j+1)^2=4A_j-3.
\label{eq:V13-SU2-dimension-Casimir}
\end{equation}

\begin{theorem}[Capacity versus output-mass retention]
\label{thm:V13-SU2-retention}
Let \(P_J\) be one resolved spin-\(J\) output copy in a tensor product
of irreducible \(SU(2)\) rows, and let
\(j_*=\max_i j_i\).  Then
\begin{equation}
\boxed{
\kappa_\otimes(P_J)
\le
\min\left\{
1,\frac{2J+1}{2j_*+1}
\right\}
=
\min\left\{
1,\sqrt{\frac{4A_J-3}{4A_{j_*}-3}}
\right\}.}
\label{eq:V13-SU2-retention}
\end{equation}
Consequently, for \(0<\theta<1\), every output satisfies the exact
dichotomy
\begin{equation}
\boxed{
\kappa_\otimes(P_J)\le\theta
\quad\hbox{or}\quad
4A_J-3>\theta^2(4A_{j_*}-3).}
\label{eq:V13-retention-dichotomy}
\end{equation}
\end{theorem}

\begin{proof}
Apply \eqref{eq:V13-one-copy} to a row of maximal dimension.  Substitute
\eqref{eq:V13-SU2-dimension-Casimir}.  If the capacity bound on the
right of \eqref{eq:V13-SU2-retention} exceeds \(\theta\), its squared
dimension ratio gives the second alternative.
\end{proof}

\begin{assessment}[Meaning of the split]
\label{ass:V13-retention-meaning}
Strongly cancelling outputs have small product overlap with at least
one large input row.  Outputs with near-unit capacity must retain a
fixed fraction of that row's Casimir scale, so their own Wilson
multiplier can potentially pay the large input mass.  This is the
correct qualitative replacement for the false claim that every
non-Cartan output is rare.  A quantitative MTA proof must still combine
the two alternatives inside the coefficient-weighted carrier.
\end{assessment}

\section{The next positive carrier}
\label{sec:V13-next-carrier}

For one local junction tree, split the non-Cartan physical carrier by
\[
\mathcal N_{\rm can}(\theta)
:=\{\nu: d_{U_\nu}\le\theta d_{R_*}\},
\qquad
\mathcal N_{\rm ret}(\theta)
:=\{\nu: d_{U_\nu}>\theta d_{R_*}\},
\]
where \(R_*\) is a largest local input row.  On the cancelling family,
\eqref{eq:V13-one-copy} gives a factor \(\theta\) per resolved copy.
On the retaining family, \eqref{eq:V13-retention-dichotomy} gives an
output Casimir floor.  Neither statement by itself supplies the two
branching denominators, but the split identifies where they may come
from without contradiction:
\begin{enumerate}[label=\textup{(\roman*)},leftmargin=4em]
\item a complete weighted trace estimate on
      \(\mathcal N_{\rm can}(\theta)\); and
\item output Wilson moments, kept inside the positive carrier, on
      \(\mathcal N_{\rm ret}(\theta)\).
\end{enumerate}

\begin{firewall}[Group scope]
\label{fire:V13-group-scope}
\Cref{thm:V13-trace-ratio-capacity} holds for every compact group and
irreducible row representations.  The explicit
dimension--Casimir identity and retention split in
\Cref{thm:V13-SU2-retention} are \(SU(2)\)-specific.  A general compact
simple group requires its Weyl-dimension analogue, with rank-dependent
wall behavior treated explicitly.
\end{firewall}

\section{V13 ledger and V14 acquisition}
\label{sec:V13-ledger}

\begin{theorem}[Integrated compact-series V13 statement]
\label{thm:V13-integrated}
Preserving compact V1--V12, V13 proves:
\begin{enumerate}[label=\textup{(V13.\arabic*)},leftmargin=6.5em]
\item scalar one-row marginals for every positive
      diagonal-invariant carrier;
\item the exact trace-ratio product-capacity upper bound
      \eqref{eq:V13-trace-ratio-capacity};
\item its one-copy physical-channel form
      \eqref{eq:V13-one-copy};
\item simultaneous recovery of the exact dual-singlet and
      \(SU(2)\) lower-channel capacities;
\item the complete weighted trace budget
      \eqref{eq:V13-weighted-trace-budget}; and
\item the exact \(SU(2)\) capacity/output-retention dichotomy
      \eqref{eq:V13-retention-dichotomy}.
\end{enumerate}
\end{theorem}

\begin{proof}
These are
\Cref{lem:V13-scalar-marginal,thm:V13-trace-ratio-capacity,
cor:V13-one-copy,cor:V13-two-extremes,
cor:V13-weighted-trace-budget,thm:V13-SU2-retention}.
\end{proof}

\begin{openproblem}[V14 mandate: one weighted cancellation band]
\label{op:V14-mandate}
\begin{enumerate}
\item Work first in \(SU(2)\), where
      \eqref{eq:V13-retention-dichotomy} is exact.
\item For a fixed dyadic dimension ratio
      \(2^{-n-1}<d_U/d_{R_*}\le2^{-n}\), keep all physical output
      copies in one positive carrier.
\item On the cancelling band, estimate the complete weighted trace
      \(\sum_\nu c_\nu\operatorname{rank}P_\nu/d_{R_*}\) without a
      multiplicity count.
\item On the retaining complement, use its output Casimir floor before
      applying a Wilson moment.
\item Preserve the V11 bankwise resolvent split and the V10 input-bank
      moments.
\item Close one dyadic non-Cartan band, or show by an explicit
      Clebsch--Gordan family which weighted trace estimate is still
      missing.
\end{enumerate}
\end{openproblem}

\begin{firewall}[Claim boundary after compact V13]
\label{fire:V13-current-boundary}
V13 supplies an exact non-Cartan capacity mechanism but does not yet
bound the complete coefficient-weighted hot carrier.  Thus non-Cartan
physical blocks, the full embedded \([4]\) sector, multiple junctions,
normalized collision aggregation, \([3,3]\), global quartic or
all-order MTA, continuum Yang--Mills existence, and a positive mass gap
remain open.
\end{firewall}

\section{Append-only record for compact V13}
\label{sec:V13-record}

V13 was appended to the complete compact V12 file.  The exact V12
parent had \(4830\) newline-delimited source records, \(164005\) bytes,
and SHA--256 digest
\[
\texttt{076cc1b4ba5a5c228c2cd8bd67d21d38a767affe971ca37951aa56a8ba9d0c64}.
\]
No compact V12 mathematical content was changed.  The sole final
\verb|\end{document}| is restored below.

% =============================================================================
% END APPENDED COMPACT VERSION V13 MATERIAL
% =============================================================================

% The compact V13 document terminator is intentionally disabled here:
% \end{document}

% =============================================================================
% BEGIN APPENDED COMPACT VERSION V14 MATERIAL
% =============================================================================

\chapter{V14: The retentive non-Cartan spectator sector}
\label{ch:V14}

\section{Retentive resolved spectator outputs}
\label{sec:V14-retentive-definition}

Fix \(0<\theta\le1\).  At an off-junction shared coordinate \(f\), let
\(U_f\) be a resolved nontrivial physical output, let
\[
b_f:=A_{U_f}>0,\qquad r_f:=q_{\alpha,U_f},
\]
and retain the input masses
\[
a_{i,f}:=1+C_2(R_{i,f}),\qquad i\in I_f.
\]
Call this resolved coordinate \(\theta\)-retentive when
\begin{equation}
\boxed{
b_f\ge\theta\max_{i\in I_f}a_{i,f}.}
\label{eq:V14-retentive-coordinate}
\end{equation}
A resolved unique-junction block is
\emph{shared-\(\theta\)-retentive} when every one of its shared
spectator coordinates is \(\theta\)-retentive.  The definition imposes
no Cartan condition at the junction or at the spectators.

\begin{remark}[Trivial outputs belong to the cancelling sector]
\label{rem:V14-trivial-cancelling}
A trivial shared output has no positive spectral gap mass and cannot
satisfy \eqref{eq:V14-retentive-coordinate} for a nontrivial active
input.  It is therefore routed to the cancelling sector left open
below, where product-state capacity rather than Wilson damping must
control it.
\end{remark}

\section{Input moments inherited from retentive outputs}
\label{sec:V14-retentive-moments}

For a fixed row \(i\), retain \(F_i,h_i\) from
\eqref{eq:V10-shared-row-mass}, and write
\[
q_S:=\prod_{\substack{f\ne e\\|I_f|\ge2}}r_f,
\qquad
H_T:=\sum_{\substack{f\ne e\\|I_f|\ge2}}b_f,
\qquad
H:=\sum_i h_i.
\]

\begin{theorem}[Retentive shared-row moments]
\label{thm:V14-retentive-row-moments}
On a shared-\(\theta\)-retentive block, for \(k=1,2,3\),
\begin{align}
\boxed{
(s_\alpha h_i)^kq_S\le C_k\theta^{-k},}
\label{eq:V14-retentive-row-moments}\\
\boxed{
(s_\alpha H)^kq_S\le C_k(4/\theta)^k.}
\label{eq:V14-retentive-total-moments}
\end{align}
Moreover
\begin{equation}
\boxed{H\le\frac4\theta H_T.}
\label{eq:V14-retentive-mass-comparison}
\end{equation}
\end{theorem}

\begin{proof}
For \(f\in F_i\), put \(x_f=s_\alpha\theta a_{i,f}\).
Retentivity gives \(x_f\le s_\alpha b_f\), and the one-link Wilson
debits give
\[
x_f^kr_f\le C_k(1-r_f),
\qquad 1\le k\le3.
\]
The V7 product-moment compiler applied on \(F_i\), followed by the
extra spectator factors in \(q_S\), proves
\eqref{eq:V14-retentive-row-moments}.

At coordinate \(f\), at most four rows are active and each
\(a_{i,f}\le b_f/\theta\).  Thus
\[
\sum_{i\in I_f}a_{i,f}\le\frac4\theta b_f.
\]
Summing proves \eqref{eq:V14-retentive-mass-comparison}.  The ordinary
moments of \(H_T\), or the same product compiler, then prove
\eqref{eq:V14-retentive-total-moments}.
\end{proof}

\begin{corollary}[Retentive output-bank resolvent moments]
\label{cor:V14-retentive-output-resolvent}
For \(k=1,2,3\),
\begin{equation}
\boxed{
\frac{(s_\alpha H_T)^kq_S}{1-q_S}\le C_k.}
\label{eq:V14-retentive-output-resolvent}
\end{equation}
\end{corollary}

\begin{proof}
This is \Cref{thm:V11-shared-output-moments}; its proof uses only the
one-link output debits and does not require Cartan outputs.
\end{proof}

\section{All-nonlocal blocks below output priority}
\label{sec:V14-below-output}

For a resolved physical tree block, let
\[
B_\tau:=\prod_i\Xi_i^{d_i-1},
\qquad
L_\tau:=\prod_i a_i^{d_i/2}.
\]
The V41 local estimate and contraction of the spectator recouplings
give
\begin{equation}
\|K_{\tau,U_e}\|\le Cs_\alpha^3L_\tau.
\label{eq:V14-local-tree-bound}
\end{equation}

\begin{theorem}[Retentive non-Cartan shared payment below
\textup{(O)}]
\label{thm:V14-retentive-below-output}
Fix \(D\ge16\).  Suppose a shared-\(\theta\)-retentive physical block
satisfies
\begin{equation}
\Xi_T\le DY,\qquad
B_\tau>DY^2,\qquad
a_i<\frac14\Xi_i\quad(1\le i\le4).
\label{eq:V14-below-output-hypotheses}
\end{equation}
Then
\begin{equation}
\boxed{
\frac{\Xi_T}{1-q_{\alpha,\boldsymbol T}}\,
q_Pq_S\|K_{\tau,U_e}\|
\le C_{D,\theta}\mathcal W_{\tau,e}.}
\label{eq:V14-retentive-below-output}
\end{equation}
\end{theorem}

\begin{proof}
Choose a maximal internal row \(i\).  The purely scalar argument of
\Cref{lem:V10-maximal-shared-row} gives
\[
h_i>\frac12\Xi_i,\qquad B_\tau<4h_i^2.
\]
The global output multiplier is a product of its coordinate
multipliers, so
\[
1-q_{\alpha,\boldsymbol T}\ge1-q_P.
\]
The bankwise private first-moment debit and \(\Xi_T\le DY\) imply
\[
\frac{\Xi_Tq_P}{1-q_{\alpha,\boldsymbol T}}
\le\frac{C_D}{s_\alpha}.
\]
Use \eqref{eq:V14-local-tree-bound} and then the quadratic retentive
row moment:
\begin{align*}
\frac{\Xi_T}{1-q_{\alpha,\boldsymbol T}}\,
q_Pq_S\|K_{\tau,U_e}\|
&\le C_Ds_\alpha^2q_SL_\tau\\
&\le\frac{C_{D,\theta}}{h_i^2}L_\tau
\le\frac{C_{D,\theta}}{B_\tau}L_\tau.
\end{align*}
Finally \(L_\tau\le\prod_i a_i^{d_i}\), which gives the marked weight.
No product-state or multiplicity estimate is used on this retentive
block.
\end{proof}

\section{All-nonlocal blocks in output priority}
\label{sec:V14-output-priority}

\begin{lemma}[Non-Cartan junction mass has the required upper bound]
\label{lem:V14-junction-upper}
For every resolved junction output \(U_e\) occurring in
\(\bigotimes_iR_{i,e}\),
\begin{equation}
A_{U_e}\le C_{\rm fus}\sum_i a_i.
\label{eq:V14-junction-upper}
\end{equation}
If \(a_i<\Xi_i/4\) for every row on a
shared-\(\theta\)-retentive block, then
\begin{equation}
\boxed{
A_{U_e}\le C_\theta(Y+H_T).}
\label{eq:V14-junction-bank-upper}
\end{equation}
\end{lemma}

\begin{proof}
The first inequality is the established fusion-Casimir upper bound
\eqref{eq:V3-fusion-output}, applied to the singleton input blocks.
The all-nonlocal condition gives
\[
3a_i<p_i+h_i\le Y+h_i.
\]
Sum in \(i\), apply
\eqref{eq:V14-retentive-mass-comparison}, and then use
\eqref{eq:V14-junction-upper}.
\end{proof}

\begin{theorem}[Retentive non-Cartan output-priority payment]
\label{thm:V14-retentive-output}
Fix \(D\ge16\).  Suppose a shared-\(\theta\)-retentive physical block
satisfies
\[
\Xi_T>DY,\qquad
a_i<\frac14\Xi_i\quad(1\le i\le4).
\]
Then every tree fibre obeys
\begin{equation}
\boxed{
\frac{\Xi_T}{1-q_{\alpha,\boldsymbol T}}\,
q_Pq_S\|K_{\tau,U_e}\|
\le C_{D,\theta}\mathcal W_{\tau,e}.}
\label{eq:V14-retentive-output}
\end{equation}
\end{theorem}

\begin{proof}
V8's nonprivate output split and the scalar allocation used in
\Cref{lem:V11-two-payer-banks} give
\begin{align}
&\frac{\Xi_T}{1-q_{\alpha,\boldsymbol T}}\,
q_Pq_S\|K_{\tau,U_e}\|\notag\\
&\quad\le C_D\left[
\frac{A_{U_e}}{1-q_{\alpha,U_e}}
q_Pq_S\|K_{\tau,U_e}\|
+
\frac{H_Tq_S}{1-q_S}
q_P\|K_{\tau,U_e}\|
\right].
\label{eq:V14-two-bank-split}
\end{align}
The tree-local junction payer proof of
\Cref{lem:V11-tree-junction-payer} is representation-independent:
\begin{equation}
\frac{A_{U_e}}{1-q_{\alpha,U_e}}\|K_{\tau,U_e}\|
\le
Cs_\alpha^2(1+s_\alpha A_{U_e})L_\tau.
\label{eq:V14-tree-junction-payer}
\end{equation}

If \(B_\tau\le DY^2\), the private quadratic moment gives
\[
s_\alpha^2q_P\le C_D/B_\tau.
\]
It pays the first term of
\eqref{eq:V14-tree-junction-payer}.  For its part containing
\(A_{U_e}\), use
\eqref{eq:V14-junction-bank-upper}.  The \(Y\)-piece is paid by the
private cubic moment; the \(H_T\)-piece is paid by the product of the
private quadratic moment and the ordinary shared-output first moment.
Explicitly,
\[
s_\alpha^3Yq_P\le C/Y^2,
\qquad
s_\alpha^3H_Tq_Pq_S\le C/Y^2.
\]
The shared-output payer in \eqref{eq:V14-two-bank-split} is bounded by
\(Cs_\alpha^2q_PL_\tau\) using
\eqref{eq:V14-retentive-output-resolvent} with \(k=1\).
Thus both payer banks are at most
\(C_DL_\tau/B_\tau\).

If \(B_\tau>DY^2\), select the maximal internal row as in
\Cref{lem:V11-nonlocal-branching}; then
\[
h_i>\Xi_i/2,\qquad B_\tau<4h_i^2.
\]
The quadratic estimate
\eqref{eq:V14-retentive-row-moments} pays the first part of
\eqref{eq:V14-tree-junction-payer}.  For the \(A_{U_e}\) part, its
\(Y\)-piece is bounded by the product of the private first moment and
the quadratic \(h_i\)-moment.  Its \(H_T\)-piece is bounded by the
ordinary cubic \(H_T\)-moment.  Hence
\[
s_\alpha^3A_{U_e}q_Pq_S
\le\frac{C_\theta}{h_i^2}.
\]
For the shared-output payer, use
\(h_i\le H\le4H_T/\theta\) and the cubic
resolvent-weighted moment:
\[
\frac{s_\alpha^3H_Th_i^2q_S}{1-q_S}
\le C_\theta.
\]
Thus both banks are again at most
\(C_\theta L_\tau/h_i^2\le
C_\theta L_\tau/B_\tau\).

In both branching cases,
\(L_\tau\le\prod_i a_i^{d_i}\), proving the theorem.
\end{proof}

\begin{corollary}[A genuine non-Cartan unique-junction sector is
closed]
\label{cor:V14-retentive-sector}
For compact connected semisimple \(G\), every
shared-\(\theta\)-retentive unique-junction physical block which is
all-nonlocal,
\[
a_i<\Xi_i/4\qquad(1\le i\le4),
\]
obeys the quartic marked-tree atlas, uniformly for fixed
\(\theta>0\).
\end{corollary}

\begin{proof}
Outside output priority, a tree either satisfies the private-dominant
conditions and is closed by V7, or satisfies
\eqref{eq:V14-below-output-hypotheses} and is closed by
\Cref{thm:V14-retentive-below-output}.  Inside output priority apply
\Cref{thm:V14-retentive-output}.  Sum the sixteen tree fibres.
\end{proof}

\begin{firewall}[The remaining non-Cartan family]
\label{fire:V14-remaining-family}
V14 does not close a block having a strongly cancelling shared
spectator output, nor the non-Cartan local-dominant junction family.
On a cancelling spectator, the multiplier \(q_{\alpha,U_f}\) need not
pay the active input mass.  That branch must use the V13
coefficient-weighted capacity before the shared projection is replaced
by an operator norm.  The local-dominant non-Cartan branch still
requires a complete local payer capacity estimate.
\end{firewall}

\section{V14 ledger and mandatory V15 audit}
\label{sec:V14-ledger}

\begin{theorem}[Integrated compact-series V14 statement]
\label{thm:V14-integrated}
Preserving compact V1--V13, V14 proves:
\begin{enumerate}[label=\textup{(V14.\arabic*)},leftmargin=6.5em]
\item the precise shared-\(\theta\)-retentive spectral sector
      \eqref{eq:V14-retentive-coordinate};
\item support-uniform input moments inherited from resolved output
      multipliers;
\item retentive payment of the all-nonlocal shared branch below output
      priority;
\item retentive payment of both nonprivate banks inside output
      priority; and
\item marked-tree closure of the complete all-nonlocal,
      shared-\(\theta\)-retentive unique-junction sector.
\end{enumerate}
\end{theorem}

\begin{proof}
These are
\Cref{thm:V14-retentive-row-moments,
thm:V14-retentive-below-output,
thm:V14-retentive-output,cor:V14-retentive-sector}.
\end{proof}

\begin{openproblem}[V15 mandate: audit and one cancelling spectator]
\label{op:V15-mandate}
\begin{enumerate}
\item Begin with the mandatory read-only audit of the newest active SM
      and NS notes.
\item Fix one shared coordinate \(f\) for which
      \(A_{U_f}<\theta\max_{i\in I_f}a_{i,f}\).
\item Keep its physical projection inside the complete positive
      carrier and apply the V13 trace-ratio theorem before taking its
      operator norm.
\item In \(SU(2)\), dyadically split the ratio
      \(d_{U_f}/\max_i d_{R_{i,f}}\) and retain the exact
      dimension--Casimir relation.
\item Determine whether the small capacity pays the missing input
      moment without an output-copy count.  If not, isolate the exact
      weighted trace or multiplicity bound which is absent.
\item Do not reuse a retentive Wilson moment on the cancelling branch.
\end{enumerate}
\end{openproblem}

\begin{firewall}[Claim boundary after compact V14]
\label{fire:V14-current-boundary}
V14 closes a non-Cartan all-nonlocal retentive sector, not the full
non-Cartan carrier.  Strongly cancelling shared outputs,
local-dominant non-Cartan blocks, the full embedded \([4]\) sector,
multiple junctions, normalized collision aggregation, \([3,3]\),
global quartic or all-order MTA, continuum Yang--Mills existence, and a
positive mass gap remain open.
\end{firewall}

\section{Append-only record for compact V14}
\label{sec:V14-record}

V14 was appended to the complete compact V13 file.  The exact V13
parent had \(5194\) newline-delimited source records, \(176259\) bytes,
and SHA--256 digest
\[
\texttt{50bc84c821c0f2918f5b7c22b97b0e2b45091a2e4efa6ca0d8ce3622bd379ec3}.
\]
No compact V13 mathematical content was changed.  The sole final
\verb|\end{document}| is restored below.

% =============================================================================
% END APPENDED COMPACT VERSION V14 MATERIAL
% =============================================================================

% The compact V14 document terminator is intentionally disabled here:
% \end{document}

% =============================================================================
% BEGIN APPENDED COMPACT VERSION V15 MATERIAL
% =============================================================================

\chapter{V15: Companion audit and the cancelling-capacity threshold}
\label{ch:V15}

\section{Mandatory companion audit}
\label{sec:V15-audit}

The active SM note advanced from V12 at the V10 audit to V18.  The NS
note remains V31.  The exact V15 read-only records are:
\begin{center}
\begin{tabular}{p{0.28\linewidth}p{0.12\linewidth}p{0.12\linewidth}
                p{0.38\linewidth}}
\toprule
file & lines & bytes & SHA--256\\
\midrule
\texttt{Renewal\_SM\_Einstein\_Continuum\_Research\_Notes\_V18.tex}
& \(6486\) & \(229765\)
& \texttt{99c6433f94fa4635f03452266195606cdb86c0cf661cc97a7a5bd85c2dbf7e09}\\
\texttt{Renewal\_NS\_Stability\_Research\_Notes\_V31.tex}
& \(23052\) & \(887537\)
& \texttt{f1121b62238ad5491bb7cf03635ea9934942edf573ed8faaa9ee79e1d38a6696}\\
\bottomrule
\end{tabular}
\end{center}

SM V18 proves that connected block averaging and removal of local
Taylor jets give gains only when the observation scale is separated
from the microscopic scale.  A smooth finite-range example retains
order-one diagonal-shell mass.  Its conclusion is that the missing
strong-source estimate must act on the diagonal output shell itself.

NS V31 remains unchanged: normalized annular marks do not control their
first moment, and the exact heat-formation lift reaches the critical
source norm rather than creating a subcritical gain.  Both companion
programmes therefore warn against treating a normalized probability or
capacity as though it paid an unbounded weight.

\begin{assessment}[V15 audit transfer decision]
\label{ass:V15-audit-decision}
For a cancelling YM output, product-state capacity is a normalized
overlap.  Before using it as a Casimir payer, its precise weighted
moment must be computed.  The calculation below shows that the basic
dual-singlet capacity pays only half of one \(SU(2)\) Casimir moment.
The direction is therefore changed from capacity-only payment to a
joint capacity/mark-routing argument.
\end{assessment}

\section{The exact dual-singlet moment threshold}
\label{sec:V15-singlet-threshold}

Let \(V_j\) be the spin-\(j\) representation of \(SU(2)\),
\[
d_j:=2j+1,\qquad a_j:=1+j(j+1),
\]
and let \(P_{\boldsymbol1}\) be the singlet projection in
\(V_j\otimes V_j^*\).  Recall
\[
\kappa_\otimes(P_{\boldsymbol1})=\frac1{d_j},
\qquad
d_j^2=4a_j-3.
\]

\begin{theorem}[Capacity pays exactly one half Casimir moment]
\label{thm:V15-half-moment}
For \(\gamma\ge0\),
\begin{equation}
\boxed{
\kappa_\otimes(a_j^\gamma P_{\boldsymbol1})
=\frac{a_j^\gamma}{\sqrt{4a_j-3}}.}
\label{eq:V15-half-moment}
\end{equation}
Consequently,
\begin{equation}
\sup_{j\in\frac12\mathbb N}
\kappa_\otimes(a_j^\gamma P_{\boldsymbol1})<\infty
\quad\Longleftrightarrow\quad
\gamma\le\frac12.
\label{eq:V15-half-moment-threshold}
\end{equation}
In particular,
\begin{equation}
\kappa_\otimes(a_jP_{\boldsymbol1})
=\frac{a_j}{2j+1}\longrightarrow\infty.
\label{eq:V15-first-moment-fails}
\end{equation}
\end{theorem}

\begin{proof}
Positive homogeneity of capacity and
\eqref{eq:V1-singlet-capacity} give the first equality.  Substitute
\(d_j=\sqrt{4a_j-3}\).  Since \(a_j\to\infty\), the right side is
asymptotic to \(\frac12a_j^{\gamma-1/2}\).  It is uniformly bounded
exactly for \(\gamma\le1/2\).
\end{proof}

\begin{corollary}[Adding harmless core identities does not improve the
threshold]
\label{cor:V15-core-identity}
Let \(C_1,\ldots,C_4\) be arbitrary finite-dimensional core row spaces,
and let the singlet spectator act between rows one and two.  On
\[
\bigotimes_{i=1}^4(C_i\otimes S_i),
\]
where \(S_1=V_j\), \(S_2=V_j^*\), and \(S_3,S_4\) are one-dimensional,
put
\[
Q_j:=I_{C_1\otimes\cdots\otimes C_4}
\otimes P_{\boldsymbol1}.
\]
Then
\begin{equation}
\boxed{\kappa_\otimes(Q_j)=1/d_j.}
\label{eq:V15-core-singlet-capacity}
\end{equation}
Thus even with an otherwise unconstrained four-row core,
\(a_jQ_j\) has unbounded capacity.
\end{corollary}

\begin{proof}
For any product state across the four original rows, trace out the core
inside each row.  The two spectator marginals remain a product density
matrix, so their singlet expectation is at most \(1/d_j\).  Conversely,
take arbitrary pure core states and spectator vectors which attain the
dual-singlet capacity.  This gives equality.
\end{proof}

\begin{firewall}[Capacity is not a first-moment debit]
\label{fire:V15-capacity-not-moment}
Theorem~\ref{thm:V15-half-moment} rules out the estimate
\[
a_j\,\kappa_\otimes(P_{\boldsymbol1})\le C.
\]
A cancelling output has \(q_{\alpha,\boldsymbol1}=1\), so no Wilson
tail factor repairs this loss.  If the MTA marking requires one or two
full inverse input masses at that row, singlet capacity alone is
insufficient.
\end{firewall}

\section{Why the missing mass may instead be a legal tree mark}
\label{sec:V15-mark-rerouting}

The capacity obstruction does not yet disprove the marked-tree atlas.
A large shared input mass appears in an actual support intersection,
so the MTA right side contains markings at that coordinate.  The
following exact regression makes the distinction visible.

Fix one junction coordinate \(e\).  Let the four junction input masses
be fixed positive numbers \(a_1,a_2,a_3,a_4\).  Add one off-junction
coordinate \(f\) shared only by rows one and two, with equal dual input
mass \(A\).  Thus
\begin{equation}
\Xi_1=A+a_1,\qquad
\Xi_2=A+a_2,\qquad
\Xi_3=a_3,\qquad
\Xi_4=a_4.
\label{eq:V15-two-row-shared-masses}
\end{equation}

\begin{proposition}[One rerouted path mark stays macroscopic]
\label{prop:V15-rerouted-path}
Let \(\sigma\) be the path with edges
\(\{12,13,24\}\).  Mark edge \(12\) by \(f\), and mark edges \(13,24\)
by \(e\).  Its exact BFA weight is
\begin{equation}
\boxed{
\mathcal W_{\sigma,\ell}
=
\frac{A^2a_1a_2a_3a_4}
     {(A+a_1)(A+a_2)}.}
\label{eq:V15-rerouted-path}
\end{equation}
If \(A\ge\max\{a_1,a_2\}\), then
\begin{equation}
\mathcal W_{\sigma,\ell}
\ge\frac14a_1a_2a_3a_4.
\label{eq:V15-rerouted-lower}
\end{equation}
By contrast, the all-junction star centered at row one has weight
\begin{equation}
\mathcal W_{\tau_1,e}
=\frac{a_1^3a_2a_3a_4}{(A+a_1)^2}
\longrightarrow0.
\label{eq:V15-all-local-star-decay}
\end{equation}
\end{proposition}

\begin{proof}
Insert \eqref{eq:V15-two-row-shared-masses} into the marked-tree
definition \eqref{eq:V1-marked-tree-weight}.  The three edge factors
are
\[
\frac{A}{\Xi_1}\frac{A}{\Xi_2},\qquad
\frac{a_1}{\Xi_1}\frac{a_3}{\Xi_3},\qquad
\frac{a_2}{\Xi_2}\frac{a_4}{\Xi_4}.
\]
Multiplication by \(\prod_i\Xi_i\) gives
\eqref{eq:V15-rerouted-path}.  If \(A\ge a_1,a_2\), then
\((A+a_1)(A+a_2)\le4A^2\), proving the lower bound.
The star formula follows from
\(\mathcal W_{\tau,e}=\prod_i a_i^{d_i}\Xi_i^{1-d_i}\).
\end{proof}

\begin{assessment}[The capacity no-go selects mark routing]
\label{ass:V15-mark-routing}
In the dual-singlet regression the all-junction marked star is too
small, and capacity supplies only \(A^{-1/2}\).  Nevertheless the legal
path which marks the actual shared coordinate has order-one weight.
Thus a proof which insists on paying every local tree fibre by its
all-\(e\) marking is stronger than the MTA requires and can fail for
the wrong reason.  The correct next question is whether every
strongly cancelling shared bank admits a bounded-incidence rerouting
to the complete sum over trees and legal coordinate markings.
\end{assessment}

\section{The exact residual after V15}
\label{sec:V15-residual}

For a general shared coordinate \(f\), cancellation means that some
active input mass is large compared with \(A_{U_f}\).  V13 provides a
raw overlap gain for its physical projection, while V15 shows that the
gain has only dimension scale and need not pay a Casimir moment.  The
MTA right side, however, contains the overlap resources
\begin{equation}
\sum_{f\in X_i\cap X_j}
\frac{A_{R_{i,f}}}{\Xi_i}
\frac{A_{R_{j,f}}}{\Xi_j}
\label{eq:V15-legal-pair-marks}
\end{equation}
on every tree edge \(ij\).  A complete proof must combine:
\begin{enumerate}[label=\textup{(\roman*)},leftmargin=4em]
\item capacity of the cancelling projection, before norming;
\item an edge \(ij\) supported by that coordinate;
\item the exact sum over all legal markings; and
\item a bounded-use assignment of local contraction-tree fibres to
      those marked trees.
\end{enumerate}
No such general routing theorem is asserted in V15.

\begin{firewall}[The example is not the routing theorem]
\label{fire:V15-example-only}
\Cref{prop:V15-rerouted-path} treats one pair-shared coordinate and one
explicit mass pattern.  It proves that legal marks can remove the
apparent counterexample; it does not prove that an arbitrary mixture
of pair-, triple-, and four-row shared banks can be routed without
support-size or multiplicity loss.
\end{firewall}

\section{V15 ledger and V16 acquisition}
\label{sec:V15-ledger}

\begin{theorem}[Integrated compact-series V15 statement]
\label{thm:V15-integrated}
Preserving compact V1--V14, V15 proves:
\begin{enumerate}[label=\textup{(V15.\arabic*)},leftmargin=6.5em]
\item the mandatory hashed audit of active SM V18 and NS V31;
\item the exact dual-singlet weighted-capacity formula
      \eqref{eq:V15-half-moment};
\item the sharp threshold that capacity alone pays at most one half
      \(SU(2)\) Casimir moment;
\item persistence of that threshold after arbitrary harmless core
      identities are inserted;
\item an exact legal path marking which stays macroscopic on the
      single dual-shared regression; and
\item replacement of the capacity-only plan by a joint
      capacity/mark-routing acquisition.
\end{enumerate}
\end{theorem}

\begin{proof}
These are
\Cref{ass:V15-audit-decision,thm:V15-half-moment,
cor:V15-core-identity,prop:V15-rerouted-path,
ass:V15-mark-routing}.
\end{proof}

\begin{openproblem}[V16 mandate: pair-shared mark routing]
\label{op:V16-mandate}
\begin{enumerate}
\item Start with shared banks of exact incidence two; postpone
      triple- and four-row shared coordinates.
\item Split each pair-shared coordinate into a retentive output part
      closed by V14 and a cancelling output part.
\item For the cancelling part on pair \(ij\), retain the physical
      projector and use the V13 row-marginal capacity bound.
\item Route one edge \(ij\) of a comparison tree to that coordinate,
      as in \Cref{prop:V15-rerouted-path}, before summing local tree
      fibres.
\item Prove bounded incidence of the map from local fibres and
      cancelling coordinates to comparison trees and legal markings.
\item If a diffuse pair bank prevents a point selection, formulate the
      exact bilinear mass sum which the capacity expectation must
      control; do not use a support cardinality bound.
\end{enumerate}
\end{openproblem}

\begin{firewall}[Claim boundary after compact V15]
\label{fire:V15-current-boundary}
V15 proves a sharp capacity no-go and a single-coordinate rerouting
regression.  It does not close general cancelling shared banks,
local-dominant non-Cartan blocks, the full embedded \([4]\) sector,
multiple junctions, normalized collision aggregation, \([3,3]\),
global quartic or all-order MTA, continuum Yang--Mills existence, or a
positive mass gap.
\end{firewall}

\section{Append-only record for compact V15}
\label{sec:V15-record}

V15 was appended to the complete compact V14 file.  The exact V14
parent had \(5602\) newline-delimited source records, \(189055\) bytes,
and SHA--256 digest
\[
\texttt{e0856774d8a230fd6853fd585971ce09cc52969b30b29e704798186a42dc8313}.
\]
No compact V14 mathematical content was changed.  The sole final
\verb|\end{document}| is restored below.

% =============================================================================
% END APPENDED COMPACT VERSION V15 MATERIAL
% =============================================================================

% The compact V15 document terminator is intentionally disabled here:
% \end{document}

% =============================================================================
% BEGIN APPENDED COMPACT VERSION V16 MATERIAL
% =============================================================================

\chapter{V16: Diffuse pair-singlet banks are paid by legal marks}
\label{ch:V16}

\section{An exact bank of dual-singlet outputs}
\label{sec:V16-bank}

Fix two distinct rows \(i,j\), and let \(F\) be a nonempty finite set
of off-junction coordinates active only in those two rows.  At
\(f\in F\), let row \(i\) carry a nontrivial irreducible \(R_f\), let
row \(j\) carry \(R_f^*\), and resolve the shared output to the trivial
representation.  Write
\[
d_f:=\dim R_f,\qquad
A_f:=1+C_2(R_f)=1+C_2(R_f^*),
\]
and let \(P_{\boldsymbol1,f}\) be the normalized singlet projection in
\(R_f\otimes R_f^*\).  Define
\begin{align}
P_F&:=\bigotimes_{f\in F}P_{\boldsymbol1,f},&
D_F&:=\prod_{f\in F}d_f,\label{eq:V16-bank-projector}\\
h_F&:=\sum_{f\in F}A_f,&
H_F&:=\sum_{f\in F}A_f^2.
\label{eq:V16-bank-masses}
\end{align}

\begin{lemma}[The whole bank is one maximally entangled projector]
\label{lem:V16-bank-maximally-entangled}
Across the two complete spectator row spaces
\[
\mathcal S_i:=\bigotimes_{f\in F}R_f,
\qquad
\mathcal S_j:=\bigotimes_{f\in F}R_f^*,
\]
\(P_F\) is the rank-one projection onto a maximally entangled vector of
Schmidt rank \(D_F\).  Consequently,
\begin{equation}
\boxed{\kappa_\otimes(P_F)=D_F^{-1}.}
\label{eq:V16-bank-capacity}
\end{equation}
The same value holds after identity factors are added to either row or
to the other two rows.
\end{lemma}

\begin{proof}
Choose orthonormal bases in each \(R_f\).  The tensor product of the
normalized singlet vectors is, after lexicographically grouping their
indices,
\[
D_F^{-1/2}\sum_{\boldsymbol a}
e_{\boldsymbol a}\otimes e_{\boldsymbol a}^*.
\]
This is the normalized maximally entangled vector between
\(\mathcal S_i\) and \(\mathcal S_j\).  Its maximum squared overlap
with a product vector is \(D_F^{-1}\), by the same Cauchy--Schwarz
calculation as in
\Cref{prop:V1-singlet-capacity}.  Mixed product states do not increase
the supremum.  Identity factors may be traced out row by row and can
also be filled by arbitrary pure states, so they do not change the
value.
\end{proof}

\section{Capacity removes the diffuse-bank counting loss}
\label{sec:V16-diffuse}

\begin{theorem}[Exact capacity--bilinear-mark inequality]
\label{thm:V16-capacity-mark}
For every dual-singlet bank above,
\begin{equation}
\boxed{
h_F^2\,\kappa_\otimes(P_F)\le H_F.}
\label{eq:V16-capacity-mark}
\end{equation}
The constant is independent of \(\lvert F\rvert\), the
representations, and their masses.
\end{theorem}

\begin{proof}
Put \(N=\lvert F\rvert\).  Cauchy--Schwarz gives
\[
h_F^2
=\left(\sum_fA_f\right)^2
\le N\sum_fA_f^2
=NH_F.
\]
For a compact connected semisimple group, every nontrivial irreducible
has dimension at least two.  Hence
\[
D_F\ge2^N,\qquad
\frac{N}{D_F}\le\frac{N}{2^N}\le1.
\]
Multiply the Cauchy--Schwarz estimate by
\(\kappa_\otimes(P_F)=D_F^{-1}\).
\end{proof}

\begin{remark}[The single-coordinate half-moment obstruction is
absorbed by the legal numerator]
\label{rem:V16-single-coordinate}
When \(F=\{f\}\), the theorem reads
\[
A_f^2/d_f\le A_f^2.
\]
Capacity still does not pay \(A_f\) by itself, as V15 proved.  Instead,
the bilinear legal edge mark supplies \(A_f^2\), and capacity controls
the normalization of its cancelling projection.  For many
coordinates, \(D_F^{-1}\) supplies exactly the factor which removes the
Cauchy--Schwarz cardinality \(N\).
\end{remark}

\section{The bilinear mass is exactly a sum of legal markings}
\label{sec:V16-legal-mark}

Let \(e\) be the four-row junction and write
\[
a_r:=1+C_2(R_{r,e}),\qquad 1\le r\le4.
\]
Choose the two remaining rows \(k,l\) so that
\(\{i,j,k,l\}=\{1,2,3,4\}\).  Let \(\sigma_{ij}\) be the path with
edges
\[
\{ij,ik,jl\}.
\]
On its edge \(ij\), allow the marking to range over \(f\in F\); mark
the other two edges by \(e\).

\begin{lemma}[Exact sum of the routed path weights]
\label{lem:V16-routed-path-sum}
The indicated subfamily of legal marked-tree weights has exact sum
\begin{equation}
\boxed{
\sum_{f\in F}\mathcal W_{\sigma_{ij},\ell_f}
=
\frac{H_F\,a_ia_ja_ka_l}{\Xi_i\Xi_j}.}
\label{eq:V16-routed-path-sum}
\end{equation}
This identity remains true when the four rows contain additional
spectator masses; those masses enter only through \(\Xi_i,\Xi_j\).
\end{lemma}

\begin{proof}
For fixed \(f\), the three edge factors in
\eqref{eq:V1-marked-tree-weight} are
\[
\frac{A_f}{\Xi_i}\frac{A_f}{\Xi_j},
\qquad
\frac{a_i}{\Xi_i}\frac{a_k}{\Xi_k},
\qquad
\frac{a_j}{\Xi_j}\frac{a_l}{\Xi_l}.
\]
Multiplication by \(\prod_r\Xi_r\) leaves
\[
\frac{A_f^2a_ia_ja_ka_l}{\Xi_i\Xi_j}.
\]
Sum over \(f\).
\end{proof}

\begin{corollary}[Bank-dominant pair routing has bounded incidence]
\label{cor:V16-bank-dominant-routing}
Assume that the same bank dominates both participating row masses:
\begin{equation}
\Xi_i\le C_0h_F,\qquad
\Xi_j\le C_0h_F.
\label{eq:V16-bank-dominance}
\end{equation}
Then
\begin{equation}
\boxed{
\kappa_\otimes(P_F)\prod_{r=1}^4a_r
\le
C_0^2
\sum_{f\in F}\mathcal W_{\sigma_{ij},\ell_f}.}
\label{eq:V16-bank-routing}
\end{equation}
\end{corollary}

\begin{proof}
By \Cref{thm:V16-capacity-mark},
\[
\kappa_\otimes(P_F)
\le H_F/h_F^2.
\]
The two inequalities in \eqref{eq:V16-bank-dominance} give
\[
\frac{H_F}{h_F^2}
\le C_0^2\frac{H_F}{\Xi_i\Xi_j}.
\]
Multiply by \(\prod_ra_r\) and use
\eqref{eq:V16-routed-path-sum}.
\end{proof}

\section{Insertion into a positive local core}
\label{sec:V16-core-insertion}

\begin{theorem}[One-sided capacity insertion for a dual-singlet bank]
\label{thm:V16-core-insertion}
Let \(Q_e\ge0\) be a positive four-row carrier on the junction core,
let \(P_F\) act on the pair-shared spectator bank, and let all remaining
spectators be positive contractions.  If
\begin{equation}
\|Q_e\|\le C_e\prod_{r=1}^4a_r
\label{eq:V16-core-op-bound}
\end{equation}
and \eqref{eq:V16-bank-dominance} holds, then the complete carrier
\(Q_{\rm full}\) satisfies
\begin{equation}
\boxed{
\kappa_\otimes(Q_{\rm full})
\le
C_eC_0^2
\sum_{f\in F}\mathcal W_{\sigma_{ij},\ell_f}.}
\label{eq:V16-core-insertion}
\end{equation}
\end{theorem}

\begin{proof}
Discard the remaining positive contractions by
\Cref{cor:V2-shared-contractions}.  The one-sided tensor law
\eqref{eq:V2-one-sided-tensor} gives
\[
\kappa_\otimes(Q_e\otimes P_F)
\le\|Q_e\|\kappa_\otimes(P_F).
\]
Apply \eqref{eq:V16-core-op-bound}, followed by
\Cref{cor:V16-bank-dominant-routing}.
\end{proof}

\begin{firewall}[Why entanglement swapping is not being assumed away]
\label{fire:V16-swapping}
The proof uses an operator norm on the junction side and capacity only
on the pair-singlet side, exactly as permitted by the V2 one-sided
tensor theorem.  It does not condition on the singlet output and then
claim that the remaining core state is product.  Such a claim would be
false by the V2 entanglement-swapping regression.
\end{firewall}

\section{What is and is not closed}
\label{sec:V16-status}

\begin{assessment}[A complete diffuse-bank subproblem is closed]
\label{ass:V16-closed-subproblem}
Equations \eqref{eq:V16-capacity-mark} and
\eqref{eq:V16-bank-routing} solve the support-uniform normalization and
mark-routing problem for a bank of exact pair-shared dual singlets
which dominates both participating rows.  Every coordinate is used
once in the legal edge-mark sum; the factor \(D_F^{-1}\) removes the
otherwise unavoidable bank cardinality.
\end{assessment}

\begin{firewall}[The local payer and general cancelling channels remain]
\label{fire:V16-remaining}
\Cref{thm:V16-core-insertion} requires the local positive carrier bound
\eqref{eq:V16-core-op-bound}.  The hot non-Cartan junction payer is not
known to satisfy that bound at the needed scale.  V16 also does not
treat nontrivial low-output cancellation channels, a pair bank which
dominates only one of its two rows, or overlapping pair banks sharing a
row.  Those cases require further allocation.
\end{firewall}

\section{V16 ledger and V17 acquisition}
\label{sec:V16-ledger}

\begin{theorem}[Integrated compact-series V16 statement]
\label{thm:V16-integrated}
Preserving compact V1--V15, V16 proves:
\begin{enumerate}[label=\textup{(V16.\arabic*)},leftmargin=6.5em]
\item exact capacity \(D_F^{-1}\) of an arbitrary diffuse
      pair-singlet bank;
\item the support-uniform joint inequality
      \(h_F^2\kappa_\otimes(P_F)\le H_F\);
\item identification of \(H_F\) as the exact sum of legal pair-edge
      markings;
\item bounded-incidence routing when the bank dominates both rows; and
\item rigorous one-sided insertion into any local positive core
      satisfying \eqref{eq:V16-core-op-bound}.
\end{enumerate}
\end{theorem}

\begin{proof}
These are
\Cref{lem:V16-bank-maximally-entangled,
thm:V16-capacity-mark,lem:V16-routed-path-sum,
cor:V16-bank-dominant-routing,thm:V16-core-insertion}.
\end{proof}

\begin{openproblem}[V17 mandate: asymmetric and overlapping pair banks]
\label{op:V17-mandate}
\begin{enumerate}
\item Keep exact dual-singlet outputs, but allow several pair labels
      \(ij\) and allow their banks to share rows.
\item Split each row's cancelling mass among its three possible
      partners with a deterministic bounded-overlap rule.
\item When a selected bank dominates only row \(i\), identify the
      second row or private/output stock which pays the comparison-path
      denominator \(\Xi_j\).
\item Sum the legal routed paths before using
      \Cref{thm:V16-capacity-mark}; do not select one coordinate by a
      support-cardinality pigeonhole.
\item Determine whether one-sided capacity can be applied bank by bank
      without entanglement swapping across a row shared by two banks.
\item If not, isolate the exact multi-bank marginal capacity needed.
\end{enumerate}
\end{openproblem}

\begin{firewall}[Claim boundary after compact V16]
\label{fire:V16-current-boundary}
V16 closes the diffuse normalization and routing subproblem for one
bank-dominant pair-singlet bank, conditionally on a local core operator
bound.  It does not close the general cancelling spectator family,
local-dominant non-Cartan blocks, the full embedded \([4]\) sector,
multiple junctions, normalized collision aggregation, \([3,3]\),
global quartic or all-order MTA, continuum Yang--Mills existence, or a
positive mass gap.
\end{firewall}

\section{Append-only record for compact V16}
\label{sec:V16-record}

V16 was appended to the complete compact V15 file.  The exact V15
parent had \(5929\) newline-delimited source records, \(201097\) bytes,
and SHA--256 digest
\[
\texttt{79940411657647982004aaeae2aa9f0badae6a9c7b5ad78425c892fe86f3f2c4}.
\]
No compact V15 mathematical content was changed.  The sole final
\verb|\end{document}| is restored below.

% =============================================================================
% END APPENDED COMPACT VERSION V16 MATERIAL
% =============================================================================

% The compact V16 document terminator is intentionally disabled here:
% \end{document}

% =============================================================================
% BEGIN APPENDED COMPACT VERSION V17 MATERIAL
% =============================================================================

\chapter{V17: Exact capacity of an overlapping singlet network}
\label{ch:V17}

\section{A graph of pair-shared singlets}
\label{sec:V17-graph}

Let \(\Gamma\) be a finite multigraph on the four row vertices.  Each
edge \(f\in E(\Gamma)\), joining distinct rows \(i(f),j(f)\), carries a
nontrivial irreducible \(R_f\) of dimension \(d_f\) at one endpoint,
its dual at the other endpoint, and the resolved trivial-output
projection \(P_{\boldsymbol1,f}\).  Put
\begin{equation}
P_\Gamma:=\bigotimes_{f\in E(\Gamma)}P_{\boldsymbol1,f},
\qquad
D_\Gamma:=\prod_{f\in E(\Gamma)}d_f.
\label{eq:V17-network-projector}
\end{equation}
The tensor factors are grouped by the original four rows.  A row may
therefore contain, and internally entangle, all half-edges incident to
that vertex.

\begin{lemma}[Finite graph Finner inequality]
\label{lem:V17-Finner}
Let \(X_f\) be independent finite random variables indexed by the
edges of a finite graph.  If
\(g_i=g_i((X_f)_{f\ni i})\), then
\begin{equation}
\boxed{
\mathbb E\prod_i|g_i|
\le\prod_i\left(\mathbb E|g_i|^2\right)^{1/2}.}
\label{eq:V17-Finner}
\end{equation}
\end{lemma}

\begin{proof}
This is the finite Finner inequality with exponent two at every graph
vertex.  For completeness, expand the expectation as an iterated sum
over the edge variables.  For the variable \(X_f\), exactly the two
endpoint functions depend on it.  Apply Cauchy--Schwarz to those two
factors in the \(X_f\)-sum.  Continue through the edge variables,
carrying the resulting squares into the remaining iterated sums.
Every endpoint function is squared exactly once for each occurrence of
one of its variables and then integrated over its full product
measure; successive square roots cancel the repeated exponents.  The
result is precisely the product of the four \(L^2\) norms displayed in
\eqref{eq:V17-Finner}.  Isolated vertices contribute constants and do
not change the argument.
\end{proof}

\begin{theorem}[Overlapping singlet capacities multiply exactly]
\label{thm:V17-network-capacity}
Across the original row grouping,
\begin{equation}
\boxed{\kappa_\otimes(P_\Gamma)=D_\Gamma^{-1}.}
\label{eq:V17-network-capacity}
\end{equation}
Thus sharing a row among several pair banks creates no
entanglement-swapping loss for the complete rank-one singlet network.
\end{theorem}

\begin{proof}
Choose an orthonormal basis on every edge representation.  A normalized
pure state in row \(i\) has coefficients
\(t_i((x_f)_{f\ni i})\) of counting-measure \(\ell^2\)-norm one.
The overlap with the tensor product of normalized singlets is
\begin{equation}
\langle \otimes_i t_i,\Omega_\Gamma\rangle
=
D_\Gamma^{-1/2}
\sum_{(x_f)}
\prod_i\overline{t_i((x_f)_{f\ni i})}.
\label{eq:V17-network-overlap}
\end{equation}
Give every \(x_f\) the uniform probability measure.  Its row-\(i\)
\(L^2\) norm is
\[
\left(\mathbb E|t_i|^2\right)^{1/2}
=
\left(\prod_{f\ni i}d_f\right)^{-1/2}.
\]
By \Cref{lem:V17-Finner}, the absolute value of the unnormalized sum
in \eqref{eq:V17-network-overlap} is at most one: multiplication by
\(\prod_fd_f\) cancels the product of row normalization factors because
every edge has two endpoints.  Hence the squared overlap is at most
\(D_\Gamma^{-1}\).

For the reverse inequality, fix one basis label on every edge and take
at each row the tensor product of its incident fixed basis vectors.
Exactly one summand in \eqref{eq:V17-network-overlap} survives, with
value \(D_\Gamma^{-1/2}\).  Convexity extends the pure-state maximum to
mixed product states.
\end{proof}

\begin{firewall}[Relation to the V2 swapping regression]
\label{fire:V17-swapping}
The theorem concerns the complete rank-one tensor network
\(P_\Gamma\) itself.  It does not condition on one singlet and assert
that an unrelated junction state remains product.  When a separate
core is present, V17 continues to use the one-sided tensor law with an
operator norm on that core.  The V2 entanglement-swapping
counterexample remains fully operative for an attempted
capacity-times-capacity argument.
\end{firewall}

\section{All pair banks and their legal edge masses}
\label{sec:V17-edge-banks}

For every unordered pair \(ij\), let \(F_{ij}\) be the corresponding
edge bank and set
\begin{align}
h_{ij}&:=\sum_{f\in F_{ij}}A_f,&
H_{ij}&:=\sum_{f\in F_{ij}}A_f^2,&
D_{ij}&:=\prod_{f\in F_{ij}}d_f.
\label{eq:V17-pair-bank-data}
\end{align}
Empty banks have \(h_{ij}=H_{ij}=0\) and \(D_{ij}=1\).  The total
pair-singlet mass in row \(i\) is
\begin{equation}
h_i^{\rm ps}:=\sum_{j\ne i}h_{ij}.
\label{eq:V17-row-pair-mass}
\end{equation}

\begin{corollary}[Every individual edge bank retains the V16
capacity--mark inequality]
\label{cor:V17-edge-bank}
For every nonempty \(F_{ij}\),
\begin{equation}
\boxed{
\frac{h_{ij}^2}{D_{ij}}\le H_{ij},
\qquad
\kappa_\otimes(P_\Gamma)\le D_{ij}^{-1}.}
\label{eq:V17-edge-bank}
\end{equation}
\end{corollary}

\begin{proof}
The first inequality is
\Cref{thm:V16-capacity-mark}.  The second follows from
\[
D_\Gamma=D_{ij}\prod_{f\notin F_{ij}}d_f\ge D_{ij}
\]
and \eqref{eq:V17-network-capacity}.
\end{proof}

\section{A deterministic dominant-bank selection}
\label{sec:V17-selection}

\begin{lemma}[One pair bank dominates two full row masses]
\label{lem:V17-dominant-pair}
Fix \(D\ge16\).  Suppose:
\begin{enumerate}[label=\textup{(\roman*)},leftmargin=4em]
\item \(a_r<\Xi_r/4\) for every row \(r\);
\item all off-junction shared mass is pair-singlet mass, so
      \(h_r=h_r^{\rm ps}\);
\item for some quartic tree,
      \(B_\tau>DY^2\).
\end{enumerate}
Choose a row \(i\) with maximal \(\Xi_i\).  Then there is a partner
\(j\ne i\) such that
\begin{equation}
\boxed{
\Xi_i<6h_{ij},\qquad
\Xi_j<6h_{ij}.}
\label{eq:V17-two-row-dominance}
\end{equation}
\end{lemma}

\begin{proof}
The failed branching condition gives an internal row \(r\) with
\(\Xi_r>\sqrt D\,Y\).  Maximality therefore gives
\(\Xi_i>\sqrt D\,Y\ge4Y\).  Hence
\[
p_i\le Y<\Xi_i/4,\qquad a_i<\Xi_i/4,
\]
so
\[
h_i^{\rm ps}=\Xi_i-a_i-p_i>\Xi_i/2.
\]
There are only three possible partners.  Choose \(j\) with
\[
h_{ij}\ge\frac13h_i^{\rm ps}>\frac16\Xi_i.
\]
This proves the first inequality in
\eqref{eq:V17-two-row-dominance}.  Since \(i\) was chosen globally
maximal, \(\Xi_j\le\Xi_i<6h_{ij}\), proving the second.
\end{proof}

\section{Bounded routing for the complete overlapping network}
\label{sec:V17-routing}

For the selected pair \(ij\), let \(\sigma_{ij}\) be the comparison
path of V16 and let \(\ell_f\) mark its \(ij\)-edge by
\(f\in F_{ij}\), with its other two edges marked by the junction.

\begin{theorem}[Overlapping pair-singlet network routing]
\label{thm:V17-network-routing}
Under the hypotheses of
\Cref{lem:V17-dominant-pair},
\begin{equation}
\boxed{
\kappa_\otimes(P_\Gamma)\prod_{r=1}^4a_r
\le
36\sum_{f\in F_{ij}}
\mathcal W_{\sigma_{ij},\ell_f}.}
\label{eq:V17-network-routing}
\end{equation}
The selected bank and comparison path depend only on the fixed input
block, not on a choice of product coefficient.
\end{theorem}

\begin{proof}
By \Cref{cor:V17-edge-bank},
\[
\kappa_\otimes(P_\Gamma)
\le D_{ij}^{-1}
\le\frac{H_{ij}}{h_{ij}^2}.
\]
The two-row dominance
\eqref{eq:V17-two-row-dominance} gives
\[
\frac{H_{ij}}{h_{ij}^2}
\le36\frac{H_{ij}}{\Xi_i\Xi_j}.
\]
Multiply by \(\prod_ra_r\) and use the exact routed path identity
\eqref{eq:V16-routed-path-sum} for \(F_{ij}\).
\end{proof}

\begin{corollary}[One-sided insertion with overlapping pair banks]
\label{cor:V17-core-insertion}
Let \(Q_e\ge0\) be a local junction carrier satisfying
\[
\|Q_e\|\le C_e\prod_ra_r,
\]
and let all nonsinglet remaining spectators be positive contractions.
Under the hypotheses of
\Cref{lem:V17-dominant-pair}, the full carrier satisfies
\begin{equation}
\boxed{
\kappa_\otimes(Q_{\rm full})
\le36C_e
\sum_{f\in F_{ij}}\mathcal W_{\sigma_{ij},\ell_f}.}
\label{eq:V17-core-insertion}
\end{equation}
\end{corollary}

\begin{proof}
The V2 one-sided tensor law gives
\[
\kappa_\otimes(Q_{\rm full})
\le\|Q_e\|\kappa_\otimes(P_\Gamma).
\]
Use \Cref{thm:V17-network-routing}.
\end{proof}

\begin{assessment}[The multi-bank marginal problem is closed for exact
singlets]
\label{ass:V17-multibank-closed}
V17 proves that arbitrary overlap of exact pair-singlet banks across
the four rows costs no additional capacity factor.  A globally maximal
all-nonlocal row selects one of its three partner banks, and that bank
dominates both endpoint row masses at a universal constant.  Its exact
bilinear mass is already the legal marking sum of one comparison path.
Thus the multi-bank normalization, diffuse incidence, and routing
parts of the exact pair-singlet problem are closed.
\end{assessment}

\begin{firewall}[What the theorem does not absorb]
\label{fire:V17-remaining}
The theorem is conditional on the local core operator bound used in
\Cref{cor:V17-core-insertion}.  It also relies on rank-one trivial
outputs at every cancelling pair coordinate.  A nontrivial low-output
channel has a higher-rank invariant projector and is not an EPR edge;
Finner's scalar contraction does not directly give its network
capacity.  Triple- and four-row cancelling outputs are likewise still
open.
\end{firewall}

\section{V17 ledger and V18 acquisition}
\label{sec:V17-ledger}

\begin{theorem}[Integrated compact-series V17 statement]
\label{thm:V17-integrated}
Preserving compact V1--V16, V17 proves:
\begin{enumerate}[label=\textup{(V17.\arabic*)},leftmargin=6.5em]
\item the graph Finner inequality
      \eqref{eq:V17-Finner};
\item exact multiplicativity
      \(\kappa_\otimes(P_\Gamma)=D_\Gamma^{-1}\) for every overlapping
      pair-singlet network;
\item a deterministic selection of one partner bank which dominates
      both endpoint row masses;
\item bounded routing of the complete network capacity to one legal
      comparison-path marking sum; and
\item conditional one-sided insertion into the local junction carrier.
\end{enumerate}
\end{theorem}

\begin{proof}
These are
\Cref{lem:V17-Finner,thm:V17-network-capacity,
lem:V17-dominant-pair,thm:V17-network-routing,
cor:V17-core-insertion}.
\end{proof}

\begin{openproblem}[V18 mandate: beyond rank-one cancelling outputs]
\label{op:V18-mandate}
\begin{enumerate}
\item At one pair-shared coordinate, replace the singlet by an arbitrary
      resolved nontrivial low-output copy \(P_{U,a}\).
\item Use the exact row marginal
      \(\Tr_{\widehat i}P_{U,a}=(d_U/d_{R_i})I\).
\item Determine whether a graph analogue of
      \eqref{eq:V17-network-capacity} holds with edge factor
      \(\min\{1,d_U/d_{R_i},d_U/d_{R_j}\}\), or construct a counterexample.
\item Keep the bilinear legal edge mass
      \(A_{R_i}A_{R_j}\) beside that capacity factor.
\item If the graph product fails, identify a fractional-cover capacity
      bound strong enough for the four-row tree atlas.
\item Continue to keep the hot local non-Cartan payer as a separate
      upstream gate.
\end{enumerate}
\end{openproblem}

\begin{firewall}[Claim boundary after compact V17]
\label{fire:V17-current-boundary}
V17 closes overlapping exact pair-singlet normalization and routing,
conditionally on a local core operator bound.  It does not close
higher-rank cancelling outputs, triple- or four-row cancelling
spectators, the hot local-dominant non-Cartan carrier, the full embedded
\([4]\) sector, multiple junctions, normalized collision aggregation,
\([3,3]\), global quartic or all-order MTA, continuum Yang--Mills
existence, or a positive mass gap.
\end{firewall}

\section{Append-only record for compact V17}
\label{sec:V17-record}

V17 was appended to the complete compact V16 file.  The exact V16
parent had \(6265\) newline-delimited source records, \(212220\) bytes,
and SHA--256 digest
\[
\texttt{397c40f9f430f8b3cd2678f67ae93e8f13ab938049b438832a6bb534877ece50}.
\]
No compact V16 mathematical content was changed.  The sole final
\verb|\end{document}| is restored below.

% =============================================================================
% END APPENDED COMPACT VERSION V17 MATERIAL
% =============================================================================

% The compact V17 document terminator is intentionally disabled here:
% \end{document}

% =============================================================================
% BEGIN APPENDED COMPACT VERSION V18 MATERIAL
% =============================================================================

\chapter{V18: Higher-rank pair cancellation and cutwise capacity}
\label{ch:V18}

\section{Naive edge-capacity multiplication is false}
\label{sec:V18-counterexample}

Before constructing a replacement for V17, one must test whether
arbitrary higher-rank edge capacities multiply after the edge
coordinates are regrouped into rows.  They do not.

\begin{proposition}[Two antisymmetric channels are supermultiplicative]
\label{prop:V18-antisymmetric-counterexample}
Let \(d\ge3\), let \(F\) be the flip on
\(\mathbb C^d\otimes\mathbb C^d\), and put
\[
P_\wedge:=\frac{I-F}{2}.
\]
For one coordinate,
\begin{equation}
\kappa_\otimes(P_\wedge)=\frac12.
\label{eq:V18-one-antisymmetric}
\end{equation}
For two parallel coordinates, regrouped as
\[
(\mathbb C^d_{A_1}\otimes\mathbb C^d_{A_2})
\otimes
(\mathbb C^d_{B_1}\otimes\mathbb C^d_{B_2}),
\]
one has
\begin{equation}
\boxed{
\kappa_\otimes(P_\wedge^{(1)}\otimes P_\wedge^{(2)})
\ge\frac12\left(1-\frac1d\right)>\frac14.}
\label{eq:V18-antisymmetric-supermultiplicative}
\end{equation}
These are invariant physical projections for the fundamental
representation of \(SU(d)\).
\end{proposition}

\begin{proof}
For one product density matrix,
\[
\Tr[P_\wedge(\rho\otimes\sigma)]
=\frac12\bigl(1-\Tr(\rho\sigma)\bigr)\le\frac12.
\]
Orthogonal pure vectors attain equality.

For two coordinates, take in each complete row the normalized
maximally entangled vector
\[
\Omega_A=d^{-1/2}\sum_a e_a^{A_1}\otimes e_a^{A_2},
\qquad
\Omega_B=d^{-1/2}\sum_a e_a^{B_1}\otimes e_a^{B_2}.
\]
This is a legal product vector across the two regrouped rows.  Expanding
the two antisymmetric projectors gives
\[
\frac14\left(
1-\langle F_1\rangle-\langle F_2\rangle
+\langle F_1F_2\rangle
\right).
\]
The reduced states at each single coordinate are \(I/d\), so
\(\langle F_1\rangle=\langle F_2\rangle=1/d\).
The double flip exchanges the two identical complete-row vectors, so
\(\langle F_1F_2\rangle=1\).  This proves
\eqref{eq:V18-antisymmetric-supermultiplicative}.
\end{proof}

\begin{firewall}[Scope of the counterexample]
\label{fire:V18-counterexample-scope}
The antisymmetric output has rank \(d(d-1)/2\), larger than either
input dimension when \(d\ge3\).  It disproves unrestricted
edge-capacity multiplication, but it is not a strongly
dimension-cancelling channel.  The replacement theorem below uses the
small one-row marginal of precisely that cancelling regime.
\end{firewall}

\section{A cutwise product theorem for one pair bank}
\label{sec:V18-cutwise}

Fix two rows \(i,j\) and a finite coordinate bank \(F\).  At coordinate
\(f\), let \(R_{i,f},R_{j,f}\) be irreducible inputs and let
\(P_f=P_{U_f,a_f}\) project onto one resolved irreducible output copy
of dimension \(d_{U_f}\).  Put
\begin{equation}
P_F:=\bigotimes_{f\in F}P_f,
\qquad
\alpha_i(F):=\prod_{f\in F}\frac{d_{U_f}}{d_{R_{i,f}}},
\qquad
\alpha_j(F):=\prod_{f\in F}\frac{d_{U_f}}{d_{R_{j,f}}}.
\label{eq:V18-cut-data}
\end{equation}

\begin{theorem}[Cutwise capacity products survive row entanglement]
\label{thm:V18-cutwise-capacity}
Across the two complete row spaces,
\begin{equation}
\boxed{
\kappa_\otimes(P_F)
\le\min\{1,\alpha_i(F),\alpha_j(F)\}.}
\label{eq:V18-cutwise-capacity}
\end{equation}
The same upper bound holds for a larger four-row spectator carrier
which contains \(P_F\) and whose remaining coordinate factors are
positive contractions.
\end{theorem}

\begin{proof}
The scalar marginal formula
\eqref{eq:V13-scalar-marginal} gives, coordinate by coordinate,
\[
\Tr_{R_{j,f}}P_f
=\frac{d_{U_f}}{d_{R_{i,f}}}I_{R_{i,f}}.
\]
Tensoring these identities yields
\[
\Tr_{\mathcal H_j}P_F
=\alpha_i(F)I_{\mathcal H_i}.
\]
For arbitrary, possibly coordinate-entangled, row density matrices
\(\rho_i,\rho_j\),
\[
\Tr[P_F(\rho_i\otimes\rho_j)]
\le\Tr[P_F(\rho_i\otimes I)]
=\alpha_i(F).
\]
Interchange \(i,j\) and include the operator-norm bound one.  In a
larger carrier, replace every other positive contraction by the
identity.  Tracing the unused factors inside each original row leaves
a product density matrix on the two selected bank spaces, so the same
bound applies.
\end{proof}

\begin{assessment}[What replaces graphwise multiplication]
\label{ass:V18-cutwise-replacement}
Arbitrary edge capacities cannot be multiplied coordinate by
coordinate, as
\Cref{prop:V18-antisymmetric-counterexample} shows.  However, all
coordinates of one selected pair bank cross the same row cut.  Their
exact scalar marginals tensor, so the dimension ratios on either
endpoint multiply even when that row is internally entangled.  A
dominant-row argument only needs one such selected bank.
\end{assessment}

\section{Strong \(SU(2)\) cancellation forces a half-dimension ratio}
\label{sec:V18-SU2}

For an \(SU(2)\) spin \(r\), write
\[
d_r:=2r+1,\qquad A_r:=1+r(r+1),
\qquad d_r^2=4A_r-3.
\]

\begin{lemma}[Mass cancellation implies dimension cancellation and a
comparable partner]
\label{lem:V18-SU2-cancellation}
Suppose spin \(J\) occurs in \(V_p\otimes V_q\), with \(p,q>0\), and
\begin{equation}
A_J<\frac17A_p.
\label{eq:V18-strong-mass-cancellation}
\end{equation}
Then
\begin{equation}
\boxed{
\frac{d_J}{d_p}<\frac12,
\qquad
A_q>\frac14A_p.}
\label{eq:V18-dimension-and-partner}
\end{equation}
\end{lemma}

\begin{proof}
The dimension--Casimir identity gives
\[
d_J^2=4A_J-3
<\frac47A_p-3
<\frac14(4A_p-3)
=\frac14d_p^2.
\]
Thus \(d_J<d_p/2\).  The Clebsch--Gordan triangle inequality
\(J\ge|p-q|\) gives \(q\ge p-J\), and hence
\[
d_q=2q+1\ge2(p-J)+1=d_p-d_J+1>\frac12d_p.
\]
Therefore
\[
A_q=\frac{d_q^2+3}{4}
>
\frac{d_p^2/4+3}{4}
\ge\frac14\,\frac{d_p^2+3}{4}
=\frac14A_p.
\]
\end{proof}

\section{The higher-rank capacity--mark inequality}
\label{sec:V18-capacity-mark}

Let \(F\) be a pair bank as above with \(SU(2)\) input spins
\((p_f,q_f)\) and resolved output spins \(J_f\).  Assume that every
coordinate is strongly cancelling relative to row \(i\):
\[
A_{J_f}<\frac17A_{p_f}.
\]
Define
\begin{equation}
h_{i,F}:=\sum_{f\in F}A_{p_f},
\qquad
H_{ij,F}:=\sum_{f\in F}A_{p_f}A_{q_f}.
\label{eq:V18-cancelling-bank-masses}
\end{equation}

\begin{theorem}[Support-uniform higher-rank cancellation payment]
\label{thm:V18-higher-rank-capacity-mark}
For the complete bank projection \(P_F\),
\begin{equation}
\boxed{
h_{i,F}^2\,\kappa_\otimes(P_F)
\le4H_{ij,F}.}
\label{eq:V18-higher-rank-capacity-mark}
\end{equation}
The same estimate holds when \(P_F\) is part of a larger four-row
positive spectator carrier.
\end{theorem}

\begin{proof}
Put \(N=\lvert F\rvert\).  By
\Cref{lem:V18-SU2-cancellation} and
\Cref{thm:V18-cutwise-capacity},
\[
\kappa_\otimes(P_F)
\le\prod_{f\in F}\frac{d_{J_f}}{d_{p_f}}
\le2^{-N}.
\]
Cauchy--Schwarz and the partner-mass comparison give
\[
h_{i,F}^2
\le N\sum_fA_{p_f}^2
<4N\sum_fA_{p_f}A_{q_f}
=4NH_{ij,F}.
\]
Since \(N2^{-N}\le1\), multiplication by the capacity proves the
claim.  The larger-carrier statement is the final assertion of
\Cref{thm:V18-cutwise-capacity}.
\end{proof}

\section{Legal path routing for a strongly cancelling bank}
\label{sec:V18-routing}

Let \(a_r\) be the four junction input masses.  For the path
\(\sigma_{ij}\) with edges \(ij,ik,jl\), mark edge \(ij\) by
\(f\in F\) and the other edges by the junction.

\begin{lemma}[Unequal-input routed path identity]
\label{lem:V18-unequal-routing}
The exact subfamily of marked weights is
\begin{equation}
\boxed{
\sum_{f\in F}\mathcal W_{\sigma_{ij},\ell_f}
=
\frac{H_{ij,F}\prod_{r=1}^4a_r}{\Xi_i\Xi_j}.}
\label{eq:V18-unequal-routing}
\end{equation}
\end{lemma}

\begin{proof}
The \(ij\)-edge contributes
\[
\frac{A_{p_f}}{\Xi_i}\frac{A_{q_f}}{\Xi_j}.
\]
The two junction-marked edges and the leading factor
\(\prod_r\Xi_r\) cancel exactly as in
\Cref{lem:V16-routed-path-sum}.  Sum over \(f\).
\end{proof}

\begin{corollary}[Bank-dominant higher-rank routing]
\label{cor:V18-bank-dominant-routing}
Suppose row \(i\) is globally maximal and
\begin{equation}
\Xi_i\le C_0h_{i,F}.
\label{eq:V18-one-row-bank-dominance}
\end{equation}
Then
\begin{equation}
\boxed{
\kappa_\otimes(P_F)\prod_ra_r
\le4C_0^2
\sum_{f\in F}\mathcal W_{\sigma_{ij},\ell_f}.}
\label{eq:V18-bank-dominant-routing}
\end{equation}
\end{corollary}

\begin{proof}
By the partner comparison in
\Cref{lem:V18-SU2-cancellation},
\[
\sum_fA_{q_f}>\frac14h_{i,F}.
\]
This is a lower bound on \(\Xi_j\), while global maximality gives
\(\Xi_j\le\Xi_i\le C_0h_{i,F}\).  Only the upper bound is needed:
\[
\Xi_i\Xi_j\le C_0^2h_{i,F}^2.
\]
Use \Cref{thm:V18-higher-rank-capacity-mark} and then
\eqref{eq:V18-unequal-routing}.
\end{proof}

\begin{firewall}[The local core remains separate]
\label{fire:V18-local-core}
The V18 routing estimate controls the complete spectator projection
by legal marked weights.  Insertion into a junction carrier is valid
with the V2 one-sided tensor law if the other side is controlled in
operator norm by \(C\prod_ra_r\), as in V16--V17.  The hot
local-dominant non-Cartan payer has not yet been proved to satisfy that
operator bound.  V18 therefore does not silently promote spectator
routing to full \([4]\) closure.
\end{firewall}

\section{V18 ledger and V19 acquisition}
\label{sec:V18-ledger}

\begin{theorem}[Integrated compact-series V18 statement]
\label{thm:V18-integrated}
Preserving compact V1--V17, V18 proves:
\begin{enumerate}[label=\textup{(V18.\arabic*)},leftmargin=6.5em]
\item an explicit \(SU(d)\) counterexample to unrestricted
      edge-capacity multiplication;
\item the exact cutwise product of one-row marginal ratios for an
      arbitrary higher-rank pair bank;
\item the \(SU(2)\) implication from \(1/7\)-mass cancellation to
      half-dimension cancellation and partner-mass comparability;
\item the support-uniform capacity--bilinear-mark estimate
      \eqref{eq:V18-higher-rank-capacity-mark}; and
\item legal comparison-path routing for a bank which dominates one
      globally maximal row.
\end{enumerate}
\end{theorem}

\begin{proof}
These are
\Cref{prop:V18-antisymmetric-counterexample,
thm:V18-cutwise-capacity,lem:V18-SU2-cancellation,
thm:V18-higher-rank-capacity-mark,
cor:V18-bank-dominant-routing}.
\end{proof}

\begin{openproblem}[V19 mandate: combine retentive and cancelling pair
mass]
\label{op:V19-mandate}
\begin{enumerate}
\item Restrict first to \(SU(2)\) and exact incidence-two shared
      coordinates.
\item In a globally maximal all-nonlocal row, split each coordinate at
      the threshold \(A_U=A_{\rm input}/7\).
\item If the retentive subbank carries a fixed fraction of the row
      mass, use its Wilson quadratic moment without requiring the other
      coordinates to be retentive.
\item If the cancelling subbank dominates, select one of the three
      partners and apply
      \Cref{thm:V18-higher-rank-capacity-mark}.
\item Insert the resulting alternative below output priority, where
      the private bank has already paid the global resolvent.
\item Track the residual local factor
      \(s_\alpha^2L_\tau/\prod_ra_r\) exactly; close it or isolate the
      star-center hot subcase which prevents closure.
\end{enumerate}
\end{openproblem}

\begin{firewall}[Claim boundary after compact V18]
\label{fire:V18-current-boundary}
V18 closes the diffuse capacity/mark comparison for strongly
cancelling higher-rank \(SU(2)\) pair banks.  It does not yet combine
that comparison with every local tree fibre, handle arbitrary compact
simple groups, triple- or four-row cancelling spectators, the hot
local-dominant non-Cartan carrier, the full embedded \([4]\) sector,
multiple junctions, normalized collision aggregation, \([3,3]\),
global quartic or all-order MTA, continuum Yang--Mills existence, or a
positive mass gap.
\end{firewall}

\section{Append-only record for compact V18}
\label{sec:V18-record}

V18 was appended to the complete compact V17 file.  The exact V17
parent had \(6622\) newline-delimited source records, \(224573\) bytes,
and SHA--256 digest
\[
\texttt{65c4d022a5cbfc69ed52d48152cfd5aa5151e8e8b7f1a09b6f86b03e048f4a22}.
\]
No compact V17 mathematical content was changed.  The sole final
\verb|\end{document}| is restored below.

% =============================================================================
% END APPENDED COMPACT VERSION V18 MATERIAL
% =============================================================================

% The compact V18 document terminator is intentionally disabled here:
% \end{document}

% =============================================================================
% BEGIN APPENDED COMPACT VERSION V19 MATERIAL
% =============================================================================

\chapter{V19: Complete dual-pair moments and output aggregation}
\label{ch:V19}

\section{Why the resolved V18 estimate must be assembled}
\label{sec:V19-assembly}

V18 controls one prescribed tuple of cancelling output copies.  The
MTA carrier, however, contains the orthogonal sum over all such tuples.
Summing the block capacities would reintroduce an output count.  V19
therefore returns upstream to the complete positive moment operator at
each pair-shared coordinate and estimates the tensor product of those
moments before resolving its outputs.

\section{One complete dual-pair Wilson moment}
\label{sec:V19-one-coordinate}

Let \(R\) be a nontrivial irreducible representation of a compact
connected semisimple group, of dimension \(d_R\), and consider the
pair \(R\otimes R^*\).  Its complete positive Wilson moment has spectral
resolution
\begin{equation}
M_R=\sum_{U,a}q_{\alpha,U}P_{U,a},
\qquad 0\le M_R\le I.
\label{eq:V19-complete-pair-moment}
\end{equation}
The trivial output occurs once; denote its rank-one projection by
\(P_{\boldsymbol1,R}\).

\begin{lemma}[Uniform separation of every nontrivial output]
\label{lem:V19-nontrivial-separation}
In the established weak-coupling range \(0<s_\alpha\le1\), there is a
group constant \(c_*>0\) such that
\begin{equation}
\boxed{
M_R\preccurlyeq
(1-\delta_\alpha)I+\delta_\alpha P_{\boldsymbol1,R},
\qquad
\delta_\alpha:=c_*s_\alpha.}
\label{eq:V19-isotropic-majorant}
\end{equation}
The constant may be decreased so that \(0<\delta_\alpha\le1\).
\end{lemma}

\begin{proof}
For every nontrivial \(U\), the one-link gap gives
\[
1-q_{\alpha,U}
\ge c_{\rm gap}\min\{1,s_\alpha A_U\}.
\]
The nontrivial spectral masses have a positive group-uniform floor.
Since \(s_\alpha\le1\), the minimum is at least a constant times
\(s_\alpha\).  Thus
\(q_{\alpha,U}\le1-c_*s_\alpha\).
On the trivial line the multiplier is one.  Apply these scalar bounds
in the orthogonal spectral resolution
\eqref{eq:V19-complete-pair-moment}.
\end{proof}

\section{A complete bank without output counting}
\label{sec:V19-bank}

Let \(F\) be a nonempty bank of coordinates shared by the same two
rows.  At \(f\), the inputs are \(R_f,R_f^*\).  Put
\begin{align}
M_F&:=\bigotimes_{f\in F}M_{R_f},&
N&:=|F|,\label{eq:V19-bank-moment}\\
A_f&:=1+C_2(R_f),&
h_F&:=\sum_{f\in F}A_f,&
H_F&:=\sum_{f\in F}A_f^2.
\label{eq:V19-bank-mass}
\end{align}

\begin{theorem}[Complete-bank product capacity]
\label{thm:V19-complete-bank-capacity}
There are constants \(c,C>0\), independent of \(F\), its
representations, and \(\alpha\), such that
\begin{align}
\boxed{
\kappa_\otimes(M_F)
\le
\prod_{f\in F}
\left(1-\delta_\alpha+\frac{\delta_\alpha}{d_{R_f}}\right)
\le e^{-cs_\alpha N},}
\label{eq:V19-complete-bank-capacity}\\
\boxed{
s_\alpha h_F^2\kappa_\otimes(M_F)
\le CH_F,}
\label{eq:V19-first-bank-mark}\\
\boxed{
s_\alpha^2h_F^2\kappa_\otimes(M_F)
\le CH_F.}
\label{eq:V19-second-bank-mark}
\end{align}
All physical output copies have already been summed inside \(M_F\).
\end{theorem}

\begin{proof}
Tensor the Loewner bounds
\eqref{eq:V19-isotropic-majorant} and expand:
\[
M_F\preccurlyeq
\sum_{S\subseteq F}
(1-\delta_\alpha)^{N-|S|}
\delta_\alpha^{|S|}
\left(\bigotimes_{f\in S}P_{\boldsymbol1,R_f}\right)
\otimes I_{F\setminus S}.
\]
Capacity is monotone and subadditive on positive operators.  By the
exact bank calculation in
\Cref{lem:V16-bank-maximally-entangled},
\[
\kappa_\otimes\!\left[
\left(\bigotimes_{f\in S}P_{\boldsymbol1,R_f}\right)
\otimes I_{F\setminus S}\right]
=\prod_{f\in S}d_{R_f}^{-1}.
\]
Summing the resulting binomial expansion proves the first product in
\eqref{eq:V19-complete-bank-capacity}.  Every \(d_{R_f}\ge2\), so each
factor is at most \(1-\delta_\alpha/2\), which is bounded by
\(e^{-\delta_\alpha/2}\).

Cauchy--Schwarz gives \(h_F^2\le NH_F\).  Therefore
\[
s_\alpha h_F^2\kappa_\otimes(M_F)
\le
(s_\alpha N)e^{-cs_\alpha N}H_F
\le CH_F.
\]
The second estimate follows because \(s_\alpha\le1\).
\end{proof}

\begin{assessment}[The output aggregation gate is closed for a dual
pair bank]
\label{ass:V19-aggregation-closed}
Unlike the V18 resolved-copy estimate,
\Cref{thm:V19-complete-bank-capacity} acts on the complete output sum.
The one-coordinate singlet line is the only undamped eigenspace, and
its exact capacity is \(1/d_R\).  Expanding only this two-level
majorant produces a Bernoulli product which simultaneously retains
output aggregation, representation normalization, and support-uniform
decay.
\end{assessment}

\section{Selecting a dominant complete bank}
\label{sec:V19-selection}

Assume all off-junction shared coordinates have incidence exactly two
and carry dual inputs on their active row pair.  For each pair \(ij\),
let \(F_{ij}\) be its complete coordinate bank and use
\(h_{ij},H_{ij}\) from \eqref{eq:V17-pair-bank-data}.

\begin{lemma}[One complete bank dominates both endpoint rows]
\label{lem:V19-complete-bank-selection}
Fix \(D\ge16\).  Suppose
\[
a_r<\Xi_r/4\quad(1\le r\le4),
\qquad
B_\tau>DY^2
\]
for at least one local tree.  Choose a globally maximal row \(i\).
Then some \(j\ne i\) satisfies
\begin{equation}
\boxed{
\Xi_i<6h_{ij},\qquad
\Xi_j<6h_{ij}.}
\label{eq:V19-complete-bank-dominance}
\end{equation}
\end{lemma}

\begin{proof}
This is the deterministic proof of
\Cref{lem:V17-dominant-pair}.  The failed branching condition supplies
a row larger than \(4Y\), hence the global maximum \(i\) satisfies
\(p_i<\Xi_i/4\).  All-nonlocality then gives
\(h_i>\Xi_i/2\).  The three pair banks partition \(h_i\), so one
\(h_{ij}> \Xi_i/6\).  Global maximality gives
\(\Xi_j\le\Xi_i\).
\end{proof}

For the selected \(ij\), let \(\sigma_{ij}\) be the comparison path
and let its \(ij\)-edge range over \(f\in F_{ij}\), while the other two
edges are marked by the junction.  Then
\begin{equation}
\sum_{f\in F_{ij}}\mathcal W_{\sigma_{ij},\ell_f}
=
\frac{H_{ij}\prod_{r=1}^4a_r}{\Xi_i\Xi_j}
\label{eq:V19-complete-bank-routing-weight}
\end{equation}
by \Cref{lem:V16-routed-path-sum}.

\section{Insertion below output priority}
\label{sec:V19-below-output}

Put
\[
A_0:=\prod_{r=1}^4a_r,
\qquad
L_\tau:=\prod_{r=1}^4a_r^{d_r/2}.
\]
Retain \(0<s_\alpha\le1\).

\begin{lemma}[Private payment leaves exactly two local powers]
\label{lem:V19-private-leaves-two}
On
\[
\Xi_T\le DY,\qquad Y>0,
\]
the normalized positive carrier of one local tree fibre, with the
selected complete bank \(M_{F_{ij}}\) retained, satisfies
\begin{equation}
\boxed{
\kappa_\otimes(\mathcal Q_{\tau,ij})
\le
C_Ds_\alpha^2L_\tau
\kappa_\otimes(M_{F_{ij}}).}
\label{eq:V19-private-leaves-two}
\end{equation}
\end{lemma}

\begin{proof}
The global denominator obeys
\(1-q_{\alpha,\boldsymbol T}\ge1-q_P\), while
\(\Xi_T\le DY\).  The retained-failure private debit gives
\[
\frac{\Xi_Tq_P}{1-q_{\alpha,\boldsymbol T}}
\le\frac{C_D}{s_\alpha}.
\]
The V41 local tree operator is bounded by
\(Cs_\alpha^3L_\tau\).  Discard all unselected spectator contractions
and use the V2 one-sided tensor law with an operator norm on the local
core.  This gives \eqref{eq:V19-private-leaves-two}.
\end{proof}

\begin{lemma}[Path fibres and balanced stars are routed]
\label{lem:V19-balanced-fibres}
Under the hypotheses of
\Cref{lem:V19-complete-bank-selection} and below output priority, a
tree fibre obeys the comparison-path marked bound whenever
\begin{equation}
s_\alpha^2L_\tau\le A_0.
\label{eq:V19-balanced-local-factor}
\end{equation}
Every path tree satisfies
\eqref{eq:V19-balanced-local-factor}.
\end{lemma}

\begin{proof}
Combine \eqref{eq:V19-private-leaves-two} with
\eqref{eq:V19-second-bank-mark} and
\eqref{eq:V19-complete-bank-dominance}:
\begin{align*}
\kappa_\otimes(\mathcal Q_{\tau,ij})
&\le
C_D L_\tau\frac{H_{ij}}{h_{ij}^2}\\
&\le
C_D L_\tau\frac{H_{ij}}{\Xi_i\Xi_j}.
\end{align*}
If \eqref{eq:V19-balanced-local-factor} holds, this is bounded by the
right side of
\eqref{eq:V19-complete-bank-routing-weight}.

For a path with internal vertices \(u,v\) and leaves \(r,t\),
\[
L_\tau=a_ua_v(a_ra_t)^{1/2}\le A_0.
\]
Since \(s_\alpha\le1\), every path is balanced.
\end{proof}

\section{Returning upstream for a hot star}
\label{sec:V19-hot-star}

\begin{lemma}[Uniform cap for the complete local fourth cumulant]
\label{lem:V19-complete-cumulant-cap}
At one coordinate, before the V41 tree split, the complete fourth
cumulant operator satisfies
\begin{equation}
\boxed{\|K_4\|\le C_4}
\label{eq:V19-complete-cumulant-cap}
\end{equation}
uniformly in all input and output representations.
\end{lemma}

\begin{proof}
The moment--cumulant formula has fifteen set-partition terms with fixed
Moebius coefficients.  Each block moment is an average of unitaries
and hence a contraction; tensor products, recouplings, and physical
compressions remain contractions.  The triangle inequality over this
fixed fifteen-term formula gives a universal constant.  This is the
complete cumulant cap used already in the proof of
\Cref{cor:V3-junction-payer-cap}.
\end{proof}

\begin{lemma}[One hot star makes the complete local product large]
\label{lem:V19-hot-star-local-product}
If a star centered at \(c\) violates
\eqref{eq:V19-balanced-local-factor}, then
\begin{equation}
\boxed{A_0>s_\alpha^{-4}\ge s_\alpha^{-2}.}
\label{eq:V19-hot-star-local-product}
\end{equation}
\end{lemma}

\begin{proof}
For that star,
\[
\frac{L_\tau}{A_0}
=
\left(\frac{a_c}{\prod_{r\ne c}a_r}\right)^{1/2}.
\]
The violation says
\(s_\alpha^4a_c>\prod_{r\ne c}a_r\).
Since every local mass is at least one,
\(a_c>s_\alpha^{-4}\), and hence
\(A_0>s_\alpha^{-4}\).  The second inequality uses
\(s_\alpha\le1\).
\end{proof}

\begin{theorem}[Complete dual-pair branch below output priority]
\label{thm:V19-dual-pair-below-output}
Fix \(D\ge16\).  Assume:
\begin{enumerate}[label=\textup{(\roman*)},leftmargin=4em]
\item all shared spectators have incidence exactly two and dual inputs
      on their active pair;
\item \(a_r<\Xi_r/4\) for all four rows;
\item \(\Xi_T\le DY\); and
\item at least one of the sixteen local trees fails the V7
      private-dominant branching inequality
      \(B_\tau\le DY^2\).
\end{enumerate}
Then the complete unique-junction physical carrier obeys the sum of
legal marked-tree weights, uniformly in all shared output
representations and multiplicities.
\end{theorem}

\begin{proof}
The failed branching condition and the first two assumptions select
one complete bank \(F_{ij}\) by
\Cref{lem:V19-complete-bank-selection}.  If every star satisfies
\eqref{eq:V19-balanced-local-factor}, then
\Cref{lem:V19-balanced-fibres} closes all sixteen local tree fibres
against the single legal comparison-path subfamily
\eqref{eq:V19-complete-bank-routing-weight}; the finite reuse costs
only a universal factor.

Otherwise one star is hot.  Do not split the local cumulant into
trees.  The private payer and
\Cref{lem:V19-complete-cumulant-cap} give a complete local core
operator bound \(C_D/s_\alpha\).  The one-sided tensor law and
\eqref{eq:V19-first-bank-mark} then give
\[
\kappa_\otimes(\mathcal Q_{\rm complete})
\le
\frac{C_D}{s_\alpha}\kappa_\otimes(M_{F_{ij}})
\le
\frac{C_DH_{ij}}{s_\alpha^2h_{ij}^2}
\le
\frac{C_DH_{ij}}{s_\alpha^2\Xi_i\Xi_j}.
\]
By \Cref{lem:V19-hot-star-local-product},
\(s_\alpha^{-2}\le A_0\).  The last expression is therefore bounded by
\eqref{eq:V19-complete-bank-routing-weight}.  All shared outputs were
kept in the complete moment \(M_{F_{ij}}\), so no output count occurs.
\end{proof}

\begin{corollary}[All-nonlocal dual-pair closure outside
\textup{(O)}]
\label{cor:V19-dual-pair-outside-output}
Under the incidence and dual-input assumptions of
\Cref{thm:V19-dual-pair-below-output}, every all-nonlocal
unique-junction block outside output priority obeys the fourth
marked-tree atlas.
\end{corollary}

\begin{proof}
If the V7 private-dominant inequalities hold, use
\Cref{thm:V7-private-MTA}.  Otherwise use
\Cref{thm:V19-dual-pair-below-output}.  The case \(Y=0\) cannot occur
outside \textup{(O)} for a nontrivial product output.
\end{proof}

\begin{firewall}[Exact scope of the complete-output closure]
\label{fire:V19-scope}
V19 closes output aggregation for complete moment banks whose two
inputs are dual, because the undamped eigenspace is the unique singlet.
It does not yet treat a mixture of nondual pair inputs, incidence-three
or incidence-four shared coordinates, output priority \textup{(O)}, or
the local-dominant non-Cartan junction family.  Those cases cannot be
inserted by analogy.
\end{firewall}

\section{V19 ledger and mandatory V20 audit}
\label{sec:V19-ledger}

\begin{theorem}[Integrated compact-series V19 statement]
\label{thm:V19-integrated}
Preserving compact V1--V18, V19 proves:
\begin{enumerate}[label=\textup{(V19.\arabic*)},leftmargin=6.5em]
\item a two-level Loewner majorant for the complete dual-pair moment;
\item output-count-free capacity decay
      \(\kappa_\otimes(M_F)\le e^{-cs_\alpha|F|}\);
\item the first- and second-order complete-bank mark estimates
      \eqref{eq:V19-first-bank-mark}--%
      \eqref{eq:V19-second-bank-mark};
\item routing of every path and balanced-star local fibre;
\item upstream closure of the remaining hot-star sum by the complete
      cumulant cap; and
\item all-nonlocal, incidence-two dual-pair closure outside output
      priority.
\end{enumerate}
\end{theorem}

\begin{proof}
These are
\Cref{lem:V19-nontrivial-separation,
thm:V19-complete-bank-capacity,lem:V19-balanced-fibres,
lem:V19-complete-cumulant-cap,
thm:V19-dual-pair-below-output,
cor:V19-dual-pair-outside-output}.
\end{proof}

\begin{openproblem}[V20 mandate: audit and nondual pair banks]
\label{op:V20-mandate}
\begin{enumerate}
\item Begin with the mandatory read-only audit of the newest active SM
      and NS notes.
\item At a nondual pair coordinate, use the absence of a trivial output
      to obtain an operator-norm gap for the complete moment.
\item Split strongly unbalanced input pairs from comparable pairs using
      the fusion lower bound on the smallest output Casimir.
\item For comparable pairs, identify the legal bilinear edge mass as in
      V18--V19.
\item Assemble a complete nondual bank before resolving its outputs,
      and prove the analogue of
      \eqref{eq:V19-first-bank-mark}--%
      \eqref{eq:V19-second-bank-mark}.
\item Insert the result below output priority and state separately what
      remains at incidence three and four.
\end{enumerate}
\end{openproblem}

\begin{firewall}[Claim boundary after compact V19]
\label{fire:V19-current-boundary}
V19 closes the complete all-nonlocal dual-pair sector outside output
priority.  It does not close nondual, triple-, or four-row shared
banks, output priority for non-Cartan blocks, the local-dominant
non-Cartan carrier, the full embedded \([4]\) sector, multiple
junctions, normalized collision aggregation, \([3,3]\), global
quartic or all-order MTA, continuum Yang--Mills existence, or a
positive mass gap.
\end{firewall}

\section{Append-only record for compact V19}
\label{sec:V19-record}

V19 was appended to the complete compact V18 file.  The exact V18
parent had \(7025\) newline-delimited source records, \(237137\) bytes,
and SHA--256 digest
\[
\texttt{7aec6ff2c8a046bf95b5da3f1eb6bce5d8814f4882a3ceca7c994df80e871328}.
\]
No compact V18 mathematical content was changed.  The sole final
\verb|\end{document}| is restored below.

% =============================================================================
% END APPENDED COMPACT VERSION V19 MATERIAL
% =============================================================================

% The compact V19 document terminator is intentionally disabled here:
% \end{document}

% =============================================================================
% BEGIN APPENDED COMPACT VERSION V20 MATERIAL
% =============================================================================

\chapter{V20: Companion audit and complete nondual pair banks}
\label{ch:V20}

\section{Mandatory companion audit}
\label{sec:V20-audit}

The active SM note advanced from V18 at the V15 audit to V24.  The NS
note remains V31.  The exact read-only records are:
\begin{center}
\begin{tabular}{p{0.28\linewidth}p{0.12\linewidth}p{0.12\linewidth}
                p{0.38\linewidth}}
\toprule
file & lines & bytes & SHA--256\\
\midrule
\texttt{Renewal\_SM\_Einstein\_Continuum\_Research\_Notes\_V24.tex}
& \(8725\) & \(304431\)
& \texttt{f212459c45ff9b4b601560027329c10ed3c36bbad10aa7a626848add6dec9fbe}\\
\texttt{Renewal\_NS\_Stability\_Research\_Notes\_V31.tex}
& \(23052\) & \(887537\)
& \texttt{f1121b62238ad5491bb7cf03635ea9934942edf573ed8faaa9ee79e1d38a6696}\\
\bottomrule
\end{tabular}
\end{center}

SM V24 proves that the translation Ward identity protects a
nonvanishing zeroth stress form factor.  It corrects a proposed
four-jet argument by splitting off that protected term rather than
assigning it nonexistent decay.  This directly supports the V19
architecture: the unique undamped trivial output is isolated exactly,
and only the orthogonal nontrivial sector receives a Wilson gap.

NS V31 remains unchanged and continues to enforce the same weighted
lesson: a normalized mark does not supply its first moment.  No NS or
SM estimate is imported below.  Their common proof-order constraint is
implemented directly in the representation-theoretic bank.

\begin{assessment}[V20 audit decision]
\label{ass:V20-audit-decision}
For a nondual \(SU(2)\) pair there is no protected trivial output.
The complete output moment may therefore be contracted as a whole.
Comparable input pairs are paid by its uniform nontrivial spectral
gap; strongly unbalanced pairs are paid by the Casimir retained in
every possible output.  These two mechanisms are proved separately
before being combined.
\end{assessment}

\section{Complete moments for unequal \(SU(2)\) spins}
\label{sec:V20-nondual-moment}

Let \(p,q>0\) be unequal half-integers, and write
\begin{equation}
M_{p,q}:=\sum_{J=|p-q|}^{p+q}
q_{\alpha,J}P_J
\label{eq:V20-complete-nondual-moment}
\end{equation}
on \(V_p\otimes V_q\).  The Clebsch--Gordan sum is
multiplicity-free.

\begin{lemma}[No protected line in a nondual pair]
\label{lem:V20-no-trivial}
If \(p\ne q\), every output in
\eqref{eq:V20-complete-nondual-moment} is nontrivial.  Hence
\begin{equation}
\boxed{
\|M_{p,q}\|\le1-c_*s_\alpha}
\label{eq:V20-uniform-nondual-gap}
\end{equation}
in the established range \(0<s_\alpha\le1\).
\end{lemma}

\begin{proof}
The trivial spin occurs in \(V_p\otimes V_q\) exactly when \(p=q\).
For all remaining spins, apply the uniform nontrivial-output gap proved
in \Cref{lem:V19-nontrivial-separation}; the operator is positive and
diagonal in the displayed Clebsch--Gordan resolution.
\end{proof}

\begin{lemma}[Comparable inputs give a bilinear legal mass]
\label{lem:V20-comparable-inputs}
If \(q\ge p/2\), then
\begin{equation}
\boxed{A_q\ge\frac14A_p.}
\label{eq:V20-comparable-mass}
\end{equation}
\end{lemma}

\begin{proof}
Monotonicity of \(r\mapsto1+r(r+1)\) gives
\[
A_q\ge1+\frac p2\left(\frac p2+1\right)
\ge\frac14\bigl(1+p(p+1)\bigr).
\]
\end{proof}

\begin{lemma}[Unbalanced inputs force output retention]
\label{lem:V20-unbalanced-retention}
If \(q<p/2\), then every output spin \(J\) in
\(V_p\otimes V_q\) satisfies
\begin{equation}
\boxed{A_J\ge\frac14A_p.}
\label{eq:V20-unbalanced-retention}
\end{equation}
\end{lemma}

\begin{proof}
The triangle inequality gives
\(J\ge p-q>p/2\).  Apply the same monotonic comparison as in
\Cref{lem:V20-comparable-inputs}.
\end{proof}

\section{A complete comparable nondual bank}
\label{sec:V20-comparable-bank}

Let \(F\) be a finite bank on one row pair \(ij\).  At coordinate \(f\),
row \(i\) has spin \(p_f\), row \(j\) has a distinct spin \(q_f\), and
\(q_f\ge p_f/2\).  Put
\begin{align}
M_F&:=\bigotimes_{f\in F}M_{p_f,q_f},&
h_{i,F}&:=\sum_fA_{p_f},&
H_{ij,F}&:=\sum_fA_{p_f}A_{q_f}.
\label{eq:V20-comparable-bank-data}
\end{align}

\begin{theorem}[Complete comparable-bank mark estimates]
\label{thm:V20-comparable-bank}
Uniformly in the bank and all its representations,
\begin{align}
\boxed{\|M_F\|\le e^{-cs_\alpha|F|},}
\label{eq:V20-comparable-bank-gap}\\
\boxed{
s_\alpha h_{i,F}^2\|M_F\|
\le C H_{ij,F},}
\label{eq:V20-comparable-first-mark}\\
\boxed{
s_\alpha^2h_{i,F}^2\|M_F\|
\le C H_{ij,F}.}
\label{eq:V20-comparable-second-mark}
\end{align}
The complete output sum is retained in every \(M_{p_f,q_f}\).
\end{theorem}

\begin{proof}
Tensor \eqref{eq:V20-uniform-nondual-gap} to obtain the exponential
operator-norm bound.  If \(N=|F|\), then Cauchy--Schwarz and
\eqref{eq:V20-comparable-mass} give
\[
h_{i,F}^2
\le N\sum_fA_{p_f}^2
\le4NH_{ij,F}.
\]
Now use
\((s_\alpha N)e^{-cs_\alpha N}\le C\).
The second-order estimate follows from \(s_\alpha\le1\).
\end{proof}

\section{A complete unbalanced nondual bank}
\label{sec:V20-unbalanced-bank}

Let \(F\) instead satisfy \(q_f<p_f/2\) at every coordinate, and retain
\[
h_{i,F}:=\sum_fA_{p_f},
\qquad
M_F:=\bigotimes_fM_{p_f,q_f}.
\]

\begin{theorem}[The output bank pays the large input row]
\label{thm:V20-unbalanced-bank}
For \(k=1,2,3\),
\begin{equation}
\boxed{
(s_\alpha h_{i,F})^k\|M_F\|\le C_k.}
\label{eq:V20-unbalanced-bank}
\end{equation}
\end{theorem}

\begin{proof}
Let \(r_f:=\|M_{p_f,q_f}\|\), and choose an output spin \(J_f\) attaining
that maximum.  By
\Cref{lem:V20-unbalanced-retention},
\(A_{J_f}\ge A_{p_f}/4\).  The one-link retained-failure debits imply
\[
(s_\alpha A_{p_f})^kr_f
\le C_k(1-r_f),
\qquad 1\le k\le3.
\]
Apply the V7 product-moment compiler to
\((s_\alpha A_{p_f},r_f)\).  Since
\(\|M_F\|=\prod_fr_f\), the result follows.
\end{proof}

\begin{assessment}[Two nondual mechanisms]
\label{ass:V20-two-mechanisms}
On a comparable bank, the partner mass converts the support-uniform
spectral-gap factor into the legal bilinear mark
\(H_{ij,F}\).  On an unbalanced bank, the smallest possible output
already retains a fixed fraction of the large input mass, so ordinary
Wilson moments pay that row directly.  No output representation is
selected in either proof.
\end{assessment}

\section{Dominant-bank selection among all incidence-two coordinates}
\label{sec:V20-selection}

Assume now that every shared spectator coordinate has incidence exactly
two, with arbitrary nontrivial \(SU(2)\) input spins.  Fix an
all-nonlocal block and suppose at least one tree fails the V7 branching
condition.  Choose a row \(i\) of globally maximal \(\Xi_i\).

\begin{lemma}[Six-bank selection]
\label{lem:V20-six-bank-selection}
Among the three partner labels \(j\ne i\) and the two types
\[
\textup{equal spin }(p_f=q_f),
\qquad
\textup{unequal spin }(p_f\ne q_f),
\]
one bank \(F\) has row-\(i\) mass
\begin{equation}
\boxed{h_{i,F}>\frac1{12}\Xi_i.}
\label{eq:V20-six-bank-selection}
\end{equation}
For its partner \(j\),
\begin{equation}
\boxed{\Xi_j\le\Xi_i<12h_{i,F}.}
\label{eq:V20-selected-two-row-upper}
\end{equation}
\end{lemma}

\begin{proof}
As in \Cref{lem:V19-complete-bank-selection}, global maximality and
failed branching give \(h_i>\Xi_i/2\).  The six displayed banks
partition \(h_i\), so one has mass greater than \(\Xi_i/12\).
The second assertion is global maximality.
\end{proof}

If the selected bank has unequal spins, split it once more into
\[
F_{\rm cmp}:=\{f:q_f\ge p_f/2\},
\qquad
F_{\rm unb}:=\{f:q_f<p_f/2\},
\]
where \(p_f\) is the row-\(i\) spin.  One subbank has mass greater than
\(\Xi_i/24\).

\section{Closure below output priority}
\label{sec:V20-below-output}

\begin{theorem}[Complete incidence-two \(SU(2)\) closure below
\textup{(O)}]
\label{thm:V20-incidence-two-below-output}
Fix \(D\ge16\).  Suppose:
\begin{enumerate}[label=\textup{(\roman*)},leftmargin=4em]
\item every shared spectator coordinate has incidence exactly two;
\item \(a_r<\Xi_r/4\) for every row;
\item \(\Xi_T\le DY\); and
\item at least one local tree fails
      \(B_\tau\le DY^2\).
\end{enumerate}
For \(G=SU(2)\), the complete unique-junction physical carrier obeys
the sum of legal marked-tree weights, uniformly in all shared output
representations and multiplicities.
\end{theorem}

\begin{proof}
Select \(F\) by
\Cref{lem:V20-six-bank-selection}.

If \(F\) has equal spins, it is a complete dual-pair bank.  The proof of
\Cref{thm:V19-dual-pair-below-output} applies with the harmless
dominance constant twelve in place of six.

Suppose \(F\) has unequal spins.  If the unbalanced subbank has mass
\(h>\Xi_i/24\), retain only its complete moments and discard all other
spectators as contractions.  The private payer leaves the factor
\(C_Ds_\alpha^2L_\tau\), exactly as in
\Cref{lem:V19-private-leaves-two}.  By
\eqref{eq:V20-unbalanced-bank} with \(k=2\),
\[
s_\alpha^2\|M_{F_{\rm unb}}\|
\le C/h^2
\le C/\Xi_i^2.
\]
For every quartic tree,
\[
B_\tau=\prod_r\Xi_r^{d_r-1}\le\Xi_i^2
\]
because \(\sum_r(d_r-1)=2\) and \(\Xi_i\) is globally maximal.
Hence the tree contribution is bounded by
\[
C_DL_\tau/B_\tau
\le C_D\mathcal W_{\tau,e}.
\]
This closes all sixteen fibres pointwise.

It remains to treat the comparable subbank with
\(h>\Xi_i/24\).  Its partner obeys
\(\Xi_j\le\Xi_i<24h\).  The exact legal comparison-path sum is
\[
\frac{H_{ij,F_{\rm cmp}}A_0}{\Xi_i\Xi_j}.
\]
The first- and second-order estimates
\eqref{eq:V20-comparable-first-mark}--%
\eqref{eq:V20-comparable-second-mark}
replace their V19 dual-bank counterparts.  Therefore every path and
every star satisfying
\(s_\alpha^2L_\tau\le A_0\) is routed exactly as in
\Cref{lem:V19-balanced-fibres}.

If a star violates that inequality, return to the complete local
cumulant.  The private payer and
\Cref{lem:V19-complete-cumulant-cap} leave \(C_D/s_\alpha\).
The first-order comparable-bank estimate gives
\[
\frac{C_D}{s_\alpha}\|M_{F_{\rm cmp}}\|
\le
\frac{C_DH_{ij,F_{\rm cmp}}}
     {s_\alpha^2h^2}
\le
\frac{C_DH_{ij,F_{\rm cmp}}}
     {s_\alpha^2\Xi_i\Xi_j}.
\]
The hot-star lemma gives \(A_0>s_\alpha^{-4}\), so this is bounded by
the legal comparison-path sum.  No output tuple was separated in any
of the three cases.
\end{proof}

\begin{corollary}[All-nonlocal incidence-two sector outside output
priority]
\label{cor:V20-incidence-two-outside-output}
For \(SU(2)\), every all-nonlocal unique-junction block whose shared
coordinates all have incidence two obeys the fourth marked-tree atlas
outside priority family \textup{(O)}.
\end{corollary}

\begin{proof}
If every tree satisfies the V7 private-dominant inequalities, apply
\Cref{thm:V7-private-MTA}.  Otherwise apply
\Cref{thm:V20-incidence-two-below-output}.  As before, a nontrivial
block outside \textup{(O)} has \(Y>0\).
\end{proof}

\begin{firewall}[What V20 does not infer from the audit]
\label{fire:V20-audit-scope}
SM V24's Ward argument motivates exact separation of protected and
gapped sectors, but none of its estimates enters
\Cref{thm:V20-incidence-two-below-output}.  That theorem is proved from
\(SU(2)\) fusion, Wilson gaps, complete moment operators, and legal
MTA markings.  It does not apply to a general compact simple group
without a replacement for the spin-order dichotomy.
\end{firewall}

\section{V20 ledger and V21 acquisition}
\label{sec:V20-ledger}

\begin{theorem}[Integrated compact-series V20 statement]
\label{thm:V20-integrated}
Preserving compact V1--V19, V20 proves:
\begin{enumerate}[label=\textup{(V20.\arabic*)},leftmargin=6.5em]
\item the mandatory hashed audit of active SM V24 and NS V31;
\item a complete-output uniform gap for every unequal-spin
      \(SU(2)\) pair;
\item complete comparable-bank bilinear mark estimates;
\item complete unbalanced-bank input-moment estimates;
\item deterministic selection among all equal/unequal pair banks; and
\item all-nonlocal incidence-two \(SU(2)\) closure outside output
      priority.
\end{enumerate}
\end{theorem}

\begin{proof}
These are
\Cref{ass:V20-audit-decision,lem:V20-no-trivial,
thm:V20-comparable-bank,thm:V20-unbalanced-bank,
lem:V20-six-bank-selection,
thm:V20-incidence-two-below-output,
cor:V20-incidence-two-outside-output}.
\end{proof}

\begin{openproblem}[V21 mandate: incidence-two output priority]
\label{op:V21-mandate}
\begin{enumerate}
\item Work on \(\Xi_T>DY\) with \(SU(2)\) incidence-two shared
      coordinates.
\item Retain the V11 split into the junction-output bank and the
      complete shared-output bank.
\item Separate equal-spin protected singlets from unequal-spin fully
      gapped moments before applying a shared-output resolvent.
\item Reuse the comparable/unbalanced dichotomy to pay the dominant
      input row.
\item Assemble all output tuples inside complete moments; do not sum
      resolved-copy capacities.
\item Close the all-nonlocal output-priority sector, or isolate the
      exact resolvent-weighted complete-bank capacity still missing.
\end{enumerate}
\end{openproblem}

\begin{firewall}[Claim boundary after compact V20]
\label{fire:V20-current-boundary}
V20 closes the complete all-nonlocal incidence-two \(SU(2)\) sector
outside output priority.  It does not close its non-Cartan
output-priority part, local-dominant non-Cartan blocks, incidence-three
or incidence-four shared coordinates, general compact simple groups,
the full embedded \([4]\) sector, multiple junctions, normalized
collision aggregation, \([3,3]\), global quartic or all-order MTA,
continuum Yang--Mills existence, or a positive mass gap.
\end{firewall}

\section{Append-only record for compact V20}
\label{sec:V20-record}

V20 was appended to the complete compact V19 file.  The exact V19
parent had \(7499\) newline-delimited source records, \(252931\) bytes,
and SHA--256 digest
\[
\texttt{dab2c84b0cdbe6e93e512896b36100926cc05525908e49068b06d3bdeb149691}.
\]
No compact V19 mathematical content was changed.  The sole final
\verb|\end{document}| is restored below.

% =============================================================================
% END APPENDED COMPACT VERSION V20 MATERIAL
% =============================================================================

% The compact V20 document terminator is intentionally disabled here:
% \end{document}

% =============================================================================
% BEGIN APPENDED COMPACT VERSION V21 MATERIAL
% =============================================================================

\chapter{V21: Protected-line separation in complete pair-bank
resolvents}
\label{ch:V21}

\section{The upstream correction forced by output priority}
\label{sec:V21-upstream}

For a Cartan shared-output bank every active output is nontrivial, so
the V11 expression
\[
   \frac{H_Tq_S}{1-q_S}
\]
has a positive denominator.  This fails on an equal-spin \(SU(2)\)
pair bank: the tensor product of all singlet outputs has
\(q_S=1\).  It is therefore not legitimate to apply the V11
shared-bank resolvent moment to the whole non-Cartan carrier.

The correct proof order is the one identified in the V20 companion
audit.  First split off the unique protected line, then estimate the
complete orthogonal output sum without resolving and counting its
copies.  Only after that split may the remaining external resolvent be
assigned to the junction, private, or other shared bank.

\section{The complete equal-spin bank}
\label{sec:V21-equal-bank}

Let \(F\ne\varnothing\) be a bank of \(N:=|F|\) coordinates shared by
two rows.  At \(f\), both inputs have spin \(p_f>0\).  Write
\[
 V_{p_f}\otimes V_{p_f}^*
   =\bigoplus_{J=0}^{2p_f}V_J,
 \qquad
 z_J:=C_2(J)=J(J+1),
 \qquad
 b_J:=1+z_J,
\]
and denote the corresponding projection by \(P_{J,f}\).  Put
\begin{align}
 M_f^{(n)}
   &:=\sum_{J=0}^{2p_f}q_{\alpha,J}^{\,n}P_{J,f},
 &
 M_F^{(n)}
   &:=\bigotimes_{f\in F}M_f^{(n)}, \label{eq:V21-power-moment}\\
 P_{0,F}
   &:=\bigotimes_{f\in F}P_{0,f}.
\label{eq:V21-protected-bank}
\end{align}
For an output tuple
\(\boldsymbol J=(J_f)_{f\in F}\), use
\[
 q_{\boldsymbol J}:=\prod_fq_{\alpha,J_f},
 \qquad
 Z_{\boldsymbol J}:=\sum_fz_{J_f},
 \qquad
 H_{\boldsymbol J}:=N+Z_{\boldsymbol J}.
\]
The only tuple with \(q_{\boldsymbol J}=1\) is
\(\boldsymbol J=\boldsymbol0\).

Define positive operators on the orthogonal complement of the
protected line by
\begin{align}
 \mathcal G_F
 &:=
 \sum_{\boldsymbol J\ne\boldsymbol0}
 \frac{q_{\boldsymbol J}}{1-q_{\boldsymbol J}}
 P_{\boldsymbol J},
\label{eq:V21-geometric-resolvent}\\
 \mathcal C_F
 &:=
 \sum_{\boldsymbol J\ne\boldsymbol0}
 \frac{Z_{\boldsymbol J}q_{\boldsymbol J}}
      {1-q_{\boldsymbol J}}
 P_{\boldsymbol J},
\label{eq:V21-Casimir-resolvent}\\
 \mathcal R_F
 &:=
 N\mathcal G_F+\mathcal C_F
 =
 \sum_{\boldsymbol J\ne\boldsymbol0}
 \frac{H_{\boldsymbol J}q_{\boldsymbol J}}
      {1-q_{\boldsymbol J}}
 P_{\boldsymbol J}.
\label{eq:V21-complete-resolvent}
\end{align}
Every output tuple and multiplicity has already been summed in these
operators.

\begin{lemma}[A geometric capacity bound with the protected line
removed]
\label{lem:V21-geometric-capacity}
There is a group-uniform constant \(C\) such that
\begin{equation}
\boxed{
 \kappa_\otimes(\mathcal G_F)
 \le \frac{C}{s_\alpha N}.}
\label{eq:V21-geometric-capacity}
\end{equation}
\end{lemma}

\begin{proof}
Let \(\delta_\alpha=c_*s_\alpha\) be the separation constant in
\Cref{lem:V19-nontrivial-separation}, decrease it if necessary so that
\(0<\delta_\alpha<1\), and put
\[
 r_n:=(1-\delta_\alpha)^n.
\]
The spectral resolution gives
\[
 M_f^{(n)}
 \preccurlyeq r_nI+(1-r_n)P_{0,f}.
\]
Since \(M_F^{(n)}\) and its majorant both equal one on \(P_{0,F}\),
\begin{equation}
 0\le M_F^{(n)}-P_{0,F}
 \preccurlyeq
 \bigotimes_f\bigl(r_nI+(1-r_n)P_{0,f}\bigr)-P_{0,F}.
\label{eq:V21-power-majorant}
\end{equation}

For \(S\subseteq F\), let
\[
 Q_S:=
 \left(\bigotimes_{f\in S}P_{0,f}\right)
 \otimes I_{F\setminus S}.
\]
Expanding the right side of
\eqref{eq:V21-power-majorant} and collecting its protected component
gives the positive identity
\[
 \bigotimes_f\bigl(r_nI+(1-r_n)P_{0,f}\bigr)-P_{0,F}
 =
 \sum_{S\subsetneq F}
 r_n^{N-|S|}(1-r_n)^{|S|}
 (Q_S-P_{0,F}).
\]
The exact maximally-entangled bank calculation in
\Cref{lem:V16-bank-maximally-entangled} and monotonicity give
\[
 \kappa_\otimes(Q_S-P_{0,F})
 \le\kappa_\otimes(Q_S)
 =\prod_{f\in S}d_{p_f}^{-1}.
\]
As every \(d_{p_f}\ge2\), and the resulting polynomial is increasing
in each \(d_{p_f}^{-1}\),
\begin{equation}
 \kappa_\otimes(M_F^{(n)}-P_{0,F})
 \le
 f_N(r_n):=
 2^{-N}\bigl[(1+r_n)^N-(1-r_n)^N\bigr].
\label{eq:V21-fN-bound}
\end{equation}

Put
\(\lambda_\alpha:=-\log(1-\delta_\alpha)\ge\delta_\alpha\).
The function \(f_N\) is increasing, so comparison on the intervals
\([e^{-\lambda_\alpha n},e^{-\lambda_\alpha(n-1)}]\) yields
\[
 \sum_{n\ge1}f_N(e^{-\lambda_\alpha n})
 \le
 \lambda_\alpha^{-1}
 \int_0^1\frac{f_N(r)}r\,dr.
\]
The odd binomial expansion and
\(k^{-1}\le2(k+1)^{-1}\) for \(k\ge1\) give
\begin{align*}
 \int_0^1\frac{f_N(r)}r\,dr
 &=
 2^{1-N}
 \sum_{\substack{1\le k\le N\\k\ {\rm odd}}}
 \binom Nk\frac1k\\
 &\le
 2^{2-N}\sum_{k=0}^N\binom Nk\frac1{k+1}
 =
 \frac{2^{2-N}(2^{N+1}-1)}{N+1}
 \le\frac8{N+1}.
\end{align*}
Finally, spectral geometric summation gives
\[
 \mathcal G_F
 =\sum_{n\ge1}(M_F^{(n)}-P_{0,F}).
\]
Capacity is subadditive on positive operators.  Combine the last four
displays, \(\lambda_\alpha\ge c_*s_\alpha\), and \(N\ge1\).
\end{proof}

\begin{lemma}[The Casimir part needs no output count]
\label{lem:V21-Casimir-resolvent}
Uniformly in \(F\) and all its spins,
\begin{equation}
\boxed{
 0\le\mathcal C_F
 \preccurlyeq\frac{C}{s_\alpha}(I-P_{0,F}),
 \qquad
 \kappa_\otimes(\mathcal C_F)\le\frac{C}{s_\alpha}.}
\label{eq:V21-Casimir-resolvent-bound}
\end{equation}
\end{lemma}

\begin{proof}
Fix \(\boldsymbol J\ne\boldsymbol0\).  Apply the retained-failure
compiler \Cref{lem:V11-retained-failure} with \(m=1\),
\[
 x_f=s_\alpha z_{J_f},
 \qquad
 r_f=q_{\alpha,J_f}.
\]
For \(J_f=0\), both \(x_f\) and \(1-r_f\) vanish.  For
\(J_f>0\), the established one-link linear Wilson debit applies.
Consequently,
\[
 s_\alpha Z_{\boldsymbol J}q_{\boldsymbol J}
 \le C(1-q_{\boldsymbol J}).
\]
This is exactly the claimed scalar coefficient bound on every
nonprotected spectral block.  The capacity statement follows from
monotonicity and \(\kappa_\otimes(I-P_{0,F})\le1\).
\end{proof}

\begin{theorem}[Complete equal-bank resolvent capacity]
\label{thm:V21-equal-resolvent}
For every nonempty equal-spin pair bank,
\begin{equation}
\boxed{
 \kappa_\otimes(\mathcal R_F)\le\frac{C}{s_\alpha}.}
\label{eq:V21-equal-resolvent}
\end{equation}
The bound is independent of the support size, input spins, and number
of output tuples.
\end{theorem}

\begin{proof}
By \eqref{eq:V21-complete-resolvent}, positivity, and
\Cref{lem:V21-geometric-capacity,lem:V21-Casimir-resolvent},
\[
 \kappa_\otimes(\mathcal R_F)
 \le
 N\kappa_\otimes(\mathcal G_F)
 +\kappa_\otimes(\mathcal C_F)
 \le\frac{C}{s_\alpha}.
\]
\end{proof}

\section{The fully gapped unequal-spin bank}
\label{sec:V21-unequal-bank}

Let \(F\) now be a bank in which the two input spins are unequal at
every coordinate.  Define \(H_{\boldsymbol J}\),
\(q_{\boldsymbol J}\), and
\[
 \mathcal R_F^{\ne}:=
 \sum_{\boldsymbol J}
 \frac{H_{\boldsymbol J}q_{\boldsymbol J}}
      {1-q_{\boldsymbol J}}P_{\boldsymbol J}
\]
with the complete Clebsch--Gordan output at every coordinate.

\begin{theorem}[Unequal banks are resolvent-gapped in operator norm]
\label{thm:V21-unequal-resolvent}
Uniformly in the complete unequal bank,
\begin{equation}
\boxed{
 \|\mathcal R_F^{\ne}\|\le\frac{C}{s_\alpha}.}
\label{eq:V21-unequal-resolvent}
\end{equation}
\end{theorem}

\begin{proof}
Every coordinate output is nontrivial by
\Cref{lem:V20-no-trivial}.  If \(N=|F|\), then
\[
 q_{\boldsymbol J}\le(1-c_*s_\alpha)^N
 =e^{-\lambda_\alpha N}.
\]
Thus
\[
 \frac{Nq_{\boldsymbol J}}{1-q_{\boldsymbol J}}
 \le\frac{N}{e^{\lambda_\alpha N}-1}
 \le\frac1{\lambda_\alpha}
 \le\frac{C}{s_\alpha}.
\]
The retained-failure argument in
\Cref{lem:V21-Casimir-resolvent} also gives
\[
 \frac{Z_{\boldsymbol J}q_{\boldsymbol J}}
      {1-q_{\boldsymbol J}}
 \le\frac{C}{s_\alpha}.
\]
Add the two scalar bounds on every spectral block.
\end{proof}

\section{Keeping the external resolvent on the protected line}
\label{sec:V21-external}

For \(0\le z<1\), define the externally damped complete equal-bank
operator
\begin{equation}
 \mathcal E_F(z):=
 \sum_{\boldsymbol J}
 \frac{H_{\boldsymbol J}zq_{\boldsymbol J}}
      {1-zq_{\boldsymbol J}}P_{\boldsymbol J}.
\label{eq:V21-external-operator}
\end{equation}
Here \(z\) is the product of the junction, private, and other shared
Wilson multipliers.  It must not be set equal to one before the
protected line is separated.

\begin{theorem}[Exact protected/gapped external split]
\label{thm:V21-external-split}
For an equal-spin bank,
\begin{equation}
\boxed{
 \mathcal E_F(z)
 \preccurlyeq
 \mathcal R_F+
 \frac{Nz}{1-z}P_{0,F},}
\label{eq:V21-external-split}
\end{equation}
and hence
\begin{equation}
\boxed{
 \kappa_\otimes(\mathcal E_F(z))
 \le
 \frac{C}{s_\alpha}
 \frac{Nz}{1-z}\prod_{f\in F}d_{p_f}^{-1}
 \le
 \frac{C}{s_\alpha}+C\frac{z}{1-z}.}
\label{eq:V21-external-capacity}
\end{equation}
For an unequal-spin bank,
\begin{equation}
\boxed{\|\mathcal E_F(z)\|\le C/s_\alpha.}
\label{eq:V21-unequal-external}
\end{equation}
\end{theorem}

\begin{proof}
If \(\boldsymbol J\ne\boldsymbol0\), then
\[
 \frac{zq_{\boldsymbol J}}{1-zq_{\boldsymbol J}}
 \le
 \frac{q_{\boldsymbol J}}{1-q_{\boldsymbol J}}.
\]
On the protected tuple,
\(H_{\boldsymbol0}=N\) and \(q_{\boldsymbol0}=1\).
This proves the operator inequality
\eqref{eq:V21-external-split}.  Apply
\Cref{thm:V21-equal-resolvent} and the exact bank capacity
\[
 \kappa_\otimes(P_{0,F})
 =\prod_fd_{p_f}^{-1}.
\]
Since \(d_{p_f}\ge2\), one has
\(N\prod_fd_{p_f}^{-1}\le N2^{-N}\le1\).
This proves \eqref{eq:V21-external-capacity}.

For an unequal bank there is no protected tuple.  Apply the same
scalar comparison and
\Cref{thm:V21-unequal-resolvent}.
\end{proof}

\begin{assessment}[What has been removed from output priority]
\label{ass:V21-output-count-removed}
The gapped part of every complete \(SU(2)\) pair bank now pays one
shared-output resolvent with cost \(C/s_\alpha\), before any output is
resolved.  No Clebsch--Gordan output count, multiplicity count, or
sum of resolved capacities remains.  Equal-spin banks have exactly
one additional term: the all-singlet projection carrying the
external factor \(z/(1-z)\).
\end{assessment}

\begin{firewall}[The protected term cannot be put back into the V11
denominator]
\label{fire:V21-protected-denominator}
On \(P_{0,F}\), the internal bank multiplier is one.  Therefore
\[
 \frac{Nq_F}{1-q_F}
\]
is infinite and has no shared-bank interpretation.  The finite
quantity is
\[
 \frac{Nz}{1-z}P_{0,F},
\]
with the external multiplier and its sole resolvent kept together.
Capacity supplies
\(N\prod_fd_{p_f}^{-1}\le C\), but V15 shows that it does not by itself
supply an arbitrary full input-Casimir moment.  Any proof which erases
\(z/(1-z)\), or replaces it by an internal shared denominator, is
circular.
\end{firewall}

\section{The exact residual incidence-two output-priority problem}
\label{sec:V21-residual}

On an all-nonlocal output-priority tree with
\(B_\tau\le DY^2\), the gapped shared payer now leaves
\[
 C_Ds_\alpha^2q_PL_\tau
 \le C_D\frac{L_\tau}{B_\tau}
\]
by the private quadratic debit.  Thus the complete gapped
shared-resolvent part of this branching family is closed.

Two coupled terms remain before one may claim the full incidence-two
output-priority sector:
\begin{enumerate}
\item the equal-bank protected carrier
      \[
      \frac{Nz}{1-z}P_{0,F}
      \]
      must be inserted while retaining the external junction/private
      payer; and
\item the junction-payer term containing
      \(s_\alpha A_{U_e}\) must route a possibly large cancelling
      shared input mass through legal edge marks.  It cannot use the
      retentive comparison \(H\lesssim H_T\), which is false on a
      singlet output.
\end{enumerate}
On a failed-branching tree \(B_\tau>DY^2\), one additionally needs a
marked version of the complete resolvent estimate which couples the
dominant incident pair bank to the legal comparison path.  These are
external-resolvent and mark-routing questions, not output-aggregation
questions.

\begin{openproblem}[V22 mandate: protected external payer and legal
mark routing]
\label{op:V22-mandate}
\begin{enumerate}
\item Keep the factor \(z/(1-z)\) on the all-singlet line and identify
      whether \(z\) is paid by the junction bank, private bank, or a
      second shared bank.
\item Use the exact singlet capacity only after that external payer is
      fixed.
\item For the \(s_\alpha A_{U_e}\) junction term, route cancelling
      input Casimir to actual pair-edge marks; do not assert
      \(H\lesssim H_T\).
\item Treat \(B_\tau\le DY^2\) and \(B_\tau>DY^2\) separately.
\item If several pair banks overlap, use a one-sided insertion or a
      cutwise positive-operator estimate.  Do not multiply their
      capacities.
\item Close the protected output-priority carrier, or state the exact
      two-bank inequality still missing.
\end{enumerate}
\end{openproblem}

\begin{firewall}[Claim boundary after compact V21]
\label{fire:V21-current-boundary}
V21 proves complete output-summed resolvent bounds for the gapped part
of every \(SU(2)\) pair bank and isolates the unique protected
all-singlet external-resolvent term.  It does not yet close that
protected term, the cancelling junction-input mark, the full
incidence-two output-priority sector, local-dominant non-Cartan
blocks, incidence-three or incidence-four shared coordinates, general
compact simple groups, the full embedded \([4]\) sector, multiple
junctions, normalized collision aggregation, \([3,3]\), global
quartic or all-order MTA, continuum Yang--Mills existence, or a
positive mass gap.
\end{firewall}

\section{Append-only record for compact V21}
\label{sec:V21-record}

V21 was appended to the complete compact V20 file.  The exact V20
parent had \(7930\) newline-delimited source records, \(267288\) bytes,
and SHA--256 digest
\[
\texttt{a4719abb5d42378fb763c83ce3d554852d71604300ae5f367303d406dce7a7b0}.
\]
No compact V20 mathematical content was changed.  The sole final
\verb|\end{document}| is restored below.

% =============================================================================
% END APPENDED COMPACT VERSION V21 MATERIAL
% =============================================================================

% The compact V21 document terminator is intentionally disabled here:
% \end{document}

% =============================================================================
% BEGIN APPENDED COMPACT VERSION V22 MATERIAL
% =============================================================================

\chapter{V22: The global all-singlet network and its external payer}
\label{ch:V22}

\section{Why the protected sector must be assembled globally}
\label{sec:V22-global-assembly}

The protected term in \Cref{thm:V21-external-split} was written for
one row pair.  In an incidence-two carrier, several equal-spin pair
banks may overlap at a row.  Their capacities cannot be multiplied:
the V18 entanglement-swapping example forbids that step.  The genuinely
singular shared output, however, has a simpler global description.  It
exists only when every shared coordinate has equal dual inputs and
every one of those coordinates is projected to its singlet.  The
resulting carrier is one tensor product of EPR projections on a
four-vertex multigraph.

V22 keeps that entire graph projection intact.  A Schmidt bound across
each vertex cut replaces any multiplication of edge capacities.

\section{Cut capacity of an EPR multigraph}
\label{sec:V22-cut-capacity}

Let \(E\ne\varnothing\) be a finite multiset of edges on the vertex set
\(\{1,2,3,4\}\).  Edge \(f=uv\) carries equal nontrivial spin \(p_f\)
at \(u\) and \(v\).  Set
\[
 d_f:=2p_f+1,
 \qquad
 A_f:=1+p_f(p_f+1),
 \qquad
 P_E:=\bigotimes_{f\in E}P_{0,f}.
\]
The row Hilbert space groups all edge factors incident to that row.
Put
\begin{align}
 N&:=|E|,
 &
 S&:=\sum_{f\in E}A_f,
 &
 H&:=\sum_{f\in E}A_f^2,
\label{eq:V22-network-data}\\
 D_r&:=\prod_{\substack{f\in E\\r\in f}}d_f,
 &
 D_E&:=\prod_{f\in E}d_f.
\label{eq:V22-cut-dimensions}
\end{align}
Empty products in \(D_r\) equal one.

\begin{lemma}[Vertex-cut EPR capacity]
\label{lem:V22-vertex-cut}
The four-row product-state capacity satisfies
\begin{equation}
\boxed{
 \kappa_\otimes(P_E)
 \le\min_{1\le r\le4}D_r^{-1}
 \le D_E^{-1/2}.}
\label{eq:V22-vertex-cut}
\end{equation}
\end{lemma}

\begin{proof}
Fix a row \(r\) and group the other three rows into one tensor factor.
Every singlet edge incident to \(r\) is maximally entangled across
this bipartition.  Their tensor product has \(D_r\) equal Schmidt
coefficients, while every edge not incident to \(r\) lies entirely on
the other side.  Hence the largest squared Schmidt coefficient is
\(D_r^{-1}\).  Four-row product vectors are a subset of product
vectors across this bipartition, proving the first inequality.

Each edge dimension occurs in exactly two of the four vertex products,
so
\[
 \prod_{r=1}^4D_r=D_E^2.
\]
Consequently \(\max_rD_r\ge D_E^{1/2}\), and minimizing \(D_r^{-1}\)
proves the second inequality.
\end{proof}

\begin{theorem}[The output-count-weighted global singlet mark]
\label{thm:V22-weighted-network}
Uniformly in the multigraph and all its spins,
\begin{align}
\boxed{N\kappa_\otimes(P_E)\le C,}
\label{eq:V22-count-capacity}\\
\boxed{
 NS^2\kappa_\otimes(P_E)\le CH.}
\label{eq:V22-weighted-network}
\end{align}
\end{theorem}

\begin{proof}
For \(SU(2)\),
\[
 d_f^2=4A_f-3\ge A_f.
\]
Since every \(p_f>0\), also \(A_f\ge7/4\).  Therefore
\[
 \kappa_\otimes(P_E)
 \le D_E^{-1/2}
 \le\prod_{f\in E}A_f^{-1/4}
 \le\rho^N,
 \qquad
 \rho:=(7/4)^{-1/4}<1.
\]
This immediately gives
\eqref{eq:V22-count-capacity}.  Cauchy--Schwarz gives
\(S^2\le NH\), whence
\[
 \frac{NS^2}{H}\kappa_\otimes(P_E)
 \le N^2\rho^N.
\]
The last quantity is uniformly bounded in \(N\), proving
\eqref{eq:V22-weighted-network}.
\end{proof}

\begin{remark}[The gain is global, not edgewise]
\label{rem:V22-global-not-edgewise}
No assertion of the form
\(\kappa_\otimes(P_E)=\prod_f\kappa_\otimes(P_{0,f})\)
is used.  Only four valid bipartite Schmidt bounds are taken, followed
by their geometric mean.  Thus
\Cref{thm:V22-weighted-network} remains valid in the overlapping-bank
configurations where naive capacity multiplication fails.
\end{remark}

\section{The legal path sum carried by the network}
\label{sec:V22-legal-routing}

For an edge \(f=uv\), let \(\sigma_{uv}\) be the comparison path whose
\(uv\)-edge is marked at \(f\), while its other two edges are marked at
the junction.  Put
\[
 A_0:=\prod_{r=1}^4a_r,
 \qquad
 \Xi_{\max}:=\max_r\Xi_r.
\]

\begin{lemma}[Global EPR marks route without a bank selection]
\label{lem:V22-global-routing}
The legal marked-tree atlas contains the sum
\begin{equation}
\boxed{
 \mathfrak P_E:=
 \sum_{f=uv\in E}\mathcal W_{\sigma_{uv},\ell_f}
 =
 A_0\sum_{f=uv\in E}\frac{A_f^2}{\Xi_u\Xi_v}
 \ge
 \frac{A_0H}{\Xi_{\max}^2}.}
\label{eq:V22-global-routing}
\end{equation}
\end{lemma}

\begin{proof}
The equality is the one-coordinate path calculation of
\Cref{lem:V16-routed-path-sum}, summed over every shared coordinate
exactly once.  The inequality follows from
\(\Xi_u\Xi_v\le\Xi_{\max}^2\).
\end{proof}

\section{Payment of the genuinely protected output-priority line}
\label{sec:V22-protected-payment}

Assume every off-junction shared coordinate has incidence two and
equal nontrivial spins on its two active rows.  Resolve every shared
coordinate to its singlet.  Then the shared multiplier is one and the
shared output mass is exactly \(N\).  Put
\[
 z:=q_{\alpha,U_e}q_P.
\]
On a retained nontrivial global output, \(z<1\).  Let \(Q_{\tau,e}\)
denote the positive local tree carrier, so that
\begin{equation}
 \|Q_{\tau,e}\|\le Cs_\alpha^3L_\tau,
 \qquad
 L_\tau:=\prod_{r=1}^4a_r^{d_r/2}.
\label{eq:V22-local-tree}
\end{equation}
The shared-output part of the normalized protected carrier is bounded,
up to the fixed V8 output-priority constant, by
\begin{equation}
 \mathfrak C_{\tau,E}:=
 \frac{Nz}{1-z}\,
 \kappa_\otimes(Q_{\tau,e}\otimes P_E).
\label{eq:V22-protected-carrier}
\end{equation}

\begin{theorem}[The global all-singlet line obeys the marked atlas]
\label{thm:V22-protected-line}
Fix \(D\ge16\).  Suppose
\[
 \Xi_T>DY,
 \qquad
 a_r<\Xi_r/4\quad(1\le r\le4).
\]
Then every globally protected incidence-two shared-output
contribution is bounded by the legal quartic marked-tree sum,
uniformly in the support and all spins.
\end{theorem}

\begin{proof}
The one-sided insertion theorem gives
\begin{equation}
 \mathfrak C_{\tau,E}
 \le
 Cs_\alpha^3L_\tau
 \frac{Nz}{1-z}\kappa_\otimes(P_E).
\label{eq:V22-one-sided-protected}
\end{equation}

First suppose
\[
 B_\tau:=\prod_r\Xi_r^{d_r-1}\le DY^2.
\]
Then \(Y>0\).  Since \(z\le q_P\),
\[
 \frac{z}{1-z}\le\frac{q_P}{1-q_P}.
\]
The \(m=2\) retained-failure product estimate on the private bank gives
\[
 \frac{(s_\alpha Y)^2q_P}{1-q_P}\le C.
\]
Use this, \eqref{eq:V22-count-capacity}, and
\eqref{eq:V22-one-sided-protected} to obtain
\[
 \mathfrak C_{\tau,E}
 \le\frac{Cs_\alpha}{Y^2}L_\tau
 \le\frac{C_D}{B_\tau}L_\tau
 \le C_D\mathcal W_{\tau,e}.
\]
The last inequality uses \(s_\alpha\le1\) and
\(L_\tau\le\prod_ra_r^{d_r}\).

Now suppose \(B_\tau>DY^2\).  Choose a globally maximal row \(i\).
The scalar proof of
\Cref{lem:V19-complete-bank-selection} gives
\[
 h_i>\Xi_i/2.
\]
Here all shared coordinates belong to \(E\), and
\[
 h_i=\sum_{\substack{f\in E\\i\in f}}A_f\le S.
\]
Thus
\begin{equation}
 S>\Xi_{\max}/2.
\label{eq:V22-S-dominates}
\end{equation}
Because the global output is nontrivial outside the protected shared
bank, at least one factor of \(z\) is nontrivial.  The uniform
nontrivial-output separation gives
\[
 \frac{z}{1-z}\le\frac{C}{s_\alpha}.
\]
Equations
\eqref{eq:V22-weighted-network},
\eqref{eq:V22-S-dominates}, and
\eqref{eq:V22-one-sided-protected} imply
\begin{equation}
 \mathfrak C_{\tau,E}
 \le
 Cs_\alpha^2L_\tau
 \frac{H}{S^2}
 \le
 Cs_\alpha^2L_\tau
 \frac{H}{\Xi_{\max}^2}.
\label{eq:V22-failed-bound}
\end{equation}

If \(\tau\) is a path, then \(L_\tau\le A_0\).  The same conclusion
holds for a star satisfying
\(s_\alpha^2L_\tau\le A_0\).  In both cases,
\eqref{eq:V22-failed-bound} is bounded by
\(\mathfrak P_E\) through
\Cref{lem:V22-global-routing}.

It remains to consider a star with
\(s_\alpha^2L_\tau>A_0\).  Return upstream to the complete positive
local fourth cumulant and use
\Cref{lem:V19-complete-cumulant-cap} instead of
\eqref{eq:V22-local-tree}.  Its norm is bounded by a constant.  The
protected contribution is then at most
\[
 \frac{C}{s_\alpha}N\kappa_\otimes(P_E)
 \le
 \frac{CH}{s_\alpha S^2}
 \le
 \frac{CH}{s_\alpha\Xi_{\max}^2}.
\]
The hot-star inequality implies \(A_0>s_\alpha^{-4}\), and hence
\(s_\alpha^{-1}\le A_0\).  Apply
\eqref{eq:V22-global-routing} once more.  The path and the two star
branches exhaust the sixteen local tree fibres.
\end{proof}

\begin{corollary}[The singular shared denominator is eliminated]
\label{cor:V22-no-singular-denominator}
On the globally all-singlet incidence-two output, the factor
\(q_S/(1-q_S)\) is never introduced.  The exact external carrier
\[
 \frac{Nz}{1-z}P_E
\]
is finite and obeys the quartic marked-tree atlas.
\end{corollary}

\begin{proof}
This is precisely
\Cref{thm:V22-weighted-network,thm:V22-protected-line}.
\end{proof}

\begin{assessment}[A genuine output-priority obstruction is closed]
\label{ass:V22-protected-closed}
The all-singlet line was the only shared incidence-two output with
internal multiplier exactly one.  V22 closes that line globally,
including overlapping pair banks, by a vertex-cut capacity rather
than an edgewise product.  The external resolvent is paid either by
the private bank in the private-dominant branch or by the uniform
nontrivial gap outside the shared bank in the failed-branching branch.
\end{assessment}

\section{What remains in the mixed gapped carrier}
\label{sec:V22-mixed-residual}

When at least one shared coordinate has a nontrivial output, the
shared multiplier is strictly below one and V21 removes the output
tuple count inside each complete pair bank.  A mixed tuple can still
place singlets on a large equal-spin subgraph while its nontrivial
output lies in a different bank.  Its baseline output mass then
couples the capacity of that singlet subgraph to the resolvent of the
gapped bank.  The full payment requires a two-bank version of
\eqref{eq:V22-weighted-network} with one retained failure.  Treating
the two factors by independent capacities would repeat the V18 error.

\begin{openproblem}[V23 mandate: one-failure EPR-network resolvent]
\label{op:V23-mandate}
\begin{enumerate}
\item Expand the complete shared moment by the set of nontrivial
      output coordinates, but keep each set as one positive operator.
\item For a fixed nonempty failure set, apply vertex-cut Schmidt bounds
      to the complementary singlet network.
\item Sum the failure sets with their Wilson multipliers before taking
      capacity; do not sum resolved irreducible copies.
\item Prove a support-uniform bound for output baseline mass times the
      one-failure resolvent.
\item Combine that bound with
      \Cref{thm:V20-comparable-bank,thm:V20-unbalanced-bank} to route
      the dominant row to legal marks.
\item Then return to the junction-payer term containing
      \(s_\alpha A_{U_e}\).
\end{enumerate}
\end{openproblem}

\begin{firewall}[Claim boundary after compact V22]
\label{fire:V22-current-boundary}
V22 closes the globally all-singlet incidence-two output-priority
line.  It does not yet close mixed singlet/nontrivial output tuples,
the cancelling junction-input mark, the full incidence-two
output-priority carrier, local-dominant non-Cartan blocks,
incidence-three or incidence-four shared coordinates, general compact
simple groups, the full embedded \([4]\) sector, multiple junctions,
normalized collision aggregation, \([3,3]\), global quartic or
all-order MTA, continuum Yang--Mills existence, or a positive mass
gap.
\end{firewall}

\section{Append-only record for compact V22}
\label{sec:V22-record}

V22 was appended to the complete compact V21 file.  The exact V21
parent had \(8412\) newline-delimited source records, \(281756\) bytes,
and SHA--256 digest
\[
\texttt{c7e6ff1152510e6a61b8c131cfb9224b8f3130d265a6e7ea7c8378c5141b5c5a}.
\]
No compact V21 mathematical content was changed.  The sole final
\verb|\end{document}| is restored below.

% =============================================================================
% END APPENDED COMPACT VERSION V22 MATERIAL
% =============================================================================

% The compact V22 document terminator is intentionally disabled here:
% \end{document}

% =============================================================================
% BEGIN APPENDED COMPACT VERSION V23 MATERIAL
% =============================================================================

\chapter{V23: One-failure EPR-network resolvents}
\label{ch:V23}

\section{Balanced incidence-two graphs}
\label{sec:V23-balanced-graphs}

Let \(E\) be the complete set of off-junction incidence-two shared
coordinates, \(N:=|E|\ge1\).  At \(f=uv\), the two nontrivial
\(SU(2)\) input spins have masses
\[
 A_{u,f}:=1+C_2(p_{u,f}),
 \qquad
 A_{v,f}:=1+C_2(p_{v,f}).
\]
Call the coordinate \emph{balanced} when
\begin{equation}
 \min(A_{u,f},A_{v,f})
 \ge\frac14\max(A_{u,f},A_{v,f}).
\label{eq:V23-balanced-edge}
\end{equation}
Equal-spin coordinates are included.  Throughout V23 every coordinate
is balanced.  Put
\begin{align}
 w_f&:=\max(A_{u,f},A_{v,f}),&
 S&:=\sum_{f\in E}w_f,&
 H&:=\sum_{f=uv\in E}A_{u,f}A_{v,f}.
\label{eq:V23-balanced-data}
\end{align}
Then
\begin{equation}
\boxed{S^2\le4NH.}
\label{eq:V23-SH}
\end{equation}
Indeed, \eqref{eq:V23-balanced-edge} gives
\(w_f^2\le4A_{u,f}A_{v,f}\), followed by Cauchy--Schwarz.

For \(t>0\), let
\[
 M_f^{(t)}:=\sum_{J}
 q_{\alpha,J}^{\,t}P_{J,f},
 \qquad
 M_E^{(t)}:=\bigotimes_{f\in E}M_f^{(t)}.
\]
At an equal-spin coordinate let \(P_{0,f}\) be the singlet
projection.  If every coordinate is equal-spin, set
\[
 P_*:=\bigotimes_{f\in E}P_{0,f};
\]
if at least one coordinate has unequal spins, set \(P_*:=0\).
Thus \(P_*\) is exactly the eigenspace on which the complete shared
multiplier equals one.

\section{A graph-level power-capacity lemma}
\label{sec:V23-power-capacity}

Fix
\[
 \rho:=2^{-1/2}<1.
\]
For \(0\le r\le1\), define
\begin{align}
 \Phi_N^{(0)}(r)
 &:=
 [\rho+(1-\rho)r]^N-[\rho(1-r)]^N,
\label{eq:V23-Phi-protected}\\
 \Phi_N^{(1)}(r)
 &:=
 r[\rho+(1-\rho)r]^{N-1}.
\label{eq:V23-Phi-failed}
\end{align}

\begin{lemma}[Power capacity with zero or one retained failure]
\label{lem:V23-power-capacity}
Let
\[
 r_t:=(1-\delta_\alpha)^t,
 \qquad
 \delta_\alpha=c_*s_\alpha
\]
with the V19 separation constant.  Then
\begin{equation}
\boxed{
 \kappa_\otimes(M_E^{(t)}-P_*)
 \le
 \begin{cases}
 \Phi_N^{(0)}(r_t),&P_*\ne0,\\
 \Phi_N^{(1)}(r_t),&P_*=0.
 \end{cases}}
\label{eq:V23-power-capacity}
\end{equation}
\end{lemma}

\begin{proof}
At an equal-spin coordinate,
\[
 M_f^{(t)}
 \preccurlyeq r_tI+(1-r_t)P_{0,f};
\]
at an unequal-spin coordinate,
\[
 M_f^{(t)}\preccurlyeq r_tI
\]
by \Cref{lem:V20-no-trivial}.

First suppose every edge has equal spins.  Expand the tensor product
of the first majorants and subtract \(P_*\), exactly as in
\eqref{eq:V21-power-majorant}.  For a proper subset
\(A\subsetneq E\), the resulting positive projection is bounded by
the projection imposing singlets on the subgraph \(A\).  The vertex-cut
argument of \Cref{lem:V22-vertex-cut} gives
\[
 \kappa_\otimes\!\left[
 \left(\bigotimes_{f\in A}P_{0,f}\right)
 \otimes I_{E\setminus A}\right]
 \le
 \left(\prod_{f\in A}d_f\right)^{-1/2}
 \le\rho^{|A|}.
\]
Summing the expansion gives
\(\Phi_N^{(0)}(r_t)\).

Now suppose \(K\ge1\) coordinates have unequal spins.  Their
majorants contribute \(r_t^K\).  Expanding the remaining
equal-spin coordinates and applying the same cut bound gives
\[
 \kappa_\otimes(M_E^{(t)})
 \le
 r_t^K[\rho+(1-\rho)r_t]^{N-K}.
\]
Since \(r_t\le\rho+(1-\rho)r_t\) and \(K\ge1\), the right side is at
most \(\Phi_N^{(1)}(r_t)\).
\end{proof}

\begin{lemma}[The analytic one-failure sum]
\label{lem:V23-analytic-sum}
For either \(\nu=0\) or \(\nu=1\),
\begin{align}
 \int_0^1\frac{\Phi_N^{(\nu)}(r)}r\,dr
 &\le\frac{C_\rho}{N},
\label{eq:V23-integral-bound}\\
 \sum_{n\ge1}\Phi_N^{(\nu)}
 \bigl((1-\delta_\alpha)^n\bigr)
 &\le
 C e^{-cs_\alpha N}
 \left(1+\frac1{s_\alpha N}\right).
\label{eq:V23-integer-sum}\\
 \sum_{n\ge0}\Phi_N^{(\nu)}
 \bigl((1-\delta_\alpha)^{n+1/2}\bigr)
 &\le
 C e^{-cs_\alpha N}
 \left(1+\frac1{s_\alpha N}\right).
\label{eq:V23-half-sum}
\end{align}
In particular, the sum in \eqref{eq:V23-integer-sum} is at most
\(C/(s_\alpha N)\).
\end{lemma}

\begin{proof}
Put \(A(r):=\rho+(1-\rho)r\) and
\(B(r):=\rho(1-r)\).  For \(\nu=1\),
\[
 \int_0^1\frac{\Phi_N^{(1)}(r)}r\,dr
 =
 \int_0^1A(r)^{N-1}\,dr
 \le\frac{1}{(1-\rho)N}.
\]
For \(\nu=0\), \(A(r)-B(r)=r\).  On \(0<r\le1/2\), the mean-value
theorem gives
\[
 \frac{A(r)^N-B(r)^N}{r}
 \le N A(1/2)^{N-1}.
\]
Its integral is \(O(N\theta^N)\), with
\(\theta:=A(1/2)<1\), and hence is \(O(N^{-1})\).  On
\(1/2\le r\le1\), use \(r^{-1}\le2\) and integrate \(2A(r)^N\).
This proves \eqref{eq:V23-integral-bound}.

Let
\(\lambda_\alpha:=-\log(1-\delta_\alpha)\), so
\(\lambda_\alpha\ge\delta_\alpha\).  Both \(\Phi\)'s are increasing.
Geometric interval comparison gives
\[
 \sum_{n\ge1}\Phi(e^{-\lambda_\alpha n})
 \le
 \Phi(e^{-\lambda_\alpha})
 +\lambda_\alpha^{-1}
 \int_0^{e^{-\lambda_\alpha}}\frac{\Phi(r)}r\,dr.
\]
Decrease \(c_*\), if needed, so that
\(\delta_\alpha\le1/2\).  Repeating the preceding split with upper
endpoint \(e^{-\lambda_\alpha}\) shows
\[
 \Phi(e^{-\lambda_\alpha})
 \le Ce^{-c\delta_\alpha N},
 \qquad
 \int_0^{e^{-\lambda_\alpha}}\frac{\Phi(r)}r\,dr
 \le\frac{C}{N}e^{-c\delta_\alpha N}.
\]
This proves \eqref{eq:V23-integer-sum}.  Starting the same geometric
comparison at \(e^{-\lambda_\alpha/2}\) proves
\eqref{eq:V23-half-sum}.  Finally,
\(\sup_{x>0}e^{-cx}(1+x^{-1})x<\infty\), giving the last assertion.
\end{proof}

\section{The complete graph resolvent}
\label{sec:V23-complete-resolvent}

For a shared output tuple \(\boldsymbol J\), put
\[
 q_{\boldsymbol J}:=\prod_{f\in E}q_{\alpha,J_f},
 \qquad
 Z_{\boldsymbol J}:=\sum_{f\in E}C_2(J_f),
 \qquad
 H_{\boldsymbol J}:=N+Z_{\boldsymbol J}.
\]
On the complement of \(P_*\), define
\begin{align}
 \mathcal G_E
 &:=
 \sum_{q_{\boldsymbol J}<1}
 \frac{q_{\boldsymbol J}}{1-q_{\boldsymbol J}}P_{\boldsymbol J},
\label{eq:V23-graph-geometric}\\
 \mathcal C_E
 &:=
 \sum_{q_{\boldsymbol J}<1}
 \frac{Z_{\boldsymbol J}q_{\boldsymbol J}}
      {1-q_{\boldsymbol J}}P_{\boldsymbol J},
\label{eq:V23-graph-Casimir}\\
 \mathcal R_E&:=N\mathcal G_E+\mathcal C_E.
\label{eq:V23-graph-resolvent}
\end{align}

\begin{theorem}[One-failure complete graph estimates]
\label{thm:V23-graph-resolvent}
There are uniform constants \(c,C>0\) such that
\begin{align}
 \kappa_\otimes(\mathcal G_E)
 &\le
 Ce^{-cs_\alpha N}
 \left(1+\frac1{s_\alpha N}\right)
 \le\frac{C}{s_\alpha N},
\label{eq:V23-geometric-bounds}\\
 \kappa_\otimes(\mathcal C_E)
 &\le
 \min\left\{
 \frac{C}{s_\alpha},
 \frac{C}{s_\alpha}e^{-cs_\alpha N}
 \left(1+\frac1{s_\alpha N}\right)
 \right\},
\label{eq:V23-Casimir-bounds}\\
 \boxed{\kappa_\otimes(\mathcal R_E)\le C/s_\alpha,}
\label{eq:V23-resolvent-unmarked}\\
 \boxed{
 s_\alpha^2S^2\kappa_\otimes(\mathcal R_E)\le CH.}
\label{eq:V23-resolvent-marked}
\end{align}
\end{theorem}

\begin{proof}
Geometric summation gives
\[
 \mathcal G_E=\sum_{n\ge1}(M_E^{(n)}-P_*).
\]
Capacity subadditivity,
\Cref{lem:V23-power-capacity,lem:V23-analytic-sum}, prove
\eqref{eq:V23-geometric-bounds}.

For a nonprotected output tuple, apply the \(m=2\)
retained-failure compiler to
\[
 x_f=s_\alpha C_2(J_f),
 \qquad
 r_f=q_{\alpha,J_f}.
\]
Writing \(X=s_\alpha Z_{\boldsymbol J}\) and
\(Q=q_{\boldsymbol J}\), it gives
\[
 X^2Q\le C(1-Q).
\]
Therefore
\[
 \frac{Z_{\boldsymbol J}Q}{1-Q}
 \le
 \frac{C}{s_\alpha}\sqrt{\frac{Q}{1-Q}}
 \le
 \frac{C}{s_\alpha}
 \left(\sum_{n\ge0}Q^{n+1/2}\right).
\]
The last inequality uses
\(\sqrt{Q/(1-Q)}\le\sqrt Q/(1-Q)\).
Consequently, in Loewner order,
\[
 \mathcal C_E
 \preccurlyeq
 \frac{C}{s_\alpha}
 \sum_{n\ge0}(M_E^{(n+1/2)}-P_*).
\]
The half-integer estimate proves the second entry in the minimum in
\eqref{eq:V23-Casimir-bounds}.  The \(m=1\) retained-failure estimate
also gives directly
\(\mathcal C_E\preccurlyeq C s_\alpha^{-1}(I-P_*)\), proving the first
entry.

Now use
\(\mathcal R_E=N\mathcal G_E+\mathcal C_E\).
The unmarked estimate follows from
\eqref{eq:V23-geometric-bounds} and the first part of
\eqref{eq:V23-Casimir-bounds}.

For the marked estimate, recall \(S^2\le4NH\), and put
\(x=s_\alpha N\).  If \(x\le1\), use
\[
 \kappa_\otimes(\mathcal G_E)\le C/(s_\alpha N),
 \qquad
 \kappa_\otimes(\mathcal C_E)\le C/s_\alpha.
\]
The two contributions to
\(s_\alpha^2S^2\kappa_\otimes(\mathcal R_E)\)
are bounded respectively by
\[
 Cs_\alpha NH=CxH,
\qquad
 Cs_\alpha NH=CxH.
\]
If \(x>1\), use the exponentially decaying entries in
\eqref{eq:V23-geometric-bounds}--%
\eqref{eq:V23-Casimir-bounds}.  The two contributions are bounded by
\[
 Cx^2e^{-cx}(1+x^{-1})H,
\qquad
 Cx e^{-cx}(1+x^{-1})H.
\]
All four displayed scalar factors are uniformly bounded.  This proves
\eqref{eq:V23-resolvent-marked}.
\end{proof}

\section{Insertion into output priority}
\label{sec:V23-output-priority}

For \(f=uv\), let \(\sigma_{uv}\) be the comparison path with its
\(uv\)-edge marked at \(f\).  The balanced-graph legal sum is
\begin{equation}
 \mathfrak P_E
 :=
 A_0\sum_{f=uv\in E}
 \frac{A_{u,f}A_{v,f}}{\Xi_u\Xi_v}
 \ge
 \frac{A_0H}{\Xi_{\max}^2}.
\label{eq:V23-balanced-routing}
\end{equation}

\begin{theorem}[Balanced gapped shared-payer closure in output
priority]
\label{thm:V23-balanced-output}
Fix \(D\ge16\).  Suppose:
\begin{enumerate}[label=\textup{(\roman*)},leftmargin=4em]
\item every shared coordinate has incidence two and satisfies
      \eqref{eq:V23-balanced-edge};
\item \(a_r<\Xi_r/4\) for all rows;
\item \(\Xi_T>DY\); and
\item the globally protected line, when present, is treated by V22.
\end{enumerate}
Then the complete gapped shared-output-resolvent payer obeys the legal
quartic marked-tree atlas.
\end{theorem}

\begin{proof}
The complete gapped output sum is exactly
\(\mathcal R_E\); all junction and private factors not assigned to
this payer are positive contractions.  The one-sided insertion
theorem and the local tree bound give
\begin{equation}
 \textup{gapped shared payer}
 \le
 Cs_\alpha^3q_PL_\tau
 \kappa_\otimes(\mathcal R_E).
\label{eq:V23-gapped-insertion}
\end{equation}

If \(B_\tau\le DY^2\), use
\eqref{eq:V23-resolvent-unmarked} and the private quadratic debit:
\[
 \eqref{eq:V23-gapped-insertion}
 \le Cs_\alpha^2q_PL_\tau
 \le C_D\frac{L_\tau}{B_\tau}
 \le C_D\mathcal W_{\tau,e}.
\]

Suppose \(B_\tau>DY^2\).  Choose a globally maximal row \(i\).  As in
\Cref{lem:V19-complete-bank-selection},
\[
 h_i>\Xi_i/2.
\]
Every incident input mass is at most the corresponding \(w_f\), so
\[
 S\ge h_i>\Xi_{\max}/2.
\]
The marked graph-resolvent estimate gives
\[
 \eqref{eq:V23-gapped-insertion}
 \le
 Cs_\alpha L_\tau\frac{H}{S^2}
\le
 Cs_\alpha L_\tau\frac{H}{\Xi_{\max}^2}.
\]
For paths, \(L_\tau\le A_0\).  For stars satisfying
\(s_\alpha L_\tau\le A_0\), the same displayed estimate is bounded by
\eqref{eq:V23-balanced-routing}.

For a hot star, return to the complete positive local fourth cumulant.
The marked estimate itself gives
\[
 \kappa_\otimes(\mathcal R_E)
 \le\frac{CH}{s_\alpha^2S^2}
 \le\frac{CH}{s_\alpha^2\Xi_{\max}^2}.
\]
For a star centered at \(c\),
\[
\frac{L_\tau}{A_0}
=
\left(\frac{a_c}{\prod_{r\ne c}a_r}\right)^{1/2}.
\]
The hot inequality \(s_\alpha L_\tau>A_0\) therefore implies
\(s_\alpha^2a_c>\prod_{r\ne c}a_r\), and hence
\(A_0>s_\alpha^{-2}\).  Thus
\eqref{eq:V23-balanced-routing} pays this final branch.
\end{proof}

\begin{assessment}[The mixed output-count problem is closed on the
balanced graph]
\label{ass:V23-balanced-closed}
V23 sums the nonempty Wilson-failure sets inside the complete graph
moment.  Vertex-cut capacity is applied only to the complementary
singlet subgraph, and the failure sets are then summed analytically.
This proves a marked complete-resolvent estimate without multiplying
overlapping pair capacities or counting resolved outputs.
\end{assessment}

\begin{firewall}[Two output-priority mechanisms remain]
\label{fire:V23-remaining}
\Cref{thm:V23-balanced-output} controls the shared-output payer, not
the junction-payer term proportional to \(s_\alpha A_{U_e}\).
It also assumes every unequal-spin edge is balanced in the sense of
\eqref{eq:V23-balanced-edge}.  Strongly unbalanced edges must use
their retained output Casimir as in
\Cref{thm:V20-unbalanced-bank}; replacing their legal edge mass by
\(w_f^2\) would be false.
\end{firewall}

\begin{openproblem}[V24 mandate: unbalanced failures and the junction
mass]
\label{op:V24-mandate}
\begin{enumerate}
\item Split the incidence-two graph into balanced and strongly
      unbalanced edges before resolving outputs.
\item Use the unbalanced output-retention estimate to pay the larger
      endpoint mass through the complete graph resolvent.
\item If the balanced graph dominates the maximal row, apply V23.
\item If the unbalanced graph dominates it, construct the direct
      \(h_i^{-2}\) payer without inserting an artificial bilinear mark.
\item Revisit the junction-payer factor
      \(s_\alpha A_{U_e}\) only after this dichotomy is complete.
\item Close the full shared-output payer for incidence two, or isolate
      the exact mixed balanced/unbalanced operator inequality.
\end{enumerate}
\end{openproblem}

\begin{firewall}[Claim boundary after compact V23]
\label{fire:V23-current-boundary}
V23 closes the complete balanced incidence-two shared-output payer in
output priority.  It does not yet close strongly unbalanced mixed
graphs, the cancelling junction-input mark, the full incidence-two
output-priority carrier, local-dominant non-Cartan blocks,
incidence-three or incidence-four shared coordinates, general compact
simple groups, the full embedded \([4]\) sector, multiple junctions,
normalized collision aggregation, \([3,3]\), global quartic or
all-order MTA, continuum Yang--Mills existence, or a positive mass
gap.
\end{firewall}

\section{Append-only record for compact V23}
\label{sec:V23-record}

V23 was appended to the complete compact V22 file.  The exact V22
parent had \(8799\) newline-delimited source records, \(294107\) bytes,
and SHA--256 digest
\[
\texttt{14bfa8a04735efc792c047acb3ac673ed62217e95e6f4fd2f1d4b2ff7cce723f}.
\]
No compact V22 mathematical content was changed.  The sole final
\verb|\end{document}| is restored below.

% =============================================================================
% END APPENDED COMPACT VERSION V23 MATERIAL
% =============================================================================

% The compact V23 document terminator is intentionally disabled here:
% \end{document}

% =============================================================================
% BEGIN APPENDED COMPACT VERSION V24 MATERIAL
% =============================================================================

\chapter{V24: Strongly unbalanced failures and the complete
incidence-two shared payer}
\label{ch:V24}

\section{The strongly unbalanced subgraph}
\label{sec:V24-unbalanced}

Let \(U\ne\varnothing\) be a set of incidence-two \(SU(2)\)
coordinates.  At \(f=uv\), put
\[
 w_f:=\max(A_{u,f},A_{v,f}),
 \qquad
 \min(A_{u,f},A_{v,f})<\frac14w_f,
 \qquad
 W:=\sum_{f\in U}w_f.
\]
For every output \(J_f\) at \(f\), put
\[
 b_{J_f}:=1+C_2(J_f),
 \qquad
 q_f:=q_{\alpha,J_f}.
\]

\begin{lemma}[Every unbalanced output is comparable to the large
input]
\label{lem:V24-output-comparison}
There are absolute constants \(c,C>0\) such that
\begin{equation}
\boxed{
 cw_f\le b_{J_f}\le Cw_f}
\label{eq:V24-output-comparison}
\end{equation}
for every allowed output.
\end{lemma}

\begin{proof}
Orient the edge so that row \(u\) carries the larger spin.  The lower
bound is \Cref{lem:V20-unbalanced-retention}.  The upper bound is the
fusion-Casimir upper estimate
\eqref{eq:V3-fusion-output}:
\[
 b_{J_f}\le C(A_{u,f}+A_{v,f})\le Cw_f.
\]
\end{proof}

Let
\[
 M_U:=\bigotimes_{f\in U}\sum_{J_f}q_fP_{J_f}.
\]
For an output tuple \(\boldsymbol J\), write
\[
 Q_{\boldsymbol J}:=\prod_fq_f,
 \qquad
 B_{\boldsymbol J}:=\sum_fb_{J_f}.
\]
Every \(Q_{\boldsymbol J}<1\).  Define
\begin{align}
 \mathcal G_U
 &:=
 \sum_{\boldsymbol J}
 \frac{Q_{\boldsymbol J}}{1-Q_{\boldsymbol J}}
 P_{\boldsymbol J},
\label{eq:V24-unbalanced-geometric}\\
 \mathcal R_U
 &:=
 \sum_{\boldsymbol J}
 \frac{B_{\boldsymbol J}Q_{\boldsymbol J}}
      {1-Q_{\boldsymbol J}}
 P_{\boldsymbol J}.
\label{eq:V24-unbalanced-resolvent}
\end{align}

\begin{theorem}[Direct unbalanced moments and resolvents]
\label{thm:V24-unbalanced-resolvent}
Uniformly in \(U\), its spins, and all complete output tuples,
\begin{align}
 (s_\alpha W)^k\|M_U\|&\le C_k,
 &&1\le k\le3,
\label{eq:V24-unbalanced-moment}\\
 \|\mathcal G_U\|&\le
 \frac{C_k}{(s_\alpha W)^k},
 &&1\le k\le3,
\label{eq:V24-unbalanced-geometric-bounds}\\
 \boxed{
 \|\mathcal R_U\|
 \le
 \min\left\{
 \frac{C}{s_\alpha},
 \frac{C}{s_\alpha^3W^2}
 \right\}.}
\label{eq:V24-unbalanced-resolvent-bounds}
\end{align}
\end{theorem}

\begin{proof}
For a fixed tuple, put \(x_f=s_\alpha w_f\).  By
\Cref{lem:V24-output-comparison} and the one-link Wilson debits,
\[
 x_f^kq_f\le C_k(1-q_f),
 \qquad 1\le k\le3.
\]
The ordinary product compiler gives
\[
 (s_\alpha W)^kQ_{\boldsymbol J}\le C_k.
\]
Taking the maximum tuple proves
\eqref{eq:V24-unbalanced-moment}.

The retained-failure compiler gives
\[
 (s_\alpha W)^kQ_{\boldsymbol J}
 \le C_k(1-Q_{\boldsymbol J}).
\]
Divide by \(1-Q_{\boldsymbol J}\) and maximize to prove
\eqref{eq:V24-unbalanced-geometric-bounds}.  Finally
\Cref{lem:V24-output-comparison} gives
\[
 cW\le B_{\boldsymbol J}\le CW.
\]
Multiply the \(k=1\) and \(k=3\) cases of
\eqref{eq:V24-unbalanced-geometric-bounds} by \(CW\).
\end{proof}

\section{A first marked moment for the balanced subgraph}
\label{sec:V24-balanced-moment}

Let \(B\ne\varnothing\) be a balanced subgraph in the sense of V23.
Use \(N_B,S_B,H_B\) from
\eqref{eq:V23-balanced-data}, restricted to \(B\), and let
\[
 M_B:=M_B^{(1)}.
\]
If every edge of \(B\) has equal spins, let \(P_B\) be its global
singlet network projection; otherwise put \(P_B=0\).

\begin{lemma}[Complete balanced-moment first mark]
\label{lem:V24-balanced-first-mark}
Uniformly in the balanced graph,
\begin{equation}
\boxed{
 s_\alpha S_B^2\kappa_\otimes(M_B)\le C H_B.}
\label{eq:V24-balanced-first-mark}
\end{equation}
Moreover, when \(P_B\ne0\),
\begin{equation}
\boxed{
 N_BS_B^2\kappa_\otimes(P_B)\le CH_B.}
\label{eq:V24-protected-balanced-mark}
\end{equation}
\end{lemma}

\begin{proof}
The second estimate is the proof of
\Cref{thm:V22-weighted-network}; the same proof applies to the
balanced subgraph because a protected line forces every one of its
edges to have equal spins.

For the first estimate, split
\[
 M_B=(M_B-P_B)+P_B.
\]
By \Cref{lem:V23-power-capacity} at \(t=1\) and the endpoint estimate
in the proof of \Cref{lem:V23-analytic-sum},
\[
 \kappa_\otimes(M_B-P_B)
 \le C\min\{1,e^{-cs_\alpha N_B}\}.
\]
Since \(S_B^2\le4N_BH_B\),
\[
 s_\alpha S_B^2\kappa_\otimes(M_B-P_B)
 \le
 C(s_\alpha N_B)
 \min\{1,e^{-cs_\alpha N_B}\}H_B
 \le CH_B.
\]
The protected part follows from
\eqref{eq:V24-protected-balanced-mark} and
\(s_\alpha/N_B\le1\).
\end{proof}

\section{The exact mixed resolvent split}
\label{sec:V24-mixed-split}

Assume both \(B\) and \(U\) are nonempty.  On the complete gapped
output of \(B\cup U\), let \(\mathcal R_{B\cup U}\) denote the shared
output-mass resolvent.  Use
\(\mathcal R_B\) for the V23 balanced resolvent, with its protected
line removed.

\begin{lemma}[Balanced/unbalanced positive-operator decomposition]
\label{lem:V24-mixed-decomposition}
In Loewner order,
\begin{equation}
\boxed{
 \mathcal R_{B\cup U}
 \preccurlyeq
 \mathcal R_B\otimes M_U
 +M_B\otimes\mathcal R_U
 +N_BP_B\otimes\mathcal G_U.}
\label{eq:V24-mixed-decomposition}
\end{equation}
The last term is zero when \(P_B=0\).
\end{lemma}

\begin{proof}
For a resolved pair of tuples write
\[
 Q=Q_BQ_U,
 \qquad
 H_{\rm out}=H_B^{\rm out}+H_U^{\rm out}.
\]
If \(Q_B<1\), then
\[
 \frac{H_B^{\rm out}Q_BQ_U}{1-Q_BQ_U}
 \le
 \frac{H_B^{\rm out}Q_B}{1-Q_B}\,Q_U.
\]
If \(Q_B=1\), its tuple is the protected balanced line and
\(H_B^{\rm out}=N_B\), producing the third term in
\eqref{eq:V24-mixed-decomposition}.  Since \(Q_U<1\),
\[
 \frac{H_U^{\rm out}Q_BQ_U}{1-Q_BQ_U}
 \le
 Q_B\frac{H_U^{\rm out}Q_U}{1-Q_U}.
\]
These scalar inequalities hold on every common orthogonal spectral
block and prove the operator statement.
\end{proof}

\begin{corollary}[Three valid one-sided insertions]
\label{cor:V24-three-insertions}
\begin{align}
 \kappa_\otimes(\mathcal R_B\otimes M_U)
 &\le
 \kappa_\otimes(\mathcal R_B)\|M_U\|,
\label{eq:V24-insertion-one}\\
 \kappa_\otimes(M_B\otimes\mathcal R_U)
 &\le
 \kappa_\otimes(M_B)\|\mathcal R_U\|,
\label{eq:V24-insertion-two}\\
 \kappa_\otimes(P_B\otimes\mathcal G_U)
 &\le
 \kappa_\otimes(P_B)\|\mathcal G_U\|.
\label{eq:V24-insertion-three}
\end{align}
\end{corollary}

\begin{proof}
Apply the V2 one-sided tensor theorem in the displayed orientation.
No product of capacities is used.
\end{proof}

\section{Closure of the complete incidence-two shared payer}
\label{sec:V24-shared-closure}

\begin{theorem}[All \(SU(2)\) incidence-two shared-output payers close]
\label{thm:V24-full-shared-payer}
Fix \(D\ge16\).  Suppose every shared coordinate has incidence two,
\[
 a_r<\Xi_r/4\quad(1\le r\le4),
 \qquad
 \Xi_T>DY.
\]
Then the complete shared-output-resolvent payer, including every
balanced and strongly unbalanced coordinate and every mixed output
tuple, obeys the legal quartic marked-tree atlas.  The globally
protected line is understood in the exact V22 sense.
\end{theorem}

\begin{proof}
Split the shared graph into \(B\) and \(U\).  If \(U=\varnothing\),
use \Cref{thm:V23-balanced-output} and
\Cref{thm:V22-protected-line}.  If \(B=\varnothing\), use
\Cref{thm:V24-unbalanced-resolvent} directly.  It remains to treat the
mixed case through
\eqref{eq:V24-mixed-decomposition}.

First suppose \(B_\tau\le DY^2\).  Equations
\eqref{eq:V24-insertion-one}--%
\eqref{eq:V24-insertion-three},
\eqref{eq:V23-resolvent-unmarked},
\eqref{eq:V24-unbalanced-resolvent-bounds}, and
\eqref{eq:V22-count-capacity} give
\[
 \kappa_\otimes(\mathcal R_{B\cup U})\le C/s_\alpha.
\]
Indeed, \(W\ge1\), so the third term is bounded using
\(\|\mathcal G_U\|\le C/(s_\alpha W)\).
After insertion of the local tree carrier, the private quadratic debit
gives
\[
 Cs_\alpha^2q_PL_\tau
 \le C_D\frac{L_\tau}{B_\tau}
 \le C_D\mathcal W_{\tau,e}.
\]

Now suppose \(B_\tau>DY^2\).  Choose a globally maximal row \(i\).
The usual failed-branching calculation gives
\[
 h_i>\Xi_i/2.
\]
Every incident balanced input is at most its \(w_f\), and the same is
true on \(U\).  Hence
\[
 S_B+W>\Xi_{\max}/2.
\label{eq:V24-dominant-split}
\]

Assume first that \(S_B>\Xi_{\max}/4\).  The three terms in
\eqref{eq:V24-mixed-decomposition}, after multiplication by the local
factor \(s_\alpha^3L_\tau\), are bounded respectively by
\[
 Cs_\alpha L_\tau\frac{H_B}{S_B^2},
 \qquad
 Cs_\alpha L_\tau\frac{H_B}{S_B^2},
 \qquad
 Cs_\alpha^2L_\tau\frac{H_B}{S_B^2}.
\]
Here use, in order,
\eqref{eq:V23-resolvent-marked},
\eqref{eq:V24-balanced-first-mark} together with
\(\|\mathcal R_U\|\le C/s_\alpha\), and
\eqref{eq:V24-protected-balanced-mark} together with
\(\|\mathcal G_U\|\le C/(s_\alpha W)\).
Thus paths and stars satisfying
\(s_\alpha L_\tau\le A_0\) are paid by the balanced legal path sum
\[
 A_0\sum_{f=uv\in B}
 \frac{A_{u,f}A_{v,f}}{\Xi_u\Xi_v}
 \ge
 \frac{A_0H_B}{\Xi_{\max}^2}.
\]
If a star is hot, return to the complete local fourth cumulant.  The
same three one-sided estimates without \(s_\alpha^3L_\tau\) are at
most
\[
 \frac{CH_B}{s_\alpha^2S_B^2}.
\]
As in V23, the hot inequality implies
\(A_0>s_\alpha^{-2}\), so the same legal path sum pays the complete
hot-star carrier.

Finally assume \(W\ge\Xi_{\max}/4\).  The first two terms of
\eqref{eq:V24-mixed-decomposition}, after local insertion, obey
\[
 Cs_\alpha^3L_\tau
 \frac{C}{s_\alpha}
 \frac{C}{s_\alpha^2W^2}
 \le\frac{CL_\tau}{W^2},
\]
and
\[
 Cs_\alpha^3L_\tau
 \frac{C}{s_\alpha^3W^2}
 \le\frac{CL_\tau}{W^2},
\]
respectively.  For the third term use
\(\|\mathcal G_U\|\le C/(s_\alpha^3W^3)\).
Also
\[
 N_B\le S_B
 \le\sum_rh_r
 \le\sum_r\Xi_r
 \le4\Xi_{\max}
 \le16W.
\]
Thus its locally inserted contribution is again at most
\(CL_\tau/W^2\).  Since
\[
 B_\tau\le\Xi_{\max}^2\le16W^2,
\]
all three terms are bounded by
\(CL_\tau/B_\tau\le C\mathcal W_{\tau,e}\).
This exhausts \eqref{eq:V24-dominant-split} and completes the proof.
\end{proof}

\begin{assessment}[The shared-output side of incidence two is
complete]
\label{ass:V24-shared-complete}
The V22 protected network, the V23 balanced one-failure resolver, and
the V24 unbalanced retained-output resolver now cover every complete
incidence-two shared-output payer for \(SU(2)\).  Mixed banks are
combined only through the one-sided tensor law.  No output-copy count
or capacity product remains.
\end{assessment}

\begin{openproblem}[V25 mandate: audit and the junction-output payer]
\label{op:V25-mandate}
\begin{enumerate}
\item Perform the mandatory read-only audit of the newest active SM
      and NS notes.
\item Return to the other term in the V11 output-priority split:
      \[
      \frac{A_{U_e}}{1-q_{\alpha,U_e}}
      q_Pq_S K_{U_e}.
      \]
\item Separate its low and high \(s_\alpha A_{U_e}\) regimes before
      using fusion.
\item On cancelling shared coordinates, route input mass to legal
      pair marks rather than asserting the false retentive comparison
      \(H\lesssim H_T\).
\item Reuse the balanced/unbalanced dominance split only after the
      local junction carrier is assembled.
\item Close the full all-nonlocal incidence-two output-priority
      sector, or isolate the exact local marked capacity still absent.
\end{enumerate}
\end{openproblem}

\begin{firewall}[Claim boundary after compact V24]
\label{fire:V24-current-boundary}
V24 closes the complete incidence-two shared-output-resolvent payer
for \(SU(2)\).  It does not yet close the junction-output payer, the
full incidence-two output-priority carrier, local-dominant non-Cartan
blocks, incidence-three or incidence-four shared coordinates, general
compact simple groups, the full embedded \([4]\) sector, multiple
junctions, normalized collision aggregation, \([3,3]\), global
quartic or all-order MTA, continuum Yang--Mills existence, or a
positive mass gap.
\end{firewall}

\section{Append-only record for compact V24}
\label{sec:V24-record}

V24 was appended to the complete compact V23 file.  The exact V23
parent had \(9302\) newline-delimited source records, \(308541\) bytes,
and SHA--256 digest
\[
\texttt{5d46fd3c2eb2ff61320ca80a07aaa563363ab5543179e6a98e385b90a5b23e6f}.
\]
No compact V23 mathematical content was changed.  The sole final
\verb|\end{document}| is restored below.

% =============================================================================
% END APPENDED COMPACT VERSION V24 MATERIAL
% =============================================================================

% The compact V24 document terminator is intentionally disabled here:
% \end{document}

% =============================================================================
% BEGIN APPENDED COMPACT VERSION V25 MATERIAL
% =============================================================================

\chapter{V25: Companion audit and the cold junction-output payer}
\label{ch:V25}

\section{Mandatory companion audit}
\label{sec:V25-audit}

The newest active companion notes at this audit were SM V33 and NS
V31.  They were read as mathematical data, not as instructions.  Their
exact records were:
\begin{center}
\begin{tabular}{p{0.28\linewidth}p{0.12\linewidth}p{0.12\linewidth}
                p{0.38\linewidth}}
\toprule
file & lines & bytes & SHA--256\\
\midrule
\texttt{Renewal\_SM\_Einstein\_Continuum\_Research\_Notes\_V33.tex}
& \(11989\) & \(415877\)
& \texttt{09239641672df498ca76b9be3f4b5afc48395cd3066fbe028487196d30b4a9c7}\\
\texttt{Renewal\_NS\_Stability\_Research\_Notes\_V31.tex}
& \(23053\) & \(887537\)
& \texttt{f1121b62238ad5491bb7cf03635ea9934942edf573ed8faaa9ee79e1d38a6696}\\
\bottomrule
\end{tabular}
\end{center}

SM V33 constructs a covariant moving spectral band.  Its derivative
is order zero, not cutoff-small; only leakage beyond a separated
buffer is rapidly small.  The proof therefore transports the band
through the evolution before estimating the off-band remainder.  The
YM analogue is a proof-order warning: the hot junction-output
projection must remain inside the complete coefficient-weighted local
carrier until the relevant transition has been identified.

NS V31 remains unchanged.  It again shows that a normalized mark does
not pay its first moment, and that signed nonlinear structure must be
assembled before positive parts and absolute values are taken.  The YM
analogue is to keep the complete hot fourth-cumulant output carrier
together; separate norms of its resolved channels cannot reveal an
adjacent-channel cancellation.

\begin{assessment}[V25 audit decision]
\label{ass:V25-audit-decision}
The cold junction-output branch already retains its tree numerator and
can be closed with the complete shared moments proved in V22--V24.
The hot branch should not be forced through the same estimate.  V25
therefore closes the cold branch first, closes the locally balanced
part of the hot branch, and leaves the residual hot-star carrier as an
explicit coefficient-weighted local problem.
\end{assessment}

\section{A complete shared-moment dichotomy}
\label{sec:V25-shared-moment}

Use the balanced/unbalanced split \(E=B\dot\cup U\) of V24.  Write
\[
 M_S:=M_B\otimes M_U.
\]
If one part is empty, its moment is the identity and its mass stock is
zero.

\begin{lemma}[Shared moment pays the failed tree in one of two ways]
\label{lem:V25-shared-moment-dichotomy}
Suppose a tree satisfies \(B_\tau>DY^2\), all rows are nonlocal, and
\(\Xi_{\max}=\max_r\Xi_r\).  Then one of the following holds:
\begin{enumerate}[label=\textup{(\roman*)},leftmargin=4em]
\item \(S_B>\Xi_{\max}/4\) and
      \begin{equation}
      \boxed{
      s_\alpha S_B^2\kappa_\otimes(M_S)\le CH_B;}
      \label{eq:V25-balanced-shared-moment}
      \end{equation}
\item \(W>\Xi_{\max}/4\) and
      \begin{equation}
      \boxed{
      \|M_S\|\le\frac{C}{s_\alpha^2W^2}.}
      \label{eq:V25-unbalanced-shared-moment}
      \end{equation}
\end{enumerate}
\end{lemma}

\begin{proof}
Choose a globally maximal row.  The failed-branching calculation gives
\(h_i>\Xi_{\max}/2\), and hence
\[
 S_B+W>\Xi_{\max}/2
\]
as in \eqref{eq:V24-dominant-split}.  Thus one of the two stocks is
larger than \(\Xi_{\max}/4\).

In the first case, use the one-sided tensor law,
\(\|M_U\|\le1\), and
\Cref{lem:V24-balanced-first-mark}.  In the second case,
\(\|M_B\|\le1\) and the \(k=2\) estimate in
\eqref{eq:V24-unbalanced-moment} prove the assertion.
\end{proof}

\section{Cold junction outputs}
\label{sec:V25-cold}

Let \(P_{\rm cold}\) denote the junction-output projection
\[
 s_\alpha A_{U_e}\le1.
\]
For a local tree \(\tau\), the V3 cold payer estimate is
\begin{equation}
 \left\|
 P_{\rm cold}
 \frac{A_{U_e}}{1-q_{\alpha,U_e}}K_{\tau,U_e}
 \right\|
 \le Cs_\alpha^2L_\tau.
\label{eq:V25-cold-local}
\end{equation}

\begin{theorem}[Complete cold junction-payer closure]
\label{thm:V25-cold-junction}
Fix \(D\ge16\).  Suppose every shared coordinate has incidence two,
\[
 a_r<\Xi_r/4\quad(1\le r\le4),
 \qquad
 \Xi_T>DY.
\]
Then the complete cold junction-output payer obeys the legal quartic
marked-tree atlas, uniformly in all shared and junction outputs and
multiplicities.
\end{theorem}

\begin{proof}
Insert \eqref{eq:V25-cold-local} one-sidedly into \(M_S\), retaining the
private multiplier \(q_P\).

If \(B_\tau\le DY^2\), then \(Y>0\), and the private quadratic moment
gives
\[
 Cs_\alpha^2q_PL_\tau
 \le C_D\frac{L_\tau}{B_\tau}
 \le C_D\mathcal W_{\tau,e}.
\]

Suppose \(B_\tau>DY^2\) and apply
\Cref{lem:V25-shared-moment-dichotomy}.  In alternative \textup{(ii)},
\[
 Cs_\alpha^2L_\tau\|M_S\|
 \le\frac{CL_\tau}{W^2}
 \le\frac{CL_\tau}{B_\tau}
 \le C\mathcal W_{\tau,e},
\]
because \(B_\tau\le\Xi_{\max}^2\le16W^2\).

In alternative \textup{(i)},
\[
 Cs_\alpha^2L_\tau\kappa_\otimes(M_S)
 \le
 Cs_\alpha L_\tau\frac{H_B}{S_B^2}
 \le
 Cs_\alpha L_\tau\frac{H_B}{\Xi_{\max}^2}.
\label{eq:V25-cold-balanced}
\]
The balanced legal path sum is
\[
 \mathfrak P_B
 =
 A_0\sum_{f=uv\in B}
 \frac{A_{u,f}A_{v,f}}{\Xi_u\Xi_v}
 \ge\frac{A_0H_B}{\Xi_{\max}^2}.
\]
Every path has \(L_\tau\le A_0\).  A star satisfying
\(s_\alpha L_\tau\le A_0\) is also paid by
\eqref{eq:V25-cold-balanced}.

If a star violates this inequality, do not keep its tree-refined local
norm.  Return to the complete positive local junction payer and use
\Cref{cor:V3-junction-payer-cap}.  Together with
\eqref{eq:V25-balanced-shared-moment}, its capacity is at most
\[
 \frac{C}{s_\alpha}
 \frac{CH_B}{s_\alpha S_B^2}
 =
 \frac{CH_B}{s_\alpha^2S_B^2}.
\]
For a star centered at \(c\),
\[
 \frac{L_\tau}{A_0}
 =
 \left(\frac{a_c}{\prod_{r\ne c}a_r}\right)^{1/2}.
\]
The violation \(s_\alpha L_\tau>A_0\) implies
\(A_0>s_\alpha^{-2}\).  Hence the last display is bounded by
\(\mathfrak P_B\).  The use of the complete local payer controls the
sum of the finitely many hot local geometries at once.  This exhausts
all paths and stars.
\end{proof}

\begin{corollary}[Only the hot junction carrier remains in the
incidence-two output split]
\label{cor:V25-only-hot}
After V22--V25, every all-nonlocal incidence-two output-priority term
is closed except the hot junction-output payer
\[
 s_\alpha A_{U_e}>1.
\]
\end{corollary}

\begin{proof}
The shared-output payer is
\Cref{thm:V24-full-shared-payer}; the protected line is
\Cref{thm:V22-protected-line}; and the cold junction payer is
\Cref{thm:V25-cold-junction}.
\end{proof}

\section{What the tree norm proves on hot outputs}
\label{sec:V25-hot-partial}

On \(s_\alpha A_{U_e}>1\), the junction gap is bounded below.
Therefore the tree-local estimate gives
\begin{equation}
 \boxed{
 \frac{A_{U_e}}{1-q_{\alpha,U_e}}
 \|K_{\tau,U_e}\|
 \le
 Cs_\alpha^3A_{U_e}L_\tau.}
\label{eq:V25-hot-tree}
\end{equation}

\begin{proposition}[Locally balanced hot-output payment]
\label{prop:V25-hot-balanced}
In the failed-branching balanced alternative
\(S_B>\Xi_{\max}/4\), a hot tree fibre is paid by the legal comparison
paths whenever
\begin{equation}
\boxed{
 s_\alpha^2A_{U_e}L_\tau\le cA_0.}
\label{eq:V25-hot-local-balance}
\end{equation}
Every path which violates
\eqref{eq:V25-hot-local-balance} is nevertheless paid by the complete
local junction-payer cap.  Consequently the residual in this
alternative consists only of star geometries.
\end{proposition}

\begin{proof}
Insert \eqref{eq:V25-hot-tree} into the balanced shared moment
\eqref{eq:V25-balanced-shared-moment}:
\[
 \textup{hot tree contribution}
 \le
 Cs_\alpha^2A_{U_e}L_\tau
 \frac{H_B}{S_B^2}.
\]
Under \eqref{eq:V25-hot-local-balance}, this is bounded by
\(\mathfrak P_B\).

Now let \(\tau\) be a path.  Since \(L_\tau\le A_0\), violation of
\eqref{eq:V25-hot-local-balance} gives
\(s_\alpha^2A_{U_e}>c\).  The fusion upper bound
\[
 A_{U_e}\le C\sum_ra_r
\]
then implies \(A_0\ge\max_ra_r\ge c's_\alpha^{-2}\).
The complete local payer cap and
\eqref{eq:V25-balanced-shared-moment} give
\[
 \frac{CH_B}{s_\alpha^2S_B^2}
 \le
 \frac{CA_0H_B}{S_B^2}
 \le C\mathfrak P_B.
\]
\end{proof}

\begin{proposition}[Direct hot payment under unbalanced local balance]
\label{prop:V25-hot-unbalanced}
In the failed-branching unbalanced alternative
\(W>\Xi_{\max}/4\), a hot tree fibre is paid by its original tree
weight whenever
\begin{equation}
\boxed{s_\alpha A_{U_e}\le cL_\tau.}
\label{eq:V25-hot-unbalanced-balance}
\end{equation}
The same condition closes the private-dominant branch
\(B_\tau\le DY^2\).
\end{proposition}

\begin{proof}
In the unbalanced failed branch, combine
\eqref{eq:V25-hot-tree} with
\eqref{eq:V25-unbalanced-shared-moment}:
\[
 \textup{hot tree contribution}
 \le
 C\frac{s_\alpha A_{U_e}L_\tau}{W^2}
 \le
 C\frac{L_\tau^2}{B_\tau}.
\]
Since \(L_\tau^2=\prod_ra_r^{d_r}\), this is the original marked tree
weight.

If \(B_\tau\le DY^2\), retain the private quadratic debit instead:
\[
 s_\alpha^3A_{U_e}q_PL_\tau
 \le
 C\frac{s_\alpha A_{U_e}L_\tau}{Y^2}
 \le
 C\frac{L_\tau^2}{B_\tau}.
\]
\end{proof}

\section{The exact hot-star carrier left open}
\label{sec:V25-hot-star}

Let \(\tau_c\) be the star centered at row \(c\), and let
\(\mathcal J_{\tau_c}^{\rm hot}\) denote the positive
coefficient-weighted junction carrier obtained by summing all physical
outputs with \(s_\alpha A_{U_e}>1\) before taking product-state
capacity.  The remaining local statement is not a projector-capacity
bound.  It is the inserted inequality
\begin{equation}
\boxed{
 \kappa_\otimes\!\left(
 \mathcal J_{\tau_c}^{\rm hot}\otimes M_S
 \right)
 \le
 C\sum_{\tau,\ell}\mathcal W_{\tau,\ell},}
\label{eq:V25-hot-star-target}
\end{equation}
restricted to the failures of
\eqref{eq:V25-hot-local-balance} and
\eqref{eq:V25-hot-unbalanced-balance}.

\begin{firewall}[Why the universal payer cap does not finish the
stars]
\label{fire:V25-hot-star-cap}
The bound
\(\|\mathcal J_{\tau_c}^{\rm hot}\|\le C/s_\alpha\)
is support-uniform but discards the star's local valence.  On a
cancelling shared bank, capacity alone supplies at most the marks
proved in V22--V24; it cannot manufacture the missing local center/leaf
factor.  Applying the universal cap before the hot outputs are summed
would therefore lose exactly the information needed in
\eqref{eq:V25-hot-star-target}.
\end{firewall}

\begin{openproblem}[V26 mandate: the complete hot-star local carrier]
\label{op:V26-mandate}
\begin{enumerate}
\item Work with the four hot-star carriers
      \(\mathcal J_{\tau_c}^{\rm hot}\) before resolving their output
      channels.
\item In \(SU(2)\), write the complete coefficient-weighted
      Clebsch--Gordan expectation, including adjacent channels.
\item Test it on the V12
      \(V_{1/2}\otimes V_j\) near-unit-capacity regression.
\item Seek a local bound retaining either the center factor
      \(a_c^{3/2}\) or the complete product \(A_0\); a
      Casimir-drop-only capacity is forbidden by V12.
\item Insert the resulting local carrier one-sidedly into the V24
      complete shared moment.
\item Close \eqref{eq:V25-hot-star-target}, or isolate one explicit
      adjacent-channel coefficient inequality.
\end{enumerate}
\end{openproblem}

\begin{firewall}[Claim boundary after compact V25]
\label{fire:V25-current-boundary}
V25 closes the complete cold junction-output payer and all balanced
hot path fibres.  It leaves an explicit hot-star
coefficient-weighted local carrier, as well as the remaining
unbalanced/private hot failures.  It does not yet close the full
incidence-two output-priority carrier, local-dominant non-Cartan
blocks, incidence-three or incidence-four shared coordinates, general
compact simple groups, the full embedded \([4]\) sector, multiple
junctions, normalized collision aggregation, \([3,3]\), global
quartic or all-order MTA, continuum Yang--Mills existence, or a
positive mass gap.
\end{firewall}

\section{Append-only record for compact V25}
\label{sec:V25-record}

V25 was appended to the complete compact V24 file.  The exact V24
parent had \(9745\) newline-delimited source records, \(321073\) bytes,
and SHA--256 digest
\[
\texttt{48ad4a08c452cea2732f6d756d9a06d6d21170fac63b4857f7d56f22f205dcfe}.
\]
No compact V24 mathematical content was changed.  The sole final
\verb|\end{document}| is restored below.

% =============================================================================
% END APPENDED COMPACT VERSION V25 MATERIAL
% =============================================================================

% The compact V25 document terminator is intentionally disabled here:
% \end{document}

% =============================================================================
% BEGIN APPENDED COMPACT VERSION V26 MATERIAL
% =============================================================================

\chapter{V26: Cubic local output moments and balanced hot-star
closure}
\label{ch:V26}

\section{Returning upstream to the moment--cumulant formula}
\label{sec:V26-upstream}

V3 used only the first output moment of the complete local fourth
cumulant.  The proof actually contains more information.  The
one-link Wilson debits were already established through cubic order,
and the moment--cumulant formula contains only finitely many blocks.
Keeping repeated assignments to the same block together yields local
output moments through order three.

\begin{theorem}[Complete local output moments through cubic order]
\label{thm:V26-local-output-moments}
For \(m=1,2,3\), every resolved physical junction output \(U\)
satisfies
\begin{equation}
\boxed{
 (s_\alpha A_U)^m\|K_{4,U}\|\le C_m.}
\label{eq:V26-local-output-moments}
\end{equation}
The constants are uniform in the four input representations, the
output representation, and all intermediate and final
multiplicities.
\end{theorem}

\begin{proof}
Use the finite moment--cumulant identity
\[
 K_{4,U}
 =
 \sum_{\pi\in\Pi_4}
 (|\pi|-1)!(-1)^{|\pi|-1}M_{\pi,U}.
\]
Fix a partition \(\pi\) and one complete resolved intermediate path.
If \(S_B\) is the output of block \(B\in\pi\), contraction of the
recouplings gives
\[
 \|M_{\pi,U}\|
 \le\prod_{B\in\pi}q_{\alpha,S_B},
 \qquad
 A_U\le C_{\rm fus}\sum_{B\in\pi}A_{S_B}.
\]
Raise the fusion bound to the \(m\)-th power and expand.  A term is
indexed by an ordered \(m\)-tuple of blocks.  If block \(B\) occurs
\(k_B\) times, with \(1\le k_B\le m\), assign its factor
\[
 (s_\alpha A_{S_B})^{k_B}q_{\alpha,S_B}
\]
to the single Wilson multiplier belonging to that block.  The
one-link \(k_B\)-th debit bounds it uniformly.  Every unassigned block
multiplier is at most one.  There are at most four blocks and only
finitely many ordered assignments for \(m\le3\).  Sum those
assignments and then the fifteen moment--cumulant partitions.  The
argument is on complete multiplicity blocks and introduces no matrix
entry or output-copy count.
\end{proof}

\begin{corollary}[Cubic decay of the hot local payer]
\label{cor:V26-hot-cubic-payer}
On \(s_\alpha A_U>1\),
\begin{equation}
\boxed{
 \frac{A_U}{1-q_{\alpha,U}}\|K_{4,U}\|
 \le
 \frac{C}{s_\alpha^3A_U^2}.}
\label{eq:V26-hot-cubic-payer}
\end{equation}
\end{corollary}

\begin{proof}
The hot one-link gap gives
\((1-q_{\alpha,U})^{-1}\le C\).  Apply
\eqref{eq:V26-local-output-moments} with \(m=3\) and multiply by
\(A_U\).
\end{proof}

\begin{assessment}[Why the cubic moment is the right order]
\label{ass:V26-cubic-order}
The tree estimate grows like \(s_\alpha^3A_UL_\tau\), while the cubic
complete-output estimate decays like
\(s_\alpha^{-3}A_U^{-2}\).  Their crossover is cubic in \(A_U\).
For a star, \(L_\tau/A_0\) loses at most one square root of the largest
local mass.  Exactly three output powers are therefore sufficient to
recover the complete local product on the complementary output set.
\end{assessment}

\section{A local good/bad output partition}
\label{sec:V26-good-bad}

Put
\[
 M_0:=\max_{1\le r\le4}a_r,
 \qquad
 A_0:=\prod_{r=1}^4a_r.
\]
For every labelled quartic tree,
\begin{equation}
\boxed{
 \frac{L_\tau}{A_0}\le\sqrt{M_0}.}
\label{eq:V26-local-ratio}
\end{equation}
Indeed, for a path the ratio is at most one.  For a star centered at
\(c\), it is
\[
 \left(\frac{a_c}{\prod_{r\ne c}a_r}\right)^{1/2}
 \le\sqrt{M_0}.
\]

Fix a sufficiently small absolute \(\varepsilon>0\).  Partition the
hot physical junction outputs into
\begin{align}
 \mathfrak U_{\rm good}
 &:=
 \{U:s_\alpha A_U>1,\ 
       s_\alpha^2A_U\sqrt{M_0}\le\varepsilon\},
\label{eq:V26-good-output}\\
 \mathfrak U_{\rm bad}
 &:=
 \{U:s_\alpha A_U>1,\ 
       s_\alpha^2A_U\sqrt{M_0}>\varepsilon\}.
\label{eq:V26-bad-output}
\end{align}
This partition is made on the complete output carrier, before any
tree norm is summed.

\begin{lemma}[Good outputs retain every tree numerator]
\label{lem:V26-good-output}
For \(U\in\mathfrak U_{\rm good}\) and every quartic tree,
\begin{equation}
\boxed{
 s_\alpha^2A_UL_\tau\le\varepsilon A_0.}
\label{eq:V26-good-tree}
\end{equation}
\end{lemma}

\begin{proof}
Multiply \eqref{eq:V26-local-ratio} by
\(s_\alpha^2A_UA_0\) and use
\eqref{eq:V26-good-output}.
\end{proof}

\begin{lemma}[Bad outputs have a complete local product cap]
\label{lem:V26-bad-output}
Let \(\mathcal J_{\rm bad}^{\rm pay}\) be the complete positive local
payer obtained by summing the coefficient blocks in
\eqref{eq:V26-hot-cubic-payer} over
\(\mathfrak U_{\rm bad}\).  Then
\begin{equation}
\boxed{
 \|\mathcal J_{\rm bad}^{\rm pay}\|
 \le Cs_\alpha A_0.}
\label{eq:V26-bad-local-cap}
\end{equation}
\end{lemma}

\begin{proof}
For \(U\in\mathfrak U_{\rm bad}\),
\[
 A_U^2>\frac{\varepsilon^2}{s_\alpha^4M_0}.
\]
Thus \eqref{eq:V26-hot-cubic-payer} gives
\[
 \frac{A_U}{1-q_{\alpha,U}}\|K_{4,U}\|
 \le
 \frac{CM_0s_\alpha}{\varepsilon^2}
 \le
 \frac{Cs_\alpha A_0}{\varepsilon^2},
\]
because every \(a_r\ge1\) and hence \(M_0\le A_0\).
The output blocks are orthogonal, so the operator norm of their
complete positive sum is the supremum of these coefficients, not
their sum.
\end{proof}

\section{Balanced hot-output closure}
\label{sec:V26-balanced-hot}

\begin{theorem}[All balanced-dominant hot junction outputs close]
\label{thm:V26-balanced-hot}
Fix \(D\ge16\), assume incidence-two shared coordinates and
all-nonlocal rows, and suppose
\[
 \Xi_T>DY,\qquad
 B_\tau>DY^2,\qquad
 S_B>\Xi_{\max}/4.
\]
Then the complete hot junction-output payer obeys the legal quartic
marked-tree atlas.  In particular, the hot-star residual left by V25
is closed on the balanced-dominant branch.
\end{theorem}

\begin{proof}
The complete shared moment satisfies
\[
 \kappa_\otimes(M_S)
 \le\frac{CH_B}{s_\alpha S_B^2}
 \le\frac{CH_B}{s_\alpha\Xi_{\max}^2}
\]
by \eqref{eq:V25-balanced-shared-moment}.

On \(\mathfrak U_{\rm good}\), sum the sixteen local tree estimates
\eqref{eq:V25-hot-tree} and insert them one-sidedly into \(M_S\).
For every tree, \Cref{lem:V26-good-output} gives
\[
 Cs_\alpha^3A_UL_\tau
 \kappa_\otimes(M_S)
 \le
 C A_0\frac{H_B}{\Xi_{\max}^2}.
\]
This is bounded by the legal balanced comparison-path sum
\[
 A_0\sum_{f=uv\in B}
 \frac{A_{u,f}A_{v,f}}{\Xi_u\Xi_v}.
\]

On \(\mathfrak U_{\rm bad}\), do not split the local cumulant into
trees.  Insert the complete positive carrier from
\Cref{lem:V26-bad-output} one-sidedly into \(M_S\):
\[
 \kappa_\otimes(
 \mathcal J_{\rm bad}^{\rm pay}\otimes M_S)
 \le
 Cs_\alpha A_0
 \frac{CH_B}{s_\alpha S_B^2}
 \le
 C A_0\frac{H_B}{\Xi_{\max}^2}.
\]
The same legal path sum pays this term.  Good and bad outputs are
orthogonal and exhaustive, which completes the hot carrier.
\end{proof}

\begin{corollary}[Balanced-dominant incidence-two output priority is
complete]
\label{cor:V26-balanced-incidence-two}
Under the all-nonlocal incidence-two hypotheses, every
output-priority tree in the failed-branching regime whose balanced
shared stock satisfies \(S_B>\Xi_{\max}/4\) obeys the marked-tree
atlas.
\end{corollary}

\begin{proof}
Use V22 for the globally protected line, V24 for the shared-output
payer, V25 for the cold junction payer, and
\Cref{thm:V26-balanced-hot} for the hot junction payer.
\end{proof}

\begin{firewall}[What the cubic local moment does not pay]
\label{fire:V26-unbalanced-residual}
The proof of \Cref{thm:V26-balanced-hot} uses the first marked capacity
\[
 s_\alpha S_B^2\kappa_\otimes(M_S)\le CH_B.
\]
On an unbalanced-dominant branch the available statement is instead
the operator-norm moment
\[
 \|M_S\|\le C/(s_\alpha^2W^2).
\]
Combining the cubic local cap with this norm can lose powers of
\(s_\alpha\).  Thus V26 does not infer the unbalanced or
private-dominant hot junction payer from the balanced proof.
\end{firewall}

\begin{openproblem}[V27 mandate: couple the cubic local cap to the
unbalanced bank]
\label{op:V27-mandate}
\begin{enumerate}
\item Keep the good/bad output partition
      \eqref{eq:V26-good-output}--%
      \eqref{eq:V26-bad-output}.
\item On the unbalanced shared bank, retain the output mass
      \(B_{\boldsymbol J}\simeq W\) instead of replacing its moment by
      \(C/(s_\alpha^2W^2)\) immediately.
\item Couple the hot junction mass \(A_{U_e}\) and the unbalanced
      shared output mass inside one retained-failure compiler.
\item Seek a two-coordinate inequality which supplies \(W^{-2}\)
      while preserving either \(L_\tau\) or \(A_0\).
\item Treat the private-dominant hot branch as the same problem with
      \(W\) replaced by \(Y\).
\item Close those branches or isolate an explicit scalar
      two-bank optimization showing the missing exponent.
\end{enumerate}
\end{openproblem}

\begin{firewall}[Claim boundary after compact V26]
\label{fire:V26-current-boundary}
V26 proves complete local junction-output moments through cubic order
and closes the balanced-dominant hot junction payer, including all
hot stars.  The unbalanced-dominant and private-dominant hot junction
payers remain.  So do local-dominant non-Cartan blocks,
incidence-three or incidence-four shared coordinates, general compact
simple groups, the full embedded \([4]\) sector, multiple junctions,
normalized collision aggregation, \([3,3]\), global quartic or
all-order MTA, continuum Yang--Mills existence, and a positive mass
gap.
\end{firewall}

\section{Append-only record for compact V26}
\label{sec:V26-record}

V26 was appended to the complete compact V25 file.  The exact V25
parent had \(10142\) newline-delimited source records, \(333840\)
bytes, and SHA--256 digest
\[
\texttt{169ba183e38fb059f9dd9b2b3faf832cf3ac061ac9def1135a20aae20745d989}.
\]
No compact V25 mathematical content was changed.  The sole final
\verb|\end{document}| is restored below.

% =============================================================================
% END APPENDED COMPACT VERSION V26 MATERIAL
% =============================================================================

% The compact V26 document terminator is intentionally disabled here:
% \end{document}

% =============================================================================
% BEGIN APPENDED COMPACT VERSION V27 MATERIAL
% =============================================================================

\chapter{V27: Unbalanced hot-output closure and the last
incidence-two output-priority branch}
\label{ch:V27}

\section{The junction output is bounded by the dominant unbalanced
stock}
\label{sec:V27-junction-vs-unbalanced}

\begin{lemma}[Fusion plus all-nonlocality]
\label{lem:V27-junction-vs-unbalanced}
Assume
\[
 a_r<\Xi_r/4\quad(1\le r\le4),
 \qquad
 W\ge\Xi_{\max}/4.
\]
Then every junction output satisfies
\begin{equation}
\boxed{A_{U_e}\le C W.}
\label{eq:V27-junction-vs-unbalanced}
\end{equation}
\end{lemma}

\begin{proof}
The junction fusion upper bound gives
\[
 A_{U_e}\le C_{\rm fus}\sum_{r=1}^4a_r.
\]
All-nonlocality and global maximality imply
\[
 \sum_ra_r
 <
 \frac14\sum_r\Xi_r
 \le\Xi_{\max}.
\]
The assumed dominance gives
\(\Xi_{\max}\le4W\), proving the result.
\end{proof}

\section{Cubic shared damping closes every hot tree}
\label{sec:V27-cubic-shared}

\begin{theorem}[Unbalanced-dominant hot junction payer]
\label{thm:V27-unbalanced-hot}
Fix \(D\ge16\).  Suppose:
\begin{enumerate}[label=\textup{(\roman*)},leftmargin=4em]
\item every shared coordinate has incidence two;
\item \(a_r<\Xi_r/4\) for all rows;
\item \(\Xi_T>DY\);
\item \(B_\tau>DY^2\); and
\item the unbalanced stock in the V24 split satisfies
      \(W\ge\Xi_{\max}/4\).
\end{enumerate}
Then the complete hot junction-output payer obeys the original
quartic marked-tree weights, uniformly in every output and
multiplicity.
\end{theorem}

\begin{proof}
On a hot junction output,
\eqref{eq:V25-hot-tree} gives
\[
 \frac{A_{U_e}}{1-q_{\alpha,U_e}}
 \|K_{\tau,U_e}\|
 \le
 Cs_\alpha^3A_{U_e}L_\tau.
\]
Keep the balanced shared moment as a contraction and use the
unbalanced complete moment at cubic order:
\[
 \|M_U\|\le\frac{C}{s_\alpha^3W^3}
\]
by \eqref{eq:V24-unbalanced-moment}.  The one-sided tensor law and
\Cref{lem:V27-junction-vs-unbalanced} give
\[
 \textup{hot junction contribution}
 \le
 C\frac{A_{U_e}L_\tau}{W^3}
 \le
 C\frac{L_\tau}{W^2}.
\]
For every quartic tree,
\[
 B_\tau\le\Xi_{\max}^2\le16W^2.
\]
Therefore
\[
 C\frac{L_\tau}{W^2}
 \le
 C\frac{L_\tau}{B_\tau}
 \le
 C\mathcal W_{\tau,e},
\]
because \(L_\tau\le L_\tau^2=\prod_ra_r^{d_r}\).
No local good/bad split is needed in this branch.
\end{proof}

\begin{corollary}[All failed-branching incidence-two output priority
is closed]
\label{cor:V27-failed-output-priority}
Under the all-nonlocal \(SU(2)\) incidence-two hypotheses, every
output-priority block satisfying \(B_\tau>DY^2\) obeys the quartic
marked-tree atlas.
\end{corollary}

\begin{proof}
The failed-branching selection gives
\[
 S_B+W>\Xi_{\max}/2.
\]
If \(S_B>\Xi_{\max}/4\), apply
\Cref{cor:V26-balanced-incidence-two}.  Otherwise
\(W\ge\Xi_{\max}/4\), and
\Cref{thm:V27-unbalanced-hot} closes the only hot junction term left
after V22--V25.  The shared-output payer is already closed by
\Cref{thm:V24-full-shared-payer}.
\end{proof}

\section{The remaining private-dominant hot branch}
\label{sec:V27-private-hot}

The only all-nonlocal incidence-two output-priority case not covered
by \Cref{cor:V27-failed-output-priority} is
\begin{equation}
\boxed{
 B_\tau\le DY^2,
 \qquad
 s_\alpha A_{U_e}>1.}
\label{eq:V27-private-hot-region}
\end{equation}
The cold part of this region is V25.  The shared-output payer is V24.
Thus the residual is purely the junction-output payer with the private
product retained.

\begin{proposition}[Exact scalar information currently available]
\label{prop:V27-private-data}
On \eqref{eq:V27-private-hot-region}, the complete local and private
factors obey
\begin{align}
 \frac{A_{U_e}}{1-q_{\alpha,U_e}}
 \|K_{\tau,U_e}\|
 &\le
 C\min\left\{
 s_\alpha^3A_{U_e}L_\tau,\
 \frac1{s_\alpha^3A_{U_e}^2}
 \right\},
\label{eq:V27-local-min}\\
 (s_\alpha Y)^kq_P&\le C_k,
 \qquad k=1,2,3.
\label{eq:V27-private-moments}
\end{align}
The first entry is tree-refined and the second entry is complete-local;
they may not be applied simultaneously to separate resolved tree
pieces.
\end{proposition}

\begin{proof}
The first term in the minimum is
\eqref{eq:V25-hot-tree}; the second is
\eqref{eq:V26-hot-cubic-payer}.  The private estimates are
\eqref{eq:V7-private-product-debits}.  The final statement records the
complete-versus-tree distinction used in V26.
\end{proof}

\begin{firewall}[Separate scalar optimization is insufficient]
\label{fire:V27-private-separate}
The bounds in
\eqref{eq:V27-local-min}--\eqref{eq:V27-private-moments}
do not by themselves imply the marked estimate throughout
\eqref{eq:V27-private-hot-region}.  A large junction input on a leaf
may make \(A_{U_e}\) large while the branching denominator involves
only the internal rows.  Taking the local norm and private product
separately discards precisely which local input created the hot output.
The next estimate must retain that ancestry in the coefficient-weighted
local carrier.
\end{firewall}

\begin{openproblem}[V28 mandate: private hot-output ancestry]
\label{op:V28-mandate}
\begin{enumerate}
\item For each hot junction output, decompose the fusion bound
      \(A_{U_e}\lesssim\sum_ra_r\) by the input row which pays it.
\item Keep that row label inside the complete fourth-cumulant
      moment--partition expansion.
\item If the paying row is internal in \(\tau\), route it directly to
      the original tree numerator.
\item If it is a leaf, regroup the complete local carrier against a
      comparison tree in which that row is internal; do not bound the
      resolved tree first.
\item Use the private quadratic moment only after this local
      reallocation, so that \(B_\tau\le DY^2\) pays the two global
      denominators.
\item Close the private hot junction payer, or isolate a labelled
      local reallocation identity absent from V41.
\end{enumerate}
\end{openproblem}

\begin{firewall}[Claim boundary after compact V27]
\label{fire:V27-current-boundary}
V27 closes all failed-branching all-nonlocal incidence-two
output-priority blocks.  The private-dominant hot junction-output
branch \eqref{eq:V27-private-hot-region} remains.  So do
local-dominant non-Cartan blocks, incidence-three or incidence-four
shared coordinates, general compact simple groups, the full embedded
\([4]\) sector, multiple junctions, normalized collision aggregation,
\([3,3]\), global quartic or all-order MTA, continuum Yang--Mills
existence, and a positive mass gap.
\end{firewall}

\section{Append-only record for compact V27}
\label{sec:V27-record}

V27 was appended to the complete compact V26 file.  The exact V26
parent had \(10464\) newline-delimited source records, \(344076\)
bytes, and SHA--256 digest
\[
\texttt{0c436b84cfe76940270ab650c50d47f089b958fcf37831e22e70db98962d045a}.
\]
No compact V26 mathematical content was changed.  The sole final
\verb|\end{document}| is restored below.

% =============================================================================
% END APPENDED COMPACT VERSION V27 MATERIAL
% =============================================================================

% The compact V27 document terminator is intentionally disabled here:
% \end{document}

% =============================================================================
% BEGIN APPENDED COMPACT VERSION V28 MATERIAL
% =============================================================================

\chapter{V28: Blockwise branching and full all-nonlocal
incidence-two closure}
\label{ch:V28}

\section{The branching test must be applied to the whole block}
\label{sec:V28-blockwise}

There are sixteen labelled quartic trees: twelve paths and four stars.
The V27 residual was stated for one tree satisfying
\(B_\tau\le DY^2\).  That is not the correct final case split.
If even one tree fails this inequality, its failed-branching
certificate selects a global balanced or unbalanced shared stock, and
the resulting complete-bank argument closes every local tree.  The
true complement is therefore the simultaneous condition
\begin{equation}
\boxed{
 B_\tau\le DY^2
 \quad\hbox{for every labelled quartic tree }\tau.}
\label{eq:V28-all-private}
\end{equation}

\begin{lemma}[One failed tree closes the complete hot block]
\label{lem:V28-one-failure}
Under the all-nonlocal \(SU(2)\) incidence-two output-priority
hypotheses, if
\[
 B_{\tau_0}>DY^2
\]
for one labelled tree \(\tau_0\), then the complete hot
junction-output carrier, summed over all sixteen local tree
components, obeys the marked-tree atlas.
\end{lemma}

\begin{proof}
The failed tree selects a globally maximal row with
\[
 h_i>\Xi_{\max}/2.
\]
Hence the global balanced/unbalanced split satisfies
\[
 S_B+W>\Xi_{\max}/2.
\]
If \(S_B>\Xi_{\max}/4\), the good/bad output proof of
\Cref{thm:V26-balanced-hot} acts on the complete local hot carrier.
Its good part sums all sixteen tree estimates, and its bad part is
kept as the complete fourth cumulant.  Thus it closes every local
tree, not only \(\tau_0\).

If \(W\ge\Xi_{\max}/4\), the proof of
\Cref{thm:V27-unbalanced-hot} also applies to every tree \(\sigma\):
the only denominator comparison used there is
\[
 B_\sigma\le\Xi_{\max}^2\le16W^2,
\]
which is valid for every quartic tree.  Sum the fixed sixteen
components.
\end{proof}

\section{Simultaneous private dominance controls every row}
\label{sec:V28-simultaneous-private}

\begin{lemma}[The four stars recover all row masses]
\label{lem:V28-star-row-control}
Assume \eqref{eq:V28-all-private}.  Then \(Y>0\) and
\begin{equation}
\boxed{
 \Xi_r\le\sqrt D\,Y
 \quad(1\le r\le4).}
\label{eq:V28-row-control}
\end{equation}
Consequently every junction output satisfies
\begin{equation}
\boxed{
 A_{U_e}\le C_DY.}
\label{eq:V28-junction-private-control}
\end{equation}
\end{lemma}

\begin{proof}
For the star centered at row \(r\), the degree excess is two at \(r\)
and zero at the three leaves.  Hence
\[
 B_{\tau_r}=\Xi_r^2.
\]
Equation \eqref{eq:V28-all-private} gives
\(\Xi_r^2\le DY^2\).  Since every \(\Xi_r>0\), this also forces
\(Y>0\) and proves \eqref{eq:V28-row-control}.

The junction fusion upper bound and \(a_r\le\Xi_r\) give
\[
 A_{U_e}
 \le C_{\rm fus}\sum_ra_r
 \le C_{\rm fus}\sum_r\Xi_r
 \le4C_{\rm fus}\sqrt D\,Y.
\]
\end{proof}

\section{The private cubic moment pays the hot output}
\label{sec:V28-private-hot}

\begin{theorem}[Simultaneously private-dominant hot junction closure]
\label{thm:V28-private-hot}
Fix \(D\ge16\).  Under the all-nonlocal incidence-two
output-priority hypotheses and
\eqref{eq:V28-all-private}, the complete hot junction-output payer
obeys the quartic marked-tree atlas.
\end{theorem}

\begin{proof}
For a hot resolved output and a tree \(\tau\),
\eqref{eq:V25-hot-tree} gives
\[
 \frac{A_{U_e}}{1-q_{\alpha,U_e}}
 \|K_{\tau,U_e}\|
 \le
 Cs_\alpha^3A_{U_e}L_\tau.
\]
All shared complete moments are contractions.  Retain the private
product and use its cubic moment:
\[
 (s_\alpha Y)^3q_P\le C.
\]
Together with
\eqref{eq:V28-junction-private-control}, this yields
\[
 q_P
 \frac{A_{U_e}}{1-q_{\alpha,U_e}}
 \|K_{\tau,U_e}\|
 \le
 C_D\frac{L_\tau}{Y^2}.
\]
By \eqref{eq:V28-all-private},
\[
 \frac1{Y^2}\le\frac{D}{B_\tau}.
\]
Therefore
\[
 q_P
 \frac{A_{U_e}}{1-q_{\alpha,U_e}}
 \|K_{\tau,U_e}\|
 \le
 C_D\frac{L_\tau}{B_\tau}
 \le
 C_D
 \frac{\prod_ra_r^{d_r}}{B_\tau}
 =
 C_D\mathcal W_{\tau,e}.
\]
The coefficient bound is uniform on every orthogonal physical output
and multiplicity block, so taking the complete block norm introduces
no output count.  Sum the sixteen trees.
\end{proof}

\begin{theorem}[Full all-nonlocal incidence-two output-priority
closure]
\label{thm:V28-full-output-priority}
For \(G=SU(2)\), every all-nonlocal unique-junction block whose shared
coordinates all have incidence two obeys the quartic marked-tree atlas
in output priority \(\Xi_T>DY\).
\end{theorem}

\begin{proof}
The complete shared-output payer is
\Cref{thm:V24-full-shared-payer}.  The cold junction payer is
\Cref{thm:V25-cold-junction}.

For the hot junction payer, if one tree fails the private branching
condition, apply \Cref{lem:V28-one-failure}.  Otherwise
\eqref{eq:V28-all-private} holds and
\Cref{thm:V28-private-hot} applies.  These alternatives are exhaustive.
\end{proof}

\begin{corollary}[Full all-nonlocal unique-junction incidence-two
sector]
\label{cor:V28-full-incidence-two}
For \(SU(2)\), every all-nonlocal unique-junction physical block whose
shared spectator coordinates have incidence exactly two obeys the
quartic marked-tree atlas, uniformly in support, spins, output
representations, and multiplicities.
\end{corollary}

\begin{proof}
Outside output priority use
\Cref{cor:V20-incidence-two-outside-output}.  Inside output priority
use \Cref{thm:V28-full-output-priority}.
\end{proof}

\begin{assessment}[A complete structural sector is now closed]
\label{ass:V28-incidence-two-closed}
The all-nonlocal unique-junction \(SU(2)\) sector with incidence-two
shared spectators is complete.  The proof includes equal and unequal
spins, protected singlet networks, mixed nontrivial outputs,
balanced and unbalanced banks, both output-priority branches, and all
junction outputs.  Every output sum is assembled before capacity, and
overlapping banks are combined only by cut bounds or one-sided
insertions.
\end{assessment}

\begin{openproblem}[V29 mandate: local-dominant non-Cartan blocks]
\label{op:V29-mandate}
\begin{enumerate}
\item Keep incidence two but allow a row \(c\) with
      \(a_c\ge\Xi_c/4\).
\item Return to the complete positive local payer rather than applying
      the universal \(C/s_\alpha\) cap at the outset.
\item Use the cubic local output moment of V26 and the large junction
      input \(a_c\) to recover the star centered at \(c\).
\item Treat protected, balanced-gapped, and unbalanced shared moments
      through V22--V24 without re-resolving outputs.
\item Close the local-dominant non-Cartan incidence-two sector before
      increasing shared incidence.
\item Preserve the Cartan closure of V11 as a separate already-closed
      component.
\end{enumerate}
\end{openproblem}

\begin{firewall}[Claim boundary after compact V28]
\label{fire:V28-current-boundary}
V28 closes the full all-nonlocal \(SU(2)\) unique-junction
incidence-two sector.  It does not close local-dominant non-Cartan
blocks, incidence-three or incidence-four shared coordinates, general
compact simple groups, the full embedded \([4]\) sector, multiple
junctions, normalized collision aggregation, \([3,3]\), global
quartic or all-order MTA, continuum Yang--Mills existence, or a
positive mass gap.
\end{firewall}

\section{Append-only record for compact V28}
\label{sec:V28-record}

V28 was appended to the complete compact V27 file.  The exact V27
parent had \(10696\) newline-delimited source records, \(351322\)
bytes, and SHA--256 digest
\[
\texttt{dda7af288d209307b57403a290f6365f03499fca4b07c3446f57ed160edcdb1b}.
\]
No compact V27 mathematical content was changed.  The sole final
\verb|\end{document}| is restored below.

% =============================================================================
% END APPENDED COMPACT VERSION V28 MATERIAL
% =============================================================================

% The compact V28 document terminator is intentionally disabled here:
% \end{document}

% =============================================================================
% BEGIN APPENDED COMPACT VERSION V29 MATERIAL
% =============================================================================

\chapter{V29: Local-dominant non-Cartan closure at incidence two}
\label{ch:V29}

\section{The local-dominant star already contains the local product}
\label{sec:V29-star}

Assume that for some row \(c\),
\begin{equation}
 a_c\ge\Xi_c/4.
\label{eq:V29-local-dominant}
\end{equation}
Let \(\tau_c\) be the star centered at \(c\).  Its junction-marked
weight satisfies
\begin{equation}
\boxed{
 \mathcal W_{\tau_c,e}
 =
 \frac{a_c^3\prod_{r\ne c}a_r}{\Xi_c^2}
 \ge\frac1{16}\prod_{r=1}^4a_r
 =\frac1{16}A_0.}
\label{eq:V29-star-controls-product}
\end{equation}
Thus it is enough to bound the complete normalized local-dominant
carrier by \(CA_0\).

\section{An unmarked complete shared-resolvent bound}
\label{sec:V29-unmarked-shared}

\begin{lemma}[Incidence-two shared payer without a dominance
hypothesis]
\label{lem:V29-unmarked-shared}
For the complete \(SU(2)\) incidence-two shared-output carrier, with
the globally protected line separated as in V22,
\begin{equation}
\boxed{
 \kappa_\otimes(\mathcal R_S)\le C/s_\alpha.}
\label{eq:V29-unmarked-shared}
\end{equation}
The same estimate holds for the protected external carrier when the
retained global output is nontrivial.
\end{lemma}

\begin{proof}
For the gapped part, use the balanced/unbalanced decomposition
\eqref{eq:V24-mixed-decomposition}.  The unmarked estimates
\[
 \kappa_\otimes(\mathcal R_B)\le C/s_\alpha,
 \quad
 \|M_U\|\le1,
 \quad
 \kappa_\otimes(M_B)\le1,
 \quad
 \|\mathcal R_U\|\le C/s_\alpha
\]
control its first two terms.  For the cross term,
\[
 N_B\kappa_\otimes(P_B)\le C,
 \qquad
 \|\mathcal G_U\|\le C/(s_\alpha W)\le C/s_\alpha.
\]
The cases with an empty subgraph follow from V21--V23 directly.

On the globally protected line the shared operator is
\[
 \frac{Nz}{1-z}P_E.
\]
The retained full output is nontrivial outside that line, so
\(z/(1-z)\le C/s_\alpha\), while
\(N\kappa_\otimes(P_E)\le C\) by
\eqref{eq:V22-count-capacity}.
\end{proof}

\section{A two-scale local product estimate}
\label{sec:V29-two-scale}

Let
\[
 M_0:=\max_ra_r.
\]
For every quartic tree,
\begin{equation}
 L_\tau\le A_0\sqrt{M_0}\le A_0^{3/2}.
\label{eq:V29-L-vs-A0}
\end{equation}

\begin{lemma}[Small local product retains the tree gain]
\label{lem:V29-small-product}
If \(A_0<s_\alpha^{-1}\), then
\begin{align}
 s_\alpha^2L_\tau&\le A_0,
\label{eq:V29-small-tree}\\
 A_{U_e}&\le C A_0,
\label{eq:V29-small-fusion}\\
 \frac{A_{U_e}}{1-q_{\alpha,U_e}}
 \|K_{\tau,U_e}\|
 &\le Cs_\alpha^2L_\tau.
\label{eq:V29-small-junction-payer}
\end{align}
\end{lemma}

\begin{proof}
Equation \eqref{eq:V29-L-vs-A0} gives
\[
 s_\alpha^2L_\tau
 \le s_\alpha^2A_0^{3/2}
 <s_\alpha^{3/2}A_0
 \le A_0.
\]
Fusion gives
\[
 A_{U_e}\le C\sum_ra_r\le4CA_0,
\]
because \(A_0\ge a_r\) for every row.  Hence
\(s_\alpha A_{U_e}\le C\) in the small-product regime.  The
representation-independent junction-payer estimate
\eqref{eq:V14-tree-junction-payer} then gives
\eqref{eq:V29-small-junction-payer}.
\end{proof}

\begin{theorem}[Complete local-dominant incidence-two payment]
\label{thm:V29-local-dominant}
For \(G=SU(2)\), every unique-junction physical block with
incidence-two shared coordinates and a row satisfying
\eqref{eq:V29-local-dominant} obeys
\begin{equation}
\boxed{
 \textup{normalized complete physical carrier}
 \le C_DA_0
 \le16C_D\mathcal W_{\tau_c,e}.}
\label{eq:V29-local-dominant-bound}
\end{equation}
This holds both inside and outside output priority.
\end{theorem}

\begin{proof}
First work in output priority.  The V8 nonprivate allocation splits the
sole output resolvent into the junction-output payer and the complete
shared-output payer.

Suppose \(A_0\ge s_\alpha^{-1}\).  The complete local junction-payer
cap \eqref{eq:V3-junction-payer-cap} is at most
\[
 C/s_\alpha\le CA_0.
\]
For the shared-output payer, use the complete local cumulant norm
\(\|K_4\|\le C\), the one-sided insertion theorem, and
\eqref{eq:V29-unmarked-shared}; this is also at most \(C/s_\alpha\le
CA_0\).  The globally protected line is included in
\Cref{lem:V29-unmarked-shared}.

Now suppose \(A_0<s_\alpha^{-1}\).  Sum the sixteen tree bounds
\eqref{eq:V29-small-junction-payer}; all shared moments are
contractions, so the junction-payer contribution is at most
\[
 Cs_\alpha^2L_\tau\le CA_0.
\]
For the shared-output payer, use the local tree estimate
\(\|K_{\tau,U_e}\|\le Cs_\alpha^3L_\tau\) and
\eqref{eq:V29-unmarked-shared}.  One-sided insertion gives
\[
 Cs_\alpha^2L_\tau\le CA_0.
\]
Again this includes the protected external carrier.  Summing the
fixed tree family proves the output-priority claim.

Outside output priority, the private bank pays the global resolvent
and leaves \(C_D/s_\alpha\) times the complete local cumulant, exactly
as in \Cref{lem:V9-private-pays-global}.  If
\(A_0\ge s_\alpha^{-1}\), use \(\|K_4\|\le C\).  If
\(A_0<s_\alpha^{-1}\), retain the tree estimate and use
\eqref{eq:V29-small-tree}.  Shared spectators are contractions in both
cases.  This gives \(C_DA_0\), and
\eqref{eq:V29-star-controls-product} finishes the proof.
\end{proof}

\begin{corollary}[Full unique-junction incidence-two sector]
\label{cor:V29-full-incidence-two}
For \(SU(2)\), every unique-junction physical block whose shared
coordinates have incidence exactly two obeys the quartic marked-tree
atlas, uniformly in support, spins, outputs, and multiplicities.
\end{corollary}

\begin{proof}
If some row is local-dominant, use
\Cref{thm:V29-local-dominant}.  Otherwise every row is all-nonlocal,
and \Cref{cor:V28-full-incidence-two} applies.
\end{proof}

\begin{assessment}[The incidence-two unique-junction programme is
complete]
\label{ass:V29-incidence-two-complete}
V29 closes the local-dominant complement of V28.  Thus the entire
\(SU(2)\) unique-junction incidence-two physical carrier is closed,
including Cartan and non-Cartan outputs, protected and gapped shared
sectors, and every output-priority branch.  The next structural
increase is shared incidence three.
\end{assessment}

\begin{openproblem}[V30 mandate: audit and incidence-three
coordinates]
\label{op:V30-mandate}
\begin{enumerate}
\item Perform the mandatory read-only audit of the newest SM and NS
      companion notes.
\item Fix one shared coordinate active on exactly three rows and keep
      its complete positive output moment.
\item Identify its protected invariant subspace and compute valid
      vertex-cut capacity bounds; do not factor it into pair
      capacities.
\item Separate fully gapped outputs from invariant outputs before
      applying a resolvent.
\item Determine the analogue of the balanced/unbalanced input-output
      dichotomy for \(SU(2)\) triple fusion.
\item Close one incidence-three bank or isolate the exact
      representation-theoretic capacity tensor still missing.
\end{enumerate}
\end{openproblem}

\begin{firewall}[Claim boundary after compact V29]
\label{fire:V29-current-boundary}
V29 closes the full \(SU(2)\) unique-junction incidence-two sector.
It does not close incidence-three or incidence-four shared
coordinates, general compact simple groups, the full embedded
\([4]\) sector, multiple junctions, normalized collision aggregation,
\([3,3]\), global quartic or all-order MTA, continuum Yang--Mills
existence, or a positive mass gap.
\end{firewall}

\section{Append-only record for compact V29}
\label{sec:V29-record}

V29 was appended to the complete compact V28 file.  The exact V28
parent had \(10946\) newline-delimited source records, \(359406\)
bytes, and SHA--256 digest
\[
\texttt{d8cde2b9dacb1fb4b57dc7a83d94cb46d63b9f351d743c4f1436a7ff8b48f350}.
\]
No compact V28 mathematical content was changed.  The sole final
\verb|\end{document}| is restored below.

% =============================================================================
% END APPENDED COMPACT VERSION V29 MATERIAL
% =============================================================================

% The compact V29 document terminator is intentionally disabled here:
% \end{document}

% =============================================================================
% BEGIN APPENDED COMPACT VERSION V30 MATERIAL
% =============================================================================

\chapter{V30: Companion audit and invariant incidence-three
hypergraphs}
\label{ch:V30}

\section{Mandatory companion audit}
\label{sec:V30-audit}

The newest active companion notes inspected at the V30 audit were SM
V36 and NS V31.  Their exact read-only records were:
\begin{center}
\begin{tabular}{p{0.28\linewidth}p{0.12\linewidth}p{0.12\linewidth}
                p{0.38\linewidth}}
\toprule
file & lines & bytes & SHA--256\\
\midrule
\texttt{Renewal\_SM\_Einstein\_Continuum\_Research\_Notes\_V36.tex}
& \(13110\) & \(454809\)
& \texttt{84dd3dbb74c2c5482faffd21c770d3b6415ec74821d18b779ea8d4d2565ecb51}\\
\texttt{Renewal\_NS\_Stability\_Research\_Notes\_V31.tex}
& \(23053\) & \(887537\)
& \texttt{f1121b62238ad5491bb7cf03635ea9934942edf573ed8faaa9ee79e1d38a6696}\\
\bottomrule
\end{tabular}
\end{center}

SM V35--V36 keeps one global buffered spectral carrier through all
gap crossings and charges an exact crossing occupation, rather than
multiplying losses at separate switches.  V36 proves a coarea formula
for the regular crossing occupation and isolates stationary branches
as a separate critical term.  The incidence-three YM analogue is
proof-order only: keep the complete invariant hypergraph projection
together and separate its protected invariant tensor before applying
a Wilson resolvent.

NS V31 is unchanged.  Its normalized-mark and continuation-stock
firewalls again prohibit treating a probability or capacity as an
unbounded first moment.  No SM or NS estimate is imported below.

\begin{assessment}[V30 audit decision]
\label{ass:V30-audit-decision}
An \(SU(2)\) triple invariant is not a product of pair singlets.
It is one three-row tensor with flat one-row marginals.  The valid
global estimate is therefore a vertex-cut Schmidt bound on the entire
incidence-three hypergraph, followed by complete Wilson-failure
summation.  This is the direct higher-incidence analogue of V22--V23.
\end{assessment}

\section{One \(SU(2)\) triple invariant}
\label{sec:V30-one-invariant}

Let \(p,q,r>0\) be half-integers and put
\[
 \mathcal H_{pqr}:=V_p\otimes V_q\otimes V_r.
\]

\begin{lemma}[Multiplicity and flat marginals of a triple invariant]
\label{lem:V30-triple-invariant}
The invariant space
\(\operatorname{Inv}(\mathcal H_{pqr})\) is nonzero exactly when
\begin{equation}
 |p-q|\le r\le p+q,
 \qquad
 p+q+r\in\mathbb Z.
\label{eq:V30-triangle}
\end{equation}
When nonzero it is one-dimensional.  If
\(\Omega_{pqr}\) is its normalized vector and
\(P_{pqr}:=|\Omega_{pqr}\rangle\langle\Omega_{pqr}|\), then
\begin{align}
 \operatorname{Tr}_{qr}P_{pqr}&=d_p^{-1}I_{V_p},
\label{eq:V30-flat-p}\\
 \operatorname{Tr}_{pr}P_{pqr}&=d_q^{-1}I_{V_q},
\label{eq:V30-flat-q}\\
 \operatorname{Tr}_{pq}P_{pqr}&=d_r^{-1}I_{V_r},
\label{eq:V30-flat-r}
\end{align}
where \(d_j=2j+1\).
\end{lemma}

\begin{proof}
The Clebsch--Gordan decomposition
\[
 V_p\otimes V_q
 =
 \bigoplus_{J=|p-q|}^{p+q}V_J
\]
is multiplicity-free.  A trivial representation occurs in
\((V_p\otimes V_q)\otimes V_r\) exactly when the intermediate spin is
\(J=r\), which is equivalent to \eqref{eq:V30-triangle}; its
multiplicity is one.

The invariant identity
\[
 (D^p(g)\otimes D^q(g)\otimes D^r(g))\Omega_{pqr}
 =\Omega_{pqr}
\]
shows that the first reduced density commutes with the irreducible
representation \(D^p\).  Schur's lemma makes it a scalar multiple of
the identity, and trace one fixes that scalar as \(d_p^{-1}\).
The other two formulas follow cyclically.
\end{proof}

\begin{corollary}[Three valid bipartite capacity bounds]
\label{cor:V30-one-invariant-capacity}
\begin{equation}
\boxed{
 \kappa_\otimes(P_{pqr})
 \le
 \min\{d_p^{-1},d_q^{-1},d_r^{-1}\}.}
\label{eq:V30-one-invariant-capacity}
\end{equation}
\end{corollary}

\begin{proof}
Across the cut \(p\,|\,qr\), the squared largest Schmidt coefficient
is \(d_p^{-1}\) by \eqref{eq:V30-flat-p}.  A three-row product vector
is a special bipartite product vector across that cut.  Repeat for the
other two cuts.
\end{proof}

\section{Invariant projections on a four-row hypergraph}
\label{sec:V30-hypergraph}

Let \(F\) be a nonempty multiset of incidence-three spectator
coordinates on four rows.  Each \(f\) has an active triple
\(I_f\subset\{1,2,3,4\}\), \(|I_f|=3\), and nontrivial spins
\((p_{i,f})_{i\in I_f}\).  On a subset \(A\subseteq F\) for which all
triple invariants exist, put
\[
 P_A:=
 \left(\bigotimes_{f\in A}P_f\right)
 \otimes I_{F\setminus A},
\qquad
 D_{i,A}:=
 \prod_{\substack{f\in A\\i\in I_f}}d_{p_{i,f}}.
\]

\begin{theorem}[Vertex-cut capacity of an invariant hypergraph]
\label{thm:V30-hypergraph-capacity}
For every such \(A\),
\begin{equation}
\boxed{
 \kappa_\otimes(P_A)
 \le
 \min_{1\le i\le4}D_{i,A}^{-1}
 \le
 \rho_3^{|A|},
 \qquad
 \rho_3:=8^{-1/4}<1.}
\label{eq:V30-hypergraph-capacity}
\end{equation}
\end{theorem}

\begin{proof}
Fix row \(i\) and group the other three rows.  At every hyperedge
containing \(i\), \Cref{lem:V30-triple-invariant} says that its
invariant tensor has \(d_{p_{i,f}}\) equal Schmidt coefficients across
this cut.  Hyperedges not containing \(i\) lie entirely on the other
side.  Tensoring gives largest squared Schmidt coefficient
\(D_{i,A}^{-1}\), and four-row product vectors form a subset of the
bipartite product vectors.

Each hyperedge contributes its three active dimensions to
\(\prod_iD_{i,A}\).  Since every active representation is nontrivial,
\[
 \prod_{i=1}^4D_{i,A}
 =
 \prod_{f\in A}\prod_{i\in I_f}d_{p_{i,f}}
 \ge8^{|A|}.
\]
Thus some \(D_{i,A}\ge8^{|A|/4}\), proving the second inequality.
\end{proof}

\begin{remark}[No trivalent tensor factorization]
\label{rem:V30-no-factorization}
The proof never writes \(P_f\) as a product of pair projections and
never multiplies capacities of overlapping hyperedges.  It uses only
the exact Schmidt spectrum across one valid row cut at a time.
\end{remark}

\section{The complete incidence-three Wilson moment}
\label{sec:V30-complete-moment}

At \(f\), let
\[
 M_f^{(t)}
 :=
 \sum_{U,a}q_{\alpha,U}^{\,t}P_{U,a},
 \qquad t>0,
\]
be the complete physical output moment on the triple tensor product.
If the triple invariant exists, denote its rank-one projection by
\(P_f\); otherwise put \(P_f=0\).  Let
\[
 M_F^{(t)}:=\bigotimes_{f\in F}M_f^{(t)}.
\]
If every \(P_f\ne0\), put \(P_*:=\bigotimes_fP_f\); otherwise
\(P_*:=0\).  Thus \(P_*\) is precisely the shared eigenspace with
multiplier one.

\begin{lemma}[Power capacity on the trivalent hypergraph]
\label{lem:V30-power-capacity}
Let \(N=|F|\),
\[
 r_t:=(1-c_*s_\alpha)^t,
\]
and define, with \(\rho=\rho_3\),
\begin{align}
 \Phi_{N,\rho}^{(0)}(r)
 &:=[\rho+(1-\rho)r]^N-[\rho(1-r)]^N,
\label{eq:V30-Phi0}\\
 \Phi_{N,\rho}^{(1)}(r)
 &:=r[\rho+(1-\rho)r]^{N-1}.
\label{eq:V30-Phi1}
\end{align}
Then
\begin{equation}
\boxed{
 \kappa_\otimes(M_F^{(t)}-P_*)
 \le
 \begin{cases}
 \Phi_{N,\rho_3}^{(0)}(r_t),&P_*\ne0,\\
 \Phi_{N,\rho_3}^{(1)}(r_t),&P_*=0.
 \end{cases}}
\label{eq:V30-power-capacity}
\end{equation}
\end{lemma}

\begin{proof}
Every nontrivial triple output obeys the uniform separation of
\Cref{lem:V19-nontrivial-separation}.  Hence
\[
 M_f^{(t)}
 \preccurlyeq
 \begin{cases}
 r_tI+(1-r_t)P_f,&P_f\ne0,\\
 r_tI,&P_f=0.
 \end{cases}
\]
Expand the tensor product.  If every invariant exists, subtract the
all-invariant projection and rewrite the result as the positive sum
over proper invariant subsets, as in
\eqref{eq:V21-power-majorant}.  Apply
\eqref{eq:V30-hypergraph-capacity} to each subset and sum the binomial
polynomial, giving \(\Phi^{(0)}\).

If at least one invariant is absent, retain one factor \(r_t\), expand
the remaining invariant coordinates, and use
\[
 r_t^K[\rho_3+(1-\rho_3)r_t]^{N-K}
 \le
 r_t[\rho_3+(1-\rho_3)r_t]^{N-1}.
\]
\end{proof}

\begin{lemma}[Trivalent one-failure sum]
\label{lem:V30-one-failure-sum}
For \(\nu=0,1\),
\begin{equation}
\boxed{
 \sum_{n\ge1}
 \Phi_{N,\rho_3}^{(\nu)}
 ((1-c_*s_\alpha)^n)
 \le
 C e^{-cs_\alpha N}
 \left(1+\frac1{s_\alpha N}\right)
 \le\frac{C}{s_\alpha N}.}
\label{eq:V30-one-failure-sum}
\end{equation}
\end{lemma}

\begin{proof}
The proof of \Cref{lem:V23-analytic-sum} uses only
\(0<\rho<1\).  For completeness, put
\(A(r)=\rho_3+(1-\rho_3)r\) and
\(B(r)=\rho_3(1-r)\).  For \(\Phi^{(1)}\),
\[
 \int_0^1\frac{\Phi^{(1)}(r)}r\,dr
 =
 \int_0^1A(r)^{N-1}\,dr\le C/N.
\]
For \(\Phi^{(0)}\), use \(A(r)-B(r)=r\), the mean-value theorem on
\((0,1/2]\), and direct integration of \(2A(r)^N\) on
\([1/2,1]\), again obtaining \(C/N\).  Geometric interval comparison
for \(r_n=(1-c_*s_\alpha)^n\), with the first interval retained,
gives the exponential refinement exactly as in
\eqref{eq:V23-integer-sum}.
\end{proof}

\section{A complete trivalent shared resolvent}
\label{sec:V30-resolvent}

For an output tuple \(\boldsymbol U\), write
\[
 q_{\boldsymbol U}:=\prod_fq_{\alpha,U_f},
 \quad
 Z_{\boldsymbol U}:=\sum_fC_2(U_f),
 \quad
 H_{\boldsymbol U}:=N+Z_{\boldsymbol U}.
\]
On the complement of \(P_*\), define
\[
 \mathcal R_F^{(3)}
 :=
 \sum_{q_{\boldsymbol U}<1}
 \frac{H_{\boldsymbol U}q_{\boldsymbol U}}
      {1-q_{\boldsymbol U}}P_{\boldsymbol U}.
\]

\begin{theorem}[Unmarked complete incidence-three resolvent]
\label{thm:V30-trivalent-resolvent}
Uniformly in the incidence-three hypergraph, its spins, all outputs,
and multiplicities,
\begin{equation}
\boxed{
 \kappa_\otimes(\mathcal R_F^{(3)})
 \le C/s_\alpha.}
\label{eq:V30-trivalent-resolvent}
\end{equation}
\end{theorem}

\begin{proof}
Split the output mass as \(H_{\boldsymbol U}=N+Z_{\boldsymbol U}\).
Geometric summation and
\Cref{lem:V30-power-capacity,lem:V30-one-failure-sum} give
\[
 \kappa_\otimes\!\left(
 \sum_{q_{\boldsymbol U}<1}
 \frac{q_{\boldsymbol U}}{1-q_{\boldsymbol U}}
 P_{\boldsymbol U}\right)
 \le\frac{C}{s_\alpha N}.
\]
Multiplication by \(N\) costs \(C/s_\alpha\).

For the Casimir part, apply the \(m=1\) retained-failure compiler to
\[
 x_f=s_\alpha C_2(U_f),
 \qquad
 r_f=q_{\alpha,U_f}.
\]
It gives
\[
 \frac{Z_{\boldsymbol U}q_{\boldsymbol U}}
      {1-q_{\boldsymbol U}}
 \le C/s_\alpha
\]
on every nonprotected spectral block.  Hence that positive operator
is bounded by \(C s_\alpha^{-1}(I-P_*)\).  Add the two capacity
bounds.
\end{proof}

\begin{corollary}[Exact external split of the invariant line]
\label{cor:V30-external-split}
For \(0\le z<1\), the complete externally damped trivalent carrier
obeys
\begin{equation}
\boxed{
 \kappa_\otimes\!\left(
 \sum_{\boldsymbol U}
 \frac{H_{\boldsymbol U}zq_{\boldsymbol U}}
      {1-zq_{\boldsymbol U}}P_{\boldsymbol U}\right)
 \le
 \frac{C}{s_\alpha}
 C\frac{z}{1-z}.}
\label{eq:V30-external-split}
\end{equation}
\end{corollary}

\begin{proof}
On nonprotected tuples,
\[
 \frac{zq}{1-zq}\le\frac q{1-q}.
\]
On \(P_*\), the coefficient is \(Nz/(1-z)\).  By
\eqref{eq:V30-hypergraph-capacity},
\[
 N\kappa_\otimes(P_*)\le N\rho_3^N\le C.
\]
Apply \Cref{thm:V30-trivalent-resolvent}.
\end{proof}

\begin{assessment}[The incidence-three aggregation gate is closed]
\label{ass:V30-aggregation-closed}
V30 proves that protected triple invariants and every gapped triple
output may be summed inside one complete hypergraph moment and one
shared resolvent.  Neither the number of hyperedges nor the number of
triple fusion outputs appears.  What remains is not aggregation but
the conversion of input row masses into legal marked-tree weights.
\end{assessment}

\begin{openproblem}[V31 mandate: triple-fusion input marks]
\label{op:V31-mandate}
\begin{enumerate}
\item At one incidence-three coordinate, order the input spins and use
      the triangle condition to identify the two necessarily
      comparable large inputs on an invariant channel.
\item For noninvariant outputs, split between retained output Casimir
      and two-input bilinear mass, as in V20.
\item Define the legal sum of pair-edge marks available at a
      three-row coordinate without using the coordinate more than once
      in a tree.
\item Prove a marked refinement of
      \eqref{eq:V30-trivalent-resolvent}.
\item Insert the result first below output priority and keep the
      incidence-two sector as an already-closed contraction.
\item Do not replace the triple invariant capacity by a product of
      pair-singlet capacities.
\end{enumerate}
\end{openproblem}

\begin{firewall}[Claim boundary after compact V30]
\label{fire:V30-current-boundary}
V30 closes complete-output aggregation and the unmarked shared
resolvent for \(SU(2)\) incidence-three spectator banks.  It does not
yet prove their marked input payment, incidence-four shared
coordinates, general compact simple groups, the full embedded
\([4]\) sector, multiple junctions, normalized collision aggregation,
\([3,3]\), global quartic or all-order MTA, continuum Yang--Mills
existence, or a positive mass gap.
\end{firewall}

\section{Append-only record for compact V30}
\label{sec:V30-record}

V30 was appended to the complete compact V29 file.  The exact V29
parent had \(11197\) newline-delimited source records, \(367571\)
bytes, and SHA--256 digest
\[
\texttt{d5026b9796913b3c93790adb0683d1f7413ed8d8f0fb4891ad3be017aca1b177}.
\]
No compact V29 mathematical content was changed.  The sole final
\verb|\end{document}| is restored below.

% =============================================================================
% END APPENDED COMPACT VERSION V30 MATERIAL
% =============================================================================

% The compact V30 document terminator is intentionally disabled here:
% \end{document}

% =============================================================================
% BEGIN APPENDED COMPACT VERSION V31 MATERIAL
% =============================================================================

\chapter{V31: Triple-fusion marks and the incidence-three shared payer}
\label{ch:V31}

\section{The trivalent balanced/unbalanced dichotomy}
\label{sec:V31-dichotomy}

At one incidence-three coordinate, order the three input spins as
\[
 j_1\ge j_2\ge j_3>0,
\qquad
 A_\nu:=1+j_\nu(j_\nu+1).
\]
Put \(w:=A_1\).

\begin{lemma}[Either a legal pair is large or every output is large]
\label{lem:V31-triple-dichotomy}
Exactly one of the following alternatives applies:
\begin{enumerate}[label=\textup{(\mathrm B)},leftmargin=4em]
\item \(j_2+j_3\ge j_1/2\).  Then
      \begin{equation}
      \boxed{
      A_2\ge\frac1{16}A_1,
      \qquad
      A_1A_2\ge\frac1{16}w^2.}
      \label{eq:V31-balanced-triple}
      \end{equation}
\end{enumerate}
or
\begin{enumerate}[label=\textup{(\mathrm U)},leftmargin=4em]
\item \(j_2+j_3<j_1/2\).  Every total output spin \(J\) in
      \(V_{j_1}\otimes V_{j_2}\otimes V_{j_3}\) then satisfies
      \begin{equation}
      \boxed{1+J(J+1)\ge\frac14w.}
      \label{eq:V31-unbalanced-triple}
      \end{equation}
\end{enumerate}
Every coordinate admitting a triple invariant belongs to
\textup{(B)}.
\end{lemma}

\begin{proof}
In case \textup{(B)}, \(j_2\ge j_3\) implies
\(2j_2\ge j_2+j_3\ge j_1/2\), hence \(j_2\ge j_1/4\).
Monotonicity and
\[
 1+\frac{j_1}{4}\left(\frac{j_1}{4}+1\right)
 \ge
 \frac1{16}\bigl(1+j_1(j_1+1)\bigr)
\]
prove \eqref{eq:V31-balanced-triple}.

In case \textup{(U)}, repeated Clebsch--Gordan triangle inequalities
give
\[
 J\ge j_1-j_2-j_3>j_1/2.
\]
The comparison used in
\Cref{lem:V20-unbalanced-retention} proves
\eqref{eq:V31-unbalanced-triple}.

Finally, an invariant triple obeys
\(j_1\le j_2+j_3\) by \eqref{eq:V30-triangle}, so it is balanced.
\end{proof}

\begin{definition}[One legal edge per balanced triple coordinate]
\label{def:V31-legal-pair}
On a balanced coordinate \(f\), let \(u_f,v_f\) be the rows carrying
the two largest spins.  Define
\[
 w_f:=A_{1,f},
 \qquad
 E_f:=A_{u_f,f}A_{v_f,f}.
\]
The coordinate is assigned only to the pair \(u_fv_f\).  Thus a
marked tree uses \(f\) at most once, even though \(f\) is active on
three rows.
\end{definition}

\section{Marked complete moments on the balanced hypergraph}
\label{sec:V31-balanced-hypergraph}

Let \(B\) be a nonempty bank of balanced incidence-three coordinates.
Put
\begin{equation}
 N_B:=|B|,
 \qquad
 S_B:=\sum_{f\in B}w_f,
 \qquad
 H_B:=\sum_{f\in B}E_f.
\label{eq:V31-balanced-stocks}
\end{equation}
By \eqref{eq:V31-balanced-triple},
\begin{equation}
\boxed{S_B^2\le16N_BH_B.}
\label{eq:V31-balanced-stock-comparison}
\end{equation}
Use the complete moment \(M_B\), protected projection \(P_B\), and
gapped resolvent \(\mathcal R_B^{(3)}\) from V30, restricted to this
bank.

\begin{lemma}[Half-power trivalent failure sum]
\label{lem:V31-half-power}
With \(\rho_3=8^{-1/4}\), the two functions in
\eqref{eq:V30-Phi0}--\eqref{eq:V30-Phi1} satisfy
\begin{equation}
\boxed{
 \sum_{n\ge0}
 \Phi_{N_B,\rho_3}^{(\nu)}
 ((1-c_*s_\alpha)^{n+1/2})
 \le
 C e^{-cs_\alpha N_B}
 \left(1+\frac1{s_\alpha N_B}\right),
 \quad \nu=0,1.}
\label{eq:V31-half-power}
\end{equation}
\end{lemma}

\begin{proof}
Repeat the geometric interval comparison in
\Cref{lem:V30-one-failure-sum}, starting at
\((1-c_*s_\alpha)^{1/2}\).  The integral estimate depends only on
\(0<\rho_3<1\), and the omitted first half interval is bounded by the
same exponential endpoint term.
\end{proof}

\begin{theorem}[Marked trivalent moment and resolvent]
\label{thm:V31-marked-trivalent}
Uniformly in the balanced hypergraph and all its representations,
\begin{align}
 \boxed{
 s_\alpha S_B^2\kappa_\otimes(M_B)\le CH_B,}
\label{eq:V31-balanced-moment}\\
 \boxed{
 N_BS_B^2\kappa_\otimes(P_B)\le CH_B,}
\label{eq:V31-balanced-protected}\\
 \boxed{
 s_\alpha^2S_B^2
 \kappa_\otimes(\mathcal R_B^{(3)})\le CH_B.}
\label{eq:V31-balanced-resolvent}
\end{align}
The middle line is interpreted as zero when \(P_B=0\).
\end{theorem}

\begin{proof}
At \(t=1\),
\Cref{lem:V30-power-capacity} and the endpoint estimate in
\Cref{lem:V30-one-failure-sum} give
\[
 \kappa_\otimes(M_B-P_B)
 \le C\min\{1,e^{-cs_\alpha N_B}\}.
\]
Together with
\eqref{eq:V31-balanced-stock-comparison},
\[
 s_\alpha S_B^2\kappa_\otimes(M_B-P_B)
 \le
 C(s_\alpha N_B)
 \min\{1,e^{-cs_\alpha N_B}\}H_B
 \le CH_B.
\]

When \(P_B\ne0\), every coordinate admits an invariant.  The
hypergraph cut bound gives
\(\kappa_\otimes(P_B)\le\rho_3^{N_B}\).  Hence
\[
 \frac{N_BS_B^2}{H_B}\kappa_\otimes(P_B)
 \le16N_B^2\rho_3^{N_B}\le C.
\]
This proves the protected estimate and completes the first estimate
because \(s_\alpha/N_B\le1\).

For the resolvent, split
\[
 \mathcal R_B^{(3)}
 =
 N_B\mathcal G_B+\mathcal C_B
\]
into its geometric and output-Casimir parts.  The proof of
\Cref{thm:V30-trivalent-resolvent} and the exponential refinement in
\eqref{eq:V30-one-failure-sum} give
\[
 \kappa_\otimes(\mathcal G_B)
 \le
 Ce^{-cs_\alpha N_B}
 \left(1+\frac1{s_\alpha N_B}\right).
\]
For a nonprotected output tuple, the \(m=2\) retained-failure compiler,
with
\[
 X=s_\alpha\sum_fC_2(U_f),
 \qquad
 Q=\prod_fq_{\alpha,U_f},
\]
gives \(X^2Q\le C(1-Q)\), and therefore
\[
 \frac{(\sum_fC_2(U_f))Q}{1-Q}
 \le
 \frac{C}{s_\alpha}
 \left(\sum_{n\ge0}Q^{n+1/2}\right).
\]
Apply
\Cref{lem:V30-power-capacity,lem:V31-half-power} to the positive
operator on the right.  Also retain the direct \(C/s_\alpha\)
operator-norm bound from the \(m=1\) compiler.

Now put \(x=s_\alpha N_B\) and use
\(S_B^2\le16N_BH_B\).  If \(x\le1\), the unmarked bounds
\[
 \kappa_\otimes(\mathcal G_B)\le C/(s_\alpha N_B),
 \qquad
 \kappa_\otimes(\mathcal C_B)\le C/s_\alpha
\]
give
\[
 s_\alpha^2S_B^2
 \kappa_\otimes(N_B\mathcal G_B+\mathcal C_B)
 \le CxH_B.
\]
If \(x>1\), the integer and half-integer exponential bounds give a
multiple of
\[
 [x^2+x]e^{-cx}(1+x^{-1})H_B.
\]
Both regimes are uniformly bounded, proving
\eqref{eq:V31-balanced-resolvent}.
\end{proof}

\section{The unbalanced trivalent bank}
\label{sec:V31-unbalanced-bank}

Let \(U\) be a bank of unbalanced incidence-three coordinates.  Put
\[
 W:=\sum_{f\in U}w_f.
\]
For an output tuple, let \(Q_{\boldsymbol U}\) be its multiplier
product and \(B_{\boldsymbol U}\) its total output mass.

\begin{theorem}[Direct unbalanced trivalent estimates]
\label{thm:V31-unbalanced-trivalent}
Uniformly in the complete unbalanced bank,
\begin{align}
 (s_\alpha W)^k\|M_U\|&\le C_k,
 &&1\le k\le3,
\label{eq:V31-unbalanced-moments}\\
 \|\mathcal G_U\|&\le\frac{C_k}{(s_\alpha W)^k},
 &&1\le k\le3,
\label{eq:V31-unbalanced-geometric}\\
 \boxed{
 \|\mathcal R_U^{(3)}\|
 \le
 \min\left\{\frac{C}{s_\alpha},
             \frac{C}{s_\alpha^3W^2}\right\}.}
\label{eq:V31-unbalanced-resolvent}
\end{align}
\end{theorem}

\begin{proof}
By \eqref{eq:V31-unbalanced-triple}, every output mass at \(f\) is at
least \(w_f/4\).  The fusion upper bound makes it at most \(Cw_f\).
Thus the proof of
\Cref{thm:V24-unbalanced-resolvent} applies verbatim with triple
outputs: ordinary product moments give the first line, retained
failure moments give the second, and
\[
 cW\le B_{\boldsymbol U}\le CW
\]
converts the \(k=1,3\) geometric estimates into the last line.
\end{proof}

\section{Legal routing of balanced triple coordinates}
\label{sec:V31-routing}

For \(f\in B\), let \(u_fv_f\) be the assigned pair in
\Cref{def:V31-legal-pair}.  Let \(\sigma_{u_fv_f}\) be the comparison
path whose \(u_fv_f\)-edge is marked at \(f\), while its other two
edges are marked at the junction.

\begin{lemma}[One-use trivalent legal path sum]
\label{lem:V31-trivalent-routing}
\begin{equation}
\boxed{
 \mathfrak P_B^{(3)}
 :=
 \sum_{f\in B}
 \mathcal W_{\sigma_{u_fv_f},\ell_f}
 =
 A_0\sum_{f\in B}
 \frac{E_f}{\Xi_{u_f}\Xi_{v_f}}
 \ge
 \frac{A_0H_B}{\Xi_{\max}^2}.}
\label{eq:V31-trivalent-routing}
\end{equation}
\end{lemma}

\begin{proof}
At coordinate \(f\), both assigned rows are active, so \(f\) is a
legal mark for their tree edge.  The coordinate occurs in exactly one
summand by construction.  The marked-weight formula gives the
equality, and
\(\Xi_{u_f}\Xi_{v_f}\le\Xi_{\max}^2\) gives the inequality.
\end{proof}

\section{Insertion below output priority}
\label{sec:V31-below-output}

\begin{theorem}[All-nonlocal incidence-three closure below
\textup{(O)}]
\label{thm:V31-below-output}
Fix \(D\ge16\).  Suppose every shared spectator coordinate has
incidence exactly three, all rows are nonlocal, and
\(\Xi_T\le DY\).  Then the complete unique-junction \(SU(2)\) carrier
obeys the quartic marked-tree atlas.
\end{theorem}

\begin{proof}
If every tree satisfies \(B_\tau\le DY^2\), use the V7 private
payment.  Otherwise one failed tree selects a globally maximal row
with
\[
 h_i>\Xi_{\max}/2.
\]
Every input mass at every coordinate is at most its \(w_f\), so
\[
 S_B+W>\Xi_{\max}/2.
\]

The private bank pays the global resolvent and leaves
\(Cs_\alpha^2L_\tau\), as in
\Cref{lem:V19-private-leaves-two}.  If
\(S_B>\Xi_{\max}/4\), insert the balanced complete moment and use
\eqref{eq:V31-balanced-moment}:
\[
 Cs_\alpha^2L_\tau\kappa_\otimes(M_B)
 \le
 Cs_\alpha L_\tau\frac{H_B}{S_B^2}.
\]
Paths and stars with \(s_\alpha L_\tau\le A_0\) are paid by
\eqref{eq:V31-trivalent-routing}.  If a star is hot, return to the
complete local cumulant.  The private payer leaves \(C/s_\alpha\);
\eqref{eq:V31-balanced-moment} gives
\[
 \frac{C}{s_\alpha}\kappa_\otimes(M_B)
 \le
 \frac{CH_B}{s_\alpha^2S_B^2}.
\]
The hot-star calculation gives \(A_0>s_\alpha^{-2}\), so the legal
path sum pays this term.

If \(W\ge\Xi_{\max}/4\), use
\eqref{eq:V31-unbalanced-moments} with \(k=2\):
\[
 Cs_\alpha^2L_\tau\|M_U\|
 \le\frac{CL_\tau}{W^2}
 \le\frac{CL_\tau}{B_\tau}
 \le C\mathcal W_{\tau,e}.
\]
The balanced part, when present, is a contraction.  The two stock
alternatives are exhaustive.
\end{proof}

\section{The output-priority shared payer}
\label{sec:V31-output-priority}

On a mixed balanced/unbalanced trivalent bank, the positive
decomposition
\begin{equation}
 \mathcal R_{B\cup U}^{(3)}
 \preccurlyeq
 \mathcal R_B^{(3)}\otimes M_U
 +M_B\otimes\mathcal R_U^{(3)}
 +N_BP_B\otimes\mathcal G_U
\label{eq:V31-mixed-resolvent}
\end{equation}
is proved by the scalar argument of
\Cref{lem:V24-mixed-decomposition}.  It uses only positivity and
\(1-Q_BQ_U\ge1-Q_B,1-Q_U\).

\begin{theorem}[Complete incidence-three shared-output payer]
\label{thm:V31-shared-output}
Under the hypotheses of
\Cref{thm:V31-below-output}, but in output priority
\(\Xi_T>DY\), the complete shared-output-resolvent payer obeys the
quartic marked-tree atlas.
\end{theorem}

\begin{proof}
First dispose of the globally protected line.  It occurs only when
\(U=\varnothing\), every balanced coordinate admits an invariant, and
every shared output is that invariant.  With
\[
 z=q_{\alpha,U_e}q_P,
\]
its exact external carrier is
\[
 \frac{N_Bz}{1-z}P_B.
\]
If \(B_\tau\le DY^2\), then \(Y>0\),
\[
 \frac{z}{1-z}\le\frac{q_P}{1-q_P},
\]
and the private quadratic retained-failure estimate together with
\(N_B\kappa_\otimes(P_B)\le C\) bounds the locally inserted carrier by
\(C_DL_\tau/B_\tau\).

If one tree fails, choose the globally maximal row.  Then
\(S_B\ge h_i>\Xi_{\max}/2\), while
\eqref{eq:V31-balanced-protected} gives
\[
 N_B\kappa_\otimes(P_B)
 \le\frac{CH_B}{S_B^2}.
\]
The retained full output is nontrivial outside the shared invariant
line, so \(z/(1-z)\le C/s_\alpha\).  The local tree estimate therefore
leaves
\[
 Cs_\alpha^2L_\tau\frac{H_B}{S_B^2}.
\]
Paths and stars satisfying \(s_\alpha^2L_\tau\le A_0\) are paid by
\eqref{eq:V31-trivalent-routing}.  For a hot star, use the complete
local cumulant norm instead.  The remaining factor is
\(CH_B/(s_\alpha S_B^2)\), while the hot inequality gives
\(A_0>s_\alpha^{-4}\ge s_\alpha^{-1}\).  Thus the same legal path sum
pays the globally protected line.

It remains to treat the gapped shared outputs.

For \(B_\tau\le DY^2\), the unmarked bounds
\[
 \kappa_\otimes(\mathcal R_B^{(3)})\le C/s_\alpha,
 \quad
 \|\mathcal R_U^{(3)}\|\le C/s_\alpha,
 \quad
 N_B\kappa_\otimes(P_B)\le C
\]
and the one-sided tensor law applied to
\eqref{eq:V31-mixed-resolvent} give
\[
 \kappa_\otimes(\mathcal R_{B\cup U}^{(3)})
 \le C/s_\alpha.
\]
The local \(s_\alpha^3L_\tau\) factor and private quadratic debit
finish this branch.

If one tree fails, then
\(S_B+W>\Xi_{\max}/2\).  When
\(S_B>\Xi_{\max}/4\), use
\eqref{eq:V31-balanced-resolvent},
\eqref{eq:V31-balanced-moment},
\eqref{eq:V31-balanced-protected}, and the unmarked unbalanced bounds
in the three terms of
\eqref{eq:V31-mixed-resolvent}.  After local insertion they are at
most
\[
 Cs_\alpha L_\tau\frac{H_B}{S_B^2},
 \quad
 Cs_\alpha L_\tau\frac{H_B}{S_B^2},
 \quad
 Cs_\alpha^2L_\tau\frac{H_B}{S_B^2}.
\]
Route paths and balanced stars through
\eqref{eq:V31-trivalent-routing}; for a hot star use the complete
local cumulant and \(A_0>s_\alpha^{-2}\), exactly as in the proof of
\Cref{thm:V24-full-shared-payer}.

When \(W\ge\Xi_{\max}/4\), use the strong unbalanced bounds
\[
 \|M_U\|\le C/(s_\alpha^2W^2),
 \quad
 \|\mathcal R_U^{(3)}\|\le C/(s_\alpha^3W^2),
 \quad
 \|\mathcal G_U\|\le C/(s_\alpha^3W^3).
\]
The three terms of \eqref{eq:V31-mixed-resolvent}, after local
insertion, are at most \(CL_\tau/W^2\).  For the cross term use
\[
 N_B\le S_B\le\sum_rh_r\le4\Xi_{\max}\le16W.
\]
Finally \(B_\tau\le\Xi_{\max}^2\le16W^2\).  This proves the marked
bound.
\end{proof}

\begin{assessment}[The trivalent shared side is closed]
\label{ass:V31-shared-closed}
For pure incidence-three \(SU(2)\) spectator banks, V31 closes the
complete carrier below output priority and closes the entire
shared-output payer inside output priority.  Protected invariants,
balanced noninvariant outputs, strongly unbalanced outputs, and mixed
tuples are all included without resolved-output counting.
\end{assessment}

\begin{openproblem}[V32 mandate: the incidence-three junction payer]
\label{op:V32-mandate}
\begin{enumerate}
\item Reuse the cold/hot junction split of V25 and the cubic complete
      local output moment of V26.
\item Insert the balanced trivalent moment
      \eqref{eq:V31-balanced-moment} into the good/bad hot-output
      partition.
\item On an unbalanced-dominant triple bank, use
      \(A_{U_e}\lesssim W\) from all-nonlocality and the cubic moment
      \eqref{eq:V31-unbalanced-moments}.
\item Apply the all-trees-private argument of V28 to the remaining
      branch.
\item Close the all-nonlocal incidence-three output-priority carrier.
\item Treat incidence-two spectators, if present, only as already
      assembled complete contractions.
\end{enumerate}
\end{openproblem}

\begin{firewall}[Claim boundary after compact V31]
\label{fire:V31-current-boundary}
V31 closes pure incidence-three blocks outside output priority and
their complete shared-output payer inside output priority.  It does
not yet close the incidence-three junction-output payer, mixtures of
incidence two and three requiring simultaneous marks, incidence-four
shared coordinates, general compact simple groups, the full embedded
\([4]\) sector, multiple junctions, normalized collision aggregation,
\([3,3]\), global quartic or all-order MTA, continuum Yang--Mills
existence, or a positive mass gap.
\end{firewall}

\section{Append-only record for compact V31}
\label{sec:V31-record}

V31 was appended to the complete compact V30 file.  The exact V30
parent had \(11641\) newline-delimited source records, \(381224\)
bytes, and SHA--256 digest
\[
\texttt{2bc9bbcdafed042a66743a156687930f2fb718dd1c4386527e46a9640235a15d}.
\]
No compact V30 mathematical content was changed.  The sole final
\verb|\end{document}| is restored below.

% =============================================================================
% END APPENDED COMPACT VERSION V31 MATERIAL
% =============================================================================

% The compact V31 document terminator is intentionally disabled here:
% \end{document}

% =============================================================================
% BEGIN APPENDED COMPACT VERSION V32 MATERIAL
% =============================================================================

\chapter{V32: Full pure incidence-three unique-junction closure}
\label{ch:V32}

\section{The trivalent moment dichotomy after a failed tree}
\label{sec:V32-moment-dichotomy}

\begin{lemma}[Balanced or unbalanced triple stock pays two row
denominators]
\label{lem:V32-triple-stock}
Assume all shared coordinates have incidence three, all rows are
nonlocal, and one tree satisfies \(B_\tau>DY^2\).  Then, with the V31
balanced/unbalanced split, either
\begin{align}
 S_B&>\Xi_{\max}/4,
 &
 s_\alpha S_B^2\kappa_\otimes(M_S)&\le CH_B,
\label{eq:V32-balanced-stock}
\end{align}
or
\begin{align}
 W&\ge\Xi_{\max}/4,
 &
 \|M_S\|&\le C/(s_\alpha^2W^2).
\label{eq:V32-unbalanced-stock}
\end{align}
Here \(M_S=M_B\otimes M_U\).
\end{lemma}

\begin{proof}
Choose a globally maximal row.  Failed branching and all-nonlocality
give \(h_i>\Xi_{\max}/2\).  Every input at a triple coordinate is at
most its \(w_f\), so
\[
 S_B+W>\Xi_{\max}/2.
\]
One stock is at least \(\Xi_{\max}/4\).  In the balanced case use
\eqref{eq:V31-balanced-moment}, the one-sided tensor law, and
\(\|M_U\|\le1\).  In the unbalanced case use
\eqref{eq:V31-unbalanced-moments} with \(k=2\) and
\(\|M_B\|\le1\).
\end{proof}

\section{The cold incidence-three junction payer}
\label{sec:V32-cold}

\begin{theorem}[Cold trivalent junction-output closure]
\label{thm:V32-cold}
Fix \(D\ge16\).  In output priority, every all-nonlocal
unique-junction \(SU(2)\) block whose shared coordinates all have
incidence three has its complete cold junction-output payer bounded by
the quartic marked-tree atlas.
\end{theorem}

\begin{proof}
On \(s_\alpha A_{U_e}\le1\), the local tree payer is
\[
 \frac{A_{U_e}}{1-q_{\alpha,U_e}}
 \|K_{\tau,U_e}\|
 \le Cs_\alpha^2L_\tau.
\]
If \(B_\tau\le DY^2\), the private quadratic moment gives
\[
 Cs_\alpha^2q_PL_\tau
 \le C_DL_\tau/B_\tau
 \le C_D\mathcal W_{\tau,e}.
\]

If one tree fails, use
\Cref{lem:V32-triple-stock}.  In the unbalanced alternative,
\[
 Cs_\alpha^2L_\tau\|M_S\|
 \le CL_\tau/W^2
 \le CL_\tau/B_\tau.
\]
The last comparison uses
\(B_\tau\le\Xi_{\max}^2\le16W^2\).

In the balanced alternative,
\[
 Cs_\alpha^2L_\tau\kappa_\otimes(M_S)
 \le
 Cs_\alpha L_\tau H_B/S_B^2.
\]
Paths and stars satisfying \(s_\alpha L_\tau\le A_0\) are paid by
\eqref{eq:V31-trivalent-routing}.  If a star is hot, return to the
complete positive local payer.  Its \(C/s_\alpha\) norm and
\eqref{eq:V32-balanced-stock} give
\[
 CH_B/(s_\alpha^2S_B^2).
\]
The hot-star inequality implies \(A_0>s_\alpha^{-2}\), so the legal
trivalent path sum pays the result.  As in V25, the complete local
carrier handles the fixed family of remaining star geometries at
once.
\end{proof}

\section{Hot outputs after one failed tree}
\label{sec:V32-hot-failed}

\begin{theorem}[Balanced-dominant hot trivalent junction payer]
\label{thm:V32-balanced-hot}
Under the hypotheses of
\Cref{thm:V32-cold}, suppose one tree fails and
\(S_B>\Xi_{\max}/4\).  Then the complete hot junction-output payer
obeys the legal marked-tree atlas.
\end{theorem}

\begin{proof}
Use the V26 complete local good/bad partition
\[
 s_\alpha^2A_{U_e}\sqrt{M_0}\le\varepsilon
 \quad\hbox{or}\quad
 s_\alpha^2A_{U_e}\sqrt{M_0}>\varepsilon,
\qquad
 M_0=\max_ra_r.
\]
Its proof is local and independent of spectator incidence.  On the
good set,
\[
 s_\alpha^2A_{U_e}L_\tau\le\varepsilon A_0
\]
for every tree.  Insert the hot tree estimate
\[
 \frac{A_{U_e}}{1-q_{\alpha,U_e}}\|K_{\tau,U_e}\|
 \le Cs_\alpha^3A_{U_e}L_\tau
\]
into \eqref{eq:V32-balanced-stock}.  The result is bounded by
\[
 C A_0H_B/\Xi_{\max}^2
\le C\mathfrak P_B^{(3)}.
\]

On the bad set, the cubic complete local output moment of V26 gives
the complete positive local payer norm
\[
 \|\mathcal J_{\rm bad}^{\rm pay}\|
 \le Cs_\alpha A_0.
\]
One-sided insertion into
\eqref{eq:V32-balanced-stock} again gives
\[
 C A_0H_B/\Xi_{\max}^2.
\]
No triple output is separately counted.
\end{proof}

\begin{lemma}[A dominant unbalanced triple stock bounds the junction
output]
\label{lem:V32-junction-vs-W}
Under all-nonlocality and \(W\ge\Xi_{\max}/4\),
\begin{equation}
\boxed{A_{U_e}\le CW.}
\label{eq:V32-junction-vs-W}
\end{equation}
\end{lemma}

\begin{proof}
Fusion and all-nonlocality give
\[
 A_{U_e}
 \le C\sum_ra_r
 <
 \frac C4\sum_r\Xi_r
 \le C\Xi_{\max}
 \le4CW.
\]
\end{proof}

\begin{theorem}[Unbalanced-dominant hot trivalent junction payer]
\label{thm:V32-unbalanced-hot}
Under the hypotheses of
\Cref{thm:V32-cold}, suppose one tree fails and
\(W\ge\Xi_{\max}/4\).  Then the complete hot junction-output payer
obeys the original marked-tree weights.
\end{theorem}

\begin{proof}
Keep the balanced moment as a contraction.  By
\eqref{eq:V31-unbalanced-moments} with \(k=3\),
\[
 \|M_U\|\le C/(s_\alpha^3W^3).
\]
Combine this with the hot local tree bound and
\eqref{eq:V32-junction-vs-W}:
\[
 \textup{hot contribution}
 \le
 Cs_\alpha^3A_{U_e}L_\tau
 \|M_U\|
 \le CL_\tau/W^2
 \le CL_\tau/B_\tau
 \le C\mathcal W_{\tau,e}.
\]
Once \(W\ge\Xi_{\max}/4\), the comparison
\(B_\sigma\le\Xi_{\max}^2\le16W^2\) holds for every tree
\(\sigma\).  Hence the one failed tree closes the full local tree
family.
\end{proof}

\section{The simultaneous private complement}
\label{sec:V32-private}

\begin{theorem}[Private-dominant hot trivalent junction payer]
\label{thm:V32-private-hot}
Suppose
\[
 B_\tau\le DY^2
\quad\hbox{for all sixteen quartic trees}.
\]
Then the complete hot junction-output payer obeys the marked-tree
atlas.
\end{theorem}

\begin{proof}
The four star trees give
\[
 \Xi_r^2\le DY^2
\quad(1\le r\le4).
\]
Thus \(Y>0\) and fusion gives
\[
 A_{U_e}\le C\sum_ra_r
 \le C\sum_r\Xi_r
 \le C_DY.
\]
Retain the private cubic moment:
\[
 (s_\alpha Y)^3q_P\le C.
\]
For every hot tree output,
\[
 q_P
 \frac{A_{U_e}}{1-q_{\alpha,U_e}}\|K_{\tau,U_e}\|
 \le
 Cs_\alpha^3A_{U_e}q_PL_\tau
 \le C_DL_\tau/Y^2
 \le C_DL_\tau/B_\tau.
\]
Finish with \(L_\tau\le\prod_ra_r^{d_r}\) and sum the fixed tree
family.
\end{proof}

\begin{theorem}[All-nonlocal pure incidence-three output priority]
\label{thm:V32-all-nonlocal-output}
Every all-nonlocal pure incidence-three \(SU(2)\) unique-junction
block in output priority obeys the quartic marked-tree atlas.
\end{theorem}

\begin{proof}
The shared-output payer is
\Cref{thm:V31-shared-output}; the cold junction payer is
\Cref{thm:V32-cold}.  For the hot junction payer, if one tree fails,
use \Cref{lem:V32-triple-stock} followed by
\Cref{thm:V32-balanced-hot} or
\Cref{thm:V32-unbalanced-hot}.  If none fails, use
\Cref{thm:V32-private-hot}.
\end{proof}

\section{The local-dominant complement}
\label{sec:V32-local-dominant}

\begin{lemma}[Unmarked pure trivalent shared payer]
\label{lem:V32-unmarked-shared}
For the complete pure incidence-three shared carrier, including the
protected external invariant line,
\begin{equation}
\boxed{\kappa_\otimes(\mathcal R_S^{(3)})\le C/s_\alpha.}
\label{eq:V32-unmarked-shared}
\end{equation}
\end{lemma}

\begin{proof}
For gapped outputs, use the mixed decomposition
\eqref{eq:V31-mixed-resolvent}, the unmarked balanced resolvent, and
\eqref{eq:V31-unbalanced-resolvent}.  The protected cross term is
bounded by
\[
 N_B\kappa_\otimes(P_B)\,
 \|\mathcal G_U\|
 \le C/s_\alpha.
\]
On the globally protected line, use
\eqref{eq:V30-external-split}; the retained full output is nontrivial
outside the invariant bank, so its external factor is \(O(1/s_\alpha)\).
\end{proof}

\begin{theorem}[Local-dominant pure incidence-three closure]
\label{thm:V32-local-dominant}
If one row \(c\) satisfies \(a_c\ge\Xi_c/4\), then the complete pure
incidence-three unique-junction physical carrier is bounded by
\[
 C_DA_0\le16C_D\mathcal W_{\tau_c,e}.
\]
\end{theorem}

\begin{proof}
The star comparison is
\eqref{eq:V29-star-controls-product}.  The two-scale proof of
\Cref{thm:V29-local-dominant} depends only on:
\begin{enumerate}
\item the complete local payer cap \(C/s_\alpha\);
\item the complete local cumulant cap;
\item the small-product tree estimates; and
\item an unmarked complete shared payer of size \(C/s_\alpha\).
\end{enumerate}
The first three are local, while the fourth is
\Cref{lem:V32-unmarked-shared}.  The V29 proof therefore applies
word for word.  Outside output priority, shared spectators are merely
contractions, so the same V29 private-payer proof applies.
\end{proof}

\begin{corollary}[Full pure incidence-three unique-junction sector]
\label{cor:V32-full-incidence-three}
For \(SU(2)\), every unique-junction physical block whose shared
spectator coordinates all have incidence exactly three obeys the
quartic marked-tree atlas.
\end{corollary}

\begin{proof}
Outside output priority, use
\Cref{thm:V31-below-output} in the all-nonlocal case and the
outside-priority part of
\Cref{thm:V32-local-dominant} otherwise.  Inside output priority, use
\Cref{thm:V32-all-nonlocal-output} or
\Cref{thm:V32-local-dominant}.
\end{proof}

\begin{assessment}[The pure trivalent structural sector is closed]
\label{ass:V32-pure-trivalent-closed}
V32 closes the complete pure incidence-three unique-junction
\(SU(2)\) carrier.  The proof includes triple invariants, all gapped
triple outputs, balanced and unbalanced input triples, cold and hot
junction outputs, all branching alternatives, and local-dominant
rows.
\end{assessment}

\begin{openproblem}[V33 mandate: mixed incidence two and three]
\label{op:V33-mandate}
\begin{enumerate}
\item Allow incidence-two and incidence-three spectator coordinates in
      the same unique-junction block.
\item Form one protected hypergraph containing both pair singlets and
      triple invariants.
\item Prove its vertex-cut capacity without multiplying the two
      incidence classes.
\item Combine their balanced legal-pair stocks and unbalanced retained
      output stocks before the failed-tree selection.
\item Reuse V24 and V31 only after this common stock is defined.
\item Close the mixed-incidence carrier before passing to incidence
      four.
\end{enumerate}
\end{openproblem}

\begin{firewall}[Claim boundary after compact V32]
\label{fire:V32-current-boundary}
V32 closes pure incidence-three unique-junction \(SU(2)\) blocks.  It
does not yet close mixtures of incidence two and three, incidence-four
shared coordinates, general compact simple groups, the full embedded
\([4]\) sector, multiple junctions, normalized collision aggregation,
\([3,3]\), global quartic or all-order MTA, continuum Yang--Mills
existence, or a positive mass gap.
\end{firewall}

\section{Append-only record for compact V32}
\label{sec:V32-record}

V32 was appended to the complete compact V31 file.  The exact V31
parent had \(12185\) newline-delimited source records, \(397330\)
bytes, and SHA--256 digest
\[
\texttt{8a589056390d04ce37c32afc7f66d013e7f71c51deeaea7993f228c7746eb7aa}.
\]
No compact V31 mathematical content was changed.  The sole final
\verb|\end{document}| is restored below.

% =============================================================================
% END APPENDED COMPACT VERSION V32 MATERIAL
% =============================================================================

% The compact V32 document terminator is intentionally disabled here:
% \end{document}

% =============================================================================
% BEGIN APPENDED COMPACT VERSION V33 MATERIAL
% =============================================================================

\chapter{V33: Mixed incidence-two/incidence-three hypergraphs}
\label{ch:V33}

\section{One protected hypergraph}
\label{sec:V33-protected-hypergraph}

Let \(F=F_2\sqcup F_3\) be a finite multiset of shared spectator
coordinates on the four rows.  A coordinate in \(F_k\) is active on
exactly \(k\) rows.  At a protected coordinate \(f\), let \(P_f\)
denote respectively the normalized pair-singlet projection or the
normalized triple-invariant projection.  Thus a protected
incidence-two coordinate necessarily has equal input spins, whereas a
protected incidence-three coordinate satisfies
\eqref{eq:V30-triangle}.  For \(A\subseteq F\) consisting of
protected coordinates, put
\[
 P_A:=
 \left(\bigotimes_{f\in A}P_f\right)
 \otimes I_{F\setminus A},
 \qquad
 D_{i,A}:=
 \prod_{\substack{f\in A\\i\text{ active at }f}}d_{i,f}.
\]

\begin{theorem}[Mixed vertex-cut capacity]
\label{thm:V33-mixed-cut}
For every protected subset \(A\),
\begin{equation}
\boxed{
 \kappa_\otimes(P_A)
 \le\min_{1\le i\le4}D_{i,A}^{-1}
 \le\rho_{23}^{|A|},
 \qquad
 \rho_{23}:=2^{-1/2}<1.}
\label{eq:V33-mixed-cut}
\end{equation}
The bound is taken on the whole mixed hypergraph; no capacity of an
incidence-two subgraph is multiplied by a capacity of an
incidence-three subgraph.
\end{theorem}

\begin{proof}
Fix a row \(i\) and group the other three rows.  A normalized
equal-spin pair singlet has reduced density \(d_{i,f}^{-1}I\) on
either active row.  A normalized triple invariant has the same flat
one-row marginal by
\Cref{lem:V30-triple-invariant}.  Therefore the largest squared
Schmidt coefficient of the tensor of all factors in \(A\), across
the cut \(i\,|\,\{1,2,3,4\}\setminus\{i\}\), equals
\(D_{i,A}^{-1}\).  Factors not active on row \(i\) lie wholly on the
other side.  A four-row product vector is a special bipartite product
vector across this cut, proving the first inequality.

Taking the product over the four choices of row gives
\[
 \prod_{i=1}^4D_{i,A}
 =
 \prod_{f\in A}\prod_{i\text{ active at }f}d_{i,f}.
\]
Every active representation is nontrivial.  Each pair coordinate
therefore contributes at least \(2^2=4\), and each triple coordinate
at least \(2^3=8\).  In particular,
\[
 \prod_{i=1}^4D_{i,A}\ge4^{|A|}.
\]
For some row \(i\), \(D_{i,A}\ge4^{|A|/4}\).  Taking its reciprocal
proves
\(\min_iD_{i,A}^{-1}\le4^{-|A|/4}=2^{-|A|/2}\).
\end{proof}

\begin{remark}[Why the incidence classes must not be separated]
\label{rem:V33-no-class-product}
The minimizing row cut can depend on the entire protected subset
\(A\).  Estimating \(F_2\) and \(F_3\) separately and multiplying
their best capacities would generally choose incompatible cuts.
\Cref{thm:V33-mixed-cut} chooses one cut only after both classes have
been assembled.
\end{remark}

\section{The mixed complete Wilson moment}
\label{sec:V33-complete-moment}

At every \(f\in F\), let
\[
 M_f^{(t)}=\sum_Uq_{\alpha,U}^{\,t}P_{U,f},
 \qquad
 M_F^{(t)}=\bigotimes_{f\in F}M_f^{(t)}.
\]
The sum is the complete multiplicity-free pair output resolution for
\(f\in F_2\), and the complete isotypic triple output resolution,
including multiplicity labels, for \(f\in F_3\).  If every coordinate
has a protected output, put \(P_*=\bigotimes_fP_f\); otherwise put
\(P_*=0\).  Let \(N=|F|\), \(r_t=(1-c_*s_\alpha)^t\), and define
\(\Phi_{N,\rho}^{(0)},\Phi_{N,\rho}^{(1)}\) as in
\eqref{eq:V30-Phi0}--\eqref{eq:V30-Phi1}.

\begin{lemma}[Mixed power capacity]
\label{lem:V33-mixed-power}
\begin{equation}
\boxed{
 \kappa_\otimes(M_F^{(t)}-P_*)
 \le
 \begin{cases}
 \Phi_{N,\rho_{23}}^{(0)}(r_t),&P_*\ne0,\\
 \Phi_{N,\rho_{23}}^{(1)}(r_t),&P_*=0.
 \end{cases}}
\label{eq:V33-mixed-power}
\end{equation}
\end{lemma}

\begin{proof}
The uniform nontrivial-output separation gives
\[
 M_f^{(t)}
 \preccurlyeq
 \begin{cases}
 r_tI+(1-r_t)P_f,&P_f\ne0,\\
 r_tI,&P_f=0.
 \end{cases}
\]
If all protected factors exist, expand the tensor product and
subtract \(P_*\).  Each proper protected subset \(A\) is bounded by
\(\rho_{23}^{|A|}\) using
\eqref{eq:V33-mixed-cut}; summing the binomial expansion gives
\(\Phi^{(0)}\).  If at least one protected factor is absent, retain
one \(r_t\) and expand all remaining protected factors.  Since
\(r_t\le\rho_{23}+(1-\rho_{23})r_t\), the resulting expression is at
most \(\Phi^{(1)}\).
\end{proof}

\begin{corollary}[Mixed one-failure sums]
\label{cor:V33-mixed-sums}
For \(\nu=0,1\),
\begin{align}
 \sum_{n\ge1}
 \Phi_{N,\rho_{23}}^{(\nu)}
 ((1-c_*s_\alpha)^n)
 &\le
 Ce^{-cs_\alpha N}
 \left(1+\frac1{s_\alpha N}\right),
\label{eq:V33-integer-sum}\\
 \sum_{n\ge0}
 \Phi_{N,\rho_{23}}^{(\nu)}
 ((1-c_*s_\alpha)^{n+1/2})
 &\le
 Ce^{-cs_\alpha N}
 \left(1+\frac1{s_\alpha N}\right).
\label{eq:V33-half-sum}
\end{align}
\end{corollary}

\begin{proof}
The analytic proofs of
\Cref{lem:V23-analytic-sum,lem:V31-half-power} require only a fixed
\(0<\rho<1\).  Apply them with \(\rho=\rho_{23}\).
\end{proof}

\section{One balanced stock for both incidences}
\label{sec:V33-balanced-stock}

Classify an incidence-two coordinate as balanced when
\[
 \min(A_{u,f},A_{v,f})
 \ge\frac14\max(A_{u,f},A_{v,f}),
\]
as in V23.  Classify an incidence-three coordinate by alternative
\textup{(B)} of \Cref{lem:V31-triple-dichotomy}.  For every balanced
coordinate choose one legal row pair:
\[
 (u_f,v_f)=
 \begin{cases}
 \text{the two active rows},&f\in F_2,\\
 \text{the rows carrying the two largest spins},&f\in F_3.
 \end{cases}
\]
Set
\[
 w_f:=\max_{i\text{ active at }f}A_{i,f},
 \qquad
 E_f:=A_{u_f,f}A_{v_f,f}.
\]
For the mixed balanced bank \(B\), define
\[
 N_B:=|B|,
 \qquad
 S_B:=\sum_{f\in B}w_f,
 \qquad
 H_B:=\sum_{f\in B}E_f.
\]

\begin{lemma}[Common balanced stock comparison and routing]
\label{lem:V33-balanced-comparison}
\begin{align}
 \boxed{S_B^2\le16N_BH_B,}
\label{eq:V33-SH}\\
 \boxed{
 \mathfrak P_B^{(2,3)}
 :=
 A_0\sum_{f\in B}
 \frac{E_f}{\Xi_{u_f}\Xi_{v_f}}
 \ge\frac{A_0H_B}{\Xi_{\max}^2}.}
\label{eq:V33-routing}
\end{align}
Every spectator coordinate occurs in exactly one summand of
\(\mathfrak P_B^{(2,3)}\).
\end{lemma}

\begin{proof}
For incidence two, balancedness gives \(E_f\ge w_f^2/4\).
For incidence three,
\eqref{eq:V31-balanced-triple} gives \(E_f\ge w_f^2/16\).
Thus \(w_f^2\le16E_f\) in both cases, and Cauchy--Schwarz proves
\eqref{eq:V33-SH}.

Both chosen rows are active at \(f\); hence \(f\) is a legal mark on
the comparison-path edge \(u_fv_f\).  Assigning exactly one pair to
each coordinate prevents repeated use.  The marked-weight formula
gives the equality in \eqref{eq:V33-routing}, and
\(\Xi_{u_f}\Xi_{v_f}\le\Xi_{\max}^2\) gives the inequality.
\end{proof}

Let \(M_B=M_B^{(1)}\), let \(P_B\) be the all-protected projection
when it exists and zero otherwise, and let
\(\mathcal R_B^{(2,3)}\) be the complete gapped shared resolvent with
output mass
\[
 H_{\boldsymbol U}^{\rm out}
 :=N_B+\sum_{f\in B}C_2(U_f).
\]

\begin{theorem}[Mixed balanced marked compiler]
\label{thm:V33-balanced-compiler}
Uniformly in the mixture, representations, and all complete output
multiplicities,
\begin{align}
 \boxed{
 s_\alpha S_B^2\kappa_\otimes(M_B)\le CH_B,}
\label{eq:V33-balanced-moment}\\
 \boxed{
 N_BS_B^2\kappa_\otimes(P_B)\le CH_B,}
\label{eq:V33-balanced-protected}\\
 \boxed{
 s_\alpha^2S_B^2
 \kappa_\otimes(\mathcal R_B^{(2,3)})\le CH_B,}
\label{eq:V33-balanced-resolvent}\\
 \boxed{
 \kappa_\otimes(\mathcal R_B^{(2,3)})\le C/s_\alpha.}
\label{eq:V33-balanced-unmarked}
\end{align}
\end{theorem}

\begin{proof}
At \(t=1\),
\Cref{lem:V33-mixed-power,cor:V33-mixed-sums} give
\[
 \kappa_\otimes(M_B-P_B)
 \le C\min\{1,e^{-cs_\alpha N_B}\}.
\]
Together with \eqref{eq:V33-SH}, this implies
\[
 s_\alpha S_B^2\kappa_\otimes(M_B-P_B)
 \le C(s_\alpha N_B)
 \min\{1,e^{-cs_\alpha N_B}\}H_B
 \le CH_B.
\]
When \(P_B\ne0\),
\eqref{eq:V33-mixed-cut} gives
\(\kappa_\otimes(P_B)\le\rho_{23}^{N_B}\).  Hence
\[
 \frac{N_BS_B^2}{H_B}\kappa_\otimes(P_B)
 \le16N_B^2\rho_{23}^{N_B}\le C.
\]
This proves the protected line and its contribution to the first
line.

Write
\[
 \mathcal R_B^{(2,3)}
 =N_B\mathcal G_B+\mathcal C_B
\]
on the complement of \(P_B\).  Geometric summation and
\eqref{eq:V33-integer-sum} give
\[
 \kappa_\otimes(\mathcal G_B)
 \le
 Ce^{-cs_\alpha N_B}
 \left(1+\frac1{s_\alpha N_B}\right)
 \le\frac{C}{s_\alpha N_B}.
\]
For a nonprotected output tuple, put
\[
 X=s_\alpha\sum_fC_2(U_f),
 \qquad Q=\prod_fq_{\alpha,U_f}.
\]
The retained-failure compiler with \(m=2\) gives
\(X^2Q\le C(1-Q)\), so
\[
 \frac{(\sum_fC_2(U_f))Q}{1-Q}
 \le\frac C{s_\alpha}\sum_{n\ge0}Q^{n+1/2}.
\]
Equations
\eqref{eq:V33-mixed-power} and \eqref{eq:V33-half-sum} bound the
capacity of the right side.  The \(m=1\) compiler also gives the
direct operator bound \(\mathcal C_B\preccurlyeq
Cs_\alpha^{-1}(I-P_B)\).

Put \(x=s_\alpha N_B\).  If \(x\le1\), use the direct bounds
\[
 \kappa_\otimes(\mathcal G_B)\le C/(s_\alpha N_B),
 \qquad
 \kappa_\otimes(\mathcal C_B)\le C/s_\alpha
\]
with \eqref{eq:V33-SH}.  If \(x>1\), use the exponential integer and
half-integer estimates.  Exactly as in
\Cref{thm:V23-graph-resolvent,thm:V31-marked-trivalent}, the scalar
remainders are bounded by
\[
 C[x^2+x]e^{-cx}(1+x^{-1}).
\]
This proves \eqref{eq:V33-balanced-resolvent}.  The two direct
unmarked estimates also give
\eqref{eq:V33-balanced-unmarked}.
\end{proof}

\section{One unbalanced stock for both incidences}
\label{sec:V33-unbalanced-stock}

Put in \(U\) every incidence-two coordinate failing the V23
balanced condition and every incidence-three coordinate satisfying
alternative \textup{(U)} of
\Cref{lem:V31-triple-dichotomy}.  Retain
\[
 w_f:=\max_{i\text{ active at }f}A_{i,f},
 \qquad
 W:=\sum_{f\in U}w_f.
\]
For an output \(U_f\), put \(b_{U_f}=1+C_2(U_f)\).

\begin{lemma}[Uniform retained-output comparison]
\label{lem:V33-unbalanced-comparison}
For every \(f\in U\) and every allowed complete output,
\begin{equation}
\boxed{cw_f\le b_{U_f}\le Cw_f.}
\label{eq:V33-unbalanced-comparison}
\end{equation}
In particular, no coordinate in \(U\) has a protected output.
\end{lemma}

\begin{proof}
For incidence two, this is
\Cref{lem:V24-output-comparison}.  For incidence three, the lower
bound is \eqref{eq:V31-unbalanced-triple} and the upper bound is the
fusion-Casimir estimate.  A protected output would have
\(b_{U_f}=1\), contradicting the lower-spin triangle argument that
places every invariant triple in the balanced class; for pairs it
would require equal spins, which are balanced.
\end{proof}

Let \(M_U\) be the complete Wilson moment and let
\(\mathcal G_U,\mathcal R_U^{(2,3)}\) be respectively the geometric
and output-mass resolvents on \(U\).

\begin{theorem}[Mixed unbalanced retained-output compiler]
\label{thm:V33-unbalanced-compiler}
\begin{align}
 (s_\alpha W)^k\|M_U\|&\le C_k,
 &&1\le k\le3,
\label{eq:V33-unbalanced-moment}\\
 \|\mathcal G_U\|&\le\frac{C_k}{(s_\alpha W)^k},
 &&1\le k\le3,
\label{eq:V33-unbalanced-geometric}\\
 \boxed{
 \|\mathcal R_U^{(2,3)}\|
 \le\min\left\{\frac C{s_\alpha},
                \frac C{s_\alpha^3W^2}\right\}.}
\label{eq:V33-unbalanced-resolvent}
\end{align}
\end{theorem}

\begin{proof}
For a resolved output tuple, write
\[
 Q=\prod_fq_{\alpha,U_f},
 \qquad
 B_{\rm out}=\sum_fb_{U_f}.
\]
By \eqref{eq:V33-unbalanced-comparison} and the one-link Wilson
debits,
\[
 (s_\alpha w_f)^kq_{\alpha,U_f}
 \le C_k(1-q_{\alpha,U_f}),
 \qquad 1\le k\le3.
\]
The ordinary product compiler yields
\((s_\alpha W)^kQ\le C_k\); maximizing proves
\eqref{eq:V33-unbalanced-moment}.  The retained-failure compiler
yields
\((s_\alpha W)^kQ\le C_k(1-Q)\); division by \(1-Q\) proves
\eqref{eq:V33-unbalanced-geometric}.  Finally,
\[
 cW\le B_{\rm out}\le CW.
\]
Multiplying the \(k=1\) and \(k=3\) geometric estimates by \(CW\)
proves \eqref{eq:V33-unbalanced-resolvent}.
\end{proof}

\section{The unified mixed resolvent}
\label{sec:V33-unified-resolvent}

On the complete gapped output of \(B\sqcup U\), positivity and the
scalar inequalities
\[
 1-Q_BQ_U\ge1-Q_B,
 \qquad
 1-Q_BQ_U\ge1-Q_U
\]
give
\begin{equation}
\boxed{
 \mathcal R_{B\sqcup U}^{(2,3)}
 \preccurlyeq
 \mathcal R_B^{(2,3)}\otimes M_U
 M_B\otimes\mathcal R_U^{(2,3)}
 N_BP_B\otimes\mathcal G_U.}
\label{eq:V33-mixed-decomposition}
\end{equation}
As in \Cref{lem:V24-mixed-decomposition}, the last term is exactly the
case \(Q_B=1\), and is zero when \(P_B=0\).  The three tensors are
inserted with the one-sided capacity theorem, never by multiplying
capacities.

\begin{theorem}[Complete mixed shared-output payer]
\label{thm:V33-shared-payer}
Fix \(D\ge16\).  Suppose all shared spectator coordinates have
incidence in \(\{2,3\}\), every row is nonlocal, and
\(\Xi_T>DY\).  Then the complete shared-output-resolvent payer obeys
the quartic marked-tree atlas.
\end{theorem}

\begin{proof}
The cases \(B=\varnothing\) and \(U=\varnothing\) are the limiting
forms of the estimates below.  Assume both are nonempty.

First consider the globally protected shared line.  It exists only
when \(U=\varnothing\) and every balanced coordinate has its
protected output.  For an external multiplier \(z<1\), its exact
coefficient is
\[
 \frac{N_Bz}{1-z}P_B.
\]
If all trees satisfy \(B_\tau\le DY^2\), use
\(z/(1-z)\le q_P/(1-q_P)\), the private quadratic debit, and
\(N_B\kappa_\otimes(P_B)\le C\).  If one tree fails, choose a
globally maximal row.  Since all coordinates are balanced,
\(S_B\ge h_i>\Xi_{\max}/2\).  The full retained output outside the
shared protected line is nontrivial, so \(z/(1-z)\le C/s_\alpha\).
Equation \eqref{eq:V33-balanced-protected} leaves
\[
 Cs_\alpha^2L_\tau H_B/S_B^2.
\]
If \(s_\alpha^2L_\tau\le A_0\), pay it with
\eqref{eq:V33-routing}; if a star violates this inequality, use the
complete local cumulant and the resulting
\(A_0>s_\alpha^{-4}\).  This is precisely the protected-line
calculation of V31, now with the mixed \(H_B\).

For gapped shared outputs, first assume
\(B_\tau\le DY^2\).  Equations
\eqref{eq:V33-balanced-unmarked},
\eqref{eq:V33-unbalanced-resolvent}, and
\[
 N_B\kappa_\otimes(P_B)\le C,
 \qquad
 \|\mathcal G_U\|\le C/(s_\alpha W)\le C/s_\alpha
\]
inserted one-sidedly into
\eqref{eq:V33-mixed-decomposition} give
\[
 \kappa_\otimes(\mathcal R_{B\sqcup U}^{(2,3)})
 \le C/s_\alpha.
\]
The local \(s_\alpha^3L_\tau\) factor and the private quadratic debit
close this branch.

Suppose instead that one tree fails.  Choose a globally maximal row.
Failed branching gives \(h_i>\Xi_{\max}/2\).  Every active input mass
is at most the corresponding \(w_f\), hence
\begin{equation}
\boxed{S_B+W>\Xi_{\max}/2.}
\label{eq:V33-dominant-split}
\end{equation}

If \(S_B>\Xi_{\max}/4\), insert
\eqref{eq:V33-balanced-resolvent},
\eqref{eq:V33-balanced-moment}, and
\eqref{eq:V33-balanced-protected} into the three terms of
\eqref{eq:V33-mixed-decomposition}, together with the unmarked
unbalanced bounds.  After multiplication by the local factor
\(s_\alpha^3L_\tau\), the terms are bounded respectively by
\[
 Cs_\alpha L_\tau\frac{H_B}{S_B^2},
 \qquad
 Cs_\alpha L_\tau\frac{H_B}{S_B^2},
 \qquad
 Cs_\alpha^2L_\tau\frac{H_B}{S_B^2}.
\]
The common legal path sum \eqref{eq:V33-routing} pays paths and
nonhot stars.  The complete local cumulant and
\(A_0>s_\alpha^{-2}\) pay the hot stars exactly as in V24 and V31.

If \(W\ge\Xi_{\max}/4\), use
\[
 \|M_U\|\le C/(s_\alpha^2W^2),\quad
 \|\mathcal R_U^{(2,3)}\|\le C/(s_\alpha^3W^2),\quad
 \|\mathcal G_U\|\le C/(s_\alpha^3W^3).
\]
The first two terms of
\eqref{eq:V33-mixed-decomposition}, after local insertion, are at
most \(CL_\tau/W^2\).  For the third, note
\[
 N_B\le S_B
 \le\sum_rh_r
 \le4\Xi_{\max}\le16W;
\]
it has the same bound.  Finally,
\(B_\tau\le\Xi_{\max}^2\le16W^2\), so every term is at most a constant
times \(L_\tau/B_\tau\).  This proves the marked-tree estimate.
\end{proof}

\begin{corollary}[Unmarked complete mixed shared payer]
\label{cor:V33-unmarked-shared}
Including the externally damped globally protected line,
\begin{equation}
\boxed{
 \kappa_\otimes(\mathcal R_S^{(2,3)})\le C/s_\alpha.}
\label{eq:V33-unmarked-shared}
\end{equation}
\end{corollary}

\begin{proof}
On gapped outputs use
\eqref{eq:V33-mixed-decomposition} and the unmarked estimates in the
preceding proof.  On the protected line,
\[
 N_B\kappa_\otimes(P_B)\frac z{1-z}\le C/s_\alpha,
\]
because the retained full output outside the shared invariant bank is
nontrivial.  This is the exact external split of
\Cref{cor:V30-external-split}, with
\eqref{eq:V33-mixed-cut} replacing the pure trivalent capacity.
\end{proof}

\section{Common stock after a failed tree}
\label{sec:V33-failed-tree}

Let \(M_S=M_B\otimes M_U\).

\begin{lemma}[Mixed stock dichotomy]
\label{lem:V33-stock-dichotomy}
Under all-nonlocality, if one tree fails, then either
\begin{align}
 S_B&>\Xi_{\max}/4,
 &
 s_\alpha S_B^2\kappa_\otimes(M_S)&\le CH_B,
\label{eq:V33-stock-balanced}
\end{align}
or
\begin{align}
 W&\ge\Xi_{\max}/4,
 &
 \|M_S\|&\le C/(s_\alpha^2W^2).
\label{eq:V33-stock-unbalanced}
\end{align}
\end{lemma}

\begin{proof}
Equation \eqref{eq:V33-dominant-split} is incidence-blind and gives
one of the two stock alternatives.  In the first use
\eqref{eq:V33-balanced-moment} and \(\|M_U\|\le1\); in the second use
\eqref{eq:V33-unbalanced-moment} with \(k=2\) and
\(\|M_B\|\le1\).
\end{proof}

\section{The mixed junction payer}
\label{sec:V33-junction}

\begin{theorem}[All-nonlocal mixed junction-output closure]
\label{thm:V33-junction}
Under the hypotheses of \Cref{thm:V33-shared-payer}, the complete cold
and hot junction-output payers obey the quartic marked-tree atlas.
\end{theorem}

\begin{proof}
On the cold set \(s_\alpha A_{U_e}\le1\), the local tree payer is at
most \(Cs_\alpha^2L_\tau\).  If
\(B_\tau\le DY^2\), the private quadratic moment closes it.  If a tree
fails, use \Cref{lem:V33-stock-dichotomy}.  The unbalanced alternative
gives
\[
 Cs_\alpha^2L_\tau\|M_S\|
 \le CL_\tau/W^2
 \le CL_\tau/B_\tau.
\]
The balanced alternative gives
\[
 Cs_\alpha L_\tau H_B/S_B^2.
\]
Equation \eqref{eq:V33-routing} pays paths and nonhot stars; for a hot
star the complete local payer has norm \(C/s_\alpha\), and
\(A_0>s_\alpha^{-2}\) pays the resulting
\(CH_B/(s_\alpha^2S_B^2)\).  This is the V32 cold proof with the
common stock substituted before any incidence-specific argument.

For hot junction outputs after a failed tree, use the V26 complete
good/bad partition
\[
 s_\alpha^2A_{U_e}\sqrt{M_0}\le\varepsilon
 \quad\hbox{or}\quad
 s_\alpha^2A_{U_e}\sqrt{M_0}>\varepsilon.
\]
In the balanced alternative, the good-set local bound
\(s_\alpha^2A_{U_e}L_\tau\le\varepsilon A_0\) and the bad-set
complete local output moment
\(\|\mathcal J_{\rm bad}^{\rm pay}\|\le Cs_\alpha A_0\), inserted in
\eqref{eq:V33-stock-balanced}, both give
\[
 C A_0H_B/\Xi_{\max}^2
 \le C\mathfrak P_B^{(2,3)}.
\]

In the unbalanced alternative, fusion and all-nonlocality give
\[
 A_{U_e}
 \le C\sum_ra_r
 <\frac C4\sum_r\Xi_r
 \le C\Xi_{\max}
 \le4CW.
\]
Using \(k=3\) in
\eqref{eq:V33-unbalanced-moment}, the hot local bound gives
\[
 Cs_\alpha^3A_{U_e}L_\tau\|M_U\|
 \le CL_\tau/W^2
 \le CL_\tau/B_\tau.
\]

It remains only the branch in which every tree satisfies
\(B_\tau\le DY^2\).  The four star inequalities imply
\(\Xi_r\le\sqrt D\,Y\) for all rows.  Hence
\(A_{U_e}\le C_DY\).  The private cubic moment
\((s_\alpha Y)^3q_P\le C\) gives, on every hot output,
\[
 q_P\frac{A_{U_e}}{1-q_{\alpha,U_e}}
 \|K_{\tau,U_e}\|
 \le C_D L_\tau/Y^2
 \le C_D L_\tau/B_\tau.
\]
The fixed tree family is therefore completely paid.
\end{proof}

\begin{theorem}[All-nonlocal mixed \(\{2,3\}\) sector]
\label{thm:V33-all-nonlocal}
Every all-nonlocal unique-junction \(SU(2)\) block whose shared
spectator incidences lie in \(\{2,3\}\) obeys the quartic marked-tree
atlas.
\end{theorem}

\begin{proof}
Below output priority, use the private payer when all branching tests
hold.  If one fails, apply
\Cref{lem:V33-stock-dichotomy}: the balanced alternative is paid by
\eqref{eq:V33-routing} after
\eqref{eq:V33-stock-balanced}, and the unbalanced alternative by
\eqref{eq:V33-stock-unbalanced} and
\(B_\tau\le\Xi_{\max}^2\le16W^2\).  This is the complete V20/V31
below-priority proof, but with selection performed only after the
two incidence classes have been combined.

Inside output priority, the shared-output term is
\Cref{thm:V33-shared-payer} and the junction-output term is
\Cref{thm:V33-junction}.
\end{proof}

\section{Local-dominant complement and closure}
\label{sec:V33-local-dominant}

\begin{theorem}[Full mixed incidence-two/incidence-three
unique-junction sector]
\label{thm:V33-full-mixed}
For \(SU(2)\), every unique-junction physical block whose shared
spectator coordinates have incidence two or three obeys the quartic
marked-tree atlas.
\end{theorem}

\begin{proof}
If every row is nonlocal, use
\Cref{thm:V33-all-nonlocal}.  Otherwise choose a row \(c\) with
\(a_c\ge\Xi_c/4\).  The two-scale local-dominant proof of
\Cref{thm:V29-local-dominant} uses only the complete local payer and
cumulant bounds, the small-product tree estimates, and an unmarked
complete shared payer of size \(C/s_\alpha\).  The first three inputs
are incidence-independent and the fourth is
\eqref{eq:V33-unmarked-shared}.  It yields
\[
 C_DA_0\le16C_D\mathcal W_{\tau_c,e}.
\]
Outside output priority the shared factors are contractions, so the
same private-payer argument applies.
\end{proof}

\begin{assessment}[The mixed incidence-\(\{2,3\}\) gate is closed]
\label{ass:V33-mixed-closed}
The protected pair and triple tensors are now aggregated in one
hypergraph capacity estimate.  Their balanced input masses are
compiled into one legal path stock, and their strongly unbalanced
input masses into one retained-output stock before failed-tree
selection.  Consequently V33 closes the full unique-junction
\(SU(2)\) sector with arbitrary mixtures of incidence-two and
incidence-three shared coordinates.
\end{assessment}

\begin{openproblem}[V34 mandate: incidence four]
\label{op:V34-mandate}
\begin{enumerate}
\item Classify the invariant multiplicity space in
      \(V_{j_1}\otimes V_{j_2}\otimes V_{j_3}\otimes V_{j_4}\);
      unlike the pair and triple cases it need not be one-dimensional.
\item Separate the canonical projector onto the full invariant
      isotypic space from a basis-dependent rank-one recoupling
      projector.
\item Prove a four-row product-capacity bound for that full invariant
      projector which is summable over repeated coordinates and does
      not pay its multiplicity by counting intermediate spins.
\item If such a capacity bound fails, identify the exact high-spin
      family and move upstream to an input-mark or local-output
      mechanism before attempting a resolvent.
\item Do not insert the already-closed mixed
      incidence-\(\{2,3\}\) sector until the incidence-four protected
      obstruction has been isolated.
\end{enumerate}
\end{openproblem}

\begin{firewall}[Claim boundary after compact V33]
\label{fire:V33-current-boundary}
V33 closes arbitrary mixtures of incidence-two and incidence-three
shared coordinates in the \(SU(2)\) unique-junction sector.  It does
not close incidence-four shared coordinates, general compact simple
groups, the full embedded \([4]\) sector, multiple junctions,
normalized collision aggregation, \([3,3]\), global quartic or
all-order MTA, continuum Yang--Mills existence, or a positive mass
gap.
\end{firewall}

\section{Append-only record for compact V33}
\label{sec:V33-record}

V33 was appended to the complete compact V32 file.  The exact V32
parent had \(12571\) newline-delimited source records, \(408679\)
bytes, and SHA--256 digest
\[
\texttt{24ef2af77d30f0bba1796fad54d9e62655c20a5f2f678adb40e86052cd05163c}.
\]
No compact V32 mathematical content was changed.  The sole final
\verb|\end{document}| is restored below.

% =============================================================================
% END APPENDED COMPACT VERSION V33 MATERIAL
% =============================================================================

% The compact V33 document terminator is intentionally disabled here:
% \end{document}

% =============================================================================
% BEGIN APPENDED COMPACT VERSION V34 MATERIAL
% =============================================================================

\chapter{V34: Full incidence-four invariant spaces and completion of
the embedded \texorpdfstring{\([4]\)}{[4]} sector}
\label{ch:V34}

\section{The canonical four-input invariant projector}
\label{sec:V34-four-invariant}

Fix nontrivial \(SU(2)\) spins \(j_1,j_2,j_3,j_4\), put
\[
 \mathcal H_{\boldsymbol j}
 :=V_{j_1}\otimes V_{j_2}\otimes V_{j_3}\otimes V_{j_4},
\]
and let \(P_{\boldsymbol j}^{\rm inv}\) denote the orthogonal
projection onto
\(\operatorname{Inv}(\mathcal H_{\boldsymbol j})\).

\begin{proposition}[Recoupling multiplicity]
\label{prop:V34-recoupling-multiplicity}
Let \(\mathcal K(\boldsymbol j)\) be the set of half-integers \(k\)
such that
\begin{align*}
 |j_1-j_2|\le k\le j_1+j_2,\qquad
 j_1+j_2+k\in\mathbb Z,\\
 |j_3-j_4|\le k\le j_3+j_4,\qquad
 j_3+j_4+k\in\mathbb Z.
\end{align*}
Then
\begin{equation}
 \dim\operatorname{Inv}(\mathcal H_{\boldsymbol j})
 =|\mathcal K(\boldsymbol j)|.
\label{eq:V34-invariant-multiplicity}
\end{equation}
For the pairing \(12\,|\,34\), normalized recoupling invariants
\(\Omega_k\), \(k\in\mathcal K(\boldsymbol j)\), form an orthonormal
basis and
\begin{equation}
\boxed{
 P_{\boldsymbol j}^{\rm inv}
 =\sum_{k\in\mathcal K(\boldsymbol j)}
 |\Omega_k\rangle\langle\Omega_k|.}
\label{eq:V34-canonical-projector}
\end{equation}
The projector is canonical although the displayed orthonormal basis
depends on the chosen pairing.
\end{proposition}

\begin{proof}
The multiplicity-free Clebsch--Gordan decompositions give
\[
 V_{j_1}\otimes V_{j_2}
 =\bigoplus_{k\in\mathcal K_{12}}V_k,
 \qquad
 V_{j_3}\otimes V_{j_4}
 =\bigoplus_{\ell\in\mathcal K_{34}}V_\ell.
\]
An invariant occurs in \(V_k\otimes V_\ell\) exactly when
\(k=\ell\), and then with multiplicity one.  This proves
\eqref{eq:V34-invariant-multiplicity}.  Choosing the normalized
singlet in every \(V_k\otimes V_k\) gives mutually orthogonal
\(\Omega_k\)'s and hence
\eqref{eq:V34-canonical-projector}.  Changing the recoupling pairing
changes this basis by a unitary matrix on the invariant space and
therefore leaves its orthogonal projector unchanged.
\end{proof}

\section{Multiplicity does not cost product capacity}
\label{sec:V34-multiplicity-capacity}

\begin{lemma}[One-row normal form of a full invariant projector]
\label{lem:V34-one-row-normal-form}
Fix \(i\in\{1,2,3,4\}\), set \(d_i=2j_i+1\), and group the other
three factors into \(\mathcal B_i\).  There is a multiplicity space
\(\mathcal M_i\) and an isometric equivariant inclusion
\[
 T_i:V_{j_i}^{*}\otimes\mathcal M_i\longrightarrow\mathcal B_i
\]
such that, on
\(V_{j_i}\otimes\mathcal B_i\),
\begin{equation}
 P_{\boldsymbol j}^{\rm inv}
 \simeq
 |\Omega_i\rangle\langle\Omega_i|
 \otimes I_{\mathcal M_i}
 \oplus0.
\label{eq:V34-one-row-normal-form}
\end{equation}
Here \(\Omega_i\) is the normalized canonical invariant in
\(V_{j_i}\otimes V_{j_i}^{*}\).  Consequently, for unit vectors
\(x_i\in V_{j_i}\) and \(y_i\in\mathcal B_i\),
\begin{equation}
\boxed{
 \langle x_i\otimes y_i,
 P_{\boldsymbol j}^{\rm inv}(x_i\otimes y_i)\rangle
 \le d_i^{-1}.}
\label{eq:V34-one-row-capacity}
\end{equation}
\end{lemma}

\begin{proof}
Decompose the \(V_{j_i}^{*}\)-isotypic component of
\(\mathcal B_i\) as
\[
 V_{j_i}^{*}\otimes\mathcal M_i,
 \qquad
 \mathcal M_i
 =\operatorname{Hom}_{SU(2)}(V_{j_i}^{*},\mathcal B_i).
\]
All other isotypic components are orthogonal to every invariant in
\(V_{j_i}\otimes\mathcal B_i\).  Schur's lemma identifies the
invariant space with
\[
 \operatorname{Inv}(V_{j_i}\otimes V_{j_i}^{*})
 \otimes\mathcal M_i.
\]
The first factor is the one-dimensional span of
\[
 \Omega_i=d_i^{-1/2}\sum_{a=1}^{d_i}e_a\otimes e_a^*.
\]
This proves \eqref{eq:V34-one-row-normal-form}.

Let \(\Pi_i\) be the projection of \(\mathcal B_i\) onto the displayed
isotypic component.  Contracting \(x_i\) against \(\Omega_i\) gives
a vector of norm \(d_i^{-1/2}\|x_i\|\).  Therefore
\[
 \langle x_i\otimes y_i,
 P_{\boldsymbol j}^{\rm inv}(x_i\otimes y_i)\rangle
 =
 \|(\langle\Omega_i|\otimes I_{\mathcal M_i})
       (x_i\otimes\Pi_i y_i)\|^2
 \le d_i^{-1}\|x_i\|^2\|y_i\|^2.
\]
No factor \(\dim\mathcal M_i\) occurs.
\end{proof}

\begin{corollary}[Four-row product capacity]
\label{cor:V34-four-product-capacity}
\begin{equation}
\boxed{
 \kappa_\otimes(P_{\boldsymbol j}^{\rm inv})
 \le\min_{1\le i\le4}d_i^{-1}.}
\label{eq:V34-four-product-capacity}
\end{equation}
\end{corollary}

\begin{proof}
A four-row product vector is a bipartite product vector across every
cut \(i\,|\,\widehat i\).  Apply
\eqref{eq:V34-one-row-capacity} and minimize over \(i\).
\end{proof}

\begin{theorem}[Tensorized mixed-incidence cut with full invariant
spaces]
\label{thm:V34-full-mixed-cut}
Let \(A\) be any protected collection of pair, triple, and four-row
coordinates.  At a four-row coordinate use the full canonical
projector \(P_f^{\rm inv}\), not a rank-one recoupling projector.
Then
\begin{equation}
\boxed{
 \kappa_\otimes(P_A)
 \le\min_iD_{i,A}^{-1}
 \le\rho_{234}^{|A|},
 \qquad
 \rho_{234}:=2^{-1/2}.}
\label{eq:V34-full-mixed-cut}
\end{equation}
\end{theorem}

\begin{proof}
Fix the row cut \(i\,|\,\widehat i\).  On every protected coordinate
active at row \(i\), the proof of
\Cref{lem:V34-one-row-normal-form} gives
\[
 P_f\simeq
 |\Omega_{i,f}\rangle\langle\Omega_{i,f}|
 \otimes I_{\mathcal M_{i,f}}\oplus0.
\]
For pair and triple coordinates the multiplicity space is
one-dimensional; for a four-row coordinate it need not be.  Tensoring
these normal forms gives one maximally entangled vector between row
\(i\)'s active tensor product, of dimension \(D_{i,A}\), and its dual
on the other side, tensored with the identity on the product
multiplicity space.  The contraction calculation in
\eqref{eq:V34-one-row-capacity} therefore gives
\(\kappa_\otimes(P_A)\le D_{i,A}^{-1}\), even though a row vector may
be entangled among its coordinate factors.

Each protected coordinate contributes at least \(4\) to
\(\prod_iD_{i,A}\): the lower bounds are \(4,8,16\) for incidences
two, three, and four.  Thus
\(\prod_iD_{i,A}\ge4^{|A|}\), and the geometric-mean argument of
\Cref{thm:V33-mixed-cut} yields
\(\min_iD_{i,A}^{-1}\le2^{-|A|/2}\).
\end{proof}

\begin{assessment}[Resolution of the multiplicity obstruction]
\label{ass:V34-multiplicity-resolution}
Summing rank-one recoupling bounds would lose
\(|\mathcal K(\boldsymbol j)|\).  The full invariant projector avoids
that loss because, across any one-row cut, its multiplicity is an
identity tensor behind a single maximally entangled representation
factor.  This is the incidence-four analogue of flat one-row
marginals, but it is a statement about a subspace rather than one
vector.
\end{assessment}

\section{Complete Wilson aggregation at incidence four}
\label{sec:V34-Wilson-aggregation}

At an incidence-four coordinate let
\[
 M_f^{(t)}
 :=\sum_{U}q_{\alpha,U}^{\,t}P_{U,f},
\]
where \(P_{U,f}\) is the full \(U\)-isotypic projection, including
all multiplicities.  Its protected projection is
\(P_f=P_f^{\rm inv}\).  For a mixed collection \(F\) of incidences
two, three, and four define \(M_F^{(t)}\) and \(P_*\) as in V33.

\begin{theorem}[Incidence-\(\{2,3,4\}\) power and resolvent compiler]
\label{thm:V34-power-resolvent}
Let \(N=|F|\), \(r_t=(1-c_*s_\alpha)^t\), and
\(\rho=\rho_{234}\).  Then
\begin{equation}
\boxed{
 \kappa_\otimes(M_F^{(t)}-P_*)
 \le
 \begin{cases}
 \Phi_{N,\rho}^{(0)}(r_t),&P_*\ne0,\\
 \Phi_{N,\rho}^{(1)}(r_t),&P_*=0.
 \end{cases}}
\label{eq:V34-power-capacity}
\end{equation}
The integer and half-integer sums of the right side obey
\eqref{eq:V33-integer-sum}--\eqref{eq:V33-half-sum}, with uniform
constants.  On gapped outputs the complete unmarked shared resolvent
obeys
\begin{equation}
\boxed{\kappa_\otimes(\mathcal R_F^{(2,3,4)})\le C/s_\alpha.}
\label{eq:V34-unmarked-resolvent}
\end{equation}
\end{theorem}

\begin{proof}
All nontrivial output representations obey the uniform Wilson
separation.  Hence
\[
 M_f^{(t)}
 \preccurlyeq r_tI+(1-r_t)P_f
\]
when an invariant space exists, and
\(M_f^{(t)}\preccurlyeq r_tI\) otherwise.  Expand the product and use
\eqref{eq:V34-full-mixed-cut} on every protected subset.  The
binomial sum is exactly the right side of
\eqref{eq:V34-power-capacity}.  The analytic sum depends only on the
fixed number \(0<\rho<1\), so
\Cref{cor:V33-mixed-sums} applies.

Geometric summation pays the \(N\)-part of the output mass by
\(C/s_\alpha\).  For its Casimir part, the retained-failure compiler
with \(m=1\) gives the operator bound
\[
 \sum_{\boldsymbol U}
 \frac{(\sum_fC_2(U_f))q_{\boldsymbol U}}
      {1-q_{\boldsymbol U}}P_{\boldsymbol U}
 \preccurlyeq
 \frac C{s_\alpha}(I-P_*).
\]
The isotypic projections already include multiplicity, so neither
step contains a recoupling-channel count.
\end{proof}

\begin{corollary}[Exact external split for full invariant spaces]
\label{cor:V34-external-split}
For \(0\le z<1\), let the external multiplier \(z\) damp the complete
mixed-incidence shared carrier.  Then
\begin{equation}
\boxed{
 \kappa_\otimes\!\left(
 \sum_{\boldsymbol U}
 \frac{H_{\boldsymbol U}^{\rm out}zq_{\boldsymbol U}}
      {1-zq_{\boldsymbol U}}P_{\boldsymbol U}\right)
 \le\frac C{s_\alpha}+C\frac z{1-z}.}
\label{eq:V34-external-split}
\end{equation}
In a physical carrier where the retained external output is
nontrivial, the right side is \(O(s_\alpha^{-1})\).
\end{corollary}

\begin{proof}
On every gapped shared tuple,
\[
 \frac{zq_{\boldsymbol U}}{1-zq_{\boldsymbol U}}
 \le
 \frac{q_{\boldsymbol U}}{1-q_{\boldsymbol U}},
 \]
so \eqref{eq:V34-unmarked-resolvent} applies.  On the all-invariant
space, \(q_{\boldsymbol U}=1\) and
\(H_{\boldsymbol U}^{\rm out}=N\).  Equation
\eqref{eq:V34-full-mixed-cut} gives
\[
 N\kappa_\otimes(P_*)\le N\rho_{234}^{N}\le C,
 \]
which proves the second term in
\eqref{eq:V34-external-split}.  If the retained external output is
nontrivial, the uniform Wilson separation gives
\(z/(1-z)\le C/s_\alpha\).
\end{proof}

\section{The four-input mark dichotomy}
\label{sec:V34-four-input-dichotomy}

At one incidence-four coordinate order the inputs as
\[
 j_1\ge j_2\ge j_3\ge j_4>0,
 \qquad
 A_\nu=1+j_\nu(j_\nu+1),
 \qquad
 w=A_1.
\]

\begin{lemma}[Either two large inputs or a retained total output]
\label{lem:V34-four-dichotomy}
Exactly one of the following alternatives holds:
\begin{enumerate}[label=\textup{(\mathrm B)},leftmargin=4em]
\item \(j_2+j_3+j_4\ge j_1/2\).  Then
      \begin{equation}
      \boxed{
      A_2\ge\frac1{36}A_1,
      \qquad
      A_1A_2\ge\frac1{36}w^2.}
      \label{eq:V34-four-balanced}
      \end{equation}
\end{enumerate}
or
\begin{enumerate}[label=\textup{(\mathrm U)},leftmargin=4em]
\item \(j_2+j_3+j_4<j_1/2\).  Every total output spin \(J\) then
      satisfies
      \begin{equation}
      \boxed{1+J(J+1)\ge\frac14w.}
      \label{eq:V34-four-unbalanced}
      \end{equation}
\end{enumerate}
Every coordinate with a nonzero invariant space is balanced.
\end{lemma}

\begin{proof}
In alternative \textup{(B)}, the ordering gives
\[
 3j_2\ge j_2+j_3+j_4\ge j_1/2,
 \qquad j_2\ge j_1/6.
\]
Therefore
\[
 1+j_2(j_2+1)
 \ge1+\frac{j_1}{6}\left(\frac{j_1}{6}+1\right)
 \ge\frac1{36}\bigl(1+j_1(j_1+1)\bigr).
\]
In alternative \textup{(U)}, repeated triangle inequalities give
\[
 J\ge j_1-j_2-j_3-j_4>j_1/2,
\]
and the monotonic comparison used in
\Cref{lem:V20-unbalanced-retention} proves
\eqref{eq:V34-four-unbalanced}.  Finally, existence of an invariant
requires \(j_1\le j_2+j_3+j_4\), so an invariant coordinate cannot
be unbalanced.
\end{proof}

On a balanced incidence-four coordinate assign its mark once, to the
pair of rows carrying \(j_1,j_2\), and put \(E_f=A_{1,f}A_{2,f}\).
On an unbalanced coordinate retain \(w_f=A_{1,f}\).

\section{The universal four-row shared compiler}
\label{sec:V34-universal-compiler}

Combine all balanced coordinates of incidences two, three, and four
into \(B\), using the legal pair assignments of V33 and the preceding
section.  Combine all unbalanced coordinates into \(U\).  Put
\[
 S_B=\sum_{f\in B}w_f,
 \qquad
 H_B=\sum_{f\in B}E_f,
 \qquad
 W=\sum_{f\in U}w_f.
\]

\begin{theorem}[Universal balanced and unbalanced estimates]
\label{thm:V34-universal-stocks}
The common stocks obey
\begin{align}
 \boxed{S_B^2\le36N_BH_B,}
\label{eq:V34-SH}\\
 \boxed{
 \mathfrak P_B^{(2,3,4)}
 :=
 A_0\sum_{f\in B}
 \frac{E_f}{\Xi_{u_f}\Xi_{v_f}}
 \ge\frac{A_0H_B}{\Xi_{\max}^2}.}
\label{eq:V34-routing}
\end{align}
Moreover,
\begin{align}
 s_\alpha S_B^2\kappa_\otimes(M_B)&\le CH_B,
\label{eq:V34-balanced-moment}\\
 N_BS_B^2\kappa_\otimes(P_B)&\le CH_B,
\label{eq:V34-balanced-protected}\\
 s_\alpha^2S_B^2
 \kappa_\otimes(\mathcal R_B^{(2,3,4)})&\le CH_B,
\label{eq:V34-balanced-resolvent}\\
 (s_\alpha W)^k\|M_U\|&\le C_k,
\quad1\le k\le3,
\label{eq:V34-unbalanced-moment}\\
 \|\mathcal R_U^{(2,3,4)}\|
 &\le\min\left\{\frac C{s_\alpha},
                 \frac C{s_\alpha^3W^2}\right\}.
\label{eq:V34-unbalanced-resolvent}
\end{align}
\end{theorem}

\begin{proof}
The balanced coordinate inequalities are
\(w_f^2\le4E_f\), \(w_f^2\le16E_f\), and
\(w_f^2\le36E_f\) for incidences two, three, and four respectively.
Cauchy--Schwarz proves \eqref{eq:V34-SH}.  Each coordinate is assigned
to one pair of active rows, so the marked-weight identity and
\(\Xi_{u_f}\Xi_{v_f}\le\Xi_{\max}^2\) prove
\eqref{eq:V34-routing}.

Use \eqref{eq:V34-power-capacity}, its integer and half-integer sums,
and \eqref{eq:V34-SH} in the proof of
\Cref{thm:V33-balanced-compiler}.  The only numerical change is
\(16\mapsto36\); all exponential scalar factors remain uniformly
bounded.  This proves
\eqref{eq:V34-balanced-moment}--%
\eqref{eq:V34-balanced-resolvent}, including the full
incidence-four invariant projector.

On every unbalanced coordinate,
\eqref{eq:V20-unbalanced-retention},
\eqref{eq:V31-unbalanced-triple}, or
\eqref{eq:V34-four-unbalanced} gives
\[
 cw_f\le1+C_2(U_f)\le Cw_f.
\]
The proof of
\Cref{thm:V33-unbalanced-compiler} uses only this comparison and the
one-link debits.  It proves
\eqref{eq:V34-unbalanced-moment}--%
\eqref{eq:V34-unbalanced-resolvent}.
\end{proof}

\begin{corollary}[Universal mixed resolvent split]
\label{cor:V34-universal-split}
On the complete gapped output,
\begin{equation}
 \mathcal R_{B\sqcup U}^{(2,3,4)}
 \preccurlyeq
 \mathcal R_B^{(2,3,4)}\otimes M_U
 M_B\otimes\mathcal R_U^{(2,3,4)}
 N_BP_B\otimes\mathcal G_U.
\label{eq:V34-universal-split}
\end{equation}
After one failed tree, either
\begin{align}
 S_B&>\Xi_{\max}/4,
 &
 s_\alpha S_B^2\kappa_\otimes(M_B\otimes M_U)&\le CH_B,
\label{eq:V34-balanced-stock}
\end{align}
or
\begin{align}
 W&\ge\Xi_{\max}/4,
 &
 \|M_B\otimes M_U\|&\le C/(s_\alpha^2W^2).
\label{eq:V34-unbalanced-stock}
\end{align}
\end{corollary}

\begin{proof}
The positive split is the scalar proof of
\eqref{eq:V33-mixed-decomposition}.  For a failed tree, choose a
globally maximal row.  All-nonlocality and failed branching give
\(h_i>\Xi_{\max}/2\), while every active input mass is at most
\(w_f\).  Hence \(S_B+W>\Xi_{\max}/2\).  Use
\eqref{eq:V34-balanced-moment} in the first alternative and
\eqref{eq:V34-unbalanced-moment} with \(k=2\) in the second.
\end{proof}

\section{Completion of the unique-junction and embedded
\texorpdfstring{\([4]\)}{[4]} sectors}
\label{sec:V34-embedded-closure}

\begin{theorem}[Full four-row unique-junction carrier]
\label{thm:V34-full-unique-junction}
For \(SU(2)\), every unique-junction physical block with arbitrary
shared spectator incidence obeys the quartic marked-tree atlas.
\end{theorem}

\begin{proof}
A shared spectator on four rows has incidence two, three, or four.
Thus
\Cref{thm:V34-universal-stocks,cor:V34-universal-split} provide every
incidence-dependent input used in the V33 payer proof.

For completeness, if all rows are nonlocal and a tree fails, use
\eqref{eq:V34-balanced-stock} or
\eqref{eq:V34-unbalanced-stock}.  In the balanced branch, the common
legal sum \eqref{eq:V34-routing} pays the shared resolver, the cold
junction payer, and both sides of the V26 hot-output partition.  In
the unbalanced branch, the quadratic and cubic estimates from
\eqref{eq:V34-unbalanced-moment} give
\[
 CL_\tau/W^2\le CL_\tau/B_\tau
\]
for the cold and hot junction terms, using
\(B_\tau\le\Xi_{\max}^2\le16W^2\).  The three terms of
\eqref{eq:V34-universal-split} obey the same bound; for its cross term
use \(N_B\le S_B\le4\Xi_{\max}\le16W\).

If all trees satisfy \(B_\tau\le DY^2\), the private quadratic and
cubic moments close the shared and junction payers.  This argument is
independent of spectator incidence.  If some row is local-dominant,
the V29 two-scale proof applies with the unmarked bound
\eqref{eq:V34-external-split}, whose physical form is
\(O(s_\alpha^{-1})\).  These cases exhaust the V8 residual atlas.
\end{proof}

\begin{corollary}[Complete normalized embedded \([4]\) MTA]
\label{cor:V34-embedded-MTA}
The coefficient-sensitive normalized quartic marked-tree estimate
holds for the complete embedded \([4]\) topology over \(SU(2)\),
uniformly in supports, representations, output multiplicities, and
the regulator parameter.
\end{corollary}

\begin{proof}
The V2 exact unique-junction factorization and one-resolvent allocation
reduce every embedded \([4]\) coefficient block to the physical
carriers covered by
\Cref{thm:V34-full-unique-junction}.  Coordinates active on only one
row belong to the already-paid private factors; every shared
coordinate is covered by the theorem.  The physical output
partition is orthogonal, so complete isotypic projectors may be kept
intact.  Finally the exact Fourier normalization
\eqref{eq:V1-exact-Fourier} and raw pinching
\eqref{eq:V1-raw-pinching} sum the coefficient blocks without a
representation-dimension or multiplicity count.
\end{proof}

\begin{assessment}[The embedded \([4]\) gate is closed]
\label{ass:V34-embedded-closed}
The first open row in the compact V1 status table is now resolved:
the complete normalized embedded \([4]\) sector satisfies quartic
MTA for \(SU(2)\).  The decisive step was not a bound on individual
recoupling invariants, but the one-row normal form of the full
invariant projector.
\end{assessment}

\begin{openproblem}[V35 mandate: audit and select the next quartic
topology]
\label{op:V35-mandate}
\begin{enumerate}
\item Perform the mandatory read-only audit of the newest active SM
      and NS notes.
\item Reconcile their newest continuation or aggregation mechanisms
      with the completed embedded-\([4]\) carrier, importing no
      estimate without proof.
\item Compare the remaining normalized collision-\([3,2]\) Boolean
      selector with the multiple-junction \([3,3]\) topology.
\item Attack first the obstruction that lies upstream of both:
      exact one-resolvent allocation across more than one junction, if
      that common obstruction is real; otherwise choose the topology
      whose coefficient partition is already orthogonal.
\item Preserve full invariant isotypic projectors throughout; do not
      return to a recoupling-channel count.
\end{enumerate}
\end{openproblem}

\begin{firewall}[Claim boundary after compact V34]
\label{fire:V34-current-boundary}
V34 closes the complete normalized embedded-\([4]\) quartic sector
for \(SU(2)\).  It does not close normalized collision aggregation in
\([3,2]\), the \([3,3]\) sector, general compact simple groups, global
quartic or all-order MTA, weighted uniform fusion control, the
continuum Yang--Mills construction, or a positive mass gap.
\end{firewall}

\section{Append-only record for compact V34}
\label{sec:V34-record}

V34 was appended to the complete compact V33 file.  The exact V33
parent had \(13346\) newline-delimited source records, \(432920\)
bytes, and SHA--256 digest
\[
\texttt{4d726ebcb2f73aadee49f66d63e4b28492ed74e7d03707ef62eddee0002e1ba4}.
\]
No compact V33 mathematical content was changed.  The sole final
\verb|\end{document}| is restored below.

% =============================================================================
% END APPENDED COMPACT VERSION V34 MATERIAL
% =============================================================================

% The compact V34 document terminator is intentionally disabled here:
% \end{document}

% =============================================================================
% BEGIN APPENDED COMPACT VERSION V35 MATERIAL
% =============================================================================

\chapter{V35: Companion audit and first-entry allocation for the
\texorpdfstring{\([3,3]\)}{[3,3]} topology}
\label{ch:V35}

\section{Mandatory companion audit}
\label{sec:V35-audit}

The newest active companion sources inspected read-only were:
\begin{center}
\begin{tabular}{p{0.42\linewidth}r r p{0.30\linewidth}}
\toprule
file & lines & bytes & SHA--256\\
\midrule
\texttt{Renewal\_SM\_Einstein\_Continuum\_Research\_Notes\_V40.tex}
& \(14698\) & \(515144\)
& \texttt{3a3b96b8a1f7115d49c3e8e2e454056f6314d6486edbd3f301695ae5a2fe9105}\\
\texttt{Renewal\_NS\_Stability\_Research\_Notes\_V31.tex}
& \(23052\) & \(887537\)
& \texttt{f1121b62238ad5491bb7cf03635ea9934942edf573ed8faaa9ee79e1d38a6696}\\
\bottomrule
\end{tabular}
\end{center}

The SM file has advanced from V36 to V40.  V37 derives an exact
Hadamard variation for scalar and bundle Laplace-type eigenvalues and
shows that lapse/shift geometry supplies a checkable, but not automatic,
transversality criterion.  V38 replaces first-order transversality by a
finite-type sublevel theorem and identifies uniform integrability as the
extra hypothesis needed for weighted crossing occupation.  V39 localizes
the cluster gap in time and makes differentiated norm-resolvent
convergence an explicit regulator input.  V40 assembles local
renormalization before the nonlocal response, proves a fractional
Volterra resolvent for an integrable causal kernel, and proves that past
history enters a continuation slab once as a known source.

The NS file is unchanged.  Its V31 endpoint proves a multiplicity-free
dyadic formation ladder, but shows that removing a flat probability mark
restores the exact missing first moment.  Its persistent branch cannot
be debited again on every recent-entry interval.

\begin{assessment}[V35 audit transfer]
\label{ass:V35-audit-transfer}
No SM or NS estimate is imported into Yang--Mills.  Two proof-order
rules transfer:
\begin{enumerate}
\item assemble all local terms belonging to one finite state before
      taking a norm; and
\item when a word first acquires a second junction type, retain its
      preceding junction bank once as an assembled history, rather than
      summing over every possible pair of witnessing coordinates.
\end{enumerate}
The Yang--Mills implementation below is an exact finite word identity,
not a causal or PDE analogy.
\end{assessment}

\section{Two signed junctions and one final resolvent}
\label{sec:V35-two-junction}

Let \(X\) be a finite coordinate set with distinct junction coordinates
\(e,f\).  On a coordinatewise orthogonal physical resolution write
\[
 P_{\boldsymbol U,\boldsymbol a}
 =\bigotimes_{g\in X}P^g_{U_g,a_g}.
\]
Assume that the norm of one resolved coordinate word has the exact
factorized majorant
\begin{equation}
 k_{\boldsymbol U,\boldsymbol a}
 \le
 k_e(U_e,a_e)k_f(U_f,a_f)
 \prod_{g\notin\{e,f\}}m_g(U_g,a_g),
 \qquad 0\le m_g\le1.
\label{eq:V35-two-junction-factorization}
\end{equation}
This assumption is valid for a fixed word with two local connected
blocks and complete positive moment spectators.  It is not asserted for
an unassembled signed topology sum.

Define positive block majorants
\begin{align*}
 J_e&:=\sum_{U,a}k_e(U,a)P^e_{U,a},
 &
 J_e^{\rm pay}
 &:=
 \sum_{U,a}\frac{A_U}{1-q_{\alpha,U}}
 k_e(U,a)P^e_{U,a},\\
 J_f&:=\sum_{U,a}k_f(U,a)P^f_{U,a},
 &
 J_f^{\rm pay}
 &:=
 \sum_{U,a}\frac{A_U}{1-q_{\alpha,U}}
 k_f(U,a)P^f_{U,a},
\end{align*}
and use \(M_g,D_g^{\rm pay}\) from
\eqref{eq:V2-spectator-payer} for the remaining coordinates.

\begin{theorem}[Two-junction one-resolvent capacity reduction]
\label{thm:V35-two-junction-reduction}
Under \eqref{eq:V35-two-junction-factorization}, the normalized
physical carrier \(\mathcal Q_{e,f}\) of the word satisfies
\begin{align}
 \kappa_\otimes(\mathcal Q_{e,f})
 \le{}&
 \min\left\{
   \kappa_\otimes(J_e^{\rm pay})\|J_f\|,
   \|J_e^{\rm pay}\|\kappa_\otimes(J_f)
 \right\}
\notag\\
&+
 \min\left\{
   \kappa_\otimes(J_f^{\rm pay})\|J_e\|,
   \|J_f^{\rm pay}\|\kappa_\otimes(J_e)
 \right\}
\notag\\
&+
 \sum_{g\notin\{e,f\}}
 \min\left\{
 \begin{aligned}
 &\kappa_\otimes(J_e)\|J_f\|\|D_g^{\rm pay}\|,\\
 &\|J_e\|\kappa_\otimes(J_f)\|D_g^{\rm pay}\|,\\
 &\|J_e\|\|J_f\|\kappa_\otimes(D_g^{\rm pay})
 \end{aligned}\right\}.
\label{eq:V35-two-junction-reduction}
\end{align}
Exactly one payer carries a resolvent denominator in every summand.
\end{theorem}

\begin{proof}
Apply the blockwise allocation
\eqref{eq:V2-carrier-allocation} to the single final denominator before
using \eqref{eq:V35-two-junction-factorization}.  When \(e\) is the
payer, the resulting positive majorant is
\[
 J_e^{\rm pay}\otimes J_f
 \otimes\bigotimes_{g\notin\{e,f\}}M_g.
\]
Discard the positive spectator contractions with
\Cref{cor:V2-shared-contractions}, then apply the one-sided tensor
capacity theorem in either orientation.  This gives the first minimum.
The \(f\)-payer gives the second.

When a spectator \(g\) is the payer, the remaining noncontractive tensor
is
\[
 J_e\otimes J_f\otimes D_g^{\rm pay}.
\]
Apply the one-sided capacity theorem twice, choosing any one of the three
factors for the capacity estimate and ordinary norms for the other two.
This gives the three-entry minimum.  Sum the positive payer
majorants.  At no stage are two product capacities multiplied.
\end{proof}

\begin{corollary}[No repeated junction denominator]
\label{cor:V35-no-double-resolvent}
A multi-junction proof may require separate first-moment estimates for
the two local connected blocks, but it may not insert
\((1-q_{\alpha,U_e})^{-1}(1-q_{\alpha,U_f})^{-1}\).
The only valid denominator in
\eqref{eq:V35-two-junction-reduction} is the one assigned from the
original final resolvent.
\end{corollary}

\begin{proof}
Every term in \eqref{eq:V35-two-junction-reduction} is produced by one
summand of \eqref{eq:V2-carrier-allocation}.  A second denominator is
absent from that identity and cannot be introduced by the one-sided
tensor estimate.
\end{proof}

\section{The surplus edge in a \([3,3]\) skeleton}
\label{sec:V35-surplus-edge}

Let \(I,J\subset\{1,2,3,4\}\) be distinct triples.  Then
\(|I\cap J|=2\) and \(I\cup J=\{1,2,3,4\}\).  Let
\(\sigma\) and \(\tau\) be spanning trees on \(I\) and \(J\),
respectively.  Their union is a connected four-edge multigraph on the
four rows.

\begin{lemma}[One normalized ternary mark is surplus]
\label{lem:V35-surplus-edge}
Assign a number \(0\le\theta_a\le1\) to each occurrence of an edge in
the multigraph \(\sigma\cup\tau\).  Then
\begin{equation}
\boxed{
 \prod_{a\in E(\sigma)\sqcup E(\tau)}\theta_a
 \le
 \sum_{\substack{\gamma\subset
                  E(\sigma)\sqcup E(\tau)\\
                  \gamma\text{ a spanning tree}}}
 \prod_{a\in\gamma}\theta_a.}
\label{eq:V35-surplus-edge}
\end{equation}
The sum contains at most four terms.
\end{lemma}

\begin{proof}
The four-edge multigraph is connected on four vertices, so it contains
a cycle, where a doubled common edge is allowed as a cycle of length
two.  Deleting any edge of that cycle leaves a spanning tree
\(\gamma\).  Since the deleted \(\theta_a\le1\),
\[
 \prod_{a\in E(\sigma)\sqcup E(\tau)}\theta_a
 \le\prod_{a\in\gamma}\theta_a.
\]
This one term occurs on the right side of
\eqref{eq:V35-surplus-edge}.  A four-edge multigraph has at most four
single-edge deletions, giving the cardinality assertion.
\end{proof}

\begin{assessment}[What local ternary control would buy]
\label{ass:V35-local-ternary-buy}
Each of the two triple blocks supplies two normalized link marks in the
already-closed ternary MTA.  Their product has four marks.  Once the
single final resolvent has been handled at the two-junction level,
\Cref{lem:V35-surplus-edge} converts those four marks into the required
three-link quartic tree atlas.  Thus the graph combinatorics of
\([3,3]\) is not the obstruction.  The obstruction is assembling all
possible junction coordinates before applying
\Cref{thm:V35-two-junction-reduction}.
\end{assessment}

\section{Why a raw sum over witness pairs is invalid}
\label{sec:V35-pair-count}

\begin{proposition}[Witness-pair labels are not orthogonal]
\label{prop:V35-pairs-not-orthogonal}
The decomposition of a \([3,3]\) topology sum by all ordered pairs
\((e,f)\) carrying distinct triple types is not, in general, an
orthogonal physical decomposition.  A word containing \(r\) occurrences
of one triple type and \(s\) occurrences of another would occur \(rs\)
times if every witness pair were selected.
\end{proposition}

\begin{proof}
The local partition chosen at a coordinate is an algebraic cumulant
letter, not a final isotypic projector.  Distinct choices of witnessing
coordinates therefore act on the same final physical block.  A fixed
word with the stated occurrences has exactly \(rs\) ordered coordinate
pairs of the two types.  Assigning it to every pair is literal
multiplicity, not orthogonality.  In the one-dimensional scalar
regression, all word operators have the same range, which rules out an
orthogonal interpretation.
\end{proof}

\section{A finite-state first-entry identity}
\label{sec:V35-first-entry}

Let \(\mathcal T=\{123,124,134,234\}\) be the four triple block types.
Order the coordinate set as \(X=\{1,\ldots,n\}\).  At coordinate \(r\),
let \(\Pi_r\subseteq\Pi_4\) be its admissible partition alphabet and
let \(L_{r,\pi}\) be the exact local cumulant letter of V130.

Define the finite state space
\[
 \mathcal Q
 :=
 \{\mathsf d,\mathsf a\}
 \cup\{S:S\subseteq\mathcal T,\ |S|\le1\}.
\]
The state \(\mathsf d\) is dead because a four-row block has appeared;
\(\mathsf a\) is the accepted \([3,3]\) state.  Starting from
\(\varnothing\), read a local partition \(\pi\) by the deterministic
transition \(\delta\):
\begin{enumerate}
\item if \(\pi\) contains a four-row block, move to \(\mathsf d\);
\item otherwise adjoin to \(S\) the unique triple block of \(\pi\), if
      present;
\item if the resulting set has at least two elements, move to
      \(\mathsf a\);
\item \(\mathsf a\) remains \(\mathsf a\) unless a later four-row block
      moves it to \(\mathsf d\), and \(\mathsf d\) remains dead.
\end{enumerate}
A set partition cannot contain two overlapping triple blocks, so its
triple type is unique.

For \(q\in\mathcal Q\), define operators recursively by
\begin{align}
 Z_0(q)&=
 \begin{cases}I,&q=\varnothing,\\0,&q\ne\varnothing,\end{cases}
\label{eq:V35-state-initial}\\
 Z_r(q)&=
 \sum_{\substack{p\in\mathcal Q,\ \pi\in\Pi_r\\
                  \delta(p,\pi)=q}}
 Z_{r-1}(p)\otimes L_{r,\pi}.
\label{eq:V35-state-recurrence}
\end{align}

\begin{theorem}[Exact one-use \([3,3]\) automaton]
\label{thm:V35-automaton}
For every \(q\in\mathcal Q\),
\begin{equation}
\boxed{
 Z_n(q)=
 \sum_{\substack{w(r)\in\Pi_r\\
                  \delta(\varnothing,w(1)\cdots w(n))=q}}
 \bigotimes_{r=1}^nL_{r,w(r)}.}
\label{eq:V35-state-word-sum}
\end{equation}
In particular, \(Z_n(\mathsf a)\) is the exact sum of all words which
contain at least two distinct triple types and no four-row block.  Every
such word is counted once.
\end{theorem}

\begin{proof}
Induct on \(r\).  The initialization is the empty word.  Appending one
letter \(\pi\) to a word ending in \(p\) tensors its operator with
\(L_{r,\pi}\) and moves its deterministic state to
\(\delta(p,\pi)\).  The pairs \((p,\pi)\) in
\eqref{eq:V35-state-recurrence} partition all words by their penultimate
state and last letter.  Thus no word is omitted or repeated.

The final state is \(\mathsf a\) exactly when two distinct triple types
have appeared and no four-block ever appears.  Two distinct triples on
four rows automatically have union all four rows, which is the
\([3,3]\) priority condition.
\end{proof}

To expose the unique first entry, split
\[
 \mathcal Q_-:=\{S\subseteq\mathcal T:|S|\le1\}.
\]
For \(1\le r\le n\), define
\begin{equation}
 F_r:=
 \sum_{\substack{p\in\mathcal Q_-,\,\pi\in\Pi_r\\
                  \delta(p,\pi)=\mathsf a}}
 Z_{r-1}(p)\otimes L_{r,\pi}.
\label{eq:V35-first-entry-source}
\end{equation}
Let \(\mathcal U_{s:r}\) denote the ordered transfer from coordinate
\(r+1\) through \(s\) restricted to trajectories which remain in
\(\mathsf a\), with a transition to \(\mathsf d\) killed.

\begin{corollary}[Discrete Volterra or first-entry expansion]
\label{cor:V35-first-entry}
\begin{equation}
\boxed{
 Z_n(\mathsf a)
 =
 \sum_{r=1}^n\mathcal U_{n:r}F_r.}
\label{eq:V35-first-entry}
\end{equation}
Each accepted word occurs in exactly one summand, indexed by the first
coordinate at which its second distinct triple type appears.
\end{corollary}

\begin{proof}
Every accepted deterministic trajectory has a unique first transition
from \(\mathcal Q_-\) to \(\mathsf a\).  Its prefix and transition
letter occur in \(F_r\); its remaining suffix occurs in
\(\mathcal U_{n:r}\).  Conversely, every term on the right is an
accepted trajectory with first-entry time \(r\).  This is a disjoint
partition of words, so the operator identity follows before taking any
norm.
\end{proof}

\begin{remark}[Finite state is not yet analytic contraction]
\label{rem:V35-state-not-contraction}
Equation \eqref{eq:V35-first-entry} removes witness-pair multiplicity
algebraically.  It does not imply
\(\|\mathcal U_{n:r}\|\le1\): the allowed local sums contain signed
cumulant letters.  Treating the suffix as a positive moment contraction
would be the same invalid shortcut ruled out by the V2 factorization
firewall.
\end{remark}

\section{The exact next analytic gate}
\label{sec:V35-next-gate}

\begin{definition}[One-type prefix and accepted transfer norms]
\label{def:V35-transfer-norms}
For a state \(S\in\mathcal Q_-\), let
\(\mathcal P_{r}(S)\) denote the complete prefix operator
\(Z_r(S)\).  Let
\(\mathcal A_{s:r}\) denote the complete accepted-state transfer
\(\mathcal U_{s:r}\), with all allowed local letters assembled at each
coordinate before a norm is taken.
\end{definition}

\begin{openproblem}[The \([3,3]\) transfer estimate]
\label{op:V35-transfer-gate}
Construct a positive or coefficient-pinched majorant for the exact
first-entry terms such that
\begin{equation}
 \sum_{r=1}^n
 \kappa_\otimes\!\left(
 \mathcal A_{n:r}F_r
 \ \text{with the single final resolvent}\right)
 \le
 C\sum_{\gamma\in\mathfrak T_4}
 \mathcal W_\gamma,
\label{eq:V35-transfer-target}
\end{equation}
with \(C\) independent of \(n\).  The prefix may use the complete
ternary MTA and the transition may use
\Cref{thm:V35-two-junction-reduction,lem:V35-surplus-edge}; the accepted
suffix must be retained once and may not be replaced by a product of
termwise cumulant norms.
\end{openproblem}

\begin{assessment}[V35 decision]
\label{ass:V35-decision}
The next target is \([3,3]\), not collision aggregation.  The embedded
\([4]\) result of V34 supplies the higher-priority exclusion, and the
already-closed ternary MTA controls a one-type prefix.  The new common
obstruction is the accepted-state transfer in
\eqref{eq:V35-first-entry}.  Solving it will also provide the correct
one-use architecture for the collision selector, whereas attacking
collision words first would mix that issue with their additional
mean/double-pair flag.
\end{assessment}

\begin{openproblem}[V36 mandate: accepted-state completion]
\label{op:V36-mandate}
\begin{enumerate}
\item Expand the no-four-block local sum into a complete moment minus
      the already-closed four-block cumulant, keeping both terms
      assembled.
\item Determine whether the accepted suffix is a contraction after the
      embedded-\([4]\) subtraction is inserted by a Duhamel telescoping
      identity.
\item If it is not, produce a sharp scalar or \(SU(2)\) regression
      before attempting \eqref{eq:V35-transfer-target}.
\item Insert the first-entry source into the two-junction one-resolvent
      reduction only once.
\item Use \Cref{lem:V35-surplus-edge} after, not before, the complete
      coefficient and output aggregation.
\end{enumerate}
\end{openproblem}

\begin{firewall}[Claim boundary after compact V35]
\label{fire:V35-current-boundary}
V35 proves the exact two-junction one-resolvent capacity reduction, the
surplus-edge conversion for two ternary trees, and a multiplicity-free
first-entry representation of the complete \([3,3]\) word sum.  It does
not prove the accepted-state transfer estimate, close \([3,3]\), close
normalized collision aggregation, prove global quartic or all-order
MTA, weighted uniform fusion control, continuum Yang--Mills existence,
or a positive mass gap.
\end{firewall}

\section{Append-only record for compact V35}
\label{sec:V35-record}

V35 was appended to the complete compact V34 file.  The exact V34
parent had \(13957\) newline-delimited source records, \(453285\)
bytes, and SHA--256 digest
\[
\texttt{3ed5ab0993c96ea9e6a624c1538261ab8cbf0f059869138351be7eaf2603ff77}.
\]
No compact V34 mathematical content was changed.  The sole final
\verb|\end{document}| is restored below.

% =============================================================================
% END APPENDED COMPACT VERSION V35 MATERIAL
% =============================================================================

% The compact V35 document terminator is intentionally disabled here:
% \end{document}

% =============================================================================
% BEGIN APPENDED COMPACT VERSION V36 MATERIAL
% =============================================================================

\chapter{V36: Restricted local transfers, a correction, and the
normalized global-tree target}
\label{ch:V36}

\section{The exact no-four-block local factor}
\label{sec:V36-no-four-factor}

At a four-row coordinate \(r\), retain the complete local partition
alphabet and write
\begin{align}
 M_r&:=\sum_{\pi\in\Pi_4}L_{r,\pi},
\label{eq:V36-complete-moment}\\
 K_r&:=L_{r,\{1234\}},
\label{eq:V36-four-letter}\\
 R_r&:=\sum_{\pi\in\Pi_4\setminus\{\{1234\}\}}L_{r,\pi}
     =M_r-K_r.
\label{eq:V36-restricted-factor}
\end{align}
The first operator is the complete four-row Wilson moment and hence a
positive contraction.  The last operator is a signed sum of all local
partitions except the four-block letter.

\begin{lemma}[Accepted-state transfer is a restricted product]
\label{lem:V36-accepted-product}
For the V35 \([3,3]\) automaton, once state \(\mathsf a\) has been
entered at coordinate \(r\), the transfer through coordinates
\(r+1,\ldots,n\), with every transition to the four-block dead state
killed, is
\begin{equation}
\boxed{
 \mathcal A_{n:r}
 =\bigotimes_{s=r+1}^nR_s.}
\label{eq:V36-accepted-product}
\end{equation}
\end{lemma}

\begin{proof}
In state \(\mathsf a\), every local partition except the four-block
letter keeps the automaton in \(\mathsf a\).  Summing those allowed
letters at coordinate \(s\) gives \(R_s\) by
\eqref{eq:V36-restricted-factor}.  Coordinate independence tensors
these sums in their fixed order.
\end{proof}

\begin{lemma}[Exact first-four-block telescope]
\label{lem:V36-four-telescope}
On any ordered interval \(a,\ldots,b\),
\begin{equation}
\boxed{
 \bigotimes_{s=a}^bM_s-\bigotimes_{s=a}^bR_s
 =
 \sum_{t=a}^b
 \left(\bigotimes_{a\le s<t}R_s\right)
 \otimes K_t\otimes
 \left(\bigotimes_{t<s\le b}M_s\right).}
\label{eq:V36-four-telescope}
\end{equation}
Every word containing at least one four-block letter occurs exactly
once, at its first such coordinate.
\end{lemma}

\begin{proof}
Use \(M_t-R_t=K_t\) and telescope consecutive partial products:
\[
 \bigotimes_{s=a}^bM_s-\bigotimes_{s=a}^bR_s
 =
 \sum_{t=a}^b
 \left(\bigotimes_{a\le s<t}R_s\right)
 \otimes(M_t-R_t)\otimes
 \left(\bigotimes_{t<s\le b}M_s\right).
\]
Expanding the prefix shows that it contains no earlier four-block,
while the suffix is unrestricted.  Hence \(t\) is the unique first
four-block coordinate of every expanded word.
\end{proof}

\section{A sharp contraction regression}
\label{sec:V36-contraction-regression}

\begin{proposition}[The restricted factor need not be a contraction]
\label{prop:V36-not-contraction}
There is no contraction theorem for \(R_r=M_r-K_r\) based only on the
moment--cumulant algebra and the assumption that the row variables are
contractions.  In a one-dimensional scalar model one can have
\begin{equation}
 M_r=1,\qquad K_r=-2,\qquad R_r=3.
\label{eq:V36-Rademacher-regression}
\end{equation}
Consequently a tensor of \(N\) such restricted factors has norm
\(3^N\).
\end{proposition}

\begin{proof}
Let \(X\) be a centered Rademacher variable and take all four row
variables equal to \(X\).  Then
\[
 \mathbb EX=0,\qquad
 \kappa_2(X,X)=1,\qquad
 \kappa_3(X,X,X)=0,
\]
and
\[
 \kappa_4(X,X,X,X)
 =\mathbb EX^4-3(\mathbb EX^2)^2
 =-2.
\]
The complete fourfold moment is \(M_r=\mathbb EX^4=1\).  Among the
non-top partition letters, only the three pair-pair partitions survive,
each with value one, so \(R_r=3\).  Tensoring scalar factors multiplies
their norms.
\end{proof}

\begin{firewall}[Scope of the scalar regression]
\label{fire:V36-regression-scope}
\Cref{prop:V36-not-contraction} is not asserted to be a Wilson-law
counterexample.  It proves that positivity of the complete moment and
finite cumulant algebra alone cannot justify
\(\|\mathcal A_{n:r}\|\le1\).  A valid Wilson proof must use its
coefficient normalization, heat-kernel debits, or the complete connected
operator before the restricted product is normed.
\end{firewall}

\section{Correction of the V34 global status}
\label{sec:V36-V34-correction}

\begin{assessment}[V34 results which remain valid]
\label{ass:V36-V34-valid}
The following V34 results are unaffected:
\begin{enumerate}
\item the canonical full incidence-four invariant projector and its
      recoupling multiplicity;
\item its one-row normal form
      \eqref{eq:V34-one-row-normal-form} and multiplicity-free product
      capacity;
\item the mixed-incidence power, marked-moment, and resolvent compilers;
      and
\item the MTA bound for a carrier with one already-selected four-row
      connected block and arbitrary complete positive moment spectators.
\end{enumerate}
These are coordinatewise statements or apply after the exact
factorization hypothesis of V2 has been verified.
\end{assessment}

\begin{theorem}[Withdrawal of the global embedded-\([4]\) corollary]
\label{thm:V36-withdraw-V34-global}
The global readings of
\Cref{cor:V34-embedded-MTA,ass:V34-embedded-closed} are withdrawn.
What V34 proves is the \emph{one-selected-junction} embedded-\([4]\)
carrier.  It does not sum the possible choices of a four-block letter
when several four-active coordinates occur.
\end{theorem}

\begin{proof}
Selecting a four-block at coordinate \(t\) and leaving complete positive
moments at all other coordinates does not form a disjoint word
partition: a word with \(m\) four-block letters would occur \(m\) times.
The exact disjoint selection by the first four-block is
\eqref{eq:V36-four-telescope}.  Its prefix is a tensor of \(R_s\)'s, not
complete positive moments.  By
\Cref{prop:V36-not-contraction}, the V2 positive-spectator contraction
argument cannot remove that prefix.  Therefore the selection aggregation
required for the global embedded-\([4]\) topology was absent from the V34
proof.
\end{proof}

\begin{remark}[Incidence four versus geometric uniqueness]
\label{rem:V36-incidence-four-uniqueness}
If \(e\) is literally the only coordinate common to all four row
supports, no off-\(e\) spectator has incidence four.  V34's
incidence-four compiler is nevertheless useful when one four-block
letter has been selected while another four-active coordinate is kept
as a complete moment.  The latter is an algebraic spectator, not a proof
that the support has only one possible junction.
\end{remark}

\begin{assessment}[Revised compact-series status]
\label{ass:V36-revised-status}
V29 closes the geometrically unique-junction incidence-two sector, V32
the pure incidence-three sector, and V33 their mixtures.  V34 adds the
full incidence-four invariant-space compiler and closes every
one-selected-junction carrier.  Multiple four-block selection and the
\([3,3]\) accepted transfer remain open and share the same restricted
factor \(R=M-K_4\).
\end{assessment}

\section{Why the global connected operator is upstream}
\label{sec:V36-global-upstream}

The frozen V135 theorem gives the complete connected quartic operator
before the final resolvent:
\[
 \|\mathbf K_I^{\rm conn}\|_{\rm op}
 \le
 C\sum_{\tau\in\mathfrak T_4}
 \prod_{ij\in E(\tau)}s_\alpha H_{ij}.
\]
Its proof retains all local partitions until after the exact connected
word sum is formed.  It therefore bypasses the exponentially unstable
restricted product of \Cref{prop:V36-not-contraction}.  Its stated
firewall is different: the global row denominators were not inserted
before the operator norm.

For nonzero row masses, define the normalized overlap numbers
\begin{equation}
 \theta_{ij}:=\frac{H_{ij}}{\Xi_i\Xi_j},
 \qquad
 H_{ij}:=\sum_eA_{i,e}A_{j,e}.
\label{eq:V36-normalized-overlap}
\end{equation}

\begin{lemma}[Normalized overlaps are probability weights]
\label{lem:V36-theta-range}
\begin{equation}
\boxed{0\le\theta_{ij}\le1.}
\label{eq:V36-theta-range}
\end{equation}
\end{lemma}

\begin{proof}
All summands are nonnegative and
\[
 H_{ij}
 =\sum_eA_{i,e}A_{j,e}
 \le
 \left(\sum_eA_{i,e}\right)
 \left(\sum_fA_{j,f}\right)
 =\Xi_i\Xi_j.
\]
Coordinates absent from a row contribute zero.
\end{proof}

For a labeled tree \(\tau\), put
\begin{equation}
 \mathfrak W_\tau^{\rm glob}
 :=
 \left(\prod_{i=1}^4\Xi_i\right)
 \prod_{ij\in E(\tau)}\theta_{ij}
 =
 \left(\prod_i\Xi_i^{1-\deg_\tau(i)}\right)
 \prod_{ij\in E(\tau)}H_{ij}.
\label{eq:V36-global-marked-tree}
\end{equation}
This is the global fixed-representation marked-tree weight whose
coefficient lift is required by the BFA seminorms.

\section{A normalized threshold compiler}
\label{sec:V36-normalized-threshold}

For a row cut \(C\mid D\), set
\[
 \Theta(C,D)
 :=\sum_{\substack{i\in C\\j\in D}}\theta_{ij}.
\]
For a partition into three nonempty row clusters, use the same notation
for quotient links.

\begin{theorem}[Finite normalized graph reduction]
\label{thm:V36-normalized-graph}
Let \(\mathcal F(\theta)\ge0\) be a quantity associated with six edge
weights \(0\le\theta_{ij}\le1\).  Suppose there are uniform constants
such that:
\begin{enumerate}[label=\textup{(\roman*)},leftmargin=4em]
\item \(\mathcal F\le C_0\);
\item for every two-cluster row partition,
      \(\mathcal F\le C_1\Theta(C,D)\);
\item for every three-cluster row partition,
      \[
      \mathcal F
      \le C_2
      \sum_{\sigma\in\mathfrak T_3}
      \prod_{CD\in E(\sigma)}\Theta(C,D);
      \]
\item if all \(\theta_{ij}<\delta_0\), then
      \[
      \mathcal F
      \le C_3\sum_{\tau\in\mathfrak T_4}
      \prod_{ij\in E(\tau)}\theta_{ij}.
      \]
\end{enumerate}
Here \(\delta_0\in(0,1)\) is fixed.  Then
\begin{equation}
\boxed{
 \mathcal F(\theta)
 \le C
 \sum_{\tau\in\mathfrak T_4}
 \prod_{ij\in E(\tau)}\theta_{ij}.}
\label{eq:V36-normalized-graph}
\end{equation}
\end{theorem}

\begin{proof}
Form the graph whose strong edges satisfy
\(\theta_{ij}\ge\delta_0\), and let \(k\) be its number of connected
components.

If \(k=1\), a strong spanning tree has weight at least
\(\delta_0^3\), so assumption \textup{(i)} is paid by the tree
polynomial.

If \(k=2\), assumption \textup{(ii)} bounds \(\mathcal F\) by the sum
of crossing edges.  Adjoin each crossing edge to fixed strong spanning
trees inside the two components.  The internal forest contains two
edges, and the corresponding original tree weight is at least
\(\delta_0^2\) times the crossing weight.

If \(k=3\), assumption \textup{(iii)} gives a sum of two-edge quotient
trees.  Expand each quotient link into original crossing edges and
adjoin the single internal strong edge.  Every term becomes an original
four-row tree with weight at least \(\delta_0\) times the quotient
monomial.

If \(k=4\), every edge is weak and assumption \textup{(iv)} applies.
The number of row edges and trees is fixed, so all expansion
multiplicities are absorbed in a numerical constant.
\end{proof}

\begin{corollary}[Exact analytic targets for normalized quartic MTA]
\label{cor:V36-four-targets}
Let \(\mathcal C_4\) be the complete coefficient-pinched, one-resolvent
quartic carrier, before any topology word is normed, and define
\begin{equation}
 \mathcal F_4
 :=
 \frac{\kappa_\otimes(\mathcal C_4)}
      {\prod_i\Xi_i}.
\label{eq:V36-normalized-carrier}
\end{equation}
If \(\mathcal F_4\) satisfies the four hypotheses of
\Cref{thm:V36-normalized-graph}, then
\begin{equation}
 \kappa_\otimes(\mathcal C_4)
 \le
 C\sum_{\tau\in\mathfrak T_4}
 \mathfrak W_\tau^{\rm glob}.
\label{eq:V36-global-MTA-target}
\end{equation}
\end{corollary}

\begin{proof}
Apply \eqref{eq:V36-normalized-graph} to
\eqref{eq:V36-normalized-carrier} and multiply by
\(\prod_i\Xi_i\).  Equation
\eqref{eq:V36-global-marked-tree} identifies every term.
\end{proof}

\begin{assessment}[What has moved upstream]
\label{ass:V36-upstream-decision}
The correct next object is not the norm of the accepted
\([3,3]\) suffix and not a sum over selected junction coordinates.  It
is the normalized complete carrier \(\mathcal F_4\) in
\eqref{eq:V36-normalized-carrier}.  The V135 threshold proof already
supplies the finite graph logic, while
\Cref{thm:V36-normalized-graph} states exactly which estimates must be
reproved with coefficient capacity and the sole resolvent still present.
\end{assessment}

\begin{openproblem}[V37 mandate: normalized weak-cut carrier]
\label{op:V37-mandate}
\begin{enumerate}
\item Return to the exact two-cluster product--cumulant identity used in
      frozen V135 before taking an operator norm.
\item Insert the exact Fourier factor, raw isotypic pinching, and the
      single final resolvent at that level.
\item Prove, first for a \(1|3\) cut and then for a \(2|2\) cut,
      \[
      \mathcal F_4\le C\Theta(C,D),
      \]
      or isolate the precise correction partition which lacks a
      normalized crossing mark.
\item Treat every lower cumulant as one complete binary or ternary MTA
      carrier; do not expand its internal coordinate words.
\item Keep the universal, three-cluster, and all-weak hypotheses of
      \Cref{thm:V36-normalized-graph} behind separate firewalls until the
      two-cluster estimate is complete.
\end{enumerate}
\end{openproblem}

\begin{firewall}[Claim boundary after compact V36]
\label{fire:V36-current-boundary}
V36 proves the exact restricted-factor and first-four-block identities,
disproves an abstract contraction shortcut, corrects the global
interpretation of V34, and reduces normalized global quartic MTA to four
capacity-level threshold estimates.  It does not yet prove the
normalized weak-cut carrier, global quartic MTA, all-order MTA, WUFC/CPM,
continuum Yang--Mills existence, or a positive mass gap.
\end{firewall}

\section{Append-only record for compact V36}
\label{sec:V36-record}

V36 was appended to the complete compact V35 file.  The exact V35
parent had \(14421\) newline-delimited source records, \(470651\)
bytes, and SHA--256 digest
\[
\texttt{4ce369192e40377c3bd13d69f3beb623bd0de4e7371bfa6e1cf2ecc7c01a7fc8}.
\]
No compact V35 mathematical content was changed.  The sole final
\verb|\end{document}| is restored below.

% =============================================================================
% END APPENDED COMPACT VERSION V36 MATERIAL
% =============================================================================

% The compact V36 document terminator is intentionally disabled here:
% \end{document}

% =============================================================================
% BEGIN APPENDED COMPACT VERSION V37 MATERIAL
% =============================================================================

\chapter{V37: A normalized coefficient-capacity weak-cut theorem}
\label{ch:V37}

\section{Lower-order carrier table}
\label{sec:V37-lower-table}

Fix the four input representations and their global masses
\(\Xi_i\).  For a nonempty row subset \(B\) of size at most three, let
\(\mathcal N_B\) denote the complete positive physical envelope of its
connected Wilson operator before a resolvent, and let
\(\mathcal P_B\) denote the corresponding envelope with its output mass
and one output resolvent inserted.  Isotypic projections and all
multiplicity spaces are retained in these envelopes.  For a singleton,
\(\mathcal N_{\{i\}}\) is the complete moment contraction.

For \(2\le|B|\le3\), put
\begin{equation}
 \mathcal T_B(\theta)
 :=
 \sum_{\sigma\in\mathfrak T(B)}
 \prod_{ij\in E(\sigma)}\theta_{ij}.
\label{eq:V37-lower-tree-polynomial}
\end{equation}

\begin{lemma}[Response-or-numerator table through order three]
\label{lem:V37-lower-table}
Uniformly in the input representations and physical output
multiplicities,
\begin{align}
 \frac{\|\mathcal N_B\|}{\prod_{i\in B}\Xi_i}
 &\le
 C s_\alpha^{|B|-1}\mathcal T_B(\theta),
 &&2\le|B|\le3,
\label{eq:V37-unpayer-table}\\
 \frac{\|\mathcal P_B\|}{\prod_{i\in B}\Xi_i}
 &\le
 C s_\alpha^{|B|-2}\mathcal T_B(\theta),
 &&2\le|B|\le3,
\label{eq:V37-payer-table}\\
 \|\mathcal N_{\{i\}}\|&\le1,
\qquad
 \frac{\|\mathcal P_{\{i\}}\|}{\Xi_i}
 \le\frac C{s_\alpha}.
\label{eq:V37-singleton-table}
\end{align}
Every displayed norm bound also bounds product-state capacity.
\end{lemma}

\begin{proof}
For two rows, the complete covariance numerator is bounded by
\(Cs_\alpha H_{ij}\); division by
\(\Xi_i\Xi_j\) gives
\(Cs_\alpha\theta_{ij}\).  Assigning its output mass to the same
complete binary block and applying the one-link first-moment/gap
estimate removes one power of \(s_\alpha\), giving
\(C\theta_{ij}\).

For three rows, the frozen V50 complete connected-operator theorem,
which is part of the inherited compact V1 status, gives
\[
 \|\mathcal N_B\|
 \le
 Cs_\alpha^2
 \sum_{\sigma\in\mathfrak T(B)}
 \left(\prod_{i\in B}\Xi_i\right)
 \prod_{ij\in E(\sigma)}\theta_{ij}.
\]
Its pointwise one-resolvent response theorem gives the same expression
with \(s_\alpha^2\) replaced by \(s_\alpha\).  Those theorems use the
complete ternary connected operator and raw isotypic pinching, so the
bounds contain no coordinate-word or output multiplicity count.

For one row, the complete Wilson moment is a contraction and
\Cref{lem:V3-spectator-payer} gives
\(\|\mathcal P_{\{i\}}\|\le C/s_\alpha\).
Since \(\Xi_i\ge1\), the last displayed normalized bound follows.
Finally \(\kappa_\otimes(Q)\le\|Q\|\) for every positive envelope.
\end{proof}

\begin{remark}[Only one payer row]
\label{rem:V37-one-payer}
Equation \eqref{eq:V37-payer-table} is used for at most one block of a
row partition.  All other connected blocks use
\eqref{eq:V37-unpayer-table}.  Using the payer table twice would insert
two resolvents and lose exactly the power of \(s_\alpha\) which the
quartic tree needs.
\end{remark}

\section{A composite covariance carrier}
\label{sec:V37-composite-covariance}

Let \(C\mid D\) be a nontrivial cut of the four rows and write
\[
 F_C:=\prod_{i\in C}F_i,\qquad
 F_D:=\prod_{j\in D}F_j
\]
in the fixed row tensor order.  Let
\(\mathcal P_{\operatorname{Cov}(C,D)}\) be the complete positive
physical envelope of
\(\operatorname{Cov}(F_C,F_D)\), with the sole final output mass and
resolvent included before the envelope is formed.

\begin{lemma}[Normalized composite binary response]
\label{lem:V37-composite-response}
\begin{equation}
\boxed{
 \frac{
 \kappa_\otimes(
 \mathcal P_{\operatorname{Cov}(C,D)})}
 {\prod_{i=1}^4\Xi_i}
 \le C\Theta(C,D).}
\label{eq:V37-composite-response}
\end{equation}
\end{lemma}

\begin{proof}
Use the operator martingale covariance formula of frozen V50 before
physical compression.  Replacing coordinate \(e\) in the product
\(F_C\) gives, by telescoping inside the product and contraction of the
remaining row factors, an influence at most the sum of the row
influences over \(i\in C\).  The same holds for \(D\).  Operator
Cauchy--Schwarz therefore gives the complete unresolvented bound
\[
 C s_\alpha
 \sum_{\substack{i\in C\\j\in D}}H_{ij}.
\]
All rows not selected in a particular \(ij\) term remain complete
moment contractions.

Resolve the two composite outputs only after this sum has been formed.
Fusion bounds the final output mass by the sum of their masses.  The
scalar one-resolvent allocation assigns the denominator once, to one
composite side.  The binary payer estimate removes the displayed
\(s_\alpha\); complete moments on the other side are contractions.
At coefficient level, exact Fourier normalization and weighted raw
pinching supply the four row masses.  The \(ij\) summand is therefore
bounded by
\[
 C\left(\prod_{r=1}^4\Xi_r\right)
 \frac{H_{ij}}{\Xi_i\Xi_j}.
\]
Sum over the crossing row pairs and apply
\Cref{thm:V1-weighted-pinching-capacity}.  This is
\eqref{eq:V37-composite-response}.
\end{proof}

\section{Partition products with one resolvent}
\label{sec:V37-partition-products}

Let \(\pi\in\Pi_4\) have \(k=|\pi|\in\{2,3\}\).  On every resolved
product of its block outputs, fusion and the scalar allocation assign
the sole final resolvent to one block \(P\in\pi\).  Define the resulting
positive envelope by using \(\mathcal P_P\) on \(P\) and
\(\mathcal N_B\) on every \(B\ne P\).

\begin{lemma}[A cut-crossing correction costs one normalized edge]
\label{lem:V37-crossing-correction}
Suppose \(\pi\vee\{C,D\}=\boldsymbol1\), so at least one block of
\(\pi\) meets both sides of the cut.  The sum over all possible payer
blocks satisfies
\begin{equation}
\boxed{
 \frac{\kappa_\otimes(\mathcal E_\pi^{\rm pay})}
      {\prod_i\Xi_i}
 \le C\Theta(C,D).}
\label{eq:V37-crossing-correction}
\end{equation}
\end{lemma}

\begin{proof}
Fix the payer \(P\).  Apply the one-sided tensor theorem to the positive
block envelopes on the disjoint row groups, using capacity on one
factor and ordinary norms on all others.  Equations
\eqref{eq:V37-unpayer-table}--%
\eqref{eq:V37-singleton-table} then multiply without any capacity
product.

If a nonsingleton block has size \(b\), its unpayed
\(s_\alpha\)-exponent is \(b-1\).  Paying a nonsingleton block lowers
that exponent by one; paying a singleton contributes exponent \(-1\).
Since \(\sum_{B\in\pi}(|B|-1)=4-k\), the total exponent after exactly
one payment is
\[
 3-k\ge0.
\]
Thus its scalar factor is bounded because
\(0<s_\alpha\le1\).

Every nonsingleton block also contributes a tree monomial in normalized
edges.  Since \(\pi\) connects the cut, one such block \(B_*\) meets
both sides.  Every spanning tree on \(B_*\) contains a crossing edge
\(ij\), while all other \(\theta\)'s lie in \([0,1]\) by
\Cref{lem:V36-theta-range}.  Discarding those other factors bounds the
monomial by \(\theta_{ij}\).  Summing its finitely many tree choices,
the finitely many blocks, and the \(k\le3\) payer choices yields
\(C\Theta(C,D)\).
\end{proof}

\section{The quartic weak-cut identity before norming}
\label{sec:V37-quartic-identity}

Let \(\rho=\{C,D\}\).  The exact product--cumulant identity is
\begin{equation}
 \operatorname{Cov}(F_C,F_D)
 =
 \kappa_4(F_1,F_2,F_3,F_4)
 +
 \sum_{\substack{\pi\in\Pi_4,\ \pi\ne\boldsymbol1\\
                  \pi\vee\rho=\boldsymbol1}}
 \prod_{B\in\pi}\kappa(F_i:i\in B).
\label{eq:V37-product-cumulant}
\end{equation}
It is an operator identity before any output resolution or norm.

\begin{theorem}[Normalized coefficient-capacity weak-cut bound]
\label{thm:V37-normalized-weak-cut}
For the complete quartic positive carrier
\(\mathcal C_4\) of \eqref{eq:V36-normalized-carrier}, and every
nontrivial cut \(C\mid D\),
\begin{equation}
\boxed{
 \mathcal F_4
 =
 \frac{\kappa_\otimes(\mathcal C_4)}
      {\prod_i\Xi_i}
 \le C\Theta(C,D).}
\label{eq:V37-normalized-weak-cut}
\end{equation}
This includes both \(1|3\) and \(2|2\) cuts.
\end{theorem}

\begin{proof}
Solve \eqref{eq:V37-product-cumulant} for the fourth cumulant.  Resolve
all physical block outputs, insert the exact final output mass and sole
resolvent, and only then form the blockwise positive envelope.  Triangle
inequality and subadditivity of capacity bound the first term by
\Cref{lem:V37-composite-response}.

Every correction partition in the sum has two or three blocks.  The
four-singleton partition does not connect the cut, and the one-block
partition is the fourth cumulant already isolated.  Fusion of the block
outputs and \Cref{lem:V2-resolvent-allocation} express each correction
as the finite sum of envelopes with exactly one payer block.  Its
normalized capacity is bounded by
\Cref{lem:V37-crossing-correction}.

There are only fifteen row partitions, independently of supports,
representations, and coordinate words.  Summing their bounds proves
\eqref{eq:V37-normalized-weak-cut}.  The exact Fourier factor and raw
pinching have already been included in
\Cref{lem:V37-composite-response,lem:V37-crossing-correction}, so no
final-channel multiplicity is added.
\end{proof}

\begin{corollary}[The two-cluster hypothesis of V36 is closed]
\label{cor:V37-two-cluster-closed}
The normalized carrier \(\mathcal F_4\) satisfies hypothesis
\textup{(ii)} of \Cref{thm:V36-normalized-graph}.
\end{corollary}

\begin{proof}
Hypothesis \textup{(ii)} is exactly
\eqref{eq:V37-normalized-weak-cut}.
\end{proof}

\begin{assessment}[Why no coordinate transfer appears]
\label{ass:V37-no-transfer}
V37 never estimates the restricted product
\(\bigotimes_r(M_r-K_r)\).  The complete product--cumulant identity is
formed first, and all correction terms are products over disjoint
\emph{row} blocks, not a sum over possible junction coordinates.
Lower-order connected carriers are inserted intact.  This is why the
V36 scalar regression does not apply to the proof.
\end{assessment}

\begin{openproblem}[V38 mandate: normalized three-cluster bound]
\label{op:V38-mandate}
\begin{enumerate}
\item Start from the exact three-cluster product--cumulant identity in
      frozen V135, with the two-row cluster kept intact.
\item Apply the complete ternary MTA to the three cluster variables,
      verifying the composite two-row cluster influence before
      coefficient pinching.
\item Bound the two triple corrections and two covariance products by
      quotient-tree polynomials using the response-or-numerator table.
\item Prove hypothesis \textup{(iii)} of
      \Cref{thm:V36-normalized-graph} without using
      \eqref{eq:V37-normalized-weak-cut} twice on the same resolvent.
\item Then isolate the universal-cap hypothesis \textup{(i)}; do not
      infer it from the V135 unresolvented operator cap.
\end{enumerate}
\end{openproblem}

\begin{firewall}[Claim boundary after compact V37]
\label{fire:V37-current-boundary}
V37 proves the normalized coefficient-capacity weak-cut estimate for
the complete one-resolvent quartic carrier.  It closes only the
two-cluster hypothesis of the V36 normalized graph compiler.  The
three-cluster, universal-cap, and all-weak carrier hypotheses, global
quartic MTA, all-order MTA, WUFC/CPM, continuum Yang--Mills existence,
and a positive mass gap remain open.
\end{firewall}

\section{Append-only record for compact V37}
\label{sec:V37-record}

V37 was appended to the complete compact V36 file.  The exact V36
parent had \(14831\) newline-delimited source records, \(484962\)
bytes, and SHA--256 digest
\[
\texttt{090039b90e3d803d1e0b352d7402d47fde773cdc75092f58d960bf00a2416570}.
\]
No compact V36 mathematical content was changed.  The sole final
\verb|\end{document}| is restored below.

% =============================================================================
% END APPENDED COMPACT VERSION V37 MATERIAL
% =============================================================================

% The compact V37 document terminator is intentionally disabled here:
% \end{document}

% =============================================================================
% BEGIN APPENDED COMPACT VERSION V38 MATERIAL
% =============================================================================

\chapter{V38: The normalized three-cluster carrier}
\label{ch:V38}

\section{Quotient-tree notation}
\label{sec:V38-quotient}

Let
\[
 C_1=\{i,j\},\qquad C_2=\{k\},\qquad C_3=\{\ell\}
\]
be a three-block partition of the four rows.  Put
\[
 \Theta_{ab}:=\Theta(C_a,C_b)
 =\sum_{\substack{u\in C_a\\v\in C_b}}\theta_{uv}
\]
and define
\begin{equation}
 \mathcal T_{\mathcal C}
 :=
 \Theta_{12}\Theta_{13}
 +\Theta_{12}\Theta_{23}
 +\Theta_{13}\Theta_{23}.
\label{eq:V38-quotient-tree}
\end{equation}
This is the spanning-tree polynomial of the three quotient vertices.

\section{Ternary response with one composite row}
\label{sec:V38-composite-ternary}

Write
\[
 G_1=F_iF_j,\qquad G_2=F_k,\qquad G_3=F_\ell
\]
in the fixed tensor order.  Let
\(\mathcal P_{\kappa_3(G_1,G_2,G_3)}\) be the complete positive
physical envelope of their third cumulant with the single final
resolvent.

\begin{lemma}[Composite-row ternary response]
\label{lem:V38-composite-ternary}
\begin{equation}
\boxed{
 \frac{
 \kappa_\otimes(
 \mathcal P_{\kappa_3(G_1,G_2,G_3)})}
 {\Xi_i\Xi_j\Xi_k\Xi_\ell}
 \le C\mathcal T_{\mathcal C}.}
\label{eq:V38-composite-ternary}
\end{equation}
\end{lemma}

\begin{proof}
Use the complete ternary connected-operator proof of frozen V50 with
the first variable left as the product \(F_iF_j\).  Replacement of one
coordinate telescopes inside this product, so its two-sided influence
is bounded by the sum of the \(i\)- and \(j\)-row influences.  Hence
the three operator link weights are bounded respectively by
\[
 Cs_\alpha
 \sum_{u\in C_1,v\in C_2}H_{uv},\qquad
 Cs_\alpha
 \sum_{u\in C_1,v\in C_3}H_{uv},\qquad
 Cs_\alpha H_{k\ell}.
\]
The V50 replicated-tree argument is valid under arbitrary spectator
amplification and therefore gives the sum of their three two-link tree
products before the final output is resolved.

Insert the exact Fourier normalization and regard the input coefficient
on the composite first variable as
\(A_i^{\rm in}\otimes A_j^{\rm in}\).  Its trace norm is the product of
the two row trace norms.  A product state on the original four rows is a
special product state on the three cluster carriers, so the
three-cluster capacity estimate is an upper bound for the required
four-row capacity.  The one-resolvent part of the V50 response theorem
removes one power of \(s_\alpha\) and supplies, for each selected row
pair, its two global row denominators.  Dividing by
\(\Xi_i\Xi_j\Xi_k\Xi_\ell\) turns the three link sums into
\(\Theta_{12},\Theta_{13},\Theta_{23}\).  The three ternary trees give
exactly \eqref{eq:V38-quotient-tree}.
\end{proof}

\begin{remark}[No composite mass denominator]
\label{rem:V38-no-composite-mass}
The proof never introduces a denominator
\((\Xi_i+\Xi_j)^{-1}\).  The composite coefficient is still the tensor
of two original row coefficients, and every link is expanded into its
original \(uv\) marks before normalization.  This is essential when both
quotient edges incident to \(C_1\) select the same original row.
\end{remark}

\section{The exact three-cluster identity}
\label{sec:V38-identity}

The Leonov--Shiryaev product--cumulant formula gives
\begin{align}
 \kappa_3(F_iF_j,F_k,F_\ell)
 ={}&
 \kappa_4(F_i,F_j,F_k,F_\ell)
\notag\\
 &+\mathbb EF_i\,\kappa_3(F_j,F_k,F_\ell)
  +\mathbb EF_j\,\kappa_3(F_i,F_k,F_\ell)
\notag\\
 &+\operatorname{Cov}(F_i,F_k)
       \operatorname{Cov}(F_j,F_\ell)
  +\operatorname{Cov}(F_i,F_\ell)
       \operatorname{Cov}(F_j,F_k).
\label{eq:V38-three-cluster-identity}
\end{align}
This is the frozen V135 identity, now retained before any norm or final
resolvent.

\begin{lemma}[The two triple corrections are quotient-tree paid]
\label{lem:V38-triple-corrections}
After resolving their complete outputs and allocating the sole final
resolvent once, the normalized positive envelopes of
\[
 \mathbb EF_i\,\kappa_3(F_j,F_k,F_\ell),
 \qquad
 \mathbb EF_j\,\kappa_3(F_i,F_k,F_\ell)
\]
are each bounded by \(C\mathcal T_{\mathcal C}\).
\end{lemma}

\begin{proof}
Consider the first term.  Fusion allocates the output payer either to
the singleton mean or to the ternary connected block.  If the ternary
block pays, use \eqref{eq:V37-payer-table} on
\(\{j,k,\ell\}\) and the singleton contraction on \(i\).  If the
singleton pays, use \eqref{eq:V37-singleton-table} and the unpayed
ternary estimate \eqref{eq:V37-unpayer-table}.  In both cases the net
power is \(s_\alpha\le1\).

Every tree on \(\{j,k,\ell\}\) maps to a quotient tree: \(jk\) is a
\(C_1C_2\) edge, \(j\ell\) a \(C_1C_3\) edge, and \(k\ell\) a
\(C_2C_3\) edge.  Its normalized monomial is therefore bounded by the
corresponding term of
\eqref{eq:V38-quotient-tree}.  The second correction is identical with
\(i\) and \(j\) interchanged.
\end{proof}

\begin{lemma}[The covariance products are quotient-tree paid]
\label{lem:V38-covariance-products}
After one-resolvent allocation, the normalized positive envelopes of
the last two products in
\eqref{eq:V38-three-cluster-identity} are bounded by
\(C\mathcal T_{\mathcal C}\).
\end{lemma}

\begin{proof}
For
\(\operatorname{Cov}(F_i,F_k)
  \operatorname{Cov}(F_j,F_\ell)\),
assign the payer to either covariance block.  Equations
\eqref{eq:V37-unpayer-table} and
\eqref{eq:V37-payer-table} give
\[
 Cs_\alpha\theta_{ik}\theta_{j\ell}
 \le C\Theta_{12}\Theta_{13}.
\]
The alternative product gives
\[
 Cs_\alpha\theta_{i\ell}\theta_{jk}
 \le C\Theta_{13}\Theta_{12}.
\]
Thus both allowed one-payer allocations have the displayed bound.  No
covariance is resolved twice.
\end{proof}

\section{The normalized three-cluster theorem}
\label{sec:V38-theorem}

\begin{theorem}[Normalized coefficient-capacity three-cluster bound]
\label{thm:V38-three-cluster}
For every three-block row partition
\(\mathcal C=\{C_1,C_2,C_3\}\),
\begin{equation}
\boxed{
 \mathcal F_4
 \le C
 \sum_{\sigma\in\mathfrak T(\mathcal C)}
 \prod_{CD\in E(\sigma)}\Theta(C,D).}
\label{eq:V38-three-cluster}
\end{equation}
\end{theorem}

\begin{proof}
Solve \eqref{eq:V38-three-cluster-identity} for the fourth cumulant.
Insert the exact final output mass and sole resolvent before forming
positive physical envelopes.  The composite ternary term is bounded by
\Cref{lem:V38-composite-ternary}, the two triple corrections by
\Cref{lem:V38-triple-corrections}, and the two covariance products by
\Cref{lem:V38-covariance-products}.  Subadditivity of capacity and the
finite five-term identity give
\eqref{eq:V38-three-cluster}.  All coefficient blocks are summed by the
same raw pinching theorem, so no cluster or multiplicity count appears.
\end{proof}

\begin{corollary}[The three-cluster hypothesis of V36 is closed]
\label{cor:V38-three-cluster-closed}
The normalized carrier \(\mathcal F_4\) satisfies hypothesis
\textup{(iii)} of \Cref{thm:V36-normalized-graph}.
\end{corollary}

\begin{proof}
Equation \eqref{eq:V38-three-cluster} is exactly that hypothesis.
\end{proof}

\begin{assessment}[One resolvent, not two weak cuts]
\label{ass:V38-one-resolvent}
The covariance-product rows were bounded by one payer table and one
unpayer table.  Applying the V37 weak-cut theorem separately to the two
covariances would place a final resolvent in both and would not prove the
displayed estimate.  V38 therefore closes the three-cluster hypothesis
without circular reuse of the two-cluster result.
\end{assessment}

\begin{openproblem}[V39 mandate: the universal normalized cap]
\label{op:V39-mandate}
\begin{enumerate}
\item Prove
      \(\mathcal F_4\le C\) with the complete final resolvent present.
\item Split low and high total output before using the universal fourth
      cumulant norm.
\item On low output, retain the order-three connected Wilson smallness;
      on high output, use the complete output first moment before
      dividing by \(\prod_i\Xi_i\).
\item Test the cap on the V1 private-mass regression and on protected
      pair, triple, and four-row invariant banks.
\item Do not use the V135 unresolvented cap as though it already paid the
      final denominator.
\end{enumerate}
\end{openproblem}

\begin{firewall}[Claim boundary after compact V38]
\label{fire:V38-current-boundary}
V38 proves the normalized coefficient-capacity three-cluster estimate,
closing hypothesis \textup{(iii)} of the V36 graph compiler.  The
universal normalized cap and all-weak normalized carrier remain open;
therefore global quartic MTA, all-order MTA, WUFC/CPM, continuum
Yang--Mills existence, and a positive mass gap remain open.
\end{firewall}

\section{Append-only record for compact V38}
\label{sec:V38-record}

V38 was appended to the complete compact V37 file.  The exact V37
parent had \(15157\) newline-delimited source records, \(497089\)
bytes, and SHA--256 digest
\[
\texttt{e03c109f0bcd6fdb395b18f6f935c95d04841e3db5506fc91c1994b1a000e24b}.
\]
No compact V37 mathematical content was changed.  The sole final
\verb|\end{document}| is restored below.

% =============================================================================
% END APPENDED COMPACT VERSION V38 MATERIAL
% =============================================================================

% The compact V38 document terminator is intentionally disabled here:
% \end{document}

% =============================================================================
% BEGIN APPENDED COMPACT VERSION V39 MATERIAL
% =============================================================================

\chapter{V39: The universal normalized quartic carrier cap}
\label{ch:V39}

\section{A global output first moment}
\label{sec:V39-output-moment}

Let \(K_\nu^{\rm conn}\) be a complete resolved physical block of the
global four-row connected Wilson operator
\(\mathbf K_I^{\rm conn}\).  The label \(\nu\) includes the final
representation, every physical outer label, and its full multiplicity
space.

\begin{lemma}[Global connected output first moment]
\label{lem:V39-global-first-moment}
Uniformly in supports, representations, and \(\nu\),
\begin{equation}
\boxed{
 s_\alpha\Xi(\boldsymbol U_\nu)
 \|K_\nu^{\rm conn}\|\le C.}
\label{eq:V39-global-first-moment}
\end{equation}
\end{lemma}

\begin{proof}
Use the fifteen-term row moment--cumulant formula before final
recoupling.  Fix one row partition \(\pi\).  At every coordinate,
resolve the output \(S_{B,e}\) of each block \(B\in\pi\).  The norm of
the resulting moment term is bounded by
\[
 Q_{\pi}:=
 \prod_{e}\prod_{B\in\pi}q_{\alpha,S_{B,e}},
\]
because all recoupling maps and final physical compressions are
contractions.

Repeated fusion gives
\[
 \Xi(\boldsymbol U_\nu)
 \le
 C\sum_e\sum_{B\in\pi}A_{S_{B,e}}.
\]
Apply the order-one product-moment compiler to the complete family
\((S_{B,e})\).  The one-link Wilson debit gives
\[
 s_\alpha
 \left(\sum_e\sum_{B\in\pi}A_{S_{B,e}}\right)Q_\pi
 \le C,
\]
independently of the number of coordinates.  Therefore each fixed
partition term satisfies \eqref{eq:V39-global-first-moment}.  Summing
the fifteen absolute Möbius coefficients proves the lemma.  Since the
estimate is made on the complete resolved multiplicity block, no
entrywise channel count is introduced.
\end{proof}

\begin{lemma}[Global fusion bound]
\label{lem:V39-global-fusion}
Every resolved final output obeys
\begin{equation}
\boxed{
 \Xi(\boldsymbol U_\nu)
 \le C_{\rm fus}\sum_{i=1}^4\Xi_i.}
\label{eq:V39-global-fusion}
\end{equation}
\end{lemma}

\begin{proof}
At coordinate \(e\), every final output is contained in the tensor
product of the active row inputs.  The Casimir fusion estimate bounds
its mass by a fixed multiple of the sum of their input masses.  Sum over
coordinates.  Each input contribution occurs in its own row mass
\(\Xi_i\).
\end{proof}

\section{High total output}
\label{sec:V39-high-output}

Let \(c_0>0\) be the fixed threshold from the two-sided product-gap
split.  Denote by \(\mathcal C_{4,\rm hi}\) the part of
\(\mathcal C_4\) on which
\[
 s_\alpha\Xi(\boldsymbol U_\nu)>c_0.
\]

\begin{theorem}[Universal normalized high-output cap]
\label{thm:V39-high-output}
\begin{equation}
\boxed{
 \frac{\kappa_\otimes(\mathcal C_{4,\rm hi})}
      {\prod_i\Xi_i}\le C.}
\label{eq:V39-high-output}
\end{equation}
\end{theorem}

\begin{proof}
The high-output product gap gives
\(1-q_{\alpha,\boldsymbol U_\nu}\ge c>0\).  Hence
\Cref{lem:V39-global-first-moment} gives, on every orthogonal block,
\[
 \frac{\Xi(\boldsymbol U_\nu)}
      {1-q_{\alpha,\boldsymbol U_\nu}}
 \|K_\nu^{\rm conn}\|
 \le\frac C{s_\alpha}.
\]
Therefore
\(\mathcal C_{4,\rm hi}\preccurlyeq Cs_\alpha^{-1}I\).

On the other hand,
\eqref{eq:V39-global-fusion} and the high-output condition imply
\[
 \max_i\Xi_i
 \ge\frac{c_0}{4C_{\rm fus}s_\alpha}.
\]
All row masses are at least one, so
\(\prod_i\Xi_i\ge\max_i\Xi_i\).  Divide the operator bound by this
lower bound and use capacity below operator norm.
\end{proof}

\section{Low output with a hot input}
\label{sec:V39-low-hot-input}

On the complementary output set,
\[
 s_\alpha\Xi(\boldsymbol U_\nu)\le c_0,
\]
the low-output gap gives
\begin{equation}
 \frac{\Xi(\boldsymbol U_\nu)}
      {1-q_{\alpha,\boldsymbol U_\nu}}
 \le\frac C{s_\alpha}.
\label{eq:V39-low-output-resolvent}
\end{equation}

\begin{lemma}[A hot input pays the coarse low-output response]
\label{lem:V39-hot-input}
Fix \(M\ge1\).  If
\(\max_i s_\alpha\Xi_i>M\), then
\begin{equation}
\boxed{
 \frac{\kappa_\otimes(\mathcal C_{4,\rm lo})}
      {\prod_i\Xi_i}
 \le\frac{C}{M}.}
\label{eq:V39-hot-input}
\end{equation}
\end{lemma}

\begin{proof}
The fifteen-term moment--cumulant formula and contraction of all row
moments give the universal block bound
\(\|K_\nu^{\rm conn}\|\le C\).  By
\eqref{eq:V39-low-output-resolvent},
\[
 \mathcal C_{4,\rm lo}\preccurlyeq C s_\alpha^{-1}I.
\]
If \(s_\alpha\Xi_c>M\), then
\(\prod_i\Xi_i\ge\Xi_c>M/s_\alpha\).  Divide and take product-state
capacity.
\end{proof}

\section{Low output with thermal inputs}
\label{sec:V39-low-thermal}

\begin{lemma}[Tree excess matches the thermal powers]
\label{lem:V39-tree-excess}
For every four-row tree \(\tau\),
\begin{equation}
 \frac{\prod_{ij\in E(\tau)}H_{ij}}
      {\prod_i\Xi_i}
 \le
 \prod_i\Xi_i^{\deg_\tau(i)-1}.
\label{eq:V39-tree-excess}
\end{equation}
If \(s_\alpha\Xi_i\le M\) for every row, then
\begin{equation}
\boxed{
 s_\alpha^2
 \frac{\prod_{ij\in E(\tau)}H_{ij}}
      {\prod_i\Xi_i}
 \le M^2.}
\label{eq:V39-thermal-tree}
\end{equation}
\end{lemma}

\begin{proof}
Use \(H_{ij}\le\Xi_i\Xi_j\) on every tree edge and collect row powers.
The tree degree identity gives
\[
 \sum_i(\deg_\tau(i)-1)=2.
\]
Under the thermal hypothesis, each factor
\(\Xi_i^{\deg_\tau(i)-1}\) is bounded by the corresponding power of
\(M/s_\alpha\).  Their total exponent is two, proving
\eqref{eq:V39-thermal-tree}.
\end{proof}

\begin{theorem}[Universal normalized low-output thermal cap]
\label{thm:V39-low-thermal}
If \(s_\alpha\Xi_i\le M\) for all four rows, then
\begin{equation}
\boxed{
 \frac{\kappa_\otimes(\mathcal C_{4,\rm lo})}
      {\prod_i\Xi_i}
 \le C_M.}
\label{eq:V39-low-thermal}
\end{equation}
\end{theorem}

\begin{proof}
The frozen V135 global numerator theorem, followed by compression to a
physical block, gives
\[
 \|K_\nu^{\rm conn}\|
 \le
 Cs_\alpha^3
 \sum_{\tau\in\mathfrak T_4}
 \prod_{ij\in E(\tau)}H_{ij}.
\]
Multiply by the low-output quotient
\eqref{eq:V39-low-output-resolvent}.  Since the physical projections
are orthogonal, the resulting positive carrier has operator norm at
most
\[
 Cs_\alpha^2
 \sum_{\tau\in\mathfrak T_4}
 \prod_{ij\in E(\tau)}H_{ij}.
\]
Divide by \(\prod_i\Xi_i\), apply
\eqref{eq:V39-thermal-tree} to each of the sixteen labeled trees, and
use capacity below operator norm.
\end{proof}

\section{The universal-cap theorem}
\label{sec:V39-universal}

\begin{theorem}[Universal normalized complete-carrier cap]
\label{thm:V39-universal-cap}
\begin{equation}
\boxed{\mathcal F_4\le C.}
\label{eq:V39-universal-cap}
\end{equation}
\end{theorem}

\begin{proof}
Split into high and low total output.  The high part is
\Cref{thm:V39-high-output}.  On the low part fix, for example, \(M=2\).
If one input is hot, use \Cref{lem:V39-hot-input}; if all are thermal,
use \Cref{thm:V39-low-thermal}.  Capacity is subadditive over the two
orthogonal output regions.
\end{proof}

\begin{corollary}[Hypothesis \textup{(i)} of the normalized graph
compiler is closed]
\label{cor:V39-universal-closed}
The normalized carrier \(\mathcal F_4\) satisfies hypothesis
\textup{(i)} of \Cref{thm:V36-normalized-graph}.
\end{corollary}

\begin{proof}
This is \eqref{eq:V39-universal-cap}.
\end{proof}

\section{Regression checks}
\label{sec:V39-regressions}

\begin{assessment}[Private-mass and invariant-bank tests]
\label{ass:V39-regressions}
In the V1 private-mass regression
\(\Xi_1=L\), \(\Xi_2=\Xi_3=\Xi_4=1\), the hot-input branch gives
\(O((s_\alpha L)^{-1})\) once \(s_\alpha L>M\); below that threshold
the thermal estimate uses exactly
\(s_\alpha^2L^2\le M^2\) for a star centered at row one.

For protected pair, triple, or four-row invariant banks, a large
cancelling input lies in the hot-input branch and is paid by the global
row product.  The proof never counts protected multiplicities.  The
sharper invariant capacities of V30--V34 remain necessary for marked
tree estimates, but no additional rarity is required for the universal
cap.
\end{assessment}

\begin{openproblem}[V40 mandate: audit and the all-weak normalized
carrier]
\label{op:V40-mandate}
\begin{enumerate}
\item Perform the mandatory read-only audit of the newest active SM and
      NS notes.
\item Return to the frozen V134 low-activity connected-word expansion
      before its operator norm.
\item Replace its unnormalized local activities by the normalized
      overlaps \(\theta_{ij}\), retaining the exact coefficient carrier
      and one final resolvent.
\item Prove hypothesis \textup{(iv)} of
      \Cref{thm:V36-normalized-graph} when all
      \(\theta_{ij}<\delta_0\), or identify the precise local letter
      whose activity is small in \(s_\alpha H_{ij}\) but not in
      \(\theta_{ij}\).
\item If hypothesis \textup{(iv)} closes, combine V37--V40 immediately
      to obtain global quartic MTA.
\end{enumerate}
\end{openproblem}

\begin{firewall}[Claim boundary after compact V39]
\label{fire:V39-current-boundary}
V39 proves the universal normalized cap for the complete
coefficient-pinched one-resolvent quartic carrier.  Together V37--V39
close hypotheses \textup{(i)--(iii)} of the V36 normalized graph
compiler.  The all-weak normalized carrier remains open, so global
quartic MTA, all-order MTA, WUFC/CPM, continuum Yang--Mills existence,
and a positive mass gap remain open.
\end{firewall}

\section{Append-only record for compact V39}
\label{sec:V39-record}

V39 was appended to the complete compact V38 file.  The exact V38
parent had \(15423\) newline-delimited source records, \(506452\)
bytes, and SHA--256 digest
\[
\texttt{f534382b9daa031feaa35bfea5bf37f1a49091254ada99cccae4cdf72180523b}.
\]
No compact V38 mathematical content was changed.  The sole final
\verb|\end{document}| is restored below.

% =============================================================================
% END APPENDED COMPACT VERSION V39 MATERIAL
% =============================================================================

% The compact V39 document terminator is intentionally disabled here:
% \end{document}

% =============================================================================
% BEGIN APPENDED COMPACT VERSION V40 MATERIAL
% =============================================================================

\chapter{V40: Companion audit and the thermal-tree part of the
all-weak carrier}
\label{ch:V40}

\section{Mandatory companion audit}
\label{sec:V40-audit}

The newest active companion sources inspected read-only were:
\begin{center}
\begin{tabular}{p{0.42\linewidth}r r p{0.30\linewidth}}
\toprule
file & lines & bytes & SHA--256\\
\midrule
\texttt{Renewal\_SM\_Einstein\_Continuum\_Research\_Notes\_V45.tex}
& \(16563\) & \(592714\)
& \texttt{a0da83031bb59350a14687dc689346a485953f959b461088bdc9634be18e104a}\\
\texttt{Renewal\_NS\_Stability\_Research\_Notes\_V31.tex}
& \(23052\) & \(887537\)
& \texttt{f1121b62238ad5491bb7cf03635ea9934942edf573ed8faaa9ee79e1d38a6696}\\
\bottomrule
\end{tabular}
\end{center}

The NS endpoint is unchanged.  Its first-moment and no-double-payment
warnings remain exactly as recorded in V35.

SM V41 derives an integrable-memory criterion from a sufficiently
subtracted spectral density and gives a nondecaying-density regression.
V42 proves the three-subtraction threshold in free scalar, Dirac, and
physical gauge sectors.  V43 permits every fixed logarithmic order but
isolates a nonperturbative Dini majorant.  V44 transfers the free result
to bounded-geometry Hadamard symbols and obtains order minus two after
three subtractions.  V45 assembles compatible renormalization forests
before taking symbol norms.  Its explicit nonorthogonal-projection
regression has tensor norm \(3^N\), independently confirming the
restricted-factor obstruction proved in YM V36.

\begin{assessment}[V40 audit transfer]
\label{ass:V40-audit-transfer}
No SM or NS inequality is imported.  The proof-order consequence is
strict: the all-weak YM branch must be estimated from the complete
connected quartic carrier.  Neither the no-four product
\(\bigotimes(M-K_4)\) nor individual topology witnesses may be assigned
independent norms.  Fixed-order companion closure also supplies no
all-order Yang--Mills summability theorem.
\end{assessment}

\section{Tree excess and the physical carrier}
\label{sec:V40-tree-excess}

For \(\tau\in\mathfrak T_4\), write
\[
 d_i=\deg_\tau(i),
 \qquad
 B_\tau:=\prod_{i=1}^4\Xi_i^{d_i-1}.
\]
Then
\begin{equation}
 \mathfrak W_\tau^{\rm glob}
 =\frac{\prod_{ij\in E(\tau)}H_{ij}}{B_\tau}.
\label{eq:V40-tree-weight-B}
\end{equation}
The degree identity gives
\begin{equation}
 \sum_i(d_i-1)=2.
\label{eq:V40-excess-two}
\end{equation}

\begin{lemma}[A thermal tree converts the V135 numerator exactly]
\label{lem:V40-thermal-tree-conversion}
If
\begin{equation}
 s_\alpha^2B_\tau\le M^2,
\label{eq:V40-tree-thermal}
\end{equation}
then
\begin{equation}
\boxed{
 s_\alpha^2\prod_{ij\in E(\tau)}H_{ij}
 \le M^2\mathfrak W_\tau^{\rm glob}.}
\label{eq:V40-thermal-conversion}
\end{equation}
\end{lemma}

\begin{proof}
Multiply \eqref{eq:V40-tree-weight-B} by
\(s_\alpha^2B_\tau\) and use
\eqref{eq:V40-tree-thermal}.
\end{proof}

\section{Complete closure when every input is thermal}
\label{sec:V40-all-input-thermal}

\begin{theorem}[Global quartic MTA on the thermal-input branch]
\label{thm:V40-thermal-input-MTA}
Fix \(M\ge1\).  If
\begin{equation}
 s_\alpha\Xi_i\le M
 \qquad(1\le i\le4),
\label{eq:V40-all-input-thermal}
\end{equation}
then the complete physical carrier satisfies
\begin{equation}
\boxed{
 \kappa_\otimes(\mathcal C_4)
 \le
 C_M\sum_{\tau\in\mathfrak T_4}
 \mathfrak W_\tau^{\rm glob}.}
\label{eq:V40-thermal-input-MTA}
\end{equation}
No assumption on the normalized overlaps \(\theta_{ij}\) is needed.
\end{theorem}

\begin{proof}
First restrict to low total output
\(s_\alpha\Xi(\boldsymbol U_\nu)\le c_0\).  The low-output resolvent
bound and the frozen V135 numerator give, blockwise,
\[
 \frac{\Xi(\boldsymbol U_\nu)}
      {1-q_{\alpha,\boldsymbol U_\nu}}
 \|K_\nu^{\rm conn}\|
 \le
 Cs_\alpha^2
 \sum_{\tau}\prod_{ij\in E(\tau)}H_{ij}.
\]

On high total output the denominator has a fixed lower bound.  By
\Cref{lem:V39-global-fusion} and
\eqref{eq:V40-all-input-thermal},
\[
 \Xi(\boldsymbol U_\nu)
 \le C\sum_i\Xi_i\le\frac{C_M}{s_\alpha}.
\]
Multiplication by the V135
\(Cs_\alpha^3\sum_\tau\prod H_{ij}\) numerator gives the same
\(C_Ms_\alpha^2\sum_\tau\prod H_{ij}\) bound.

By \eqref{eq:V40-excess-two}, the thermal-input hypothesis implies
\[
 s_\alpha^2B_\tau
 =
 \prod_i(s_\alpha\Xi_i)^{d_i-1}
 \le M^2
\]
for every tree.  Apply
\Cref{lem:V40-thermal-tree-conversion}.  The physical blocks are
orthogonal, so the pointwise scalar majorant is an operator majorant of
\(\mathcal C_4\).  Taking product-state capacity proves
\eqref{eq:V40-thermal-input-MTA}.
\end{proof}

\begin{corollary}[All-weak thermal-input closure]
\label{cor:V40-all-weak-thermal}
Hypothesis \textup{(iv)} of
\Cref{thm:V36-normalized-graph} holds on the subfamily satisfying
\eqref{eq:V40-all-input-thermal}, with no loss as
\(\max_{ij}\theta_{ij}\to0\).
\end{corollary}

\begin{proof}
Divide \eqref{eq:V40-thermal-input-MTA} by
\(\prod_i\Xi_i\) and use
\eqref{eq:V36-global-marked-tree}.
\end{proof}

\section{A treewise low-output refinement}
\label{sec:V40-treewise-low}

Even when a leaf row is hot, a fixed low-output tree is closed whenever
its two internal excess powers are thermal.

\begin{theorem}[Low-output thermal-tree closure]
\label{thm:V40-low-thermal-tree}
On the low-output sector, the part of the V135 scalar tree majorant
indexed by any \(\tau\) satisfying
\eqref{eq:V40-tree-thermal} is bounded by
\[
 C_M\mathfrak W_\tau^{\rm glob}.
\]
Consequently all low-output thermal-tree pieces satisfy their marked
tree weights even if one or more degree-one rows have
\(s_\alpha\Xi_i>M\).
\end{theorem}

\begin{proof}
The low-output quotient always turns the V135 tree numerator into
\[
 Cs_\alpha^2\prod_{ij\in E(\tau)}H_{ij}.
\]
Apply \eqref{eq:V40-thermal-conversion}.  A leaf has exponent
\(d_i-1=0\), so its mass is absent from
\eqref{eq:V40-tree-thermal}.
\end{proof}

\begin{lemma}[Failure selects a hot internal row]
\label{lem:V40-hot-internal-row}
If a tree fails \eqref{eq:V40-tree-thermal}, then some internal vertex
\(i\), \(d_i\ge2\), satisfies
\begin{equation}
\boxed{s_\alpha\Xi_i>M.}
\label{eq:V40-hot-internal-row}
\end{equation}
For a star this is its center.  For a path it is at least one of its
two internal vertices.
\end{lemma}

\begin{proof}
If every internal row satisfied
\(s_\alpha\Xi_i\le M\), then
\[
 s_\alpha^2B_\tau
 =
 \prod_i(s_\alpha\Xi_i)^{d_i-1}
 \le M^{\sum_i(d_i-1)}=M^2,
\]
contrary to failure.
\end{proof}

\section{The narrowed all-weak obstruction}
\label{sec:V40-narrowed-obstruction}

\begin{assessment}[What remains after V40]
\label{ass:V40-narrowed}
In the all-weak regime \(\theta_{ij}<\delta_0\), the last unproved part
of hypothesis \textup{(iv)} has the following sharper form:
\begin{enumerate}
\item on low output, a selected tree has a hot internal row as in
      \eqref{eq:V40-hot-internal-row};
\item on high output, at least one input is hot, since the all-input
      thermal branch is already closed.
\end{enumerate}
Thus the remaining issue is not weak overlap alone.  It is conversion of
a hot global row mass into its missing repeated normalized marks while
the complete connected operator stays assembled.
\end{assessment}

For a row \(i\), split its mass into
\[
 \Xi_i=p_i+h_i,
\]
where \(p_i\) is supported on coordinates private to row \(i\) and
\(h_i\) on coordinates shared with at least one other row.  Unlike the
one-selected-junction setting, no local junction mass is split off:
every coordinate letter remains inside the complete connected operator.

\begin{openproblem}[V41 mandate: hot internal-row dichotomy]
\label{op:V41-mandate}
\begin{enumerate}
\item Prove the exact multilinear factorization of the complete global
      fourth cumulant by the private Wilson multiplier of a selected hot
      row.
\item If \(p_i\ge\Xi_i/2\), use first or second private product moments
      according to \(d_i-1\), while retaining the global connected
      operator and sole final resolvent.
\item If \(h_i>\Xi_i/2\), split its shared mass by partner and incidence,
      but form one common stock before selecting a coordinate bank.
\item Use the all-weak inequalities
      \(H_{ij}<\delta_0\Xi_i\Xi_j\) only after the private/shared
      dichotomy; do not infer \(h_i\ll\Xi_i\) without controlling the
      partner masses.
\item Close the private-dominant hot-row branch and identify the exact
      shared-dominant residual for V42.
\end{enumerate}
\end{openproblem}

\begin{firewall}[Claim boundary after compact V40]
\label{fire:V40-current-boundary}
V40 completes the mandatory companion audit, proves global quartic MTA
when all four inputs are thermal, and closes every low-output tree whose
two excess input powers are thermal.  It does not close the hot
internal-row all-weak residual.  Hence hypothesis \textup{(iv)} of the
V36 compiler, global quartic MTA, all-order MTA, WUFC/CPM, continuum
Yang--Mills existence, and a positive mass gap remain open.
\end{firewall}

\section{Append-only record for compact V40}
\label{sec:V40-record}

V40 was appended to the complete compact V39 file.  The exact V39
parent had \(15757\) newline-delimited source records, \(516585\)
bytes, and SHA--256 digest
\[
\texttt{756300928db8e5671a3440fd50fe0a565829cac1bb3b6ea426aaab2912f09c89}.
\]
No compact V39 mathematical content was changed.  The sole final
\verb|\end{document}| is restored below.

% =============================================================================
% END APPENDED COMPACT VERSION V40 MATERIAL
% =============================================================================

% The compact V40 document terminator is intentionally disabled here:
% \end{document}

% =============================================================================
% BEGIN APPENDED COMPACT VERSION V41 MATERIAL
% =============================================================================

\chapter{V41: Exact private factorization and the hot-tree dichotomy}
\label{ch:V41}

\section{Purpose and proof boundary}
\label{sec:V41-purpose}

V40 reduced the unproved all-weak quartic carrier to hot global row
masses.  This chapter performs the first exact split of such a mass.
Coordinates private to one row are removed algebraically from the
complete fourth cumulant, before norms and before the output resolvent.
Their product moments then pay the two excess powers of a quartic tree.
On high output the same calculation also pays the part of the output
mass inherited from private coordinates.

This is not a return to the one-selected-junction carrier.  There is no
selected coordinate, and no term of the global cumulant is discarded.
All moment partitions remain assembled throughout the factorization.

\section{Private banks and exact moment factorization}
\label{sec:V41-private-factorization}

Work in one fixed physical input and output block of the global
four-row carrier.  Let \(P_i\) be the set of coordinate letters that
occur in row \(i\) and in no other row.  At \(e\in P_i\), denote the
normalized Wilson expectation in the input representation by
\(q_{i,e}\).  Put
\begin{equation}
 p_i:=\sum_{e\in P_i}\bigl(1+C_2(R_{i,e})\bigr),
 \qquad
 q_i:=\prod_{e\in P_i}q_{i,e},
 \qquad
 q_{\rm priv}:=\prod_{i=1}^4q_i.
\label{eq:V41-private-data}
\end{equation}
Empty products equal one.  Thus the V40 row split is
\[
 \Xi_i=p_i+h_i,
\]
where \(h_i\) contains every nonprivate coordinate letter of row \(i\).
In the established heat-kernel range,
\begin{equation}
 0\le q_i\le1,
 \qquad
 (s_\alpha p_i)^kq_i\le C_k
 \quad (k=1,2,3)
\label{eq:V41-row-private-moments}
\end{equation}
whenever \(p_i>0\).  These are
\Cref{cor:V7-private-products} applied separately to the four disjoint
private banks.

Let \(F_i\) be the row insertion in the fixed block and let
\(F_i^{\rm sh}\) be the same insertion with all coordinates in \(P_i\)
integrated out.  Since the coordinate measure is a product and a
private coordinate occurs in only one row, for every nonempty
\(A\subseteq[4]\),
\begin{equation}
 \mathbb E\!\left[\prod_{i\in A}F_i\right]
 =
 \left(\prod_{i\in A}q_i\right)
 \mathbb E\!\left[\prod_{i\in A}F_i^{\rm sh}\right].
\label{eq:V41-submoment-factorization}
\end{equation}
The equality is an equality of the remaining output-block operators;
the factors \(q_i\) are scalars.

\begin{lemma}[Exact private multiplier of the complete fourth
cumulant]
\label{lem:V41-exact-private-factor}
In every fixed physical block,
\begin{equation}
\boxed{
 \kappa_4(F_1,F_2,F_3,F_4)
 =
 q_{\rm priv}\,
 \kappa_4(F_1^{\rm sh},F_2^{\rm sh},
          F_3^{\rm sh},F_4^{\rm sh}).}
\label{eq:V41-exact-private-factor}
\end{equation}
The equality holds before resolving the final output norm and before
using a tree estimate.
\end{lemma}

\begin{proof}
Use the moment--partition identity
\[
 \kappa_4(F_1,F_2,F_3,F_4)
 =
 \sum_{\pi\in\Pi([4])}
 (|\pi|-1)!(-1)^{|\pi|-1}
 \prod_{B\in\pi}
 \mathbb E\!\left[\prod_{i\in B}F_i\right].
\]
Insert \eqref{eq:V41-submoment-factorization}.  The blocks of a
partition are disjoint and cover \([4]\), hence
\[
 \prod_{B\in\pi}\prod_{i\in B}q_i
 =\prod_{i=1}^4q_i=q_{\rm priv}
\]
for every \(\pi\).  This common scalar factors out of the complete
partition sum.  What remains is exactly the moment--partition formula
for the four private-free rows.  No triangle inequality has been used.
\end{proof}

\begin{corollary}[Private-refined frozen tree numerator]
\label{cor:V41-private-refined-numerator}
Let
\[
 L_\tau:=\prod_{ij\in E(\tau)}H_{ij}.
\]
The frozen V135 connected estimate may be retained in the form
\begin{equation}
\boxed{
 \|K_\nu^{\rm conn}\|
 \le
 Cs_\alpha^3q_{\rm priv}
 \sum_{\tau\in\mathfrak T_4}L_\tau.}
\label{eq:V41-private-refined-numerator}
\end{equation}
For a low-output block its coefficient-weighted version is
\begin{equation}
 \frac{\Xi(\boldsymbol U_\nu)}
      {1-q_{\alpha,\boldsymbol U_\nu}}
 \|K_\nu^{\rm conn}\|
 \le
 Cs_\alpha^2q_{\rm priv}
 \sum_{\tau\in\mathfrak T_4}L_\tau.
\label{eq:V41-private-refined-low}
\end{equation}
\end{corollary}

\begin{proof}
Apply the V135 connected tree estimate to the private-free cumulant in
\eqref{eq:V41-exact-private-factor}.  Private coordinates cannot
contribute to \(H_{ij}\), since they occur in only one row, so the
overlap numerators are unchanged.  This proves
\eqref{eq:V41-private-refined-numerator}.  The low-output resolvent
estimate used in V40 removes one factor of \(s_\alpha\), proving
\eqref{eq:V41-private-refined-low}.
\end{proof}

\begin{firewall}[What exact factorization does not permit]
\label{fire:V41-factorization-boundary}
Equation \eqref{eq:V41-exact-private-factor} permits the common private
multiplier to be retained on the complete carrier.  It does not permit
one to attach different private factors to separate cumulant
partitions, output blocks, or topology witnesses.  In particular the
nonprivate cumulant on the right remains the complete connected
operator.
\end{firewall}

\section{A treewise private-moment compiler}
\label{sec:V41-tree-compiler}

For \(\tau\in\mathfrak T_4\), put
\[
 e_i(\tau):=d_i(\tau)-1,
 \qquad
 B_\tau=\prod_i\Xi_i^{e_i(\tau)}.
\]
Recall that \(e_i\in\{0,1,2\}\) and
\(\sum_i e_i=2\).  Fix \(M\ge1\).

\begin{definition}[Privately controlled tree]
\label{def:V41-private-controlled}
A labelled quartic tree \(\tau\) is \(M\)-privately controlled if,
for every \(i\) with \(e_i(\tau)>0\), at least one of
\begin{equation}
 s_\alpha\Xi_i\le M,
 \qquad
 p_i\ge\frac12\Xi_i
\label{eq:V41-controlled-alternatives}
\end{equation}
holds.
\end{definition}

\begin{lemma}[Two excess powers]
\label{lem:V41-two-excess-powers}
If \(\tau\) is \(M\)-privately controlled, then
\begin{equation}
\boxed{
 s_\alpha^2q_{\rm priv}B_\tau\le C_M.}
\label{eq:V41-two-excess-powers}
\end{equation}
\end{lemma}

\begin{proof}
There are only two degree patterns.

If \(\tau\) is a star with center \(c\), then
\(B_\tau=\Xi_c^2\).  In the thermal alternative,
\[
 s_\alpha^2q_{\rm priv}B_\tau
 \le(s_\alpha\Xi_c)^2\le M^2.
\]
In the private alternative \(\Xi_c\le2p_c\), and
\eqref{eq:V41-row-private-moments} with \(k=2\) gives
\[
 s_\alpha^2q_{\rm priv}B_\tau
 \le4(s_\alpha p_c)^2q_c\prod_{i\ne c}q_i
 \le4C_2.
\]

If \(\tau\) is a path with internal vertices \(u,v\), then
\(B_\tau=\Xi_u\Xi_v\).  A thermal internal factor contributes at most
\(M\).  A private-dominant internal factor satisfies
\[
 s_\alpha\Xi_iq_i\le2s_\alpha p_iq_i\le2C_1.
\]
If both internal vertices are private-dominant, apply the last
inequality independently to their disjoint banks.  The remaining
\(q_i\)'s are at most one.  This proves the result in all four
thermal/private combinations.
\end{proof}

\begin{lemma}[A third private payer]
\label{lem:V41-third-private-payer}
If \(\tau\) is \(M\)-privately controlled, then
\begin{equation}
\boxed{
 s_\alpha^3q_{\rm priv}
 \left(\sum_{r=1}^4p_r\right)B_\tau
 \le C_M.}
\label{eq:V41-third-private-payer}
\end{equation}
\end{lemma}

\begin{proof}
It suffices to bound the summand with payer \(r\):
\begin{equation}
 q_{\rm priv}(s_\alpha p_r)
 \prod_i(s_\alpha\Xi_i)^{e_i(\tau)}.
\label{eq:V41-labelled-private-payer}
\end{equation}
For each internal row in the thermal alternative, replace its excess
factor by \(M^{e_i}\).  For each internal row in the private
alternative, use \(\Xi_i\le2p_i\).  If \(r\) is itself thermal and
internal, one may instead use
\(s_\alpha p_r\le s_\alpha\Xi_r\le M\).
All remaining powers of \(s_\alpha p_i\) are charged to \(q_i\) using
\eqref{eq:V41-row-private-moments}.

The exponent charged to any one private bank is at most three:
for a star it is at most \(2+1\), and for a path it is at most
\(1+1\).  If \(r\) is a leaf, its exponent is one and is disjoint from
the internal charges.  Thus only \(C_1,C_2,C_3\) occur.  The number of
rows and the number of labelled payer choices are fixed, so summing
over \(r\) proves \eqref{eq:V41-third-private-payer}.
\end{proof}

\section{Closed low-output and private-output branches}
\label{sec:V41-closed-branches}

\begin{theorem}[Low-output closure for every privately controlled
tree]
\label{thm:V41-low-private-controlled}
On the low-output sector, the scalar tree contribution indexed by an
\(M\)-privately controlled \(\tau\) obeys
\begin{equation}
\boxed{
 Cs_\alpha^2q_{\rm priv}L_\tau
 \le C_M\mathfrak W_\tau^{\rm glob}.}
\label{eq:V41-low-private-controlled}
\end{equation}
Thus the contribution satisfies the required global marked-tree
weight, even when an internal row is hot, provided every hot internal
row is private-dominant.
\end{theorem}

\begin{proof}
By \eqref{eq:V40-tree-weight-B},
\[
 s_\alpha^2q_{\rm priv}L_\tau
 =
 \bigl(s_\alpha^2q_{\rm priv}B_\tau\bigr)
 \mathfrak W_\tau^{\rm glob}.
\]
Apply \Cref{lem:V41-two-excess-powers}.  Since the common multiplier
was extracted from the complete cumulant in
\Cref{lem:V41-exact-private-factor}, this scalar conversion is made
only after the connected carrier has been assembled.
\end{proof}

For high output the private part of the output mass can also be
identified exactly.  Private coordinates have only one input
representation and no other row with which to fuse.  Their output
representation therefore equals that input representation.  Hence,
in every physical output block \(\nu\),
\begin{equation}
 \Xi(\boldsymbol U_\nu)
 =
 P+Z_\nu,
 \qquad
 P:=\sum_{i=1}^4p_i,
 \qquad
 Z_\nu\ge0,
\label{eq:V41-output-mass-split}
\end{equation}
where \(Z_\nu\) is the total mass on nonprivate coordinates.

\begin{theorem}[The private part of every high-output payer]
\label{thm:V41-high-private-output}
On the high-output sector, fix an \(M\)-privately controlled tree
\(\tau\).  The portion of the scalar coefficient majorant obtained by
using \(P\) in \eqref{eq:V41-output-mass-split} satisfies
\begin{equation}
\boxed{
 Cs_\alpha^3Pq_{\rm priv}L_\tau
 \le C_M\mathfrak W_\tau^{\rm glob}.}
\label{eq:V41-high-private-output}
\end{equation}
\end{theorem}

\begin{proof}
The high-output denominator has a fixed positive lower bound.
Equation \eqref{eq:V41-private-refined-numerator}, followed only now by
the exact scalar split \eqref{eq:V41-output-mass-split}, gives the
left side of \eqref{eq:V41-high-private-output}.  Since
\[
 s_\alpha^3Pq_{\rm priv}L_\tau
 =
 \bigl(s_\alpha^3Pq_{\rm priv}B_\tau\bigr)
 \mathfrak W_\tau^{\rm glob},
\]
\Cref{lem:V41-third-private-payer} proves the claim.
\end{proof}

\begin{remark}[The output split is bookkeeping, not a new operator
decomposition]
\label{rem:V41-output-bookkeeping}
The two nonnegative scalars \(P\) and \(Z_\nu\) split the coefficient
\(\Xi(\boldsymbol U_\nu)\) multiplying the same connected block.
No assertion is made that \(K_\nu^{\rm conn}\) itself decomposes into
private-output and shared-output operators.  Consequently the theorem
does not duplicate the final resolvent.
\end{remark}

\section{Exact shared-dominant residual}
\label{sec:V41-shared-residual}

\begin{lemma}[Failure of private control selects shared hot mass]
\label{lem:V41-shared-hot}
If a tree \(\tau\) is not \(M\)-privately controlled, then some
internal row \(i\) satisfies
\begin{equation}
\boxed{
 s_\alpha\Xi_i>M,
 \qquad
 h_i>\frac12\Xi_i.}
\label{eq:V41-shared-hot}
\end{equation}
\end{lemma}

\begin{proof}
Failure of \Cref{def:V41-private-controlled} supplies an internal row
for which both alternatives in
\eqref{eq:V41-controlled-alternatives} fail.  Thus
\(s_\alpha\Xi_i>M\) and \(p_i<\Xi_i/2\).  Since
\(\Xi_i=p_i+h_i\), the second inequality is exactly
\(h_i>\Xi_i/2\).
\end{proof}

\begin{proposition}[Exhaustive post-V41 all-weak ledger]
\label{prop:V41-exhaustive-ledger}
After V40 and the estimates above, every still-unproved scalar piece
of the complete all-weak quartic carrier belongs to at least one of
the following two classes:
\begin{enumerate}
\item a low-output tree with a shared-dominant hot internal row
      satisfying \eqref{eq:V41-shared-hot};
\item a high-output contribution carrying the nonprivate output mass
      \(Z_\nu\), or a private-output tree with a shared-dominant hot
      internal row satisfying \eqref{eq:V41-shared-hot}.
\end{enumerate}
In particular there is no remaining low-output tree whose internal
rows are all thermal or private-dominant, and there is no remaining
high-output \(P\)-payer on such a tree.
\end{proposition}

\begin{proof}
For low output, a privately controlled tree is closed by
\Cref{thm:V41-low-private-controlled}; failure gives
\Cref{lem:V41-shared-hot}.  On high output, split the coefficient by
\eqref{eq:V41-output-mass-split}.  The \(P\)-part is closed by
\Cref{thm:V41-high-private-output} whenever the tree is privately
controlled, while failure again gives
\Cref{lem:V41-shared-hot}.  The estimate has not yet treated the
\(Z_\nu\)-part, so it is retained explicitly.  These alternatives are
exhaustive.
\end{proof}

\begin{assessment}[Direction forced for V42]
\label{ass:V41-V42-direction}
The residual is now genuinely nonprivate.  For a selected row \(i\)
in \eqref{eq:V41-shared-hot}, write its shared mass as
\[
 h_i=\sum_{\varnothing\ne A\subseteq[4]\setminus\{i\}}h_{i;A},
\]
where \(h_{i;A}\) is carried by coordinates whose other active rows
are exactly \(A\).  The next step must combine all pair, triple, and
quadruple incidence classes into one common stock before choosing a
bank.  On high output it must simultaneously retain the ancestry of
\(Z_\nu\).  Applying separate incidence-class norms would repeat the
restricted-factor error ruled out in V36.
\end{assessment}

\begin{openproblem}[V42 mandate: common shared stock]
\label{op:V42-mandate}
\begin{enumerate}
\item For the shared-dominant row selected by
      \eqref{eq:V41-shared-hot}, form a single finite incidence
      hypergraph containing all supports of size two, three, and four.
\item Prove a common balanced/unbalanced stock alternative without
      assigning separate norms to overlapping banks.
\item In the balanced branch, use the protected invariant-projector
      capacities of V33--V34 and preserve one global denominator.
\item In the unbalanced branch, retain the output representation
      carried by the dominant row; this must pay both its repeated
      tree incidence and the \(Z_\nu\)-payer when present.
\item Convert the result to the original global marked-tree weights
      before invoking the finite graph reduction of V36.
\end{enumerate}
\end{openproblem}

\begin{firewall}[Claim boundary after compact V41]
\label{fire:V41-current-boundary}
V41 proves the exact private-coordinate factorization of the complete
global fourth cumulant, closes every low-output tree whose internal
rows are thermal or private-dominant, and closes the private-mass part
of every corresponding high-output payer.  It does not close the
shared-dominant hot-row residual or the nonprivate high-output payer.
Therefore hypothesis \textup{(iv)} of the V36 compiler, global
quartic MTA, all-order MTA, WUFC/CPM, continuum Yang--Mills existence,
and a positive mass gap remain open.
\end{firewall}

\section{Append-only record for compact V41}
\label{sec:V41-record}

V41 was appended to the complete compact V40 file.  The exact V40
parent had \(16055\) newline-delimited source records, \(526431\)
bytes, and SHA--256 digest
\[
\texttt{5cd708f12a9617255105a60c2bc48152dc794c82bbb2f85c97bc2c0ab334dcc1}.
\]
No compact V40 mathematical content was changed.  The sole final
\verb|\end{document}| is restored below.

% =============================================================================
% END APPENDED COMPACT VERSION V41 MATERIAL
% =============================================================================

% The compact V41 document terminator is intentionally disabled here:
% \end{document}

% =============================================================================
% BEGIN APPENDED COMPACT VERSION V42 MATERIAL
% =============================================================================

\chapter{V42: The global partition-state carrier and common shared
stock}
\label{ch:V42}

\section{Why the global carrier needs a finite-state formulation}
\label{sec:V42-purpose}

V34 already constructed one balanced/unbalanced stock containing
incidences two, three, and four.  Its later global interpretation was
withdrawn in V36 because the proof there had first selected a unique
junction.  The stock inequalities themselves remain valid, but their
use in the complete global cumulant requires an algebraic carrier
which does not expand an independent partition choice at every
coordinate.

This chapter supplies that carrier.  A fourth cumulant has only the
fifteen global states in the partition lattice \(\Pi_4\).  Coordinate
moments act diagonally on those same fifteen states.  Thus an arbitrary
number of shared coordinates does not create \(3^N\), \(15^N\), or
any other support-dependent family of histories.

\section{The exact fifteen-state representation}
\label{sec:V42-fifteen-state}

Let \(\Pi_4\) be the set of partitions of \([4]\), and define
\[
 \mu(\pi):=(|\pi|-1)!(-1)^{|\pi|-1}.
\]
After the private factorization of V41, let \(E_{\rm sh}\) be the
finite set of all remaining coordinate letters.  At \(e\in E_{\rm sh}\)
let \(I_e\subseteq[4]\), \(2\le|I_e|\le4\), be its active rows.  For
\(A\subseteq[4]\), denote its local block moment by
\[
 m_e(A):=
 \mathbb E_e\!\left[
   \prod_{i\in A\cap I_e}X_{i,e}\right],
 \qquad
 m_e(A)=1\quad\hbox{if }A\cap I_e=\varnothing .
\]
Products are read in the fixed physical intertwiner block.  For
\(\pi\in\Pi_4\), put
\begin{equation}
 T_e(\pi):=\bigotimes_{A\in\pi}m_e(A).
\label{eq:V42-local-state}
\end{equation}
Factors belonging to distinct coordinates act on distinct tensor
factors.

Let \(\mathscr V\) be the space of operator-valued functions on
\(\Pi_4\).  Define the diagonal coordinate map and cumulant functional
by
\begin{align}
 (\mathsf T_ev)(\pi)&:=T_e(\pi)v(\pi),
\label{eq:V42-diagonal-map}\\
 \Lambda(v)&:=\sum_{\pi\in\Pi_4}\mu(\pi)v(\pi).
\label{eq:V42-cumulant-functional}
\end{align}
Write \(\boldsymbol1\) for the constant state.

\begin{theorem}[Exact global partition-state carrier]
\label{thm:V42-exact-state-carrier}
In every fixed physical output block,
\begin{equation}
\boxed{
 K_\nu^{\rm conn}
 =
 q_{\rm priv}\,
 \Lambda\!\left(
   \prod_{e\in E_{\rm sh}}\mathsf T_e\boldsymbol1
 \right).}
\label{eq:V42-exact-state-carrier}
\end{equation}
The product is independent of the coordinate ordering.  It contains
exactly fifteen global partition states, irrespective of
\(|E_{\rm sh}|\).
\end{theorem}

\begin{proof}
For a fixed \(\pi\), independence of the coordinate measures gives
\[
 \prod_{A\in\pi}
 \mathbb E\!\left[\prod_{i\in A}F_i^{\rm sh}\right]
 =
 \bigotimes_{e\in E_{\rm sh}}T_e(\pi).
\]
Insert this identity into the moment--partition formula used in
\Cref{lem:V41-exact-private-factor}, and then restore its common
factor \(q_{\rm priv}\).  This is exactly
\eqref{eq:V42-exact-state-carrier}.  The maps are diagonal in the same
global variable \(\pi\), and coordinate tensor factors commute, so the
ordering is immaterial.
\end{proof}

\begin{corollary}[Support-independent state norm]
\label{cor:V42-state-norm}
For every \(v\in\mathscr V\),
\begin{equation}
 \|\Lambda(v)\|
 \le26\max_{\pi\in\Pi_4}\|v(\pi)\|.
\label{eq:V42-state-norm}
\end{equation}
Consequently passing from a statewise bound to the complete cumulant
costs a fixed constant, not an exponential in the support size.
\end{corollary}

\begin{proof}
The numbers of partitions with \(1,2,3,4\) blocks are
\(1,7,6,1\).  Therefore
\[
 \sum_{\pi\in\Pi_4}|\mu(\pi)|
 =1+7+2\cdot6+6=26.
\]
Apply the triangle inequality once, after the coordinate product has
been formed.
\end{proof}

\section{A one-use response identity}
\label{sec:V42-one-use}

Set
\[
 r_e:=\bigotimes_{i\in I_e}m_e(\{i\}),
 \qquad
 \mathsf R_e:=r_e\,\mathsf I.
\]
Thus \(\mathsf R_e\boldsymbol1\) is constant in the partition state.
Let \(\mathsf D_e:=\mathsf T_e-\mathsf R_e\).  Since
\(\sum_{\pi}\mu(\pi)=0\), the fully row-separated reference has zero
fourth cumulant.

\begin{lemma}[First-defect telescope]
\label{lem:V42-first-defect}
For any ordering \(e_1,\dots,e_N\) of \(E_{\rm sh}\),
\begin{equation}
\boxed{
 K_\nu^{\rm conn}
 =
 q_{\rm priv}\sum_{k=1}^N
 \Lambda\!\left[
  \left(\prod_{\ell<k}\mathsf R_{e_\ell}\right)
  \mathsf D_{e_k}
  \left(\prod_{\ell>k}\mathsf T_{e_\ell}\right)
  \boldsymbol1
 \right].}
\label{eq:V42-first-defect}
\end{equation}
Every coordinate appears exactly once in every summand.
\end{lemma}

\begin{proof}
The commuting-product identity
\[
 \prod_{k=1}^N\mathsf T_{e_k}
 -
 \prod_{k=1}^N\mathsf R_{e_k}
 =
 \sum_{k=1}^N
 \left(\prod_{\ell<k}\mathsf R_{e_\ell}\right)
 (\mathsf T_{e_k}-\mathsf R_{e_k})
 \left(\prod_{\ell>k}\mathsf T_{e_\ell}\right)
\]
is a telescoping equality.  Apply \(\Lambda\) to its action on
\(\boldsymbol1\).  The second full product is constant in \(\pi\), so
its image under \(\Lambda\) vanishes.  Use
\Cref{thm:V42-exact-state-carrier}.
\end{proof}

\begin{assessment}[Resolution of the V36 regression mechanism]
\label{ass:V42-regression-resolved}
The forbidden tensor
\(\bigotimes_e(M_e-K_{4,e})\) creates an independent binary choice at
each coordinate and can have norm \(3^N\).  Formula
\eqref{eq:V42-first-defect} creates no such choice: the only state is
the single global partition \(\pi\), and a coordinate is used once in
the ordered telescope.  This resolves the algebraic, but not yet the
analytic, part of the global shared-stock problem.
\end{assessment}

\section{One common incidence-\texorpdfstring{\(\{2,3,4\}\)}{{2,3,4}}
stock}
\label{sec:V42-common-stock}

At a shared coordinate \(e\), let
\[
 a_{i,e}:=1+C_2(R_{i,e}),
 \qquad
 w_e:=\max_{i\in I_e}a_{i,e}.
\]
Use exactly the representation-theoretic balanced/unbalanced
classification of V20, V31, and V34.  If \(e\) is balanced, assign it
once to the legal pair \(u_e,v_e\) and set
\(E_e=a_{u_e,e}a_{v_e,e}\).  If it is unbalanced, retain \(w_e\).
Globally, without first selecting an incidence or a junction, put
\begin{equation}
 S_B:=\sum_{e\in B}w_e,
 \qquad
 H_B:=\sum_{e\in B}E_e,
 \qquad
 W:=\sum_{e\in U}w_e.
\label{eq:V42-global-stocks}
\end{equation}

\begin{lemma}[Global common-stock inequalities]
\label{lem:V42-global-common-stock}
The definitions in \eqref{eq:V42-global-stocks} obey
\begin{align}
 S_B^2&\le36N_BH_B,
\label{eq:V42-global-SH}\\
 A_0\sum_{e\in B}
 \frac{E_e}{\Xi_{u_e}\Xi_{v_e}}
 &\ge\frac{A_0H_B}{\Xi_{\max}^2}.
\label{eq:V42-global-routing}
\end{align}
If row \(i\) is shared-dominant, then
\begin{equation}
 S_B+W\ge h_i>\frac12\Xi_i.
\label{eq:V42-row-stock}
\end{equation}
Hence either \(S_B>\Xi_i/4\) or \(W>\Xi_i/4\).
\end{lemma}

\begin{proof}
For balanced coordinates of incidences two, three, and four, the
pointwise inequalities are respectively
\[
 w_e^2\le4E_e,\qquad
 w_e^2\le16E_e,\qquad
 w_e^2\le36E_e.
\]
Cauchy--Schwarz proves \eqref{eq:V42-global-SH}.  Every balanced
coordinate is assigned to one legal active pair, so
\(\Xi_{u_e}\Xi_{v_e}\le\Xi_{\max}^2\), proving
\eqref{eq:V42-global-routing}.

For a fixed row,
\[
 h_i=\sum_{e:\,i\in I_e}a_{i,e}
 \le\sum_{e:\,i\in I_e}w_e
 \le S_B+W.
\]
Combine this with \(h_i>\Xi_i/2\).
\end{proof}

\begin{remark}[No incidence pigeonhole]
\label{rem:V42-no-incidence-pigeonhole}
The alternative in \eqref{eq:V42-row-stock} is formed after summing all
incidences.  There is no loss by the seven possible partner sets and
no overlapping collection of incidence-specific banks.
\end{remark}

\section{State-uniform decay of the unbalanced stock}
\label{sec:V42-unbalanced-state}

\begin{lemma}[Unbalanced retention in every global partition state]
\label{lem:V42-statewise-unbalanced}
There are universal \(c,C>0\) such that, for every
\(\pi\in\Pi_4\),
\begin{equation}
 \left\|
 \bigotimes_{e\in U}T_e(\pi)
 \right\|
 \le C e^{-cs_\alpha W}.
\label{eq:V42-statewise-unbalanced}
\end{equation}
Consequently, for \(k=1,2,3\),
\begin{equation}
\boxed{
 (s_\alpha W)^k
 \max_{\pi\in\Pi_4}
 \left\|\bigotimes_{e\in U}T_e(\pi)\right\|
 \le C_k.}
\label{eq:V42-statewise-unbalanced-moments}
\end{equation}
\end{lemma}

\begin{proof}
Fix \(e\in U\), order its active spins so that \(j_1\) is dominant,
and let \(A\in\pi\) be the unique partition block containing row \(1\).
Only a subset of the other active spins can occur in the same local
moment.  The unbalanced alternative gives
\[
 J\ge j_1-\sum_{r\ne1}j_r>\frac12j_1
\]
for every output spin \(J\) in that block.  The heat multiplier on
this output is therefore at most
\(Ce^{-cs_\alpha w_e}\); all other block moments are contractions.
This argument is independent of \(\pi\).  Tensoring over \(e\in U\)
proves \eqref{eq:V42-statewise-unbalanced}.  The scalar supremum
\(\sup_{x\ge0}x^ke^{-cx}<\infty\) proves
\eqref{eq:V42-statewise-unbalanced-moments}.
\end{proof}

\begin{corollary}[No partition-count loss in the unbalanced tail]
\label{cor:V42-unbalanced-tail}
If a complete carrier term retains the unbalanced coordinate moments
statewise, then
\[
 (s_\alpha W)^k
 \left\|
 \Lambda\!\left[
 \left(\bigotimes_{e\in U}T_e\right)v
 \right]\right\|
 \le
 26C_k\max_{\pi}\|v(\pi)\|,
 \qquad k=1,2,3.
\]
\end{corollary}

\begin{proof}
Apply \Cref{cor:V42-state-norm} and
\eqref{eq:V42-statewise-unbalanced-moments}.
\end{proof}

\section{Tree marks and the reserve dichotomy}
\label{sec:V42-reserve}

Each monomial in
\[
 L_\tau=\prod_{ij\in E(\tau)}H_{ij}
\]
chooses one coordinate for each of the three tree edges.  Let
\(Q\subseteq E_{\rm sh}\) be the set of distinct chosen coordinates.
Then \(|Q|\le3\).  Define the unbalanced reserve
\[
 W_Q:=\sum_{e\in U\setminus Q}w_e.
\]

\begin{lemma}[Reserve or one heavy marked coordinate]
\label{lem:V42-reserve-heavy}
For every tree monomial, either
\begin{equation}
 W_Q\ge\frac12W,
\label{eq:V42-large-reserve}
\end{equation}
or some marked unbalanced coordinate \(e\in Q\cap U\) satisfies
\begin{equation}
 w_e>\frac16W.
\label{eq:V42-heavy-mark}
\end{equation}
\end{lemma}

\begin{proof}
If \eqref{eq:V42-large-reserve} fails, the at most three coordinates
of \(Q\cap U\) carry total mass \(W-W_Q>W/2\).  One of them carries
more than \(W/6\).
\end{proof}

\begin{proposition}[Retained-tail refinement of the frozen tree
expansion]
\label{prop:V42-retained-tail}
The V135 connected tree expansion can be organized so that, in the
monomial indexed by \(Q\), every coordinate in
\(E_{\rm sh}\setminus Q\) remains inside its exact state multiplier
\(T_e(\pi)\).  Passing to a statewise norm costs at most the fixed
factor \(26\).
\end{proposition}

\begin{proof}
Expand each edge sum \(H_{ij}\) only after the three row-connecting
differences in the frozen proof have been chosen.  A difference
selecting the edge \(ij\) acts on one coordinate supporting both
rows; thus three edges distinguish at most three coordinates.
Multilinearity leaves every other coordinate moment unchanged.
Before the final norm, collect the terms with the same global
moment-partition state \(\pi\).  There are precisely the fifteen
states in \eqref{eq:V42-exact-state-carrier}, not independent states
at the unmarked coordinates.  Applying
\eqref{eq:V42-state-norm} gives the asserted constant.  This is the
same frozen tree argument with contractions on unmarked coordinates
postponed, so its edge coefficients and \(L_\tau\) are unchanged.
\end{proof}

\begin{definition}[Unbalanced-controlled tree]
\label{def:V42-unbalanced-controlled}
A tree with a selected internal row \(i\) is
\((M,W)\)-unbalanced-controlled if \(W\ge\Xi_i/4\) and every other
internal excess factor is either thermal, private-dominant, or bounded
by \(4W\).
\end{definition}

\begin{theorem}[Low-output unbalanced reserve closure]
\label{thm:V42-unbalanced-reserve-closure}
Consider a low-output tree monomial whose marked set \(Q\) satisfies
\eqref{eq:V42-large-reserve}.  If the tree is
\((M,W)\)-unbalanced-controlled, then that monomial is bounded by its
global marked-tree weight with a constant depending only on \(M\).
\end{theorem}

\begin{proof}
Retain the tail in
\Cref{prop:V42-retained-tail}.  By
\Cref{lem:V42-statewise-unbalanced},
\[
 \max_\pi\left\|
 \bigotimes_{e\in U\setminus Q}T_e(\pi)
 \right\|
 \le\frac{C_k}{(s_\alpha W_Q)^k}
 \le\frac{C_k'}{(s_\alpha W)^k}
\]
for \(k=1,2\).

For a star, its center is the selected row and
\(B_\tau=\Xi_i^2\le16W^2\); use \(k=2\).  For a path,
\(B_\tau=\Xi_i\Xi_j\).  The selected factor obeys
\(\Xi_i\le4W\).  If the other factor is thermal, use \(k=1\) and
\(s_\alpha\Xi_j\le M\).  If it is private-dominant, use \(k=1\) and
the separate first private moment of
\eqref{eq:V41-row-private-moments}.  If \(\Xi_j\le4W\), use \(k=2\).
Thus in every case the retained private and unbalanced moments give
\[
 s_\alpha^2q_{\rm priv}B_\tau
 \max_\pi\left\|
 \bigotimes_{e\in U\setminus Q}T_e(\pi)
 \right\|
 \le C_M.
\]
Multiply by \(L_\tau/B_\tau\).  The low-output coefficient is the
V41 factor \(Cs_\alpha^2\), and the fixed state loss is absorbed into
\(C_M\).
\end{proof}

\begin{assessment}[What V42 has and has not closed]
\label{ass:V42-result}
V42 proves that the global shared carrier has only fifteen partition
states, forms the incidence-\(\{2,3,4\}\) stock before any bank
selection, and gives state-uniform unbalanced decay.  It closes the
low-output unbalanced branch whenever the three tree marks leave at
least half of the unbalanced stock in reserve and the other internal
tree factor has one of the listed payers.

Two concrete branches remain upstream of the V41 residual:
\begin{enumerate}
\item the heavy marked-coordinate alternative
      \eqref{eq:V42-heavy-mark}, where the coordinate that should
      provide heat decay has already been used by a tree difference;
\item the balanced-stock alternative, where protected invariant
      sectors require a capacity estimate on the fifteen-state tail,
      not merely the selected-junction estimate of V34.
\end{enumerate}
\end{assessment}

\begin{openproblem}[V43 mandate: marked-heavy absorption and balanced
state capacity]
\label{op:V43-mandate}
\begin{enumerate}
\item At a heavy marked unbalanced coordinate, retain the output-spin
      lower bound inside the marked difference instead of replacing
      that difference immediately by \(H_{ij}\).
\item Prove the resulting marked-and-decaying local estimate, with no
      second use of the coordinate.
\item For balanced coordinates, lift the V34 full invariant-projector
      normal form to each of the fifteen global partition states.
\item Tensor those statewise cut capacities before applying the fixed
      factor \(26\).
\item Revisit the high-output \(Z_\nu\)-payer only after these two
      common-stock branches are available.
\end{enumerate}
\end{openproblem}

\begin{firewall}[Claim boundary after compact V42]
\label{fire:V42-current-boundary}
V42 does not prove the complete all-weak estimate.  In particular,
heavy marked unbalanced coordinates, balanced protected state tails,
and the nonprivate high-output payer remain open.  Hence global
quartic MTA, all-order MTA, WUFC/CPM, continuum Yang--Mills existence,
and a positive mass gap remain open.
\end{firewall}

\section{Append-only record for compact V42}
\label{sec:V42-record}

V42 was appended to the complete compact V41 file.  The exact V41
parent had \(16514\) newline-delimited source records, \(542399\)
bytes, and SHA--256 digest
\[
\texttt{df28a0fc9d662861904ce728d85a6332a12027beed3b1602cd644bad5aa20c44}.
\]
No compact V41 mathematical content was changed.  The sole final
\verb|\end{document}| is restored below.

% =============================================================================
% END APPENDED COMPACT VERSION V42 MATERIAL
% =============================================================================

% The compact V42 document terminator is intentionally disabled here:
% \end{document}

% =============================================================================
% BEGIN APPENDED COMPACT VERSION V43 MATERIAL
% =============================================================================

\chapter{V43: Heavy unbalanced marks and the cut-connected balanced
quotient}
\label{ch:V43}

\section{The marked unbalanced coordinate retains its heat decay}
\label{sec:V43-heavy-mark}

The second branch of
\Cref{lem:V42-reserve-heavy} is not a loss of the unbalanced moment.
The coordinate has been marked by a row-connecting difference, but
every term of that difference still contains the dominant input row.
Its output-spin lower bound can therefore be retained together with
the tree mark.

\begin{lemma}[Marked-and-decaying unbalanced coordinate]
\label{lem:V43-marked-decay}
Let \(e\) be unbalanced, with dominant active spin \(j_1\) and mass
\(w_e\).  Let \(\mathsf A_e\) be any local coordinate operator arising
after one or more of the three row-connecting differences in the
frozen V135 tree extraction have marked \(e\).  After its displayed
Casimir edge factors have been removed, its residual state multiplier
obeys
\begin{equation}
\boxed{
 \max_{\pi\in\Pi_4}\|\mathsf A_e(\pi)\|
 \le C e^{-cs_\alpha w_e}.}
\label{eq:V43-marked-decay}
\end{equation}
The constants are independent of the incidence and of how many of the
three tree edges use \(e\).
\end{lemma}

\begin{proof}
Each row-connecting difference is a subtraction between two local
moment products.  Expanding at this single coordinate produces a
fixed linear combination, with coefficients depending only on the
four-row partition lattice, of terms of the form
\[
 \bigotimes_{A\in\pi}
 \mathbb E_e\!\left[
   \prod_{i\in A\cap I_e}X_{i,e}\right].
\]
In every such term the unique block containing the dominant row
contains only a subset of the remaining active rows.  Therefore every
output spin \(J\) in that block satisfies
\[
 J\ge j_1-\sum_{r\ne1}j_r>\frac12j_1.
\]
The corresponding heat multiplier is bounded by
\(Ce^{-cs_\alpha w_e}\); all other local block moments are
contractions.  The number of local partition terms and connecting
differences is bounded in terms of four rows only.  Summing them proves
\eqref{eq:V43-marked-decay}.

The Casimir factors already displayed in the selected tree monomial
are not estimated again.  The exponential in
\eqref{eq:V43-marked-decay} is the heat multiplier remaining in the
same exact local coefficient, so no coordinate is duplicated.
\end{proof}

\begin{corollary}[Heavy mark replaces the missing reserve]
\label{cor:V43-heavy-replaces-reserve}
If a tree monomial has a marked unbalanced coordinate satisfying
\(w_e>W/6\), then its marked local multiplier together with all
unmarked unbalanced coordinates obeys
\begin{equation}
 (s_\alpha W)^k
 \max_{\pi\in\Pi_4}
 \left\|
 \mathsf A_e(\pi)
 \bigotimes_{f\in U\setminus Q}T_f(\pi)
 \right\|
 \le C_k,
 \qquad k=1,2,3.
\label{eq:V43-heavy-full-stock}
\end{equation}
\end{corollary}

\begin{proof}
By \Cref{lem:V43-marked-decay} and the contraction bound on all
unmarked coordinate moments, the norm is at most
\(Ce^{-cs_\alpha w_e}\le Ce^{-(c/6)s_\alpha W}\).  Multiply by
\((s_\alpha W)^k\) and optimize the scalar exponential.
\end{proof}

\begin{theorem}[Complete low-output unbalanced-stock closure]
\label{thm:V43-full-unbalanced-low}
Every low-output tree monomial which is
\((M,W)\)-unbalanced-controlled in the sense of
\Cref{def:V42-unbalanced-controlled} obeys its global marked-tree
weight, with no reserve hypothesis.
\end{theorem}

\begin{proof}
Apply \Cref{lem:V42-reserve-heavy}.  In the reserve alternative use
\Cref{thm:V42-unbalanced-reserve-closure}.  In the heavy alternative
replace the reserve moment in that proof by
\eqref{eq:V43-heavy-full-stock}.  The same star/path case split pays
\(s_\alpha^2B_\tau\).  The state sum costs only the fixed factor \(26\)
from \Cref{cor:V42-state-norm}.
\end{proof}

\begin{theorem}[Private high-output payer with unbalanced shared
control]
\label{thm:V43-unbalanced-private-high}
Let a high-output tree be \((M,W)\)-unbalanced-controlled.  The part
of its coefficient majorant carrying the private output mass
\(P=\sum_rp_r\) obeys its global marked-tree weight.
\end{theorem}

\begin{proof}
Expand \(P\) into its four labelled payers.  Use the reserve/heavy
dichotomy exactly as in
\Cref{thm:V43-full-unbalanced-low}.  One or two powers of \(W\) pay
the shared-controlled internal excess factors.  Thermal internal
factors cost \(M\).  Private-dominant internal factors and the labelled
output payer are paid from the disjoint row banks by
\eqref{eq:V41-row-private-moments}.

The total polynomial degree is three:
two degrees come from \(\sum_i(d_i-1)=2\), and one from \(p_r\).
No single private bank receives exponent greater than three, and
\eqref{eq:V43-heavy-full-stock} is available through \(k=3\).
Thus
\[
 s_\alpha^3Pq_{\rm priv}B_\tau
 \times(\hbox{retained unbalanced state multiplier})
 \le C_M.
\]
Multiply by \(L_\tau/B_\tau\), as in
\Cref{thm:V41-high-private-output}.
\end{proof}

\section{A failed balanced shortcut}
\label{sec:V43-balanced-regression}

The V42 mandate proposed a capacity estimate on each of the fifteen
partition states.  That statement is false, even though the complete
connected carrier can still satisfy the desired cut estimate.

\begin{proposition}[Statewise protected cut capacity is false]
\label{prop:V43-statewise-false}
Consider the cut \(\{1,2\}\mid\{3,4\}\) and the partition
\(\pi=\{\{1,2\},\{3,4\}\}\).  There are balanced incidence-four input
blocks for which the protected part of \(T_e(\pi)\) has product-state
capacity one across this cut.  Hence no estimate
\[
 \kappa_\otimes(T_e(\pi))\le d^{-1}
\]
with \(d\to\infty\) can hold uniformly state by state.
\end{proposition}

\begin{proof}
Take rows \(1,2\) in mutually dual spin-\(j\) representations and rows
\(3,4\) in mutually dual spin-\(k\) representations.  In the indicated
partition state, choose the singlet output in each two-row block.  The
protected operator is
\[
 P_{12}^{\rm inv}\otimes P_{34}^{\rm inv}.
\]
It is a tensor product across
\(\{1,2\}\mid\{3,4\}\), has norm one, and is attained on the product
of its normalized left and right singlets.  Its product-state capacity
across the cut is therefore one, independently of \(j,k\).
\end{proof}

\begin{assessment}[Why this does not contradict V34]
\label{ass:V43-no-V34-contradiction}
V34 bounds the full invariant projector when it joins the two sides of
the tested cut.  The operator in
\Cref{prop:V43-statewise-false} consists of two separately invariant
projectors and is disconnected across that cut.  It cancels only
after the product--cumulant identity is assembled.  Thus V34 remains
valid, while the proposed statewise lift in the V42 mandate must be
replaced.
\end{assessment}

\section{The cut-connected quotient before norms}
\label{sec:V43-cut-connected}

Let \(C\mid D\) be a nontrivial row cut and
\(\rho=\{C,D\}\).  V37 established the exact identity
\begin{equation}
 \operatorname{Cov}(F_C,F_D)
 =
 \kappa_4(F_1,F_2,F_3,F_4)
 +
 \sum_{\substack{\pi\in\Pi_4,\ \pi\ne\boldsymbol1\\
                  \pi\vee\rho=\boldsymbol1}}
 \prod_{A\in\pi}\kappa(F_i:i\in A).
\label{eq:V43-product-cumulant}
\end{equation}
It is essential that \eqref{eq:V43-product-cumulant} is formed before
capacity.  Every term on its right other than \(\kappa_4\) is a product
of lower connected carriers whose row blocks collectively cross the
cut.

\begin{definition}[Cut-connected protected tail]
\label{def:V43-cut-tail}
For a balanced coordinate bank \(B'\), its
\(C\mid D\) cut-connected protected tail is the complete operator
obtained by:
\begin{enumerate}
\item forming \eqref{eq:V43-product-cumulant} with all coordinate
      moments on \(B'\) retained;
\item subtracting its displayed lower connected correction products;
\item restricting only then to zero heat output on \(B'\).
\end{enumerate}
No individual partition state is normed in this definition.
\end{definition}

\begin{lemma}[Disconnected protected products vanish before
capacity]
\label{lem:V43-disconnected-cancel}
Every protected coordinate contribution which factors into a left
operator on \(C\) and a right operator on \(D\), including the
two-singlet operator of \Cref{prop:V43-statewise-false}, cancels from
the cut-connected protected tail.
\end{lemma}

\begin{proof}
Replace the coordinate variable on the two sides by independent
copies.  Then \(F_C\) and \(F_D\) are independent, so the left side of
\eqref{eq:V43-product-cumulant} is zero.  The product--cumulant
identity is polynomial in the coordinate moments, hence the sum of all
right-side terms having this left/right factorization is identically
zero.  The cancellation is operator-valued and precedes the
restriction to protected outputs.  Equivalently, it is the Möbius
identity on \(\Pi_4\) after restriction to the two sides of \(\rho\).
\end{proof}

\begin{lemma}[Every surviving protected block crosses the cut]
\label{lem:V43-surviving-crosses}
After the cancellation in
\Cref{lem:V43-disconnected-cancel}, each surviving protected summand
contains at least one invariant projector whose active row block meets
both \(C\) and \(D\).
\end{lemma}

\begin{proof}
If every protected invariant block were supported wholly on one side,
their tensor product would factor as a left operator times a right
operator.  Such a term vanishes by
\Cref{lem:V43-disconnected-cancel}.  Therefore a surviving term has a
crossing invariant block.
\end{proof}

\begin{proposition}[Single-coordinate crossing capacity]
\label{prop:V43-single-crossing-capacity}
Let a surviving protected summand contain a crossing invariant block
at a balanced coordinate \(e\).  If \(D_{e,C}\) is the dimension of
the smaller nontrivial representation leg exposed by the induced cut,
then its product-state capacity across \(C\mid D\) is at most
\begin{equation}
 \frac1{D_{e,C}}.
\label{eq:V43-single-crossing-capacity}
\end{equation}
All multiplicity spaces are included in this estimate.
\end{proposition}

\begin{proof}
Regroup the active tensor factors of that invariant block according to
the two sides of the cut.  The V34 full-projector normal form is
\[
 P_{\rm inv}\simeq
 |\Omega\rangle\langle\Omega|\otimes I_{\mathcal M}\oplus0,
\]
where \(\Omega\) is maximally entangled on the exposed irreducible
legs.  Its overlap with a product vector across the cut is at most
\(D_{e,C}^{-1}\).  The identity on the complete multiplicity space
does not create a channel count.  All other protected blocks and
gapped moments in the summand are contractions, so one-sided tensor
submultiplicativity proves \eqref{eq:V43-single-crossing-capacity}.
\end{proof}

\section{The precise balanced analytic gate}
\label{sec:V43-balanced-gate}

For a cut \(C\mid D\), let \(B_{C|D}\) be the common balanced
coordinates whose assigned legal pair \(u_e,v_e\) is separated by the
cut.  Put
\[
 S_{C|D}:=\sum_{e\in B_{C|D}}w_e,
 \qquad
 H_{C|D}:=\sum_{e\in B_{C|D}}E_e.
\]
For a tree monomial with marked set \(Q\), use the same notation with
\(B_{C|D}\setminus Q\).

\begin{conjecture}[Cut-connected balanced-tail estimate]
\label{conj:V43-balanced-tail}
The cut-connected protected tail of
\Cref{def:V43-cut-tail}, together with its gapped complement, should
satisfy
\begin{equation}
\boxed{
 s_\alpha S_{C|D,Q}^{\,2}
 \kappa_\otimes(\mathcal B_{C|D,Q})
 \le C H_{C|D,Q}}
\label{eq:V43-balanced-tail}
\end{equation}
uniformly in \(Q\), provided \(S_{C|D,Q}\) retains a fixed fraction of
\(S_{C|D}\).  The capacity is taken only after the cut-connected
identity has been assembled.
\end{conjecture}

\begin{proposition}[Finite reduction of the balanced V41 residual]
\label{prop:V43-balanced-reduction}
Assume \eqref{eq:V43-balanced-tail} for the seven nontrivial row cuts,
and assume the marked-heavy analogue in which one removed coordinate
retains its crossing capacity as in
\Cref{prop:V43-single-crossing-capacity}.  Then every low-output
balanced-stock tree whose other internal excess factor is thermal,
private-dominant, or controlled by the same cut stock obeys its global
marked-tree weight.  The same is true for its private high-output
part.
\end{proposition}

\begin{proof}
Choose the cut obtained by deleting the tree edge whose legal marked
pair is supplied by the common routing sum
\eqref{eq:V42-global-routing}.  The cut stock pays the shared-dominant
internal mass by the same star/path algebra as
\Cref{thm:V43-full-unbalanced-low}, with
\eqref{eq:V43-balanced-tail} replacing the unbalanced exponential
moment.  If the reserve is small, one of the at most three marked
coordinates is heavy; use the assumed marked-heavy version.
Thermal and private factors are paid as in V41.  For private high
output, add the single labelled private payer and use moments only
through total degree three.  There are seven cuts and sixteen labelled
trees, so all losses are fixed.
\end{proof}

\begin{openproblem}[V44 mandate: prove the balanced cut-tail gate]
\label{op:V44-mandate}
\begin{enumerate}
\item Expand the cut-connected identity
      \eqref{eq:V43-product-cumulant} into its finite crossing-block
      types, not into coordinatewise protected/nonprotected choices.
\item Tensor the crossing projector estimate
      \eqref{eq:V43-single-crossing-capacity} over a common reserve.
\item Combine that capacity decay with
      \(S_{C|D}^2\le36N_{C|D}H_{C|D}\).
\item Treat gapped output by the same scalar heat tail, and treat a
      heavy marked protected coordinate without using it twice.
\item Either prove \eqref{eq:V43-balanced-tail} or exhibit a finite
      crossing-block type for which the proposed normalization fails.
\end{enumerate}
\end{openproblem}

\begin{firewall}[Claim boundary after compact V43]
\label{fire:V43-current-boundary}
V43 closes the reserve and heavy-mark alternatives for the
unbalanced-stock low-output branch and its private high-output payer.
It disproves the false statewise balanced-capacity shortcut and
replaces it by the cut-connected gate
\eqref{eq:V43-balanced-tail}.  That gate, the nonprivate high-output
mass, and some mixed payer configurations remain open.  Therefore
global quartic MTA, all-order MTA, WUFC/CPM, continuum Yang--Mills
existence, and a positive mass gap remain open.
\end{firewall}

\section{Append-only record for compact V43}
\label{sec:V43-record}

V43 was appended to the complete compact V42 file.  The exact V42
parent had \(17005\) newline-delimited source records, \(558513\)
bytes, and SHA--256 digest
\[
\texttt{a376e6c33d593aeb98389113527109492a7b9ad0eab360b77d5a3483b26671d6}.
\]
No compact V42 mathematical content was changed.  The sole final
\verb|\end{document}| is restored below.

% =============================================================================
% END APPENDED COMPACT VERSION V43 MATERIAL
% =============================================================================

% The compact V43 document terminator is intentionally disabled here:
% \end{document}

% =============================================================================
% BEGIN APPENDED COMPACT VERSION V44 MATERIAL
% =============================================================================

\chapter{V44: Balanced cut-tail capacity and finite mass transport}
\label{ch:V44}

\section{Local alternatives after cut-connected cancellation}
\label{sec:V44-local-cut}

Fix a nontrivial cut \(C\mid D\).  Consider one of the finitely many
crossing-block types produced by the exact product--cumulant identity
\eqref{eq:V43-product-cumulant}.  At a balanced coordinate \(e\) whose
chosen legal pair is separated by the cut, retain all local physical
outputs until the cut-connected cancellations of
\Cref{lem:V43-disconnected-cancel} have been made.

\begin{lemma}[Gapped-or-crossing-protected alternative]
\label{lem:V44-local-alternative}
Every surviving local output at such a coordinate has one of the
following two properties:
\begin{enumerate}
\item some local block output is nontrivial and its heat multiplier
      has norm at most \(r_\alpha:=1-c_*s_\alpha<1\);
\item a protected invariant block crosses \(C\mid D\), and its complete
      projector has product-state capacity at most
      \(D_{e,C}^{-1}\).
\end{enumerate}
If the legal pair is split between distinct protected blocks, either
one of those blocks crosses the cut or the full local contribution is
cut-disconnected and has already cancelled.
\end{lemma}

\begin{proof}
If a local block output is nontrivial, the established balanced
spectral separation gives its heat multiplier the uniform bound
\(1-c_*s_\alpha\).  No comparison with \(w_e\) is asserted: balanced
large spins may fuse to a small nonzero output.

Suppose instead that every local block output is trivial.  Each
nonzero block is then an invariant tensor.  If every such block lies
on one side of the cut, their product factors into a \(C\)-operator
and a \(D\)-operator and vanishes by
\Cref{lem:V43-disconnected-cancel}.  Otherwise one invariant block
crosses the cut, and
\Cref{prop:V43-single-crossing-capacity} gives
\(D_{e,C}^{-1}\).  The last assertion is the same dichotomy applied to
the two blocks containing the endpoints of the legal pair.
\end{proof}

\begin{lemma}[Uniform protected contraction]
\label{lem:V44-uniform-protected}
There is a universal \(\rho<1\) such that the protected alternative in
\Cref{lem:V44-local-alternative} has product-state capacity at most
\(\rho\).
\end{lemma}

\begin{proof}
The crossing invariant block exposes a nontrivial \(SU(2)\)
representation across the cut.  Its dimension is at least two, so
\Cref{prop:V43-single-crossing-capacity} gives a bound at most
\(1/2\).  Enlarging this fixed bound if required by the one-sided
tensor convention gives a universal \(\rho<1\).
\end{proof}

\section{The balanced cut-tail theorem}
\label{sec:V44-cut-tail}

For \(B_{C|D}\), \(S_{C|D}\), and \(H_{C|D}\), use the notation of
V43.  If a tree monomial marks \(Q\), analytic information at a marked
coordinate is retained as in V43: its displayed edge Casimir factor is
removed, but its residual heat multiplier or invariant-projector
capacity is not replaced by one.

\begin{theorem}[Cut-connected balanced-tail estimate]
\label{thm:V44-balanced-tail}
The conjectured estimate
\eqref{eq:V43-balanced-tail} holds:
\begin{equation}
\boxed{
 s_\alpha S_{C|D}^{\,2}
 \kappa_\otimes(\mathcal B_{C|D})
 \le C H_{C|D}.}
\label{eq:V44-balanced-tail}
\end{equation}
The same estimate holds for a tree-resolved tail with at most three
marked coordinates, provided their residual analytic multipliers are
retained.  No constant depends on the support size, spins,
multiplicities, or the chosen global partition type.
\end{theorem}

\begin{proof}
First fix one crossing-block type from
\eqref{eq:V43-product-cumulant}.  There are only finitely many such
types.  At every coordinate of \(B_{C|D}\), apply
\Cref{lem:V44-local-alternative}.  The gapped output has the same
uniform spectral factor \(r_\alpha=1-c_*s_\alpha\) as the gapped
summand in V34.  The protected output has the uniform capacity
\(\rho<1\) from \Cref{lem:V44-uniform-protected}.

Now repeat the scalar power-capacity summation in the proof of
\Cref{thm:V34-universal-stocks}, without changing its order:
orthogonal gapped/protected output sectors are summed first,
coordinate tensor products second, and product-state capacity last.
The only scalar inputs of that proof are
\[
 w_e^2\le36E_e,
 \qquad
 S_{C|D}^2\le36N_{C|D}H_{C|D},
\]
the gapped heat tail, and the protected full-projector capacity.  All
four inputs have just been verified for the cut-connected type.
Consequently its complete gapped-plus-protected tail obeys
\[
 s_\alpha S_{C|D}^2\kappa_\otimes
 \le C H_{C|D}.
\]

Sum only after this estimate over the finite crossing-block types in
\eqref{eq:V43-product-cumulant}.  The lower connected correction
products use the one-payer rule of
\Cref{rem:V37-one-payer}; all other factors are contractions.  Thus
the constant is multiplied by a fixed number no larger than the
partition-state constant \(26\).

For a marked coordinate, the frozen tree extraction leaves a fixed
linear combination of the same local physical outputs.  If it is
gapped, retain its uniform factor \(r_\alpha\).  If it is protected,
retain the crossing capacity from
\Cref{prop:V43-single-crossing-capacity}.  Therefore the preceding
power-capacity sum includes marked and unmarked coordinates in one
stock; no reserve split or second use of a coordinate is required.
\end{proof}

\begin{corollary}[The V43 analytic gate is closed]
\label{cor:V44-gate-closed}
The cut-connected balanced tail, rather than an individual partition
state, is the correct global extension of the V34 balanced compiler.
In particular the analytic assumption in
\Cref{prop:V43-balanced-reduction} is now a theorem.
\end{corollary}

\section{A row-rooted stock}
\label{sec:V44-row-rooted}

The global stock \eqref{eq:V42-global-stocks} detects a
shared-dominant row but does not say which legal edge carries that
row's mass.  The following refinement is made only after the common
balanced/unbalanced classification; it does not repartition
coordinates by incidence.

For a row \(i\) and a balanced coordinate \(e\ni i\), choose
\(v_i(e)\ne i\) with maximal active mass among the other rows.  Put
\begin{align}
 S_{ij}&:=
 \sum_{\substack{e\in B:\ i\in I_e\\v_i(e)=j}}a_{i,e},
 &
 H_{ij}^{(i)}&:=
 \sum_{\substack{e\in B:\ i\in I_e\\v_i(e)=j}}
 a_{i,e}a_{j,e},
\label{eq:V44-row-pair-stocks}\\
 S_i&:=\sum_{j\ne i}S_{ij},
 &
 U_i&:=\sum_{\substack{e\in U\\i\in I_e}}a_{i,e}.
\label{eq:V44-row-stocks}
\end{align}
Then
\begin{equation}
 h_i=S_i+U_i.
\label{eq:V44-row-stock-exact}
\end{equation}

\begin{lemma}[Row-rooted balanced power inequality]
\label{lem:V44-row-power}
For every \(i,j\),
\begin{equation}
\boxed{
 S_{ij}^2\le36N_{ij}H_{ij}^{(i)},}
\label{eq:V44-row-power}
\end{equation}
where \(N_{ij}\) is the number of coordinates in
\eqref{eq:V44-row-pair-stocks}.  Moreover
\begin{equation}
 H_{ij}^{(i)}\le H_{ij},
\label{eq:V44-row-H-domination}
\end{equation}
where \(H_{ij}\) is the full raw overlap of rows \(i,j\).
\end{lemma}

\begin{proof}
If \(a_{i,e}\) is not maximal, then
\(a_{j,e}\ge a_{i,e}\).  If it is maximal, the balanced
incidence-two/three/four alternatives imply
\(a_{j,e}\ge a_{i,e}/36\).  Thus
\[
 a_{i,e}^2\le36a_{i,e}a_{j,e}.
\]
Cauchy--Schwarz over the \(N_{ij}\) coordinates gives
\eqref{eq:V44-row-power}.  Every summand of
\(H_{ij}^{(i)}\) is one of the nonnegative coordinate summands in the
full overlap \(H_{ij}\), proving
\eqref{eq:V44-row-H-domination}.
\end{proof}

\begin{lemma}[Shared-hot row alternative]
\label{lem:V44-row-alternative}
If \(h_i>\Xi_i/2\), then either
\begin{equation}
 U_i>\frac14\Xi_i
\label{eq:V44-row-unbalanced}
\end{equation}
or
\begin{equation}
 S_i>\frac14\Xi_i.
\label{eq:V44-row-balanced}
\end{equation}
In the second case there is a partner \(j\ne i\) such that
\begin{equation}
 S_{ij}>\frac1{12}\Xi_i.
\label{eq:V44-partner-stock}
\end{equation}
\end{lemma}

\begin{proof}
Use \eqref{eq:V44-row-stock-exact}.  If neither
\eqref{eq:V44-row-unbalanced} nor
\eqref{eq:V44-row-balanced} held, then \(h_i\le\Xi_i/2\).
There are three possible partners, so
\eqref{eq:V44-row-balanced} implies
\eqref{eq:V44-partner-stock} for at least one of them.
\end{proof}

\section{All-weak mass transport}
\label{sec:V44-mass-transport}

Assume now the all-weak inequalities
\begin{equation}
 H_{ij}\le\delta\,\Xi_i\Xi_j
 \qquad(i\ne j).
\label{eq:V44-all-weak}
\end{equation}

\begin{theorem}[Balanced stock transports mass to one partner]
\label{thm:V44-mass-transport}
If \eqref{eq:V44-partner-stock} holds, then
\begin{equation}
\boxed{
 \Xi_i\le5184\,\delta\,N_{ij}\Xi_j.}
\label{eq:V44-mass-transport}
\end{equation}
Consequently, for every \(L\ge1\), either
\begin{equation}
 N_{ij}\ge L
\label{eq:V44-dense-bank}
\end{equation}
or
\begin{equation}
 \Xi_j>
 \frac{\Xi_i}{5184\,\delta L}.
\label{eq:V44-mass-escalation}
\end{equation}
\end{theorem}

\begin{proof}
By \eqref{eq:V44-partner-stock},
\[
 \frac{\Xi_i^2}{144}<S_{ij}^2.
\]
Use \Cref{lem:V44-row-power},
\eqref{eq:V44-row-H-domination}, and
\eqref{eq:V44-all-weak}:
\[
 \frac{\Xi_i^2}{144}
 <
 36N_{ij}H_{ij}^{(i)}
 \le36\delta N_{ij}\Xi_i\Xi_j.
\]
Cancel the positive factor \(\Xi_i\).  The two alternatives follow by
splitting at \(N_{ij}=L\).
\end{proof}

\begin{corollary}[Finite escalation cannot cycle]
\label{cor:V44-no-cycle}
Choose \(L\ge1\) and
\(\delta<(10368L)^{-1}\).  Along every sparse balanced transport edge
\eqref{eq:V44-mass-escalation},
\begin{equation}
 \Xi_j>2\Xi_i.
\label{eq:V44-doubling}
\end{equation}
Hence a chain of sparse balanced transports among four rows has
length at most three.  It terminates at a row with a dense partner
bank, a private-dominant row, or an unbalanced row stock greater than
one quarter of that row mass.
\end{corollary}

\begin{proof}
The choice of \(\delta\) turns
\eqref{eq:V44-mass-escalation} into
\eqref{eq:V44-doubling}.  A repeated vertex would imply that its
positive mass is strictly larger than twice itself, so no vertex
repeats.  There are four vertices.

At a terminal row which is not dense, the sparse balanced rule can
continue whenever its balanced row stock exceeds one quarter of its
mass.  If it cannot continue, then \(S_i\le\Xi_i/4\).  Since
\(\Xi_i=p_i+S_i+U_i\), either \(p_i\ge\Xi_i/2\), or
\[
 U_i=\Xi_i-p_i-S_i>\frac14\Xi_i.
\]
These are respectively the private and unbalanced terminal
alternatives.
\end{proof}

\begin{corollary}[Quantitative cut-tail transfer factor]
\label{cor:V44-cut-transfer-factor}
On the partner bank selected by
\eqref{eq:V44-partner-stock}, the balanced cut-tail theorem gives
\begin{equation}
 \kappa_\otimes(\mathcal B_{i|j})
 \le
 \frac{C H_{ij}^{(i)}}{s_\alpha S_{ij}^2}
 \le
 \frac{C\delta\,\Xi_j}{s_\alpha\Xi_i}.
\label{eq:V44-cut-transfer-factor}
\end{equation}
\end{corollary}

\begin{proof}
Apply \Cref{thm:V44-balanced-tail} to the cut separating \(i\) from
the chosen partner \(j\), retaining only the row-rooted bank.  Then
use \eqref{eq:V44-partner-stock},
\eqref{eq:V44-row-H-domination}, and
\eqref{eq:V44-all-weak}.
\end{proof}

\begin{assessment}[The balanced obstruction is now finite]
\label{ass:V44-finite-obstruction}
There is no longer an uncontrolled support-size or partition-state
problem.  A shared-hot internal row is paid by unbalanced decay, or it
creates a balanced partner edge.  A sparse sequence of such edges
strictly increases row mass and stops within three steps; a dense edge
has the cut-tail estimate \eqref{eq:V44-balanced-tail}.  What remains
is a finite graph-allocation problem: insert the at most three
transport edges into the original tree, choose a comparison spanning
tree, and ensure that the unique final resolvent and a possible
nonprivate output payer are charged only once.
\end{assessment}

\begin{openproblem}[V45 mandate: audit and transport-tree compiler]
\label{op:V45-mandate}
\begin{enumerate}
\item Perform the mandatory read-only audit of the newest active SM
      and NS notes before adding new mathematics.
\item Import only proof-order or bookkeeping lessons relevant to the
      finite transport chain; do not import an unproved estimate.
\item Form the connected multigraph consisting of the original
      quartic tree and the at most three transport edges of
      \Cref{cor:V44-no-cycle}.
\item Prove a comparison-tree selection which handles repeated edges
      and makes the terminal private/unbalanced payer internal when
      its mass is needed.
\item Allocate the sole output resolvent before applying
      \eqref{eq:V44-cut-transfer-factor}.
\item Treat the nonprivate high-output mass \(Z_\nu\) by labelled
      ancestry on this same finite graph.
\end{enumerate}
\end{openproblem}

\begin{firewall}[Claim boundary after compact V44]
\label{fire:V44-current-boundary}
V44 proves the cut-connected balanced-tail estimate and a finite
all-weak mass-transport alternative.  It does not yet prove that the
resulting transport edges can be converted to one quartic marked tree
with one output payer.  The nonprivate high-output branch remains
open.  Therefore global quartic MTA, all-order MTA, WUFC/CPM,
continuum Yang--Mills existence, and a positive mass gap remain open.
\end{firewall}

\section{Append-only record for compact V44}
\label{sec:V44-record}

V44 was appended to the complete compact V43 file.  The exact V43
parent had \(17390\) newline-delimited source records, \(573510\)
bytes, and SHA--256 digest
\[
\texttt{7d62d41fbb18b2f209782782bc0c69698ad1f2cb9e94c4d4da74cf35399d2553}.
\]
No compact V43 mathematical content was changed.  The sole final
\verb|\end{document}| is restored below.

% =============================================================================
% END APPENDED COMPACT VERSION V44 MATERIAL
% =============================================================================

% The compact V44 document terminator is intentionally disabled here:
% \end{document}

% =============================================================================
% BEGIN APPENDED COMPACT VERSION V45 MATERIAL
% =============================================================================

\chapter{V45: Companion audit, shared-output ancestry, and the finite
transport graph}
\label{ch:V45}

\section{Mandatory SM/NS companion audit}
\label{sec:V45-audit}

The audit was performed before the V45 argument.  During the audit the
parallel SM programme advanced from V49 to V50, so the newer file was
used.  Its exact identity was
\[
\begin{array}{ll}
\text{file:}&
\texttt{Renewal\_SM\_Einstein\_Continuum\_Research\_Notes\_V50.tex},\\
\text{records:}&18451,\\
\text{bytes:}&667784,\\
\text{SHA--256:}&
\texttt{e2fceb86428998694275ab9cade2a9de86e5e87c86a0a94a8638cac3e9bbb2e2}.
\end{array}
\]
SM V47 shows that perturbative coefficients and finitely many local
normalizations do not by themselves select an exponentially flat
nonlocal completion.  V48 uses reflection positivity only after the
complete Euclidean observable is assembled.  V49 separates
fixed-background tightness from uniqueness and from the dynamical
conformal problem.  V50 proves that faster \(p^{-4}\) decay of a local
fourth-order scalar covariance is incompatible with a nonzero positive
K\"all\'en--Lehmann measure.  Its relevant proof-order lesson is that
improved decay of signed components is not a substitute for the
positive weighted stock required by the final physical carrier.

The newest NS companion remained
\[
\begin{array}{ll}
\text{file:}&
\texttt{Renewal\_NS\_Stability\_Research\_Notes\_V31.tex},\\
\text{records:}&23052,\\
\text{bytes:}&887537,\\
\text{SHA--256:}&
\texttt{f1121b62238ad5491bb7cf03635ea9934942edf573ed8faaa9ee79e1d38a6696}.
\end{array}
\]
NS V31 proves an exact heat-formation ladder and a one-use geometric
parent sum.  It also proves that replacing the critical source norm by
a weaker energy-paid norm merely returns the missing powers through
the shell weights.  Its relevant lesson is that a mark and the
weighted moment needed to remove that mark are distinct debits; the
same transport edge may not be charged once as an overlap and again as
an independent output payer.

\begin{assessment}[V45 audit transfer]
\label{ass:V45-audit-transfer}
No SM or NS estimate is imported.  Two proof-order rules are retained:
\begin{enumerate}
\item keep the gapped \(s_\alpha N\) factor and the protected
      \(N\)-capacity factor separate until the final balanced
      optimization;
\item insert every transport mark and every output payer into one
      globally assembled finite graph before taking capacity or
      discarding a normalized edge.
\end{enumerate}
\end{assessment}

\section{The balanced tail has two different currencies}
\label{sec:V45-two-currencies}

The combined estimate \eqref{eq:V44-balanced-tail} is convenient but
conceals the distinction emphasized by the audit.  Let
\(\mathcal G_{C|D}\) and \(\mathcal P_{C|D}\) be respectively the
gapped and protected parts of the cut-connected balanced tail, formed
after the cancellations of V43.

\begin{proposition}[Separated gapped and protected cut ledgers]
\label{prop:V45-separated-cut-ledgers}
For a cut stock of size \(N=N_{C|D}\), mass \(S=S_{C|D}\), and mark
\(H=H_{C|D}\),
\begin{align}
 \kappa_\otimes(\mathcal G_{C|D})
 &\le C\min\{1,e^{-cs_\alpha N}\},
\label{eq:V45-gapped-endpoint}\\
 s_\alpha S^2
 \kappa_\otimes(\mathcal G_{C|D})
 &\le CH,
\label{eq:V45-gapped-ledger}\\
 N S^2
 \kappa_\otimes(\mathcal P_{C|D})
 &\le CH.
\label{eq:V45-protected-ledger}
\end{align}
The same estimates hold when at most three tree-marked coordinates
retain their residual analytic multipliers.
\end{proposition}

\begin{proof}
In the proof of \Cref{thm:V44-balanced-tail}, do not combine the two
lines of the local alternative.  The nonprotected line has the
power-capacity endpoint
\[
 C\min\{1,e^{-cs_\alpha N}\}
\]
from the V23/V30 scalar functions.  Since
\(S^2\le36NH\),
\[
 s_\alpha S^2
 \min\{1,e^{-cs_\alpha N}\}
 \le
 C(s_\alpha N)\min\{1,e^{-cs_\alpha N}\}H
 \le CH.
\]
This proves \eqref{eq:V45-gapped-endpoint} and
\eqref{eq:V45-gapped-ledger}.

On the protected line, every coordinate has a crossing invariant
projector of capacity at most a universal \(\rho<1\).  Hence
\[
 \kappa_\otimes(\mathcal P_{C|D})\le C\rho^N.
\]
Together with \(S^2\le36NH\),
\[
 \frac{NS^2}{H}
 \kappa_\otimes(\mathcal P_{C|D})
 \le C N^2\rho^N\le C,
\]
which is \eqref{eq:V45-protected-ledger}.  The marked-coordinate
statement follows because V44 retained, rather than discarded, the
same gapped or protected multiplier at every mark.
\end{proof}

\begin{remark}[Why the currencies are not interchangeable]
\label{rem:V45-currencies}
The gapped line pays with \(s_\alpha N\); the protected line pays with
tensor capacity \(\rho^N\) and yields the stronger explicit factor
\(N^{-1}\).  Replacing both by the unmarked norm bound one loses the
only mechanism that cancels a dense transport multiplicity.
\end{remark}

\section{Exact ancestry of the nonprivate output mass}
\label{sec:V45-output-ancestry}

Recall
\(\Xi(\boldsymbol U_\nu)=P+Z_\nu\) from
\eqref{eq:V41-output-mass-split}.  At a nonprivate coordinate \(e\),
the output spin is contained in the tensor product of at most four
active inputs.  The \(SU(2)\) triangle inequality and Cauchy--Schwarz
give
\begin{equation}
 1+C_2(U_{\nu,e})
 \le C\sum_{i\in I_e}a_{i,e}.
\label{eq:V45-coordinate-output-ancestry}
\end{equation}

\begin{lemma}[Labelled shared-output ancestry]
\label{lem:V45-shared-output-ancestry}
In every physical block,
\begin{equation}
\boxed{
 Z_\nu
 \le C\sum_{i=1}^4 h_i
 =
 C\sum_{i=1}^4(S_i+U_i).}
\label{eq:V45-shared-output-ancestry}
\end{equation}
Thus the nonprivate high-output coefficient has a finite positive
majorant labelled by a row and by balanced or unbalanced stock.
\end{lemma}

\begin{proof}
If \(U_{\nu,e}\) has spin \(J\), then
\(J\le\sum_{i\in I_e}j_{i,e}\).  Since \(|I_e|\le4\),
\[
 1+J(J+1)
 \le C\sum_{i\in I_e}\bigl(1+j_{i,e}(j_{i,e}+1)\bigr),
\]
which is \eqref{eq:V45-coordinate-output-ancestry}.  Sum over all
nonprivate coordinates and interchange the finite row and coordinate
sums.  Equation \eqref{eq:V44-row-stock-exact} gives the last equality.
\end{proof}

\begin{theorem}[Unbalanced part of the nonprivate high-output payer]
\label{thm:V45-unbalanced-shared-output}
Let a high-output tree be \((M,W)\)-unbalanced-controlled.  The part of
the coefficient majorant in
\eqref{eq:V45-shared-output-ancestry} carrying
\(\sum_iU_i\) obeys the global marked-tree weight.
\end{theorem}

\begin{proof}
Since every unbalanced row contribution is bounded by the total
unbalanced stock,
\[
 \sum_iU_i\le4W.
\]
For a fixed tree monomial, apply the reserve/heavy dichotomy of
\Cref{lem:V42-reserve-heavy}.  In the reserve branch retain the
statewise moment
\eqref{eq:V42-statewise-unbalanced-moments}; in the heavy branch use
\eqref{eq:V43-heavy-full-stock}.  The output payer costs one factor
\(s_\alpha W\).  The two tree-excess degrees are paid, as in
\Cref{thm:V43-unbalanced-private-high}, by further powers of \(W\),
thermal constants, or disjoint private moments.

The total degree is three, so no moment beyond those already proved in
V7 and V42--V43 is used.  Hence
\[
 s_\alpha^3q_{\rm priv}
 \left(\sum_iU_i\right)B_\tau
 \times(\hbox{retained unbalanced multiplier})
 \le C_M.
\]
Multiplication by \(L_\tau/B_\tau\) proves the claim.  The sole final
resolvent was used only to declare the high-output denominator bounded;
it is not inserted into the stock.
\end{proof}

\begin{corollary}[Remaining high-output ancestry is balanced]
\label{cor:V45-high-balanced-only}
For an unbalanced-controlled tree, both its private output part \(P\)
and its unbalanced nonprivate output part \(\sum_iU_i\) are closed.
Its only unclosed output ancestry is therefore
\[
 \sum_iS_i.
\]
\end{corollary}

\begin{proof}
Use \Cref{thm:V43-unbalanced-private-high} for \(P\) and
\Cref{thm:V45-unbalanced-shared-output} for \(\sum_iU_i\).
\end{proof}

\section{The transport multigraph}
\label{sec:V45-transport-graph}

Let \(\tau\) be the original labelled quartic tree.  A sparse balanced
transport chain from \Cref{cor:V44-no-cycle} is
\[
 i_0i_1,\ i_1i_2,\ldots,i_{m-1}i_m,
 \qquad 0\le m\le3.
\]
Adjoin these edge occurrences to \(\tau\), allowing a transport edge
to be parallel to an original or earlier edge.  Denote the resulting
connected multigraph by \(\mathcal G\).  Give every edge occurrence its
normalized overlap mark \(\theta_e\in[0,1]\).

\begin{lemma}[Spanning-tree domination with one-use marks]
\label{lem:V45-spanning-tree-domination}
For every spanning tree \(\tau'\) contained in the underlying
multigraph,
\begin{equation}
\boxed{
 \prod_{e\in E(\mathcal G)}\theta_e
 \le
 \prod_{e\in E(\tau')}\theta_e.}
\label{eq:V45-spanning-tree-domination}
\end{equation}
Each occurrence, including a parallel occurrence, is used at most
once on the left.
\end{lemma}

\begin{proof}
Delete from \(\mathcal G\) all edge occurrences not belonging to
\(\tau'\).  Every deleted factor lies in \([0,1]\), so deletion can
only increase the product.  Parallel occurrences cause no exception.
\end{proof}

\begin{lemma}[When a terminal vertex can be made internal]
\label{lem:V45-terminal-internal}
Let \(v\) be a vertex of a connected finite multigraph.  There exists a
spanning tree in which \(v\) has degree at least two if and only if
\(v\) has at least two distinct neighbours in the underlying simple
graph.
\end{lemma}

\begin{proof}
Necessity is immediate.  For sufficiency, choose edges from \(v\) to
two distinct neighbours.  These two edges form a forest.  Every forest
in a connected graph extends to a spanning tree by repeatedly adding
an edge joining two different components.  The resulting tree retains
both selected edges at \(v\).
\end{proof}

\begin{corollary}[Exact repeated-edge obstruction]
\label{cor:V45-repeated-edge-obstruction}
Let \(i_m\) be the terminal vertex of a transport chain.  If it cannot
be made internal in a comparison spanning tree of \(\mathcal G\), then
it is a leaf of the original tree and every incident transport edge is
parallel to its unique original neighbour.  Conversely this
configuration cannot make \(i_m\) internal.
\end{corollary}

\begin{proof}
The original tree already supplies at least one neighbour.  By
\Cref{lem:V45-terminal-internal}, failure means that no added edge
supplies a distinct second neighbour, which is exactly the stated
parallel-edge configuration.
\end{proof}

\begin{proposition}[Finite comparison-tree reduction]
\label{prop:V45-finite-tree-reduction}
Suppose an analytic branch of the complete carrier has been bounded by
\[
 C\left(\prod_i\Xi_i\right)
 \prod_{e\in E(\mathcal G)}\theta_e
\]
for the original tree plus its transport-chain marks.  Then it obeys
the quartic marked-tree atlas.  If its terminal payer must be internal,
the same conclusion follows unless
\Cref{cor:V45-repeated-edge-obstruction} occurs.
\end{proposition}

\begin{proof}
Choose any spanning tree \(\tau'\subseteq\mathcal G\), or one with the
terminal vertex internal using
\Cref{lem:V45-terminal-internal}.  Apply
\eqref{eq:V45-spanning-tree-domination}.  The right side is precisely
the normalized marked tree associated with \(\tau'\).  There are only
six edge occurrences at most, sixteen labelled quartic trees, and
fixed transport length at most three, so summation costs a universal
constant.
\end{proof}

\begin{firewall}[The graph lemma does not manufacture a mark]
\label{fire:V45-no-manufactured-mark}
\Cref{prop:V45-finite-tree-reduction} applies only after the analytic
estimate has produced every edge occurrence on its left.  In
particular one may not use \(H_{ij}\) in
\eqref{eq:V44-cut-transfer-factor} to transport row mass and then
insert a second copy of \(\theta_{ij}\) into the multigraph.  The
transport edge is that same one-use mark.
\end{firewall}

\section{The exact post-audit obstruction}
\label{sec:V45-obstruction}

\begin{assessment}[What remains after the V45 compiler]
\label{ass:V45-obstruction}
All support-dependent combinatorics have been removed.  The unresolved
balanced branch has at most three transport edges and one of two
endpoints:
\begin{enumerate}
\item a dense cut bank, where the gapped and protected currencies in
      \eqref{eq:V45-gapped-ledger}--%
      \eqref{eq:V45-protected-ledger} must be combined with the
      original two tree-excess powers;
\item a private or unbalanced terminal payer, where the comparison-tree
      lemma applies except for the explicit repeated-edge geometry of
      \Cref{cor:V45-repeated-edge-obstruction}.
\end{enumerate}
For high output, the only additional term is the balanced ancestry
\(\sum_iS_i\).  The private and unbalanced ancestry terms are closed.
\end{assessment}

\begin{openproblem}[V46 mandate: scalar transport ledger]
\label{op:V46-mandate}
\begin{enumerate}
\item Write the balanced tail bound with its transport edge normalized
      immediately as \(H_{ij}/(\Xi_i\Xi_j)\).
\item Track separately the protected \(N^{-1}\) gain and the gapped
      \(s_\alpha\)-gain; do not replace both by
      \eqref{eq:V44-balanced-tail}.
\item Along a sparse chain, telescope the row-mass ratios before
      applying any tree norm.
\item In the dense endpoint, optimize \(s_\alpha N\) against the two
      tree-excess powers.
\item In the repeated-edge endpoint, determine whether the second
      occurrence supplies a genuine square mark or merely repeats the
      same coordinate bank.
\item Add the balanced output payer \(\sum_iS_i\) to this single
      scalar ledger and use the sole final resolvent once.
\end{enumerate}
\end{openproblem}

\begin{firewall}[Claim boundary after compact V45]
\label{fire:V45-current-boundary}
V45 completes the mandatory companion audit, separates the balanced
gapped/protected currencies, closes the unbalanced part of the
nonprivate high-output payer, and proves the finite transport-graph
compiler.  It does not close the scalar balanced transport ledger or
the repeated-edge endpoint.  Therefore the all-weak hypothesis of the
V36 compiler, global quartic MTA, all-order MTA, WUFC/CPM, continuum
Yang--Mills existence, and a positive mass gap remain open.
\end{firewall}

\section{Append-only record for compact V45}
\label{sec:V45-record}

V45 was appended to the complete compact V44 file.  The exact V44
parent had \(17788\) newline-delimited source records, \(587539\)
bytes, and SHA--256 digest
\[
\texttt{8119f8b7a2553957be1ca93988c018b057458b4a552b7278c60d2a8b4608af9e}.
\]
No compact V44 mathematical content was changed.  The sole final
\verb|\end{document}| is restored below.

% =============================================================================
% END APPENDED COMPACT VERSION V45 MATERIAL
% =============================================================================

% The compact V45 document terminator is intentionally disabled here:
% \end{document}

% =============================================================================
% BEGIN APPENDED COMPACT VERSION V46 MATERIAL
% =============================================================================

\chapter{V46: Telescoping protected transport and the overlap gate}
\label{ch:V46}

\section{Normalize a balanced transport at once}
\label{sec:V46-normalized-step}

Fix a row-rooted partner bank from V44 and abbreviate
\[
 S=S_{ij},\qquad
 H=H_{ij}^{(i)},\qquad
 N=N_{ij},\qquad
 \vartheta_{ij}:=\frac{H}{\Xi_i\Xi_j}.
\]
Then \(0\le\vartheta_{ij}\le\theta_{ij}\le1\).  If
\(S>\Xi_i/12\), the two cut currencies in V45 take a particularly
transparent form.

\begin{lemma}[Normalized protected and gapped transport factors]
\label{lem:V46-normalized-factors}
For the row-rooted bank selected by
\eqref{eq:V44-partner-stock},
\begin{align}
 \kappa_\otimes(\mathcal P_{ij})
 &\le
 \frac{C}{N}\,
 \vartheta_{ij}\frac{\Xi_j}{\Xi_i},
\label{eq:V46-protected-factor}\\
 \kappa_\otimes(\mathcal G_{ij})
 &\le
 C\min\left\{
 e^{-cs_\alpha N},
 \frac1{s_\alpha}\,
 \vartheta_{ij}\frac{\Xi_j}{\Xi_i}
 \right\}.
\label{eq:V46-gapped-factor}
\end{align}
\end{lemma}

\begin{proof}
By \eqref{eq:V45-protected-ledger},
\[
 \kappa_\otimes(\mathcal P_{ij})
 \le\frac{CH}{NS^2}
 \le
 \frac{144C}{N}\frac{H}{\Xi_i^2}
 =
 \frac{144C}{N}\,
 \vartheta_{ij}\frac{\Xi_j}{\Xi_i}.
\]
Absorb \(144\) into \(C\).  The second entry of the minimum follows in
the same way from \eqref{eq:V45-gapped-ledger}; its first entry is
\eqref{eq:V45-gapped-endpoint}.
\end{proof}

\begin{remark}[Direction of the mass ratio]
\label{rem:V46-ratio-direction}
The factor \(\Xi_j/\Xi_i\) is not an error and may be large on a sparse
transport edge.  It is harmless only if ratios from the complete
transport chain are multiplied before the terminal moment is used.
Bounding each step separately would lose the telescope.
\end{remark}

\section{Terminal moment certificates}
\label{sec:V46-terminal}

\begin{definition}[Moment-controlled terminal row]
\label{def:V46-terminal}
A row \(v\) has a terminal moment certificate through order three if
the complete carrier retains a scalar or positive tail \(Q_v\) such
that
\begin{equation}
 0\le Q_v\le1,
 \qquad
 (s_\alpha\Xi_v)^kQ_v\le C_k,
 \quad k=1,2,3.
\label{eq:V46-terminal-certificate}
\end{equation}
\end{definition}

\begin{lemma}[Private and unbalanced terminals]
\label{lem:V46-terminal-types}
Each of the terminal alternatives in
\Cref{cor:V44-no-cycle} has a certificate:
\begin{enumerate}
\item if \(p_v\ge\Xi_v/2\), take the retained private factor \(q_v\);
\item if \(U_v>\Xi_v/4\), take the retained common unbalanced tail.
\end{enumerate}
\end{lemma}

\begin{proof}
In the private case,
\(\Xi_v\le2p_v\), so
\eqref{eq:V41-row-private-moments} proves
\eqref{eq:V46-terminal-certificate}.  In the unbalanced case,
\[
 W\ge U_v>\Xi_v/4.
\]
The reserve/heavy estimates
\eqref{eq:V42-statewise-unbalanced-moments} and
\eqref{eq:V43-heavy-full-stock} prove the same three moment bounds.
Balanced and unbalanced coordinate banks are disjoint, so this
terminal tail is not one of the protected transport factors below.
\end{proof}

\section{The scalar telescope}
\label{sec:V46-scalar-telescope}

Let
\[
 i_0\longrightarrow i_1\longrightarrow\cdots
 \longrightarrow i_m,
 \qquad 1\le m\le3,
\]
be a sparse balanced transport chain.  Put
\[
 R_\ell:=\frac{\Xi_{i_{\ell+1}}}{\Xi_{i_\ell}},
 \qquad
 \vartheta_\ell:=\vartheta_{i_\ell i_{\ell+1}},
 \qquad
 N_\ell:=N_{i_\ell i_{\ell+1}}.
\]
By construction \(R_\ell>2\).

\begin{lemma}[Exact ratio telescope]
\label{lem:V46-ratio-telescope}
For \(k=1,2,3\),
\begin{equation}
\boxed{
 (s_\alpha\Xi_{i_0})^k
 \left[
 \prod_{\ell=0}^{m-1}
 \left(
  \frac{\vartheta_\ell R_\ell}{N_\ell}
 \right)
 \right]
 \frac1{(s_\alpha\Xi_{i_m})^k}
 =
 \left(\prod_{\ell=0}^{m-1}
 \frac{\vartheta_\ell}{N_\ell}\right)
 \left(\frac{\Xi_{i_0}}{\Xi_{i_m}}\right)^{k-1}.}
\label{eq:V46-ratio-telescope}
\end{equation}
In particular the left side is at most
\(\prod_\ell\vartheta_\ell\).
\end{lemma}

\begin{proof}
The row ratios telescope once:
\[
 \prod_{\ell=0}^{m-1}R_\ell
 =\frac{\Xi_{i_m}}{\Xi_{i_0}}.
\]
Insert this identity and cancel powers.  Since
\(\Xi_{i_m}>\Xi_{i_0}\) and every \(N_\ell\ge1\), the final bound
follows.
\end{proof}

\begin{theorem}[Protected sparse-chain scalar closure]
\label{thm:V46-protected-scalar}
Suppose a carrier branch contains:
\begin{enumerate}
\item \(k\in\{1,2,3\}\) unpaid factors of
      \(s_\alpha\Xi_{i_0}\), consisting of zero, one, or two tree
      excess powers and at most one labelled balanced-output payer;
\item the protected factor
      \eqref{eq:V46-protected-factor} from every edge of a sparse
      transport chain;
\item a terminal certificate
      \eqref{eq:V46-terminal-certificate} at \(i_m\).
\end{enumerate}
Then its scalar coefficient is bounded by the product of the transport
marks \(\prod_\ell\vartheta_\ell\), up to a universal constant.
\end{theorem}

\begin{proof}
Multiply the \(m\) instances of
\eqref{eq:V46-protected-factor}.  Use the order-\(k\) terminal
certificate and then \Cref{lem:V46-ratio-telescope}.  The factors
\(\prod N_\ell^{-1}\) may be discarded because they are at most one.
No fourth moment is needed: a quartic tree has two excess powers and
the high-output ancestry contributes at most one more.
\end{proof}

\begin{corollary}[Parallel terminal edges are harmless in the scalar
protected branch]
\label{cor:V46-parallel-harmless}
The repeated-edge geometry of
\Cref{cor:V45-repeated-edge-obstruction} does not obstruct
\Cref{thm:V46-protected-scalar}.  After the terminal moment has been
used, parallel normalized transport marks can be removed by
\Cref{lem:V45-spanning-tree-domination}.
\end{corollary}

\begin{proof}
The private multiplier or unbalanced tail in
\eqref{eq:V46-terminal-certificate} is global and does not require the
terminal row to be internal in the comparison tree.  Thus the scalar
telescope is valid before a spanning tree is selected.  Every parallel
mark lies in \([0,1]\), so the finite multigraph compiler then removes
surplus occurrences.
\end{proof}

\section{When the analytic factors may not yet be multiplied}
\label{sec:V46-compatible-banks}

The scalar calculation is not by itself permission to multiply cut
capacities.  Even coordinate-disjoint banks retain the same global
partition variable \(\pi\) from V42, and each cut-connected
cancellation is performed by the same functional \(\Lambda\).

\begin{proposition}[Disjoint support is not sufficient]
\label{prop:V46-disjoint-not-sufficient}
There is no algebraic implication
\[
 \Lambda(x)=0,\quad\Lambda(y)=0
 \quad\Longrightarrow\quad
 \Lambda(xy)=0
\]
for partition-state functions \(x,y\), even when \(x\) and \(y\)
represent disjoint coordinate banks.
\end{proposition}

\begin{proof}
Choose a one-block partition \(\pi_1\), for which
\(\mu(\pi_1)=1\), and a three-block partition \(\pi_3\), for which
\(\mu(\pi_3)=2\).  Let \(x=y\) vanish off these two states and set
\[
 x(\pi_1)=1,\qquad x(\pi_3)=-\frac12.
\]
Then
\[
 \Lambda(x)=1+2(-1/2)=0,
 \qquad
 \Lambda(xy)=1+2(1/4)=\frac32.
\]
Coordinate disjointness makes the local tensor factors commute, but it
does not replace the single global state sum by a product of two state
sums.
\end{proof}

\begin{proposition}[Conditional compatible protected-chain compiler]
\label{prop:V46-compatible-protected}
Suppose a single global cut-forest identity, formed before norms,
produces the product of protected factors
\eqref{eq:V46-protected-factor} along a sparse transport chain, with
each coordinate multiplier used once.  If the terminal row is
private- or unbalanced-controlled and the total degree \(k\) is at
most three, then that branch obeys the global marked-tree atlas.
\end{proposition}

\begin{proof}
The assumed identity is exactly the missing operator-level
authorization to use the product appearing in
\Cref{thm:V46-protected-scalar}.  Apply that theorem.  The result
contains the original three tree marks and the transport marks once,
so \Cref{prop:V45-finite-tree-reduction} gives a quartic comparison
tree.  A labelled balanced output contributes the optional third
degree already included in the scalar theorem.
\end{proof}

\begin{firewall}[One global state sum, one capacity]
\label{fire:V46-no-overlap-product}
Separate cut-connected estimates may not be multiplied, whether or
not their coordinate banks overlap.  If banks overlap, there is also a
local same-tensor-factor obstruction.  If they are disjoint, the
regression in \Cref{prop:V46-disjoint-not-sufficient} shows that the
common partition-state cancellation is still an obstruction.  The
needed object is one simultaneous cut-forest identity followed by one
capacity estimate.
\end{firewall}

\section{Exact classification of bank overlap}
\label{sec:V46-overlap}

\begin{lemma}[Overlap forces a higher-incidence chain coordinate]
\label{lem:V46-overlap-incidence}
Suppose a coordinate belongs to the row-rooted banks of two distinct
edges of a simple transport chain.  If the edges are adjacent, the
coordinate has incidence at least three.  If they are disjoint, it has
incidence four.
\end{lemma}

\begin{proof}
Membership in a row-rooted \(ij\)-bank means that both \(i\) and \(j\)
are active at the coordinate.  Two adjacent chain edges have three
distinct endpoints; two disjoint chain edges have four.  All those
rows are active at the common coordinate.
\end{proof}

\begin{corollary}[There are only two coordinate-overlap geometries]
\label{cor:V46-two-overlaps}
Every actual coordinate overlap between distinct transport banks is
supported on:
\begin{enumerate}
\item an incidence-three coordinate shared by two adjacent transport
      edges; or
\item an incidence-four coordinate shared by two or three transport
      edges.
\end{enumerate}
No incidence-two coordinate can create a transport-bank overlap.
\end{corollary}

\begin{proof}
A chain has at most three edges.  Apply
\Cref{lem:V46-overlap-incidence} to every overlapping pair.
\end{proof}

\section{The genuinely gapped residual}
\label{sec:V46-gapped-residual}

\begin{proposition}[Why the gapped line cannot be substituted into the
protected telescope]
\label{prop:V46-gapped-not-protected}
Replacing a protected factor in
\eqref{eq:V46-ratio-telescope} by the second entry of
\eqref{eq:V46-gapped-factor} inserts one extra
\(s_\alpha^{-1}\).  With up to three transport steps, the available
terminal moments through order three do not uniformly pay all such
losses in addition to two tree-excess powers and an output payer.
\end{proposition}

\begin{proof}
The protected factor is
\(\vartheta R/N\), whereas the marked gapped estimate is
\(\vartheta R/s_\alpha\).  Their ratio is \(N/s_\alpha\).  If this
replacement is made independently at \(g\) steps, it introduces
\(s_\alpha^{-g}\) before the endpoint estimate
\(e^{-cs_\alpha N}\) has been combined across the complete carrier.
The scalar data allow \(g=3\), while the tree and output can already
require three terminal degrees.  The known compiler provides no
moments of order six.  Thus separate substitution is not a proof.
\end{proof}

\begin{assessment}[Post-V46 reduction]
\label{ass:V46-reduction}
The protected sparse-chain scalar ledger is complete, parallel
transport edges are harmless, and
\Cref{prop:V46-disjoint-not-sufficient} prevents an invalid
operator-level conclusion.  The remaining balanced obstruction has
three precise components:
\begin{enumerate}
\item a simultaneous cut-forest identity for the single global
      partition-state carrier, even when its banks are disjoint;
\item overlapping incidence-three/four protected banks inside that
      identity;
\item gapped transport states, whose \(s_\alpha N\) endpoint factors
      must be assembled globally rather than charged step by step.
\end{enumerate}
The balanced high-output ancestry is included in the same two
components; it is not a separate unlabelled output sum.
\end{assessment}

\begin{openproblem}[V47 mandate: simultaneous multi-cut carrier]
\label{op:V47-mandate}
\begin{enumerate}
\item Construct one row-partition cut-forest identity whose leaves are
      the complete protected transport chain; do not multiply
      separately cancelled cut responses.
\item For one incidence-three coordinate supporting two adjacent
      transport edges, compute the full invariant-projector normal
      form across both cuts at once.
\item Do the same for an incidence-four coordinate supporting a
      three-edge chain.
\item Bound the simultaneous product-state capacity without
      multiplying two capacities on the same tensor factor.
\item For the gapped complement, combine all bank counts into one
      \(N_{\rm gap}\) before using \(e^{-cs_\alpha N_{\rm gap}}\).
\item Insert the resulting multi-cut marks into the finite graph of
      V45 and retain the terminal telescope of V46.
\end{enumerate}
\end{openproblem}

\begin{firewall}[Claim boundary after compact V46]
\label{fire:V46-current-boundary}
V46 closes the scalar protected transport ledger with terminal
private/unbalanced control and proves that disjoint coordinate support
does not by itself lift that ledger to the operator carrier.  The
simultaneous cut-forest identity, overlapping higher-incidence
projectors, and globally gapped transport states remain open.
Consequently the complete all-weak estimate, global quartic MTA,
all-order MTA, WUFC/CPM, continuum Yang--Mills existence, and a
positive mass gap remain open.
\end{firewall}

\section{Append-only record for compact V46}
\label{sec:V46-record}

V46 was appended to the complete compact V45 file.  The exact V45
parent had \(18193\) newline-delimited source records, \(602507\)
bytes, and SHA--256 digest
\[
\texttt{75ce35573c027ed615c5bbdb208f4affc78e94173ece4f472ddd9eaeb5459b85}.
\]
No compact V45 mathematical content was changed.  The sole final
\verb|\end{document}| is restored below.

% =============================================================================
% END APPENDED COMPACT VERSION V46 MATERIAL
% =============================================================================

% The compact V46 document terminator is intentionally disabled here:
% \end{document}

% =============================================================================
% BEGIN APPENDED COMPACT VERSION V47 MATERIAL
% =============================================================================

\chapter{V47: An exact connected-graph identity on the four-row
partition lattice}
\label{ch:V47}

\section{Merge differences on \texorpdfstring{\(\Pi_4\)}{Pi4}}
\label{sec:V47-merge}

Let
\[
 \widehat0:=\{\{1\},\{2\},\{3\},\{4\}\},
 \qquad
 \widehat1:=\{\{1,2,3,4\}\}.
\]
For an unordered row edge \(a=\{i,j\}\), define the join operator on a
function \(M:\Pi_4\to\mathcal X\) by
\[
 (\mathsf J_aM)(\pi):=M(\pi\vee\{\{i,j\}\}),
 \qquad
 \nabla_a:=\mathsf J_a-\mathsf I.
\]
The \(\mathsf J_a\)'s commute and are idempotent.  For a simple graph
\(G\) on \([4]\), put
\begin{equation}
 \mathfrak D_GM
 :=
 \left(\prod_{a\in E(G)}\nabla_a\right)M(\widehat0).
\label{eq:V47-graph-difference}
\end{equation}
Equivalently,
\begin{equation}
 \mathfrak D_GM
 =
 \sum_{F\subseteq E(G)}
 (-1)^{|E(G)|-|F|}M(\pi(F)),
\label{eq:V47-forest-expansion}
\end{equation}
where \(\pi(F)\) is the partition into connected components of the
spanning subgraph \(([4],F)\).

\section{Why trees alone cannot represent the cumulant}
\label{sec:V47-tree-rank}

\begin{theorem}[Exact obstruction to an unweighted tree-difference
identity]
\label{thm:V47-tree-rank-obstruction}
Let \(\boldsymbol\mu\in\mathbb R^{\Pi_4}\) be the fourth-cumulant
coefficient vector
\[
 \boldsymbol\mu(\pi)=(|\pi|-1)!(-1)^{|\pi|-1}.
\]
Then \(\boldsymbol\mu\) is not in the linear span of the sixteen
coefficient vectors of \(\mathfrak D_TM\), with \(T\) ranging over
labelled spanning trees.
\end{theorem}

\begin{proof}
Define a linear functional \(\mathscr Y\) on partition coefficient
vectors by assigning weight \(3\) to \(\widehat1\), weight \(1\) to
each of the seven two-block partitions, and weight zero to all other
partitions.

For a spanning tree \(T\), the full edge set contributes
\(+M(\widehat1)\) to \eqref{eq:V47-forest-expansion}.  Its three
two-edge subforests give three distinct two-block partitions, each
with coefficient \(-1\).  Subforests with at most one edge have at
least three blocks.  Therefore
\[
 \mathscr Y(\mathfrak D_T)=3-3=0
\]
for every \(T\).

On the other hand, the cumulant vector has coefficient \(1\) on
\(\widehat1\) and coefficient \(-1\) on each two-block partition.
Hence
\[
 \mathscr Y(\boldsymbol\mu)=3-7=-4\ne0.
\]
Thus no linear combination of tree-difference vectors equals
\(\boldsymbol\mu\).
\end{proof}

\begin{assessment}[Meaning of the obstruction]
\label{ass:V47-tree-obstruction-meaning}
The frozen tree numerator is a valid estimate, but the exact global
M\"obius cancellation is not an unweighted sum of three merge
differences.  Consequently the V46 protected factors cannot be
authorized merely by choosing one of the sixteen trees.  Either
partition-dependent weights or a fixed family with surplus graph
edges is necessary.
\end{assessment}

\section{A symmetric connected-graph repair}
\label{sec:V47-graph-repair}

Let \(\mathfrak S\), \(\mathfrak P\), \(\mathfrak C\), and
\(\mathfrak A\) be respectively the labelled graph classes consisting
of:
\begin{enumerate}
\item the four three-edge stars;
\item the twelve three-edge paths;
\item the three four-cycles;
\item the twelve four-edge paws, namely a triangle with one pendant
      edge.
\end{enumerate}
The letters refer only to graph classes and not to protected
projectors.

\begin{theorem}[Exact connected-graph cumulant identity]
\label{thm:V47-connected-graph-identity}
For every operator-valued function \(M\) on \(\Pi_4\),
\begin{equation}
\boxed{
 \sum_{\pi\in\Pi_4}\boldsymbol\mu(\pi)M(\pi)
 =
 \sum_{G\in\mathfrak S}\mathfrak D_GM
 +\frac23\sum_{G\in\mathfrak P}\mathfrak D_GM
 -\frac13\sum_{G\in\mathfrak C}\mathfrak D_GM
 +\frac7{12}\sum_{G\in\mathfrak A}\mathfrak D_GM.}
\label{eq:V47-connected-graph-identity}
\end{equation}
Thus four-edge connected graphs repair the tree-rank obstruction; no
five- or six-edge graph is required.
\end{theorem}

\begin{proof}
Both sides are invariant under permutation of the four row labels.
It suffices to compare the coefficient of \(M(\pi)\) for the five
integer partition types of \(4\).  Direct counting of subgraphs in
\eqref{eq:V47-forest-expansion} gives the following table.  Each graph
column is already summed over all labelled graphs in that class.
\[
\begin{array}{c|r|rrrr}
\text{block sizes of }\pi&
\boldsymbol\mu(\pi)&
\mathfrak S&\mathfrak P&\mathfrak C&\mathfrak A\\ \hline
4       & 1  & 4  & 12 & -9 & -24\\
3+1     &-1  &-3  & -6 &  3 &  12\\
2+2     &-1  & 0  & -4 &  2 &   4\\
2+1+1   & 2  & 2  &  6 & -2 &  -8\\
1+1+1+1 &-6  &-4  &-12 &  3 &  12
\end{array}
\]
For example, the star entry in the first row is four because there are
four labelled stars and the full edge set of each has sign \(+1\).
All other entries follow by choosing the subgraph whose connected
components have the displayed sizes.

Multiply the four graph columns by
\[
 1,\qquad\frac23,\qquad-\frac13,\qquad\frac7{12}
\]
respectively.  Row by row the results are
\(1,-1,-1,2,-6\), which are exactly the five values of
\(\boldsymbol\mu\).  This proves the identity for every partition and
hence for operator-valued \(M\).
\end{proof}

\begin{corollary}[Exact graph carrier for the Yang--Mills block]
\label{cor:V47-YM-graph-carrier}
With
\[
 M_\nu(\pi):=
 q_{\rm priv}\bigotimes_{e\in E_{\rm sh}}T_e(\pi),
\]
the complete physical fourth-cumulant block is the right side of
\eqref{eq:V47-connected-graph-identity} applied to \(M_\nu\).
All merge differences are therefore formed before the output norm,
capacity, and sole final resolvent.
\end{corollary}

\begin{proof}
The left side is
\eqref{eq:V42-exact-state-carrier}.  Apply
\Cref{thm:V47-connected-graph-identity}.
\end{proof}

\section{Coordinate resolution of a graph difference}
\label{sec:V47-coordinate-resolution}

\begin{lemma}[At most four marked coordinates]
\label{lem:V47-four-marks}
After expanding the coordinate-product differences in one
\(\mathfrak D_GM_\nu\), every resolved monomial distinguishes at most
\(|E(G)|\le4\) coordinates.  Every undisinguished coordinate remains
inside its exact global partition-state multiplier.
\end{lemma}

\begin{proof}
Apply the product telescope once for each factor
\(\nabla_a\) in \eqref{eq:V47-graph-difference}.  A row-merge
difference chooses one coordinate on which the two row blocks interact;
subsequent differences may choose the same coordinate or another one.
Thus at most one new coordinate is distinguished per graph edge.
Multilinearity leaves all other coordinate factors untouched.  Since
the largest graphs in
\eqref{eq:V47-connected-graph-identity} have four edges, the bound is
four.
\end{proof}

\begin{lemma}[The stock compilers extend from three to four marks]
\label{lem:V47-four-mark-stock}
All unbalanced and balanced retained-tail estimates of V43--V45 remain
valid for graph-resolved monomials with four marked coordinates.
\end{lemma}

\begin{proof}
For the unbalanced stock, if the unmarked reserve is below \(W/2\),
one of at most four marked coordinates has mass greater than \(W/8\).
\Cref{lem:V43-marked-decay} then supplies
\(e^{-cs_\alpha W/8}\), and moments through order three follow exactly
as in \Cref{cor:V43-heavy-replaces-reserve}.  In the reserve branch
nothing changes.

For the balanced stock, the proof of
\Cref{thm:V44-balanced-tail} already retains at a marked coordinate
the same uniform gapped factor or crossing protected capacity used at
an unmarked coordinate.  Its power-capacity product is independent of
whether a fixed coordinate also displays an edge Casimir factor.
Changing the fixed maximum number of marks from three to four changes
only a universal combinatorial constant.
\end{proof}

\begin{corollary}[Surplus graph edges cost no new normalized power]
\label{cor:V47-surplus-edge}
Every graph in
\eqref{eq:V47-connected-graph-identity} contains a spanning tree.  If
a resolved analytic estimate supplies all its normalized edge marks,
then its product is bounded by a quartic marked-tree monomial.
\end{corollary}

\begin{proof}
Stars and paths already are trees.  A four-cycle or paw is connected
with four edges.  Delete any cycle edge while preserving connectivity.
Because every normalized mark is at most one, apply
\Cref{lem:V45-spanning-tree-domination}.
\end{proof}

\section{The simultaneous protected branch}
\label{sec:V47-simultaneous-protected}

\begin{definition}[Edge-compatible protected resolution]
\label{def:V47-edge-compatible}
A resolved term of \(\mathfrak D_GM_\nu\) is edge-compatible with a
transport chain if every chain edge is one of the merge edges of
\(G\), its row-rooted protected stock is retained in the corresponding
merge difference, and no coordinate multiplier is used in two
separate capacity estimates.
\end{definition}

\begin{theorem}[Operator authorization of the disjoint protected
chain]
\label{thm:V47-disjoint-protected-authorized}
Suppose an edge-compatible protected resolution has pairwise
coordinate-disjoint transport banks and a private- or
unbalanced-controlled terminal row.  Then it obeys the global
marked-tree atlas, including one labelled balanced-output payer when
the total terminal degree is at most three.
\end{theorem}

\begin{proof}
Unlike the invalid product of separately cancelled responses,
\(\mathfrak D_GM_\nu\) is one exact global partition-lattice
difference formed by
\Cref{cor:V47-YM-graph-carrier}.  Within the fixed resolved term,
disjoint bank multipliers act on disjoint coordinate tensor factors.
One-sided capacity submultiplicativity now gives the product of the
protected factors \eqref{eq:V46-protected-factor}; there is no second
application of \(\Lambda\).

Apply \Cref{thm:V46-protected-scalar}.  The graph marks occur once in
the same \(\mathfrak D_G\) term.  Finally use
\Cref{cor:V47-surplus-edge}.  The labelled output payer is the optional
third terminal degree already covered by V46.
\end{proof}

\begin{assessment}[What the connected-graph repair accomplishes]
\label{ass:V47-repair-accomplishes}
\Cref{thm:V47-connected-graph-identity} supplies the operator-level
authorization which was explicitly absent in V46.  It does not
multiply separately cancelled cut responses.  Instead it forms one
global graph difference and takes capacity once.  As a result, the
pairwise disjoint protected transport branch is now genuinely closed.
\end{assessment}

\section{One globally gapped stock}
\label{sec:V47-global-gap}

\begin{lemma}[Gapped factors are assembled once]
\label{lem:V47-one-global-gap}
In a fixed resolved graph-difference term, let \(N_{\rm gap}\) be the
number of balanced coordinate multipliers lying in a nonprotected
output sector after all merge differences have been formed.  Their
joint norm is bounded by
\begin{equation}
\boxed{
 C_Ge^{-cs_\alpha N_{\rm gap}},}
\label{eq:V47-one-global-gap}
\end{equation}
where \(C_G\) is universal because \(|E(G)|\le4\).
\end{lemma}

\begin{proof}
Every such coordinate contributes the uniform spectral factor
\(r_\alpha=1-c_*s_\alpha\).  Undistinguished factors are contractions,
and at most four marked local differences contribute a fixed
coefficient sum.  Hence the product is bounded by
\[
 C_G r_\alpha^{N_{\rm gap}}
 \le C_Ge^{-c s_\alpha N_{\rm gap}}.
\]
Only one product norm is taken.
\end{proof}

\begin{firewall}[The global gap is not yet a mass moment]
\label{fire:V47-gap-not-mass}
\eqref{eq:V47-one-global-gap} removes the false factor
\(s_\alpha^{-g}\) from treating \(g\) gapped transport steps
separately.  It controls the count \(N_{\rm gap}\), not the row masses
carried by protected coordinates.  Converting the remaining tree
excess powers requires one joint scalar optimization with
\(S^2\le36NH\); it is not automatic from exponential count decay.
\end{firewall}

\section{The narrowed multi-cut residual}
\label{sec:V47-residual}

\begin{assessment}[Post-V47 obstruction]
\label{ass:V47-obstruction}
The exact global M\"obius assembly and the disjoint protected sparse
chain are closed.  The unresolved graph-difference terms are now:
\begin{enumerate}
\item an incidence-three or incidence-four coordinate used by two or
      more merge edges, requiring a simultaneous local multi-cut
      capacity rather than a product of capacities;
\item mixed protected/gapped resolutions for which the single global
      gap \eqref{eq:V47-one-global-gap} must be optimized together with
      the row-rooted masses and terminal moment.
\end{enumerate}
There are four connected graph classes and at most four marked
coordinates.  Thus the residual is finite and support-independent.
\end{assessment}

\begin{openproblem}[V48 mandate: local multi-cut normal forms]
\label{op:V48-mandate}
\begin{enumerate}
\item For incidence three, decompose the complete invariant projector
      relative to two adjacent merge edges and compute its fully
      row-product capacity.
\item For incidence four, do the same for path, star, cycle, and paw
      edge sets.
\item Compare the result with the product of the row-rooted factors
      required by \eqref{eq:V46-ratio-telescope}; record any loss
      explicitly rather than assuming independent cuts.
\item Combine the protected result with the one global gapped factor
      \eqref{eq:V47-one-global-gap}.
\item If a simultaneous projector supplies only one capacity power,
      return upstream to the all-weak overlap marks at that same
      coordinate instead of iterating a false cut estimate.
\end{enumerate}
\end{openproblem}

\begin{firewall}[Claim boundary after compact V47]
\label{fire:V47-current-boundary}
V47 proves the exact four-row connected-graph difference identity,
closes its disjoint protected transport branch, and assembles all
gapped coordinate counts into one exponential.  It does not close
higher-incidence multi-cut projectors or the final mixed scalar
optimization.  Hence the complete all-weak estimate, global quartic
MTA, all-order MTA, WUFC/CPM, continuum Yang--Mills existence, and a
positive mass gap remain open.
\end{firewall}

\section{Append-only record for compact V47}
\label{sec:V47-record}

V47 was appended to the complete compact V46 file.  The exact V46
parent had \(18600\) newline-delimited source records, \(616723\)
bytes, and SHA--256 digest
\[
\texttt{70308b19b67cccb51847cdba23be1a7511ab26dcedc4b9479d34883b4e30444c}.
\]
No compact V46 mathematical content was changed.  The sole final
\verb|\end{document}| is restored below.

% =============================================================================
% END APPENDED COMPACT VERSION V47 MATERIAL
% =============================================================================

% The compact V47 document terminator is intentionally disabled here:
% \end{document}

% =============================================================================
% BEGIN APPENDED COMPACT VERSION V48 MATERIAL
% =============================================================================

\chapter{V48: Sharp multi-cut capacity and coordinate-word
compression}
\label{ch:V48}

\section{The trivalent invariant supplies only one capacity power}
\label{sec:V48-trivalent-capacity}

Let
\[
 \iota_{j_1j_2j_3}
 \in V_{j_1}\otimes V_{j_2}\otimes V_{j_3}
\]
be the normalized \(SU(2)\)-invariant vector when it exists.  In the
standard magnetic basis its coefficients are Wigner \(3j\)-symbols.

\begin{lemma}[Full row-product capacity of a trivalent invariant]
\label{lem:V48-trivalent-upper}
If \(d_r=2j_r+1\), then
\begin{equation}
\boxed{
 \sup_{\substack{\|v_r\|=1\\r=1,2,3}}
 |\langle v_1\otimes v_2\otimes v_3,
          \iota_{j_1j_2j_3}\rangle|^2
 \le
 \min_r\frac1{d_r}
 =
 \frac1{\max_r d_r}.}
\label{eq:V48-trivalent-upper}
\end{equation}
\end{lemma}

\begin{proof}
Across the cut \(r\mid\{1,2,3\}\setminus\{r\}\), invariance and
Schur's lemma make the reduced density matrix on \(V_{j_r}\) equal to
\(d_r^{-1}I\).  The largest squared Schmidt coefficient across that
cut is therefore \(d_r^{-1}\).  A fully row-product vector is a
special case of a product vector across the cut, so its overlap is at
most \(d_r^{-1}\).  Apply this for all three \(r\).
\end{proof}

\begin{proposition}[The one-power bound is sharp]
\label{prop:V48-trivalent-sharp}
For integer or half-integer \(j\ge1\), the invariant with spins
\((j,j,1)\) has a row-product basis coefficient satisfying
\begin{equation}
\boxed{
 \left|
 \begin{pmatrix}
  j&j&1\\ j&-j&0
 \end{pmatrix}
 \right|^2
 =
 \frac{j}{(j+1)(2j+1)}
 \asymp\frac1{2j+1}.}
\label{eq:V48-trivalent-sharp}
\end{equation}
Consequently no uniform estimate by
\(C[(2j+1)^{-1}]^2\) is possible.
\end{proposition}

\begin{proof}
The standard \(3j\)-identity
\[
 \begin{pmatrix}
  j&j&1\\ m&-m&0
 \end{pmatrix}
 =
 (-1)^{j-m}
 \frac{m}{\sqrt{j(j+1)(2j+1)}}
\]
with \(m=j\) gives the equality.  Multiplication by
\((2j+1)^2\) makes the displayed quantity grow linearly in \(j\), so
the proposed squared-capacity bound fails.
\end{proof}

\begin{firewall}[Two cuts at one coordinate are not two capacities]
\label{fire:V48-one-coordinate-one-capacity}
An incidence-three coordinate may support two adjacent merge edges,
but its protected invariant projector supplies only the one capacity
power in \eqref{eq:V48-trivalent-upper}.  Multiplying the two
single-cut bounds would contradict
\Cref{prop:V48-trivalent-sharp}.  The same coordinate must instead be
treated as one multi-edge local block.
\end{firewall}

\section{Normalized local marks}
\label{sec:V48-local-marks}

For an active row \(i\) at coordinate \(e\), put
\[
 x_i(e):=\frac{a_{i,e}}{\Xi_i}\in[0,1].
\]
The local contribution of \(e\) to the normalized overlap of rows
\(i,j\) is
\begin{equation}
 \vartheta_{ij}(e):=x_i(e)x_j(e)
 =
 \frac{a_{i,e}a_{j,e}}{\Xi_i\Xi_j}.
\label{eq:V48-local-mark}
\end{equation}
It obeys
\[
 0\le\vartheta_{ij}(e)\le\theta_{ij}\le1.
\]

\begin{lemma}[Same-coordinate path compression]
\label{lem:V48-path-compression}
Let \(v_0,v_1,\ldots,v_r\) be distinct rows active at the same
coordinate \(e\), with \(1\le r\le3\).  Then
\begin{equation}
\boxed{
 \prod_{\ell=0}^{r-1}
 \vartheta_{v_\ell v_{\ell+1}}(e)
 =
 \vartheta_{v_0v_r}(e)
 \prod_{\ell=1}^{r-1}x_{v_\ell}(e)^2
 \le
 \vartheta_{v_0v_r}(e).}
\label{eq:V48-path-compression}
\end{equation}
\end{lemma}

\begin{proof}
Expand every mark using
\eqref{eq:V48-local-mark}.  Each endpoint factor occurs once and each
internal row factor twice.  Since every \(x_i(e)\le1\), discard the
internal squares.
\end{proof}

\begin{corollary}[Mass ratios compress at the same time]
\label{cor:V48-ratio-compression}
For the same path,
\begin{equation}
 \prod_{\ell=0}^{r-1}
 \left[
 \vartheta_{v_\ell v_{\ell+1}}(e)
 \frac{\Xi_{v_{\ell+1}}}{\Xi_{v_\ell}}
 \right]
 \le
 \vartheta_{v_0v_r}(e)
 \frac{\Xi_{v_r}}{\Xi_{v_0}}.
\label{eq:V48-ratio-compression}
\end{equation}
\end{corollary}

\begin{proof}
The mass ratios telescope, and
\Cref{lem:V48-path-compression} bounds the mark product.
\end{proof}

\begin{assessment}[Geometric meaning]
\label{ass:V48-compression-meaning}
When one higher-incidence coordinate carries consecutive transport
edges, the correct replacement is not a product of cut capacities.
It is one shortcut edge between the endpoints, one local protected
capacity, and the same telescoped endpoint mass ratio.  This matches
the sharp one-power result of
\Cref{prop:V48-trivalent-sharp}.
\end{assessment}

\section{Compression of a transport coordinate word}
\label{sec:V48-word}

Let a simple transport chain have edge-coordinate word
\[
 (e_0,e_1,\ldots,e_{m-1}),
 \qquad m\le3,
\]
where \(e_\ell\) supports
\(i_\ell i_{\ell+1}\).

\begin{definition}[Run compression]
\label{def:V48-run-compression}
Replace every maximal consecutive run
\[
 e_a=e_{a+1}=\cdots=e_{b-1}=e
\]
by the single shortcut edge \(i_ai_b\) carrying the local mark
\(\vartheta_{i_ai_b}(e)\).  The resulting shorter chain and coordinate
word are called the run compression.
\end{definition}

\begin{lemma}[Run compression dominates the original word]
\label{lem:V48-run-word}
The product of local normalized marks and row-mass ratios of the
original word is at most the corresponding product for its run
compression.
\end{lemma}

\begin{proof}
Apply \Cref{cor:V48-ratio-compression} independently on the disjoint
maximal runs.
\end{proof}

\begin{lemma}[Only one repeated-coordinate word survives]
\label{lem:V48-only-ABA}
After run compression, either all coordinate letters are distinct, or
the chain has length three and its word is
\begin{equation}
 (e,f,e),\qquad e\ne f.
\label{eq:V48-ABA}
\end{equation}
In the latter case \(e\) is an incidence-four coordinate.
\end{lemma}

\begin{proof}
A reduced word has no equal adjacent letters and has length at most
three.  The only possible repetition is therefore the first letter in
the third position, giving \eqref{eq:V48-ABA}.  That coordinate
supports the disjoint chain edges \(i_0i_1\) and \(i_2i_3\), so all four
rows are active.
\end{proof}

\begin{theorem}[All adjacent protected overlaps reduce to distinct
coordinate factors]
\label{thm:V48-adjacent-overlap-closed}
In a protected transport resolution of the connected-graph carrier,
compress every same-coordinate consecutive run.  If the reduced word
does not have the crossed-repeat form \eqref{eq:V48-ABA}, its
operator-level protected factors act on distinct coordinate tensor
factors and its scalar mark-ratio product is bounded by the compressed
chain.  Hence the protected terminal-moment argument of V46--V47
applies.
\end{theorem}

\begin{proof}
By \Cref{lem:V48-run-word}, compression only enlarges the normalized
mark-ratio majorant.  Each run is kept as one local multi-edge block
and is assigned only the single protected capacity permitted by
\Cref{fire:V48-one-coordinate-one-capacity}.  By
\Cref{lem:V48-only-ABA}, absence of \eqref{eq:V48-ABA} makes the
remaining coordinate letters pairwise distinct.  The exact
connected-graph identity of V47 has already assembled their merge
differences before norms.  Product-state capacity may therefore be
tensored across these distinct coordinate factors, followed by the
scalar telescope \eqref{eq:V46-ratio-telescope} and the surplus-edge
reduction \eqref{eq:V45-spanning-tree-domination}.
\end{proof}

\section{The crossed incidence-four repeat}
\label{sec:V48-crossed-repeat}

In the remaining word \((e,f,e)\), write the chain as
\[
 0\mathbin{-}1\mathbin{-}2\mathbin{-}3.
\]
The two \(e\)-marks obey a second elementary compression:
\begin{equation}
 \vartheta_{01}(e)\vartheta_{23}(e)
 =
 x_0(e)x_1(e)x_2(e)x_3(e)
 \le
 \vartheta_{03}(e),
\label{eq:V48-crossed-mark-compression}
\end{equation}
because \(x_1(e)x_2(e)\le1\).  Unlike
\eqref{eq:V48-path-compression}, the repeated edges are disjoint.
Their marks must therefore be compressed only after all three chain
mass ratios have been multiplied.

\begin{proposition}[Exact crossed-repeat compression]
\label{prop:V48-crossed-loss}
After replacing the two \(e\)-marks by the shortcut mark
\(\vartheta_{03}(e)\), the complete mark-ratio product of the word
\((e,f,e)\) is bounded by
\begin{equation}
 \vartheta_{03}(e)\vartheta_{12}(f)
 \frac{\Xi_3}{\Xi_0}.
\label{eq:V48-crossed-loss}
\end{equation}
Thus the \(e\)-block carries the endpoint transport factor, while the
middle \(f\)-edge remains only as a surplus normalized mark.
\end{proposition}

\begin{proof}
The three mass ratios telescope:
\[
 \frac{\Xi_1}{\Xi_0}
 \frac{\Xi_2}{\Xi_1}
 \frac{\Xi_3}{\Xi_2}
 =
 \frac{\Xi_3}{\Xi_0}
\].
Multiply this identity by the three original local marks and use
\eqref{eq:V48-crossed-mark-compression}.
\end{proof}

\begin{corollary}[The crossed repeat has no scalar loss]
\label{cor:V48-crossed-scalar}
The crossed-repeat scalar ledger is bounded by one endpoint transport
factor \(\vartheta_{03}(e)\Xi_3/\Xi_0\), times the surplus mark
\(\vartheta_{12}(f)\le1\).
\end{corollary}

\begin{proof}
This is \eqref{eq:V48-crossed-loss}.
\end{proof}

\begin{firewall}[The local incidence-four capacity is still single]
\label{fire:V48-crossed-capacity}
\Cref{cor:V48-crossed-scalar} proves that there is no scalar
middle-ratio loss.  It does not prove that the incidence-four
\(e\)-projector supplies two
capacity factors.  The shortcut uses one \(e\)-capacity and one
normalized \(03\)-mark.  A direct simultaneous normal form is still
needed to verify that this one-capacity shortcut is retained by the
resolved paw/cycle graph difference.
\end{firewall}

\section{The narrowed local gate}
\label{sec:V48-narrowed}

\begin{assessment}[Post-V48 protected status]
\label{ass:V48-status}
\begin{enumerate}
\item Two independent cut capacities at one trivalent coordinate are
      rigorously impossible.
\item Every consecutive same-coordinate transport run compresses to
      one endpoint edge and one mass ratio.
\item After compression, all protected coordinate factors are
      distinct except the unique incidence-four word \((e,f,e)\).
\item That word also has a favorable scalar compression; only the
      operator statement that its \(e\)-block retains the shortcut
      capacity inside a paw/cycle difference remains.
\end{enumerate}
The separate global-gapped optimization from V47 is unchanged.
\end{assessment}

\begin{openproblem}[V49 mandate: incidence-four shortcut normal form]
\label{op:V49-mandate}
\begin{enumerate}
\item In the word \((e,f,e)\), write the full incidence-four
      invariant projector in an \(03\mid12\) coupled basis.
\item Keep the complete multiplicity identity and prove one
      \(03\)-shortcut capacity, not a product of the \(01\) and \(23\)
      cut capacities.
\item Verify that the paw and four-cycle merge differences preserve
      this shortcut after their signed partition terms are assembled.
\item Combine the result with
      \Cref{cor:V48-crossed-scalar} and the terminal moment telescope.
\item Then perform the one-stock optimization of the gapped
      complement using \eqref{eq:V47-one-global-gap}.
\end{enumerate}
\end{openproblem}

\begin{firewall}[Claim boundary after compact V48]
\label{fire:V48-current-boundary}
V48 proves the sharp one-power trivalent capacity and closes all
adjacent same-coordinate protected overlaps by shortcut compression.
It does not yet prove the incidence-four crossed-repeat shortcut inside
the signed graph difference or close the global-gapped scalar branch.
Thus the complete all-weak estimate, global quartic MTA, all-order MTA,
WUFC/CPM, continuum Yang--Mills existence, and a positive mass gap
remain open.
\end{firewall}

\section{Append-only record for compact V48}
\label{sec:V48-record}

V48 was appended to the complete compact V47 file.  The exact V47
parent had \(18995\) newline-delimited source records, \(631510\)
bytes, and SHA--256 digest
\[
\texttt{fce2ad491f805923b38155df557fdfd8114ae5baa99e48aed6a4ff9715dc1d79}.
\]
No compact V47 mathematical content was changed.  The sole final
\verb|\end{document}| is restored below.

% =============================================================================
% END APPENDED COMPACT VERSION V48 MATERIAL
% =============================================================================

% The compact V48 document terminator is intentionally disabled here:
% \end{document}

% =============================================================================
% BEGIN APPENDED COMPACT VERSION V49 MATERIAL
% =============================================================================

\chapter{V49: The incidence-four recoupling cut and the valid
one-row shortcut capacity}
\label{ch:V49}

\section{Exact capacity across a two-pair recoupling cut}
\label{sec:V49-two-pair}

Relabel the crossed-repeat rows as \(0,1,2,3\) and group the
incidence-four coordinate as
\[
 \mathcal A:=V_{j_0}\otimes V_{j_3},
 \qquad
 \mathcal B:=V_{j_1}\otimes V_{j_2}.
\]
The multiplicity-free decompositions are
\[
 \mathcal A=\bigoplus_{k\in\mathcal K_{03}}V_k,
 \qquad
 \mathcal B=\bigoplus_{\ell\in\mathcal K_{12}}V_\ell.
\]
Let
\(\mathcal K=\mathcal K_{03}\cap\mathcal K_{12}\).  As in V34,
\[
 P_{\boldsymbol j}^{\rm inv}
 =
 \sum_{k\in\mathcal K}|\Omega_k\rangle\langle\Omega_k|,
\]
where \(\Omega_k\) is the normalized singlet in
\(V_k\otimes V_k\).

\begin{theorem}[Exact composite-cut capacity]
\label{thm:V49-composite-capacity}
The product-state capacity of the full invariant projector across the
composite cut \(03\mid12\) is
\begin{equation}
\boxed{
 \sup_{\substack{\|x\|_{\mathcal A}=1\\
                 \|y\|_{\mathcal B}=1}}
 \langle x\otimes y,
 P_{\boldsymbol j}^{\rm inv}(x\otimes y)\rangle
 =
 \max_{k\in\mathcal K}\frac1{2k+1}.}
\label{eq:V49-composite-capacity}
\end{equation}
In particular this capacity equals one whenever \(0\in\mathcal K\).
\end{theorem}

\begin{proof}
Write \(x=\bigoplus_kx_k\) and \(y=\bigoplus_ky_k\) in the two
Clebsch--Gordan decompositions.  Since \(\Omega_k\) is maximally
entangled on \(V_k\otimes V_k\),
\[
 |\langle\Omega_k,x_k\otimes y_k\rangle|^2
 \le\frac1{2k+1}\|x_k\|^2\|y_k\|^2.
\]
Orthogonality of the recoupling channels gives
\[
 \langle x\otimes y,
 P_{\boldsymbol j}^{\rm inv}(x\otimes y)\rangle
 \le
 \left(\max_{k\in\mathcal K}\frac1{2k+1}\right)
 \sum_k\|x_k\|^2\|y_k\|^2
 \le
 \max_{k\in\mathcal K}\frac1{2k+1}.
\]
For equality, choose a channel \(k\) attaining the maximum and take
\(x_k,y_k\) to be matching Schmidt basis vectors of \(\Omega_k\), with
all other components zero.
\end{proof}

\begin{corollary}[The proposed \(03\mid12\) shortcut contraction is
false]
\label{cor:V49-composite-shortcut-false}
If \(j_0=j_3\) and \(j_1=j_2\), then \(0\in\mathcal K\), and the
protected projector has composite-cut capacity one.  Hence no
spin-uniform strict contraction across \(03\mid12\) can justify the
crossed-repeat shortcut.
\end{corollary}

\begin{proof}
Equal spins couple to spin zero on each side.  Apply
\Cref{thm:V49-composite-capacity}.
\end{proof}

\begin{assessment}[Correction to the V48 acquisition route]
\label{ass:V49-route-correction}
The scalar mark compression
\eqref{eq:V48-crossed-loss} remains exact.  What fails is the proposed
analytic interpretation of its \(03\) shortcut as a uniformly
contracting two-pair cut.  The \(k=0\) channel is a product of a
\(03\)-singlet and a \(12\)-singlet across that composite cut.  V49
therefore returns to capacities against individual row factors, where
Schur's lemma gives a genuine dimension loss.
\end{assessment}

\section{Protected partition terms retain one-row capacity}
\label{sec:V49-one-row}

Fix one incidence-four coordinate and one global row partition
\(\pi\).  Restrict every block moment in \(T_e(\pi)\) to trivial heat
output.  A nonzero protected block is its complete invariant
projector; a singleton block containing a nontrivial active
representation has zero protected part.

\begin{lemma}[One-row capacity in every protected partition state]
\label{lem:V49-state-one-row}
Let \(T_e^{\rm prot}(\pi)\) be a nonzero protected partition-state
operator.  For every active row \(i\),
\begin{equation}
 \kappa_{\otimes,i}
 \bigl(T_e^{\rm prot}(\pi)\bigr)
 \le\frac1{d_i},
\label{eq:V49-state-one-row}
\end{equation}
where the capacity tests a product between row \(i\) and all other row
factors.  The estimate includes the full invariant multiplicity of
every block.
\end{lemma}

\begin{proof}
The partition block \(A\in\pi\) containing \(i\) cannot be a protected
singleton because \(R_{i,e}\) is nontrivial.  On that block, group row
\(i\) against \(A\setminus\{i\}\).  The Schur normal form from
\Cref{lem:V34-one-row-normal-form}, applied to this smaller invariant
tensor product, gives capacity \(d_i^{-1}\) and an identity on its full
multiplicity space.  Protected projectors from all other partition
blocks are contractions acting away from row \(i\).  One-sided tensor
submultiplicativity proves \eqref{eq:V49-state-one-row}.
\end{proof}

\begin{proposition}[Signed graph differences preserve one-row
capacity]
\label{prop:V49-signed-one-row}
Let \(G\) be any star, path, four-cycle, or paw in
\eqref{eq:V47-connected-graph-identity}.  The protected part of the
single-coordinate signed difference
\(\mathfrak D_GT_e\) satisfies
\begin{equation}
\boxed{
 \kappa_{\otimes,i}
 \bigl((\mathfrak D_GT_e)^{\rm prot}\bigr)
 \le\frac{C_G}{d_i}}
\label{eq:V49-signed-one-row}
\end{equation}
for every active row \(i\), with \(C_G\le2^{|E(G)|}\le16\).
\end{proposition}

\begin{proof}
By \eqref{eq:V47-forest-expansion},
\(\mathfrak D_GT_e\) is a signed sum of at most
\(2^{|E(G)|}\) protected partition-state operators.  Apply capacity
subadditivity only after this exact signed graph difference has been
formed, and use \Cref{lem:V49-state-one-row} on every nonzero term.
\end{proof}

\begin{corollary}[The crossed \(e\)-block has exactly one valid
capacity]
\label{cor:V49-crossed-one-capacity}
In the crossed word \((e,f,e)\), the incidence-four \(e\)-block inside
a resolved paw or four-cycle difference retains one one-row capacity
factor \(C/d_i\) for any chosen endpoint \(i\).  It does not retain two
independent cut capacities.
\end{corollary}

\begin{proof}
Use \Cref{prop:V49-signed-one-row} for the positive assertion and
\Cref{cor:V49-composite-shortcut-false} for the forbidden stronger
interpretation.
\end{proof}

\section{Comparison with the local Casimir mass}
\label{sec:V49-Casimir}

\begin{lemma}[Dimension is one square root of Casimir mass]
\label{lem:V49-dimension-Casimir}
For \(a_i=1+j_i(j_i+1)\) and \(d_i=2j_i+1\),
\begin{equation}
 c\,a_i^{1/2}\le d_i\le C\,a_i^{1/2}.
\label{eq:V49-dimension-Casimir}
\end{equation}
Hence the valid local protected gain in
\eqref{eq:V49-signed-one-row} is of order \(a_i^{-1/2}\), not
\(a_i^{-1}\) or \(a_i^{-2}\).
\end{lemma}

\begin{proof}
Both \(d_i^2=4j_i(j_i+1)+1\) and \(a_i=1+j_i(j_i+1)\) are bounded
above and below by fixed positive multiples of one another.
\end{proof}

\begin{proposition}[A single-coordinate estimate cannot pay the global
row mass]
\label{prop:V49-local-not-global}
There is no uniform implication from
\eqref{eq:V49-signed-one-row} to a negative power of the full row mass
\(\Xi_i\).  In particular, keeping \(j_i\) fixed at the selected
coordinate while adding arbitrarily many other shared coordinates to
row \(i\) leaves \(d_i^{-1}\) fixed and makes \(\Xi_i\) arbitrarily
large.
\end{proposition}

\begin{proof}
The one-coordinate representation dimension depends only on
\(a_{i,e}\).  The global mass is the sum of \(a_{i,f}\) over every
coordinate active in row \(i\).  Adding coordinates \(f\ne e\) changes
the latter sum without changing \(d_{i,e}\).  Thus no bound
\(d_{i,e}^{-1}\le C\Xi_i^{-\varepsilon}\), \(\varepsilon>0\), follows
from the local normal form alone.
\end{proof}

\begin{assessment}[The correct incidence-four use]
\label{ass:V49-correct-use}
The crossed incidence-four coordinate supplies:
\begin{enumerate}
\item the exact scalar shortcut marks of V48;
\item one row-factor capacity \(d_{i,e}^{-1}\) retained through the
      complete signed graph difference;
\item no uniform \(03\mid12\) contraction and no second capacity
      power.
\end{enumerate}
Therefore it must be returned to the common row-rooted stock and
combined with all other coordinates before a global mass denominator
is claimed.  A selected-coordinate shortcut by itself would be
circular.
\end{assessment}

\section{A stock-level one-row product}
\label{sec:V49-stock-product}

Let \(E_i^{\rm prot}\) be any protected incidence-three/four bank in
which row \(i\) is active at every coordinate.  Put
\[
 D_{i,E}:=\prod_{e\in E_i^{\rm prot}}d_{i,e}.
\]

\begin{lemma}[Tensorized one-row capacity survives signed local
differences]
\label{lem:V49-tensor-one-row}
If the complete graph difference is formed first and its protected
coordinates in \(E_i^{\rm prot}\) are then tensored, their one-row
capacity is bounded by
\begin{equation}
 \frac{C_G}{D_{i,E}}.
\label{eq:V49-tensor-one-row}
\end{equation}
No recoupling multiplicity factor occurs.
\end{lemma}

\begin{proof}
At each coordinate use the Schur normal form with its full
multiplicity identity, as in
\Cref{prop:V49-signed-one-row}.  The tested vector is a product in the
coordinate tensor factors of row \(i\) and its complement, so
one-sided capacity submultiplicativity multiplies the factors
\(d_{i,e}^{-1}\).  The signed graph coefficient sum is taken once and
is bounded by the fixed \(C_G\), rather than once per coordinate.
\end{proof}

\begin{corollary}[Support-uniform exponential capacity]
\label{cor:V49-stock-exponential}
Since every active \(d_{i,e}\ge2\),
\begin{equation}
 \kappa_{\otimes,i}
 \le C_G2^{-|E_i^{\rm prot}|}.
\label{eq:V49-stock-exponential}
\end{equation}
\end{corollary}

\begin{proof}
Apply \eqref{eq:V49-tensor-one-row}.
\end{proof}

\begin{firewall}[The remaining conversion is weighted]
\label{fire:V49-weighted-conversion}
\eqref{eq:V49-stock-exponential} pays coordinate count.  The tree
requires powers of the weighted sum \(\Xi_i\).  As in the V45 audit,
count decay cannot be called a weighted mass moment until it is
combined with
\(S_{ij}^2\le36N_{ij}H_{ij}^{(i)}\) and the all-weak marks in the same
graph-difference term.
\end{firewall}

\begin{openproblem}[V50 mandate: audit and weighted shortcut stock]
\label{op:V50-mandate}
\begin{enumerate}
\item Perform the mandatory read-only audit of the newest SM and NS
      notes.
\item In the crossed-repeat branch, keep the one-row product
      \eqref{eq:V49-tensor-one-row} over the entire row-rooted bank,
      not just the selected incidence-four coordinate.
\item Combine it with the exact shortcut marks
      \eqref{eq:V48-crossed-loss} and the row-power inequality
      \eqref{eq:V44-row-power}.
\item Treat the gapped complement with the single global factor
      \eqref{eq:V47-one-global-gap}.
\item Determine whether these two stock-level estimates close the
      final weighted scalar optimization or leave a sharp
      concentration counterexample.
\end{enumerate}
\end{openproblem}

\begin{firewall}[Claim boundary after compact V49]
\label{fire:V49-current-boundary}
V49 computes the exact two-pair recoupling-cut capacity, disproves a
uniform composite-cut shortcut, and proves that one valid row-factor
capacity survives the complete signed incidence-four graph
difference.  The remaining task is the weighted common-stock
conversion, together with the globally gapped complement.  Thus the
complete all-weak estimate, global quartic MTA, all-order MTA,
WUFC/CPM, continuum Yang--Mills existence, and a positive mass gap
remain open.
\end{firewall}

\section{Append-only record for compact V49}
\label{sec:V49-record}

V49 was appended to the complete compact V48 file.  The exact V48
parent had \(19374\) newline-delimited source records, \(643955\)
bytes, and SHA--256 digest
\[
\texttt{91720b37083b85fdb60603f7f078d2f7c1e9d229250f673a920dfc839a05e162}.
\]
No compact V48 mathematical content was changed.  The sole final
\verb|\end{document}| is restored below.

% =============================================================================
% END APPENDED COMPACT VERSION V49 MATERIAL
% =============================================================================

% The compact V49 document terminator is intentionally disabled here:
% \end{document}

% =============================================================================
% BEGIN APPENDED COMPACT VERSION V50 MATERIAL
% =============================================================================

\chapter{V50: Companion audit and closure of the crossed protected
stock}
\label{ch:V50}

\section{Mandatory SM/NS companion audit}
\label{sec:V50-audit}

The newest SM companion inspected before this argument was
\[
\begin{array}{ll}
\text{file:}&
\texttt{Renewal\_SM\_Einstein\_Continuum\_Research\_Notes\_V54.tex},\\
\text{records:}&20300,\\
\text{bytes:}&734041,\\
\text{SHA--256:}&
\texttt{f238e35d9d422cae67d1100e213bcd9cc1bd1ac082c80937290bde0f3dadf1f8}.
\end{array}
\]
SM V51 proves fixed-background Lorentzian state compactness only in a
topology containing the composite stress insertions and uniform
microlocal bounds.  V52--V53 identify a genuinely new
background-comparison modulus and generate it from relative Cauchy
response.  V54 constructs one common physical algebra chart from an
absolutely summable ordered product of shell derivations.  It proves
the product by one noncommutative telescope and shows that shellwise
decay without summability may fail through a harmonic-series phase.
BRST descent is imposed before the infinite product.

The NS companion remained
\[
\begin{array}{ll}
\text{file:}&
\texttt{Renewal\_NS\_Stability\_Research\_Notes\_V31.tex},\\
\text{records:}&23052,\\
\text{bytes:}&887537,\\
\text{SHA--256:}&
\texttt{f1121b62238ad5491bb7cf03635ea9934942edf573ed8faaa9ee79e1d38a6696}.
\end{array}
\]
Its exact formation ladder and packet regressions continue to say that
a weak norm does not pay the weighted endpoint moment, and that the
one-use persistence ledger may not be charged again on recent-entry
intervals.

\begin{assessment}[V50 audit transfer]
\label{ass:V50-audit-transfer}
No companion estimate is imported.  The proof-order transfer is:
\begin{enumerate}
\item disjointize or combine overlapping transport banks before
      multiplying analytic factors;
\item use one exact graph-difference telescope, as in V47, rather than
      a product of marginal cancellations;
\item retain the weighted row stock explicitly; pointwise
      representation decay is not a substitute for its sum.
\end{enumerate}
\end{assessment}

\section{The only residual protected coordinate word}
\label{sec:V50-residual-word}

After the run compression of V48, the only repeated-coordinate
transport word is
\[
 (e,f,e)
\]
on the simple chain
\(0-1-2-3\).  Let \(B_{01}\) and \(B_{23}\) be the two row-rooted
balanced banks for the exterior transports, after every adjacent-bank
overlap has already been compressed.  Put
\begin{align}
 I&:=B_{01}\cap B_{23},
\label{eq:V50-intersection}\\
 S_{01}&:=\sum_{g\in B_{01}}a_{0,g},
 &
 S_{23}&:=\sum_{g\in B_{23}}a_{2,g},
\label{eq:V50-exterior-stocks}\\
 S_{0I}&:=\sum_{g\in I}a_{0,g},
 &
 S_{2I}&:=\sum_{g\in I}a_{2,g}.
\label{eq:V50-intersection-stocks}
\end{align}
The transport selection gives
\begin{equation}
 S_{01}>\frac1{12}\Xi_0,
 \qquad
 S_{23}>\frac1{12}\Xi_2.
\label{eq:V50-exterior-large}
\end{equation}

\begin{lemma}[Light intersection permits disjointization]
\label{lem:V50-light-disjoint}
If
\[
 S_{0I}\le\frac12S_{01}
\quad\hbox{or}\quad
 S_{2I}\le\frac12S_{23},
\]
then the two exterior banks can be replaced by disjoint banks while
retaining at least half of the selected row stock in each bank.
\end{lemma}

\begin{proof}
If \(S_{0I}\le S_{01}/2\), replace \(B_{01}\) by
\(B_{01}\setminus I\) and leave \(B_{23}\) unchanged.  The new banks
are disjoint, the first retains at least \(S_{01}/2\), and the second
retains all its stock.  The other alternative is symmetric.
\end{proof}

\begin{corollary}[Light crossed overlap is closed]
\label{cor:V50-light-closed}
Every protected crossed-repeat branch satisfying
\Cref{lem:V50-light-disjoint} obeys the global marked-tree atlas.
\end{corollary}

\begin{proof}
The factor-two stock loss changes only the constant in
\eqref{eq:V46-protected-factor}.  The exact connected-graph difference
has already been formed, and the remaining protected banks are
coordinate-disjoint.  Apply
\Cref{thm:V47-disjoint-protected-authorized}, followed by the scalar
compression in \Cref{prop:V48-crossed-loss}.
\end{proof}

\section{A heavy intersection creates a legal shortcut bank}
\label{sec:V50-heavy-shortcut}

It remains to assume
\begin{equation}
 S_{0I}>\frac12S_{01},
 \qquad
 S_{2I}>\frac12S_{23}.
\label{eq:V50-double-heavy}
\end{equation}

\begin{lemma}[Maximum-partner rigidity on the intersection]
\label{lem:V50-max-rigidity}
At every \(g\in I\),
\begin{equation}
 a_{1,g}=a_{3,g}
 \ge\max\{a_{0,g},a_{2,g}\}.
\label{eq:V50-max-rigidity}
\end{equation}
Consequently row \(0\) may reassign the tied maximal partner from row
\(1\) to row \(3\) at every coordinate of \(I\).
\end{lemma}

\begin{proof}
Membership in \(B_{01}\) means that row \(1\) was chosen maximal among
the active partners of row \(0\).  Hence
\[
 a_{1,g}\ge a_{2,g},a_{3,g}.
\]
Membership in \(B_{23}\) means that row \(3\) was chosen maximal among
the active partners of row \(2\).  Hence
\[
 a_{3,g}\ge a_{0,g},a_{1,g}.
\]
Thus \(a_{1,g}=a_{3,g}\), and both dominate the other two entries.
The partner rule allows either member of a tie.
\end{proof}

Define the shortcut stock
\begin{equation}
 S_{03}^{I}:=\sum_{g\in I}a_{0,g}=S_{0I},
 \qquad
 H_{03}^{I}:=\sum_{g\in I}a_{0,g}a_{3,g},
 \qquad
 N_I:=|I|.
\label{eq:V50-shortcut-stock}
\end{equation}

\begin{lemma}[Shortcut stock power and size]
\label{lem:V50-shortcut-power}
Under \eqref{eq:V50-double-heavy},
\begin{align}
 S_{03}^{I}&>\frac1{24}\Xi_0,
\label{eq:V50-shortcut-large}\\
 (S_{03}^{I})^2&\le N_IH_{03}^{I}.
\label{eq:V50-shortcut-power}
\end{align}
\end{lemma}

\begin{proof}
The first line follows from
\eqref{eq:V50-double-heavy} and
\eqref{eq:V50-exterior-large}.  By
\eqref{eq:V50-max-rigidity},
\(a_{3,g}\ge a_{0,g}\), so
\[
 a_{0,g}^2\le a_{0,g}a_{3,g}.
\]
Cauchy--Schwarz over \(I\) proves the second line.
\end{proof}

\begin{lemma}[Middle mass ratio is uniformly bounded]
\label{lem:V50-middle-ratio}
Under \eqref{eq:V50-double-heavy},
\begin{equation}
\boxed{\Xi_2<24\Xi_1.}
\label{eq:V50-middle-ratio}
\end{equation}
\end{lemma}

\begin{proof}
By \eqref{eq:V50-max-rigidity},
\[
 S_{2I}
 =\sum_{g\in I}a_{2,g}
 \le\sum_{g\in I}a_{1,g}
 \le\Xi_1.
\]
On the other hand,
\eqref{eq:V50-double-heavy} and
\eqref{eq:V50-exterior-large} give
\(S_{2I}>\Xi_2/24\).
\end{proof}

\section{The protected shortcut factor}
\label{sec:V50-shortcut-factor}

Let
\[
 \vartheta_{03}^{I}
 :=
 \frac{H_{03}^{I}}{\Xi_0\Xi_3}
 \le\theta_{03}.
\]

\begin{proposition}[One stock-level shortcut capacity]
\label{prop:V50-shortcut-capacity}
Inside the complete protected paw/four-cycle graph difference, the
heavy intersection bank \(I\) obeys
\begin{equation}
\boxed{
 \kappa_{\otimes,0}(\mathcal P_I)
 \le
 \frac{C}{N_I}
 \vartheta_{03}^{I}
 \frac{\Xi_3}{\Xi_0}.}
\label{eq:V50-shortcut-capacity}
\end{equation}
Only one row-factor capacity is used at each incidence-four
coordinate.
\end{proposition}

\begin{proof}
By \Cref{lem:V50-max-rigidity}, the \(03\) pairing is a legal
row-rooted partner assignment on \(I\).  The signed one-row normal form
of \Cref{prop:V49-signed-one-row} and its tensorization
\Cref{lem:V49-tensor-one-row} supply the protected capacity input to
the stock compiler.  Repeating the protected line of
\Cref{prop:V45-separated-cut-ledgers}, now with
\eqref{eq:V50-shortcut-power}, gives
\[
 N_I(S_{03}^{I})^2
 \kappa_{\otimes,0}(\mathcal P_I)
 \le C H_{03}^{I}.
\]
Use \eqref{eq:V50-shortcut-large}:
\[
 \kappa_{\otimes,0}(\mathcal P_I)
 \le
 \frac{C H_{03}^{I}}{N_I\Xi_0^2}
 =
 \frac{C}{N_I}
 \vartheta_{03}^{I}\frac{\Xi_3}{\Xi_0}.
\]
The full multiplicity space stays inside the one-row projector normal
form, so no recoupling-channel count is inserted.
\end{proof}

\begin{theorem}[Closure of the heavy crossed protected branch]
\label{thm:V50-heavy-crossed-closed}
Assume the crossed word \((e,f,e)\) lies in the double-heavy branch
\eqref{eq:V50-double-heavy} and its terminal row \(3\) is private- or
unbalanced-controlled.  Then its complete protected contribution,
including at most one labelled balanced-output payer and total degree
at most three, obeys the global marked-tree atlas.
\end{theorem}

\begin{proof}
Use \eqref{eq:V50-shortcut-capacity} as the single transport factor
from row \(0\) to row \(3\).  The middle \(12\)-edge remains in the
same paw or four-cycle graph difference as a surplus normalized mark;
it is not assigned a second mass ratio.  This is precisely the scalar
compression proved in
\Cref{prop:V48-crossed-loss,cor:V48-crossed-scalar}.

Apply the terminal moment certificate of
\Cref{lem:V46-terminal-types} at row \(3\), with degree equal to the
one or two unpaid tree powers plus the optional output payer.
The one-step case of
\Cref{lem:V46-ratio-telescope} closes the scalar coefficient.
Finally discard surplus graph edges using
\Cref{cor:V47-surplus-edge}.
\end{proof}

\begin{theorem}[All protected transport-bank overlaps are closed]
\label{thm:V50-all-protected-overlaps}
After the exact connected-graph assembly of V47, every protected
sparse transport chain terminating in a private- or
unbalanced-controlled row obeys the global marked-tree atlas.  This
includes all incidence-three and incidence-four coordinate overlaps.
\end{theorem}

\begin{proof}
Compress adjacent repeated-coordinate runs by
\Cref{thm:V48-adjacent-overlap-closed}.  The only residual word is
\((e,f,e)\) by \Cref{lem:V48-only-ABA}.  Its light-intersection branch
is \Cref{cor:V50-light-closed}; its double-heavy branch is
\Cref{thm:V50-heavy-crossed-closed}.  These alternatives are
exhaustive.
\end{proof}

\section{The remaining gapped degree ledger}
\label{sec:V50-gapped-ledger}

The protected branch is now complete, but the global-gapped factor
\eqref{eq:V47-one-global-gap} has a different scalar cost.  If a
transport bank is estimated through
\eqref{eq:V46-gapped-factor}, its marked entry contains
\(s_\alpha^{-1}\) where the protected entry contains \(N^{-1}\).

\begin{lemma}[Finite terminal-degree criterion for gapped steps]
\label{lem:V50-gapped-degree}
Let \(k\le3\) be the number of tree-excess and labelled-output powers
at the starting row, and let \(g\) be the number of transport steps
whose mark is obtained from the gapped ledger.  If a terminal moment
certificate is available through order \(k+g\), then the same
mass-ratio telescope closes the scalar chain.
\end{lemma}

\begin{proof}
Relative to the protected product, each gapped marked factor replaces
\(N_\ell^{-1}\) by \(s_\alpha^{-1}\).  Move these \(g\) factors to the
terminal debit.  The calculation
\eqref{eq:V46-ratio-telescope} then uses order \(k+g\) instead of
order \(k\).  All normalized marks and row ratios telescope in the
same way.
\end{proof}

\begin{firewall}[Only cubic terminal moments have been established]
\label{fire:V50-cubic-only}
The current private and unbalanced terminal certificates were proved
only through order three.  Therefore
\Cref{lem:V50-gapped-degree} closes only configurations with
\(k+g\le3\).  Invoking higher moments because the chain length is
finite would require new one-link debit estimates; they are plausible
for heat multipliers but have not been assumed in the compact
programme.
\end{firewall}

\begin{assessment}[Post-V50 obstruction]
\label{ass:V50-obstruction}
All protected overlap geometries are closed, including the crossed
incidence-four word.  The sole remaining part of the global all-weak
quartic carrier is the gapped transport degree \(k+g>3\), together
with any branch whose terminal is not yet private/unbalanced but ends
at a dense gapped bank.  This is now a finite scalar analytic problem,
not a partition, multiplicity, or coordinate-overlap problem.
\end{assessment}

\begin{openproblem}[V51 mandate: higher fixed moments or direct global
gap optimization]
\label{op:V51-mandate}
\begin{enumerate}
\item Return to the exact one-link Wilson multiplier and determine
      whether the debit
      \((s_\alpha A)^rq_{\alpha,A}\le C_r\) is available for
      \(4\le r\le6\).
\item If so, extend the Bernoulli product compiler of V7 and close all
      finite gapped degrees by
      \Cref{lem:V50-gapped-degree}.
\item If only cubic information is legitimate, optimize the single
      global exponential
      \(e^{-cs_\alpha N_{\rm gap}}\) before extracting marked
      \(s_\alpha^{-1}\) factors.
\item Test the dense one-coordinate/high-spin sector separately:
      coordinate-count decay alone must not be used as row-mass decay.
\item On success, insert the resulting all-weak estimate into the V36
      finite graph compiler and state exactly whether global quartic
      MTA is complete.
\end{enumerate}
\end{openproblem}

\begin{firewall}[Claim boundary after compact V50]
\label{fire:V50-current-boundary}
V50 completes the mandatory companion audit and closes every protected
transport-bank overlap by a light-intersection disjointization or a
heavy-intersection shortcut stock.  It does not close gapped transport
degrees beyond the established cubic terminal moments.  Hence the
complete all-weak estimate, global quartic MTA, all-order MTA,
WUFC/CPM, continuum Yang--Mills existence, and a positive mass gap
remain open.
\end{firewall}

\section{Append-only record for compact V50}
\label{sec:V50-record}

V50 was appended to the complete compact V49 file.  The exact V49
parent had \(19710\) newline-delimited source records, \(655939\)
bytes, and SHA--256 digest
\[
\texttt{4b5dc6ad53e07dfb6ebed176df3209e87b3bc2ae23022aee5a8eefc705a38132}.
\]
No compact V49 mathematical content was changed.  The sole final
\verb|\end{document}| is restored below.

% =============================================================================
% END APPENDED COMPACT VERSION V50 MATERIAL
% =============================================================================

% The compact V50 document terminator is intentionally disabled here:
% \end{document}

% =============================================================================
% BEGIN APPENDED COMPACT VERSION V51 MATERIAL
% =============================================================================

\chapter{V51: Arbitrary unbalanced moments and the sharp private-tail
gate}
\label{ch:V51}

\section{The unbalanced terminal has every fixed moment}
\label{sec:V51-unbalanced-moments}

V42 stated unbalanced stock moments only through order three because
that was the largest degree then needed.  Its proof gives every fixed
order without a new representation-theoretic input.

\begin{theorem}[All fixed moments of the unbalanced stock]
\label{thm:V51-all-unbalanced-moments}
For every fixed integer \(m\ge0\),
\begin{equation}
\boxed{
 (s_\alpha W)^m
 \max_{\pi\in\Pi_4}
 \left\|\bigotimes_{e\in U}T_e(\pi)\right\|
 \le C_m.}
\label{eq:V51-all-unbalanced-moments}
\end{equation}
The same estimate holds in a graph-resolved term with at most four
marked coordinates when their residual local heat multipliers are
retained.
\end{theorem}

\begin{proof}
\Cref{lem:V42-statewise-unbalanced} gives
\[
 \max_\pi
 \left\|\bigotimes_{e\in U}T_e(\pi)\right\|
 \le Ce^{-cs_\alpha W}.
\]
The scalar supremum
\(\sup_{x\ge0}x^me^{-cx}\) is finite for every fixed \(m\), proving the
first assertion.

For a graph-resolved term, either the unmarked reserve is at least
\(W/2\), or one of at most four marked coordinates has mass at least
\(W/8\).  In the first case apply the same exponential to the reserve.
In the second use the marked-and-decaying estimate
\eqref{eq:V43-marked-decay}.  Both give \(Ce^{-c's_\alpha W}\), and
the same scalar optimization applies.
\end{proof}

\begin{corollary}[Unbalanced terminal certificates of arbitrary
finite order]
\label{cor:V51-unbalanced-terminal-all-orders}
If \(U_v>\Xi_v/4\), the common unbalanced tail satisfies
\begin{equation}
 (s_\alpha\Xi_v)^m Q_v\le C_m
\label{eq:V51-unbalanced-terminal-all-orders}
\end{equation}
for every fixed \(m\).
\end{corollary}

\begin{proof}
The total stock obeys
\(W\ge U_v>\Xi_v/4\).  Apply
\Cref{thm:V51-all-unbalanced-moments}.
\end{proof}

\begin{theorem}[All finite gapped degrees with an unbalanced terminal]
\label{thm:V51-gapped-unbalanced-terminal}
Every graph-resolved gapped transport chain ending at an
unbalanced-controlled row obeys the global marked-tree atlas.  This
includes two tree-excess powers and one labelled balanced-output payer.
\end{theorem}

\begin{proof}
A transport chain has at most three steps.  Hence the degree required
by \Cref{lem:V50-gapped-degree} is at most
\[
 k+g\le3+3=6.
\]
Use the order-\((k+g)\) certificate from
\Cref{cor:V51-unbalanced-terminal-all-orders}, perform the exact
mass-ratio telescope, and reduce the connected graph to a tree by
\Cref{cor:V47-surplus-edge}.
\end{proof}

\section{Why cubic private data do not imply a fourth moment}
\label{sec:V51-private-regression}

The private multiplier is governed in the compact programme by the
one-link gap and by explicit debits through cubic order.  These
hypotheses alone do not imply the next debit.

\begin{proposition}[Sharp scalar regression at degree four]
\label{prop:V51-private-fourth-regression}
Define, for \(x\ge0\),
\begin{equation}
 q(x):=
 \begin{cases}
 e^{-x},&0\le x\le1,\\
 e^{-1}x^{-3},&x\ge1.
 \end{cases}
\label{eq:V51-regression-multiplier}
\end{equation}
Then \(0<q(x)\le1\), and there are constants \(c,C_1,C_2,C_3>0\)
such that
\begin{align}
 1-q(x)&\ge c\min\{1,x\},
\label{eq:V51-regression-gap}\\
 x^rq(x)&\le C_r,\qquad r=1,2,3.
\label{eq:V51-regression-cubic}
\end{align}
Nevertheless
\begin{equation}
 \sup_{x\ge0}x^4q(x)=\infty.
\label{eq:V51-regression-fourth-fails}
\end{equation}
\end{proposition}

\begin{proof}
The two definitions agree at \(x=1\).  On \([0,1]\),
\[
 1-e^{-x}\ge(1-e^{-1})x.
\]
On \([1,\infty)\), \(q(x)\le e^{-1}\), so
\[
 1-q(x)\ge1-e^{-1}.
\]
This proves \eqref{eq:V51-regression-gap}.  For \(x\le1\), all three
displayed moments are at most one.  For \(x\ge1\),
\[
 x^rq(x)=e^{-1}x^{r-3},
\]
which is bounded for \(r\le3\), while
\(x^4q(x)=e^{-1}x\to\infty\).
\end{proof}

\begin{corollary}[No formal extension of the private product compiler]
\label{cor:V51-no-formal-private-extension}
The standing one-link gap and private debits through order three do
not imply a support-uniform fourth private moment.  Consequently the
V7 Bernoulli compiler cannot be invoked at orders four through six
without an additional one-link hypothesis.
\end{corollary}

\begin{proof}
Use a private bank consisting of one coordinate and set
\(x=s_\alpha p\).  The multiplier
\eqref{eq:V51-regression-multiplier} satisfies every stated scalar
hypothesis through cubic order but violates the fourth-order
conclusion.
\end{proof}

\begin{firewall}[Scope of the regression]
\label{fire:V51-regression-scope}
\Cref{prop:V51-private-fourth-regression} does not assert that the
actual Wilson multiplier has a polynomial tail.  It proves a logical
point: the compact assumptions recorded so far do not exclude that
tail.  A higher private moment must be derived from the exact one-link
measure or avoided by a different global-gap optimization.
\end{firewall}

\section{The exact residual degree classes}
\label{sec:V51-degree-classes}

Let \(k\in\{0,1,2,3\}\) be the number of starting row-mass degrees:
two at most from the tree and one at most from a labelled output payer.
Let \(g\in\{0,1,2,3\}\) be the number of gapped transport steps.

\begin{proposition}[Finite private-terminal residual]
\label{prop:V51-private-residual}
For a private-controlled terminal row, every gapped transport branch
with \(k+g\le3\) is closed by the existing cubic certificate.  The only
unclosed degree classes are
\begin{equation}
\boxed{
 (k,g)\ \hbox{with}\ k+g\in\{4,5,6\}.}
\label{eq:V51-private-residual}
\end{equation}
Equivalently, the minimal unresolved configurations are
\[
 (3,1),\qquad(2,2),\qquad(1,3).
\]
\end{proposition}

\begin{proof}
Apply \Cref{lem:V50-gapped-degree}.  The cubic private certificate
closes exactly \(k+g\le3\).  Since \(k,g\le3\), every failure has
degree between four and six.  Reducing either \(k\) or \(g\) gives one
of the three displayed minimal pairs.
\end{proof}

\section{A large global gap can repay missing degrees}
\label{sec:V51-large-gap}

Write \(N_{\rm gap}\) for the single count in
\eqref{eq:V47-one-global-gap}, and assume
\(0<s_\alpha\le1\).

\begin{lemma}[Half of the global gap survives mark extraction]
\label{lem:V51-half-gap}
Fix \(g\le3\).  In one connected-graph term, the polynomial in the
gapped bank counts used to extract \(g\) marked transport factors can
be absorbed by half of the global exponential:
\begin{equation}
 (s_\alpha N_{\rm gap})^g
 e^{-cs_\alpha N_{\rm gap}}
 \le
 C_g e^{-(c/2)s_\alpha N_{\rm gap}}.
\label{eq:V51-half-gap}
\end{equation}
Thus the marked gapped ledger can be formed while a single residual
global exponential remains.
\end{lemma}

\begin{proof}
Put \(x=s_\alpha N_{\rm gap}\).  The function
\(x^ge^{-(c/2)x}\) is bounded on \([0,\infty)\), so
\[
 x^ge^{-cx}
 =
 \bigl(x^ge^{-(c/2)x}\bigr)e^{-(c/2)x}
 \le C_ge^{-(c/2)x}.
\]
All individual gapped bank counts are bounded by \(N_{\rm gap}\), and
there are at most three transport steps.  Hence their count polynomial
is bounded by the left side.
\end{proof}

\begin{lemma}[Large-count repayment]
\label{lem:V51-large-count}
For every integer \(r\ge1\), if
\begin{equation}
 N_{\rm gap}
 \ge
 \frac{2r}{c\,s_\alpha}
 \log\frac1{s_\alpha},
\label{eq:V51-large-count-threshold}
\end{equation}
then
\begin{equation}
 e^{-cs_\alpha N_{\rm gap}}\le s_\alpha^{2r}.
\label{eq:V51-large-count-repayment}
\end{equation}
\end{lemma}

\begin{proof}
The threshold gives
\[
 cs_\alpha N_{\rm gap}
 \ge2r\log(1/s_\alpha).
\]
Exponentiate its negative.
\end{proof}

\begin{corollary}[Large-count private gapped branches]
\label{cor:V51-large-count-private}
If a private-terminal branch is missing \(r=k+g-3\) powers of
\(s_\alpha\) after its cubic debit and satisfies
\eqref{eq:V51-large-count-threshold}, the single global gapped factor
repays those powers.
\end{corollary}

\begin{proof}
The deficit is at most \(s_\alpha^{-r}\).  Use
\Cref{lem:V51-half-gap} when the gapped marks are extracted, and
decrease the constant \(c\) in
\eqref{eq:V51-large-count-threshold} to the residual exponent
constant.  Applying \eqref{eq:V51-large-count-repayment} to that
remaining exponential leaves the favorable factor
\(s_\alpha^r\le1\).  The two uses concern disjoint halves of the same
exponent, so the gap is not charged twice.
\end{proof}

\begin{assessment}[The residual is sparse-gapped and private]
\label{ass:V51-sparse-private}
After V51, an unclosed gapped transport branch must:
\begin{enumerate}
\item terminate at a private-dominant row;
\item have total degree \(k+g>3\);
\item have
\[
 N_{\rm gap}
 <
 \frac{2(k+g-3)}{c\,s_\alpha}
 \log\frac1{s_\alpha}.
\]
\end{enumerate}
All protected branches, all unbalanced-terminal gapped branches, and
all sufficiently large-count private-terminal gapped branches are
closed.  The remaining problem is a finite sparse-gapped/private
interaction.
\end{assessment}

\begin{openproblem}[V52 mandate: sparse-gap/private interaction]
\label{op:V52-mandate}
\begin{enumerate}
\item Keep the exact private multiplier \(q_v\) and the single global
      gapped exponential in one scalar optimization; do not extract
      \(g\) separate \(s_\alpha^{-1}\) factors.
\item Use the sparse count bound from
      \Cref{ass:V51-sparse-private} together with the all-weak mass
      escalation at each transport step.
\item Test whether the actual one-link Wilson integral supplies an
      exponential representation tail stronger than the compact
      cubic assumptions.
\item If not, build a scalar concentration example satisfying all
      proven stock inequalities and decide whether an additional
      theorem is genuinely necessary.
\item Close the sparse-private branch before invoking the V36 global
      quartic compiler.
\end{enumerate}
\end{openproblem}

\begin{firewall}[Claim boundary after compact V51]
\label{fire:V51-current-boundary}
V51 proves unbalanced moments of every fixed order and closes every
gapped transport chain with an unbalanced terminal.  It also proves
that the existing private assumptions cannot be extended beyond cubic
order by logic alone, while a sufficiently large global gapped count
repays the missing degrees.  The sparse-gapped/private-terminal branch
remains open.  Therefore the complete all-weak estimate, global
quartic MTA, all-order MTA, WUFC/CPM, continuum Yang--Mills existence,
and a positive mass gap remain open.
\end{firewall}

\section{Append-only record for compact V51}
\label{sec:V51-record}

V51 was appended to the complete compact V50 file.  The exact V50
parent had \(20129\) newline-delimited source records, \(669980\)
bytes, and SHA--256 digest
\[
\texttt{c08dae1f8fafa7f51125e09edf2db15dd2692cdca4aaf38394cec4339ca197a8}.
\]
No compact V50 mathematical content was changed.  The sole final
\verb|\end{document}| is restored below.

% =============================================================================
% END APPENDED COMPACT VERSION V51 MATERIAL
% =============================================================================

% The compact V51 document terminator is intentionally disabled here:
% \end{document}

% =============================================================================
% BEGIN APPENDED COMPACT VERSION V52 MATERIAL
% =============================================================================

\chapter{V52: All-order Wilson tails and closure of global quartic MTA}
\label{ch:V52}

\section{Purpose and upstream normalization check}
\label{sec:V52-purpose}

V51 isolated one residual scalar branch: a gapped transport chain ending
at a private-dominant row and requiring a terminal moment of degree four,
five, or six.  The abstract gap and cubic-debit assumptions do not imply
such a moment.  This chapter returns to the exact one-link Wilson density
and proves every fixed Casimir moment by self-adjoint integration by parts.
The argument is not a heat-kernel substitution and does not require an
explicit special-function formula.

The clock normalization was also checked against the frozen f1 archive.
There
\[
 s_\alpha=-\frac{\log\eta_F(\alpha)}{C_F},
\]
and legacy V122 corrects the Wilson-well asymptotic to
\begin{equation}
 s_\alpha=\frac3\alpha+O(\alpha^{-2}).
\label{eq:V52-clock}
\end{equation}
Only the consequence
\begin{equation}
 0<s_\alpha\le \frac{C_s}{\alpha}
 \qquad(\alpha\ge\alpha_s)
\label{eq:V52-clock-upper}
\end{equation}
is used below.  Thus the proof is insensitive to the convention-dependent
leading constant.

\section{Normalized derivatives of the exact Wilson density}
\label{sec:V52-density-derivatives}

Let (G) be the standing compact connected semisimple gauge group, with
its fixed bi-invariant metric and normalized Haar measure.  Write the
one-link law as
\begin{equation}
 q_\alpha(U)\,\mathrm dU
 =Z_\alpha^{-1}e^{-\alpha V(U)}\,\mathrm dU,
 \qquad
 V(U)=\frac13\left(
 \operatorname{ReTr}\mathbf 1-\operatorname{ReTr}U\right).
\label{eq:V52-Wilson-density}
\end{equation}
The additive constant makes (V(1)=0) and cancels from the normalized
law.  The argument below uses
the Wilson-well facts already established in the legacy one-link analysis:
(V\ge0), the identity is its unique minimum, and the Hessian of (V) at
the identity is positive definite.  These facts hold for the standard
faithful Wilson trace on (SU(2)) and (SU(3)).

Fix the geodesic distance
\[
 r(U):=d_G(U,1).
\]

\begin{lemma}[Uniform Laplace moments of the Wilson well]
\label{lem:V52-Laplace-moments}
For every integer \(\ell\ge0\), there are constants
\(C_\ell<\infty\) and \(\alpha_\ell\) such that
\begin{equation}
 \boxed{
 \int_G r(U)^\ell q_\alpha(U)\,\mathrm dU
 \le C_\ell\alpha^{-\ell/2}}
 \qquad(\alpha\ge\alpha_\ell).
\label{eq:V52-Laplace-moments}
\end{equation}
Moreover there are (0<c_Z<C_Z<\infty) such that, if
\(d=\dim G\),
\begin{equation}
 c_Z\alpha^{-d/2}\le Z_\alpha\le C_Z\alpha^{-d/2}
\label{eq:V52-partition-scale}
\end{equation}
for all sufficiently large \(\alpha\).
\end{lemma}

\begin{proof}
Choose a normal ball (B_\varepsilon(1)) on which the Haar Jacobian is
bounded above and below and
\[
 c_0r(U)^2\le V(U)\le C_0r(U)^2.
\]
The lower bound for (Z_\alpha) follows by integrating on the ball of
radius \(\alpha^{-1/2}\).  The upper bound on the normal ball follows
from the Gaussian integral after the change of variables
\(x=\sqrt\alpha X\).  On its compact complement, uniqueness of the
minimum gives (V\ge c_1>0), whose contribution is (O(e^{-c_1\alpha})).
This proves \eqref{eq:V52-partition-scale}.

The same decomposition, with the factor (r^\ell), gives on the normal
ball
\[
 \int r^\ell e^{-\alpha V}\,\mathrm dU
 \le C\alpha^{-(d+\ell)/2}.
\]
The complement contributes (O(e^{-c_1\alpha})).  Divide by the lower
bound for (Z_\alpha); the exponential remainder is bounded by a constant
times \(\alpha^{-\ell/2}\).  This proves
\eqref{eq:V52-Laplace-moments}.
\end{proof}

Let (X_1,\ldots,X_d) be a fixed orthonormal left-invariant frame and use
the nonnegative Casimir convention
\[
 \mathcal L_G:=-\sum_{a=1}^dX_a^2.
\]

\begin{lemma}[Derivative monomials have parabolic weight at most their
half-order]
\label{lem:V52-derivative-weight}
Fix a word (D=X_{i_1}\cdots X_{i_r}).  On a sufficiently small normal
ball, (e^{\alpha V}D(e^{-\alpha V})) is a finite sum of terms bounded by
\begin{equation}
 C_D\alpha^t r(U)^\ell,
 \qquad
 t-\frac\ell2\le\frac r2.
\label{eq:V52-derivative-weight}
\end{equation}
The constants and number of terms depend on (D), but not on \(\alpha\).
\end{lemma}

\begin{proof}
Repeated use of
\[
 Xe^{-\alpha V}=-\alpha(XV)e^{-\alpha V}
\]
shows that every term is \(\alpha^t\) times a product of (t)
derivatives of (V), of positive orders whose sum is at most (r).
Let \(\ell\) be the number of first derivatives among those (t) factors.
Every other factor uses at least two derivatives, so
\[
 2t-\ell\le r.
\]
Because (\mathrm dV(1)=0), every first derivative satisfies
\(|XV(U)|\le Cr(U)\) on the normal ball; all derivatives of order at least
two are bounded there.  Hence the term is bounded by
\(C_D\alpha^tr^\ell\), and the displayed order inequality is equivalent to
\eqref{eq:V52-derivative-weight}.
\end{proof}

\begin{theorem}[All normalized Laplacian derivatives]
\label{thm:V52-density-all-derivatives}
For every integer (m\ge0), there are (C_m<\infty) and \(\alpha_m\)
such that
\begin{equation}
 \boxed{
 \|\mathcal L_G^m q_\alpha\|_{L^1(G)}
 \le C_m\alpha^m}
 \qquad(\alpha\ge\alpha_m).
\label{eq:V52-density-all-derivatives}
\end{equation}
\end{theorem}

\begin{proof}
Expand \(\mathcal L_G^m\) into finitely many words of length (2m).
On the normal ball, every term in
\Cref{lem:V52-derivative-weight} satisfies
\[
 \alpha^t\mathbb E_{q_\alpha}r^\ell
 \le C_\ell\alpha^{t-\ell/2}
 \le C_\ell\alpha^m
\]
by \Cref{lem:V52-Laplace-moments}.  On the complement of a smaller
normal ball, all derivatives of (V) are bounded, so the same expansion
is at most a fixed polynomial of degree (2m) in \(\alpha\), multiplied
by (e^{-c\alpha}/Z_\alpha).  Equation
\eqref{eq:V52-partition-scale} and exponential domination bound this by
\(C_m\alpha^m\).  Sum the finitely many terms.  The normalizing factor
(Z_\alpha^{-1}) is constant in (U), so the result is exactly
\eqref{eq:V52-density-all-derivatives}.
\end{proof}

\section{All-order Casimir and failure debits}
\label{sec:V52-all-order-debits}

For an irreducible unitary representation (R), let (d_R) be its
dimension, (C_2(R)) its Casimir eigenvalue, and
\[
 \eta_R(\alpha)
 =\frac1{d_R}\int_G\chi_R(U)q_\alpha(U)\,\mathrm dU
\]
its central Wilson multiplier.

\begin{theorem}[Every fixed one-link Casimir moment]
\label{thm:V52-all-Casimir-moments}
For every integer (m\ge0),
\begin{equation}
 \boxed{
 C_2(R)^m|\eta_R(\alpha)|
 \le C_m\alpha^m}
\label{eq:V52-all-Casimir-moments}
\end{equation}
uniformly in (R) and in sufficiently large \(\alpha\).  Consequently,
for every nontrivial (R),
\begin{equation}
 \boxed{
 \bigl(s_\alpha(1+C_2(R))\bigr)^m
 \eta_R(\alpha)\le M_m.}
\label{eq:V52-scaled-Casimir-moments}
\end{equation}
\end{theorem}

\begin{proof}
The character is a Casimir eigenfunction:
\[
 \mathcal L_G^m\chi_R=C_2(R)^m\chi_R.
\]
Self-adjointness on Haar measure, followed by
\(|\chi_R(U)|\le d_R\), gives
\begin{align*}
 C_2(R)^m|\eta_R(\alpha)|
 &=\left|\frac1{d_R}
   \int_G\chi_R\mathcal L_G^mq_\alpha\,\mathrm dU\right|\\
 &\le\|\mathcal L_G^mq_\alpha\|_{L^1(G)}.
\end{align*}
Apply \Cref{thm:V52-density-all-derivatives}.

Because (G) is fixed and semisimple, the nonzero Casimir spectrum has a
positive floor (c_G).  Hence (1+C_2(R)\le K_GC_2(R)) for nontrivial
(R).  Combine this with \eqref{eq:V52-clock-upper} and
\eqref{eq:V52-all-Casimir-moments}.  Positivity of the standing Wilson
multipliers removes the absolute value and proves
\eqref{eq:V52-scaled-Casimir-moments}.
\end{proof}

The product compiler needs a failure debit, rather than merely a bounded
moment.  The established one-link gap supplies it.

\begin{corollary}[All-order one-site failure debits]
\label{cor:V52-all-failure-debits}
For every fixed integer (m\ge1), there are constants (D_k),
(1\le k\le m), such that, for every nontrivial (R),
\begin{equation}
 \boxed{
 x_R^k\eta_R(\alpha)
 \le D_k\bigl(1-\eta_R(\alpha)\bigr),
 \qquad
 x_R:=s_\alpha(1+C_2(R)).}
\label{eq:V52-all-failure-debits}
\end{equation}
The constants are uniform in (R) and sufficiently large \(\alpha\).
\end{corollary}

\begin{proof}
Put (y_R=s_\alpha C_2(R)).  The Casimir floor gives
(x_R\le K_Gy_R).  If (y_R\le1), then, using (0\le\eta_R\le1),
\[
 x_R^k\eta_R\le K_G^ky_R^k\le K_G^ky_R.
\]
The one-link gap \eqref{eq:V3-one-link-gap} bounds the last expression by
a constant times (1-\eta_R).  If (y_R>1), then the same gap gives
(1-\eta_R\ge c_{\rm gap}), while
\Cref{thm:V52-all-Casimir-moments} bounds (x_R^k\eta_R) uniformly.
This proves \eqref{eq:V52-all-failure-debits}.
\end{proof}

\begin{corollary}[Private Wilson products of every fixed order]
\label{cor:V52-private-products-all-orders}
Let (P) be a finite exact private coordinate bank and put
\[
 Y:=\sum_{f\in P}(1+C_2(R_f)),
 \qquad
 q_P:=\prod_{f\in P}\eta_{R_f}(\alpha).
\]
For every fixed integer (m\ge0),
\begin{equation}
 \boxed{(s_\alpha Y)^mq_P\le C_m^{\rm priv}}
\label{eq:V52-private-products-all-orders}
\end{equation}
uniformly in the size of (P), all representations, and sufficiently
large \(\alpha\).
\end{corollary}

\begin{proof}
For (m=0), use (0\le q_P\le1).  For (m\ge1), apply
\Cref{lem:V7-product-moment-compiler} with
\[
 x_f=s_\alpha(1+C_2(R_f)),
 \qquad r_f=\eta_{R_f}(\alpha),
\]
and use \Cref{cor:V52-all-failure-debits} for every (1\le k\le m).
The Bernoulli sum in that compiler is a probability, so no support count
is introduced.
\end{proof}

\begin{firewall}[Fixed order versus an all-order atlas constant]
\label{fire:V52-fixed-order}
The constants (C_m^{\rm priv}) are finite for each fixed (m), but this
chapter does not prove (C_m^{\rm priv}\le C^m m!), or the stronger
per-tree exponential bound needed by the all-order marked-tree atlas.
Only degrees at most six are used in the quartic transport closure below.
Thus \Cref{cor:V52-private-products-all-orders} must not be read as
all-order MTA.
\end{firewall}

\section{Closure of the private terminal and the all-weak carrier}
\label{sec:V52-private-closure}

\begin{theorem}[Every finite gapped degree with a private terminal closes]
\label{thm:V52-gapped-private-terminal}
Every graph-resolved gapped transport chain ending at a private-dominant
row obeys the global quartic marked-tree atlas.  In particular, all degree
classes (k+g\in\{4,5,6\}) isolated in
\eqref{eq:V51-private-residual} are closed.
\end{theorem}

\begin{proof}
Let (v) be the terminal row.  Private dominance gives a private bank
(P_v) whose mass (Y_v) satisfies (Y_v\ge\Xi_v/2), and the exact
private-coordinate factorization retains its multiplier (q_{P_v}).
A quartic transport chain has at most three starting degrees and at most
three gapped steps, so (d:=k+g\le6).  By
\Cref{cor:V52-private-products-all-orders},
\[
 (s_\alpha\Xi_v)^dq_{P_v}
 \le2^d(s_\alpha Y_v)^dq_{P_v}
 \le C_d.
\]
This is precisely the terminal moment certificate required in
\Cref{lem:V50-gapped-degree}.  Apply that lemma, telescope all row-mass
ratios as in V46, and discard surplus graph edges only after the complete
connected-graph difference has been formed.  No second cut capacity and no
second output resolvent is used.
\end{proof}

\begin{theorem}[Complete all-weak normalized quartic carrier]
\label{thm:V52-all-weak-closed}
When all normalized overlaps satisfy
\(
 \theta_{ij}<\delta_0
\), the complete coefficient-pinched, one-resolvent quartic carrier obeys
\begin{equation}
 \boxed{
 \mathcal F_4
 \le C\sum_{\tau\in\mathfrak T_4}
       \prod_{ij\in E(\tau)}\theta_{ij}.}
\label{eq:V52-all-weak-closed}
\end{equation}
Thus hypothesis \textup{(iv)} of
\Cref{thm:V36-normalized-graph} is proved.
\end{theorem}

\begin{proof}
The mass-transport alternative of
\Cref{cor:V44-no-cycle} is exhaustive and terminates after at most three
steps.  Protected terminal banks, including every incidence-three or
incidence-four overlap and the crossed word ((e,f,e)), are closed by
\Cref{thm:V50-all-protected-overlaps}.  Unbalanced terminals of arbitrary
finite degree are closed by
\Cref{thm:V51-gapped-unbalanced-terminal}.  Private terminals are closed by
\Cref{thm:V52-gapped-private-terminal}.  The dense gapped ledger was reduced
in V50--V51 to these same finite terminal alternatives; large global gap
count is additionally covered by
\Cref{cor:V51-large-count-private}.  These cases exhaust the post-V41
all-weak ledger and preserve the one global connected-graph difference.
Summing its fixed finite graph types proves
\eqref{eq:V52-all-weak-closed}.
\end{proof}

\begin{corollary}[Global quartic marked-tree atlas]
\label{cor:V52-global-quartic-MTA}
The complete compact-series quartic carrier satisfies
\begin{equation}
 \boxed{
 \kappa_\otimes(\mathcal C_4)
 \le C\sum_{\tau\in\mathfrak T_4}
 \left(\prod_{i=1}^4\Xi_i\right)
 \prod_{ij\in E(\tau)}\theta_{ij}.}
\label{eq:V52-global-quartic-MTA}
\end{equation}
Equivalently, the global (m=4) instance of MTA is closed for the
complete coefficient-pinched carrier with its single final resolvent.
\end{corollary}

\begin{proof}
V37 proves the two-cluster hypothesis \textup{(ii)} of
\Cref{thm:V36-normalized-graph}; V38 proves the three-cluster hypothesis
\textup{(iii)}; V39 proves the universal hypothesis \textup{(i)}; and
\Cref{thm:V52-all-weak-closed} proves hypothesis \textup{(iv)}.  Apply
\Cref{thm:V36-normalized-graph}, then multiply its normalized conclusion by
\(\prod_i\Xi_i\) as in \Cref{cor:V36-four-targets}.
\end{proof}

\section{V52 ledger and V53 mandate}
\label{sec:V52-ledger}

\begin{theorem}[Integrated compact-series V52 statement]
\label{thm:V52-integrated}
Preserving compact V1--V51, V52 proves:
\begin{enumerate}[label=\textup{(V52.\arabic*)},leftmargin=6.7em]
\item normalized Wilson-well moments at every fixed radial order;
\item the derivative estimate
      \(\|\mathcal L_G^mq_\alpha\|_1\le C_m\alpha^m\);
\item every fixed one-link Casimir moment and failure debit;
\item every fixed private-product moment, uniformly in support size;
\item closure of all private-terminal gapped degrees through six;
\item the complete all-weak normalized quartic estimate; and
\item global quartic MTA by the V36 finite normalized graph compiler.
\end{enumerate}
\end{theorem}

\begin{proof}
The assertions are
\Cref{lem:V52-Laplace-moments,
thm:V52-density-all-derivatives,
thm:V52-all-Casimir-moments,
cor:V52-all-failure-debits,
cor:V52-private-products-all-orders,
thm:V52-gapped-private-terminal,
thm:V52-all-weak-closed,
cor:V52-global-quartic-MTA}.
\end{proof}

\begin{openproblem}[V53 mandate: the first post-quartic tree fibre]
\label{op:V53-mandate}
\begin{enumerate}
\item Return to the legacy V39 connected replica-partition formula and
      its canonical marked-tree extraction; do not replace it by a new
      partition expansion.
\item Refine every local hyperedge by contraction trees before taking
      absolute values, as required by the legacy V40 high-weight
      regression and the quartic construction now completed.
\item Treat (m=5) first.  Keep the complete Wilson cumulant, all
      Peter--Weyl output blocks, and the single final resolvent assembled.
\item Separate analytic growth from combinatorial growth.  The fixed-order
      constants in \Cref{thm:V52-density-all-derivatives} do not yet have
      the per-edge exponential bound required for MTA.
\item Prove a (C_0^4) bound for each contraction-refined quintic tree
      fibre, or return the first resolved support/representation family
      for which that fibre estimate fails.
\item Only after a uniform leaf- or hyperedge-deletion step is proved may
      the all-order MTA compiler, WUFC, or CPM be invoked.
\end{enumerate}
\end{openproblem}

\begin{firewall}[Claim boundary after compact V52]
\label{fire:V52-current-boundary}
V52 closes the last scalar branch of the global quartic proof and thereby
proves the complete (m=4) marked-tree atlas stated in
\eqref{eq:V52-global-quartic-MTA}.  It does not prove the contraction-refined
quintic fibre, an order-uniform MTA constant, WUFC, CPM, the nonlinear
Perron/source packet, finite-edge margins, capacity, protected-sector
floors, a nontrivial reflection-positive continuum state, or a positive
physical mass gap.  The full Yang--Mills mass gap remains open.
\end{firewall}

\section{Append-only record for compact V52}
\label{sec:V52-record}

V52 was appended to the complete compact V51 file.  The exact V51 parent
had (20474) newline-delimited source records, (681339) bytes, and
SHA--256 digest
\[
\texttt{86a7e85c9c2d9109533d85955a785685d699242ef88436a604cfd290638d40db}.
\]
No compact V51 mathematical content was changed.  The sole final
\verb|\end{document}| is restored below.

% =============================================================================
% END APPENDED COMPACT VERSION V52 MATERIAL
% =============================================================================

% The compact V52 document terminator is intentionally disabled here:
% \end{document}

% =============================================================================
% BEGIN APPENDED COMPACT VERSION V53 MATERIAL
% =============================================================================

\chapter{V53: The contraction-refined quintic coefficient}
\label{ch:V53}

\section{Purpose and fixed-bank normal form}
\label{sec:V53-purpose}

V52 closes global quartic MTA.  The first genuinely new order is therefore
(m=5).  The frozen f1 analysis already supplies two indispensable pieces:
the exact connected replica-partition formula and a deterministic projection
of every connected coordinate partition to a marked spanning tree.  It also
shows, at order four, that assigning a whole local hyperedge to an arbitrary
tree before resolving its contractions gives the wrong high-weight valence.
This chapter does not repeat those results.  It computes the first quintic
coefficient and identifies the correct contraction-refined tree index.

Put (s=s_\alpha).  On a fixed finite Peter--Weyl bank, the legacy
Wilson-well normal form can be written
\begin{equation}
 \mathrm d\nu_s(Z)
 =\phi(Z)\left(1+s a_1(Z)+O_{\mathscr W}(s^2)\right)\mathrm dZ,
 \qquad
 \mathbb E_\phi a_1=0,
\label{eq:V53-density-normal-form}
\end{equation}
after (U=\exp(\sqrt{s},Z)).  Here \(\phi\) is the centered nondegenerate
Gaussian on the Lie algebra, (a_1) is even and has polynomial degree at
most four, and the remainder, together with its fixed-bank moments, is
uniform.  Inversion invariance makes every displayed density correction
even.

For a smooth matrix coefficient (g_i), let
\[
 D_i^{(r)}:=\mathrm d^rg_i(1)
\]
be its symmetric (r)-linear differential.  Let
\(
 \Gamma(\ell,\ell')=\mathbb E_\phi[\ell(Z)\ell'(Z)]
\)
be the complex bilinear Gaussian covariance.

\begin{lemma}[First variation of a cumulant]
\label{lem:V53-cumulant-variation}
Let (a) be integrable with \(\mathbb E_\phi a=0\), and let
\(
 \mathbb E_\varepsilon F
 =\mathbb E_\phi[(1+\varepsilon a)F]
\).
For variables (Y_1,\ldots,Y_m) having all required moments,
\begin{equation}
 \boxed{
 \left.\frac{\mathrm d}{\mathrm d\varepsilon}\right|_{0}
 \kappa_m^{(\varepsilon)}(Y_1,\ldots,Y_m)
 =\kappa_{m+1}^{\phi}(a,Y_1,\ldots,Y_m).}
\label{eq:V53-cumulant-variation}
\end{equation}
\end{lemma}

\begin{proof}
For formal source variables (t_i), set
\[
 M_\varepsilon(t)
 =\mathbb E_\phi\left[(1+\varepsilon a)
   \exp\left(\sum_it_iY_i\right)\right].
\]
Then
\[
 \left.\partial_\varepsilon\right|_0\log M_\varepsilon(t)
 =\frac{\mathbb E_\phi[a e^{\sum_it_iY_i}]}
        {\mathbb E_\phi[e^{\sum_it_iY_i}]}.
\]
The coefficient of (t_1\cdots t_m) on the right is, by the defining
logarithmic moment generating function, the joint cumulant of
(a,Y_1,\ldots,Y_m).  Equality of these coefficients proves
\eqref{eq:V53-cumulant-variation}.
\end{proof}

\section{Connected Gaussian pairings at five sources}
\label{sec:V53-Gaussian-trees}

For a labeled tree \(\tau\in\mathfrak T_5\), define
\begin{equation}
 \mathcal W_\tau(g_1,\ldots,g_5)
 :=\operatorname{Contr}_{\Gamma,\tau}
   \bigotimes_{i=1}^5D_i^{(\deg_\tau(i))},
\label{eq:V53-Wick-tree}
\end{equation}
contracting one differential slot at each endpoint of every tree edge.
For (1\le r\le5), define the density-correction coefficient
\begin{equation}
 \mathcal U_{5,r}
 :=\kappa_6^\phi\left(
 a_1,
 D_r^{(2)}[Z,Z],
 (D_i^{(1)}[Z])_{i\ne r}
 \right).
\label{eq:V53-density-correction}
\end{equation}

\begin{lemma}[Connected-pairing cancellation]
\label{lem:V53-connected-pairings}
Let (P_i(Z)) be homogeneous polynomials of positive degrees (r_i).
In their Gaussian joint cumulant, a Wick pairing survives if and only if
the multigraph induced on the five polynomial labels is connected.  Hence
the cumulant vanishes when \(\sum_i r_i<8\).  If
\(\sum_i r_i=8\), every surviving pairing has no self-pair and its induced
graph is a labeled tree.
\end{lemma}

\begin{proof}
Expand every Gaussian moment in the moment--cumulant formula by Wick's
rule.  Fix one pairing of all scalar Gaussian slots and let \(\pi\) be the
partition of the five polynomial labels into the connected components of
its induced graph.  The coefficient of this pairing is
\[
 \sum_{\rho\ge\pi}\mu(\rho,\boldsymbol1_5),
\]
which equals one if \(\pi=\boldsymbol1_5\) and zero otherwise by Möbius
inversion.  A connected graph on five vertices needs at least four
cross-pairs, hence at least eight scalar slots.  At exactly eight slots
there are exactly four pairs, all cross-pairs, and the connected induced
graph has five vertices and four edges, so it is a tree.
\end{proof}

\begin{theorem}[Exact leading quintic Wilson coefficient]
\label{thm:V53-leading-quintic}
On every fixed finite Peter--Weyl bank, uniformly for (g_i) in a bounded
set,
\begin{align}
 &\kappa_5^{\nu_s}(g_1,\ldots,g_5)
 \notag\\
 &\quad=s^4\left[
 \sum_{\tau\in\mathfrak T_5}\mathcal W_\tau(g_1,\ldots,g_5)
 +\frac12\sum_{r=1}^5\mathcal U_{5,r}
 \right]
 +O_{\mathscr W}(s^5).
\label{eq:V53-leading-quintic}
\end{align}
In particular, the pure Gaussian part has exactly
\(
 |\mathfrak T_5|=5^{5-2}=125
\)
contraction trees.
\end{theorem}

\begin{proof}
Cumulants of order at least two are invariant under adding constants, so
replace every (g_i) by (g_i-g_i(1)).  On the fixed bank,
\begin{equation}
 g_i(\exp(\sqrt{s}Z))-g_i(1)
 =\sum_{r\ge1}\frac{s^{r/2}}{r!}D_i^{(r)}[Z^{\otimes r}],
\label{eq:V53-source-Taylor}
\end{equation}
with a Gaussian-dominated remainder through the finite order used here.

First use the Gaussian part of \eqref{eq:V53-density-normal-form}.
By \Cref{lem:V53-connected-pairings}, all source terms of total degree
below eight cancel, and odd total degrees vanish.  At degree eight, the
surviving pairing graph is a tree.  For a fixed labeled tree \(\tau\), the
degree at source (i) is \(\deg_\tau(i)\).  Its incident edges can be
assigned to the symmetric differential slots in
\(\deg_\tau(i)!\) ways, cancelling the Taylor denominator in
\eqref{eq:V53-source-Taylor}.  The contribution is therefore exactly
\(\mathcal W_\tau\), and Cayley's formula gives 125 trees.

Now take the coefficient (s a_1) in the density.  By
\Cref{lem:V53-cumulant-variation}, its contribution is the Gaussian joint
cumulant with (a_1) as one additional variable.  Five linear source
terms have odd total degree and vanish because (a_1) is even.  The next
possibility has total source degree six: exactly one source is quadratic
and the other four are linear.  Its Taylor coefficient is (1/2), and
its value is \(\mathcal U_{5,r}\).  Multiplication by the density factor
(s) and the source factor (s^{6/2}) gives (s^4/2).

All remaining Gaussian source terms have total degree at least ten.  The
remaining first-density terms have total source degree at least eight, and
higher density corrections carry at least (s^2).  Parity removes the
intervening half-integer orders.  Fixed-bank Gaussian domination therefore
places their sum in (O_{\mathscr W}(s^5)), proving
\eqref{eq:V53-leading-quintic}.
\end{proof}

\section{Tree-valence bounds and the sharp regression}
\label{sec:V53-valence}

Let (A_i:=1+C_2(R_i)) for normalized matrix coefficients in irreducible
representations (R_i).

\begin{proposition}[Every leading term admits a quintic tree valence]
\label{prop:V53-leading-valence}
There is a fixed-bank constant (C_{\mathscr W}) such that
\begin{align}
 |\mathcal W_\tau(g_1,\ldots,g_5)|
 &\le C_{\mathscr W}
 \prod_{i=1}^5 A_i^{\deg_\tau(i)/2},
\label{eq:V53-Wick-valence-bound}\\
 |\mathcal U_{5,r}|
 &\le C_{\mathscr W}
 A_r\prod_{i\ne r}A_i^{1/2}.
\label{eq:V53-density-valence-bound}
\end{align}
For each (r), the right side of
\eqref{eq:V53-density-valence-bound} is bounded by the valence monomial of
any labeled tree in which (r) has degree at least two.  Consequently the
whole coefficient of (s^4) in
\eqref{eq:V53-leading-quintic} is bounded by
\begin{equation}
 C_{\mathscr W}\sum_{\tau\in\mathfrak T_5}
 \prod_{i=1}^5A_i^{\deg_\tau(i)/2}.
\label{eq:V53-leading-tree-bound}
\end{equation}
\end{proposition}

\begin{proof}
Repeated differentiation of a normalized matrix coefficient and the
derived-representation Casimir identity give, for each fixed (k),
\[
 \|D_i^{(k)}\|\le C_kA_i^{k/2}.
\]
Gaussian contraction along the four edges of \(\tau\) proves
\eqref{eq:V53-Wick-valence-bound}.  The polynomial (a_1) is fixed and
Gaussian-integrable, so Hölder's inequality and the (k=1,2) derivative
bounds prove \eqref{eq:V53-density-valence-bound}.

If \(\deg_\tau(r)\ge2\), its valence monomial contains at least (A_r).
Every other vertex of a tree has degree at least one and hence contributes
at least (A_i^{1/2}).  Since all (A_i\ge1), the density monomial is no
larger.  Choose one such tree for each (r), and enlarge the constant to
obtain \eqref{eq:V53-leading-tree-bound}.
\end{proof}

\begin{proposition}[A quintic one-Casimir-per-input bound is false]
\label{prop:V53-bare-quintic-false}
For (SU(3)), no constant uniform in (s) and the five irreducible inputs
can satisfy
\begin{equation}
 \|\mathbf K_{12345}\|_{\rm op}
 \le C\min\left\{
 s^4\prod_{i=1}^5(1+C_2(R_i)),1\right\},
\label{eq:V53-false-bare-quintic}
\end{equation}
where \(\mathbf K_{12345}\) is the complete one-link fifth cumulant.
The failure already occurs for one coherent input
\(R_N=(N,0)\) and four fundamental inputs.
\end{proposition}

\begin{proof}
Let
\[
 g_N(U)=(U_{11})^N,
 \qquad h(U)=U_{11},
\]
and take ((g_1,\ldots,g_5)=(g_N,h,h,h,h)).  If
\(\ell(X)=X_{11}\), then
\[
 \mathrm d^kg_N(1)=N^k\ell^{\otimes k}+O_k(N^{k-1}).
\]
The unique labeled star centered at the coherent input contributes
\(c^4N^4+O(N^3)\) to the tree sum, where
\(c=\Gamma(\ell,\ell)\ne0\).  Every other tree has coherent-input degree at
most three and contributes (O(N^3)).  The density terms are (O(N^2)):
at most the quadratic differential of (g_N) occurs in
\eqref{eq:V53-density-correction}.  Therefore
\[
 \lim_{s\downarrow0}s^{-4}
 \kappa_5^{\nu_s}(g_N,h,h,h,h)
 =c^4N^4+O(N^3).
\]
This scalar is a matrix entry of \(\mathbf K_{12345}\).
But (1+C_2(R_N)\asymp N^2), while the four fundamental Casimirs are
fixed.  Dividing \eqref{eq:V53-false-bare-quintic} by (s^4), sending
(s\downarrow0) at fixed (N), and then sending (N\to\infty) would
force (N^4\lesssim N^2), a contradiction.
\end{proof}

\begin{assessment}[What V53 changes in the all-order fibre]
\label{ass:V53-fibre-index}
At a five-row coordinate, the pure Gaussian coefficient must be split into
125 source trees before absolute values.  The first density correction must
remain assembled as the five coefficients
\(\mathcal U_{5,r}\), each routed only to a tree in which (r) is internal.
An arbitrary lexicographic star projection loses the high-weight valence,
while one Casimir per input is disproved by
\Cref{prop:V53-bare-quintic-false}.  This is an indexing correction, not yet
a finite-(s), mixed-support estimate.
\end{assessment}

\section{V53 ledger and V54 mandate}
\label{sec:V53-ledger}

\begin{theorem}[Integrated compact-series V53 statement]
\label{thm:V53-integrated}
Preserving compact V1--V52, V53 proves:
\begin{enumerate}[label=\textup{(V53.\arabic*)},leftmargin=6.7em]
\item the exact first-variation identity for cumulants;
\item cancellation of every disconnected Gaussian pairing at five rows;
\item the exact leading quintic coefficient, consisting of 125 source
      trees and five assembled first-density terms;
\item tree-valence bounds for every leading coefficient; and
\item a coherent high-weight regression disproving the bare
      one-Casimir-per-input quintic estimate.
\end{enumerate}
\end{theorem}

\begin{proof}
These are
\Cref{lem:V53-cumulant-variation,
lem:V53-connected-pairings,
thm:V53-leading-quintic,
prop:V53-leading-valence,
prop:V53-bare-quintic-false}.
\end{proof}

\begin{openproblem}[V54 mandate: finite-(s) quintic support fibres]
\label{op:V54-mandate}
\begin{enumerate}
\item Lift the contraction-refined local index of
      \Cref{ass:V53-fibre-index} to the exact finite-(s) Wilson operator.
      Fixed-bank asymptotics may test the result but may not bound growing
      representations.
\item In the connected replica formula, classify the inclusion-minimal
      five-row coordinate hypergraphs by a deterministic leaf or core
      reduction.  Retain repeated coordinates explicitly.
\item Treat first a single five-row hyperedge embedded in larger supports.
      Keep all private multipliers and the single final resolvent with the
      local coefficient before applying any norm.
\item Route every density-type term only through an internal-vertex tree
      allowed by \eqref{eq:V53-density-valence-bound}; do not send the whole
      hyperedge to one lexicographic star.
\item Test the finite-(s) statement on the coherent family in
      \Cref{prop:V53-bare-quintic-false} and on a dual-singlet cancellation
      bank.
\item Isolate one exact normalized quintic carrier, or produce a resolved
      counterexample and go upstream to its local multiplier identity.
\end{enumerate}
\end{openproblem}

\begin{firewall}[Claim boundary after compact V53]
\label{fire:V53-current-boundary}
V53 determines and bounds the complete leading quintic coefficient on each
fixed Peter--Weyl bank.  It does not provide a finite-(s) estimate uniform
in representation height, close an embedded five-row hyperedge, prove the
global quintic or all-order MTA, WUFC, CPM, continuum Yang--Mills existence,
or a positive mass gap.  The full Yang--Mills mass gap remains open.
\end{firewall}

\section{Append-only record for compact V53}
\label{sec:V53-record}

V53 was appended to the complete compact V52 file.  The exact V52 parent
had (20946) newline-delimited source records, (698394) bytes, and
SHA--256 digest
\[
\texttt{1714d80b4b1bbf099e60ffd3f62bbadd46d713c738a058bbb2c2fd2a8ea4d905}.
\]
No compact V52 mathematical content was changed.  The sole final
\verb|\end{document}| is restored below.

% =============================================================================
% END APPENDED COMPACT VERSION V53 MATERIAL
% =============================================================================

% The compact V53 document terminator is intentionally disabled here:
% \end{document}

% =============================================================================
% BEGIN APPENDED COMPACT VERSION V54 MATERIAL
% =============================================================================

\chapter{V54: A fixed-order private-terminal compiler}
\label{ch:V54}

\section{Purpose}
\label{sec:V54-purpose}

V53 identifies the contraction-refined tree valence of the leading
five-row Wilson coefficient, but fixed-bank asymptotics do not control the
finite-(s_\alpha) remainder at growing representation height.  Before
attacking that local operator problem, this chapter removes a separate
global issue: private spectators and the final output resolvent.

The result is a fixed-order compiler.  At every order (m), once one
selected local connected block has a finite-(s_\alpha) tree-valence bound,
an exact private coordinate bank pays all repeated global row marks with no
support count.  At (m=5), the required private moments have orders three
and four and are already supplied by V52.  Thus the embedded quintic
private-dominant branch has no further scalar or resolvent obstruction.

\section{The local tree-card interface}
\label{sec:V54-local-interface}

Fix (m\ge2), nonempty exact row supports (X_1,\ldots,X_m), and a
coordinate (e\in\bigcap_iX_i).  Let
\(
 \tau\in\mathfrak T_m
\)
have degrees (d_i=\deg_\tau(i)), and put
\begin{equation}
 a_i:=1+C_2(R_{i,e}),
 \qquad
 w_{\tau,e}^{(m)}
 :=s_\alpha^{m-1}\prod_{i=1}^ma_i^{d_i/2}.
\label{eq:V54-local-tree-weight}
\end{equation}
The tree identity gives
\begin{equation}
 \sum_{i=1}^m(d_i-1)=m-2.
\label{eq:V54-tree-excess}
\end{equation}

Let (P_i\subseteq X_i) be the coordinates private to row (i), put
\[
 P:=\mathop{\dot\bigcup}_{i=1}^mP_i,
 \qquad
 Y:=\sum_{f\in P}(1+C_2(R_f)),
 \qquad
 q_P:=\prod_{f\in P}\eta_{R_f}(\alpha),
\]
and let \(\Xi_i\) be the full mass of row (i).

\begin{definition}[Finite-(s\) local tree card]
\label{def:V54-local-tree-card}
A resolved connected local fibre \(\mathbf K_{\tau,e}^{(m)}\) satisfies
the local tree card with constant (L_m) if
\begin{equation}
 \boxed{
 \|\mathbf K_{\tau,e}^{(m)}\|_{\rm op}
 \le L_m w_{\tau,e}^{(m)}.}
\label{eq:V54-local-tree-card}
\end{equation}
All local cumulant partitions and final multiplicity spaces belonging to
this fibre must be assembled before the norm is taken.
\end{definition}

\begin{lemma}[Exact private factorization at arbitrary fixed order]
\label{lem:V54-private-factorization}
Suppose one connected local fibre has been selected at (e), every
coordinate in (P) is private, and all other coordinates carry complete
moment operators.  Then its exact operator has the form
\begin{equation}
 \mathbf J_{\tau,e}^{(m)}
 =q_P\mathbf K_{\tau,e}^{(m)}
 \otimes\bigotimes_{f\notin P\cup\{e\}}\mathbf Q_f,
 \qquad
 \|\mathbf Q_f\|_{\rm op}\le1.
\label{eq:V54-private-factorization}
\end{equation}
Consequently, under \eqref{eq:V54-local-tree-card},
\begin{equation}
 \boxed{
 \|\mathbf J_{\tau,e}^{(m)}\|_{\rm op}
 \le L_mq_Pw_{\tau,e}^{(m)}.}
\label{eq:V54-private-factor-bound}
\end{equation}
\end{lemma}

\begin{proof}
At a private coordinate exactly one row is nontrivial.  Its complete local
moment is therefore the scalar central multiplier
\(\eta_{R_f}(\alpha)\); independence tensors their product into (q_P).
At any remaining spectator coordinate, the complete moment is an average
of tensor-product unitaries under the one-link Wilson probability law and
hence is a contraction.  Coordinate independence gives
\eqref{eq:V54-private-factorization}; submultiplicativity and the local
tree card give \eqref{eq:V54-private-factor-bound}.
\end{proof}

\section{The two resolvent branches}
\label{sec:V54-resolvent-branches}

Let \(\boldsymbol T\) be a nontrivial resolved final output, with mass
\(\Xi_T\).  Retain the established product-gap split
\begin{equation}
 \frac{\Xi_T}{1-q_{\alpha,\boldsymbol T}}
 \le C_{\rm out}
 \begin{cases}
 s_\alpha^{-1},&s_\alpha\Xi_T\le c_0,\\
 \Xi_T,&s_\alpha\Xi_T>c_0.
 \end{cases}
\label{eq:V54-output-split}
\end{equation}

\begin{theorem}[Fixed-order private-terminal scalar compiler]
\label{thm:V54-private-scalar-compiler}
Fix (m\ge2), (D\ge1), and a labeled tree \(\tau\).  Assume (Y>0),
the local tree card \eqref{eq:V54-local-tree-card}, and
\begin{equation}
 \Xi_T\le DY,
 \qquad
 \prod_{i=1}^m\Xi_i^{d_i-1}\le DY^{m-2}.
\label{eq:V54-private-dominance}
\end{equation}
Then
\begin{equation}
 \boxed{
 \frac{\Xi_T}{1-q_{\alpha,\boldsymbol T}}
 \|\mathbf J_{\tau,e}^{(m)}\|_{\rm op}
 \le C_{m,D}L_m
 \prod_{i=1}^ma_i^{d_i}\Xi_i^{1-d_i}.}
\label{eq:V54-private-scalar-compiler}
\end{equation}
The constant is independent of all supports, representations, output
multiplicities, and volume.
\end{theorem}

\begin{proof}
On the low-output branch of \eqref{eq:V54-output-split}, use
\eqref{eq:V54-private-factor-bound} and
\Cref{cor:V52-private-products-all-orders} at order (m-2):
\begin{align*}
 \frac{\Xi_T}{1-q_{\alpha,\boldsymbol T}}
 \|\mathbf J_{\tau,e}^{(m)}\|_{\rm op}
 &\le C L_m q_Ps_\alpha^{m-2}
       \prod_i a_i^{d_i/2}\\
 &=\frac{CL_m}{Y^{m-2}}
   (s_\alpha Y)^{m-2}q_P
   \prod_i a_i^{d_i/2}\\
 &\le\frac{C_mL_m}{Y^{m-2}}
   \prod_i a_i^{d_i/2}.
\end{align*}
The second inequality in \eqref{eq:V54-private-dominance} gives
\[
 Y^{-(m-2)}
 \le D\prod_i\Xi_i^{1-d_i}.
\]
Since (a_i\ge1), one has
\(a_i^{d_i/2}\le a_i^{d_i}\), proving the desired low-output bound.

On the high-output branch, use the first inequality in
\eqref{eq:V54-private-dominance} and the private product moment of order
(m-1):
\begin{align*}
 \frac{\Xi_T}{1-q_{\alpha,\boldsymbol T}}
 \|\mathbf J_{\tau,e}^{(m)}\|_{\rm op}
 &\le CL_m\Xi_Tq_Ps_\alpha^{m-1}
       \prod_i a_i^{d_i/2}\\
 &\le CDL_mYq_Ps_\alpha^{m-1}
       \prod_i a_i^{d_i/2}\\
 &=\frac{CDL_m}{Y^{m-2}}
   (s_\alpha Y)^{m-1}q_P
   \prod_i a_i^{d_i/2}\\
 &\le\frac{C_{m,D}L_m}{Y^{m-2}}
   \prod_i a_i^{d_i/2}.
\end{align*}
Apply the same branching-mass inequality and local-Casimir comparison.
This proves \eqref{eq:V54-private-scalar-compiler} in both cases.
\end{proof}

\begin{remark}[Exact origin of the moment orders]
\label{rem:V54-moment-orders}
The sole final resolvent removes one power of (s_\alpha) on low output,
so the tree's (m-1) powers leave the private moment (m-2).  On high
output the resolvent is bounded by \(\Xi_T\), and private dominance changes
that mass into one factor (Y), requiring moment (m-1).  No moment of
order (m) is used, and the resolvent is not duplicated.
\end{remark}

\begin{corollary}[The quintic private-dominant scalar row is closed]
\label{cor:V54-quintic-private-row}
At (m=5), every selected five-row fibre satisfying the finite-(s) local
tree card and \eqref{eq:V54-private-dominance} obeys its normalized quintic
marked-tree term.  Only the private moments of orders three and four are
used.
\end{corollary}

\begin{proof}
Set (m=5) in \Cref{thm:V54-private-scalar-compiler}.  The two moment
orders are (m-2=3) and (m-1=4), both supplied support-uniformly by
\Cref{cor:V52-private-products-all-orders}.  The right side of
\eqref{eq:V54-private-scalar-compiler} is the resolved fixed-input tree
weight, since
\[
 \prod_i a_i^{d_i}\Xi_i^{1-d_i}
 =\left(\prod_i\Xi_i\right)
  \prod_{ij\in E(\tau)}
  \frac{a_i}{\Xi_i}\frac{a_j}{\Xi_j}.
\]
The exact Fourier factor and raw isotypic pinching then lift the resolved
row exactly as in the already-proved quartic capacity criterion.
\end{proof}

\begin{firewall}[What remains local]
\label{fire:V54-local-remains}
The local tree card \eqref{eq:V54-local-tree-card} is an explicit
hypothesis, not a conclusion of V53.  V53 proves its leading fixed-bank
coefficient only.  Accordingly
\Cref{cor:V54-quintic-private-row} removes the private-spectator, repeated
row-mass, and final-resolvent obstructions conditional on the exact
finite-(s) local operator card; it does not prove that card or global
quintic MTA.
\end{firewall}

\section{Constant-growth audit}
\label{sec:V54-constant-growth}

\begin{proposition}[Fixed-order closure does not imply the all-order
constant]
\label{prop:V54-constant-growth}
Even if \eqref{eq:V54-local-tree-card} is proved for every fixed (m),
the preceding compiler yields the all-order MTA constant only if the
combined local and private constants satisfy an exponential-per-edge
bound.  Finiteness of (L_m) and (C_m^{\rm priv}) separately for each
(m) is insufficient.
\end{proposition}

\begin{proof}
The definition of the all-order marked-tree atlas places a factor
(K_0^{m-1}) before the sum over labeled trees.  The proof of
\Cref{thm:V54-private-scalar-compiler} contributes (L_m) times the
order-(m-2) or order-(m-1) private-product constant.  Therefore a
single atlas constant can follow only from a bound of the form
\[
 L_m\max\{C_{m-2}^{\rm priv},C_{m-1}^{\rm priv}\}
 \le K_0^{m-1}
\]
up to fixed geometric constants.  V52 proves each factor finite at fixed
order but does not prove this growth estimate.  Logical quantification
over each (m) separately cannot supply one uniform (K_0).
\end{proof}

\section{V54 ledger and V55 mandate}
\label{sec:V54-ledger}

\begin{theorem}[Integrated compact-series V54 statement]
\label{thm:V54-integrated}
Preserving compact V1--V53, V54 proves:
\begin{enumerate}[label=\textup{(V54.\arabic*)},leftmargin=6.7em]
\item exact private-coordinate factorization at arbitrary fixed order;
\item a two-branch scalar compiler using private moments (m-2) and
      (m-1), with the final resolvent charged exactly once;
\item conditional closure of the private-dominant quintic row once its
      finite-(s) local tree card is supplied; and
\item the precise exponential constant-growth requirement separating
      fixed-order closure from all-order MTA.
\end{enumerate}
\end{theorem}

\begin{proof}
These are
\Cref{lem:V54-private-factorization,
thm:V54-private-scalar-compiler,
cor:V54-quintic-private-row,
prop:V54-constant-growth}.
\end{proof}

\begin{openproblem}[V55 mandate: audit and the exact quintic local card]
\label{op:V55-mandate}
\begin{enumerate}
\item Perform the mandatory read-only audit of the newest active
      Standard-Model--Einstein and Navier--Stokes notes before choosing the
      proof direction.  Import only type-compatible proved estimates.
\item Return to the exact one-link fifth cumulant operator, not merely the
      fixed-bank expansion \eqref{eq:V53-leading-quintic}.
\item Seek an interpolation or integration-by-parts formula indexed by the
      125 contraction trees and the five density terms of V53, uniform in
      representation height.
\item Prove \eqref{eq:V54-local-tree-card} at (m=5), or isolate one
      resolved high-weight channel for which the proposed finite-(s) tree
      card fails.
\item Test the coherent star and dual-singlet families before inserting
      private spectators through \Cref{thm:V54-private-scalar-compiler}.
\item Keep the all-order constant-growth row behind
      \Cref{prop:V54-constant-growth}; a quintic theorem alone does not
      activate WUFC or CPM.
\end{enumerate}
\end{openproblem}

\begin{firewall}[Claim boundary after compact V54]
\label{fire:V54-current-boundary}
V54 proves a support-uniform fixed-order private-terminal and resolvent
compiler and removes those scalar issues from the embedded quintic
private-dominant branch.  It does not prove the exact finite-(s) quintic
local tree card, close the nonprivate quintic branches, establish an
order-uniform MTA constant, WUFC, CPM, continuum Yang--Mills existence, or
a positive mass gap.  The full Yang--Mills mass gap remains open.
\end{firewall}

\section{Append-only record for compact V54}
\label{sec:V54-record}

V54 was appended to the complete compact V53 file.  The exact V53 parent
had (21314) newline-delimited source records, (712568) bytes, and
SHA--256 digest
\[
\texttt{07a9326115997d5915ce020b0d4dd0ba3c928d1186cf5983eeff3ed6ce775c29}.
\]
No compact V53 mathematical content was changed.  The sole final
\verb|\end{document}| is restored below.

% =============================================================================
% END APPENDED COMPACT VERSION V54 MATERIAL
% =============================================================================

% The compact V54 document terminator is intentionally disabled here:
% \end{document}

% =============================================================================
% BEGIN APPENDED COMPACT VERSION V55 MATERIAL
% =============================================================================

\chapter{V55: Companion audit and the bounded-thermal quintic card}
\label{ch:V55}

\section{Mandatory companion audit}
\label{sec:V55-audit}

Before choosing the V55 direction, the newest active companion notes in
\texttt{latex\_notes} were inspected read-only.  The snapshots were
\begin{center}
\begin{tabular}{p{0.56\textwidth}p{0.36\textwidth}}
\texttt{Renewal\_SM\_Einstein\_Continuum\_Research\_Notes\_V60.tex}
& (22792) lines, (811882) bytes\\
&\scriptsize\ttfamily
8b37f4a70d008bc51daffba63f10d357\newline
52b21b5aff6f66f9c1ebb8d0974cd425\\[0.7em]
\texttt{Renewal\_NS\_Stability\_Research\_Notes\_V31.tex}
& (23052) lines, (887537) bytes\\
&\scriptsize\ttfamily
f1121b62238ad5491bb7cf03635ea9934\newline
942edf573ed8faaa9ee79e1d38a6696
\end{tabular}
\end{center}

SM V60 proves an exact Banach-valued Taylor remainder only after the
complete Ward/species row is assembled.  It explicitly leaves the required
uniform fifth derivative as a hypothesis.  The type-correct lesson for the
present quintic problem is to extract a finite connected Taylor degree after
the complete Wilson cumulant is formed; its unproved low--low derivative
estimate is not a Wilson operator bound and is not imported.

NS V31 proves that flat dyadic probability marks remove stopping-time
multiplicity, but unmarking them costs an exact first radial moment.  Its
sharp forced-mode and packet tests show that a missing moment cannot be
created by heat smoothing alone.  The type-correct lesson is negative but
useful: the repeated link marks in MTA must remain in the estimate until an
actual Casimir moment pays them.  No helical flux, coarea, or Navier--Stokes
source inequality is imported.

\begin{firewall}[Cross-program type boundary]
\label{fire:V55-audit-types}
The audit changes proof order only.  A Banach external-momentum derivative
is not a derivative of a Wilson representation, and a dyadic radial mark is
not a Peter--Weyl link mark.  V55 therefore proves its local estimate from
the Wilson normal form, derived-representation bounds, and connected
Gaussian diagrams already present in the Yang--Mills programme.
\end{firewall}

\section{Thermal variables and a degree-allocation lemma}
\label{sec:V55-thermal-variables}

Fix (L\ge1) and five nontrivial irreducible representations
(R_1,\ldots,R_5).  Put
\begin{equation}
 A_i:=1+C_2(R_i),
 \qquad
 b_i:=\sqrt{s_\alpha A_i},
 \qquad
 b_i^2\le L.
\label{eq:V55-thermal-variables}
\end{equation}
The last inequalities define the bounded-thermal window.

\begin{lemma}[Every connected quintic degree contains a tree degree]
\label{lem:V55-degree-allocation}
Let (t_1,\ldots,t_5\) be positive integers with
\(\sum_it_i\ge8\).  There is a labeled tree
\(\tau\in\mathfrak T_5\) such that
\begin{equation}
 1\le d_i:=\deg_\tau(i)\le t_i
 \qquad(1\le i\le5).
\label{eq:V55-degree-allocation}
\end{equation}
If (0<b_i\le\sqrt L), then
\begin{equation}
 \prod_i b_i^{t_i}
 \le C_{L,\boldsymbol t}
 \prod_i b_i^{d_i}.
\label{eq:V55-monomial-tree-domination}
\end{equation}
At any fixed diagram order, the constant is uniform in the
representations and (s_\alpha).
\end{lemma}

\begin{proof}
The available excess above one is
\(
 \sum_i(t_i-1)\ge3
\).
Capping the contribution of one vertex at three does not reduce the total
below three: if it did, the uncapped total would also be below three unless
one term alone exceeded three, in which case its cap already equals three.
Choose integers
\[
 0\le e_i\le\min\{t_i-1,3\},
 \qquad
 \sum_ie_i=3,
\]
and put (d_i=1+e_i).  Then (1\le d_i\le4),
\(
 \sum_id_i=8=2(5-1)
\), and (d_i\le t_i).  By the Pruefer degree-sequence theorem, these are
the degrees of a labeled tree on five vertices.  Finally,
\[
 \prod_i b_i^{t_i-d_i}
 \le \max\{1,L^{(\sum_i(t_i-d_i))/2}\},
\]
which proves \eqref{eq:V55-monomial-tree-domination}.
\end{proof}

\section{A connected-diagram majorant on a thermal window}
\label{sec:V55-diagram-majorant}

Let
\[
 \mathbf Y_i(s,Z)
 :=D^{R_i}(\exp(\sqrt{s}\,Z))-I_{R_i}.
\]
Constants do not affect a fifth cumulant, so the exact one-link operator is
\begin{equation}
 \mathbf K_5^\alpha(R_1,\ldots,R_5)
 =\kappa_5^{\nu_\alpha}
  (\mathbf Y_1,\ldots,\mathbf Y_5).
\label{eq:V55-exact-local-cumulant}
\end{equation}

\begin{lemma}[Uniform source-series majorant]
\label{lem:V55-source-majorant}
For every (r\ge1), the degree-(r) Taylor coefficient of
\(\mathbf Y_i(s,Z)\) has operator norm at most
\begin{equation}
 \frac{(C b_i\|Z\|)^r}{r!}.
\label{eq:V55-source-majorant}
\end{equation}
Consequently the five source series are absolutely dominated, on the
window \eqref{eq:V55-thermal-variables}, by
\(
 \exp(C_L\|Z\|)
\), uniformly in representation dimensions.
\end{lemma}

\begin{proof}
For an orthonormal Lie-algebra basis (X_a), the Casimir identity is
\[
 \sum_a \mathrm dR_i(X_a)^*\mathrm dR_i(X_a)
 =C_2(R_i)I.
\]
Thus every generator has norm at most \(C_2(R_i)^{1/2}\), and every
ordered product of (r) generators has norm at most
\(C^rA_i^{r/2}\).  Multiplication by the exponential-coordinate Taylor
factor (s^{r/2}/r!) gives
\eqref{eq:V55-source-majorant}.  Sum the exponential series and use
(b_i\le\sqrt L).
\end{proof}

In the Wilson normal form, a nonquadratic action vertex of even valence
(2\ell\) carries (s^{\ell-1}); Haar-Jacobian and normalization vertices
carry at least the same parabolic weight.  Split products in the exponential
into separate unmarked vertices.

\begin{lemma}[Thermal domination of every connected diagram]
\label{lem:V55-connected-diagram}
Consider a connected Gaussian contraction diagram meeting all five source
vertices.  Suppose its source Taylor degrees are (r_i\ge1), and the total
explicit small-noise power contributed by its unmarked density vertices is
(s^p), (p\in\mathbb N_0).  Then
\begin{equation}
 p+\frac12\sum_i r_i\ge4.
\label{eq:V55-connected-degree}
\end{equation}
On the thermal window, the norm of the diagram is bounded by
\begin{equation}
 C_L\prod_i b_i^{\deg_\tau(i)}
\label{eq:V55-diagram-tree-bound}
\end{equation}
for at least one labeled tree 
\(\tau\in\mathfrak T_5\).  The tree can be chosen deterministically from
the diagram.
\end{lemma}

\begin{proof}
If an unmarked action vertex has valence (2\ell), charge it
(s^{\ell-1}); a better Jacobian power may be weakened to this one.  If
the connected diagram has (E) Gaussian contraction edges and (V_u)
unmarked vertices, its parabolic degree (D) satisfies
\[
 D=p+\frac12\sum_i r_i\ge E-V_u.
\]
The graph has (5+V_u) vertices and is connected, so
\(E\ge5+V_u-1\).  This proves
\eqref{eq:V55-connected-degree}.

Choose nonnegative integers (q_i) with
\(
 \sum_iq_i=2p
\).  Since (A_i\ge1), one has (b_i\ge\sqrt{s}), and hence
\[
 s^p\prod_i b_i^{r_i}
 \le\prod_i b_i^{r_i+q_i}.
\]
The integers (t_i=r_i+q_i) are positive and have total at least eight by
\eqref{eq:V55-connected-degree}.  Apply
\Cref{lem:V55-degree-allocation}; the surplus bounded thermal powers are
absorbed into (C_L).  The Gaussian contractions and fixed density-vertex
coefficients contribute only another (L)-dependent constant, proving
\eqref{eq:V55-diagram-tree-bound}.  Ordering all labeled trees once makes
the first admissible choice deterministic.
\end{proof}

\begin{theorem}[Bounded-thermal finite-(s) quintic tree card]
\label{thm:V55-thermal-quintic-card}
For every (L\ge1), there are (C_L<\infty) and \(\alpha_L\) such that,
whenever \eqref{eq:V55-thermal-variables} holds and
\(\alpha\ge\alpha_L\), the exact local fifth-cumulant operator admits a
decomposition
\begin{equation}
 \boxed{
 \mathbf K_5^\alpha
 =\sum_{\tau\in\mathfrak T_5}\mathbf K_{5,\tau}^\alpha,
 \qquad
 \|\mathbf K_{5,\tau}^\alpha\|_{\rm op}
 \le C_Ls_\alpha^4
 \prod_{i=1}^5A_i^{\deg_\tau(i)/2}.}
\label{eq:V55-thermal-quintic-card}
\end{equation}
The bound is uniform in representation dimensions and all final
Peter--Weyl multiplicities.
\end{theorem}

\begin{proof}
Use the Wilson Laplace normal form on a fixed exponential neighbourhood,
to depth five.  Expand the five source operators by their absolutely
convergent series from \Cref{lem:V55-source-majorant}, expand each density
coefficient into its finitely many Gaussian polynomial vertices, and form
the complete fifth cumulant before taking norms.  Wick expansion followed
by the cumulant Möbius sum cancels every diagram whose induced graph on the
five marked sources is disconnected.  Assign each surviving connected
diagram to the deterministic tree supplied by
\Cref{lem:V55-connected-diagram}.

The source exponential majorant is integrable against a slightly weakened
Gaussian, uniformly for (b_i\le\sqrt L).  The normal-form remainder has a
fixed polynomial times the same weakened Gaussian.  Consequently the
absolute sum of all diagrams assigned to one tree is bounded by one
constant (C_L), with no representation-dimension factor.  Since
\[
 \prod_i b_i^{\deg_\tau(i)}
 =s_\alpha^{\frac12\sum_i\deg_\tau(i)}
  \prod_iA_i^{\deg_\tau(i)/2}
 =s_\alpha^4\prod_iA_i^{\deg_\tau(i)/2},
\]
the local part gives the stated bound.

On the complement of the exponential neighbourhood, the Wilson-well gap
gives (Ce^{-c/s_\alpha}) times the universal fifth-cumulant norm.  Every
tree monomial on the right is at least (s_\alpha^4), because (A_i\ge1).
Thus exponential domination absorbs the complement into the same estimate.
The assignment partitions the complete connected diagram sum, proving the
decomposition.  Orthogonal output projection and multiplicity compression
are contractions, so every resolved final block inherits the bound without
an output count.
\end{proof}

\begin{assessment}[Consistency with the exact V53 coefficient]
\label{ass:V55-V53-consistency}
At (s_\alpha=0), the degree-four diagrams in the proof of
\Cref{thm:V55-thermal-quintic-card} are exactly the 125 Gaussian source
trees and five density terms of \eqref{eq:V53-leading-quintic}.  Higher
diagrams are not estimated by fixed-bank remainder notation; their surplus
thermal powers are retained and absorbed only through
\(s_\alpha A_i\le L\).  Thus the proof is uniform in representation height
inside the stated window and passes the coherent valence regression.
\end{assessment}

\begin{corollary}[Bounded-thermal embedded private quintic closure]
\label{cor:V55-thermal-private-quintic}
For every fixed (L,D\), a selected five-row junction satisfying
\eqref{eq:V55-thermal-variables} and the private-dominance conditions
\eqref{eq:V54-private-dominance} obeys its complete normalized quintic
marked-tree contribution, uniformly in supports, representations,
multiplicities, and volume.
\end{corollary}

\begin{proof}
Theorem~\ref{thm:V55-thermal-quintic-card} supplies the local tree card
\eqref{eq:V54-local-tree-card} with (m=5) and (L_5=C_L).  Apply
\Cref{thm:V54-private-scalar-compiler} to each of its 125 tree fibres, then
use exact Fourier normalization and raw multiplicity pinching.
\end{proof}

\section{V55 ledger and V56 mandate}
\label{sec:V55-ledger}

\begin{theorem}[Integrated compact-series V55 statement]
\label{thm:V55-integrated}
Preserving compact V1--V54, V55 proves:
\begin{enumerate}[label=\textup{(V55.\arabic*)},leftmargin=6.7em]
\item the scheduled hashed SM V60 and NS V31 audit and its no-import
      decision;
\item the quintic degree-allocation lemma;
\item a representation-dimension-free thermal source-series bound;
\item deterministic domination of every connected quintic diagram by one
      tree-valence monomial;
\item the exact finite-(s) quintic local tree card on every fixed thermal
      window; and
\item the corresponding embedded private-dominant quintic closure.
\end{enumerate}
\end{theorem}

\begin{proof}
These are
\Cref{sec:V55-audit,
lem:V55-degree-allocation,
lem:V55-source-majorant,
lem:V55-connected-diagram,
thm:V55-thermal-quintic-card,
cor:V55-thermal-private-quintic}.
\end{proof}

\begin{openproblem}[V56 mandate: thermally escaping quintic inputs]
\label{op:V56-mandate}
\begin{enumerate}
\item Split the complement of \eqref{eq:V55-thermal-variables} by the set
      of rows with (s_\alpha A_i>L\).  Keep the complete fifth cumulant and
      final output resolution assembled.
\item Use the all-order one-link failure debits of V52 only on multipliers
      actually present in a resolved moment term; do not transfer damping
      across a recoupling matrix without proof.
\item Prove that a single hot input either remains hot in the final output
      or forces a paired high-to-low fusion.  In the latter case retain the
      complete formation multiplicity before summing outputs.
\item Test the hot star on the coherent family of
      \Cref{prop:V53-bare-quintic-false}; its center needs four valence
      powers, not one input Casimir.
\item Close one hot-output or paired-cancellation branch, or return the
      first exact recoupling/formation ratio which lacks a tree mark.
\item Do not invoke global quintic MTA until all minimal connected support
      hypergraphs, not only a single embedded five-row junction, are
      controlled.
\end{enumerate}
\end{openproblem}

\begin{firewall}[Claim boundary after compact V55]
\label{fire:V55-current-boundary}
V55 proves the exact finite-(s) local fifth-cumulant tree card on every
fixed thermal window and closes its embedded private-dominant branch.  It
does not control thermally escaping quintic representations, all mixed
five-row support skeletons, global quintic or all-order MTA, WUFC, CPM,
continuum Yang--Mills existence, or a positive mass gap.  The full
Yang--Mills mass gap remains open.
\end{firewall}

\section{Append-only record for compact V55}
\label{sec:V55-record}

V55 was appended to the complete compact V54 file.  The exact V54 parent
had (21645) newline-delimited source records, (724790) bytes, and
SHA--256 digest
\[
\texttt{8b92f8f137d007cb0479e99e7a574a42ea09956833cbbd3ce05c92ae6ce469c2}.
\]
No compact V54 mathematical content was changed.  The sole final
\verb|\end{document}| is restored below.

% =============================================================================
% END APPENDED COMPACT VERSION V55 MATERIAL
% =============================================================================

% The compact V55 document terminator is intentionally disabled here:
% \end{document}

% =============================================================================
% BEGIN APPENDED COMPACT VERSION V56 MATERIAL
% =============================================================================

\chapter{V56: Hot-star displacement and the uniform one-link quintic card}
\label{ch:V56}

\section{The exact displacement factor}
\label{sec:V56-displacement}

V55 proves the local fifth-cumulant tree card when every thermal variable
\(b_i=\sqrt{s_\alpha A_i}\) is bounded by one.  Its complement contains a
hot row (i) with (b_i\ge1).  Instead of resolving every hot fusion
channel, one may keep the complete cumulant and use its invariance under
constants.  Four exact representation displacements then pay the four
leaves of the star centered at the hot row.

For a nontrivial irreducible (R), set
\begin{equation}
 b_R:=\sqrt{s_\alpha(1+C_2(R))},
 \qquad
 c_R:=\min\{1,b_R\}.
\label{eq:V56-bc}
\end{equation}

\begin{lemma}[One-link representation displacement]
\label{lem:V56-representation-displacement}
There are (C<\infty) and \(\alpha_0\) such that, for every nontrivial
irreducible (R),
\begin{equation}
 \boxed{
 \mathbb E_{\nu_\alpha}
 \|D^R(U)-I_R\|_{\rm op}^k
 \le C_k c_R^k}
\label{eq:V56-displacement-moment}
\end{equation}
for each fixed integer (k\ge1).  More generally, for finitely many
representations (R_j),
\begin{equation}
 \mathbb E_{\nu_\alpha}
 \prod_{j\in J}\|D^{R_j}(U)-I_{R_j}\|_{\rm op}
 \le C_{|J|}\prod_{j\in J}c_{R_j}.
\label{eq:V56-mixed-displacement}
\end{equation}
The constants are independent of representation dimensions.
\end{lemma}

\begin{proof}
On a fixed exponential neighbourhood write
\(U=\exp(\sqrt{s_\alpha}Z)\).  Since the derived representation is
skew-adjoint,
\[
 \|D^R(U)-I_R\|_{\rm op}
 \le\min\{2,\sqrt{s_\alpha}\|\mathrm dR(Z)\|_{\rm op}\}
 \le\min\{2,Cb_R\|Z\|\}.
\]
The last inequality follows from the Casimir identity exactly as in
\Cref{lem:V55-source-majorant}.  For (b\le1), the right side is at most
\(Cb(1+\|Z\|)\); for (b>1), it is at most
\(2=2c_R\).  Hence, uniformly in (R),
\begin{equation}
 \|D^R(U)-I_R\|_{\rm op}
 \le Cc_R(1+\|Z\|)
\label{eq:V56-pointwise-displacement}
\end{equation}
on the normal neighbourhood.  The rescaled radial moments are uniformly
finite by \Cref{lem:V52-Laplace-moments}, so Hölder proves both displayed
estimates there.

Outside the neighbourhood the displacement is at most two and the
normalized Wilson mass is (O(e^{-c/s_\alpha})).  Because every nontrivial
Casimir has a positive floor and (s_\alpha\le1), one has
\(c_R\ge c_G\sqrt{s_\alpha}\).  Exponential domination therefore gives
\[
 e^{-c/s_\alpha}
 \le C_k\prod_{j\in J}c_{R_j}
\]
for every fixed nonempty (J).  This absorbs the complement and proves the
lemma.
\end{proof}

\section{Cumulant products retain four leaf displacements}
\label{sec:V56-cumulant-products}

Let
\[
 \mathbf X_i(U):=D^{R_i}(U),
 \qquad
 \mathbf Y_i(U):=\mathbf X_i(U)-I_{R_i},
 \qquad 1\le i\le5,
\]
with every operator placed on its own tensor leg.

\begin{lemma}[Hot-centered complete cumulant bound]
\label{lem:V56-hot-cumulant}
Fix an index (i\).  The complete one-link fifth cumulant satisfies
\begin{equation}
 \boxed{
 \|\kappa_5^{\nu_\alpha}
   (\mathbf X_1,\ldots,\mathbf X_5)\|_{\rm op}
 \le C\prod_{j\ne i}c_{R_j}.}
\label{eq:V56-hot-cumulant}
\end{equation}
The bound is uniform in all representations and multiplicities.
\end{lemma}

\begin{proof}
A joint cumulant of order at least two is invariant under adding a constant
to any one argument.  Thus the left side equals
\(
 \kappa_5(\mathbf Y_1,\ldots,\mathbf Y_5)
\).
Expand it by the moment--cumulant formula:
\begin{equation}
 \kappa_5(\mathbf Y_1,\ldots,\mathbf Y_5)
 =\sum_{\pi\in\Pi_5}
 (|\pi|-1)!(-1)^{|\pi|-1}
 \bigotimes_{B\in\pi}
 \mathbb E_{\nu_\alpha}\bigotimes_{j\in B}\mathbf Y_j.
\label{eq:V56-moment-cumulant}
\end{equation}
For a block (B), submultiplicativity and
\Cref{lem:V56-representation-displacement} give
\[
 \left\|\mathbb E\bigotimes_{j\in B}\mathbf Y_j\right\|_{\rm op}
 \le C_B\prod_{j\in B\setminus\{i\}}c_{R_j};
\]
if (i\notin B), the set difference does nothing, while if (B=\{i\})
the empty product is one.  Multiplying over the blocks of \(\pi\) uses
each (j\ne i) exactly once.  The fifth partition lattice is finite, so
the sum of its numerical Möbius coefficients is absorbed into (C).
No output decomposition has been taken, hence no multiplicity count occurs.
\end{proof}

\section{The thermally escaping star}
\label{sec:V56-hot-star}

For a labeled tree \(\tau\) on five vertices, write
\begin{equation}
 \mathfrak w_\tau
 :=s_\alpha^4\prod_{j=1}^5A_j^{\deg_\tau(j)/2}
 =\prod_{j=1}^5b_j^{\deg_\tau(j)}.
\label{eq:V56-b-tree-weight}
\end{equation}

\begin{theorem}[Hot-row quintic star card]
\label{thm:V56-hot-star-card}
If (b_i\ge1) for at least one row (i), then the exact local fifth
cumulant obeys
\begin{equation}
 \boxed{
 \|\mathbf K_5^\alpha(R_1,\ldots,R_5)\|_{\rm op}
 \le C\mathfrak w_{\tau_i},}
\label{eq:V56-hot-star-card}
\end{equation}
where \(\tau_i\) is the labeled star centered at (i).  Thus the entire
hot contribution may be assigned to one legitimate contraction-refined
tree.
\end{theorem}

\begin{proof}
By \Cref{lem:V56-hot-cumulant}, the left side is at most
\(C\prod_{j\ne i}c_{R_j}\).  Since (b_i\ge1) and
\(c_{R_j}\le b_j\),
\[
 \prod_{j\ne i}c_{R_j}
 \le b_i^4\prod_{j\ne i}b_j
 =\mathfrak w_{\tau_i}.
\]
This proves \eqref{eq:V56-hot-star-card}.  The complete operator is assigned
only after its cumulant sum \eqref{eq:V56-moment-cumulant} has been formed,
so the construction does not multiply separately cancelled capacities.
\end{proof}

\begin{theorem}[Representation-uniform one-link quintic tree card]
\label{thm:V56-uniform-one-link-quintic}
There are (C_5<\infty) and \(\alpha_5\) such that, for all five-tuples of
nontrivial irreducible representations and all
\(\alpha\ge\alpha_5\),
\begin{equation}
 \boxed{
 \mathbf K_5^\alpha
 =\sum_{\tau\in\mathfrak T_5}\mathbf K_{5,\tau}^\alpha,
 \qquad
 \|\mathbf K_{5,\tau}^\alpha\|_{\rm op}
 \le C_5s_\alpha^4
 \prod_iA_i^{\deg_\tau(i)/2}.}
\label{eq:V56-uniform-one-link-quintic}
\end{equation}
Every resolved output compression inherits the same estimate with no
dimension or multiplicity factor.
\end{theorem}

\begin{proof}
If all (b_i<1), apply
\Cref{thm:V55-thermal-quintic-card} with (L=1).  If some (b_i\ge1),
choose the least such label, assign the complete cumulant to its star by
\Cref{thm:V56-hot-star-card}, and set all other tree fibres to zero.  These
two alternatives are exhaustive.  Orthogonal final-channel projection and
the established partial-trace compression are contractions, proving the
last assertion.
\end{proof}

\begin{corollary}[Uniform embedded private-dominant five-junction]
\label{cor:V56-uniform-private-five-junction}
A selected five-row coordinate junction satisfying the private-dominance
conditions \eqref{eq:V54-private-dominance} obeys its complete normalized
quintic marked-tree contribution for arbitrary representation heights.
The bound is uniform in supports, final outputs, multiplicities, and volume.
\end{corollary}

\begin{proof}
Use \Cref{thm:V56-uniform-one-link-quintic} as the local tree card in
\Cref{thm:V54-private-scalar-compiler}, then apply exact Fourier
normalization and raw isotypic pinching.  V52 supplies the required private
moments of orders three and four.
\end{proof}

\begin{assessment}[Why no hot fusion tail was needed]
\label{ass:V56-no-tail-needed}
The route proposed at the end of V55 would have tracked which fused output
retains a hot Casimir.  That is unnecessary for the one-coordinate local
card.  Once one row is hot, the four leaf displacements already fit under
the centered star weight.  This avoids a recoupling-dependent transfer of
Wilson damping and is stronger in type: the estimate is taken on the
complete cumulant operator before any output resolution.
\end{assessment}

\section{V56 ledger and V57 mandate}
\label{sec:V56-ledger}

\begin{theorem}[Integrated compact-series V56 statement]
\label{thm:V56-integrated}
Preserving compact V1--V55, V56 proves:
\begin{enumerate}[label=\textup{(V56.\arabic*)},leftmargin=6.7em]
\item uniform mixed moments of exact representation displacements;
\item a complete fifth-cumulant bound retaining four chosen leaf
      displacements;
\item the hot-row star card for every thermally escaping tuple;
\item the representation-uniform one-link quintic tree card, by joining
      the hot theorem to V55 at (s_\alpha A_i=1); and
\item the corresponding arbitrary-height embedded private-dominant
      five-junction closure.
\end{enumerate}
\end{theorem}

\begin{proof}
These are
\Cref{lem:V56-representation-displacement,
lem:V56-hot-cumulant,
thm:V56-hot-star-card,
thm:V56-uniform-one-link-quintic,
cor:V56-uniform-private-five-junction}.
\end{proof}

\begin{openproblem}[V57 mandate: mixed-support quintic skeletons]
\label{op:V57-mandate}
\begin{enumerate}
\item Return to the exact connected replica-partition formula.  Classify
      inclusion-minimal connected hypergraphs on five row labels without
      expanding their local Wilson cumulants into absolute partition sums.
\item Prove a deterministic leaf-hyperedge or two-core reduction which
      exposes either one local five-row block, a lower-order connected block
      joined by a pair edge, or an irreducible overlapping core.
\item Insert \Cref{thm:V56-uniform-one-link-quintic} for a five-row block,
      the global binary/ternary/quartic MTA theorems for lower blocks, and
      keep the final resolvent on exactly one factor.
\item Treat repeated-coordinate overlaps before multiplying capacities;
      the quartic same-coordinate regressions remain binding.
\item Close one exhaustive mixed-support class or isolate the first
      five-row core for which lower-order MTA does not tensor through the
      single resolvent.
\item Keep the constant-growth requirement of
      \Cref{prop:V54-constant-growth} explicit; quintic closure would still
      not prove all-order MTA.
\end{enumerate}
\end{openproblem}

\begin{firewall}[Claim boundary after compact V56]
\label{fire:V56-current-boundary}
V56 proves the local fifth-cumulant tree card uniformly in all
representations and closes the selected private-dominant five-junction.
It does not yet assemble all mixed-support quintic skeletons, prove global
quintic or all-order MTA, WUFC, CPM, continuum Yang--Mills existence, or a
positive mass gap.  The full Yang--Mills mass gap remains open.
\end{firewall}

\section{Append-only record for compact V56}
\label{sec:V56-record}

V56 was appended to the complete compact V55 file.  The exact V55 parent
had (22019) newline-delimited source records, (739185) bytes, and
SHA--256 digest
\[
\texttt{318588551bfca5c272f622fc571dc087539f751d70712a09705205dabb0568e4}.
\]
No compact V55 mathematical content was changed.  The sole final
\verb|\end{document}| is restored below.

% =============================================================================
% END APPENDED COMPACT VERSION V56 MATERIAL
% =============================================================================

% The compact V56 document terminator is intentionally disabled here:
% \end{document}

% =============================================================================
% BEGIN APPENDED COMPACT VERSION V57 MATERIAL
% =============================================================================

\chapter{V57: The quintic one-payer forest compiler}
\label{ch:V57}

\section{Purpose}
\label{sec:V57-purpose}

V56 closes the complete one-coordinate fifth cumulant.  A mixed-support
quintic replica term, however, may be a product of connected blocks on a
row partition.  Applying a response theorem to every block would insert
several output resolvents.  This chapter proves the exact replacement:
one block pays, all other blocks use numerator bounds, and the missing
connections form a quotient tree.  The internal forests and quotient tree
then lift to one ordinary five-row tree with the correct power of
(s_\alpha).

\section{Response-or-numerator table through order four}
\label{sec:V57-table}

For a nonempty row set (B\subseteq[5]), let \(\mathcal N_B\) denote the
complete positive envelope of its connected Wilson numerator before a
resolvent, and let \(\mathcal P_B\) denote the corresponding envelope with
one output mass and its resolvent.  All output projections and multiplicity
spaces are retained.  For (2\le|B|\le4), put
\begin{equation}
 \mathcal T_B(\theta)
 :=\sum_{\tau\in\mathfrak T(B)}
   \prod_{ij\in E(\tau)}\theta_{ij}.
\label{eq:V57-block-tree-polynomial}
\end{equation}

\begin{theorem}[Lower-order payer table through quartic order]
\label{thm:V57-lower-table}
Uniformly in supports, representations, output multiplicities, and volume,
\begin{align}
 \frac{\|\mathcal N_B\|}{\prod_{i\in B}\Xi_i}
 &\le C^{|B|-1}s_\alpha^{|B|-1}\mathcal T_B(\theta),
 &&2\le|B|\le4,
\label{eq:V57-unpayer-table}\\
 \frac{\|\mathcal P_B\|}{\prod_{i\in B}\Xi_i}
 &\le C^{|B|-1}s_\alpha^{|B|-2}\mathcal T_B(\theta),
 &&2\le|B|\le4.
\label{eq:V57-payer-table}
\end{align}
For a singleton,
\begin{equation}
 \|\mathcal N_{\{i\}}\|\le1,
 \qquad
 \frac{\|\mathcal P_{\{i\}}\|}{\Xi_i}\le C/s_\alpha.
\label{eq:V57-singleton-table}
\end{equation}
Every norm estimate also bounds the product-state capacity used in a
one-sided tensor product.
\end{theorem}

\begin{proof}
For block sizes two and three these are exactly the response-or-numerator
bounds proved in V37.  For size four, the unpayed estimate is the complete
connected quartic numerator bound retained from the frozen f1 archive,
after exact Fourier normalization and division by the four row masses.
Its three normalized links give the factor
\(s_\alpha^3\mathcal T_B\).  The paid estimate is the global quartic MTA
proved in \Cref{cor:V52-global-quartic-MTA}; the sole output resolvent
removes one power of (s_\alpha).  The singleton statements are complete
moment contraction and \Cref{lem:V3-spectator-payer}.  One-sided tensor
capacity is no larger than operator norm on all unselected factors.
\end{proof}

\section{Exactly one payer on a row partition}
\label{sec:V57-one-payer}

Let \(\pi\in\Pi_5\) have (k=|\pi|\), with (2\le k\le4).  Resolve the
product of the connected outputs of its blocks.  Fusion assigns the final
output mass to one block (P\in\pi); use \(\mathcal P_P\) there and
\(\mathcal N_B\) on every (B\ne P).  Sum the resulting positive envelopes
over all possible payer blocks and call the result
\(\mathcal E_\pi^{\rm pay}\).

\begin{lemma}[One-payer partition-product bound]
\label{lem:V57-one-payer-partition}
For (2\le|\pi|=k\le4),
\begin{equation}
 \boxed{
 \frac{\kappa_\otimes(\mathcal E_\pi^{\rm pay})}
      {\prod_{i=1}^5\Xi_i}
 \le C^4s_\alpha^{4-k}
 \prod_{\substack{B\in\pi\\|B|\ge2}}
 \mathcal T_B(\theta).}
\label{eq:V57-one-payer-partition}
\end{equation}
No capacity product and no second resolvent occurs.
\end{lemma}

\begin{proof}
Fix the payer block.  Apply the one-sided tensor theorem to the positive
block envelopes: capacity is used on the payer and ordinary operator norms
on all other blocks.  If the payer is nonsingleton, the total exponent of
(s_\alpha) is
\[
 (|P|-2)+\sum_{B\ne P}(|B|-1)
 =\sum_{B\in\pi}(|B|-1)-1
 =(5-k)-1=4-k.
\]
If the payer is a singleton, its exponent is (-1), and the identical
calculation holds.  The row masses multiply to \(\prod_i\Xi_i\), because
the blocks partition the row set.  Insert
\Cref{thm:V57-lower-table} and sum the at most four payer choices.  This
proves \eqref{eq:V57-one-payer-partition}.
\end{proof}

\begin{firewall}[The discrete partition remains assembled]
\label{fire:V57-discrete-partition}
The case (k=5) is deliberately absent.  Five singleton means with one
payer would carry (s_\alpha^{-1}) and no spatial edge.  In a connected
fifth cumulant that term cancels against the other moment partitions.  It
must remain inside the complete cumulant and may not be normed as an
independent positive correction.
\end{firewall}

\section{Lifting a quotient tree to a five-row tree}
\label{sec:V57-quotient-lift}

For two distinct blocks (B,C\in\pi), define
\begin{equation}
 \Theta_{BC}:=\sum_{i\in B,\,j\in C}\theta_{ij}.
\label{eq:V57-quotient-edge}
\end{equation}

\begin{lemma}[Forest plus quotient tree is an ordinary tree]
\label{lem:V57-forest-quotient-tree}
Fix an internal labeled tree \(\tau_B\) on every nonsingleton block
(B\in\pi), and a labeled tree \(\zeta\) on the (k) quotient vertices
(B\in\pi).  Expand every quotient factor
\(\Theta_{BC}\) by choosing one original edge (ij), (i\in B,j\in C).
For every such choice,
\begin{equation}
 \left(\bigcup_{B\in\pi}E(\tau_B)\right)
 \cup\{ij:BC\in E(\zeta)\}
\label{eq:V57-lifted-edge-set}
\end{equation}
is a labeled tree on ([5]).  Consequently,
\begin{align}
 &\left(
 \prod_{\substack{B\in\pi\\|B|\ge2}}
 \mathcal T_B(\theta)
 \right)
 \left(
 \sum_{\zeta\in\mathfrak T(\pi)}
 \prod_{BC\in E(\zeta)}\Theta_{BC}
 \right)
\notag\\
 &\qquad\le C_5
 \sum_{\tau\in\mathfrak T_5}
 \prod_{ij\in E(\tau)}\theta_{ij}.
\label{eq:V57-forest-quotient-lift}
\end{align}
\end{lemma}

\begin{proof}
The internal block trees form a forest with (k) connected components and
\(
 \sum_B(|B|-1)=5-k
\)
edges.  Each expanded quotient edge joins two different components, and
the quotient edges form a tree, so their (k-1) lifts connect the forest
without a cycle.  The union is connected and has
\((5-k)+(k-1)=4\) edges on five vertices; hence it is a tree.  Expanding
the two finite sums gives only lifted five-row tree monomials.  A fixed
five-row tree can have only a finite number, depending on five alone, of
preimages under these choices.  Absorb that number into (C_5).
\end{proof}

\begin{theorem}[Quintic one-payer forest composition]
\label{thm:V57-forest-composition}
Suppose a connected quintic identity produces, for a row partition
\(\pi\) with (2\le k\le4), the normalized quotient-tree factor
\begin{equation}
 C^{k-1}s_\alpha^{k-1}
 \sum_{\zeta\in\mathfrak T(\pi)}
 \prod_{BC\in E(\zeta)}\Theta_{BC}
\label{eq:V57-quotient-tree-factor}
\end{equation}
before its block outputs are normed.  After exact one-resolvent allocation,
all lower-order block products belonging to this identity obey
\begin{equation}
 \boxed{
 C^4s_\alpha^3
 \sum_{\tau\in\mathfrak T_5}
 \prod_{ij\in E(\tau)}\theta_{ij}.}
\label{eq:V57-composed-quintic-tree}
\end{equation}
Thus the final resolvent and every lower-order partition product are paid
with exactly the three (s_\alpha)-powers required after a fifth-order
connected numerator loses one power to response.
\end{theorem}

\begin{proof}
Multiply \eqref{eq:V57-one-payer-partition} by
\eqref{eq:V57-quotient-tree-factor}.  The small-noise exponent is
\[
 (4-k)+(k-1)=3.
\]
Apply \Cref{lem:V57-forest-quotient-tree} to the normalized edge factors.
Only one block used \(\mathcal P_B\), so the final resolvent occurs once.
All other factors used numerator bounds.  This proves
\eqref{eq:V57-composed-quintic-tree}.
\end{proof}

\begin{assessment}[The remaining mixed-support gate]
\label{ass:V57-remaining-gate}
V57 removes the final-resolvent and lower-order partition-product issues
from quintic assembly.  What remains is now exact: for every connected
mixed-support replica fibre, one must expose the quotient-tree factor
\eqref{eq:V57-quotient-tree-factor} while the local coordinate cumulants
are still assembled.  Once that factor is present, the internal binary,
ternary, and quartic blocks compile automatically by
\Cref{thm:V57-forest-composition}.  The all-singleton term cannot be
separated and remains part of this connected extraction.
\end{assessment}

\section{V57 ledger and V58 mandate}
\label{sec:V57-ledger}

\begin{theorem}[Integrated compact-series V57 statement]
\label{thm:V57-integrated}
Preserving compact V1--V56, V57 proves:
\begin{enumerate}[label=\textup{(V57.\arabic*)},leftmargin=6.7em]
\item the response-or-numerator table through global quartic order;
\item exact allocation of one and only one output payer on every
      nontrivial row partition of five;
\item the exponent identity (s_\alpha^{4-k}s_\alpha^{k-1}=s_\alpha^3);
\item the forest-plus-quotient-tree lift to an ordinary five-row tree; and
\item the quintic one-payer forest compiler
      \eqref{eq:V57-composed-quintic-tree}.
\end{enumerate}
\end{theorem}

\begin{proof}
These are
\Cref{thm:V57-lower-table,
lem:V57-one-payer-partition,
lem:V57-forest-quotient-tree,
thm:V57-forest-composition}.
\end{proof}

\begin{openproblem}[V58 mandate: connected quotient-tree extraction]
\label{op:V58-mandate}
\begin{enumerate}
\item Work in the exact five-row incidence replica formula and keep the
      all-singleton cancellation assembled.
\item Contract every already-controlled local block of size at most four
      to one quotient vertex, but retain its coordinate and output carrier.
\item Prove the quotient connected operator has a marked-tree expansion
      with the factor \eqref{eq:V57-quotient-tree-factor}; do not obtain it
      by multiplying separately normed weak cuts.
\item Insert the uniform local five-row card of
      \Cref{thm:V56-uniform-one-link-quintic} whenever one coordinate block
      meets all five rows.
\item Classify the residual overlapping cores in which no five-row block
      occurs and no lower block is a leaf in the quotient incidence graph.
\item Close the leaf-reducible class with
      \Cref{thm:V57-forest-composition}, or isolate the first irreducible
      core and its exact missing normalized mark.
\end{enumerate}
\end{openproblem}

\begin{firewall}[Claim boundary after compact V57]
\label{fire:V57-current-boundary}
V57 proves the complete lower-order one-payer compiler required for
mixed-support quintic assembly.  It does not yet prove the connected
quotient-tree factor for every replica fibre, global quintic or all-order
MTA, WUFC, CPM, continuum Yang--Mills existence, or a positive mass gap.
The full Yang--Mills mass gap remains open.
\end{firewall}

\section{Append-only record for compact V57}
\label{sec:V57-record}

V57 was appended to the complete compact V56 file.  The exact V56 parent
had (22325) newline-delimited source records, (750396) bytes, and
SHA--256 digest
\[
\texttt{db63334bca4065b108e05ba080e467c4ddb4ec807ac75694c7db1254cd9a82e9}.
\]
No compact V56 mathematical content was changed.  The sole final
\verb|\end{document}| is restored below.

% =============================================================================
% END APPENDED COMPACT VERSION V57 MATERIAL
% =============================================================================

% The compact V57 document terminator is intentionally disabled here:
% \end{document}

% =============================================================================
% BEGIN APPENDED COMPACT VERSION V58 MATERIAL
% =============================================================================

\chapter{V58: Exact five-row graph-difference algebra}
\label{ch:V58}

\section{Purpose and upstream correction}
\label{sec:V58-purpose}

The V57 compiler is conditional on exposing a quotient-tree factor while
the complete fifth cumulant is still assembled.  There is a finite
algebraic question upstream of every Wilson estimate: can the fifth
cumulant itself be written as a signed sum of tree merge-differences?
The four-row calculation in the frozen programme warned that the answer
need not be yes.  This chapter settles the five-row question exactly.

Trees alone again fail.  Nevertheless, the complete fifth cumulant is a
linear combination of six symmetric graph-difference orbits: the three
five-vertex tree types and three unicyclic types.  Thus every term exposes
four connectivity differences before any norm is taken, and the only
algebra beyond a tree is one surplus edge.  This is the precise row-level
replacement for the false tree-only ansatz.

\section{Partition-lattice merge differences}
\label{sec:V58-partition-algebra}

Write \(\Pi_5\) for the partition lattice of \([5]\), with discrete
partition \(\widehat 0\) and one-block partition \(\widehat 1\).  For an
edge \(a=\{i,j\}\), let \(\sigma_a\) be the partition whose only
nonsingleton block is \(\{i,j\}\), and set
\[
 J_a\pi:=\pi\vee\sigma_a,
 \qquad \nabla_a:=J_a-I.
\]
Equivalently, \(J_a\pi\) merges the two blocks of \(\pi\) containing
\(i\) and \(j\), doing nothing when they are already the same block.
The maps \(J_a\), and hence the differences \(\nabla_a\), commute.

Let \(M:\Pi_5\to\mathcal V\) be an arbitrary moment datum with values in
a rational vector space.  If \(G\) is a simple graph on \([5]\), define
\begin{equation}
 D_GM:=\prod_{a\in E(G)}\nabla_aM(\widehat0)
 =\sum_{F\subseteq E(G)}(-1)^{|E(G)|-|F|}M(\pi(F)),
\label{eq:V58-graph-difference}
\end{equation}
where \(\pi(F)\) is the component partition of the spanning graph
\(([5],F)\).  Define the fifth cumulant functional by
\begin{equation}
 \Lambda_5M
 :=\sum_{\pi\in\Pi_5}
 (-1)^{|\pi|-1}(|\pi|-1)!\,M(\pi).
\label{eq:V58-fifth-cumulant-functional}
\end{equation}

For an unlabeled graph type \(H\), let \(\mathcal O(H)\) be its orbit of
distinct labeled graphs on \([5]\), and put
\begin{equation}
 \mathfrak D_HM:=\sum_{G\in\mathcal O(H)}D_GM.
\label{eq:V58-orbit-difference}
\end{equation}
We use the following six types; the displayed edges specify a
representative and therefore remove every naming ambiguity:
\begin{align*}
 S&: \{12,13,14,15\},
 &|\mathcal O(S)|&=5,
 &\deg S&=(4,1,1,1,1),\\
 Y&: \{12,14,15,23\},
 &|\mathcal O(Y)|&=60,
 &\deg Y&=(3,2,1,1,1),\\
 P&: \{13,15,23,24\},
 &|\mathcal O(P)|&=60,
 &\deg P&=(2,2,2,1,1),\\
 U_\Delta&: \{12,13,14,15,23\},
 &|\mathcal O(U_\Delta)|&=30,
 &\deg U_\Delta&=(4,2,2,1,1),\\
 U_4&: \{13,14,15,23,24\},
 &|\mathcal O(U_4)|&=60,
 &\deg U_4&=(3,2,2,2,1),\\
 C_5&: \{14,15,23,25,34\},
 &|\mathcal O(C_5)|&=12,
 &\deg C_5&=(2,2,2,2,2).
\end{align*}
Here \(U_\Delta\) is a triangle with two pendant edges at the same
triangle vertex, and \(U_4\) is a four-cycle with one pendant edge.

\section{A reproducible coefficient count}
\label{sec:V58-coefficient-count}

The next lemma makes the finite calculation auditable without appealing
to a symbolic package.

\begin{lemma}[Induced-block formula for a graph-difference coefficient]
\label{lem:V58-induced-block-formula}
For a finite graph \(K\), set
\begin{equation}
 r(K):=\sum_{A\subseteq E(K)\,:,(V(K),A)\ {\rm connected}}
       (-1)^{|E(K)|-|A|},
\label{eq:V58-rK}
\end{equation}
with \(r(K)=1\) when \(K\) has one vertex.  Fix
\(\pi=\{B_1,\ldots,B_k\}\in\Pi_5\).  The coefficient of \(M(\pi)\) in
\(D_GM\) is
\begin{equation}
 (-1)^{e_\pi(G)}\prod_{q=1}^k r(G[B_q]),
\label{eq:V58-induced-block-coefficient}
\end{equation}
where \(e_\pi(G)\) is the number of edges of \(G\) joining different
blocks of \(\pi\).  In particular it is zero if some induced graph
\(G[B_q]\) is disconnected.
\end{lemma}

\begin{proof}
The equality \(\pi(F)=\pi\) holds exactly when \(F\) has no cross-block
edge and each \(F\cap E(G[B_q])\) is connected and spanning on \(B_q\).
For such an \(F\),
\[
 |E(G)|-|F|
 =e_\pi(G)+\sum_q\bigl(|E(G[B_q])|-|F\cap E(G[B_q])|\bigr).
\]
Insert this in \eqref{eq:V58-graph-difference}; the independent sums over
the blocks factor and give \eqref{eq:V58-induced-block-coefficient}.
\end{proof}

Since an orbit sum is invariant under relabeling, its coefficient depends
only on the integer partition formed by the block sizes of \(\pi\).
Applying \Cref{lem:V58-induced-block-formula} to the six displayed edge
sets and their relabelings gives the following complete table.  The last
column is the coefficient in \(\Lambda_5\).
\begin{equation}
\begin{array}{c|rrrrrr|r}
 \operatorname{type}(\pi)
 &\mathfrak D_S&\mathfrak D_Y&\mathfrak D_P
 &\mathfrak D_{U_\Delta}&\mathfrak D_{U_4}&\mathfrak D_{C_5}
 &\mu(\pi,\widehat1)\\ \hline
 5       & 5& 60& 60&-60&-180&-48& 1\\
 4+1     &-4&-36&-24& 36&  72& 12&-1\\
 3+2     & 0& -6&-12&  3&  18&  6&-1\\
 3+1+1   & 3& 24& 18&-21& -36& -6& 2\\
 2+2+1   & 0&  8& 12& -4& -16& -4& 2\\
 2+1+1+1 &-2&-24&-24& 15&  30&  6&-6\\
 1+1+1+1+1&5&60&60&-30&-60&-12&24
\end{array}
\label{eq:V58-coefficient-table}
\end{equation}

For illustration, a one-block partition receives coefficient one from
each labeled tree, so the first entry is \(5\).  For a fixed \(4+1\)
partition, precisely the four stars centred in the four-block contribute,
each with the negative sign of its unique cross edge, giving \(-4\).
For a labeled five-cycle, the connected spanning subsets are the full
cycle and the five paths obtained by deleting one edge, hence
\(r(C_5)=1-5=-4\); multiplying by its orbit size twelve gives \(-48\).
All remaining entries follow from the same formula on induced graphs of
at most five vertices.  Thus the table is a finite application of
\eqref{eq:V58-induced-block-coefficient}, not a numerical fit.

\section{The exact identity and the tree-only obstruction}
\label{sec:V58-exact-identity}

\begin{theorem}[Exact symmetric five-row graph-difference identity]
\label{thm:V58-exact-graph-identity}
For every moment datum \(M:\Pi_5\to\mathcal V\),
\begin{equation}
 \boxed{
 \Lambda_5M
 =\mathfrak D_SM+\frac32\mathfrak D_YM-\frac7{12}\mathfrak D_PM
  +\frac53\mathfrak D_{U_\Delta}M-\frac14\mathfrak D_{U_4}M
  +\frac1{12}\mathfrak D_{C_5}M.}
\label{eq:V58-exact-graph-identity}
\end{equation}
Every graph on the right is connected and has either four or five edges.
\end{theorem}

\begin{proof}
Multiply the six graph-orbit columns of
\eqref{eq:V58-coefficient-table} by
\[
 \left(1,\frac32,-\frac7{12},\frac53,-\frac14,\frac1{12}\right).
\]
The resulting seven entries, in the displayed row order, are
\[
 (1,-1,-1,2,2,-6,24).
\]
These are exactly
\((-1)^{|\pi|-1}(|\pi|-1)!\) for partitions with respectively one
through five blocks.  Equality therefore holds coefficient by coefficient
for all fifty-two partitions in \(\Pi_5\), proving
\eqref{eq:V58-exact-graph-identity}.
\end{proof}

\begin{proposition}[Trees alone cannot represent \(\Lambda_5\)]
\label{prop:V58-tree-only-obstruction}
There are no rational numbers \(a,b,c\) for which
\(
 \Lambda_5=a\mathfrak D_S+b\mathfrak D_Y+c\mathfrak D_P
\)
as functionals on arbitrary moment data.
\end{proposition}

\begin{proof}
The \(3+2\) and \(2+2+1\) rows of
\eqref{eq:V58-coefficient-table} would give
\[
 -6b-12c=-1,
 \qquad 8b+12c=2,
\]
so \(b=1/2\) and \(c=-1/6\).  The one-block row would then force
\(a=-19/5\).  But the \(4+1\) row would have coefficient
\[
 -4a-36b-24c=\frac65\ne-1,
\]
a contradiction.
\end{proof}

\begin{firewall}[Why a positive tree expansion may not be inserted]
\label{fire:V58-no-tree-insertion}
The obstruction is an identity-level statement on arbitrary moment data.
Consequently one may not replace the complete fifth cumulant by a sum of
tree differences merely because a desired final norm has tree form.  The
three unicyclic corrections in
\eqref{eq:V58-exact-graph-identity} carry genuine cancellation.  They must
be retained until their one surplus merge mark has been controlled.
\end{firewall}

\section{Quotient marks before norms}
\label{sec:V58-quotient-marks}

Fix \(\pi\in\Pi_5\) with \(k=|\pi|\).  Contract every block of \(\pi\)
in a connected graph \(G\), retaining parallel cross edges and deleting
loops; call the resulting multigraph \(G/\pi\).

\begin{lemma}[Exact quotient-tree factorization of a selected graph]
\label{lem:V58-quotient-factorization}
The multigraph \(G/\pi\) is connected.  Choose a quotient spanning tree
and one original edge above each of its edges; let \(L_\pi(G)\subseteq
E(G)\) be the resulting lift.  Then \(|L_\pi(G)|=k-1\), every edge in
\(L_\pi(G)\) joins two different blocks, and
\begin{equation}
 D_GM=left(\prod_{a\in L_\pi(G)}\nabla_a\right)
       \left(\prod_{a\in E(G)\setminus L_\pi(G)}\nabla_a\right)
       M(\widehat0).
\label{eq:V58-quotient-factorization}
\end{equation}
For a tree type the second product has \(5-k\) factors; for a unicyclic
type it has \(6-k\) factors.
\end{lemma}

\begin{proof}
The image of a path in \(G\) is a walk between the corresponding quotient
vertices, so contraction preserves connectedness.  A finite connected
multigraph has a spanning tree.  Lifting its edges gives the asserted set
of \(k-1\) cross-block edges.  Equation
\eqref{eq:V58-quotient-factorization} follows from commutativity of the
merge differences.  The residual counts are
\(4-(k-1)=5-k\) and \(5-(k-1)=6-k\), respectively.
\end{proof}

\begin{corollary}[Exact status of the V57 quotient request]
\label{cor:V58-status-V57-request}
At the partition-lattice level, every term in the complete fifth cumulant
now contains a quotient spanning tree of \(k-1\) cross-block differences
before norms are taken.  The discrete term \(M(\widehat0)\) remains
inside each full alternating difference \(D_GM\).  Relative to a tree
budget, the exact identity has at most one surplus merge difference.
\end{corollary}

\begin{proof}
Apply \Cref{lem:V58-quotient-factorization} termwise to the six orbit sums
in \Cref{thm:V58-exact-graph-identity}.  The empty-subgraph term in
\eqref{eq:V58-graph-difference} is never separated, which preserves the
all-singleton cancellation required by
\Cref{fire:V57-discrete-partition}.  The edge counts give the last claim.
\end{proof}

The corollary is deliberately algebraic.  To obtain the normalized
quotient factor of V57, a cross-block difference must still be realized
on the exact Wilson coordinate carrier and shown to produce its
\(s_\alpha\Theta_{BC}\) mark.  Moreover, for a unicyclic orbit the one
surplus difference must be bounded in the same one-payer product state;
bounding a separately expanded moment is forbidden.

\begin{theorem}[Conditional completion criterion for quintic assembly]
\label{thm:V58-conditional-quintic-criterion}
Suppose that, for each of the six graph types in
\eqref{eq:V58-exact-graph-identity}, the assembled Wilson carrier admits
the factorization \eqref{eq:V58-quotient-factorization} statewise, with:
\begin{enumerate}[label=\textup{(\alph*)}]
\item each selected quotient edge delivering the normalized factor
      \(Cs_\alpha\theta_{ij}\);
\item the residual lower blocks obeying the one-payer capacity of V57;
      and
\item the sole surplus factor in a unicyclic graph costing at most a
      universal constant in that same capacity.
\end{enumerate}
Then the complete mixed-support fifth response obeys the global quintic
MTA bound
\begin{equation}
 C^4s_\alpha^3
 \sum_{\tau\in\mathfrak T_5}
 \prod_{ij\in E(\tau)}\theta_{ij}.
\label{eq:V58-conditional-global-quintic}
\end{equation}
\end{theorem}

\begin{proof}
Use \Cref{thm:V58-exact-graph-identity} before taking any norm.  For each
row partition choose the quotient lifts from
\Cref{lem:V58-quotient-factorization}.  Assumption (a) supplies the
quotient-tree marks required in V57.  Assumption (b) invokes the exact
one-payer forest composition theorem
\Cref{thm:V57-forest-composition}; its exponent is
\(s_\alpha^{4-k}s_\alpha^{k-1}=s_\alpha^3\).
For a unicyclic type, assumption (c) absorbs the only additional edge.
The six coefficients and all orbit multiplicities are numerical constants
depending only on five.  Summing them and applying the forest-plus-quotient
lift gives \eqref{eq:V58-conditional-global-quintic}.
\end{proof}

\section{V58 ledger and V59 mandate}
\label{sec:V58-ledger}

\begin{theorem}[Integrated compact-series V58 statement]
\label{thm:V58-integrated}
Preserving compact V1--V57, V58 proves:
\begin{enumerate}[label=\textup{(V58.\arabic*)},leftmargin=6.7em]
\item the exact induced-block coefficient formula for arbitrary graph
      merge-differences;
\item the complete seven-by-six orbit coefficient table;
\item the exact fifth-cumulant identity
      \eqref{eq:V58-exact-graph-identity};
\item a direct proof that the three tree orbits alone cannot represent
      the fifth cumulant;
\item exact quotient spanning-tree extraction before norms, with at most
      one surplus difference; and
\item a precise conditional criterion reducing global quintic MTA to a
      statewise Wilson estimate for six graph orbits.
\end{enumerate}
\end{theorem}

\begin{proof}
These are
\Cref{lem:V58-induced-block-formula,
thm:V58-exact-graph-identity,
prop:V58-tree-only-obstruction,
lem:V58-quotient-factorization,
thm:V58-conditional-quintic-criterion}.
\end{proof}

\begin{openproblem}[V59 mandate: lift merge differences to Wilson carriers]
\label{op:V59-mandate}
\begin{enumerate}
\item Realize each row merge difference \(\nabla_{ij}\) inside the exact
      incidence-replica Wilson formula, before Fourier and operator norms.
\item Prove that a selected quotient lift delivers one normalized
      \(s_\alpha\theta_{ij}\) factor in the one-payer product state.
\item Close the three tree orbits first using V57, without introducing a
      second resolvent.
\item For \(U_\Delta,U_4,C_5\), delete one cycle edge, keep its difference
      assembled, and prove a universal one-surplus capacity bound.
\item If that bound fails, classify the first local coordinate block that
      meets the cycle in two non-leaf rows and state its exact missing
      numerator estimate; do not recycle a failed global argument.
\end{enumerate}
\end{openproblem}

\begin{firewall}[Claim boundary after compact V58]
\label{fire:V58-current-boundary}
V58 closes the exact five-row partition algebra and proves that the
remaining signed correction is only unicyclic.  It does not yet prove the
statewise Wilson realization of every merge mark, the one-surplus
capacity estimate, global quintic or all-order MTA, WUFC, CPM, continuum
Yang--Mills existence, or a positive mass gap.  The full Yang--Mills mass
gap remains open.
\end{firewall}

\section{Append-only record for compact V58}
\label{sec:V58-record}

V58 was appended to the complete compact V57 file.  The exact V57 parent
had (22622) newline-delimited source records, (761879) bytes, and
SHA--256 digest
\[
\texttt{4935205bb613fa7f60b4f290e028f479da6b545423fc0f5ebde635b4c5900897}.
\]
No compact V57 mathematical content was changed.  The sole final
\verb|\end{document}| is restored below.

% =============================================================================
% END APPENDED COMPACT VERSION V58 MATERIAL
% =============================================================================

% The compact V58 document terminator is intentionally disabled here:
% \end{document}

% =============================================================================
% BEGIN APPENDED COMPACT VERSION V59 MATERIAL
% =============================================================================

\chapter{V59: Graphwise leakage and the assembled cluster repair}
\label{ch:V59}

\section{Purpose and correction}
\label{sec:V59-purpose}

V58 found an exact six-orbit graph-difference identity.  Its next mandate
was to attach one normalized Wilson mark to every displayed graph edge.
That step is false for an individual graph difference.  Once an earlier
merge has joined two row blocks, a later merge can correlate rows other
than the two labels on its graph edge.  A path difference therefore
contains a covariance of its endpoints, multiplied by the means of its
internal vertices.  It can have only one small-noise power regardless of
the length of the path.

This chapter proves the leakage exactly, gives a mixed-support Wilson
counterexample, and records the consequent scope correction.  It also
identifies the cancellation between graph orbits which V58 must preserve.
The replacement route is the assembled product--cumulant identity on row
clusters, followed by a five-row normalized threshold compiler.  No norm
of an individual graph orbit is used in that route.

\section{A graph difference in the cumulant basis}
\label{sec:V59-cumulant-basis}

Let \(X_1,\ldots,X_n\) be scalar random variables with all moments used
below.  For \(\pi,\sigma\in\Pi_n\), put
\begin{align}
 M_X(\pi)&:=\prod_{B\in\pi}\mathbb E\prod_{i\in B}X_i,
\label{eq:V59-moment-partition}\
 K_X(\sigma)&:=\prod_{B\in\sigma}
                  \kappa(X_i:i\in B).
\label{eq:V59-cumulant-partition}
\end{align}
Then the moment--cumulant formula is
\begin{equation}
 M_X(\pi)=\sum_{\sigma\le\pi}K_X(\sigma).
\label{eq:V59-moment-cumulant-refinement}
\end{equation}

\begin{lemma}[Exact cumulant-basis coefficient of \(D_G\)]
\label{lem:V59-cumulant-basis-coefficient}
For a graph \(G\) on \([n]\), the coefficient of \(K_X(\sigma)\) in
\(D_GM_X\) is
\begin{equation}
 c_G(\sigma)
 :=\sum_{F\subseteq E(G)\,:\,\sigma\le\pi(F)}
       (-1)^{|E(G)|-|F|}.
\label{eq:V59-cumulant-basis-coefficient}
\end{equation}
In particular, connectedness of \(G\) does not imply that only the
one-block cumulant survives.
\end{lemma}

\begin{proof}
Insert \eqref{eq:V59-moment-cumulant-refinement} into
\eqref{eq:V58-graph-difference} and interchange the two finite sums.  A
fixed \(K_X(\sigma)\) occurs in \(M_X(\pi(F))\) exactly when
\(\sigma\le\pi(F)\), which gives
\eqref{eq:V59-cumulant-basis-coefficient}.
\end{proof}

\begin{proposition}[Endpoint covariance leakage on every path]
\label{prop:V59-path-leakage}
Let \(P_n\) be the path
\(1-2-\cdots-n\), and set
\[
 \sigma_*:=\{\{1,n\},\{2\},\ldots,\{n-1\}\}.
\]
Then
\begin{equation}
 c_{P_n}(\sigma_*)=1.
\label{eq:V59-path-pair-coefficient}
\end{equation}
Consequently \(D_{P_n}M_X\) contains the term
\begin{equation}
 \kappa(X_1,X_n)\prod_{j=2}^{n-1}\mathbb EX_j
\label{eq:V59-path-endpoint-term}
\end{equation}
with coefficient one.
\end{proposition}

\begin{proof}
The relation \(\sigma_*\le\pi(F)\) holds exactly when vertices \(1\)
and \(n\) lie in the same component of \(([n],F)\).  A subgraph of a
path connects its endpoints only if it contains every path edge.  Thus
the only set in \eqref{eq:V59-cumulant-basis-coefficient} is
\(F=E(P_n)\), and its sign is positive.
\end{proof}

For \(n=3\), direct expansion gives the particularly transparent
identity
\begin{equation}
 D_{1-2-3}M_X
 =\kappa(X_1,X_2,X_3)
  +\kappa(X_1,X_3)\mathbb EX_2.
\label{eq:V59-three-path-identity}
\end{equation}
The second term is not associated with either displayed path edge.

\section{An exact mixed-support Wilson counterexample}
\label{sec:V59-Wilson-counterexample}

\begin{theorem}[A path graph difference has no path-edge marked bound]
\label{thm:V59-Wilson-path-counterexample}
Fix sufficiently small \(\alpha\), so that a chosen nontrivial one-link
Fourier multiplier is nonzero.  There is a five-row Wilson product
measure and admissible finite Peter--Weyl row observables such that, for
the labeled path \(P_5=1-2-3-4-5\),
\begin{equation}
 D_{P_5}M_X\ne0,
 \qquad
 H_{12}=H_{23}=H_{34}=H_{45}=0.
\label{eq:V59-Wilson-path-counterexample}
\end{equation}
Hence no support-uniform estimate of this individual graph difference
by a product of its four path-edge marks can hold.  The analogous
statement already fails for the four-row path.
\end{theorem}

\begin{proof}
Take one Wilson coordinate \(e\) shared only by rows \(1\) and \(5\).
Choose a nonconstant finite Peter--Weyl function \(f\) and use
\(X_1=f(U_e)\), \(X_5=\overline{f(U_e)}\).  The one-link Wilson density
is positive on the compact group, so strict Cauchy--Schwarz gives
\begin{equation}
 \kappa(X_1,X_5)
 =\mathbb E|f|^2-|\mathbb Ef|^2>0.
\label{eq:V59-positive-Wilson-variance}
\end{equation}

Give each of rows \(2,3,4\) its own private coordinate, disjoint from
all the others, and choose a diagonal coefficient \(X_j\) whose mean
is the corresponding nonzero central Fourier multiplier.  These three
variables are mutually independent and independent of the pair
\((X_1,X_5)\).  Therefore every mixed cumulant vanishes except the
endpoint covariance and singleton means.  By
\Cref{prop:V59-path-leakage},
\begin{equation}
 D_{P_5}M_X
 =\kappa(X_1,X_5)
  \prod_{j=2}^4\mathbb EX_j\ne0.
\label{eq:V59-exact-Wilson-path-value}
\end{equation}

No adjacent pair in the displayed path shares a coordinate.  Hence all
four raw overlaps in \eqref{eq:V59-Wilson-path-counterexample} vanish.
Deleting row \(3\) and one private factor gives the identical four-row
path counterexample.
\end{proof}

\begin{corollary}[Failure of the graph-edge witness assertion]
\label{cor:V59-witness-fails}
In an iterated merge difference, it is false that a difference labeled
\(ij\) can only select a coordinate supporting both rows \(i\) and
\(j\).  It can select a coordinate joining any row in the current block
of \(i\) to any row in the current block of \(j\).
\end{corollary}

\begin{proof}
In \eqref{eq:V59-three-path-identity}, after the \(23\) merge the next
\(12\) merge joins row \(1\) to the block \(\{2,3\}\).  It may therefore
detect a coordinate shared by rows \(1\) and \(3\), even when row \(2\)
is absent from that coordinate.  The Wilson construction above realizes
this possibility exactly.
\end{proof}

\section{The cancellation that must remain assembled}
\label{sec:V59-orbit-cancellation}

The leakage is not a contradiction to the exact identities of V47 or
V58.  Those identities cancel it between different graph orbits.

\begin{proposition}[Pair-cumulant cancellation in the V58 identity]
\label{prop:V59-quintic-pair-cancellation}
Fix
\(
 \sigma=\{\{1,2\},\{3\},\{4\},\{5\}\}
\).
The coefficients of \(K_X(\sigma)\) in the six orbit sums of V58 are
\begin{equation}
\begin{array}{c|rrrrrr}
 &\mathfrak D_S&\mathfrak D_Y&\mathfrak D_P
 &\mathfrak D_{U_\Delta}&\mathfrak D_{U_4}&\mathfrak D_{C_5}\\ \hline
 [2,1,1,1]&0&0&6&0&-18&-12.
\end{array}
\label{eq:V59-quintic-pair-row}
\end{equation}
With the coefficients of \eqref{eq:V58-exact-graph-identity}, their
sum is zero:
\begin{equation}
 6\left(-\frac7{12}\right)
 -18\left(-\frac14\right)
 -12\left(\frac1{12}\right)=0.
\label{eq:V59-quintic-pair-cancels}
\end{equation}
\end{proposition}

\begin{proof}
Apply the finite coefficient formula
\eqref{eq:V59-cumulant-basis-coefficient} to a representative of every
orbit and sum over its distinct relabelings.  This gives the displayed
row.  The arithmetic in
\eqref{eq:V59-quintic-pair-cancels} is exact.  It must vanish because
\(\Lambda_5M_X=\kappa(X_1,\ldots,X_5)\) has no proper cumulant-product
term.
\end{proof}

The same audit at four rows gives, in the V47 column order
\((\mathfrak S,\mathfrak P,\mathfrak C,\mathfrak A)\),
\begin{equation}
 (0,2,-3,-4)
\label{eq:V59-quartic-pair-row}
\end{equation}
for a fixed \([2,1,1]\) cumulant product.  The V47 weights cancel it:
\begin{equation}
 2\left(\frac23\right)
 -3\left(-\frac13\right)
 -4\left(\frac7{12}\right)=0.
\label{eq:V59-quartic-pair-cancels}
\end{equation}

\begin{firewall}[Graph orbits may not be normed separately]
\label{fire:V59-no-graphwise-norm}
Equations \eqref{eq:V59-quintic-pair-cancels} and
\eqref{eq:V59-quartic-pair-cancels} occur across graph classes.  Taking
an absolute value, operator norm, or positive envelope of each graph or
each orbit before summing destroys them.  Connectedness of the graph is
not a substitute for connectedness of the cumulant functional.
\end{firewall}

\section{Dependency correction}
\label{sec:V59-dependency-correction}

\begin{theorem}[Exact scope of the graphwise regression]
\label{thm:V59-scope-correction}
The following earlier results remain valid and are not changed:
\begin{enumerate}[label=\textup{(\alph*)}]
\item the partition-lattice identities
      \eqref{eq:V47-connected-graph-identity} and
      \eqref{eq:V58-exact-graph-identity};
\item the V37--V39 assembled weak-cut, three-cluster, and universal
      normalized estimates;
\item the exact one-link moment and private-tail theorems of V52;
\item the uniform one-link fifth-cumulant card of V56; and
\item the binary and ternary entries of the V57 lower-order table.
\end{enumerate}
The recorded proofs do \emph{not} establish any conclusion that requires
norming separate graph differences and assigning a mark to every graph
edge.  In particular, the all-weak quartic conclusion
\eqref{eq:V52-all-weak-closed} and therefore the global quartic conclusion
\eqref{eq:V52-global-quartic-MTA} are withdrawn pending an assembled
repair.  Consequently the paid quartic entry of
\eqref{eq:V57-payer-table}, and V57 compositions using a four-row block,
are conditional.  The V58 identity remains exact, but it may not be
inserted graph orbit by graph orbit into
\Cref{thm:V58-conditional-quintic-criterion}.
\end{theorem}

\begin{proof}
The positive assertions have proofs independent of the false witness
claim in \Cref{cor:V59-witness-fails}.  By contrast, V47--V52 resolve and
bound the path, cycle, and paw graph differences as separate analytic
branches when proving the all-weak carrier.  The exact Wilson example in
\Cref{thm:V59-Wilson-path-counterexample} disproves the required
graph-edge estimate already on the path branch.  Thus those proofs do not
imply their all-weak and global conclusions.  V57 cites the latter global
quartic response for its paid size-four entry, so precisely that entry
and its uses inherit the conditional status.
\end{proof}

\begin{remark}[Withdrawal is not a counterexample to quartic MTA]
\label{rem:V59-no-quartic-counterexample}
The counterexample concerns a summand introduced by an auxiliary graph
identity.  The complete cumulant cancels that summand.  It therefore does
not disprove the global quartic or quintic marked-tree estimates; it
shows that the recorded graphwise proof route cannot establish them.
\end{remark}

\section{The assembled cluster identity}
\label{sec:V59-cluster-identity}

Let \(\rho=\{C_1,\ldots,C_k\}\) be a partition of \([5]\), and put
\(Y_a:=\prod_{i\in C_a}X_i\) in the fixed row order.

\begin{theorem}[Product--cumulant identity on arbitrary row clusters]
\label{thm:V59-product-cumulant}
For \(2\le k\le5\),
\begin{equation}
 \boxed{
 \kappa(Y_1,\ldots,Y_k)
 =\sum_{\substack{\pi\in\Pi_5\\\pi\vee\rho=\widehat1}}
   \prod_{B\in\pi}\kappa(X_i:i\in B).}
\label{eq:V59-product-cumulant}
\end{equation}
The term \(\pi=\widehat1\) is the complete fifth cumulant.  Every other
term is a product of lower connected blocks whose supports collectively
connect the quotient clusters.
\end{theorem}

\begin{proof}
Expand the left side by the cumulant M\"obius formula on \(\Pi_k\).
For a partition \(\eta\in\Pi_k\), each moment
\(
 \prod_{D\in\eta}\mathbb E\prod_{a\in D}Y_a
\)
is then expanded into row cumulants.  A fixed row partition \(\pi\)
appears precisely for those \(\eta\) whose lifted blocks coarsen
\(\pi\vee\rho\).  Its coefficient is therefore
\[
 \sum_{\eta\,:\,\eta\ge \pi\vee\rho}\mu(\eta,\widehat1).
\]
The defining M\"obius cancellation makes this sum one when
\(\pi\vee\rho=\widehat1\) and zero otherwise.  This proves
\eqref{eq:V59-product-cumulant}.  The same finite polynomial identity
holds coefficientwise for the programme's fixed tensor ordering.
\end{proof}

Unlike a graph difference, the left side of
\eqref{eq:V59-product-cumulant} is a genuinely connected cumulant of the
\(k\) composite variables.  It may be bounded once by the known order-
\(k\) response theorem.  The lower products are precisely the objects
for which the one-payer rule was designed.  No cancellation between
auxiliary graph classes is lost.

\section{A five-row normalized threshold compiler}
\label{sec:V59-five-row-compiler}

For a row partition
\(\rho=\{C_1,\ldots,C_k\}\), set
\begin{align}
 \Theta(C_a,C_b)&:=
   \sum_{i\in C_a,\,j\in C_b}\theta_{ij},
\label{eq:V59-quotient-overlap}\
 \mathcal T_\rho(\theta)&:=
   \sum_{\zeta\in\mathfrak T(\rho)}
   \prod_{C_aC_b\in E(\zeta)}\Theta(C_a,C_b).
\label{eq:V59-quotient-tree-polynomial}
\end{align}

\begin{theorem}[Finite normalized graph reduction on five rows]
\label{thm:V59-five-row-normalized-graph}
Let \(0\le\theta_{ij}\le1\), fix \(\delta\in(0,1)\), and let
\(\mathcal F_5(\theta)\ge0\).  Suppose:
\begin{enumerate}[label=\textup{(\roman*)},leftmargin=4em]
\item \(\mathcal F_5\le C_1\);
\item for every row partition \(\rho\) with
      \(2\le|\rho|\le4\),
      \(\mathcal F_5\le C_{|\rho|}\mathcal T_\rho(\theta)\);
\item if every \(\theta_{ij}<\delta\), then
      \[
       \mathcal F_5\le C_5
       \sum_{\tau\in\mathfrak T_5}
       \prod_{ij\in E(\tau)}\theta_{ij}.
      \]
\end{enumerate}
Then
\begin{equation}
 \boxed{
 \mathcal F_5(\theta)
 \le C
 \sum_{\tau\in\mathfrak T_5}
 \prod_{ij\in E(\tau)}\theta_{ij},}
\label{eq:V59-five-row-normalized-graph}
\end{equation}
where \(C\) depends only on \(\delta\) and the displayed constants.
\end{theorem}

\begin{proof}
Form the graph of strong edges \(ij\) with
\(\theta_{ij}\ge\delta\), and let its connected components be
\(C_1,\ldots,C_k\).

If \(k=1\), choose a strong spanning tree.  Its monomial is at least
\(\delta^4\), so hypothesis \textup{(i)} is bounded by the full tree
polynomial times \(C_1\delta^{-4}\).

If \(2\le k\le4\), choose a strong spanning tree inside each component.
Their union is a forest with
\(
 \sum_a(|C_a|-1)=5-k
\)
edges and weight at least \(\delta^{5-k}\).  Apply hypothesis
\textup{(ii)} to the component partition and expand every quotient edge
sum in \eqref{eq:V59-quotient-tree-polynomial}.  Adjoining the chosen
internal forest to each lifted quotient tree gives a labeled tree on
\([5]\).  A labeled tree has only a fixed finite number of such
preimages.  Thus
\[
 \mathcal T_\rho(\theta)
 \le C\delta^{-(5-k)}
 \sum_{\tau\in\mathfrak T_5}
 \prod_{ij\in E(\tau)}\theta_{ij}.
\]

If \(k=5\), all edges are weak and hypothesis \textup{(iii)} applies
directly.  These cases exhaust the strong graph and prove
\eqref{eq:V59-five-row-normalized-graph}.
\end{proof}

\begin{assessment}[The repaired order of proof]
\label{ass:V59-repaired-order}
The safe route is now:
\begin{enumerate}
\item restore the quartic all-weak carrier with its complete cumulant
      assembled, so the paid size-four response used by V57 is valid;
\item apply \Cref{thm:V59-product-cumulant} for quotient partitions with
      two, three, and four clusters, and use exactly one payer on every
      lower-block correction;
\item prove the genuinely all-weak five-row estimate without taking
      graph-orbit norms; and
\item invoke \Cref{thm:V59-five-row-normalized-graph}.
\end{enumerate}
The V58 graph identity remains an exact diagnostic for cancellations,
not an authorized positive decomposition.
\end{assessment}

\section{V59 ledger and V60 mandate}
\label{sec:V59-ledger}

\begin{theorem}[Integrated compact-series V59 statement]
\label{thm:V59-integrated}
Preserving compact V1--V58, V59 proves:
\begin{enumerate}[label=\textup{(V59.\arabic*)},leftmargin=6.7em]
\item the exact cumulant-basis coefficient formula for a graph
      difference;
\item endpoint-covariance leakage on paths of every length;
\item an exact mixed-support Wilson counterexample to graph-edge
      marking of an individual path difference;
\item the pair-cumulant cancellation between the V47 and V58 graph
      orbits;
\item the precise withdrawal and dependency correction forced by this
      regression;
\item the arbitrary-cluster product--cumulant identity; and
\item a finite normalized five-row threshold compiler independent of
      graph-difference norming.
\end{enumerate}
\end{theorem}

\begin{proof}
These are
\Cref{lem:V59-cumulant-basis-coefficient,
prop:V59-path-leakage,
thm:V59-Wilson-path-counterexample,
prop:V59-quintic-pair-cancellation,
thm:V59-scope-correction,
thm:V59-product-cumulant,
thm:V59-five-row-normalized-graph}.
\end{proof}

\begin{openproblem}[V60 mandate: audit and assembled quartic repair]
\label{op:V60-mandate}
\begin{enumerate}
\item Perform the mandatory read-only audit of the newest active SM and
      NS notes before selecting the repair.
\item Return to the complete quartic cumulant carrier used in V36--V40,
      before the V47 graph decomposition and before positive envelopes.
\item Preserve the valid two-cluster, three-cluster, universal, thermal,
      and one-link tail estimates, but do not cite the withdrawn V52
      all-weak conclusion.
\item Determine whether the hot-row stock estimates can be applied to one
      assembled cut-connected product--cumulant identity.  If not,
      isolate the exact simultaneous protected capacity still missing.
\item Restore the paid quartic response before using V57 on a partition
      containing a four-row block.
\end{enumerate}
\end{openproblem}

\begin{firewall}[Claim boundary after compact V59]
\label{fire:V59-current-boundary}
V59 repairs a false graphwise inference and supplies the exact assembled
cluster route.  The global quartic MTA is no longer claimed as proved;
its weak-cut, three-cluster, universal, thermal, and Wilson-tail inputs
remain available.  Global quintic and all-order MTA, WUFC, CPM, continuum
Yang--Mills existence, and a positive mass gap remain open.  The full
Yang--Mills mass gap remains open.
\end{firewall}

\section{Append-only record for compact V59}
\label{sec:V59-record}

V59 was appended to the complete compact V58 file.  The exact V58 parent
had (23018) newline-delimited source records, (777554) bytes, and
SHA--256 digest
\[
\texttt{25ac0e70cc11c0f6faf25a3bdd74dfcabed456889e068f55966a07e4899f2c8f}.
\]
No compact V58 mathematical content was changed.  The sole final
\verb|\end{document}| is restored below.

% =============================================================================
% END APPENDED COMPACT VERSION V59 MATERIAL
% =============================================================================

% The compact V59 document terminator is intentionally disabled here:
% \end{document}

% =============================================================================
% BEGIN APPENDED COMPACT VERSION V60 MATERIAL
% =============================================================================

\chapter{V60: Companion audit and the topology-transfer Green gate}
\label{ch:V60}

\section{Mandatory SM/NS companion audit}
\label{sec:V60-audit}

The newest active companion sources inspected read-only were:
\begin{center}
\begin{tabular}{p{0.42\linewidth}r r p{0.30\linewidth}}
\toprule
file & lines & bytes & SHA--256\\
\midrule
\texttt{Renewal\_SM\_Einstein\_Continuum\_Research\_Notes\_V68.tex}
& \(26090\) & \(914156\)
& \texttt{249e261736eb0a2ca7d6c6e0777b7a4375636b2bc01e4ad3b8e75a22f57614f8}\\
\texttt{Renewal\_NS\_Stability\_Research\_Notes\_V31.tex}
& \(23052\) & \(887537\)
& \texttt{f1121b62238ad5491bb7cf03635ea9934942edf573ed8faaa9ee79e1d38a6696}\\
\bottomrule
\end{tabular}
\end{center}

SM V65 places its fifth-derivative ledger on complete Ward words and
proves that an offending factorwise allocation may disappear only in an
exact packet formed before the norm.  V66--V67 then classify complete
packets rather than individual graph drawings.  V68 proves an exact Hodge
contraction at each finite regulator, but also proves that a collapsing
BRST Laplacian gap makes the contracting homotopy unbounded.  Thus finite
algebraic exactness does not by itself give a regulator-uniform analytic
descent.

NS V31 is unchanged.  Its dyadic probability mark can be removed only by
paying the exact first radial moment, and its persistent intervals may not
be charged again in the formation ledger.  The signed Navier--Stokes
source has to remain assembled if it is to improve the heat-only bound.

\begin{assessment}[V60 audit decision]
\label{ass:V60-audit-decision}
No SM or NS estimate is imported.  The proof-order lessons agree exactly
with V59: keep the full cumulant packet through cancellation, charge the
single Yang--Mills output payer once, and demand a quantitative uniform
bound on the remaining topology transfer.  A finite partition automaton
or an exact quotient identity alone does not supply that bound.
\end{assessment}

\begin{firewall}[Cross-program type boundary]
\label{fire:V60-audit-boundary}
A BRST Hodge gap is not a Wilson spectral gap, and an NS radial first
moment is not a Yang--Mills row mark.  Only the logical distinction between
exact finite decomposition and uniform analytic control is transferred.
\end{firewall}

\section{Relation to the legacy V133 regression}
\label{sec:V60-V133}

The frozen f1 archive had already found the local ancestor of the V59
problem.  Legacy V133 proved that two Bernoulli bridge derivatives which
become parallel after contraction collapse to one merger rank; a product
tax for both nominal bridges is then false.  V59 is not a repetition of
that result.  It proves that the later V47 merge-difference route
reintroduced the same defect globally, identifies the endpoint-covariance
term in the cumulant basis, gives an actual mixed-support Wilson
counterexample, and computes the cancellation between the V47/V58 graph
orbits.

\begin{assessment}[Surviving upstream objects]
\label{ass:V60-surviving-objects}
The safe objects after both regressions are:
\begin{enumerate}
\item the exact connected coordinate-word formula and the low-activity
      canonical-skeleton theorem of legacy V134;
\item the complete weak-cut/cluster threshold numerator of legacy V135
      and compact V36--V39;
\item the exact topology priority partition and local physical capacity
      cards through compact V35;
\item the one-link Wilson tail moments of V52 and the quintic local card
      of V56; and
\item the assembled cluster identity of V59.
\end{enumerate}
What is missing is a support-uniform way to propagate a topology after its
first connected source has occurred, without replacing the restricted
signed suffix by a positive product.
\end{assessment}

\section{Exact first-entry transfer in block form}
\label{sec:V60-block-transfer}

The V35 \([3,3]\) automaton is the first unresolved instance, but the
following statement applies to every finite topology priority automaton.
Let \(P\) denote the direct sum of all pre-entry states and \(A\) the
accepted state.  At ordered coordinate \(r\), write the exact local
operator-valued transfer in triangular block form
\begin{equation}
 \mathbb T_r=
 \begin{pmatrix}
  B_r&0\\ F_r&A_r
 \end{pmatrix}.
\label{eq:V60-triangular-transfer}
\end{equation}
The zero block records that an accepted word cannot return to a pre-entry
state.  A higher-priority dead state is omitted, since its contribution is
assigned to its own already separated topology packet.

Starting with \((p_0,0)\), let
\((p_r,a_r)=\mathbb T_r(p_{r-1},a_{r-1})\).  For \(s\ge r\), put
\begin{equation}
 G_{s:r}:=A_sA_{s-1}\cdots A_r,
 \qquad G_{s:s+1}:=I.
\label{eq:V60-accepted-propagator}
\end{equation}

\begin{theorem}[Exact discrete Duhamel formula for first entry]
\label{thm:V60-first-entry-Duhamel}
For every \(n\),
\begin{align}
 p_n&=B_nB_{n-1}\cdots B_1p_0,
\label{eq:V60-pre-entry-product}\
 a_n&=\sum_{r=1}^nG_{n:r+1}F_rp_{r-1}.
\label{eq:V60-first-entry-Duhamel}
\end{align}
Every accepted coordinate word occurs exactly once, indexed by its first
transition into \(A\).
\end{theorem}

\begin{proof}
The upper equation in
\eqref{eq:V60-triangular-transfer} gives
\eqref{eq:V60-pre-entry-product}.  The lower equation is
\[
 a_r=F_rp_{r-1}+A_ra_{r-1}.
\]
Iterating from \(a_0=0\) gives
\eqref{eq:V60-first-entry-Duhamel}.  Determinism of the state transition
gives uniqueness of the first-entry index, exactly as in the V35
automaton proof.
\end{proof}

\section{The capacity-compatible Green stock}
\label{sec:V60-Green-stock}

Let \(\|cdot\|_*\) be any seminorm on accepted carriers in which exact
signed packets are assembled before a positive envelope is taken.  It is
required to dominate the final coefficient product-state capacity after
the sole output payer has been allocated.  Denote the induced transfer
norm by the same symbol.

\begin{definition}[Accepted-transfer Green constants]
\label{def:V60-Green-constants}
Define
\begin{align}
 \mathfrak G_0
 &:=\sup_{n\ge r\ge1}\|G_{n:r+1}\|_* ,
\label{eq:V60-G0}\
 \mathfrak G(\omega)
 &:=\sup_{n\ge1}
   \sum_{r=1}^n
   \omega_r\|G_{n:r+1}\|_* ,
\label{eq:V60-Gomega}
\end{align}
for a nonnegative scalar stock \(\omega=(\omega_r)\).
\end{definition}

\begin{lemma}[Two exact aggregation alternatives]
\label{lem:V60-Green-aggregation}
Suppose the source terms in
\eqref{eq:V60-first-entry-Duhamel} satisfy either
\begin{align}
 \sum_{r=1}^n\|F_rp_{r-1}\|_*&\le W_n,
\label{eq:V60-source-l1}
\end{align}
or, for one common majorant \(W\),
\begin{align}
 \|F_rp_{r-1}\|_*&\le\omega_rW.
\label{eq:V60-source-weighted}
\end{align}
Then respectively
\begin{align}
 \|a_n\|_*&\le\mathfrak G_0W_n,
\label{eq:V60-G0-bound}\
 \|a_n\|_*&\le\mathfrak G(\omega)W.
\label{eq:V60-Gomega-bound}
\end{align}
\end{lemma}

\begin{proof}
Apply the triangle inequality to the exact Duhamel formula
\eqref{eq:V60-first-entry-Duhamel}, followed by the induced transfer
norm.  Use \eqref{eq:V60-G0} with
\eqref{eq:V60-source-l1}, or sum
\eqref{eq:V60-source-weighted} with
\eqref{eq:V60-Gomega}.
\end{proof}

The first alternative is appropriate when the global normalized tree
sums already pay every possible first-entry location.  The second is the
renewal alternative when a source occurrence and later accepted transfer
share one decay stock.

\section{A weighted hazard theorem}
\label{sec:V60-hazard}

\begin{theorem}[Local hazard implies a uniform Green first moment]
\label{thm:V60-hazard-Green}
Let \(0\le\omega_r\le1\).  Suppose there is \(\gamma>0\) and a common
accepted-carrier seminorm such that
\begin{equation}
 \|A_rz\|_*
 \le e^{-\gamma\omega_r}\|z\|_*
\label{eq:V60-local-hazard}
\end{equation}
for every accepted packet \(z\).  Then
\begin{equation}
 \boxed{
 \mathfrak G(\omega)
 \le\frac1{1-e^{-\gamma}}.}
\label{eq:V60-hazard-Green}
\end{equation}
Consequently sources satisfying
\eqref{eq:V60-source-weighted} aggregate uniformly in the support size.
\end{theorem}

\begin{proof}
Submultiplicativity gives
\[
 \|G_{n:r+1}\|_*
 \le\exp\left(-\gamma\sum_{t=r+1}^n\omega_t\right).
\]
Put \(S_r:=\sum_{t=r}^n\omega_t\).  Concavity of
\(1-e^{-\gamma x}\) on \([0,1]\) gives
\[
 1-e^{-\gamma\omega_r}
 \ge(1-e^{-\gamma})\omega_r.
\]
Therefore
\begin{align*}
 \omega_r e^{-\gamma S_{r+1}}
 &\le\frac{e^{-\gamma S_{r+1}}-e^{-\gamma S_r}}
             {1-e^{-\gamma}}.
\end{align*}
Sum over \(r\).  The numerator telescopes to
\(1-e^{-\gamma S_1}\le1\), proving
\eqref{eq:V60-hazard-Green}.
\end{proof}

\begin{corollary}[Gapped and protected hazards may be combined once]
\label{cor:V60-combined-hazard}
Suppose a local accepted transfer decomposes into orthogonal physical
outputs, and after the complete local packet is formed each output has
either a gapped norm factor \(e^{-c s_\alpha}\) or a protected
product-state capacity factor at most \(e^{-c}\).  If the output sum is
performed before the accepted-carrier seminorm, then
\eqref{eq:V60-local-hazard} follows with
\(\omega_r\) equal to the corresponding gapped/protected occupation
weight, provided the local packet contains no cut-disconnected neutral
sector.
\end{corollary}

\begin{proof}
On each orthogonal output use the stated contraction.  The gapped factor
is bounded by \(e^{-c'\omega_r}\) when its occupation is normalized by
\(s_\alpha\); the protected factor has the same form after setting its
occupation to one.  Sum the orthogonal outputs once.  A neutral sector
would have factor one and is therefore explicitly excluded.
\end{proof}

\begin{firewall}[The neutral-sector exclusion is the YM gate]
\label{fire:V60-neutral-sector}
V44 proves a gapped-or-crossing-protected alternative only after one
cut-connected identity has been assembled.  It has not been proved for
the entire accepted topology transfer simultaneously across all required
cuts.  Thus \Cref{cor:V60-combined-hazard} is a criterion, not a proof
that the Yang--Mills \(A_r\) satisfies a local hazard.
\end{firewall}

\section{Sharp no-gap regressions}
\label{sec:V60-no-gap}

\begin{proposition}[Finite state and exact first entry do not bound the
accepted transfer]
\label{prop:V60-finite-state-no-Green}
On the one-dimensional accepted space, take
\(
 A_r=I
\)
and scalar sources \(F_rp_{r-1}=w>0\).  Then
\begin{equation}
 a_n=nw.
\label{eq:V60-linear-growth}
\end{equation}
If instead \(A_r=1+\varepsilon\) for fixed \(\varepsilon>0\), then
\begin{equation}
 a_n=w\frac{(1+\varepsilon)^n-1}{\varepsilon}.
\label{eq:V60-exponential-growth}
\end{equation}
Thus neither a finite automaton, unique first entry, nor a fixed bound on
one-step transfer coefficients implies a support-uniform Green bound.
\end{proposition}

\begin{proof}
Insert the scalar transfers into
\eqref{eq:V60-first-entry-Duhamel}.  The first sum has \(n\) equal terms.
The second is the elementary geometric sum.
\end{proof}

\begin{proposition}[A collapsing hazard loses its exact first moment]
\label{prop:V60-collapsing-hazard}
Let \(A_r=e^{-\gamma_n}\), \(F_rp_{r-1}=\gamma_nW\), and take
\(n=\lfloor\gamma_n^{-1}\rfloor\) with \(\gamma_n\downarrow0\).  Then
\(a_n\) remains of order \(W\), whereas a bound which discards the
factor \(\gamma_n\) from each source would grow like
\(\gamma_n^{-1}W\).  Hence the source occupation and the transfer hazard
must use the same normalized first-moment stock.
\end{proposition}

\begin{proof}
The exact sum is
\[
 \gamma_nW\sum_{j=0}^{n-1}e^{-j\gamma_n}
 =\gamma_nW\frac{1-e^{-n\gamma_n}}{1-e^{-\gamma_n}},
\]
which is bounded above and below by fixed positive multiples of \(W\).
Removing \(\gamma_n\) multiplies this expression by
\(\gamma_n^{-1}\).
\end{proof}

\section{The exact quartic repair criterion}
\label{sec:V60-quartic-criterion}

Let \(\mathcal A_{s:r}^{33}\) be the accepted \([3,3]\) transfer from
V35, with every allowed local moment/cumulant letter retained, and let
\(\|cdot\|_{\rm pay}\) denote the coefficient-capacity seminorm after
one final payer allocation.  Define
\begin{equation}
 \mathfrak G_{33}
 :=\sup_{s\ge r}
 \|\mathcal A_{s:r}^{33}\|_{{\rm pay}\to{\rm pay}}.
\label{eq:V60-G33}
\end{equation}
Define analogous constants for the multiple-\([4]\) and collision
accepted states of the exact legacy priority automaton.

\begin{theorem}[Topology-transfer criterion for restoring quartic MTA]
\label{thm:V60-quartic-repair-criterion}
Assume:
\begin{enumerate}[label=\textup{(\alph*)}]
\item the first-entry source sums in each unresolved quartic topology
      obey their already specified normalized marked-tree bounds with the
      final resolvent allocated exactly once;
\item the three accepted-transfer Green constants are bounded uniformly
      in support, representations, physical multiplicities, and volume;
      and
\item the topology priority packets are summed before their positive
      coefficient envelopes are taken.
\end{enumerate}
Then the all-weak normalized quartic carrier
\eqref{eq:V52-all-weak-closed} is valid.  Together with the unaffected
V36--V39 hypotheses, this restores the global quartic marked-tree atlas
\eqref{eq:V52-global-quartic-MTA}.
\end{theorem}

\begin{proof}
Apply \Cref{thm:V60-first-entry-Duhamel} in each topology state and
\Cref{lem:V60-Green-aggregation} with the source alternative in
assumption (a).  Assumption (b) makes every accepted suffix
support-uniform.  The exact priority partition counts each coordinate
word once; assumption (c) preserves the cancellations before capacity.
The embedded \([4]\), two-junction \([3,3]\), collision, noncollision,
and exclusive-pair local source estimates then sum over only a fixed
finite family of row trees.  This proves the all-weak hypothesis of the
normalized threshold theorem.  The weak-cut, three-cluster, and universal
hypotheses were not affected by V59, so the V36 finite normalized graph
reduction gives the global conclusion.
\end{proof}

\begin{assessment}[The first concrete analytic target]
\label{ass:V60-first-target}
The earliest unresolved Green constant is \(\mathfrak G_{33}\).  V35
already gives its exact first-entry source, two-junction one-payer
allocation, and surplus-edge tree conversion.  Proving
\begin{equation}
 \boxed{\mathfrak G_{33}\le C}
\label{eq:V60-G33-target}
\end{equation}
would close the complete \([3,3]\) aggregation without graphwise merge
marks.  The required proof must work on the accepted packet itself.  A
bound on every local letter separately is insufficient by
\Cref{prop:V60-finite-state-no-Green}.
\end{assessment}

\section{V60 ledger and V61 mandate}
\label{sec:V60-ledger}

\begin{theorem}[Integrated compact-series V60 statement]
\label{thm:V60-integrated}
Preserving compact V1--V59, V60 proves:
\begin{enumerate}[label=\textup{(V60.\arabic*)},leftmargin=6.7em]
\item the mandatory SM V68/NS V31 audit and its strict type firewall;
\item the precise relation between the legacy V133 rank-collapse warning
      and the new V59 graphwise regression;
\item the exact block-transfer Duhamel formula for every first-entry
      topology automaton;
\item unweighted and weighted accepted-transfer Green criteria;
\item a telescoping local-hazard theorem with a support-uniform constant;
\item sharp finite-state and collapsing-hazard regressions; and
\item an exact criterion reducing restoration of quartic MTA to three
      accepted-transfer Green bounds, beginning with
      \eqref{eq:V60-G33-target}.
\end{enumerate}
\end{theorem}

\begin{proof}
These are
\Cref{ass:V60-audit-decision,
thm:V60-first-entry-Duhamel,
lem:V60-Green-aggregation,
thm:V60-hazard-Green,
prop:V60-finite-state-no-Green,
prop:V60-collapsing-hazard,
thm:V60-quartic-repair-criterion}.
\end{proof}

\begin{openproblem}[V61 mandate: resolve the \([3,3]\) accepted transfer]
\label{op:V61-mandate}
\begin{enumerate}
\item Write one coordinate step of \(\mathcal A^{33}\) in the complete
      partition-state basis, not as a product of separately normed
      cumulant letters.
\item Cancel every cut-disconnected protected sector before capacity and
      identify the remaining gapped, crossing-protected, and neutral
      blocks.
\item Test whether those blocks satisfy a common Lyapunov seminorm and the
      local hazard \eqref{eq:V60-local-hazard}.
\item If a neutral block survives, construct its smallest scalar or
      \(SU(2)\) realization and determine whether it cancels only after a
      larger topology packet is added.
\item Prove \eqref{eq:V60-G33-target}, or replace it by the exact weighted
      Green stock which the first-entry source can pay.
\end{enumerate}
\end{openproblem}

\begin{firewall}[Claim boundary after compact V60]
\label{fire:V60-current-boundary}
V60 performs the required companion audit and converts the former vague
accepted-suffix issue into the quantitative Green target
\eqref{eq:V60-G33-target}.  It does not prove that target or restore the
withdrawn global quartic MTA.  Global quintic and all-order MTA, WUFC,
CPM, continuum Yang--Mills existence, and a positive mass gap remain
open.  The full Yang--Mills mass gap remains open.
\end{firewall}

\section{Append-only record for compact V60}
\label{sec:V60-record}

V60 was appended to the complete compact V59 file.  The exact V59 parent
had (23519) newline-delimited source records, (796599) bytes, and
SHA--256 digest
\[
\texttt{64ab78ab9ceed172909a3bebeedfe6e4dd50e09529286982125b72a86728f1d0}.
\]
No compact V59 mathematical content was changed.  The sole final
\verb|\end{document}| is restored below.

% =============================================================================
% END APPENDED COMPACT VERSION V60 MATERIAL
% =============================================================================

% =============================================================================
% BEGIN APPENDED COMPACT VERSION V61 MATERIAL
% =============================================================================

\chapter{First-connectivity regrouping: complete suffixes replace the
accepted-transfer Green problem}
\label{chap:V61-first-connectivity}

\section{Purpose and upstream correction}
\label{sec:V61-purpose}

V60 formulated the correct Green problem for the \emph{old} topology
priority automaton: after a word first entered the accepted \([3,3]\)
state, every later four-block letter was killed, leaving the restricted
factor \(R_s=M_s-K_s\).  The present note goes upstream of that
restriction.  A globally connected coordinate word has a unique first
coordinate at which the join of all row partitions seen so far becomes
the one-block partition.  Classifying by that coordinate leaves every
later letter unrestricted.  The entire suffix therefore sums to the
complete moment \(M_s\), not to \(R_s\).

This does not prove the missing quartic source estimate.  It does prove
that the source need not be propagated by the unproved Green operator
\(\mathcal A^{33}\).  The remaining question is a finite family of
first-connectivity source estimates, together with the ordinary case in
which the sole final payer lies in a complete-moment suffix.  This is a
strict reduction of the V60 gate.

\begin{assessment}[Relation to the frozen programme]
\label{ass:V61-nonduplication}
The frozen V134 first-skeleton construction selected a first local
skeleton and then used product-envelope estimates.  The construction
below instead selects the first \emph{cumulative join equal to
\(\widehat1\)}, retains the exact prefix join, and sums the suffix to
complete local moments before any norm or capacity is taken.  The
linked-word identity is implicit in the replica-cumulant formula of
frozen V39, but the first-connectivity disintegration, exact prefix
factorization, and its removal of the V36 restricted product are new in
this compact continuation.
\end{assessment}

\section{The local partition alphabet}
\label{sec:V61-alphabet}

Fix a finite row set \(I\), \(|I|=m\), and a finite ordered coordinate
set
\[
 E=\{1<\cdots<n\}.
\]
Write \(\Pi(I)\) for the partition lattice, ordered by refinement,
with discrete and one-block partitions \(\widehat0\) and
\(\widehat1\).  At coordinate \(e\), let
\(k_e(B)\) be the joint local cumulant carrier on the nonempty row set
\(B\subseteq I\).  For \(\pi\in\Pi(I)\), put
\begin{equation}
 L_{e,\pi}:=\bigotimes_{B\in\pi}k_e(B),
 \qquad
 M_e:=\sum_{\pi\in\Pi(I)}L_{e,\pi}.
\label{eq:V61-local-letters}
\end{equation}
Thus \(M_e\) is the complete local moment carrier.  Tensor signs below
refer to the coordinate tensor product; within a partition letter the
blocks are on disjoint replica-row factors.

For a word \(\boldsymbol\pi=(\pi_e)_{e\in E}\), define
\begin{equation}
 W(\boldsymbol\pi):=\bigotimes_{e\in E}L_{e,\pi_e},
 \qquad
 j_r(\boldsymbol\pi):=\bigvee_{e\le r}\pi_e,
 \qquad
 j_0(\boldsymbol\pi):=\widehat0.
\label{eq:V61-word-join}
\end{equation}
All identities in this chapter are finite algebraic identities.  No
absolute value, physical-output sum, or analytic limit is used in their
proof.

\begin{lemma}[Linked-word identity]
\label{lem:V61-linked-word}
The connected \(I\)-row carrier of the complete coordinate product is
\begin{equation}
 \boxed{
 K_I^{E,\mathrm{conn}}
 =\sum_{\substack{\boldsymbol\pi\in\Pi(I)^E\\
                   \bigvee_{e\in E}\pi_e=\widehat1}}
   W(\boldsymbol\pi).}
\label{eq:V61-linked-word}
\end{equation}
In particular, a coordinate word occurs with coefficient one if its
local row partitions jointly connect \(I\), and with coefficient zero
otherwise.
\end{lemma}

\begin{proof}
For \(\tau\in\Pi(I)\), the moment in which different blocks of
\(\tau\) use independent replicas is
\begin{equation}
 M_E(\tau)
 =\sum_{\substack{\boldsymbol\pi\in\Pi(I)^E\\
                   \pi_e\le\tau\ \forall e}}
   W(\boldsymbol\pi).
\label{eq:V61-replica-moment}
\end{equation}
Indeed, at every coordinate the moment--cumulant formula allows exactly
the local cumulant blocks which refine \(\tau\), and coordinate
independence then tensors the local expansions.  M\"obius inversion on
\(\Pi(I)\) gives
\[
 K_I^{E,\mathrm{conn}}
 =\sum_{\tau\in\Pi(I)}
   \mu(\tau,\widehat1)M_E(\tau).
\]
The coefficient of a fixed word with
\(\sigma=\bigvee_e\pi_e\) is therefore
\[
 \sum_{\tau:\,\sigma\le\tau\le\widehat1}
       \mu(\tau,\widehat1)
 =\mathbf1_{\{\sigma=\widehat1\}},
\]
which is the defining cancellation identity for the M\"obius function.
This proves \eqref{eq:V61-linked-word}.
\end{proof}

\section{Exact disintegration at the first connected join}
\label{sec:V61-first-time}

For \(r\ge1\) and \(\rho\in\Pi(I)\), define the exact prefix-join
carrier
\begin{equation}
 C_{<r}(\rho)
 :=\sum_{\substack{(\pi_e)_{e<r}\\
                    \bigvee_{e<r}\pi_e=\rho}}
     \bigotimes_{e<r}L_{e,\pi_e},
\label{eq:V61-prefix-carrier}
\end{equation}
where the empty prefix equals the identity operator \(\mathbf 1\) when \(r=1\) and
\(\rho=\widehat0\), and equals zero otherwise.  Define the local
joining letter and complete suffix by
\begin{align}
 J_r(\rho)
 &:=\sum_{\substack{\pi\in\Pi(I)\\
                     \rho\vee\pi=\widehat1}}
       L_{r,\pi},
\label{eq:V61-joining-letter}\\
 M_{>r}&:=\bigotimes_{e>r}M_e.
\label{eq:V61-complete-suffix}
\end{align}

\begin{theorem}[First-connectivity formula]
\label{thm:V61-first-connectivity}
For every finite ordered coordinate set,
\begin{equation}
 \boxed{
 K_I^{E,\mathrm{conn}}
 =\sum_{r=1}^n
   \sum_{\rho<\widehat1}
   C_{<r}(\rho)\otimes J_r(\rho)\otimes M_{>r}.}
\label{eq:V61-first-connectivity}
\end{equation}
Every connected coordinate word appears exactly once.  Its index \(r\)
is the first coordinate for which \(j_r=\widehat1\), and \(\rho\) is
its exact prefix join \(j_{r-1}\).
\end{theorem}

\begin{proof}
Take a word in \eqref{eq:V61-linked-word}.  Since
\(j_0=\widehat0<\widehat1\) and \(j_n=\widehat1\), the monotone chain
\(j_0\le\cdots\le j_n\) has a unique first index \(r\) at which it
equals \(\widehat1\).  Then
\(\rho=j_{r-1}<\widehat1\) and
\(\rho\vee\pi_r=\widehat1\).  The prefix of the word occurs once in
\(C_{<r}(\rho)\), and its transition letter occurs once in
\(J_r(\rho)\).  After time \(r\), every local partition is allowed;
expanding
\[
 \bigotimes_{e>r}M_e
 =\bigotimes_{e>r}\sum_{\pi_e}L_{e,\pi_e}
\]
recovers its suffix once.  Conversely, every word on the right has
prefix join \(\rho<\widehat1\), becomes connected at \(r\), and stays
connected.  Hence the displayed classes are disjoint and exhaust
\eqref{eq:V61-linked-word}.
\end{proof}

The formula is independent of the chosen coordinate order: changing the
order changes only which coordinate is designated as the first joining
coordinate, not the total operator on the left.

\section{Exact factorization of an unconnected prefix}
\label{sec:V61-prefix-factorization}

For a nonempty block \(B\subseteq I\), let
\(K_B^{<r,\mathrm{conn}}\) denote the connected \(B\)-row cumulant
carrier formed from the coordinates \(e<r\).  For \(|B|=1\), this means
the complete one-row prefix moment, rather than a centered variable.

\begin{theorem}[Prefix-join factorization]
\label{thm:V61-prefix-factorization}
For every \(\rho<\widehat1\),
\begin{equation}
 \boxed{
 C_{<r}(\rho)
 =\bigotimes_{B\in\rho}K_B^{<r,\mathrm{conn}}.}
\label{eq:V61-prefix-factorization}
\end{equation}
Thus an exact prefix-join carrier contains no hidden sum over ways of
connecting different blocks of \(\rho\).
\end{theorem}

\begin{proof}
If \(\bigvee_{e<r}\pi_e=\rho\), then every \(\pi_e\le\rho\).  Write
\(\pi_e|_B\) for its restriction to \(B\in\rho\).  From
\eqref{eq:V61-local-letters},
\[
 L_{e,\pi_e}
 =\bigotimes_{B\in\rho}L^B_{e,\pi_e|_B}.
\]
Moreover,
\[
 \bigvee_{e<r}\pi_e=\rho
 \quad\Longleftrightarrow\quad
 \bigvee_{e<r}(\pi_e|_B)=\widehat1_B
 \quad\hbox{for every }B\in\rho.
\]
Consequently the choices of restricted partition words factor over the
blocks \(B\).  Applying \Cref{lem:V61-linked-word} separately on each
block gives exactly \eqref{eq:V61-prefix-factorization}.  The singleton
case is the ordinary one-row moment expansion and is included in the
same calculation.
\end{proof}

\begin{corollary}[No prefix support multiplicity]
\label{cor:V61-no-prefix-multiplicity}
Suppose connected marked-tree bounds are available through order
\(m-1\).  Every proper-prefix carrier in
\eqref{eq:V61-first-connectivity} is a tensor product of those complete
lower-order connected carriers and singleton moment carriers.  In
particular, it must be estimated after the exact join \(\rho\) is fixed;
there is no additional count of witnessing coordinates or witnessing
edges.
\end{corollary}

\begin{proof}
Every block of a proper partition \(\rho<\widehat1\) has size at most
\(m-1\).  Apply \Cref{thm:V61-prefix-factorization}.  Exact joins and the
unique first time in \Cref{thm:V61-first-connectivity} already perform
the only combinatorial classification.
\end{proof}

\section{The complete suffix is contractive}
\label{sec:V61-suffix}

\begin{lemma}[Complete moment suffix contraction]
\label{lem:V61-suffix-contraction}
On the raw replica coefficient Hilbert space, each complete local Wilson
moment \(M_e\) is the expectation of a tensor product of unitary matrix
coefficients and hence
\begin{equation}
 \|M_e\|_{\mathrm{op}}\le1,
 \qquad
 \|M_{>r}\|_{\mathrm{op}}\le1.
\label{eq:V61-suffix-contraction}
\end{equation}
For every source carrier \(Q\), as long as the sole final payer has not
been allocated to the suffix,
\begin{equation}
 |Q\otimes M_{>r}|
 =|Q|\otimes|M_{>r}|
 \le |Q|\otimes I.
\label{eq:V61-positive-suffix}
\end{equation}
The same inequality survives the previously used completely positive
isotypic pinching and a one-sided capacity allocation whose payer lies
in \(Q\).  If the payer lies at a suffix coordinate \(s>r\), the exact
V2 allocation instead leaves one factor \(D_s^{\rm pay}\) and complete
moment contractions at all other suffix coordinates; that paid factor
must be bounded by the one-link first-moment compiler and is not covered
by \eqref{eq:V61-positive-suffix}.
\end{lemma}

\begin{proof}
A tensor product of unitary representation matrices has operator norm
one.  Haar expectation is a unital completely positive contraction, so
its expectation \(M_e\) has norm at most one.  Coordinate independence
gives the second inequality in \eqref{eq:V61-suffix-contraction}.  The
absolute value of a tensor product factorizes, and
\(|M_{>r}|\le I\), proving \eqref{eq:V61-positive-suffix}.  Completely
positive pinching is order preserving.  Finally, if the V2 allocation
places the paid weight in \(Q\), it uses operator norms on the suffix
spectators, so their contraction cannot increase the result.  If it
places the payer at \(s>r\), the corresponding tensor factor is
\(D_s^{\rm pay}\), not \(M_s\); the last assertion merely records this
separate exact alternative and makes no contraction claim for it.
\end{proof}

\begin{firewall}[Where contraction is used]
\label{fire:V61-contraction-order}
The contraction in \Cref{lem:V61-suffix-contraction} applies only after
the complete sum \(M_e=\sum_\pi L_{e,\pi}\) has been formed.  It gives no
termwise estimate on \(L_{e,\pi}\), and it does not assert that
\(R_e=M_e-K_e\) is a contraction.  Thus it does not repeat the invalid
step isolated in V36 or the graphwise merge-mark regression of V59.
Nor may a paid suffix carrier \(D_s^{\rm pay}\) be replaced by \(M_s\).
\end{firewall}

\section{Elimination of the restricted accepted transfer}
\label{sec:V61-eliminate-transfer}

\begin{theorem}[Exact replacement of the V36 accepted suffix]
\label{thm:V61-replace-accepted-suffix}
Group the full connected carrier by first global connectivity as in
\eqref{eq:V61-first-connectivity}.  For every fixed source type
\((\rho,\pi)\) satisfying
\begin{equation}
 \rho<\widehat1,
 \qquad
 \rho\vee\pi=\widehat1,
\label{eq:V61-source-admissible}
\end{equation}
its exact contribution is
\begin{equation}
 \boxed{
 \mathcal S_{\rho,\pi}
 =\sum_{r=1}^n
 C_{<r}(\rho)\otimes L_{r,\pi}\otimes M_{>r}.}
\label{eq:V61-source-packet}
\end{equation}
No higher-priority letter is deleted from the suffix.  In particular,
later four-row cumulant letters are already included once inside
\(M_{>r}\), and no accepted-transfer Green constant occurs.
\end{theorem}

\begin{proof}
Expand \(J_r(\rho)\) in
\eqref{eq:V61-first-connectivity} and select the fixed letter \(\pi\).
This gives \eqref{eq:V61-source-packet}.  A later local partition cannot
change the fact that the cumulative join is already
\(\widehat1\), so every later partition belongs to this unique first-time
class.  Their complete sum is \(M_{>r}\).  Disjointness in
\Cref{thm:V61-first-connectivity} proves that none is double counted.
\end{proof}

\begin{corollary}[The V60 Green target is not intrinsic]
\label{cor:V61-G33-not-needed}
The bound \(\mathfrak G_{33}\le C\) in
\eqref{eq:V60-G33-target} is sufficient for the old priority
decomposition but is unnecessary for the first-connectivity
decomposition.  The latter requires only a support-uniform bound for the
once-paid source packet in \eqref{eq:V61-source-packet}.  Every nonpayer
suffix factor has norm at most one by
\Cref{lem:V61-suffix-contraction}; a suffix payer is one ordinary
\(D_s^{\rm pay}\) factor and never a restricted accepted transfer.
\end{corollary}

\begin{proof}
The old automaton remained in its accepted state only while excluding a
later four-block letter, and hence propagated by
\(\bigotimes_{s>r}(M_s-K_s)\).  The new class imposes no condition after
the first globally connected join and propagates by
\(\bigotimes_{s>r}M_s\).  Apply
\Cref{lem:V61-suffix-contraction} to every nonpayer suffix factor and the
exact V2 allocation to the unique payer factor.  This removes the
operator \(\mathcal A^{33}\) rather than proving a bound for it.
\end{proof}

\begin{assessment}[Exact telescoping interpretation]
\label{ass:V61-telescope}
Expanding each later \(M_s=R_s+K_s\) shows exactly what the old priority
packet omitted: after a \([3,3]\) first-connectivity source, the term
with only \(R\)'s is the V36 accepted suffix, while every term whose
first later \(K_s\) occurs at \(s\) is the corresponding completion
term.  Their algebraic sum is the full product of the \(M_s\)'s.  The
completion must occur before capacity; estimating the terms separately
would recreate the Green problem.
\end{assessment}

\section{The \([3,3]\) first-connectivity source}
\label{sec:V61-33-source}

Now take \(I=\{1,2,3,4\}\).  For \(a\in I\), write
\(\rho_a=(I\setminus\{a\})|\{a\}\) for the triple--singleton
partition.  Distinct \(a,b\) satisfy
\(\rho_a\vee\rho_b=\widehat1\).

\begin{proposition}[Exact \([3,3]\) first-connectivity packet]
\label{prop:V61-33-packet}
The part of the connected four-row carrier whose prefix join is
\(\rho_a\) and whose first joining letter is \(\rho_b\), with
\(a\ne b\), is
\begin{equation}
 \boxed{
 \mathcal S^{33}_{a,b}
 =\sum_{r=1}^n
 \left(
 K_{I\setminus\{a\}}^{<r,\mathrm{conn}}
 \otimes M_{\{a\}}^{<r}
 \right)
 \otimes
 \left(
 k_r(I\setminus\{b\})\otimes k_r(\{b\})
 \right)
 \otimes M_{>r}.}
\label{eq:V61-33-packet}
\end{equation}
It is counted without a choice of an earlier witnessing coordinate and
without a restricted suffix.
\end{proposition}

\begin{proof}
Use \Cref{thm:V61-prefix-factorization} with
\(\rho=\rho_a\); its two blocks give the two factors in the first
parenthesis.  Formula \eqref{eq:V61-local-letters} gives the second
parenthesis for \(L_{r,\rho_b}\).  Then apply
\Cref{thm:V61-replace-accepted-suffix}.  The exact prefix join already
sums every possible earlier realization of the first triple type, and
the first-connectivity time \(r\) is unique.
\end{proof}

\begin{proposition}[The graph part of the \([3,3]\) source is already
closed]
\label{prop:V61-33-graph-closed}
Assume the two ternary factors in a fixed summand of
\eqref{eq:V61-33-packet} have been converted by the established ternary
MTA into trees on the two distinct triples.  After the single final
resolvent is allocated by
\Cref{thm:V35-two-junction-reduction}, their four normalized marks are
bounded by a sum of at most four quartic spanning-tree weights.  The
complete nonpayer suffix costs no mark and no support factor; a suffix
payer remains the single paid-moment alternative described in
\Cref{lem:V61-suffix-contraction}.
\end{proposition}

\begin{proof}
The two triples have two common rows and their union is \(I\).  The
union of their two tree edges is therefore a connected four-edge
multigraph on four vertices.  Apply \Cref{lem:V35-surplus-edge}; one edge
is surplus and deletion leaves one of at most four spanning trees.  The
one-resolvent assertion is exactly
\Cref{thm:V35-two-junction-reduction}.  Finally use
\Cref{lem:V61-suffix-contraction} on nonpayer suffix factors.  The
suffix-payer alternative retains exactly one \(D_s^{\rm pay}\) and does
not change the row graph.
\end{proof}

\begin{firewall}[The remaining \([3,3]\) source estimate]
\label{fire:V61-33-source-gate}
\Cref{prop:V61-33-graph-closed} is not yet a bound for the sum over
\(r\) in \eqref{eq:V61-33-packet}.  Such a bound must show, before
absolute values, that the exact local transition stock and the complete
ternary prefix stock share the row masses and the one final payer without
losing a count of possible \(r\).  V35 identifies the correct
two-junction allocation, but does not prove this support-uniform source
summation.  This finite source normalization, not accepted-state
propagation, is the current \([3,3]\) gate.
\end{firewall}

\section{A finite first-connectivity criterion for quartic MTA}
\label{sec:V61-quartic-criterion}

For \(\rho<\widehat1\) and
\(\pi\in\Pi_4\) with \(\rho\vee\pi=\widehat1\), let
\(\mathfrak S_{\rho,\pi}\) denote the normalized capacity of the full
packet \eqref{eq:V61-source-packet} after the exact V2 allocation of the
sole final payer.  Thus source-payer terms use complete contractive
suffixes, while a suffix-payer term contains exactly one
\(D_s^{\rm pay}\) and complete moments at all other suffix coordinates.
There are only finitely many ordered partition-pair types, although each
packet still contains its support-uniform joining-coordinate and payer
sums.

\begin{theorem}[First-connectivity source criterion]
\label{thm:V61-source-criterion}
Suppose that for every admissible quartic pair \((\rho,\pi)\), the exact
source sum satisfies
\begin{equation}
 \mathfrak S_{\rho,\pi}
 \le C_{\rho,\pi}
 \sum_{T\in\mathfrak T_4}\mathcal W_T,
\label{eq:V61-source-criterion}
\end{equation}
uniformly in coordinate support, representations, physical
multiplicities, and volume.  Assume the source packet is formed before
positive coefficient envelopes and the final resolvent is allocated
exactly once.  Then the all-weak normalized quartic carrier is bounded by
\begin{equation}
 \left(\sum_{\substack{\rho<\widehat1,\ \pi\in\Pi_4\\
                       \rho\vee\pi=\widehat1}}
 C_{\rho,\pi}\right)
 \sum_{T\in\mathfrak T_4}\mathcal W_T.
\label{eq:V61-quartic-sum}
\end{equation}
Consequently the unaffected weak-cut and large-link branches of the V52
threshold compiler would restore global quartic MTA.
\end{theorem}

\begin{proof}
Use \Cref{thm:V61-first-connectivity} and expand each joining letter
into the finitely many admissible \(\pi\)'s.  Apply the exact V2 payer
allocation and then the assumed bound to the resulting full packet.  By
\Cref{lem:V61-suffix-contraction}, all nonpayer complete moments may be
discarded after the packet has been formed; the unique paid suffix factor
is retained in the definition of \(\mathfrak S_{\rho,\pi}\).  Sum the
finite partition-pair constants.  The first-time disintegration ensures
that no word or coordinate is counted more than once before the finite
triangle inequality.  The last assertion is the normalized quartic
threshold theorem, whose non-all-weak branches were not withdrawn in
V59.
\end{proof}

\begin{assessment}[Why this criterion is smaller than V60]
\label{ass:V61-smaller-gate}
V60 required both first-entry source bounds and support-uniform Green
bounds for accepted transfers.  \Cref{thm:V61-source-criterion} requires
only finitely many first-connectivity once-paid packet bounds.  Every
nonpayer complete suffix contraction is a theorem, not an additional
hypothesis; the one possible paid suffix moment is part of the packet
bound.  The price is that topology priority is abandoned: all local
partition types must be audited together according to the pair
\((\rho,\pi)\).
\end{assessment}

\section{V61 ledger and V62 mandate}
\label{sec:V61-ledger}

\begin{theorem}[Integrated compact-series V61 statement]
\label{thm:V61-integrated}
Preserving compact V1--V60, V61 proves:
\begin{enumerate}[label=\textup{(V61.\arabic*)},leftmargin=6.7em]
\item the linked-word formula selecting exactly those local partition
      words whose cumulative join is connected;
\item the multiplicity-free disintegration at the unique first globally
      connected coordinate;
\item exact factorization of every proper prefix-join carrier into
      lower-order connected cumulants;
\item contraction of every nonpayer factor in the unrestricted complete
      moment suffix and exact isolation of the single suffix-payer
      alternative;
\item elimination, rather than estimation, of the V36 restricted
      accepted transfer and the V60 Green target;
\item the exact \([3,3]\) first-connectivity packet and closure of its
      graph/surplus-edge part; and
\item a finite source criterion for restoring quartic MTA.
\end{enumerate}
\end{theorem}

\begin{proof}
These are \Cref{lem:V61-linked-word,
thm:V61-first-connectivity,
thm:V61-prefix-factorization,
lem:V61-suffix-contraction,
thm:V61-replace-accepted-suffix,
prop:V61-33-packet,
prop:V61-33-graph-closed,
thm:V61-source-criterion}.
\end{proof}

\begin{openproblem}[V62 mandate: enumerate and compile the quartic
first-connectivity sources]
\label{op:V62-mandate}
\begin{enumerate}
\item Enumerate the finite \(S_4\)-orbits of admissible pairs
      \((\rho,\pi)\) with
      \(\rho<\widehat1\) and
      \(\rho\vee\pi=\widehat1\).
\item For every orbit, insert
      \eqref{eq:V61-prefix-factorization} and identify which established
      binary, ternary, embedded-four, collision, or private-tail compiler
      applies without a graphwise merge mark.
\item For the \((3+1,3+1)\) orbit, prove or regress the support-uniform
      sum over the unique joining coordinate in
      \eqref{eq:V61-33-packet}, preserving the shared row masses and the
      single final payer.
\item State the smallest source orbit not closed by the existing
      compilers.  If a source estimate becomes circular, move upstream
      to its local cumulant/M\"obius identity before introducing any new
      norm.
\end{enumerate}
\end{openproblem}

\begin{firewall}[Claim boundary after compact V61]
\label{fire:V61-current-boundary}
V61 removes the accepted-suffix Green obstruction by an exact algebraic
regrouping.  It does not yet prove all of the finite quartic source
bounds \eqref{eq:V61-source-criterion}; in particular, the support-uniform
\([3,3]\) source sum remains open.  Therefore global quartic and quintic
MTA, all-order MTA, WUFC, CPM, continuum Yang--Mills existence, and a
positive mass gap remain open.  The full Yang--Mills mass gap remains
open.
\end{firewall}

\section{Append-only record for compact V61}
\label{sec:V61-record}

V61 was appended to the complete compact V60 file.  The exact V60 parent
had (23992) newline-delimited source records, (814488) bytes, and
SHA--256 digest
\[
\texttt{91ecee09fee6b521ac8301e258e4775239112fbc5810ec023e81524a3069b099}.
\]
No compact V60 mathematical content was changed.  The sole final
\verb|\end{document}| is restored below.

% =============================================================================
% END APPENDED COMPACT VERSION V61 MATERIAL
% =============================================================================

% =============================================================================
% BEGIN APPENDED COMPACT VERSION V62 MATERIAL
% =============================================================================

\chapter{The twelve quartic first-connectivity orbits and the
ordered-simplex compiler}
\label{chap:V62-orbits}

\section{Purpose}
\label{sec:V62-purpose}

V61 replaced every restricted accepted suffix by a complete-moment
suffix and reduced quartic MTA to the once-paid first-connectivity
packets \(\mathcal S_{\rho,\pi}\).  This chapter performs the finite
partition-lattice audit demanded there.  It proves that there are twelve,
and only twelve, simultaneous-\(S_4\) source orbits.  It also proves two
general facts:
\begin{enumerate}[label=\textup{(\roman*)}]
\item the union of block trees for every admissible pair is already a
      connected four-row multigraph with at most two surplus edges; and
\item every ordering condition on joining and payer coordinates is a
      subset of a Cartesian product and therefore costs no support
      factor once the coefficient estimate has been resolved into
      nonnegative \(\ell^1\) coordinate stocks.
\end{enumerate}
Consequently neither row-graph topology nor ordered support summation is
the remaining quartic obstruction.  What remains is to prove that the
exact once-paid coefficient carrier admits the required compatible
stock resolution before product-state capacity is taken.

\section{The complete orbit census}
\label{sec:V62-census}

For a partition \(\sigma\in\Pi_4\), write \(t(\sigma)\) for the
nonincreasing list of its block sizes and \(b(\sigma)\) for its number
of blocks.  Two admissible ordered pairs \((\rho,\pi)\) are equivalent
when one simultaneous permutation of the four rows sends both
partitions to the other pair.

\begin{theorem}[Twelve-orbit census]
\label{thm:V62-twelve-orbits}
There are \(104\) ordered pairs
\begin{equation}
 (\rho,\pi)\in
 \bigl(\Pi_4\setminus\{\widehat1\}\bigr)\times\Pi_4,
 \qquad
 \rho\vee\pi=\widehat1,
\label{eq:V62-admissible-pairs}
\end{equation}
and they form exactly the following twelve simultaneous-\(S_4\)
orbits.  The multiplicity is the cardinality of the orbit.  The integers
\[
 p=4-b(\rho),\qquad q=4-b(\pi),\qquad
 s=p+q-3
\]
are respectively the number of edges in block-spanning forests for the
prefix and transition and their surplus over a quartic tree.
\begin{center}
\small
\begin{tabular}{c c c l l r r r r}
\toprule
\#&\(t(\rho)\)&\(t(\pi)\)&\(\rho\)&\(\pi\)&mult.&\(p\)&\(q\)&\(s\)\\
\midrule
1&\(1+1+1+1\)&\(4\)&\(1|2|3|4\)&\(1234\)&1&0&3&0\\
2&\(2+1+1\)&\(2+2\)&\(12|3|4\)&\(13|24\)&12&1&2&0\\
3&\(2+1+1\)&\(3+1\)&\(12|3|4\)&\(134|2\)&12&1&2&0\\
4&\(2+1+1\)&\(4\)&\(12|3|4\)&\(1234\)&6&1&3&1\\
5&\(2+2\)&\(2+1+1\)&\(12|34\)&\(13|2|4\)&12&2&1&0\\
6&\(2+2\)&\(2+2\)&\(12|34\)&\(13|24\)&6&2&2&1\\
7&\(2+2\)&\(3+1\)&\(12|34\)&\(123|4\)&12&2&2&1\\
8&\(2+2\)&\(4\)&\(12|34\)&\(1234\)&3&2&3&2\\
9&\(3+1\)&\(2+1+1\)&\(123|4\)&\(14|2|3\)&12&2&1&0\\
10&\(3+1\)&\(2+2\)&\(123|4\)&\(12|34\)&12&2&2&1\\
11&\(3+1\)&\(3+1\)&\(123|4\)&\(124|3\)&12&2&2&1\\
12&\(3+1\)&\(4\)&\(123|4\)&\(1234\)&4&2&3&2\\
\bottomrule
\end{tabular}
\end{center}
The orbit multiplicities sum to \(104\).
\end{theorem}

\begin{proof}
If \(\rho=\widehat0\), no proper transition connects four discrete
blocks, so \(\pi=\widehat1\).  This gives orbit 1.

There are six partitions \(\rho\) of type \(2+1+1\).  Fix
\(\rho=12|3|4\).  To connect its three blocks, a transition of type
\(2+2\) must be \(13|24\) or \(14|23\); a transition of type \(3+1\)
must contain \(3,4\) and one of \(1,2\); and the one-block transition is
always allowed.  Thus this prefix type contributes
\[
 6(2+2+1)=30
\]
pairs, represented by orbits 2--4, with multiplicities \(12,12,6\).

There are three partitions \(\rho\) of type \(2+2\).  Fix
\(\rho=12|34\).  A transition of type \(2+1+1\) must use one of the four
cross pairs; a transition of type \(2+2\) must be one of the two other
perfect matchings; every one of the four triple--singleton partitions
joins the two prefix blocks; and the one-block transition is allowed.
This contributes
\[
 3(4+2+4+1)=33
\]
pairs, represented by orbits 5--8, with multiplicities
\(12,6,12,3\).

There are four partitions \(\rho\) of type \(3+1\).  Fix
\(\rho=123|4\).  A transition of type \(2+1+1\) must pair row \(4\)
with one of \(1,2,3\); each of the three perfect matchings contains such
a cross pair; a transition of type \(3+1\) may be any of the three
triple partitions other than \(\rho\); and the one-block transition is
allowed.  This contributes
\[
 4(3+3+3+1)=40
\]
pairs, represented by orbits 9--12, with multiplicities
\(12,12,12,4\).  Hence the total is
\(1+30+33+40=104\).

Within each fixed pair of block-size types, the displayed incidence
pattern is unique up to simultaneous row permutation: it is determined
by the nonzero intersection sizes between the blocks of \(\rho\) and
\(\pi\), up to permuting equal-sized blocks.  The representatives above
have distinct ordered block-size types or distinct intersection
patterns.  They are therefore all the orbits, and the displayed
multiplicities, whose sum is \(104\), prove that none is missing.
\end{proof}

\begin{remark}[Ordered, not symmetric, source types]
\label{rem:V62-ordered}
Interchanging \(\rho\) and \(\pi\) is not an allowed equivalence:
\(\rho\) is a cumulative prefix join, whereas \(\pi\) is one local
transition partition.  Thus orbits 2 and 5, and orbits 3 and 9, are
genuinely different analytic source packets even though their
unlabelled union graphs have the same edge count.
\end{remark}

\section{Exact source factors in the twelve rows}
\label{sec:V62-source-factors}

For a block \(B\), abbreviate
\[
 K_B=K_B^{<r,\mathrm{conn}},
 \qquad
 M_i=M_{\{i\}}^{<r},
 \qquad
 k_r(B)=\text{the local connected \(B\)-row letter at \(r\)}.
\]
Then \Cref{thm:V61-prefix-factorization} turns the twelve
representatives into the following source cores; every row is followed
by the same complete suffix \(M_{>r}\).
\begin{center}
\small
\begin{tabular}{c l l}
\toprule
\#&prefix carrier \(C_{<r}(\rho)\)&transition \(L_{r,\pi}\)\\
\midrule
1&\(M_1M_2M_3M_4\)&\(k_r(1234)\)\\
2&\(K_{12}M_3M_4\)&\(k_r(13)k_r(24)\)\\
3&\(K_{12}M_3M_4\)&\(k_r(134)k_r(2)\)\\
4&\(K_{12}M_3M_4\)&\(k_r(1234)\)\\
5&\(K_{12}K_{34}\)&\(k_r(13)k_r(2)k_r(4)\)\\
6&\(K_{12}K_{34}\)&\(k_r(13)k_r(24)\)\\
7&\(K_{12}K_{34}\)&\(k_r(123)k_r(4)\)\\
8&\(K_{12}K_{34}\)&\(k_r(1234)\)\\
9&\(K_{123}M_4\)&\(k_r(14)k_r(2)k_r(3)\)\\
10&\(K_{123}M_4\)&\(k_r(12)k_r(34)\)\\
11&\(K_{123}M_4\)&\(k_r(124)k_r(3)\)\\
12&\(K_{123}M_4\)&\(k_r(1234)\)\\
\bottomrule
\end{tabular}
\end{center}
Products in this table are tensor products on disjoint row blocks.
Singleton factors are displayed because they retain complete Wilson
moments and can carry the sole payer; they contribute no row edge.

\begin{corollary}[No higher-than-ternary prefix input]
\label{cor:V62-prefix-order}
Every proper prefix in the quartic first-connectivity formula uses only
the already established binary and ternary connected carriers, together
with singleton moments.  A four-row connected local carrier can occur
only in the transition column, namely in orbits \(1,4,8,12\).
\end{corollary}

\begin{proof}
This is immediate from the block sizes of a proper partition of four
rows and the source table.
\end{proof}

\section{The union-forest theorem}
\label{sec:V62-union-forest}

For a partition \(\sigma\), choose independently a spanning tree on each
nonsingleton block and let \(F(\sigma)\) be their disjoint union.
Its connected-component partition is exactly \(\sigma\), and
\[
 |E(F(\sigma))|=\sum_{B\in\sigma}(|B|-1)=4-b(\sigma).
\]

\begin{theorem}[Join equals connectivity of block forests]
\label{thm:V62-join-forest}
For any \(\rho,\pi\in\Pi_4\), the multigraph
\begin{equation}
 G_{\rho,\pi}:=F(\rho)\sqcup F(\pi)
\label{eq:V62-union-graph}
\end{equation}
has connected-component partition \(\rho\vee\pi\).  Hence every
admissible pair in \eqref{eq:V62-admissible-pairs} produces a connected
multigraph with
\begin{equation}
 |E(G_{\rho,\pi})|=p+q=3+s,
 \qquad 0\le s\le2.
\label{eq:V62-surplus-count}
\end{equation}
\end{theorem}

\begin{proof}
Two vertices are connected in \(F(\rho)\) exactly when they lie in one
block of \(\rho\), and similarly for \(\pi\).  A path in the union graph
is therefore a chain of identifications, each licensed by a block of
\(\rho\) or of \(\pi\).  The equivalence relation generated by precisely
those identifications is the lattice join \(\rho\vee\pi\).  This proves
the component assertion.

The edge count is the sum of the two forest edge counts.  A connected
four-vertex multigraph has at least three edges, so \(s\ge0\).  Since
\(\rho\ne\widehat1\), \(p\le2\), while \(q\le3\), hence \(s\le2\).
\end{proof}

\begin{corollary}[Uniform surplus deletion for all twelve orbits]
\label{cor:V62-surplus-deletion}
Attach a normalized mark \(0\le\theta_a\le1\) to every edge occurrence
of \(G_{\rho,\pi}\).  Then
\begin{equation}
 \prod_{a\in E(G_{\rho,\pi})}\theta_a
 \le
 \sum_{\substack{T\subseteq G_{\rho,\pi}\\
                  T\ {\rm a\ spanning\ tree}}}
 \prod_{a\in E(T)}\theta_a,
\label{eq:V62-surplus-deletion}
\end{equation}
and the sum contains at most
\(\binom53=10\) trees, counting parallel edge occurrences separately.
\end{corollary}

\begin{proof}
The union graph is connected by
\Cref{thm:V62-join-forest}, so it contains a spanning tree.  For every
such tree, deleting the other \(s\) edge occurrences can only increase
the product because all marks are at most one.  One displayed tree
would suffice for the inequality; summing all of them remains valid.
There are at most five edge occurrences and a spanning tree chooses
three, giving the stated uniform bound.
\end{proof}

\begin{assessment}[Graph topology is closed orbitwise]
\label{ass:V62-graph-closed}
The V35 surplus-edge lemma was the \((3+1,3+1)\) instance of
\Cref{cor:V62-surplus-deletion}.  The general statement closes the graph
conversion for every row of the census at once.  It does not manufacture
the normalized marks: those must still come from valid binary, ternary,
and selected-four coefficient estimates before this deletion is used.
\end{assessment}

\section{Ordered-simplex domination}
\label{sec:V62-simplex}

The unique joining coordinate and the optional later payer coordinate
introduce order constraints.  The following elementary lemma shows that
such constraints never create a support cardinality.

\begin{lemma}[Cartesian domination of an ordered coordinate sum]
\label{lem:V62-cartesian}
Let \(E\) be finite and let \(A\) be a finite set of labelled stock
occurrences.  For every \(a\in A\), let \(x_{a,e}\ge0\) and
\(X_a=\sum_{e\in E}x_{a,e}\).  For an arbitrary set
\(\Omega\subseteq E^A\),
\begin{equation}
 \boxed{
 \sum_{(e_a)_{a\in A}\in\Omega}
 \prod_{a\in A}x_{a,e_a}
 \le\prod_{a\in A}X_a.}
\label{eq:V62-cartesian}
\end{equation}
The occurrences in \(A\) may use identical sequences; they are still
treated as separately labelled stocks.
\end{lemma}

\begin{proof}
All summands are nonnegative, so enlarge \(\Omega\) to the full Cartesian
product \(E^A\).  The resulting finite sum factorizes:
\[
 \sum_{(e_a)\in E^A}\prod_{a\in A}x_{a,e_a}
 =\prod_{a\in A}\sum_{e_a\in E}x_{a,e_a}
 =\prod_{a\in A}X_a.
\]
\end{proof}

For a stock \(x\), write
\[
 X^{<r}=\sum_{e<r}x_e,\qquad
 X^{>r}=\sum_{e>r}x_e.
\]

\begin{corollary}[First-connectivity stock sum]
\label{cor:V62-first-stock}
For labelled prefix stocks \(a\in A\) and nonempty labelled transition
stocks \(b\in B\),
\begin{equation}
 \sum_{r\in E}
 \prod_{a\in A}X_a^{<r}
 \prod_{b\in B}x_{b,r}
 \le
 \prod_{a\in A}X_a\prod_{b\in B}X_b.
\label{eq:V62-first-stock}
\end{equation}
If the sole payer contributes a nonnegative density \(d_s\), then both
the source-payer case \(s\le r\) and suffix-payer case \(r<s\) obey the
same bound with the additional factor \(D=\sum_sd_s\).
\begin{equation}
 \sum_{\substack{r,s\in E\\r<s}}
 \prod_{a\in A}X_a^{<r}
 \prod_{b\in B}x_{b,r}\,d_s
 \le
 \prod_{a\in A}X_a\prod_{b\in B}X_b\,D.
\label{eq:V62-suffix-payer-stock}
\end{equation}
\end{corollary}

\begin{proof}
Expand each prefix sum in
\eqref{eq:V62-first-stock}.  The resulting coordinate tuples satisfy
\(e_a<r\), while all transition occurrences have coordinate \(r\).
They form a subset of the Cartesian product in
\Cref{lem:V62-cartesian}.  For
\eqref{eq:V62-suffix-payer-stock}, add the labelled payer occurrence and
the constraint \(r<s\).  The source-payer alternatives are identical
with their actual order constraints.  Discarding any of these constraints
gives the stated product of total stocks.
\end{proof}

\begin{firewall}[What the simplex lemma does not prove]
\label{fire:V62-simplex-boundary}
\Cref{lem:V62-cartesian} applies after the operator/coefficient packet
has been bounded by nonnegative coordinate densities.  It cannot justify
taking absolute values of individual cumulant letters, multiplying two
product-state capacities, or assigning one graph-edge mark to a
graph-difference operator.  Those are precisely the invalid operations
isolated in V2 and V59.
\end{firewall}

\section{Compiler audit and the remaining interface}
\label{sec:V62-compiler-audit}

The twelve rows fall into four transition families:
\begin{center}
\begin{tabular}{c l l}
\toprule
transition type&orbits&available fixed-block input\\
\midrule
\(4\)&\(1,4,8,12\)&V34 selected full four-row carrier\\
\(3+1\)&\(3,7,11\)&local ternary carrier; V35 for orbit 11\\
\(2+2\)&\(2,6,10\)&two disjoint local binary carriers\\
\(2+1+1\)&\(5,9\)&one local binary carrier\\
\bottomrule
\end{tabular}
\end{center}
All prefix factors have order at most three by
\Cref{cor:V62-prefix-order}.  Thus the individual block estimates needed
by the table already exist.  That statement alone does not allow their
capacities to be multiplied.

\begin{definition}[Compatible once-paid density interface]
\label{def:V62-density-interface}
An orbit has a compatible once-paid density interface if, after its
exact source core from \Cref{sec:V62-source-factors} and the V2
one-resolvent allocation have been formed, every positive resolved
carrier is bounded by a finite sum of terms with the following
properties:
\begin{enumerate}[label=\textup{(\alph*)}]
\item exactly one factor is estimated by product-state capacity and all
      other noncontractive factors by operator norm;
\item each required prefix-tree edge, transition-tree edge, and paid
      first moment is represented by a nonnegative coordinate stock with
      a support-uniform \(\ell^1\) total;
\item the coordinate restrictions are only the actual first-connectivity
      and payer-order restrictions; and
\item no row mass and no resolvent denominator is used twice.
\end{enumerate}
\end{definition}

\begin{theorem}[Orbitwise compilation from the density interface]
\label{thm:V62-interface-compiler}
If an orbit in \Cref{thm:V62-twelve-orbits} has the compatible once-paid
density interface, then its full packet
\(\mathfrak S_{\rho,\pi}\) satisfies the quartic source bound
\eqref{eq:V61-source-criterion}, with a constant independent of support.
At most ten quartic tree cards are produced for each resolved block-tree
choice.
\end{theorem}

\begin{proof}
Apply the one-sided capacity/norm choice in
\Cref{def:V62-density-interface}(a).  Sum the coordinate densities by
\Cref{cor:V62-first-stock}; a suffix payer uses
\eqref{eq:V62-suffix-payer-stock}.  No coordinate cardinality appears.
The block-tree edge occurrences form \(G_{\rho,\pi}\) by
\Cref{thm:V62-join-forest}.  Delete its at most two surplus normalized
marks using \Cref{cor:V62-surplus-deletion}.  The remaining three marks
are a quartic spanning-tree weight.  Conditions (c)--(d) preserve the
unique first-connectivity classification and the single final payer.
\end{proof}

\begin{assessment}[Exact status of the twelve rows]
\label{ass:V62-orbit-status}
For all twelve orbits, the finite lattice census, prefix factorization,
row-graph connectivity, surplus deletion, and ordered scalar stock
summation are closed.  What has not been proved is
\Cref{def:V62-density-interface} for any whole signed once-paid source
packet merely by citing separate lower-order capacities.  The first
useful test is orbit 9:
\[
 K_{123}^{<r,\mathrm{conn}}\otimes
 k_r(14)\otimes M_{>r},
\]
with singleton moments restored.  It has exactly three row edges and no
surplus, uses only an established ternary prefix and one binary
transition, and avoids the local four-row carrier and the two-triple
interaction.  Failure there would identify a coefficient/payer
obstruction upstream of every more complicated orbit.
\end{assessment}

\begin{proposition}[Unpaid scalar skeleton of orbit 9]
\label{prop:V62-orbit9-scalar}
Let \(h_{ij,e}\ge0\) be normalized local link densities and
\(H_{ij}=\sum_eh_{ij,e}\).  For either tree
\(T\in\mathfrak T(\{1,2,3\})\),
\begin{equation}
 \sum_r
 \prod_{ij\in E(T)}H_{ij}^{<r}\,h_{14,r}
 \le
 \prod_{ij\in E(T)}H_{ij}\,H_{14}.
\label{eq:V62-orbit9-scalar}
\end{equation}
The three edges on the right form a quartic spanning tree.  Hence orbit
9 has no scalar mark or support-count obstruction.
\end{proposition}

\begin{proof}
Use \Cref{cor:V62-first-stock} with the two prefix edge stocks of \(T\)
and the transition stock \(h_{14,r}\).  Since \(T\) connects rows
\(1,2,3\), adjoining edge \(14\) connects row \(4\) and creates no
cycle.
\end{proof}

\begin{firewall}[Current quartic bottleneck]
\label{fire:V62-current-gate}
\Cref{prop:V62-orbit9-scalar} is not yet the normalized operator
estimate for orbit 9.  The missing statement is that, after the exact
physical output resolution and one-resolvent allocation, the assembled
ternary-prefix/binary-transition carrier has a positive majorant
satisfying \Cref{def:V62-density-interface}.  In particular one must not
multiply the ternary and binary product-state capacities.  This is now
the smallest quartic source gate.
\end{firewall}

\section{V62 ledger and V63 mandate}
\label{sec:V62-ledger}

\begin{theorem}[Integrated compact-series V62 statement]
\label{thm:V62-integrated}
Preserving compact V1--V61, V62 proves:
\begin{enumerate}[label=\textup{(V62.\arabic*)},leftmargin=6.7em]
\item the exact census of \(104\) admissible quartic source pairs in
      twelve simultaneous-\(S_4\) orbits;
\item the exact lower-order prefix and local-transition factor in each
      orbit;
\item the join--union-forest theorem and a uniform deletion of at most
      two surplus marks;
\item Cartesian domination of every ordered joining/payer coordinate
      sum;
\item a sufficient compatible once-paid density interface for every
      orbit; and
\item removal of graph and scalar support-count obstructions from the
      minimal orbit-9 source.
\end{enumerate}
\end{theorem}

\begin{proof}
These are \Cref{thm:V62-twelve-orbits,
cor:V62-prefix-order,
thm:V62-join-forest,
cor:V62-surplus-deletion,
lem:V62-cartesian,
cor:V62-first-stock,
thm:V62-interface-compiler,
prop:V62-orbit9-scalar}.
\end{proof}

\begin{openproblem}[V63 mandate: the orbit-9 coefficient interface]
\label{op:V63-mandate}
\begin{enumerate}
\item Write the exact resolved carrier for
      \(K_{123}^{<r,\mathrm{conn}}\otimes k_r(14)\), with singleton
      moments and the complete suffix retained.
\item Apply the V2 blockwise one-resolvent allocation once and split only
      by the actual payer location: ternary prefix, binary transition, or
      complete-moment spectator.
\item In each branch, use product-state capacity on exactly one factor
      and operator norm on the other; prove the compatible coordinate
      densities required by \Cref{def:V62-density-interface}.
\item Sum joining and suffix-payer coordinates only through
      \Cref{cor:V62-first-stock}; do not introduce a raw pair count.
\item If the one-sided norm loses a needed row mark, move upstream to the
      common physical output pinching and test the smallest \(SU(2)\)
      representation before treating any more complicated orbit.
\end{enumerate}
\end{openproblem}

\begin{firewall}[Claim boundary after compact V62]
\label{fire:V62-current-boundary}
V62 proves a complete finite orbit census and closes graph conversion and
ordered scalar support summation for all quartic first-connectivity
sources.  It does not prove the compatible once-paid coefficient
interface, even for orbit 9.  Therefore the hypotheses of
\Cref{thm:V61-source-criterion} are not yet established.  Global
quartic and quintic MTA, all-order MTA, WUFC, CPM, continuum
Yang--Mills existence, and a positive mass gap remain open.  The full
Yang--Mills mass gap remains open.
\end{firewall}

\section{Append-only record for compact V62}
\label{sec:V62-record}

V62 was appended to the complete compact V61 file.  The exact V61 parent
had (24597) newline-delimited source records, (838615) bytes, and
SHA--256 digest
\[
\texttt{265b2ed092845cd4e853b45d271ef50bdca52ee1c94894c927b4a676722da4d5}.
\]
No compact V61 mathematical content was changed.  The sole final
\verb|\end{document}| is restored below.

% =============================================================================
% END APPENDED COMPACT VERSION V62 MATERIAL
% =============================================================================

% =============================================================================
% BEGIN APPENDED COMPACT VERSION V63 MATERIAL
% =============================================================================

\chapter{Assembled quotient-cluster joining and the repeated-row
regression}
\label{chap:V63-quotient-cluster}

\section{Purpose and correction of analytic granularity}
\label{sec:V63-purpose}

The twelve-orbit census of V62 is exact and remains useful as a finite
algebraic audit.  It is too fine for the next analytic estimate.
Already in orbit 9, multiplying the valid ternary-prefix and
binary-transition bounds loses one full mass of their common row.  The
loss is unbounded and is not repaired by the ordered-simplex lemma.

This chapter returns to the unexpanded joining letter
\(J_r(\rho)\) in V61.  The Leonov--Shiryaev identity identifies it
exactly as the connected cumulant of the block-products indexed by the
prefix partition \(\rho\).  Thus the twelve local partition orbits must
be assembled into four quotient-cluster families before capacity.  This
is the same proof-order principle which kept the V37--V39 cluster
estimates valid after the V59 graphwise regression.

\section{The joining letter is a block-product cumulant}
\label{sec:V63-joining-cumulant}

At a fixed coordinate \(r\), write \(X_{i,r}\) for the \(i\)-th row
variable in the fixed tensor order.  For a block \(B\subseteq I\), put
\[
 Z_{B,r}:=\prod_{i\in B}X_{i,r}.
\]
For \(\rho\in\Pi(I)\), let
\[
 \kappa_r^\rho
 :=
 \kappa\bigl(Z_{B,r}:B\in\rho\bigr)
\]
be the connected cumulant of the \(|\rho|\) composite block variables.
All products and cumulants here are coefficient carriers before a norm
or physical positive envelope is taken.

\begin{theorem}[Assembled joining-letter identity]
\label{thm:V63-joining-cumulant}
For every \(\rho<\widehat1\),
\begin{equation}
 \boxed{
 J_r(\rho)
 =
 \sum_{\substack{\pi\in\Pi(I)\\
                  \rho\vee\pi=\widehat1}}L_{r,\pi}
 =
 \kappa_r^\rho.}
\label{eq:V63-joining-cumulant}
\end{equation}
Consequently the first-connectivity formula has the exact unexpanded
form
\begin{equation}
 \boxed{
 K_I^{E,\mathrm{conn}}
 =
 \sum_{r=1}^n\sum_{\rho<\widehat1}
 \left(\bigotimes_{B\in\rho}
 K_B^{<r,\mathrm{conn}}\right)
 \otimes\kappa_r^\rho\otimes M_{>r}.}
\label{eq:V63-four-family-formula}
\end{equation}
\end{theorem}

\begin{proof}
Apply the Leonov--Shiryaev product-cumulant identity proved in V59 to
the family of products \(Z_{B,r}\), \(B\in\rho\).  A partition
\(\pi\) of the underlying row variables contributes if and only if
\(\pi\vee\rho=\widehat1\), and its contribution is
\[
 \prod_{D\in\pi}\kappa(X_{i,r}:i\in D)
 =L_{r,\pi}.
\]
This proves \eqref{eq:V63-joining-cumulant}.  Substitute it and the
prefix factorization
\eqref{eq:V61-prefix-factorization} into
\eqref{eq:V61-first-connectivity}.
\end{proof}

\begin{corollary}[Four analytic source families]
\label{cor:V63-four-families}
Up to row permutation, the proper prefix partition \(\rho\) has one of
four types:
\begin{center}
\begin{tabular}{c l l c c}
\toprule
\(t(\rho)\)&prefix connected factors&joining carrier&
internal edges&quotient edges\\
\midrule
\(1+1+1+1\)&four singleton moments&
local fourth cumulant&0&3\\
\(2+1+1\)&one binary cumulant and two moments&
ternary cumulant of three block-products&1&2\\
\(2+2\)&two binary cumulants&
covariance of two block-products&2&1\\
\(3+1\)&one ternary cumulant and one moment&
covariance of two block-products&2&1\\
\bottomrule
\end{tabular}
\end{center}
The twelve V62 orbits are the partition-letter expansion of these four
assembled joining carriers.  In particular, V62 orbits \(9\)--\(12\)
must be summed before estimating
\(\operatorname{Cov}_r(Z_{123,r},X_{4,r})\).
\end{corollary}

\begin{proof}
The four proper partition types are exhaustive.  Formula
\eqref{eq:V63-four-family-formula} gives the prefix column.  The order of
\(\kappa_r^\rho\) is \(b(\rho)\), giving the joining column.  A connected
block of size \(d\) uses \(d-1\) internal tree edges, so the prefix has
\(\sum_{B\in\rho}(|B|-1)=4-b(\rho)\) edges.  A tree on the
\(b(\rho)\) quotient blocks has \(b(\rho)-1\) edges.
\end{proof}

\begin{corollary}[Two-block moment-difference form]
\label{cor:V63-two-block-difference}
If \(\rho=\{C,D\}\) has two blocks, then
\begin{equation}
 \boxed{
 J_r(\rho)
 =
 \operatorname{Cov}_r(Z_{C,r},Z_{D,r})
 =
 M_r-M_r(\rho),}
\label{eq:V63-two-block-difference}
\end{equation}
where \(M_r(\rho)=
\mathbb E Z_{C,r}\otimes\mathbb E Z_{D,r}\) is the local moment with the
two block replicas independent.
\end{corollary}

\begin{proof}
The second cumulant is covariance.  Its moment--cumulant formula is the
difference of the joint moment and the product of the two marginal
moments.  Equivalently, when \(\rho\) has two blocks, a row partition
\(\pi\) fails to join them precisely when \(\pi\le\rho\); subtracting
the sum of those letters from \(M_r=\sum_\pi L_{r,\pi}\) gives
\eqref{eq:V63-two-block-difference}.
\end{proof}

\section{Internal forests plus a quotient tree}
\label{sec:V63-quotient-tree}

Let \(\rho\) be any proper partition of \(I=\{1,2,3,4\}\).  Choose a
spanning tree \(T_B\) inside each nonsingleton block \(B\in\rho\).
Choose a tree \(\zeta\) on the quotient vertex set \(\rho\), and for
each quotient edge \(BC\in E(\zeta)\) choose one row pair
\(i(BC)j(BC)\) with \(i(BC)\in B\) and \(j(BC)\in C\).

\begin{theorem}[Exact forest--quotient-tree lift at order four]
\label{thm:V63-forest-quotient}
The row graph
\begin{equation}
 T(\rho,\zeta)
 :=
 \left(\bigsqcup_{B\in\rho}T_B\right)
 \sqcup
 \{i(BC)j(BC):BC\in E(\zeta)\}
\label{eq:V63-lifted-tree}
\end{equation}
is a spanning tree on \(I\).  It has exactly
\begin{equation}
 (4-b(\rho))+(b(\rho)-1)=3
\label{eq:V63-edge-balance}
\end{equation}
edges.  Hence the assembled quotient-cluster formulation needs no
surplus-edge deletion.
\end{theorem}

\begin{proof}
Every \(T_B\) connects the rows within \(B\).  Contracting each such
tree to one vertex turns \eqref{eq:V63-lifted-tree} into \(\zeta\), which
is connected.  Therefore the lifted row graph is connected.  Its edge
count is \eqref{eq:V63-edge-balance}.  A connected graph on four vertices
with three edges is a tree.
\end{proof}

\begin{assessment}[Why quotient edges are the correct marks]
\label{ass:V63-quotient-marks}
The block cumulants in the prefix provide the internal forests.  The
assembled carrier \(\kappa_r^\rho\) must provide a connected tree on the
quotient blocks, with each quotient edge expanded as a sum of actual
crossing row influences.  This produces exactly the missing three-edge
quartic tree.  Expanding \(\kappa_r^\rho\) into its row-partition letters
before the crossing influences are formed destroys this normalization
and recreates the repeated-row deficit below.
\end{assessment}

\section{The orbit-9 repeated-row deficit}
\label{sec:V63-orbit9-deficit}

Take \(\rho=123|4\) and isolate the orbit-9 letter
\[
 K_{123}^{<r,\mathrm{conn}}\otimes k_r(14),
\]
with its singleton contractions understood.  Fix a ternary tree
\(T\in\mathfrak T(\{1,2,3\})\).  To give the separated-block method its
strongest possible input, suppose the V37 response-or-numerator table is
available with the prefix-restricted and one-coordinate densities:
\begin{align}
 \|\mathcal N_{123}^{<r}(T)\|
 &\le
 Cs_\alpha^2\Xi_1\Xi_2\Xi_3
 \prod_{ij\in E(T)}\theta_{ij}^{<r},
\label{eq:V63-ternary-separated}\\
 \|\mathcal N_{14,r}\|
 &\le
 Cs_\alpha\Xi_1\Xi_4\theta_{14,r}.
\label{eq:V63-binary-separated}
\end{align}

\begin{proposition}[Exact duplicated attachment mass]
\label{prop:V63-duplicated-mass}
Multiplying
\eqref{eq:V63-ternary-separated} and
\eqref{eq:V63-binary-separated} and dividing by the four input masses
gives
\begin{equation}
 \frac{
 \|\mathcal N_{123}^{<r}(T)\|
 \|\mathcal N_{14,r}\|}
 {\Xi_1\Xi_2\Xi_3\Xi_4}
 \le
 Cs_\alpha^3\,
 \boxed{\Xi_1}\,
 \prod_{ij\in E(T)}\theta_{ij}^{<r}\,
 \theta_{14,r}.
\label{eq:V63-extra-Xi}
\end{equation}
The desired quartic tree normalization has the same product without the
boxed factor \(\Xi_1\).  Paying either nontrivial block removes one
power of \(s_\alpha\) but leaves the factor \(\Xi_1\).  Paying a complete
spectator contributes the ordinary \(O(s_\alpha^{-1})\) first moment
and again leaves \(\Xi_1\).
\end{proposition}

\begin{proof}
The two separate block bounds each contain the full mass of their common
row \(1\).  Their product therefore contains \(\Xi_1^2\), whereas the
four-row normalization cancels only one copy.  This proves
\eqref{eq:V63-extra-Xi}.  The paid bounds in
\eqref{eq:V37-payer-table} have the same products of input row masses as
their unpayed counterparts, and differ only by one power of
\(s_\alpha\).  The singleton payer bound
\eqref{eq:V37-singleton-table} likewise changes the small-noise power,
not the duplicated row mass.
\end{proof}

\begin{corollary}[Ordered summation cannot repair the mass deficit]
\label{cor:V63-simplex-no-repair}
Applying \Cref{cor:V62-first-stock} to
\eqref{eq:V63-extra-Xi} removes the joining-coordinate count but retains
the factor \(\Xi_1\).  Thus support-uniform scalar summation and
quartic row-mass normalization are logically different gates.
\end{corollary}

\begin{proof}
The Cartesian domination lemma multiplies the total edge stocks and
does not alter any prefactor independent of the coordinate indices.
\end{proof}

\begin{proposition}[Sharp abstract regression for separate block tables]
\label{prop:V63-scalar-regression}
No argument using only the two separate V37 norm/capacity bounds and the
V2 one-sided tensor inequality can remove the factor \(\Xi_1\) in
\eqref{eq:V63-extra-Xi}.  More precisely, there are one-dimensional
positive carriers satisfying those abstract bounds for which the ratio
between the separated product and the desired four-row tree mass equals
\(\Xi_1\).
\end{proposition}

\begin{proof}
Set \(s_\alpha=1\), take every Hilbert space one-dimensional, put
\(\Xi_1=L\) and \(\Xi_2=\Xi_3=\Xi_4=1\), and set the three displayed
normalized edge factors to one.  Choose the ternary and binary carrier
scalars equal to the right sides of
\eqref{eq:V63-ternary-separated} and
\eqref{eq:V63-binary-separated}, with the harmless fixed constant
removed.  Operator norm and product-state capacity then both equal the
scalar itself, and the one-sided tensor inequality is equality.  After
division by \(\prod_i\Xi_i\), the separated product is \(L\), whereas
the normalized target tree monomial is one.  Letting \(L\to\infty\)
proves the claim.
\end{proof}

\begin{firewall}[Scope of the regression]
\label{fire:V63-regression-scope}
\Cref{prop:V63-scalar-regression} is not a Wilson-law counterexample and
does not disprove quartic MTA.  It proves that the separate lower-block
tables, positivity, tensor factorization, and one-sided capacity law do
not contain the missing attachment-row normalization.  A valid proof
must use the assembled block-product cumulant, a stronger
Wilson-specific rooted estimate, or cancellation with the other local
partition letters.
\end{firewall}

\section{Assembly of the V62 orbit rows}
\label{sec:V63-orbit-assembly}

For the two-block family \(123|4\),
\Cref{cor:V63-two-block-difference} reads
\begin{equation}
 \operatorname{Cov}_r(Z_{123,r},X_{4,r})
 =
 \sum_{\substack{\pi\in\Pi_4\\
                  \pi\vee(123|4)=\widehat1}}L_{r,\pi}.
\label{eq:V63-assemble-9-12}
\end{equation}
The right side is exactly the sum of the row-labelled instances of V62
orbits 9--12 at this prefix.  Similarly:
\begin{align}
 \rho=12|34&:
 &J_r(\rho)&=\operatorname{Cov}_r(Z_{12,r},Z_{34,r}),
\label{eq:V63-assemble-5-8}\\
 \rho=12|3|4&:
 &J_r(\rho)&=\kappa_3(Z_{12,r},X_{3,r},X_{4,r}),
\label{eq:V63-assemble-2-4}\\
 \rho=\widehat0&:
 &J_r(\rho)&=\kappa_4(X_{1,r},X_{2,r},X_{3,r},X_{4,r}).
\label{eq:V63-assemble-1}
\end{align}

\begin{theorem}[Cancellation-preserving four-family replacement]
\label{thm:V63-four-family-replacement}
The exact V61 connected carrier may be classified by the four prefix
types in \Cref{cor:V63-four-families}, with joining carriers
\eqref{eq:V63-assemble-9-12}--\eqref{eq:V63-assemble-1}.  This
classification counts every connected coordinate word once and retains
all cancellations internal to the local composite cumulant.  Expanding
those four joining carriers into the twelve V62 orbits is optional
algebraically but is not authorized before a positive envelope.
\end{theorem}

\begin{proof}
Exactness and unique word counting follow from
\Cref{thm:V63-joining-cumulant} and
\Cref{thm:V61-first-connectivity}.  The four prefix types are disjoint
and exhaustive.  Each displayed joining carrier is one connected
cumulant, so its moment--cumulant cancellations occur before any norm.
\Cref{prop:V63-scalar-regression} shows that the available analytic
inequalities do not justify the finer expansion.
\end{proof}

\begin{assessment}[Relation to the valid V37--V39 results]
\label{ass:V63-V37-V39}
V59 explicitly preserved the assembled V37 weak-cut, V38 three-cluster,
and V39 universal estimates.  Their leading carriers are respectively a
composite covariance, a composite ternary cumulant, and the complete
quartic carrier---the same cluster orders appearing in
\eqref{eq:V63-assemble-9-12}--\eqref{eq:V63-assemble-1}.  Those theorems
bound the complete assembled carriers, but their recorded statements do
not provide a density conditioned simultaneously on the exact prefix
join and its first joining coordinate.  They therefore identify the
correct analytic machinery without, by themselves, proving the
first-connectivity source bounds.
\end{assessment}

\section{The density-refined quotient-cluster gate}
\label{sec:V63-density-gate}

\begin{definition}[Density-refined quotient-cluster estimate]
\label{def:V63-DRQC}
For a proper partition \(\rho\), the DRQC estimate is the following
statement.  Form the exact source
\[
 \sum_r
 \left(\bigotimes_{B\in\rho}K_B^{<r,\mathrm{conn}}\right)
 \otimes\kappa_r^\rho\otimes M_{>r}
\]
and allocate the sole final payer before taking a positive envelope.
The resulting carrier is bounded by a finite sum of terms consisting of:
\begin{enumerate}[label=\textup{(\alph*)}]
\item one valid capacity factor and ordinary norms on every other
      tensor factor;
\item one internal tree on every nonsingleton prefix block;
\item one tree on the quotient blocks, with every quotient edge expanded
      into an actual crossing row influence;
\item nonnegative coordinate densities compatible with the actual
      first-time and payer ordering; and
\item no repeated use of a row mass, coordinate mark, or resolvent.
\end{enumerate}
\end{definition}

\begin{theorem}[Four DRQC estimates imply quartic MTA]
\label{thm:V63-DRQC-criterion}
If the DRQC estimate holds for one representative of each of the four
prefix types in \Cref{cor:V63-four-families}, uniformly under row
permutation, then the all-weak normalized quartic carrier satisfies
\eqref{eq:V52-all-weak-closed}.  Together with the unaffected V37--V39
branches and the V36 threshold compiler, this restores
\eqref{eq:V52-global-quartic-MTA}.
\end{theorem}

\begin{proof}
Use the exact four-family formula
\eqref{eq:V63-four-family-formula}.  In each DRQC term, lift the prefix
forests and quotient tree to one row tree by
\Cref{thm:V63-forest-quotient}.  Sum the ordered coordinate densities by
\Cref{lem:V62-cartesian}; the actual ordering constraints only shrink
the Cartesian domain.  There are four prefix types and finitely many
row permutations and tree choices, so their constants remain uniform.
This proves the all-weak tree sum.  V59 did not withdraw the weak-cut,
three-cluster, or universal hypotheses, so the stated V36 conclusion
follows.
\end{proof}

\begin{assessment}[Smallest current gate]
\label{ass:V63-smallest-gate}
The first DRQC target should be the two-block partition
\(\rho=123|4\).  Its joining carrier is the exact local composite
covariance in \eqref{eq:V63-assemble-9-12}.  A density-refined version of
the operator martingale covariance proof used in
\Cref{lem:V37-composite-response} must be formed while the ternary prefix
remains intact.  The required gain is exactly the missing
\(\Xi_1^{-1}\) exposed by \Cref{prop:V63-duplicated-mass}; it must arise
from the common coefficient pinching/row influence allocation, not from
a second capacity or a discarded suffix.
\end{assessment}

\section{V63 ledger and V64 mandate}
\label{sec:V63-ledger}

\begin{theorem}[Integrated compact-series V63 statement]
\label{thm:V63-integrated}
Preserving compact V1--V62, V63 proves:
\begin{enumerate}[label=\textup{(V63.\arabic*)},leftmargin=6.7em]
\item the exact identity between the V61 joining letter and the
      Leonov--Shiryaev cumulant of prefix block-products;
\item reduction of twelve analytic orbit candidates to four assembled
      quotient-cluster families;
\item the exact internal-forest plus quotient-tree lift with three and
      only three quartic edges;
\item the unbounded duplicated attachment-row mass in the separated
      orbit-9 estimate, for every payer location;
\item a sharp abstract regression proving that the V2/V37 separate-block
      inequalities cannot repair that loss; and
\item a four-family DRQC criterion sufficient to restore quartic MTA.
\end{enumerate}
\end{theorem}

\begin{proof}
These are \Cref{thm:V63-joining-cumulant,
cor:V63-four-families,
thm:V63-forest-quotient,
prop:V63-duplicated-mass,
cor:V63-simplex-no-repair,
prop:V63-scalar-regression,
thm:V63-four-family-replacement,
thm:V63-DRQC-criterion}.
\end{proof}

\begin{openproblem}[V64 mandate: density-refined composite covariance]
\label{op:V64-mandate}
\begin{enumerate}
\item Return to the operator martingale covariance formula underlying
      \Cref{lem:V37-composite-response}, before physical compression.
\item Insert the exact prefix factor
      \(K_{123}^{<r,\mathrm{conn}}\) and the local composite covariance
      \(\operatorname{Cov}_r(Z_{123,r},X_{4,r})\) without expanding
      either into coordinate words or row partitions.
\item Resolve a common row-1 coefficient carrier once, and test whether
      its martingale influence supplies the missing normalized
      attachment mark \(\Xi_1^{-1}\).
\item Treat prefix-, transition-, and suffix-payer branches with one
      blockwise V2 allocation and use
      \eqref{eq:V62-suffix-payer-stock} only after the common carrier has
      been formed.
\item If the missing inverse mass is false, construct the smallest
      \(SU(2)\) representation/support family realizing the failure and
      enlarge the assembled packet before taking any norm.
\item The next mandatory SM/NS audit is compact V65, after V64.
\end{enumerate}
\end{openproblem}

\begin{firewall}[Claim boundary after compact V63]
\label{fire:V63-current-boundary}
V63 proves that the twelve individual transition orbits are too fine for
the available analytic bounds and replaces them exactly by four
assembled quotient-cluster families.  It does not prove the
density-refined composite covariance, any of the four DRQC estimates,
or the restored all-weak quartic theorem.  Global quartic and quintic
MTA, all-order MTA, WUFC, CPM, continuum Yang--Mills existence, and a
positive mass gap remain open.  The full Yang--Mills mass gap remains
open.
\end{firewall}

\section{Append-only record for compact V63}
\label{sec:V63-record}

V63 was appended to the complete compact V62 file.  The exact V62 parent
had (25145) newline-delimited source records, (859900) bytes, and
SHA--256 digest
\[
\texttt{310bd016b12551033751fb2056f73ae5f88ec0eac2da67908668157a42e76827}.
\]
No compact V62 mathematical content was changed.  The sole final
\verb|\end{document}| is restored below.

% =============================================================================
% END APPENDED COMPACT VERSION V63 MATERIAL
% =============================================================================

% =============================================================================
% BEGIN APPENDED COMPACT VERSION V64 MATERIAL
% =============================================================================

\chapter{Thermal duplicated-row accounting and the five compulsory
assembled branches}
\label{chap:V64-thermal-ledger}

\section{Purpose}
\label{sec:V64-purpose}

V63 found one duplicated attachment mass in the separated orbit-9
estimate.  The Wilson clock can pay such a mass when the row is thermal:
\(s_\alpha\Xi_i\) is then bounded.  This chapter performs that accounting
for all twelve V62 orbits and for every possible location of the sole
final payer.

The outcome is exact.  Seven orbit types close by separated block
estimates in the fully thermal sector for every payer location.  Four
more close except when a singleton block pays.  Orbit 6 remains one
power short even when no singleton can pay.  Thus only five thermal
branches compel assembled quotient-cluster cancellation.  This is a
strictly smaller target than either the twelve-orbit list or the full
four-family DRQC criterion.

\section{Block incidence and duplicated mass}
\label{sec:V64-incidence}

Fix one representative \((\rho,\pi)\) from
\Cref{thm:V62-twelve-orbits}.  Its nonsingleton block occurrences are
the nonsingleton blocks of both \(\rho\) and \(\pi\); the two copies are
kept distinct even if their row sets agree.  Define
\begin{align}
 m_i(\rho,\pi)
 &:=
 \#\{B:\ B\text{ is a nonsingleton block occurrence and }i\in B\},
\label{eq:V64-row-multiplicity}\\
 E(\rho,\pi)
 &:=
 \sum_{B}(|B|-1),
\label{eq:V64-edge-order}\\
 d(\rho,\pi)
 &:=
 \sum_{i=1}^4\bigl(m_i(\rho,\pi)-1\bigr).
\label{eq:V64-duplication-degree}
\end{align}
Admissibility implies \(m_i\ge1\): a row singleton in both partitions
would be isolated in their join.  Thus \(d\) is the total exponent of
duplicated input masses after the product of separate block tables is
divided by \(\prod_i\Xi_i\).

\begin{lemma}[Separated once-paid mass factor]
\label{lem:V64-separated-factor}
Choose one tree for every nonsingleton block occurrence.  Apply the
valid V37 response-or-numerator tables to block sizes two and three and
the valid one-selected V34 local table to a transition block of size
four.  Before surplus marks are discarded, the normalized separated
majorant has the following scalar prefactor.
\begin{enumerate}[label=\textup{(\alph*)}]
\item If a nonsingleton block is the payer, or if the payer is an
      ordinary complete-moment spectator coordinate, then
      \begin{equation}
       C s_\alpha^{E-1}
       \prod_{i=1}^4\Xi_i^{m_i-1}.
      \label{eq:V64-nonsingleton-factor}
      \end{equation}
\item If a singleton block factor on row \(j\) is the payer, then
      \begin{equation}
       C s_\alpha^{E-1}\Xi_j
       \prod_{i=1}^4\Xi_i^{m_i-1}.
      \label{eq:V64-singleton-factor}
      \end{equation}
\end{enumerate}
Each prefactor multiplies the product of the selected normalized edge
marks.  It is independent of support cardinality.
\end{lemma}

\begin{proof}
An unpayed connected block \(B\) contributes
\[
 C s_\alpha^{|B|-1}
 \left(\prod_{i\in B}\Xi_i\right)
 \prod_{uv\in E(T_B)}\theta_{uv}
\]
by \eqref{eq:V37-unpayer-table}; the selected local four-row statement
has the same three-edge numerator scaling.  Multiplying all
nonsingleton occurrences gives the power \(s_\alpha^E\) and row mass
\(\prod_i\Xi_i^{m_i}\).  Division by the four-row input normalization
leaves \(\prod_i\Xi_i^{m_i-1}\).

If a nonsingleton block pays, its response table lowers the
\(s_\alpha\)-power by one and does not change its input-mass product.
If a complete-moment spectator coordinate pays, the ordinary
one-link moment bound contributes \(Cs_\alpha^{-1}\) and no additional
input row mass.  These are
\eqref{eq:V64-nonsingleton-factor}.

An unpayed singleton has norm at most one.  If it pays, the last line of
\Cref{lem:V37-lower-table} replaces that contraction by
\(C\Xi_j/s_\alpha\), proving
\eqref{eq:V64-singleton-factor}.  In every case V2 allocates only one
resolvent.  Coordinate-refined marks and the payer density are summed by
\Cref{lem:V62-cartesian}, so no support count is present.
\end{proof}

\begin{remark}[Why a paid singleton is different]
\label{rem:V64-singleton-extra}
A singleton which is unpayed supplies no input-mass factor.  When it
pays it supplies \(\Xi_j/s_\alpha\).  Since admissibility forces row
\(j\) to belong to a nonsingleton block of the other partition, this is
one additional duplicated row mass.  It cannot be included silently in
the spectator-payer case.
\end{remark}

\section{Thermal exponents}
\label{sec:V64-thermal-exponents}

Fix \(M<\infty\) and call the input fully \(M\)-thermal when
\begin{equation}
 s_\alpha\Xi_i\le M,\qquad 1\le i\le4.
\label{eq:V64-thermal}
\end{equation}
Define the remaining small-noise exponents
\begin{equation}
 g_{\rm ns}:=E-1-d,
 \qquad
 g_{\rm sing}:=E-2-d.
\label{eq:V64-gaps}
\end{equation}
The second is relevant only when \(\rho\) or \(\pi\) has a singleton
block which may receive the payer.

\begin{lemma}[Thermal absorption criterion]
\label{lem:V64-thermal-absorption}
Under \eqref{eq:V64-thermal}, the prefactor
\eqref{eq:V64-nonsingleton-factor} is bounded by
\[
 C M^d s_\alpha^{g_{\rm ns}}
\]
and the prefactor \eqref{eq:V64-singleton-factor} is bounded by
\[
 C M^{d+1}s_\alpha^{g_{\rm sing}}.
\]
Consequently the corresponding separated branch is uniform for
\(0<s_\alpha\le1\) whenever the relevant \(g\) is nonnegative.
\end{lemma}

\begin{proof}
Rewrite every duplicated factor as
\(\Xi_i^{m_i-1}
=s_\alpha^{-(m_i-1)}(s_\alpha\Xi_i)^{m_i-1}\)
and use \eqref{eq:V64-thermal}.  In the singleton-payer case do the same
with the extra \(\Xi_j\).  The total negative powers are \(d\) and
\(d+1\), respectively, which gives \eqref{eq:V64-gaps}.
\end{proof}

\section{The complete twelve-row thermal ledger}
\label{sec:V64-table}

\begin{theorem}[Thermal duplication table]
\label{thm:V64-thermal-table}
For the twelve representatives in
\Cref{thm:V62-twelve-orbits}, the exact separated-table data are:
\begin{center}
\small
\begin{tabular}{c c c c c c c}
\toprule
\#&\(E\)&\((m_1,m_2,m_3,m_4)\)&\(d\)&\(g_{\rm ns}\)&
singleton locations&\(g_{\rm sing}\)\\
\midrule
1&3&(1,1,1,1)&0&2&\(\rho:1,2,3,4\)&1\\
2&3&(2,2,1,1)&2&0&\(\rho:3,4\)&\(-1\)\\
3&3&(2,1,1,1)&1&1&\(\rho:3,4;\ \pi:2\)&0\\
4&4&(2,2,1,1)&2&1&\(\rho:3,4\)&0\\
5&3&(2,1,2,1)&2&0&\(\pi:2,4\)&\(-1\)\\
6&4&(2,2,2,2)&4&\(-1\)&none&--\\
7&4&(2,2,2,1)&3&0&\(\pi:4\)&\(-1\)\\
8&5&(2,2,2,2)&4&0&none&--\\
9&3&(2,1,1,1)&1&1&\(\rho:4;\ \pi:2,3\)&0\\
10&4&(2,2,2,1)&3&0&\(\rho:4\)&\(-1\)\\
11&4&(2,2,1,1)&2&1&\(\rho:4;\ \pi:3\)&0\\
12&5&(2,2,2,1)&3&1&\(\rho:4\)&0\\
\bottomrule
\end{tabular}
\end{center}
The five negative entries are exactly:
\begin{equation}
 \boxed{
 \begin{gathered}
 \text{singleton-payer branches of orbits }2,5,7,10,\\
 \text{and every separated branch of orbit }6.
 \end{gathered}}
\label{eq:V64-five-bad-branches}
\end{equation}
\end{theorem}

\begin{proof}
Read the nonsingleton block occurrences from the two source columns of
\Cref{thm:V62-twelve-orbits}.  For example, orbit 10 has prefix block
\(123\) and transition blocks \(12,34\).  Hence
\[
 (m_1,m_2,m_3,m_4)=(2,2,2,1),\quad
 E=2+1+1=4,\quad d=3,
\]
so \(g_{\rm ns}=0\) and \(g_{\rm sing}=-1\); the only singleton is row
\(4\) in the prefix.  The other rows are identical finite counts.
Substitution in \eqref{eq:V64-gaps} gives the table.
\end{proof}

\begin{corollary}[Seven fully thermal orbit types close separately]
\label{cor:V64-seven-close}
Under the exact factorization and coordinate-density hypotheses of the
valid V34/V37 block tables, the fully \(M\)-thermal once-paid packets of
orbits
\begin{equation}
 \boxed{1,3,4,8,9,11,12}
\label{eq:V64-seven-orbits}
\end{equation}
obey the quartic marked-tree bound for every payer location.  For orbits
\(2,5,7,10\), all nonsingleton-block and ordinary spectator-payer
branches obey it; only the singleton-payer branches remain.  No
separated branch of orbit 6 is certified.
\end{corollary}

\begin{proof}
Apply \Cref{lem:V64-thermal-absorption} to the nonnegative entries of
\Cref{thm:V64-thermal-table}.  Sum joining and payer coordinates through
\Cref{lem:V62-cartesian}.  The chosen block forests have at most two
surplus normalized marks; apply
\Cref{cor:V62-surplus-deletion} to obtain a finite sum of quartic tree
weights.  The local size-four block is used only in its valid
one-selected form, and V61 supplies complete nonpayer suffix
contractions.
\end{proof}

\begin{corollary}[Thermal orbit 9 satisfies DRQC after finite expansion]
\label{cor:V64-orbit9-thermal}
In the fully thermal sector, the orbit-9 packet satisfies the
once-paid density interface of \Cref{def:V62-density-interface} and the
quartic source bound.  The factor \(\Xi_1\) in
\eqref{eq:V63-extra-Xi} is paid by one of the two available powers of
\(s_\alpha\), and a singleton payer on row \(2\), \(3\), or \(4\) uses
the second.
\end{corollary}

\begin{proof}
Orbit 9 has \(d=1\), \(g_{\rm ns}=1\), and
\(g_{\rm sing}=0\).  The assertion follows from
\Cref{lem:V64-thermal-absorption} and the orbit-9 tree identity
\eqref{eq:V62-orbit9-scalar}.
\end{proof}

\section{Sharpness of the five failures}
\label{sec:V64-sharpness}

\begin{proposition}[Thermal scaling regression for orbit 6]
\label{prop:V64-orbit6-regression}
The separate block tables cannot close orbit 6 even under the strongest
uniform thermal scaling \(s_\alpha\Xi_i=1\) for all four rows.  Their
once-paid prefactor is then \(Cs_\alpha^{-1}\).
\end{proposition}

\begin{proof}
Orbit 6 consists of four binary block occurrences:
\[
 (12),(34)\quad\hbox{in the prefix},\qquad
 (13),(24)\quad\hbox{at the transition}.
\]
Thus \(E=4\) and \(m_i=2\) for all \(i\), so the prefactor in
\eqref{eq:V64-nonsingleton-factor} is
\[
 Cs_\alpha^3\Xi_1\Xi_2\Xi_3\Xi_4.
\]
Putting \(\Xi_i=s_\alpha^{-1}\) gives
\(Cs_\alpha^{-1}\).  There is no singleton block and hence no alternate
singleton accounting.
\end{proof}

\begin{proposition}[Thermal scaling regression for a bad singleton payer]
\label{prop:V64-singleton-regression}
For each of orbits \(2,5,7,10\), the separated singleton-payer
prefactor can scale as \(Cs_\alpha^{-1}\) while
\(s_\alpha\Xi_i=1\) on every duplicated row and on the paying singleton
row.
\end{proposition}

\begin{proof}
For these four rows of
\Cref{thm:V64-thermal-table}, \(g_{\rm sing}=-1\).  Set every mass
appearing in \eqref{eq:V64-singleton-factor} equal to
\(s_\alpha^{-1}\).  The exact factorization in
\Cref{lem:V64-thermal-absorption} then leaves
\(s_\alpha^{-1}\).  As in
\Cref{prop:V63-scalar-regression}, one-dimensional positive carriers
may saturate the abstract block inequalities, so no consequence of
those inequalities alone removes the divergence.
\end{proof}

\begin{firewall}[Not a Wilson counterexample]
\label{fire:V64-regression-scope}
The two regressions prove insufficiency of separated response-or-numerator
tables.  They do not assert that Wilson cumulants saturate all four
binary blocks or the paid singleton simultaneously.  The exact
block-product cumulants may cancel these scalings, which is why the
corresponding joining letters must now remain assembled.
\end{firewall}

\section{Where the five failures assemble}
\label{sec:V64-assembly-map}

\begin{theorem}[Assembly map for all thermal failures]
\label{thm:V64-assembly-map}
Every branch in \eqref{eq:V64-five-bad-branches} lies in one of the
following three assembled joining carriers:
\begin{align}
 \rho=12|3|4:\quad&
 \kappa_3(Z_{12,r},X_{3,r},X_{4,r})
 &&\text{contains the bad orbit-2 singleton branches},
\label{eq:V64-family-three}\\
 \rho=12|34:\quad&
 \operatorname{Cov}(Z_{12,r},Z_{34,r})
 &&\text{contains bad orbit 5, orbit 6, and bad orbit 7},
\label{eq:V64-family-22}\\
 \rho=123|4:\quad&
 \operatorname{Cov}(Z_{123,r},X_{4,r})
 &&\text{contains the bad orbit-10 singleton branch}.
\label{eq:V64-family-31}
\end{align}
The discrete-prefix local-four family has no thermal separated deficit.
\end{theorem}

\begin{proof}
Use the orbit-to-prefix table in
\Cref{thm:V62-twelve-orbits}.  Orbits \(2\)--\(4\) have prefix type
\(2+1+1\), orbits \(5\)--\(8\) have prefix type \(2+2\), and orbits
\(9\)--\(12\) have prefix type \(3+1\).  Apply the assembled identities
\eqref{eq:V63-assemble-2-4},
\eqref{eq:V63-assemble-5-8}, and
\eqref{eq:V63-assemble-9-12}.  Orbit 1 is the sole discrete-prefix
orbit and has positive gaps in the thermal table.
\end{proof}

\begin{assessment}[Revised V65 target]
\label{ass:V64-next-target}
The first assembled estimate should be
\eqref{eq:V64-family-22}.  It contains the unique no-singleton thermal
failure, orbit 6, and its exact two-block form is the local moment
difference
\[
 M_r-M_r(12|34).
\]
The composite covariance must be pinched before its four
pair-partition letters are separated.  A successful density estimate
would simultaneously close the orbit-5 and orbit-7 singleton-payer
failures.  If it fails, the four-binary \(SU(2)\) sector is the smallest
place to construct a genuine Wilson regression.
\end{assessment}

\section{V64 ledger and V65 mandate}
\label{sec:V64-ledger}

\begin{theorem}[Integrated compact-series V64 statement]
\label{thm:V64-integrated}
Preserving compact V1--V63, V64 proves:
\begin{enumerate}[label=\textup{(V64.\arabic*)},leftmargin=6.7em]
\item the exact duplicated-row mass and small-noise factor for every
      separated once-paid source;
\item the distinct extra mass created by a singleton-block payer;
\item the complete twelve-orbit thermal exponent table;
\item closure of seven fully thermal orbit types for every payer
      location and partial closure of four further types;
\item sharp separated-table regressions for orbit 6 and the four bad
      singleton-payer branches; and
\item the exact assembly of all five failures into three quotient-cluster
      cumulants, beginning with the \(2+2\) composite covariance.
\end{enumerate}
\end{theorem}

\begin{proof}
These are \Cref{lem:V64-separated-factor,
lem:V64-thermal-absorption,
thm:V64-thermal-table,
cor:V64-seven-close,
cor:V64-orbit9-thermal,
prop:V64-orbit6-regression,
prop:V64-singleton-regression,
thm:V64-assembly-map}.
\end{proof}

\begin{openproblem}[V65 mandate: audit and assembled \(2+2\)
covariance]
\label{op:V65-mandate}
\begin{enumerate}
\item Perform the mandatory read-only audit of the newest active SM and
      NS notes before changing direction.
\item Write the coordinate martingale difference of
      \(\operatorname{Cov}(Z_{12},Z_{34})\) with the exact prefix
      \(K_{12}^{<r,\mathrm{conn}}\otimes
       K_{34}^{<r,\mathrm{conn}}\) retained.
\item Keep
      \(J_r(12|34)=M_r-M_r(12|34)\) assembled through physical
      pinching and the one-resolvent allocation.
\item Prove a quotient crossing influence which supplies one edge
      between the two prefix blocks without the orbit-6 factor
      \(s_\alpha^{-1}\), or construct the smallest \(SU(2)\) Wilson
      regression.
\item Treat the orbit-5 and orbit-7 singleton payers inside the same
      assembled carrier; do not estimate them as separate corrections.
\end{enumerate}
\end{openproblem}

\begin{firewall}[Claim boundary after compact V64]
\label{fire:V64-current-boundary}
V64 closes the separated fully thermal branches listed in
\eqref{eq:V64-seven-orbits} and identifies exactly five branches which
require assembled quotient-cluster control.  It does not close any hot
duplicated-row branch, the assembled \(2+2\) covariance, all four DRQC
families, or the restored all-weak quartic theorem.  Global quartic and
quintic MTA, all-order MTA, WUFC, CPM, continuum Yang--Mills existence,
and a positive mass gap remain open.  The full Yang--Mills mass gap
remains open.
\end{firewall}

\section{Append-only record for compact V64}
\label{sec:V64-record}

V64 was appended to the complete compact V63 file.  The exact V63 parent
had (25654) newline-delimited source records, (879791) bytes, and
SHA--256 digest
\[
\texttt{8132a4d396d99bdd3a1393c8698fd1249f7faa7fa53472417b169b5aeac85824}.
\]
No compact V63 mathematical content was changed.  The sole final
\verb|\end{document}| is restored below.

% =============================================================================
% END APPENDED COMPACT VERSION V64 MATERIAL
% =============================================================================

% =============================================================================
% BEGIN APPENDED COMPACT VERSION V65 MATERIAL
% =============================================================================

\chapter{Mandatory audit and the assembled thermal
\texorpdfstring{\(2+2\)}{2+2} covariance}
\label{chap:V65-assembled-22}

\section{Mandatory read-only companion audit}
\label{sec:V65-audit}

The newest active companion files present during this audit were:
\begin{center}
\small
\begin{tabular}{p{0.47\linewidth}r r p{0.27\linewidth}}
\toprule
file&lines&bytes&SHA--256\\
\midrule
\texttt{Renewal\_SM\_Einstein\_Continuum\_Research\_Notes\_V79.tex}
&\(30454\)&\(1043768\)&
\texttt{4379f13aac813eb0f9bb6f280e693ada9cc351b99278d37918ab633a4ef91f40}\\
\texttt{Renewal\_NS\_Stability\_Research\_Notes\_V31.tex}
&\(23052\)&\(887537\)&
\texttt{f1121b62238ad5491bb7cf03635ea9934942edf573ed8faaa9ee79e1d38a6696}\\
\bottomrule
\end{tabular}
\end{center}
The SM programme advanced from V78 to V79 while the audit was being
performed, so the later file was used.  Both companion files were read
only.

SM V75--V79 first obtains a complete high-tail inverse for order-zero
Dirac perturbations, then treats metric principal-symbol changes by a
form cone.  It removes a massive shift by a reduced inverse on the
complement of a moving kernel, uses Kato transport to put that complement
on a fixed scale, and finally enlarges the infrared sector to one
fixed-rank exact Riesz cluster containing every zero crossing.  The
external complement is exactly reducing and uniformly invertible; every
remaining singularity is confined to the finite Hermitian cluster
matrix.  No Schur remainder is introduced for an exact spectral
projection.

NS V31 is unchanged.  Its dyadic probability mark is multiplicity-free,
but removing the mark restores the exact radial first moment.  Heat
lifting merely transfers that first moment to the critical source
action, and the parent formation bound remains at continuation-class
size.  The persistent branch may not be charged a second time.

\begin{assessment}[V65 audit decision]
\label{ass:V65-audit-decision}
The type-correct common lesson is proof order.  A potentially singular or
mass-deficient collection must be enlarged to the exact reducing packet
before inversion, unmarking, or positive majorization.  In YM this means
keeping
\[
 J_r(12|34)=M_r-M_r(12|34)
\]
as one composite covariance through physical pinching and the sole
resolvent allocation.  Separating its orbit-5, orbit-6, and orbit-7
letters would recreate the V64 deficits.  No SM resolvent estimate or NS
formation estimate is imported.
\end{assessment}

\begin{firewall}[Cross-program type boundary]
\label{fire:V65-audit-types}
An exact Riesz cluster is a spectral reducing subspace; a YM
block-product cumulant is a replica-partition identity; and an NS dyadic
mark is a probability measure on cutoffs.  Only the instruction to
assemble before applying a singular operation transfers.  All YM
inequalities below are proved from Wilson moment and cumulant estimates.
\end{firewall}

\section{A replica-difference formula for the local composite covariance}
\label{sec:V65-double-difference}

Fix a coordinate \(r\) and let \(\mu_r\) be its Wilson probability law.
For two independent variables \(U,V\sim\mu_r\), put
\[
 A(U)=X_{1,r}(U)\otimes X_{2,r}(U),
 \qquad
 B(U)=X_{3,r}(U)\otimes X_{4,r}(U).
\]
The two displayed operators act on disjoint row tensor factors.

\begin{lemma}[Exact double-replica covariance]
\label{lem:V65-double-replica}
\begin{equation}
 \boxed{
 \operatorname{Cov}_{\mu_r}(A,B)
 =
 \frac12
 \mathbb E_{U,V}
 \bigl[(A(U)-A(V))\otimes(B(U)-B(V))\bigr].}
\label{eq:V65-double-replica}
\end{equation}
Moreover,
\begin{align}
 A(U)-A(V)
 ={}&
 (X_{1,r}(U)-X_{1,r}(V))\otimes X_{2,r}(U)
\notag\\
 &+
 X_{1,r}(V)\otimes
 (X_{2,r}(U)-X_{2,r}(V)),
\label{eq:V65-A-telescope}\\
 B(U)-B(V)
 ={}&
 (X_{3,r}(U)-X_{3,r}(V))\otimes X_{4,r}(U)
\notag\\
 &+
 X_{3,r}(V)\otimes
 (X_{4,r}(U)-X_{4,r}(V)).
\label{eq:V65-B-telescope}
\end{align}
Thus the assembled joining covariance is a sum of four terms, each with
one row difference in \(\{1,2\}\) and one in \(\{3,4\}\), while all
other factors remain unitary contractions.
\end{lemma}

\begin{proof}
Expanding the right side of
\eqref{eq:V65-double-replica}, the two diagonal terms are equal to
\(\mathbb E[A(U)\otimes B(U)]\), while independence makes each mixed
term equal to
\(\mathbb EA\otimes\mathbb EB\).  The factor \(1/2\) gives covariance.
Equations \eqref{eq:V65-A-telescope} and
\eqref{eq:V65-B-telescope} are the ordered product-difference identity
\[
 Y_1Y_2-Z_1Z_2=(Y_1-Z_1)Y_2+Z_1(Y_2-Z_2).
\]
\end{proof}

Put
\[
 a_{i,r}:=1+C_2(R_{i,r}),
 \qquad
 h_{ij,r}:=a_{i,r}a_{j,r},
 \qquad
 H_{ij}^{<r}:=\sum_{e<r}h_{ij,e}.
\]
Inactive rows have \(a_{i,r}=0\).  These are the local densities in
\eqref{eq:V36-normalized-overlap}.

\begin{lemma}[One quotient edge for the assembled local covariance]
\label{lem:V65-local-crossing}
Let \(\mathcal N_{r}^{12|34}\) be the complete positive physical
envelope of
\(J_r(12|34)=\operatorname{Cov}_r(A,B)\) before a resolvent, and let
\(\mathcal P_{r}^{12|34}\) be the corresponding envelope when the sole
output mass and resolvent are assigned to this joining carrier.  Then
\begin{align}
 \|\mathcal N_{r}^{12|34}\|
 &\le
 C s_\alpha
 \sum_{\substack{i\in\{1,2\}\\j\in\{3,4\}}}h_{ij,r},
\label{eq:V65-local-unpaid-crossing}\\
 \|\mathcal P_{r}^{12|34}\|
 &\le
 C
 \sum_{\substack{i\in\{1,2\}\\j\in\{3,4\}}}h_{ij,r}.
\label{eq:V65-local-paid-crossing}
\end{align}
The inequalities are formed after the covariance difference
\eqref{eq:V65-double-replica}, not after its eleven partition letters
are separately enveloped.
\end{lemma}

\begin{proof}
Insert \eqref{eq:V65-A-telescope}--%
\eqref{eq:V65-B-telescope} into
\eqref{eq:V65-double-replica}.  Each of the four resulting terms has
one Wilson difference on row \(i\in\{1,2\}\), one on
\(j\in\{3,4\}\), and only unitary factors on the other two rows.  The
one-link two-difference estimate used in the operator martingale proof
of \Cref{lem:V37-composite-response} gives
\(Cs_\alpha a_{i,r}a_{j,r}\) before payment.  Sum the four crossing
pairs only after their common physical output is resolved.  This proves
\eqref{eq:V65-local-unpaid-crossing}.

If the joining carrier pays, fusion assigns the final output mass to one
of its two composite sides and the scalar allocation introduces one,
not two, resolvent denominators.  The binary first-moment/gap estimate in
the same V37 proof removes the single factor \(s_\alpha\).  This gives
\eqref{eq:V65-local-paid-crossing}.  Exact Fourier normalization and raw
isotypic pinching are the ones already used in V37, so no matrix-entry or
multiplicity count is added.
\end{proof}

\begin{firewall}[Why orbit 6 has disappeared]
\label{fire:V65-orbit6-disappears}
The right sides of
\eqref{eq:V65-local-unpaid-crossing}--%
\eqref{eq:V65-local-paid-crossing} contain one crossing edge, not the
two local pair blocks of orbit 6.  The second apparent crossing edge is
an artifact of expanding
\(M_r-M_r(12|34)\) into partition cumulants before using the
double-replica difference.  No individual orbit-6 envelope is bounded
in this proof.
\end{firewall}

\section{The assembled first-connectivity source}
\label{sec:V65-source}

Define
\begin{equation}
 \mathcal S_{22}
 :=
 \sum_{r=1}^n
 K_{12}^{<r,\mathrm{conn}}
 \otimes K_{34}^{<r,\mathrm{conn}}
 \otimes J_r(12|34)
 \otimes M_{>r}.
\label{eq:V65-S22}
\end{equation}
This is the exact \(2+2\) prefix family in
\eqref{eq:V63-four-family-formula}.  It includes all row-labelled
instances of V62 orbits 5--8 and counts every such connected word once.

\begin{lemma}[Four payer locations have the same thermal scale]
\label{lem:V65-four-payers}
After the exact V2 allocation of the sole final payer, the positive
envelope of \(\mathcal S_{22}\) is a sum of four kinds of terms:
\begin{enumerate}[label=\textup{(\roman*)}]
\item the \(12\)-prefix covariance pays;
\item the \(34\)-prefix covariance pays;
\item the local joining covariance pays; or
\item one complete-moment suffix coordinate pays.
\end{enumerate}
For every crossing pair
\(i\in\{1,2\}\), \(j\in\{3,4\}\), each kind is bounded, after summing its
coordinate indices, by
\begin{equation}
 C s_\alpha^2 H_{12}H_{34}H_{ij}.
\label{eq:V65-common-payer-scale}
\end{equation}
There is no singleton-block payer.
\end{lemma}

\begin{proof}
The binary entries of \Cref{lem:V37-lower-table}, restricted to the
prefix coordinates, give
\begin{align*}
 \|\mathcal N_{12}^{<r}\|&\le Cs_\alpha H_{12}^{<r},
 &
 \|\mathcal P_{12}^{<r}\|&\le CH_{12}^{<r},\\
 \|\mathcal N_{34}^{<r}\|&\le Cs_\alpha H_{34}^{<r},
 &
 \|\mathcal P_{34}^{<r}\|&\le CH_{34}^{<r}.
\end{align*}
Combine these with
\eqref{eq:V65-local-unpaid-crossing} or
\eqref{eq:V65-local-paid-crossing}.  If either prefix or the joining
carrier pays, exactly one of the three factors loses its one power of
\(s_\alpha\), leaving \(s_\alpha^2\).

If a suffix coordinate pays, all three source factors are unpayed and
give \(s_\alpha^3\).  The exact complete-product payer sum is
\(O(s_\alpha^{-1})\), uniformly in the number of suffix coordinates, by
the product-moment compiler together with
\Cref{cor:V52-all-failure-debits}.  It again leaves
\(s_\alpha^2\).  Every other suffix factor is a contraction by
\Cref{lem:V61-suffix-contraction}.

Finally,
\[
 \sum_r
 H_{12}^{<r}H_{34}^{<r}h_{ij,r}
 \le H_{12}H_{34}H_{ij}
\]
by \Cref{cor:V62-first-stock}.  A later suffix-payer coordinate merely
adds its genuine ordering restriction and is covered by
\eqref{eq:V62-suffix-payer-stock}.  This proves
\eqref{eq:V65-common-payer-scale}.  Because both prefix blocks and both
joining sides have two rows, no singleton block exists after the
assembled covariance is formed.
\end{proof}

\begin{theorem}[Fully thermal assembled \(2+2\) DRQC]
\label{thm:V65-thermal-22}
Assume the fully \(M\)-thermal condition
\eqref{eq:V64-thermal}.  Then the complete once-paid packet
\(\mathcal S_{22}\) satisfies
\begin{equation}
 \boxed{
 \frac{\kappa_\otimes(\mathcal S_{22}^{\rm pay})}
      {\Xi_1\Xi_2\Xi_3\Xi_4}
 \le
 C_M
 \sum_{\substack{i\in\{1,2\}\\j\in\{3,4\}}}
 \theta_{12}\theta_{34}\theta_{ij}.}
\label{eq:V65-thermal-22}
\end{equation}
Each monomial on the right is a quartic spanning-tree weight.
Consequently the \(2+2\) prefix type satisfies the DRQC estimate in the
fully thermal sector.
\end{theorem}

\begin{proof}
Use \Cref{lem:V65-four-payers}, divide
\eqref{eq:V65-common-payer-scale} by
\(\prod_k\Xi_k\), and observe that
\begin{align*}
 \frac{s_\alpha^2H_{12}H_{34}H_{ij}}
      {\Xi_1\Xi_2\Xi_3\Xi_4}
 &=
 (s_\alpha^2\Xi_i\Xi_j)
 \theta_{12}\theta_{34}\theta_{ij}\\
 &\le
 M^2\theta_{12}\theta_{34}\theta_{ij}.
\end{align*}
The three edges \(12,34,ij\) connect the two internal pairs by one
crossing edge and hence form a spanning tree.  Positive capacity is at
most the ordinary norm on the assembled envelope; all physical output
and payer alternatives were formed before that envelope in
\Cref{lem:V65-local-crossing,lem:V65-four-payers}.
\end{proof}

\begin{corollary}[Three V64 failures close simultaneously]
\label{cor:V65-three-failures-close}
The fully thermal orbit-5 singleton-payer branches, every orbit-6 branch,
and the orbit-7 singleton-payer branches are bounded after their exact
assembly in \(J_r(12|34)\).  No separate estimate of those orbit
letters is required.
\end{corollary}

\begin{proof}
They are exactly the bad \(2+2\)-prefix rows identified in
\eqref{eq:V64-family-22}.  Apply
\Cref{thm:V65-thermal-22}.
\end{proof}

\begin{assessment}[Thermal ledger after V65]
\label{ass:V65-thermal-ledger}
Among the five compulsory thermal failures in
\eqref{eq:V64-five-bad-branches}, only two remain:
\begin{enumerate}[label=\textup{(\alph*)}]
\item the orbit-2 singleton-payer branch inside
      \(\kappa_3(Z_{12},X_3,X_4)\); and
\item the orbit-10 singleton-payer branch inside
      \(\operatorname{Cov}(Z_{123},X_4)\).
\end{enumerate}
The all-hot or mixed hot/thermal duplicated-row sectors remain open and
are not changed by this thermal theorem.
\end{assessment}

\section{V65 ledger and V66 mandate}
\label{sec:V65-ledger}

\begin{theorem}[Integrated compact-series V65 statement]
\label{thm:V65-integrated}
Preserving compact V1--V64, V65 proves:
\begin{enumerate}[label=\textup{(V65.\arabic*)},leftmargin=6.7em]
\item the mandatory SM V79/NS V31 read-only audit and its strict type
      firewall;
\item the exact double-replica difference for the local \(2+2\)
      composite covariance;
\item a one-quotient-edge unpaid and paid local crossing bound formed
      before the eleven partition letters are separated;
\item a four-location one-payer ledger with common scale
      \(s_\alpha^2H_{12}H_{34}H_{ij}\);
\item the fully thermal assembled \(2+2\) DRQC estimate; and
\item simultaneous closure of the orbit-5, orbit-6, and orbit-7 thermal
      deficits isolated in V64.
\end{enumerate}
\end{theorem}

\begin{proof}
These are \Cref{ass:V65-audit-decision,
lem:V65-double-replica,
lem:V65-local-crossing,
lem:V65-four-payers,
thm:V65-thermal-22,
cor:V65-three-failures-close}.
\end{proof}

\begin{openproblem}[V66 mandate: assembled \(3+1\) covariance]
\label{op:V66-mandate}
\begin{enumerate}
\item Apply the same double-replica covariance identity to
      \(A=Z_{123}\), \(B=X_4\), telescoping the three-row product before
      physical pinching.
\item Keep the ternary prefix
      \(K_{123}^{<r,\mathrm{conn}}\) intact and produce exactly one
      quotient edge \(i4\), \(i\in\{1,2,3\}\).
\item Prove the fully thermal once-paid packet including the formerly bad
      orbit-10 singleton-payer branch; do not estimate orbits 9--12
      separately.
\item Record the remaining \(2+1+1\) composite-ternary singleton branch.
\item Only after the thermal quotient-cluster ledger closes, split the
      first hot duplicated row and test whether V52 Wilson moment debits
      or protected invariant capacity pays it.
\end{enumerate}
\end{openproblem}

\begin{firewall}[Claim boundary after compact V65]
\label{fire:V65-current-boundary}
V65 closes the fully thermal assembled \(2+2\) prefix family and three
of the five thermal failures from V64.  It does not close the assembled
\(3+1\) or \(2+1+1\) thermal families, any mixed hot/thermal family, all
four DRQC estimates globally, or the restored all-weak quartic theorem.
Global quartic and quintic MTA, all-order MTA, WUFC, CPM, continuum
Yang--Mills existence, and a positive mass gap remain open.  The full
Yang--Mills mass gap remains open.
\end{firewall}

\section{Append-only record for compact V65}
\label{sec:V65-record}

V65 was appended to the complete compact V64 file.  The exact V64 parent
had (26093) newline-delimited source records, (896171) bytes, and
SHA--256 digest
\[
\texttt{568479b7f8f22755276ad53fc38fb62a13ae4869de39dc4ad63a06018a7baad5}.
\]
No compact V64 mathematical content was changed.  The sole final
\verb|\end{document}| is restored below.

% =============================================================================
% END APPENDED COMPACT VERSION V65 MATERIAL
% =============================================================================

% =============================================================================
% BEGIN APPENDED COMPACT VERSION V66 MATERIAL
% =============================================================================

\chapter{The assembled fully thermal
\texorpdfstring{\(3+1\)}{3+1} covariance}
\label{chap:V66-assembled-31}

\section{Purpose}
\label{sec:V66-purpose}

V65 closed the fully thermal \(2+2\) prefix family by applying an exact
double-replica difference to its composite covariance before expanding
the local partition letters.  The \(3+1\) family has the same quotient
order and a simpler second composite variable.  Its prefix contains one
complete ternary cumulant and one singleton moment; its joining carrier
is
\[
 J_r(123|4)=\operatorname{Cov}_r(Z_{123,r},X_{4,r}).
\]
This chapter proves the corresponding thermal DRQC estimate, including
the prefix-singleton payer which the separated orbit-10 table could not
control.

\section{The three-row product difference}
\label{sec:V66-product-difference}

Let \(U,V\) be independent with the local Wilson law at coordinate \(r\)
and put
\[
 A(U)=X_{1,r}(U)\otimes X_{2,r}(U)\otimes X_{3,r}(U),
 \qquad B(U)=X_{4,r}(U).
\]

\begin{lemma}[Ordered three-factor telescope]
\label{lem:V66-three-telescope}
In the fixed row order,
\begin{align}
 A(U)-A(V)
 ={}&
 (X_1(U)-X_1(V))\otimes X_2(U)\otimes X_3(U)
\notag\\
 &+X_1(V)\otimes(X_2(U)-X_2(V))\otimes X_3(U)
\notag\\
 &+X_1(V)\otimes X_2(V)\otimes(X_3(U)-X_3(V)).
\label{eq:V66-three-telescope}
\end{align}
Consequently
\begin{equation}
 \boxed{
 J_r(123|4)
 =
 \frac12\mathbb E_{U,V}
 [(A(U)-A(V))\otimes(X_4(U)-X_4(V))]}
\label{eq:V66-composite-covariance}
\end{equation}
is a sum of three exact influence terms.  Each contains one row
difference \(i\in\{1,2,3\}\), one row-4 difference, and contractions on
the other two rows.
\end{lemma}

\begin{proof}
The product telescope is obtained by inserting successively
\[
 X_1(V)X_2(U)X_3(U),\qquad
 X_1(V)X_2(V)X_3(U)
\]
between \(A(U)\) and \(A(V)\).  Formula
\eqref{eq:V66-composite-covariance} is
\Cref{lem:V65-double-replica} with the present \(A,B\).  Substitute the
telescope.
\end{proof}

\begin{lemma}[One quotient edge for the \(3+1\) joining covariance]
\label{lem:V66-local-crossing}
Let \(\mathcal N_r^{123|4}\) and
\(\mathcal P_r^{123|4}\) be the complete positive physical envelopes of
\(J_r(123|4)\), respectively before payment and with the sole final
output mass and resolvent assigned to the joining carrier.  Then
\begin{align}
 \|\mathcal N_r^{123|4}\|
 &\le
 Cs_\alpha\sum_{i=1}^3h_{i4,r},
\label{eq:V66-unpaid-crossing}\\
 \|\mathcal P_r^{123|4}\|
 &\le
 C\sum_{i=1}^3h_{i4,r}.
\label{eq:V66-paid-crossing}
\end{align}
\end{lemma}

\begin{proof}
Use \eqref{eq:V66-three-telescope} inside
\eqref{eq:V66-composite-covariance}.  The two-difference Wilson estimate
used in \Cref{lem:V65-local-crossing} bounds the \(i,4\) influence by
\(Cs_\alpha a_{i,r}a_{4,r}\); all other representation factors are
unitary contractions.  Sum the three terms after their common physical
pinching.  If the joining carrier pays, allocate the one final output
mass to a composite side and use the binary first-moment/gap estimate to
remove exactly one \(s_\alpha\).  This gives the two displayed bounds
without expanding V62 orbits 9--12.
\end{proof}

\section{The exact \(3+1\) source and its payer ledger}
\label{sec:V66-source}

Define
\begin{equation}
 \mathcal S_{31}
 :=
 \sum_{r=1}^n
 K_{123}^{<r,\mathrm{conn}}
 \otimes M_{\{4\}}^{<r}
 \otimes J_r(123|4)
 \otimes M_{>r}.
\label{eq:V66-S31}
\end{equation}
For a ternary tree \(T\in\mathfrak T(\{1,2,3\})\), let
\(\mathcal N_{123}^{<r}(T)\) and
\(\mathcal P_{123}^{<r}(T)\) denote its numerator and payer envelopes.
The V37 table gives
\begin{align}
 \|\mathcal N_{123}^{<r}(T)\|
 &\le
 Cs_\alpha^2\Xi_1\Xi_2\Xi_3
 \prod_{uv\in E(T)}\theta_{uv}^{<r},
\label{eq:V66-prefix-unpaid}\\
 \|\mathcal P_{123}^{<r}(T)\|
 &\le
 Cs_\alpha\Xi_1\Xi_2\Xi_3
 \prod_{uv\in E(T)}\theta_{uv}^{<r}.
\label{eq:V66-prefix-paid}
\end{align}

\begin{lemma}[Complete \(3+1\) payer ledger]
\label{lem:V66-payer-ledger}
The exact V2 allocation of the sole final payer in
\eqref{eq:V66-S31} has four kinds of terms:
\begin{enumerate}[label=\textup{(\roman*)}]
\item the ternary prefix pays;
\item the prefix singleton moment on row \(4\) pays;
\item the local joining covariance pays; or
\item one complete-moment suffix coordinate pays.
\end{enumerate}
After summing the joining coordinate, the non-singleton and suffix-payer
branches associated with \(T\) and \(i\in\{1,2,3\}\) are bounded by
\begin{equation}
 C s_\alpha^2\Xi_1\Xi_2\Xi_3
 \left(\prod_{uv\in E(T)}\theta_{uv}\right)H_{i4}.
\label{eq:V66-ordinary-payer-scale}
\end{equation}
The row-4 singleton-payer branch is bounded by the same expression with
one additional factor \(\Xi_4\):
\begin{equation}
 C s_\alpha^2\Xi_1\Xi_2\Xi_3\Xi_4
 \left(\prod_{uv\in E(T)}\theta_{uv}\right)H_{i4}.
\label{eq:V66-singleton-payer-scale}
\end{equation}
\end{lemma}

\begin{proof}
For a ternary-prefix payer, multiply
\eqref{eq:V66-prefix-paid} by
\eqref{eq:V66-unpaid-crossing}; for a joining payer, multiply
\eqref{eq:V66-prefix-unpaid} by
\eqref{eq:V66-paid-crossing}.  Both products have
\(s_\alpha^2\).  An ordinary suffix payer contributes
\(Cs_\alpha^{-1}\) to the three unpayed powers, again leaving
\(s_\alpha^2\), by the product-moment and failure-debit argument in
\Cref{lem:V65-four-payers}.

The unpayed row-4 prefix moment is a contraction.  If it pays,
\eqref{eq:V37-singleton-table} supplies
\(C\Xi_4/s_\alpha\).  Multiplication by the unpayed ternary and joining
carriers leaves \(s_\alpha^2\) and adds \(\Xi_4\), proving
\eqref{eq:V66-singleton-payer-scale}.

Finally apply \Cref{cor:V62-first-stock} to the two prefix edge stocks
and \(h_{i4,r}\).  A suffix payer uses the ordered form
\eqref{eq:V62-suffix-payer-stock}.  This replaces partial quantities by
the totals displayed above without a support count.
\end{proof}

\begin{theorem}[Fully thermal assembled \(3+1\) DRQC]
\label{thm:V66-thermal-31}
Under \eqref{eq:V64-thermal},
\begin{equation}
 \boxed{
 \frac{\kappa_\otimes(\mathcal S_{31}^{\rm pay})}
      {\Xi_1\Xi_2\Xi_3\Xi_4}
 \le
 C_M
 \sum_{T\in\mathfrak T(\{1,2,3\})}
 \sum_{i=1}^3
 \left(\prod_{uv\in E(T)}\theta_{uv}\right)\theta_{i4}.}
\label{eq:V66-thermal-31}
\end{equation}
Every monomial on the right is a quartic spanning-tree weight.  Hence
the full \(3+1\) prefix family satisfies DRQC in the fully thermal
sector.
\end{theorem}

\begin{proof}
For \eqref{eq:V66-ordinary-payer-scale}, division by
\(\prod_k\Xi_k\) gives
\[
 (s_\alpha^2\Xi_i)
 \left(\prod_{uv\in E(T)}\theta_{uv}\right)\theta_{i4}.
\]
This is at most \(M\) times the tree monomial because
\(s_\alpha^2\Xi_i\le Ms_\alpha\le M\).
For \eqref{eq:V66-singleton-payer-scale}, division gives
\[
 (s_\alpha^2\Xi_i\Xi_4)
 \left(\prod_{uv\in E(T)}\theta_{uv}\right)\theta_{i4}
 \le
 M^2
 \left(\prod_{uv\in E(T)}\theta_{uv}\right)\theta_{i4}.
\]
The two edges of \(T\) connect rows \(1,2,3\), and \(i4\) attaches row
\(4\), so their union is a spanning tree.  The exact joining covariance
was pinched and paid before capacity in
\Cref{lem:V66-local-crossing,lem:V66-payer-ledger}.
\end{proof}

\begin{corollary}[The orbit-10 thermal deficit is closed]
\label{cor:V66-orbit10-close}
The formerly deficient prefix-singleton payer in orbit 10 is bounded as
part of the assembled orbits 9--12 joining covariance.  It must not be
reintroduced as a separate positive orbit envelope.
\end{corollary}

\begin{proof}
Orbit 10 is one partition-letter term in
\eqref{eq:V63-assemble-9-12}, and its only V64 deficit was the row-4
prefix singleton payer.  That payer is
\eqref{eq:V66-singleton-payer-scale}, bounded in
\Cref{thm:V66-thermal-31}.
\end{proof}

\begin{assessment}[Last fully thermal quotient-cluster gate]
\label{ass:V66-last-thermal}
After V65--V66, the only thermal deficit from
\eqref{eq:V64-five-bad-branches} is the orbit-2 singleton payer inside
\[
 J_r(12|3|4)
 =
 \kappa_3(Z_{12,r},X_{3,r},X_{4,r}).
\]
Its correct treatment is a three-replica third-difference formula for
the composite variables, producing a two-edge tree on the quotient
blocks.  Together with the prefix edge \(12\), that would close the
fully thermal four-family ledger.
\end{assessment}

\section{V66 ledger and V67 mandate}
\label{sec:V66-ledger}

\begin{theorem}[Integrated compact-series V66 statement]
\label{thm:V66-integrated}
Preserving compact V1--V65, V66 proves:
\begin{enumerate}[label=\textup{(V66.\arabic*)},leftmargin=6.7em]
\item the exact three-factor product difference inside the assembled
      \(3+1\) joining covariance;
\item unpaid and paid one-quotient-edge bounds without expanding orbits
      9--12;
\item the complete four-location payer ledger, including the prefix
      singleton;
\item the fully thermal \(3+1\) DRQC estimate; and
\item closure of the orbit-10 thermal singleton-payer deficit.
\end{enumerate}
\end{theorem}

\begin{proof}
These are \Cref{lem:V66-three-telescope,
lem:V66-local-crossing,
lem:V66-payer-ledger,
thm:V66-thermal-31,
cor:V66-orbit10-close}.
\end{proof}

\begin{openproblem}[V67 mandate: assembled composite ternary joining]
\label{op:V67-mandate}
\begin{enumerate}
\item Write a three-replica polarization identity for
      \(\kappa_3(Z_{12},X_3,X_4)\), or use the operator martingale
      ternary identity underlying \Cref{lem:V38-composite-ternary}.
\item Telescope \(Z_{12}\) before physical pinching and expose a tree of
      two quotient crossing influences among blocks \(12,3,4\).
\item Retain the prefix covariance \(K_{12}^{<r,\mathrm{conn}}\) and
      allocate the sole final payer once, including both prefix singleton
      moments.
\item Prove the fully thermal DRQC estimate and close the orbit-2
      singleton-payer deficit without separately enveloping orbits 2--4.
\item Then begin the mixed hot/thermal ledger by selecting the first hot
      duplicated row only after the assembled joining cumulant is formed.
\end{enumerate}
\end{openproblem}

\begin{firewall}[Claim boundary after compact V66]
\label{fire:V66-current-boundary}
V66 closes the fully thermal assembled \(3+1\) prefix family.  The
assembled \(2+1+1\) thermal family and every mixed hot/thermal DRQC
family remain open.  Thus the global all-weak quartic theorem is not
restored.  Global quartic and quintic MTA, all-order MTA, WUFC, CPM,
continuum Yang--Mills existence, and a positive mass gap remain open.
The full Yang--Mills mass gap remains open.
\end{firewall}

\section{Append-only record for compact V66}
\label{sec:V66-record}

V66 was appended to the complete compact V65 file.  The exact V65 parent
had (26514) newline-delimited source records, (911410) bytes, and
SHA--256 digest
\[
\texttt{a3cb11c5d655fae161e4d282c68d5c8c89cfb563a0d8263f74a3a44bad03829f}.
\]
No compact V65 mathematical content was changed.  The sole final
\verb|\end{document}| is restored below.

% =============================================================================
% END APPENDED COMPACT VERSION V66 MATERIAL
% =============================================================================

% =============================================================================
% BEGIN APPENDED COMPACT VERSION V67 MATERIAL
% =============================================================================

\chapter{The assembled composite-ternary family and completion of the
fully thermal quartic ledger}
\label{chap:V67-composite-ternary}

\section{Purpose and the singleton-payer correction}
\label{sec:V67-purpose}

V64 used the pointwise singleton response bound
\(\|\mathcal P_{\{i\}}\|\le C\Xi_i/s_\alpha\).  That estimate is valid
but unnecessarily coarse when the singleton factor is a complete Wilson
bank.  The bank must be paid before its coordinates are separated.
The exact Bernoulli failure debit removes the factor \(\Xi_i\).

After this correction, the assembled composite ternary joining carrier
\[
 J_r(12|3|4)=\kappa_3(Z_{12,r},X_{3,r},X_{4,r})
\]
has the same two quotient-edge scaling for every payer location.  This
closes the last fully thermal prefix family.  The remaining quartic
problem is thereby restricted to sectors with at least one hot input
row.

\section{A complete Wilson bank pays once}
\label{sec:V67-bank-payer}

Let \(P\) be a finite coordinate bank on one row.  Put
\[
 Y_P:=\sum_{e\in P}a_e,\qquad
 q_P:=\prod_{e\in P}\eta_e,
\]
where \(a_e=1+C_2(R_e)\), \(0\le\eta_e\le1\), and the one-link failure
debit is
\begin{equation}
 s_\alpha a_e\eta_e\le D(1-\eta_e).
\label{eq:V67-one-link-debit}
\end{equation}

\begin{lemma}[Bankwise singleton payer]
\label{lem:V67-bank-payer}
If \(P\ne\varnothing\), then
\begin{equation}
 \boxed{
 \frac{Y_Pq_P}{1-q_P}\le\frac D{s_\alpha}.}
\label{eq:V67-bank-payer}
\end{equation}
The constant is independent of \(|P|\), \(Y_P\), and the
representations.  The same bound holds for every initial or terminal
subbank.
\end{lemma}

\begin{proof}
Multiply \eqref{eq:V67-one-link-debit} by
\(\prod_{f\ne e}\eta_f\) and sum:
\[
 s_\alpha Y_Pq_P
 \le
 D\sum_{e\in P}(1-\eta_e)\prod_{f\ne e}\eta_f.
\]
The sum on the right is the probability of exactly one failure in
independent Bernoulli trials with success probabilities \(\eta_e\).  It
is bounded by the probability of at least one failure,
\(1-\prod_e\eta_e=1-q_P\).  Divide by
\(s_\alpha(1-q_P)\).  Restricting the index set gives the same proof for
any subbank.
\end{proof}

\begin{corollary}[The coarse singleton mass is avoidable]
\label{cor:V67-no-Xi-singleton}
When the sole final output mass and resolvent are allocated to a complete
prefix singleton bank, its entire paid multiplier costs
\(C/s_\alpha\), not \(C\Xi_i/s_\alpha\).  Splitting the bank into
coordinate payers before using \eqref{eq:V67-bank-payer} can lose its
support size and is not authorized.
\end{corollary}

\begin{proof}
The complete bank output has mass at most a fixed multiple of \(Y_P\)
and Wilson multiplier \(q_P\).  Apply
\Cref{lem:V67-bank-payer}.  This is the bankwise version of the product
failure-debit observation in V8 and uses the sole denominator only once.
\end{proof}

\section{The local composite ternary tree}
\label{sec:V67-local-ternary}

For the quotient blocks
\[
 A=\{1,2\},\qquad B=\{3\},\qquad C=\{4\},
\]
define the three local crossing stocks
\begin{align}
 b_{AB,r}&:=h_{13,r}+h_{23,r},&
 b_{AC,r}&:=h_{14,r}+h_{24,r},&
 b_{BC,r}&:=h_{34,r},
\label{eq:V67-cluster-stocks}
\end{align}
and the quotient-tree polynomial
\begin{equation}
 Q_r^{12|3|4}
 :=
 b_{AB,r}b_{AC,r}
 +b_{AB,r}b_{BC,r}
 +b_{AC,r}b_{BC,r}.
\label{eq:V67-local-Q}
\end{equation}

\begin{lemma}[Local composite-ternary numerator and response]
\label{lem:V67-local-ternary}
Let \(\mathcal N_r^{12|3|4}\) be the complete positive physical envelope
of \(J_r(12|3|4)\) before payment, and let
\(\mathcal P_r^{12|3|4}\) be its envelope when it carries the sole final
output mass and resolvent.  Then
\begin{align}
 \|\mathcal N_r^{12|3|4}\|
 &\le Cs_\alpha^2Q_r^{12|3|4},
\label{eq:V67-local-ternary-unpaid}\\
 \|\mathcal P_r^{12|3|4}\|
 &\le Cs_\alpha Q_r^{12|3|4}.
\label{eq:V67-local-ternary-paid}
\end{align}
The three terms of \(Q_r^{12|3|4}\) are the three trees on the quotient
blocks.  Each quotient edge is expanded into the sum of its actual row
crossings before physical pinching.
\end{lemma}

\begin{proof}
Use the operator martingale ternary identity underlying
\Cref{lem:V38-composite-ternary}, restricted to the one-coordinate
family at \(r\).  Replacing the composite variable
\(Z_{12,r}=X_{1,r}X_{2,r}\) telescopes in the fixed order, so its
influence toward row \(3\) is bounded by
\(h_{13,r}+h_{23,r}\), and its influence toward row \(4\) by
\(h_{14,r}+h_{24,r}\).  The third influence is \(h_{34,r}\).
The replicated ternary tree formula is the sum of the three products of
two such influences, exactly \eqref{eq:V67-local-Q}, and contributes
\(s_\alpha^2\).  This proves
\eqref{eq:V67-local-ternary-unpaid}.

Resolve the complete ternary output after the three tree terms have been
formed.  Fusion and the scalar one-resolvent allocation assign the output
mass once, and the ternary response estimate used in V38 removes one
power of \(s_\alpha\).  This proves
\eqref{eq:V67-local-ternary-paid}.  No partition-letter or output
multiplicity is counted.
\end{proof}

\begin{firewall}[The local joining carrier remains assembled]
\label{fire:V67-no-orbit-expansion}
Equation \eqref{eq:V67-local-Q} is an influence-tree expansion of one
composite third cumulant.  It is not the positive envelope of V62 orbits
2, 3, and 4 separately.  In particular, the bad orbit-2 singleton term
is never isolated.
\end{firewall}

\section{The exact \(2+1+1\) source}
\label{sec:V67-source}

Define
\begin{equation}
 \mathcal S_{211}
 :=
 \sum_{r=1}^n
 K_{12}^{<r,\mathrm{conn}}
 \otimes M_{\{3\}}^{<r}
 \otimes M_{\{4\}}^{<r}
 \otimes J_r(12|3|4)
 \otimes M_{>r}.
\label{eq:V67-S211}
\end{equation}
The prefix binary envelopes obey
\begin{equation}
 \|\mathcal N_{12}^{<r}\|\le Cs_\alpha H_{12}^{<r},
 \qquad
 \|\mathcal P_{12}^{<r}\|\le CH_{12}^{<r}.
\label{eq:V67-prefix-binary}
\end{equation}

\begin{lemma}[Every payer branch has the same two-power scale]
\label{lem:V67-payer-ledger}
Allocate the sole final payer in
\eqref{eq:V67-S211} to one of:
\begin{enumerate}[label=\textup{(\roman*)}]
\item the binary prefix;
\item the composite ternary joining carrier;
\item the complete prefix singleton bank on row \(3\) or row \(4\); or
\item one complete-moment suffix bank.
\end{enumerate}
After exact bankwise aggregation, every branch is bounded by
\begin{equation}
 C s_\alpha^2 H_{12}
 \left(
 B_{AB}B_{AC}
 +B_{AB}B_{BC}
 +B_{AC}B_{BC}
 \right),
\label{eq:V67-common-scale}
\end{equation}
where
\begin{align*}
 B_{AB}&:=H_{13}+H_{23},&
 B_{AC}&:=H_{14}+H_{24},&
 B_{BC}&:=H_{34}.
\end{align*}
\end{lemma}

\begin{proof}
If the binary prefix pays, use its paid bound in
\eqref{eq:V67-prefix-binary} and the unpayed joining bound
\eqref{eq:V67-local-ternary-unpaid}.  If the joining carrier pays, use
the unpayed prefix and
\eqref{eq:V67-local-ternary-paid}.  Both cases contain
\(s_\alpha^2\).

If a prefix singleton bank pays, keep its whole Wilson product intact and
apply \Cref{lem:V67-bank-payer}; its cost is \(C/s_\alpha\), with no
extra row mass.  Multiplication by the unpayed binary
\(s_\alpha H_{12}^{<r}\) and unpayed composite ternary
\(s_\alpha^2Q_r^{12|3|4}\) again leaves \(s_\alpha^2\).  The identical
complete-product argument treats a suffix payer.  Every other singleton
or suffix moment is a contraction.

Expand the three terms of \(Q_r^{12|3|4}\).  Apply
\Cref{lem:V62-cartesian} to the prefix edge coordinate, the joining
coordinate shared by the two quotient influences, and the possible
payer-bank coordinate.  Dropping the actual order constraints bounds
their sums by the three products of total stocks in
\eqref{eq:V67-common-scale}.  No raw coordinate count occurs.
\end{proof}

\begin{theorem}[Fully thermal assembled \(2+1+1\) DRQC]
\label{thm:V67-thermal-211}
Under \eqref{eq:V64-thermal},
\begin{equation}
 \boxed{
 \frac{\kappa_\otimes(\mathcal S_{211}^{\rm pay})}
      {\Xi_1\Xi_2\Xi_3\Xi_4}
 \le
 C_M
 \sum_{\substack{T\in\mathfrak T_4\\12\in E(T)}}
 \prod_{uv\in E(T)}\theta_{uv}.}
\label{eq:V67-thermal-211}
\end{equation}
Thus the \(2+1+1\) prefix family satisfies DRQC throughout the fully
thermal sector.
\end{theorem}

\begin{proof}
Expand, for example,
\[
 H_{12}B_{AB}B_{AC}
 =
 \sum_{i\in\{1,2\}}\sum_{j\in\{1,2\}}
 H_{12}H_{i3}H_{j4}.
\]
The edges \(12,i3,j4\) form a spanning tree.  The other two quotient
trees similarly give \(12,i3,34\) and \(12,i4,34\), also spanning
trees.

For any resulting tree \(T\),
\begin{equation}
 \frac{s_\alpha^2\prod_{uv\in E(T)}H_{uv}}
      {\prod_{k=1}^4\Xi_k}
 =
 s_\alpha^2
 \left(\prod_{k=1}^4\Xi_k^{\deg_T(k)-1}\right)
 \prod_{uv\in E(T)}\theta_{uv}.
\label{eq:V67-degree-factor}
\end{equation}
The degree-excess identity for a four-vertex tree is
\[
 \sum_{k=1}^4(\deg_T(k)-1)=2.
\]
Therefore \eqref{eq:V64-thermal} bounds the middle factor in
\eqref{eq:V67-degree-factor} by \(M^2\).  Apply
\Cref{lem:V67-payer-ledger} and sum the finite tree choices.
\end{proof}

\begin{corollary}[The last V64 thermal deficit is closed]
\label{cor:V67-orbit2-close}
The orbit-2 singleton-payer deficit is bounded inside the assembled
composite ternary joining carrier.  Hence every one of the five thermal
failures in \eqref{eq:V64-five-bad-branches} has now been closed by V65,
V66, or V67.
\end{corollary}

\begin{proof}
The orbit-2 branch belongs to
\eqref{eq:V64-family-three}.  Its complete prefix singleton bank is paid
by \Cref{lem:V67-bank-payer}, and the full joining cumulant is bounded by
\Cref{thm:V67-thermal-211}.
\end{proof}

\section{Completion of the fully thermal first-connectivity atlas}
\label{sec:V67-thermal-completion}

\begin{theorem}[All fully thermal quartic first-connectivity sources]
\label{thm:V67-all-thermal}
Assume \(s_\alpha\Xi_i\le M\) for all four input rows.  Then every source
family in the exact first-connectivity formula
\eqref{eq:V63-four-family-formula} satisfies its normalized once-paid
quartic marked-tree bound, uniformly in support, representations,
physical multiplicities, and volume.
\end{theorem}

\begin{proof}
For the discrete prefix, V64 orbit 1 uses the valid one-selected local
four-row carrier and has positive thermal exponent for every payer
location.  The \(2+2\) prefix family is
\Cref{thm:V65-thermal-22}; the \(3+1\) family is
\Cref{thm:V66-thermal-31}; and the \(2+1+1\) family is
\Cref{thm:V67-thermal-211}.  These four prefix types are disjoint and
exhaustive by \Cref{cor:V63-four-families}.  Each proof keeps its
composite joining cumulant assembled, allocates one payer, and uses
complete suffix contractions.
\end{proof}

\begin{corollary}[The remaining quartic source gate is hot]
\label{cor:V67-hot-only}
Any counterexample to the four DRQC source bounds, and hence any
remaining obstruction to restoring the all-weak quartic theorem through
\Cref{thm:V63-DRQC-criterion}, must lie in a sector with
\begin{equation}
 \max_{1\le i\le4}s_\alpha\Xi_i>M.
\label{eq:V67-hot-sector}
\end{equation}
\end{corollary}

\begin{proof}
The complement of \eqref{eq:V67-hot-sector} is exactly the sector closed
by \Cref{thm:V67-all-thermal}.
\end{proof}

\begin{assessment}[What the thermal theorem does not use]
\label{ass:V67-thermal-scope}
The proof does not invoke the withdrawn V52 all-weak theorem, a
graph-difference mark, or an accepted-transfer Green bound.  It uses
only the exact first-connectivity regrouping, lower-order and
one-selected local response tables, bankwise Wilson failure debits, and
assembled composite covariance/ternary influence trees.  Hot input rows
require additional Wilson damping or protected-capacity information and
remain open.
\end{assessment}

\section{V67 ledger and V68 mandate}
\label{sec:V67-ledger}

\begin{theorem}[Integrated compact-series V67 statement]
\label{thm:V67-integrated}
Preserving compact V1--V66, V67 proves:
\begin{enumerate}[label=\textup{(V67.\arabic*)},leftmargin=6.7em]
\item the support-uniform bankwise singleton payer
      \eqref{eq:V67-bank-payer};
\item the local assembled composite-ternary quotient-tree numerator and
      response bounds;
\item a common \(s_\alpha^2\) ledger for every payer location in the
      \(2+1+1\) source;
\item the fully thermal \(2+1+1\) DRQC estimate;
\item closure of the final orbit-2 thermal deficit; and
\item the complete fully thermal quartic first-connectivity atlas,
      reducing the source gate to sectors with a hot input row.
\end{enumerate}
\end{theorem}

\begin{proof}
These are \Cref{lem:V67-bank-payer,
lem:V67-local-ternary,
lem:V67-payer-ledger,
thm:V67-thermal-211,
cor:V67-orbit2-close,
thm:V67-all-thermal,
cor:V67-hot-only}.
\end{proof}

\begin{openproblem}[V68 mandate: the first hot duplicated row]
\label{op:V68-mandate}
\begin{enumerate}
\item In each of the four assembled source families, select a row with
      \(s_\alpha\Xi_i>M\) only after the composite joining cumulant and
      its quotient tree have been formed.
\item Split that row's mass into the coordinates actually used by its
      internal and quotient edges, complete private banks, and unused
      spectators; preserve the exact row probability budget.
\item On a private-dominant hot row, use the all-order Wilson product
      moments of \Cref{cor:V52-private-products-all-orders}.
\item On a shared-dominant hot row, keep the complete protected/gapped
      output packet and test whether one crossing influence earns the
      repeated row mark without invoking a graphwise difference.
\item Treat the \(2+2\) composite covariance first, because its quotient
      edge is explicit and its hot endpoints are the only repeated
      masses in \eqref{eq:V65-thermal-22}.
\item If neither branch pays the hot mass, construct the smallest
      \(SU(2)\) source packet realizing the failure before enlarging the
      packet further.
\end{enumerate}
\end{openproblem}

\begin{firewall}[Claim boundary after compact V67]
\label{fire:V67-current-boundary}
V67 completes the fully thermal quartic first-connectivity source atlas.
It does not control mixed hot/thermal or all-hot sources, restore the
all-weak quartic theorem globally, or restore global quartic MTA.  Global
quintic and all-order MTA, WUFC, CPM, continuum Yang--Mills existence,
and a positive mass gap remain open.  The full Yang--Mills mass gap
remains open.
\end{firewall}

\section{Append-only record for compact V67}
\label{sec:V67-record}

V67 was appended to the complete compact V66 file.  The exact V66 parent
had (26839) newline-delimited source records, (922561) bytes, and
SHA--256 digest
\[
\texttt{f4047a07920b55fc157b5007941a3b5d4a7e55c5e9cdcb7079e8e1de3b2fcc7f}.
\]
No compact V66 mathematical content was changed.  The sole final
\verb|\end{document}| is restored below.

% =============================================================================
% END APPENDED COMPACT VERSION V67 MATERIAL
% =============================================================================

\clearpage

% =============================================================================
% BEGIN APPENDED COMPACT VERSION V68 MATERIAL
% =============================================================================

\chapter{Exact private multipliers for first-connectivity packets}
\label{chap:V68}

\section{Context and scope}
\label{sec:V68-context}

V61--V67 replaced the twelve separately enveloped quartic partition
orbits by four exact first-connectivity source families and closed all
four in the fully thermal sector.  The next issue is whether the private
Wilson-product damping of V41 survives that new regrouping.  This is not
automatic from the complete-cumulant identity alone: an arbitrary
positive subdivision of a cumulant need not inherit a factor that is
visible only after cancellation.

This note proves the stronger, packetwise statement.  A private
coordinate cannot be a first-connection coordinate, and at every other
position it is a scalar spectator independent of the partition word.
Consequently each of the four exact source families, before norms and
before the final output split, has the same exact private multiplier as
the complete fourth cumulant.  The result is then applied to the
assembled \(2+2\) covariance source.  It closes every low-output term
and every private-output term whose two repeated tree rows are thermal
or private-dominant.  What remains is stated exactly: a nonprivate
high-output payer or a genuinely shared-dominant hot internal row.

\section{Private coordinates cannot create first connectivity}
\label{sec:V68-private-spectators}

Keep the ordered coordinate model of V61.  For a coordinate \(e\), let
\[
 A_e:=\{i\in[4]:\hbox{row \(i\) is nontrivial at \(e\)}\}.
\]
Thus \(e\in P_i\) precisely when \(A_e=\{i\}\).  Delete all private
coordinates and preserve the relative order of the remaining
coordinates; a superscript \({\rm sh}\) denotes the resulting
private-free carrier.  The products \(q_i\) and \(q_{\rm priv}\) remain
those of \eqref{eq:V41-private-data}.

\begin{lemma}[A private letter is the discrete partition letter]
\label{lem:V68-private-discrete}
If \(e\in P_i\), its local moment--cumulant expansion has only the
discrete row-partition carrier.  Equivalently, after inactive identity
rows are suppressed, its local operator is the scalar
\(q_{i,e}\), independent of the incoming prefix partition.
In particular, if \(\rho<\widehat1\), then
\begin{equation}
 J_e(\rho)=0 .
\label{eq:V68-private-not-joining}
\end{equation}
\end{lemma}

\begin{proof}
At \(e\) only the random factor in row \(i\) is nonconstant.  Every
joint cumulant involving that factor and a nonempty set of the other
three, constant, row factors is zero.  The moment--cumulant expansion
therefore has the single singleton cumulant
\(\mathbb E F_{i,e}=q_{i,e}\), with the three inactive rows contributing
identities.  Its induced row partition is \(\widehat0\).  Since
\(\rho\vee\widehat0=\rho<\widehat1\), no term in the defining sum
\[
 J_e(\rho)
 =
 \sum_{\pi:\,\rho\vee\pi=\widehat1}L_{e,\pi}
\]
survives.  This proves \eqref{eq:V68-private-not-joining}.
\end{proof}

\begin{theorem}[Exact private multiplier for every
first-connectivity family]
\label{thm:V68-packetwise-private-factor}
Let \(\mathcal S_\rho\) be the sum of all terms in the V61
first-connectivity formula whose prefix partition immediately before
the first-connection coordinate is the fixed
\(\rho<\widehat1\).  Then, in every fixed physical output block,
\begin{equation}
\boxed{
 \mathcal S_\rho
 =
 q_{\rm priv}\,\mathcal S_\rho^{\rm sh}.}
\label{eq:V68-packetwise-private-factor}
\end{equation}
Hence the same identity holds for each of the four assembled families
in \eqref{eq:V63-four-family-formula}: the discrete, \(2+1+1\),
\(2+2\), and \(3+1\) prefix families.  The identity is before a
triangle inequality, before physical pinching, and before the final
output mass is split.
\end{theorem}

\begin{proof}
Use the exact V61 formula
\[
 \mathcal S_\rho
 =
 \sum_r
 \left(\bigotimes_{B\in\rho}
 K_B^{<r,\mathrm{conn}}\right)
 \otimes J_r(\rho)\otimes M_{>r}.
\]
By \Cref{lem:V68-private-discrete}, a nonzero summand has a shared
first-connection coordinate \(r\).

Fix such an \(r\).  If \(e<r\) is private to row \(i\), it occurs in the
unique connected prefix factor whose block \(B\in\rho\) contains \(i\).
The submoment factorization
\eqref{eq:V41-submoment-factorization}, now applied to that block,
extracts precisely \(q_{i,e}\).  If \(e>r\), the suffix is a complete
moment and its private local factor is again precisely \(q_{i,e}\).
There is no third case because \(r\) is shared.  The blocks of \(\rho\)
are disjoint and cover all four rows, so every private coordinate is
extracted once and only once.  Their product is \(q_{\rm priv}\).

After deleting those scalar letters, the remaining prefix, joining,
and suffix words have exactly the same relative order and the same
cumulative joins.  They are therefore the corresponding summand of
\(\mathcal S_\rho^{\rm sh}\).  This correspondence is bijective after
reindexing the shared coordinates.  Summing over \(r\) proves
\eqref{eq:V68-packetwise-private-factor}.  The four-family assertion
follows from the exact regrouping of
\Cref{cor:V63-four-families}.
\end{proof}

\begin{firewall}[No factor is assigned to an orbit letter]
\label{fire:V68-no-private-orbit-factor}
The proof factors a whole first-connectivity family, not an individual
V62 orbit envelope.  In particular,
\eqref{eq:V68-packetwise-private-factor} does not authorize different
private multipliers for different moment partitions inside an
assembled covariance or composite third cumulant.
\end{firewall}

\section{Private-refined assembled \(2+2\) estimates}
\label{sec:V68-private-22}

For \(i\in\{1,2\}\) and \(j\in\{3,4\}\), let
\begin{equation}
 T_{ij}:=\{12,34,ij\}.
\label{eq:V68-Tij}
\end{equation}
This is a path on the four labelled rows.  Its internal vertices are
exactly \(i,j\), and
\begin{equation}
 B_{T_{ij}}
 =\prod_{k=1}^4\Xi_k^{\deg_{T_{ij}}(k)-1}
 =\Xi_i\Xi_j.
\label{eq:V68-path-excess}
\end{equation}

In a fixed output block \(\nu\), retain the exact split
\[
 \Xi(\boldsymbol U_\nu)=P+Z_\nu,
 \qquad
 P=\sum_{k=1}^4p_k,
\]
from \eqref{eq:V41-output-mass-split}.  The terms low output and high
output refer to the same disjoint spectral sectors as in V40--V41.

\begin{proposition}[Three exact private-refined \(2+2\) scales]
\label{prop:V68-private-22-scales}
For every crossing choice \(i\in\{1,2\}\),
\(j\in\{3,4\}\), the assembled source
\(\mathcal S_{22}\) of \eqref{eq:V65-S22} obeys the following
blockwise positive-envelope estimates:
\begin{align}
 \mathcal N_{22}(T_{ij})
 &\le
 C s_\alpha^3q_{\rm priv}H_{12}H_{34}H_{ij},
\label{eq:V68-22-unpaid}\\
 \mathcal L_{22}(T_{ij})
 &\le
 C s_\alpha^2q_{\rm priv}H_{12}H_{34}H_{ij},
\label{eq:V68-22-low}\\
 \mathcal H_{22}^{P}(T_{ij})
 &\le
 C s_\alpha^3Pq_{\rm priv}H_{12}H_{34}H_{ij}.
\label{eq:V68-22-private-high}
\end{align}
Here \(\mathcal N\) is before the final payer,
\(\mathcal L\) is the coefficient-weighted low-output block, and
\(\mathcal H^P\) is the \(P\)-part of the high-output coefficient.
No assertion about the \(Z_\nu\)-part is included.
\end{proposition}

\begin{proof}
Before payment, the two connected binary prefixes contribute
\(Cs_\alpha H_{12}\) and \(Cs_\alpha H_{34}\), while the assembled
joining covariance contributes \(Cs_\alpha H_{ij}\), by
\Cref{lem:V37-lower-table,lem:V65-local-crossing}.  The ordered
coordinate sums are bounded by the Cartesian stock argument of
\Cref{lem:V62-cartesian,cor:V62-first-stock}.  This gives
\(Cs_\alpha^3H_{12}H_{34}H_{ij}\) for the private-free packet.
Apply the exact factor
\eqref{eq:V68-packetwise-private-factor} before taking its envelope to
obtain \eqref{eq:V68-22-unpaid}.

On the low-output sector, the blockwise one-resolvent estimate used in
\eqref{eq:V41-private-refined-low} removes one factor of
\(s_\alpha\) and leaves the already extracted scalar
\(q_{\rm priv}\) untouched.  It is a positive estimate on each
resolved physical block, so it applies to the present assembled source
packet and yields \eqref{eq:V68-22-low}.

On high output the denominator is bounded below uniformly.  Split only
the nonnegative scalar output mass, not the operator, according to
\eqref{eq:V41-output-mass-split}.  Multiplying
\eqref{eq:V68-22-unpaid} by its private part \(P\) proves
\eqref{eq:V68-22-private-high}.  This is the same bookkeeping as
\Cref{rem:V41-output-bookkeeping}; it introduces no second resolvent.
\end{proof}

\begin{theorem}[Hot private-dominant closure for the assembled
\(2+2\) path]
\label{thm:V68-hot-private-22}
Fix \(M\ge1\).  Suppose that, for each of the two internal vertices
\(k\in\{i,j\}\) of \(T_{ij}\), at least one of
\begin{equation}
 s_\alpha\Xi_k\le M,
 \qquad
 p_k\ge\frac12\Xi_k
\label{eq:V68-22-control}
\end{equation}
holds.  Then
\begin{align}
 \frac{\mathcal L_{22}(T_{ij})}
      {\Xi_1\Xi_2\Xi_3\Xi_4}
 &\le
 C_M\theta_{12}\theta_{34}\theta_{ij},
\label{eq:V68-22-low-DRQC}\\
 \frac{\mathcal H_{22}^{P}(T_{ij})}
      {\Xi_1\Xi_2\Xi_3\Xi_4}
 &\le
 C_M\theta_{12}\theta_{34}\theta_{ij}.
\label{eq:V68-22-private-high-DRQC}
\end{align}
In particular, these two sectors are controlled even if one or both
internal rows are hot, provided every hot internal row is
private-dominant.
\end{theorem}

\begin{proof}
For either scalar envelope,
\[
 \frac{H_{12}H_{34}H_{ij}}
      {\Xi_1\Xi_2\Xi_3\Xi_4}
 =
 B_{T_{ij}}\,
 \theta_{12}\theta_{34}\theta_{ij}
 =
 \Xi_i\Xi_j\,
 \theta_{12}\theta_{34}\theta_{ij}.
\]
The hypothesis says exactly that \(T_{ij}\) is
\(M\)-privately controlled in the sense of
\Cref{def:V41-private-controlled}.  Therefore
\Cref{lem:V41-two-excess-powers} bounds
\(s_\alpha^2q_{\rm priv}B_{T_{ij}}\), proving
\eqref{eq:V68-22-low-DRQC}.  Likewise
\Cref{lem:V41-third-private-payer} bounds
\(s_\alpha^3Pq_{\rm priv}B_{T_{ij}}\), proving
\eqref{eq:V68-22-private-high-DRQC}.  These are precisely the
coefficients in \eqref{eq:V68-22-low} and
\eqref{eq:V68-22-private-high}.
\end{proof}

\begin{remark}[The first genuinely hot \(2+2\) rows]
\label{rem:V68-genuine-hot-22}
Because \(T_{ij}\) is a path, the only duplicated row masses are those
of the endpoints \(i,j\) of the assembled quotient edge.  Thus V68
does not merely select an arbitrary hot row: it identifies the two
rows on which a hot estimate can actually fail.
\end{remark}

\section{Exact residual and the next upstream move}
\label{sec:V68-residual}

\begin{theorem}[Exhaustive post-V68 \(2+2\) residual]
\label{thm:V68-exhaustive-residual}
Combine the fully thermal theorem
\Cref{thm:V65-thermal-22} with
\Cref{thm:V68-hot-private-22}.  Every still-unproved normalized
\(2+2\) source term belongs to at least one of the following classes:
\begin{enumerate}
\item it is a high-output term carrying the nonprivate output mass
      \(Z_\nu\);
\item for one of the two internal vertices \(k\in\{i,j\}\),
      \begin{equation}
       s_\alpha\Xi_k>M,
       \qquad
       h_k>\frac12\Xi_k.
      \label{eq:V68-shared-hot-endpoint}
      \end{equation}
\end{enumerate}
The second alternative selects a shared-dominant hot endpoint of the
actual crossing edge of the assembled covariance tree.
\end{theorem}

\begin{proof}
If the source is fully thermal, apply
\Cref{thm:V65-thermal-22}.  Otherwise consider each of its path
monomials \(T_{ij}\).  On low output, or on the \(P\)-part of high
output, \Cref{thm:V68-hot-private-22} applies whenever the tree is
privately controlled.  If it is not, the definition of private control
gives an internal vertex \(k\) for which both alternatives in
\eqref{eq:V68-22-control} fail.  Since \(\Xi_k=p_k+h_k\), this is
exactly \eqref{eq:V68-shared-hot-endpoint}.  The only coefficient not
covered by that dichotomy is the \(Z_\nu\)-part of high output, which
is class~1.
\end{proof}

\begin{firewall}[Why V41 is not being used circularly]
\label{fire:V68-no-circular-V41}
V68 uses from V41 only the exact scalar product moments
\eqref{eq:V41-row-private-moments} and the two algebraic compilers
\Cref{lem:V41-two-excess-powers,lem:V41-third-private-payer}.
The source numerator is supplied independently by the exact
first-connectivity decomposition and the assembled V65 covariance
estimate.  No withdrawn global all-weak theorem and no claimed
graphwise cancellation is used.
\end{firewall}

\begin{assessment}[Upstream target selected by the residual]
\label{ass:V68-upstream-target}
The two residuals in
\Cref{thm:V68-exhaustive-residual} have the same nonprivate origin.
At a shared-dominant internal endpoint \(k\), the repeated factor
\(\Xi_k\) must be tied to the actual incident stocks \(H_{12}\) or
\(H_{34}\) and \(H_{ij}\).  The high-output mass \(Z_\nu\) must
simultaneously be traced to the shared coordinate that creates it.
The next note should therefore keep, for one endpoint \(k\), the
joint packet consisting of its internal-pair influence, its quotient
crossing influence, and its physical output ancestry.  Separating any
of those three objects would recreate the duplicated-row loss.
\end{assessment}

\section{V68 ledger and V69 mandate}
\label{sec:V68-ledger}

\begin{theorem}[Integrated compact-series V68 statement]
\label{thm:V68-integrated}
Preserving compact V1--V67, V68 proves:
\begin{enumerate}[label=\textup{(V68.\arabic*)},leftmargin=6.7em]
\item private coordinates cannot be first-connection coordinates;
\item every exact first-connectivity family has the common
      packetwise multiplier \(q_{\rm priv}\);
\item the assembled \(2+2\) source has private-refined unpaid,
      low-output, and private-high-output scales;
\item all low-output and private-high-output \(2+2\) paths with
      thermal or private-dominant internal vertices satisfy DRQC; and
\item every remaining \(2+2\) obstruction is either a nonprivate
      high-output payer or has a shared-dominant hot endpoint on the
      actual assembled path.
\end{enumerate}
\end{theorem}

\begin{proof}
These are
\Cref{lem:V68-private-discrete,
thm:V68-packetwise-private-factor,
prop:V68-private-22-scales,
thm:V68-hot-private-22,
thm:V68-exhaustive-residual}.
\end{proof}

\begin{openproblem}[V69 mandate: shared ancestry at one hot endpoint]
\label{op:V69-mandate}
\begin{enumerate}
\item Fix \(T_{ij}=\{12,34,ij\}\) and a shared-dominant hot internal
      endpoint \(k\in\{i,j\}\) selected only after the assembled
      covariance tree has been formed.
\item Partition the shared coordinates of row \(k\) by their exact
      companion set, but retain their union as one probability bank.
\item Couple the internal-pair and quotient-edge differences through
      a common replica of row \(k\); do not apply two independent
      Cauchy--Schwarz estimates to the duplicated row.
\item Track the physical output representation of that same shared
      coordinate and test whether its gap pays either the repeated
      \(\Xi_k\) or the nonprivate mass \(Z_\nu\).
\item If this joint packet has no uniform bound, exhibit the smallest
      \(SU(2)\) representation/support configuration for which the
      normalized ratio diverges.  That counterpacket, rather than a
      coarse norm estimate, will determine the required upstream
      lemma.
\end{enumerate}
\end{openproblem}

\begin{firewall}[Claim boundary after compact V68]
\label{fire:V68-current-boundary}
V68 proves an exact private factor for every first-connectivity family
and closes the stated privately controlled parts of the assembled
\(2+2\) source.  It does not control the nonprivate high-output payer,
a shared-dominant hot internal row, the corresponding hot sectors of
the other three source families, or the full all-weak quartic theorem.
Global quartic and higher MTA, WUFC, CPM, continuum Yang--Mills
existence, and a positive mass gap remain open.  The full Yang--Mills
mass gap remains open.
\end{firewall}

\section{Append-only record for compact V68}
\label{sec:V68-record}

V68 was appended to the complete compact V67 file.  The exact V67
parent had \(27265\) newline-delimited source records, \(937562\)
bytes, and SHA--256 digest
\[
\texttt{b0a5d67a67b51e182a9e79cf51b8af8a75da046bfcf2cc5bf98d049b098af21f}.
\]
No compact V67 mathematical content was changed.  The sole final
\verb|\end{document}| is restored below.

% =============================================================================
% END APPENDED COMPACT VERSION V68 MATERIAL
% =============================================================================

\clearpage

% =============================================================================
% BEGIN APPENDED COMPACT VERSION V69 MATERIAL
% =============================================================================

\chapter{Ordered hot endpoints and the unbalanced source tail}
\label{chap:V69}

\section{Context: return to the coordinate-resolved source}
\label{sec:V69-context}

V68 reduced the assembled \(2+2\) source to a nonprivate high-output
payer or a shared-dominant hot endpoint of the path
\(T_{ij}=\{12,34,ij\}\).  Older V42--V50 arguments also studied
shared hot rows, but their final global use passed through separately
normed graph differences and was withdrawn in V59.  The local
balanced/unbalanced spin alternative itself was not the source of that
regression.

This note rebuilds only the part needed by the exact
first-connectivity source.  It first restores the genuine coordinate
ordering hidden by the total-stock bound.  The two edge marks incident
to either internal row occur at distinct coordinates.  It then retains
the elementary unbalanced heat tail on every unmarked coordinate and
on every unbalanced marked coordinate.  Those tails have moments of
every fixed degree, so they pay both repeated endpoint masses and the
unbalanced part of the nonprivate output ancestry.  The surviving
obstruction is balanced: a balanced-dominant hot endpoint or balanced
high-output ancestry.

\section{The two incident marks are necessarily punctured}
\label{sec:V69-ordered-puncture}

For a coordinate \(u\), retain
\[
 a_{k,u}:=
 \begin{cases}
  1+C_2(R_{k,u}),&k\in A_u,\\
  0,&k\notin A_u,
 \end{cases}
 \qquad
 h_{k\ell,u}:=a_{k,u}a_{\ell,u}.
\]
Before the Cartesian enlargement in V65, the scalar stock attached to
the crossing choice \(i\in\{1,2\}\), \(j\in\{3,4\}\) is
\begin{equation}
 \Omega_{ij}
 :=
 \sum_r
 \left(\sum_{e<r}h_{12,e}\right)
 \left(\sum_{f<r}h_{34,f}\right)
 h_{ij,r}.
\label{eq:V69-ordered-stock}
\end{equation}
Thus
\begin{equation}
 \Omega_{ij}
 =
 \sum_{\substack{e,f,r\\e<r,\ f<r}}
 h_{12,e}h_{34,f}h_{ij,r}.
\label{eq:V69-ordered-expanded}
\end{equation}

\begin{lemma}[Distinct incident coordinates at each internal endpoint]
\label{lem:V69-distinct-incident}
In every monomial of \eqref{eq:V69-ordered-expanded}, the two marked
coordinates incident to row \(i\) are \(e\) and \(r\), and the two
marked coordinates incident to row \(j\) are \(f\) and \(r\).  They
satisfy
\begin{equation}
 e\ne r,\qquad f\ne r.
\label{eq:V69-punctures}
\end{equation}
The two prefix coordinates \(e,f\) may coincide.
\end{lemma}

\begin{proof}
The mark \(h_{12,e}\) belongs to the connected \(12\)-prefix and
\(h_{34,f}\) to the connected \(34\)-prefix.  Both prefixes end
strictly before the first-connection coordinate \(r\).  The crossing
mark \(h_{ij,r}\) is produced by the assembled local covariance at
that coordinate.  Hence \(e<r\) and \(f<r\), proving
\eqref{eq:V69-punctures}.  Nothing orders \(e\) relative to \(f\), so
their equality is permitted, as it must be for an incidence-four
prefix coordinate with local partition \(12|34\).
\end{proof}

For row \(i\), define the punctured mass
\[
 R_i(e,r):=\sum_{g\notin\{e,r\}}a_{i,g}.
\]
Define \(R_j(f,r)\) analogously.

\begin{lemma}[Exact nine-cell endpoint-payer decomposition]
\label{lem:V69-nine-cell}
For every monomial indexed by \(e,f,r\) in
\eqref{eq:V69-ordered-expanded},
\begin{align}
 \Xi_i&=a_{i,e}+a_{i,r}+R_i(e,r),
\label{eq:V69-i-payer-split}\\
 \Xi_j&=a_{j,f}+a_{j,r}+R_j(f,r).
\label{eq:V69-j-payer-split}
\end{align}
Consequently \(\Xi_i\Xi_j\) is the sum of exactly nine nonnegative
terms, indexed by
\begin{equation}
 \{e,r,\mathrm{res}\}\times\{f,r,\mathrm{res}\}.
\label{eq:V69-nine-cells}
\end{equation}
This decomposition introduces neither a support count nor a second
copy of either row mass.
\end{lemma}

\begin{proof}
The coordinates active in row \(i\) are the disjoint union of
\(\{e\}\), \(\{r\}\), and their complement; an inactive marked
coordinate simply has \(a_{i,u}=0\).  This proves
\eqref{eq:V69-i-payer-split}; the proof of
\eqref{eq:V69-j-payer-split} is identical.  Multiplying the two
three-term identities gives \eqref{eq:V69-nine-cells}.  It is an
identity for each already selected source monomial, so no pigeonhole
or coordinate summation has been estimated.
\end{proof}

\begin{firewall}[A higher-incidence one-link estimate cannot replace
the puncture]
\label{fire:V69-no-one-link-replacement}
Even if \(r\) has incidence three or four, it cannot be the coordinate
of the other edge incident to \(i\) or \(j\).  Thus a local ternary or
quartic estimate at \(r\) alone cannot pay the complete duplicated
mass of an internal endpoint.  The relevant analytic object is a
two-puncture row packet.
\end{firewall}

\section{The scalar envelope really is insufficient}
\label{sec:V69-scalar-regression}

\begin{proposition}[No stock-only hot \(2+2\) conversion]
\label{prop:V69-no-stock-only}
There is no constant \(C\), uniform in nonnegative row masses and
overlap stocks, for which
\begin{equation}
 \frac{s_\alpha^2H_{12}H_{34}H_{13}}
      {\Xi_1\Xi_2\Xi_3\Xi_4}
 \le
 C\,\theta_{12}\theta_{34}\theta_{13}
\label{eq:V69-false-stock-conversion}
\end{equation}
follows only from
\(0\le H_{k\ell}\le\Xi_k\Xi_\ell\).
The failure occurs on data realized by three ordered shared
coordinates and with row \(1\) shared-dominant and arbitrarily hot.
\end{proposition}

\begin{proof}
Fix \(s_\alpha>0\), put
\(b:=1+C_2(\tfrac12)>1\), choose an \(SU(2)\) spin \(J\) with
\(L:=1+C_2(J)\), and use coordinates \(e,f<r\).  Set
\[
 a_{1,e}=L,\quad a_{2,e}=b,\qquad
 a_{3,f}=a_{4,f}=b,\qquad
 a_{1,r}=a_{3,r}=b,
\]
with all other entries zero.  Then
\[
 \Xi_1=L+b,\quad \Xi_2=b,\quad
 \Xi_3=2b,\quad \Xi_4=b,
\]
and
\[
 H_{12}=Lb,\qquad H_{34}=b^2,\qquad H_{13}=b^2.
\]
These are exactly the three ordered edge stocks in
\eqref{eq:V69-ordered-expanded}.  Division of the two sides of
\eqref{eq:V69-false-stock-conversion} by their common positive tree
monomial shows that the asserted inequality would require
\[
 s_\alpha^2\Xi_1\Xi_3
 =2b\,s_\alpha^2(L+b)\le C.
\]
This fails along \(J\to\infty\).  Row \(1\) has no private mass, so
\(h_1=\Xi_1\), and it is hot once
\(L+b>M/s_\alpha\).  All displayed nonzero entries are genuine
nontrivial \(SU(2)\) representation masses.  Therefore the missing
factor must come from a retained analytic multiplier in the physical
source packet; no rearrangement of \(\Xi\), \(H\), and \(\theta\)
alone can supply it.
\end{proof}

\begin{remark}[This is not a Yang--Mills counterexample]
\label{rem:V69-not-YM-counterexample}
\Cref{prop:V69-no-stock-only} disproves only the coarse positive scalar
conversion.  The actual Wilson packet has heat multipliers and
invariant-projector capacities that were discarded in reaching
\eqref{eq:V65-common-payer-scale}.  The proposition says precisely
that at least one of those analytic factors must be retained.
\end{remark}

\section{Sourcewise unbalanced tails}
\label{sec:V69-unbalanced-tail}

At a shared coordinate \(u\), order the active \(SU(2)\) spins and call
the coordinate unbalanced if one spin \(j_*(u)\) satisfies the
quantitative alternative
\begin{equation}
 \sum_{\ell\ne *}j_\ell(u)\le\frac12j_*(u).
\label{eq:V69-unbalanced}
\end{equation}
Otherwise call it balanced.  Put
\[
 w_u:=\max_{k\in A_u}a_{k,u}.
\]
For a row \(k\), split its shared mass exactly as
\begin{align}
 B_k&:=\sum_{\substack{u\ {\rm balanced}\\k\in A_u}}a_{k,u},
 &
 U_k&:=\sum_{\substack{u\ {\rm unbalanced}\\k\in A_u}}a_{k,u},
\label{eq:V69-row-BU}\\
 h_k&=B_k+U_k.
\label{eq:V69-row-BU-exact}
\end{align}
For a coordinate-resolved source monomial let
\(Q:=\{e,f,r\}\), with repetitions removed, and define
\begin{equation}
 W_Q:=\sum_{\substack{u\ {\rm unbalanced}\\u\notin Q}}w_u.
\label{eq:V69-unmarked-W}
\end{equation}

\begin{lemma}[Uniform unbalanced output separation]
\label{lem:V69-unbalanced-separation}
There are universal \(c,C>0\) such that every local physical summand
of a partition moment at an unbalanced coordinate \(u\) contains a
block with
\begin{equation}
 1+C_2(U_{\rm out})\ge c\,w_u.
\label{eq:V69-output-separation}
\end{equation}
Its heat multiplier is therefore at most
\begin{equation}
 C e^{-c s_\alpha w_u}.
\label{eq:V69-unbalanced-exponential}
\end{equation}
The assertion is uniform in the global row partition.
\end{lemma}

\begin{proof}
Consider the unique partition block containing the dominant input
spin \(j_*\).  Whatever subset of the other active inputs belongs to
that block, every Clebsch--Gordan output spin \(J\) satisfies
\[
 J\ge j_*-\sum_{\ell\ne *}j_\ell\ge\frac12j_*.
\]
Since
\(1+j_*(j_*+1)\) is comparable to \(w_u\), this proves
\eqref{eq:V69-output-separation}.  The established one-link heat
upper bound then gives \eqref{eq:V69-unbalanced-exponential}.  Other
partition blocks are contractions, and the argument did not depend on
which subset accompanied \(j_*\).
\end{proof}

\begin{lemma}[The exact \(2+2\) packet retains unbalanced moments]
\label{lem:V69-source-retains-unbalanced}
Resolve the two binary prefixes and the assembled joining covariance
into their finite global partition states, but do not norm separate
moment-partition letters.  For every fixed integer \(m\ge0\), their
coordinate-resolved \(2+2\) source envelope can be chosen so that
\begin{equation}
 \boxed{
 (s_\alpha W_Q)^m\,
 \mathcal R_Q\le C_m,}
\label{eq:V69-reserve-moments}
\end{equation}
where \(\mathcal R_Q\) is the product of the retained local physical
multipliers at unbalanced coordinates outside \(Q\).

If \(q\in Q\) is itself unbalanced, its marked local source envelope
retains the factor \eqref{eq:V69-unbalanced-exponential}; hence, for
every fixed \(m\),
\begin{equation}
 (s_\alpha w_q)^m
 e^{-c s_\alpha w_q}\le C_m.
\label{eq:V69-marked-moments}
\end{equation}
All constants are independent of the number of coordinates.
\end{lemma}

\begin{proof}
Each binary prefix is a covariance and hence has only the two global
partition states on its two rows.  The joining carrier is one complete
covariance of the two composite variables.  Their tensor product
therefore has a fixed number of global states, independent of support.
At every unmarked unbalanced coordinate,
\Cref{lem:V69-unbalanced-separation} supplies a factor
\(Ce^{-cs_\alpha w_u}\) in every such state.  Tensoring coordinates
first gives
\[
 \mathcal R_Q
 \le C e^{-cs_\alpha W_Q}.
\]
The constant from summing the fixed state set is universal.  The bound
\(\sup_{x\ge0}x^me^{-cx}<\infty\) proves
\eqref{eq:V69-reserve-moments}.

At a marked coordinate, the binary martingale difference or the
assembled composite covariance is a fixed finite signed combination
of the same local physical moments.  Under
\eqref{eq:V69-unbalanced}, the block containing the dominant spin is
nontrivial in every summand, so its heat multiplier is still
\eqref{eq:V69-unbalanced-exponential}.  The displayed edge Casimir
factors have already been placed in \(h_{12,e}\), \(h_{34,f}\), or
\(h_{ij,r}\); retaining the remaining exponential proves
\eqref{eq:V69-marked-moments}.  No graph-difference expansion occurs.
\end{proof}

\begin{firewall}[Balanced cancellation is not estimated here]
\label{fire:V69-balanced-not-estimated}
The preceding proof uses only the nonzero output gap forced by an
unbalanced dominant spin.  At a balanced coordinate, large inputs can
fuse to an invariant output and the exponential may equal one.
No balanced protected projector is separated or normed in
\Cref{lem:V69-source-retains-unbalanced}.
\end{firewall}

\section{A support-free row alternative at the punctures}
\label{sec:V69-row-alternative}

\begin{lemma}[Balanced, reserve, or heavy marked coordinate]
\label{lem:V69-BRH}
Let \(k\in\{i,j\}\) be a shared-dominant hot internal row:
\[
 s_\alpha\Xi_k>M,\qquad h_k>\frac12\Xi_k.
\]
For the marked set \(Q=\{e,f,r\}\), at least one of the following holds:
\begin{enumerate}[label=\textup{(\roman*)}]
\item \(B_k\ge\Xi_k/4\);
\item \(W_Q\ge\Xi_k/8\);
\item some unbalanced \(q\in Q\) satisfies
      \(a_{k,q}\ge\Xi_k/24\).
\end{enumerate}
\label{enum:V69-BRH}
\end{lemma}

\begin{proof}
If (i) fails, then
\eqref{eq:V69-row-BU-exact} and
\(h_k>\Xi_k/2\) give \(U_k>\Xi_k/4\).  The contribution to \(U_k\)
from unbalanced coordinates outside \(Q\) is at most \(W_Q\).  If (ii)
also fails, the contribution from unbalanced coordinates in \(Q\) is
larger than \(\Xi_k/8\).  There are at most three distinct coordinates
in \(Q\), so one of their row-\(k\) contributions is at least
\(\Xi_k/24\), proving (iii).
\end{proof}

\begin{lemma}[Mixed private/unbalanced fixed-degree compiler]
\label{lem:V69-mixed-compiler}
Take any product of at most three factors, each of which is bounded by
a fixed multiple of one of:
\[
 s_\alpha p_k,\qquad
 s_\alpha W_Q,\qquad
 s_\alpha w_q\quad(q\in Q\hbox{ unbalanced}),\qquad
 M.
\]
Multiply it by the exact private scalar \(q_{\rm priv}\), by
\(\mathcal R_Q\), and by the marked exponentials belonging to the
listed \(q\)''s.  The result is bounded by a constant depending only on
the total degree and \(M\).
\end{lemma}

\begin{proof}
Private and shared coordinate sets are disjoint.  On each private row
bank use the fixed moments
\eqref{eq:V52-private-products-all-orders}.  On the unmarked
unbalanced bank use \eqref{eq:V69-reserve-moments}; repeated uses of
the same \(W_Q\) are paid by the corresponding higher moment.  On a
marked unbalanced coordinate use
\eqref{eq:V69-marked-moments}, again with its total exponent if it
pays more than one listed factor.  The remaining factors are at most
\(M\).  A multinomial expansion has a number of terms depending only
on degree three, not on support.
\end{proof}

\section{Closure of every unbalanced hot endpoint}
\label{sec:V69-unbalanced-closure}

Use the nonprivate output ancestry
\eqref{eq:V45-shared-output-ancestry} only as the elementary
Clebsch--Gordan inequality
\[
 Z_\nu\le C\sum_{k=1}^4(B_k+U_k).
\]
Write
\begin{equation}
 Z_\nu^{B}:=C\sum_kB_k,\qquad
 Z_\nu^{U}:=C\sum_kU_k,
\qquad
 Z_\nu\le Z_\nu^B+Z_\nu^U.
\label{eq:V69-output-BU}
\end{equation}
Since each unbalanced coordinate has at most four active rows,
\begin{equation}
 Z_\nu^U\le4C\sum_{u\ {\rm unbalanced}}w_u.
\label{eq:V69-output-U-W}
\end{equation}

\begin{theorem}[Sourcewise unbalanced \(2+2\) closure]
\label{thm:V69-unbalanced-22}
Fix a coordinate-resolved path monomial of the assembled \(2+2\)
source.  Assume that each internal endpoint \(k\in\{i,j\}\) satisfies
at least one of:
\begin{enumerate}[label=\textup{(\alph*)}]
\item \(s_\alpha\Xi_k\le M\);
\item \(p_k\ge\Xi_k/2\);
\item \(k\) is shared-dominant and alternatives
      \textup{(ii)} or \textup{(iii)} of
      \Cref{lem:V69-BRH} hold.
\end{enumerate}
\label{enum:V69-controlled}
Then its low-output part and its private high-output part satisfy the
normalized tree bound
\begin{equation}
 C_M\theta_{12}\theta_{34}\theta_{ij}.
\label{eq:V69-low-P-bound}
\end{equation}
Under the same endpoint assumptions, the \(Z_\nu^U\)-part of its
nonprivate high-output payer also satisfies
\eqref{eq:V69-low-P-bound}.
\end{theorem}

\begin{proof}
After division by the four row masses, the low-output coefficient in
\eqref{eq:V68-22-low} is
\[
 s_\alpha^2q_{\rm priv}\Xi_i\Xi_j
 \theta_{12}\theta_{34}\theta_{ij},
\]
with the unbalanced multipliers of
\Cref{lem:V69-source-retains-unbalanced} retained before the envelope.
A thermal endpoint contributes at most \(M\).  A private-dominant
endpoint is paid by its private Wilson product.  For a shared-dominant
endpoint in case (ii), \(\Xi_k\le8W_Q\); in case (iii),
\(\Xi_k\le24a_{k,q}\le24w_q\).  Thus the two endpoint factors are a
mixed product of total degree two covered by
\Cref{lem:V69-mixed-compiler}.  This proves the low-output assertion.

The private high-output coefficient has one additional factor
\(s_\alpha P\), by \eqref{eq:V68-22-private-high}.  Allocate \(P\)
among the four private row banks and apply the same compiler at total
degree three.  This proves the private-output assertion.

For \(Z_\nu^U\), split the total unbalanced stock in
\eqref{eq:V69-output-U-W} into \(W_Q\) and the at most three marked
weights \(w_q\).  Each summand is paid by its retained reserve or marked
exponential.  Together with the two endpoint factors the total degree
is three, so \Cref{lem:V69-mixed-compiler} applies once more.  The
high-output denominator has already been bounded below and is not
duplicated.
\end{proof}

\begin{corollary}[The exact post-V69 \(2+2\) obstruction is balanced]
\label{cor:V69-balanced-residual}
After V65, V68, and V69, every still-unproved \(2+2\) source term has
at least one of the following properties:
\begin{enumerate}
\item an internal row \(k\in\{i,j\}\) is hot and
      \begin{equation}
       B_k\ge\frac14\Xi_k;
      \label{eq:V69-balanced-hot}
      \end{equation}
\item it is a high-output term carrying the balanced ancestry
      \(Z_\nu^B\) from \eqref{eq:V69-output-BU}.
\end{enumerate}
Thus neither private mass nor unbalanced shared mass remains as an
independent \(2+2\) obstruction.
\end{corollary}

\begin{proof}
Thermal terms were closed in V65.  V68 closes thermal/private-controlled
endpoints and the private output coefficient.  If a remaining endpoint
is shared-dominant, \Cref{lem:V69-BRH} is exhaustive.  V69 closes its
reserve and heavy-mark alternatives, leaving only
\eqref{eq:V69-balanced-hot}.  Finally
\eqref{eq:V69-output-BU} and
\Cref{thm:V69-unbalanced-22} remove the unbalanced part of the
nonprivate output majorant.  The balanced part is not estimated.
\end{proof}

\begin{firewall}[No withdrawn transport theorem is used]
\label{fire:V69-no-withdrawn-transport}
The proof uses the exact first-connectivity source, the assembled V65
covariance, the elementary Clebsch--Gordan lower bound in an unbalanced
coordinate, and the valid V52 fixed one-link moments.  It does not use
the V47 graph carrier, graphwise marks, a separately normed transport
chain, or the withdrawn V52 global quartic conclusion.
\end{firewall}

\section{The balanced gate exposed by V69}
\label{sec:V69-balanced-gate}

At a balanced coordinate, a large input can fuse to the trivial output.
The invariant projector across one row has product-state capacity
\((2j+1)^{-1}\), which is only one square root of the Casimir mass.
For completeness, the source-relevant sharpness is recorded without
using any graph expansion.

\begin{proposition}[One protected coordinate cannot pay one hot row
mass]
\label{prop:V69-one-protected-insufficient}
Let \(P_j\) be the normalized projector onto the invariant line in
\(V_j\otimes V_j^*\).  Its product-state capacity across the two
factors is exactly
\begin{equation}
 \kappa_\otimes(P_j)=\frac1{2j+1}.
\label{eq:V69-singlet-capacity}
\end{equation}
Consequently
\begin{equation}
 \bigl(1+j(j+1)\bigr)\kappa_\otimes(P_j)\longrightarrow\infty
 \qquad(j\to\infty).
\label{eq:V69-one-capacity-fails}
\end{equation}
Thus a balanced hot endpoint cannot be closed by assigning one local
protected capacity to one selected incident coordinate.
\end{proposition}

\begin{proof}
In orthonormal bases the invariant unit vector is maximally entangled:
\[
 \Omega_j=(2j+1)^{-1/2}
 \sum_{m=-j}^{j}e_m\otimes e_m^*.
\]
The largest squared overlap of \(\Omega_j\) with a product unit vector
is the square of its largest Schmidt coefficient,
\((2j+1)^{-1}\), proving \eqref{eq:V69-singlet-capacity}.  The
left side of \eqref{eq:V69-one-capacity-fails} grows like \(j/2\).
\end{proof}

\begin{assessment}[The needed source-specific balanced packet]
\label{ass:V69-needed-balanced-packet}
The ordered source supplies exactly the structure missing from a
one-coordinate bound: the two incident marks of a hot internal row
lie at distinct coordinates.  A successful balanced estimate must
keep the complete signed binary-prefix covariance, the assembled
joining covariance, and all balanced spectator outputs until their
common physical block is fixed.  It must then obtain either two
independent one-row capacities, or one capacity together with a
nontrivial heat output, at the stock level.  V59 forbids obtaining
those factors by norming graph or partition orbits separately.
\end{assessment}

\section{V69 ledger and V70 mandate}
\label{sec:V69-ledger}

\begin{theorem}[Integrated compact-series V69 statement]
\label{thm:V69-integrated}
Preserving compact V1--V68, V69 proves:
\begin{enumerate}[label=\textup{(V69.\arabic*)},leftmargin=6.7em]
\item the two incident marks at either internal endpoint of an
      ordered \(2+2\) source monomial are at distinct coordinates;
\item the two excess row masses have an exact nine-cell
      marked/residual payer decomposition;
\item the coarse stock envelope alone cannot close a hot endpoint;
\item the exact first-connectivity packet retains all fixed
      unbalanced reserve and marked-coordinate heat moments;
\item every unbalanced hot endpoint and the unbalanced part of
      nonprivate output ancestry satisfy the \(2+2\) DRQC bound; and
\item the remaining \(2+2\) obstruction is entirely balanced.
\end{enumerate}
\end{theorem}

\begin{proof}
These are
\Cref{lem:V69-distinct-incident,
lem:V69-nine-cell,
prop:V69-no-stock-only,
lem:V69-source-retains-unbalanced,
thm:V69-unbalanced-22,
cor:V69-balanced-residual}.
\end{proof}

\begin{openproblem}[V70 mandate: audit and balanced two-puncture
capacity]
\label{op:V70-mandate}
\begin{enumerate}
\item Perform the mandatory five-version audit of the newest SM and NS
      notes, importing only proof-order information of the correct
      mathematical type.
\item For the ordered \(2+2\) monomial, fix one balanced-hot internal
      row and retain its two distinct incident coordinates together
      with the complete balanced spectator bank.
\item Form the common physical packet before any positive envelope.
      Prove a two-puncture alternative: two independent row capacities,
      or one row capacity plus a nontrivial output heat factor.
\item Include the balanced output ancestry \(Z_\nu^B\) in the same
      packet so the sole high-output payer is used once.
\item Test the statement first on the minimal two-coordinate \(SU(2)\)
      packet.  If the two capacities collapse to one, record the exact
      recoupling channel causing the collapse and enlarge the packet
      upstream.
\end{enumerate}
\end{openproblem}

\begin{firewall}[Claim boundary after compact V69]
\label{fire:V69-current-boundary}
V69 closes the unbalanced portions of the assembled \(2+2\) hot
source.  It does not prove the balanced two-puncture capacity estimate,
close balanced nonprivate output ancestry, extend the result to all
four source families, or restore the full all-weak quartic theorem.
Global quartic and higher MTA, WUFC, CPM, continuum Yang--Mills
existence, and a positive mass gap remain open.  The full Yang--Mills
mass gap remains open.
\end{firewall}

\section{Append-only record for compact V69}
\label{sec:V69-record}

V69 was appended to the complete compact V68 file.  The exact V68
parent had \(27683\) newline-delimited source records, \(954053\)
bytes, and SHA--256 digest
\[
\texttt{ccc5d377dbf1d461e214df771be416fb484d419350076c2af4690ad220bb004e}.
\]
No compact V68 mathematical content was changed.  The sole final
\verb|\end{document}| is restored below.

% =============================================================================
% END APPENDED COMPACT VERSION V69 MATERIAL
% =============================================================================

\clearpage

% =============================================================================
% BEGIN APPENDED COMPACT VERSION V70 MATERIAL
% =============================================================================

\chapter{Companion audit and the balanced sourcewise transfer
certificate}
\label{chap:V70}

\section{Mandatory five-version companion audit}
\label{sec:V70-audit}

The newest Standard Model and Navier--Stokes continuations in the shared
\texttt{latex\_notes} folder were inspected read-only:
\begin{center}
\small
\begin{tabular}{p{0.48\linewidth}r r p{0.31\linewidth}}
\toprule
file&lines&bytes&SHA--256\\
\midrule
\texttt{Renewal\_SM\_Einstein\_Continuum\_Research\_Notes\_V86.tex}
&33298&1133175&
\texttt{5122e7bba445951621febc604533d3eda4c5771c5a8917daca31d1b619a1981e}
\\
\texttt{Renewal\_NS\_Stability\_Research\_Notes\_V31.tex}
&23052&887537&
\texttt{f1121b62238ad5491bb7cf03635ea9934942edf573ed8faaa9ee79e1d38a6696}
\\
\bottomrule
\end{tabular}
\end{center}
The NS file is unchanged from the V65 Yang--Mills audit.

\begin{assessment}[Standard Model proof-order transfer]
\label{ass:V70-SM-transfer}
SM V86 isolates two simple singular eigenspaces, proves the
complementary resolvents uniformly regular, and compresses the complete
interaction numerator between the two on-shell projections before
taking its norm.  Its warning is directly relevant to the balanced
Yang--Mills gate: protected and gapped local outputs should be split
inside one assembled physical packet, and a possible zero of the
projected numerator must not be replaced by norms of its summands.
No pole, BRST, or microlocal estimate is imported into Yang--Mills.
\end{assessment}

\begin{assessment}[Navier--Stokes proof-order transfer]
\label{ass:V70-NS-transfer}
NS V31 remains at the exact first-moment formation gate.  A flat mark
can be removed only by paying its precise first moment, and a
persistent time interval may not be charged again in the recent-entry
ledger.  For the present source this reinforces two requirements:
the balanced row mass is transferred once through one retained bank,
and the final output payer is not charged a second time during tree
routing.  No NS coarea or Duhamel inequality is imported.
\end{assessment}

\begin{firewall}[Cross-program type boundary]
\label{fire:V70-audit-boundary}
The companion notes contribute proof order only.  A Hermitian
on-shell projection is not a Wilson invariant projector, and a dyadic
formation first moment is not a Casimir row mass.  All estimates below
are proved directly in the Yang--Mills finite-row carrier.
\end{firewall}

\section{Row-rooted balanced banks}
\label{sec:V70-row-banks}

Use the quantitative balanced/unbalanced split of V69.  Fix a row
\(k\) and a balanced coordinate \(u\) active in that row.  Choose
\(v_k(u)\ne k\) whose active spin is largest among the other rows.
For \(\ell\ne k\) and a coordinate set \(D\), put
\begin{align}
 D_{k\ell}
 &:=
 \{u\in D:u\hbox{ is balanced},\ k\in A_u,\
             v_k(u)=\ell\},
\label{eq:V70-partner-bank}\\
 S_{k\ell}(D)
 &:=
 \sum_{u\in D_{k\ell}}a_{k,u},
 &
 H_{k\ell}^{(k)}(D)
 &:=
 \sum_{u\in D_{k\ell}}a_{k,u}a_{\ell,u}.
\label{eq:V70-bank-stocks}
\end{align}

\begin{lemma}[Row-rooted power inequality]
\label{lem:V70-row-power}
There is a universal \(C\) such that, for
\(N_{k\ell}(D):=|D_{k\ell}|\),
\begin{equation}
 \boxed{
 S_{k\ell}(D)^2
 \le
 C N_{k\ell}(D)H_{k\ell}^{(k)}(D).}
\label{eq:V70-row-power}
\end{equation}
Moreover
\begin{equation}
 H_{k\ell}^{(k)}(D)\le H_{k\ell}.
\label{eq:V70-H-domination}
\end{equation}
\end{lemma}

\begin{proof}
If row \(k\) does not carry the largest active spin at \(u\), then the
chosen partner has spin at least \(j_{k,u}\).  If it does carry the
largest spin, failure of the unbalanced condition
\eqref{eq:V69-unbalanced} gives
\[
 \sum_{m\ne k}j_{m,u}>\frac12j_{k,u}.
\]
There are at most three other rows, so the chosen partner has
\(j_{\ell,u}\ge j_{k,u}/6\).  The Casimir formula and the positive
nontrivial spin floor therefore give
\[
 a_{k,u}^2\le C a_{k,u}a_{\ell,u}.
\]
Cauchy--Schwarz over \(D_{k\ell}\) proves
\eqref{eq:V70-row-power}.  The second assertion follows because every
summand on its left is one of the nonnegative summands of the full
overlap \(H_{k\ell}\).
\end{proof}

\begin{lemma}[One partner bank carries a balanced-hot row]
\label{lem:V70-large-partner-bank}
If
\begin{equation}
 B_k\ge\frac14\Xi_k,
\label{eq:V70-balanced-hot}
\end{equation}
then some \(\ell\ne k\) satisfies
\begin{equation}
 S_{k\ell}(E)\ge\frac1{12}\Xi_k,
\label{eq:V70-large-bank}
\end{equation}
where \(E\) is the full coordinate set.
\end{lemma}

\begin{proof}
The three partner banks partition the balanced coordinates active in
row \(k\), and their \(S\)-stocks sum to \(B_k\).  One of them is at
least \(B_k/3\), which together with
\eqref{eq:V70-balanced-hot} proves
\eqref{eq:V70-large-bank}.
\end{proof}

\section{Protected/gapped splitting inside the assembled packet}
\label{sec:V70-spectral-split}

Fix the row cut
\[
 \{k\}\mid[4]\setminus\{k\}.
\]
At a coordinate \(u\in D_{k\ell}\), first resolve the complete local
physical output of the relevant global partition state.  Let
\(\mathsf P_u\) denote the sum of sectors in which every local output
block is trivial, and \(\mathsf G_u\) the sum of the remaining,
nontrivial-output sectors.  This is an orthogonal output
decomposition of the complete local moment:
\begin{equation}
 \mathsf T_u=\mathsf P_u+\mathsf G_u.
\label{eq:V70-PG-split}
\end{equation}
For a marked prefix or joining coordinate the same notation is applied
only after its complete local covariance difference has been formed.

\begin{lemma}[Balanced local protected/gapped alternative]
\label{lem:V70-local-PG}
There are universal \(c>0\) and \(\rho<1\) such that
\begin{align}
 \|\mathsf G_u\|
 &\le1-cs_\alpha,
\label{eq:V70-G-gap}\\
 \kappa_\otimes^{\,k|k^c}(\mathsf P_u)
 &\le\rho.
\label{eq:V70-P-capacity}
\end{align}
The same bounds hold for the residual physical multiplier at a marked
coordinate after its displayed edge Casimir factors have been
extracted.
\end{lemma}

\begin{proof}
Every sector in \(\mathsf G_u\) contains a nontrivial output
representation.  The positive Casimir floor and the standing one-link
heat estimate give \eqref{eq:V70-G-gap}.

In a sector of \(\mathsf P_u\), the nontrivial input in row \(k\)
cannot form a trivial block by itself.  Hence the invariant block
containing it also contains at least one row across the displayed cut.
Regroup that block as \(V_{k,u}\otimes W\).  The normalized invariant
projector has Schmidt coefficients at most
\((\dim V_{k,u})^{-1/2}\), so its product-state capacity across the
cut is at most \((\dim V_{k,u})^{-1}\le1/2\).  Other blocks are
contractions.  Thus \eqref{eq:V70-P-capacity} holds with a universal
\(\rho<1\).

A marked local covariance is a fixed finite signed combination of the
same physical output blocks.  Form that combination first and then
compress it with the orthogonal protected and gapped projections.
The two bounds survive up to a fixed local constant already absorbed
in the source card.  No moment-partition term is separately assigned a
capacity.
\end{proof}

\begin{firewall}[The split is spectral, not combinatorial]
\label{fire:V70-no-PG-word-expansion}
Equation \eqref{eq:V70-PG-split} is not expanded into \(2^N\)
protected/gapped words over a bank of \(N\) coordinates.  The complete
local physical blocks are orthogonal, and the bank estimate below is
taken after their tensor product.  This is the Yang--Mills analogue of
isolating a finite singular sector and its complement; it is not an
imported SM resolvent argument.
\end{firewall}

\section{The balanced bank tail}
\label{sec:V70-bank-tail}

Let \(D_{k\ell}\) be disjoint from the source mark set
\(Q=\{e,f,r\}\).  Denote by \(\mathcal B_{k\ell}(D)\) the complete
tensor product of its retained local physical multipliers in a fixed
global partition state, followed by the one-sided product-state
capacity across \(k|k^c\).

\begin{theorem}[Sourcewise balanced-bank tail]
\label{thm:V70-balanced-tail}
Uniformly in the number of coordinates, representations, and the
fixed global state,
\begin{equation}
 \boxed{
 s_\alpha S_{k\ell}(D)^2
 \kappa_\otimes^{\,k|k^c}
 \bigl(\mathcal B_{k\ell}(D)\bigr)
 \le
 C H_{k\ell}^{(k)}(D).}
\label{eq:V70-balanced-tail}
\end{equation}
The estimate remains valid after the finite global-state sum defining
the two binary prefixes and their assembled joining covariance.
\end{theorem}

\begin{proof}
Put \(N=N_{k\ell}(D)\) and \(r=1-cs_\alpha\).  At one coordinate,
write the complete positive physical block as
\(\mathsf P_u+\mathsf G_u\), with the two output sectors orthogonal.
For a product unit vector \(v\) across \(k|k^c\), set
\(p=\langle v,\mathsf P_uv\rangle\).  Then \(p\le\rho\), while
\(\mathsf G_u\le r(I-\mathsf P_u)\) on the complementary sector.
Consequently
\[
 \langle v,\mathsf T_uv\rangle
 \le p+r(1-p)
 \le r+(1-r)\rho
 \le1-c''s_\alpha.
\]
One-sided tensor submultiplicativity across the same fixed row cut now
gives
\[
 \kappa_\otimes^{\,k|k^c}
 \bigl(\mathcal B_{k\ell}(D)\bigr)
 \le
 C(1-c''s_\alpha)^N
 \le
 C e^{-c''s_\alpha N}.
\]
Here \(0<s_\alpha\le1\) is the established weak-link range; constants
are enlarged for its fixed compact remainder.

Apply \Cref{lem:V70-row-power}:
\[
 s_\alpha S_{k\ell}(D)^2 e^{-c''s_\alpha N}
 \le
 C(s_\alpha N)e^{-c''s_\alpha N}
 H_{k\ell}^{(k)}(D)
 \le
 C H_{k\ell}^{(k)}(D).
\]
The binary-prefix moment formula and the assembled local covariance
contain only a fixed number of global partition states.  Sum those
states after the bank estimate; this changes only the universal
constant.  Since the proof never introduces a graph difference, the
V59 graphwise regression is inapplicable.
\end{proof}

\begin{proposition}[Separated protected and gapped bank currencies]
\label{prop:V70-separated-currencies}
Let \(N=N_{k\ell}(D)\ge1\), let
\(\mathcal P_{k\ell}(D)\) be the sector in which every coordinate of
the bank is protected, and let \(\mathcal G_{k\ell}(D)\) be its
orthogonal complement inside the complete bank.  Then
\begin{align}
 \kappa_\otimes^{\,k|k^c}
 \bigl(\mathcal P_{k\ell}(D)\bigr)
 &\le
 \frac{C}{N}\,
 \frac{H_{k\ell}^{(k)}(D)}
      {S_{k\ell}(D)^2},
\label{eq:V70-protected-currency}\\
 \kappa_\otimes^{\,k|k^c}
 \bigl(\mathcal G_{k\ell}(D)\bigr)
 &\le
 C\min\left\{
 e^{-cs_\alpha N},
 \frac{H_{k\ell}^{(k)}(D)}
      {s_\alpha S_{k\ell}(D)^2}
 \right\}.
\label{eq:V70-gapped-currency}
\end{align}
If \(S_{k\ell}(D)\ge c_0\Xi_k\), these imply
\begin{align}
 \kappa_\otimes(\mathcal P_{k\ell}(D))
 &\le
 \frac{C_{c_0}}{N}\,
 \theta_{k\ell}\frac{\Xi_\ell}{\Xi_k},
\label{eq:V70-protected-ratio}\\
 \kappa_\otimes(\mathcal G_{k\ell}(D))
 &\le
 C_{c_0}\min\left\{
 e^{-cs_\alpha N},
 \frac1{s_\alpha}\,
 \theta_{k\ell}\frac{\Xi_\ell}{\Xi_k}
 \right\}.
\label{eq:V70-gapped-ratio}
\end{align}
\end{proposition}

\begin{proof}
Tensoring \eqref{eq:V70-P-capacity} over the distinct coordinates
gives
\[
 \kappa_\otimes(\mathcal P_{k\ell}(D))\le C\rho^N.
\]
Since \(\sup_{N\ge1}N^2\rho^N<\infty\) and
\eqref{eq:V70-row-power} gives
\(H_{k\ell}^{(k)}(D)/S_{k\ell}(D)^2\ge c/N\), this proves
\eqref{eq:V70-protected-currency}.

The complementary sector is a positive compression of the complete
bank, so the affine local argument in
\Cref{thm:V70-balanced-tail} gives its first bound in
\eqref{eq:V70-gapped-currency}.  The second follows by dividing
\eqref{eq:V70-balanced-tail} by \(s_\alpha S_{k\ell}(D)^2\).

If \(S_{k\ell}(D)\ge c_0\Xi_k\), then
\[
 \frac{H_{k\ell}^{(k)}(D)}
      {S_{k\ell}(D)^2}
 \le
 C_{c_0}
 \frac{H_{k\ell}}
      {\Xi_k^2}
 =
 C_{c_0}\theta_{k\ell}\frac{\Xi_\ell}{\Xi_k}.
\]
Insert this into the first two estimates.
\end{proof}

\begin{lemma}[A heavy marked balanced coordinate has the same card]
\label{lem:V70-heavy-mark-card}
Let \(q\in Q\cap D_{k\ell}\).  For the residual physical multiplier
\(\mathcal B_q\) of that marked coordinate,
\begin{equation}
 \boxed{
 s_\alpha a_{k,q}^2
 \kappa_\otimes^{\,k|k^c}(\mathcal B_q)
 \le
 C a_{k,q}a_{\ell,q}.}
\label{eq:V70-heavy-mark-card}
\end{equation}
\end{lemma}

\begin{proof}
The residual multiplier is a contraction after the fixed local source
constant is removed.  Balance and the maximal-partner choice in the
proof of \Cref{lem:V70-row-power} give
\(a_{k,q}\le C a_{\ell,q}\).  Since \(s_\alpha\le1\),
\[
 s_\alpha a_{k,q}^2
 \kappa_\otimes(\mathcal B_q)
 \le a_{k,q}^2
 \le C a_{k,q}a_{\ell,q}.
\]
Keeping the sharper alternatives of
\Cref{lem:V70-local-PG} is permitted but not needed for this
single-coordinate inequality.
\end{proof}

\section{A one-use transfer certificate for every balanced-hot row}
\label{sec:V70-transfer}

For a nonempty bank \(D_{k\ell}\), define
\begin{equation}
 \vartheta_{k\to\ell}(D)
 :=
 \frac{H_{k\ell}^{(k)}(D)}
      {S_{k\ell}(D)\Xi_\ell}.
\label{eq:V70-transfer-factor}
\end{equation}
It obeys
\begin{equation}
 0\le\vartheta_{k\to\ell}(D)\le1,
\label{eq:V70-transfer-contraction}
\end{equation}
because \(a_{\ell,u}\le\Xi_\ell\) termwise.

\begin{theorem}[Balanced sourcewise transfer certificate]
\label{thm:V70-transfer-certificate}
Let \(k\) be a balanced-hot internal endpoint satisfying
\eqref{eq:V70-balanced-hot}, and let \(Q\) be the at most three marked
coordinates of one ordered \(2+2\) source monomial.  Then there exist
a partner \(\ell\ne k\) and either
\begin{enumerate}[label=\textup{(\roman*)}]
\item an unmarked reserve bank
      \(D\subseteq E\setminus Q\), or
\item one marked coordinate \(q\in Q\),
\end{enumerate}
whose retained physical multiplier \(\mathcal C\) satisfies
\begin{equation}
 \boxed{
 s_\alpha\Xi_k\,
 \kappa_\otimes^{\,k|k^c}(\mathcal C)
 \le
 C\Xi_\ell\,
 \vartheta_{k\to\ell},}
\label{eq:V70-transfer-certificate}
\end{equation}
where the right side uses the stock of the selected reserve bank or
the singleton stock at \(q\).  Moreover
\begin{equation}
 \vartheta_{k\to\ell}
 \le C\theta_{k\ell}.
\label{eq:V70-transfer-to-tree-mark}
\end{equation}
The certificate is obtained inside the assembled first-connectivity
packet and uses no final output resolvent.
\end{theorem}

\begin{proof}
Choose \(\ell\) by \Cref{lem:V70-large-partner-bank} and abbreviate the
full selected stock by \(S\), so \(S\ge\Xi_k/12\).  Split it into the
unmarked part \(S^\circ\) and the at most three marked contributions.

If \(S^\circ\ge S/2\), take \(D=D_{k\ell}\setminus Q\).
Then \(S^\circ\ge\Xi_k/24\).  Divide
\eqref{eq:V70-balanced-tail} by \(S^\circ\):
\[
 s_\alpha\Xi_k\kappa_\otimes(\mathcal C)
 \le
 24s_\alpha S^\circ\kappa_\otimes(\mathcal C)
 \le
 C\frac{H_{k\ell}^{(k)}(D)}{S^\circ}
 =
 C\Xi_\ell\vartheta_{k\to\ell}(D).
\]

If \(S^\circ<S/2\), the marked contribution exceeds \(S/2\).  Some
\(q\in Q\) has
\[
 a_{k,q}\ge S/6\ge\Xi_k/72.
\]
Divide \eqref{eq:V70-heavy-mark-card} by \(a_{k,q}\) and use this
comparison to obtain
\[
 s_\alpha\Xi_k\kappa_\otimes(\mathcal B_q)
 \le
 C a_{\ell,q}
 =
 C\Xi_\ell
 \frac{a_{k,q}a_{\ell,q}}{a_{k,q}\Xi_\ell}.
\]
This is \eqref{eq:V70-transfer-certificate} for the singleton bank.

Finally, for either selected bank,
\[
 \vartheta_{k\to\ell}
 =
 \frac{\Xi_k}{S_{k\ell}}
 \frac{H_{k\ell}^{(k)}}{\Xi_k\Xi_\ell}
 \le
 C\theta_{k\ell},
\]
because the selected bank stock is at least a fixed multiple of
\(\Xi_k\) and
\eqref{eq:V70-H-domination} holds.  The multiplier \(\mathcal C\)
comes from spectator physical blocks or from one already marked local
block, not from the final payer, proving the last assertion.
\end{proof}

\begin{remark}[What the certificate transfers]
\label{rem:V70-transfer-meaning}
The left side of \eqref{eq:V70-transfer-certificate} is one hot
endpoint factor.  The right side replaces it by one normalized legal
edge mark and the partner mass \(\Xi_\ell\).  Thus the analytic
balanced problem is closed at one row, but a finite four-vertex
topology problem remains: the partner mass must become a legal excess
degree of a comparison spanning tree or be paid by another retained
bank.
\end{remark}

\begin{firewall}[No manufactured tree edge]
\label{fire:V70-no-manufactured-edge}
The factor \(\theta_{k\ell}\) in
\eqref{eq:V70-transfer-to-tree-mark} is backed by the actual overlap
stock \(H_{k\ell}^{(k)}\) of the retained balanced bank.  It is not
assigned to a graph edge introduced by a cumulant identity.  Tree
comparison is postponed until after the complete source packet has
supplied this physical edge certificate.
\end{firewall}

\section{Post-audit reduction and next acquisition}
\label{sec:V70-reduction}

\begin{theorem}[Post-V70 balanced reduction]
\label{thm:V70-reduction}
For every balanced-hot internal endpoint in the assembled \(2+2\)
source, V70 supplies a sourcewise one-use transfer
\[
 k\longrightarrow\ell,
 \qquad
 s_\alpha\Xi_k\mapsto
 C\Xi_\ell\theta_{k\ell}.
\]
Consequently the unresolved low/private-output scalar problem is
finite: starting from one of the four paths
\(T_{ij}=\{12,34,ij\}\), determine whether the certified edge
\(k\ell\) permits a comparison spanning tree whose degree-excess
masses absorb the transported \(\Xi_\ell\).  Parallel-edge and
two-cycle cases must be treated without counting the same certificate
twice.  Balanced high-output ancestry adds at most one further
labelled transfer and must use the same sole payer.
\end{theorem}

\begin{proof}
\Cref{cor:V69-balanced-residual} leaves only balanced-hot endpoints and
balanced output ancestry.  Apply
\Cref{thm:V70-transfer-certificate} to each selected endpoint, after
the ordered source monomial and its physical bank have been fixed.
There are four choices of \(T_{ij}\), four vertices, and at most three
partner choices per selected row; hence the remaining routing state
space is finite and support-independent.  The theorem makes no claim
that every state closes: it identifies exactly the scalar comparison
still required after the analytic bank estimate.
\end{proof}

\begin{assessment}[Effect of the companion audit]
\label{ass:V70-audit-effect}
The audit changes the proof order, not the mathematical currency.
The balanced bank is first split into its complete protected and gapped
physical sectors, like the singular/complementary spectral split in SM
V86.  Its mass is then transferred once, matching the one-use lesson
of NS V31.  Only after both operations is the normalized tree topology
examined.  This avoids both an orbitwise triangle inequality and a
second charge of the final resolvent.
\end{assessment}

\section{V70 ledger and V71 mandate}
\label{sec:V70-ledger}

\begin{theorem}[Integrated compact-series V70 statement]
\label{thm:V70-integrated}
Preserving compact V1--V69, V70 proves:
\begin{enumerate}[label=\textup{(V70.\arabic*)},leftmargin=6.7em]
\item the mandatory SM V86 and NS V31 audit with explicit fingerprints;
\item a row-rooted balanced power inequality and a large-partner bank;
\item an assembled protected/gapped spectral alternative at every
      balanced coordinate;
\item a support-uniform balanced-bank tail inside the exact \(2+2\)
      source;
\item separated protected and gapped bank currencies, including the
      protected \(N^{-1}\) mass-ratio gain;
\item the corresponding heavy marked-coordinate card; and
\item a one-use transfer certificate for every balanced-hot internal
      endpoint, reducing the remaining analytic problem to finite
      four-vertex routing plus balanced output ancestry.
\end{enumerate}
\end{theorem}

\begin{proof}
These are
\Cref{ass:V70-SM-transfer,
ass:V70-NS-transfer,
lem:V70-row-power,
lem:V70-large-partner-bank,
lem:V70-local-PG,
thm:V70-balanced-tail,
prop:V70-separated-currencies,
lem:V70-heavy-mark-card,
thm:V70-transfer-certificate,
thm:V70-reduction}.
\end{proof}

\begin{openproblem}[V71 mandate: enumerate the certified routing
states]
\label{op:V71-mandate}
\begin{enumerate}
\item For each \(T_{ij}=\{12,34,ij\}\), each internal endpoint
      \(k\in\{i,j\}\), and each certified partner \(\ell\ne k\),
      enumerate the simple graph \(T_{ij}\cup\{k\ell\}\), retaining
      parallel occurrences.
\item Determine exactly when the transported mass \(\Xi_\ell\) is a
      degree-excess factor of a spanning tree contained in that graph.
\item In every failure state, apply a second row-rooted certificate
      only after verifying that it uses a disjoint bank or a higher
      moment of the same retained bank.
\item Prove termination on four vertices or identify the unique
      directed cycle/parallel state that prevents it.
\item Add one labelled balanced-output ancestry transfer and verify
      that the final payer is still used exactly once.
\item Do not invoke V47 graph differences: the graph in this
      enumeration contains only the original source edges and the
      physical certified overlap edges.
\end{enumerate}
\end{openproblem}

\begin{firewall}[Claim boundary after compact V70]
\label{fire:V70-current-boundary}
V70 proves the balanced bank tail and a sourcewise transfer certificate.
It does not yet prove that every finite transfer state closes, pay all
balanced high-output ancestry, extend the hot result to the other three
source families, or restore the full all-weak quartic theorem.  Global
quartic and higher MTA, WUFC, CPM, continuum Yang--Mills existence, and
a positive mass gap remain open.  The full Yang--Mills mass gap remains
open.
\end{firewall}

\section{Append-only record for compact V70}
\label{sec:V70-record}

V70 was appended to the complete compact V69 file.  The exact V69
parent had \(28313\) newline-delimited source records, \(977542\)
bytes, and SHA--256 digest
\[
\texttt{d4d6a94dca78d998cf8cc9f9ca80cde0bd080029f18c1d6bddae2c99f16e31d0}.
\]
No compact V69 mathematical content was changed.  The sole final
\verb|\end{document}| is restored below.

% =============================================================================
% END APPENDED COMPACT VERSION V70 MATERIAL
% =============================================================================

\clearpage

% =============================================================================
% BEGIN APPENDED COMPACT VERSION V71 MATERIAL
% =============================================================================

\chapter{Clean sparse protected chains in the assembled \(2+2\)
source}
\label{chap:V71}

\section{Context and branch definition}
\label{sec:V71-context}

V70 converted a balanced-hot row into two separate analytic
currencies.  The all-protected sector supplies a normalized
mass-ratio factor with no \(s_\alpha^{-1}\) loss, whereas the
complementary sector supplies a gapped count tail.  They must not be
merged during routing.

This note closes the clean sparse all-protected branch.  The proof is
source-specific: all protected banks are coordinate-disjoint, are
disjoint from the three source marks, and are estimated in each fixed
global partition state of the two binary prefixes and assembled
joining covariance.  Under those hypotheses their four-row
product-state capacities tensorize before the finite state sum.  A
strict row-mass escalation prevents cycles, and a private or
unbalanced terminal pays the remaining moment.  The additional
physical transfer marks are contractions, so after the scalar
telescope they may be deleted while retaining the original \(2+2\)
spanning tree.

Work in the all-weak sector
\begin{equation}
 \theta_{uv}\le\delta
 \qquad(u\ne v),
\label{eq:V71-all-weak}
\end{equation}
the only quartic threshold sector still at issue.  Fix
\[
 T=a-i-j-b,
\]
one of the paths \(T_{ij}=\{12,34,ij\}\), with internal rows \(i,j\)
and leaves \(a,b\).

\begin{definition}[Clean sparse protected transfer]
\label{def:V71-clean-transfer}
A transfer \(u\to v\) is clean, sparse, and protected if:
\begin{enumerate}[label=\textup{(\roman*)}]
\item row \(u\) has a row-rooted balanced bank
      \(D_{uv}\) with
      \(S_{uv}\ge c_0\Xi_u\), where \(c_0>0\) is fixed;
\item \(D_{uv}\) is disjoint from the source mark set
      \(Q=\{e,f,r\}\) and from every earlier transfer bank;
\item \(1\le N_{uv}:=|D_{uv}|<L\), where \(L\ge2\) is fixed; and
\item the complete bank lies in the all-protected sector
      \(\mathcal P_{uv}(D_{uv})\) of
      \Cref{prop:V70-separated-currencies}.
\end{enumerate}
Its physical mark and mass ratio are
\begin{equation}
 \vartheta_{uv}
 :=
 \frac{H_{uv}^{(u)}(D_{uv})}{\Xi_u\Xi_v}
 \le\theta_{uv},
 \qquad
 R_{uv}:=\frac{\Xi_v}{\Xi_u}.
\label{eq:V71-transfer-data}
\end{equation}
\end{definition}

\section{Strict mass escalation on a sparse edge}
\label{sec:V71-escalation}

\begin{lemma}[Quantitative protected transfer factor]
\label{lem:V71-protected-factor}
Every transfer in \Cref{def:V71-clean-transfer} has the statewise
capacity bound
\begin{equation}
 \boxed{
 \kappa_\otimes(\mathcal P_{uv}(D_{uv}))
 \le
 \frac{C_{c_0}}{N_{uv}}\,
 \vartheta_{uv}R_{uv}.}
\label{eq:V71-protected-factor}
\end{equation}
\end{lemma}

\begin{proof}
This is \eqref{eq:V70-protected-ratio}, with the subbank mark
\(\vartheta_{uv}\) retained instead of enlarging it immediately to
\(\theta_{uv}\).
\end{proof}

\begin{lemma}[Sparse transfers more than double row mass]
\label{lem:V71-doubling}
There is a universal \(C_0\) such that every balanced partner bank in
\Cref{def:V71-clean-transfer} satisfies
\begin{equation}
 \Xi_u\le C_0\delta N_{uv}\Xi_v.
\label{eq:V71-mass-transport}
\end{equation}
Consequently, if
\begin{equation}
 \delta<\frac1{2C_0L},
\label{eq:V71-delta-choice}
\end{equation}
then every sparse transfer obeys
\begin{equation}
 \Xi_v>2\Xi_u.
\label{eq:V71-doubling}
\end{equation}
\end{lemma}

\begin{proof}
The lower stock bound and \Cref{lem:V70-row-power} give
\[
 c_0^2\Xi_u^2
 \le S_{uv}^2
 \le C N_{uv}H_{uv}^{(u)}.
\]
By \eqref{eq:V71-all-weak} and
\eqref{eq:V70-H-domination},
\[
 H_{uv}^{(u)}
 \le H_{uv}
 \le\delta\Xi_u\Xi_v.
\]
Cancel \(\Xi_u>0\) and absorb \(C/c_0^2\) into \(C_0\), proving
\eqref{eq:V71-mass-transport}.  If \(N_{uv}<L\), the choices in
\eqref{eq:V71-delta-choice} give
\(\Xi_v>2\Xi_u\).
\end{proof}

\begin{corollary}[A clean sparse chain has length at most three]
\label{cor:V71-chain-length}
A sequence
\[
 i_0\longrightarrow i_1\longrightarrow\cdots
 \longrightarrow i_m
\]
of clean sparse transfers among four rows has \(m\le3\) and has no
repeated vertex.
\end{corollary}

\begin{proof}
Along the chain, \(\Xi_{i_{\ell+1}}>2\Xi_{i_\ell}\).
A repeated vertex would have mass strictly larger than itself.
Therefore all vertices are distinct, and four vertices permit at most
three edges.
\end{proof}

\section{Authorization to multiply clean sourcewise capacities}
\label{sec:V71-authorization}

\begin{lemma}[Statewise tensorization over disjoint transfer banks]
\label{lem:V71-statewise-tensorization}
Fix one global moment-partition state in each of the two binary
prefixes and in the assembled local joining covariance.  For a clean
chain with pairwise coordinate-disjoint protected banks,
\begin{equation}
 \kappa_{\otimes,4}\!\left(
 \bigotimes_{\ell=0}^{m-1}
 \mathcal P_{i_\ell i_{\ell+1}}
 \right)
 \le
 \prod_{\ell=0}^{m-1}
 \kappa_{\otimes,4}\!\left(
 \mathcal P_{i_\ell i_{\ell+1}}
 \right).
\label{eq:V71-statewise-tensorization}
\end{equation}
The same product bound holds after the fixed finite global-state sum
of the exact \(2+2\) first-connectivity source.
\end{lemma}

\begin{proof}
Each protected bank acts on a disjoint tensor product of coordinate
Hilbert spaces.  The four-row product-state capacity is the injective
crossnorm with respect to the four row factors.  Its one-sided tensor
submultiplicativity gives
\eqref{eq:V71-statewise-tensorization}, even though consecutive
transfers use different row cuts: a four-row product vector is a
product vector for each of those cuts.

The two connected binary prefixes have two global partition states
each, and the assembled joining covariance has two composite states.
Thus the source uses a fixed finite state set.  Apply the statewise
bound first and sum that set only afterward.  No assertion of the form
\(\Lambda(xy)=\Lambda(x)\Lambda(y)\) is made, and no cancellation is
assigned separately to two banks.
\end{proof}

\begin{firewall}[Why the V46 disjoint-support warning is respected]
\label{fire:V71-V46-warning}
Coordinate disjointness alone would not let two independently
cut-connected cumulant functionals be multiplied.  V71 uses a
different statement: the exact \(2+2\) source has already been fixed
in one common finite global state, and each V70 protected estimate is
valid statewise for that same source.  Capacities tensorize inside the
state; the common state functional is applied once at the end.
\end{firewall}

\section{Terminal moment certificates}
\label{sec:V71-terminal}

\begin{definition}[Admissible clean terminal]
\label{def:V71-terminal}
The final row \(i_m\) of a clean chain is an admissible terminal if
the exact source packet retains a nonnegative factor
\(\mathcal Q_{i_m}\le1\) such that
\begin{equation}
 (s_\alpha\Xi_{i_m})^q\mathcal Q_{i_m}
 \le C_q,
 \qquad q=1,2,3.
\label{eq:V71-terminal-moments}
\end{equation}
\end{definition}

\begin{lemma}[Private and unbalanced rows are admissible terminals]
\label{lem:V71-terminal-types}
If
\[
 p_{i_m}\ge\frac12\Xi_{i_m},
\]
then one may take the exact private Wilson product of that row.
If the unbalanced shared stock of row \(i_m\) is at least a fixed
positive fraction of \(\Xi_{i_m}\), one may take the retained
unbalanced reserve/marked tail of V69.  In either case
\eqref{eq:V71-terminal-moments} holds.
\end{lemma}

\begin{proof}
For the private alternative use
\eqref{eq:V52-private-products-all-orders}; its factor is part of the
exact packetwise multiplier of
\Cref{thm:V68-packetwise-private-factor}.  For the unbalanced
alternative use the reserve/heavy dichotomy
\Cref{lem:V69-BRH} and the moments
\eqref{eq:V69-reserve-moments}--%
\eqref{eq:V69-marked-moments}.  Balanced transfer banks and
unbalanced coordinate banks are disjoint, so none of the protected
capacities used along the chain is spent again at the terminal.
\end{proof}

\section{The protected scalar telescope}
\label{sec:V71-telescope}

For a clean chain write
\[
 R_\ell:=\frac{\Xi_{i_{\ell+1}}}{\Xi_{i_\ell}},
 \qquad
 \vartheta_\ell:=
 \vartheta_{i_\ell i_{\ell+1}},
 \qquad
 N_\ell:=N_{i_\ell i_{\ell+1}}.
\]

\begin{lemma}[Exact \(q\)-moment ratio telescope]
\label{lem:V71-ratio-telescope}
For \(q\in\{1,2,3\}\),
\begin{equation}
\boxed{
 (s_\alpha\Xi_{i_0})^q
 \left[
  \prod_{\ell=0}^{m-1}
  \frac{\vartheta_\ell R_\ell}{N_\ell}
 \right]
 \frac1{(s_\alpha\Xi_{i_m})^q}
 =
 \left(\prod_{\ell=0}^{m-1}
       \frac{\vartheta_\ell}{N_\ell}\right)
 \left(\frac{\Xi_{i_0}}{\Xi_{i_m}}\right)^{q-1}.}
\label{eq:V71-ratio-telescope}
\end{equation}
For a nonempty chain its right side is at most
\(\prod_\ell\vartheta_\ell\).
\end{lemma}

\begin{proof}
The mass ratios telescope:
\[
 \prod_{\ell=0}^{m-1}R_\ell
 =
 \frac{\Xi_{i_m}}{\Xi_{i_0}}.
\]
Substitution gives the equality.  Every \(N_\ell\ge1\), and
\Cref{lem:V71-doubling} gives
\(\Xi_{i_m}>\Xi_{i_0}\), proving the final assertion.
\end{proof}

\begin{theorem}[Clean sparse protected-chain closure]
\label{thm:V71-clean-chain-closure}
Suppose a branch of the exact assembled \(2+2\) source contains:
\begin{enumerate}[label=\textup{(\roman*)}]
\item at most \(q\le3\) unpaid factors
      \(s_\alpha\Xi_{i_0}\);
\item a clean sparse protected chain from \(i_0\) to an admissible
      terminal \(i_m\); and
\item no use of a transfer bank by the final output payer.
\end{enumerate}
Then its scalar coefficient, including all transfer-bank capacities,
is bounded by
\begin{equation}
 C\prod_{\ell=0}^{m-1}\vartheta_\ell.
\label{eq:V71-chain-bound}
\end{equation}
The constant is uniform in support, representations, and volume.
\end{theorem}

\begin{proof}
Use \Cref{lem:V71-statewise-tensorization} and multiply the protected
factors \eqref{eq:V71-protected-factor}.  The terminal estimate
\eqref{eq:V71-terminal-moments} contributes
\(C_q(s_\alpha\Xi_{i_m})^{-q}\).  The remaining scalar is exactly the
left side of \eqref{eq:V71-ratio-telescope}, so
\eqref{eq:V71-chain-bound} follows.  There are at most three transfer
banks and a fixed number of source states, hence all suppressed
constants are universal.
\end{proof}

\section{Application to the \(2+2\) path}
\label{sec:V71-application}

\begin{lemma}[Largest internal row roots the low-output moment]
\label{lem:V71-largest-root}
Let \(i_0\) be whichever of the two internal vertices \(i,j\) has
larger row mass.  Then
\begin{equation}
 s_\alpha^2\Xi_i\Xi_j
 \le
 (s_\alpha\Xi_{i_0})^2.
\label{eq:V71-root-two}
\end{equation}
For the private part \(P=\sum_t p_t\) of a high-output payer, its
labelled summands require total private/protected terminal degree at
most three.
\end{lemma}

\begin{proof}
The first assertion is immediate from
\(\Xi_i\Xi_j\le\max\{\Xi_i,\Xi_j\}^2\).  A labelled private-output
summand adds one factor \(s_\alpha p_t\) to the two endpoint factors.
The exact private bank and terminal factors have all fixed moments by
V52 and V69; if they coincide, their exponents add to at most three.
\end{proof}

\begin{theorem}[Closed clean protected part of the assembled
\(2+2\) source]
\label{thm:V71-protected-22}
Choose the larger internal row \(i_0\) of each ordered path monomial.
If its balanced routing consists of a clean sparse all-protected chain
ending at a private or unbalanced terminal, then both:
\begin{enumerate}[label=\textup{(\alph*)}]
\item the low-output contribution; and
\item the private part of the high-output contribution
\end{enumerate}
satisfy the original DRQC tree bound
\begin{equation}
 C\,\theta_{12}\theta_{34}\theta_{ij}.
\label{eq:V71-original-tree-bound}
\end{equation}
\end{theorem}

\begin{proof}
For low output, use \eqref{eq:V71-root-two} and
\Cref{thm:V71-clean-chain-closure} with \(q=2\).  For a labelled
private-output summand, use the same telescope together with the
additional private factor; the total moment degree is at most three by
\Cref{lem:V71-largest-root}.

Both cases leave the original source mark
\(\theta_{12}\theta_{34}\theta_{ij}\) multiplied by
\(\prod_\ell\vartheta_\ell\).  Every transfer factor lies in
\([0,1]\), so
\[
 \theta_{12}\theta_{34}\theta_{ij}
 \prod_\ell\vartheta_\ell
 \le
 \theta_{12}\theta_{34}\theta_{ij}.
\]
Thus no comparison graph or alternative spanning tree is needed after
the scalar coefficient has been paid.  All transfer marks came from
actual retained overlap banks.
\end{proof}

\begin{proposition}[Complete one-edge routing enumeration]
\label{prop:V71-one-edge-enumeration}
Write the original path as \(a-i-j-b\).  A transfer from \(i\) has
exactly three partner types:
\begin{enumerate}[label=\textup{(\roman*)}]
\item \(i\to a\), parallel to the leaf edge;
\item \(i\to j\), parallel to the central edge;
\item \(i\to b\), the unique nonparallel chord.
\end{enumerate}
The reflected list holds from \(j\).  All three types are harmless in
\Cref{thm:V71-protected-22}, including parallel occurrences.
\end{proposition}

\begin{proof}
There are only the other three vertices.  The first two are already
neighbours of \(i\) in \(T\), while the opposite leaf \(b\) is not.
After the scalar telescope, each occurrence contributes only a factor
\(\vartheta\le1\), so it may be deleted regardless of whether its
underlying simple edge was already present.  This proves the last
assertion without manufacturing a tree degree.
\end{proof}

\begin{firewall}[The protected theorem does not cover a mixed bank]
\label{fire:V71-protected-only}
The proof applies only when every selected transfer bank is in its
all-protected sector.  If a bank contains the complementary gapped
sector, its available factor is
\eqref{eq:V70-gapped-ratio}, which carries a different
\(s_\alpha^{-1}\) alternative and must be combined with its exponential
count tail before routing.  Substituting it into the protected
telescope would be invalid.
\end{firewall}

\section{Exact residual after the protected-chain theorem}
\label{sec:V71-residual}

\begin{theorem}[Post-V71 \(2+2\) residual]
\label{thm:V71-residual}
Relative to the balanced residual of
\Cref{cor:V69-balanced-residual}, the clean sparse all-protected branch
with a private or unbalanced terminal is closed.  Every remaining
\(2+2\) term has at least one of:
\begin{enumerate}
\item a selected balanced transfer bank in its gapped complementary
      sector;
\item a dense selected bank, \(N\ge L\);
\item a selected bank concentrated on one of the at most three source
      marks;
\item two required transfer banks sharing a coordinate;
\item no private or unbalanced terminal before one of the preceding
      cases occurs; or
\item the balanced part \(Z_\nu^B\) of the high-output payer.
\end{enumerate}
\label{enum:V71-residual}
\end{theorem}

\begin{proof}
V69 leaves balanced internal/output ancestry.  V70 gives, for every
balanced-hot row, either an unmarked reserve or a heavy source mark and
splits every reserve into protected and gapped sectors.  If successive
reserve banks are disjoint, sparse, protected, and terminate in a
private/unbalanced row, apply
\Cref{thm:V71-protected-22}.  Failure of one of those hypotheses gives
items 1--5.  V71 has not allocated \(Z_\nu^B\), giving item 6.
The list is exhaustive by construction.
\end{proof}

\begin{assessment}[Next upstream target]
\label{ass:V71-next-target}
The gapped complementary sector should be treated before coordinate
overlaps.  Its retained factor has two simultaneous descriptions:
\[
 e^{-cs_\alpha N},
 \qquad
 \frac1{s_\alpha}\theta_{uv}\frac{\Xi_v}{\Xi_u}.
\]
The first pays powers of the bank count; the second transports row
mass but loses one \(s_\alpha\).  The next note must split dense and
sparse counts at a scale determined by the number \(q\le3\) of unpaid
row/output factors, then combine both descriptions before taking the
minimum.  Treating either description alone cannot close all counts.
\end{assessment}

\section{V71 ledger and V72 mandate}
\label{sec:V71-ledger}

\begin{theorem}[Integrated compact-series V71 statement]
\label{thm:V71-integrated}
Preserving compact V1--V70, V71 proves:
\begin{enumerate}[label=\textup{(V71.\arabic*)},leftmargin=6.7em]
\item sparse balanced transfers strictly more than double row mass and
      hence form chains of length at most three;
\item clean coordinate-disjoint protected capacities tensorize
      statewise inside the exact \(2+2\) source;
\item private and unbalanced terminal moments close the exact
      \(q\le3\) scalar ratio telescope;
\item every clean sparse protected low/private-output \(2+2\) branch
      obeys its original spanning-tree mark; and
\item the remaining balanced cases are the explicit six classes in
      \eqref{enum:V71-residual}.
\end{enumerate}
\end{theorem}

\begin{proof}
These are
\Cref{lem:V71-doubling,
cor:V71-chain-length,
lem:V71-statewise-tensorization,
lem:V71-terminal-types,
lem:V71-ratio-telescope,
thm:V71-protected-22,
thm:V71-residual}.
\end{proof}

\begin{openproblem}[V72 mandate: gapped count/mass interpolation]
\label{op:V72-mandate}
\begin{enumerate}
\item Keep the two entries of
      \eqref{eq:V70-gapped-ratio} on the same selected bank.
\item For \(q=2\) endpoint factors and \(q=3\) with one labelled output
      payer, determine the sharp split in \(s_\alpha N\) at which the
      exponential count tail pays the inverse-scale loss of the
      transport entry.
\item Preserve all transfer ratios until a private or unbalanced
      terminal moment is applied; do not estimate ratios edgewise.
\item Treat a dense bank directly from the exponential entry and a
      sparse bank by mass escalation, recording any intermediate count
      power exactly.
\item Only after the scalar coefficient is bounded, delete the
      physical transfer marks and retain the original \(2+2\) tree.
\item Leave marked-heavy and overlapping-bank cases separate unless
      the same interpolation genuinely retains their local physical
      multipliers.
\end{enumerate}
\end{openproblem}

\begin{firewall}[Claim boundary after compact V71]
\label{fire:V71-current-boundary}
V71 closes only the clean sparse all-protected chain branch of the
assembled \(2+2\) source.  The six residual classes in
\eqref{enum:V71-residual}, the other three first-connectivity families,
global quartic and higher MTA, WUFC, CPM, continuum Yang--Mills
existence, and a positive mass gap remain open.  The full Yang--Mills
mass gap remains open.
\end{firewall}

\section{Append-only record for compact V71}
\label{sec:V71-record}

V71 was appended to the complete compact V70 file.  The exact V70
parent had \(28957\) newline-delimited source records, \(999512\)
bytes, and SHA--256 digest
\[
\texttt{19e6931c1a9a39d0b640ab668b807ab1083781a3540744c523d69ca2eb000871}.
\]
No compact V70 mathematical content was changed.  The sole final
\verb|\end{document}| is restored below.

% =============================================================================
% END APPENDED COMPACT VERSION V71 MATERIAL
% =============================================================================

\clearpage

% =============================================================================
% BEGIN APPENDED COMPACT VERSION V72 MATERIAL
% =============================================================================

\chapter{Sparse mixed protected/gapped chains without an
inverse-scale loss}
\label{chap:V72}

\section{Context and correction of the proposed interpolation}
\label{sec:V72-context}

V71 closed a clean sparse chain when every selected balanced bank lies
in its all-protected sector.  V70 also gave, for the complementary
gapped sector, the pair of bounds
\[
 e^{-cs_\alpha N},
 \qquad
 \frac1{s_\alpha}\theta_{uv}\frac{\Xi_v}{\Xi_u}.
\]
The V71 mandate proposed interpolating those two entries.  That is
unnecessary and obscures a simpler fact.  Sparse all-weak mass
transport already gives
\[
 R_{uv}:=\frac{\Xi_v}{\Xi_u}>2.
\]
Therefore the exponential entry itself may be written as a transport
factor:
\[
 e^{-cs_\alpha N}\le
 R_{uv}e^{-cs_\alpha N}.
\]
The row ratio needed by the terminal telescope is retained, while the
remaining exponential is a contraction.  No \(s_\alpha^{-1}\) is
created, no higher terminal moment is required, and no new graph mark
is claimed.  This note implements that observation inside the exact
\(2+2\) source and closes every clean sparse chain with an arbitrary
mixture of protected and gapped selected banks.

\section{A common ratio form for both spectral sectors}
\label{sec:V72-common-ratio}

\begin{definition}[Clean sparse spectral chain]
\label{def:V72-spectral-chain}
A clean sparse spectral chain is a V71 clean chain
\[
 i_0\longrightarrow i_1\longrightarrow\cdots
 \longrightarrow i_m
\]
with pairwise coordinate-disjoint, unmarked selected banks, except
that each bank is allowed to lie either in:
\begin{enumerate}[label=\textup{(\roman*)}]
\item its all-protected sector \(\mathcal P_\ell\), or
\item its gapped complementary sector \(\mathcal G_\ell\).
\end{enumerate}
The all-weak and sparse choices are
\eqref{eq:V71-all-weak} and \eqref{eq:V71-delta-choice}.
\end{definition}

For the \(\ell\)-th edge put
\[
 R_\ell:=\frac{\Xi_{i_{\ell+1}}}{\Xi_{i_\ell}},
 \qquad
 N_\ell:=N_{i_\ell i_{\ell+1}},
 \qquad
 \vartheta_\ell:=
 \frac{H_{i_\ell i_{\ell+1}}^{(i_\ell)}}
      {\Xi_{i_\ell}\Xi_{i_{\ell+1}}}.
\]

\begin{lemma}[Sector factors in a common transport form]
\label{lem:V72-common-transport}
For each edge of a clean sparse spectral chain, its retained sector
factor \(\mathcal A_\ell\) satisfies
\begin{equation}
 \boxed{
 \kappa_{\otimes,4}(\mathcal A_\ell)
 \le
 C R_\ell\gamma_\ell,}
\label{eq:V72-common-transport}
\end{equation}
where
\begin{equation}
 \gamma_\ell:=
 \begin{cases}
  \vartheta_\ell/N_\ell,
       &\mathcal A_\ell=\mathcal P_\ell,\\[2mm]
  e^{-cs_\alpha N_\ell},
       &\mathcal A_\ell=\mathcal G_\ell.
 \end{cases}
\label{eq:V72-gamma}
\end{equation}
In both cases \(0\le\gamma_\ell\le1\).
\end{lemma}

\begin{proof}
For a protected edge, use
\eqref{eq:V71-protected-factor}.  For a gapped edge,
\eqref{eq:V70-gapped-ratio} gives
\[
 \kappa_{\otimes,4}(\mathcal G_\ell)
 \le C e^{-cs_\alpha N_\ell}.
\]
By \Cref{lem:V71-doubling}, \(R_\ell>2>1\), so the last bound is at
most
\(CR_\ell e^{-cs_\alpha N_\ell}\).
Finally \(\vartheta_\ell\le1\), \(N_\ell\ge1\), and the exponential is
at most one.
\end{proof}

\begin{firewall}[No gapped edge mark is manufactured]
\label{fire:V72-no-gapped-mark}
On a gapped bank, \(\gamma_\ell\) is an analytic heat contraction, not
an overlap mark.  V72 never replaces it by a graph edge.  This is
legitimate because the original three \(2+2\) source marks already
form a spanning tree; every additional factor in the final scalar
bound only has to be at most one.
\end{firewall}

\section{Statewise multiplication of mixed sectors}
\label{sec:V72-statewise}

\begin{lemma}[Mixed clean-bank tensorization]
\label{lem:V72-mixed-tensorization}
In each fixed global state of the exact assembled \(2+2\) source,
\begin{equation}
 \kappa_{\otimes,4}\!\left(
  \bigotimes_{\ell=0}^{m-1}\mathcal A_\ell
 \right)
 \le
 C^m\prod_{\ell=0}^{m-1}R_\ell\gamma_\ell.
\label{eq:V72-mixed-tensorization}
\end{equation}
The same estimate, with a universal constant, holds after the fixed
finite source-state sum.
\end{lemma}

\begin{proof}
Selected banks are pairwise coordinate-disjoint.  Protected sectors
are positive physical compressions, and gapped complementary sectors
are orthogonal positive physical compressions of the same complete
bank.  The four-row injective capacity is one-sided tensor
submultiplicative over disjoint coordinate tensor factors.  Apply
\Cref{lem:V72-common-transport} to each bank.  Since
\Cref{cor:V71-chain-length} gives \(m\le3\), the factor \(C^m\) is
universal.  As in \Cref{lem:V71-statewise-tensorization}, sum the fixed
global source states only after this product estimate.
\end{proof}

\begin{firewall}[The global gapped factor is used once]
\label{fire:V72-one-gap-use}
If several selected banks lie in gapped sectors, their coordinate sets
are disjoint and their local exponentials multiply to the exponential
of the union count.  V72 neither extracts the same gap twice nor first
uses it to create a mark and later uses it as a terminal moment.
\end{firewall}

\section{The mixed scalar telescope}
\label{sec:V72-telescope}

\begin{lemma}[Sparse mixed ratio telescope]
\label{lem:V72-mixed-telescope}
Let \(q\in\{1,2,3\}\) and let the terminal row \(i_m\) have the moment
certificate \eqref{eq:V71-terminal-moments}.  Then
\begin{equation}
 \boxed{
 (s_\alpha\Xi_{i_0})^q
 \left(\prod_{\ell=0}^{m-1}R_\ell\gamma_\ell\right)
 \mathcal Q_{i_m}
 \le
 C_q\prod_{\ell=0}^{m-1}\gamma_\ell
 \le C_q.}
\label{eq:V72-mixed-telescope}
\end{equation}
\end{lemma}

\begin{proof}
The terminal certificate gives
\(\mathcal Q_{i_m}\le
C_q(s_\alpha\Xi_{i_m})^{-q}\).
The row ratios telescope once:
\[
 \prod_{\ell=0}^{m-1}R_\ell
 =
 \frac{\Xi_{i_m}}{\Xi_{i_0}}.
\]
Therefore the left side of
\eqref{eq:V72-mixed-telescope} is at most
\[
 C_q
 \left(\frac{\Xi_{i_0}}{\Xi_{i_m}}\right)^{q-1}
 \prod_\ell\gamma_\ell.
\]
The chain is strictly increasing, so the mass ratio is at most one,
and every \(\gamma_\ell\le1\).
\end{proof}

\begin{theorem}[Clean sparse mixed-chain closure]
\label{thm:V72-clean-mixed-chain}
Every branch of the exact assembled \(2+2\) source with:
\begin{enumerate}[label=\textup{(\roman*)}]
\item at most three unpaid row/private-output factors rooted at
      \(i_0\);
\item a clean sparse spectral chain;
\item a private or unbalanced admissible terminal; and
\item no transfer bank reused by the sole final payer,
\end{enumerate}
has a support-uniform scalar coefficient.  The bound is independent of
the number and positions of gapped banks in the chain.
\end{theorem}

\begin{proof}
Combine
\Cref{lem:V72-mixed-tensorization,
lem:V72-mixed-telescope}.
Private high-output terms are treated as in
\Cref{lem:V71-largest-root}: allocate the labelled private factor
inside the exact \(q_{\rm priv}\) product, so the total fixed moment
degree is at most three.  The selected balanced banks, private banks,
and unbalanced terminal banks are disjoint analytic resources.
\end{proof}

\section{The \(2+2\) DRQC consequence}
\label{sec:V72-DRQC}

\begin{theorem}[Closed clean sparse spectral part of the assembled
\(2+2\) source]
\label{thm:V72-spectral-22}
Choose the larger internal row of every ordered \(2+2\) path monomial.
If its balanced routing is a clean sparse spectral chain ending at a
private or unbalanced terminal, then its low-output and
private-high-output contributions satisfy
\begin{equation}
 \boxed{
 \frac{\kappa_\otimes(\mathcal S_{22}^{\rm branch})}
      {\Xi_1\Xi_2\Xi_3\Xi_4}
 \le
 C\,\theta_{12}\theta_{34}\theta_{ij}.}
\label{eq:V72-spectral-22}
\end{equation}
\end{theorem}

\begin{proof}
The normalized source coefficient before the retained transfer sectors
is bounded by
\[
 (s_\alpha\Xi_{i_0})^2
 \theta_{12}\theta_{34}\theta_{ij}
\]
on low output, by \Cref{lem:V71-largest-root}.  The private-output
part adds one labelled private factor but no fourth moment.
Apply \Cref{thm:V72-clean-mixed-chain}.  What remains is the original
tree monomial times
\(\prod_\ell\gamma_\ell\le1\).  Protected
\(\gamma_\ell\)''s contain genuine extra overlap marks; gapped
\(\gamma_\ell\)''s contain heat contractions.  Both may be discarded
only after the scalar telescope, leaving
\eqref{eq:V72-spectral-22}.
\end{proof}

\begin{corollary}[The inverse-scale gapped ledger is unnecessary on
sparse escalating chains]
\label{cor:V72-no-inverse-s}
On the branch of \Cref{thm:V72-spectral-22}, no factor
\(s_\alpha^{-1}\) from
\eqref{eq:V70-gapped-ratio} is used.  Hence a gapped transfer does not
raise the required terminal moment order.
\end{corollary}

\begin{proof}
The proof uses the exponential entry and the free inequality
\(1<R_\ell\) in \Cref{lem:V72-common-transport}.  The terminal
telescope remains of order two, or total order three after a private
output payer, regardless of the number of gapped steps.
\end{proof}

\section{Residual after sparse spectral closure}
\label{sec:V72-residual}

\begin{theorem}[Post-V72 balanced \(2+2\) residual]
\label{thm:V72-residual}
All clean sparse chains, with arbitrary protected/gapped sector choices
and a private or unbalanced terminal, are closed.  Every still-unproved
balanced \(2+2\) term therefore has at least one of:
\begin{enumerate}
\item a selected row-rooted bank with \(N\ge L\);
\item a selected bank whose large stock is concentrated on one of the
      at most three source marks;
\item two required selected banks sharing a coordinate;
\item a routing sequence that reaches no private/unbalanced terminal
      before one of items 1--3 occurs; or
\item the balanced nonprivate output ancestry \(Z_\nu^B\).
\end{enumerate}
\label{enum:V72-residual}
\end{theorem}

\begin{proof}
Start from the exhaustive V71 list
\eqref{enum:V71-residual}.  V72 removes every gapped-sector item when
the bank remains clean and sparse, so protected and gapped sector
choices no longer distinguish such a chain.  The hypotheses excluded
from the clean sparse theorem are exactly items 1--4 above.  Balanced
output ancestry was not used and remains item 5.
\end{proof}

\begin{assessment}[The dense branch is not merely large count]
\label{ass:V72-dense-warning}
For a dense bank, the count \(N\) may be large relative to the fixed
threshold \(L\) but small relative to \(s_\alpha^{-1}\).  Then
\(e^{-cs_\alpha N}\) need not be small.  Conversely, the row mass may
be concentrated in one arbitrarily large representation, so count
decay alone cannot pay it.  The next step must keep both the
row-rooted power inequality and the protected/gapped sector split;
declaring \(N\ge L\) to be a terminal moment would be false.
\end{assessment}

\section{V72 ledger and V73 mandate}
\label{sec:V72-ledger}

\begin{theorem}[Integrated compact-series V72 statement]
\label{thm:V72-integrated}
Preserving compact V1--V71, V72 proves:
\begin{enumerate}[label=\textup{(V72.\arabic*)},leftmargin=6.7em]
\item protected and gapped sectors of a sparse escalating bank share
      the common ratio form \(R_\ell\gamma_\ell\);
\item mixed sector factors tensorize statewise over clean disjoint
      banks;
\item the exact terminal mass-ratio telescope needs no
      \(s_\alpha^{-1}\) extraction and no higher moment degree;
\item every clean sparse mixed spectral chain satisfies the assembled
      \(2+2\) DRQC bound; and
\item the balanced residual is reduced to the five classes in
      \eqref{enum:V72-residual}.
\end{enumerate}
\end{theorem}

\begin{proof}
These are
\Cref{lem:V72-common-transport,
lem:V72-mixed-tensorization,
lem:V72-mixed-telescope,
thm:V72-spectral-22,
thm:V72-residual}.
\end{proof}

\begin{openproblem}[V73 mandate: dense-bank stopping rule]
\label{op:V73-mandate}
\begin{enumerate}
\item Fix the first selected bank with \(N\ge L\) and keep its exact
      protected/gapped spectral split.
\item In the protected sector, combine the \(\rho^N\) capacity with the
      actual row-mass ratio before deciding whether to stop or
      transport.
\item In the gapped sector, combine \(e^{-cs_\alpha N}\) with
      \(S^2\le CNH\); do not treat count as Casimir mass.
\item Split according to the dimensionless variables
      \(s_\alpha N\) and \(\Xi_v/\Xi_u\), and prove a stopping moment or
      a new strictly increasing transfer.
\item Test the one-coordinate/high-spin limit separately; it is sparse
      in count and belongs to the marked-heavy or clean branch, not to
      a dense-count argument.
\item Keep marked-heavy, coordinate-overlap, and balanced-output
      branches outside the dense proof unless their resources are
      explicitly retained.
\end{enumerate}
\end{openproblem}

\begin{firewall}[Claim boundary after compact V72]
\label{fire:V72-current-boundary}
V72 closes clean sparse mixed protected/gapped chains only.  It does
not close the five residual classes in \eqref{enum:V72-residual}, the
other three first-connectivity families, global quartic and higher MTA,
WUFC, CPM, continuum Yang--Mills existence, or a positive mass gap.
The full Yang--Mills mass gap remains open.
\end{firewall}

\section{Append-only record for compact V72}
\label{sec:V72-record}

V72 was appended to the complete compact V71 file.  The exact V71
parent had \(29504\) newline-delimited source records, \(1018595\)
bytes, and SHA--256 digest
\[
\texttt{2f3e051f78675f7bc61fcec6ce7f1d4fb677f1977937bcda31c3d1a59778580e}.
\]
No compact V71 mathematical content was changed.  The sole final
\verb|\end{document}| is restored below.

% =============================================================================
% END APPENDED COMPACT VERSION V72 MATERIAL
% =============================================================================

\clearpage

% =============================================================================
% BEGIN APPENDED COMPACT VERSION V73 MATERIAL
% =============================================================================

\chapter{Diffuse protected terminals and removal of the dense
protected obstruction}
\label{chap:V73}

\section{Context and the correct dense split}
\label{sec:V73-context}

V72 closed every clean sparse spectral chain.  A selected bank with
large coordinate count was left open because count alone is not
Casimir mass.  In the all-protected sector, however, one has more than
count decay: every balanced coordinate exposes one inverse
representation dimension across the row cut.

The all-weak condition forces a non-escalating balanced bank to be
diffuse.  No coordinate can carry a fixed fraction of the root row
mass, since its comparable partner mass would create a large normalized
overlap.  Diffuseness allows the coordinates to be divided into six
mass groups.  The product of the protected one-row capacities over
each pair of groups pays one full row mass.  Six groups therefore pay
every moment through degree three.  This turns a non-escalating
all-protected bank into a terminal moment certificate.  If the partner
mass does escalate, the usual ratio telescope continues without any
sparsity hypothesis.

\section{A scalar product-capacity lemma}
\label{sec:V73-scalar-capacity}

For a nontrivial \(SU(2)\) spin \(j\), write
\[
 a(j):=1+j(j+1),
 \qquad
 d(j):=2j+1.
\]
Then
\begin{equation}
 d(j)^2=4a(j)-3\ge a(j),
\qquad
 a(j)\ge a_*:=\frac74.
\label{eq:V73-dimension-mass}
\end{equation}

\begin{lemma}[One group supplies one square-root moment]
\label{lem:V73-one-group}
For every nonempty finite set \(G\) of nontrivial spins,
\begin{equation}
 \boxed{
 \prod_{u\in G}\frac1{d(j_u)}
 \le
 C_*
 \left(\sum_{u\in G}a(j_u)\right)^{-1/2}.}
\label{eq:V73-one-group}
\end{equation}
The constant is independent of \(|G|\) and of the spins.
\end{lemma}

\begin{proof}
By \eqref{eq:V73-dimension-mass},
\[
 \prod_{u\in G}\frac1{d(j_u)}
 \le
 \left(\prod_{u\in G}a(j_u)\right)^{-1/2}.
\]
If \(n=|G|\), then
\[
 \frac{\sum_{u\in G}a(j_u)}
      {\prod_{u\in G}a(j_u)}
 =
 \sum_{u\in G}
 \frac1{\prod_{v\in G\setminus\{u\}}a(j_v)}
 \le
 n a_*^{-(n-1)}.
\]
The last expression has a finite supremum over \(n\ge1\).  Taking
square roots proves \eqref{eq:V73-one-group}.
\end{proof}

\begin{lemma}[Diffuse weights admit \(2q\) large groups]
\label{lem:V73-grouping}
Let \(x_1,\ldots,x_N>0\), put \(S=\sum_nx_n\), and fix an integer
\(q\ge1\).  If
\begin{equation}
 \max_nx_n\le\frac{S}{4q},
\label{eq:V73-diffuse}
\end{equation}
then the indices can be partitioned into \(2q\) nonempty groups
\(G_1,\ldots,G_{2q}\) such that
\begin{equation}
 \sum_{n\in G_t}x_n\ge\frac{S}{4q}
 \qquad(1\le t\le2q).
\label{eq:V73-group-mass}
\end{equation}
\end{lemma}

\begin{proof}
Read the indices in any order.  Form \(G_1\) by stopping the first time
its accumulated mass reaches \(S/(4q)\), then repeat for
\(G_2,\ldots,G_{2q-1}\).  By
\eqref{eq:V73-diffuse}, every completed group has mass at most
\(S/(2q)\).  Hence the first \(2q-1\) groups consume at most
\[
 \frac{2q-1}{2q}S,
\]
leaving at least \(S/(2q)\) for the last group.  This proves
\eqref{eq:V73-group-mass}.
\end{proof}

\begin{theorem}[Diffuse protected moments through degree three]
\label{thm:V73-diffuse-protected-moments}
Let \(D_{uv}\) be a row-rooted balanced bank and let
\(\mathcal P_{uv}(D)\) be its all-protected sector.  If
\begin{equation}
 \max_{w\in D}a_{u,w}\le\frac{S_{uv}(D)}{12},
\label{eq:V73-six-group-diffuse}
\end{equation}
then, for \(q=1,2,3\),
\begin{equation}
 \boxed{
 S_{uv}(D)^q\,
 \kappa_{\otimes,4}\bigl(\mathcal P_{uv}(D)\bigr)
 \le C_q.}
\label{eq:V73-protected-moments}
\end{equation}
Consequently the same bounds hold with
\((s_\alpha S_{uv}(D))^q\) in place of \(S_{uv}(D)^q\).
\end{theorem}

\begin{proof}
At each protected coordinate the invariant block crossing
\(u|u^c\) has product-state capacity at most
\(d(j_{u,w})^{-1}\), as in
\Cref{lem:V70-local-PG}.  Tensorization over distinct coordinates
gives
\[
 \kappa_{\otimes,4}(\mathcal P_{uv}(D))
 \le
 \prod_{w\in D}d(j_{u,w})^{-1}.
\]

For a fixed \(q\le3\), condition
\eqref{eq:V73-six-group-diffuse} implies
\eqref{eq:V73-diffuse}.  Partition \(D\) into \(2q\) groups using
\Cref{lem:V73-grouping}.  Apply
\Cref{lem:V73-one-group} to each group:
\[
 \prod_{w\in D}d(j_{u,w})^{-1}
 \le
 C_*^{2q}
 \prod_{t=1}^{2q}
 \left(\sum_{w\in G_t}a_{u,w}\right)^{-1/2}
 \le
 C_q S_{uv}(D)^{-q}.
\]
This proves \eqref{eq:V73-protected-moments}.  The scaled assertion
uses \(0<s_\alpha\le1\).
\end{proof}

\section{All-weak non-escalation forces diffuseness}
\label{sec:V73-forced-diffuse}

\begin{lemma}[A non-escalating large partner bank is diffuse]
\label{lem:V73-nonescalating-diffuse}
Let \(D_{uv}\) satisfy
\begin{equation}
 S_{uv}(D)\ge c_0\Xi_u
\label{eq:V73-large-D}
\end{equation}
and the all-weak bound
\(\theta_{uv}\le\delta\).  If
\begin{equation}
 \Xi_v\le2\Xi_u,
\label{eq:V73-nonescalating}
\end{equation}
then
\begin{equation}
 \max_{w\in D}a_{u,w}
 \le
 C_{c_0}\sqrt{\delta}\,S_{uv}(D).
\label{eq:V73-forced-diffuse}
\end{equation}
In particular, after fixing
\begin{equation}
 \delta\le\delta_*(c_0)
 :=
 \frac1{144C_{c_0}^2},
\label{eq:V73-delta-star}
\end{equation}
condition \eqref{eq:V73-six-group-diffuse} holds.
\end{lemma}

\begin{proof}
The maximal-partner comparison in
\Cref{lem:V70-row-power} gives
\[
 a_{u,w}^2\le C a_{u,w}a_{v,w}
\]
at every \(w\in D\).  Hence
\[
 \max_{w\in D}a_{u,w}^2
 \le
 C H_{uv}^{(u)}(D)
 \le
 C H_{uv}
 \le
 C\delta\Xi_u\Xi_v
 \le
 2C\delta\Xi_u^2.
\]
Use \eqref{eq:V73-large-D} to replace
\(\Xi_u\) by \(c_0^{-1}S_{uv}(D)\), then take square roots.  Enlarge
\(C_{c_0}\) to absorb \(2C\).  The choice
\eqref{eq:V73-delta-star} gives the final assertion.
\end{proof}

\begin{corollary}[A non-escalating all-protected bank is a terminal]
\label{cor:V73-protected-terminal}
Under
\eqref{eq:V73-large-D},
\eqref{eq:V73-nonescalating}, and
\eqref{eq:V73-delta-star}, the retained all-protected sector satisfies
\begin{equation}
 (s_\alpha\Xi_u)^q
 \kappa_{\otimes,4}\bigl(\mathcal P_{uv}(D)\bigr)
 \le C_{q,c_0},
 \qquad q=1,2,3.
\label{eq:V73-terminal-card}
\end{equation}
\end{corollary}

\begin{proof}
By \eqref{eq:V73-large-D},
\(\Xi_u\le c_0^{-1}S_{uv}(D)\).  Apply
\Cref{lem:V73-nonescalating-diffuse,
thm:V73-diffuse-protected-moments}.
\end{proof}

\begin{remark}[Dense count is a consequence, not the payer]
\label{rem:V73-count-consequence}
Diffuseness implies at least twelve active coordinates, so this branch
is dense in a fixed sense.  The proof does not use
\(e^{-cs_\alpha N}\) to pay Casimir mass.  It uses the product of the
actual inverse representation dimensions and the all-weak exclusion
of a dominant balanced coordinate.
\end{remark}

\section{Clean protected routing without a count threshold}
\label{sec:V73-routing}

\begin{definition}[Clean protected step]
\label{def:V73-protected-step}
A clean protected step \(u\to v\) has the same unmarked,
coordinate-disjoint, large-bank hypotheses as V71, lies in the
all-protected sector, and has no restriction on \(N_{uv}\).
If \(\Xi_v>2\Xi_u\), call it escalating; otherwise call it terminal.
\end{definition}

\begin{lemma}[Every clean protected step either transports or pays]
\label{lem:V73-transport-or-terminal}
For \(\delta\le\delta_*(c_0)\), an escalating clean protected step has
\begin{equation}
 \kappa_\otimes(\mathcal P_{uv})
 \le
 \frac{C}{N_{uv}}\,
 \vartheta_{uv}\frac{\Xi_v}{\Xi_u},
\label{eq:V73-escalating-factor}
\end{equation}
and a terminal clean protected step has the moment certificate
\eqref{eq:V73-terminal-card} through degree three.
\end{lemma}

\begin{proof}
The first assertion is
\eqref{eq:V70-protected-ratio}; it never required sparsity.
The second is \Cref{cor:V73-protected-terminal}.
\end{proof}

\begin{theorem}[Clean all-protected chain closure without a density
exception]
\label{thm:V73-all-protected-chain}
Start from at most three unpaid factors rooted at one row.  A clean
coordinate-disjoint all-protected routing sequence closes if at every
step the V70 large partner bank is available.  No sparse/dense count
split is required.
\end{theorem}

\begin{proof}
At a step \(u\to v\), use
\Cref{lem:V73-transport-or-terminal}.  If it is terminal, apply its
degree-\(q\le3\) moment card.  If it escalates, retain the factor
\eqref{eq:V73-escalating-factor} and continue from \(v\).  Escalating
steps more than double row mass, so they cannot revisit a vertex and
there are at most three.  If no private or unbalanced terminal occurs,
the first non-escalating protected step supplies the terminal card.

Multiply coordinate-disjoint capacities statewise as in
\Cref{lem:V71-statewise-tensorization}.  The exact ratio telescope
\eqref{eq:V71-ratio-telescope} applies to all escalating steps; the
last terminal card pays the remaining moment.  Every factor
\(\vartheta/N\) is at most one.
\end{proof}

\begin{theorem}[Dense all-protected \(2+2\) branch is closed]
\label{thm:V73-dense-protected-22}
In the exact assembled \(2+2\) source, every low-output or
private-high-output branch whose selected balanced banks are clean,
coordinate-disjoint, and all-protected satisfies its original DRQC
tree bound, regardless of the bank counts.
\end{theorem}

\begin{proof}
Choose the largest internal endpoint as in
\Cref{lem:V71-largest-root}.  Apply
\Cref{thm:V73-all-protected-chain}; the private output factor raises
the fixed total degree only to three.  Once the scalar coefficient is
paid, delete the additional physical transfer contractions exactly as
in \Cref{thm:V71-protected-22}.  The original
\(\theta_{12}\theta_{34}\theta_{ij}\) remains.
\end{proof}

\section{Refined mixed-chain consequence}
\label{sec:V73-mixed}

\begin{corollary}[Only a non-escalating gapped bank can stop a clean
mixed chain]
\label{cor:V73-only-gapped-stop}
Consider a clean coordinate-disjoint mixed protected/gapped routing
sequence.  Every selected bank with partner mass
\(\Xi_v>2\Xi_u\) is a valid transport step: use
\eqref{eq:V73-escalating-factor} in the protected sector and
\Cref{lem:V72-common-transport} in the gapped sector.  If
\(\Xi_v\le2\Xi_u\), an all-protected sector is a terminal by V73.
Therefore the only clean bank not yet controlled is a
non-escalating gapped complementary sector.
\end{corollary}

\begin{proof}
The escalating protected and gapped assertions are the cited ratio
forms and use no count restriction once the mass inequality is assumed.
The non-escalating protected assertion is
\Cref{cor:V73-protected-terminal}.  These cases are exhaustive for the
two spectral sectors.
\end{proof}

\begin{firewall}[No claim about a low-spin gapped projection]
\label{fire:V73-gapped-remains}
A non-escalating diffuse bank may contain large balanced inputs fusing
to small nonzero outputs.  Its heat multiplier then has only the
uniform \(1-cs_\alpha\) gap, and the present proof has not shown that
its physical projection retains the inverse-dimension products used in
the protected sector.  Hence the gapped terminal is not inferred from
the protected terminal.
\end{firewall}

\section{Post-V73 residual and next acquisition}
\label{sec:V73-residual}

\begin{theorem}[Post-V73 balanced \(2+2\) residual]
\label{thm:V73-residual}
The dense all-protected exception in
\eqref{enum:V72-residual} is closed.  Every remaining balanced
\(2+2\) term has at least one of:
\begin{enumerate}
\item a clean non-escalating gapped balanced bank;
\item a large balanced bank concentrated on a source mark;
\item two required balanced banks sharing a coordinate;
\item failure to reach a private, unbalanced, or diffuse protected
      terminal because one of items 1--3 occurs; or
\item balanced nonprivate output ancestry \(Z_\nu^B\).
\end{enumerate}
\label{enum:V73-residual}
\end{theorem}

\begin{proof}
\Cref{thm:V73-dense-protected-22} removes every all-protected bank
irrespective of count.  By
\Cref{cor:V73-only-gapped-stop}, the only clean spectral stopping case
left is a non-escalating gapped bank.  Mark concentration,
coordinate overlap, and balanced output ancestry were excluded from
the clean-bank theorem and remain exactly as listed.
\end{proof}

\begin{assessment}[The next local normal form]
\label{ass:V73-next-normal-form}
At a non-escalating gapped coordinate, let \(j\) be the row-root spin
and \(J\ne0\) the output spin in its block.  The needed estimate is a
capacity/heat bound for the complete \(J\)-isotypic projection,
schematically
\[
 \kappa_\otimes(P_J)\,e^{-s_\alpha C_2(J)}
 \le
 \frac{C}{d_j}\,
 \Phi\!\left(s_\alpha C_2(J),\frac{d_J}{d_j}\right),
\]
with a summable right side after the physical \(J\)-sum.  Small \(J\)
must recover a one-row inverse dimension; large \(J\) may use heat
decay.  The statement must be proved for the whole isotypic projector,
including multiplicity, before it is tensorized over the diffuse bank.
\end{assessment}

\section{V73 ledger and V74 mandate}
\label{sec:V73-ledger}

\begin{theorem}[Integrated compact-series V73 statement]
\label{thm:V73-integrated}
Preserving compact V1--V72, V73 proves:
\begin{enumerate}[label=\textup{(V73.\arabic*)},leftmargin=6.7em]
\item a product of protected one-row capacities pays one square-root
      moment on every coordinate group;
\item a diffuse bank can be split into six groups and hence pays every
      moment through degree three;
\item all-weak non-escalation forces the selected balanced bank to be
      diffuse;
\item every clean all-protected routing chain closes without a
      sparse/dense split;
\item the dense all-protected \(2+2\) branch satisfies DRQC; and
\item the only clean spectral stopping obstruction is a
      non-escalating gapped bank.
\end{enumerate}
\end{theorem}

\begin{proof}
These are
\Cref{lem:V73-one-group,
lem:V73-grouping,
thm:V73-diffuse-protected-moments,
lem:V73-nonescalating-diffuse,
thm:V73-all-protected-chain,
thm:V73-dense-protected-22,
cor:V73-only-gapped-stop}.
\end{proof}

\begin{openproblem}[V74 mandate: gapped isotypic capacity]
\label{op:V74-mandate}
\begin{enumerate}
\item For a fixed balanced coordinate and row cut, decompose the
      complete nontrivial-output block into \(SU(2)\) output
      isotypic projectors \(P_J\), retaining multiplicity spaces.
\item Prove the exact product-state capacity of \(P_J\), or a sharp
      upper bound uniform in input spins and multiplicity.
\item Multiply that capacity by the Wilson heat factor
      \(e^{-s_\alpha C_2(J)}\) before summing \(J\).
\item Show that the resulting complete gapped block supplies one
      inverse root-row dimension uniformly, or construct the minimal
      recoupling counterexample.
\item If the one-coordinate estimate succeeds, tensor it over the
      diffuse non-escalating bank and repeat the six-group moment
      argument.
\item Keep marked-heavy, overlapping-bank, and balanced-output cases
      outside this local calculation.
\end{enumerate}
\end{openproblem}

\begin{firewall}[Claim boundary after compact V73]
\label{fire:V73-current-boundary}
V73 closes dense and sparse clean all-protected routing, but not the
gapped isotypic-capacity gate or the other four residual classes in
\eqref{enum:V73-residual}.  The other three first-connectivity
families, global quartic and higher MTA, WUFC, CPM, continuum
Yang--Mills existence, and a positive mass gap remain open.  The full
Yang--Mills mass gap remains open.
\end{firewall}

\section{Append-only record for compact V73}
\label{sec:V73-record}

V73 was appended to the complete compact V72 file.  The exact V72
parent had \(29897\) newline-delimited source records, \(1032389\)
bytes, and SHA--256 digest
\[
\texttt{1e9d582281b45d20ed570d7772f3abc7dac6b201d5ec9f9880f08ae90ff3d529}.
\]
No compact V72 mathematical content was changed.  The sole final
\verb|\end{document}| is restored below.

% =============================================================================
% END APPENDED COMPACT VERSION V73 MATERIAL
% =============================================================================

\clearpage

% =============================================================================
% BEGIN APPENDED COMPACT VERSION V74 MATERIAL
% =============================================================================

\chapter{Anti-aligned coherent-state regression for the gapped
terminal}
\label{chap:V74}

\section{Context: test the proposed local normal form first}
\label{sec:V74-context}

V73 reduced the clean balanced obstruction to a non-escalating gapped
bank.  Its mandate proposed proving that the complete nontrivial-output
isotypic block retains one inverse root-row dimension after heat
weighting.  This note tests that assertion in the smallest possible
model: two equal \(SU(2)\) spins at one coordinate.

The proposed bound is false.  Two anti-aligned spin-coherent product
vectors have total Casimir expectation \(2j\), rather than order
\(j^2\).  At heat time \(s_\alpha\), choosing
\(j\asymp s_\alpha^{-1}\) therefore leaves an order-one expectation of
the nontrivial-output heat block, while the input Casimir mass is order
\(s_\alpha^{-2}\).  Repeating a fixed number of such coordinates makes
the normalized overlap arbitrarily weak and the bank diffuse, but
still does not create a scaled row-moment tail.

This is a regression of a positive statewise route, not a
counterexample to the complete first-connectivity source.  It forces
the programme upstream: the gapped spectator bank must remain
correlated with the signed binary-prefix and joining covariance state
sum.

\section{The exact two-spin heat block}
\label{sec:V74-two-spin}

Let \(V_j\) be the spin-\(j\) irreducible representation and
\[
 V_j\otimes V_j
 =
 \bigoplus_{J=0}^{2j}V_J.
\]
Write \(P_J\) for the orthogonal projection onto \(V_J\).  In the exact
heat-kernel model define
\begin{align}
 \mathsf T_{s,j}
 &:=
 \sum_{J=0}^{2j}e^{-sJ(J+1)}P_J,
\label{eq:V74-full-heat}\\
 \mathsf G_{s,j}
 &:=
 \sum_{J=1}^{2j}e^{-sJ(J+1)}P_J
 =
 \mathsf T_{s,j}-P_0.
\label{eq:V74-gapped-heat}
\end{align}
The second sum means all nontrivial \(J\); for half-integral notation
the lower limit is interpreted as the smallest positive admissible
output.

Let
\[
 \psi_j:=|j,j\rangle\otimes|j,-j\rangle.
\]
This is a product vector across the two input rows.

\begin{lemma}[Anti-aligned total Casimir]
\label{lem:V74-mean-Casimir}
If \(\mathbf J=\mathbf J_1+\mathbf J_2\) is the total angular momentum,
then
\begin{equation}
 \boxed{
 \langle\psi_j,\mathbf J^2\psi_j\rangle=2j.}
\label{eq:V74-mean-Casimir}
\end{equation}
\end{lemma}

\begin{proof}
The two factors have
\[
 \langle\mathbf J_1\rangle=(0,0,j),
 \qquad
 \langle\mathbf J_2\rangle=(0,0,-j).
\]
Because \(\psi_j\) is a product vector,
\[
 \langle\mathbf J_1\cdot\mathbf J_2\rangle
 =
 \langle\mathbf J_1\rangle\cdot
 \langle\mathbf J_2\rangle
 =-j^2.
\]
Using
\(\mathbf J^2=\mathbf J_1^2+\mathbf J_2^2+
2\mathbf J_1\cdot\mathbf J_2\) gives
\[
 2j(j+1)-2j^2=2j.
\]
\end{proof}

\begin{lemma}[Exact singlet weight]
\label{lem:V74-singlet-weight}
The spin-zero probability in \(\psi_j\) is
\begin{equation}
 \boxed{
 \langle\psi_j,P_0\psi_j\rangle=\frac1{2j+1}.}
\label{eq:V74-singlet-weight}
\end{equation}
\end{lemma}

\begin{proof}
The normalized invariant vector in
\(V_j\otimes V_j\) is
\[
 \Omega_j
 =
 \frac1{\sqrt{2j+1}}
 \sum_{m=-j}^{j}(-1)^{j-m}
 |j,m\rangle\otimes|j,-m\rangle.
\]
Its coefficient on
\(|j,j\rangle\otimes|j,-j\rangle\) has modulus
\((2j+1)^{-1/2}\).  Squaring proves
\eqref{eq:V74-singlet-weight}.
\end{proof}

\begin{theorem}[Order-one gapped expectation at a hot input]
\label{thm:V74-order-one-gapped}
There are universal \(s_0,c_0>0\) such that, for every
\(0<s\le s_0\), one may choose a half-integer \(j_s\) satisfying
\begin{equation}
 \frac1{2s}\le j_s\le\frac1s
\label{eq:V74-j-choice}
\end{equation}
and
\begin{equation}
 \boxed{
 \langle\psi_{j_s},
 \mathsf G_{s,j_s}\psi_{j_s}\rangle
 \ge c_0.}
\label{eq:V74-order-one-gapped}
\end{equation}
At the same time,
\begin{equation}
 a_s:=1+j_s(j_s+1)\asymp s^{-2},
 \qquad
 sa_s\asymp s^{-1}\longrightarrow\infty.
\label{eq:V74-hot-input}
\end{equation}
\end{theorem}

\begin{proof}
Let \(p_J=\langle\psi_j,P_J\psi_j\rangle\).  These numbers form a
probability distribution, and
\Cref{lem:V74-mean-Casimir} gives
\[
 \sum_Jp_JJ(J+1)=2j.
\]
The function \(x\mapsto e^{-sx}\) is convex.  Jensen''s inequality
therefore yields
\[
 \langle\psi_j,\mathsf T_{s,j}\psi_j\rangle
 =
 \sum_Jp_Je^{-sJ(J+1)}
 \ge e^{-2sj}.
\]
Subtract the singlet term using
\Cref{lem:V74-singlet-weight}:
\begin{equation}
 \langle\psi_j,\mathsf G_{s,j}\psi_j\rangle
 \ge
 e^{-2sj}-\frac1{2j+1}.
\label{eq:V74-lower-prelimit}
\end{equation}
Choose any half-integer in the interval
\eqref{eq:V74-j-choice}, which is nonempty for small \(s\).
Then \(e^{-2sj_s}\ge e^{-2}\), while
\((2j_s+1)^{-1}\to0\).  After decreasing \(s_0\), the right side of
\eqref{eq:V74-lower-prelimit} is at least
\(c_0=e^{-2}/2\).  The estimates
\eqref{eq:V74-hot-input} follow directly from
\eqref{eq:V74-j-choice}.
\end{proof}

\begin{corollary}[The desired inverse-dimension gapped card is false]
\label{cor:V74-no-inverse-dimension}
There is no universal \(C\) such that
\begin{equation}
 \kappa_\otimes(\mathsf G_{s,j})
 \le\frac{C}{2j+1}
\label{eq:V74-false-dimension-card}
\end{equation}
for all small \(s\) and all spins \(j\).  More generally, no upper
bound tending to zero as a function of \(sa(j)\) alone can hold for
the complete gapped block.
\end{corollary}

\begin{proof}
Product-state capacity is at least the expectation on the product
vector \(\psi_{j_s}\).  By
\Cref{thm:V74-order-one-gapped}, that expectation is at least \(c_0\),
whereas \((2j_s+1)^{-1}\to0\) and
\(sa(j_s)\to\infty\).
\end{proof}

\section{A diffuse all-weak bank countertest}
\label{sec:V74-bank-test}

Fix an integer
\begin{equation}
 N\ge\max\{12,\lceil\delta^{-1}\rceil\}.
\label{eq:V74-N-choice}
\end{equation}
Use \(N\) distinct coordinates, with the two rows carrying spin
\(j_s\) at every coordinate.  Let the row-\(u\) vector be the tensor
product of all highest-weight vectors and the row-\(v\) vector the
tensor product of all lowest-weight vectors.  This remains one product
vector across the two rows.

\begin{proposition}[Diffuse weak overlap does not give a gapped
terminal moment]
\label{prop:V74-diffuse-gapped-fails}
For the bank just described,
\begin{align}
 \Xi_u=\Xi_v&=Na_s,
\label{eq:V74-bank-masses}\\
 H_{uv}&=Na_s^2,
\label{eq:V74-bank-overlap}\\
 \theta_{uv}&=\frac1N\le\delta,
\label{eq:V74-bank-theta}\\
 \frac{\max_w a_{u,w}}{\sum_w a_{u,w}}
 &=\frac1N\le\frac1{12}.
\label{eq:V74-bank-diffuse}
\end{align}
Nevertheless the sector in which all \(N\) coordinates have
nontrivial output obeys
\begin{equation}
 \kappa_\otimes
 \left(\mathsf G_{s,j_s}^{\otimes N}\right)
 \ge c_0^N,
\label{eq:V74-bank-lower}
\end{equation}
and hence, for every fixed \(q\ge1\),
\begin{equation}
 \boxed{
 (s\Xi_u)^q
 \kappa_\otimes
 \left(\mathsf G_{s,j_s}^{\otimes N}\right)
 \longrightarrow\infty
 \quad(s\downarrow0).}
\label{eq:V74-no-gapped-moment}
\end{equation}
\end{proposition}

\begin{proof}
The mass and overlap identities follow by summing \(N\) identical
coordinate contributions.  Their quotient is \(1/N\), and
\eqref{eq:V74-N-choice} gives
\eqref{eq:V74-bank-theta}--%
\eqref{eq:V74-bank-diffuse}.

Evaluate the tensor product on the tensor product of the \(N\) vectors
\(\psi_{j_s}\).  Expectations factor across coordinates, so
\Cref{thm:V74-order-one-gapped} gives
\eqref{eq:V74-bank-lower}.  Finally
\[
 s\Xi_u=sNa_s\asymp\frac{N}{s}\longrightarrow\infty,
\]
while \(c_0^N>0\) is fixed once \(\delta\) is fixed.  This proves
\eqref{eq:V74-no-gapped-moment}.
\end{proof}

\begin{firewall}[What the bank countertest disproves]
\label{fire:V74-scope}
\Cref{prop:V74-diffuse-gapped-fails} disproves a statewise positive
terminal-moment estimate for the complete gapped spectator bank, even
under balanced diffuseness, non-escalation, and arbitrarily weak
normalized overlap.  It does not show that the complete signed
first-connectivity source is large on the same vectors.  The binary
prefix and joining covariance state sum may cancel precisely the
anti-aligned sector retained by the positive envelope.
\end{firewall}

\section{Why heat and dimension cannot be optimized separately}
\label{sec:V74-mechanism}

\begin{assessment}[The anti-aligned window]
\label{ass:V74-window}
The obstruction uses three scales:
\[
 j_s\asymp s^{-1},
 \qquad
 C_2(\hbox{input})\asymp s^{-2},
 \qquad
 C_2(\hbox{typical total output})\asymp j_s\asymp s^{-1}.
\]
Thus the input row is strongly hot, but the coupled output lies at the
parabolic heat scale and is not exponentially suppressed.  The
singlet sector has only probability \(O(s)\), so removing it does not
change the order-one gapped expectation.  An estimate based only on
input dimension controls the singlet but misses the broad
\(J\asymp\sqrt j\) band; an estimate based only on output heat regards
that band as order one.
\end{assessment}

\begin{proposition}[The V73 protected argument cannot be transferred
to the gapped block]
\label{prop:V74-no-V73-transfer}
The six-group proof of
\Cref{thm:V73-diffuse-protected-moments} has no analogue obtained by
replacing each protected factor \(d_j^{-1}\) with the complete gapped
heat block.  The failure already occurs for a fixed number
\(N\ge12\) of identical diffuse coordinates.
\end{proposition}

\begin{proof}
The protected proof obtains one factor \(d_j^{-1}\) at every
coordinate and tensors those factors before grouping.  If a gapped
replacement supplied the same conclusion, its degree-one consequence
would bound
\[
 s\Xi_u\,
 \kappa_\otimes(\mathsf G_{s,j_s}^{\otimes N})
\]
uniformly.  This is the case \(q=1\) contradicted by
\eqref{eq:V74-no-gapped-moment}.
\end{proof}

\section{The forced upstream carrier}
\label{sec:V74-upstream}

\begin{theorem}[Positive statewise routing cannot close the gapped
residual]
\label{thm:V74-positive-route-fails}
Any proof of the remaining non-escalating gapped \(2+2\) branch which
uses only:
\begin{enumerate}[label=\textup{(\roman*)}]
\item row masses and overlap stocks;
\item the all-weak inequality;
\item a positive envelope in each fixed global moment-partition state;
      and
\item the complete nontrivial-output heat block at each spectator
      coordinate
\end{enumerate}
cannot produce the required uniform terminal moment.  Additional
signed correlation with the exact source state sum is necessary.
\end{theorem}

\begin{proof}
The bank in
\Cref{prop:V74-diffuse-gapped-fails} satisfies the scalar and
all-weak hypotheses, is balanced and diffuse, and lies wholly in the
nontrivial-output sector used by the proposed statewise envelope.  Its
product-state expectation violates every positive terminal moment.
Therefore no implication from items (i)--(iv) to that moment can hold.
\end{proof}

\begin{assessment}[New acquisition direction]
\label{ass:V74-new-direction}
Return to the exact packet
\[
 K_{12}^{<r,\mathrm{conn}}
 \otimes K_{34}^{<r,\mathrm{conn}}
 \otimes\operatorname{Cov}_r(Z_{12,r},Z_{34,r})
 \otimes M_{>r}.
\]
At a candidate gapped spectator bank, do not first fix and norm a
global partition state.  Instead compress the complete signed binary
prefix and joining covariance functional to the anti-aligned
\(J\asymp\sqrt j\) output band.  The next question is whether the
four covariance signs annihilate its leading semiclassical product
symbol.  If they do, the remainder may gain a derivative or inverse
dimension.  If they do not, the resulting two-prefix packet is the
minimal counterpacket to the present DRQC formulation.
\end{assessment}

\section{V74 ledger and V75 mandate}
\label{sec:V74-ledger}

\begin{theorem}[Integrated compact-series V74 statement]
\label{thm:V74-integrated}
Preserving compact V1--V73, V74 proves:
\begin{enumerate}[label=\textup{(V74.\arabic*)},leftmargin=6.7em]
\item an anti-aligned spin-\(j\) product vector has total Casimir mean
      \(2j\) and singlet probability \((2j+1)^{-1}\);
\item for \(j\asymp s_\alpha^{-1}\), the complete gapped heat block
      has order-one product-state expectation although the input is
      strongly hot;
\item a fixed diffuse all-weak bank of such coordinates has no scaled
      gapped terminal moment of any positive order;
\item the proposed inverse-dimension gapped isotypic card and a direct
      transfer of the V73 six-group proof are false; and
\item the remaining branch must be returned to the complete signed
      first-connectivity source state sum.
\end{enumerate}
\end{theorem}

\begin{proof}
These are
\Cref{lem:V74-mean-Casimir,
lem:V74-singlet-weight,
thm:V74-order-one-gapped,
prop:V74-diffuse-gapped-fails,
cor:V74-no-inverse-dimension,
prop:V74-no-V73-transfer,
thm:V74-positive-route-fails}.
\end{proof}

\begin{openproblem}[V75 mandate: audit and signed gapped covariance
symbol]
\label{op:V75-mandate}
\begin{enumerate}
\item Perform the mandatory five-version audit of the newest SM and NS
      files, recording fingerprints and importing proof order only.
\item Write the exact four-term double-replica formula for the two
      binary prefixes together with the assembled joining covariance,
      before any statewise norm.
\item Compress that complete signed packet to the anti-aligned
      \(J\asymp\sqrt j\) spectator output band of V74.
\item Compute its leading coherent-state symbol and test whether the
      covariance signs cancel it.
\item If the leading symbol cancels, quantify the first surviving
      derivative/inverse-dimension factor uniformly over the diffuse
      bank.
\item If it does not cancel, preserve the smallest explicit
      two-prefix/one-joining \(SU(2)\) packet as a counterexample to
      the current DRQC source bound and revise the target upstream.
\end{enumerate}
\end{openproblem}

\begin{firewall}[Claim boundary after compact V74]
\label{fire:V74-current-boundary}
V74 disproves a proposed positive gapped-terminal route but does not
disprove the complete \(2+2\) first-connectivity source bound.  The
signed gapped covariance packet, marked-heavy and overlapping-bank
branches, balanced output ancestry, the other three source families,
global quartic and higher MTA, WUFC, CPM, continuum Yang--Mills
existence, and a positive mass gap remain open.  The full Yang--Mills
mass gap remains open.
\end{firewall}

\section{Append-only record for compact V74}
\label{sec:V74-record}

V74 was appended to the complete compact V73 file.  The exact V73
parent had \(30380\) newline-delimited source records, \(1048518\)
bytes, and SHA--256 digest
\[
\texttt{cdd8f5baf1ab882cf76fe84ad12bc5f41b7ea52bfe6e91faa7d09e22ee63445c}.
\]
No compact V73 mathematical content was changed.  The sole final
\verb|\end{document}| is restored below.

% =============================================================================
% END APPENDED COMPACT VERSION V74 MATERIAL
% =============================================================================

\clearpage

% =============================================================================
% BEGIN APPENDED COMPACT VERSION V75 MATERIAL
% =============================================================================

\chapter{Companion audit, order covariance, and the hot binary-prefix
repair}
\label{chap:V75}

\section{Mandatory five-version companion audit}
\label{sec:V75-audit}

The newest companion files in the shared notes folder were inspected
read-only:
\begin{center}
\small
\begin{tabular}{p{0.48\linewidth}r r p{0.31\linewidth}}
\toprule
file&lines&bytes&SHA--256\\
\midrule
\texttt{Renewal\_SM\_Einstein\_Continuum\_Research\_Notes\_V91.tex}
&35543&1202031&
\texttt{1d7427fff226c1ba0274ddbbf78d8610a88749875d80657a599fe36fedef601c}
\\
\texttt{Renewal\_NS\_Stability\_Research\_Notes\_V31.tex}
&23052&887537&
\texttt{f1121b62238ad5491bb7cf03635ea9934942edf573ed8faaa9ee79e1d38a6696}
\\
\bottomrule
\end{tabular}
\end{center}
The NS file is unchanged from the V70 audit.

\begin{assessment}[Standard Model transfer]
\label{ass:V75-SM-transfer}
SM V90 proves that a finite obstruction must be projected onto the
range/cokernel split of the complete map with a uniform singular-value
gap; SM V91 then computes the complete representation trace before
declaring the local anomaly coefficient zero.  The proof-order lesson
for V75 is to test the complete binary covariance, not a protected or
gapped summand by itself, and to distinguish an actual kernel from a
small but nonzero projected carrier.  No anomaly or counterterm
identity is imported.
\end{assessment}

\begin{assessment}[Navier--Stokes transfer]
\label{ass:V75-NS-transfer}
NS V31 continues to show that the marked-to-unmarked conversion is
governed by its exact first moment and by correlation with the actual
source/exit packet.  In the present programme a large spectator bank
must therefore be moved into an exact prefix response before its row
mass is estimated.  It may not be charged once as a suffix tail and
again as a prefix edge.  No NS estimate is imported.
\end{assessment}

\begin{firewall}[Audit type boundary]
\label{fire:V75-audit-boundary}
Only proof order is transferred.  The finite Ward cokernel is not a
Wilson covariance state space, and an NS formation mark is not an
\(SU(2)\) overlap stock.
\end{firewall}

\section{Correction of the V74 acquisition order}
\label{sec:V75-correction}

V74 correctly disproved a positive gapped terminal moment.  It then
proposed testing signed cancellation first.  That proposal was too
narrow: failure of a one-row terminal moment does not imply failure of
the two-endpoint source coefficient, and the complete binary covariance
has a second, elementary contraction bound absent from the tree
envelope.

\begin{lemma}[One-row terminal failure need not be a source failure]
\label{lem:V75-terminal-not-source}
Let a V74 bank have
\[
 X_s=Na_s\asymp Ns^{-2},
 \qquad
 \Gamma_s:=
 \kappa_\otimes(\mathsf G_{s,j_s}^{\otimes N})\le1.
\]
Although \(sX_s\Gamma_s\) need not be bounded, for every fixed
companion mass \(Y\),
\begin{equation}
 s^2X_sY\Gamma_s\le C_NY.
\label{eq:V75-cold-companion}
\end{equation}
Thus replacing two endpoint factors by the square of the larger one
can create a false obstruction.
\end{lemma}

\begin{proof}
Use \(s^2X_s\asymp N\) and \(\Gamma_s\le1\).
\end{proof}

\begin{assessment}[Revised order]
\label{ass:V75-revised-order}
The next estimate must retain the exact product
\((s_\alpha\Xi_i)(s_\alpha\Xi_j)\) and the two complete binary-prefix
responses.  A one-row terminal or a root-square domination may be used
only on branches where it is quantitatively favorable.  Signed
cancellation is tested after, not before, the complete binary
covariance contraction.
\end{assessment}

\section{A suffix bank cannot see earlier covariance signs}
\label{sec:V75-suffix}

\begin{lemma}[Common suffix-factor identity]
\label{lem:V75-common-suffix}
Fix a first-connection coordinate \(r\), and let \(D\) be a coordinate
bank lying strictly after \(r\).  In every term of the exact
first-connectivity source, its local moment product is a common right
factor:
\begin{equation}
 \mathcal S_{\rho,r}
 =
 \mathcal A_{\rho,r}\otimes M_D\otimes M_{\rm rest}.
\label{eq:V75-common-suffix}
\end{equation}
For any physical spectral projection \(Q_D\) of that bank,
\begin{equation}
 (I\otimes Q_D)\mathcal S_{\rho,r}(I\otimes Q_D)
 =
 \mathcal A_{\rho,r}\otimes
 (Q_DM_DQ_D)\otimes M_{\rm rest}.
\label{eq:V75-compressed-suffix}
\end{equation}
Consequently covariance signs formed at or before \(r\) cannot cancel
a nonzero compressed suffix multiplier.
\end{lemma}

\begin{proof}
This is the suffix factor in the exact V61 formula.  Once the
cumulative join has reached \(\widehat1\) at \(r\), every later
coordinate belongs to \(M_{>r}\), independent of the earlier
moment-partition letter.  Tensor factors at distinct coordinates
commute, giving \eqref{eq:V75-common-suffix}.  Compression on the
suffix coordinate Hilbert space gives
\eqref{eq:V75-compressed-suffix}.  A scalar or operator common tensor
factor is not altered by a signed sum in
\(\mathcal A_{\rho,r}\).
\end{proof}

\begin{firewall}[Why the V74 signed proposal cannot treat a pure
suffix]
\label{fire:V75-no-suffix-cancellation}
If the V74 anti-aligned bank lies wholly after the first-connection
coordinate, projecting the earlier covariance functional cannot
remove it.  One must either use the bank as a physical transfer stock
or choose a different deterministic coordinate order in which it
belongs to a prefix response.
\end{firewall}

\section{Order covariance of first connectivity}
\label{sec:V75-order}

\begin{theorem}[Exact first-connectivity formula for every coordinate
order]
\label{thm:V75-order-covariance}
Let \(\sigma\) be any deterministic permutation of the finite
coordinate set, chosen from the input supports and representation
labels before integration.  Replacing the original order by
\(\sigma(1),\ldots,\sigma(n)\) in the V61 construction gives
\begin{equation}
 K^{\rm conn}
 =
 \sum_t\sum_{\rho<\widehat1}
 C_{<t}^{\,\sigma}(\rho)
 \otimes J_{\sigma(t)}(\rho)
 \otimes M_{>t}^{\,\sigma}.
\label{eq:V75-reordered-first-connection}
\end{equation}
Every connected coordinate word is counted exactly once, and the sum
is the same connected cumulant as in the original order.
\end{theorem}

\begin{proof}
Coordinate operators act on distinct tensor factors, so permuting
their displayed product does not change any row moment or the final
moment--partition cumulant.  Apply
\Cref{thm:V61-first-connectivity} to the permuted list.  Along any
connected word, the cumulative partition join is monotone and has a
unique first index \(t\) at which it equals \(\widehat1\).  This proves
unique counting.  Since both ordered formulas equal the same
moment--partition cumulant, their sums agree.
\end{proof}

\begin{corollary}[Pure-pair priority order]
\label{cor:V75-pair-priority}
Let \(D_{12}\) and \(D_{34}\) be any deterministic banks whose active
row sets are exactly \(\{1,2\}\) and \(\{3,4\}\).  There is an
authorized first-connectivity order placing
\(D_{12}\cup D_{34}\) before every coordinate which joins the two
pair clusters.  In every resulting \(2+2\) source term, the two banks
belong to the exact binary prefixes, not to the positive suffix.
\end{corollary}

\begin{proof}
Place the two pure-pair banks first, in any internal order, followed by
all remaining coordinates.  No coordinate in their union joins
\(12|34\) to \(\widehat1\).  Therefore any first full connection occurs
later.  Apply \Cref{thm:V75-order-covariance}.  Prefix factorization
\Cref{thm:V61-prefix-factorization} puts the banks in
\(K_{12}^{<r,\rm conn}\) and
\(K_{34}^{<r,\rm conn}\).
\end{proof}

\begin{remark}[The order is deterministic, not outcome-adaptive]
\label{rem:V75-order-deterministic}
The permutation depends only on fixed representation/support data.
It is not chosen from a sampled group variable or from a physical
output label.  Hence it changes neither the probability law nor the
one-payer allocation.
\end{remark}

\section{The complete pure-pair binary covariance}
\label{sec:V75-binary}

Let \(D\) be a bank shared by precisely two rows \(u,v\).  At
coordinate \(w\in D\), write
\[
 T_w:=\mathbb E_w[X_{u,w}\otimes X_{v,w}],
 \qquad
 A_w:=\mathbb E_wX_{u,w},
 \qquad
 B_w:=\mathbb E_wX_{v,w}.
\]

\begin{theorem}[Exact two-term bank covariance]
\label{thm:V75-bank-covariance}
The complete binary connected prefix on \(D\) is
\begin{equation}
 \boxed{
 K_{uv}^{\rm conn}(D)
 =
 \bigotimes_{w\in D}T_w
 -
 \left(\bigotimes_{w\in D}A_w\right)
 \otimes
 \left(\bigotimes_{w\in D}B_w\right).}
\label{eq:V75-bank-covariance}
\end{equation}
Consequently
\begin{equation}
 \boxed{
 \kappa_\otimes
 \bigl(K_{uv}^{\rm conn}(D)\bigr)
 \le
 C\min\{1,s_\alpha H_{uv}(D)\}.}
\label{eq:V75-binary-min}
\end{equation}
\end{theorem}

\begin{proof}
Independence across coordinates gives
\[
 \mathbb E\!\left[
  \bigotimes_{w\in D}X_{u,w}\otimes
  \bigotimes_{w\in D}X_{v,w}\right]
 =
 \bigotimes_{w\in D}T_w,
\]
while the product of the two row expectations is the second term in
\eqref{eq:V75-bank-covariance}.  Their difference is the definition
of binary covariance.

Both complete moments are contractions, so the capacity of their
difference is at most a universal constant.  Independently, the valid
binary entry of \Cref{lem:V37-lower-table} gives
\[
 \kappa_\otimes(K_{uv}^{\rm conn}(D))
 \le Cs_\alpha H_{uv}(D).
\]
Take the smaller of the two bounds.
\end{proof}

\begin{firewall}[The minimum is taken after the covariance]
\label{fire:V75-min-after-covariance}
The contraction entry in \eqref{eq:V75-binary-min} applies to the
complete two-term covariance
\eqref{eq:V75-bank-covariance}.  It is not a coordinatewise minimum
and creates no \(2^{|D|}\) expansion.
\end{firewall}

\section{The V74 anti-aligned bank inside the exact prefix}
\label{sec:V75-anti-prefix}

For the identical-spin heat-kernel bank of V74,
\[
 T_w=\mathsf T_{s,j_s},
\qquad
 A_w=B_w=e^{-s j_s(j_s+1)}I
\]
up to the harmless fixed heat-time convention.  Compress to the
all-gapped tensor sector and evaluate on the anti-aligned product
vector.

\begin{proposition}[The separated covariance term does not cancel the
anti-aligned band]
\label{prop:V75-no-leading-cancel}
For fixed \(N\) and all sufficiently small \(s\),
\begin{equation}
 \left\langle
 \psi_{j_s}^{\otimes N},
 K_{uv}^{\rm conn}(D)
 \psi_{j_s}^{\otimes N}
 \right\rangle_{\rm all\text{-}gapped}
 \ge
 \frac12c_0^N.
\label{eq:V75-no-leading-cancel}
\end{equation}
\end{proposition}

\begin{proof}
The joint term has expectation at least \(c_0^N\) by
\eqref{eq:V74-bank-lower}.  The separated term has norm
\[
 e^{-2sNj_s(j_s+1)}.
\]
Since \(j_s\asymp s^{-1}\), its exponent tends to
\(-\infty\) like \(-cN/s\).  It is therefore at most
\(c_0^N/2\) for small \(s\).  Subtract it from the joint lower bound.
\end{proof}

\begin{proposition}[Nevertheless the normalized hot prefix is small]
\label{prop:V75-normalized-prefix-small}
For the same bank,
\begin{equation}
 \frac{\kappa_\otimes(K_{uv}^{\rm conn}(D))}
      {\Xi_u\Xi_v}
 \le
 \frac{C}{N^2a_s^2},
\label{eq:V75-prefix-normalized}
\end{equation}
whereas \(\theta_{uv}=1/N\).  Hence
\begin{equation}
 \frac{\kappa_\otimes(K_{uv}^{\rm conn}(D))}
      {\Xi_u\Xi_v\,\theta_{uv}}
 \le
 \frac{C}{Na_s^2}
 \longrightarrow0
\quad(s\downarrow0).
\label{eq:V75-prefix-vs-mark}
\end{equation}
\end{proposition}

\begin{proof}
Use the contraction entry of
\eqref{eq:V75-binary-min} and the exact masses
\(\Xi_u=\Xi_v=Na_s\) from
\eqref{eq:V74-bank-masses}.  Divide once more by
\(\theta_{uv}=1/N\).
\end{proof}

\begin{assessment}[Meaning of the noncancellation]
\label{ass:V75-noncancellation-meaning}
The anti-aligned gapped band is not in the kernel of the binary
covariance; the joint moment survives and the separated term is
exponentially smaller.  Even so, the exact prefix is harmless after
row normalization because its contraction bound is vastly sharper
than the coarse \(s_\alpha H_{uv}\) tree numerator in this hot regime.
Thus V74 found a failure of the terminal-tail parametrization, not a
failure of the binary connected carrier.
\end{assessment}

\section{The hot binary-prefix gate}
\label{sec:V75-hot-prefix-gate}

\begin{definition}[Hot binary minimum ratio]
\label{def:V75-beta}
For a two-row prefix bank define
\begin{equation}
 \beta_{uv}
 :=
 \frac{\min\{1,s_\alpha H_{uv}\}}{H_{uv}}
 =
 \min\left\{\frac1{H_{uv}},s_\alpha\right\}.
\label{eq:V75-beta}
\end{equation}
when \(H_{uv}>0\), and set \(\beta_{uv}=0\) when \(H_{uv}=0\).
In the latter case the binary covariance vanishes.
The exact covariance estimate is
\[
 \kappa_\otimes(K_{uv}^{\rm conn})
 \le C\beta_{uv}H_{uv}.
\]
\end{definition}

\begin{proposition}[Refined joining-payer coefficient]
\label{prop:V75-joining-payer-refined}
For the \(2+2\) source branch in which the assembled local joining
covariance carries the sole payer, the crossing choice
\(i\in\{1,2\}\), \(j\in\{3,4\}\) has normalized coefficient at most
\begin{equation}
 C\,
 \Xi_i\Xi_j\,
 \beta_{12}\beta_{34}\,
 \theta_{12}\theta_{34}\theta_{ij}.
\label{eq:V75-refined-joining-payer}
\end{equation}
\end{proposition}

\begin{proof}
Use the two binary minimum bounds for the prefixes and the paid local
joining estimate
\eqref{eq:V65-local-paid-crossing}.  The numerator is
\[
 C\beta_{12}\beta_{34}H_{12}H_{34}H_{ij}.
\]
Divide by the four row masses.  Comparing with the path tree mark
introduces precisely the degree-excess factor \(\Xi_i\Xi_j\), proving
\eqref{eq:V75-refined-joining-payer}.
\end{proof}

\begin{corollary}[The double anti-aligned test closes in the refined
joining-payer ledger]
\label{cor:V75-anti-test-closes}
If the \(12\) and \(34\) prefixes are the identical V74 banks, then
the coefficient multiplying the tree mark in
\eqref{eq:V75-refined-joining-payer} tends to zero as
\(s_\alpha\downarrow0\), for fixed \(N\).
\end{corollary}

\begin{proof}
Here \(\Xi_i=Na_s\), \(H_{12}=H_{34}=Na_s^2\), and
\(s_\alpha H_{12}\to\infty\).  Thus
\[
 \Xi_i\beta_{12}
 =
 \frac{Na_s}{Na_s^2}
 =
 \frac1{a_s},
\]
and similarly for \(j,34\).  Their product is \(a_s^{-2}\to0\).
\end{proof}

\begin{assessment}[What remains of the hot binary problem]
\label{ass:V75-hot-binary-residual}
The contraction/response minimum closes the extremal anti-aligned
test but not yet every representation distribution.  The required
uniform scalar statement for the joining-payer branch is
\[
 \Xi_i\Xi_j\beta_{12}\beta_{34}\le C
\]
after the exact row-rooted private, unbalanced, and balanced
alternatives are inserted.  Prefix-payer and suffix-payer branches
require corresponding paid minimum cards; they are not consequences
of the unpaid contraction bound.
\end{assessment}

\section{V75 ledger and V76 mandate}
\label{sec:V75-ledger}

\begin{theorem}[Integrated compact-series V75 statement]
\label{thm:V75-integrated}
Preserving compact V1--V74, V75 proves:
\begin{enumerate}[label=\textup{(V75.\arabic*)},leftmargin=6.7em]
\item the mandatory SM V91 and NS V31 audit;
\item a suffix physical bank is a common factor invisible to earlier
      covariance signs;
\item the exact first-connectivity formula is valid for every
      deterministic coordinate order, permitting pure-pair priority;
\item a pure-pair binary prefix has the exact two-term covariance and
      the bound \(C\min\{1,s_\alpha H\}\);
\item the V74 anti-aligned band survives the covariance signs but is
      harmless after exact prefix normalization; and
\item the joining-payer \(2+2\) branch has the refined coefficient
      \eqref{eq:V75-refined-joining-payer}.
\end{enumerate}
\end{theorem}

\begin{proof}
These are
\Cref{lem:V75-common-suffix,
thm:V75-order-covariance,
cor:V75-pair-priority,
thm:V75-bank-covariance,
prop:V75-no-leading-cancel,
prop:V75-normalized-prefix-small,
prop:V75-joining-payer-refined}.
\end{proof}

\begin{openproblem}[V76 mandate: uniform hot binary minimum compiler]
\label{op:V76-mandate}
\begin{enumerate}
\item Keep the two-endpoint product
      \(\Xi_i\Xi_j\beta_{12}\beta_{34}\); do not dominate it by the
      square of one row.
\item Prove the joining-payer bound first, using the exact
      private/unbalanced/balanced row alternatives and the fact that
      each internal endpoint belongs to its own binary prefix.
\item Derive a paid binary minimum card by combining the V37 response
      estimate with the complete covariance contraction and the single
      output-mass split.
\item Rebuild the two prefix-payer branches with one paid minimum and
      one unpaid minimum.
\item For a suffix payer, retain the exact suffix Wilson failure debit
      without discarding the two prefix minima.
\item If a large shared bank lies after first connectivity, use order
      covariance to move a deterministic pure-pair component into the
      prefix; treat higher-incidence components separately because
      they may change the prefix family.
\end{enumerate}
\end{openproblem}

\begin{firewall}[Claim boundary after compact V75]
\label{fire:V75-current-boundary}
V75 repairs the V74 acquisition order and closes its extremal
anti-aligned test in the joining-payer ledger.  It does not yet prove
the uniform hot binary minimum compiler, the other payer branches,
higher-incidence order priority, marked-heavy/overlapping balanced
banks, or balanced output ancestry.  The other first-connectivity
families, global quartic and higher MTA, WUFC, CPM, continuum
Yang--Mills existence, and a positive mass gap remain open.  The full
Yang--Mills mass gap remains open.
\end{firewall}

\section{Append-only record for compact V75}
\label{sec:V75-record}

V75 was appended to the complete compact V74 file.  The exact V74
parent had \(30838\) newline-delimited source records, \(1063469\)
bytes, and SHA--256 digest
\[
\texttt{07f972a6c6966e40b2136cba916a407fa1c46342224bd4dff930d7423e919479}.
\]
No compact V74 mathematical content was changed.  The sole final
\verb|\end{document}| is restored below.

% =============================================================================
% END APPENDED COMPACT VERSION V75 MATERIAL
% =============================================================================

\clearpage

% =============================================================================
% BEGIN APPENDED COMPACT VERSION V76 MATERIAL
% =============================================================================

\chapter{Prefix-dominant hot binary minima}
\label{chap:V76}

\section{Context and coordinate-resolved prefix masses}
\label{sec:V76-context}

V75 proved the complete binary covariance estimate
\[
 \kappa_\otimes(K_{uv}^{\rm conn})
 \le C\min\{1,s_\alpha H_{uv}\}.
\]
This note identifies the sector in which that minimum pays the repeated
row mass directly.  The observation is elementary but decisive: on
every coordinate shared by two nontrivial rows, the partner Casimir
weight has a positive group-dependent floor.  Hence the overlap stock
dominates the one-row mass carried by that binary bank.  In the
small-\(sH\) branch the response factor \(s\) pays the row; in the
large-\(sH\) branch the covariance contraction pays it through
\(H^{-1}\).

Fix a \(2+2\) first-connection term with joining coordinate \(r\).
For its \(12\)-prefix put
\begin{align}
 X_{1|2}^{<r}
 &:=
 \sum_{\substack{e<r\\1,2\in A_e}}a_{1,e},
 &
 X_{2|1}^{<r}
 &:=
 \sum_{\substack{e<r\\1,2\in A_e}}a_{2,e},
\label{eq:V76-12-row-stocks}\\
 H_{12}^{<r}
 &:=
 \sum_{e<r}a_{1,e}a_{2,e}.
\label{eq:V76-12-overlap}
\end{align}
Define \(X_{3|4}^{<r}\), \(X_{4|3}^{<r}\), and
\(H_{34}^{<r}\) analogously.  Inactive entries are zero.

\section{The binary minimum pays one bank-carried row}
\label{sec:V76-scalar-minimum}

\begin{lemma}[Overlap dominates either endpoint bank mass]
\label{lem:V76-overlap-dominates}
There is a fixed \(a_*>0\), depending only on the compact gauge group,
such that
\begin{align}
 H_{12}^{<r}
 &\ge a_*X_{1|2}^{<r},
 &
 H_{12}^{<r}
 &\ge a_*X_{2|1}^{<r},
\label{eq:V76-H-dominates-X}
\end{align}
and the analogous inequalities hold for the \(34\)-prefix.
\end{lemma}

\begin{proof}
At a summand contributing to \(X_{1|2}^{<r}\), row \(2\) carries a
nontrivial irreducible representation.  The nonzero Casimir spectrum
of the fixed compact semisimple group has a positive floor, so
\[
 a_{2,e}=1+C_2(R_{2,e})\ge a_*.
\]
Thus \(a_{1,e}a_{2,e}\ge a_*a_{1,e}\).  Sum over \(e<r\).
Interchanging the rows gives the second inequality, and the \(34\)
case is identical.
\end{proof}

\begin{lemma}[Scalar binary-minimum compiler]
\label{lem:V76-scalar-compiler}
Let \(X,H\ge0\) satisfy \(H\ge a_*X\).  Then
\begin{equation}
 \boxed{
 X\min\left\{\frac1H,s_\alpha\right\}
 \le\frac1{a_*},}
\label{eq:V76-scalar-compiler}
\end{equation}
with the left side interpreted as zero when \(H=0\).
\end{lemma}

\begin{proof}
If \(s_\alpha H\le1\), then
\[
 s_\alpha X\le a_*^{-1}s_\alpha H\le a_*^{-1}.
\]
If \(s_\alpha H\ge1\), then
\[
 X/H\le a_*^{-1}.
\]
These are the two entries of the minimum.
\end{proof}

\begin{definition}[Prefix-dominant endpoint]
\label{def:V76-prefix-dominant}
Fix \(0<\kappa\le1\).  An endpoint \(u\) of a binary prefix
\(uv\) is \(\kappa\)-prefix-dominant at \(r\) if
\begin{equation}
 X_{u|v}^{<r}\ge\kappa\Xi_u.
\label{eq:V76-prefix-dominant}
\end{equation}
\end{definition}

\begin{corollary}[One prefix minimum pays one repeated endpoint]
\label{cor:V76-one-endpoint}
If \(u\) is \(\kappa\)-prefix-dominant in the \(uv\)-prefix, then
\begin{equation}
 \boxed{
 \min\{1,s_\alpha H_{uv}^{<r}\}
 \le
 C_\kappa\frac{H_{uv}^{<r}}{\Xi_u}.}
\label{eq:V76-one-endpoint}
\end{equation}
No thermal, private, or output assumption is required.
\end{corollary}

\begin{proof}
For \(H_{uv}^{<r}>0\), write the left side as
\[
 H_{uv}^{<r}
 \min\left\{\frac1{H_{uv}^{<r}},s_\alpha\right\}.
\]
By \eqref{eq:V76-prefix-dominant},
\(\Xi_u\le\kappa^{-1}X_{u|v}^{<r}\).  Apply
\Cref{lem:V76-overlap-dominates,
lem:V76-scalar-compiler}.  If \(H=0\), both sides vanish because the
binary covariance is zero.
\end{proof}

\begin{firewall}[No private tail is hidden in the minimum]
\label{fire:V76-no-hidden-tail}
\Cref{cor:V76-one-endpoint} uses only mass carried by coordinates
where both prefix rows are nontrivial.  Mass from a private coordinate,
a different companion row, the joining coordinate, or the suffix is
not included in \(X_{u|v}^{<r}\) and is not silently paid.
\end{firewall}

\section{Uniform joining-payer closure in the prefix-dominant sector}
\label{sec:V76-joining-payer}

\begin{theorem}[Prefix-dominant joining-payer \(2+2\) DRQC]
\label{thm:V76-joining-payer}
Fix \(i\in\{1,2\}\) and \(j\in\{3,4\}\).  Consider the ordered
\(2+2\) source terms whose first full connection occurs at \(r\) and
whose local joining covariance carries the sole final payer.  If
\begin{equation}
 X_{i|\{1,2\}\setminus\{i\}}^{<r}
 \ge\kappa\Xi_i,
 \qquad
 X_{j|\{3,4\}\setminus\{j\}}^{<r}
 \ge\kappa\Xi_j,
\label{eq:V76-double-prefix-dominant}
\end{equation}
then their normalized sum satisfies
\begin{equation}
 \boxed{
 \frac{\kappa_\otimes(\mathcal S_{22}^{\rm join\text{-}pay,
 \,prefix\text{-}dom})}
      {\Xi_1\Xi_2\Xi_3\Xi_4}
 \le
 C_\kappa\,
 \theta_{12}\theta_{34}\theta_{ij}.}
\label{eq:V76-joining-payer-DRQC}
\end{equation}
The estimate is uniform in how hot the two internal rows are.
\end{theorem}

\begin{proof}
At fixed \(r\), apply the two complete binary covariance bounds
\eqref{eq:V75-binary-min} with the prefix stocks
\(H_{12}^{<r}\), \(H_{34}^{<r}\).  Apply the paid local joining bound
\eqref{eq:V65-local-paid-crossing}.  The resulting numerator is at
most
\[
 C
 \min\{1,s_\alpha H_{12}^{<r}\}
 \min\{1,s_\alpha H_{34}^{<r}\}
 h_{ij,r}.
\]
Use \Cref{cor:V76-one-endpoint} on rows \(i\) and \(j\).  After
division by the four row masses this is bounded by
\[
 C_\kappa
 \frac{H_{12}^{<r}H_{34}^{<r}h_{ij,r}}
 {\Xi_1\Xi_2\Xi_3\Xi_4\Xi_i\Xi_j}.
\]
The denominator is exactly the product of row masses with powers equal
to the degrees of the path \(\{12,34,ij\}\).  Finally sum \(r\) and
use the ordered stock inequality
\[
 \sum_r
 H_{12}^{<r}H_{34}^{<r}h_{ij,r}
 \le
 H_{12}H_{34}H_{ij}.
\]
This gives \eqref{eq:V76-joining-payer-DRQC}.
\end{proof}

\begin{corollary}[All pure-pair-dominant hot tests close after priority
ordering]
\label{cor:V76-pure-pair-close}
Suppose deterministic pure-pair banks \(D_{12}\), \(D_{34}\) carry at
least a fraction \(\kappa\) of the chosen internal row masses \(i,j\).
Use the pair-priority order of
\Cref{cor:V75-pair-priority}.  Then every subsequent \(2+2\)
first-connection term with a joining payer satisfies
\eqref{eq:V76-joining-payer-DRQC}.  This includes the single and
double anti-aligned V74 test banks.
\end{corollary}

\begin{proof}
The priority order puts the entire banks before every coordinate that
joins \(12|34\).  Hence their row stocks are included in the two
quantities on the left of
\eqref{eq:V76-double-prefix-dominant}.  Apply
\Cref{thm:V76-joining-payer}.  The V74 banks are pure-pair and carry
all but fixed lower-order additions to their large row masses.
\end{proof}

\begin{remark}[The proof uses the exact two-endpoint product]
\label{rem:V76-two-endpoints}
Each binary minimum pays its own internal endpoint.  At no stage is
\(\Xi_i\Xi_j\) enlarged to
\(\max\{\Xi_i,\Xi_j\}^2\).  Thus the cold-companion loss identified in
\Cref{lem:V75-terminal-not-source} cannot occur.
\end{remark}

\section{A mixed endpoint version}
\label{sec:V76-mixed}

\begin{proposition}[One prefix-dominant endpoint plus one certified
endpoint]
\label{prop:V76-mixed-endpoint}
Assume row \(i\) is \(\kappa\)-prefix-dominant in the \(12\)-prefix.
Suppose the complete \(34\)-prefix packet, together with retained
private/unbalanced/balanced factors, has a sourcewise estimate
\begin{equation}
 \min\{1,s_\alpha H_{34}^{<r}\}\,\mathcal C_j
 \le
 C\frac{H_{34}^{<r}}{\Xi_j},
\label{eq:V76-other-certificate}
\end{equation}
where \(\mathcal C_j\) is not used by the final joining payer.  Then
the same DRQC conclusion
\eqref{eq:V76-joining-payer-DRQC} holds for that mixed branch.
\end{proposition}

\begin{proof}
Apply \Cref{cor:V76-one-endpoint} to the \(12\)-prefix and
\eqref{eq:V76-other-certificate} to the \(34\)-prefix.  The rest of the
proof is identical to \Cref{thm:V76-joining-payer}.
\end{proof}

\begin{assessment}[How V69--V73 feed the mixed certificate]
\label{ass:V76-mixed-inputs}
The factor \(\mathcal C_j\) may be supplied by a private Wilson product,
an unbalanced reserve/marked heat tail, a clean escalating spectral
chain with terminal moment, or a diffuse all-protected terminal.
What is not yet authorized is an overlapping bank, a marked-heavy
balanced transfer used twice, or the V74 non-escalating gapped
terminal.  The proposition records the exact interface rather than
claiming those residual cases.
\end{assessment}

\section{Exact residual inside the joining-payer branch}
\label{sec:V76-residual}

For the \(12\)-prefix decompose row \(i\)''s mass as
\begin{equation}
 \Xi_i=X_{i|i''}^{<r}+Y_i^{(r)},
 \qquad i''=\{1,2\}\setminus\{i\},
\label{eq:V76-prefix-residual-split}
\end{equation}
and analogously
\(\Xi_j=X_{j|j''}^{<r}+Y_j^{(r)}\) in the \(34\)-prefix.

\begin{corollary}[Every unclosed joining-payer term has a large
nonprefix endpoint residual]
\label{cor:V76-large-nonprefix-residual}
Fix \(0<\kappa<1\).  If a joining-payer term is not covered by
\Cref{thm:V76-joining-payer}, then at least one internal endpoint
\(k\in\{i,j\}\) satisfies
\begin{equation}
 Y_k^{(r)}>(1-\kappa)\Xi_k.
\label{eq:V76-large-nonprefix-residual}
\end{equation}
That mass lies at coordinates which are not shared with the endpoint''s
own prefix partner before \(r\).
\end{corollary}

\begin{proof}
Failure of
\eqref{eq:V76-double-prefix-dominant} means
\(X_{k|k''}^{<r}<\kappa\Xi_k\) for one endpoint.  Subtract from
\eqref{eq:V76-prefix-residual-split}.
\end{proof}

\begin{assessment}[The next source-local split]
\label{ass:V76-next-split}
The large residual in
\eqref{eq:V76-large-nonprefix-residual} is now sharply localized:
it consists of pair-private prefix coordinates, coordinates involving a
different row and hence capable of changing the prefix partition, the
joining coordinate, and the suffix.  Private and unbalanced pieces are
already controlled.  The remaining balanced piece must be divided
between:
\begin{enumerate}[label=\textup{(\alph*)}]
\item higher-incidence coordinates which should be moved earlier and
      reclassified into the \(2+1+1\), \(3+1\), or discrete source
      family; and
\item genuinely later coordinates, for which
      \Cref{lem:V75-common-suffix} proves that no earlier covariance
      cancellation is available.
\end{enumerate}
This is the correct upstream interface for order optimization.
\end{assessment}

\section{V76 ledger and V77 mandate}
\label{sec:V76-ledger}

\begin{theorem}[Integrated compact-series V76 statement]
\label{thm:V76-integrated}
Preserving compact V1--V75, V76 proves:
\begin{enumerate}[label=\textup{(V76.\arabic*)},leftmargin=6.7em]
\item a binary overlap stock dominates either endpoint mass carried by
      its shared prefix bank;
\item the complete covariance minimum pays any
      prefix-dominant repeated row uniformly;
\item the full joining-payer \(2+2\) source satisfies DRQC whenever
      both internal endpoints are prefix-dominant;
\item pure-pair-dominant hot banks, including the V74 tests, close
      after deterministic pair-priority ordering; and
\item every remaining joining-payer term has a quantitatively large
      nonprefix residual at an internal endpoint.
\end{enumerate}
\end{theorem}

\begin{proof}
These are
\Cref{lem:V76-overlap-dominates,
lem:V76-scalar-compiler,
cor:V76-one-endpoint,
thm:V76-joining-payer,
cor:V76-pure-pair-close,
cor:V76-large-nonprefix-residual}.
\end{proof}

\begin{openproblem}[V77 mandate: higher-incidence priority compiler]
\label{op:V77-mandate}
\begin{enumerate}
\item Split the large nonprefix residual
      \eqref{eq:V76-large-nonprefix-residual} into private,
      unbalanced, balanced higher-incidence, joining, and suffix
      stocks without overlapping coordinate charges.
\item Move the deterministic balanced higher-incidence stock before
      the first connection using
      \Cref{thm:V75-order-covariance}.
\item Record exactly which of the other three source families the new
      first-connection word enters; do not force it back into
      \(2+2\).
\item Apply the corresponding assembled covariance/ternary/local-four
      minimum before taking norms.
\item Treat the joining-coordinate-heavy case by a local paid minimum,
      and retain the suffix case as a separate obstruction because of
      \Cref{lem:V75-common-suffix}.
\item Only after the joining-payer branch is complete, derive the paid
      binary minimum needed for prefix-payer branches.
\end{enumerate}
\end{openproblem}

\begin{firewall}[Claim boundary after compact V76]
\label{fire:V76-current-boundary}
V76 closes the prefix-dominant joining-payer sector of the assembled
\(2+2\) source.  It does not close large balanced nonprefix residuals,
prefix- or suffix-payer branches, marked-heavy/overlapping banks,
balanced output ancestry, or the other three hot source families.
Global quartic and higher MTA, WUFC, CPM, continuum Yang--Mills
existence, and a positive mass gap remain open.  The full Yang--Mills
mass gap remains open.
\end{firewall}

\section{Append-only record for compact V76}
\label{sec:V76-record}

V76 was appended to the complete compact V75 file.  The exact V75
parent had \(31373\) newline-delimited source records, \(1081672\)
bytes, and SHA--256 digest
\[
\texttt{9975e4cac06aaacaf22dfa508d10ca350229fead0d449b4e1420c225ad5473e8}.
\]
No compact V75 mathematical content was changed.  The sole final
\verb|\end{document}| is restored below.

% =============================================================================
% END APPENDED COMPACT VERSION V76 MATERIAL
% =============================================================================

% =============================================================================
% BEGIN APPENDED COMPACT VERSION V77 MATERIAL
% =============================================================================

\clearpage
\chapter{Bankwise first-connection resummation}
\label{chap:V77}

\section{Context: why coordinate priority is insufficient}
\label{sec:V77-context}

V76 reduced every still-unclosed joining-payer \(2+2\) term to a large
endpoint mass outside its own pure-pair prefix.  One part of that
residual may be carried by coordinates incident to three or four rows.
The proposed V76 mandate was to move such coordinates earlier.  That
proposal is not valid in its naive form.  A higher-incidence coordinate
which sees both blocks of an already formed cut may itself be the first
full connection.  All later coordinates then occur in the complete
suffix of the V61 formula, regardless of how large their combined row
mass is.

This note replaces the ordering proposal by an exact resummation.  Once
a proper two-block partition \(C\mid D\) has been formed, sum over
\emph{every} possible first coordinate which joins \(C\) to \(D\).
The resulting packet is a single composite covariance over the entire
remaining coordinate bank.  Thus the apparent suffix coordinates of
the bank are not discarded: they are part of the assembled covariance
before any norm is taken.  The construction works for pure-pair,
three-row, and four-row coordinates at once.

The main analytic consequence proved here is an unpaid
contraction/response minimum for that complete bank.  It pays one
cut-exposed row whenever a fixed fraction of that row's mass is carried
by the bank, including when the mass is carried by higher-incidence
coordinates.  This is a genuine extension of V76.  It does not yet give
the analogous \emph{paid} minimum required when the same bank carries
the sole final resolvent.

\section{A fixed two-block carrier and its residual bank}
\label{sec:V77-fixed-cut}

Let
\[
 \rho=\{C,D\},\qquad
 C\mathbin{\dot\cup}D=\{1,2,3,4\},
\]
be a proper two-block partition.  We condition algebraically on a fixed
earlier V61 prefix carrier whose cumulative partition is exactly
\(\rho\).  Thus the rows in \(C\) are already connected to one another,
the rows in \(D\) are already connected to one another, and the two
blocks have not yet been joined.  Let
\[
 B=(e_1,\ldots,e_m)
\]
be the complete ordered set of coordinates not used by that fixed
carrier.  Coordinates which see only one side of the cut are retained
in \(B\); their cut-joining letter will simply be zero.

At \(e\in B\), define the cut-replicated local moment by
\[
 M_e(\rho):=
 \sum_{\substack{\pi\in\Pi_4\\\pi\le\rho}}L_{e,\pi},
 \qquad
 M_e(\widehat1)=M_e
\]
in the notation of \eqref{eq:V61-local-letters}, and put
\begin{align}
 T_e&:=M_e(\widehat 1),&
 R_e^\rho&:=M_e(\rho),&
 \Delta_e^\rho&:=T_e-R_e^\rho .
\label{eq:V77-local-TRD}
\end{align}
Here \(T_e\) uses the common coordinate variable for all active rows,
whereas \(R_e^\rho\) uses independent replicas on the two blocks.
Because \(\rho\) has exactly two blocks, a local partition fails to
refine \(\rho\) exactly when it joins those two blocks.  Accordingly,
\(\Delta_e^\rho=J_e(\rho)\) is precisely the assembled local two-block
joining letter in \eqref{eq:V61-joining-letter}.  If \(e\) has no
active row on one side of the cut, then \(T_e=R_e^\rho\) and
\(\Delta_e^\rho=0\).

\begin{definition}[Conditioned first-join packet]
\label{def:V77-first-join-packet}
For the fixed cut carrier and residual bank \(B\), define
\begin{equation}
 \mathfrak J_B^\rho
 :=
 \sum_{r=1}^m
 \left(\bigotimes_{q<r}R_{e_q}^\rho\right)
 \otimes\Delta_{e_r}^\rho
 \otimes
 \left(\bigotimes_{q>r}T_{e_q}\right).
\label{eq:V77-first-join-packet}
\end{equation}
Empty tensor products are identities.  Terms with
\(\Delta_{e_r}^\rho=0\) vanish.
\end{definition}

The convention in \eqref{eq:V77-first-join-packet} is forced by first
connectivity.  Before the first full join, the two blocks are replicated
separately and hence contribute \(R^\rho\).  After the first full join,
the V61 suffix is the complete moment and hence contributes \(T\).

\section{The exact bank telescope}
\label{sec:V77-telescope}

\begin{theorem}[Bankwise first-connection telescope]
\label{thm:V77-bank-telescope}
For every finite residual bank \(B=(e_1,\ldots,e_m)\),
\begin{equation}
 \boxed{
 \mathfrak J_B^\rho
 =
 \bigotimes_{q=1}^m T_{e_q}
 -
 \bigotimes_{q=1}^m R_{e_q}^\rho .}
\label{eq:V77-bank-telescope}
\end{equation}
In particular, the sum of the conditioned V61 terms over all possible
first full-connection coordinates in \(B\) has only two endpoint
moments.
\end{theorem}

\begin{proof}
Put
\[
 P_r:=
 \left(\bigotimes_{q<r}R_{e_q}^\rho\right)
 \otimes
 \left(\bigotimes_{q\ge r}T_{e_q}\right),
 \qquad 1\le r\le m+1,
\]
where
\[
 P_1=\bigotimes_{q=1}^mT_{e_q},
 \qquad
 P_{m+1}=\bigotimes_{q=1}^mR_{e_q}^\rho.
\]
Tensor factors belonging to distinct coordinates commute.  Therefore
\begin{align*}
 P_r-P_{r+1}
 &=
 \left(\bigotimes_{q<r}R_{e_q}^\rho\right)
 \otimes(T_{e_r}-R_{e_r}^\rho)
 \otimes
 \left(\bigotimes_{q>r}T_{e_q}\right)\\
 &=
 \left(\bigotimes_{q<r}R_{e_q}^\rho\right)
 \otimes\Delta_{e_r}^\rho
 \otimes
 \left(\bigotimes_{q>r}T_{e_q}\right).
\end{align*}
Summing \(P_r-P_{r+1}\) from \(r=1\) to \(m\) telescopes to
\(P_1-P_{m+1}\), which is
\eqref{eq:V77-bank-telescope}.
\end{proof}

\begin{corollary}[Exact compatibility with V61 first connectivity]
\label{cor:V77-V61-packet}
Multiplying \(\mathfrak J_B^\rho\) by the fixed earlier prefix carrier
gives exactly the sum of all terms in
\eqref{eq:V61-first-connectivity} which have that carrier, enter the
bank with partition \(\rho\), and first reach
\(\widehat1\) at a coordinate of \(B\).  No connected word is omitted
or counted twice.
\end{corollary}

\begin{proof}
For a term whose first join occurs at \(e_r\), V61 gives the separated
letters \(R_{e_q}^\rho\) for \(q<r\), the joining letter
\(\Delta_{e_r}^\rho=J_{e_r}(\rho)\), and the complete suffix letters
\(T_{e_q}\) for \(q>r\).  This is the \(r\)-th summand of
\eqref{eq:V77-first-join-packet}.  Monotonicity of the cumulative
partition join gives a unique first joining index.  Summing those
unique indices proves the assertion.
\end{proof}

\begin{corollary}[Order independence of the conditioned packet]
\label{cor:V77-packet-order}
The value of \(\mathfrak J_B^\rho\) is unchanged by any deterministic
permutation of \(B\).
\end{corollary}

\begin{proof}
The two endpoint tensor products on the right of
\eqref{eq:V77-bank-telescope} are unchanged by permuting distinct
coordinate factors.  This also follows by applying
\Cref{thm:V75-order-covariance} to the fixed conditioned packet.
\end{proof}

\begin{proposition}[Priority alone cannot expose a whole joining bank]
\label{prop:V77-priority-no-go}
Suppose at least two coordinates \(e,f\in B\) have nonzero joining
letters for \(\rho\).  In every deterministic order, each individual
V61 first-connection term has exactly one of these coordinates as its
first joining letter; every later cut-joining coordinate lies in its
complete suffix.  Consequently no deterministic priority rule can turn
the sum of the local joining stocks of all of \(B\) into the stock of
one local first-joining letter.
\end{proposition}

\begin{proof}
In a fixed connected word the cumulative partition is monotone and has
a unique first index at which it becomes \(\widehat1\), by
\Cref{thm:V61-first-connectivity}.  All later coordinates therefore
belong to \(M_{>r}\), even when they also see both sides of \(\rho\).
A permutation changes which coordinate is first, but it does not change
the uniqueness statement.  Only the sum over all possible first
indices, evaluated by \Cref{thm:V77-bank-telescope}, retains the whole
bank in one assembled object.
\end{proof}

\begin{firewall}[What has and has not been moved]
\label{fire:V77-no-formal-move}
No higher-incidence coordinate is declared to be an internal prefix
coordinate.  Such a declaration would be false when the coordinate
already joins \(C\) to \(D\).  Instead, the full residual bank is left
in its original conditioned source packet and the first-index sum is
performed exactly.  The result is an algebraic regrouping of the V61
formula, not an outcome-adaptive ordering and not a new probabilistic
independence assertion.
\end{firewall}

\section{Identification as a composite covariance}
\label{sec:V77-composite}

Let \(F_C^B\) and \(F_D^B\) be the ordered products of the row Wilson
factors belonging respectively to the blocks \(C\) and \(D\), restricted
to the coordinates in \(B\).  Expectations use the same coordinate
variables on the two sides unless a prime is displayed.

\begin{theorem}[The bank telescope is the complete composite covariance]
\label{thm:V77-composite-covariance}
\begin{equation}
 \boxed{
 \mathfrak J_B^\rho
 =
 \mathbb E_B\!\left[F_C^B\otimes F_D^B\right]
 -
 \mathbb E_BF_C^B\otimes\mathbb E'_BF_D^B
 =
 \operatorname{Cov}_B(F_C^B,F_D^B).}
\label{eq:V77-composite-covariance}
\end{equation}
The equality is at the signed operator level, before physical
enveloping or a one-resolvent allocation.
\end{theorem}

\begin{proof}
Coordinate independence gives
\[
 \mathbb E_B[F_C^B\otimes F_D^B]
 =
 \bigotimes_{e\in B}T_e.
\]
Using independent replicas on the two blocks gives
\[
 \mathbb E_BF_C^B\otimes\mathbb E'_BF_D^B
 =
 \bigotimes_{e\in B}R_e^\rho.
\]
Subtract and apply \Cref{thm:V77-bank-telescope}.
\end{proof}

\begin{remark}[Relation to the V75 pure-pair identity]
\label{rem:V77-extends-V75}
When \(C=\{u\}\), \(D=\{v\}\), and every coordinate of \(B\) is active
on precisely \(u,v\), \eqref{eq:V77-composite-covariance} is exactly
\eqref{eq:V75-bank-covariance}.  V77 adds two features not present in
that special case: either side may already be a composite connected
cluster, and \(B\) may contain three-row and four-row coordinates.
\end{remark}

\section{The composite bank minimum}
\label{sec:V77-bank-minimum}

As in V65--V76, inactive row entries are zero and
\[
 h_{ij,e}:=a_{i,e}a_{j,e}.
\]
For the cut \(\rho=C\mid D\), define its residual crossing stock
\begin{equation}
 H_B(C,D)
 :=
 \sum_{e\in B}
 \sum_{\substack{i\in C\\j\in D}}
 h_{ij,e}.
\label{eq:V77-cross-stock}
\end{equation}

\begin{theorem}[Unpaid composite bank minimum]
\label{thm:V77-composite-minimum}
The product-state capacity of the assembled signed bank covariance
satisfies
\begin{equation}
 \boxed{
 \kappa_\otimes(\mathfrak J_B^\rho)
 \le
 C\min\{1,s_\alpha H_B(C,D)\}.}
\label{eq:V77-composite-minimum}
\end{equation}
The constant is independent of \(|B|\), the coordinate order, the
representations, and physical output multiplicities.
\end{theorem}

\begin{proof}
Both endpoint moments in
\eqref{eq:V77-composite-covariance} are complete contractions.
The contraction argument used in
\Cref{thm:V75-bank-covariance} therefore gives
\begin{equation}
 \kappa_\otimes(\mathfrak J_B^\rho)\le C.
\label{eq:V77-contraction-entry}
\end{equation}
Independently, apply the operator martingale covariance formula to the
two composite products \(F_C^B,F_D^B\), exactly as in the proof of
\Cref{lem:V37-composite-response}.  Replacing coordinate \(e\) in
\(F_C^B\) costs at most the sum of the row influences
\(\sum_{i\in C}a_{i,e}\), and the corresponding cost on the other side
is at most \(\sum_{j\in D}a_{j,e}\).  The one-link
two-difference estimate and operator Cauchy--Schwarz give
\begin{align}
 \kappa_\otimes(\mathfrak J_B^\rho)
 &\le
 C s_\alpha
 \sum_{e\in B}
 \left(\sum_{i\in C}a_{i,e}\right)
 \left(\sum_{j\in D}a_{j,e}\right)\notag\\
 &=C s_\alpha H_B(C,D).
\label{eq:V77-response-entry}
\end{align}
The equality uses the zero convention for inactive rows.  This is a
complete composite estimate: the coordinate sum is formed inside the
martingale covariance and no coordinate-word expansion or output
multiplicity count is introduced.  Taking the smaller of
\eqref{eq:V77-contraction-entry} and
\eqref{eq:V77-response-entry} proves
\eqref{eq:V77-composite-minimum}.
\end{proof}

\begin{firewall}[Unpaid means no final resolvent]
\label{fire:V77-unpaid-only}
The contraction entry in
\eqref{eq:V77-composite-minimum} concerns the unresolvented assembled
covariance.  Multiplication by an output mass and a resolvent is not a
contraction and cannot be inserted after this estimate.  V37 supplies
a paid composite response estimate, but combining it with the
contraction entry into a useful paid minimum requires a separate
one-payer argument.  No such paid minimum is asserted in V77.
\end{firewall}

\section{A cut-exposed row compiler}
\label{sec:V77-row-compiler}

For \(k\in C\), define the part of row \(k\)'s mass in \(B\) which is
exposed to the opposite block by
\begin{equation}
 X_{k\to D}(B)
 :=
 \sum_{\substack{e\in B:\ a_{k,e}>0\\
                         \sum_{j\in D}a_{j,e}>0}}
 a_{k,e}.
\label{eq:V77-exposed-row-stock}
\end{equation}
For \(k\in D\), define \(X_{k\to C}(B)\) symmetrically.

\begin{lemma}[Cross stock dominates every exposed row stock]
\label{lem:V77-cross-dominates-row}
With the group-dependent nontrivial Casimir floor \(a_*>0\) from
\Cref{lem:V76-overlap-dominates},
\begin{align}
 H_B(C,D)&\ge a_*X_{k\to D}(B),
 &&k\in C,
\label{eq:V77-cross-dominates-C}\\
 H_B(C,D)&\ge a_*X_{k\to C}(B),
 &&k\in D.
\label{eq:V77-cross-dominates-D}
\end{align}
\end{lemma}

\begin{proof}
Fix \(k\in C\).  At every coordinate counted by
\(X_{k\to D}(B)\), at least one row \(j\in D\) is nontrivial, so
\[
 \sum_{j\in D}a_{j,e}\ge a_*.
\]
The terms of \eqref{eq:V77-cross-stock} involving row \(k\) therefore
satisfy
\[
 a_{k,e}\sum_{j\in D}a_{j,e}\ge a_*a_{k,e}.
\]
Sum over the indicated coordinates and discard all other nonnegative
terms.  The proof for \(k\in D\) is identical.
\end{proof}

\begin{definition}[Bank-dominant cut exposure]
\label{def:V77-bank-dominant}
For \(0<\kappa\le1\), a row \(k\in C\) is
\(\kappa\)-bank-dominant across \(C\mid D\) if
\[
 X_{k\to D}(B)\ge\kappa\Xi_k.
\]
The definition for \(k\in D\) uses \(X_{k\to C}(B)\).
\end{definition}

\begin{corollary}[The unpaid bank minimum pays one exposed row]
\label{cor:V77-bank-pays-row}
If \(k\) is \(\kappa\)-bank-dominant across \(C\mid D\), then
\begin{equation}
 \boxed{
 \kappa_\otimes(\mathfrak J_B^\rho)
 \le
 C_\kappa\frac{H_B(C,D)}{\Xi_k}.}
\label{eq:V77-bank-pays-row}
\end{equation}
\end{corollary}

\begin{proof}
Write
\[
 \min\{1,s_\alpha H_B(C,D)\}
 =
 H_B(C,D)
 \min\{H_B(C,D)^{-1},s_\alpha\}
\]
when \(H_B(C,D)>0\).  By bank dominance,
\(\Xi_k\le\kappa^{-1}X_{k\to D}(B)\) or its symmetric analogue.
Apply
\Cref{lem:V77-cross-dominates-row,
lem:V76-scalar-compiler} and then
\Cref{thm:V77-composite-minimum}.  If \(H_B(C,D)=0\), the covariance
vanishes by \eqref{eq:V77-response-entry}.
\end{proof}

\begin{corollary}[Higher-incidence mass is admissible]
\label{cor:V77-higher-incidence}
Let \(B_{\ge3}(k;C,D)\) be the coordinates \(e\in B\) such that
\(a_{k,e}>0\), at least one row on the opposite side of the cut is
active, and the total active row set \(A_e\) has cardinality at least
three.  If
\begin{equation}
 \sum_{e\in B_{\ge3}(k;C,D)}a_{k,e}
 \ge\kappa\Xi_k,
\label{eq:V77-higher-incidence-dominance}
\end{equation}
then \eqref{eq:V77-bank-pays-row} holds.
\end{corollary}

\begin{proof}
Every coordinate counted on the left of
\eqref{eq:V77-higher-incidence-dominance} is also counted in the
cut-exposed stock of \eqref{eq:V77-exposed-row-stock}.  Thus \(k\) is
\(\kappa\)-bank-dominant and
\Cref{cor:V77-bank-pays-row} applies.
\end{proof}

\begin{assessment}[Graph-theoretic meaning]
\label{ass:V77-graph-meaning}
Expanding \(H_B(C,D)\) chooses one crossing pair
\(i\in C,j\in D\) and one coordinate.  Hence it supplies exactly one
quotient-tree edge across the cut.  Division by \(\Xi_k\) in
\eqref{eq:V77-bank-pays-row} removes one occurrence of a row mass which
would otherwise be repeated when the crossing edge is multiplied by an
earlier connected carrier.  No edge is manufactured from a coordinate
which lies wholly inside one block.
\end{assessment}

\section{Insertion into first-connectivity source families}
\label{sec:V77-source-insertion}

\begin{proposition}[One-sided tensor insertion]
\label{prop:V77-tensor-insertion}
Let \(\mathcal A_\rho\) be the complete physical envelope of the fixed
earlier carrier at partition \(\rho=C\mid D\), acting on tensor slots
disjoint from the residual bank \(B\).  If row \(k\) is
\(\kappa\)-bank-dominant and the bank is unpaid, then
\begin{equation}
 \kappa_\otimes
 \bigl(\mathcal A_\rho\otimes\mathfrak J_B^\rho\bigr)
 \le
 C_\kappa
 \|\mathcal A_\rho\|
 \frac{H_B(C,D)}{\Xi_k}.
\label{eq:V77-tensor-insertion}
\end{equation}
The same statement holds with a product-state capacity bound on the
earlier carrier and an ordinary norm bound on the bank whenever the
opposite one-sided orientation is the admissible one.
\end{proposition}

\begin{proof}
Apply the inherited one-sided tensor capacity inequality with capacity
on the bank factor and ordinary norm on
\(\mathcal A_\rho\).  Then use
\Cref{cor:V77-bank-pays-row}.  The alternative orientation is the same
inequality with the two tensor factors interchanged.
\end{proof}

\begin{corollary}[Exact two-block source-family coverage]
\label{cor:V77-source-families}
The resummation applies without changing its proof to the two V63
families whose partition immediately before the first full join has
exactly two blocks:
\begin{enumerate}[label=\textup{(\alph*)}]
\item a \(2+2\) carrier, with \(C\) and \(D\) the two connected pairs;
\item a \(3+1\) carrier, with one connected triple and one singleton.
\end{enumerate}
Three-row and four-row coordinates of \(B\) require no separate local
partition enumeration.
\end{corollary}

\begin{proof}
After conditioning on the indicated earlier carrier, each case has a
proper two-block cumulative partition \(\rho\).  A local partition
\(\pi\) fails to refine such a \(\rho\) if and only if it has a block
meeting both \(C\) and \(D\); because \(\rho\) has only two blocks, this
is equivalent to \(\rho\vee\pi=\widehat1\).  Hence the V61 joining
letter is exactly \(M_e(\widehat1)-M_e(\rho)\).  Therefore
\Cref{thm:V77-bank-telescope} applies verbatim.  The analytic estimate
depends only on the crossing influence sum
\eqref{eq:V77-cross-stock}.
\end{proof}

\begin{firewall}[Why the multi-block families are not included]
\label{fire:V77-source-not-closed}
For a prefix partition with three or four blocks,
\(M_e(\widehat1)-M_e(\rho)\) includes local partitions which coarsen
\(\rho\) without reaching \(\widehat1\).  It is therefore strictly
larger than the V61 joining letter in general.  The telescope does not
directly apply to the \(2+1+1\) or discrete first-join family.  If a
refined word first reaches a genuine two-block stage, the telescope may
be applied from that stage onward, but its eventual first-full-join
term is then classified by its actual two-block prefix.  No incompatible
source families may be merged inside one norm.
\end{firewall}

\section{What this repairs in the V76 residual}
\label{sec:V77-V76-repair}

Fix a \(2+2\) cut \(12\mid34\) and one of the large nonprefix endpoint
residuals \(Y_k^{(r)}\) from
\eqref{eq:V76-large-nonprefix-residual}.  After removing the already
controlled private and unbalanced coordinates, denote the remaining
balanced stock by \(Y_{k,\mathrm{rem}}^{(r)}\le Y_k^{(r)}\), and split
it into
\begin{equation}
 Y_{k,\mathrm{rem}}^{(r)}
 =
 Y_{k,\mathrm{cross}}^{(r)}
 +
 Y_{k,\mathrm{internal}}^{(r)},
\label{eq:V77-residual-cut-split}
\end{equation}
where the first stock is carried by residual coordinates which see the
opposite block of \(12\mid34\), and the second is carried by coordinates
which do not.

\begin{proposition}[Cut-crossing higher-incidence residuals no longer
require priority]
\label{prop:V77-cross-residual-repair}
Suppose a fixed conditioned \(12\mid34\) packet has an unpaid residual
bank and
\[
 Y_{k,\mathrm{cross}}^{(r)}\ge\kappa\Xi_k
\]
for an endpoint \(k\).  Then the entire first-joining-position sum over
that bank satisfies the one-row gain
\eqref{eq:V77-bank-pays-row}.  This remains true if all of
\(Y_{k,\mathrm{cross}}^{(r)}\) is carried by three-row or four-row
coordinates.
\end{proposition}

\begin{proof}
The crossing residual is bounded above by the appropriate exposed
stock in \eqref{eq:V77-exposed-row-stock}.  Apply
\Cref{cor:V77-bank-pays-row}.  Exact recovery of the whole
first-position sum, including its apparent suffix factors, follows from
\Cref{cor:V77-V61-packet}.  No priority permutation is used.
\end{proof}

\begin{assessment}[The sharper remaining split]
\label{ass:V77-sharper-residual}
V76 grouped all higher-incidence coordinates together.  V77 shows that
incidence number is not the decisive distinction.  The correct split
is:
\begin{enumerate}[label=\textup{(\roman*)}]
\item coordinates crossing the current two-block cut, which belong to
      the exact composite covariance packet and admit
      \eqref{eq:V77-composite-minimum};
\item coordinates internal to one current block, which must be absorbed
      into or estimated with that block's connected carrier; and
\item coordinates in a packet carrying the sole final payer, for which
      the unpaid contraction entry is unavailable until a paid minimum
      is proved.
\end{enumerate}
Thus the ordering circularity in the V76 mandate is removed, but the
one-payer bottleneck is not.
\end{assessment}

\section{Failure modes and the next upstream target}
\label{sec:V77-failure-modes}

\begin{proposition}[Why the ordinary spectator estimate does not prove
the paid bank minimum]
\label{prop:V77-no-spectator-shortcut}
The inequality
\[
 \kappa_\otimes(\mathfrak J_B^\rho)
 \le C\min\{1,s_\alpha H_B(C,D)\}
\]
and the ordinary spectator estimate
\eqref{eq:V3-spectator-payer} do not, by themselves, imply a paid bound
for the same signed covariance bank.
\end{proposition}

\begin{proof}
The spectator estimate applies to a positive complete Wilson moment
whose output coefficient contains the heat multiplier
\(q_{\alpha,U}\).  In a signed composite covariance, the output-mass
allocation must be made after the two endpoint moments have been
assembled and physically resolved.  It may land on either composite
side, and no common scalar \(q_{\alpha,U}\) can be factored from the
signed difference before that resolution.  Consequently multiplying
the unpaid capacity bound by \(C/s_\alpha\) is not authorized by
\Cref{lem:V3-spectator-payer}.  The valid paid estimate must instead be
derived from the one-payer allocation, as in
\Cref{lem:V37-composite-response}.
\end{proof}

\begin{assessment}[Candidate paid minimum, not yet a theorem]
\label{ass:V77-paid-candidate}
The two proved endpoint bounds suggest that the complete paid bank
should obey a minimum combining:
\begin{enumerate}[label=\textup{(\alph*)}]
\item the V37 composite response, which removes one factor
      \(s_\alpha\) from the martingale covariance; and
\item a high-output contraction branch obtained only after the final
      output mass is split between the two composite sides.
\end{enumerate}
The useful form must retain a factor capable of paying one
cut-exposed row.  Establishing that high-output branch without
duplicating the sole resolvent is the next upstream problem.
\end{assessment}

\section{V77 ledger and V78 mandate}
\label{sec:V77-ledger}

\begin{theorem}[Integrated compact-series V77 statement]
\label{thm:V77-integrated}
Preserving compact V1--V76, V77 proves:
\begin{enumerate}[label=\textup{(V77.\arabic*)},leftmargin=6.7em]
\item an exact telescope which resums every possible first joining
      coordinate after a fixed two-block carrier;
\item equality of that packet with one complete composite covariance;
\item order independence of the resummed packet and a proof that
      priority alone cannot expose an entire joining bank;
\item the support-uniform unpaid bound
      \(C\min\{1,s_\alpha H_B(C,D)\}\);
\item a one-row compiler for every bank-dominant cut-exposed mass; and
\item coverage of cut-crossing three-row and four-row residuals in the
      \(2+2\) and \(3+1\) families without reassigning them artificially
      to an internal prefix.
\end{enumerate}
\end{theorem}

\begin{proof}
These assertions are
\Cref{thm:V77-bank-telescope,
thm:V77-composite-covariance,
cor:V77-packet-order,
prop:V77-priority-no-go,
thm:V77-composite-minimum,
cor:V77-bank-pays-row,
cor:V77-higher-incidence,
prop:V77-cross-residual-repair}.
\end{proof}

\begin{openproblem}[V78 mandate: paid composite-bank minimum]
\label{op:V78-mandate}
\begin{enumerate}
\item Write the exact physical output decomposition of the two endpoint
      moments in \eqref{eq:V77-composite-covariance} before taking an
      envelope.
\item Apply the sole final output-mass/resolvent allocation once, between
      the two composite sides; do not insert a denominator at every
      possible first joining coordinate.
\item Combine the V37 paid response with a genuine high-output branch,
      retaining the bank covariance signs until the output split.
\item Determine whether the resulting minimum pays one
      \(\kappa\)-bank-dominant row in the sense of
      \Cref{def:V77-bank-dominant}.
\item If it does, combine that gain with the remaining V76 binary-prefix
      minimum to close the one-residual joining-payer \(2+2\) sector.
\item If it does not, construct a representation-level counterexample
      and move upstream to the exact output ancestry rather than
      repeating the unpaid argument.
\end{enumerate}
\end{openproblem}

\begin{firewall}[Claim boundary after compact V77]
\label{fire:V77-current-boundary}
V77 repairs the higher-incidence ordering step and closes the
bank-dominant cut-crossing residual whenever the resummed bank is
unpaid.  It does not prove the paid composite-bank minimum and therefore
does not yet close the joining-payer \(2+2\) residual to which V76 was
directed.  Cut-internal balanced residuals, prefix- and suffix-payer
branches, marked-heavy/overlapping banks, balanced output ancestry, and
the other source-family estimates remain open.  Global quartic and
higher MTA, WUFC, CPM, continuum Yang--Mills existence, and a positive
mass gap remain open.  The full Yang--Mills mass gap remains open.
\end{firewall}

\section{Append-only record for compact V77}
\label{sec:V77-record}

V77 was appended to the complete compact V76 file.  The exact V76
parent had \(31776\) newline-delimited source records, \(1095484\)
bytes, and SHA--256 digest
\[
\texttt{f5da5bc0e3d9a0174f7e94aa0e212ee9205539423539b36d11db910e9c7b72ed}.
\]
No compact V76 mathematical content was changed.  The sole final
\verb|\end{document}| is restored below.

% =============================================================================
% END APPENDED COMPACT VERSION V77 MATERIAL
% =============================================================================

% =============================================================================
% BEGIN APPENDED COMPACT VERSION V78 MATERIAL
% =============================================================================

\clearpage
\chapter{Paid composite-bank minima and one-residual closure}
\label{chap:V78}

\section{Context and proof strategy}
\label{sec:V78-context}

V77 resummed the possible first joining positions after a fixed
two-block carrier into the exact composite covariance
\[
 \mathfrak J_B^{C|D}
 =
 \operatorname{Cov}_B(F_C^B,F_D^B)
\]
and proved the unpaid minimum
\[
 \kappa_\otimes(\mathfrak J_B^{C|D})
 \le C\min\{1,s_\alpha H_B(C,D)\}.
\]
The contraction entry in that minimum cannot simply be multiplied by a
spectator payer.  The payer belongs to the physically resolved output
of the signed covariance itself.

This note resolves that obstruction by estimating the \emph{whole}
assembled cut covariance in two independent ways after the sole output
mass and resolvent have been inserted:
\[
 CH_B(C,D)
 \qquad\hbox{and}\qquad
 C/s_\alpha.
\]
The first is the V37 paid response.  The second follows from a complete
output first moment, not from the unpaid covariance contraction.
Their minimum is therefore legitimate.  Only after this operator
estimate is complete do we split its \emph{scalar} crossing stock into
the finitely many row pairs.  This distinction is essential: an
operator row-pair decomposition of a multicoordinate covariance need
not be localized by \(H_{ij}\), because spectator rows can correlate
through other coordinates.  The scalar refinement pays one
bank-dominant row at cost \(s_\alpha^{-1}\).  That cost is exactly
canceled by an unpaid response factor elsewhere in a quartic
first-connectivity tree.

\section{Why the operator carrier must remain assembled}
\label{sec:V78-no-pair-operator}

It is tempting to telescope the two composite row products and declare
the resulting \((i,j)\) term to be controlled by \(H_{ij}(B)\).  That
declaration is false for a multicoordinate bank.  For example, a term
selected by differences on rows \(1\) and \(3\) may still contain
undifferenced rows \(2\) and \(4\) evaluated on the same replica.
Correlations along overlaps \(14\) and \(23\) can then survive even when
\(H_{13}(B)=0\).

\begin{proposition}[No row-pair operator localization from the row
telescope]
\label{prop:V78-no-pair-operator}
The whole-row double-replica identity gives a finite algebraic
decomposition of
\(\operatorname{Cov}_B(F_C^B,F_D^B)\), but its nominal
\((i,j)\) summand is not, in general, bounded by
\(Cs_\alpha H_{ij}(B)\).  Hence such summands may not be separately
physically enveloped as row-pair response carriers.
\end{proposition}

\begin{proof}
It already fails for scalar contractions on a two-coordinate product
space.  Let \(X,Y\) be independent symmetric Rademacher variables,
attach \(X\) to rows \(1,4\) at coordinate \(e\), and attach \(Y\) to
rows \(2,3\) at coordinate \(f\).  Thus
\[
 W_1=W_4=X,\qquad W_2=W_3=Y,
\]
and no coordinate is shared by rows \(1,3\), so \(H_{13}(B)=0\).
For independent copies
\((X_U,Y_U)\), \((X_V,Y_V)\), the nominal \((1,3)\) whole-row
telescope summand is
\[
 \frac12\mathbb E\!
 \left[
 (X_U-X_V)Y_U\,(Y_U-Y_V)X_U
 \right]
 =\frac12.
\]
Indeed, after expansion only
\(\mathbb E[X_U^2Y_U^2]=1\) survives.  A bound by
\(Cs_\alpha H_{13}(B)\) would force this nonzero term to vanish.  The
example is not used as a Yang--Mills model; it proves that the proposed
localization does not follow from the double-replica algebra and
contraction alone.  Only the complete signed sum has the required
covariance cancellations.
\end{proof}

\begin{firewall}[The valid refinement is scalar]
\label{fire:V78-scalar-only}
V78 keeps the exact two-endpoint operator
\(\mathfrak J_B^{C|D}\) assembled through payment and physical
pinching.  The individual \(H_{ij}(B)\) are introduced only after a
single scalar bound for that operator has been proved.  No
\((i,j)\)-labelled signed operator is asserted.
\end{firewall}

\section{Physical payment of the complete cut covariance}
\label{sec:V78-payment}

Resolve \(\mathfrak J_B^{C|D}\) into complete orthogonal physical
output blocks \(\nu\).  Write
\[
 \Lambda_\nu:=\Xi(\boldsymbol U_\nu),
 \qquad
 Q_\nu:=q_{\alpha,\boldsymbol U_\nu}
\]
for the final output mass and its Wilson product multiplier.  Let
\(\mathfrak P_B^{C|D}\) denote the positive physical envelope obtained
by multiplying the absolute carrier in each nontrivial block once by
\begin{equation}
 \frac{\Lambda_\nu}{1-Q_\nu}.
\label{eq:V78-sole-payer}
\end{equation}
This is the sole final payer.  The definition contains no
coordinatewise denominator.

Recall the crossing stock
\[
 H_B(C,D)
 =
 \sum_{\substack{i\in C\\j\in D}}H_{ij}(B),
 \qquad
 H_{ij}(B):=\sum_{e\in B}a_{i,e}a_{j,e}.
\]

\begin{lemma}[Paid complete-cut response]
\label{lem:V78-paid-response}
\begin{equation}
 \boxed{
 \kappa_\otimes(\mathfrak P_B^{C|D})
 \le C H_B(C,D).}
\label{eq:V78-paid-response}
\end{equation}
The estimate is uniform in \(|B|\), the representations, and all
physical multiplicities.
\end{lemma}

\begin{proof}
Apply the operator martingale covariance formula directly to the
assembled covariance
\eqref{eq:V77-composite-covariance}.  Replacing coordinate \(e\) in
the \(C\)-product costs at most
\(\sum_{i\in C}a_{i,e}\), and replacing it in the \(D\)-product costs
at most \(\sum_{j\in D}a_{j,e}\).  The two-difference one-link estimate
and operator Cauchy--Schwarz therefore give the unpaid response
\[
 C s_\alpha
 \sum_{e\in B}
 \left(\sum_{i\in C}a_{i,e}\right)
 \left(\sum_{j\in D}a_{j,e}\right)
 =
 C s_\alpha H_B(C,D).
\]

Resolve the two composite outputs only after this covariance sum has
been formed.  Fusion assigns \(\Lambda_\nu\) to the sum of the two
composite masses.  The scalar one-resolvent allocation assigns
\((1-Q_\nu)^{-1}\) once to one composite side.  The binary
first-moment/gap estimate in the proof of
\Cref{lem:V37-composite-response} removes exactly the displayed factor
\(s_\alpha\).  Complete weighted pinching is performed after the
finite crossing sum and therefore introduces neither a support-size
factor nor an output-multiplicity count.  This proves
\eqref{eq:V78-paid-response}.
\end{proof}

\begin{lemma}[Output first moment of the complete cut covariance]
\label{lem:V78-output-first-moment}
For every resolved physical block \(\nu\),
\begin{equation}
 \boxed{
 s_\alpha\Lambda_\nu
 \left\|(\mathfrak J_B^{C|D})_\nu\right\|
 \le C.}
\label{eq:V78-output-first-moment}
\end{equation}
\end{lemma}

\begin{proof}
Use the two endpoint moments in
\eqref{eq:V77-composite-covariance}; do not expand the V77
first-joining-position sum.  On a fully resolved path in either
endpoint moment, all recouplings and physical compressions are
contractions and the norm is bounded by
\[
 Q_{\rm path}
 =
 \prod_{e\in B}\prod_{S\in\mathcal B_e}q_{\alpha,S},
\]
where \(\mathcal B_e\) is the local block-output family for that
endpoint moment.  Repeated fusion gives
\[
 \Lambda_\nu
 \le
 C_{\rm fus}
 \sum_{e\in B}\sum_{S\in\mathcal B_e}A_S.
\]
The complete-product first-moment compiler used in
\Cref{lem:V39-global-first-moment} gives
\[
 s_\alpha
 \left(
 \sum_{e\in B}\sum_{S\in\mathcal B_e}A_S
 \right)Q_{\rm path}
 \le C
\]
without a coordinate count.  Thus each endpoint moment obeys the
displayed first-moment estimate.  Their difference changes only the
constant.  The estimate is on a complete orthogonal multiplicity block,
so physical pinching adds no channel count.
\end{proof}

\begin{lemma}[Universal paid cap for the cut covariance]
\label{lem:V78-universal-paid-cap}
\begin{equation}
 \boxed{
 \left\|\mathfrak P_B^{C|D}\right\|
 \le\frac C{s_\alpha}.}
\label{eq:V78-universal-paid-cap}
\end{equation}
\end{lemma}

\begin{proof}
Use the V39 low/high output split.  If
\(s_\alpha\Lambda_\nu\le c_0\), the low-output product gap gives
\[
 \frac{\Lambda_\nu}{1-Q_\nu}\le\frac C{s_\alpha}.
\]
Both endpoint moments in
\eqref{eq:V77-composite-covariance} are contractions, so their
difference has uniformly bounded norm.  The paid block is therefore
at most \(C/s_\alpha\).

If \(s_\alpha\Lambda_\nu>c_0\), the product gap is bounded below by a
positive constant.  Hence
\Cref{lem:V78-output-first-moment} gives
\[
 \frac{\Lambda_\nu}{1-Q_\nu}
 \left\|(\mathfrak J_B^{C|D})_\nu\right\|
 \le\frac C{s_\alpha}.
\]
The physical blocks are orthogonal.  Their complete positive envelope
has norm equal to the supremum of the block norms, and product-state
capacity is no larger than that norm.
\end{proof}

\begin{theorem}[Paid composite-bank minimum]
\label{thm:V78-paid-cut-minimum}
\begin{equation}
 \boxed{
 \kappa_\otimes(\mathfrak P_B^{C|D})
 \le
 C\min\left\{H_B(C,D),\frac1{s_\alpha}\right\}.}
\label{eq:V78-paid-cut-minimum}
\end{equation}
\end{theorem}

\begin{proof}
Take the smaller of
\eqref{eq:V78-paid-response} and
\eqref{eq:V78-universal-paid-cap}.  Both estimates concern the same
assembled carrier with the same single payer
\eqref{eq:V78-sole-payer}.
\end{proof}

\begin{lemma}[Finite scalar row-pair refinement]
\label{lem:V78-scalar-pair-refinement}
For nonnegative \(x_1,\ldots,x_N\) and \(c>0\),
\begin{equation}
 \min\left\{\sum_{\ell=1}^Nx_\ell,c\right\}
 \le
 \sum_{\ell=1}^N\min\{x_\ell,c\}.
\label{eq:V78-scalar-pair-refinement}
\end{equation}
Consequently,
\begin{equation}
 \boxed{
 \kappa_\otimes(\mathfrak P_B^{C|D})
 \le
 C\sum_{\substack{i\in C\\j\in D}}
 \min\left\{H_{ij}(B),\frac1{s_\alpha}\right\}.}
\label{eq:V78-paid-pair-cards}
\end{equation}
The summands on the right are scalar bookkeeping cards, not separately
enveloped operators.
\end{lemma}

\begin{proof}
If \(\sum_\ell x_\ell\le c\), both sides of
\eqref{eq:V78-scalar-pair-refinement} equal or exceed
\(\sum_\ell x_\ell\).  If \(\sum_\ell x_\ell>c\), either some
\(x_\ell\ge c\), in which case the right side is at least \(c\), or
every \(x_\ell<c\), in which case the right side is
\(\sum_\ell x_\ell>c\).  Apply the scalar inequality to
\eqref{eq:V78-paid-cut-minimum} with
\(x_\ell=H_{ij}(B)\) and \(c=s_\alpha^{-1}\).
\end{proof}

\begin{corollary}[Paid binary bank minimum]
\label{cor:V78-paid-binary-minimum}
For a complete binary covariance bank \(D_{uv}\),
\begin{equation}
 \boxed{
 \max\left\{
 \left\|\mathcal P_{uv}(D_{uv})\right\|,
 \kappa_\otimes(\mathcal P_{uv}(D_{uv}))
 \right\}
 \le
 C\min\left\{H_{uv}(D_{uv}),\frac1{s_\alpha}\right\}.}
\label{eq:V78-paid-binary-minimum}
\end{equation}
\end{corollary}

\begin{proof}
For the binary cut, the V77 bank covariance is exactly
\eqref{eq:V75-bank-covariance}.  The paid binary entry of
\Cref{lem:V37-lower-table}, applied to the complete bank, gives the
ordinary-norm response bound \(CH_{uv}(D_{uv})\).  The proof of
\Cref{lem:V78-universal-paid-cap} is also an ordinary-norm proof and
applies to this binary covariance.  Take the smaller of these two norm
bounds.  Product-state capacity is no larger than ordinary norm.
\end{proof}

\begin{lemma}[Binary minimum compilers in ordinary norm]
\label{lem:V78-binary-norm-compilers}
Let \(D_{uv}\) be a complete binary bank.  If
\[
 X_{u\to v}(D_{uv})\ge\kappa\Xi_u,
\]
then
\begin{align}
 \left\|\mathcal N_{uv}(D_{uv})\right\|
 &\le
 C_\kappa\frac{H_{uv}(D_{uv})}{\Xi_u},
\label{eq:V78-unpaid-binary-norm-compiler}\\
 \left\|\mathcal P_{uv}(D_{uv})\right\|
 &\le
 C_\kappa\frac{H_{uv}(D_{uv})}
 {s_\alpha\Xi_u}.
\label{eq:V78-paid-binary-norm-compiler}
\end{align}
Both inequalities also hold with product-state capacity in place of
ordinary norm.
\end{lemma}

\begin{proof}
For the unpaid bank, the two endpoint moments in
\eqref{eq:V75-bank-covariance} give the ordinary-norm contraction
bound, while the binary entry of
\Cref{lem:V37-lower-table} gives the ordinary-norm response bound.
Thus the minimum in \eqref{eq:V75-binary-min} is valid in ordinary
norm as well as capacity.  Apply
\Cref{lem:V76-overlap-dominates,
lem:V76-scalar-compiler} to obtain
\eqref{eq:V78-unpaid-binary-norm-compiler}.

For the paid bank use the ordinary-norm statement in
\eqref{eq:V78-paid-binary-minimum} and the identity
\[
 \min\{H,s_\alpha^{-1}\}
 =
 s_\alpha^{-1}\min\{s_\alpha H,1\}.
\]
Apply the same compiler to the last minimum.  Capacity is bounded by
ordinary norm.
\end{proof}
\section{The paid minimum pays a directed row}
\label{sec:V78-directed-compiler}

For a fixed crossing pair \(i,j\), define the directed row stock
\begin{equation}
 X_{j\to i}(B)
 :=
 \sum_{\substack{e\in B\\a_{i,e}a_{j,e}>0}}a_{j,e}.
\label{eq:V78-directed-stock}
\end{equation}

\begin{lemma}[Pair stock dominates its directed row stock]
\label{lem:V78-pair-dominates}
\begin{equation}
 H_{ij}(B)\ge a_*X_{j\to i}(B),
\label{eq:V78-pair-dominates}
\end{equation}
and symmetrically
\(H_{ij}(B)\ge a_*X_{i\to j}(B)\).
\end{lemma}

\begin{proof}
At a coordinate counted by \(X_{j\to i}(B)\), row \(i\) is nontrivial
and hence \(a_{i,e}\ge a_*\).  Thus
\(a_{i,e}a_{j,e}\ge a_*a_{j,e}\).  Sum the inequality.  Interchange
\(i,j\) for the symmetric statement.
\end{proof}

\begin{corollary}[One paid scalar pair card pays one bank-dominant row]
\label{cor:V78-paid-pays-row}
If
\begin{equation}
 X_{j\to i}(B)\ge\kappa\Xi_j,
\label{eq:V78-directed-dominance}
\end{equation}
then
\begin{equation}
 \boxed{
 \min\left\{H_{ij}(B),\frac1{s_\alpha}\right\}
 \le
 C_\kappa
 \frac{H_{ij}(B)}{s_\alpha\Xi_j}.}
\label{eq:V78-paid-pays-row}
\end{equation}
The symmetric conclusion holds when row \(i\) is dominant toward
row \(j\).  For a general cut this estimates the \((i,j)\) scalar card
on the right of \eqref{eq:V78-paid-pair-cards}; it does not define a
separate operator.  For a binary cut it bounds the entire paid
covariance through \eqref{eq:V78-paid-binary-minimum}.
\end{corollary}

\begin{proof}
By \Cref{lem:V78-pair-dominates}, the hypotheses of
\Cref{lem:V76-scalar-compiler} hold with
\(X=X_{j\to i}(B)\) and \(H=H_{ij}(B)\).  Therefore
\[
 \min\{1,s_\alpha H_{ij}(B)\}
 \le
 C_\kappa\frac{H_{ij}(B)}{\Xi_j}.
\]
Divide by \(s_\alpha\) and use
\[
 \min\{H_{ij}(B),s_\alpha^{-1}\}
 =
 s_\alpha^{-1}\min\{s_\alpha H_{ij}(B),1\}
\]
This is the displayed scalar card.  The final binary-cut assertion
follows because its crossing sum has only the pair \(ij\).
\end{proof}

\begin{remark}[The \(s_\alpha^{-1}\) is necessary and useful]
\label{rem:V78-paid-s-cost}
The paid compiler differs from the unpaid compiler by exactly
\(s_\alpha^{-1}\).  This is sharp already for a low-output spectator,
by \Cref{rem:V3-payer-sharp}.  In a once-paid quartic tree it is not a
loss: one of the other two unpayed connected factors supplies its
response power \(s_\alpha\).
\end{remark}

\section{A new joining-payer residual sector}
\label{sec:V78-one-residual}

Consider a conditioned \(2+2\) carrier with internal pairs \(12\) and
\(34\), followed by its complete residual bank \(B\).  Fix a crossing
pair \(i\in\{1,2\}\), \(j\in\{3,4\}\).  After applying
\eqref{eq:V78-paid-pair-cards} to the paid residual bank, denote by
\(\mathfrak b_{22,ij}^{\rm bank\text{-}pay}\) the positive scalar
majorant obtained from its \((i,j)\) card and the ordinary norms of both
binary prefixes.  The sum of these at most four cards
majorizes the physical source envelope by the one-sided tensor
inequality.  No individual card is declared to be a signed operator.

\begin{theorem}[Complementary prefix/bank dominance closes one paid
residual]
\label{thm:V78-one-residual-22}
Assume
\begin{align}
 X_{i|i'}^{\rm pre}&\ge\kappa\Xi_i,
 &i'&=\{1,2\}\setminus\{i\},
\label{eq:V78-i-prefix-dominant}\\
 X_{j\to i}(B)&\ge\kappa\Xi_j.
\label{eq:V78-j-bank-dominant}
\end{align}
Then this resummed joining-payer majorant card obeys
\begin{equation}
 \boxed{
 \frac{\mathfrak b_{22,ij}^{\rm bank\text{-}pay}}
 {\Xi_1\Xi_2\Xi_3\Xi_4}
 \le
 C_\kappa\theta_{12}\theta_{34}\theta_{ij}.}
\label{eq:V78-one-residual-DRQC}
\end{equation}
The symmetric statement holds with \(12,i\) and \(34,j\) interchanged.
\end{theorem}

\begin{proof}
The unpaid \(12\)-prefix minimum and
\eqref{eq:V78-i-prefix-dominant} give, by
\Cref{lem:V78-binary-norm-compilers},
\[
 \left\|\mathcal N_{12}^{\rm pre}\right\|
 \le C_\kappa\frac{H_{12}^{\rm pre}}{\Xi_i}.
\]
Use the ordinary unpaid response on the other binary prefix:
\[
 \|\mathcal N_{34}^{\rm pre}\|
 \le Cs_\alpha H_{34}^{\rm pre}.
\]
Finally, the \((i,j)\) card in
\eqref{eq:V78-paid-pair-cards}, together with
\Cref{cor:V78-paid-pays-row} and
\eqref{eq:V78-j-bank-dominant} give
\[
 \min\left\{H_{ij}(B),\frac1{s_\alpha}\right\}
 \le
 C_\kappa\frac{H_{ij}(B)}{s_\alpha\Xi_j}.
\]
Use the one-sided tensor capacity bound with capacity only on the paid
bank and ordinary norms on both prefix factors.  The factors
\(s_\alpha\) and
\(s_\alpha^{-1}\) cancel, leaving
\[
 C_\kappa
 \frac{
 H_{12}^{\rm pre}H_{34}^{\rm pre}H_{ij}(B)}
 {\Xi_i\Xi_j}.
\]
The V62 ordered stock inequality bounds the sum over fixed prefix
carriers and residual banks by
\[
 C_\kappa
 \frac{H_{12}H_{34}H_{ij}}{\Xi_i\Xi_j}.
\]
After division by \(\prod_k\Xi_k\), this is exactly the path-tree mark
\(\theta_{12}\theta_{34}\theta_{ij}\).
V77 supplies the exact equality between the residual bank covariance
and the sum over all first joining positions, so no suffix bank is
discarded.
\end{proof}

\begin{corollary}[Higher-incidence paid residuals are included]
\label{cor:V78-higher-incidence-paid}
Theorem \ref{thm:V78-one-residual-22} remains valid when the stock in
\eqref{eq:V78-j-bank-dominant} is carried entirely by three-row or
four-row coordinates.
\end{corollary}

\begin{proof}
The scalar card uses only simultaneous activity of rows \(i,j\) and the
pair stock \(H_{ij}(B)\).  The assembled operator estimate preceding
the scalar split retains all other row factors.  Neither step requires
the active set of a coordinate to equal \(\{i,j\}\).
\end{proof}

\section{All payer locations in the double-prefix-dominant sector}
\label{sec:V78-all-payers}

Return to the original local first-connection presentation
\eqref{eq:V65-S22}.  For a crossing pair \(i,j\), assume both endpoints
are \(\kappa\)-prefix-dominant in their respective binary prefixes, as
in \eqref{eq:V76-double-prefix-dominant}.

\begin{proposition}[A prefix payer has the tree scale]
\label{prop:V78-prefix-payer}
If the \(12\)-prefix carries the sole payer, its \((i,j)\) component is
bounded after normalization by
\[
 C_\kappa\theta_{12}\theta_{34}\theta_{ij}.
\]
The same conclusion holds when the \(34\)-prefix pays.
\end{proposition}

\begin{proof}
For a \(12\)-prefix payer, use
\Cref{lem:V78-binary-norm-compilers} on endpoint \(i\):
\[
 \|\mathcal P_{12}^{<r}\|
 \le
 C_\kappa\frac{H_{12}^{<r}}{s_\alpha\Xi_i}.
\]
Use the unpaid entry of
\Cref{lem:V78-binary-norm-compilers} on endpoint \(j\) of the
\(34\)-prefix:
\[
 \|\mathcal N_{34}^{<r}\|
 \le C_\kappa\frac{H_{34}^{<r}}{\Xi_j}.
\]
The local joining covariance is unpaid and contributes
\(Cs_\alpha h_{ij,r}\) by
\eqref{eq:V65-local-unpaid-crossing}.  Multiplication cancels
\(s_\alpha^{-1}s_\alpha\) and yields
\[
 C_\kappa
 \frac{H_{12}^{<r}H_{34}^{<r}h_{ij,r}}
 {\Xi_i\Xi_j}.
\]
Sum \(r\) with \Cref{cor:V62-first-stock} and divide by the four row
masses.  Interchanging the two prefixes proves the other case.
\end{proof}

\begin{proposition}[A suffix payer has the tree scale]
\label{prop:V78-suffix-payer}
Under the same double-prefix-dominance hypothesis, every component in
which a complete-moment suffix coordinate carries the sole payer is
bounded after normalization by
\[
 C_\kappa\theta_{12}\theta_{34}\theta_{ij}.
\]
\end{proposition}

\begin{proof}
Use the unpaid ordinary-norm entry of
\Cref{lem:V78-binary-norm-compilers} on both binary prefixes, giving
\[
 C_\kappa
 \frac{H_{12}^{<r}H_{34}^{<r}}{\Xi_i\Xi_j}.
\]
The unpaid local joining covariance contributes
\(Cs_\alpha h_{ij,r}\).  The exact complete-product suffix payer sum is
at most \(C/s_\alpha\), by the product moment compiler and
\Cref{cor:V52-all-failure-debits}, while every other suffix moment is a
contraction.  The powers of \(s_\alpha\) cancel.  The ordered
first-stock and suffix-payer inequalities used in
\Cref{lem:V65-four-payers} give
\[
 C_\kappa
 \frac{H_{12}H_{34}H_{ij}}{\Xi_i\Xi_j}.
\]
Divide by \(\prod_k\Xi_k\).
\end{proof}

\begin{theorem}[Full payer-location closure for double-prefix-dominant
\(2+2\)]
\label{thm:V78-all-payers-prefix-dominant}
For every crossing pair \(i,j\), the complete \(2+2\)
first-connectivity source restricted by
\eqref{eq:V76-double-prefix-dominant} satisfies
\begin{equation}
 \boxed{
 \frac{
 \kappa_\otimes(
 \mathcal S_{22,ij}^{\rm prefix\text{-}dom,pay})}
 {\Xi_1\Xi_2\Xi_3\Xi_4}
 \le
 C_\kappa\theta_{12}\theta_{34}\theta_{ij},}
\label{eq:V78-all-payers-prefix-dominant}
\end{equation}
after summing all four payer locations.
\end{theorem}

\begin{proof}
The local joining-payer location is
\Cref{thm:V76-joining-payer}.  The two prefix-payer locations are
\Cref{prop:V78-prefix-payer}.  The suffix-payer location is
\Cref{prop:V78-suffix-payer}.  These are precisely the four alternatives
in \Cref{lem:V65-four-payers}, produced by the exact allocation of one
final payer.  Their finite sum proves the display.
\end{proof}

\begin{assessment}[Progress relative to V76--V77]
\label{ass:V78-progress}
Two previously distinct obstructions are now removed:
\begin{enumerate}[label=\textup{(\alph*)}]
\item the double-prefix-dominant \(2+2\) sector is controlled for every
      payer location, not merely when the local join pays; and
\item a term with one prefix-dominant endpoint and the opposite endpoint
      dominant in the paid residual bank is controlled after exact
      first-position resummation.
\end{enumerate}
The second statement is pair-directed.  If a large row residual is
spread over both opposite partners, its scalar row-pair cards must
remain separate until the matching prefix endpoint is identified.
\end{assessment}

\section{Residual cases and V79 direction}
\label{sec:V78-residual}

\begin{firewall}[What the paid minimum does not pay]
\label{fire:V78-one-row-only}
Each scalar minimum in
\eqref{eq:V78-paid-pair-cards} pays at most one factor
\(\Xi_i\) or \(\Xi_j\).  It cannot pay both endpoints of the same
crossing edge.  Thus it does not close a pair card for which both
crossing endpoints have large nonprefix residuals and neither has an
independent private, unbalanced, spectral, or prefix certificate.
\end{firewall}

\begin{assessment}[Exact remaining \(2+2\) pair alternatives]
\label{ass:V78-remaining-pairs}
After the V69--V73 certificates and
\Cref{thm:V78-all-payers-prefix-dominant,
thm:V78-one-residual-22}, a still-unclosed scalar pair card \(ij\) must
have one of the following forms:
\begin{enumerate}[label=\textup{(\roman*)}]
\item both \(i\) and \(j\) have large balanced nonprefix residuals;
\item the residual of one endpoint crosses the cut only through the
      other opposite row, whose own prefix minimum cannot pay the
      corresponding repeated mass;
\item a large residual is cut-internal rather than cut-crossing;
\item the necessary certificate overlaps a marked factor already used
      elsewhere; or
\item balanced nonprivate output ancestry carries the only unused
      payment resource.
\end{enumerate}
These are structural alternatives, not a claim that they are nonempty.
\end{assessment}

\begin{proposition}[Pigeonhole localization of a cut-crossing residual]
\label{prop:V78-cross-pigeonhole}
Let \(j\in D\).  If
\[
 X_{j\to C}(B)\ge\kappa\Xi_j,
\]
then for some \(i\in C\),
\[
 X_{j\to i}(B)\ge\frac{\kappa}{|C|}\Xi_j.
\]
\end{proposition}

\begin{proof}
Every coordinate counted by \(X_{j\to C}(B)\) has at least one active
partner in \(C\).  Assign it deterministically to the least such
partner.  The assigned directed stocks are disjoint, their sum is
\(X_{j\to C}(B)\), and there are \(|C|\) of them.  One has at least the
average.
\end{proof}

\begin{remark}[Why pigeonhole is not yet exhaustion]
\label{rem:V78-pigeonhole-limit}
The partner \(i\) produced by
\Cref{prop:V78-cross-pigeonhole} need not be prefix-dominant.  Replacing
it by a different row would also replace the crossing tree edge and is
not allowed.  V79 must therefore analyze the two-endpoint residual
graph jointly, rather than applying independent row pigeonholes.
\end{remark}

\section{V78 ledger and V79 mandate}
\label{sec:V78-ledger}

\begin{theorem}[Integrated compact-series V78 statement]
\label{thm:V78-integrated}
Preserving compact V1--V77, V78 proves:
\begin{enumerate}[label=\textup{(V78.\arabic*)},leftmargin=6.7em]
\item a proof that whole-row telescoping does not justify a row-pair
      operator localization in a multicoordinate bank;
\item a support-uniform output first moment for the complete assembled
      cut covariance;
\item the paid composite-bank minimum
      \(C\min\{H_B(C,D),s_\alpha^{-1}\}\), with the sole resolvent used
      once, and its finite scalar row-pair refinement;
\item the corresponding paid binary-prefix minimum;
\item a directed one-row paid compiler;
\item DRQC for one-prefix/one-residual complementary dominance,
      including higher-incidence residual banks; and
\item all four payer locations in the double-prefix-dominant \(2+2\)
      sector.
\end{enumerate}
\end{theorem}

\begin{proof}
These are
\Cref{prop:V78-no-pair-operator,
lem:V78-output-first-moment,
thm:V78-paid-cut-minimum,
lem:V78-scalar-pair-refinement,
cor:V78-paid-binary-minimum,
lem:V78-binary-norm-compilers,
cor:V78-paid-pays-row,
thm:V78-one-residual-22,
cor:V78-higher-incidence-paid,
thm:V78-all-payers-prefix-dominant}.
\end{proof}

\begin{openproblem}[V79 mandate: the two-endpoint residual graph]
\label{op:V79-mandate}
\begin{enumerate}
\item For each crossing pair \(ij\), form the exact disjoint allocation
      of the two endpoint masses among own-prefix, directed cross-bank,
      cut-internal, private, unbalanced, and balanced output-ancestry
      stocks.
\item Use a finite bipartite matching table, not two independent
      pigeonholes, to determine which scalar pair cards are covered by
      complementary dominance.
\item For every unmatched card, prove that both endpoints carry a
      fixed large balanced residual in specified disjoint coordinate
      banks.
\item Test whether the two paid/unpaid minima can be oriented on
      different factors to pay those two residuals without using the
      sole payer twice.
\item If a component remains unmatched, move upstream to its exact
      balanced output-ancestry graph and retain the first unused
      coordinate mark.
\item Keep the \(3+1\), \(2+1+1\), and discrete families separate; the
      V77 telescope is only a two-block identity.
\end{enumerate}
\end{openproblem}

\begin{firewall}[Claim boundary after compact V78]
\label{fire:V78-current-boundary}
V78 proves the paid assembled cut-bank and binary-bank minima, together
with a scalar row-pair refinement, and closes the
specified double-prefix and complementary one-residual \(2+2\)
sectors.  It does not close scalar pair cards with two unresolved
balanced endpoints, cut-internal or overlapping balanced residuals, or
balanced output ancestry.  Nor does it complete the remaining
\(3+1\), \(2+1+1\), and discrete hot source families.  Global quartic
and higher MTA, WUFC, CPM, continuum Yang--Mills existence, and a
positive mass gap remain open.  The full Yang--Mills mass gap remains
open.
\end{firewall}

\section{Append-only record for compact V78}
\label{sec:V78-record}

V78 was appended to the complete compact V77 file.  The exact V77
parent had \(32504\) newline-delimited source records, \(1122398\)
bytes, and SHA--256 digest
\[
\texttt{9cc550091efa82e0b3485fd11c1564c90d5c8d84073139d4100328f1ca1bf441}.
\]
No compact V77 mathematical content was changed.  The sole final
\verb|\end{document}| is restored below.

% =============================================================================
% END APPENDED COMPACT VERSION V78 MATERIAL
% =============================================================================

% =============================================================================
% BEGIN APPENDED COMPACT VERSION V79 MATERIAL
% =============================================================================

\clearpage
\chapter{The two-endpoint matching ledger}
\label{chap:V79}

\section{Context: replacing two independent pigeonholes}
\label{sec:V79-context}

V78 proved a paid minimum for the complete residual covariance bank and
only then refined its scalar crossing stock into the four possible
row-pair cards.  A card \(ij\) has two repeated endpoint masses.  One
may be paid by an internal binary-prefix minimum and the other by the
paid residual-bank minimum.  The two choices cannot be made
independently: the residual bank must pay the endpoint on the actual
crossing card, while the prefix minimum must pay the other endpoint.

This note constructs an exact disjoint incidence ledger for a fixed
\(2+2\) carrier, gives the complete finite matching table, and proves a
column-grouped paid-bank theorem.  The grouped theorem repairs the
\emph{wrong-partner} failure of an individual pigeonhole whenever both
rows on the opposite side have dominant binary prefixes.

The analysis also corrects one assertion in the V79 mandate.  An
unmatched card need not leave large residuals at both of its own
endpoints.  It may instead create a three-vertex obstruction: one
endpoint has a good prefix, the other endpoint's large crossing mass
points to the good endpoint's partner, and that partner has a poor
prefix.  An explicit finite incidence table proves this.  The surviving
mass is then localized to balanced foreign-prefix or cut-internal
coordinates unless a previously proved certificate applies.

\section{A disjoint row-incidence partition}
\label{sec:V79-incidence-partition}

Fix a coordinate-resolved \(2+2\) carrier with cut
\[
 C=\{1,2\},\qquad D=\{3,4\}.
\]
Fix its earlier two-block carrier and let \(B\) be the complete residual
bank of V77.  For \(k\in C\), write \(k^\circ\) for the other row of
\(C\); for \(k\in D\), let \(k^\circ\) be the other row of \(D\).
Let \(k^c\) denote the opposite block.

We partition \emph{row-coordinate incidences} \((k,e)\), not merely
coordinates.  A coordinate active on several rows may therefore have
different labels on different rows, but each fixed incidence receives
exactly one label.  Use the following deterministic priority.

\begin{enumerate}[label=\textup{(\roman*)}]
\item If \(e\) is private to row \(k\), label \((k,e)\) private.
\item Otherwise, if \(e\) is unbalanced in the V69 sense, label
      \((k,e)\) unbalanced.
\item Assume now that \(e\) is balanced.  If \(e\) belongs to the fixed
      earlier carrier and row \(k^\circ\) is active at \(e\), label the
      incidence own-prefix.
\item If \(e\) belongs to the earlier carrier but \(k^\circ\) is
      inactive, label the incidence foreign-prefix.
\item If \(e\in B\) and at least one row of \(k^c\) is active, assign
      \((k,e)\) to the least active row
      \(\ell\in k^c\) in the fixed row order and label it
      \(k\to\ell\).
\item If \(e\in B\) and no row of \(k^c\) is active, label the incidence
      cut-internal.
\end{enumerate}

Define the corresponding nonnegative row stocks
\begin{align}
 P_k&:=\sum_{(k,e)\ {\rm private}}a_{k,e},
 &
 U_k&:=\sum_{(k,e)\ {\rm unbalanced}}a_{k,e},
\label{eq:V79-PU}\\
 A_k&:=\sum_{(k,e)\ {\rm own\text{-}prefix}}a_{k,e},
 &
 F_k&:=\sum_{(k,e)\ {\rm foreign\text{-}prefix}}a_{k,e},
\label{eq:V79-AF}\\
 D_{k\to\ell}
 &:=\sum_{(k,e)\ {\rm assigned}\ k\to\ell}a_{k,e},
 &
 D_k&:=\sum_{\ell\in k^c}D_{k\to\ell},
\label{eq:V79-directed-D}\\
 I_k&:=\sum_{(k,e)\ {\rm cut\text{-}internal}}a_{k,e}.
\label{eq:V79-I}
\end{align}

\begin{lemma}[Exact disjoint incidence identity]
\label{lem:V79-incidence-identity}
For every row \(k\),
\begin{equation}
 \boxed{
 \Xi_k
 =
 P_k+U_k+A_k+F_k+D_k+I_k.}
\label{eq:V79-incidence-identity}
\end{equation}
Moreover, for every \(\ell\in k^c\),
\begin{equation}
 D_{k\to\ell}
 \le X_{k\to\ell}(B),
 \qquad
 D_k\le X_{k\to k^c}(B).
\label{eq:V79-assigned-below-actual}
\end{equation}
\end{lemma}

\begin{proof}
Private, unbalanced shared, and balanced shared incidences are disjoint
and exhaustive.  The fixed earlier-carrier coordinate set and the
residual set \(B\) are disjoint and exhaustive.  On a balanced earlier
incidence, the own-block partner is either active or inactive.  On a
balanced residual incidence, the opposite active set is either empty
or nonempty; in the nonempty case the least-partner rule makes exactly
one assignment.  Thus the six displayed classes partition every
nonzero incidence \((k,e)\).  Summing its weight \(a_{k,e}\) proves
\eqref{eq:V79-incidence-identity}.  An assigned incidence has the
simultaneous activity required by the corresponding V78 directed
stock, which proves \eqref{eq:V79-assigned-below-actual}.
\end{proof}

\begin{firewall}[Foreign prefix is not residual bank mass]
\label{fire:V79-foreign-prefix}
A balanced coordinate in the earlier carrier may be active on row
\(k\) and on an opposite-block row while the local partition letter
still refines \(C|D\).  If \(k^\circ\) is inactive, its row-\(k\)
incidence is not counted by the own binary-prefix overlap and it is
also not in the residual bank \(B\).  This is the foreign-prefix stock
\(F_k\).  Merging it into \(A_k\) or \(D_k\) would make either the V76
prefix compiler or the V78 paid-bank compiler circular.
\end{firewall}

\begin{lemma}[Finite dominant-cell alternative]
\label{lem:V79-dominant-cell}
Fix \(0<\eta<1/4\).  If
\begin{equation}
 P_k<\eta\Xi_k,\quad
 U_k<\eta\Xi_k,\quad
 A_k<\eta\Xi_k,\quad
 D_k<\eta\Xi_k,
\label{eq:V79-four-small}
\end{equation}
then
\begin{equation}
 \boxed{
 F_k+I_k>(1-4\eta)\Xi_k.}
\label{eq:V79-FI-large}
\end{equation}
In particular, either
\[
 F_k>\frac{1-4\eta}{2}\Xi_k
 \quad\hbox{or}\quad
 I_k>\frac{1-4\eta}{2}\Xi_k.
\]
\end{lemma}

\begin{proof}
Subtract the four strict inequalities in
\eqref{eq:V79-four-small} from the exact identity
\eqref{eq:V79-incidence-identity}.  The final assertion is the
two-cell pigeonhole principle.
\end{proof}

\begin{assessment}[Previously certified cells]
\label{ass:V79-certified-cells}
On a hot endpoint branch, a fixed-fraction private cell is controlled
by the complete private Wilson product, and a fixed-fraction
unbalanced cell is controlled by the V69 reserve/heavy alternative.
The constants in those fixed-degree moment arguments may depend on
\(\eta\), but not on support or representations.  Hence after the
V68--V69 certified branches are removed, failure of both own-prefix and
residual-cross dominance forces the quantitative
foreign-prefix/cut-internal alternative
\eqref{eq:V79-FI-large}.  This assessment invokes only the already
proved sourcewise certificates; it does not declare \(F_k\) or \(I_k\)
controlled.
\end{assessment}

\section{The exact two-endpoint matching table}
\label{sec:V79-matching-table}

Call row \(k\) prefix-good if
\begin{equation}
 A_k\ge\eta\Xi_k.
\label{eq:V79-prefix-good}
\end{equation}
Call the directed incidence \(k\to\ell\) bank-good if
\begin{equation}
 D_{k\to\ell}\ge\frac{\eta}{2}\Xi_k.
\label{eq:V79-directed-good}
\end{equation}
Call row \(k\) balanced-cross-good if
\begin{equation}
 D_k\ge\eta\Xi_k.
\label{eq:V79-cross-good}
\end{equation}
Both are sufficient, rather than necessary, conditions for the V76 and
V78 compilers because of
\eqref{eq:V79-assigned-below-actual}.

Fix a paid residual-bank scalar card \(ij\), with \(i\in C,j\in D\).
Call it directly matched if at least one of
\begin{align}
 &i\ {\rm prefix\text{-}good}
 \quad\hbox{and}\quad
 j\ {\rm prefix\text{-}good},
\label{eq:V79-match-PP}\\
 &i\ {\rm prefix\text{-}good}
 \quad\hbox{and}\quad
 j\to i\ {\rm bank\text{-}good},
\label{eq:V79-match-PD}\\
 &j\ {\rm prefix\text{-}good}
 \quad\hbox{and}\quad
 i\to j\ {\rm bank\text{-}good}
\label{eq:V79-match-DP}
\end{align}
holds.

\begin{theorem}[Complete direct-matching truth table]
\label{thm:V79-matching-table}
The direct status of the card \(ij\) is exactly:
\begin{center}
\begin{tabular}{c c c}
\toprule
\(i\) prefix-good & \(j\) prefix-good & condition for direct match\\
\midrule
yes & yes & automatic\\
yes & no  & \(j\to i\) bank-good\\
no  & yes & \(i\to j\) bank-good\\
no  & no  & impossible\\
\bottomrule
\end{tabular}
\end{center}
Every directly matched card satisfies the normalized path-tree bound
of \eqref{eq:V78-one-residual-DRQC}.
\end{theorem}

\begin{proof}
The first three rows are exactly
\eqref{eq:V79-match-PP}--%
\eqref{eq:V79-match-DP}; when neither endpoint is prefix-good none of
those conditions can hold.  In the first row use
\Cref{thm:V78-all-payers-prefix-dominant}.  In the second row use
\Cref{thm:V78-one-residual-22} with endpoint \(i\) paid by its prefix
and endpoint \(j\) paid by the directed bank card.  The third row is
the symmetric V78 theorem.
\end{proof}

\begin{proposition}[The two-endpoint-residual claim is false]
\label{prop:V79-not-two-endpoints}
An unmatched card need not have a large nonprefix residual at both of
its endpoints.
\end{proposition}

\begin{proof}
Let \(i'\) be the other row of \(C\).  Consider the finite incidence
table with \(0<\varepsilon<\eta/4\):
\[
 A_i=(1-\varepsilon)\Xi_i,\qquad
 D_{i\to j}=\varepsilon\Xi_i,
\]
and
\[
 A_j=0,\qquad
 D_{j\to i}=\varepsilon\Xi_j,\qquad
 D_{j\to i'}=(1-\varepsilon)\Xi_j,
\]
with all other cells of rows \(i,j\) zero.
It obeys the exact row identities and can be realized at the support
shared predominantly by \(j,i'\), together with one small residual
coordinate shared by \(i,j\).  The latter makes the \(ij\) crossing
card nonzero.  The card is unmatched by
\Cref{thm:V79-matching-table}: row \(j\) is not prefix-good and its
directed mass toward \(i\) is below the bank-good threshold.
Row \(i\)'s entire nonprefix residual is only
\(\varepsilon\Xi_i\), which can be made arbitrarily small.  Thus item 3
of the V79 mandate was too strong.
\end{proof}

\begin{firewall}[This is a matching counterexample, not a Yang--Mills
counterexample]
\label{fire:V79-not-YM-counterexample}
\Cref{prop:V79-not-two-endpoints} refutes only the proposed
combinatorial conclusion about where residual row masses must lie.  It
does not show that the physical source card diverges.  Indeed the next
section proves that the wrong-partner configuration closes once the
other row \(i'\) is also prefix-good.
\end{firewall}

\section{Grouping the paid minimum by an endpoint}
\label{sec:V79-column-grouping}

For \(q\in D\), put
\begin{equation}
 H_{Cq}(B):=\sum_{p\in C}H_{pq}(B),
\qquad
 X_{q\to C}(B):=
 \sum_{\substack{e\in B:\ a_{q,e}>0\\
                         \sum_{p\in C}a_{p,e}>0}}a_{q,e}.
\label{eq:V79-column-stocks}
\end{equation}

\begin{lemma}[Column minimum]
\label{lem:V79-column-minimum}
For \(C=\{p,p'\}\),
\begin{equation}
 \sum_{r\in C}
 \min\left\{H_{rq}(B),\frac1{s_\alpha}\right\}
 \le
 2\min\left\{H_{Cq}(B),\frac1{s_\alpha}\right\}.
\label{eq:V79-column-minimum}
\end{equation}
If \(X_{q\to C}(B)\ge\eta\Xi_q\), then
\begin{equation}
 \boxed{
 \sum_{r\in C}
 \min\left\{H_{rq}(B),\frac1{s_\alpha}\right\}
 \le
 C_\eta
 \frac{H_{Cq}(B)}{s_\alpha\Xi_q}.}
\label{eq:V79-column-pays-q}
\end{equation}
\end{lemma}

\begin{proof}
Each \(H_{rq}(B)\le H_{Cq}(B)\), so each minimum on the left of
\eqref{eq:V79-column-minimum} is at most the minimum on the right;
there are two rows.  At every coordinate counted by
\(X_{q\to C}(B)\), some nontrivial row in \(C\) has mass at least
\(a_*\).  Hence
\[
 H_{Cq}(B)\ge a_*X_{q\to C}(B).
\]
Apply \Cref{lem:V76-scalar-compiler} to
\(\min\{1,s_\alpha H_{Cq}(B)\}\), divide by \(s_\alpha\), and combine
with \eqref{eq:V79-column-minimum}.
\end{proof}

\begin{theorem}[Prefix-saturated side closes a cross-dominant column]
\label{thm:V79-column-closure}
Suppose both rows \(p,p'\in C\) are prefix-good.  Fix \(q\in D\) and
assume
\begin{equation}
 X_{q\to C}(B)\ge\eta\Xi_q.
\label{eq:V79-q-cross-dominant}
\end{equation}
Then the sum of the two \(pq,p'q\) scalar cards in the paid residual
bank satisfies
\begin{equation}
 \boxed{
 \frac{\mathfrak b_{Cq}^{\rm bank\text{-}pay}}
 {\Xi_1\Xi_2\Xi_3\Xi_4}
 \le
 C_\eta\,
 \theta_{12}\theta_{34}
 \sum_{r\in C}\theta_{rq}.}
\label{eq:V79-column-DRQC}
\end{equation}
Here \(\mathfrak b_{Cq}^{\rm bank\text{-}pay}\) is the corresponding
positive scalar majorant after applying
\eqref{eq:V78-paid-pair-cards}.  The symmetric row-grouped statement
holds with \(C,D\) interchanged.
\end{theorem}

\begin{proof}
Use the unpaid response bound on the binary prefix in \(D\):
\[
 \|\mathcal N_D^{\rm pre}\|
 \le Cs_\alpha H_{34}^{\rm pre}.
\]
By \Cref{lem:V79-column-minimum}, the \(q\)-column of the paid-bank
majorant is at most
\[
 C_\eta\frac{H_{Cq}(B)}{s_\alpha\Xi_q}
 =
 \sum_{r\in C}
 C_\eta\frac{H_{rq}(B)}{s_\alpha\Xi_q}.
\]
Distribute this finite scalar sum.  Since both \(r=p,p'\) are
prefix-good, the unpaid ordinary-norm binary compiler
\eqref{eq:V78-unpaid-binary-norm-compiler} may be oriented separately
on row \(r\) in the \(r\)-th scalar term:
\[
 \|\mathcal N_C^{\rm pre}\|
 \le
 C_\eta\frac{H_{12}^{\rm pre}}{\Xi_r}.
\]
Use capacity only on the paid residual bank and ordinary norms on the
two prefixes.  The factors \(s_\alpha\) cancel.  The \(r\)-th term is
bounded by
\[
 C_\eta
 \frac{
 H_{12}^{\rm pre}H_{34}^{\rm pre}H_{rq}(B)}
 {\Xi_r\Xi_q}.
\]
The V62 ordered stock inequality permits replacement by the
corresponding global stocks after summing fixed carriers.  Division by
\(\prod_k\Xi_k\) gives
\(\theta_{12}\theta_{34}\theta_{rq}\).  Sum the two values of \(r\).
\end{proof}

\begin{corollary}[One-side saturation closes the entire paid bank]
\label{cor:V79-side-saturation}
If both rows in \(C\) are prefix-good and both rows \(q\in D\) satisfy
\eqref{eq:V79-q-cross-dominant}, then every scalar card in the paid
residual-bank majorant satisfies the quartic tree atlas after summing
the two columns.  The same holds with \(C,D\) interchanged.
\end{corollary}

\begin{proof}
Apply \Cref{thm:V79-column-closure} to the two rows of \(D\).  The four
resulting monomials are precisely the four path trees
\(\{12,34,rq\}\).
\end{proof}

\section{The three-vertex wrong-partner obstruction}
\label{sec:V79-wrong-partner}

\begin{proposition}[Wrong partner under an untyped cross hypothesis]
\label{prop:V79-wrong-partner}
Let \(i,i'\) be the two rows of \(C\), and let \(j\in D\).  Assume:
\begin{enumerate}[label=\textup{(\alph*)}]
\item \(i\) is prefix-good;
\item \(X_{j\to C}(B)\ge\eta\Xi_j\);
\item \(D_{j\to i}<\eta\Xi_j/2\); and
\item the \(j\)-column has not been closed by
      \Cref{thm:V79-column-closure}.
\end{enumerate}
These assumptions do not imply
\begin{equation}
 D_{j\to i'}\ge\frac{\eta}{2}\Xi_j.
\label{eq:V79-wrong-partner-large}
\end{equation}
\end{proposition}

\begin{proof}
By construction,
\[
 D_j=D_{j\to i}+D_{j\to i'}.
\]
However, the untyped V78 stock \(X_{j\to C}(B)\) also counts residual
crossing incidences which were classified unbalanced before the
balanced assignment \(D_j\).  Thus
\(D_j\le X_{j\to C}(B)\), with no reverse inequality.  The lower bound
on \(X_{j\to C}(B)\) therefore gives no lower bound on \(D_j\), and
\eqref{eq:V79-wrong-partner-large} does not follow.  Explicitly, take
all row-\(j\) residual crossing mass in an unbalanced bank shared with
\(i'\), so that
\[
 U_j=X_{j\to C}(B)=\Xi_j,
 \qquad D_j=0.
\]
Take \(i\) prefix-good and \(i'\) prefix-poor.  Then
\textup{(a)}--\textup{(d)} hold while the left side of
\eqref{eq:V79-wrong-partner-large} is zero.
\end{proof}

\begin{assessment}[The first matching regression]
\label{ass:V79-matching-regression}
\Cref{prop:V79-wrong-partner} records a failed inference rather than a
positive closure theorem.  A theorem about the assigned balanced wrong
partner must assume
\[
 D_j\ge\eta\Xi_j,
\]
not merely \(X_{j\to C}(B)\ge\eta\Xi_j\).  With that corrected
hypothesis,
\[
 D_{j\to i'}=D_j-D_{j\to i}\ge\eta\Xi_j/2.
\]
If \(i'\) were prefix-good,
\Cref{thm:V79-column-closure} would close the column because
\eqref{eq:V79-assigned-below-actual} gives
\(X_{j\to C}(B)\ge D_j\).  Hence a surviving wrong-partner state has
the three-vertex form
\[
 i\ {\rm prefix\text{-}good},
 \qquad
 j\longrightarrow i'\ {\rm balanced\ bank\text{-}good},
 \qquad
 i'\ {\rm prefix\text{-}poor}.
\]
Replacing an assigned balanced stock by the larger untyped cross stock
would undo the disjoint ledger.
\end{assessment}

\begin{corollary}[Corrected wrong-partner theorem]
\label{cor:V79-corrected-wrong-partner}
In \Cref{prop:V79-wrong-partner}, replace assumption \textup{(b)} by
\[
 D_j\ge\eta\Xi_j.
\]
Then \eqref{eq:V79-wrong-partner-large} holds and \(i'\) is not
prefix-good.
\end{corollary}

\begin{proof}
This is the calculation in
\Cref{ass:V79-matching-regression}.
\end{proof}

\section{Balanced output ancestry uses the same finite cells}
\label{sec:V79-output-ancestry}

The balanced shared row stock \(B_k\) of V69 is, in the present
incidence partition,
\begin{equation}
 B_k=A_k+F_k+D_k+I_k.
\label{eq:V79-balanced-stock-cells}
\end{equation}
Private and unbalanced incidences are absent by definition.

\begin{lemma}[Sixteen-cell output-ancestry localization]
\label{lem:V79-output-cells}
In every physical output block, the V69 balanced nonprivate output
majorant obeys
\begin{equation}
 \boxed{
 Z_\nu^B
 \le
 C\sum_{k=1}^4(A_k+F_k+D_k+I_k).}
\label{eq:V79-output-cells}
\end{equation}
Consequently its positive majorant is the sum of at most sixteen
row-labelled incidence cells.  If the majorant is nonzero, one cell is
at least \(1/16\) of that displayed majorant.
\end{lemma}

\begin{proof}
By \eqref{eq:V69-output-BU},
\[
 Z_\nu^B=C\sum_k B_k.
\]
Insert \eqref{eq:V79-balanced-stock-cells}.  There are four rows and
four cells per row.  The last assertion is the finite pigeonhole
principle applied to this positive majorant.
\end{proof}

\begin{firewall}[Output ancestry is not a second copy of row mass]
\label{fire:V79-output-not-input}
Equation \eqref{eq:V79-output-cells} is a Clebsch--Gordan upper
majorant for the scalar output mass.  It is not added to
\eqref{eq:V79-incidence-identity}.  When an ancestry cell is used, its
underlying row-coordinate stock and its physical multiplier must be
used once in the same source packet.  Treating \(Z_\nu^B\) as an
additional disjoint input bank would double-count the incidence.
\end{firewall}

\begin{assessment}[Ancestry cells already routed and cells still open]
\label{ass:V79-output-routing}
An \(A_k\)-labelled ancestry cell is attached to an own-prefix bank and
a \(D_k\)-labelled cell to a residual crossing bank.  The V78--V79
minima provide the correct candidate routes for those cells, subject
to the sole-payer and matching conditions already stated.  The
\(F_k\)- and \(I_k\)-labelled ancestry cells have no such theorem.
Thus balanced output ancestry is no longer an undifferentiated scalar:
its unreached part is localized to the same foreign-prefix and
cut-internal banks that occur in
\eqref{eq:V79-FI-large}.
\end{assessment}

\section{Post-V79 reduction}
\label{sec:V79-reduction}

\begin{theorem}[Exact residual after the finite matching ledger]
\label{thm:V79-residual}
After applying the V68--V69 private/unbalanced certificates, the V78
direct matches, and the V79 prefix-saturated column theorem, every
still-unproved paid-residual-bank \(2+2\) card has at least one of:
\begin{enumerate}[label=\textup{(\roman*)}]
\item a three-vertex wrong-partner state from
      \Cref{cor:V79-corrected-wrong-partner};
\item two prefix-poor endpoints whose balanced residual crossing stocks
      are both large, while the sole paid minimum can pay only one;
\item an endpoint with a fixed large \(F_k+I_k\) stock as in
      \eqref{eq:V79-FI-large};
\item a marked-heavy or overlapping balanced bank excluded from the
      clean V71--V73 compilers; or
\item a balanced output-ancestry cell labelled by \(F_k\) or \(I_k\).
\end{enumerate}
The alternatives may overlap, but every stock named inside one
alternative is a cell of the disjoint incidence partition.
\end{theorem}

\begin{proof}
For a fixed card, use
\Cref{thm:V79-matching-table}.  A directly matched card closes by V78.
If neither endpoint is prefix-good but both have large assigned
balanced crossing stocks, item \textup{(ii)} occurs.
If a prefix-good endpoint faces a large assigned balanced cross stock
pointing to the other partner, either the opposite side is
prefix-saturated and
\Cref{thm:V79-column-closure} closes the column, or
\Cref{cor:V79-corrected-wrong-partner} gives item \textup{(i)}.

For every endpoint which is neither prefix-good nor
balanced-cross-good, the exact identity
\eqref{eq:V79-incidence-identity} applies.  Private and unbalanced
dominance was already removed.  Thus
\Cref{lem:V79-dominant-cell} gives item \textup{(iii)}, except when the
needed bank is one of the previously recorded marked/overlap
exceptions, which is item \textup{(iv)}.  Finally
\Cref{lem:V79-output-cells} and
\Cref{ass:V79-output-routing} give item \textup{(v)} for any
unrouted balanced ancestry majorant.
\end{proof}

\begin{assessment}[The next upstream object]
\label{ass:V79-next-object}
The matching problem is no longer a free four-vertex graph search.
Apart from the explicit two-cross/no-prefix state in
\Cref{thm:V79-residual}, every survivor contains a prefix-poor row with
a concrete balanced bank of one of two types:
\begin{enumerate}[label=\textup{(\alph*)}]
\item a foreign-prefix bank, occurring before first connectivity but
      outside that row's own binary overlap; or
\item a cut-internal residual bank, occurring after the fixed carrier
      but not crossing its cut.
\end{enumerate}
Neither bank can be paid by the existing crossing minima.  The
two-cross/no-prefix state instead asks for two independent gains from
one residual covariance bank and is the two-puncture obstruction in a
new form.  V80 should first perform the mandatory companion audit and
then derive the exact
conditional covariance identity which moves one such internal bank
into its own block carrier without destroying the V77 residual-bank
telescope.
\end{assessment}

\section{V79 ledger and V80 mandate}
\label{sec:V79-ledger}

\begin{theorem}[Integrated compact-series V79 statement]
\label{thm:V79-integrated}
Preserving compact V1--V78, V79 proves:
\begin{enumerate}[label=\textup{(V79.\arabic*)},leftmargin=6.7em]
\item an exact six-cell disjoint partition of every row's coordinate
      mass relative to a fixed \(2+2\) carrier;
\item the complete direct-matching truth table for paid residual-bank
      scalar cards;
\item a counterexample to the proposed claim that every unmatched card
      leaves large residuals at both of its own endpoints;
\item a column-grouped paid-minimum theorem which closes wrong-partner
      cards whenever the opposite side is prefix-saturated;
\item the corrected three-vertex wrong-partner obstruction;
\item a sixteen-cell localization of balanced output ancestry without
      duplicating input mass; and
\item reduction of the residual to foreign-prefix, cut-internal,
      two-cross/no-prefix, marked/overlap, or correspondingly labelled
      ancestry banks.
\end{enumerate}
\end{theorem}

\begin{proof}
These are
\Cref{lem:V79-incidence-identity,
thm:V79-matching-table,
prop:V79-not-two-endpoints,
thm:V79-column-closure,
cor:V79-side-saturation,
cor:V79-corrected-wrong-partner,
lem:V79-output-cells,
thm:V79-residual}.
\end{proof}

\begin{openproblem}[V80 mandate: audit and internal-bank conditioning]
\label{op:V80-mandate}
\begin{enumerate}
\item Perform the mandatory read-only audit of the newest SM and NS
      notes and record exact fingerprints.
\item For a foreign-prefix cell \(F_k\), retain its actual local
      partition letter and determine which connected carrier of the
      current block contains its Wilson moment.
\item For a cut-internal residual cell \(I_k\), condition the V77
      composite covariance on that internal coordinate bank and prove
      an exact tower/covariance identity before taking a norm.
\item Test whether the conditioned internal bank has an unpaid or paid
      binary minimum capable of paying the prefix-poor row.
\item Keep marked-heavy and overlapping banks separate unless the
      conditional identity proves that their physical multipliers are
      used only once.
\item Insert the \(F_k\)- and \(I_k\)-labelled output-ancestry cells
      only after the corresponding input-bank identity is proved; do
      not count the ancestry majorant as new mass.
\end{enumerate}
\end{openproblem}

\begin{firewall}[Claim boundary after compact V79]
\label{fire:V79-current-boundary}
V79 proves the disjoint endpoint ledger, the exact matching table, and
the prefix-saturated column closure.  It does not control the remaining
two-cross/no-prefix state, foreign-prefix or cut-internal balanced
banks, their labelled output ancestry, or the
marked-heavy/overlapping exceptions.  It also does not
complete the hot \(3+1\), \(2+1+1\), or discrete source families.
Global quartic and higher MTA, WUFC, CPM, continuum Yang--Mills
existence, and a positive mass gap remain open.  The full Yang--Mills
mass gap remains open.
\end{firewall}

\section{Append-only record for compact V79}
\label{sec:V79-record}

V79 was appended to the complete compact V78 file.  The exact V78
parent had \(33327\) newline-delimited source records, \(1150336\)
bytes, and SHA--256 digest
\[
\texttt{176cba166a5cb6efd9fb9f641ab71ad92511240efbf1e2c904cb8150c0467b5a}.
\]
No compact V78 mathematical content was changed.  The sole final
\verb|\end{document}| is restored below.

% =============================================================================
% END APPENDED COMPACT VERSION V79 MATERIAL
% =============================================================================

% =============================================================================
% BEGIN APPENDED COMPACT VERSION V80 MATERIAL
% =============================================================================

\clearpage
\chapter{Foreign-prefix debits and conditional internal-bank
factorization}
\label{chap:V80}

\section{Context and acquisition order}
\label{sec:V80-context}

V79 replaced the informal two-endpoint pigeonhole by an exact six-cell
row-incidence partition relative to a fixed \(2+2\) carrier.  Once the
private, unbalanced, own-prefix, and residual-cross cells were routed,
the new row-local obstruction was
\[
 F_k+I_k\gtrsim\Xi_k.
\]
Here \(F_k\) is balanced mass in the earlier carrier at coordinates
where row \(k\) is active but its own block partner is not, and \(I_k\)
is balanced mass in the residual bank at coordinates which see only
the block containing \(k\).

Those two cells have different algebraic locations.  A foreign-prefix
coordinate is not private globally, but its restriction to the current
two-row block is a singleton.  Its one-row Wilson multiplier is
therefore an exact scalar factor of the actual prefix letter.  A
cut-internal coordinate, by contrast, has zero joining letter for the
fixed cut.  It factors out of the V77 cross covariance as a complete
same-side moment.  That complete moment must then be split exactly into
its internal binary covariance and its replicated one-row product.

This chapter carries out those two operations before taking a norm.  It
proves a foreign-prefix row debit and an unpaid internal-bank
replacement compiler.  The latter says that a large clean
cut-internal bank may replace the old binary prefix as the carrier of
the same quotient-tree edge.  It does \emph{not} say that an internal
bank which also carries the sole final payer is already controlled.
That payer-overlap, the two-cross/no-prefix state, and the other
first-connectivity families remain separate.

\section{Mandatory V80 companion audit}
\label{sec:V80-audit}

The newest active companion files actually present at the end of the
read-only audit were:
\begin{center}
\small
\begin{tabular}{p{0.47\linewidth}r r p{0.27\linewidth}}
\toprule
file&lines&bytes&SHA--256\\
\midrule
\texttt{Renewal\_SM\_Einstein\_Continuum\_Research\_Notes\_V3.tex}
&\(1750\)&\(61527\)&
\texttt{132f08c23177877e0baf500423c67157a8233ab423916a8457c35cbfe7397a95}\\
\texttt{Renewal\_NS\_Stability\_Research\_Notes\_V31.tex}
&\(23052\)&\(887537\)&
\texttt{f1121b62238ad5491bb7cf03635ea9934942edf573ed8faaa9ee79e1d38a6696}\\
\bottomrule
\end{tabular}
\end{center}
The SM programme advanced from third-volume V1 to V3 during the audit,
so V3 was used for the final fingerprint.  The audit was read only.

SM V1--V3 first identifies the defect of an ordinary boundary lift as
a cocycle in the zero-trace kernel.  It repairs the defect only after
separating harmonic cohomology, constructs a tangential-plus-normal
Hodge--Poisson collar lift as one exact chain packet, and then proves
that the global obstruction is the connecting homomorphism of
\((M,\partial M)\).  A finite-dimensional harmonic correction is
performed before the relative Green contraction; nonextendable
harmonic classes remain an explicit shared cokernel.

NS V31 proves exact flat dyadic marking and its first-moment
unmarking cost.  Its heat/Duhamel ladder reaches the critical
formation moment only from the critical source norm.  A
balanced-helicity plane wave is an exact null vector for the signed
hinge while its positive total critical coarea remains nonzero.

\begin{assessment}[V80 audit decision]
\label{ass:V80-audit-decision}
Only proof order transfers.  The SM calculation requires the defect to
be placed in its exact kernel before a contraction is applied.  The NS
nullspace test forbids inferring control of a positive common factor
from a signed defect which does not see it.  Accordingly, the YM
residual bank is first conditioned into cut-internal and cut-crossing
coordinate algebras.  The internal complete moment is then split into
its replicated and connected parts before either part is estimated.
No SM boundary estimate or NS formation estimate is imported.
\end{assessment}

\begin{firewall}[Cross-program type boundary]
\label{fire:V80-audit-types}
The SM zero-trace kernel is a cochain complex, the NS hinge nullspace is
a Fourier-polarization phenomenon, and the YM object below is a
replica covariance over independent coordinate algebras.  Their
operators and estimates are not interchangeable.  The only common
statement is the logical requirement to identify the exact
kernel/common factor before applying a norm.
\end{firewall}

\section{The actual foreign-prefix letter}
\label{sec:V80-foreign-letter}

Fix the cut
\[
 \rho=C|D,\qquad C=\{k,k^\circ\},
\]
and a coordinate-resolved word in the V61 prefix carrier with exact
join \(\rho\).  Let \(E_{<}\) be its prefix coordinate set.  Define
\(G_k\subseteq E_{<}\) to be the coordinates whose row-\(k\)
incidence has the V79 foreign-prefix label.  Thus row \(k\) is
nontrivial, row \(k^\circ\) is trivial, the coordinate is balanced,
and its actual local partition \(\pi_e\) refines \(\rho\).  Activity
on \(D\) is retained and is not replaced by a private-coordinate
model.

For \(e\in G_k\), put
\begin{equation}
 m_{k,e}:=\mathbb E_e X_{k,e}
 =\eta_{R_{k,e}}(\alpha)I_{R_{k,e}},
 \qquad
 q_k(G_k):=\prod_{e\in G_k}\eta_{R_{k,e}}(\alpha).
\label{eq:V80-foreign-q}
\end{equation}
Centrality of the Wilson law gives the scalar identity in
\eqref{eq:V80-foreign-q}.

\begin{lemma}[Wordwise location of a foreign-prefix incidence]
\label{lem:V80-foreign-word}
For every \(e\in G_k\), the actual local prefix letter factors as
\begin{equation}
 \boxed{
 L_{e,\pi_e}
 =
 m_{k,e}\otimes L^{D}_{e,\pi_e|_D},}
\label{eq:V80-foreign-local-factor}
\end{equation}
up to identity factors on inactive rows.  In particular, the
opposite-block restriction \(\pi_e|_D\) is unchanged.
\end{lemma}

\begin{proof}
Because \(\pi_e\le\rho\), every block of \(\pi_e\) lies wholly in
\(C\) or wholly in \(D\).  Within \(C\), only row \(k\) is
nontrivial.  The restriction of the moment--cumulant letter to its
active rows is therefore the singleton cumulant
\(k_e(\{k\})=\mathbb E_eX_{k,e}=m_{k,e}\).  Formula
\eqref{eq:V61-local-letters} tensors this singleton with the untouched
restriction to \(D\), proving
\eqref{eq:V80-foreign-local-factor}.  No assertion is made that the
coordinate is private in the four-row probability space.
\end{proof}

\begin{theorem}[Exact block factorization of the foreign bank]
\label{thm:V80-foreign-factor}
Let \(K_C^{E_<,\mathrm{conn}}\) be the complete connected binary
carrier on \(C\) in the fixed prefix.  Then
\begin{equation}
 \boxed{
 K_C^{E_<,\mathrm{conn}}
 =
 q_k(G_k)\,
 K_C^{E_<\setminus G_k,\mathrm{conn}}.}
\label{eq:V80-foreign-block-factor}
\end{equation}
Consequently the full exact prefix carrier satisfies
\begin{equation}
 C_<(\rho)
 =
 q_k(G_k)
 \left(
 K_C^{E_<\setminus G_k,\mathrm{conn}}
 \otimes K_D^{E_<,\mathrm{conn}}
 \right).
\label{eq:V80-foreign-full-prefix}
\end{equation}
The same factorization may be performed simultaneously for every
foreign-prefix row; the resulting scalar factors commute.
\end{theorem}

\begin{proof}
Use the two-endpoint representation of the binary connected carrier:
\[
 K_C^{E_<,\mathrm{conn}}
 =
 \mathbb E\!\left[W_k^{E_<}\otimes W_{k^\circ}^{E_<}\right]
 -
 \mathbb EW_k^{E_<}\otimes
 \mathbb EW_{k^\circ}^{E_<}.
\]
At every coordinate of \(G_k\), row \(k^\circ\) is trivial, while
row \(k\) contributes the one-row moment \(m_{k,e}\).
Coordinate independence therefore factors
\(\prod_{e\in G_k}m_{k,e}=q_k(G_k)I\) from both endpoint moments.
A singleton coordinate cannot connect \(k\) to \(k^\circ\), so every
word which connects \(C\) remains connected after the coordinates
\(G_k\) are removed.  This proves
\eqref{eq:V80-foreign-block-factor}.  Apply the exact prefix
factorization \eqref{eq:V61-prefix-factorization} to obtain
\eqref{eq:V80-foreign-full-prefix}.  Disjoint singleton row factors
are central, so all rowwise factorizations commute.
\end{proof}

\begin{firewall}[Why this is not the forbidden merge
\(F_k\subset A_k\)]
\label{fire:V80-foreign-not-own}
The foreign coordinate supplies no \(kk^\circ\) overlap and no
binary response stock:
\[
 a_{k,e}a_{k^\circ,e}=0,\qquad e\in G_k.
\]
It is not inserted into the own-prefix mass \(A_k\).  Instead, its
one-row scalar is factored from the complete block carrier, while the
actual opposite-block letter remains exactly where V61 placed it.
\end{firewall}

\section{Foreign-prefix Wilson debits}
\label{sec:V80-foreign-debit}

By the V79 definition,
\begin{equation}
 F_k=\sum_{e\in G_k}a_{k,e},
 \qquad a_{k,e}=1+C_2(R_{k,e}).
\label{eq:V80-F-stock}
\end{equation}

\begin{lemma}[All fixed foreign-prefix moments]
\label{lem:V80-foreign-moments}
For every fixed integer \(m\ge0\),
\begin{equation}
 \boxed{
 (s_\alpha F_k)^m q_k(G_k)\le C_m.}
\label{eq:V80-foreign-moments}
\end{equation}
The constant is independent of \(|G_k|\), the representations, and
the activity or local partition on the opposite block.
\end{lemma}

\begin{proof}
The scalar in \eqref{eq:V80-foreign-q} is a product of the same
one-row Wilson multipliers used in
\Cref{cor:V52-private-products-all-orders}.  The proof of that
corollary uses only the one-link failure debits
\eqref{eq:V52-all-failure-debits} and the finite Bernoulli product
compiler.  It does not use global privateness of the underlying
coordinate.  Apply that proof to the row-\(k\) scalar factors in
\eqref{eq:V80-foreign-local-factor}.  The opposite-block tensor
factors do not enter the scalar product.
\end{proof}

\begin{corollary}[A foreign-dominant row has a prefix certificate]
\label{cor:V80-foreign-row-certificate}
Suppose
\begin{equation}
 F_k\ge\kappa\Xi_k,\qquad 0<\kappa\le1.
\label{eq:V80-foreign-dominance}
\end{equation}
Let
\[
 H_C^{<}:=\sum_{e\in E_<}a_{k,e}a_{k^\circ,e}.
\]
Then the complete binary block carrier obeys
\begin{equation}
 \boxed{
 \max\left\{
 \|K_C^{E_<,\mathrm{conn}}\|,
 \kappa_\otimes(K_C^{E_<,\mathrm{conn}})
 \right\}
 \le
 C_\kappa\frac{H_C^{<}}{\Xi_k}.}
\label{eq:V80-foreign-prefix-certificate}
\end{equation}
\end{corollary}

\begin{proof}
The coordinates of \(G_k\) contribute zero to \(H_C^{<}\).
The ordinary-norm binary response estimate on the remaining bank gives
\[
 \|K_C^{E_<\setminus G_k,\mathrm{conn}}\|
 \le Cs_\alpha H_C^{<}.
\]
By \Cref{lem:V80-foreign-moments} with \(m=1\) and
\eqref{eq:V80-foreign-dominance},
\[
 q_k(G_k)\le\frac{C_\kappa}{s_\alpha\Xi_k}.
\]
Multiply these estimates in the exact identity
\eqref{eq:V80-foreign-block-factor}.  Product-state capacity is no
larger than ordinary norm.
\end{proof}

\begin{corollary}[Foreign-labelled ancestry is not new input mass]
\label{cor:V80-foreign-ancestry}
Under the foreign-dominance hypothesis
\eqref{eq:V80-foreign-dominance}, for all integers \(r,\ell\ge0\)
with \(r+\ell\le3\),
\begin{equation}
 \boxed{
 (s_\alpha\Xi_k)^r(s_\alpha F_k)^\ell q_k(G_k)
 \le C_{\kappa,r,\ell}.}
\label{eq:V80-foreign-ancestry-moments}
\end{equation}
Thus an existing coordinate-resolved \(2+2\) source compiler may
assign one endpoint debit and one \(F_k\)-labelled
output-ancestry debit to the same retained scalar, provided its total
degree is at most three.  The statement supplies only this scalar
debit; it does not estimate unrelated graph factors.
\end{corollary}

\begin{proof}
By \eqref{eq:V80-foreign-dominance},
\(\Xi_k^r\le\kappa^{-r}F_k^r\).  Hence the left side of
\eqref{eq:V80-foreign-ancestry-moments} is at most
\[
 \kappa^{-r}(s_\alpha F_k)^{r+\ell}q_k(G_k),
\]
which is bounded by \eqref{eq:V80-foreign-moments}.  This is one
higher moment of one exact scalar, not several copies of the
coordinate bank.
\end{proof}

\begin{remark}[Opposite-block source marks]
\label{rem:V80-foreign-opposite-mark}
A coordinate of \(G_k\) may also be a marked prefix coordinate for the
opposite block.  Equation \eqref{eq:V80-foreign-local-factor} proves
that the row-\(k\) scalar and the actual opposite-block connected
letter are separate tensor factors in that replica letter.  Hence the
foreign debit survives such a coincidence.  This does not authorize a
second use of the opposite-block response or of the sole final payer.
\end{remark}

\section{Conditioning the residual bank on its internal algebra}
\label{sec:V80-conditioning}

Return to the fixed V77 residual bank \(B\) for
\(\rho=C|D\).  Split it exactly as
\begin{align}
 B_0
 &:=
 \{e\in B:\varnothing\ne A_e\subseteq C\ \hbox{or}\
                     \varnothing\ne A_e\subseteq D\},
\label{eq:V80-B0}\\
 B_\times
 &:=
 \{e\in B:A_e\cap C\ne\varnothing,\
                 A_e\cap D\ne\varnothing\}.
\label{eq:V80-Bcross}
\end{align}
Also write \(B_0=B_C\mathbin{\dot\cup}B_D\) according to the side
seen by the coordinate.  For \(e\in B_0\),
\begin{equation}
 T_e=R_e^\rho,\qquad \Delta_e^\rho=0.
\label{eq:V80-internal-zero-letter}
\end{equation}
Put
\begin{equation}
 \mathcal M_0:=\bigotimes_{e\in B_0}T_e.
\label{eq:V80-M0}
\end{equation}

\begin{theorem}[Exact conditional covariance factorization]
\label{thm:V80-conditional-factor}
The complete V77 first-join packet satisfies
\begin{equation}
 \boxed{
 \mathfrak J_B^\rho
 =
 \mathcal M_0\otimes
 \mathfrak J_{B_\times}^\rho.}
\label{eq:V80-conditional-factor}
\end{equation}
Equivalently, if \(\mathcal A_0\) is the coordinate algebra of \(B_0\),
then the operator law of total covariance gives
\begin{align}
 \operatorname{Cov}_{B}(F_C^B,F_D^B)
 &=
 \mathbb E_{B_0}
 \operatorname{Cov}_{B_\times}
 (F_C^B,F_D^B\mid\mathcal A_0)
\notag\\
 &\quad+
 \operatorname{Cov}_{B_0}
 \left(
  \mathbb E_{B_\times}(F_C^B\mid\mathcal A_0),
  \mathbb E_{B_\times}(F_D^B\mid\mathcal A_0)
 \right),
\label{eq:V80-total-covariance}\\
 \operatorname{Cov}_{B_0}
 \left(
  \mathbb E_{B_\times}(F_C^B\mid\mathcal A_0),
  \mathbb E_{B_\times}(F_D^B\mid\mathcal A_0)
 \right)
 &=0.
\label{eq:V80-zero-conditional-covariance}
\end{align}
\end{theorem}

\begin{proof}
By \eqref{eq:V80-internal-zero-letter}, the \(B_0\) factors in the
two endpoints of the V77 telescope are identical.  Order invariance
\Cref{cor:V77-packet-order} allows them to be placed first.  Thus
\begin{align*}
 \mathfrak J_B^\rho
 &=
 \left(\bigotimes_{e\in B_0}T_e\right)
 \otimes
 \left[
  \bigotimes_{e\in B_\times}T_e
  -
  \bigotimes_{e\in B_\times}R_e^\rho
 \right]\\
 &=\mathcal M_0\otimes\mathfrak J_{B_\times}^\rho,
\end{align*}
by \eqref{eq:V77-bank-telescope}.

For the tower form, write
\[
 F_C^B=G_C^{B_C}\widetilde F_C^{B_\times},
 \qquad
 F_D^B=G_D^{B_D}\widetilde F_D^{B_\times}.
\]
The three coordinate algebras \(B_C,B_D,B_\times\) are independent.
Conditioning on \(B_0\), the first covariance in
\eqref{eq:V80-total-covariance} is the crossing covariance multiplied
by \(G_C^{B_C}\otimes G_D^{B_D}\).  Its expectation gives
\(\mathcal M_0\otimes\mathfrak J_{B_\times}^\rho\).
The two conditional means in the second covariance have internal
random factors depending respectively only on \(B_C\) and \(B_D\).
Those algebras are independent, so their covariance is zero.  This
proves \eqref{eq:V80-zero-conditional-covariance} and the theorem.
\end{proof}

\begin{corollary}[The crossing minima cannot see internal row mass]
\label{cor:V80-cross-minima-blind}
One has
\begin{equation}
 H_B(C,D)=H_{B_\times}(C,D).
\label{eq:V80-cross-stock-only}
\end{equation}
Consequently the unpaid and paid V77--V78 cut minima remain
\begin{align}
 \kappa_\otimes(\mathfrak J_B^\rho)
 &\le C\min\{1,s_\alpha H_{B_\times}(C,D)\},
\label{eq:V80-unpaid-cross-only}\\
 \kappa_\otimes(\mathfrak P_B^\rho)
 &\le C\min\{H_{B_\times}(C,D),s_\alpha^{-1}\}.
\label{eq:V80-paid-cross-only}
\end{align}
Neither right side contains \(I_k\).
\end{corollary}

\begin{proof}
A coordinate in \(B_0\) has no active pair across the cut and hence
contributes zero to \eqref{eq:V77-cross-stock}.  This proves
\eqref{eq:V80-cross-stock-only}.  Apply
\Cref{thm:V77-composite-minimum,thm:V78-paid-cut-minimum}.
For the paid statement, this direct application is essential: the
global output mass and resolvent are not factored through
\eqref{eq:V80-conditional-factor}.
\end{proof}

\begin{firewall}[The zero conditional defect is not decay]
\label{fire:V80-zero-not-decay}
If \(B_\times=\varnothing\), then
\(\mathfrak J_B^\rho=0\) for every size of the internal masses.  If a
small crossing bank is added, the entire internal complete moment is a
common factor of that crossing covariance.  Thus no inequality using
only the signed cross defect can manufacture a decreasing function of
\(I_k\).  The internal factor must be split and used on its own block.
\end{firewall}

\section{The balanced internal binary split}
\label{sec:V80-internal-split}

Fix \(C=\{k,k^\circ\}\).  Let \(D_C\subseteq B_C\) be a clean
balanced cut-internal coordinate bank: both \(k\) and \(k^\circ\) are
nontrivial at every coordinate of \(D_C\), no opposite-block row is
active, and the bank is disjoint from the selected source marks and
from the coordinate carrying the sole final payer.  The V79 priority
implies that every incidence counted by \(I_k\), after private and
unbalanced branches have been removed, has precisely this two-row
form.

Put
\begin{align}
 I_k(D_C)&:=\sum_{e\in D_C}a_{k,e},
 &
 I_{k^\circ}(D_C)&:=\sum_{e\in D_C}a_{k^\circ,e},
\label{eq:V80-I-stocks}\\
 H_C(D_C)&:=\sum_{e\in D_C}a_{k,e}a_{k^\circ,e}.
\label{eq:V80-HI}
\end{align}
Let
\begin{align}
 T_C(D_C)
 &:=
 \mathbb E_{D_C}
 \left[W_k^{D_C}\otimes W_{k^\circ}^{D_C}\right],
\label{eq:V80-internal-joint}\\
 Q_k(D_C)&:=\mathbb E_{D_C}W_k^{D_C}
 =q_k(D_C)I,
\label{eq:V80-internal-Qk}\\
 Q_{k^\circ}(D_C)&:=\mathbb E_{D_C}W_{k^\circ}^{D_C}
 =q_{k^\circ}(D_C)I.
\label{eq:V80-internal-Qkcirc}
\end{align}

\begin{lemma}[Exact connected/replicated split]
\label{lem:V80-internal-binary-split}
The complete internal moment has the exact two-term decomposition
\begin{equation}
 \boxed{
 T_C(D_C)
 =
 K_C^{D_C,\mathrm{conn}}
 +
 Q_k(D_C)\otimes Q_{k^\circ}(D_C),}
\label{eq:V80-internal-binary-split}
\end{equation}
where
\begin{equation}
 K_C^{D_C,\mathrm{conn}}
 =
 T_C(D_C)
 -
 Q_k(D_C)\otimes Q_{k^\circ}(D_C).
\label{eq:V80-internal-K}
\end{equation}
Inserted into the conditioned residual packet, this gives
\begin{align}
 \mathfrak J_B^\rho
 &=
 \mathcal M_{\rm rest}
 \otimes K_C^{D_C,\mathrm{conn}}
 \otimes\mathfrak J_{B_\times}^\rho
\notag\\
 &\quad+
 \mathcal M_{\rm rest}
 \otimes Q_k(D_C)\otimes Q_{k^\circ}(D_C)
 \otimes\mathfrak J_{B_\times}^\rho,
\label{eq:V80-conditioned-two-branch}
\end{align}
where \(\mathcal M_{\rm rest}\) is the product of the remaining
cut-internal complete moments.
\end{lemma}

\begin{proof}
Equation \eqref{eq:V80-internal-binary-split} is the defining
moment--cumulant identity for two composite row variables, equivalently
\Cref{thm:V75-bank-covariance}.  Substitute it for the \(D_C\) factor
of \(\mathcal M_0\) in
\eqref{eq:V80-conditional-factor}.  Coordinate tensor products
distribute over the two-term sum, proving
\eqref{eq:V80-conditioned-two-branch}.
\end{proof}

\begin{lemma}[The replicated internal branch has row moments]
\label{lem:V80-internal-replicated-moments}
For every fixed integer \(m\ge0\),
\begin{align}
 (s_\alpha I_k(D_C))^m q_k(D_C)&\le C_m,
\label{eq:V80-internal-k-moments}\\
 (s_\alpha I_{k^\circ}(D_C))^m q_{k^\circ}(D_C)&\le C_m.
\label{eq:V80-internal-kcirc-moments}
\end{align}
\end{lemma}

\begin{proof}
After the replicated split, each \(Q\) is exactly a product of one-row
Wilson multipliers.  Apply the one-link failure debits and the product
compiler exactly as in
\Cref{cor:V52-private-products-all-orders}.  The original coordinates
are shared, but the two terms in
\eqref{eq:V80-internal-binary-split} have already separated their
replica blocks algebraically; no probabilistic independence beyond
that exact identity is asserted.
\end{proof}

\section{The unpaid internal replacement compiler}
\label{sec:V80-internal-compiler}

Let \(E_C\) be the earlier coordinate bank of the connected
\(C\)-prefix and write
\[
 K_C^{\rm pre}:=K_C^{E_C,\mathrm{conn}},
 \qquad
 H_C^{\rm pre}:=\sum_{e\in E_C}a_{k,e}a_{k^\circ,e}.
\]
The banks \(E_C\) and \(D_C\) are disjoint.

\begin{theorem}[Internal complete moment replaces a poor prefix]
\label{thm:V80-internal-replacement}
Assume
\begin{equation}
 I_k(D_C)\ge\kappa\Xi_k
\label{eq:V80-I-dominance}
\end{equation}
and that \(D_C\) is unpaid in the current one-resolvent allocation.
Then
\begin{equation}
 \boxed{
 \max\left\{
 \|K_C^{\rm pre}\otimes T_C(D_C)\|,
 \kappa_\otimes(K_C^{\rm pre}\otimes T_C(D_C))
 \right\}
 \le
 C_\kappa
 \frac{H_C^{\rm pre}+H_C(D_C)}{\Xi_k}.}
\label{eq:V80-internal-replacement}
\end{equation}
Thus the internal bank supplies the same ordinary-norm row certificate
as a prefix-dominant binary bank, although its coordinates occurred
after the fixed two-block carrier.
\end{theorem}

\begin{proof}
Use the exact split
\eqref{eq:V80-internal-binary-split}.  For its connected branch,
\(\|K_C^{\rm pre}\|\le C\) by the two endpoint moments, while
\Cref{lem:V78-binary-norm-compilers} and
\eqref{eq:V80-I-dominance} give
\[
 \|K_C^{D_C,\mathrm{conn}}\|
 \le C_\kappa\frac{H_C(D_C)}{\Xi_k}.
\]
Hence
\begin{equation}
 \|K_C^{\rm pre}\otimes K_C^{D_C,\mathrm{conn}}\|
 \le
 C_\kappa\frac{H_C(D_C)}{\Xi_k}.
\label{eq:V80-connected-branch-bound}
\end{equation}

For the replicated branch use the ordinary-norm response estimate
\[
 \|K_C^{\rm pre}\|\le Cs_\alpha H_C^{\rm pre}.
\]
The factor \(Q_{k^\circ}(D_C)\) is a contraction.  By
\Cref{lem:V80-internal-replicated-moments} with \(m=1\),
\[
 \|Q_k(D_C)\|
 =q_k(D_C)
 \le
 \frac{C}{s_\alpha I_k(D_C)}
 \le
 \frac{C_\kappa}{s_\alpha\Xi_k}.
\]
Therefore
\begin{equation}
 \|K_C^{\rm pre}\otimes
 Q_k(D_C)\otimes Q_{k^\circ}(D_C)\|
 \le
 C_\kappa\frac{H_C^{\rm pre}}{\Xi_k}.
\label{eq:V80-replicated-branch-bound}
\end{equation}
Add
\eqref{eq:V80-connected-branch-bound} and
\eqref{eq:V80-replicated-branch-bound}.  Product-state capacity is
bounded by ordinary norm.
\end{proof}

\begin{assessment}[What the two branches pay]
\label{ass:V80-two-branch-payment}
The two summands in
\eqref{eq:V80-conditioned-two-branch} use different mechanisms.
The connected internal covariance is bounded by its
contraction/response minimum; it supplies the \(kk^\circ\) tree edge,
while the old prefix is used only as a contraction.  The replicated
summand keeps the old prefix response edge and pays row \(k\) with the
one-row multiplier \(q_k(D_C)\).  Adding the two scalar edge stocks is
legal because
\[
 H_C^{\rm pre}+H_C(D_C)\le H_{kk^\circ}.
\]
No new graph edge and no second coordinate copy is introduced.
\end{assessment}

\begin{corollary}[The clean \(F_k+I_k\) input alternative closes
rowwise]
\label{cor:V80-FI-rowwise}
Fix \(0<\eta<1/4\).  In the V79 finite-cell alternative, suppose
\[
 F_k+I_k>(1-4\eta)\Xi_k.
\]
If the \(I_k\)-bank is clean and unpaid, then the current block
carrier has an ordinary-norm certificate
\begin{equation}
 \boxed{
 \|\mathcal B_C\|
 \le
 C_\eta\frac{H_{kk^\circ}}{\Xi_k},}
\label{eq:V80-FI-block-certificate}
\end{equation}
where \(\mathcal B_C\) is the exact earlier \(C\)-prefix multiplied by
all unpaid cut-internal complete moments on \(C\).
\end{corollary}

\begin{proof}
Either
\[
 F_k>\frac{1-4\eta}{2}\Xi_k
\quad\hbox{or}\quad
 I_k>\frac{1-4\eta}{2}\Xi_k.
\]
In the first case apply
\Cref{cor:V80-foreign-row-certificate} and use contraction on all
cut-internal moments.  In the second case take \(D_C\) to be the
balanced two-row coordinates supporting \(I_k\), apply
\Cref{thm:V80-internal-replacement}, and use contraction on the
remaining internal factors.  The two internal edge stocks are bounded
by the global \(H_{kk^\circ}\).
\end{proof}

\begin{corollary}[Clean wrong-partner columns admit the missing block
orientation]
\label{cor:V80-wrong-partner-clean}
In the V79 three-vertex wrong-partner state, suppose the prefix-poor
row \(i'\) satisfies the V79 \(F_{i'}+I_{i'}\) alternative and its
internal bank is clean and unpaid.  Then the \(i'\)-term of the
column-grouped proof may use
\eqref{eq:V80-FI-block-certificate} in place of
\eqref{eq:V78-unpaid-binary-norm-compiler}.  Hence that
wrong-partner term obeys the same path-tree majorant as
\eqref{eq:V79-column-DRQC}.
\end{corollary}

\begin{proof}
The proof of \Cref{thm:V79-column-closure} uses the block-\(C\)
ordinary norm only through a bound of the form
\[
 \|\mathcal N_C\|
 \le C\frac{H_{12}}{\Xi_r}
\]
for the row \(r\) assigned to the current scalar column card.
Equation \eqref{eq:V80-FI-block-certificate} is exactly that interface
with \(r=i'\), and its numerator is a genuine subset of the global
\(H_{12}\).  The paid column minimum and the opposite block are
unchanged.  Repeating the scalar calculation of V79 proves the result.
\end{proof}

\section{Why the internal minimum is not a terminal moment}
\label{sec:V80-internal-no-terminal}

The replacement theorem must not be simplified to a decay statement
for the internal connected covariance alone.

\begin{proposition}[No hot-row terminal debit from the internal
covariance]
\label{prop:V80-no-internal-terminal}
There is no universal \(C_q\) such that every balanced internal binary
bank satisfies
\begin{equation}
 (s_\alpha I_k)^q
 \kappa_\otimes(K_C^{D_C,\mathrm{conn}})
 \le C_q
\label{eq:V80-false-internal-terminal}
\end{equation}
for any fixed \(q\ge1\).  The failure persists for a fixed diffuse
bank with arbitrarily weak normalized overlap.
\end{proposition}

\begin{proof}
Use the \(N\)-coordinate anti-aligned heat-kernel bank of V74 as an
internal \(kk^\circ\) bank.  Proposition
\ref{prop:V75-no-leading-cancel} gives
\[
 \kappa_\otimes(K_C^{D_C,\mathrm{conn}})
 \ge \frac12c_0^N
\]
on the all-gapped product vector, while
\[
 s_\alpha I_k=s_\alpha Na_s\asymp\frac{N}{s_\alpha}
 \longrightarrow\infty.
\]
This contradicts \eqref{eq:V80-false-internal-terminal}.  The same
example has \(\theta_{kk^\circ}=1/N\) and is diffuse by
\eqref{eq:V74-bank-theta}--\eqref{eq:V74-bank-diffuse}.
\end{proof}

\begin{assessment}[The paid binary minimum does not give the unpaid
replacement for free]
\label{ass:V80-paid-test}
If the sole payer is assigned to the internal connected covariance,
\Cref{cor:V78-paid-binary-minimum,cor:V78-paid-pays-row} give
\begin{equation}
 \|\mathcal P_C(D_C)\|
 \le
 C_\kappa\frac{H_C(D_C)}
 {s_\alpha\Xi_k}.
\label{eq:V80-paid-internal-card}
\end{equation}
Multiplying by the response bound of the earlier prefix produces two
parallel pair stocks,
\[
 s_\alpha H_C^{\rm pre}\,
 \frac{H_C(D_C)}{s_\alpha\Xi_k}
 =
 \frac{H_C^{\rm pre}H_C(D_C)}{\Xi_k},
\]
not the single-edge certificate in
\eqref{eq:V80-internal-replacement}.  Using the prefix only as a
contraction leaves the explicit \(s_\alpha^{-1}\) in
\eqref{eq:V80-paid-internal-card}.  Either expression may become useful
after the complete source scale and output ancestry are fixed, but
neither is a source-independent paid replacement theorem.  The sole
payer must not also be assigned to the crossing covariance.
\end{assessment}

\begin{firewall}[Anti-alignment delimits the claim rather than
defeating it]
\label{fire:V80-anti-aligned-boundary}
\Cref{prop:V80-no-internal-terminal} does not contradict
\Cref{thm:V80-internal-replacement}.  In the anti-aligned regime the
connected covariance remains order one, but the replacement theorem
uses it as the actual normalized \(kk^\circ\) tree carrier and uses
the old prefix only as a contraction.  V75 already proved that this is
the harmless normalization of the anti-aligned bank.  What fails is
only an attempt to use the same covariance as a graph-free terminal
moment.
\end{firewall}

\section{Marked, payer, and ancestry boundaries}
\label{sec:V80-boundaries}

\begin{theorem}[Exact input-bank routing achieved in V80]
\label{thm:V80-input-routing}
For a conditioned \(2+2\) source packet:
\begin{enumerate}[label=\textup{(\roman*)}]
\item every foreign-prefix bank is an exact scalar Wilson debit inside
      its actual block carrier, including when the same coordinate is
      a source mark for the opposite block;
\item every clean unpaid cut-internal bank has the two-branch
      factorization \eqref{eq:V80-conditioned-two-branch} and the row
      certificate \eqref{eq:V80-internal-replacement}; and
\item the V79 clean \(F_k+I_k\) input alternative therefore no longer
      remains an undifferentiated obstruction.
\end{enumerate}
\end{theorem}

\begin{proof}
Items \textup{(i)--(ii)} are
\Cref{thm:V80-foreign-factor,
lem:V80-foreign-moments,
thm:V80-conditional-factor,
lem:V80-internal-binary-split,
thm:V80-internal-replacement}.  Item \textup{(iii)} is
\Cref{cor:V80-FI-rowwise}.
\end{proof}

\begin{firewall}[What remains outside the clean compiler]
\label{fire:V80-marked-payer-boundary}
The internal replacement theorem assumes that \(D_C\) is not a source
mark and does not carry the sole final payer.  If an internal
coordinate is the paid suffix/output coordinate, the global output
mass is coupled to the internal physical block and
\eqref{eq:V80-paid-internal-card} must be inserted into the complete
one-payer source before a surplus parallel edge is discarded.  If two
selected transfer banks share an internal coordinate, statewise
tensorization is also unavailable until a one-use local identity is
proved.  These are precisely the marked/payer-overlap exclusions, not
exceptions hidden inside \eqref{eq:V80-internal-replacement}.
\end{firewall}

\begin{assessment}[Output ancestry after the input identity]
\label{ass:V80-output-after-input}
The \(F_k\)-labelled output cell is now routed by the same exact scalar
factor as its input incidence, with total moment degree at most three.
For an \(I_k\)-labelled output cell, the replicated branch of
\eqref{eq:V80-conditioned-two-branch} has the analogous one-row
moments.  The connected branch, however, attaches the output payer to
the very bank used as the replacement edge.  It is governed by the
paid card \eqref{eq:V80-paid-internal-card} and remains open pending a
complete one-payer parallel-edge calculation.  The ancestry majorant
has not been counted as a second input bank.
\end{assessment}

\section{Post-V80 reduction and V81 mandate}
\label{sec:V80-reduction}

\begin{theorem}[Integrated compact-series V80 statement]
\label{thm:V80-integrated}
Preserving compact V1--V79, V80 proves:
\begin{enumerate}[label=\textup{(V80.\arabic*)},leftmargin=6.7em]
\item the mandatory read-only audit of the newest SM V3 and NS V31
      companions, with exact fingerprints and a strict type firewall;
\item the wordwise and blockwise location of every foreign-prefix
      one-row Wilson multiplier in its actual V61 prefix letter;
\item support-uniform fixed moments and an ordinary-norm row
      certificate for a foreign-dominant prefix;
\item the exact law-of-total-covariance factorization of the residual
      bank into a common cut-internal moment and the complete crossing
      covariance;
\item the exact connected/replicated split of a balanced internal
      binary moment;
\item an unpaid internal replacement compiler which pays a
      prefix-poor row while retaining a genuine \(kk^\circ\) tree
      edge;
\item rowwise closure of the clean \(F_k+I_k\) input alternative and
      the corresponding clean wrong-partner orientation;
\item an anti-aligned proof that the connected internal bank is not a
      stand-alone terminal moment; and
\item localization of the remaining \(I_k\)-labelled output problem to
      the paid connected internal branch.
\end{enumerate}
\end{theorem}

\begin{proof}
These are
\Cref{ass:V80-audit-decision,
lem:V80-foreign-word,
thm:V80-foreign-factor,
lem:V80-foreign-moments,
cor:V80-foreign-row-certificate,
thm:V80-conditional-factor,
cor:V80-cross-minima-blind,
lem:V80-internal-binary-split,
thm:V80-internal-replacement,
cor:V80-FI-rowwise,
cor:V80-wrong-partner-clean,
prop:V80-no-internal-terminal,
ass:V80-output-after-input}.
\end{proof}

\begin{openproblem}[V81 mandate: paid internal parallel-edge ledger]
\label{op:V81-mandate}
\begin{enumerate}
\item Start from the exact split
      \eqref{eq:V80-conditioned-two-branch}; do not merge the connected
      and replicated internal terms after physical resolution.
\item Apply the V2 sole-payer ancestry allocation to the connected
      internal branch and write the complete scalar coefficient when
      that branch, the crossing bank, or either earlier prefix pays.
\item In the internally paid case, retain both parallel
      \(kk^\circ\) stocks until division by all four row masses.
      Determine whether one normalized parallel edge is automatically
      surplus or whether a hot partner factor remains.
\item Route the \(I_k\)-labelled output-ancestry debit in the same
      calculation.  Do not add \(I_k\) a second time as input mass.
\item Test the resulting paid coefficient on the V74--V75
      anti-aligned bank and on a one-large/one-small mass pair.
\item Return to the two-cross/no-prefix state only after the internal
      paid branch is settled.  It requires two endpoint gains from one
      crossing covariance and may need a genuine two-puncture
      difference rather than another scalar pigeonhole.
\item Keep marked-heavy and coordinate-overlap branches separate until
      an exact one-use identity is proved.
\item Do not extend the two-block conditional identity to \(3+1\),
      \(2+1+1\), or the discrete source family by analogy; derive each
      source's exact partition formula first.
\end{enumerate}
\end{openproblem}

\begin{firewall}[Claim boundary after compact V80]
\label{fire:V80-current-boundary}
V80 proves the foreign-prefix debit, the exact internal/cross
conditioning identity, and the clean unpaid internal replacement
compiler.  It does not close an internally paid or overlapping
\(I_k\)-branch, the two-cross/no-prefix state, all marked-heavy
branches, or the remaining \(3+1\), \(2+1+1\), and discrete hot source
families.  Global quartic and higher MTA, WUFC, CPM, constructive
continuum Yang--Mills existence, Osterwalder--Schrader reconstruction,
and a positive physical mass gap remain open.  The full Yang--Mills
mass gap remains open.
\end{firewall}

\section{Append-only record for compact V80}
\label{sec:V80-record}

V80 was appended to the complete compact V79 file.  The exact V79
parent had \(34044\) newline-delimited source records, \(1176151\)
bytes, and SHA--256 digest
\[
\texttt{b2c82c8acae1e28cfb348f00fa13aeeffd25884697df52a2cbfbf915ec317cbe}.
\]
No compact V79 mathematical content was changed.  The sole final
\verb|\end{document}| is restored below.

% =============================================================================
% END APPENDED COMPACT VERSION V80 MATERIAL
% =============================================================================

% =============================================================================
% BEGIN APPENDED COMPACT VERSION V81 MATERIAL
% =============================================================================

\clearpage
\chapter{Paid augmented blocks and the one-payer three-carrier ledger}
\label{chap:V81}

\section{Context: contract the unused parallel carrier}
\label{sec:V81-context}

V80 conditioned the residual \(2+2\) covariance into a crossing bank
and two same-side complete-moment banks.  On one side
\(C=\{k,k^\circ\}\), it proved the exact split
\[
 K_C^{\rm pre}\otimes T_C(D_C)
 =
 K_C^{\rm pre}\otimes K_C^{D_C,\rm conn}
 +
 K_C^{\rm pre}\otimes Q_k(D_C)\otimes Q_{k^\circ}(D_C).
\]
When \(D_C\) is unpaid and carries a fixed fraction of \(\Xi_k\), the
first summand uses the internal covariance as the \(kk^\circ\) tree
edge and contracts the old prefix; the second uses the old prefix edge
and pays row \(k\) with \(Q_k\).

The remaining question is what happens when the sole final payer lies
inside this augmented block.  The apparent obstruction in V80 came
from applying response estimates to both parallel \(kk^\circ\)
carriers.  That is unnecessary.  In every payer location one of the
parallel carriers can be used only as a contraction.  The other
carrier supplies the unique \(kk^\circ\) tree edge and the missing row
denominator.

This note proves the resulting paid augmented-block certificate.  It
then groups the complete conditioned \(2+2\) packet into three
carriers: the augmented \(C\)-block, the augmented \(D\)-block, and
the crossing covariance.  The V2 allocation puts the sole payer in
exactly one of them.  All three alternatives have the same path-tree
scale.  In particular, the clean \(I_k\)-labelled output-ancestry
branch is closed.  Coordinate overlaps between independently selected
transfer banks and the two-cross/no-prefix state are not included.

\section{A paid product-moment lemma}
\label{sec:V81-product-payer}

Let \(G\) be a finite bank of nontrivial one-row Wilson factors.  Put
\begin{equation}
 Y_G:=\sum_{e\in G}a_e,
 \qquad
 q_G:=\prod_{e\in G}\eta_e,
 \qquad
 x_G:=s_\alpha Y_G.
\label{eq:V81-product-data}
\end{equation}
For \(G\ne\varnothing\), define its complete scalar payer by
\begin{equation}
 \mathscr D_G^{\rm pay}
 :=
 \frac{Y_Gq_G}{1-q_G}.
\label{eq:V81-product-payer}
\end{equation}
This is the output-mass/resolvent multiplier of the complete one-row
product, not a sum of independently resolvented coordinates.

\begin{lemma}[Product gap at the total row scale]
\label{lem:V81-product-gap}
There is \(c>0\) such that
\begin{equation}
 \boxed{
 1-q_G\ge c\min\{1,x_G\}.}
\label{eq:V81-product-gap}
\end{equation}
\end{lemma}

\begin{proof}
For \(x_e=s_\alpha a_e\), the one-link gap gives
\[
 \eta_e\le1-c_0\min\{1,x_e\}
 \le\exp[-c_0\min\{1,x_e\}].
\]
Hence
\[
 q_G\le
 \exp\left[-c_0\sum_{e\in G}\min\{1,x_e\}\right].
\]
If every \(x_e<1\), the sum in the exponent is \(x_G\).  If some
\(x_e\ge1\), it is at least one.  Thus
\[
 1-q_G\ge
 1-e^{-c_0\min\{1,x_G\}},
\]
which is bounded below by a fixed multiple of
\(\min\{1,x_G\}\).
\end{proof}

\begin{theorem}[Every fixed paid product moment]
\label{thm:V81-paid-product-moments}
For every fixed integer \(m\ge0\),
\begin{equation}
 \boxed{
 x_G^m\,s_\alpha\mathscr D_G^{\rm pay}
 =
 x_G^m\frac{x_Gq_G}{1-q_G}
 \le C_m.}
\label{eq:V81-paid-product-moments}
\end{equation}
In particular,
\begin{align}
 \mathscr D_G^{\rm pay}
 &\le\frac C{s_\alpha},
\label{eq:V81-product-paid-zero}\\
 \mathscr D_G^{\rm pay}
 &\le\frac C{s_\alpha^2Y_G}.
\label{eq:V81-product-paid-one}
\end{align}
\end{theorem}

\begin{proof}
If \(x_G\le1\), use
\eqref{eq:V81-product-gap}:
\[
 x_G^m\frac{x_Gq_G}{1-q_G}
 \le Cx_G^m\le C.
\]
If \(x_G\ge1\), the denominator is bounded below by a positive
constant.  The product moment compiler
\eqref{eq:V52-private-products-all-orders}, whose proof uses only the
one-row scalar multipliers, gives
\[
 x_G^{m+1}q_G\le C_{m+1}.
\]
This proves \eqref{eq:V81-paid-product-moments}.  Set \(m=0\) and
\(m=1\) for
\eqref{eq:V81-product-paid-zero}--%
\eqref{eq:V81-product-paid-one}.
\end{proof}

\begin{firewall}[One product denominator, not one per coordinate]
\label{fire:V81-one-product-denominator}
Applying
\(\mathscr D_{\{e\}}^{\rm pay}\le C/s_\alpha\) separately and summing
over \(e\) would introduce \(|G|\).  Theorem
\ref{thm:V81-paid-product-moments} first forms the complete multiplier
\(q_G\) and its single denominator \(1-q_G\).  It is therefore
support-uniform and compatible with the V2 sole-payer allocation.
\end{firewall}

\section{The internally dominant paid block}
\label{sec:V81-internal-block}

In this section all foreign-prefix one-row scalars have first been
factored by \Cref{thm:V80-foreign-factor}; \(E_C\) denotes the
remaining pair-connecting prefix bank.  The foreign factors and their
possible payer locations are reinserted in the next section.

Retain the V80 notation
\[
 K_C^{\rm pre}=K_C^{E_C,\rm conn},
 \qquad
 T_C(D_C)=K_C^{D_C,\rm conn}
 +Q_k(D_C)\otimes Q_{k^\circ}(D_C),
\]
where \(E_C\cap D_C=\varnothing\).  Put
\[
 H_E:=H_C^{\rm pre},
 \qquad
 H_I:=H_C(D_C),
 \qquad
 I_k:=I_k(D_C).
\]
Let \(\mathcal N_C^{E,I}\) be the positive unpayed envelope of this
exact two-branch carrier.  Let \(\mathcal P_C^{E,I}\) be the sum of
the positive envelopes obtained after applying the V2 scalar
allocation to the output coordinates in \(E_C\cup D_C\), with exactly
one final denominator.

\begin{theorem}[Paired internal augmented-block certificate]
\label{thm:V81-paid-internal-block}
If
\begin{equation}
 I_k\ge\kappa\Xi_k,
\label{eq:V81-internal-dominance}
\end{equation}
then
\begin{align}
 \boxed{
 \|\mathcal N_C^{E,I}\|
 \le
 C_\kappa\frac{H_E+H_I}{\Xi_k},}
\label{eq:V81-internal-unpaid-certificate}\\
 \boxed{
 \|\mathcal P_C^{E,I}\|
 \le
 C_\kappa\frac{H_E+H_I}
 {s_\alpha\Xi_k}.}
\label{eq:V81-internal-paid-certificate}
\end{align}
Both inequalities also hold with product-state capacity in place of
ordinary norm.
\end{theorem}

\begin{proof}
The unpaid statement is
\Cref{thm:V80-internal-replacement}.  For the paid statement, keep the
two V80 branches separate.

On the connected branch
\[
 K_C^{\rm pre}\otimes K_C^{D_C,\rm conn},
\]
the payer lies in exactly one binary factor.  If it lies in \(E_C\),
the universal paid binary cap and the internal unpaid compiler give
\[
 \|\mathcal P_C(E_C)\|\,
 \|K_C^{D_C,\rm conn}\|
 \le
 \frac C{s_\alpha}\,
 C_\kappa\frac{H_I}{\Xi_k}.
\]
If it lies in \(D_C\), use contraction on the old prefix and the paid
internal compiler:
\[
 \|K_C^{\rm pre}\|\,
 \|\mathcal P_C(D_C)\|
 \le
 C\,C_\kappa\frac{H_I}{s_\alpha\Xi_k}.
\]
Thus the complete connected branch is bounded by
\begin{equation}
 C_\kappa\frac{H_I}{s_\alpha\Xi_k}.
\label{eq:V81-connected-paid-branch}
\end{equation}

On the replicated branch
\[
 K_C^{\rm pre}\otimes Q_k(D_C)\otimes Q_{k^\circ}(D_C),
\]
there are three payer locations after grouping coordinate allocations.
If the old prefix pays, its paid response is \(CH_E\), while
\eqref{eq:V80-internal-k-moments} and
\eqref{eq:V81-internal-dominance} give
\[
 q_k(D_C)\le\frac{C_\kappa}{s_\alpha\Xi_k}.
\]
The \(k^\circ\)-factor is a contraction, so this case is at most
\[
 C_\kappa\frac{H_E}{s_\alpha\Xi_k}.
\]

If the \(k\)-row product pays, use
\eqref{eq:V81-product-paid-one} with \(Y_G=I_k\), the response
\(\|K_C^{\rm pre}\|\le Cs_\alpha H_E\), and contraction on
\(Q_{k^\circ}\):
\[
 Cs_\alpha H_E\,
 \frac{C}{s_\alpha^2I_k}
 \le
 C_\kappa\frac{H_E}{s_\alpha\Xi_k}.
\]
If the \(k^\circ\)-row product pays, use
\eqref{eq:V81-product-paid-zero} on that row and the unpaid first
moment on \(Q_k\):
\[
 Cs_\alpha H_E\,
 \frac{C_\kappa}{s_\alpha\Xi_k}\,
 \frac C{s_\alpha}
 \le
 C_\kappa\frac{H_E}{s_\alpha\Xi_k}.
\]
These are all payer locations in the replicated two-row product.
Adding them gives
\begin{equation}
 C_\kappa\frac{H_E}{s_\alpha\Xi_k}.
\label{eq:V81-replicated-paid-branch}
\end{equation}
Add
\eqref{eq:V81-connected-paid-branch} and
\eqref{eq:V81-replicated-paid-branch}.  The V2 one-sided tensor law
uses ordinary norm on the factors displayed above, so product-state
capacity obeys the same bound.
\end{proof}

\begin{assessment}[The parallel edge is genuinely surplus]
\label{ass:V81-parallel-surplus}
In the connected branch, if the internal covariance pays, the old
prefix is a contraction; if the old prefix pays, the internal
covariance supplies the unpaid \(kk^\circ\) edge.  In the replicated
branch, exactly one row-product payer is combined with the response
edge of the old prefix.  Hence \(H_EH_I\) never appears in the final
certificate.  The paid block contains the single edge stock
\(H_E+H_I\), with the expected single loss \(s_\alpha^{-1}\).
\end{assessment}

\section{Foreign payers and the unified block interface}
\label{sec:V81-foreign-paid}

Let \(G_k\) be the foreign-prefix bank of row \(k\), let
\[
 q_F:=q_k(G_k),\qquad F_k=\sum_{e\in G_k}a_{k,e},
\]
and factor
\begin{equation}
 K_C^{\rm pre}
 =
 q_F K_C^{R,\rm conn}
\label{eq:V81-foreign-reduced-prefix}
\end{equation}
as in \eqref{eq:V80-foreign-block-factor}.  The reduced bank \(R\)
contains every coordinate which can contribute to the
\(kk^\circ\) response stock; write that stock as \(H_R\).

\begin{theorem}[Paired foreign-prefix certificate]
\label{thm:V81-paid-foreign-block}
Assume
\begin{equation}
 F_k\ge\kappa\Xi_k.
\label{eq:V81-foreign-dominance}
\end{equation}
Let \(\mathcal S_C\) be any product of complete same-side moment
spectators on coordinate banks disjoint from \(G_k\cup R\).
Let \(\mathcal N_C^{F}\) be the unpayed envelope of
\[
 q_FK_C^{R,\rm conn}\otimes\mathcal S_C,
\]
and let \(\mathcal P_C^{F}\) be the exact V2 sum in which the sole
payer lies in \(G_k\), \(R\), or one spectator bank of
\(\mathcal S_C\).  Then
\begin{align}
 \boxed{
 \|\mathcal N_C^{F}\|
 \le C_\kappa\frac{H_R}{\Xi_k},}
\label{eq:V81-foreign-unpaid}\\
 \boxed{
 \|\mathcal P_C^{F}\|
 \le C_\kappa\frac{H_R}{s_\alpha\Xi_k}.}
\label{eq:V81-foreign-paid}
\end{align}
The same bounds hold for product-state capacity.
\end{theorem}

\begin{proof}
The unpaid bound is
\Cref{cor:V80-foreign-row-certificate}; all factors of
\(\mathcal S_C\) are contractions.

For the paid bound, first suppose the payer lies in the reduced
binary prefix.  Its paid response bound is \(CH_R\).  The foreign
first moment gives
\[
 q_F\le\frac{C_\kappa}{s_\alpha\Xi_k},
\]
so this contribution is at most
\(C_\kappa H_R/(s_\alpha\Xi_k)\).

If the payer lies in the foreign product, use
\eqref{eq:V81-product-paid-one} with \(Y_G=F_k\), together with the
unpaid response \(Cs_\alpha H_R\):
\[
 Cs_\alpha H_R\,
 \frac{C}{s_\alpha^2F_k}
 \le
 C_\kappa\frac{H_R}{s_\alpha\Xi_k}.
\]

Finally suppose a complete moment spectator pays.  Its complete
one-payer first-moment bound is \(C/s_\alpha\), while all other
spectators are contractions.  Use the reduced-prefix response and
the unpayed foreign first moment:
\[
 Cs_\alpha H_R\,
 \frac{C_\kappa}{s_\alpha\Xi_k}\,
 \frac C{s_\alpha}
 \le
 C_\kappa\frac{H_R}{s_\alpha\Xi_k}.
\]
The V2 allocation groups these mutually exclusive payer locations
before the one-sided tensor inequality.  Product-state capacity is
therefore bounded by the same ordinary-norm estimates.
\end{proof}

\begin{remark}[A foreign coordinate marked on the opposite block]
\label{rem:V81-foreign-marked-paid}
If \(e\in G_k\) is also a marked coordinate for the opposite block,
the local identity \eqref{eq:V80-foreign-local-factor} assigns the
row-\(k\) scalar and the opposite-block letter to separate replica
factors.  A payer assigned to the row-\(k\) scalar is included in the
second case of the preceding proof; a payer assigned to the
opposite-block letter belongs to the opposite augmented block.  It is
never counted in both.
\end{remark}

\begin{definition}[Paired block certificate]
\label{def:V81-paired-block}
For \(C=\{k,k^\circ\}\), an exact augmented block carrier
\(\mathcal B_C\) has a paired \(k\)-certificate with stock
\(\widehat H_C\) if its unpayed and internally paid envelopes satisfy
\begin{align}
 \|\mathcal N(\mathcal B_C)\|
 &\le C\frac{\widehat H_C}{\Xi_k},
\label{eq:V81-paired-unpaid}\\
 \|\mathcal P(\mathcal B_C)\|
 &\le C\frac{\widehat H_C}{s_\alpha\Xi_k},
\label{eq:V81-paired-paid}
\end{align}
and the same inequalities hold with capacity in place of norm.
Here
\(\widehat H_C\le H_{kk^\circ}\) is a sum of genuine disjoint
\(kk^\circ\) overlap stocks.
\end{definition}

\begin{corollary}[Own-prefix, foreign-prefix, and internal dominance
are paired]
\label{cor:V81-AFI-paired}
Suppose one of the following holds for row \(k\):
\begin{align}
 A_k&\ge\kappa\Xi_k,
\label{eq:V81-A-dom}\\
 F_k&\ge\kappa\Xi_k,
\label{eq:V81-F-dom}\\
 I_k&\ge\kappa\Xi_k.
\label{eq:V81-I-dom}
\end{align}
Assume that the selected coordinate banks are disjoint and that no
second transfer certificate uses one of them.  Then the complete
augmented \(C\)-block has a paired \(k\)-certificate with
\begin{equation}
 \widehat H_C\le H_{kk^\circ}.
\label{eq:V81-AFI-stock}
\end{equation}
\end{corollary}

\begin{proof}
Under \eqref{eq:V81-A-dom}, the unpaid and paid binary minimum
compilers
\eqref{eq:V78-unpaid-binary-norm-compiler}--%
\eqref{eq:V78-paid-binary-norm-compiler}
give the two certificates when the pair prefix itself is unpayed or
paid.  If a disjoint complete spectator pays, use its \(C/s_\alpha\)
cap together with the unpaid prefix certificate.

Under \eqref{eq:V81-F-dom}, apply
\Cref{thm:V81-paid-foreign-block}.  Under
\eqref{eq:V81-I-dom}, apply
\Cref{thm:V81-paid-internal-block}.  Foreign scalars not selected as
payers are contractions.  If a foreign scalar pays in the
internally dominant case, use its \(C/s_\alpha\) paid cap and the
unpaid internal certificate.  If a disjoint internal complete moment
pays in the foreign-dominant case, it is already covered by the last
case of Theorem~\ref{thm:V81-paid-foreign-block}.  In every case the
displayed overlap stock is a sum of subsets of the global pair stock,
which proves \eqref{eq:V81-AFI-stock}.
\end{proof}

\begin{corollary}[The V79 \(F_k+I_k\) alternative is paired]
\label{cor:V81-FI-paired}
If
\[
 F_k+I_k>(1-4\eta)\Xi_k
\]
on a clean disjoint branch, then the augmented block has a paired
\(k\)-certificate with a constant depending only on \(\eta\).
\end{corollary}

\begin{proof}
One of \(F_k,I_k\) is larger than
\((1-4\eta)\Xi_k/2\).  Apply
\Cref{cor:V81-AFI-paired}.
\end{proof}

\begin{firewall}[What paired does not mean]
\label{fire:V81-paired-scope}
The two estimates in a paired certificate are alternatives selected
by the unique V2 payer location.  They are not multiplied, and no
source term contains two resolvent denominators.  The disjointness
hypothesis excludes the V71 overlap branch in which two independently
chosen transfer certificates attempt to use the same coordinate.
\end{firewall}

\section{The exact three-carrier payer ledger}
\label{sec:V81-three-carrier}

Fix \(C=\{1,2\}\), \(D=\{3,4\}\), and use the exact V80
conditioning.  After choosing one connected/replicated branch in each
same-side complete moment, the source has the factor form
\begin{equation}
 \mathcal B_C\otimes\mathcal B_D
 \otimes\mathfrak J_{B_\times}^{C|D}.
\label{eq:V81-three-carrier-factor}
\end{equation}
Here \(\mathcal B_C,\mathcal B_D\) are the complete augmented block
carriers and \(\mathfrak J_{B_\times}^{C|D}\) is the assembled
crossing covariance.  No norm has been used in
\eqref{eq:V81-three-carrier-factor}.

Choose \(i\in C\), \(j\in D\).  Assume the two augmented blocks have
paired certificates
\begin{align}
 \|\mathcal N_C\|
 &\le C\frac{\widehat H_C}{\Xi_i},
 &
 \|\mathcal P_C\|
 &\le C\frac{\widehat H_C}{s_\alpha\Xi_i},
\label{eq:V81-C-paired}\\
 \|\mathcal N_D\|
 &\le C\frac{\widehat H_D}{\Xi_j},
 &
 \|\mathcal P_D\|
 &\le C\frac{\widehat H_D}{s_\alpha\Xi_j},
\label{eq:V81-D-paired}
\end{align}
where
\begin{equation}
 \widehat H_C\le H_{12},
 \qquad
 \widehat H_D\le H_{34}.
\label{eq:V81-block-stock-inclusion}
\end{equation}

For the crossing covariance, the unpaid composite response and paid
scalar refinement give the positive scalar cards
\begin{align}
 \mathfrak n_{ij}^{\times}
 &:=Cs_\alpha H_{ij}(B_\times),
\label{eq:V81-cross-unpaid-card}\\
 \mathfrak p_{ij}^{\times}
 &:=CH_{ij}(B_\times).
\label{eq:V81-cross-paid-card}
\end{align}
These cards are introduced only after the complete signed crossing
operator has been estimated.

\begin{lemma}[Three mutually exclusive payer locations]
\label{lem:V81-three-payer-locations}
After the V2 scalar allocation, the \((i,j)\) scalar majorant of
\eqref{eq:V81-three-carrier-factor} is at most
\begin{align}
 \mathfrak b_{ij}^{\rm pay}
 \le{}&
 \|\mathcal P_C\|\,\|\mathcal N_D\|\,
 \mathfrak n_{ij}^{\times}
\notag\\
 &+
 \|\mathcal N_C\|\,\|\mathcal P_D\|\,
 \mathfrak n_{ij}^{\times}
\notag\\
 &+
 \|\mathcal N_C\|\,\|\mathcal N_D\|\,
 \mathfrak p_{ij}^{\times}.
\label{eq:V81-three-payer-ledger}
\end{align}
The three lines mean respectively that the sole payer lies in the
\(C\)-block, the \(D\)-block, or the crossing bank.
\end{lemma}

\begin{proof}
Resolve a fixed coordinate word and its local partition blocks before
positive majorization.  Repeated fusion bounds the final output mass
by a fixed multiple of the sum of the output masses belonging to the
three factors in
\eqref{eq:V81-three-carrier-factor}.  Apply the scalar allocation
\eqref{eq:V2-resolvent-allocation}; exactly one of those factor masses
receives the unique denominator.

If a same physical coordinate has replica blocks in both \(C\) and
\(D\), the local fusion inequality first splits its output mass
between those replica blocks.  The actual local letter remains the
tensor product furnished by
\eqref{eq:V61-prefix-factorization}; no coordinate is declared to be
independent twice.

For a block payer use ordinary norm on both augmented blocks and
product-state capacity on the assembled unpayed crossing covariance.
For a crossing payer use its paid capacity and ordinary norms on the
two blocks.  This is the one-sided tensor law
\eqref{eq:V2-one-sided-tensor}; no product of capacities occurs.
Finally apply the finite scalar row-pair refinement only to the
already estimated crossing operator.  This gives
\eqref{eq:V81-three-payer-ledger}.
\end{proof}

\begin{theorem}[Every paired \(2+2\) card has the path-tree scale]
\label{thm:V81-paired-22}
Under
\eqref{eq:V81-C-paired}--%
\eqref{eq:V81-block-stock-inclusion},
\begin{equation}
 \boxed{
 \frac{\mathfrak b_{ij}^{\rm pay}}
 {\Xi_1\Xi_2\Xi_3\Xi_4}
 \le
 C\,\theta_{12}\theta_{34}\theta_{ij}.}
\label{eq:V81-paired-22}
\end{equation}
The constant is uniform in the coordinate support, representations,
physical multiplicities, and the location of the sole payer.
\end{theorem}

\begin{proof}
For a \(C\)-block payer, substitute
\eqref{eq:V81-C-paired}--%
\eqref{eq:V81-cross-unpaid-card}:
\[
 \|\mathcal P_C\|\,
 \|\mathcal N_D\|\,
 \mathfrak n_{ij}^{\times}
 \le
 C
 \frac{\widehat H_C\widehat H_DH_{ij}(B_\times)}
 {\Xi_i\Xi_j}.
\]
The factors \(s_\alpha^{-1}\) and \(s_\alpha\) cancel.  The
\(D\)-block payer gives the identical expression.  For a crossing
payer,
\eqref{eq:V81-cross-paid-card} gives that expression directly.
Therefore all three lines of
\eqref{eq:V81-three-payer-ledger} have the common bound
\begin{equation}
 C
 \frac{\widehat H_C\widehat H_DH_{ij}(B_\times)}
 {\Xi_i\Xi_j}.
\label{eq:V81-common-tree-numerator}
\end{equation}
Divide by the four row masses and use
\eqref{eq:V81-block-stock-inclusion} together with
\(H_{ij}(B_\times)\le H_{ij}\).  The result is
\[
 C
 \frac{H_{12}H_{34}H_{ij}}
 {\Xi_1\Xi_2\Xi_3\Xi_4\Xi_i\Xi_j}
 =
 C\theta_{12}\theta_{34}\theta_{ij}.
\]
\end{proof}

\begin{firewall}[The crossing cards remain scalar]
\label{fire:V81-cross-cards-scalar}
Neither
\(\mathfrak n_{ij}^{\times}\) nor
\(\mathfrak p_{ij}^{\times}\) is an
\((i,j)\)-labelled signed operator.  The assembled crossing covariance
is kept intact through its capacity estimate.  Only its finite
nonnegative stock sum is distributed among the four row-pair cards,
exactly as in V78.
\end{firewall}

\section{Closure of the clean internally paid ancestry}
\label{sec:V81-I-ancestry}

\begin{theorem}[The \(I_k\)-labelled payer is already inside the paid
block]
\label{thm:V81-I-ancestry}
Suppose the V79 balanced output-ancestry allocation labels the final
output mass by a clean cut-internal cell \(I_k\).  Suppose the
opposite selected endpoint \(j\) has a paired block certificate and
that
\[
 I_k\ge\kappa\Xi_k.
\]
Then both the connected and replicated internal branches satisfy the
path-tree estimate \eqref{eq:V81-paired-22}.  No additional factor
\(I_k\) is inserted after the paid block estimate.
\end{theorem}

\begin{proof}
An \(I_k\)-labelled output term is, by the V2 ancestry allocation, a
term whose unique output mass and denominator are assigned to the
augmented block containing that internal coordinate.  On the
connected internal branch it is one of the two payer locations
estimated in
\eqref{eq:V81-connected-paid-branch}.  On the replicated branch it is
one of the row-product payer locations estimated using
\eqref{eq:V81-product-paid-zero}--%
\eqref{eq:V81-product-paid-one}.  Hence the entire \(k\)-block obeys
\eqref{eq:V81-internal-paid-certificate}.  The opposite block has
\eqref{eq:V81-D-paired}, and the crossing bank is unpayed, so the first
or second line of
\eqref{eq:V81-three-payer-ledger} applies.  Invoke
\Cref{thm:V81-paired-22}.

The scalar \(I_k\) was used only to establish
\eqref{eq:V81-internal-dominance} and, on the replicated branch, the
degree-one paid product moment.  The output mass is the numerator of
that same paid factor.  Multiplying by another \(I_k\) would count the
ancestry twice.
\end{proof}

\begin{corollary}[Clean foreign/internal wrong-partner states close]
\label{cor:V81-clean-wrong-partner}
In the V79 three-vertex wrong-partner configuration, suppose every
prefix-poor row needed by the column grouping satisfies a clean
\(F+I\) alternative, and the opposite selected endpoint has one of
the paired \(A,F,I\) certificates.  Then every payer location,
including an \(F\)- or \(I\)-labelled output payer, satisfies the
quartic path-tree atlas.
\end{corollary}

\begin{proof}
Use \Cref{cor:V81-FI-paired} on every prefix-poor row and
\Cref{cor:V81-AFI-paired} on the opposite endpoint.  Distribute the
finite column stock exactly as in
\Cref{thm:V79-column-closure}; for each row-pair card apply
\Cref{thm:V81-paired-22}.  Foreign-labelled ancestry uses
\eqref{eq:V80-foreign-ancestry-moments}; internal-labelled ancestry
uses \Cref{thm:V81-I-ancestry}.
\end{proof}

\section{Regression tests}
\label{sec:V81-tests}

\begin{proposition}[The anti-aligned internal bank passes the paid
ledger]
\label{prop:V81-anti-aligned-test}
Take the fixed \(N\)-coordinate anti-aligned \(SU(2)\) bank of
V74--V75 as \(D_C\).  Then
\[
 \Xi_k=\Xi_{k^\circ}=Na_s,\qquad
 H_I=Na_s^2,\qquad s_\alpha H_I\longrightarrow\infty.
\]
The connected covariance remains order one on the anti-aligned
product vector, but every line of the V81 paid block proof is bounded
by
\begin{equation}
 C\frac{H_I}{s_\alpha\Xi_k}
 =C\frac{a_s}{s_\alpha}.
\label{eq:V81-anti-paid-scale}
\end{equation}
After insertion into
\eqref{eq:V81-three-payer-ledger}, the factor
\(s_\alpha^{-1}\) in \eqref{eq:V81-anti-paid-scale} is cancelled
exactly once by the unpayed crossing response.  No false terminal
decay is used.
\end{proposition}

\begin{proof}
The identities are
\eqref{eq:V74-bank-masses}--%
\eqref{eq:V74-bank-overlap}.  The nondecay is
\Cref{prop:V75-no-leading-cancel}.  Because the bank is
\(k\)-dominant, the paid binary compiler gives
\eqref{eq:V81-anti-paid-scale}.  The first line of
\eqref{eq:V81-three-payer-ledger} contains the complementary factor
\(s_\alpha H_{ij}(B_\times)\), so only one \(s_\alpha\) cancels.
The remaining stocks form
\eqref{eq:V81-common-tree-numerator}.
\end{proof}

\begin{proposition}[One-large/one-small partner regression]
\label{prop:V81-large-small-test}
At one internal coordinate let
\[
 a_{k}=L,\qquad a_{k^\circ}=b,\qquad L\ge b\ge a_*.
\]
Then
\[
 I_k=L,\qquad H_I=Lb,
\]
and the paired compiler oriented on the large row gives
\begin{align}
 \|\mathcal N_C^{E,I}\|
 &\le Cb+C\frac{H_E}{L},
\label{eq:V81-large-small-unpaid}\\
 \|\mathcal P_C^{E,I}\|
 &\le \frac C{s_\alpha}
 \left(b+\frac{H_E}{L}\right).
\label{eq:V81-large-small-paid}
\end{align}
In particular, no factor \(L/b\) is created.
\end{proposition}

\begin{proof}
Substitute \(H_I/I_k=b\) into
\eqref{eq:V81-internal-unpaid-certificate}--%
\eqref{eq:V81-internal-paid-certificate}.  The Casimir floor makes
the orientation legal uniformly in \(b\).  If \(L/b\) is too large
for the coordinate to be V69-balanced, its unbalanced heat tail gives
an additional certificate; the displayed algebraic regression does
not rely on that stronger route.
\end{proof}

\begin{assessment}[What V81 removes]
\label{ass:V81-removal}
V80 left open the possibility that a cut-internal bank could pay its
row only while unpayed, and that assigning the final resolvent to the
same bank would force two parallel \(kk^\circ\) stocks.  Theorem
\ref{thm:V81-paid-internal-block} removes both concerns.  The unused
parallel carrier is contractive, and the connected and replicated
payer branches each yield the single paid certificate
\(\widehat H_C/(s_\alpha\Xi_k)\).  Theorem
\ref{thm:V81-paired-22} then closes every clean \(2+2\) card for which
the chosen endpoint on each side has a paired block certificate.
\end{assessment}

\section{Allocation firewall and the exact next obstruction}
\label{sec:V81-final-reduction}

\begin{firewall}[The paid block is a once-only envelope]
\label{fire:V81-paid-block-once}
The quantities
\(\mathcal P_C^{E,I}\), \(\mathcal P_C\), and \(\mathcal P_D\)
are the positive envelopes obtained after the exact V2 allocation of
the unique resolvent numerator and denominator.  They are not extra
operators to be multiplied into an already paid source.  Likewise,
foreign complete products are factored before the reduced
pair-connecting prefix is selected.  Consequently
\Cref{thm:V81-paid-internal-block} is used once per source term and
cannot be iterated on the two replicas of the same physical
coordinate.
\end{firewall}

\begin{theorem}[Integrated clean \(2+2\) reduction]
\label{thm:V81-integrated-clean-22}
Consider a balanced clean \(2+2\) source term after the V68--V69
unbalanced tails and the V76--V79 overlap, column, and wrong-partner
reductions have been removed.  Fix a crossing row-pair card
\((i,j)\).  If the selected endpoint on each side admits one of the
paired \(A,F,I\) certificates of
\Cref{cor:V81-AFI-paired}, with disjoint selected banks, then its
entire one-payer contribution satisfies
\[
 \frac{\mathfrak b_{ij}^{\rm pay}}
 {\Xi_1\Xi_2\Xi_3\Xi_4}
 \le C\theta_{12}\theta_{34}\theta_{ij}.
\]
This includes the connected and replicated cut-internal branches and
all three possible locations of the unique denominator.
\end{theorem}

\begin{proof}
The connected and replicated internal blocks obey the same unpaid and
paid paired certificates by
\Cref{thm:V81-paid-internal-block}.
The foreign and prefix alternatives obey those certificates by
\Cref{thm:V81-paid-foreign-block,cor:V81-AFI-paired}.  The exact V2
one-payer allocation then has precisely the three terms in
\eqref{eq:V81-three-payer-ledger}; there is no fourth term because
there is only one resolvent denominator.  Apply
\Cref{thm:V81-paired-22}.  The disjointness hypothesis and
\Cref{fire:V81-paid-block-once,fire:V81-cross-cards-scalar} exclude a
second use of either a response stock or a capacity denominator.
\end{proof}

\begin{reduction}[What is not covered by the integrated theorem]
\label{red:V81-residual-crossing}
Within the balanced clean \(2+2\) family, a card not covered by
\Cref{thm:V81-integrated-clean-22} must fail to provide a paired
certificate at at least one selected endpoint.  After the established
prefix, foreign, internal, and unbalanced alternatives are removed,
the only stock still capable of supplying the missing selected-row
mass is crossing stock.  Thus the next algebraic problem is a
crossing-dominant endpoint problem.  Marked-heavy transfers and
overlapping selected banks remain separate exceptional families and
are not included in this reduction.
\end{reduction}

\begin{proof}
The paired certificate list is exhaustive only for the already
isolated \(A,F,I\) alternatives.  Failure of all three says that none
of those disjoint local banks supplies the required fraction of the
selected row mass.  The sourcewise row-mass decomposition used in
V68--V79 then assigns that fraction to the residual crossing stock,
unless the term belongs to one of the marked or overlapping
exceptions explicitly excluded in the statement.  This is a
classification of the remaining stock location, not an estimate of
its signed operator.
\end{proof}

\begin{openproblem}[V82: two-endpoint crossing payer]
\label{open:V81-V82-mandate}
Construct an exact sourcewise allocation for the crossing-dominant
case which pays both selected endpoint denominators before the
row-pair card is formed.  At the scalar-card level, the target
certificates are
\begin{align}
 \mathfrak n_{ij}^{\times,2}
 &\le
 C\,\frac{H_{ij}(B_\times)}
 {s_\alpha\Xi_i\Xi_j},
\label{eq:V81-double-cross-unpaid-target}\\
 \mathfrak p_{ij}^{\times,2}
 &\le
 C\,\frac{H_{ij}(B_\times)}
 {s_\alpha^2\Xi_i\Xi_j}.
\label{eq:V81-double-cross-paid-target}
\end{align}
Indeed, multiplication by two raw side responses
\(s_\alpha H_{12}\) and \(s_\alpha H_{34}\) turns
\eqref{eq:V81-double-cross-paid-target} into the common tree
numerator, while one paid side and one raw side require
\eqref{eq:V81-double-cross-unpaid-target}.

These formulas are a target, not a theorem.  V82 must derive the
relevant two-puncture or double-replica identity from the exact
complete Wilson source, keep the unique payer visible, and test it on:
\begin{enumerate}
\item a single \(SU(2)\) crossing edge;
\item the V74--V75 anti-aligned product bank;
\item unequal large/small representation masses;
\item a Rademacher sign model in which pairwise cancellation is
maximally misleading.
\end{enumerate}
If either target fails, the programme must go upstream to the
unallocated complete signed source and identify the replacement
multilinear cancellation; it must not promote a scalar card to a
signed two-row operator.
\end{openproblem}

\section{Status and claim boundary}
\label{sec:V81-status}

\begin{conclusion}[V81 conclusion]
\label{con:V81}
The paid cut-internal branch is no longer a gap in the clean
\(2+2\) ledger.  Exact complete-product moments, contraction of the
unused parallel carrier, and the one-payer three-carrier identity give
the required path-tree scale whenever both selected endpoints possess
paired certificates.  The remaining clean algebraic obstruction is
now localized to the crossing-dominant endpoint allocation
\eqref{eq:V81-double-cross-unpaid-target}--%
\eqref{eq:V81-double-cross-paid-target}; exceptional marked and
overlap families, the other source partitions, the continuum limit,
and the Osterwalder--Schrader spectral-gap passage remain outside the
present theorem.
\end{conclusion}

\begin{firewall}[No mass-gap claim]
\label{fire:V81-no-mass-gap}
This note proves a finite-cutoff sourcewise reduction.  It does not
construct four-dimensional continuum Yang--Mills theory and does not
prove a positive mass gap.  Such a claim would additionally require
closure of every residual source family, uniform cutoff removal,
reflection positivity, reconstruction, and a uniform positive
spectral gap.
\end{firewall}

\begin{record}[Parent and continuation record]
\label{rec:V81-parent}
This note was copied byte-for-byte from
\texttt{Renewal\_Yang\_Mills\_Mass\_Gap\_Research\_Notes\_V80.tex}
before the present chapter was appended.  The parent had \(35027\)
lines, \(1211024\) bytes, and SHA--256 digest
\[
\texttt{44266d80d20d153f2de7cb70af3e53610b7aabf0119a515c2c77661d7a83104f}.
\]
The next active version is V82 and its first task is
\Cref{open:V81-V82-mandate}.  The next compulsory cross-program audit
of the current SM and NS notes remains V85.
\end{record}

% =============================================================================
% BEGIN APPENDED COMPACT VERSION V82 MATERIAL
% =============================================================================

\clearpage
\chapter{Two-endpoint payment from a common crossing subbank}
\label{chap:V82}

\section{Context and correction of the naive target}
\label{sec:V82-context}

V81 reduced the clean balanced \(2+2\) obstruction to a crossing
bank which must sometimes pay two selected endpoint masses.  It
recorded the scalar targets
\eqref{eq:V81-double-cross-unpaid-target}--%
\eqref{eq:V81-double-cross-paid-target}, but deliberately did not
assert them.

There are two distinct issues.  First, V78 already proves that a
whole-row double-replica summand cannot be localized as an
\((i,j)\)-labelled signed operator.  Thus crossing dominance at two
rows does not authorize attaching both gains to an arbitrary scalar
card.  Second, when a \emph{common physical subbank} is carried
precisely by rows \(i,j\), the older complete-bank estimates of
V23--V24 do contain the two required powers.  The missing step is to
insert that common subbank into the exact V77 covariance without
using either it or the sole payer twice.

This note performs that insertion.  It treats balanced, strongly
unbalanced, and mixed \(SU(2)\) representation banks.  The proof is
not a new estimate of the V78 row-pair telescope: it selects an actual
coordinate subbank, estimates its complete physical moment, and keeps
the rest of the crossing covariance assembled.

\section{A pure-pair crossing bank and its two scales}
\label{sec:V82-pure-pair}

Fix distinct rows \(i,j\) on opposite sides of a two-block cut.  Let
\(D\) be a nonempty coordinate bank such that, among the four active
rows of the source, every coordinate \(e\in D\) is active precisely
on \(i,j\).  Put
\begin{align}
 x_e&:=a_{i,e},&
 y_e&:=a_{j,e},&
 w_e&:=\max\{x_e,y_e\},
\label{eq:V82-local-masses}\\
 T_D&:=\sum_{e\in D}w_e,&
 H_D&:=\sum_{e\in D}x_ey_e=H_{ij}(D).
\label{eq:V82-bank-scales}
\end{align}
Both \(x_e,y_e\) are nontrivial representation masses, and hence are
at least the group-dependent floor \(a_*>0\).

Split \(D=D_{\rm b}\mathbin{\dot\cup}D_{\rm u}\), where
\begin{align}
 e\in D_{\rm b}
 &\quad\Longleftrightarrow\quad
 \min\{x_e,y_e\}\ge\frac14w_e,
\label{eq:V82-balanced-split}\\
 e\in D_{\rm u}
 &\quad\Longleftrightarrow\quad
 \min\{x_e,y_e\}<\frac14w_e.
\label{eq:V82-unbalanced-split}
\end{align}
Write
\[
 S:=\sum_{e\in D_{\rm b}}w_e,\qquad
 W:=\sum_{e\in D_{\rm u}}w_e,\qquad
 H_{\rm b}:=\sum_{e\in D_{\rm b}}x_ey_e,\qquad
 H_{\rm u}:=\sum_{e\in D_{\rm u}}x_ey_e.
\]
Thus
\begin{equation}
 T_D=S+W,\qquad H_D=H_{\rm b}+H_{\rm u}.
\label{eq:V82-SW-split}
\end{equation}
An empty subbank contributes zero to its displayed quantities.

\begin{lemma}[The two elementary bank comparisons]
\label{lem:V82-bank-comparisons}
If \(N_{\rm b}=|D_{\rm b}|\), then
\begin{align}
 S^2&\le4N_{\rm b}H_{\rm b},
\label{eq:V82-balanced-power}\\
 H_{\rm u}&\ge a_*W,
\label{eq:V82-unbalanced-linear}\\
 H_D&\ge a_*T_D.
\label{eq:V82-total-linear}
\end{align}
\end{lemma}

\begin{proof}
On \(D_{\rm b}\),
\(w_e^2\le4x_ey_e\); Cauchy--Schwarz gives
\eqref{eq:V82-balanced-power}.  On every coordinate of \(D\), the
smaller of \(x_e,y_e\) is at least \(a_*\), so
\[
 x_ey_e=w_e\min\{x_e,y_e\}\ge a_*w_e.
\]
Summing over \(D_{\rm u}\) and over all of \(D\) proves the other two
claims.
\end{proof}

Let
\begin{align}
 \mathcal M_D
 &:=
 \mathbb E_D\!\left[F_i^D\otimes F_j^D\right],
\label{eq:V82-joint-moment}\\
 \mathcal R_D
 &:=
 \mathbb E_DF_i^D\otimes\mathbb E'_DF_j^D,
\label{eq:V82-separated-moment}\\
 \mathcal J_D&:=\mathcal M_D-\mathcal R_D.
\label{eq:V82-binary-covariance}
\end{align}
The prime denotes an independent row replica.  Thus
\(\mathcal J_D\) is the exact binary covariance
\eqref{eq:V75-bank-covariance}.

The balanced/unbalanced spectral estimates invoked below use the
\(SU(2)\) Clebsch--Gordan calculus of V23--V24.  For a complete
physical moment \(\mathcal A\), write
\(\mathcal P(\mathcal A)\) for its positive physical envelope after
the output mass and the unique resolvent denominator have both been
assigned to that moment.  The globally protected output, on which
the heat multiplier is one and the output Casimir is zero, is omitted
from the resolvent exactly as in V22--V24.

\section{The complete moment has two endpoint powers}
\label{sec:V82-moment-estimates}

\begin{lemma}[Unpaid joint-moment two-endpoint estimate]
\label{lem:V82-joint-unpaid}
In the standing range \(0<s_\alpha\le1\),
\begin{equation}
 \boxed{
 \kappa_\otimes^{\,i|j}(\mathcal M_D)
 \le
 C\,\frac{H_D}{s_\alpha T_D^2}.}
\label{eq:V82-joint-unpaid}
\end{equation}
\end{lemma}

\begin{proof}
Coordinate independence gives
\(\mathcal M_D=\mathcal M_{D_{\rm b}}\otimes
\mathcal M_{D_{\rm u}}\).
If \(S\ge T_D/2\), use the complete balanced first-mark estimate
\eqref{eq:V24-balanced-first-mark} and contraction on the
unbalanced factor:
\[
 \kappa_\otimes(\mathcal M_D)
 \le
 C\frac{H_{\rm b}}{s_\alpha S^2}
 \le
 C\frac{H_D}{s_\alpha T_D^2}.
\]
If \(W\ge T_D/2\), use ordinary norm on the balanced factor and the
\(k=1\) unbalanced product moment
\eqref{eq:V24-unbalanced-moment}:
\[
 \kappa_\otimes(\mathcal M_D)
 \le\|\mathcal M_{D_{\rm u}}\|
 \le\frac{C}{s_\alpha W}
 \le C\frac{H_{\rm u}}{s_\alpha W^2}
 \le C\frac{H_D}{s_\alpha T_D^2}.
\]
The middle inequality is
\eqref{eq:V82-unbalanced-linear}.  One of the two alternatives always
holds, including when one subbank is empty.
\end{proof}

\begin{lemma}[Paid joint-moment two-endpoint estimate]
\label{lem:V82-joint-paid}
The complete joint-moment payer satisfies
\begin{equation}
 \boxed{
 \kappa_\otimes^{\,i|j}
 \bigl(\mathcal P(\mathcal M_D)\bigr)
 \le
 C\,\frac{H_D}{s_\alpha^2T_D^2}.}
\label{eq:V82-joint-paid}
\end{equation}
\end{lemma}

\begin{proof}
First suppose \(D=D_{\rm b}\).  The claimed estimate is precisely the
marked complete graph-resolvent estimate
\eqref{eq:V23-resolvent-marked}, with
\(S=T_D\) and \(H=H_D\).  If \(D=D_{\rm u}\), then
\eqref{eq:V24-unbalanced-resolvent-bounds} gives
\[
 \bigl\|\mathcal P(\mathcal M_D)\bigr\|
 \le
 C\min\left\{\frac1{s_\alpha},
              \frac1{s_\alpha^3W^2}\right\}.
\]
For \(z=s_\alpha W\),
\[
 \min\{1,z^{-2}\}\le\frac{C}{z}.
\]
Together with \eqref{eq:V82-unbalanced-linear}, this gives
\[
 \bigl\|\mathcal P(\mathcal M_D)\bigr\|
 \le\frac{C}{s_\alpha^2W}
 \le C\frac{H_{\rm u}}{s_\alpha^2W^2},
\]
which is the result in the purely unbalanced case.

It remains to treat nonempty \(D_{\rm b},D_{\rm u}\).  Let
\(\mathscr Z_{\rm b}\) and \(\mathscr Z_{\rm u}\) be their complete
output-mass resolvents, let \(\mathsf P_{\rm b}\) be the globally
protected balanced projection, and let
\(\mathscr G_{\rm u}\) be the unbalanced geometric resolver.  The
exact positive-operator split
\eqref{eq:V24-mixed-decomposition} says
\begin{equation}
 \mathcal P(\mathcal M_D)
 \preccurlyeq
 \mathscr Z_{\rm b}\otimes\mathcal M_{D_{\rm u}}
 +\mathcal M_{D_{\rm b}}\otimes\mathscr Z_{\rm u}
 +N_{\rm b}\mathsf P_{\rm b}\otimes\mathscr G_{\rm u}.
\label{eq:V82-mixed-paid-split}
\end{equation}

Suppose first that \(S\ge T_D/2\).  For the first term use
\eqref{eq:V23-resolvent-marked}; for the second use
\eqref{eq:V24-balanced-first-mark} and
\(\|\mathscr Z_{\rm u}\|\le C/s_\alpha\).  They are bounded by
\[
 C\frac{H_{\rm b}}{s_\alpha^2S^2}.
\]
For the third term,
\eqref{eq:V24-protected-balanced-mark} and the \(k=1\) unbalanced
geometric estimate give
\[
 N_{\rm b}\kappa_\otimes(\mathsf P_{\rm b})
 \|\mathscr G_{\rm u}\|
 \le
 C\frac{H_{\rm b}}{S^2}\frac1{s_\alpha W}
 \le
 C\frac{H_{\rm b}}{s_\alpha^2S^2}.
\]
Here \(W\ge a_*\) and \(s_\alpha\le1\).  Hence all three terms are at
most
\(CH_D/(s_\alpha^2T_D^2)\).

Suppose instead that \(W\ge T_D/2\).  Use the unmarked balanced
resolver bound
\eqref{eq:V23-resolvent-unmarked} and the \(k=1\) unbalanced moment
on the first term:
\[
 \kappa_\otimes(\mathscr Z_{\rm b})
 \|\mathcal M_{D_{\rm u}}\|
 \le\frac{C}{s_\alpha^2W}
 \le C\frac{H_{\rm u}}{s_\alpha^2W^2}.
\]
The second term is bounded by the purely unbalanced resolver estimate
proved above.  Finally,
\(N_{\rm b}\kappa_\otimes(\mathsf P_{\rm b})\le C\) by the V22
vertex-cut estimate, whereas the \(k=2\) unbalanced geometric moment
gives
\[
 \|\mathscr G_{\rm u}\|\le\frac{C}{s_\alpha^2W^2}
 \le C\frac{H_{\rm u}}{s_\alpha^2W^2};
\]
the last comparison uses \(H_{\rm u}\ge a_*W\) and \(W\ge a_*\).
Apply the V2 one-sided tensor inequality in each displayed
orientation.  Since \(W\ge T_D/2\), the three terms in
\eqref{eq:V82-mixed-paid-split} have the asserted bound.
\end{proof}

\begin{lemma}[The separated moment has the same two powers]
\label{lem:V82-separated-two-powers}
The separated product in \eqref{eq:V82-separated-moment} obeys
\begin{align}
 \|\mathcal R_D\|
 &\le
 C\,\frac{H_D}{s_\alpha T_D^2},
\label{eq:V82-separated-unpaid}\\
 \bigl\|\mathcal P(\mathcal R_D)\bigr\|
 &\le
 C\,\frac{H_D}{s_\alpha^2T_D^2}.
\label{eq:V82-separated-paid}
\end{align}
\end{lemma}

\begin{proof}
The separated moment is a scalar contraction
\[
 \mathcal R_D=q_iq_j I.
\]
Regard the two row products as one complete family \(G\), with
\[
 Y_G=\sum_{e\in D}(x_e+y_e),\qquad
 q_G=q_iq_j.
\]
Then
\begin{equation}
 T_D\le Y_G\le2T_D.
\label{eq:V82-YT-comparison}
\end{equation}
The order-one complete product moment from
\eqref{eq:V52-private-products-all-orders} gives
\[
 q_G\le\frac{C}{s_\alpha Y_G}
 \le\frac{C}{s_\alpha T_D}.
\]
Now use \(H_D\ge a_*T_D\) to obtain
\eqref{eq:V82-separated-unpaid}.

For the paid moment, repeated fusion bounds its output mass by
\(CY_G\), and the unique denominator is \(1-q_G\).  Therefore its
positive envelope is bounded by a fixed multiple of the complete
product payer \(\mathscr D_G^{\rm pay}\) from
\eqref{eq:V81-product-payer}.  The order-one paid product moment
\eqref{eq:V81-product-paid-one} gives
\[
 \bigl\|\mathcal P(\mathcal R_D)\bigr\|
 \le\frac{C}{s_\alpha^2Y_G}
 \le\frac{C}{s_\alpha^2T_D}
 \le C\frac{H_D}{s_\alpha^2T_D^2}.
\]
This proves the second estimate without assigning independent
denominators to rows \(i,j\).
\end{proof}

\begin{theorem}[Pure-pair bank pays both endpoints]
\label{thm:V82-pure-pair-double}
The exact binary covariance \(\mathcal J_D\) and its complete paid
envelope satisfy
\begin{align}
 \boxed{
 \kappa_\otimes^{\,i|j}(\mathcal J_D)
 \le
 C\,\frac{H_D}{s_\alpha T_D^2},}
\label{eq:V82-covariance-unpaid}\\
 \boxed{
 \kappa_\otimes^{\,i|j}
 \bigl(\mathcal P(\mathcal J_D)\bigr)
 \le
 C\,\frac{H_D}{s_\alpha^2T_D^2}.}
\label{eq:V82-covariance-paid}
\end{align}
\end{theorem}

\begin{proof}
For the unpaid estimate use
\(\mathcal J_D=\mathcal M_D-\mathcal R_D\), capacity
subadditivity, \Cref{lem:V82-joint-unpaid}, and
\eqref{eq:V82-separated-unpaid}.  For the paid estimate, resolve the
common physical output blocks first.  On every block, the positive
output-mass/resolvent weight times
\(|\mathcal M_D-\mathcal R_D|\) is bounded by the corresponding two
weighted positive envelopes.  Hence
\[
 \mathcal P(\mathcal J_D)
 \preccurlyeq_{\rm env}
 \mathcal P(\mathcal M_D)+\mathcal P(\mathcal R_D).
\]
Apply
\Cref{lem:V82-joint-paid} and
\eqref{eq:V82-separated-paid}.  The globally protected line cannot
be a selected-bank payer; it is retained as an unpaid factor, or the
payer is assigned to the remainder, consistently in both terms.
\end{proof}

\begin{firewall}[This is a physical subbank estimate]
\label{fire:V82-not-card-operator}
Theorem \ref{thm:V82-pure-pair-double} does not construct the forbidden
\((i,j)\)-summand of the whole-row double-replica telescope.  The bank
\(D\) is an actual set of coordinates supported precisely on rows
\(i,j\), and \(\mathcal J_D\) is its complete two-term covariance.
Only this assembled physical operator is enveloped.  The numerator
\(H_D\) is introduced after that estimate.
\end{firewall}

\section{Extraction from the full crossing covariance}
\label{sec:V82-extraction}

Let \(\rho=C_0|C_1\) be the fixed two-block cut, let
\(i\in C_0\), \(j\in C_1\), and let \(B\) be its complete V77
residual bank.  Suppose \(D\subseteq B\) is a pure \(ij\)-bank in the
sense of \Cref{sec:V82-pure-pair}, and put \(R:=B\setminus D\).
No assumption is made on the incidences in \(R\).

Denote the common-variable and cut-replicated endpoint moments on a
bank \(A\) by \(\mathcal M_A^\rho\) and
\(\mathcal R_A^\rho\), respectively, and put
\[
 \mathcal J_A^\rho:=\mathcal M_A^\rho-\mathcal R_A^\rho.
\]
On the pure pair bank \(D\), these are exactly
\(\mathcal M_D,\mathcal R_D,\mathcal J_D\) from
\eqref{eq:V82-joint-moment}--%
\eqref{eq:V82-binary-covariance}.

\begin{lemma}[Exact selected-subbank covariance extraction]
\label{lem:V82-subbank-extraction}
At the signed operator level,
\begin{equation}
 \boxed{
 \mathcal J_B^\rho
 =
 \mathcal J_D\otimes\mathcal M_R^\rho
 +
 \mathcal R_D\otimes\mathcal J_R^\rho.}
\label{eq:V82-subbank-extraction}
\end{equation}
This is an identity before positive enveloping and before the
one-payer allocation.
\end{lemma}

\begin{proof}
Coordinate independence and \(B=D\mathbin{\dot\cup}R\) give
\[
 \mathcal M_B^\rho=\mathcal M_D\otimes\mathcal M_R^\rho,
 \qquad
 \mathcal R_B^\rho=\mathcal R_D\otimes\mathcal R_R^\rho.
\]
Insert and subtract
\(\mathcal R_D\otimes\mathcal M_R^\rho\):
\begin{align*}
 \mathcal M_B^\rho-\mathcal R_B^\rho
 &=
 (\mathcal M_D-\mathcal R_D)\otimes\mathcal M_R^\rho
 +\mathcal R_D\otimes
   (\mathcal M_R^\rho-\mathcal R_R^\rho),
\end{align*}
which is \eqref{eq:V82-subbank-extraction}.
\end{proof}

\begin{theorem}[The selected subbank pays both endpoints of the full
crossing bank]
\label{thm:V82-full-bank-double}
Let \(\mathcal N_B^{[D]}\) be the positive unpaid envelope obtained
from \eqref{eq:V82-subbank-extraction}.  Let
\(\mathcal P_B^{[D]}\) be the exact V2 positive envelope in which the
single output-mass/resolvent payer is assigned once among the selected
factor and the remainder factor in the two terms of that identity.
Then
\begin{align}
 \boxed{
 \kappa_\otimes(\mathcal N_B^{[D]})
 \le
 C\,\frac{H_D}{s_\alpha T_D^2},}
\label{eq:V82-full-bank-unpaid}\\
 \boxed{
 \kappa_\otimes(\mathcal P_B^{[D]})
 \le
 C\,\frac{H_D}{s_\alpha^2T_D^2}.}
\label{eq:V82-full-bank-paid}
\end{align}
\end{theorem}

\begin{proof}
The complete moments
\(\mathcal M_R^\rho,\mathcal R_D\) are contractions, and the
difference \(\mathcal J_R^\rho\) has capacity at most a fixed
constant.  Apply the V2 one-sided tensor inequality to the two terms
of \eqref{eq:V82-subbank-extraction}, using capacity on
\(\mathcal J_D\) in the first and on \(\mathcal J_R^\rho\) in the
second.  Equations
\eqref{eq:V82-covariance-unpaid} and
\eqref{eq:V82-separated-unpaid} prove
\eqref{eq:V82-full-bank-unpaid}.

For the paid envelope there are four mutually exclusive locations:
\begin{enumerate}[label=\textup{(\roman*)}]
\item \(\mathcal J_D\) pays and \(\mathcal M_R^\rho\) is a
contraction;
\item \(\mathcal J_D\) is unpaid and
\(\mathcal M_R^\rho\) pays;
\item \(\mathcal R_D\) pays and \(\mathcal J_R^\rho\) is unpaid;
\item \(\mathcal R_D\) is unpaid and
\(\mathcal J_R^\rho\) pays.
\end{enumerate}
Cases \textup{(i)} and \textup{(iii)} are bounded by
\eqref{eq:V82-covariance-paid} and
\eqref{eq:V82-separated-paid}.  In case \textup{(ii)}, the complete
moment payer has the universal \(C/s_\alpha\) cap from
\eqref{eq:V3-spectator-carrier-cap} and the complete-product form
\eqref{eq:V81-product-paid-zero}; combine it with
\eqref{eq:V82-covariance-unpaid}.  In case \textup{(iv)}, use the
universal paid covariance cap
\eqref{eq:V78-universal-paid-cap} and
\eqref{eq:V82-separated-unpaid}.  The V2 allocation forms this
four-term sum before the one-sided tensor bound.  Each case is at most
\(CH_D/(s_\alpha^2T_D^2)\), proving the claim.
\end{proof}

\begin{corollary}[Conditional resolution of the V81 double-cross
targets]
\label{cor:V82-V81-targets}
Define
\[
 X_i(D):=\sum_{e\in D}a_{i,e},
 \qquad
 X_j(D):=\sum_{e\in D}a_{j,e}.
\]
If, for some \(\kappa>0\),
\begin{equation}
 X_i(D)\ge\kappa\Xi_i,
 \qquad
 X_j(D)\ge\kappa\Xi_j,
\label{eq:V82-common-dominance}
\end{equation}
then
\begin{align}
 \boxed{
 \kappa_\otimes(\mathcal N_B^{[D]})
 \le
 C_\kappa
 \frac{H_{ij}(D)}
 {s_\alpha\Xi_i\Xi_j},}
\label{eq:V82-double-cross-unpaid}\\
 \boxed{
 \kappa_\otimes(\mathcal P_B^{[D]})
 \le
 C_\kappa
 \frac{H_{ij}(D)}
 {s_\alpha^2\Xi_i\Xi_j}.}
\label{eq:V82-double-cross-paid}
\end{align}
Thus
\eqref{eq:V81-double-cross-unpaid-target}--%
\eqref{eq:V81-double-cross-paid-target} are valid when the two row
gains are backed by one common pure-pair physical subbank.  They are
not asserted for an arbitrary scalar card of the full bank.
\end{corollary}

\begin{proof}
Since \(T_D\ge X_i(D)\) and \(T_D\ge X_j(D)\),
\[
 T_D^2
 \ge X_i(D)X_j(D)
 \ge\kappa^2\Xi_i\Xi_j.
\]
Insert this in
\eqref{eq:V82-full-bank-unpaid}--%
\eqref{eq:V82-full-bank-paid}, and recall \(H_D=H_{ij}(D)\).
\end{proof}

\begin{firewall}[No duplicated crossing stock]
\label{fire:V82-one-subbank-use}
The overlap \(H_{ij}(D)\) in
\eqref{eq:V82-double-cross-unpaid}--%
\eqref{eq:V82-double-cross-paid} is the response numerator of the
complete selected subbank estimate.  It is not multiplied by a second
\(ij\) response extracted from the same coordinates.  The remainder
\(R\) is used only as a contraction or as the unique payer in the
corresponding V2 alternative.
\end{firewall}

\section{Insertion into the clean \(2+2\) source}
\label{sec:V82-source}

Let \(\widehat H_C\le H_{12}\) and
\(\widehat H_{C'}\le H_{34}\) be the genuine disjoint response stocks
of the two same-side augmented blocks.  In the present branch no
one-row paired certificate is assumed.  Retain only their raw
unpaid/paid response estimates
\begin{align}
 \|\mathcal N_C\|&\le Cs_\alpha\widehat H_C,
 &
 \|\mathcal P_C\|&\le C\widehat H_C,
\label{eq:V82-C-raw}\\
 \|\mathcal N_{C'}\|&\le Cs_\alpha\widehat H_{C'},
 &
 \|\mathcal P_{C'}\|&\le C\widehat H_{C'}.
\label{eq:V82-D-raw}
\end{align}
These are the response entries of the V78 binary compilers, before a
row-dominance minimum is taken.

\begin{theorem}[Common-bank two-cross/no-prefix closure]
\label{thm:V82-common-bank-22}
Suppose the selected crossing endpoints \(i\in\{1,2\}\),
\(j\in\{3,4\}\) obey
\eqref{eq:V82-common-dominance} on a pure \(ij\)-subbank
\(D\subseteq B_\times\) disjoint from the two same-side selected
stocks.  Then all three possible payer locations in
\[
 \mathcal B_C\otimes\mathcal B_{C'}
 \otimes\mathcal J_{B_\times}^{12|34}
\]
satisfy
\begin{equation}
 \boxed{
 \frac{\mathfrak b_{ij}^{\rm common}}
 {\Xi_1\Xi_2\Xi_3\Xi_4}
 \le
 C_\kappa\theta_{12}\theta_{34}\theta_{ij}.}
\label{eq:V82-common-bank-tree}
\end{equation}
\end{theorem}

\begin{proof}
If the \(C\)-block pays, use its paid raw response, the unpaid raw
response of \(C'\), and
\eqref{eq:V82-double-cross-unpaid}:
\[
 C\widehat H_C\,
 Cs_\alpha\widehat H_{C'}\,
 C_\kappa
 \frac{H_{ij}(D)}{s_\alpha\Xi_i\Xi_j}
 \le
 C_\kappa
 \frac{H_{12}H_{34}H_{ij}}{\Xi_i\Xi_j}.
\]
The \(C'\)-payer case is symmetric.  If the crossing bank pays, use
both unpaid raw responses and
\eqref{eq:V82-double-cross-paid}; the two factors
\(s_\alpha\) cancel the two factors
\(s_\alpha^{-1}\).  This gives the same numerator.

These are exactly the three alternatives of the V81 one-payer ledger.
Divide by the four row masses.  The result is
\(\theta_{12}\theta_{34}\theta_{ij}\).
\end{proof}

\begin{assessment}[Progress on the two-cross state]
\label{ass:V82-two-cross-progress}
The two-cross/no-prefix obstruction is closed whenever both selected
endpoint masses are carried by one common incidence-two crossing
subbank.  Unlike the V78 complementary theorem, neither same-side
block must supply a row denominator.  Unlike the proposed unrestricted
V81 card, the two gains come from one complete physical bank estimate
and survive every location of the sole payer.
\end{assessment}

\section{Regression tests}
\label{sec:V82-tests}

\begin{proposition}[Single-edge and protected-channel regression]
\label{prop:V82-single-edge}
Let \(D=\{e\}\).  Then
\[
 T_D=\max\{x_e,y_e\},
 \qquad
 H_D=x_ey_e,
\]
and
\begin{align}
 \kappa_\otimes(\mathcal J_D)
 &\le
 \frac{C}{s_\alpha}
 \frac{\min\{x_e,y_e\}}{\max\{x_e,y_e\}},
\label{eq:V82-single-unpaid}\\
 \kappa_\otimes(\mathcal P(\mathcal J_D))
 &\le
 \frac{C}{s_\alpha^2}
 \frac{\min\{x_e,y_e\}}{\max\{x_e,y_e\}}.
\label{eq:V82-single-paid}
\end{align}
For equal \(SU(2)\) spins, these bounds do not falsely suppress the
singlet channel.  For a strongly unequal pair they are supported by
the retained nontrivial-output heat separation.
\end{proposition}

\begin{proof}
The displays are
\eqref{eq:V82-covariance-unpaid}--%
\eqref{eq:V82-covariance-paid} with one coordinate.  In the equal-spin
case their right sides are \(C/s_\alpha\) and
\(C/s_\alpha^2\), both at least a fixed constant because
\(s_\alpha\le1\).  Thus the invariant protected component is not
claimed to decay.  In the unequal case, if the ratio is below
\(1/4\), \Cref{lem:V24-output-comparison} forces every joint output
mass to be comparable to the larger input mass.  The moment and
resolvent estimates in
\Cref{thm:V24-unbalanced-resolvent} then give exactly the inverse
ratio after the elementary minimum used in
\Cref{lem:V82-joint-paid}.
\end{proof}

\begin{proposition}[The V74--V75 anti-aligned bank is consistent]
\label{prop:V82-anti-aligned}
For the \(N\)-coordinate equal-spin anti-aligned bank of V74--V75,
\[
 T_D=Na_s,\qquad H_D=Na_s^2,
\]
and therefore
\begin{align}
 \kappa_\otimes(\mathcal J_D)
 &\le\frac{C}{s_\alpha N},
\label{eq:V82-anti-unpaid}\\
 \kappa_\otimes(\mathcal P(\mathcal J_D))
 &\le\frac{C}{s_\alpha^2N}.
\label{eq:V82-anti-paid}
\end{align}
The order-one lower bound
\eqref{eq:V75-no-leading-cancel} for fixed \(N\) does not contradict
\eqref{eq:V82-anti-unpaid}.
\end{proposition}

\begin{proof}
Insert
\eqref{eq:V74-bank-masses}--%
\eqref{eq:V74-bank-overlap} into
\Cref{thm:V82-pure-pair-double}.  When \(N\) is fixed and
\(s_\alpha\downarrow0\), the right side of
\eqref{eq:V82-anti-unpaid} grows, whereas
\eqref{eq:V75-no-leading-cancel} stays bounded below.  When
\(s_\alpha N\) is large, the V23 protected-network capacity and
one-failure sum supply the decay used in the proof.  Thus the theorem
does not assume the false terminal moment excluded in V74.
\end{proof}

\begin{proposition}[Alternating large/small representations]
\label{prop:V82-alternating}
Take two pure \(ij\)-coordinates with masses
\[
 (x_{e_1},y_{e_1})=(L,b),
 \qquad
 (x_{e_2},y_{e_2})=(b,L),
 \qquad L\ge4b\ge4a_*.
\]
Then both endpoint row masses equal \(L+b\), while
\[
 T_D=2L,\qquad H_D=2Lb.
\]
Consequently
\begin{align}
 \kappa_\otimes(\mathcal J_D)
 &\le C\frac{b}{s_\alpha L},
\label{eq:V82-alternating-unpaid}\\
 \kappa_\otimes(\mathcal P(\mathcal J_D))
 &\le C\frac{b}{s_\alpha^2L}.
\label{eq:V82-alternating-paid}
\end{align}
No ratio \(L/b\) is lost, despite the two endpoint dominances being
carried in opposite orientations.
\end{proposition}

\begin{proof}
Both coordinates lie in \(D_{\rm u}\).  Substitute their exact
stocks into \Cref{thm:V82-pure-pair-double}.  Analytically, each local
joint output is separated at the scale of its larger input by
\Cref{lem:V24-output-comparison}; the order-one unbalanced moment and
the two alternative resolvent moments sum those two orientations
before the physical envelope is taken.  The separated term is paid by
the complete product moment of
\Cref{lem:V82-separated-two-powers}.  Hence the displayed small ratio
comes from retained heat, not from a false reverse Cauchy--Schwarz
inequality.
\end{proof}

\begin{proposition}[The Rademacher wrong-card test is excluded]
\label{prop:V82-rademacher}
In the V78 Rademacher model with coordinates supported on row pairs
\(14\) and \(23\), the nominal whole-row \((1,3)\) telescope summand
is nonzero although \(H_{13}=0\).  This does not contradict any V82
estimate: there is no nonempty physical pure \(13\)-subbank \(D\), so
\eqref{eq:V82-common-dominance} fails.
\end{proposition}

\begin{proof}
The nonzero nominal summand and \(H_{13}=0\) are
\Cref{prop:V78-no-pair-operator}.  V82 never envelopes that summand.
Its selected object is the complete covariance on an actual
coordinate bank supported on the claimed pair.  Here that bank is
empty, so
\Cref{fire:V82-not-card-operator} forbids the proposed application.
\end{proof}

\section{Exact residual and V83 direction}
\label{sec:V82-residual}

\begin{reduction}[Residual inside the two-cross/no-prefix sector]
\label{red:V82-two-cross-residual}
Apply V81 to every endpoint carrying a paired \(A,F,I\) certificate,
V78--V79 to every complementary or column match, and
\Cref{thm:V82-common-bank-22} to every common dominant pure-pair
subbank.  A still-unproved clean two-cross/no-prefix card must then
have at least one of:
\begin{enumerate}[label=\textup{(\roman*)}]
\item the two large crossing endpoint stocks are carried by different
pure-pair partner banks;
\item a fixed fraction of at least one unresolved endpoint mass lies
on incidence-three or incidence-four crossing coordinates;
\item a common pure-pair bank is large for only one of the two
endpoints;
\item the bank needed for one of the preceding selections meets a
same-side selected stock or a marked-heavy coordinate; or
\item the remaining output-ancestry allocation is attached to one of
those non-disjoint banks.
\end{enumerate}
This list concerns the two-cross/no-prefix sector only; it does not
include the other first-connectivity source partitions.
\end{reduction}

\begin{proof}
Split each residual crossing incidence by its actual active-row set.
Incidence-two coordinates then split disjointly by their unique
opposite partner, whereas all remaining crossing coordinates have
incidence three or four.  If one pure pair bank is a fixed-fraction
bank for both selected endpoints, V82 closes it.  If not, the two
dominant endpoint stocks either choose different pure partners, only
one endpoint is dominant on their common bank, or higher incidence
carries the missing fraction.  These are items \textup{(i)}--%
\textup{(iii)}.  Failure of the clean disjoint selection gives
\textup{(iv)}, and the V79 ancestry firewall gives \textup{(v)}.
All earlier paired and complementary alternatives were removed before
this split.
\end{proof}

\begin{assessment}[The useful upstream lesson]
\label{ass:V82-upstream-lesson}
The missing two-endpoint power is not produced by differentiating the
V77 covariance twice.  It is a complete-bank capacity effect.  The
balanced protected network supplies inverse bank size, the gapped
sector supplies a one-failure heat sum, and strongly unequal inputs
supply retained output Casimir.  The exact extraction
\eqref{eq:V82-subbank-extraction} then transports that capacity
through the signed crossing covariance.  This is why the result
survives the anti-aligned and Rademacher tests.
\end{assessment}

\begin{openproblem}[V83: higher-incidence square gain and tree
exchange]
\label{open:V82-V83-mandate}
Treat items \textup{(i)}--\textup{(iii)} of
\Cref{red:V82-two-cross-residual} without reverting to a scalar-card
operator.  More precisely:
\begin{enumerate}
\item choose the larger unresolved endpoint \(k\), so that
      \(\Xi_k^2\ge\Xi_i\Xi_j\);
\item partition its balanced higher-incidence crossing bank by the
      V70 maximal-partner rule and select a fixed-fraction bank
      \(E_{k\ell}\);
\item extend
      \eqref{eq:V82-subbank-extraction} to the complete
      higher-incidence physical moment on \(E_{k\ell}\);
\item prove both unpaid and paid square-tail estimates
      \[
      \kappa_\otimes(\mathcal N_{E_{k\ell}})
      \le
      C\frac{H_{k\ell}(E_{k\ell})}
      {s_\alpha\Xi_k^2},
      \qquad
      \kappa_\otimes(\mathcal P_{E_{k\ell}})
      \le
      C\frac{H_{k\ell}(E_{k\ell})}
      {s_\alpha^2\Xi_k^2};
      \]
\item enumerate the finite tree exchange which replaces the nominal
      crossing card \(ij\) by the actual certified edge \(k\ell\),
      keeping every same-side edge and selected bank disjoint; and
\item test the construction on a four-row coordinate with two large
      and two small representations and on the V78 Rademacher
      wrong-card pattern.
\end{enumerate}
If the paid square tail fails for a higher-incidence local moment,
return to the V24 mixed-resolvent decomposition with the full
four-row physical output; do not multiply two row-cut capacities.
\end{openproblem}

\section{V82 conclusion and claim boundary}
\label{sec:V82-conclusion}

\begin{conclusion}[V82 conclusion]
\label{con:V82}
For a genuine common incidence-two crossing subbank, the V81
two-endpoint targets are now theorems.  The proof covers balanced,
unbalanced, and mixed representation banks, extracts the selected
bank exactly from the full V77 covariance, and preserves the unique
payer.  Consequently every clean \(2+2\) two-cross/no-prefix term
with a common dominant pure \(ij\)-bank has the path-tree scale for
all three payer locations.

The remaining crossing obstruction is no longer an undifferentiated
double debit.  It is the higher-incidence or mismatched-partner
problem in \Cref{red:V82-two-cross-residual}, whose next exact
interface is the square-gain/tree-exchange problem of
\Cref{open:V82-V83-mandate}.
\end{conclusion}

\begin{firewall}[No mass-gap claim]
\label{fire:V82-no-mass-gap}
V82 is a finite-cutoff sourcewise theorem for one residual sector.  It
does not close higher-incidence/overlapping crossing banks, the
\(3+1\), \(2+1+1\), or discrete source families, uniform cutoff
removal, Osterwalder--Schrader reconstruction, or a positive continuum
spectral gap.  The full Yang--Mills mass gap remains open.
\end{firewall}

\begin{record}[Parent and continuation record]
\label{rec:V82-parent}
This file was copied byte-for-byte from
\texttt{Renewal\_Yang\_Mills\_Mass\_Gap\_Research\_Notes\_V81.tex}
before the present chapter was appended.  The parent had \(35976\)
lines, \(1243239\) bytes, and SHA--256 digest
\[
\texttt{1610676a7cabac1c83ce05cf1defad8a01fc463ad0f11e6f735aee6d6a4d9c34}.
\]
The next active version is V83 and begins with
\Cref{open:V82-V83-mandate}.  The next compulsory read-only audit of
the current SM and NS notes remains V85.
\end{record}

% =============================================================================
% BEGIN APPENDED COMPACT VERSION V83 MATERIAL
% =============================================================================

\clearpage
\chapter{Higher-incidence square tails and the honest tree-exchange
gate}
\label{chap:V83}

\section{Context: lift the one-failure resolver, then compare masses}
\label{sec:V83-context}

V82 proved the two required endpoint powers when a common physical
incidence-two \(ij\)-bank carries both unresolved endpoint masses.
Its next mandate proposed using the V70 maximal-partner bank at the
larger endpoint when the crossing mass instead lies on
incidence-three or incidence-four coordinates.

The analytic part of that proposal is valid.  The V23 one-failure
resolvent proof uses only two local facts: the all-trivial output has
product-state capacity strictly below one across a chosen row cut,
and every complementary output has a uniform heat gap.  V70 proved
exactly those facts for an arbitrary balanced higher-incidence
physical moment.  Consequently both the unpaid and paid square tails
extend to a V70 row-rooted bank.

There is nevertheless a separate graph issue.  The analytic numerator
is the actual maximal-partner stock \(H_{k\ell}\), and its denominator
is \(\Xi_k^2\).  The legal tree edge \(k\ell\) is normalized by
\(\Xi_k\Xi_\ell\).  Thus a tree exchange is immediate only if
\(\Xi_\ell\lesssim\Xi_k\).  If the opposite partner is much heavier,
the proof has found a genuine mass-ascending transfer, not a completed
tree estimate.  If the maximal partner lies on the same side of the
cut, the new edge is parallel to an existing same-side edge and does
not connect the cut.  V83 keeps both failures visible.

\section{A row-rooted balanced physical bank}
\label{sec:V83-row-bank}

Fix a two-block cut \(\rho=C_0|C_1\), a row \(k\), and a nonempty
bank \(E\) of balanced residual coordinates which cross \(\rho\).
Throughout the analytic statements below, \(E\) is a clean retained
bank disjoint from the finite source-mark set and from every previously
selected transfer bank.  Marked or overlapping cells are retained in
the residual theorem rather than inserted into the tensor product.
At each \(e\in E\), choose a row
\[
 v_k(e)\in[4]\setminus\{k\}
\]
whose active representation mass is maximal among the other active
rows.  Fix \(\ell\ne k\) and restrict to the cell
\[
 E=E_{k\ell}:=\{e:v_k(e)=\ell\}.
\]
Use the V70 stocks
\begin{equation}
 S:=\sum_{e\in E}a_{k,e},
 \qquad
 H:=\sum_{e\in E}a_{k,e}a_{\ell,e}
 =H_{k\ell}^{(k)}(E),
 \qquad N:=|E|.
\label{eq:V83-SHN}
\end{equation}
Then
\begin{equation}
 \boxed{S^2\le C_0NH}
\label{eq:V83-row-power}
\end{equation}
by \eqref{eq:V70-row-power}.  Here and below \(C_0\) is a fixed
group/row constant, not the left block of the cut.

Fix one global partition state of either endpoint moment in the V77
covariance.  At coordinate \(e\), resolve its complete local physical
output and, for \(t>0\), raise every Wilson multiplier to the power
\(t\).  Denote the resulting positive contraction by
\(\mathsf M_e^{(t)}\).  Let \(\mathsf P_e\) be the projection on which
every local output block is trivial and put
\[
 \mathsf G_e^{(t)}:=
 \mathsf M_e^{(t)}-\mathsf P_e.
\]
If the partition state admits no all-trivial output, set
\(\mathsf P_e=0\).

\begin{lemma}[Uniform powered protected/gapped input]
\label{lem:V83-powered-PG}
There are \(c>0\) and \(0<\varrho<1\), independent of \(e\), the
representations, and the chosen global state, such that
\begin{align}
 0\preccurlyeq\mathsf G_e^{(t)}
 &\preccurlyeq
 r_t(I-\mathsf P_e),
 &
 r_t&:=(1-cs_\alpha)^t,
\label{eq:V83-powered-gap}\\
 \kappa_\otimes^{\,k|k^c}(\mathsf P_e)
 &\le\varrho.
\label{eq:V83-powered-protected}
\end{align}
The standing weak-link range is \(0<s_\alpha\le1\), after decreasing
\(c\) so that \(cs_\alpha<1\).
\end{lemma}

\begin{proof}
Every block in \(\mathsf G_e^{(t)}\) has at least one nontrivial
output.  The positive Casimir floor and the one-link heat bound give
\(q^t\le(1-cs_\alpha)^t\), which proves
\eqref{eq:V83-powered-gap}.

On \(\mathsf P_e\), the nontrivial input on row \(k\) cannot form a
trivial representation by itself.  Its invariant block therefore
contains at least one other row.  Across the cut
\(\{k\}|k^c\), the normalized invariant projector has largest squared
Schmidt coefficient at most
\((\dim V_{k,e})^{-1}\le1/2\).  All other blocks are contractions.
This is the proof of \Cref{lem:V70-local-PG}, and gives
\eqref{eq:V83-powered-protected} with a fixed
\(\varrho<1\).  If \(\mathsf P_e=0\), the assertion is immediate.
\end{proof}

Put
\[
 \mathsf M_E^{(t)}:=\bigotimes_{e\in E}\mathsf M_e^{(t)}.
\]
If every \(\mathsf P_e\ne0\), set
\(\mathsf P_*:=\bigotimes_e\mathsf P_e\); otherwise set
\(\mathsf P_*=0\).

\begin{lemma}[The V23 one-failure sum lifts to the row cut]
\label{lem:V83-one-failure-lift}
For the functions
\(\Phi_N^{(0)},\Phi_N^{(1)}\) of
\eqref{eq:V23-Phi-protected}--%
\eqref{eq:V23-Phi-failed}, with their fixed
\(\rho\) enlarged if necessary so that \(\rho\ge\varrho\),
\begin{equation}
 \kappa_\otimes^{\,k|k^c}
 \bigl(\mathsf M_E^{(t)}-\mathsf P_*\bigr)
 \le
 \begin{cases}
 \Phi_N^{(0)}(r_t),&\mathsf P_*\ne0,\\
 \Phi_N^{(1)}(r_t),&\mathsf P_*=0.
 \end{cases}
\label{eq:V83-one-failure-lift}
\end{equation}
Consequently the integer and half-integer sums in
\eqref{eq:V23-integer-sum}--%
\eqref{eq:V23-half-sum} hold for this higher-incidence bank.
\end{lemma}

\begin{proof}
Expand the product of the upper bounds
\[
 \mathsf P_e+r_t(I-\mathsf P_e)
\]
and, when present, subtract the all-protected product.  For every
selected protected subset \(A\subseteq E\), tensor
submultiplicativity across the single fixed cut \(k|k^c\) gives
\[
 \kappa_\otimes^{\,k|k^c}
 \left(\bigotimes_{e\in A}\mathsf P_e
       \otimes I_{E\setminus A}\right)
 \le\varrho^{|A|}
 \le\rho^{|A|}.
\]
The binomial sum is therefore exactly the proof of
\Cref{lem:V23-power-capacity}; it gives the corresponding
\(\Phi_N\).  The analytic proof of
\Cref{lem:V23-analytic-sum} depends only on
\(r_t=(1-cs_\alpha)^t\) and \(\rho<1\), so it applies after changing
fixed constants.
\end{proof}

\section{The paid square tail}
\label{sec:V83-paid-tail}

Let \(\mathcal A_E\) be the complete physical moment of the fixed
global state at \(t=1\).  Resolve its output tuple
\(\boldsymbol U\), and write
\[
 Q_{\boldsymbol U}:=\prod q_{\alpha,U},
 \qquad
 Z_{\boldsymbol U}:=\sum C_2(U).
\]
Repeated fusion bounds the component of the final output mass assigned
to this bank by
\begin{equation}
 \Lambda_{\boldsymbol U}
 \le C(N+Z_{\boldsymbol U}).
\label{eq:V83-output-fusion}
\end{equation}
On the complement of \(\mathsf P_*\), let
\(\mathcal P_E(\mathcal A)\) denote the positive physical envelope
with multiplier
\[
 \frac{\Lambda_{\boldsymbol U}Q_{\boldsymbol U}}
 {1-Q_{\boldsymbol U}}.
\]
The all-protected line is unpaid or sends the payer to a different
factor; it is not included in \(\mathcal P_E(\mathcal A)\).

\begin{theorem}[Higher-incidence paid row-rooted square tail]
\label{thm:V83-paid-square-tail}
Uniformly in the coordinate support, representations, and fixed
global partition state,
\begin{align}
 \kappa_\otimes^{\,k|k^c}
 \bigl(\mathcal P_E(\mathcal A)\bigr)
 &\le\frac{C}{s_\alpha},
\label{eq:V83-paid-universal}\\
 \boxed{
 s_\alpha^2S^2
 \kappa_\otimes^{\,k|k^c}
 \bigl(\mathcal P_E(\mathcal A)\bigr)
 \le CH.}
\label{eq:V83-paid-square-tail}
\end{align}
\end{theorem}

\begin{proof}
Define the geometric and output-Casimir resolvers
\begin{align*}
 \mathscr G_E
 &:=
 \sum_{Q_{\boldsymbol U}<1}
 \frac{Q_{\boldsymbol U}}{1-Q_{\boldsymbol U}}
 P_{\boldsymbol U},\\
 \mathscr C_E
 &:=
 \sum_{Q_{\boldsymbol U}<1}
 \frac{Z_{\boldsymbol U}Q_{\boldsymbol U}}
 {1-Q_{\boldsymbol U}}
 P_{\boldsymbol U}.
\end{align*}
By \eqref{eq:V83-output-fusion},
\begin{equation}
 \mathcal P_E(\mathcal A)
 \preccurlyeq C(N\mathscr G_E+\mathscr C_E).
\label{eq:V83-resolver-majorant}
\end{equation}

Geometric summation and
\Cref{lem:V83-one-failure-lift} give
\begin{equation}
 \kappa_\otimes(\mathscr G_E)
 \le
 Ce^{-cs_\alpha N}
 \left(1+\frac1{s_\alpha N}\right)
 \le\frac{C}{s_\alpha N}.
\label{eq:V83-geometric-resolver}
\end{equation}
For the Casimir resolver, the retained-failure compiler with exponent
two gives, exactly as in the proof of
\eqref{eq:V23-Casimir-bounds},
\begin{equation}
 \kappa_\otimes(\mathscr C_E)
 \le
 \min\left\{
 \frac{C}{s_\alpha},
 \frac{C}{s_\alpha}
 e^{-cs_\alpha N}
 \left(1+\frac1{s_\alpha N}\right)
 \right\}.
\label{eq:V83-Casimir-resolver}
\end{equation}
Equations
\eqref{eq:V83-resolver-majorant}--%
\eqref{eq:V83-Casimir-resolver} first give
\eqref{eq:V83-paid-universal}.

For the marked estimate put \(x=s_\alpha N\) and use
\eqref{eq:V83-row-power}.  If \(x\le1\), the last entries of
\eqref{eq:V83-geometric-resolver}--%
\eqref{eq:V83-Casimir-resolver} give
\begin{align*}
 s_\alpha^2S^2N\,
 \kappa_\otimes(\mathscr G_E)
 &\le CxH,\\
 s_\alpha^2S^2\,
 \kappa_\otimes(\mathscr C_E)
 &\le CxH.
\end{align*}
If \(x>1\), use their exponentially decaying entries instead; the
two right sides become
\[
 Cx^2e^{-cx}(1+x^{-1})H,
 \qquad
 Cxe^{-cx}(1+x^{-1})H.
\]
All four scalar factors are bounded.  Insert them in
\eqref{eq:V83-resolver-majorant} to prove
\eqref{eq:V83-paid-square-tail}.
\end{proof}

\begin{firewall}[Why this is not a second proof of V23]
\label{fire:V83-not-duplicate-V23}
The integer/half-integer analytic sums and retained-failure compiler
are imported unchanged from V23.  The new assertion is that V70's
higher-incidence local protected/gapped decomposition satisfies the
same abstract hypotheses across the row cut \(k|k^c\), including for
the endpoint moments of the V77 covariance.  No incidence-two graph
or edgewise singlet factorization is assumed in
\Cref{thm:V83-paid-square-tail}.
\end{firewall}

\section{Both endpoint moments and the selected covariance}
\label{sec:V83-endpoints}

For \(\sigma\in\{+,\rho\}\), let
\(\mathcal A_E^\sigma\) denote respectively the common-variable and
cut-replicated endpoint moment restricted to \(E\).  Each is first
written as the fixed finite sum of global row-partition states and
only then physically resolved.  Let
\[
 \mathcal J_E^\rho
 :=\mathcal A_E^+-\mathcal A_E^\rho.
\]

\begin{theorem}[Higher-incidence selected-bank certificates]
\label{thm:V83-selected-bank}
The common endpoint, replicated endpoint, and their covariance obey
\begin{align}
 \max_{\sigma\in\{+,\rho\}}
 \kappa_\otimes^{\,k|k^c}(\mathcal A_E^\sigma),
 \quad
 \kappa_\otimes^{\,k|k^c}(\mathcal J_E^\rho)
 &\le
 C\frac{H}{s_\alpha S^2},
\label{eq:V83-selected-unpaid}\\
 \max_{\sigma\in\{+,\rho\}}
 \kappa_\otimes^{\,k|k^c}
 \bigl(\mathcal P_E(\mathcal A_E^\sigma)\bigr),
 \quad
 \kappa_\otimes^{\,k|k^c}
 \bigl(\mathcal P_E(\mathcal J_E^\rho)\bigr)
 &\le
 C\frac{H}{s_\alpha^2S^2}.
\label{eq:V83-selected-paid}
\end{align}
\end{theorem}

\begin{proof}
For either endpoint and each fixed global state, the unpaid estimate
is \Cref{thm:V70-balanced-tail}; the proof there already permits
higher-incidence coordinates.  The paid estimate is
\Cref{thm:V83-paid-square-tail}.  The number of global row partitions
is bounded in terms of four rows and is independent of support, so
their positive envelope sum changes only the constant.

The covariance is the difference of the two endpoints.  Capacity
subadditivity proves its unpaid estimate.  For its paid estimate,
resolve common physical blocks and apply the positive
output-mass/resolvent weight after taking the two-term absolute
envelope, exactly as in
\Cref{thm:V82-pure-pair-double}.  The all-protected component is not
assigned to this payer.
\end{proof}

\begin{firewall}[One row cut, not two capacities]
\label{fire:V83-one-row-cut}
All capacity estimates in
\eqref{eq:V83-selected-unpaid}--%
\eqref{eq:V83-selected-paid} use the single cut
\(\{k\}|k^c\).  No estimate across an \(\ell\)-cut is multiplied into
it.  The square \(S^2\) comes from the bank-count power inequality and
the one-failure sum, not from a product of two endpoint capacities.
\end{firewall}

\section{Extraction from the V77 bank}
\label{sec:V83-extraction}

Let \(B\) be the complete residual bank of the fixed cut and let
\(R:=B\setminus E\).  Coordinate factorization gives
\[
 \mathcal A_B^+
 =\mathcal A_E^+\otimes\mathcal A_R^+,
 \qquad
 \mathcal A_B^\rho
 =\mathcal A_E^\rho\otimes\mathcal A_R^\rho.
\]

\begin{lemma}[Higher-incidence selected-subbank identity]
\label{lem:V83-subbank-identity}
The full V77 covariance satisfies
\begin{equation}
 \boxed{
 \mathcal J_B^\rho
 =
 \mathcal J_E^\rho\otimes\mathcal A_R^+
 +
 \mathcal A_E^\rho\otimes\mathcal J_R^\rho.}
\label{eq:V83-subbank-identity}
\end{equation}
\end{lemma}

\begin{proof}
Insert and subtract
\(\mathcal A_E^\rho\otimes\mathcal A_R^+\) in
\[
 \mathcal J_B^\rho
 =
 \mathcal A_E^+\otimes\mathcal A_R^+
 -
 \mathcal A_E^\rho\otimes\mathcal A_R^\rho.
\]
This is the algebra of
\Cref{lem:V82-subbank-extraction}; no incidence assumption enters it.
\end{proof}

\begin{theorem}[The full crossing bank inherits the selected square
tail]
\label{thm:V83-full-bank-square}
Let \(\mathcal N_B^{[E]}\) be the unpaid positive envelope of
\eqref{eq:V83-subbank-identity}.  Let
\(\mathcal P_B^{[E]}\) be its exact V2 one-payer envelope.  Then
\begin{align}
 \boxed{
 \kappa_\otimes^{\,k|k^c}(\mathcal N_B^{[E]})
 \le C\frac{H}{s_\alpha S^2},}
\label{eq:V83-full-unpaid}\\
 \boxed{
 \kappa_\otimes^{\,k|k^c}(\mathcal P_B^{[E]})
 \le C\frac{H}{s_\alpha^2S^2}.}
\label{eq:V83-full-paid}
\end{align}
\end{theorem}

\begin{proof}
For the unpaid envelope, use
\eqref{eq:V83-selected-unpaid} on
\(\mathcal J_E^\rho\) in the first term and on
\(\mathcal A_E^\rho\) in the second.  The remaining common endpoint is
an ordinary contraction and the remaining covariance has ordinary
norm at most two, as a difference of two contractions.  Apply the
one-sided tensor inequality.

For the paid envelope, the unique payer lies in the selected or
remaining factor of one of the two terms.  If it lies in the selected
factor, use \eqref{eq:V83-selected-paid} and contract the remainder.
If it lies in the remainder, use
\eqref{eq:V83-selected-unpaid} and its universal complete-moment or
complete-covariance paid cap \(C/s_\alpha\), from
\eqref{eq:V3-spectator-carrier-cap} or
\eqref{eq:V78-universal-paid-cap}.  These four alternatives are the
exact V2 scalar allocation, so no second denominator or capacity is
introduced.
\end{proof}

\begin{corollary}[A dominant row gives the square denominator]
\label{cor:V83-dominant-square}
If
\begin{equation}
 S\ge\kappa\Xi_k,
\label{eq:V83-row-dominance}
\end{equation}
then
\begin{align}
 \kappa_\otimes(\mathcal N_B^{[E]})
 &\le
 C_\kappa\frac{H_{k\ell}(E)}
 {s_\alpha\Xi_k^2},
\label{eq:V83-dominant-unpaid}\\
 \kappa_\otimes(\mathcal P_B^{[E]})
 &\le
 C_\kappa\frac{H_{k\ell}(E)}
 {s_\alpha^2\Xi_k^2}.
\label{eq:V83-dominant-paid}
\end{align}
If \(k\) is the larger of two unresolved endpoints \(i,j\), then
\[
 \Xi_k^2\ge\Xi_i\Xi_j,
\]
so the same estimates also pay their two endpoint denominators.
\end{corollary}

\begin{proof}
Insert \eqref{eq:V83-row-dominance} in
\eqref{eq:V83-full-unpaid}--%
\eqref{eq:V83-full-paid}.  The last statement is the definition of
the larger endpoint.
\end{proof}

\section{The tree-exchange table}
\label{sec:V83-tree-exchange}

Write the nominal \(2+2\) path as
\[
 a-k-m-b,
\]
where \(ak\) and \(mb\) are the two same-side edges and \(km\) is the
nominal crossing edge.  Thus \(k\) is one unresolved endpoint and
\(m\) the other.  Suppose \(k\) is chosen so that
\(\Xi_k\ge\Xi_m\), and let \(\ell\) be the maximal partner of its
dominant bank \(E_{k\ell}\).

\begin{proposition}[Complete partner table at the larger endpoint]
\label{prop:V83-partner-table}
There are exactly three partner positions:
\begin{center}
\begin{tabular}{c c c}
\toprule
partner \(\ell\) & graph
\(\{ak,mb,k\ell\}\) & outcome\\
\midrule
\(a\) & \(ak\) repeated & nonspanning parallel state\\
\(m\) & original path & closes automatically\\
\(b\) & exchanged path \(a-k-b-m\) &
closes if \(\Xi_b\le L\Xi_k\)\\
\bottomrule
\end{tabular}
\end{center}
In the last row, failure of the displayed comparison is the strict
mass ascent \(k\to b\), \(\Xi_b>L\Xi_k\).
\end{proposition}

\begin{proof}
The other three vertices are \(a,m,b\).  If \(\ell=a\), the certified
edge is parallel to the same-side edge and the graph has two
components.  If \(\ell=m\), the original three-edge path is recovered.
Moreover \(\Xi_m\le\Xi_k\), so
\[
 \frac{H_{km}}{\Xi_k^2}
 \le
 \frac{H_{km}}{\Xi_k\Xi_m}.
\]
If \(\ell=b\), the edges \(ak,kb,bm\) form the displayed path.  Under
\(\Xi_b\le L\Xi_k\),
\[
 \frac{H_{kb}}{\Xi_k^2}
 \le
 L\frac{H_{kb}}{\Xi_k\Xi_b}.
\]
If that mass comparison fails, the direction \(k\to b\) is strictly
ascending.  The list is exhaustive.
\end{proof}

\begin{theorem}[Mass-compatible higher-incidence tree exchange]
\label{thm:V83-tree-exchange}
Assume:
\begin{enumerate}[label=\textup{(\roman*)}]
\item \(E_{k\ell}\) is a clean selected bank disjoint from the two
same-side response stocks and from all marked payer banks;
\item \(S\ge\kappa\Xi_k\);
\item \(k\) is a larger unresolved endpoint;
\item \(\ell\) lies in the opposite block of \(\rho\); and
\item \(\Xi_\ell\le L\Xi_k\).
\end{enumerate}
Then every payer location of the conditioned \(2+2\) packet satisfies
\begin{equation}
 \boxed{
 \frac{\mathfrak b_{k\ell}^{\rm exch}}
 {\Xi_1\Xi_2\Xi_3\Xi_4}
 \le
 C_{\kappa,L}\,
 \theta_{12}\theta_{34}\theta_{k\ell}.}
\label{eq:V83-tree-exchange}
\end{equation}
\end{theorem}

\begin{proof}
Use the raw same-side response entries
\eqref{eq:V82-C-raw}--%
\eqref{eq:V82-D-raw}.  If one same-side block pays, its paid response
has no \(s_\alpha\), the other contributes one
\(s_\alpha\), and
\eqref{eq:V83-dominant-unpaid} contributes
\(s_\alpha^{-1}\).  If the crossing bank pays, the two same-side
responses contribute \(s_\alpha^2\), cancelled by
\eqref{eq:V83-dominant-paid}.  Every case is bounded by
\[
 C_\kappa
 \frac{H_{12}H_{34}H_{k\ell}}{\Xi_k^2}.
\]
The mass comparison gives
\[
 \frac1{\Xi_k^2}
 \le\frac{L}{\Xi_k\Xi_\ell}.
\]
After division by the four row masses, the three genuine edges
\(\{12,34,k\ell\}\) give the legal path mark in
\eqref{eq:V83-tree-exchange}.  The payer was allocated before the
one-sided tensor estimate by
\Cref{thm:V83-full-bank-square}.
\end{proof}

\begin{corollary}[Finite maximal-partner selection]
\label{cor:V83-partner-pigeonhole}
Let \(E_k^{\rm bal}\) be a clean balanced crossing bank at row \(k\)
with
\[
 \sum_{e\in E_k^{\rm bal}}a_{k,e}\ge\kappa\Xi_k.
\]
Partition it by the V70 maximal-partner rule.  Some
\(E_{k\ell}\) satisfies
\[
 \sum_{e\in E_{k\ell}}a_{k,e}
 \ge\frac{\kappa}{3}\Xi_k.
\]
If one such dominant cell has an opposite-block, mass-compatible
partner, \Cref{thm:V83-tree-exchange} closes that allocated cell.
Every unclosed allocated cell has either a same-side maximal partner
or a strict opposite-block mass ascent.
\end{corollary}

\begin{proof}
There are only three possible partners and their cells partition
\(E_k^{\rm bal}\).  Apply the finite pigeonhole principle and then
\Cref{prop:V83-partner-table,thm:V83-tree-exchange}.
\end{proof}

\begin{firewall}[The nominal card is not separately bounded]
\label{fire:V83-no-nominal-card}
The edge \(k\ell\) is selected from the actual V70 physical bank
before the full crossing covariance is enveloped.  V83 does not
assert that the nominal \(km\) whole-row telescope summand is bounded
by \(H_{k\ell}\), or even by \(H_{km}\).  Once
\eqref{eq:V83-full-paid} or
\eqref{eq:V83-full-unpaid} has been proved for the assembled operator,
the resulting nonnegative numerator is compared with the legal
\(k\ell\) tree.
\end{firewall}

\section{Regression and termination tests}
\label{sec:V83-tests}

\begin{proposition}[Two-large/two-small four-row test]
\label{prop:V83-four-row-test}
At one residual coordinate for the cut \(12|34\), assign masses
\[
 a_1=a_2=L,\qquad a_3=a_4=b,\qquad L\gg b\ge a_*.
\]
This coordinate is balanced in the V69 sense.  Rooting at row \(1\)
selects row \(2\) as maximal partner.  The square-tail estimate is
valid with \(S=L\), \(H=L^2\), but its certified edge is \(12\), so
the tree exchange fails as a parallel state.

If instead
\[
 a_1=a_3=L,\qquad a_2=a_4=b,
\]
rooting at row \(1\) selects the opposite-block partner \(3\).
The masses are compatible, and the certified edge \(13\) gives a
legal path.
\end{proposition}

\begin{proof}
With two inputs of size \(L\), no single largest spin dominates the
sum of the other active spins by the V69 factor two, so both patterns
are balanced.  The maximal-partner and stock statements follow
directly from \eqref{eq:V83-SHN}.  Apply
\Cref{prop:V83-partner-table}: the first pattern is its parallel row,
and the second its opposite-block mass-compatible row.  This test
separates validity of the analytic square tail from legality of its
graph numerator.
\end{proof}

\begin{proposition}[Rademacher wrong-card routing test]
\label{prop:V83-rademacher-routing}
In the V78 Rademacher pattern with physical coordinate banks on
\(14\) and \(23\), the nonzero nominal \(13\) whole-row summand is
never bounded by \(H_{13}=0\).  With equal row masses, a selection
rooted at row \(1\) uses its actual opposite partner \(4\), and the
edge set
\[
 \{12,34,14\}
\]
is a mass-compatible path.  The symmetric selection uses the actual
edge \(23\).
\end{proposition}

\begin{proof}
The false \(13\) localization is
\Cref{prop:V78-no-pair-operator}.  The partner table for \(k=1\) has
\(a=2,m=3,b=4\), so partner \(4\) is the exchanged-path row.
Equal masses satisfy the comparison with \(L=1\).  This proposition
tests only the exact allocation and graph comparison: the Rademacher
variables are not substituted for Wilson heat multipliers in
\Cref{thm:V83-paid-square-tail}.  Replacing the two physical banks by
equal-mass balanced Wilson banks supplies the analytic hypotheses.
\end{proof}

\begin{lemma}[Strict mass ascent cannot cycle]
\label{lem:V83-ascent-terminates}
Fix \(L>1\).  Any directed chain of distinct transfer states satisfying
\[
 \Xi_{v_{r+1}}>L\Xi_{v_r}
\label{eq:V83-strict-ascent}
\]
at every step has length at most three on four rows.  In particular,
no directed cycle can consist entirely of failed mass-compatible
opposite-block exchanges.
\end{lemma}

\begin{proof}
Along such a chain the positive row masses strictly increase.
Therefore no vertex can recur.  A simple directed path on four
vertices has at most three edges.
\end{proof}

\begin{firewall}[Termination is not acquisition]
\label{fire:V83-termination-not-acquisition}
\Cref{lem:V83-ascent-terminates} says that already available strict
transfers cannot cycle.  It does not prove that the row reached after
an ascent possesses a new clean dominant bank.  That acquisition must
come from its exact V79 incidence identity, with previously selected
coordinates removed.
\end{firewall}

\section{Post-V83 reduction}
\label{sec:V83-reduction}

\begin{theorem}[Residual higher-incidence crossing states]
\label{thm:V83-residual}
Within the clean balanced two-cross/no-prefix sector, apply V82 to
common dominant incidence-two banks and V83 to every dominant
row-rooted bank with an opposite-block mass-compatible maximal
partner.  Every remaining higher-incidence selection has at least one
of:
\begin{enumerate}[label=\textup{(\roman*)}]
\item a dominant cell whose maximal partner lies in the same block,
so its square-tail numerator is parallel to a same-side tree edge;
\item a dominant cell whose opposite-block maximal partner satisfies
the strict mass ascent \eqref{eq:V83-strict-ascent};
\item no clean dominant cell after removing previously selected and
marked coordinates;
\item an output-ancestry payer attached to one of those overlapping
cells; or
\item a local representation pattern outside the \(SU(2)\)
protected/gapped calculus used in V70 and V83.
\end{enumerate}
\end{theorem}

\begin{proof}
The maximal-partner cells partition the balanced crossing incidences
at the chosen row.  A fixed-fraction clean cell is handled by
\Cref{cor:V83-partner-pigeonhole}.  If its partner is across the cut,
\Cref{prop:V83-partner-table} leaves only the mass-compatible and
strictly ascending alternatives; the former closes by
\Cref{thm:V83-tree-exchange}.  A same-side partner is the parallel
row of that table.  Failure to select a disjoint fixed-fraction cell
is item \textup{(iii)}, and the existing sole-payer ancestry
allocation gives item \textup{(iv)}.  The local powered
protected/gapped lemma was proved using the \(SU(2)\) invariant-block
input, so no wider group statement is included.
\end{proof}

\begin{assessment}[What V83 genuinely removes]
\label{ass:V83-progress}
Higher incidence does not destroy the paid two-endpoint mechanism.
The V70 row-cut bank has the same one-failure resolver as the V23
incidence-two graph, and its exact extraction from V77 preserves both
square powers under every payer location.  The clean
mass-compatible opposite-partner branch is therefore closed.

The obstruction has moved to a sharper place: when the maximal local
partner is same-side, the available square-tail numerator is not a
crossing edge; when it is an excessively heavy opposite partner, the
tree comparison leaves a mass ratio.  Neither defect is analytic
support growth, and neither may be erased by changing the nominal
row-pair card.
\end{assessment}

\begin{openproblem}[V84: cross-retained parallel repair and ascending
acquisition]
\label{open:V83-V84-mandate}
\begin{enumerate}
\item For a same-side-maximal bank \(E_{kk^\circ}\) whose coordinates
      nevertheless cross \(\rho\), return to the exact double-replica
      covariance and retain one opposite-block difference inside the
      V83 one-failure resolver.
\item Test the genuinely crossing square-tail targets
      \[
      \kappa_\otimes(\mathcal N_E^\rho)
      \le
      C\frac{H_{k,C_1}(E)}
      {s_\alpha S_k(E)^2},
      \qquad
      \kappa_\otimes(\mathcal P_E^\rho)
      \le
      C\frac{H_{k,C_1}(E)}
      {s_\alpha^2S_k(E)^2}
      \]
      when \(k\in C_0\), with
      \(H_{k,C_1}=\sum_{\ell\in C_1}H_{k\ell}\).
      Do not infer them by replacing \(H_{kk^\circ}\) in V83.
\item Run the two-large-same-side/two-small-opposite test of
      \Cref{prop:V83-four-row-test}; if the targets fail, identify the
      weakest additional cross-output moment which is true.
\item For every strict opposite-block ascent, apply the exact V79
      incidence identity at the new row after deleting used banks.
      Either acquire a fresh paired or row-rooted certificate, or
      record the precise depleted/overlap cell.
\item Use \Cref{lem:V83-ascent-terminates} only after acquisition has
      been proved at every step.
\item Keep the finite tree edge produced by the final certificate;
      do not restore the original nominal card.
\end{enumerate}
\end{openproblem}

\section{V83 conclusion and claim boundary}
\label{sec:V83-conclusion}

\begin{conclusion}[V83 conclusion]
\label{con:V83}
The higher-incidence paid square tail and its exact V77 extraction are
now proved.  They close every clean \(2+2\) two-cross/no-prefix packet
whose larger unresolved endpoint has a dominant balanced
maximal-partner bank with an opposite-block partner of comparable or
smaller row mass.  The proof is uniform in bank size, representations,
and payer location within the \(SU(2)\) calculus.

The next bottleneck is the same-side-maximal parallel state, together
with acquisition along the finite strict-ascent chain.  V84 must keep
an actual cross-cut difference inside the one-failure resolver rather
than substituting the same-side square-tail stock.
\end{conclusion}

\begin{firewall}[No mass-gap claim]
\label{fire:V83-no-mass-gap}
V83 remains a finite-cutoff sourcewise reduction.  It does not close
the parallel/ascending/overlap states, the other first-connectivity
families, the extension beyond the stated local representation
calculus, continuum construction, reflection-positive
reconstruction, or a uniform positive spectral gap.  The full
Yang--Mills mass gap remains open.
\end{firewall}

\begin{record}[Parent and continuation record]
\label{rec:V83-parent}
This file was copied byte-for-byte from
\texttt{Renewal\_Yang\_Mills\_Mass\_Gap\_Research\_Notes\_V82.tex}
before the present chapter was appended.  The parent had \(36878\)
lines, \(1273518\) bytes, and SHA--256 digest
\[
\texttt{69729185d276766616f80a8fbb336d49bccf7c2c1b30bdef1f3e2d124bedc2a6}.
\]
The next active version is V84 and begins with
\Cref{open:V83-V84-mandate}.  The next compulsory read-only audit of
the current SM and NS notes remains V85.
\end{record}

% =============================================================================
% BEGIN APPENDED COMPACT VERSION V84 MATERIAL
% =============================================================================

\clearpage
\chapter{The same-side-maximal no-go and the intermediate-output
currency}
\label{chap:V84}

\section{Context: test the cross-retained square tail exactly}
\label{sec:V84-context}

V83 proved a higher-incidence square tail with numerator
\(H_{k\ell}\), where \(\ell\) is the maximal local partner of the
root row \(k\).  If \(\ell\) lies across the \(2+2\) cut and has
compatible row mass, that numerator is a legal replacement crossing
edge.  The unresolved parallel state occurs when \(\ell=k^\circ\)
lies on the same side even though every coordinate of the bank also
sees the opposite block.

The V84 mandate proposed retaining an opposite-block difference
inside the one-failure resolver and replacing the same-side numerator
by the crossing stock \(H_{k,C_1}\).  This note first tests that
proposal on one exact \(SU(2)\) heat-kernel coordinate.  Both proposed
square-tail inequalities are false.  The obstruction is a low
nontrivial intermediate output of the large same-side pair: its
product-state capacity is inverse representation dimension, not
inverse Casimir.

This is not a failure of the V83 square tail, whose numerator remains
the large same-side overlap.  It shows that the parallel state must be
resolved one recoupling level earlier, retaining the intermediate
same-side representation before converting to input overlap stocks.
The second part of the note constructs that exact output-resolved
currency and proves the finite fresh-cell acquisition alternative for
strict mass ascents.

\section{A three-row one-coordinate heat covariance}
\label{sec:V84-model}

Use the cut
\[
 C_0=\{1,2\},\qquad C_1=\{3,4\}.
\]
At one residual coordinate let rows \(1,2\) carry \(V_j,V_j\), let
row \(3\) carry \(V_1\), and let row \(4\) be inactive.  Put
\[
 q_J:=e^{-sJ(J+1)},\qquad 0<s\le s_0.
\]
On \(V_j\otimes V_j\), let \(P_J^{12}\) be the spin-\(J\)
projection.  On
\((V_j\otimes V_j)\otimes V_1\), let \(P_{J,K}\) be the joint
projection on pair spin \(J\) and total spin \(K\).  The two Casimirs
commute, so these are orthogonal physical blocks.

The common-variable heat moment and the cut-replicated endpoint are
\begin{align}
 \mathsf M_s^+
 &:=
 \sum_{\substack{0\le J\le2j\\K\subset J\otimes1}}
 q_KP_{J,K},
\label{eq:V84-common-heat}\\
 \mathsf M_s^\rho
 &:=
 \sum_{\substack{0\le J\le2j\\K\subset J\otimes1}}
 q_Jq_1P_{J,K}.
\label{eq:V84-replicated-heat}
\end{align}
Thus the exact one-coordinate cut covariance is
\begin{equation}
 \boxed{
 \mathsf J_{s,j}
 :=
 \mathsf M_s^+-\mathsf M_s^\rho
 =
 \sum_{J,K}(q_K-q_Jq_1)P_{J,K}.}
\label{eq:V84-exact-covariance}
\end{equation}
Its unpaid positive envelope is
\begin{equation}
 \mathsf N_{s,j}:=
 |\mathsf J_{s,j}|
 =
 \sum_{J,K}|q_K-q_Jq_1|P_{J,K}.
\label{eq:V84-unpaid-envelope}
\end{equation}

\begin{proof}
The Fourier transform of the heat kernel on a total spin-\(K\)
irreducible is \(q_KI\), which gives
\eqref{eq:V84-common-heat}.  Replicating the two sides of the cut
makes the heat expectations independent.  The pair-\(J\) multiplier
is \(q_J\), the row-\(3\) multiplier is \(q_1\), and recoupling their
tensor product to total \(K\) does not change the scalar product
\(q_Jq_1\).  Subtraction gives
\eqref{eq:V84-exact-covariance}, which is the local instance of
\eqref{eq:V65-double-replica}.
\end{proof}

Let
\[
 \psi_j:=|j,j\rangle\otimes|j,-j\rangle.
\]
Its pair magnetic number is zero.  Write
\[
 p_{j,J}:=
 \langle\psi_j,P_J^{12}\psi_j\rangle.
\]

\begin{lemma}[Exact spin-one weight of the anti-aligned pair]
\label{lem:V84-spin-one-weight}
For every \(j\ge\frac12\),
\begin{equation}
 \boxed{
 p_{j,1}
 =
 3\frac{[(2j)!]^2}{(2j-1)!(2j+2)!}
 =
 \frac{6j}{(2j+1)(2j+2)}.}
\label{eq:V84-spin-one-weight}
\end{equation}
In particular, there are constants \(c,C>0\) such that
\[
 \frac{c}{j}\le p_{j,1}\le\frac{C}{j}
\qquad(j\ge1).
\]
\end{lemma}

\begin{proof}
The standard extremal Clebsch--Gordan coefficient is
\[
 \left|
 \langle j,j;j,-j\mid J,0\rangle
 \right|^2
 =
 (2J+1)
 \frac{[(2j)!]^2}
 {(2j-J)!(2j+J+1)!}.
\]
Set \(J=1\) and cancel factorials.  The two-sided estimate follows
immediately.
\end{proof}

\begin{proposition}[Unpaid lower bound on the \(J=1,K=0\) block]
\label{prop:V84-unpaid-lower}
For the product vector
\[
 \Psi_j^{(0)}:=\psi_j\otimes|1,0\rangle
\]
one has
\begin{equation}
 \left\langle\Psi_j^{(0)},
 P_{1,0}\mathsf N_{s,j}P_{1,0}
 \Psi_j^{(0)}\right\rangle
 =
 \frac{p_{j,1}}3(1-q_1^2)
 \ge c\frac{s}{j}.
\label{eq:V84-unpaid-lower}
\end{equation}
\end{proposition}

\begin{proof}
Only pair magnetic number \(M=0\) occurs in the spin-one component of
\(\psi_j\).  The coefficient
\[
 |\langle1,0;1,0\mid0,0\rangle|^2=\frac13.
\]
On the block \(J=1,K=0\),
\eqref{eq:V84-exact-covariance} has coefficient
\[
 q_0-q_1^2=1-e^{-4s}.
\]
This is at least \(cs\) for \(0<s\le s_0\).  Combine with
\Cref{lem:V84-spin-one-weight}.
\end{proof}

\begin{proposition}[Paid lower bound on the \(J=1,K=1\) block]
\label{prop:V84-paid-lower}
Let \(\mathsf P_{s,j}\) be the positive physical envelope of
\(\mathsf J_{s,j}\) after the sole payer is assigned to its nontrivial
total output.  For
\[
 \Psi_j^{(1)}:=\psi_j\otimes|1,1\rangle
\]
one has
\begin{equation}
 \left\langle\Psi_j^{(1)},
 P_{1,1}\mathsf P_{s,j}P_{1,1}
 \Psi_j^{(1)}\right\rangle
 =
 \frac{3q_1p_{j,1}}2
 \ge\frac{c}{j}.
\label{eq:V84-paid-lower}
\end{equation}
\end{proposition}

\begin{proof}
The relevant coupling coefficient is
\[
 |\langle1,0;1,1\mid1,1\rangle|^2=\frac12.
\]
On \(J=1,K=1\), the signed covariance coefficient is
\[
 q_1-q_1^2=q_1(1-q_1)>0.
\]
The final output mass is
\(A_1=1+C_2(1)=3\), and the payer multiplier is
\(A_1/(1-q_1)\).  Their product with the covariance coefficient is
\(3q_1\).  Multiply by the pair weight and the displayed coupling
coefficient.  Since \(q_1=e^{-2s}\) is bounded below for
\(s\le s_0\), the last inequality follows from
\Cref{lem:V84-spin-one-weight}.
\end{proof}

\section{The proposed crossing square tails are false}
\label{sec:V84-no-go}

Put
\[
 a_j:=1+j(j+1).
\]
For the one-coordinate bank above, rooted at row \(1\),
\begin{equation}
 S_1=a_j,\qquad
 H_{12}=a_j^2,\qquad
 H_{1,C_1}=a_ja_3=3a_j.
\label{eq:V84-input-stocks}
\end{equation}
Row \(2\) is the same-side maximal partner.  The coordinate is
balanced because the two largest spins are equal.

\begin{theorem}[No uniform cross-stock square tail in the parallel
state]
\label{thm:V84-cross-square-false}
There is no constant \(C\), uniform in \(0<s\le s_0\) and \(j\), for
which either of the proposed bounds
\begin{align}
 \kappa_\otimes(\mathsf N_{s,j})
 &\le
 C\frac{H_{1,C_1}}{sS_1^2},
\label{eq:V84-false-unpaid}\\
 \kappa_\otimes(\mathsf P_{s,j})
 &\le
 C\frac{H_{1,C_1}}{s^2S_1^2}
\label{eq:V84-false-paid}
\end{align}
holds.  Hence the two targets in
\Cref{open:V83-V84-mandate} are false even for one clean
incidence-three \(SU(2)\) coordinate.
\end{theorem}

\begin{proof}
Choose an integer \(j_s\) with
\[
 s^{-3}\le j_s\le2s^{-3}
\]
for all sufficiently small \(s\).  By
\Cref{prop:V84-unpaid-lower},
\[
 \kappa_\otimes(\mathsf N_{s,j_s})
 \ge c\frac{s}{j_s}\ge c's^4.
\]
Indeed, \(\mathsf N_{s,j_s}\) is positive and dominates its
\((J,K)=(1,0)\) compression, while \(\Psi_{j_s}^{(0)}\) is a product
vector across the input rows.
On the other hand,
\eqref{eq:V84-input-stocks} makes the right side of
\eqref{eq:V84-false-unpaid} at most
\[
 \frac{C_1}{sa_{j_s}}\le C_2s^5.
\]
Their ratio tends to infinity.

Likewise \Cref{prop:V84-paid-lower} gives
\[
 \kappa_\otimes(\mathsf P_{s,j_s})
 \ge\frac{c}{j_s}\ge c's^3,
\]
whereas the proposed paid right side is at most
\[
 \frac{C_1}{s^2a_{j_s}}\le C_2s^4.
\]
This ratio also tends to infinity.  The paid envelope is positive and
dominates its \((J,K)=(1,1)\) compression; both test vectors are
products across the input rows.  Thus the lower bounds apply to the
positive product-state capacities used in the programme.
\end{proof}

\begin{firewall}[What the counterexample does and does not refute]
\label{fire:V84-counterexample-scope}
The theorem refutes only the replacement of the V83 same-side
numerator \(H_{12}\) by the much smaller crossing stock
\(H_{1,C_1}\) while retaining the full square \(S_1^2\).  It does not
refute the valid V65 crossing response, the V77 assembled covariance,
or the V83 estimate with numerator \(H_{12}\).  It is a counterexample
to a proposed intermediate inequality, not to Yang--Mills theory.
\end{firewall}

\section{Exact protected cancellation and the missing currency}
\label{sec:V84-output-currency}

\begin{lemma}[The same-side singlet cancels exactly]
\label{lem:V84-singlet-cancels}
Let \(P_0^{12}\) be the invariant projection in
\(V_j\otimes V_j\).  For every opposite-block Wilson product \(B(U)\),
\begin{equation}
 (P_0^{12}\otimes I)
 \operatorname{Cov}
 \bigl(R_j(U)\otimes R_j(U),B(U)\bigr)
 (P_0^{12}\otimes I)
 =0.
\label{eq:V84-singlet-cancels}
\end{equation}
\end{lemma}

\begin{proof}
The diagonal action is trivial on the invariant line:
\[
 P_0^{12}
 (R_j(U)\otimes R_j(U))
 P_0^{12}=P_0^{12}
\]
for every \(U\).  Therefore the common expectation compressed to this
line is
\[
 P_0^{12}\otimes\mathbb EB,
\]
which is exactly the compressed product of the two separate
expectations.  Their difference vanishes.
\end{proof}

\begin{corollary}[General same-side protected cancellation]
\label{cor:V84-general-protected-cancel}
Let \(A(U)\) be any complete same-side representation block and let
\(P_{\mathrm{inv}}\) be the projection onto its invariant subspace,
with arbitrary finite multiplicity.  For every opposite-block Wilson
product \(B(U)\),
\begin{equation}
 (P_{\mathrm{inv}}\otimes I)
 \operatorname{Cov}\bigl(A(U),B(U)\bigr)
 (P_{\mathrm{inv}}\otimes I)=0.
\label{eq:V84-general-protected-cancel}
\end{equation}
\end{corollary}

\begin{proof}
The group acts as the identity on the entire invariant subspace, so
\[
 P_{\mathrm{inv}}A(U)P_{\mathrm{inv}}=P_{\mathrm{inv}}
\]
for every \(U\).  Compressing the common expectation therefore gives
\(P_{\mathrm{inv}}\otimes\mathbb E B\), exactly the compression of the
product of the two separate expectations.  Their difference is zero.
The argument acts as the identity on the multiplicity space and hence
does not require multiplicity one.
\end{proof}

\begin{lemma}[Intermediate-spin dimension capacity]
\label{lem:V84-intermediate-dimension}
Let \(P_J^{12}\) be the multiplicity-one spin-\(J\) projection in
\(V_j\otimes V_j\).  Then
\begin{equation}
 \boxed{
 \kappa_\otimes^{\,1|2}(P_J^{12})
 \le
 \min\left\{1,\frac{2J+1}{2j+1}\right\}.}
\label{eq:V84-intermediate-dimension}
\end{equation}
If \(P_{J,K}\) is any further recoupling projection after tensoring
with other rows, the same upper bound holds across
\(1\,|\,\{2,3,4\}\).
\end{lemma}

\begin{proof}
The partial trace over row \(2\) commutes with the irreducible
spin-\(j\) action.  By Schur's lemma and
\(\operatorname{Tr}P_J^{12}=2J+1\),
\[
 \operatorname{Tr}_2P_J^{12}
 =
 \frac{2J+1}{2j+1}I_{V_j}.
\]
For a unit vector \(y\in V_j\), extend \(y\) to an orthonormal basis.
Positivity gives
\[
 (I\otimes\langle y|)P_J^{12}(I\otimes|y\rangle)
 \preccurlyeq
 \operatorname{Tr}_2P_J^{12}.
\]
Taking an expectation on a unit \(x\in V_j\) proves the first bound,
with the trivial cap one added when the ratio exceeds one.

For a vector on rows \(2,3,4\), take its reduced density matrix on
row \(2\) and use convexity of the preceding estimate.
Since
\(P_{J,K}\preccurlyeq P_J^{12}\otimes I_{34}\), the refined
projection has no larger capacity.
\end{proof}

\begin{definition}[Cross-output covariance currencies]
\label{def:V84-output-currencies}
For one same-side pair output \(J\), opposite-block output \(R\), and
final recoupled output \(K\subset J\otimes R\), define the complete
local unpaid and paid positive carriers
\begin{align}
 \mathscr C_e^{\rm out,N}
 &:=
 \sum_{\substack{J>0,R,K}}
 |q_K-q_Jq_R|\,P_{J,R,K},
\label{eq:V84-output-currency-unpaid}\\
 \mathscr C_e^{\rm out,P}
 &:=
 \sum_{\substack{J>0,R,K\\K\ne0}}
 \frac{A_K}{1-q_K}
 |q_K-q_Jq_R|\,P_{J,R,K}.
\label{eq:V84-output-currency-paid}
\end{align}
Here the sum is taken only after the complete local cut covariance is
formed; multiplicity spaces remain inside \(P_{J,R,K}\).
\end{definition}

\begin{proposition}[The output-resolved capacity bound]
\label{prop:V84-output-capacity}
For a two-row same-side block \(V_j\otimes V_j\),
\begin{align}
 \kappa_\otimes(\mathscr C_e^{\rm out,N})
 &\le
 \sum_{\substack{J>0,R,K}}
 \min\left\{1,\frac{2J+1}{2j+1}\right\}
 |q_K-q_Jq_R|,
\label{eq:V84-output-capacity-unpaid}\\
 \kappa_\otimes(\mathscr C_e^{\rm out,P})
 &\le
 \sum_{\substack{J>0,R,K\\K\ne0}}
 \min\left\{1,\frac{2J+1}{2j+1}\right\}
 \frac{A_K|q_K-q_Jq_R|}{1-q_K}.
\label{eq:V84-output-capacity-paid}
\end{align}
The \(J=0\) line is absent by
\Cref{lem:V84-singlet-cancels}.  In the counterexample, the \(J=1\)
summand has order \(s/j\) before payment and \(1/j\) after payment, so
the dimension factor has the sharp order required by
\eqref{eq:V84-unpaid-lower}--%
\eqref{eq:V84-paid-lower}.
\end{proposition}

\begin{proof}
The projections in
\eqref{eq:V84-output-currency-unpaid}--%
\eqref{eq:V84-output-currency-paid} are positive and orthogonal.
Apply capacity subadditivity and
\Cref{lem:V84-intermediate-dimension} to each physical block.  Exact
singlet cancellation removes \(J=0\).  For \(J=1\), the ratio
\((2J+1)/(2j+1)\) is comparable to \(1/j\); the coefficient
\(1-q_1^2\) is comparable to \(s\), while payment on \(K=1\) turns
\(q_1(1-q_1)\) into \(A_1q_1\), a fixed positive quantity.
\end{proof}

\begin{assessment}[The correct upstream object]
\label{ass:V84-correct-object}
Input stocks alone do not distinguish the low intermediate output
which defeats
\eqref{eq:V84-false-unpaid}--%
\eqref{eq:V84-false-paid}.  The smallest exact object exposed by the
counterexample is the same-side intermediate representation \(J\),
together with its dimension ratio and the heat difference
\(q_K-q_Jq_R\).  These data must be summed at the complete-bank level
before they are compared with legal crossing trees.  Replacing them
immediately by \(H_{k,C_1}\) loses one factor of order \(sj\) on the
test sequence.
\end{assessment}

\section{Fresh-cell acquisition after a strict ascent}
\label{sec:V84-acquisition}

Let \(\mathcal U\) be the union of coordinate banks already used by a
finite transfer path, and let
\[
 X_r(\mathcal U):=
 \sum_{e\in\mathcal U}a_{r,e}.
\]
Restrict each of the six V79 incidence cells
\[
 P_r,\ U_r,\ A_r,\ F_r,\ D_r,\ I_r
\]
to coordinates outside \(\mathcal U\), and denote the restricted
stocks by a superscript \({\rm fr}\).

\begin{lemma}[Depleted incidence identity]
\label{lem:V84-depleted-identity}
For every row \(r\),
\begin{equation}
 \boxed{
 \Xi_r
 =
 X_r(\mathcal U)
 +P_r^{\rm fr}+U_r^{\rm fr}
 +A_r^{\rm fr}+F_r^{\rm fr}
 +D_r^{\rm fr}+I_r^{\rm fr}.}
\label{eq:V84-depleted-identity}
\end{equation}
\end{lemma}

\begin{proof}
The six V79 cells partition every nonzero row-coordinate incidence by
\Cref{lem:V79-incidence-identity}.  Intersect that partition with
\(\mathcal U\) and its complement.  The weights in the first part sum
to \(X_r(\mathcal U)\); those in the complement give the six displayed
restricted stocks.
\end{proof}

\begin{theorem}[Fresh-cell or overlap acquisition]
\label{thm:V84-fresh-cell}
Fix \(0<\delta<1\).  At every row \(r\), at least one of the following
alternatives is available:
\begin{enumerate}[label=\textup{(\roman*)}]
\item the already used banks carry
      \[
      X_r(\mathcal U)\ge\delta\Xi_r;
      \tag{overlap}
      \]
\item some fresh V79 incidence cell
      \(Z_r^{\rm fr}\in
      \{P_r^{\rm fr},U_r^{\rm fr},A_r^{\rm fr},
        F_r^{\rm fr},D_r^{\rm fr},I_r^{\rm fr}\}\)
      satisfies
      \begin{equation}
      \boxed{
      Z_r^{\rm fr}
      \ge\frac{1-\delta}{6}\Xi_r.}
      \label{eq:V84-fresh-cell}
      \end{equation}
\end{enumerate}
In alternative \textup{(ii)}, the selected cell is coordinate-disjoint
from every bank in \(\mathcal U\).
\end{theorem}

\begin{proof}
If \textup{(i)} fails, subtract it from
\eqref{eq:V84-depleted-identity}.  The sum of the six fresh cells is
larger than \((1-\delta)\Xi_r\), so one is at least one sixth of that
quantity.  Disjointness holds by the definition of the fresh
restriction.
\end{proof}

\begin{corollary}[Noncircular acquisition along a mass ascent]
\label{cor:V84-ascent-acquisition}
Suppose a failed V83 tree exchange produces a strict ascent
\(k\to r\).  Apply
\Cref{thm:V84-fresh-cell} at \(r\).
\begin{enumerate}[label=\textup{(\alph*)}]
\item Alternative \textup{(i)} is an explicit same-bank/used-bank
overlap state; it is not counted as a new transfer.
\item A fresh \(P\)- or \(U\)-cell enters the established
private/unbalanced moment compiler.
\item A fresh \(A\)-, \(F\)-, or \(I\)-cell enters the V81 paired
block interface.
\item A fresh \(D\)-cell may be partitioned by actual crossing partner
and then by the V70 maximal-partner rule, after which V82--V83 apply
whenever their graph and mass gates hold.
\end{enumerate}
If every step produces a fresh crossing cell and another strict mass
ascent, the process terminates after at most three transfers.
\end{corollary}

\begin{proof}
The routing statements are precisely the definitions of the six
disjoint V79 cells and the certificates already proved for them.
Freshness permits one-sided tensor insertion without reusing an
earlier bank.  If the new crossing certificate again fails only by a
strict ascent, apply
\Cref{lem:V83-ascent-terminates}.  That lemma is now used after
acquisition, as required: every continued arrow is backed by a newly
selected disjoint cell.
\end{proof}

\begin{firewall}[The overlap alternative remains analytic]
\label{fire:V84-overlap-open}
When \(X_r(\mathcal U)\ge\delta\Xi_r\), the row mass needed at \(r\)
is carried by banks already used earlier in the path.  Reapplying
their first-moment or capacity estimates would double charge them.
This branch requires a higher moment of the same complete physical
bank or a joint multirow estimate; the finite acquisition theorem does
not close it.
\end{firewall}

\section{Post-V84 reduction}
\label{sec:V84-reduction}

\begin{theorem}[Exact status after the parallel test]
\label{thm:V84-status}
In the clean balanced \(2+2\) two-cross/no-prefix sector:
\begin{enumerate}[label=\textup{(\roman*)}]
\item common incidence-two double dominance is closed by V82;
\item higher-incidence, opposite-block, mass-compatible maximal
      partners are closed by V83;
\item strict mass ascents admit the fresh-cell/overlap alternative of
      \Cref{thm:V84-fresh-cell} and cannot continue through fresh cells
      for more than three transfers;
\item the proposed input-stock square-tail repair of a same-side
      maximal partner is false by
      \Cref{thm:V84-cross-square-false}; and
\item the unresolved same-side branch is localized to the complete
      output currencies
      \eqref{eq:V84-output-currency-unpaid}--%
      \eqref{eq:V84-output-currency-paid}, together with the explicit
      overlap alternative
      \Cref{fire:V84-overlap-open}.
\end{enumerate}
\end{theorem}

\begin{proof}
The first two items cite the stated theorems.  The third is
\Cref{cor:V84-ascent-acquisition}.  The one-coordinate lower bounds
\eqref{eq:V84-unpaid-lower}--%
\eqref{eq:V84-paid-lower} prove item \textup{(iv)}.  General protected
cancellation \Cref{cor:V84-general-protected-cancel} and the
intermediate-spin block decomposition prove that every surviving local
parallel covariance lies outside the invariant same-side output,
yielding item \textup{(v)}.  No input-stock estimate stronger than the
counterexample is asserted.
\end{proof}

\begin{openproblem}[V85: audited intermediate-output bank compiler]
\label{open:V84-V85-mandate}
\begin{enumerate}
\item Perform the compulsory read-only audit of the newest active SM
      and NS notes before changing the mathematical direction.
\item Form the complete-bank versions of
      \(\mathscr C_e^{\rm out,N}\) and
      \(\mathscr C_e^{\rm out,P}\) without summing individual
      multiplicity copies.
\item Use Clebsch--Gordan orthogonality or a weighted partial-trace
      identity to sum the factors
      \[
      \min\left\{1,\frac{2J+1}{2j+1}\right\}
      |q_K-q_Jq_R|
      \]
      without a representation-channel count.
\item Retain the \(J=0\) cancellation before absolute values and test
      the resulting compiler on the sequence \(j_s\asymp s^{-3}\);
      it must permit the sharp orders \(s/j_s\) unpaid and
      \(1/j_s\) paid.
\item Insert the compiled output currency into the two same-side raw
      responses and determine whether it is bounded by a finite legal
      tree atlas.  If a dimension/Casimir conversion is still missing,
      state it as an output-resolved inequality rather than restoring
      the false cross-stock square tail.
\item Treat the overlap alternative of
      \Cref{thm:V84-fresh-cell} only with a higher complete-bank moment
      or joint capacity; do not reuse a first-moment certificate.
\end{enumerate}
\end{openproblem}

\section{V84 conclusion and claim boundary}
\label{sec:V84-conclusion}

\begin{conclusion}[V84 conclusion]
\label{con:V84}
The same-side-maximal crossing square tail proposed in V83 is false.
An exact three-row \(SU(2)\) heat block disproves both its unpaid and
paid forms.  The obstruction is a low nontrivial same-side
intermediate spin whose capacity is inverse representation dimension,
not inverse Casimir.

V84 goes upstream rather than forcing the failed estimate.  It proves
exact cancellation of the intermediate singlet, derives the
dimension-weighted \(J>0\) output currency, and supplies a disjoint
fresh-cell acquisition theorem for strict mass ascents.  The next
analytic task is the support-uniform complete-bank summation of that
output currency; V85 begins with the compulsory companion audit.
\end{conclusion}

\begin{firewall}[No mass-gap claim]
\label{fire:V84-no-mass-gap}
V84 corrects a false intermediate target and proves a finite
acquisition reduction.  It does not close the intermediate-output or
overlap banks, the other first-connectivity source families, the
general compact-group extension, continuum construction,
Osterwalder--Schrader reconstruction, or a uniform positive spectral
gap.  The full Yang--Mills mass gap remains open.
\end{firewall}

\begin{record}[Parent and continuation record]
\label{rec:V84-parent}
This file was copied byte-for-byte from
\texttt{Renewal\_Yang\_Mills\_Mass\_Gap\_Research\_Notes\_V83.tex}
before the present chapter was appended.  The parent had \(37719\)
lines, \(1302259\) bytes, and SHA--256 digest
\[
\texttt{1405ff091de241e9cdc6839668f1c76f23d0d975bb8c636aaed9183fdf58afcc}.
\]
The next active version is V85 and begins with the compulsory audit
and \Cref{open:V84-V85-mandate}.  No companion result is imported
without a separate Yang--Mills proof.
\end{record}

% =============================================================================
% BEGIN APPENDED COMPACT VERSION V85 MATERIAL
% =============================================================================

\clearpage
\chapter{The audited output first moment and the signed
incidence-three obstruction}
\label{chap:V85}

\section{Checkpoint context and companion audit}
\label{sec:V85-audit}

V84 isolated a same-side intermediate-output currency and proposed
summing its isotypic blocks at complete-bank level.  Before attempting
that sum, this compulsory five-version checkpoint reads the newest active
companion notes and compares the proposed operation with the older YM
anti-aligned tests.

\begin{record}[Companion snapshots used by V85]
\label{rec:V85-companion-snapshots}
At the time of the audit the active companion files were:
\begin{enumerate}
\item
\(\texttt{Renewal\_SM\_Einstein\_Continuum\_Research\_Notes\_V19.tex}\),
with \(9231\) lines, \(316742\) bytes, and SHA--256
\[
\texttt{f4b1bcc5882d2f1718f2b0ccd3d5766002438cd3fd5e41c4f3c4c03db7838398};
\]
\item
\(\texttt{Renewal\_NS\_Terminal\_Source\_Concentration\_and\_Singular\_Thread\_Closure\_Programme\_V31.tex}\),
with \(23052\) lines, \(887537\) bytes, and SHA--256
\[
\texttt{f1121b62238ad5491bb7cf03635ea9934942edf573ed8faaa9ee79e1d38a6696}.
\]
\end{enumerate}
The files were read only.
\end{record}

\begin{assessment}[Type-correct audit transfer]
\label{ass:V85-audit-transfer}
The SM V18--V19 calculation cancels an acyclic quartet only after the
bosonic and fermionic operators, domains, and cutoffs have been paired.
Its relative heat calculation also shows that failure of a raw spectral
sum may represent genuine low-order data rather than a removable channel
count.  The NS V31 calculation removes an artificial dyadic multiplicity
by orthogonality and geometric parent summation, but the surviving first
moment returns exactly the critical source action.  Thus heat smoothing
does not make a genuine first moment subcritical.

The type-correct YM consequence is procedural.  Form the complete
isotypic covariance before choosing multiplicity bases, and test whether
the resulting output first moment is actually small.  If it is not,
return to the complete signed first-connectivity source before taking a
positive envelope.  No SM determinant theorem or NS energy estimate is
imported into the Yang--Mills proof.
\end{assessment}

\begin{record}[Late SM drift during the V85 calculation]
\label{rec:V85-late-SM-drift}
Before V85 was finalized, the parallel SM programme advanced to
\(\texttt{Renewal\_SM\_Einstein\_Continuum\_Research\_Notes\_V22.tex}\).
The new file was audited read only.  It had \(10507\) lines,
\(362531\) bytes, and SHA--256
\[
\texttt{d16a2c889f4188f3ad8f15245049660d2fb19a36ccbce49c944a07ccd0225648}.
\]
V22 proves that the first lower Dirichlet-to-Neumann symbol on a
nonproduct collar is direction dependent: angular averaging retains
only its mean-curvature part and loses the tracefree extrinsic data.
It also records that positivity of the full operator does not make
each lower symbol positive.

The corresponding YM proof-order lesson strengthens, but does not
change, the V85 direction: retain the complete signed output-channel
symbol before replacing it by an averaged scalar or positive carrier.
No Dirichlet-to-Neumann or Schatten estimate is imported.
\end{record}

\begin{assessment}[Internal nonduplication audit]
\label{ass:V85-internal-audit}
V74 already proves that a coordinatewise positive gapped terminal route
fails on a broad anti-aligned heat window, and V75 proves that a pure
two-row bank must instead be kept inside its exact two-term binary
covariance.  V84 concerns a different object: an incidence-three
same-side-pair/opposite-row joining covariance.  V85 therefore does not
repeat the V74 gapped-bank counterexample or the V75 pure-binary
calculation.  It sums the total-output channels of the V84
incidence-three operator and tests both its signed and positive forms.
\end{assessment}

\section{Complete isotypic carriers before multiplicity resolution}
\label{sec:V85-complete-carrier}

At one coordinate \(e\), group the two same-side rows into a representation
\(\mathcal H_e^A\) and the opposite rows into
\(\mathcal H_e^B\).  Let
\(\Pi_{J,R,K}^{e}\) denote the projection onto the complete simultaneous
isotypic subspace on which the \(A\)-output is \(J\), the \(B\)-output is
\(R\), and the combined output is \(K\).  The projection includes the
identity on every multiplicity space; it is not split into individual
Clebsch--Gordan copies.  Put
\[
 q_L=e^{-sC_2(L)}.
\]

\begin{definition}[Complete local signed and positive output carriers]
\label{def:V85-local-complete-carriers}
Define
\begin{align}
 \mathcal D_e
 &:=
 \sum_{J,R,K}(q_K-q_Jq_R)\Pi_{J,R,K}^{e},
\label{eq:V85-local-signed}\\
 \mathcal N_e^{\rm out}
 &:=
 |\mathcal D_e|_{\rm phys}
 =
 \sum_{\substack{J>0,R,K}}
 |q_K-q_Jq_R|\Pi_{J,R,K}^{e},
\label{eq:V85-local-unpaid}\\
 \mathcal P_e^{\rm out}
 &:=
 \sum_{\substack{J>0,R,K\\K\ne0}}
 \frac{A_K}{1-q_K}
 |q_K-q_Jq_R|\Pi_{J,R,K}^{e}.
\label{eq:V85-local-paid}
\end{align}
The physical absolute value is taken after simultaneous block
resolution, but before any multiplicity basis is chosen.
\end{definition}

\begin{lemma}[Invariant cancellation occurs before absolute values]
\label{lem:V85-invariant-before-absolute}
On the complete \(J=0\) isotypic subspace,
\[
 \mathcal D_e=0.
\]
Consequently the restrictions \(J>0\) in
\eqref{eq:V85-local-unpaid}--\eqref{eq:V85-local-paid} discard an exact
zero operator, including all invariant multiplicities.
\end{lemma}

\begin{proof}
If \(J=0\), then \(K=R\).  Hence
\[
 q_K-q_Jq_R=q_R-q_0q_R=0.
\]
Equivalently, apply
\Cref{cor:V84-general-protected-cancel}.  Since
\(\Pi_{0,R,R}^{e}\) is a complete isotypic projection, the argument is
the identity on its multiplicity space.
\end{proof}

Let \(E=\{e_1,\ldots,e_N\}\) be an ordered coordinate bank.  Write
\(\mathcal M_e^+\) and \(\mathcal M_e^\rho\) for its common and
cut-replicated positive heat endpoints.  They are commuting positive
contractions on every physical block.  Put
\[
 \mathcal M_E^+:=\bigotimes_{e\in E}\mathcal M_e^+,
 \qquad
 \mathcal M_E^\rho:=\bigotimes_{e\in E}\mathcal M_e^\rho.
\]

\begin{definition}[Complete-bank output carriers]
\label{def:V85-bank-output-carriers}
The exact signed bank covariance and its telescopic unpaid carrier are
\begin{align}
 \mathcal D_E
 &:=
 \mathcal M_E^+-\mathcal M_E^\rho,
\label{eq:V85-bank-signed}\\
 \mathcal N_E^{\rm out}
 &:=
 \sum_{a=1}^{N}
 \left(\bigotimes_{b<a}\mathcal M_{e_b}^+\right)
 \otimes\mathcal N_{e_a}^{\rm out}
 \otimes
 \left(\bigotimes_{b>a}\mathcal M_{e_b}^\rho\right).
\label{eq:V85-bank-unpaid}
\end{align}
Resolve \(\mathcal D_E\) only into its complete product-isotypic
projections \(\Pi_{\boldsymbol\nu}\), and write
\[
 \mathcal D_E=\sum_{\boldsymbol\nu}
 d_E(\boldsymbol\nu)\Pi_{\boldsymbol\nu}.
\]
Write
\(\boldsymbol\nu=((J_e,R_e,K_e))_{e\in E}\).  The V2 allocated paid
majorant is
\begin{equation}
 \mathcal P_E^{\rm out}
 :=
 \sum_{a=1}^{N}
 \sum_{\substack{\boldsymbol\nu\\K_{e_a}\ne0}}
 \frac{A_{K_{e_a}}}{1-q_{K_{e_a}}}
 |d_E(\boldsymbol\nu)|\Pi_{\boldsymbol\nu},
\label{eq:V85-bank-paid}
\end{equation}
where a zero output is omitted from the corresponding payer-position
sum.
\end{definition}

\begin{theorem}[The bank construction has no multiplicity-copy sum]
\label{thm:V85-no-multiplicity-sum}
The definitions
\eqref{eq:V85-bank-unpaid}--\eqref{eq:V85-bank-paid} never resolve a
multiplicity basis.  The unpaid carrier uses complete isotypic
projections, and the paid carrier has one copy of such a projection for
each coordinate payer position, with the identity on its multiplicity
space.
Moreover,
\begin{equation}
 |\mathcal D_E|_{\rm phys}
 \preccurlyeq \mathcal N_E^{\rm out}.
\label{eq:V85-bank-majorization}
\end{equation}
For \(E=\{e\}\), the bank carriers reduce exactly to
\(\mathcal N_e^{\rm out}\) and \(\mathcal P_e^{\rm out}\).
\end{theorem}

\begin{proof}
The tensor-product difference has the exact telescope
\[
 \mathcal D_E
 =
 \sum_{a=1}^{N}
 \left(\bigotimes_{b<a}\mathcal M_{e_b}^+\right)
 \otimes\mathcal D_{e_a}
 \otimes
 \left(\bigotimes_{b>a}\mathcal M_{e_b}^\rho\right).
\]
All displayed factors are functions of commuting physical Casimirs on
distinct coordinates.  Simultaneous block resolution and the scalar
triangle inequality give
\eqref{eq:V85-bank-majorization}.  Apply
\Cref{cor:V2-carrier-allocation} on each complete product-isotypic
block.  It produces the finite sum over coordinate payer positions in
\eqref{eq:V85-bank-paid}, not a product of resolvents, and neither
selects nor sums multiplicity vectors.
The singleton assertion is immediate.
\end{proof}

\begin{firewall}[Removal of a basis count is not an analytic estimate]
\label{fire:V85-no-count-not-small}
\Cref{thm:V85-no-multiplicity-sum} removes an artificial
multiplicity-basis expansion.  It does not say that the remaining
isotypic first moment is uniformly small.  Orthogonal output sectors add
inside the expectation of a positive carrier, and that addition can be a
genuine heat-scale effect.
\end{firewall}

\section{A signed product gate and the complete spin-one channel sum}
\label{sec:V85-signed-gate}

\begin{definition}[Signed product-state radius]
\label{def:V85-signed-radius}
For a self-adjoint four-row coefficient \(T\), set
\begin{equation}
 \omega_\otimes(T)
 :=
 \sup_{\substack{\rho_i\ge0\\\Tr\rho_i=1}}
 \left|
 \Tr\!\left[T(\rho_1\otimes\rho_2\otimes\rho_3\otimes\rho_4)\right]
 \right|.
\label{eq:V85-signed-radius}
\end{equation}
\end{definition}

\begin{lemma}[Raw signed coefficient gate]
\label{lem:V85-signed-gate}
For arbitrary trace-class row coefficients \(A_i\),
\begin{equation}
 \left|
 \Tr\!\left[T(A_1\otimes A_2\otimes A_3\otimes A_4)\right]
 \right|
 \le
 16\,\omega_\otimes(T)\prod_{i=1}^4\|A_i\|_1.
\label{eq:V85-signed-gate}
\end{equation}
If \(T\) has a physical spectral absolute value, then
\begin{equation}
 \omega_\otimes(T)\le\kappa_\otimes(|T|_{\rm phys}).
\label{eq:V85-signed-vs-positive}
\end{equation}
\end{lemma}

\begin{proof}
Use \Cref{lem:V1-positive-decomposition} on each \(A_i\).  Expanding
the tensor product gives products of four positive operators.  After
normalizing their traces, the absolute value of every expectation is
at most \(\omega_\otimes(T)\), while the product of the four sums of
positive-part traces is at most
\(2^4\prod_i\|A_i\|_1\).  This proves
\eqref{eq:V85-signed-gate}.  The spectral inequalities
\(-|T|_{\rm phys}\preccurlyeq T\preccurlyeq|T|_{\rm phys}\) prove
\eqref{eq:V85-signed-vs-positive}.
\end{proof}

Return to the V84 coordinate with same-side spins \(j,j\), opposite
spin \(1\), and product vector
\[
 \Psi_j:=\psi_j\otimes|1,0\rangle.
\]
The inactive fourth row is suppressed.  Put
\begin{equation}
 p_{j,J}
 =
 (2J+1)
 \frac{[(2j)!]^2}{(2j-J)!(2j+J+1)!},
 \qquad 0\le J\le2j.
\label{eq:V85-pjJ}
\end{equation}
For \(J\ge1\), define
\begin{align}
 u_J&:=\frac{J}{2J+1},
 &
 v_J&:=\frac{J+1}{2J+1},
\label{eq:V85-uv}\\
 d_J^-&:=q_{J-1}-q_Jq_1,
 &
 d_J^+&:=q_Jq_1-q_{J+1}.
\label{eq:V85-dpm}
\end{align}
Both differences in \eqref{eq:V85-dpm} are positive.

\begin{lemma}[Exact complete-channel expectations]
\label{lem:V85-exact-channel-expectations}
Let
\begin{align}
 \mathcal E_{s,j}^{\rm S}
 &:=
 \langle\Psi_j,\mathsf J_{s,j}\Psi_j\rangle,
\label{eq:V85-ES}\\
 \mathcal E_{s,j}^{\rm N}
 &:=
 \langle\Psi_j,\mathsf N_{s,j}\Psi_j\rangle,
\label{eq:V85-EN}\\
 \mathcal E_{s,j}^{\rm P}
 &:=
 \langle\Psi_j,\mathsf P_{s,j}\Psi_j\rangle.
\label{eq:V85-EP}
\end{align}
Then
\begin{align}
 \mathcal E_{s,j}^{\rm S}
 &=
 \sum_{J=1}^{2j}p_{j,J}
 \bigl(u_Jd_J^- -v_Jd_J^+\bigr),
\label{eq:V85-exact-S}\\
 \mathcal E_{s,j}^{\rm N}
 &=
 \sum_{J=1}^{2j}p_{j,J}
 \bigl(u_Jd_J^- +v_Jd_J^+\bigr),
\label{eq:V85-exact-N}\\
 \mathcal E_{s,j}^{\rm P}
 &=
 \sum_{J=2}^{2j}p_{j,J}u_J
 \frac{A_{J-1}}{1-q_{J-1}}d_J^-
 {}+
 \sum_{J=1}^{2j}p_{j,J}v_J
 \frac{A_{J+1}}{1-q_{J+1}}d_J^+.
\label{eq:V85-exact-P}
\end{align}
In particular, these are expectations of the complete isotypic
operators; no multiplicity or output-channel triangle inequality has
been used.
\end{lemma}

\begin{proof}
The spin-\(J\) component of \(\psi_j\) has magnetic number zero and
weight \eqref{eq:V85-pjJ}.  The only nonzero couplings with
\(|1,0\rangle\) are
\begin{align*}
 |\langle J,0;1,0\mid J-1,0\rangle|^2&=u_J,\\
 |\langle J,0;1,0\mid J+1,0\rangle|^2&=v_J,
\end{align*}
whereas
\(|\langle J,0;1,0\mid J,0\rangle|^2=0\).
The \(K=J-1\) covariance coefficient is \(d_J^-\), and the
\(K=J+1\) coefficient is \(-d_J^+\).  This proves
\eqref{eq:V85-exact-S} and, after physical absolute value,
\eqref{eq:V85-exact-N}.  Applying the unique payer factor to every
nonzero \(K\) gives \eqref{eq:V85-exact-P}; the \(J=1,K=0\) term is
correctly omitted.  At \(J=0\), the sole \(K=1\) coefficient is zero
by \Cref{lem:V85-invariant-before-absolute}.
\end{proof}

\begin{lemma}[Uniform anti-aligned weights on the heat window]
\label{lem:V85-window-weights}
There are universal \(c,C>0\) such that
\begin{equation}
 p_{j,J}\le C\frac{J}{j}
 \qquad(1\le J\le2j).
\label{eq:V85-p-upper}
\end{equation}
For
\begin{equation}
 \mathcal I_s
 :=
 \left\{J\in\mathbb N:
 \frac1{2\sqrt s}\le J\le\frac1{\sqrt s}\right\},
\label{eq:V85-heat-window}
\end{equation}
if \(0<s\le s_0\) and \(j\ge16s^{-1}\), then
\begin{equation}
 c\frac{J}{j}\le p_{j,J}\le C\frac{J}{j}
 \qquad(J\in\mathcal I_s).
\label{eq:V85-p-window}
\end{equation}
\end{lemma}

\begin{proof}
Set \(n=2j\).  Factor \eqref{eq:V85-pjJ} as
\begin{equation}
 p_{j,J}
 =
 \frac{2J+1}{n+J+1}
 \prod_{r=0}^{J-1}\frac{n-r}{n+r+1}.
\label{eq:V85-p-product}
\end{equation}
The product is at most one, which proves
\eqref{eq:V85-p-upper}.  When \(J^2\le n/16\), elementary logarithmic
bounds give
\[
 \sum_{r=0}^{J-1}
 \log\frac{n-r}{n+r+1}
 \ge-C\frac{J^2}{n},
\]
so the product in \eqref{eq:V85-p-product} is bounded below by a
universal positive constant.  Under the stated hypotheses,
\(J\in\mathcal I_s\) implies \(J^2\le s^{-1}\le n/32\).
The prefactor is comparable to \(J/n\), proving
\eqref{eq:V85-p-window}.
\end{proof}

\begin{lemma}[First-order channel cancellation and positive first moments]
\label{lem:V85-channel-scales}
Set
\begin{align}
 \Delta_J^{\rm S}
 &:=
 u_Jd_J^- -v_Jd_J^+,
\label{eq:V85-DeltaS}\\
 \Delta_J^{\rm N}
 &:=
 u_Jd_J^- +v_Jd_J^+,
\label{eq:V85-DeltaN}
\end{align}
and, for \(J\ge2\),
\begin{equation}
 \Delta_J^{\rm P}
 :=
 u_J\frac{A_{J-1}}{1-q_{J-1}}d_J^-
 +v_J\frac{A_{J+1}}{1-q_{J+1}}d_J^+.
\label{eq:V85-DeltaP}
\end{equation}
For \(J=1\), let \(\Delta_1^{\rm P}\) be only the second term in
\eqref{eq:V85-DeltaP}.  Uniformly for \(0<s\le s_0\) and \(J\ge1\),
\begin{align}
 0\le\Delta_J^{\rm S}
 &\le
 Cs^2(J+1)^2e^{-cs(J-1)^2},
\label{eq:V85-signed-global}\\
 \Delta_J^{\rm N}
 &\le
 Cs(J+1)e^{-cs(J-1)^2},
\label{eq:V85-unpaid-global}\\
 \Delta_J^{\rm P}
 &\le
 C\bigl(s(J+1)^3+(J+1)\bigr)
 e^{-cs(J-1)^2}.
\label{eq:V85-paid-global}
\end{align}
On the heat window \(J\in\mathcal I_s\),
\begin{align}
 cs\le\Delta_J^{\rm S}&\le Cs,
\label{eq:V85-signed-window}\\
 c\sqrt s\le\Delta_J^{\rm N}&\le C\sqrt s,
\label{eq:V85-unpaid-window}\\
 cs^{-1/2}\le\Delta_J^{\rm P}&\le Cs^{-1/2}.
\label{eq:V85-paid-window}
\end{align}
\end{lemma}

\begin{proof}
The signed conditional coefficient has the exact form
\begin{equation}
 \Delta_J^{\rm S}
 =
 q_J\left[
 u_Je^{2sJ}
 +v_Je^{-2s(J+1)}
 -e^{-2s}\right].
\label{eq:V85-Jensen-gap}
\end{equation}
Let \(X_J\) take the two values
\[
 2J,\qquad -2(J+1)
\]
with probabilities \(u_J,v_J\).  Direct calculation gives
\begin{equation}
 \mathbb EX_J=-2,
 \qquad
 \operatorname{Var}(X_J)=4J(J+1).
\label{eq:V85-X-moments}
\end{equation}
Convexity of the exponential applied to
\eqref{eq:V85-Jensen-gap} proves
\(\Delta_J^{\rm S}\ge0\).  The centered Taylor remainder
\[
 e^y-1-y\le\frac12y^2e^{|y|}
\]
and \eqref{eq:V85-X-moments}, followed by multiplication by \(q_J\),
give
\[
 \Delta_J^{\rm S}
 \le
 Cs^2(J+1)^2e^{-sJ(J-1)}.
\]
This is \eqref{eq:V85-signed-global}.  On
\(\mathcal I_s\), all exponents in
\eqref{eq:V85-Jensen-gap} lie in a fixed compact interval and
\(q_J\asymp1\).  The exponential has a uniform strictly positive
second derivative there, so its quantitative Jensen remainder and
\eqref{eq:V85-X-moments} give
\(\Delta_J^{\rm S}\asymp s^2J(J+1)\asymp s\).

The two unsigned differences can be written
\begin{align*}
 d_J^-
 &=
 q_{J-1}\bigl(1-e^{-2s(J+1)}\bigr),\\
 d_J^+
 &=
 q_Jq_1\bigl(1-e^{-2sJ}\bigr).
\end{align*}
Using \(1-e^{-x}\le\min\{1,x\}\) proves
\eqref{eq:V85-unpaid-global}; on
\(\mathcal I_s\), both differences are comparable to \(\sqrt s\),
which proves \eqref{eq:V85-unpaid-window}.

For every nonzero spin \(K\),
\begin{equation}
 \frac{A_K}{1-q_K}
 \le C\bigl(K^2+s^{-1}\bigr),
\label{eq:V85-payer-elementary}
\end{equation}
because \(1-e^{-sK(K+1)}\) is comparable to
\(\min\{1,sK(K+1)\}\).  Combine
\eqref{eq:V85-payer-elementary} with the preceding bounds on
\(d_J^\pm\) to obtain \eqref{eq:V85-paid-global}.  On the heat
window, \(A_{J\pm1}\asymp s^{-1}\), the two nonzero denominators are
bounded above and below away from zero, and
\(d_J^\pm\asymp\sqrt s\).  This proves
\eqref{eq:V85-paid-window}.
\end{proof}

\begin{theorem}[Sharp complete-channel accumulation on one product vector]
\label{thm:V85-sharp-channel-accumulation}
There are universal \(c,C,s_0>0\) such that, whenever
\[
 0<s\le s_0,
 \qquad
 j\ge16s^{-1},
\]
the complete expectations of
\Cref{lem:V85-exact-channel-expectations} satisfy
\begin{align}
 \frac{c}{j}
 \le\mathcal E_{s,j}^{\rm S}
 &\le\frac{C}{j},
\label{eq:V85-sharp-S}\\
 \frac{c}{j\sqrt s}
 \le\mathcal E_{s,j}^{\rm N}
 &\le\frac{C}{j\sqrt s},
\label{eq:V85-sharp-N}\\
 \frac{c}{js^{3/2}}
 \le\mathcal E_{s,j}^{\rm P}
 &\le\frac{C}{js^{3/2}}.
\label{eq:V85-sharp-P}
\end{align}
Thus absolute value before the total-output sum loses exactly one
factor \(s^{-1/2}\) on this family, and the output-mass/resolvent payer
costs one further factor \(s^{-1}\).
\end{theorem}

\begin{proof}
For every integer \(m\ge0\),
\begin{equation}
 \sum_{J\ge1}J^m e^{-cs(J-1)^2}
 \le C_m s^{-(m+1)/2}.
\label{eq:V85-Gaussian-sum}
\end{equation}
Use \eqref{eq:V85-p-upper} and
\eqref{eq:V85-signed-global} to get
\[
 \mathcal E_{s,j}^{\rm S}
 \le
 \frac{Cs^2}{j}\sum_{J\ge1}J(J+1)^2e^{-cs(J-1)^2}
 \le\frac Cj.
\]
Similarly,
\eqref{eq:V85-unpaid-global} gives
\[
 \mathcal E_{s,j}^{\rm N}
 \le
 \frac{Cs}{j}\sum_{J\ge1}J(J+1)e^{-cs(J-1)^2}
 \le\frac{C}{j\sqrt s},
\]
and \eqref{eq:V85-paid-global} gives
\[
 \mathcal E_{s,j}^{\rm P}
 \le
 \frac Cj
 \sum_{J\ge1}
 \bigl(sJ(J+1)^3+J(J+1)\bigr)e^{-cs(J-1)^2}
 \le\frac{C}{js^{3/2}}.
\]

For the lower bounds, restrict the exact positive sums to
\(\mathcal I_s\).  By
\eqref{eq:V85-p-window} and
\eqref{eq:V85-signed-window},
\[
 \mathcal E_{s,j}^{\rm S}
 \ge
 \frac{cs}{j}\sum_{J\in\mathcal I_s}J
 \ge\frac cj.
\]
The same calculation with
\eqref{eq:V85-unpaid-window} gives
\[
 \mathcal E_{s,j}^{\rm N}
 \ge
 \frac{c\sqrt s}{j}\sum_{J\in\mathcal I_s}J
 \ge\frac{c}{j\sqrt s},
\]
and \eqref{eq:V85-paid-window} gives
\[
 \mathcal E_{s,j}^{\rm P}
 \ge
 \frac{c}{j\sqrt s}\sum_{J\in\mathcal I_s}J
 \ge\frac{c}{js^{3/2}}.
\]
The window contains a fixed multiple of \(s^{-1/2}\) integers and
its sum of \(J\)'s is comparable to \(s^{-1}\).  This completes all
three estimates.
\end{proof}

\begin{corollary}[The channel sum is not a multiplicity artefact]
\label{cor:V85-not-multiplicity-artifact}
Choose integers \(j_s\) with
\[
 s^{-3}\le j_s\le2s^{-3}.
\]
Then
\begin{align}
 \omega_\otimes(\mathsf J_{s,j_s})
 &\ge c s^3,
\label{eq:V85-raw-lower-sequence}\\
 \kappa_\otimes(\mathsf N_{s,j_s})
 &\ge c s^{5/2},
\label{eq:V85-unpaid-lower-sequence}\\
 \kappa_\otimes(\mathsf P_{s,j_s})
 &\ge c s^{3/2}.
\label{eq:V85-paid-lower-sequence}
\end{align}
The same input stocks satisfy
\begin{align}
 \frac{H_{1,C_1}}{sS_1^2}&\asymp s^5,
 &
 \frac{H_{1,C_1}}{s^2S_1^2}&\asymp s^4.
\label{eq:V85-false-target-scales}
\end{align}
Consequently even the signed local covariance cannot satisfy the
unpaid crossing square-tail target, while the complete positive
unpaid and paid carriers are larger by still greater powers.
\end{corollary}

\begin{proof}
The three lower bounds follow from
\Cref{thm:V85-sharp-channel-accumulation}, because \(\Psi_{j_s}\) is
a product vector and the two positive carriers dominate their
expectations.  The signed expectation is positive by
\Cref{lem:V85-channel-scales}, so it lower-bounds
\(\omega_\otimes\).  Finally use
\eqref{eq:V84-input-stocks} and \(a_{j_s}\asymp j_s^2\).
\end{proof}

\begin{firewall}[Strengthening, not repeating, V84]
\label{fire:V85-strengthens-V84}
V84 used one fixed \(J=1\) block and obtained the sufficient lower
orders \(s/j\) unpaid and \(1/j\) paid.  V85 sums the complete physical
operator on the same product vector and proves the sharper orders
\(1/(j\sqrt s)\) and \(1/(js^{3/2})\).  It also quantifies the
first-order cancellation retained by the signed \(K=J-1,J+1\) pair.
Therefore the new result is not another proof of the V84
counterexample; it identifies the actual heat-window obstruction.
\end{firewall}

\section{The forced upstream packet and a finite overlap witness}
\label{sec:V85-upstream}

\begin{theorem}[No local signed repair of the crossing square tail]
\label{thm:V85-no-local-signed-repair}
There is no universal constant \(C\) such that the signed local
incidence-three covariance obeys
\begin{equation}
 \omega_\otimes(\mathsf J_{s,j})
 \le
 C\frac{H_{1,C_1}}{sS_1^2}
\label{eq:V85-false-signed-tail}
\end{equation}
for all small \(s\) and all \(j\).  Hence retaining the two
total-output signs inside the isolated local joining covariance does
not repair the V84 terminal estimate.
\end{theorem}

\begin{proof}
Use \(j=j_s\) from
\Cref{cor:V85-not-multiplicity-artifact}.  The left side of
\eqref{eq:V85-false-signed-tail} is at least \(cs^3\), while its
right side is at most \(Cs^5\).  Their ratio tends to infinity.
\end{proof}

\begin{corollary}[The next cancellation must precede the local gate]
\label{cor:V85-prelocal-cancellation}
Any proof of the unresolved same-side-maximal branch which first
replaces the incidence-three joining coordinate by either
\[
 \omega_\otimes(\mathcal D_e),\qquad
 \kappa_\otimes(\mathcal N_e^{\rm out}),\qquad\hbox{or}\qquad
 \kappa_\otimes(\mathcal P_e^{\rm out})
\]
and then seeks the V84 normalization
\(CH_{1,C_1}/(sS_1^2)\), or its paid \(s^{-2}\) analogue, cannot
establish that crossing square-tail card without an additional
currency.  Within the present route, the joining coordinate must
therefore remain inside the complete signed first-connectivity source
together with at least one prefix or partition-state difference.
\end{corollary}

\begin{proof}
The first replacement is ruled out by
\Cref{thm:V85-no-local-signed-repair}; the two positive replacements
are ruled out, with stronger lower bounds, by
\Cref{cor:V85-not-multiplicity-artifact}.  The only source factors
discarded before these three local gates are the signed
first-connectivity prefix and partition-state factors.  Therefore any
additional cancellation must be sought before that discard.  This is
a necessity statement, not an assertion that the larger packet
already closes.
\end{proof}

\begin{assessment}[Quantitative gain still required]
\label{ass:V85-required-gain}
On the sequence \(j_s\asymp s^{-3}\), the signed local expectation is
of order \(j_s^{-1}\), whereas the crossing square-tail target is of
order \((sj_s^2)^{-1}\).  The missing packet must therefore gain at
least
\begin{equation}
 \frac1{sj_s}
\label{eq:V85-required-gain}
\end{equation}
relative to the isolated signed joining covariance.  A mere
constant-factor cancellation, or the already realized
\(\sqrt s\)-gain between positive and signed output sums, is
insufficient.
\end{assessment}

We also sharpen the overlap alternative of
\Cref{fire:V84-overlap-open}.  Let
\(\mathcal B_1,\ldots,\mathcal B_m\) be the pairwise
coordinate-disjoint banks acquired along one fresh strict-ascent path,
and put
\(\mathcal U=\bigcup_{a=1}^m\mathcal B_a\).

\begin{theorem}[Finite one-bank witness for every ascent overlap]
\label{thm:V85-overlap-witness}
If
\begin{equation}
 X_r(\mathcal U)\ge\delta\Xi_r,
\label{eq:V85-overlap-hypothesis}
\end{equation}
then some previously acquired bank satisfies
\begin{equation}
 \boxed{
 X_r(\mathcal B_a)\ge\frac{\delta}{m}\Xi_r.}
\label{eq:V85-overlap-witness}
\end{equation}
Along a four-row strict-ascent path one may take \(m\le4\), and hence
the right side is at least \(\delta\Xi_r/4\).
\end{theorem}

\begin{proof}
Disjointness of the acquired banks gives
\[
 X_r(\mathcal U)
 =
 \sum_{a=1}^mX_r(\mathcal B_a).
\]
If every summand were smaller than
\(\delta\Xi_r/m\), their sum would contradict
\eqref{eq:V85-overlap-hypothesis}.  By
\Cref{lem:V83-ascent-terminates}, a strict ascent visits at most four
rows and has at most three edges.  Counting the initial bank and every
newly acquired bank gives \(m\le4\).
\end{proof}

\begin{corollary}[The overlap analytic gate is one-bank and joint]
\label{cor:V85-overlap-joint-gate}
The overlap branch does not require a sum over arbitrary subsets of
the path history.  After a loss of at most \(4/\delta\), it reduces to
one identified old bank \(\mathcal B_a\) which is large at the current
row and occurs in both the earlier and current source factors.  Its
admissible estimate must be a joint complete-bank moment or joint
product-state radius; two separate applications of its first-moment
card remain forbidden.
\end{corollary}

\begin{proof}
Choose the bank supplied by
\Cref{thm:V85-overlap-witness}.  Its lower mass fraction is fixed, so
all other history banks may be discarded in the scalar acquisition
step.  The same physical coordinate bank is present at two locations
of the source construction, so separate first-moment estimates would
charge it twice.  Keeping the two appearances in one joint carrier is
exactly the alternative allowed by
\Cref{fire:V84-overlap-open}.
\end{proof}

\begin{firewall}[The witness is not the missing joint estimate]
\label{fire:V85-overlap-not-closed}
\Cref{thm:V85-overlap-witness} removes combinatorial uncertainty from
the overlap branch but does not bound the repeated-bank joint
carrier.  In particular, the square of
\(X_r(\mathcal B_a)\) is not declared to be paid.  A higher
complete-bank moment, with the two source occurrences assembled before
capacity, remains necessary.
\end{firewall}

\section{Post-V85 reduction and next acquisition order}
\label{sec:V85-reduction}

\begin{theorem}[Exact status after the audited channel sum]
\label{thm:V85-status}
Relative to V1--V84, V85 proves:
\begin{enumerate}[label=\textup{(V85.\arabic*)},leftmargin=6.5em]
\item the V85 companion audit was performed on the initial SM V19 and
      NS V31 snapshots and refreshed after the late SM drift to V22;
      only proof order was transferred;
\item complete local and bank output carriers can be formed with
      complete isotypic projections and no multiplicity-copy sum; the
      paid carrier has only the finite V2 sum over coordinate payer
      positions;
\item the \(J=0\) isotypic block cancels before physical absolute value,
      including arbitrary invariant multiplicity;
\item for the exact spin-one incidence-three test, the signed,
      positive unpaid, and positive paid complete-channel expectations
      have the sharp orders
      \[
      j^{-1},\qquad (j\sqrt s)^{-1},\qquad
      (js^{3/2})^{-1};
      \]
\item neither multiplicity-free positive aggregation nor the signed
      isolated local covariance can satisfy the crossing-stock
      square-tail target;
\item the same-side-maximal branch must be returned to the complete
      signed first-connectivity packet before its local joining
      coordinate is normed; and
\item every ascent-overlap branch has one witness among at most four
      previously acquired banks, reducing its analytic obstruction to
      a single repeated-bank joint carrier.
\end{enumerate}
\end{theorem}

\begin{proof}
Items \textup{(V85.1)--(V85.3)} are
\Cref{rec:V85-companion-snapshots,rec:V85-late-SM-drift,
thm:V85-no-multiplicity-sum,
lem:V85-invariant-before-absolute}.  Item \textup{(V85.4)} is
\Cref{thm:V85-sharp-channel-accumulation}.  Items
\textup{(V85.5)--(V85.6)} are
\Cref{cor:V85-not-multiplicity-artifact,
thm:V85-no-local-signed-repair,
cor:V85-prelocal-cancellation}.  Item \textup{(V85.7)} is
\Cref{thm:V85-overlap-witness,cor:V85-overlap-joint-gate}.
\end{proof}

\begin{openproblem}[V86: complete source-packet cancellation]
\label{open:V85-V86-mandate}
\begin{enumerate}
\item Return to the exact V61 first-connectivity formula for the
      same-side-maximal \(2+2\) branch.  Retain the two binary prefixes,
      all joining partition-state signs, the incidence-three
      coordinate \(\mathcal D_e\), and the common suffix before taking
      a product-state gate.
\item Use the deterministic order freedom of
      \Cref{thm:V75-order-covariance} to put every pure same-side bank
      into its binary prefix.  Do not misclassify the incidence-three
      coordinate itself as a pure-pair bank.
\item Compress the complete signed packet, not the isolated joining
      coordinate, to the heat window
      \(\mathcal I_s\).  Compute whether the partition-state sum gains
      the factor \((sj)^{-1}\) required by
      \eqref{eq:V85-required-gain}.
\item If the gain fails, preserve the smallest exact
      prefix/joining/suffix product vector as a counterpacket and move
      upstream to the assembled cumulant identity before
      first-connectivity splitting.
\item For the overlap branch, insert the one bank from
      \Cref{thm:V85-overlap-witness} into a two-occurrence joint carrier
      and seek a second-moment or joint signed-radius estimate.  Never
      multiply two first-moment bounds for that bank.
\item Test every proposed estimate on both the V74 broad
      anti-aligned bank and the V85 spin-one incidence-three sequence.
      The next compulsory companion audit is V90.
\end{enumerate}
\end{openproblem}

\section{V85 conclusion and claim boundary}
\label{sec:V85-conclusion}

\begin{conclusion}[V85 conclusion]
\label{con:V85}
The complete-bank output carriers can be written without expanding
multiplicity copies, but that bookkeeping improvement does not close
the same-side branch.  The complete spin-one channel sum has a genuine
heat-window first moment.  Exact output-channel signs cancel its
first-order component and improve the unpaid scale by \(\sqrt s\), yet
even the remaining signed covariance is too large for the proposed
crossing square tail.

The programme therefore moves one level upstream: V86 must retain the
incidence-three coordinate inside the full signed first-connectivity
packet.  Separately, the ascent-overlap branch is now localized to one
of at most four old banks, but the corresponding joint repeated-bank
estimate remains open.
\end{conclusion}

\begin{firewall}[No mass-gap claim]
\label{fire:V85-no-mass-gap}
V85 proves a sharp local obstruction and two finite reductions.  It
does not prove the required complete source-packet cancellation, the
repeated-bank joint estimate, the remaining source families, a
general compact-group bound, continuum construction,
Osterwalder--Schrader reconstruction, or a uniform positive spectral
gap.  The full Yang--Mills mass gap remains open.
\end{firewall}

\begin{record}[Parent and continuation record]
\label{rec:V85-parent}
This file was copied byte-for-byte from
\texttt{Renewal\_Yang\_Mills\_Mass\_Gap\_Research\_Notes\_V84.tex}
before the present chapter was appended.  The parent had \(38416\)
lines, \(1325375\) bytes, and SHA--256 digest
\[
\texttt{01c79c9bd6cd741ec512c456115cb37df1a5faf042eba1eef9f35eaf72fc37ab}.
\]
The next active version is V86 and starts from
\Cref{open:V85-V86-mandate}.  No companion theorem is imported without
an independent Yang--Mills proof.
\end{record}

% =============================================================================
% BEGIN APPENDED COMPACT VERSION V86 MATERIAL
% =============================================================================

\clearpage
\chapter{The complete \(2+2\) packet neutralizes the local heat-window
test}
\label{chap:V86}

\section{Context: restore both binary prefixes}
\label{sec:V86-context}

V85 proves that the isolated incidence-three joining covariance remains
too large even when its two total-output signs are retained.  The V61
first-connectivity formula never presents that coordinate in isolation.
For prefix join
\[
 \rho=12|34,
\]
\Cref{thm:V61-prefix-factorization,lem:V65-double-replica} give the exact
source summand
\begin{equation}
 \mathcal Q_r
 =
 K_{12}^{<r,\mathrm{conn}}
 \otimes K_{34}^{<r,\mathrm{conn}}
 \otimes
 \operatorname{Cov}_r(A_r,B_r)
 \otimes M_{>r}.
\label{eq:V86-exact-S22-summand}
\end{equation}
The two binary prefixes are part of the same packet as the joining
coordinate.  Dropping them before the local capacity test created the
V85 obstruction.

\begin{lemma}[Deterministic order puts every pure pair into the prefix]
\label{lem:V86-pure-prefix-order}
Fix pure-pair banks \(D_{12}\) and \(D_{34}\), and an
incidence-three coordinate \(e\) meeting rows \(1,2\) and at least one
row of \(\{3,4\}\).  There is a deterministic first-connectivity order
in which every coordinate of \(D_{12}\cup D_{34}\) precedes \(e\).
If their cumulative prefix join is \(12|34\) and \(e\) is the first
coordinate joining the two blocks, then its exact V61 packet is
\eqref{eq:V86-exact-S22-summand}.  The coordinate \(e\) itself is not
included in either pure-pair prefix.
\end{lemma}

\begin{proof}
Use \Cref{thm:V75-order-covariance} with
\(D_{12}\cup D_{34}\) first, followed by all higher-incidence
coordinates in a fixed support order.  Pure-pair coordinates refine
\(12|34\), so none can connect the two blocks.  At the stated first
joining coordinate, the prefix carrier factors as
\[
 C_{<e}(12|34)
 =
 K_{12}^{<e,\mathrm{conn}}
 \otimes K_{34}^{<e,\mathrm{conn}}
\]
by \Cref{thm:V61-prefix-factorization}.  The joining letter is the
assembled composite covariance by
\Cref{lem:V65-double-replica}; the unrestricted later coordinates form
\(M_{>e}\).  These are precisely the factors in
\eqref{eq:V86-exact-S22-summand}.
\end{proof}

\begin{firewall}[There is no cross-coordinate sign cancellation]
\label{fire:V86-no-cross-coordinate-sign}
For a fixed first joining coordinate, the prefix, joining, and suffix
operators act on distinct coordinate tensor factors.  Their
expectations factor on row-product vectors which are tensor products
within each row.  Thus V86 does not claim that a sign at an earlier
coordinate cancels a sign at the joining coordinate.  The missing gain
must be an actual small prefix response, retained multiplicatively
inside the complete packet.
\end{firewall}

\section{A sharp equal-spin binary-prefix heat card}
\label{sec:V86-binary-heat-card}

On \(V_j\otimes V_j\), let
\begin{align}
 \mathsf T_{s,j}
 &:=
 \sum_{J=0}^{2j}q_JP_J,
\label{eq:V86-pair-heat}\\
 \mathsf K_{s,j}
 &:=
 \mathsf T_{s,j}-q_j^2I.
\label{eq:V86-pair-covariance}
\end{align}
The second operator is the exact one-coordinate binary covariance of
two spin-\(j\) rows.

\begin{theorem}[Equal-spin binary-prefix heat capacity]
\label{thm:V86-binary-heat-capacity}
Uniformly for \(0<s\le s_0\) and \(j\ge\frac12\),
\begin{equation}
 \boxed{
 \omega_\otimes(\mathsf K_{s,j})
 \le
 C\min\left\{1,\frac1{s(2j+1)}\right\}.}
\label{eq:V86-binary-heat-capacity}
\end{equation}
The scale is sharp: if \(s\downarrow0\) and \(sj\to\infty\), then on
the anti-aligned vector \(\psi_j\),
\begin{equation}
 \langle\psi_j,\mathsf K_{s,j}\psi_j\rangle
 \asymp\frac1{sj}.
\label{eq:V86-binary-heat-sharp}
\end{equation}
\end{theorem}

\begin{proof}
By \Cref{lem:V84-intermediate-dimension},
\[
 \kappa_\otimes(\mathsf T_{s,j})
 \le
 \frac1{2j+1}
 \sum_{J=0}^{2j}(2J+1)e^{-sJ(J+1)}.
\]
Comparison with the radial Gaussian integral gives
\[
 \sum_{J\ge0}(2J+1)e^{-sJ(J+1)}
 \le\frac Cs.
\]
Both \(\mathsf T_{s,j}\) and \(q_j^2I\) are contractions, while
\[
 q_j^2=e^{-2sj(j+1)}
 \le
 C\min\left\{1,\frac1{s(2j+1)}\right\}.
\]
The triangle inequality for \(\omega_\otimes\) proves
\eqref{eq:V86-binary-heat-capacity}.

For sharpness, use the exact weights \eqref{eq:V85-pjJ}.
On the heat range \(J\asymp s^{-1/2}\), the hypothesis \(sj\to\infty\)
gives \(J^2/j\to0\), so
\[
 p_{j,J}\asymp\frac{J}{j}.
\]
Summing \(p_{j,J}q_J\) over that range gives the lower bound
\(c/(sj)\); \eqref{eq:V86-binary-heat-capacity} gives the matching
upper bound.  Finally \(q_j^2=o((sj)^{-1})\), so subtracting the
separated endpoint does not alter the order.
\end{proof}

\begin{assessment}[What the new card adds to V75]
\label{ass:V86-binary-card-new}
The generic V75 covariance card
\(\min\{1,sH_{12}\}\) is representation-blind and becomes merely
order one in a hot equal-spin coordinate.  The output trace in
\Cref{thm:V86-binary-heat-capacity} uses the complete intermediate-spin
resolution and recovers the sharper factor \((sj)^{-1}\).  This is the
factor identified as missing in
\eqref{eq:V85-required-gain}; it is available when a mass-compatible
hot equal-spin bank occurs in the pure binary prefix.
\end{assessment}

\section{The signed once-paid joining profile}
\label{sec:V86-signed-paid}

For \(K\ne0\), put
\[
 W_K:=\frac{A_K}{1-q_K}.
\]
Before taking a physical absolute value, the joining-payer operator on
the V85 incidence-three coordinate is
\begin{equation}
 \mathsf R_{s,j}
 :=
 \sum_{\substack{J,K\\K\ne0}}
 W_K(q_K-q_Jq_1)P_{J,K}.
\label{eq:V86-raw-paid-joining}
\end{equation}
It is self-adjoint but not positive.

\begin{lemma}[Exact signed paid expectation]
\label{lem:V86-raw-paid-expectation}
On \(\Psi_j=\psi_j\otimes|1,0\rangle\),
\begin{align}
 \mathcal E_{s,j}^{\rm R}
 &:=
 \langle\Psi_j,\mathsf R_{s,j}\Psi_j\rangle
\notag\\
 &=
 \sum_{J=2}^{2j}p_{j,J}u_JW_{J-1}d_J^-
 -
 \sum_{J=1}^{2j}p_{j,J}v_JW_{J+1}d_J^+.
\label{eq:V86-exact-raw-paid}
\end{align}
\end{lemma}

\begin{proof}
Use the same Clebsch--Gordan weights as in
\Cref{lem:V85-exact-channel-expectations}, but retain the sign of the
\(K=J+1\) covariance coefficient.  The \(J=1,K=0\) term cannot carry
the payer and is absent.  This gives
\eqref{eq:V86-exact-raw-paid}.
\end{proof}

For \(J\ge2\), denote the conditional summand in brackets by
\begin{equation}
 \Delta_{s,J}^{\rm R}
 :=
 u_JW_{J-1}d_J^-
 -
 v_JW_{J+1}d_J^+.
\label{eq:V86-DeltaR}
\end{equation}

\begin{lemma}[Parabolic signed payer profile]
\label{lem:V86-paid-profile}
Let \(s\downarrow0\) and let
\(\sqrt s\,J\to y\in(0,\infty)\).  Then
\begin{equation}
 \Delta_{s,J}^{\rm R}
 \longrightarrow
 r(y^2),
\label{eq:V86-paid-profile-limit}
\end{equation}
where
\begin{equation}
 r(x)
 :=
 2x
 \frac{(x+2)-(2-x)e^x}{(e^x-1)^2},
 \qquad x>0.
\label{eq:V86-paid-profile}
\end{equation}
The profile is strictly positive.  Moreover, for \(J\ge2\),
\begin{equation}
 |\Delta_{s,J}^{\rm R}|
 \le
 C(1+sJ^2)^3e^{-csJ^2},
\label{eq:V86-paid-profile-dominant}
\end{equation}
and the exceptional \(J=1\) conditional summand in
\eqref{eq:V86-exact-raw-paid} is uniformly bounded.
\end{lemma}

\begin{proof}
Set
\[
 x_-:=sJ(J-1),\qquad
 x_+:=s(J+1)(J+2),
\qquad
 G_s(x):=\frac{x+s}{e^x-1}.
\]
The two paid magnitudes have the exact forms
\begin{align}
 W_{J-1}d_J^-
 &=
 \frac1sG_s(x_-)
 \bigl(1-e^{-2s(J+1)}\bigr),
\label{eq:V86-paid-minus-exact}\\
 W_{J+1}d_J^+
 &=
 \frac1sG_s(x_+)
 \bigl(e^{2sJ}-1\bigr).
\label{eq:V86-paid-plus-exact}
\end{align}
Write \(h=\sqrt s\) and \(y_s=hJ\), so that \(y_s\to y\).
Then
\[
 x_-=y_s^2-hy_s,
 \qquad
 x_+=y_s^2+3hy_s+2h^2.
\]
Taylor expansion of
\eqref{eq:V86-paid-minus-exact}--%
\eqref{eq:V86-paid-plus-exact}, including
\[
 u_J=\frac12-\frac{h}{4y_s}+O(h^2),
 \qquad
 v_J=\frac12+\frac{h}{4y_s}+O(h^2),
\]
shows that the terms of order \(h^{-1}\) cancel and the order-one
remainder is
\[
 -2y_s^2\bigl(F(y_s^2)+2F'(y_s^2)\bigr)+O(h),
 \qquad
 F(x):=\frac{x}{e^x-1}.
\]
Putting the limiting expression over the common denominator
\((e^x-1)^2\) gives \eqref{eq:V86-paid-profile}.
All remainders are uniform when \(y\) lies in a compact subset of
\((0,\infty)\).

Let
\[
 H(x):=(x+2)-(2-x)e^x.
\]
Then \(H(0)=0\) and
\[
 H'(x)=1-(1-x)e^x>0
\]
for \(x>0\): for \(0<x<1\), use
\(\log(1-x)+x<0\), while for \(x\ge1\) the second term is nonpositive.
Thus \(r(x)>0\).

For completeness, split \(sJ^2\le1\) and \(sJ^2>1\) in the exact
formulae.  In the first range, the power series for
\((e^z-1)^{-1}\) and the cancellation of the common linear term give
a uniform bound by a cubic polynomial in \(sJ^2\).  In the second,
the mean-value theorem applied to \(G_s\), together with
\[
 |G_s^{(a)}(z)|
 \le C_a(1+z)^{a+1}e^{-cz}
 \qquad(z\ge1/4,\ a=0,1,2),
\]
gives the same polynomial times \(e^{-csJ^2}\).  Enlarging the
constant covers the transition range and proves
\eqref{eq:V86-paid-profile-dominant}.  The \(J=1\) term follows
directly from
\[
 v_1W_2d_1^+
 =
 \frac23\frac{A_2}{1-q_2}(q_1^2-q_2),
\]
whose limit is finite.
\end{proof}

\begin{theorem}[Sharp signed joining-payer accumulation]
\label{thm:V86-signed-paid-accumulation}
Let \(s\downarrow0\) and choose half-integers \(j_s\) with
\[
 sj_s\longrightarrow\infty.
\]
Then
\begin{equation}
 j_ss\,\mathcal E_{s,j_s}^{\rm R}
 \longrightarrow
 c_{\rm R}
 :=
 \frac12\int_0^\infty r(x)\,\dd x
 \in(0,\infty).
\label{eq:V86-paid-Riemann-limit}
\end{equation}
In particular,
\begin{equation}
 \mathcal E_{s,j_s}^{\rm R}
 \asymp\frac1{j_ss}.
\label{eq:V86-signed-paid-scale}
\end{equation}
\end{theorem}

\begin{proof}
For \(J=y/\sqrt s\) with \(y\) in a fixed compact subset of
\((0,\infty)\), the product formula
\eqref{eq:V85-p-product} and \(sj_s\to\infty\) give
\[
 p_{j_s,J}
 =
 (1+o(1))\frac{J}{j_s}.
\]
Combine this with \Cref{lem:V86-paid-profile}.  The scaled sum becomes
\[
 j_ss\sum_{J\ge2}p_{j_s,J}\Delta_{s,J}^{\rm R}
 =
 s\sum_{J\ge2}J\,r(sJ^2)+o(1).
\]
The dominant bound
\eqref{eq:V86-paid-profile-dominant} and
\eqref{eq:V85-p-upper} justify passage to the Riemann integral:
\[
 s\sum_{J\ge2}J\,r(sJ^2)
 \longrightarrow
 \int_0^\infty y\,r(y^2)\,\dd y
 =
 \frac12\int_0^\infty r(x)\,\dd x.
\]
The \(J=1\) term contributes \(O(s)\) after this scaling.
The profile behaves as \(O(x^2)\) at zero and
\(O(x^2e^{-x})\) at infinity, so the integral is finite; its strict
positivity follows from \eqref{eq:V86-paid-profile}.
\end{proof}

\begin{assessment}[The payer cancellation is scale-critical]
\label{ass:V86-paid-cancellation}
The positive joining-payer expectation in V85 has order
\((js^{3/2})^{-1}\).  Keeping the \(K=J-1\) and \(K=J+1\) signs through
the payer yields \((js)^{-1}\), a gain of \(\sqrt s\).  This signed
gain is necessary for the complete minimal packet below; taking the
positive joining envelope first would still fail its paid target by
\(s^{-1/2}\).
\end{assessment}

\section{The exact three-coordinate \(2+2\) packet}
\label{sec:V86-minimal-packet}

We now restore the two proper-prefix covariances which were absent
from the isolated joining tests of V84--V85.  Use three independent
coordinates, ordered \(p<q<e\).  At \(p\), rows \(1,2\) carry
spin \(\frac12\) and rows \(3,4\) are trivial.  At \(q\), rows
\(3,4\) carry spin \(\frac12\) and rows \(1,2\) are trivial.  At
\(e\), rows \(1,2\) carry spins \(j,j\), row \(3\) carries spin
\(1\), and row \(4\) is trivial.  Thus the two pure pair banks are
already connected before \(e\), and the coordinate \(e\) first joins
the partition \(12\mid34\).  There is no nontrivial suffix.

Let
\[
 \chi
 :=
 \left|\tfrac12,\tfrac12\right\rangle
 \otimes
 \left|\tfrac12,-\tfrac12\right\rangle .
\]
On \(V_{1/2}\otimes V_{1/2}=V_0\oplus V_1\), define the connected
binary heat operator
\begin{equation}
 \mathsf B_s
 :=
 P_0+q_1P_1-q_{1/2}^{\,2}I.
\label{eq:V86-binary-prefix}
\end{equation}
Its expectation on \(\chi\) is
\begin{equation}
 b_s
 :=
 \langle\chi,\mathsf B_s\chi\rangle
 =
 \frac12(1+e^{-2s})-e^{-3s/2}
 =
 \frac s2+O(s^2).
\label{eq:V86-binary-prefix-expectation}
\end{equation}
In particular, \(b_s>0\) for all sufficiently small \(s>0\).

When the unique resolvent payer is assigned to this prefix, the
zero-output channel is omitted and the resulting signed operator is
\begin{equation}
 \mathsf B_s^{\rm pay}
 :=
 \frac{A_1}{1-q_1}
 \bigl(q_1-q_{1/2}^{\,2}\bigr)P_1.
\label{eq:V86-binary-prefix-paid}
\end{equation}
Since \(A_1=3\), its expectation has magnitude
\begin{equation}
 b_s^{\rm pay}
 :=
 \left|
 \langle\chi,\mathsf B_s^{\rm pay}\chi\rangle
 \right|
 =
 \frac32
 \frac{e^{-3s/2}-e^{-2s}}{1-e^{-2s}}
 =
 \frac38+O(s).
\label{eq:V86-binary-prefix-paid-expectation}
\end{equation}

Let
\begin{equation}
 \Omega_{s,j}
 :=
 \chi^{(p)}\otimes\chi^{(q)}\otimes\Psi_j^{(e)}.
\label{eq:V86-packet-vector}
\end{equation}
The superscripts indicate independent coordinates; after regrouping
the tensor factors, \(\Omega_{s,j}\) is a product vector across the
four input rows.  On these same coordinates define
\begin{align}
 \mathsf Q_{s,j}^{\rm U}
 &:=
 \mathsf B_s^{(p)}
 \otimes\mathsf B_s^{(q)}
 \otimes\mathsf J_{s,j}^{(e)},
\label{eq:V86-packet-unpaid}\\
 \mathsf Q_{s,j}^{\rm P,e}
 &:=
 \mathsf B_s^{(p)}
 \otimes\mathsf B_s^{(q)}
 \otimes\mathsf R_{s,j}^{(e)},
\label{eq:V86-packet-paid-e}\\
 \mathsf Q_{s,j}^{\rm P,p}
 &:=
 (\mathsf B_s^{\rm pay})^{(p)}
 \otimes\mathsf B_s^{(q)}
 \otimes\mathsf J_{s,j}^{(e)},
\label{eq:V86-packet-paid-p}\\
 \mathsf Q_{s,j}^{\rm P,q}
 &:=
 \mathsf B_s^{(p)}
 \otimes(\mathsf B_s^{\rm pay})^{(q)}
 \otimes\mathsf J_{s,j}^{(e)}.
\label{eq:V86-packet-paid-q}
\end{align}
These are not surrogate products.  They are respectively the raw
unpaid \(S_{22}\) summand and its three nonzero once-paid alternatives
on the present no-suffix packet.

\begin{theorem}[The complete minimal packet exactly restores the
missing powers]
\label{thm:V86-minimal-packet}
Let \(s\downarrow0\) and let \(j_s\) be half-integral with
\[
 j_s\ge16s^{-1},
 \qquad
 sj_s\longrightarrow\infty.
\]
Then
\begin{align}
 \langle\Omega_{s,j_s},
 \mathsf Q_{s,j_s}^{\rm N}\Omega_{s,j_s}\rangle
 &\asymp
 \frac{s^2}{j_s},
\label{eq:V86-packet-unpaid-scale}\\
 \left|
 \langle\Omega_{s,j_s},
 \mathsf Q_{s,j_s}^{\rm P,e}\Omega_{s,j_s}\rangle
 \right|
 &\asymp
 \frac{s}{j_s},
\label{eq:V86-packet-paid-e-scale}\\
 \left|
 \langle\Omega_{s,j_s},
 \mathsf Q_{s,j_s}^{\rm P,p}\Omega_{s,j_s}\rangle
 \right|
 +
 \left|
 \langle\Omega_{s,j_s},
 \mathsf Q_{s,j_s}^{\rm P,q}\Omega_{s,j_s}\rangle
 \right|
 &\asymp
 \frac{s}{j_s}.
\label{eq:V86-packet-paid-prefix-scale}
\end{align}
For \(s^{-3}\le j_s\le2s^{-3}\), the respective orders are
\[
 s^5,\qquad s^4,\qquad s^4.
\]
\end{theorem}

\begin{proof}
Independence of the three coordinate Haar variables and the tensor
form of \(\Omega_{s,j}\) give the exact identities
\begin{align}
 \langle\Omega_{s,j},
 \mathsf Q_{s,j}^{\rm U}\Omega_{s,j}\rangle
 &=
 b_s^2\mathcal E_{s,j}^{\rm S},
\label{eq:V86-packet-factor-U}\\
 \langle\Omega_{s,j},
 \mathsf Q_{s,j}^{\rm P,e}\Omega_{s,j}\rangle
 &=
 b_s^2\mathcal E_{s,j}^{\rm R},
\label{eq:V86-packet-factor-Pe}\\
 \left|
 \langle\Omega_{s,j},
 \mathsf Q_{s,j}^{\rm P,p}\Omega_{s,j}\rangle
 \right|
 &=
 b_s^{\rm pay}b_s\mathcal E_{s,j}^{\rm S},
\label{eq:V86-packet-factor-Pp}\\
 \left|
 \langle\Omega_{s,j},
 \mathsf Q_{s,j}^{\rm P,q}\Omega_{s,j}\rangle
 \right|
 &=
 b_sb_s^{\rm pay}\mathcal E_{s,j}^{\rm S}.
\label{eq:V86-packet-factor-Pq}
\end{align}
The quantity \(\mathcal E_{s,j}^{\rm S}\) is positive by
\eqref{eq:V85-Jensen-gap}.  Equations
\eqref{eq:V86-binary-prefix-expectation} and
\eqref{eq:V86-binary-prefix-paid-expectation},
\Cref{thm:V85-sharp-channel-accumulation}, and
\Cref{thm:V86-signed-paid-accumulation} now prove
\eqref{eq:V86-packet-unpaid-scale}--%
\eqref{eq:V86-packet-paid-prefix-scale}.  Substitution of
\(j_s\asymp s^{-3}\) gives the final orders.
\end{proof}

\begin{corollary}[The V84 adversarial sequence is neutralized, not
discarded]
\label{cor:V86-V84-neutralized}
On the packet of \Cref{thm:V86-minimal-packet}, the row stocks obey
\begin{equation}
 S_1=a_j+O(1),
 \qquad
 H_{12}=a_j^2+O(a_j),
 \qquad
 H_{1,C_1}=3a_j+O(1).
\label{eq:V86-packet-stocks}
\end{equation}
Consequently, if \(j_s\asymp s^{-3}\), then
\begin{align}
 \frac{H_{1,C_1}}{sS_1^2}
 &\asymp s^5,
\label{eq:V86-unpaid-target-saturated}\\
 \frac{H_{1,C_1}}{s^2S_1^2}
 &\asymp s^4.
\label{eq:V86-paid-target-saturated}
\end{align}
The unpaid test coefficient and the three nonzero once-paid test
coefficients in \Cref{thm:V86-minimal-packet} therefore have exactly
the orders required by the two crossing square-tail targets which the
isolated joining coordinate violated in V84.  Thus this particular
product-vector family is no longer a counterexample after its exact
prefix packet is restored; no capacity upper bound is inferred from
that fact.
\end{corollary}

\begin{proof}
Only the coordinate \(e\) contributes an unbounded spin.  Its
row-\(1\) and row-\(2\) masses are \(a_j\), its row-\(3\) mass is
\(3\), and the two spin-\(\frac12\) prefix coordinates contribute
only constants.  This gives \eqref{eq:V86-packet-stocks}.  Since
\(a_j\asymp j^2\), the two displayed targets follow.  Compare them
with
\eqref{eq:V86-packet-unpaid-scale}--%
\eqref{eq:V86-packet-paid-prefix-scale}.
\end{proof}

\begin{firewall}[What the minimal packet theorem proves]
\label{fire:V86-minimal-packet-scope}
\Cref{thm:V86-minimal-packet} proves that the concrete obstruction
constructed in V84--V85 was caused by deleting the two mandatory
proper-prefix connected factors.  It is a sharp test of one
three-coordinate family, not a uniform \(S_{22}\) estimate.  A
uniform proof still has to show, for every representation assignment
and every product input, that either a mass-compatible hot prefix
supplies its inverse heat-window factor or that both prefixes are in
their perturbative \(s\)-regime.  No such scalar dichotomy is assumed
here.
\end{firewall}

\section{A joint moment card for an identified overlap bank}
\label{sec:V86-overlap-card}

The V85 ascent argument reduced every overlap branch to one bank
chosen from a list of at most four, but deliberately left its two
source occurrences inside one carrier.  The following elementary
operator inequality records exactly what can be gained once those
occurrences are adjacent.

\begin{lemma}[Adjacent repeated-bank joint card]
\label{lem:V86-adjacent-joint-card}
Let \(A,B\) be bounded self-adjoint operators on the four-row input
space, and let \(\varphi\) be any product state across the rows.  Then
\begin{equation}
 |\varphi(AB)|
 \le
 \frac12\left(
 \|A\|\varphi(|A|)
 +
 \|B\|\varphi(|B|)
 \right).
\label{eq:V86-adjacent-joint-card}
\end{equation}
Consequently,
\begin{equation}
 \sup_{\varphi\ {\rm product}}|\varphi(AB)|
 \le
 \frac12\left(
 \|A\|\kappa_\otimes(|A|)
 +
 \|B\|\kappa_\otimes(|B|)
 \right).
\label{eq:V86-adjacent-joint-capacity}
\end{equation}
If \(A\) and \(B\) commute, then the stronger operator statement
\begin{equation}
 |AB|
 \preccurlyeq
 \frac12(A^2+B^2)
 \preccurlyeq
 \frac12\bigl(\|A\||A|+\|B\||B|\bigr)
\label{eq:V86-commuting-joint-order}
\end{equation}
holds, and hence \eqref{eq:V86-adjacent-joint-capacity} also controls
the positive product-state capacity of the physical absolute value
\(|AB|\).
\end{lemma}

\begin{proof}
The Cauchy--Schwarz inequality for the positive functional
\(\varphi\) gives
\[
 |\varphi(AB)|^2
 \le
 \varphi(A^2)\varphi(B^2).
\]
Functional calculus gives
\[
 A^2\preccurlyeq\|A\||A|,
 \qquad
 B^2\preccurlyeq\|B\||B|.
\]
Apply \(2\sqrt{xy}\le x+y\) with
\(x=\|A\|\varphi(|A|)\) and
\(y=\|B\|\varphi(|B|)\).  This proves
\eqref{eq:V86-adjacent-joint-card}; taking the supremum proves
\eqref{eq:V86-adjacent-joint-capacity}.

If \(A\) and \(B\) commute, their joint spectral calculus reduces the
first inequality in \eqref{eq:V86-commuting-joint-order} to
\(2|ab|\le a^2+b^2\).  The second is the preceding one-variable
functional-calculus inequality.
\end{proof}

\begin{corollary}[Conditional analytic closure of the V85 witness]
\label{cor:V86-overlap-conditional}
Let \(\mathcal B_a\) be the old bank furnished by
\Cref{thm:V85-overlap-witness}.  Suppose that, after the exact
first-connectivity packet is assembled, its earlier and current
occurrences become adjacent self-adjoint complete-bank operators
\(A_a,B_a\).  Then their joint product-state radius is charged
linearly, rather than by a product of two first-moment bounds, through
\eqref{eq:V86-adjacent-joint-capacity}.  If the occurrences are
commuting functions of the same physical block, the same statement
holds after physical absolute value.
\end{corollary}

\begin{proof}
The combinatorial selection and its loss of at most \(4/\delta\) are
already contained in
\Cref{thm:V85-overlap-witness,cor:V85-overlap-joint-gate}.  Apply
\Cref{lem:V86-adjacent-joint-card} to the two retained occurrences.
No second selection or independent use of the bank stock is made.
\end{proof}

\begin{firewall}[The actual overlap carrier has not yet been
identified with \(AB\)]
\label{fire:V86-overlap-card-scope}
The hypothesis of \Cref{cor:V86-overlap-conditional} is substantive.
In the full V61 source packet, the two occurrences can be separated
by joining, suffix, or replica-embedding factors.  Cyclicity of a
trace cannot be used inside an arbitrary product-state matrix
coefficient, and noncommuting interposed factors cannot be silently
removed.  Thus \Cref{lem:V86-adjacent-joint-card} closes the analytic
gate only after an exact algebraic reduction to an adjacent pair.
Otherwise the required object remains a noncommutative second moment.
\end{firewall}

\section{Post-V86 status and the next upstream attack}
\label{sec:V86-status}

\begin{theorem}[Exact status after restoring the minimal packet]
\label{thm:V86-status}
Relative to V1--V85, the present note proves:
\begin{enumerate}[label=\textup{(V86.\arabic*)},leftmargin=6.5em]
\item the deterministic \(S_{22}\) first-joining summand contains two
      proper-prefix connected banks before its joining covariance;
\item the equal-spin binary prefix has the sharp product-state heat
      capacity
      \[
      \omega_\otimes(\mathsf K_{s,j})
      \lesssim
      \min\left\{1,\frac1{s(2j+1)}\right\},
      \]
      and the inverse heat-window scale is attained by the
      anti-aligned input;
\item the signed joining-payer response on the V85
      incidence-three family has the exact parabolic profile
      \eqref{eq:V86-paid-profile} and accumulated order
      \((js)^{-1}\), rather than the positive-envelope order
      \((js^{3/2})^{-1}\);
\item on the exact three-coordinate \(2+2\) packet, the two prefix
      covariances supply the full missing powers: the unpaid response
      is \(s^2/j\), while paying either prefix or the joining
      coordinate gives \(s/j\);
\item the sequence \(j\asymp s^{-3}\), which disproved the isolated
      V84 crossing tails, saturates both desired tails after the
      complete prefix packet is restored;
\item an identified repeated overlap bank has the linear adjacent
      joint card \eqref{eq:V86-adjacent-joint-capacity}; commuting
      occurrences also admit the physical-absolute operator order
      \eqref{eq:V86-commuting-joint-order}.
\end{enumerate}
None of these statements proves the uniform \(S_{22}\) estimate, the
remaining source-partition estimates, Osterwalder--Schrader
reconstruction, or a positive continuum mass gap.
\end{theorem}

\begin{proof}
Items \textup{(V86.1)}--\textup{(V86.2)} are the packet reconstruction
and binary-prefix theorem of the first part of this chapter.
Item \textup{(V86.3)} is
\Cref{lem:V86-paid-profile,thm:V86-signed-paid-accumulation}.
Items \textup{(V86.4)}--\textup{(V86.5)} are
\Cref{thm:V86-minimal-packet,cor:V86-V84-neutralized}.
Item \textup{(V86.6)} is
\Cref{lem:V86-adjacent-joint-card}.  The final sentence lists gates
not addressed by those results.
\end{proof}

\begin{openproblem}[Uniform prefix--joining scalar dichotomy]
\label{open:V86-uniform-dichotomy}
For every same-side \(S_{22}\) representation assignment and every
product input, prove one of the following two alternatives with
constants independent of the support:
\begin{enumerate}[label=\textup{(\alph*)},leftmargin=4em]
\item a mass-compatible equal-spin or balanced proper-prefix bank is
      hot and supplies its sharp inverse factor
      \((sJ)^{-1}\) through a complete output-channel heat card;
\item both proper-prefix banks lie in their perturbative regimes and
      their signed connected expectations supply \(s^2\).
\end{enumerate}
The alternative must be expressed in the V79 row-cell stocks and
must survive each of the four general V2 payer locations.  It may not
infer product-state capacity from trace cancellation.
\end{openproblem}

\begin{programme}[Mandate for V87]
\label{prog:V86-next}
\begin{enumerate}[label=\textup{(\arabic*)},leftmargin=4em]
\item Start from the full V61 \(S_{22}\) packet, retain both binary
      prefixes through capacity, and formulate the dichotomy in
      \Cref{open:V86-uniform-dichotomy} as a scalar inequality in the
      V79 row masses and balanced cells.
\item Prove the hot-prefix branch with the sharp complete-channel
      heat capacity, and the cold-prefix branch with two first-order
      connected estimates; carry separately all four V2 payer
      positions, including a nontrivial suffix.
\item Test every proposed uniform inequality on both the V74
      deterministic nondecay family and the V84--V86
      \(j\asymp s^{-3}\) family before inserting it into the compiler.
\item In the overlap branch, write the two occurrences of the V85
      witness bank inside the exact V61 carrier.  Apply
      \Cref{lem:V86-adjacent-joint-card} only if algebra proves
      adjacency; otherwise retain and estimate the intervening
      noncommutative second moment.
\item Do not perform another companion audit before V90.  At V90,
      inspect the newest Standard-Model and Navier--Stokes notes
      before choosing the V91 direction.
\end{enumerate}
\end{programme}

\begin{assessment}[Current bottleneck]
\label{ass:V86-bottleneck}
The isolated joining-coordinate counterexample is no longer the
bottleneck: the exact minimal packet neutralizes it sharply.  The
remaining obstruction is upstream and uniform.  One must extract
from arbitrary row-cell data either the hot-prefix inverse
heat-window factor or two perturbative prefix factors, without
discarding the complete output signs and without charging an overlap
bank twice.
\end{assessment}

\paragraph{Parent-version record.}
This chapter was appended to
\texttt{Renewal\_Yang\_Mills\_Mass\_Gap\_Research\_Notes\_V85.tex},
whose final record was \(39380\) lines, \(1357382\) bytes, and SHA256
\[
 \texttt{2d60e74fe8f6a16e26a7527f1f4f86dbd09fdaeef3c00365836010cbc7d32f94}.
\]
No companion audit is due at V86; the next scheduled audit is V90.

% =============================================================================
% BEGIN APPENDED COMPACT VERSION V87 MATERIAL
% =============================================================================

\clearpage
\chapter{Three binary heat regimes and paired prefix cards}
\label{chap:V87}

\section{Context: separate the carrier gate from the scalar gate}
\label{sec:V87-context}

V86 showed on one coordinate that an equal-spin binary prefix can
have product-state heat capacity of order \((sj)^{-1}\), and that the
two low-spin prefix responses exactly neutralize the V84--V85 test
packet.  Two uniformity questions were still conflated:
\begin{enumerate}[label=\textup{(\roman*)}]
\item whether a complete binary bank admits compatible unpaid and
      once-paid heat cards; and
\item whether the resulting scalar cards pay the two repeated row
      masses in the V79 incidence ledger.
\end{enumerate}
This chapter treats the first question before returning to the
second.  It also corrects the proposed two-regime language in
\Cref{open:V86-uniform-dichotomy}: input-Casimir hotness
\(s\,j(j+1)\gg1\) does not by itself make the inverse heat-dimension
card small.  There is an unavoidable mesoscopic interval.

\section{Unequal-spin complete-channel heat cards}
\label{sec:V87-unequal-heat}

For half-integers \(j,k\ge\frac12\), set
\[
 c_j:=j(j+1),\qquad
 m:=\max\{j,k\},\qquad
 \ell:=|j-k|,\qquad
 d_m:=2m+1.
\]
On
\[
 V_j\otimes V_k
 =
 \bigoplus_{J=\ell}^{j+k}V_J,
\]
write \(P_J^{j,k}\) for the complete spin-\(J\) projection and put
\begin{align}
 \mathsf T_{s;j,k}
 &:=
 \sum_{J=\ell}^{j+k}e^{-sJ(J+1)}P_J^{j,k},
\label{eq:V87-two-spin-heat}\\
 r_{s;j,k}
 &:=
 e^{-s(c_j+c_k)},
\label{eq:V87-separated-scalar}\\
 \mathsf K_{s;j,k}
 &:=
 \mathsf T_{s;j,k}-r_{s;j,k}I.
\label{eq:V87-two-spin-covariance}
\end{align}

\begin{lemma}[Complete projection product bound]
\label{lem:V87-projection-product-bound}
For every product density matrix \(\rho\otimes\sigma\) on
\(V_j\otimes V_k\),
\begin{equation}
 \operatorname{Tr}
 \bigl(P_J^{j,k}(\rho\otimes\sigma)\bigr)
 \le
 \frac{2J+1}{d_m}.
\label{eq:V87-projection-product-bound}
\end{equation}
\end{lemma}

\begin{proof}
Diagonal \(SU(2)\)-invariance and Schur's lemma give
\[
 \operatorname{Tr}_{V_k}P_J^{j,k}
 =
 \frac{2J+1}{2j+1}I_{V_j},
 \qquad
 \operatorname{Tr}_{V_j}P_J^{j,k}
 =
 \frac{2J+1}{2k+1}I_{V_k}.
\]
Since a density matrix is bounded above by the identity,
\begin{align*}
 \operatorname{Tr}
 \bigl(P_J^{j,k}(\rho\otimes\sigma)\bigr)
 &\le
 \operatorname{Tr}
 \bigl(P_J^{j,k}(\rho\otimes I)\bigr)
 =
 \frac{2J+1}{2j+1},\\
 \operatorname{Tr}
 \bigl(P_J^{j,k}(\rho\otimes\sigma)\bigr)
 &\le
 \frac{2J+1}{2k+1}.
\end{align*}
Take the smaller right side.
\end{proof}

\begin{lemma}[Radial heat tails]
\label{lem:V87-radial-heat-tails}
There are universal \(c,C>0\) such that, for
\(0<s\le s_0\) and every integer or half-integer \(L\ge0\),
\begin{align}
 \sum_{J\ge L}(2J+1)e^{-sJ(J+1)}
 &\le
 \frac Cs e^{-csL(L+1)},
\label{eq:V87-radial-tail-zero}\\
 \sum_{J\ge\max\{L,1\}}
 (2J+1)
 \frac{(1+J(J+1))e^{-sJ(J+1)}}
 {1-e^{-sJ(J+1)}}
 &\le
 \frac C{s^2}e^{-csL(L+1)}.
\label{eq:V87-radial-tail-paid}
\end{align}
\end{lemma}

\begin{proof}
The derivative of \(x(x+1)\) is \(2x+1\).  Comparing the sums with
the corresponding integrals on consecutive unit intervals gives
\[
 \int_L^\infty(2x+1)e^{-csx(x+1)}\,\dd x
 =
 \frac1{cs}e^{-csL(L+1)}.
\]
Changing \(c\) absorbs the first interval and proves
\eqref{eq:V87-radial-tail-zero}.

For \(J\ge1\), the elementary gap bound gives
\[
 \frac{(1+J(J+1))e^{-sJ(J+1)}}
 {1-e^{-sJ(J+1)}}
 \le
 \frac Cs
 \bigl(1+sJ(J+1)\bigr)e^{-csJ(J+1)}.
\]
The same substitution, now with one polynomial factor, gives
\(Cs^{-2}(1+sL(L+1))e^{-csL(L+1)}\).  Absorb that polynomial by
decreasing \(c\).
\end{proof}

\begin{theorem}[Physical-absolute unequal-spin heat card]
\label{thm:V87-unequal-unpaid-card}
Uniformly for \(0<s\le s_0\),
\begin{equation}
 \boxed{
 \kappa_\otimes\bigl(|\mathsf K_{s;j,k}|\bigr)
 \le
 C\min\left\{
 1,\
 \frac{e^{-cs\ell(\ell+1)}}{s(2m+1)}
 +
 e^{-s(c_j+c_k)}
 \right\}.}
\label{eq:V87-unequal-unpaid-card}
\end{equation}
In particular, if \(sm\ge1\), then
\begin{equation}
 \kappa_\otimes\bigl(|\mathsf K_{s;j,k}|\bigr)
 \le\frac C{sm}.
\label{eq:V87-superhot-unpaid}
\end{equation}
\end{theorem}

\begin{proof}
Because \(\mathsf T_{s;j,k}\) is positive,
\[
 |\mathsf T_{s;j,k}-r_{s;j,k}I|
 \preccurlyeq
 \mathsf T_{s;j,k}+r_{s;j,k}I.
\]
Apply \Cref{lem:V87-projection-product-bound} and then
\eqref{eq:V87-radial-tail-zero}.  The constant bound follows from
\(\|\mathsf T_{s;j,k}\|\le1\) and \(r_{s;j,k}\le1\).
If \(sm\ge1\), discard the mismatch exponential and note that
\[
 e^{-s(c_j+c_k)}
 \le e^{-csm^2}
 \le\frac C{sm}.
\]
\end{proof}

For nonzero \(J\), let
\[
 W_J:=\frac{1+J(J+1)}{1-e^{-sJ(J+1)}}.
\]
The one-coordinate positive once-paid covariance envelope is
\begin{equation}
 \mathsf P_{s;j,k}
 :=
 \sum_{\substack{\ell\le J\le j+k\\J>0}}
 W_J
 \left|e^{-sJ(J+1)}-r_{s;j,k}\right|
 P_J^{j,k}.
\label{eq:V87-local-paid-envelope}
\end{equation}

\begin{theorem}[Superhot paid heat card]
\label{thm:V87-superhot-paid-card}
If \(sm\ge1\), then
\begin{equation}
 \boxed{
 \kappa_\otimes(\mathsf P_{s;j,k})
 \le\frac C{s^2m}.}
\label{eq:V87-superhot-paid-card}
\end{equation}
\end{theorem}

\begin{proof}
Use
\[
 W_J|q_J-r_{s;j,k}|
 \le W_Jq_J+r_{s;j,k}W_J,
 \qquad q_J=e^{-sJ(J+1)}.
\]
By
\Cref{lem:V87-projection-product-bound,lem:V87-radial-heat-tails},
the product-state capacity of the first sum is at most
\(C/(s^2d_m)\).

Also
\[
 W_J\le C\bigl(s^{-1}+1+J(J+1)\bigr).
\]
Since \(J\le j+k\le2m\), the second sum has capacity at most
\[
 \frac{Cr_{s;j,k}}{d_m}
 \left(\frac{(m+1)^2}{s}+(m+1)^4\right).
\]
Put \(x=sm^2\).  Under \(sm\ge1\), one has \(m\asymp m+1\) and
\[
 x e^{-cx}+x^2e^{-cx}\le C.
\]
After division by \(d_m\asymp m\), the last display is therefore at
most \(C/(s^2m)\).
\end{proof}

\section{The mesoscopic band is a genuine third regime}
\label{sec:V87-mesoscopic}

\begin{theorem}[Input-hot but unsuppressed binary covariance]
\label{thm:V87-mesoscopic-countertest}
Fix \(\alpha\in(\frac12,1)\), and choose half-integers
\(j_s\asymp s^{-\alpha}\).  Then
\begin{align}
 s\,j_s(j_s+1)&\longrightarrow\infty,
\label{eq:V87-mesoscopic-input-hot}\\
 sj_s&\longrightarrow0,
\label{eq:V87-mesoscopic-output-cold}\\
 \langle\psi_{j_s},
 \mathsf K_{s;j_s,j_s}\psi_{j_s}\rangle
 &\longrightarrow1.
\label{eq:V87-mesoscopic-response-one}
\end{align}
Thus no binary prefix card which tends to zero merely from
\(s\,j(j+1)\to\infty\) can hold.
\end{theorem}

\begin{proof}
The first two limits follow from
\(j_s\asymp s^{-\alpha}\).  By
\Cref{lem:V74-mean-Casimir} and convexity,
\[
 \langle\psi_{j_s},
 \mathsf T_{s;j_s,j_s}\psi_{j_s}\rangle
 \ge e^{-2sj_s}\longrightarrow1.
\]
The separated scalar obeys
\[
 e^{-2sj_s(j_s+1)}\longrightarrow0.
\]
Since the heat operator is a contraction, these two facts squeeze
the covariance expectation to one.
\end{proof}

\begin{firewall}[Correction of the V86 dichotomy]
\label{fire:V87-three-regimes}
The heat parameter relevant to input Casimir is \(sj^2\), whereas the
anti-aligned output window is governed by \(sj\).  Hence one must
distinguish:
\[
 sj^2\lesssim1
 \quad\hbox{(input-cold)},\qquad
 sj^2\gg1\ \hbox{and}\ sj\ll1
 \quad\hbox{(mesoscopic)},\qquad
 sj\gtrsim1
 \quad\hbox{(superhot)}.
\]
\Cref{thm:V87-mesoscopic-countertest} shows that the middle regime
cannot be assigned a vanishing heat card.  This does not contradict
\Cref{thm:V86-binary-heat-capacity}: its sharp expression is capped
by one there.  The scalar compiler must treat the mesoscopic band
separately.
\end{firewall}

\section{A superhot coordinate annihilates a complete binary bank}
\label{sec:V87-bank-annihilator}

Let \(D\) be a nonempty pair-connecting bank in one binary prefix
block \(uv\), after its one-row factors have been removed as in V80.
At \(e\in D\), rows \(u,v\) carry spins \(j_e,k_e\).  Other physical
rows may be active at \(e\); the prefix partition places their block
moments in independent replica factors and does not alter the
\(uv\)-block heat operator below.  Put
\[
 m_e:=\max\{j_e,k_e\},
 \qquad
 \mathsf T_e:=\mathsf T_{s;j_e,k_e},
 \qquad
 r_e:=r_{s;j_e,k_e}.
\]
Coordinate independence gives
\begin{align}
 \mathsf T_D&:=\bigotimes_{e\in D}\mathsf T_e,
\label{eq:V87-bank-joint-moment}\\
 r_D&:=\prod_{e\in D}r_e,
\label{eq:V87-bank-separated-moment}\\
 \mathsf K_D&:=\mathsf T_D-r_DI.
\label{eq:V87-bank-covariance}
\end{align}
Let \(\mathcal N_D:=|\mathsf K_D|\), and let
\(\mathcal P_D\) denote the exact V2 positive once-paid envelope of
this complete binary covariance.

\begin{lemma}[One-coordinate positive restriction]
\label{lem:V87-one-coordinate-restriction}
For every \(e\in D\),
\begin{equation}
 \kappa_\otimes(\mathsf T_D)
 \le
 \kappa_\otimes(\mathsf T_e).
\label{eq:V87-one-coordinate-restriction}
\end{equation}
This remains true although each row input may be entangled among its
own coordinate factors.
\end{lemma}

\begin{proof}
Write
\(\mathsf T_D=\mathsf T_e\otimes\mathsf T_{D\setminus\{e\}}\).
Every heat moment is a positive contraction, so
\[
 0\preccurlyeq\mathsf T_D
 \preccurlyeq\mathsf T_e\otimes I.
\]
On a product state across rows \(u,v\), tracing out all other
coordinate factors leaves a product of two mixed states at \(e\).
The supremum over product mixed states equals the supremum over
product pure states by separate convexity.  This proves the display.
\end{proof}

\begin{theorem}[Unpaid superhot bank witness]
\label{thm:V87-unpaid-bank-witness}
If \(e\in D\) and \(sm_e\ge1\), then
\begin{equation}
 \boxed{
 \kappa_\otimes(\mathcal N_D)
 \le\frac C{sm_e}.}
\label{eq:V87-unpaid-bank-witness}
\end{equation}
\end{theorem}

\begin{proof}
Since \(\mathsf T_D\) is positive,
\[
 |\mathsf T_D-r_DI|
 \preccurlyeq
 \mathsf T_D+r_DI.
\]
Use
\Cref{lem:V87-one-coordinate-restriction} and
\(r_D\le r_e\), followed by
\Cref{thm:V87-unequal-unpaid-card}.  Both resulting terms are
bounded by \(C/(sm_e)\).
\end{proof}

\begin{theorem}[Paid superhot bank witness]
\label{thm:V87-paid-bank-witness}
If \(e\in D\) and \(sm_e\ge1\), then
\begin{equation}
 \boxed{
 \kappa_\otimes(\mathcal P_D)
 \le\frac C{s^2m_e}.}
\label{eq:V87-paid-bank-witness}
\end{equation}
\end{theorem}

\begin{proof}
Put \(R=D\setminus\{e\}\).  At the signed level,
\begin{equation}
 \mathsf K_D
 =
 \mathsf K_e\otimes\mathsf T_R
 +
 r_eI\otimes\mathsf K_R.
\label{eq:V87-selected-coordinate-extraction}
\end{equation}
This is the one-coordinate instance of the exact selected-subbank
identity \eqref{eq:V82-subbank-extraction}.  The V2 allocation of the
sole payer gives four positive terms.

If \(\mathsf K_e\) pays, its capacity is
\(C/(s^2m_e)\) by
\Cref{thm:V87-superhot-paid-card}, and \(\mathsf T_R\) is a
contraction.  If \(\mathsf K_e\) is unpaid and the complete moment
\(\mathsf T_R\) pays, use
\eqref{eq:V87-superhot-unpaid} and the support-uniform complete-product
payer bound \(C/s\) from
\eqref{eq:V81-product-paid-zero}; the product is again
\(C/(s^2m_e)\).

For the second summand of
\eqref{eq:V87-selected-coordinate-extraction}, suppose first that
the separated scalar at \(e\) pays.  Fusion bounds its payer mass by
\(C(1+c_{j_e}+c_{k_e})\), while its Wilson multiplier is \(r_e\).
Because \(sm_e\ge1\),
\[
 \frac{(1+c_{j_e}+c_{k_e})r_e}{1-r_e}
 \le
 C m_e^2e^{-csm_e^2}
 \le
 \frac C{s^2m_e}.
\]
The unpaid remainder covariance has capacity at most a constant.
Finally, if the separated scalar is unpaid and \(\mathsf K_R\) pays,
use the universal paid covariance cap \(C/s\) and
\[
 \frac{r_e}{s}
 \le
 \frac C{s^2m_e}.
\]
Positive order permits the one-coordinate capacity bounds to be
multiplied by the ordinary norms of the remaining positive factors.
The four terms prove the theorem.  Empty-remainder terms are simply
absent.
\end{proof}

Define the paired binary card
\begin{equation}
 \mathfrak b_s(D)
 :=
 \min\left\{
 1,\ sH_{uv}(D),\
 \inf_{\substack{e\in D\\sm_e\ge1}}\frac1{sm_e}
 \right\},
\label{eq:V87-paired-binary-card}
\end{equation}
where the infimum over the empty set is \(+\infty\).

\begin{corollary}[Compatible unpaid and paid bank cards]
\label{cor:V87-paired-bank-cards}
Every complete pair-connecting binary bank satisfies
\begin{align}
 \boxed{
 \kappa_\otimes(\mathcal N_D)
 \le C\mathfrak b_s(D),}
\label{eq:V87-paired-unpaid}\\
 \boxed{
 \kappa_\otimes(\mathcal P_D)
 \le \frac C{s}\mathfrak b_s(D).}
\label{eq:V87-paired-paid}
\end{align}
\end{corollary}

\begin{proof}
The generic complete-bank response/contraction bounds are
\[
 \kappa_\otimes(\mathcal N_D)
 \le C\min\{sH_{uv}(D),1\}
\]
and, by \eqref{eq:V78-paid-binary-minimum},
\[
 \kappa_\otimes(\mathcal P_D)
 \le
 C\min\{H_{uv}(D),s^{-1}\}
 =
 \frac Cs\min\{sH_{uv}(D),1\}.
\]
For every superhot witness \(e\), use
\Cref{thm:V87-unpaid-bank-witness,thm:V87-paid-bank-witness}.
Take the infimum of all valid bounds.
\end{proof}

\begin{corollary}[Completion by factored singleton prefix banks]
\label{cor:V87-prefix-spectator-completion}
Suppose the complete binary block prefix has the exact V80 form
\begin{equation}
 K_{uv}^{\rm pre}
 =
 q_G\,\mathsf K_D,
\label{eq:V87-prefix-spectator-form}
\end{equation}
where \(D\) is its pair-connecting bank and \(q_G\) is the product of
all factored one-row prefix moments.  Then its unpaid envelope is
bounded by \(C\mathfrak b_s(D)\), and the sum of its two internal
prefix-payer locations---payer in \(D\) or payer in \(G\)---is
bounded by \(C\mathfrak b_s(D)/s\).
\end{corollary}

\begin{proof}
Equation \eqref{eq:V87-prefix-spectator-form} is
\Cref{thm:V80-foreign-factor}, applied to every factored singleton
row bank.  Unpaid, \(q_G\le1\).  If the pair-connecting covariance
pays, use \eqref{eq:V87-paired-paid}.  If the complete singleton
product pays, use \eqref{eq:V87-paired-unpaid} and the
support-uniform \(C/s\) product-payer bound
\eqref{eq:V81-product-paid-zero}.  The V2 allocation uses these two
locations once, and their finite sum has the asserted bound.
\end{proof}

\begin{firewall}[The witness is coordinate-resolved]
\label{fire:V87-witness-coordinate-resolved}
The third entry of \eqref{eq:V87-paired-binary-card} is not a
function of the aggregate stock \(H_{uv}(D)\).  It records an actual
coordinate and its two input representations.  Collapsing the bank
to a single total Casimir before selecting the witness would erase
the intermediate-spin dimension in
\Cref{lem:V87-projection-product-bound}.  The witness can be tagged
inside a V79 incidence cell, but cannot be reconstructed from the
cell's scalar sum alone.
\end{firewall}

\section{Two disjoint prefix capacities may be multiplied}
\label{sec:V87-disjoint-prefixes}

\begin{lemma}[Disjoint-row positive tensorization]
\label{lem:V87-disjoint-row-tensorization}
Let \(A_{12}\ge0\) act on designated coordinate factors of rows
\(1,2\), let \(B_{34}\ge0\) act on designated disjoint coordinate
factors of rows \(3,4\), and let \(R\ge0\) act on all remaining
factors.  If \(\|R\|\le L\), then
\begin{equation}
 \boxed{
 \kappa_\otimes^{\,1|2|3|4}
 (A_{12}\otimes B_{34}\otimes R)
 \le
 L\,
 \kappa_\otimes^{\,1|2}(A_{12})
 \kappa_\otimes^{\,3|4}(B_{34}).}
\label{eq:V87-disjoint-row-tensorization}
\end{equation}
Arbitrary entanglement among coordinate factors within each single
row is allowed.
\end{lemma}

\begin{proof}
Positive order gives
\[
 A_{12}\otimes B_{34}\otimes R
 \preccurlyeq
 L\,A_{12}\otimes B_{34}\otimes I.
\]
For a four-row product state, tracing out the remaining coordinates
leaves
\[
 (\rho_1\otimes\rho_2)\otimes(\rho_3\otimes\rho_4)
\]
on the two prefix systems.  The expectation therefore factors into
an \(A_{12}\) expectation and a \(B_{34}\) expectation.  Separate
convexity reduces each mixed product state to the corresponding pure
product-state capacity.
\end{proof}

\begin{remark}[Why ordinary norm is not required on both prefixes]
\label{rem:V87-no-two-prefix-norm}
For equal superhot spins, the binary covariance has ordinary norm of
order one because of its almost-invariant output, while its physical
product-state capacity is \(O((sj)^{-1})\).  The positivity and
disjoint row sets in
\Cref{lem:V87-disjoint-row-tensorization} allow both prefix capacities
to be retained simultaneously.  This is not a general
multiplicativity assertion for overlapping banks.
\end{remark}

\section{Insertion into the complete \(S_{22}\) payer ledger}
\label{sec:V87-S22-insertion}

Fix a first-joining position \(r\), crossing endpoints
\(i\in\{1,2\}\), \(j\in\{3,4\}\), and the corresponding scalar local
crossing stock \(h_{ij,r}\).  After the exact singleton
factorizations of \Cref{thm:V80-foreign-factor}, let
\(D_{12}^{<r}\) and \(D_{34}^{<r}\) be the two pair-connecting
binary prefix banks, and abbreviate
\[
 b_{12}^{<r}:=\mathfrak b_s(D_{12}^{<r}),
 \qquad
 b_{34}^{<r}:=\mathfrak b_s(D_{34}^{<r}).
\]

\begin{theorem}[One paired formula covers all four payer locations]
\label{thm:V87-four-payer-packet}
The positive unpaid envelope of the fixed \(S_{22}\) packet satisfies
\begin{equation}
 \kappa_\otimes(\mathcal Q_{r,ij}^{\rm unpay})
 \le
 Cs\,b_{12}^{<r}b_{34}^{<r}h_{ij,r}.
\label{eq:V87-unpaid-packet-card}
\end{equation}
After the exact V2 allocation of the sole payer among the
\(12\)-prefix, the \(34\)-prefix, the joining covariance, and the
complete suffix, the sum of the four positive envelopes satisfies
\begin{equation}
 \boxed{
 \kappa_\otimes(\mathcal Q_{r,ij}^{\rm pay})
 \le
 C\,b_{12}^{<r}b_{34}^{<r}h_{ij,r}.}
\label{eq:V87-four-payer-packet}
\end{equation}
In the suffix case, the left side includes the complete sum over
possible suffix-payer coordinates.
\end{theorem}

\begin{proof}
Use
\Cref{cor:V87-paired-bank-cards,
cor:V87-prefix-spectator-completion} on the two complete prefix
envelopes and
\Cref{lem:V87-disjoint-row-tensorization} to retain both capacities.
The unpaid local joining covariance has ordinary norm at most
\(Csh_{ij,r}\), while its paid envelope has norm at most
\(Ch_{ij,r}\), by
\eqref{eq:V65-local-unpaid-crossing}--%
\eqref{eq:V65-local-paid-crossing}.  Every nonpayer suffix factor is
a contraction.

If the joining covariance pays, the product is
\[
 Cb_{12}^{<r}b_{34}^{<r}h_{ij,r}.
\]
If the \(12\)-prefix pays, the factors are
\[
 \frac C{s}b_{12}^{<r},\qquad
 Cb_{34}^{<r},\qquad
 Csh_{ij,r};
\]
the powers of \(s\) cancel.  The \(34\)-prefix case is symmetric.
If a suffix coordinate pays, both prefixes are unpaid, the joining
factor contributes \(Csh_{ij,r}\), and the support-uniform complete
suffix-payer sum contributes \(C/s\).  These are exactly the four
locations in \Cref{lem:V65-four-payers}.  With no payer, retain the
joining factor \(Csh_{ij,r}\), proving
\eqref{eq:V87-unpaid-packet-card}.
\end{proof}

\begin{firewall}[Carrier closure is not yet scalar closure]
\label{fire:V87-carrier-not-scalar}
Equation \eqref{eq:V87-four-payer-packet} is a support-uniform
operator estimate for the exact payer ledger.  To obtain a normalized
tree bound, the two dimensionless prefix cards still have to pay the
two repeated endpoint masses.  V86's proposed dichotomy is therefore
replaced by a precise scalar-certificate question for
\(b_{12}^{<r}\) and \(b_{34}^{<r}\).
\end{firewall}

\section{Ordinary and heat-witness endpoint certificates}
\label{sec:V87-endpoint-certificates}

For a binary prefix bank \(D_{uv}\), set
\[
 H_D:=H_{uv}(D_{uv}),
 \qquad
 \gamma_s(D_{uv})
 :=
 \begin{cases}
 \mathfrak b_s(D_{uv})/H_D,&H_D>0,\\
 0,&H_D=0.
 \end{cases}
\]

\begin{definition}[Certified prefix endpoint]
\label{def:V87-certified-endpoint}
Fix \(0<\eta\le1\).  A row \(u\) is ordinarily certified by
\(D_{uv}\) if
\begin{equation}
 H_D\ge\eta\Xi_u.
\label{eq:V87-ordinary-certificate}
\end{equation}
It is heat-witness certified if some \(e\in D_{uv}\) satisfies
\begin{equation}
 sm_e\ge1,
 \qquad
 sm_eH_D\ge\eta\Xi_u.
\label{eq:V87-heat-witness-certificate}
\end{equation}
\end{definition}

\begin{lemma}[Either certificate pays one endpoint]
\label{lem:V87-certificate-pays-endpoint}
If row \(u\) has either certificate in
\Cref{def:V87-certified-endpoint}, then
\begin{equation}
 \boxed{
 \gamma_s(D_{uv})
 \le\frac{C_\eta}{\Xi_u}.}
\label{eq:V87-gamma-pays-endpoint}
\end{equation}
The V79 own-prefix condition \(A_u\ge\eta\Xi_u\) implies the ordinary
certificate after changing \(\eta\) by the fixed nontrivial mass
\(a_*\).
\end{lemma}

\begin{proof}
Under \eqref{eq:V87-ordinary-certificate}, use
\(\mathfrak b_s(D_{uv})\le1\):
\[
 \gamma_s(D_{uv})
 \le\frac1{H_D}
 \le\frac1{\eta\Xi_u}.
\]
Under \eqref{eq:V87-heat-witness-certificate}, use the witness entry
of \eqref{eq:V87-paired-binary-card}:
\[
 \gamma_s(D_{uv})
 \le
 \frac1{sm_eH_D}
 \le\frac1{\eta\Xi_u}.
\]
Finally, every own-prefix incidence has both rows nontrivial, so
\[
 H_D
 =
 \sum_{e\in D_{uv}}a_{u,e}a_{v,e}
 \ge
 a_*\sum_{e\in D_{uv}}a_{u,e}
 =
 a_*A_u.
\]
\end{proof}

\begin{theorem}[Certified all-payer \(S_{22}\) tree sector]
\label{thm:V87-certified-S22-sector}
Restrict the crossing-\((i,j)\) \(S_{22}\) source to first-joining
summands for which \(D_{12}^{<r}\) certifies endpoint \(i\) and
\(D_{34}^{<r}\) certifies endpoint \(j\), each either ordinarily or
by a superhot witness with fixed parameter \(\eta\).  Then, after
summing all four payer locations and all such \(r\),
\begin{equation}
 \boxed{
 \frac{
 \kappa_\otimes(
 \mathcal S_{22,ij}^{\rm cert,pay})}
 {\Xi_1\Xi_2\Xi_3\Xi_4}
 \le
 C_\eta\,
 \theta_{12}\theta_{34}\theta_{ij}.}
\label{eq:V87-certified-S22-sector}
\end{equation}
\end{theorem}

\begin{proof}
For a nonzero prefix stock, write
\[
 b_{uv}^{<r}
 =
 \gamma_s(D_{uv}^{<r})H_{uv}^{<r}.
\]
By \Cref{lem:V87-certificate-pays-endpoint},
\[
 \gamma_s(D_{12}^{<r})\le\frac{C_\eta}{\Xi_i},
 \qquad
 \gamma_s(D_{34}^{<r})\le\frac{C_\eta}{\Xi_j}.
\]
Insert these bounds into
\eqref{eq:V87-four-payer-packet}, sum the restricted positions, and
use the ordered first-stock inequality
\[
 \sum_r
 H_{12}^{<r}H_{34}^{<r}h_{ij,r}
 \le
 H_{12}H_{34}H_{ij}.
\]
Thus the unnormalized capacity is at most
\[
 C_\eta
 \frac{H_{12}H_{34}H_{ij}}{\Xi_i\Xi_j}.
\]
Division by the four row masses is exactly the product of the three
tree weights in \eqref{eq:V87-certified-S22-sector}.
\end{proof}

\begin{assessment}[Strict enlargement of the V78 sector]
\label{ass:V87-enlarges-V78}
The ordinary certificate recovers the double-prefix-dominant sector
of \Cref{thm:V78-all-payers-prefix-dominant}.  The heat-witness
certificate is new: it may pay an endpoint even when the V79
own-prefix row stock is a small fraction of that row's total mass,
provided one actual prefix coordinate is superhot and
\(sm_eH_D\) reaches the missing row mass.  The same certificate works
for a prefix payer because the unpaid and paid bank cards were paired
before the scalar comparison.
\end{assessment}

\section{Regression tests and remaining scalar state}
\label{sec:V87-regression}

\begin{proposition}[The V74 and V86 tests select different branches]
\label{prop:V87-regression-tests}
\begin{enumerate}[label=\textup{(\roman*)}]
\item For the V74 bank with \(N\) equal coordinates and
      \(j_s\asymp s^{-1}\), the paired card is of order one, but both
      endpoints are ordinarily certified because the bank contains
      their full row masses.
\item For an equal-spin prefix coordinate with \(j_s\asymp s^{-3}\),
      the superhot witness gives
      \[
      \mathfrak b_s(D)\lesssim s^2,
      \qquad
      s^{-1}\mathfrak b_s(D)\lesssim s,
      \]
      the unpaid and paid scales required in the V85 heat-window
      comparison.
\item For the exact V86 three-coordinate packet, the two
      spin-\(\frac12\) prefixes use the response entry
      \(\mathfrak b_s(D)\lesssim s\); their product is \(s^2\), in
      agreement with
      \Cref{thm:V86-minimal-packet}.
\end{enumerate}
\end{proposition}

\begin{proof}
In \textup{(i)}, \(H_{uv}=Na_s^2\) and
\(\Xi_u=\Xi_v=Na_s\), so
\(H_{uv}\ge a_*\Xi_u,a_*\Xi_v\).
The witness value \(1/(sj_s)\) is only order one at
\(j_s\asymp s^{-1}\).
In \textup{(ii)}, \(1/(sj_s)\asymp s^2\), and
\Cref{cor:V87-paired-bank-cards} supplies the paid value after one
factor \(s^{-1}\).
In \textup{(iii)}, the spin-\(\frac12\) pair stock is a fixed
constant, so the response entry \(sH\) is \(O(s)\).
\end{proof}

\begin{firewall}[Why the minimal packet is still only a regression
test]
\label{fire:V87-minimal-still-test}
The low-spin prefixes in item \textup{(iii)} need not certify the
large-spin joining rows through
\Cref{def:V87-certified-endpoint}.  Their \(s^2\) product neutralizes
the particular coordinate-factorized matrix coefficient computed in
V86, but \Cref{thm:V87-certified-S22-sector} does not declare the
corresponding arbitrary-row-state capacity closed.  That remaining
case belongs to the scalar mesoscopic/response ledger, not to the
carrier theorem.
\end{firewall}

\begin{theorem}[Exact status after V87]
\label{thm:V87-status}
V87 establishes:
\begin{enumerate}[label=\textup{(V87.\arabic*)},leftmargin=6.5em]
\item unequal-spin physical-absolute unpaid and paid one-coordinate
      heat cards;
\item a mesoscopic family which is input-Casimir hot but has binary
      covariance expectation tending to one;
\item a superhot witness card for an arbitrary complete binary bank,
      paired compatibly across unpaid and once-paid envelopes;
\item positive tensorization of the two disjoint binary-prefix
      capacities despite arbitrary within-row coordinate
      entanglement;
\item the all-four-payer packet reduction
      \eqref{eq:V87-four-payer-packet}; and
\item a normalized spanning-tree bound whenever the two repeated
      crossing endpoints receive ordinary or superhot-witness prefix
      certificates.
\end{enumerate}
The remaining \(S_{22}\) terms are precisely those in which at least
one repeated endpoint fails the ordinary inequality
\eqref{eq:V87-ordinary-certificate} and every superhot coordinate in
its binary prefix fails
\eqref{eq:V87-heat-witness-certificate}.  This is a scalar
row-incidence condition; it is not a failure of the paired carrier
bound.
\end{theorem}

\begin{proof}
Items \textup{(V87.1)}--\textup{(V87.2)} are
\Cref{thm:V87-unequal-unpaid-card,
thm:V87-superhot-paid-card,
thm:V87-mesoscopic-countertest}.
Item \textup{(V87.3)} is
\Cref{cor:V87-paired-bank-cards}; item \textup{(V87.4)} is
\Cref{lem:V87-disjoint-row-tensorization}; and items
\textup{(V87.5)}--\textup{(V87.6)} are
\Cref{thm:V87-four-payer-packet,
thm:V87-certified-S22-sector}.  Negating the two alternatives in
\Cref{def:V87-certified-endpoint} gives the final description.
\end{proof}

\begin{programme}[Mandate for V88]
\label{prog:V87-next}
\begin{enumerate}[label=\textup{(\arabic*)},leftmargin=4em]
\item Refine each V79 own-prefix incidence cell into input-cold,
      mesoscopic, and superhot coordinate-resolved subcells without
      duplicating row mass.
\item On an endpoint which fails both V87 certificates, retain the
      inequalities
      \(H_D<\eta\Xi_u\) and
      \(sm_eH_D<\eta\Xi_u\) for every superhot witness; combine them
      with the exact six-cell identity
      \eqref{eq:V79-incidence-identity}.
\item Determine whether the mesoscopic subcell is paid by ordinary
      contraction after row normalization.  Test first on
      \(j_s\asymp s^{-\alpha}\),
      \(\frac12<\alpha<1\), and on diffuse mixtures of such
      coordinates.  If the scalar claim fails, preserve the smallest
      incidence-table counterexample.
\item Treat the response-only low-spin prefixes of the V86 minimal
      packet at arbitrary row-product capacity, rather than only on
      its coordinate-factorized test vector.
\item Write the V85 repeated-bank witness inside the exact V61
      operator carrier.  Use the V86 adjacent joint card only after
      proving adjacency or commutation.
\item No companion audit is due in V88.  The next mandatory audit
      remains V90.
\end{enumerate}
\end{programme}

\begin{firewall}[Claim boundary after V87]
\label{fire:V87-boundary}
The paired carrier and certified \(S_{22}\) sector are proved.  The
mesoscopic/response-only uncertified sector, nonadjacent overlap
carrier, other first-connectivity families, global quartic and higher
MTA, WUFC, CPM, continuum construction, and positive mass gap remain
open.  In particular, the full Yang--Mills mass gap remains open.
\end{firewall}

\paragraph{Parent-version record.}
V87 was copied byte-for-byte from
\texttt{Renewal\_Yang\_Mills\_Mass\_Gap\_Research\_Notes\_V86.tex}
before this chapter was appended.  The parent had \(40232\) source
records, \(1384322\) bytes, and SHA256
\[
 \texttt{1183c7c7670e4fc8ec05af3dc1db48ab749e0fee1e59907dddf2c6ca81050115}.
\]
No companion audit is due at V87.

% =============================================================================
% BEGIN APPENDED COMPACT VERSION V88 MATERIAL
% =============================================================================

\clearpage
\chapter{Mesoscopic incidence ledger and the positive-envelope
regression}
\label{chap:V88}

\section{Context: resolve the uncertified endpoint before proceeding}
\label{sec:V88-context}

V87 proved compatible unpaid and paid heat cards for every binary
prefix and inserted them into all four \(S_{22}\) payer locations.
It closed the sector in which each repeated crossing endpoint is paid
either by ordinary prefix overlap or by a superhot coordinate
witness.  The remaining scalar question is not whether the carrier
has a bound: it is whether its dimensionless paired card pays a row
mass which lies mostly outside that prefix.

This chapter first refines the V79 own-prefix cell into the three heat
regimes without duplicating an incidence.  It proves that failure of
all V87 endpoint certificates forces a fixed fraction of the row mass
into the other five V79 cells.  It then tests the remaining
response-only route on the exact V86 packet.  The physical positive
envelope loses a sharp factor \(s^{-1/2}\) which the complete signed
joining-channel sum retains.  Thus the next estimate must move
upstream again: physical absolute value cannot be taken separately on
the joining covariance in this sector.

\section{A coordinate-resolved three-regime refinement}
\label{sec:V88-three-regime-cells}

Let \(D_{uv}\) be the pair-connecting bank of a binary prefix after
the V80 one-row factors have been removed.  At \(e\in D_{uv}\), put
\[
 m_e:=\max\{j_{u,e},j_{v,e}\}.
\]
Partition the coordinates, with the displayed boundary convention,
as
\begin{align}
 D_{uv}^{\rm c}
 &:=
 \{e\in D_{uv}:sm_e<1,\ sm_e^2\le1\},
\label{eq:V88-cold-coordinates}\\
 D_{uv}^{\rm m}
 &:=
 \{e\in D_{uv}:sm_e^2>1,\ sm_e<1\},
\label{eq:V88-mesoscopic-coordinates}\\
 D_{uv}^{\rm h}
 &:=
 \{e\in D_{uv}:sm_e\ge1\}.
\label{eq:V88-superhot-coordinates}
\end{align}
For \(\tau\in\{\mathrm c,\mathrm m,\mathrm h\}\), define
\begin{align}
 A_u^\tau(D_{uv})
 &:=
 \sum_{e\in D_{uv}^\tau}a_{u,e},
\label{eq:V88-regime-row-stock}\\
 H_{uv}^\tau(D_{uv})
 &:=
 \sum_{e\in D_{uv}^\tau}a_{u,e}a_{v,e}.
\label{eq:V88-regime-pair-stock}
\end{align}

\begin{lemma}[Exact regime refinement of the own-prefix cell]
\label{lem:V88-exact-regime-refinement}
The three coordinate sets in
\eqref{eq:V88-cold-coordinates}--%
\eqref{eq:V88-superhot-coordinates} are disjoint and exhaustive.
Consequently,
\begin{align}
 A_u(D_{uv})
 &=
 A_u^{\rm c}(D_{uv})
 +A_u^{\rm m}(D_{uv})
 +A_u^{\rm h}(D_{uv}),
\label{eq:V88-row-regime-identity}\\
 H_{uv}(D_{uv})
 &=
 H_{uv}^{\rm c}(D_{uv})
 +H_{uv}^{\rm m}(D_{uv})
 +H_{uv}^{\rm h}(D_{uv}).
\label{eq:V88-pair-regime-identity}
\end{align}
Replacing the V79 own-prefix incidence \(A_u\) by these three
subcells preserves the exact six-cell row identity and introduces no
new row mass.
\end{lemma}

\begin{proof}
If \(sm_e\ge1\), the coordinate is in \(D^{\rm h}\).  Otherwise
\(sm_e<1\), and exactly one of \(sm_e^2\le1\) or \(sm_e^2>1\)
holds.  Thus the sets form a partition.  Summing the same
nonnegative incidence weights over that partition proves both
identities.  Their sum is the original V79 stock \(A_u\), so
\eqref{eq:V79-incidence-identity} is unchanged.
\end{proof}

\begin{firewall}[Regime labels are not three copies of a coordinate]
\label{fire:V88-no-regime-duplication}
The label is assigned from the maximum of the two representations in
the binary block, but the row stocks remain row-labelled:
\(a_{u,e}\) is entered only in the \(u\)-cell and \(a_{v,e}\) only in
the \(v\)-cell.  The pair stock
\(a_{u,e}a_{v,e}\) is a scalar overlap numerator, not an additional
row incidence.
\end{firewall}

\section{What failure of both V87 certificates forces}
\label{sec:V88-uncertified-endpoint}

Retain
\[
 \gamma_s(D_{uv})
 =
 \frac{\mathfrak b_s(D_{uv})}{H_{uv}(D_{uv})}
\]
when the denominator is nonzero.

\begin{lemma}[Three immediate endpoint compilers]
\label{lem:V88-three-endpoint-compilers}
Fix \(\Lambda\ge1\) and \(0<\eta<1\).  Each of the following
conditions implies
\begin{equation}
 \gamma_s(D_{uv})
 \le\frac{C_{\Lambda,\eta}}{\Xi_u}:
\label{eq:V88-endpoint-compiler}
\end{equation}
\begin{enumerate}[label=\textup{(\roman*)}]
\item \(s\Xi_u\le\Lambda\);
\item \(H_{uv}(D_{uv})\ge\eta\Xi_u\);
\item some \(e\in D_{uv}^{\rm h}\) satisfies
      \(sm_eH_{uv}(D_{uv})\ge\eta\Xi_u\).
\end{enumerate}
\end{lemma}

\begin{proof}
In case \textup{(i)}, the response entry in
\eqref{eq:V87-paired-binary-card} gives
\[
 \gamma_s(D_{uv})
 \le s\le\frac{\Lambda}{\Xi_u}.
\]
Cases \textup{(ii)} and \textup{(iii)} are respectively the ordinary
and heat-witness arguments of
\Cref{lem:V87-certificate-pays-endpoint}.
\end{proof}

\begin{corollary}[Exact scalar state of an uncertified hot endpoint]
\label{cor:V88-uncertified-state}
Choose \(0<\eta<a_*/4\).  If none of the three conditions in
\Cref{lem:V88-three-endpoint-compilers} holds, then
\begin{align}
 s\Xi_u&>\Lambda,
\label{eq:V88-row-hot}\\
 H_{uv}(D_{uv})&<\eta\Xi_u,
\label{eq:V88-overlap-poor}\\
 sm_eH_{uv}(D_{uv})&<\eta\Xi_u
 \quad(e\in D_{uv}^{\rm h}),
\label{eq:V88-no-hot-witness}\\
 A_u(D_{uv})&<\frac{\eta}{a_*}\Xi_u.
\label{eq:V88-own-prefix-small}
\end{align}
In the V79 row identity, the other five incidence cells therefore
satisfy
\begin{equation}
 P_u+U_u+F_u+D_u+I_u
 >
 \left(1-\frac{\eta}{a_*}\right)\Xi_u.
\label{eq:V88-outside-prefix-large}
\end{equation}
\end{corollary}

\begin{proof}
The first three inequalities are the negations of the three compiler
hypotheses.  Since both rows are nontrivial at every coordinate of
the pair-connecting bank,
\[
 H_{uv}(D_{uv})
 \ge a_*A_u(D_{uv}).
\]
This and \eqref{eq:V88-overlap-poor} give
\eqref{eq:V88-own-prefix-small}.  Subtract it from
\eqref{eq:V79-incidence-identity}.
\end{proof}

\begin{assessment}[The residual is forced outside the binary prefix]
\label{ass:V88-forced-outside}
After the private and unbalanced certificates are applied,
\eqref{eq:V88-outside-prefix-large} forces a fixed fraction into the
foreign-prefix, directed-cross, or cut-internal cells, unless the
term belongs to a marked/overlap exception.  Thus a failed paired
prefix card is not repaired by renaming its mesoscopic mass.  The
next acquisition must use the actual outside-prefix cell in the same
source packet.
\end{assessment}

\begin{proposition}[Mesoscopic contraction alone cannot pay an
external row mass]
\label{prop:V88-mesoscopic-scalar-no-go}
Fix \(\alpha\in(\frac12,1)\).  There are scalar incidence tables with
one equal-spin mesoscopic prefix coordinate
\(j_s\asymp s^{-\alpha}\) and additional nonprefix row mass for which
\begin{equation}
 \Xi_u\,\gamma_s(D_{uv})\longrightarrow\infty.
\label{eq:V88-mesoscopic-scalar-no-go}
\end{equation}
Hence the ordinary contraction of that prefix cannot, by itself,
pay an arbitrarily larger row mass outside the prefix.
\end{proposition}

\begin{proof}
For the one-coordinate prefix,
\[
 a_s:=1+j_s(j_s+1)\asymp s^{-2\alpha},
 \qquad
 H_{uv}(D_{uv})=a_s^2\asymp s^{-4\alpha}.
\]
It is mesoscopic because \(sj_s^2\to\infty\) and \(sj_s\to0\).
There is no superhot witness, and
\(sH_{uv}(D_{uv})\to\infty\); hence
\(\mathfrak b_s(D_{uv})=1\) and
\(\gamma_s(D_{uv})=H_{uv}(D_{uv})^{-1}\).
Add a disjoint nonprefix row incidence of mass
\[
 L_s=\lambda_sH_{uv}(D_{uv}),
 \qquad \lambda_s\longrightarrow\infty.
\]
Such masses are realized, up to fixed factors, by choosing an
auxiliary half-integral representation with the corresponding
Casimir.  Then
\[
 \Xi_u\ge L_s,
 \qquad
 \Xi_u\gamma_s(D_{uv})\ge\lambda_s\longrightarrow\infty.
\]
\end{proof}

\begin{firewall}[The scalar no-go is not a source divergence]
\label{fire:V88-scalar-no-go-scope}
The added mass belongs to another V79 incidence cell and may carry
its own valid source certificate.  The proposition refutes only the
attempt to pay that external mass with the mesoscopic prefix
contraction.  It does not assert that the full Yang--Mills source
term is unbounded.
\end{firewall}

\section{The exact minimal packet after physical absolute value}
\label{sec:V88-positive-minimal-packet}

Return to the spin-\(\frac12\) prefix operator
\(\mathsf B_s\) of \eqref{eq:V86-binary-prefix}.  Its two spectral
coefficients have opposite signs:
\[
 1-e^{-3s/2}>0,
 \qquad
 e^{-2s}-e^{-3s/2}<0.
\]
Since \(\chi\) has weight \(1/2\) in both \(V_0\) and \(V_1\), put
\begin{equation}
 \widetilde b_s
 :=
 \langle\chi,|\mathsf B_s|\chi\rangle
 =
 \frac12(1-e^{-2s})
 =
 s+O(s^2).
\label{eq:V88-absolute-prefix-expectation}
\end{equation}
The corresponding paid prefix expectation remains
\[
 \langle\chi,|\mathsf B_s^{\rm pay}|\chi\rangle
 =
 b_s^{\rm pay}
 =
 \frac38+O(s)
\]
by \eqref{eq:V86-binary-prefix-paid-expectation}.

On the same vector \(\Omega_{s,j}\) as in
\eqref{eq:V86-packet-vector}, define the no-suffix positive packet
envelopes
\begin{align}
 \mathsf Q_{s,j}^{|{\rm U}|}
 &:=
 |\mathsf B_s|^{(p)}
 \otimes|\mathsf B_s|^{(q)}
 \otimes\mathsf N_{s,j}^{(e)},
\label{eq:V88-positive-packet-unpaid}\\
 \mathsf Q_{s,j}^{|{\rm P},e|}
 &:=
 |\mathsf B_s|^{(p)}
 \otimes|\mathsf B_s|^{(q)}
 \otimes\mathsf P_{s,j}^{(e)},
\label{eq:V88-positive-packet-paid-e}\\
 \mathsf Q_{s,j}^{|{\rm P},p|}
 &:=
 |\mathsf B_s^{\rm pay}|^{(p)}
 \otimes|\mathsf B_s|^{(q)}
 \otimes\mathsf N_{s,j}^{(e)}.
\label{eq:V88-positive-packet-paid-p}
\end{align}
The \(q\)-prefix-payer operator is obtained by interchanging \(p,q\).

\begin{theorem}[Sharp half-power loss of the local positive envelope]
\label{thm:V88-positive-packet-loss}
Uniformly for \(0<s\le s_0\) and \(j\ge16s^{-1}\),
\begin{align}
 \langle\Omega_{s,j},
 \mathsf Q_{s,j}^{|{\rm U}|}\Omega_{s,j}\rangle
 &\asymp
 \frac{s^{3/2}}{j},
\label{eq:V88-positive-unpaid-scale}\\
 \langle\Omega_{s,j},
 \mathsf Q_{s,j}^{|{\rm P},e|}\Omega_{s,j}\rangle
 &\asymp
 \frac{s^{1/2}}{j},
\label{eq:V88-positive-paid-e-scale}\\
 \langle\Omega_{s,j},
 \mathsf Q_{s,j}^{|{\rm P},p|}\Omega_{s,j}\rangle
 &\asymp
 \frac{s^{1/2}}{j}.
\label{eq:V88-positive-paid-p-scale}
\end{align}
The same estimate as \eqref{eq:V88-positive-paid-p-scale} holds for
the \(q\)-prefix payer.
\end{theorem}

\begin{proof}
Coordinate factorization gives the exact identities
\begin{align*}
 \langle\Omega_{s,j},
 \mathsf Q_{s,j}^{|{\rm U}|}\Omega_{s,j}\rangle
 &=
 \widetilde b_s^2\mathcal E_{s,j}^{\rm N},\\
 \langle\Omega_{s,j},
 \mathsf Q_{s,j}^{|{\rm P},e|}\Omega_{s,j}\rangle
 &=
 \widetilde b_s^2\mathcal E_{s,j}^{\rm P},\\
 \langle\Omega_{s,j},
 \mathsf Q_{s,j}^{|{\rm P},p|}\Omega_{s,j}\rangle
 &=
 b_s^{\rm pay}\widetilde b_s\mathcal E_{s,j}^{\rm N}.
\end{align*}
Now use \eqref{eq:V88-absolute-prefix-expectation},
\eqref{eq:V86-binary-prefix-paid-expectation}, and the sharp V85
estimates
\eqref{eq:V85-sharp-N}--\eqref{eq:V85-sharp-P}.
\end{proof}

\begin{corollary}[The positive route still fails the V84 square-tail
scales]
\label{cor:V88-positive-route-fails}
Choose \(j_s\asymp s^{-3}\).  The three positive expectations in
\Cref{thm:V88-positive-packet-loss} have orders
\begin{equation}
 s^{9/2},\qquad s^{7/2},\qquad s^{7/2},
\label{eq:V88-positive-sequence-orders}
\end{equation}
whereas the corresponding signed complete-packet expectations in
\Cref{thm:V86-minimal-packet} have orders
\begin{equation}
 s^5,\qquad s^4,\qquad s^4.
\label{eq:V88-signed-sequence-orders}
\end{equation}
Thus taking the physical absolute value separately on the joining
covariance loses exactly \(s^{-1/2}\) on every payer placement in
this family.
\end{corollary}

\begin{proof}
Substitute \(j_s\asymp s^{-3}\) into
\eqref{eq:V88-positive-unpaid-scale}--%
\eqref{eq:V88-positive-paid-p-scale} and compare with
\eqref{eq:V86-packet-unpaid-scale}--%
\eqref{eq:V86-packet-paid-prefix-scale}.
\end{proof}

\begin{theorem}[No separately enveloped joining proof of the sharp
response-only target]
\label{thm:V88-no-separate-joining-envelope}
There is no constant \(C\), uniform in small \(s\), which bounds the
three positive packet envelopes above by the sharp signed scales
\[
 C\frac{s^2}{j}
 \quad\hbox{unpaid},\qquad
 C\frac{s}{j}
 \quad\hbox{once paid}.
\]
Consequently the response-only uncertified branch cannot be closed by
first replacing the joining covariance by
\(\mathsf N_{s,j}\) or \(\mathsf P_{s,j}\) and then applying any
scalar refinement of the two prefix cards.
\end{theorem}

\begin{proof}
The ratios of
\eqref{eq:V88-positive-unpaid-scale} and
\eqref{eq:V88-positive-paid-e-scale}--%
\eqref{eq:V88-positive-paid-p-scale} to the proposed right sides are
all comparable to \(s^{-1/2}\).  They diverge.  Since
\(\Omega_{s,j}\) is a product vector across the four input rows,
these expectations are lower bounds for the corresponding positive
product-state capacities.  Subsequent scalar bookkeeping cannot
decrease an already formed positive envelope.
\end{proof}

\begin{firewall}[This is not a counterexample to the signed source]
\label{fire:V88-positive-only}
\Cref{thm:V88-no-separate-joining-envelope} refutes a proof order:
physical absolute value at the joining factor before the complete
packet.  It does not refute the raw signed first-connectivity source.
Indeed
\Cref{thm:V86-minimal-packet} proves that the signed channel sum has
exactly the desired powers on the same product vector.  Nor does the
theorem contradict the certified positive sector
\Cref{thm:V87-certified-S22-sector}, whose endpoint hypotheses fail
for these low-spin prefixes and high joining-row masses.
\end{firewall}

\section{The upstream signed complete-channel gate}
\label{sec:V88-signed-gate}

Let
\[
 \mathcal Q_r^{\rm raw}
 =
 K_{12}^{<r,\mathrm{conn}}
 \otimes K_{34}^{<r,\mathrm{conn}}
 \otimes J_r(12|34)
 \otimes M_{>r}
\]
be the exact signed summand of \eqref{eq:V65-S22}, with the unique
payer inserted before absolute value when present.

\begin{openproblem}[Signed response-only \(S_{22}\) capacity]
\label{open:V88-signed-response-gate}
On the branch in which one crossing endpoint satisfies
\eqref{eq:V88-row-hot}--\eqref{eq:V88-no-hot-witness}, prove a
support-uniform estimate for
\[
 \sup_{\varphi\ {\rm row\text{-}product}}
 |\varphi(\mathcal Q_r^{\rm raw})|
\]
at the required crossing square-tail scale, keeping the complete
joining output-channel sum inside the matrix coefficient.  The paid
version must keep the payer on its actual output before the signed
sum and may take physical absolute value only after the complete
packet estimate.
\end{openproblem}

\begin{assessment}[What the minimal packet says about the new gate]
\label{ass:V88-new-gate}
On the coordinate-factorized vector \(\Omega_{s,j}\), the necessary
gain is exactly the signed cancellation
\[
 \mathcal E_{s,j}^{\rm N}
 \asymp\frac1{j\sqrt s}
 \quad\longrightarrow\quad
 \mathcal E_{s,j}^{\rm S}
 \asymp\frac1j
\]
unpaid, and
\[
 \mathcal E_{s,j}^{\rm P}
 \asymp\frac1{js^{3/2}}
 \quad\longrightarrow\quad
 \mathcal E_{s,j}^{\rm R}
 \asymp\frac1{js}
\]
when the joining coordinate pays.  The next proof must decide whether
these half-power gains are uniform over all row-product inputs or
whether a tilted coherent input destroys them.  A single-vector
calculation cannot settle that capacity question.
\end{assessment}

\section{Post-V88 status and the next test}
\label{sec:V88-status}

\begin{theorem}[Exact status after the mesoscopic regression]
\label{thm:V88-status}
Relative to V1--V87, V88 proves:
\begin{enumerate}[label=\textup{(V88.\arabic*)},leftmargin=6.5em]
\item the V79 own-prefix incidence has an exact, nonduplicating
      input-cold/mesoscopic/superhot coordinate refinement;
\item a row-cold response card, an ordinary prefix certificate, or a
      superhot witness each pays one repeated endpoint;
\item failure of all three alternatives forces a fixed fraction of
      that row mass into the other five V79 incidence cells;
\item a mesoscopic prefix contraction alone cannot pay an arbitrarily
      larger nonprefix row mass;
\item on the exact V86 three-coordinate packet, separately enveloping
      the joining covariance loses precisely \(s^{-1/2}\), both
      unpaid and at every nonzero payer location; and
\item the response-only residual must be estimated as a complete
      signed prefix--joining packet before physical absolute value.
\end{enumerate}
Items \textup{(V88.1)}--\textup{(V88.4)} are scalar and incidence
statements.  Items \textup{(V88.5)}--\textup{(V88.6)} are operator
proof-order statements.  Neither group proves the open signed
capacity estimate in \Cref{open:V88-signed-response-gate}.
\end{theorem}

\begin{proof}
Items \textup{(V88.1)}--\textup{(V88.3)} are
\Cref{lem:V88-exact-regime-refinement,
lem:V88-three-endpoint-compilers,
cor:V88-uncertified-state}.
Item \textup{(V88.4)} is
\Cref{prop:V88-mesoscopic-scalar-no-go}.
Items \textup{(V88.5)}--\textup{(V88.6)} are
\Cref{thm:V88-positive-packet-loss,
cor:V88-positive-route-fails,
thm:V88-no-separate-joining-envelope}.
\end{proof}

\begin{programme}[Mandate for V89]
\label{prog:V88-next}
\begin{enumerate}[label=\textup{(\arabic*)},leftmargin=4em]
\item For the \(SU(2)\) incidence-three joining operator, parameterize
      arbitrary product inputs in the two equal large-spin rows and
      the spin-one row.  Do not assume the anti-aligned magnetic-zero
      vector.
\item Prove or refute a uniform signed complete-channel gain of
      \(s^{1/2}\) relative to the positive envelope in the regime
      \(sj\to\infty\).  Test tilted coherent pairs whose classical
      resultant lies at the output heat scale \(J\asymp s^{-1/2}\).
\item Repeat the test with the payer retained on the joining output.
      The target scales are \(j^{-1}\) unpaid and \((js)^{-1}\)
      paid.  A positive result must control arbitrary row-product
      inputs, not one Clebsch--Gordan column.
\item If either uniform target fails, preserve the smallest raw
      signed complete-packet counterexample and move upstream to the
      sum over first-joining positions or partition states.  Do not
      return to a separately positive joining envelope.
\item Insert the forced outside-prefix alternative
      \eqref{eq:V88-outside-prefix-large} into the V79 routing table,
      keeping foreign-prefix, directed-cross, cut-internal, and
      overlap cells distinct.
\item Write the V85 repeated-bank witness in the exact V61 operator
      order.  The V86 adjacent joint card may be used only if the
      intervening factors are proved removable.
\item No companion audit is due at V89.  V90 must inspect the newest
      Standard-Model and Navier--Stokes notes before the V91
      direction is chosen.
\end{enumerate}
\end{programme}

\begin{firewall}[Claim boundary after V88]
\label{fire:V88-boundary}
V88 proves that the remaining positive joining-envelope route is
false at the sharp response-only scale; it does not prove that the
raw signed source fails.  Uniform signed joining capacity, the
outside-prefix routing completion, the nonadjacent overlap carrier,
other first-connectivity families, global quartic and higher MTA,
WUFC, CPM, continuum construction, and a positive mass gap remain
open.  The full Yang--Mills mass gap remains open.
\end{firewall}

\paragraph{Parent-version record.}
V88 was copied byte-for-byte from
\texttt{Renewal\_Yang\_Mills\_Mass\_Gap\_Research\_Notes\_V87.tex}
before this chapter was appended.  The parent had \(41169\) source
records, \(1413054\) bytes, and SHA256
\[
 \texttt{ebafd9961174d65073ae8db02229454b6667109eb6e7475a119cea9dbe3a2606}.
\]
No companion audit is due at V88.

% =============================================================================
% BEGIN APPENDED COMPACT VERSION V89 MATERIAL
% =============================================================================

\clearpage
\chapter{The signed joining capacity reduces to one vector heat
moment}
\label{chap:V89}

\section{Context: isolate the first-order obstruction exactly}
\label{sec:V89-context}

V88 proved that the response-only \(S_{22}\) branch cannot take
physical absolute value separately on the incidence-three joining
covariance: that proof order loses \(s^{-1/2}\).  The next question is
whether the corresponding gain survives uniformly over all product
inputs, rather than only on the anti-aligned vector used in V85.

For the \(SU(2)\) joining coordinate, the answer reduces to one
operator-valued heat moment on the two equal large-spin rows.  The
quadratic Taylor remainder already has the target \(j^{-1}\) scale.
Only the first-order vector component remains.  This chapter proves
that reduction in both directions and closes the entire
zero-total-magnetic family.

\section{Exact linear--quadratic decomposition}
\label{sec:V89-decomposition}

On \(V_j\otimes V_j\), let
\[
 L_a:=J_{1,a}+J_{2,a},
 \qquad
 L^2:=\sum_{a=1}^3L_a^2.
\]
On the third-row spin-one space let \(S_a\) be the generators, so
\(S^2=2I\), and put
\begin{equation}
 X:=2L\mathbin{\cdot}S
 =2\sum_{a=1}^3L_a\otimes S_a.
\label{eq:V89-interaction}
\end{equation}
The complete incidence-three joining covariance is
\begin{equation}
 \mathsf J_{s,j}
 =
 e^{-s(L+S)^2}
 -
 e^{-sL^2}\otimes e^{-2s}I.
\label{eq:V89-joining-generator-form}
\end{equation}

\begin{lemma}[Exact commuting exponential form]
\label{lem:V89-commuting-form}
The operators \(L^2\), \(S^2\), and \(X\) commute, and
\begin{equation}
 \boxed{
 \mathsf J_{s,j}
 =
 e^{-s(L^2+2)}
 \bigl(e^{-sX}-I\bigr).}
\label{eq:V89-commuting-form}
\end{equation}
On the block with pair spin \(J\) and final spin \(K\),
\begin{equation}
 X
 =
 x_{J,K}I,
 \qquad
 x_{J,K}
 =
 K(K+1)-J(J+1)-2,
\label{eq:V89-X-spectrum}
\end{equation}
and
\begin{equation}
 |x_{J,K}|\le2(J+1).
\label{eq:V89-X-bound}
\end{equation}
\end{lemma}

\begin{proof}
The Casimir \(L^2\) commutes with every \(L_a\), and the two tensor
factors commute with each other.  Hence all three displayed
operators commute.  Since
\[
 (L+S)^2=L^2+2+X,
\]
factorization gives \eqref{eq:V89-commuting-form}.  On a coupled
output block,
\((L+S)^2=K(K+1)\), proving
\eqref{eq:V89-X-spectrum}.  The allowed outputs are
\(K=J-1,J,J+1\), for which \(x_{J,K}\) is respectively
\(-2(J+1),-2,2J\); unavailable negative spins are omitted.
\end{proof}

Taylor's formula with integral remainder gives
\begin{align}
 \mathsf J_{s,j}
 &=
 \mathsf L_{s,j}+\mathsf R_{s,j}^{(2)},
\label{eq:V89-linear-quadratic}\\
 \mathsf L_{s,j}
 &:=
 -s e^{-s(L^2+2)}X,
\label{eq:V89-linear-part}\\
 \mathsf R_{s,j}^{(2)}
 &:=
 s^2e^{-s(L^2+2)}X^2
 \int_0^1(1-t)e^{-tsX}\,\dd t.
\label{eq:V89-quadratic-remainder}
\end{align}

\begin{lemma}[Blockwise Gaussian remainder]
\label{lem:V89-block-remainder}
On every \((J,K)\) block,
\begin{equation}
 \left|
 \mathsf R_{s,j}^{(2)}
 \right|_{J,K}
 \le
 Cs^2(J+1)^2e^{-cs(J-1)^2}.
\label{eq:V89-block-remainder}
\end{equation}
\end{lemma}

\begin{proof}
On that block, the exponential inside
\eqref{eq:V89-quadratic-remainder} has exponent
\[
 (1-t)\bigl(J(J+1)+2\bigr)+tK(K+1).
\]
For \(0\le t\le1\), this is at least the smaller endpoint, hence at
least \(J(J-1)\).  Combine this with
\eqref{eq:V89-X-bound} and
\(\int_0^1(1-t)\,\dd t=1/2\).  Enlarging the constant covers
\(J=0,1\); at \(J=0\), \(X\) vanishes.
\end{proof}

\begin{theorem}[The quadratic remainder already has the target scale]
\label{thm:V89-quadratic-target}
Uniformly for \(0<s\le s_0\) and \(j\ge\frac12\),
\begin{equation}
 \boxed{
 \kappa_\otimes
 \bigl(|\mathsf R_{s,j}^{(2)}|\bigr)
 \le\frac C{2j+1}.}
\label{eq:V89-quadratic-target}
\end{equation}
The product-state capacity is across the three input rows.
\end{theorem}

\begin{proof}
For any product state of rows \(1,2\), the complete pair-spin
probability obeys
\[
 \operatorname{Tr}(P_J^{j,j}\rho_1\otimes\rho_2)
 \le\frac{2J+1}{2j+1}
\]
by \Cref{lem:V87-projection-product-bound}.  Conditional on \(J\),
the three final-output probabilities sum to one for every state of
the spin-one row.  Therefore
\[
 \kappa_\otimes
 \bigl(|\mathsf R_{s,j}^{(2)}|\bigr)
 \le
 \frac{Cs^2}{2j+1}
 \sum_{J\ge0}
 (2J+1)(J+1)^2e^{-cs(J-1)^2}.
\]
The Gaussian sum is \(O(s^{-2})\), which proves the result.
\end{proof}

\section{Equivalence with a two-row vector heat moment}
\label{sec:V89-vector-moment}

For a product unit vector \(x\otimes y\in V_j\otimes V_j\), define
\begin{equation}
 \mathbf v_{s,j}(x,y)
 :=
 \left(
 \langle x\otimes y,e^{-sL^2}L_a(x\otimes y)\rangle
 \right)_{a=1}^3
\label{eq:V89-vector-heat-moment}
\end{equation}
and
\begin{equation}
 \mathcal V_{s,j}
 :=
 \sup_{\|x\|=\|y\|=1}
 |\mathbf v_{s,j}(x,y)|.
\label{eq:V89-vector-radius}
\end{equation}

\begin{theorem}[Signed joining capacity is equivalent to the vector
heat estimate]
\label{thm:V89-vector-equivalence}
Let
\[
 \omega_\otimes(\mathsf J_{s,j})
 :=
 \sup_{\substack{\|x\|=\|y\|=\|z\|=1}}
 \left|
 \langle x\otimes y\otimes z,
 \mathsf J_{s,j}
 x\otimes y\otimes z\rangle
 \right|.
\]
Then
\begin{equation}
 2se^{-2s}\mathcal V_{s,j}
 -\frac C{2j+1}
 \le
 \omega_\otimes(\mathsf J_{s,j})
 \le
 2se^{-2s}\mathcal V_{s,j}
 +\frac C{2j+1}.
\label{eq:V89-vector-equivalence}
\end{equation}
Consequently,
\begin{equation}
 \boxed{
 \omega_\otimes(\mathsf J_{s,j})
 \le\frac C{2j+1}}
\label{eq:V89-desired-signed-card}
\end{equation}
is equivalent, up to universal constants, to
\begin{equation}
 \boxed{
 \mathcal V_{s,j}
 \le\frac C{s(2j+1)}.}
\label{eq:V89-vector-gate}
\end{equation}
\end{theorem}

\begin{proof}
For a product vector across the pair and the spin-one row,
\eqref{eq:V89-linear-part} factorizes:
\begin{equation}
 \langle x\otimes y\otimes z,
 \mathsf L_{s,j}x\otimes y\otimes z\rangle
 =
 -2se^{-2s}
 \mathbf v_{s,j}(x,y)
 \mathbin{\cdot}
 \langle z,\mathbf S z\rangle.
\label{eq:V89-linear-factorization}
\end{equation}
Every spin-one state has
\(|\langle z,\mathbf S z\rangle|\le1\), proving the upper bound after
\Cref{thm:V89-quadratic-target}.

Conversely, for every nonzero real vector \(\mathbf v\), a spin-one
coherent state polarized along \(\mathbf v\) has
\(\langle\mathbf S\rangle=\mathbf v/|\mathbf v|\).
Choose this state in \eqref{eq:V89-linear-factorization}, use the
reverse triangle inequality, and take a maximizing sequence for
\(\mathcal V_{s,j}\).  The lower bound follows from
\Cref{thm:V89-quadratic-target}.  Since \(e^{-2s}\) is bounded above
and below on \(0<s\le s_0\), the two boxed estimates are equivalent.
\end{proof}

\begin{firewall}[The remaining gate is rank one, not a positive heat
tail]
\label{fire:V89-rank-one-gate}
The positive estimate
\[
 \kappa_\otimes(e^{-sL^2}|L|)
 \lesssim\frac1{j s^{3/2}}
\]
obtained from complete projection traces is too weak by
\(s^{-1/2}\).  The required improvement in
\eqref{eq:V89-vector-gate} concerns the signed rank-one vector moment
before its Euclidean magnitude is taken.  Replacing \(L_a\) by
\(|L|\), or summing magnetic components absolutely, reproduces the
V88 loss and cannot prove the gate.
\end{firewall}

\section{Magnetic columns and a closed extremal ladder}
\label{sec:V89-magnetic-columns}

Let \(|j,m\rangle\) denote the usual \(J_3\)-eigenbasis.

\begin{lemma}[Exact vector moment on a magnetic product column]
\label{lem:V89-magnetic-vector}
For
\[
 x=|j,m_1\rangle,\qquad y=|j,m_2\rangle,\qquad M=m_1+m_2,
\]
put
\begin{equation}
 h_{s,j}(m_1,m_2)
 :=
 \langle x\otimes y,e^{-sL^2}x\otimes y\rangle.
\label{eq:V89-magnetic-heat}
\end{equation}
Then
\begin{equation}
 \boxed{
 \mathbf v_{s,j}(x,y)
 =
 (0,0,Mh_{s,j}(m_1,m_2)).}
\label{eq:V89-magnetic-vector}
\end{equation}
\end{lemma}

\begin{proof}
The product vector is an eigenvector of \(L_3\) with eigenvalue
\(M\), and \(e^{-sL^2}\) commutes with every \(L_a\).  Hence the
third component is \(M h_{s,j}\).  The operators
\(L_1,L_2\) are linear combinations of \(L_+\) and \(L_-\), which
map the total magnetic subspace \(M\) into the orthogonal subspaces
\(M+1\) and \(M-1\).  Their diagonal expectations vanish.
\end{proof}

\begin{corollary}[All opposite magnetic pairs have the signed target]
\label{cor:V89-zero-magnetic-family}
For every \(m\in\{-j,-j+1,\ldots,j\}\), every spin-one unit vector
\(z\), and every \(0<s\le s_0\),
\begin{equation}
 \left|
 \langle j,m;j,-m;z,
 \mathsf J_{s,j}
 j,m;j,-m;z\rangle
 \right|
 \le\frac C{2j+1}.
\label{eq:V89-zero-magnetic-family}
\end{equation}
\end{corollary}

\begin{proof}
Here \(M=0\), so
\Cref{lem:V89-magnetic-vector} makes the linear term vanish.
Apply \Cref{thm:V89-quadratic-target} to the remainder.
\end{proof}

\begin{remark}[A precise magnetic-column subproblem]
\label{rem:V89-magnetic-subproblem}
For a general magnetic product column, the complete projection bound
gives only
\[
 h_{s,j}(m_1,m_2)\le\frac C{s(2j+1)}
\]
and hence
\[
 |\mathbf v_{s,j}|
 \le
 \frac{C|M|}{s(2j+1)}.
\]
The desired vector gate follows on all magnetic columns from the
sharper estimate
\begin{equation}
 h_{s,j}(m_1,m_2)
 \le
 \frac C{s(2j+1)(1+|M|)}.
\label{eq:V89-magnetic-column-target}
\end{equation}
This estimate cannot be obtained by summing complete projection
dimensions alone; it is a Clebsch--Gordan column bound.
\end{remark}

For \(0\le n\le2j\), define the extremal ladder state
\begin{equation}
 \psi_{j,n}
 :=
 |j,j\rangle\otimes|j,-j+n\rangle.
\label{eq:V89-extremal-ladder}
\end{equation}
It has total magnetic number \(n\).

\begin{lemma}[Exact extremal-ladder pair weights]
\label{lem:V89-extremal-weights}
For \(n\le J\le2j\),
\begin{equation}
 p_{j,J}^{(n)}
 :=
 \left|
 \langle j,j;j,-j+n\mid J,n\rangle
 \right|^2
 =
 (2J+1)
 \frac{
 (J+n)!(2j)!(2j-n)!}
 {(J-n)!\,n!\,(2j-J)!(2j+J+1)!}.
\label{eq:V89-extremal-weights}
\end{equation}
For \(n\le j\),
\begin{equation}
 p_{j,J}^{(n)}
 \le
 C\frac{J+1}{j+1}
 \frac1{n!}
 \left(\frac{C(J+1)^2}{j+1}\right)^n.
\label{eq:V89-extremal-weight-upper}
\end{equation}
\end{lemma}

\begin{proof}
In the Racah formula for the Clebsch--Gordan coefficient, the
extremal value \(m_1=j\) forces the summation index to be zero.
Squaring the remaining factorial prefactor gives
\eqref{eq:V89-extremal-weights}; at \(n=0\) it reduces to
\eqref{eq:V85-pjJ}.

Divide \eqref{eq:V89-extremal-weights} by its \(n=0\) value:
\begin{equation}
 \frac{p_{j,J}^{(n)}}{p_{j,J}^{(0)}}
 =
 \frac{(J+n)!}{(J-n)!\,n!}
 \frac{(2j-n)!}{(2j)!}.
\label{eq:V89-extremal-ratio}
\end{equation}
If \(n\le j\) and \(J\ge n\), the first factorial quotient is at
most \((2J+1)^{2n}\), while the second is at most
\((j+1)^{-n}\).  Combine this with
\eqref{eq:V85-p-upper}, enlarging universal constants to include
small \(j,J\).
\end{proof}

\begin{lemma}[Heat moment of the extremal ladder]
\label{lem:V89-extremal-heat}
There are universal \(C_0,C>0\) such that, if \(sj\ge C_0\), then
for every \(0<s\le s_0\) and \(0\le n\le2j\),
\begin{equation}
 \boxed{
 n\,
 \langle\psi_{j,n},e^{-sL^2}\psi_{j,n}\rangle
 \le
 \frac C{s(2j+1)}.}
\label{eq:V89-extremal-heat}
\end{equation}
The left side is interpreted as zero at \(n=0\).
\end{lemma}

\begin{proof}
First suppose \(1\le n\le j\).  By
\eqref{eq:V89-extremal-weight-upper} and the Gaussian moment bound,
\begin{align*}
 \langle\psi_{j,n},e^{-sL^2}\psi_{j,n}\rangle
 &\le
 \frac C{j+1}
 \frac{(C/(j+1))^n}{n!}
 \sum_{J\ge0}(J+1)^{2n+1}e^{-sJ(J+1)}\\
 &\le
 \frac C{s(j+1)}
 \left(\frac C{s(j+1)}\right)^n.
\end{align*}
Indeed, comparison with the radial integral gives
\[
 \sum_{J\ge0}(J+1)^{2n+1}e^{-sJ(J+1)}
 \le C^{n+1}n!s^{-(n+1)}.
\]
Choose \(C_0\) so that \(C/(s(j+1))\le1/2\).  Then
\(n2^{-n}\le1\), proving the claim in this range.

If \(j<n\le2j\), every pair output has \(J\ge n\), so
\[
 \langle\psi_{j,n},e^{-sL^2}\psi_{j,n}\rangle
 \le e^{-sn(n+1)}.
\]
Since \(sj\ge C_0\), the function
\(n\mapsto ne^{-sn(n+1)}\) is decreasing for \(n\ge j\), after
enlarging \(C_0\).  At \(n=j\),
\[
 j e^{-sj(j+1)}
 \le
 \frac C{sj},
\]
because \(xe^{-x}\) is bounded for \(x=sj(j+1)\), and
\[
 j e^{-sj(j+1)}
 =
 \frac{sj(j+1)e^{-sj(j+1)}}{s(j+1)}
 \le \frac C{s(2j+1)}.
\]
\end{proof}

\begin{theorem}[The whole extremal magnetic ladder satisfies the
uniform signed joining card]
\label{thm:V89-extremal-ladder-closed}
There are universal \(C_0,C>0\) such that, if \(sj\ge C_0\), then for
every \(0<s\le s_0\), \(0\le n\le2j\), and every spin-one unit
vector \(z\),
\begin{equation}
 \boxed{
 \left|
 \langle\psi_{j,n}\otimes z,
 \mathsf J_{s,j}\psi_{j,n}\otimes z\rangle
 \right|
 \le\frac C{2j+1}.}
\label{eq:V89-extremal-ladder-closed}
\end{equation}
The symmetric ladders obtained by exchanging the two large-spin rows
or reversing all magnetic numbers obey the same estimate.
\end{theorem}

\begin{proof}
By \Cref{lem:V89-magnetic-vector}, the vector heat moment has magnitude
\[
 n\,
 \langle\psi_{j,n},e^{-sL^2}\psi_{j,n}\rangle.
\]
Apply \Cref{lem:V89-extremal-heat} in
\eqref{eq:V89-linear-factorization}, then add the quadratic remainder
from \Cref{thm:V89-quadratic-target}.  Row exchange and simultaneous
magnetic reversal are unitary symmetries of \(L^2\) and
\mathsf J_{s,j}\).
\end{proof}

\section{The paid quadratic remainder}
\label{sec:V89-paid-remainder}

On nontrivial final spin \(K\), put
\[
 \mathsf W
 :=
 \frac{1+(L+S)^2}{1-e^{-s(L+S)^2}},
\]
and set it to zero on \(K=0\).  The signed joining-payer operator is
\[
 \mathsf J_{s,j}^{\rm pay}:=\mathsf W\mathsf J_{s,j}.
\]
Multiplying \eqref{eq:V89-linear-quadratic} by \(\mathsf W\) gives
\[
 \mathsf J_{s,j}^{\rm pay}
 =
 \mathsf L_{s,j}^{\rm pay}
 +
 \mathsf R_{s,j}^{\rm pay,(2)},
\]
where
\[
 \mathsf L_{s,j}^{\rm pay}
 :=
 -s\mathsf W e^{-s(L^2+2)}X,
 \qquad
 \mathsf R_{s,j}^{\rm pay,(2)}
 :=
 \mathsf W\mathsf R_{s,j}^{(2)}.
\]

\begin{theorem}[The paid quadratic remainder has the paid target]
\label{thm:V89-paid-remainder}
Uniformly for \(0<s\le s_0\) and \(j\ge\frac12\),
\begin{equation}
 \boxed{
 \kappa_\otimes
 \bigl(|\mathsf R_{s,j}^{\rm pay,(2)}|\bigr)
 \le
 \frac C{s(2j+1)}.}
\label{eq:V89-paid-remainder}
\end{equation}
\end{theorem}

\begin{proof}
On a nonzero \(K\)-block,
\[
 \frac{1+K(K+1)}{1-e^{-sK(K+1)}}
 \le
 C\bigl(s^{-1}+(K+1)^2\bigr)
 \le
 C\bigl(s^{-1}+(J+2)^2\bigr).
\]
Combine this with \eqref{eq:V89-block-remainder}.  As in
\Cref{thm:V89-quadratic-target}, the complete pair projection costs
at most \((2J+1)/(2j+1)\), and the conditional final-output
probabilities sum to one.  Hence the capacity is at most
\[
 \frac{Cs^2}{2j+1}
 \sum_{J\ge0}
 (2J+1)(J+1)^2
 \bigl(s^{-1}+(J+2)^2\bigr)
 e^{-cs(J-1)^2}.
\]
Here the \(s^{-1}\)-weighted cubic Gaussian moment and the quintic
Gaussian moment are both \(O(s^{-3})\).  Therefore the display is
\(O((s(2j+1))^{-1})\).
\end{proof}

\begin{openproblem}[Paid linear joining gate]
\label{open:V89-paid-linear-gate}
Prove
\begin{equation}
 \sup_{\varphi\ {\rm row\text{-}product}}
 |\varphi(\mathsf L_{s,j}^{\rm pay})|
 \le
 \frac C{s(2j+1)}
\label{eq:V89-paid-linear-gate}
\end{equation}
for \(sj\ge C_0\), or construct a product-state counterexample.
\Cref{thm:V89-paid-remainder} shows that this is the sole remaining
local joining-payer term.  Because \(\mathsf W\) depends on the final
coupled output \(K\), it cannot be factored into the unpayed vector
moment \eqref{eq:V89-vector-heat-moment} without a new three-channel
interpolation estimate.
\end{openproblem}

\section{Post-V89 status and the next test}
\label{sec:V89-status}

\begin{theorem}[Exact status after the first-order joining reduction]
\label{thm:V89-status}
Relative to V1--V88, V89 proves:
\begin{enumerate}[label=\textup{(V89.\arabic*)},leftmargin=6.5em]
\item the complete \(SU(2)\) incidence-three joining covariance has
      an exact commuting first-order plus quadratic decomposition;
\item the absolute product-state capacity of the quadratic remainder
      is \(O((2j+1)^{-1})\);
\item the full signed joining card is equivalent, up to that already
      controlled remainder, to the two-row vector heat estimate
      \eqref{eq:V89-vector-gate};
\item every opposite magnetic pair, and the full extremal magnetic
      ladder together with its row and magnetic symmetries, satisfies
      the sharp signed joining card when \(sj\ge C_0\);
\item after the final-output payer is inserted, the paid quadratic
      remainder is \(O((s(2j+1))^{-1})\); and
\item only the paid first-order term remains in the local
      joining-payer card.
\end{enumerate}
In particular, the former positive-envelope bottleneck has been
replaced by two precise signed first-order questions:
\eqref{eq:V89-vector-gate} for arbitrary pair-product vectors and
\eqref{eq:V89-paid-linear-gate} with the final channel retained.
\end{theorem}

\begin{proof}
Item \textup{(V89.1)} is
\Cref{lem:V89-commuting-form,lem:V89-block-remainder}.
Item \textup{(V89.2)} is \Cref{thm:V89-quadratic-target}.
Item \textup{(V89.3)} is \Cref{thm:V89-vector-equivalence}.
Item \textup{(V89.4)} is
\Cref{cor:V89-zero-magnetic-family,thm:V89-extremal-ladder-closed}.
Items \textup{(V89.5)}--\textup{(V89.6)} are
\Cref{thm:V89-paid-remainder,open:V89-paid-linear-gate}.
\end{proof}

\begin{openproblem}[Uniform product-column vector heat gate]
\label{open:V89-uniform-vector-gate}
For \(0<s\le s_0\), \(sj\ge C_0\), and arbitrary unit vectors
\(x,y\in V_j\), prove
\[
 \left|
 \left(
 \langle x\otimes y,e^{-sL^2}L_a(x\otimes y)\rangle
 \right)_{a=1}^3
 \right|
 \le \frac C{s(2j+1)},
\]
or exhibit a product-state counterexample.  Magnetic columns do not
settle this problem because off-diagonal coherences between total
magnetic sectors contribute to \(L_1\) and \(L_2\).  Any proof must
control those coherences before taking the Euclidean norm.
\end{openproblem}

\begin{programme}[Mandate for V90]
\label{prog:V89-next}
\begin{enumerate}[label=\textup{(\arabic*)},leftmargin=4em]
\item Before choosing the V91 direction, perform the compulsory
      read-only audit of the newest Standard-Model and
      Navier--Stokes notes.  Import only proved mechanisms whose
      hypotheses match the present row-product and heat-window
      geometry.
\item Attack \Cref{open:V89-uniform-vector-gate} by decomposing the
      rank-one pair product into adjacent magnetic columns.  Seek a
      vector-valued Clebsch--Gordan or Berezin estimate that retains
      cancellation between the three components.
\item Test the genuinely missing inputs: tilted coherent pairs,
      superpositions of separated magnetic columns, and states whose
      classical resultant lies at \(J\asymp s^{-1/2}\).  A numerical
      test may guide the proof but cannot replace a uniform estimate.
\item For the paid gate
      \eqref{eq:V89-paid-linear-gate}, retain the three final channels
      \(K=J-1,J,J+1\) and seek an interpolation identity before
      absolute value.  Do not dominate the payer by a scalar maximum
      over \(K\), which would recreate the V88 half-power loss.
\item If either signed gate is false, preserve an exact smallest
      counterexample and move upstream to the sum over first-joining
      positions or to partition-state cancellation.  Do not return
      to a separately positive joining envelope.
\item If both local gates close, insert them into the V79
      cold/mesoscopic/superhot routing table, then resume the
      outside-prefix and nonadjacent-overlap programme without
      conflating their incidence cells.
\item At V100, freeze the then-current active note with suffix
      \texttt{\_f2}, restart the active sequence at V1, and begin that
      compact note with context and a proof-indexed summary of the
      established results.
\end{enumerate}
\end{programme}

\begin{firewall}[Claim boundary after V89]
\label{fire:V89-boundary}
V89 does not prove the uniform vector heat gate, the paid linear
joining gate, or the full incidence-source estimate.  The
outside-prefix routing completion, nonadjacent overlap carrier,
remaining first-connectivity families, global quartic and higher
MTA, WUFC, CPM, continuum construction, and positive spectral gap
also remain open.  Consequently, the full four-dimensional quantum
Yang--Mills existence and mass-gap problem remains open.
\end{firewall}

\paragraph{Parent-version record.}
V89 was copied byte-for-byte from
\texttt{Renewal\_Yang\_Mills\_Mass\_Gap\_Research\_Notes\_V88.tex}
before this chapter was appended.  The parent had \(41740\) source
records, \(1432651\) bytes, and SHA256
\[
 \texttt{b8ec74846518680a94b310e7b0e7d79d8601acd37097bd4c05d799c5dea0310c}.
\]
No companion audit is due at V89.  The audit is compulsory at V90.

% =============================================================================
% BEGIN APPENDED COMPACT VERSION V90 MATERIAL
% =============================================================================

\clearpage
\chapter{A coherent-tilt counterexample to the local signed joining
card}
\label{chap:V90}

\section{Context and compulsory companion audit}
\label{sec:V90-context}

V89 reduced the complete incidence-three \(SU(2)\) joining covariance
to a first-order vector heat moment plus a quadratic remainder already
at the desired \(j^{-1}\) scale.  It left open whether the vector moment
has an additional \(s^{1/2}\) cancellation uniformly over all
row-product inputs.  V90 first performs the scheduled companion audit
and then tests that uniform assertion on the tilted coherent family
explicitly requested in V88--V89.

\begin{assessment}[Exact companion files audited at V90]
\label{ass:V90-companion-identity}
The read-only audit used the following stable file identities:
\begin{align*}
 &\texttt{Renewal\_SM\_Einstein\_Continuum\_Research\_Notes\_V39.tex},
 \\
 &\qquad
 \text{SHA256 }\
 \texttt{c1623ae802b7bbbe3f01242cf00b80f57a55bfe906055bcfaa829b39d0b7c871},
 \\
 &\texttt{Renewal\_NS\_Stability\_Research\_Notes\_V31.tex},
 \\
 &\qquad
 \text{SHA256 }\
 \texttt{f1121b62238ad5491bb7cf03635ea9934942edf573ed8faaa9ee79e1d38a6696}.
\end{align*}
The SM file advanced during the initial inventory; V39, not the stale
inventory entry, was read and hashed before and after inspection.
\end{assessment}

\begin{proposition}[What the companion audit transfers]
\label{prop:V90-companion-transfer}
The following proof-order conclusions transfer to the present
programme, but no analytic estimate is imported across models.
\begin{enumerate}[label=\textup{(\roman*)},leftmargin=3.7em]
\item SM V39 obtains its tracefree boundary margin only after the
      physical material, multiplier, and bath variables are assembled
      in one coupled block with the energy-compatible sign.  SM V37's
      preceding quartet calculation shows that an intermediate
      positive square in a removable gauge range cancels after complete
      elimination.
\item NS V31 proves that the heat/Duhamel first-moment exponent is sharp
      on a forced mode and that a signed polarization hinge has a
      nontrivial balanced-sector nullspace.  A missing scale cannot be
      recovered by renaming a weaker positive heat norm, and a signed
      observable need not control the positive total sector.
\item Therefore a Yang--Mills joining estimate may be used locally only
      after it is tested on the complete product-state family.  If it
      fails, the next admissible location for cancellation is the
      coupled sum over first-joining positions or the partition-state
      packet, before separate absolute values.
\end{enumerate}
\end{proposition}

\begin{proof}
Items \textup{(i)} and \textup{(ii)} are direct logical consequences of
the stated finite-dimensional cancellation, coupled energy, sharp
forced-mode, and nullspace theorems in the audited notes.  Their
hypotheses do not identify the Einstein boundary symbol or the
Navier--Stokes Duhamel operator with the present finite-dimensional
\(SU(2)\) heat operator; hence only the proof-order implication in
\textup{(iii)} is transferred.  The remaining assertion is precisely
the alternative already required by \Cref{prog:V89-next}.
\end{proof}

\begin{firewall}[No cross-program theorem substitution]
\label{fire:V90-no-cross-substitution}
Neither companion note supplies a Clebsch--Gordan estimate or a
Yang--Mills measure construction.  The audit changes the acquisition
order, not the hypotheses of the current representation-theoretic
problem.  All estimates below are proved directly in \(V_j\otimes V_j\).
\end{firewall}

\section{Exact scalarization into total-magnetic columns}
\label{sec:V90-column-scalarization}

Write
\[
 x=\sum_{m=-j}^{j}x_m|j,m\rangle,
 \qquad
 y=\sum_{n=-j}^{j}y_n|j,n\rangle
\]
in the \(J_3\)-basis.  For every integer \(-2j\le M\le2j\), let
\[
 \mathcal I_{j,M}
 :=
 \{m\in\{-j,-j+1,\ldots,j\}:M-m\in[-j,j]\}
\]
and define the product column
\begin{equation}
 w_M(m):=x_m y_{M-m},
 \qquad m\in\mathcal I_{j,M}.
\label{eq:V90-product-column}
\end{equation}

\begin{lemma}[Exact Jacobi matrix in one magnetic column]
\label{lem:V90-column-Jacobi}
The restriction of \(L^2\) to total magnetic number \(M\), in the
ordered basis
\(\{|j,m\rangle\otimes|j,M-m\rangle:m\in\mathcal I_{j,M}\}\),
is the real symmetric Jacobi matrix \(H_{j,M}\) with
\begin{align}
 H_{j,M}(m,m)
 &=
 2j(j+1)+2m(M-m),
\label{eq:V90-Jacobi-diagonal}\\
 H_{j,M}(m,m+1)
 &=
 \bigl[
 (j-m)(j+m+1)
 (j+M-m)(j-M+m+1)
 \bigr]^{1/2}.
\label{eq:V90-Jacobi-offdiagonal}
\end{align}
All other entries vanish apart from the symmetric transpose.
\end{lemma}

\begin{proof}
Use
\[
 L^2
 =
 2j(j+1)+2J_{1,3}J_{2,3}
 +J_{1,+}J_{2,-}+J_{1,-}J_{2,+}.
\]
The diagonal term gives \eqref{eq:V90-Jacobi-diagonal}.  The standard
ladder coefficients give \eqref{eq:V90-Jacobi-offdiagonal}; both ladder
terms preserve \(m+(M-m)=M\).
\end{proof}

\begin{theorem}[The vector radius is one signed column functional]
\label{thm:V90-vector-column-scalarization}
For \(0<s\le s_0\),
\begin{equation}
 \boxed{
 \mathcal V_{s,j}
 =
 \sup_{\|x\|=\|y\|=1}
 \left|
 \sum_{M=-2j}^{2j}
 M\,
 \langle w_M,e^{-sH_{j,M}}w_M\rangle
 \right|.}
\label{eq:V90-vector-column-scalarization}
\end{equation}
Each scalar
\begin{equation}
 Q_{s,j,M}(x,y)
 :=
 \langle w_M,e^{-sH_{j,M}}w_M\rangle
\label{eq:V90-column-heat-mass}
\end{equation}
is nonnegative.  Thus there are no coherences between different
total-\(M\) columns after the vector is aligned with the third axis;
the unresolved coherences lie inside each product column
\eqref{eq:V90-product-column}.
\end{theorem}

\begin{proof}
Because \(L_3\) and \(L^2\) commute, the restriction of
\(e^{-sL^2}L_3\) to the \(M\)-column is
\(M e^{-sH_{j,M}}\).  Orthogonality of the columns gives
\[
 \langle x\otimes y,e^{-sL^2}L_3(x\otimes y)\rangle
 =
 \sum_M M Q_{s,j,M}(x,y).
\]
Positivity of the heat operator proves \(Q_{s,j,M}\ge0\).

Under a common rotation \(g\), the three components of
\(\mathbf v_{s,j}(x,y)\) transform by the corresponding rotation in
\(\mathbb R^3\).  For every nonzero vector choose \(g\) which sends it
to the third axis.  The rotated input is still a row-product vector,
so the supremum of the third component equals the supremum of the
Euclidean magnitude.  This proves
\eqref{eq:V90-vector-column-scalarization}.
\end{proof}

\begin{remark}[Correction of the V89 acquisition language]
\label{rem:V90-no-adjacent-column-gate}
Direct formulas for \(L_1,L_2\) contain adjacent magnetic columns, but
they are not an additional gate for the rotationally invariant
supremum \(\mathcal V_{s,j}\).  The exact gate is the signed first
moment of the nonnegative masses \eqref{eq:V90-column-heat-mass}.
Replacing that signed first moment by
\(\sum_M|M|Q_{s,j,M}\) is again the positive-envelope route and need
not be sharp.
\end{remark}

\section{Heat susceptibility of a tilted antipodal coherent pair}
\label{sec:V90-coherent-susceptibility}

Let \(|j,\mathbf e_x\rangle\) and
\(|j,-\mathbf e_x\rangle\) be the spin-\(j\) coherent states of
maximal and minimal single-spin \(J_x\), and put
\begin{equation}
 \psi_j^0
 :=
 |j,\mathbf e_x\rangle\otimes|j,-\mathbf e_x\rangle.
\label{eq:V90-antipodal-base}
\end{equation}
This is the V85 anti-aligned pair after a common rotation.  Its
spin-\(J\) weights are
\begin{equation}
 p_{j,J}
 =
 (2J+1)
 \frac{[(2j)!]^2}{(2j-J)!(2j+J+1)!}.
\label{eq:V90-antipodal-weights}
\end{equation}

\begin{lemma}[The projected base state has zero \(x\)-magnetic number]
\label{lem:V90-base-projection}
For every \(0\le J\le2j\), there is a phase choice such that
\begin{equation}
 P_J^{j,j}\psi_j^0
 =
 \sqrt{p_{j,J}}\,|J,0\rangle_x.
\label{eq:V90-base-projection}
\end{equation}
Moreover,
\begin{equation}
 {}_x\langle J,0|L_3^2|J,0\rangle_x
 =
 \frac12J(J+1).
\label{eq:V90-transverse-second-moment}
\end{equation}
\end{lemma}

\begin{proof}
The base vector obeys \(L_x\psi_j^0=0\), and \(P_J^{j,j}\) commutes
with \(L_x\).  The zero-\(x\)-magnetic line in the multiplicity-free
spin-\(J\) summand is one dimensional, giving
\eqref{eq:V90-base-projection}; its norm is the weight
\eqref{eq:V85-pjJ}.  In the state \(|J,0\rangle_x\),
\(\langle L_x^2\rangle=0\), while rotation by \(\pi/2\) about the
\(x\)-axis interchanges the second moments of \(L_2\) and \(L_3\).
Their sum is \(J(J+1)\), proving
\eqref{eq:V90-transverse-second-moment}.
\end{proof}

For \(h\ge0\), define the normalized row-product vector
\begin{align}
 x_{j,h}
 &:=
 \cosh(h/2)^{-j}
 e^{hJ_3/2}|j,\mathbf e_x\rangle,
\label{eq:V90-tilted-x}\\
 y_{j,h}
 &:=
 \cosh(h/2)^{-j}
 e^{hJ_3/2}|j,-\mathbf e_x\rangle,
\label{eq:V90-tilted-y}\\
 \psi_{j,h}
 &:=
 x_{j,h}\otimes y_{j,h}.
\label{eq:V90-tilted-product}
\end{align}
Also put
\begin{align}
 Z_{s,j}(h)
 &:=
 \langle\psi_j^0,
 e^{-sL^2}e^{hL_3}\psi_j^0\rangle,
\label{eq:V90-tilted-partition}\\
 N_j(h)
 &:=
 \langle\psi_j^0,e^{hL_3}\psi_j^0\rangle.
\label{eq:V90-tilted-normalization}
\end{align}

\begin{lemma}[Exact tilt derivative and its positive lower bound]
\label{lem:V90-tilt-derivative}
For every \(h\ge0\),
\begin{align}
 N_j(h)
 &=
 \cosh(h/2)^{4j},
\label{eq:V90-normalization-exact}\\
 \langle\psi_{j,h},
 e^{-sL^2}L_3\psi_{j,h}\rangle
 &=
 \frac{Z_{s,j}'(h)}{N_j(h)},
\label{eq:V90-vector-as-derivative}\\
 Z_{s,j}'(h)
 &\ge
 \frac h2
 \sum_{J=0}^{2j}
 p_{j,J}e^{-sJ(J+1)}J(J+1).
\label{eq:V90-derivative-lower}
\end{align}
\end{lemma}

\begin{proof}
A spin-\(j\) coherent state is the symmetric tensor power of \(2j\)
identical spin-\(\frac12\) vectors.  For either sign of the \(x\)-axis
state,
\[
 \langle j,\pm\mathbf e_x|
 e^{hJ_3}|j,\pm\mathbf e_x\rangle
 =
 \cosh(h/2)^{2j}.
\]
Multiplication of the two row factors proves
\eqref{eq:V90-normalization-exact} and the normalizations in
\eqref{eq:V90-tilted-x}--\eqref{eq:V90-tilted-product}.
Because \(L_3\) commutes with \(L^2\), differentiation of
\eqref{eq:V90-tilted-partition} proves
\eqref{eq:V90-vector-as-derivative}.

Use \Cref{lem:V90-base-projection} to write
\[
 Z_{s,j}'(h)
 =
 \sum_J p_{j,J}e^{-sJ(J+1)}
 {}_x\langle J,0|L_3e^{hL_3}|J,0\rangle_x.
\]
Rotation by \(\pi\) about the \(x\)-axis shows that the \(L_3\)
spectral weights of \(|J,0\rangle_x\) at \(M\) and \(-M\) are equal.
Pairing them and using \(\sinh(hM)\ge hM\) for \(M,h\ge0\) gives
\[
 {}_x\langle J,0|L_3e^{hL_3}|J,0\rangle_x
 \ge
 h\,{}_x\langle J,0|L_3^2|J,0\rangle_x.
\]
Now apply \eqref{eq:V90-transverse-second-moment}.
\end{proof}

\begin{lemma}[The heat-window susceptibility is of order
\(j^{-1}s^{-2}\)]
\label{lem:V90-heat-susceptibility}
There is a universal \(c>0\) such that, whenever \(0<s\le s_0\) and
\(j\ge16s^{-1}\),
\begin{equation}
 \sum_{J=0}^{2j}
 p_{j,J}e^{-sJ(J+1)}J(J+1)
 \ge
 \frac{c}{j s^2}.
\label{eq:V90-heat-susceptibility}
\end{equation}
\end{lemma}

\begin{proof}
Restrict the sum to the heat window
\(\mathcal I_s\) in \eqref{eq:V85-heat-window}.  On that window,
\Cref{lem:V85-window-weights} gives
\(p_{j,J}\ge cJ/j\), the heat factor is bounded below by a universal
positive constant, and \(J(J+1)\ge c/s\).  The window contains at
least \(c/\sqrt s\) integers, each with \(J\ge c/\sqrt s\).
Equivalently,
\[
 \sum_{J\in\mathcal I_s}
 \frac{J}{j}J(J+1)
 \ge
 \frac c j
 \sum_{J\in\mathcal I_s}J^3
 \ge\frac{c}{js^2}.
\]
\end{proof}

\begin{theorem}[Tilted coherent products refute the uniform vector
gate]
\label{thm:V90-coherent-counterexample}
For integers \(j\) tending to infinity, set
\begin{equation}
 s_j:=j^{-3/4},
 \qquad
 h_j:=j^{-1/2}.
\label{eq:V90-counterexample-scales}
\end{equation}
Then the vectors \(x_{j,h_j}\) and \(y_{j,h_j}\) are spin coherent
states, \(s_jj\to\infty\), and there is a universal \(c>0\) such that
\begin{align}
 \mathcal V_{s_j,j}
 &\ge
 \langle\psi_{j,h_j},
 e^{-s_jL^2}L_3\psi_{j,h_j}\rangle
 \ge c,
\label{eq:V90-vector-counterexample}\\
 s_j(2j+1)\mathcal V_{s_j,j}
 &\longrightarrow\infty.
\label{eq:V90-vector-gate-divergence}
\end{align}
Consequently the uniform estimate
\eqref{eq:V89-vector-gate} is false.
\end{theorem}

\begin{proof}
The normalized spin-\(\frac12\) vectors underlying
\eqref{eq:V90-tilted-x}--\eqref{eq:V90-tilted-y} have Bloch directions
\[
 \mathbf n_\pm(h)
 =
 \bigl(\pm\operatorname{sech}(h/2),0,\tanh(h/2)\bigr).
\]
Their symmetric powers are therefore spin-\(j\) coherent states.

The elementary inequality
\(\log\cosh u\le u^2/2\) gives
\[
 N_j(h_j)
 \le
 \exp\!\left(\frac{jh_j^2}{2}\right)
 =e^{1/2}.
\]
For all sufficiently large \(j\), the hypotheses of
\Cref{lem:V90-heat-susceptibility} hold because
\(j s_j=j^{1/4}\to\infty\).  Hence
\Cref{lem:V90-tilt-derivative,lem:V90-heat-susceptibility} gives
\[
 \langle\psi_{j,h_j},
 e^{-s_jL^2}L_3\psi_{j,h_j}\rangle
 \ge
 \frac{c h_j}{j s_j^2}
 =
 c.
\]
Finally
\(s_j(2j+1)\asymp j^{1/4}\), proving
\eqref{eq:V90-vector-gate-divergence}.
\end{proof}

\begin{corollary}[The complete local signed joining card is false]
\label{cor:V90-local-joining-counterexample}
For the sequence \eqref{eq:V90-counterexample-scales},
\begin{equation}
 \boxed{
 \omega_\otimes(\mathsf J_{s_j,j})
 \ge c j^{-3/4}.}
\label{eq:V90-local-joining-lower}
\end{equation}
In particular,
\begin{equation}
 (2j+1)\omega_\otimes(\mathsf J_{s_j,j})
 \longrightarrow\infty,
\label{eq:V90-local-card-divergence}
\end{equation}
so the uniform local signed card
\eqref{eq:V89-desired-signed-card} is false.
\end{corollary}

\begin{proof}
The lower inequality in \eqref{eq:V89-vector-equivalence} and
\Cref{thm:V90-coherent-counterexample} give
\[
 \omega_\otimes(\mathsf J_{s_j,j})
 \ge
 2s_je^{-2s_j}c-\frac C{2j+1}.
\]
Since \(s_j=j^{-3/4}\gg j^{-1}\), the right side is at least
\(c'j^{-3/4}\) for all sufficiently large \(j\).  Equivalently, in
the first-order factorization
\eqref{eq:V89-linear-factorization}, choose the spin-one row coherent
state polarized opposite to the positive third component in
\eqref{eq:V90-vector-counterexample}; the quadratic error is
\(O(j^{-1})\) by \Cref{thm:V89-quadratic-target}.
\end{proof}

\begin{assessment}[Mechanism of the failure]
\label{ass:V90-counterexample-mechanism}
At \(h=0\), the antipodal coherent product has zero signed vector
moment, but its heat-window components
\(J\asymp s^{-1/2}\) have transverse magnetic variance
\(J(J+1)/2\).  A common real tilt \(h=j^{-1/2}\) costs only the bounded
normalization \(N_j(h)\le e^{1/2}\), while converting that variance
into a first moment of size
\[
 \frac{h}{j s^2}.
\]
At \(s=j^{-3/4}\) this is order one.  Cancellation on the exactly
antipodal state is therefore not stable uniformly under coherent
tilts.  This is a raw signed complete-channel effect; no separate
positive output-channel envelope entered the proof.
\end{assessment}

\begin{firewall}[What the counterexample does and does not refute]
\label{fire:V90-counterexample-scope}
\Cref{cor:V90-local-joining-counterexample} refutes the use of a
uniform \(O(j^{-1})\) signed radius for one isolated incidence-three
joining coordinate.  It does not prove that the complete
first-joining-position sum has the same sign, nor that the
partition-state source estimate is false.  The final-output paid
linear gate \eqref{eq:V89-paid-linear-gate} is not decided here,
because its \(K\)-dependent multiplier changes the three-channel
combination.  It cannot, however, repair the unweighted local card
which the source compiler also requires.
\end{firewall}

\section{Post-V90 status and upstream acquisition order}
\label{sec:V90-status}

\begin{theorem}[Exact status after the coherent-tilt test]
\label{thm:V90-status}
Relative to V1--V89, V90 proves:
\begin{enumerate}[label=\textup{(V90.\arabic*)},leftmargin=6.5em]
\item the compulsory SM/NS companion audit was performed on stable,
      hash-identified active notes, with only compatible proof-order
      lessons transferred;
\item rotational covariance reduces the arbitrary vector heat radius
      exactly to the signed first moment of nonnegative
      total-magnetic column heat masses;
\item the antipodal coherent pair has a heat-window transverse
      susceptibility bounded below by \(c/(js^2)\);
\item the normalized coherent tilt
      \(h=j^{-1/2}\), \(s=j^{-3/4}\) has vector heat moment bounded
      below by a positive universal constant;
\item consequently the V89 vector gate and the isolated complete
      local signed joining card are false; and
\item any surviving cancellation must be sought upstream of the
      isolated joining coordinate.
\end{enumerate}
\end{theorem}

\begin{proof}
Item \textup{(V90.1)} is
\Cref{ass:V90-companion-identity,prop:V90-companion-transfer}.
Item \textup{(V90.2)} is
\Cref{lem:V90-column-Jacobi,thm:V90-vector-column-scalarization}.
Item \textup{(V90.3)} is
\Cref{lem:V90-base-projection,lem:V90-tilt-derivative,
lem:V90-heat-susceptibility}.
Item \textup{(V90.4)} is
\Cref{thm:V90-coherent-counterexample}.
Item \textup{(V90.5)} is
\Cref{cor:V90-local-joining-counterexample}.
Item \textup{(V90.6)} follows from
\Cref{fire:V90-counterexample-scope} and the exact alternative in
\Cref{prog:V89-next}.
\end{proof}

\begin{programme}[Mandate for V91]
\label{prog:V90-next}
\begin{enumerate}[label=\textup{(\arabic*)},leftmargin=4em]
\item Return to the exact V61 operator order and write the complete
      signed sum over first-joining positions before any local
      product-state supremum or physical absolute value is taken.
\item Insert the coherent-tilt family
      \eqref{eq:V90-tilted-x}--\eqref{eq:V90-counterexample-scales}
      into one joining coordinate and compute the signs of every term
      generated at adjacent first-joining positions.  Determine
      whether the local susceptibility is a telescoping range term or
      survives in the coupled packet.
\item Following the applicable SM audit lesson, separate the algebraic
      range of the first-position difference map from its cokernel.
      A positive residual may be claimed only on the cokernel after
      the complete range cancellation is performed.
\item Following the applicable NS audit lesson, keep the heat-window
      first moment correlated with its actual joining-position
      coefficient.  Do not sum isolated local radii or replace the
      signed packet by a weaker positive heat norm.
\item If the coherent tilt survives the first-position sum, extend it
      to the smallest admissible multi-coordinate row-product state
      and record an exact global-source counterexample.  Then move
      further upstream to partition states rather than weakening the
      target.
\item If it cancels, prove the cancellation as an operator identity
      and quantify the remaining boundary terms at both unweighted
      and paid locations before returning to the V79 routing table.
\item Keep the outside-prefix, directed-cross, cut-internal, and
      nonadjacent-overlap cells distinct.  No companion audit is due
      at V91; the next compulsory audit is V95.
\item At V100, freeze the current active note with suffix
      \texttt{\_f2} and restart at a compact, contextual V1 as already
      required.
\end{enumerate}
\end{programme}

\begin{firewall}[Claim boundary after V90]
\label{fire:V90-boundary}
V90 disproves one proposed local signed estimate; it does not prove
the complete source bound.  The first-joining-position cancellation,
partition-state alternative, outside-prefix routing, nonadjacent
overlap carrier, remaining connectivity families, global quartic and
higher MTA, WUFC, CPM, continuum construction, and positive spectral
gap remain open.  The full four-dimensional quantum Yang--Mills
existence and mass-gap problem remains open.
\end{firewall}

\paragraph{Parent-version record.}
The complete V89 mathematical body was retained before this chapter
was appended.  V89 had \(42450\) source records, \(1453371\) bytes,
and SHA256
\[
 \texttt{0631bcf94afcd677abc65da5b907d546a4dee6991902896c270ed3a5699de703}.
\]
The required companion audit was performed at V90.

% =============================================================================
% BEGIN APPENDED COMPACT VERSION V91 MATERIAL
% =============================================================================

\clearpage
\chapter{Mesoscopic prefixes transmit the coherent tilt}
\label{chap:V91}

\section{Context: exhaust the first-position telescope}
\label{sec:V91-context}

V90 gives a coherent-product sequence on which one isolated
incidence-three joining covariance has signed radius
\(\gtrsim j^{-3/4}\), rather than \(O(j^{-1})\).  The mandated upstream
question is whether the sum over possible first-joining positions
removes that susceptibility.  V77 has already identified this sum
exactly as a two-endpoint composite covariance.  The first result
below records the consequence which matters for the new
counterexample: a singleton crossing bank has no adjacent
first-position term with which to cancel.

\begin{proposition}[A singleton crossing bank survives the complete
first-position sum]
\label{prop:V91-singleton-survives}
Fix a two-block prefix partition \(\rho=C|D\).  If the residual
crossing bank is \(B=\{e\}\), then
\begin{equation}
 \mathfrak J_{\{e\}}^\rho
 =
 T_e-R_e^\rho
 =
 \Delta_e^\rho.
\label{eq:V91-singleton-telescope}
\end{equation}
This is both the unique V61 first-joining term and the complete V77
first-position sum.  Adding coordinates internal to \(C\) or \(D\)
only tensors the common complete moment from V80 and cannot create a
second signed crossing term.
\end{proposition}

\begin{proof}
Apply \eqref{eq:V77-bank-telescope} with \(m=1\).  The sole summand in
\eqref{eq:V77-first-join-packet} is
\(\Delta_e^\rho\), proving \eqref{eq:V91-singleton-telescope}.
For internal coordinates, use the exact conditional factorization
\eqref{eq:V80-conditional-factor}; their common factor multiplies
\(\mathfrak J_{\{e\}}^\rho\) and does not alter the endpoint
difference.
\end{proof}

\begin{corollary}[First-position resummation cannot repair the V90
local card uniformly]
\label{cor:V91-first-position-no-repair}
There is no uniform \(O(j^{-1})\) product-state card for the whole
first-position sum over an arbitrary crossing bank.  Indeed, the
one-coordinate bank in
\Cref{cor:V90-local-joining-counterexample} is an admissible singleton
instance of \eqref{eq:V91-singleton-telescope}.
\end{corollary}

\begin{proof}
A bound uniform over finite banks must hold when the bank has one
coordinate.  In that case the summed operator is the local signed
joining difference itself by
\Cref{prop:V91-singleton-survives}, and V90 gives the contradictory
lower bound \eqref{eq:V90-local-card-divergence}.
\end{proof}

\section{A mesoscopic binary prefix is asymptotically transparent}
\label{sec:V91-mesoscopic-prefix}

On \(V_k\otimes V_k\), put
\begin{align}
 \mathsf T_{s,k}
 &:=
 e^{-s(J^{(1)}+J^{(2)})^2},
\label{eq:V91-prefix-heat}\\
 \mathsf K_{s,k}
 &:=
 \mathsf T_{s,k}
 -
 e^{-2sk(k+1)}I.
\label{eq:V91-prefix-covariance}
\end{align}
Let
\begin{equation}
 \eta_k
 :=
 |k,k\rangle\otimes|k,-k\rangle,
 \qquad
 \beta_{s,k}
 :=
 \langle\eta_k,\mathsf K_{s,k}\eta_k\rangle.
\label{eq:V91-prefix-test}
\end{equation}

\begin{lemma}[Mesoscopic prefix transmission]
\label{lem:V91-mesoscopic-transmission}
If \(s\downarrow0\), \(sk\to0\), and \(sk^2\to\infty\), then
\begin{equation}
 \beta_{s,k}\longrightarrow1.
\label{eq:V91-prefix-tends-one}
\end{equation}
In particular, the connected binary prefix supplies no vanishing
factor on this product vector.
\end{lemma}

\begin{proof}
On \(\eta_k\),
\begin{equation}
 \langle (J^{(1)}+J^{(2)})^2\rangle
 =
 2k(k+1)-2k^2
 =
 2k.
\label{eq:V91-prefix-mean-Casimir}
\end{equation}
Since \(x\mapsto e^{-sx}\) is convex, Jensen's inequality gives
\[
 \langle\eta_k,\mathsf T_{s,k}\eta_k\rangle
 \ge e^{-2sk}\longrightarrow1.
\]
The heat operator is a contraction, so this expectation is squeezed
to one.  Meanwhile
\[
 e^{-2sk(k+1)}\longrightarrow0
\]
because \(sk^2\to\infty\).  Subtraction proves
\eqref{eq:V91-prefix-tends-one}.
\end{proof}

\begin{remark}[This is the V87 mesoscopic regime]
\label{rem:V91-mesoscopic-not-new}
\Cref{lem:V91-mesoscopic-transmission} is the precise scalar input
from \Cref{thm:V87-mesoscopic-countertest}, restated here only to fix
the operator and test vector used in the new three-coordinate packet.
The new result will be its coupling to the V90 coherent tilt.
\end{remark}

\section{The exact mesoscopic-prefix/coherent-join packet}
\label{sec:V91-three-coordinate-packet}

Fix once and for all
\begin{equation}
 \frac12<\alpha<\frac23.
\label{eq:V91-alpha-range}
\end{equation}
For integers \(j\to\infty\), retain
\[
 s_j=j^{-3/4},
 \qquad
 h_j=j^{-1/2},
\]
and choose a half-integer \(k_j\asymp s_j^{-\alpha}\).  Then
\begin{align}
 s_jk_j&\longrightarrow0,
 &
 s_jk_j^2&\longrightarrow\infty,
\label{eq:V91-mesoscopic-scales}\\
 \frac{k_j^2}{j^2}&\longrightarrow0,
 &
 \frac{k_j^4}{j^2}&\longrightarrow0.
\label{eq:V91-prefix-stock-small}
\end{align}
The last limit uses \(\alpha<2/3\).

Use three independent coordinates \(p<q<e\).  At \(p\), rows \(1,2\)
carry \(k_j,k_j\); at \(q\), rows \(3,4\) carry \(k_j,k_j\); and at
\(e\), rows \(1,2,3,4\) carry \(j,j,1,0\).  Thus \(p,q\) form the
proper prefix \(12|34\), while \(e\) is its unique first full join and
there is no suffix.

Let \(|1,-1\rangle\) be the third-axis lowest-weight vector in the
spin-one row and put
\begin{equation}
 \Theta_j
 :=
 \psi_{j,h_j}\otimes|1,-1\rangle,
 \qquad
 \varepsilon_j
 :=
 \langle\Theta_j,\mathsf J_{s_j,j}\Theta_j\rangle.
\label{eq:V91-joining-test}
\end{equation}

\begin{lemma}[The coherent joining expectation is positive of order
\(s_j\)]
\label{lem:V91-coherent-joining-lower}
For all sufficiently large \(j\),
\begin{equation}
 \varepsilon_j\ge c s_j.
\label{eq:V91-coherent-joining-lower}
\end{equation}
\end{lemma}

\begin{proof}
The spin-one vector has
\(\langle\mathbf S\rangle=-\mathbf e_3\).  Hence
\eqref{eq:V89-linear-factorization} and
\eqref{eq:V90-vector-counterexample} give a first-order expectation
at least \(2cs_je^{-2s_j}\).  The quadratic remainder has absolute
expectation at most \(C/(2j+1)\) by
\Cref{thm:V89-quadratic-target}.  Since
\(j^{-1}=o(s_j)\), the reverse triangle inequality proves
\eqref{eq:V91-coherent-joining-lower}.
\end{proof}

Define the exact raw no-suffix \(S_{22}\) operator
\begin{equation}
 \mathsf Q_j^{\rm raw}
 :=
 \mathsf K_{s_j,k_j}^{(p)}
 \otimes
 \mathsf K_{s_j,k_j}^{(q)}
 \otimes
 \mathsf J_{s_j,j}^{(e)}
\label{eq:V91-raw-packet}
\end{equation}
and the coordinate-factorized vector
\begin{equation}
 \Omega_j
 :=
 \eta_{k_j}^{(p)}
 \otimes
 \eta_{k_j}^{(q)}
 \otimes
 \Theta_j^{(e)}.
\label{eq:V91-packet-vector}
\end{equation}
After regrouping coordinate factors by rows, \(\Omega_j\) is a unit
product vector across all four input rows.

\begin{theorem}[Mesoscopic prefixes transmit the coherent joining
obstruction]
\label{thm:V91-raw-packet-lower}
The operator \(\mathsf Q_j^{\rm raw}\) is the exact V61
\(12|34\) first-joining summand for this no-suffix support, and
\begin{equation}
 \boxed{
 \langle\Omega_j,\mathsf Q_j^{\rm raw}\Omega_j\rangle
 =
 \beta_{s_j,k_j}^{\,2}\varepsilon_j
 \ge c s_j.}
\label{eq:V91-raw-packet-lower}
\end{equation}
Thus neither restoration of the two mandatory prefixes nor summation
over possible first positions supplies a vanishing factor on this
family.
\end{theorem}

\begin{proof}
The prefix-join factorization
\eqref{eq:V61-prefix-factorization} gives the two binary connected
operators at \(p,q\).  The V65 double-replica identity gives
\(\mathsf J_{s_j,j}\) as the joining letter at \(e\), exactly as in
\eqref{eq:V86-exact-S22-summand}.  There is no later coordinate.
Since \(e\) is the only crossing coordinate after the fixed prefix,
\Cref{prop:V91-singleton-survives} also shows that this is the complete
first-position sum for the packet.

Independence of the coordinate factors and the tensor form of
\(\Omega_j\) give the displayed equality.  By
\eqref{eq:V91-mesoscopic-scales} and
\Cref{lem:V91-mesoscopic-transmission},
\(\beta_{s_j,k_j}\to1\); combine this with
\Cref{lem:V91-coherent-joining-lower}.
\end{proof}

\begin{firewall}[Why the V86 minimal packet does not contradict V91]
\label{fire:V91-v86-distinction}
The V86 test used fixed spin-\(\frac12\) prefix coordinates, whose
connected expectations are \(O(s)\) and together supply \(s^2\).
Here the prefixes obey
\(s_jk_j^2\to\infty\) but \(s_jk_j\to0\); their connected
expectations tend to one.  The representation assignments therefore
belong to different heat regimes, and
\Cref{thm:V86-minimal-packet} makes no uniform-capacity assertion
covering the present family.
\end{firewall}

\section{Going upstream of the first-connectivity split}
\label{sec:V91-full-connected}

Write \(m_r(i):=k_r(\{i\})\) for a singleton local moment, and
suppress identity factors on inactive rows.  At \(p\), define the
complete two-row moment
\begin{equation}
 M_p^{12}
 :=
 k_p(12)+m_p(1)\otimes m_p(2).
\label{eq:V91-complete-p-moment}
\end{equation}
Let \(\mathcal K_j^{\rm conn}\) be the complete connected four-row
carrier on the three-coordinate support \(\{p,q,e\}\), before
first-connectivity is split by prefix state or first position.

\begin{theorem}[Exact unconditioned connected-carrier expansion]
\label{thm:V91-full-connected-expansion}
For the activity pattern in
\Cref{sec:V91-three-coordinate-packet},
\begin{align}
 \mathcal K_j^{\rm conn}
 &=
 k_q(34)
 \otimes
 \Bigl[
 M_p^{12}\otimes k_e(123)
\notag\\
 &\hspace{5.5em}
 {}+
 k_p(12)\otimes k_e(13)\otimes m_e(2)
 +
 k_p(12)\otimes k_e(23)\otimes m_e(1)
 \Bigr].
\label{eq:V91-full-connected-expansion}
\end{align}
There are no omitted connected local-partition words.
\end{theorem}

\begin{proof}
Apply the linked-word identity \eqref{eq:V61-linked-word}.  Coordinate
\(q\) is the only coordinate containing row \(4\).  Therefore a
globally connected word must use the block \(34\) at \(q\), giving
\(k_q(34)\).

It remains to join rows \(1,2,3\) using \(p,e\).  If the local
partition at \(e\) contains the block \(123\), either partition at
\(p\) is allowed, and their sum is \(M_p^{12}\).  If the partition at
\(e\) is \(13|2\) or \(23|1\), the block \(12\) is forced at \(p\).
The remaining local partitions at \(e\), namely \(12|3\) and the
discrete partition, do not connect row \(3\) to rows \(1,2\).
These cases are disjoint and give exactly
\eqref{eq:V91-full-connected-expansion}.
\end{proof}

At the joining coordinate, let \(M_e(A)\) denote the complete local
moment on the indicated active row set.  The joining covariance is
\[
 \mathsf J_{s_j,j}
 =
 M_e(123)-M_e(12)\otimes m_e(3).
\]
The ordinary three-row moment--cumulant identity therefore gives
\begin{align}
 k_e(123)
 &=
 \mathsf J_{s_j,j}
 -
 M_e(13)\otimes m_e(2)
 -
 M_e(23)\otimes m_e(1)
\notag\\
 &\quad+
 2m_e(1)\otimes m_e(2)\otimes m_e(3).
\label{eq:V91-triple-vs-joining}
\end{align}
Here
\begin{equation}
 m_e(1)=m_e(2)=q_{s_j,j}I,
 \qquad
 q_{s_j,j}:=e^{-s_j\,j(j+1)},
 \qquad
 m_e(3)=e^{-2s_j}I.
\label{eq:V91-singleton-suppression}
\end{equation}

\begin{lemma}[All competing connected words are exponentially
suppressed on the counterpacket]
\label{lem:V91-competing-words-small}
On the joining-coordinate factor \(\Theta_j\),
\begin{align}
 \left|
 \langle k_e(123)\rangle_{\Theta_j}
 -
 \varepsilon_j
 \right|
 &\le
 2q_{s_j,j}+2q_{s_j,j}^2,
\label{eq:V91-triple-error}\\
 \left|
 \langle k_e(13)\otimes m_e(2)\rangle_{\Theta_j}
 \right|
 +
 \left|
 \langle k_e(23)\otimes m_e(1)\rangle_{\Theta_j}
 \right|
 &\le
 4q_{s_j,j}.
\label{eq:V91-pair-partition-error}
\end{align}
Moreover,
\begin{equation}
 q_{s_j,j}
 \le e^{-c j^{5/4}}
 =
 o(s_j).
\label{eq:V91-singleton-exponential}
\end{equation}
\end{lemma}

\begin{proof}
Every complete local Wilson moment is a contraction.  Hence
\eqref{eq:V91-triple-vs-joining} and
\eqref{eq:V91-singleton-suppression} give
\eqref{eq:V91-triple-error}.  A binary cumulant is the difference of
two complete contractions and has norm at most two; multiplication
by the opposite large-spin singleton gives
\eqref{eq:V91-pair-partition-error}.  Finally,
\[
 s_j\,j(j+1)\ge c j^{5/4},
\]
which proves \eqref{eq:V91-singleton-exponential}.
\end{proof}

\begin{theorem}[The fully unconditioned connected carrier still has
order \(s_j\)]
\label{thm:V91-full-connected-lower}
On the same row-product vector \(\Omega_j\),
\begin{equation}
 \boxed{
 \langle\Omega_j,
 \mathcal K_j^{\rm conn}\Omega_j\rangle
 \ge c s_j}
\label{eq:V91-full-connected-lower}
\end{equation}
for all sufficiently large \(j\).
\end{theorem}

\begin{proof}
Let
\[
 \tau_{s,k}
 :=
 \langle\eta_k,\mathsf T_{s,k}\eta_k\rangle.
\]
The proof of \Cref{lem:V91-mesoscopic-transmission} gives
\(\tau_{s_j,k_j}\to1\), while
\(\beta_{s_j,k_j}\to1\).  Evaluate
\eqref{eq:V91-full-connected-expansion} on \(\Omega_j\).
The main term is
\[
 \beta_{s_j,k_j}\,
 \tau_{s_j,k_j}\,
 \langle k_e(123)\rangle_{\Theta_j}.
\]
By \Cref{lem:V91-coherent-joining-lower,
lem:V91-competing-words-small}, it is at least
\(cs_j-o(s_j)\).  The sum of the other two terms has magnitude at
most \(4q_{s_j,j}\), because the two prefix expectations have
magnitude at most one.  Apply
\eqref{eq:V91-singleton-exponential} once more.
\end{proof}

\section{Failure of the crossing-only square-tail normalization}
\label{sec:V91-crossing-tail-fails}

Put \(a_r:=1+r(r+1)\).  On the three-coordinate representation
assignment above, the row-\(1\) mass and its crossing numerator are
\begin{equation}
 S_1=a_j+a_{k_j},
 \qquad
 H_{1,C_1}=a_ja_1=3a_j.
\label{eq:V91-counterpacket-stocks}
\end{equation}
Indeed, \(p\) is internal to \(12\), \(q\) is supported on \(34\),
and \(e\) is the sole coordinate where row \(1\) meets the opposite
block.

\begin{theorem}[The response-only \(S_{22}\) crossing square tail is
false even before first-connectivity splitting]
\label{thm:V91-crossing-square-counterexample}
There is no universal \(C\) such that the fully connected carrier in
this family satisfies
\begin{equation}
 \omega_\otimes(\mathcal K_j^{\rm conn})
 \le
 C\frac{H_{1,C_1}}{s_jS_1^2}.
\label{eq:V91-false-crossing-square}
\end{equation}
More precisely,
\begin{equation}
 \frac{
 \langle\Omega_j,\mathcal K_j^{\rm conn}\Omega_j\rangle}
 {H_{1,C_1}/(s_jS_1^2)}
 \ge
 c s_j^2j^2
 =
 c j^{1/2}
 \longrightarrow\infty.
\label{eq:V91-crossing-square-divergence}
\end{equation}
\end{theorem}

\begin{proof}
By \eqref{eq:V91-prefix-stock-small},
\(a_{k_j}=o(a_j)\).  Therefore
\[
 \frac{H_{1,C_1}}{s_jS_1^2}
 \asymp
 \frac1{s_j\,j^2}.
\]
The numerator in
\eqref{eq:V91-crossing-square-divergence} is at least \(cs_j\) by
\Cref{thm:V91-full-connected-lower}.  Their ratio is bounded below
by \(cs_j^2j^2=cj^{1/2}\).  Since \(\Omega_j\) is a row-product
vector, its expectation is a lower bound for
\(\omega_\otimes(\mathcal K_j^{\rm conn})\).
\end{proof}

\begin{firewall}[The no-go concerns one currency, not quartic MTA]
\label{fire:V91-crossing-currency-only}
\Cref{thm:V91-crossing-square-counterexample} proves that the
same-side mass \(a_j^2\) cannot be discarded and replaced solely by
the much smaller crossing numerator \(3a_j\) while the full
row-square \(S_1^2\) remains in the denominator.  It does not refute
a bound using the same-side numerator \(H_{12}\), a sum over all
spanning-tree currencies, or the full quartic MTA.  The next
regression must test those complete alternatives on this exact
counterpacket instead of weakening the failed estimate.
\end{firewall}

\section{The legal marked-tree currency absorbs the counterpacket}
\label{sec:V91-legal-rerouting}

Let \(\sigma\) be the path with edges
\[
 12,\qquad13,\qquad34.
\]
Mark these edges respectively by \(p,e,q\), and denote this legal
marking by \(\ell\).  The remaining row masses are
\begin{equation}
 S_2=a_j+a_{k_j},
 \qquad
 S_3=a_{k_j}+3,
 \qquad
 S_4=a_{k_j}.
\label{eq:V91-all-row-stocks}
\end{equation}

\begin{proposition}[Exact rerouted path weight]
\label{prop:V91-rerouted-path}
The marked-tree weight of \((\sigma,\ell)\) is
\begin{equation}
 \boxed{
 \mathcal W_{\sigma,\ell}
 =
 \frac{3a_j a_{k_j}^{\,4}}
 {(a_j+a_{k_j})(a_{k_j}+3)}.}
\label{eq:V91-rerouted-path-weight}
\end{equation}
Consequently,
\begin{equation}
 \mathcal W_{\sigma,\ell}
 \ge c a_{k_j}^{\,3}
\label{eq:V91-rerouted-path-lower}
\end{equation}
for all sufficiently large \(j\).
\end{proposition}

\begin{proof}
The three marked edge factors are
\[
 \frac{a_{k_j}}{S_1}\frac{a_{k_j}}{S_2},
 \qquad
 \frac{a_j}{S_1}\frac{3}{S_3},
 \qquad
 \frac{a_{k_j}}{S_3}\frac{a_{k_j}}{S_4}.
\]
Multiply their product by \(\prod_{i=1}^4S_i\), as in
\eqref{eq:V1-marked-tree-weight}, and use
\eqref{eq:V91-counterpacket-stocks}--%
\eqref{eq:V91-all-row-stocks}.  This gives
\eqref{eq:V91-rerouted-path-weight}.  The limits in
\eqref{eq:V91-prefix-stock-small} imply
\[
 \frac{a_j}{a_j+a_{k_j}}\ge\frac12,
 \qquad
 \frac{a_{k_j}}{a_{k_j}+3}\ge\frac12
\]
eventually, proving \eqref{eq:V91-rerouted-path-lower}.
\end{proof}

\begin{corollary}[The complete carrier passes this legal MTA test]
\label{cor:V91-counterpacket-MTA-paid}
Uniformly in \(j\),
\begin{equation}
 \omega_\otimes(\mathcal K_j^{\rm conn})
 \le28.
\label{eq:V91-connected-norm-bound}
\end{equation}
Hence, after enlarging a universal constant,
\begin{equation}
 \boxed{
 \omega_\otimes(\mathcal K_j^{\rm conn})
 \le
 C\mathcal W_{\sigma,\ell}.}
\label{eq:V91-counterpacket-MTA-paid}
\end{equation}
Thus the V91 family refutes the crossing-only square tail but is not a
counterexample to the complete marked-tree right side.
\end{corollary}

\begin{proof}
Complete moments and singleton moments are contractions, a binary
cumulant has norm at most two, and the ordinary three-row cumulant
formula gives
\(\|k_e(123)\|\le1+3+2=6\).  Therefore the bracket in
\eqref{eq:V91-full-connected-expansion} has norm at most
\[
 6+2\cdot(2\cdot2)=14.
\]
Multiplication by \(\|k_q(34)\|\le2\) proves
\eqref{eq:V91-connected-norm-bound}.  The lower bound
\eqref{eq:V91-rerouted-path-lower} and \(a_{k_j}\ge1\) prove
\eqref{eq:V91-counterpacket-MTA-paid}.
\end{proof}

\begin{assessment}[The regression selects mark routing]
\label{ass:V91-routing-selection}
The failed currency
\(H_{1,C_1}/(sS_1^2)\) insists that the joining coordinate alone mark
the repeated large row.  The legal path
\eqref{eq:V91-rerouted-path-weight} instead keeps the two actual
prefix incidences as the \(12\) and \(34\) edge marks.  This is the
same structural repair seen in V15, now on the exact V90 coherent
counterpacket and after the complete cumulant has been restored.  The
remaining problem is a uniform bounded-incidence routing theorem, not
a stronger local joining card.
\end{assessment}

\section{Post-V91 status and next acquisition order}
\label{sec:V91-status}

\begin{theorem}[Exact status after the first-position and cumulant
regressions]
\label{thm:V91-status}
Relative to V1--V90, V91 proves:
\begin{enumerate}[label=\textup{(V91.\arabic*)},leftmargin=6.5em]
\item a singleton crossing bank survives the complete V77
      first-position telescope exactly, so that resummation alone
      cannot repair the V90 local card;
\item equal-spin mesoscopic binary prefixes have connected
      expectations tending to one on anti-aligned product vectors;
\item two such mandatory prefixes transmit the V90 tilted-coherent
      joining expectation, giving an exact raw \(S_{22}\) packet of
      order \(s_j\);
\item the full unconditioned three-coordinate connected carrier has
      the exact three-term expansion
      \eqref{eq:V91-full-connected-expansion}, and every term which
      might cancel the joining susceptibility is exponentially
      suppressed;
\item the crossing-only square-tail normalization
      \(H_{1,C_1}/(sS_1^2)\) fails on this full connected carrier by a
      factor at least \(cj^{1/2}\); and
\item the same carrier is controlled by one explicit legal marked
      path, so the counterpacket selects mark rerouting rather than
      refuting the complete marked-tree atlas.
\end{enumerate}
\end{theorem}

\begin{proof}
Item \textup{(V91.1)} is
\Cref{prop:V91-singleton-survives,
cor:V91-first-position-no-repair}.
Item \textup{(V91.2)} is
\Cref{lem:V91-mesoscopic-transmission}.
Item \textup{(V91.3)} is
\Cref{lem:V91-coherent-joining-lower,thm:V91-raw-packet-lower}.
Item \textup{(V91.4)} is
\Cref{thm:V91-full-connected-expansion,
lem:V91-competing-words-small,thm:V91-full-connected-lower}.
Item \textup{(V91.5)} is
\Cref{thm:V91-crossing-square-counterexample}.
Item \textup{(V91.6)} is
\Cref{prop:V91-rerouted-path,
cor:V91-counterpacket-MTA-paid}.
\end{proof}

\begin{programme}[Mandate for V92]
\label{prog:V91-next}
\begin{enumerate}[label=\textup{(\arabic*)},leftmargin=4em]
\item Abandon the false crossing-only response card
      \eqref{eq:V91-false-crossing-square}.  Start from the complete
      marked-tree sum and allow edges to be marked at the actual two
      prefix banks, as in
      \eqref{eq:V91-rerouted-path-weight}.
\item Insert the forced outside-prefix alternative
      \eqref{eq:V88-outside-prefix-large} into a scalar mark-routing
      table.  Treat foreign-prefix, directed-cross, cut-internal, and
      overlap incidences separately; do not rename them all as
      crossing stock.
\item Close at least one of those cells by assigning a legal
      comparison-tree edge and proving bounded incidence of the map
      from source packets to marked trees.  The operator carrier must
      remain assembled until its capacity bound is proved; only then
      may its scalar row mass choose a mark.
\item On a repeated old bank, retain the two occurrences in the exact
      V61 carrier.  Use
      \eqref{eq:V86-adjacent-joint-capacity} only after proving the
      intervening factors removable or contractive in the required
      operator order.
\item Carry the sole V2 payer separately at a prefix, joining, and
      suffix location.  A legal unpaid rerouting is not automatically
      a paid rerouting.
\item Test every proposed assignment on the V15 dual-shared
      regression, the V86 fixed-low-prefix packet, and the V91
      mesoscopic-prefix/coherent-join packet.
\item If no bounded-incidence routing exists, preserve the smallest
      conflicting family and move upstream to the full marked-tree
      sum before source-family splitting.  Do not weaken MTA to the
      false local normalization.
\item No companion audit is due at V92.  The next compulsory audit is
      V95.  At V100 freeze the active note with suffix
      \texttt{\_f2} and restart at a contextual, proof-indexed V1.
\end{enumerate}
\end{programme}

\begin{firewall}[Claim boundary after V91]
\label{fire:V91-boundary}
V91 repairs the interpretation of the V90 counterexample and proves
one exact legal rerouting for that family.  It does not prove the
uniform bounded-incidence routing required for all V79 cells, the
paid rerouting, the nonadjacent repeated-bank carrier, the remaining
first-connectivity families, global quartic or higher MTA, WUFC, CPM,
continuum construction, or a positive spectral gap.  The full
four-dimensional quantum Yang--Mills existence and mass-gap problem
remains open.
\end{firewall}

\paragraph{Parent-version record.}
The complete V90 mathematical body was retained before this chapter
was appended.  V90 had \(43047\) source records, \(1473459\) bytes,
and SHA256
\[
 \texttt{934ba585dffd8cbb86d6d24fbcd2d4953aaa1af565d80980f7d792e879ff5667}.
\]
No companion audit is due at V91.

% =============================================================================
% BEGIN APPENDED COMPACT VERSION V92 MATERIAL
% =============================================================================

\clearpage
\chapter{Exact mark expansion and contraction-sandwiched overlap}
\label{chap:V92}

\section{Context: route only after the carrier has been assembled}
\label{sec:V92-context}

V91 ruled out the response-only crossing-square estimate on an exact
connected \(S_{22}\) packet.  The same packet was nevertheless paid
by a legal path whose two same-side edges were marked at the actual
prefix coordinates.  Thus the failed local card must not be weakened:
the proof has to retain the complete marked-tree sum.

This chapter supplies the missing bookkeeping interface.  First, an
aggregate product of three genuine edge stocks is expanded
\emph{exactly} into legal coordinate markings.  For the assembled
\(S_{22}\) carrier, the map from the connected/replicated and
one-payer branches to those markings has a universal fibre bound.
This inserts the V88 outside-prefix alternative into a four-cell
routing table without merging foreign-prefix, directed-cross,
cut-internal, and overlap incidences.

Second, the V86 adjacent repeated-bank estimate is replaced by a
noncommutative sandwich inequality.  An intervening contraction is
not deleted: it is retained inside a Schatten estimate.  A weighted
pinched version identifies exactly what is still needed for the
physical positive envelope.  Finally, all possible payer locations
of the explicit V91 counterpacket are checked against its legal path
weight.

\section{A bank product is exactly a sum of legal markings}
\label{sec:V92-mark-expansion}

Retain the fixed-input notation
\[
 \Xi_r>0,\qquad
 a_{r,e}=1+C_2(R_{r,e}),\qquad 1\le r\le4.
\]
For a pair \(g=uv\) and a coordinate bank
\(E_g\subseteq\operatorname{supp}R_u\cap
\operatorname{supp}R_v\), put
\begin{equation}
 H_g(E_g):=\sum_{e\in E_g}a_{u,e}a_{v,e}.
\label{eq:V92-edge-bank-stock}
\end{equation}
If \(\tau\) is a labelled tree on \([4]\), let
\(\operatorname{Mark}(\tau;\mathbf E)\) be the set of functions
\(\ell\) satisfying \(\ell(g)\in E_g\) for every
\(g\in E(\tau)\).

\begin{proposition}[Exact bank-to-mark identity]
\label{prop:V92-bank-to-mark}
For every labelled quartic tree \(\tau\) and every family of finite
edge banks \(\mathbf E=(E_g)_{g\in E(\tau)}\),
\begin{equation}
 \boxed{
 \sum_{\ell\in\operatorname{Mark}(\tau;\mathbf E)}
 \mathcal W_{\tau,\ell}
 =
 \left(\prod_{r=1}^4\Xi_r\right)
 \prod_{uv\in E(\tau)}
 \frac{H_{uv}(E_{uv})}{\Xi_u\Xi_v}.}
\label{eq:V92-bank-to-mark}
\end{equation}
The map
\[
 (e_g)_{g\in E(\tau)}
 \longmapsto
 \bigl(\tau,\ell(g)=e_g\bigr)
\]
is a bijection.  In particular, no support-cardinality factor occurs.
\end{proposition}

\begin{proof}
Expand the finite product of the three sums
\[
 \prod_{uv\in E(\tau)}
 \left(\sum_{e\in E_{uv}}a_{u,e}a_{v,e}\right).
\]
A term is indexed by one and only one coordinate \(e_g\) for each
labelled edge \(g\).  After multiplication by
\(\prod_r\Xi_r\), its normalized factors are
\[
 \left(\prod_{r=1}^4\Xi_r\right)
 \prod_{uv\in E(\tau)}
 \frac{a_{u,e_{uv}}}{\Xi_u}
 \frac{a_{v,e_{uv}}}{\Xi_v},
\]
which is precisely \(\mathcal W_{\tau,\ell}\) from
\eqref{eq:V1-marked-tree-weight}.  Edge labels recover every
coordinate \(e_g\), proving bijectivity.
\end{proof}

\begin{firewall}[Algebraic repetition is not an analytic second use]
\label{fire:V92-mark-overlap}
Proposition~\ref{prop:V92-bank-to-mark} remains an algebraic identity
if the same physical coordinate belongs to two edge banks: the
corresponding legal marking uses that coordinate on two edges.
It does not prove that a source operator supplies two independent
first-moment estimates at that coordinate.  Such a diagonal term
belongs to the overlap cell and must be kept as one joint carrier.
\end{firewall}

For the cut \(C=\{1,2\}\), \(D=\{3,4\}\), and
\(i\in C,j\in D\), let \(\tau_{ij}\) be the path with edges
\[
 \{12,34,ij\}.
\]

\begin{corollary}[Exact \(S_{22}\) path expansion]
\label{cor:V92-S22-path-expansion}
For arbitrary edge banks \(E_C,E_D,E_{ij}\),
\begin{equation}
 \boxed{
 \sum_{\substack{p\in E_C,\ q\in E_D\\ e\in E_{ij}}}
 \mathcal W_{\tau_{ij},\ell_{pqe}}
 =
 \frac{
 H_{12}(E_C)H_{34}(E_D)H_{ij}(E_{ij})
 }{\Xi_i\Xi_j},}
\label{eq:V92-S22-path-expansion}
\end{equation}
where
\(\ell_{pqe}(12)=p\),
\(\ell_{pqe}(34)=q\), and
\(\ell_{pqe}(ij)=e\).
\end{corollary}

\begin{proof}
Apply \Cref{prop:V92-bank-to-mark}.  The four row denominators
belonging to the edges \(12\) and \(34\) cancel
\(\prod_r\Xi_r\), leaving only \(\Xi_i\Xi_j\).
\end{proof}

\begin{remark}[The V81 numerator is already the marked sum]
\label{rem:V92-V81-numerator}
The right side of \eqref{eq:V92-S22-path-expansion} is exactly the
common three-carrier numerator in
\eqref{eq:V81-common-tree-numerator}.  Thus no post hoc choice of one
large coordinate is needed.  The complete edge-stock product is
itself the sum over all legal coordinate markings of the selected
path.
\end{remark}

\section{Finite fibres for the assembled clean \(S_{22}\) carrier}
\label{sec:V92-finite-fibre}

The word \emph{assembled} is essential here.  The prefix join has
already been formed as in
\eqref{eq:V61-prefix-factorization}; the crossing bank has already
been telescoped as in the V77 composite covariance; and complete
nonpayer suffixes have not been expanded into partition words.

\begin{proposition}[Universal core routing fibre]
\label{prop:V92-core-routing-fibre}
Fix a deterministic V79 cell allocation, a cut \(12|34\), and a
crossing card \(ij\).  In the clean V80--V81 conditioned carrier,
map an assembled core branch to the marked path
\((\tau_{ij},\ell_{pqe})\) supplied by
\eqref{eq:V92-S22-path-expansion}.  Then every marked path has at most
\begin{equation}
 2\cdot2\cdot3=12
\label{eq:V92-core-fibre}
\end{equation}
preimages.  The factors count the connected/replicated choice on
the \(C\)-side, the corresponding choice on the \(D\)-side, and the
three mutually exclusive V2 payer locations
\[
 C,\qquad D,\qquad \times .
\]
The bound is independent of coordinate support and representations.
\end{proposition}

\begin{proof}
Once the marked tree is given, its unique cross edge determines
\((i,j)\); its two leaf edges determine the cut.  Their three marks
recover \(p,q,e\).  The deterministic cell priority recovers the
selected row-cell type before this fixed-allocation map is formed, so
that type is not a fibre label.  The only labels not visible in the
marked tree are the two binary connected/replicated choices and the
core payer location, of which there are at most \(2,2,3\), respectively.

There is no hidden choice of a first witnessing prefix coordinate by
\Cref{cor:V61-no-prefix-multiplicity}.  There is no choice of a first
crossing coordinate after the complete V77 telescope has been
assembled.  Hence the displayed list exhausts the fibre.
\end{proof}

\begin{firewall}[The suffix payer is a fourth, separate location]
\label{fire:V92-suffix-payer-separate}
A nonpayer complete suffix costs no fibre and is contractive by
\Cref{lem:V61-suffix-contraction}.  If the unique payer is at a suffix
coordinate, it is the distinct V61 factor \(D_t^{\rm pay}\), not a
fourth copy of one of the three core lines.  If an analytic suffix
estimate produces the same path numerator, adding that location
changes \(12\) to at most \(16\).  No such general suffix-payer
estimate is inferred merely from the finite-fibre count.
\end{firewall}

\begin{theorem}[Clean paired blocks route into the literal MTA sum]
\label{thm:V92-clean-paired-routing}
Assume the exact V81 three-carrier factorization for a clean
\(S_{22}\) source.  Suppose the selected \(C\)- and \(D\)-block
certificates have stocks
\[
 \widehat H_C=H_{12}(E_C),\qquad
 \widehat H_D=H_{34}(E_D),
\]
where each bank is the disjoint union of the genuine prefix and
internal banks used in its certificate.  Let \(E_{ij}\) be the bank
supporting the already assembled crossing card.  Then the sum of the
three core payer locations satisfies
\begin{equation}
 \boxed{
 \mathfrak b_{ij}^{\rm pay}
 \le
 C
 \sum_{\substack{p\in E_C,\ q\in E_D\\e\in E_{ij}}}
 \mathcal W_{\tau_{ij},\ell_{pqe}}.}
\label{eq:V92-clean-paired-routing}
\end{equation}
After summing the three possible paired certificate types
\(A,F,I\) on each side, each marked path is used at most
\begin{equation}
 3\cdot3\cdot12=108
\label{eq:V92-all-clean-fibre}
\end{equation}
times.
\end{theorem}

\begin{proof}
\Cref{lem:V81-three-payer-locations,thm:V81-paired-22} give, before
division by the four row masses,
\[
 \mathfrak b_{ij}^{\rm pay}
 \le
 C\frac{
 \widehat H_C\widehat H_DH_{ij}(E_{ij})
 }{\Xi_i\Xi_j}.
\]
Apply \eqref{eq:V92-S22-path-expansion}.  If a block certificate uses
a sum of disjoint genuine subbanks, distribute that sum first; their
union gives exactly \(E_C\) or \(E_D\), with no repeated incidence.
\Cref{prop:V92-core-routing-fibre} gives \(12\) for each fixed
allocation.  There are at most three paired certificate types
\(A,F,I\) on either side, giving
\eqref{eq:V92-all-clean-fibre}.
\end{proof}

\begin{corollary}[The clean foreign and cut-internal cells have legal
marks]
\label{cor:V92-FI-legal-marks}
Suppose a selected row \(k\) satisfies either
\[
 F_k\ge\kappa\Xi_k
 \qquad\hbox{or}\qquad
 I_k\ge\kappa\Xi_k,
\]
and the clean disjoint hypotheses of
\Cref{cor:V81-AFI-paired} hold.  If the opposite block has a paired
certificate, every core payer location is bounded by a literal sum of
legal path markings.  In the foreign case the \(kk^\circ\) edge is
marked in the reduced prefix response bank \(R\); in the internal
case it is marked in the disjoint union of the old prefix bank and
the connected internal bank.
\end{corollary}

\begin{proof}
For \(F_k\)-dominance use
\Cref{thm:V81-paid-foreign-block}; its genuine edge numerator is
\(H_R\), while the foreign bank supplies the scalar row debit and is
not falsely renamed as \(kk^\circ\) overlap.  For \(I_k\)-dominance
use \Cref{thm:V81-paid-internal-block}; its edge numerator is the sum
of the old-prefix and internal connected stocks.  These banks are
disjoint by hypothesis.  In either case
\Cref{thm:V92-clean-paired-routing} converts the aggregate numerator
to legal markings.
\end{proof}

\section{The V88 outside-prefix routing table}
\label{sec:V92-cell-table}

Let
\begin{equation}
 \delta_\eta:=1-\frac{\eta}{a_*}>0.
\label{eq:V92-delta}
\end{equation}
On the uncertified V88 endpoint,
\[
 P_k+U_k+F_k+D_k+I_k>\delta_\eta\Xi_k.
\]
The four balanced cells relevant after the earlier private and
unbalanced compilers have the following distinct routes.

\begin{center}
\small
\begin{tabular}{p{0.13\textwidth}p{0.25\textwidth}
                p{0.20\textwidth}p{0.29\textwidth}}
\toprule
cell & exact carrier location & legal edge & proved status\\
\midrule
\(F_k\) &
one-row scalar \(q_k(G_k)\) inside the actual prefix letter &
\(kk^\circ\), marked in the reduced prefix bank &
clean core branch closed by
\Cref{cor:V92-FI-legal-marks}\\
\(D_{k\to\ell}\) &
complete assembled crossing covariance &
\(k\ell\), marked in the selected residual bank &
direct matches and V82 pure common banks close; V83
mass-compatible opposite partners close; the V91 same-side-maximal
branch is not assigned a false cross-only card\\
\(I_k\) &
complete cut-internal moment, split into connected and replicated
branches &
\(kk^\circ\), marked in the actual internal or old-prefix bank &
clean core branch closed by
\Cref{cor:V92-FI-legal-marks}\\
overlap &
one identified old bank occurring twice in the exact source carrier &
the edge carried by that one bank, used once analytically &
the separator is treated in
\Cref{thm:V92-pinched-sandwich}; the final pinched one-occurrence
absolute card remains to be certified\\
\bottomrule
\end{tabular}
\end{center}

\begin{firewall}[Directed-cross is not synonymous with crossing-only]
\label{fire:V92-directed-not-cross-only}
The directed cell records an actual residual incidence
\(k\to\ell\).  In the same-side-maximal configuration of V91, the
largest analytic partner is \(k^\circ\), so its square-tail numerator
is parallel to the already present same-side edge.  The failed V84
cross-retained estimate cannot be restored by relabelling that
numerator.  The correct legal path may instead use both proper-prefix
marks, exactly as V91 demonstrated.
\end{firewall}

\begin{theorem}[Quantitative residual after clean cell routing]
\label{thm:V92-forced-cell-reduction}
Fix an uncertified V88 endpoint \(k\).  Suppose the previously proved
private and unbalanced branches have been removed at threshold
\(\delta_\eta/5\).  Apply
\Cref{cor:V92-FI-legal-marks} whenever a clean disjoint foreign or
cut-internal cell has mass at least
\(\delta_\eta\Xi_k/5\).  Then every remaining endpoint has one of:
\begin{enumerate}[label=\textup{(\roman*)},leftmargin=4em]
\item
\begin{equation}
 D_k>\frac{\delta_\eta}{5}\Xi_k;
\label{eq:V92-directed-residual}
\end{equation}
\item a foreign or cut-internal dominant bank which intersects a
selected source, payer, or transfer bank and hence belongs to the
overlap cell; or
\item a nonclean local representation state explicitly excluded from
the V80--V83 \(SU(2)\) bank compilers.
\end{enumerate}
Thus, within the clean \(S_{22}\) calculus, the V88
outside-prefix alternative reduces to directed-cross or overlap
routing.
\end{theorem}

\begin{proof}
If neither \textup{(ii)} nor \textup{(iii)} occurs, every dominant
\(F_k\)- or \(I_k\)-bank is clean and eligible for
\Cref{cor:V92-FI-legal-marks}.  On a surviving term,
\[
 P_k,\ U_k,\ F_k,\ I_k
 \le\frac{\delta_\eta}{5}\Xi_k.
\]
Subtract these four bounds from
\[
 P_k+U_k+F_k+D_k+I_k>\delta_\eta\Xi_k
\]
to obtain \eqref{eq:V92-directed-residual}.  If a selected bank is
not disjoint, the deterministic V79 incidence is retained but its
analytic use is classified as overlap, giving \textup{(ii)}.
\end{proof}

\section{A noncommutative sandwich for a repeated bank}
\label{sec:V92-sandwich}

For an operator \(X\) on the four-row input space, define the
product trace radius
\begin{equation}
 \omega_{\otimes,1}(X)
 :=
 \sup_{\rho=\rho_1\otimes\cdots\otimes\rho_4}
 \left\|\rho^{1/2}X\rho^{1/2}\right\|_1,
\label{eq:V92-product-trace-radius}
\end{equation}
where each \(\rho_r\) is a density matrix.  If \(X\ge0\), this is
exactly \(\kappa_\otimes(X)\), and for general \(X\) it dominates
\(\sup_\rho|\Tr(\rho X)|\).

\begin{lemma}[Contraction-sandwiched joint card]
\label{lem:V92-sandwiched-joint-card}
Let \(A,B\) be bounded self-adjoint operators and let
\(\|C\|\le L\).  Then
\begin{align}
 \left\|\rho^{1/2}ACB\rho^{1/2}\right\|_1
 &\le
 L\sqrt{\Tr(\rho A^2)\Tr(\rho B^2)}
\label{eq:V92-sandwich-CS}\\
 &\le
 \frac L2\left(
 \|A\|\Tr(\rho|A|)
 +
 \|B\|\Tr(\rho|B|)
 \right)
\label{eq:V92-sandwich-linear}
\end{align}
for every density matrix \(\rho\).  Consequently,
\begin{equation}
 \boxed{
 \omega_{\otimes,1}(ACB)
 \le
 \frac L2\left(
 \|A\|\kappa_\otimes(|A|)
 +
 \|B\|\kappa_\otimes(|B|)
 \right).}
\label{eq:V92-sandwich-capacity}
\end{equation}
\end{lemma}

\begin{proof}
Schatten H\"older gives
\begin{align*}
 \left\|\rho^{1/2}ACB\rho^{1/2}\right\|_1
 &\le
 \|\rho^{1/2}A\|_2\,\|C\|\,
 \|B\rho^{1/2}\|_2\\
 &=
 L\sqrt{\Tr(\rho A^2)\Tr(\rho B^2)}.
\end{align*}
Functional calculus gives
\[
 A^2\preccurlyeq\|A\||A|,
 \qquad
 B^2\preccurlyeq\|B\||B|.
\]
Use \(2\sqrt{xy}\le x+y\) to obtain
\eqref{eq:V92-sandwich-linear}, then take the supremum over product
density matrices.
\end{proof}

\begin{remark}[What has improved over adjacency]
\label{rem:V92-no-adjacency}
The operator \(C\) need not commute with \(A\) or \(B\), and it is not
removed from the word.  When \(C=I\),
\eqref{eq:V92-sandwich-capacity} recovers the V86 adjacent signed
card.  Thus noncommutativity of a uniformly bounded separator is not
by itself a second-moment obstruction at the product trace-radius
level.
\end{remark}

The physical proof uses weighted isotypic compressions.  The following
version keeps those compressions in their exact order.

\begin{theorem}[Weighted pinched sandwich]
\label{thm:V92-pinched-sandwich}
Let \((P_\nu)\) be mutually orthogonal projections, let \(c_\nu\ge0\),
and define the positive pinching map
\[
 \mathscr P_c(X):=\sum_\nu c_\nu P_\nu XP_\nu.
\]
For self-adjoint \(A,B\), \(\|C\|\le L\), and every density matrix
\(\rho\),
\begin{align}
 &\sum_\nu c_\nu
 \left\|
 \rho^{1/2}P_\nu ACBP_\nu\rho^{1/2}
 \right\|_1
\notag\\
 &\qquad\le
 L\sqrt{
 \Tr\!\left(\rho\mathscr P_c(A^2)\right)
 \Tr\!\left(\rho\mathscr P_c(B^2)\right)}
\label{eq:V92-pinched-geometric}\\
 &\qquad\le
 \frac L2\left[
 \|A\|\Tr\!\left(\rho\mathscr P_c(|A|)\right)
 +
 \|B\|\Tr\!\left(\rho\mathscr P_c(|B|)\right)
 \right].
\label{eq:V92-pinched-linear}
\end{align}
In particular, after taking the supremum over product \(\rho\), the
repeated carrier is reduced to the two positive one-occurrence
pinched absolute cards, with the explicit bound
\begin{align}
 &\sup_{\rho\ {\rm product}}
 \sum_\nu c_\nu
 \left\|
 \rho^{1/2}P_\nu ACBP_\nu\rho^{1/2}
 \right\|_1
\notag\\
 &\qquad\le
 \frac L2\left[
 \|A\|\kappa_\otimes\!\left(\mathscr P_c(|A|)\right)
 +
 \|B\|\kappa_\otimes\!\left(\mathscr P_c(|B|)\right)
 \right].
\label{eq:V92-pinched-capacity}
\end{align}
Only one linear use of the identified bank is required.
\end{theorem}

\begin{proof}
For each \(\nu\), Schatten H\"older gives
\begin{align*}
 &\left\|
 \rho^{1/2}P_\nu ACBP_\nu\rho^{1/2}
 \right\|_1\\
 &\qquad\le
 L
 \|\rho^{1/2}P_\nu A\|_2
 \|BP_\nu\rho^{1/2}\|_2.
\end{align*}
Apply Cauchy--Schwarz to the sum over \(\nu\) with weights \(c_\nu\).
Cyclicity of the ordinary trace gives
\begin{align*}
 \sum_\nu c_\nu\|\rho^{1/2}P_\nu A\|_2^2
 &=
 \Tr\!\left(\rho\sum_\nu c_\nu P_\nu A^2P_\nu\right),\\
 \sum_\nu c_\nu\|BP_\nu\rho^{1/2}\|_2^2
 &=
 \Tr\!\left(\rho\sum_\nu c_\nu P_\nu B^2P_\nu\right).
\end{align*}
This proves \eqref{eq:V92-pinched-geometric}.  Compress the
functional-calculus inequalities
\(A^2\preccurlyeq\|A\||A|\) and
\(B^2\preccurlyeq\|B\||B|\), sum with weights \(c_\nu\), and use
the arithmetic--geometric mean inequality.
\end{proof}

\begin{lemma}[Uniform size of a V61 separator]
\label{lem:V92-separator-size}
Let \(k_m\) be any \(m\)-row local cumulant letter with \(m\le4\).
Then
\begin{equation}
 \|k_m\|
 \le c_m,
 \qquad
 (c_1,c_2,c_3,c_4)=(1,2,6,26).
\label{eq:V92-cumulant-caps}
\end{equation}
Every local partition letter on four rows has norm at most \(26\),
and every V61 joining letter has norm at most
\begin{equation}
 L_4:=15\cdot26=390.
\label{eq:V92-joining-cap}
\end{equation}
A separator made of one joining letter, complete moment strings, and
replica regrouping is therefore uniformly bounded by \(L_4\).
\end{lemma}

\begin{proof}
The moment--cumulant formula and contraction of every complete moment
give
\[
 \|k_m\|
 \le
 \sum_{\pi\in\Pi_m}(|\pi|-1)!
 =
 \sum_{r=1}^m S(m,r)(r-1)!,
\]
where \(S(m,r)\) is a Stirling number of the second kind.  Direct
evaluation for \(m=1,2,3,4\) gives \(1,2,6,26\).
A local partition letter is a tensor product of cumulants on disjoint
blocks; among partitions of four its norm is at most \(26\).
There are at most \(B_4=15\) partition letters in a joining sum.
Complete moment strings are contractions by
\Cref{lem:V61-suffix-contraction}, while regrouping unitaries and
replica identifications have norm one.
\end{proof}

\begin{corollary}[Exact reduction of a contractively separated
overlap]
\label{cor:V92-overlap-reduction}
Let the V85 witness bank occur through self-adjoint complete-bank
operators \(A_a,B_a\).  Suppose that, after the V61 packet and all
nonpayer complete suffixes have been assembled, its compressed source
core has the exact form
\[
 P_\nu A_aC_aB_aP_\nu,
 \qquad
 \|C_a\|\le L_4.
\]
Then its weighted product trace majorant obeys
\eqref{eq:V92-pinched-linear}.  Hence the intervening joining,
complete-moment, and replica factors cause only a universal constant;
the remaining analytic gate is the linear capacity of the two
one-occurrence pinched absolute carriers.
\end{corollary}

\begin{proof}
Apply
\Cref{thm:V92-pinched-sandwich,lem:V92-separator-size}.
The witness selection has multiplicity at most four by
\Cref{thm:V85-overlap-witness}; no second selection is performed.
\end{proof}

\begin{firewall}[The physical pinched absolute card is still a theorem,
not a notation]
\label{fire:V92-pinched-absolute-open}
Neither \(\mathscr P_c(|A_a|)\) nor
\(\mathscr P_c(|B_a|)\) may be replaced automatically by the
unpinched capacities in
\eqref{eq:V86-adjacent-joint-capacity}.  Physical isotypic pinching
can correlate the four input rows.  Thus
\Cref{cor:V92-overlap-reduction} removes the noncommutative separator
as the obstruction, but it does not assert the still-missing
weighted pinched first-moment bound.
\end{firewall}

\section{The sole payer remains visible in the sandwich}
\label{sec:V92-payer-ledger}

\begin{proposition}[Four-location repeated-bank ledger]
\label{prop:V92-four-payer-ledger}
In the setting of
\Cref{thm:V92-pinched-sandwich}, the three core payer locations are
the separate substitutions below, whenever the displayed paid
replacement is self-adjoint in the resolved physical carrier:
\begin{align}
 (A,C,B)&\longmapsto(A^{\rm pay},C,B),
\label{eq:V92-payer-A}\\
 (A,C,B)&\longmapsto(A,C^{\rm pay},B),
\label{eq:V92-payer-C}\\
 (A,C,B)&\longmapsto(A,C,B^{\rm pay}).
\label{eq:V92-payer-B}
\end{align}
For \eqref{eq:V92-payer-A} and
\eqref{eq:V92-payer-B}, the corresponding paid occurrence replaces
one of the two pinched absolute cards.  For
\eqref{eq:V92-payer-C}, the factor \(L\) is replaced by
\(\|C^{\rm pay}\|\).  A suffix payer is a fourth tensor factor and is
estimated by the V2 one-sided tensor law; it is not inserted into any
of \eqref{eq:V92-payer-A}--%
\eqref{eq:V92-payer-B}.
\end{proposition}

\begin{proof}
The V2 scalar allocation assigns the unique output numerator and
resolvent to exactly one coordinate block.  Apply
\Cref{thm:V92-pinched-sandwich} to each of the first three resulting
operators separately.  For a suffix payer,
\Cref{lem:V61-suffix-contraction} leaves one
\(D_t^{\rm pay}\) and contracts every other suffix factor; the
one-sided tensor law then uses its ordinary norm.  At no point are two
paid substitutions multiplied.
\end{proof}

\begin{assessment}[The overlap bottleneck is now sharply typed]
\label{ass:V92-overlap-bottleneck}
The old description that the two occurrences are nonadjacent was too
coarse.  Uniformly bounded separation is harmless for the
trace-sandwiched carrier.  What remains is one of:
\begin{enumerate}[label=\textup{(\alph*)},leftmargin=4em]
\item a weighted pinched absolute estimate for an unpaid occurrence;
\item the analogous paid pinched absolute estimate when the payer lies
inside that occurrence;
\item a paid separator norm when the joining block pays; or
\item the ordinary one-link suffix payer combined with the same
one-use bank mark.
\end{enumerate}
These are four different analytic cards.  None is a product of two
first-moment estimates for the witness bank.
\end{assessment}

\section{Regression tests and the V91 payer completion}
\label{sec:V92-regressions}

\begin{proposition}[The exact marking map passes the V15 and V86
tests]
\label{prop:V92-V15-V86-tests}
\begin{enumerate}[label=\textup{(\roman*)},leftmargin=4em]
\item In the V15 dual-shared regression, take the path
\(\{12,13,24\}\), the bank \(\{f\}\) on edge \(12\), and
\(\{e\}\) on the other two edges.  Then
\eqref{eq:V92-bank-to-mark} is exactly
\eqref{eq:V15-rerouted-path}; the large dual-shared mass is retained
as a legal mark and is not demanded from singlet capacity.
\item In the V86 fixed-low-prefix packet, take
\[
 E_{12}=\{p\},\qquad E_{34}=\{q\},\qquad E_{13}=\{e\}.
\]
The exact path weight stays bounded below by a positive constant as
\(s\downarrow0\) and \(j\asymp s^{-3}\).  The unpaid and once-paid
test coefficients, of orders \(s^5\) and \(s^4\), are therefore
compatible with the assigned marking.
\end{enumerate}
\end{proposition}

\begin{proof}
Item \textup{(i)} follows by substituting the three singleton banks
into \eqref{eq:V92-bank-to-mark}; the resulting formula is
\[
 \frac{A^2a_1a_2a_3a_4}{(A+a_1)(A+a_2)}.
\]
For \textup{(ii)}, write \(a_{1/2}=1+C_2(1/2)\).  The selected
stocks are
\[
 H_{12}(\{p\})=a_{1/2}^2,\quad
 H_{34}(\{q\})=a_{1/2}^2,\quad
 H_{13}(\{e\})=3a_j.
\]
The denominator in \eqref{eq:V92-S22-path-expansion} is
\((a_j+a_{1/2})(3+a_{1/2})\), so the path weight tends to a positive
constant.  Compare with
\eqref{eq:V86-packet-unpaid-scale}--%
\eqref{eq:V86-packet-paid-prefix-scale}.
\end{proof}

\begin{lemma}[Finite-order local payer cap]
\label{lem:V92-local-payer-cap}
Every physically resolved moment--cumulant letter on at most four
rows, with the unique nontrivial output numerator and resolvent
assigned to that coordinate, has operator norm at most
\begin{equation}
 \frac{C}{s_\alpha}.
\label{eq:V92-local-payer-cap}
\end{equation}
The same is true for a fixed finite sum of such letters.
\end{lemma}

\begin{proof}
The proof of
\Cref{thm:V3-four-output-moment,cor:V3-junction-payer-cap}
uses only the finite moment--cumulant formula, contraction of
recoupling maps, fusion ancestry of the output Casimir, the one-link
first-moment debit, and the one-link gap.  The identical argument
with \(\Pi_m\), \(m\le4\), gives
\[
 s_\alpha A_U\|K_{m,U}\|\le C_m
\]
and then \eqref{eq:V92-local-payer-cap}.  Products belonging to a
local partition letter are handled by assigning the output Casimir
to one block at a time, as in
\eqref{eq:V3-fusion-output}.  There are at most fifteen letters, so a
fixed joining sum changes only the constant.
\end{proof}

\begin{theorem}[Every payer location of the V91 counterpacket is
legally marked]
\label{thm:V92-V91-paid}
Retain the V91 family
\[
 s=s_j,\qquad
 k=k_j\asymp s^{-\alpha},\qquad
 \frac12<\alpha<\frac23.
\]
For \(r\in\{p,q,e\}\), let
\(\mathcal Q_j^{{\rm pay},r}\) be the positive physical envelope of
the exact no-suffix carrier obtained by placing the sole V2 payer at
\(r\) in the full connected three-coordinate carrier of
\Cref{thm:V91-full-connected-lower}.  Then
\begin{equation}
 \boxed{
 \sum_{r\in\{p,q,e\}}
 \kappa_\otimes(\mathcal Q_j^{{\rm pay},r})
 \le
 C\,
 \mathcal W_{\sigma,\ell}.}
\label{eq:V92-V91-paid}
\end{equation}
Thus the V91 family is not a counterexample at any of its three
nonzero payer locations.
\end{theorem}

\begin{proof}
The exact connected carrier is the finite three-term expansion
\eqref{eq:V91-full-connected-expansion}.  Every unpaid binary or
ternary cumulant factor has uniformly bounded norm by
\eqref{eq:V92-cumulant-caps}.  If the payer is at \(p,q\), or \(e\),
\Cref{lem:V92-local-payer-cap} bounds the paid local factor by
\(C/s\).  There is no suffix, and the remaining factors have
uniform norm.  The finite triangle inequality therefore gives
\[
 \sum_{r\in\{p,q,e\}}
 \kappa_\otimes(\mathcal Q_j^{{\rm pay},r})
 \le\frac C{s}.
\]

On the other hand,
\Cref{prop:V91-rerouted-path} gives
\[
 \mathcal W_{\sigma,\ell}\ge c\,a_k^3.
\]
Since \(a_k\asymp k^2\asymp s^{-2\alpha}\),
\[
 \mathcal W_{\sigma,\ell}
 \ge c s^{-6\alpha}
 \ge \frac c{s}
\]
for all sufficiently small \(s\).  Absorb the finite initial range
into the constant.
\end{proof}

\begin{firewall}[A regression closure is not the uniform
same-side theorem]
\label{fire:V92-V91-not-uniform}
\Cref{thm:V92-V91-paid} uses the explicit three-coordinate V91
factorization and the very large legal prefix marks of that family.
It does not prove that every same-side-maximal higher-incidence bank
contains an \(s^{-1}\)-heavy marked path, nor does it sum arbitrary
first-connectivity positions.
\end{firewall}

\section{V92 status and the next upstream gate}
\label{sec:V92-status}

\begin{theorem}[Integrated V92 statement]
\label{thm:V92-status}
Preserving V1--V91, this chapter proves:
\begin{enumerate}[label=\textup{(V92.\arabic*)},leftmargin=6.8em]
\item every product of genuine edge-bank stocks is exactly the sum of
the corresponding legal coordinate markings;
\item the assembled clean \(S_{22}\) core has routing fibre at most
\(12\), with no support multiplicity;
\item the V81 aggregate paired-block estimate therefore lies in the
literal complete marked-tree sum;
\item clean foreign-prefix and cut-internal dominance have explicit
legal same-side marks;
\item after those routes and the earlier private/unbalanced routes,
the V88 forced outside-prefix alternative reduces quantitatively to
directed-cross or overlap cells;
\item a noncommutative contraction-sandwich estimate replaces the V86
adjacency requirement at the product trace-radius level;
\item its weighted pinched version reduces the physical overlap
problem to one-occurrence pinched absolute cards, while keeping the
sole payer locations distinct; and
\item the V15, V86, and V91 regression families pass the routing map,
with every nonzero payer location of the V91 counterpacket controlled
by its explicit legal path.
\end{enumerate}
\end{theorem}

\begin{proof}
The assertions are respectively
\Cref{prop:V92-bank-to-mark,
prop:V92-core-routing-fibre,
thm:V92-clean-paired-routing,
cor:V92-FI-legal-marks,
thm:V92-forced-cell-reduction,
lem:V92-sandwiched-joint-card,
thm:V92-pinched-sandwich,
prop:V92-four-payer-ledger,
prop:V92-V15-V86-tests,
thm:V92-V91-paid}.
\end{proof}

\begin{programme}[Mandate for V93]
\label{prog:V92-next}
\begin{enumerate}[label=\textup{(\arabic*)},leftmargin=4em]
\item Start from the weighted pinched carriers
\(\mathscr P_c(|A_a|)\) and
\(\mathscr P_c(|B_a|)\) in
\eqref{eq:V92-pinched-linear}.  Prove their one-bank linear bounds,
or construct a product-density counterexample.  Do not replace
pinched product capacity by ordinary operator norm.
\item Insert the V85 witness bank into the actual V61 source core and
record which of the four payer lines in
\Cref{prop:V92-four-payer-ledger} occurs.  Use the separator cap only
after the exact factor \(A_aC_aB_a\) has been identified.
\item For the directed residual
\eqref{eq:V92-directed-residual}, keep the V83
mass-compatible opposite-partner routes, but send a
same-side-maximal branch directly to the complete marked-tree sum.
Do not reuse the false V84 crossing-only square tail.
\item Treat the suffix payer as a fourth line.  Its one-link
\(C/s_\alpha\) norm must be combined with a genuine marked numerator
before the sum over suffix locations.
\item Keep the operator carrier assembled until the pinched or
one-sided capacity estimate is complete.  Only its resulting
nonnegative bank stocks may be expanded using
\Cref{prop:V92-bank-to-mark}.
\item Recheck every new bound on the V15, V86, and V91 families.
If a pinched bound fails, preserve the smallest exact physical block
and move upstream to the complete output-channel sum.
\item No companion audit is due at V93.  The next compulsory SM/NS
audit is V95.  At V100 freeze the active note with suffix
\texttt{\_f2} and restart at a contextual proof-indexed V1.
\end{enumerate}
\end{programme}

\begin{firewall}[Claim boundary after V92]
\label{fire:V92-boundary}
V92 proves exact coordinate-level mark routing for the clean paired
\(S_{22}\) core and removes bounded noncommutative separation as a
signed trace-radius obstruction.  It does not prove the required
weighted pinched absolute bank bound, the general same-side-maximal
directed-cross source estimate, the support-uniform suffix-payer sum,
or the remaining \(3+1\), \(2+1+1\), and discrete
first-connectivity families.  Global quartic and higher MTA, WUFC,
CPM, constructive continuum Yang--Mills theory,
Osterwalder--Schrader reconstruction, and a positive spectral gap
remain open.  The full Yang--Mills mass gap remains open.
\end{firewall}

\paragraph{Parent-version record.}
The complete V91 mathematical body was retained before this chapter
was appended.  V91 had \(43785\) source records, \(1496919\) bytes,
and SHA256
\[
 \texttt{354cb2844cee45e572e9fdc2806b07e130bec0b3252d70bb489e45d33f686cae}.
\]
No companion audit is due at V92; the next required audit remains V95.

% =============================================================================
% BEGIN APPENDED COMPACT VERSION V93 MATERIAL
% =============================================================================

\clearpage
\chapter{Pinching instability and outer-multiplicity coarsening}
\label{chap:V93}

\section{Context and the upstream correction}
\label{sec:V93-context}

V92 reduced every non-adjacent repeated-bank overlap to a weighted
pinched absolute one-occurrence card.  That was a useful reduction,
but it deliberately stopped before asserting that product capacity is
stable under pinching.  The first result below shows that such a
stability assertion is false even for an unweighted three-block
pinching of a positive rank-one operator.  The loss is unbounded with
the row dimension.  This is the operator-theoretic form of the
entanglement-swapping warning in V2.

The failure does not apply indiscriminately to the physical output
resolution.  There are two different operations which had been
conflated:
\begin{enumerate}
\item coarse projection onto a total \(G\)-isotypic output, which
commutes with every equivariant complete-bank operator; and
\item fine projection onto an outer multiplicity label, which need
not commute with that operator.
\end{enumerate}
We prove that the first operation leaves the unpaid V92 overlap card
unchanged.  We also prove an exact trace-norm coarsening theorem which
removes the second operation \emph{before} absolute values are taken,
provided the scalar output weight depends only on the isomorphism
class \(U\), as the physical mass and resolvent factors do.  Thus the
next proof target is no longer an abstract pinching-stability
principle.  It is a coarse-isotypic raw trace estimate in which the
outer-label source multiplier has not been replaced by a separate
operator norm.

\section{A dimension-growing obstruction to product-capacity
pinching}
\label{sec:V93-pinch-obstruction}

For a positive operator \(T\) on a row tensor product, write
\[
 \kappa_\otimes(T)
 :=
 \sup_{\substack{\rho=\rho_1\otimes\cdots\otimes\rho_m\\
                  \rho_i\ {\rm density}}}
 \Tr(\rho T).
\]
This agrees with the positive specialization of the product trace
radius used in V92.

\begin{theorem}[Unweighted pinching may amplify product capacity]
\label{thm:V93-pinch-counterexample}
For every \(d\ge2\) there are a positive rank-one operator \(A_d\) on
\(\mathbb C^d\otimes\mathbb C^d\) and an orthogonal resolution
\(P_+,P_-,P_0\) such that, for
\[
 {\cal E}_d(X):=P_+XP_++P_-XP_-+P_0XP_0,
\]
one has
\[
 \|A_d\|=1,\qquad
 \kappa_\otimes(A_d)=\frac1d,\qquad
 \kappa_\otimes({\cal E}_d(A_d))\ge\frac12.
\]
Consequently no dimension-independent constant \(C\) can satisfy
\[
 \kappa_\otimes({\cal E}(|A|))
 \le C\,\kappa_\otimes(|A|)
 \tag{V93.1}\label{eq:V93-false-pinch-card}
\]
for all positive \(A\) and all orthogonal pinchings \({\cal E}\).
The same conclusion holds for four rows after adjoining two
one-dimensional rows.
\end{theorem}

\begin{proof}
Let \(e_1,\ldots,e_d\) be the standard basis and put
\[
 \Omega_d=d^{-1/2}\sum_{r=1}^d e_r\otimes e_r,
 \qquad
 x=e_1\otimes e_2,
 \qquad
 A_d=|\Omega_d\rangle\langle\Omega_d|.
\]
The vectors \(\Omega_d\) and \(x\) are orthonormal.  The largest
Schmidt coefficient of \(\Omega_d\) is \(d^{-1/2}\); hence the
maximum squared overlap with a product unit vector is \(d^{-1}\).
Convexity then gives
\(\kappa_\otimes(A_d)=d^{-1}\).

Set
\[
 u_\pm=2^{-1/2}(\Omega_d\pm x),\qquad
 P_\pm=|u_\pm\rangle\langle u_\pm|,\qquad
 P_0=I-P_+-P_-.
\]
Since
\(\langle u_\pm,\Omega_d\rangle=2^{-1/2}\) and
\(P_0\Omega_d=0\),
\[
 {\cal E}_d(A_d)
 =\frac12(P_++P_-)
 =\frac12\bigl(
     |\Omega_d\rangle\langle\Omega_d|
     +|x\rangle\langle x|
   \bigr).
\]
The product vector \(x\) therefore has expectation \(1/2\) in
\({\cal E}_d(A_d)\).  The amplification ratio is at least \(d/2\),
which proves all assertions.
\end{proof}

\begin{remark}[Precise scope of the obstruction]
\label{rem:V93-obstruction-scope}
The projectors \(P_\pm\) above are not asserted to be Yang--Mills
isotypic projectors.  The theorem instead rules out a purely
operator-theoretic completion of the V92 programme: any successful
argument must use the centrality of the physical coarse output
projection or the exact order in which outer multiplicities are
summed.  Notice also that \(\|A_d\|=1\), so multiplying the
right-hand side of \eqref{eq:V93-false-pinch-card} by any fixed power
of the ordinary norm cannot repair the estimate.
\end{remark}

\section{Outer labels can be coarsened before the trace norm}
\label{sec:V93-coarsening}

\begin{theorem}[Exact outer-multiplicity coarsening]
\label{thm:V93-outer-coarsening}
Let \(P_U\) be pairwise orthogonal finite-rank projections and, for
each \(U\), let
\[
 P_U=\sum_{a\in{\cal A}_U}P_{U,a}
\]
be an orthogonal refinement.  Then every trace-class operator \(X\)
satisfies
\[
 \sum_{a\in{\cal A}_U}
   \|P_{U,a}XP_{U,a}\|_1
 \le
 \|P_UXP_U\|_1.
\tag{V93.2}\label{eq:V93-coarsen-one-U}
\]
Consequently, for arbitrary numbers \(w_U\ge0\),
\[
 \sum_U\sum_{a\in{\cal A}_U}
 w_U\|P_{U,a}XP_{U,a}\|_1
 \le
 \sum_Uw_U\|P_UXP_U\|_1.
\tag{V93.3}\label{eq:V93-coarsen-weighted}
\]
There is no factor involving \(|{\cal A}_U|\).
\end{theorem}

\begin{proof}
Work first on \(P_U{\cal H}\).  The map
\[
 {\cal D}_U(Y)=\sum_{a\in{\cal A}_U}P_{U,a}YP_{U,a}
\]
is an average of unitary conjugations.  Indeed, if the phases
\(\theta_a\) are integrated independently on the circle and
\(V_\theta=\sum_a e^{i\theta_a}P_{U,a}\), then
\[
 {\cal D}_U(Y)
 =\int V_\theta YV_\theta^*\,d\theta.
\]
Unitary invariance and convexity of the trace norm imply
\(\|{\cal D}_U(Y)\|_1\le\|Y\|_1\).  On the other hand
\({\cal D}_U(Y)\) is block diagonal, so its trace norm is the sum of
the trace norms of its diagonal blocks.  Taking \(Y=P_UXP_U\) proves
\eqref{eq:V93-coarsen-one-U}; multiplication by \(w_U\) and summation
give \eqref{eq:V93-coarsen-weighted}.
\end{proof}

\begin{corollary}[The source multiplier must stay inside the block]
\label{cor:V93-keep-multiplier}
Suppose the physical scalar output weight is \(w_U\), independent of
the outer label \(a\), and let \(K\) be the exact source multiplier.
For every trace-class input coefficient \(Z\),
\[
 \sum_{U,a}w_U
 \|P_{U,a}KZP_{U,a}\|_1
 \le
 \sum_Uw_U\|P_UKZP_U\|_1.
\tag{V93.4}\label{eq:V93-raw-coarse-carrier}
\]
Thus one may remove the outer-label sum without counting
multiplicities.  Replacing \(P_{U,a}K P_{U,a}\) first by its ordinary
operator norm generally creates \(a\)-dependent weights and destroys
the hypothesis needed for this coarsening.
\end{corollary}

\begin{proof}
Apply \Cref{thm:V93-outer-coarsening} to \(X=KZ\).
\end{proof}

\section{Equivariance makes coarse isotypic pinching harmless}
\label{sec:V93-equivariance}

\begin{lemma}[Intertwiners commute with coarse isotypic projections]
\label{lem:V93-intertwiner-centrality}
Let \(\pi\) be a finite-dimensional unitary representation of a
compact group \(G\), and let \(P_U\) denote the projection onto the
full \(U\)-isotypic summand.  If
\[
 T\pi(g)=\pi(g)T\qquad(g\in G),
\]
then
\[
 [T,P_U]=0\qquad\hbox{for every \(U\)}.
\tag{V93.5}\label{eq:V93-intertwiner-commutes}
\]
If \(T=T^*\), then \(|T|\), \(T^2\), and every bounded Borel function
of \(T\) also commute with \(P_U\).
\end{lemma}

\begin{proof}
The \(U\)-isotypic subspace is the sum of all invariant subspaces
isomorphic to \(U\).  An intertwiner maps each irreducible invariant
subspace either to zero or to an irreducible subspace of the same
isomorphism class.  It therefore preserves every full isotypic
subspace.  The same argument applied to \(T^*\) shows that its
orthogonal complement is preserved, which is equivalent to
\eqref{eq:V93-intertwiner-commutes}.  The final assertion follows
from functional calculus.
\end{proof}

\begin{lemma}[Complete Wilson bank letters are intertwiners]
\label{lem:V93-bank-intertwiners}
At fixed input labels, every complete one-coordinate Wilson moment,
every cumulant obtained from such moments, and every finite product
of complete bank letters commutes with the diagonal \(G\)-action.
Consequently all such operators commute with the coarse projections
\(P_U\).  This conclusion need not hold for a chosen outer
multiplicity projection \(P_{U,a}\).
\end{lemma}

\begin{proof}
The heat-kernel law and Haar law used in a complete coordinate bank
are invariant under conjugation.  Conjugating the tensor
representation inside the defining group integral and changing the
integration variable therefore leaves the moment operator unchanged.
Thus every complete moment is an intertwiner.  Cumulants are finite
linear combinations of products and tensor products of moments, so
they remain intertwiners; the same is true of finite products of bank
letters.  Apply \Cref{lem:V93-intertwiner-centrality}.  An outer
multiplicity projection belongs to a noncentral matrix algebra in the
commutant and is not forced to commute with another element of that
algebra.
\end{proof}

\begin{proposition}[Coarse weighted absolute card]
\label{prop:V93-coarse-weighted-card}
Let \(A=A^*\) be an equivariant complete-bank operator, let
\(\{P_U\}\) be the complete coarse isotypic resolution, and let
\[
 Q_w:=\sum_Uw_UP_U,\qquad w_U\ge0.
\]
Then
\[
 \sum_Uw_UP_U|A|P_U=Q_w|A|=|A|Q_w
\tag{V93.6}\label{eq:V93-weighted-commuting-identity}
\]
and
\[
 \kappa_\otimes(Q_w|A|)
 \le
 \min\bigl\{
   \|Q_w\|\,\kappa_\otimes(|A|),
   \|A\|\,\kappa_\otimes(Q_w)
 \bigr\}.
\tag{V93.7}\label{eq:V93-weighted-card-min}
\]
In particular, for \(w_U=1\),
\[
 \sum_UP_U|A|P_U=|A|,
\qquad
 \kappa_\otimes\!\left(\sum_UP_U|A|P_U\right)
 =\kappa_\otimes(|A|).
\tag{V93.8}\label{eq:V93-unpaid-coarse-exact}
\]
\end{proposition}

\begin{proof}
By \Cref{lem:V93-bank-intertwiners}, \(|A|\) commutes with every
\(P_U\), which proves \eqref{eq:V93-weighted-commuting-identity}.
The commuting operators \(Q_w\) and \(|A|\) are positive.  Hence
\[
 0\le Q_w|A|\le\|Q_w\||A|,
 \qquad
 0\le Q_w|A|\le\|A\|Q_w.
\]
Taking expectations in product density matrices proves
\eqref{eq:V93-weighted-card-min}.  Completeness gives \(Q_w=I\) when
all \(w_U=1\), proving \eqref{eq:V93-unpaid-coarse-exact}.
\end{proof}

\begin{theorem}[Closure of the unpaid coarse repeated-bank overlap]
\label{thm:V93-unpaid-coarse-overlap}
Let \(A=A^*\) and \(B=B^*\) be equivariant complete-bank operators,
let \(\|C\|\le L\), and let \(\{P_U\}\) be the complete coarse
isotypic resolution.  For every product density matrix
\(\rho=\rho_1\otimes\cdots\otimes\rho_m\),
\[
 \sum_U
 \bigl\|
  \rho^{1/2}P_UACBP_U\rho^{1/2}
 \bigr\|_1
 \le
 L\sqrt{\Tr(\rho A^2)\Tr(\rho B^2)}
\tag{V93.9}\label{eq:V93-unpaid-coarse-overlap}
\]
and therefore
\[
 \sum_U
 \bigl\|
  \rho^{1/2}P_UACBP_U\rho^{1/2}
 \bigr\|_1
 \le
 \frac L2\bigl(
   \|A\|\Tr(\rho|A|)
   +\|B\|\Tr(\rho|B|)
 \bigr).
\tag{V93.10}\label{eq:V93-unpaid-coarse-linear}
\]
Thus the V92 non-adjacent sandwich has no extra product-capacity loss
from the \emph{coarse} unpaid output resolution.
\end{theorem}

\begin{proof}
Apply the weighted pinched sandwich theorem
\Cref{thm:V92-pinched-sandwich} with \(c_U=1\).  Equivariance and
\Cref{lem:V93-intertwiner-centrality} give
\[
 \sum_UP_UA^2P_U=A^2,\qquad
 \sum_UP_UB^2P_U=B^2.
\]
The square-root estimate in
\eqref{eq:V93-unpaid-coarse-overlap} follows.  Since
\(A^2\le\|A\||A|\) and \(B^2\le\|B\||B|\), the arithmetic--geometric
mean inequality gives \eqref{eq:V93-unpaid-coarse-linear}.
\end{proof}

\begin{remark}[Raw coarsening and sandwiched coarsening are distinct]
\label{rem:V93-order-firewall}
Equation \eqref{eq:V93-raw-coarse-carrier} is an assertion about the
raw trace-class blocks \(P_{U,a}KZP_{U,a}\).  It does not by itself
imply
\[
 \sum_a\|
 \rho^{1/2}P_{U,a}KP_{U,a}\rho^{1/2}
 \|_1
 \le
 \|\rho^{1/2}P_UKP_U\rho^{1/2}\|_1,
\]
because a product density \(\rho\) need not commute with the
projectors.  The latter inequality is not used here.  The raw Fourier
coefficient must first be put in the form covered by
\Cref{thm:V93-outer-coarsening}; alternatively one must remain at the
coarse sandwiched expression of
\Cref{thm:V93-unpaid-coarse-overlap}.
\end{remark}

\section{Paid letters and the corrected overlap ledger}
\label{sec:V93-paid-ledger}

\begin{corollary}[Coarse resolution of a paid local letter]
\label{cor:V93-paid-local-coarse}
Let \(\widetilde A=\widetilde A^*\) be a complete equivariant paid
local letter covered by \Cref{lem:V92-local-payer-cap}; in particular,
the unique output numerator and resolvent have already been assigned
to this letter.  Then
\[
 \kappa_\otimes\!\left(
   \sum_UP_U|\widetilde A|P_U
 \right)
 =
 \kappa_\otimes(|\widetilde A|)
 \le \|\widetilde A\|
 \le \frac C{s_\alpha}.
\tag{V93.11}\label{eq:V93-paid-local-coarse}
\]
Thus coarse output resolution creates no additional loss beyond the
existing local payer debit.
\end{corollary}

\begin{proof}
The equality is \eqref{eq:V93-unpaid-coarse-exact} applied to the
already paid operator \(\widetilde A\).  Product capacity is bounded
by ordinary operator norm, and the final bound is
\eqref{eq:V92-local-payer-cap}.
\end{proof}

\begin{assessment}[What has and has not been removed]
\label{ass:V93-corrected-ledger}
The overlap ledger must now be read as follows.
\begin{enumerate}[label=\textup{(\roman*)},leftmargin=4em]
\item If the repeated complete banks are unpaid and only the coarse
isotypic output is resolved, the overlap is closed by
\Cref{thm:V93-unpaid-coarse-overlap}.
\item If the output numerator and resolvent have already been
incorporated in one complete local bank letter, its coarse resolution
is controlled by \Cref{cor:V93-paid-local-coarse}; the resulting
\(s_\alpha^{-1}\) debit must still be paid by the legal marked tree.
\item If the exact coefficient is first split into outer labels, but
the scalar weight is \(w_U\) rather than \(w_{U,a}\), keep the source
multiplier inside the raw trace-class block and apply
\Cref{cor:V93-keep-multiplier}.  No outer multiplicity count is then
present.
\item If a source multiplier is replaced by
\(\|P_{U,a}KP_{U,a}\|\) before the \(a\)-sum, or if one insists on a
fine sandwiched expression with a noncommuting product density, none
of the preceding theorems applies.  The counterexample
\Cref{thm:V93-pinch-counterexample} forbids filling this gap by a
general dimension-free pinching lemma.
\item A payer in the bounded separator, in the ordinary suffix, or in
a different complete bank remains a separate analytic card, exactly
as recorded in \Cref{ass:V92-overlap-bottleneck}.
\end{enumerate}
\end{assessment}

\begin{proposition}[The V92 routing identities survive the correction]
\label{prop:V93-routing-survives}
The exact bank-to-mark identity
\eqref{eq:V92-bank-to-mark}, the \(S_{22}\) path expansion
\eqref{eq:V92-S22-path-expansion}, and the universal clean routing
fibres of V92 are unchanged by the replacement of fine pinched
capacities with coarse raw blocks.  In particular, coarsening the
outer output label neither changes a coordinate bank
\(E_{uv}\) nor sums over a new coordinate marking.
\end{proposition}

\begin{proof}
The markings in V92 are chosen from the tensor-factor coordinates
appearing in complete edge banks.  An outer label \(a\) indexes a
basis component inside one fixed output isotypic multiplicity space;
it is not an additional tensor-factor coordinate.  The coarsening
map of \Cref{thm:V93-outer-coarsening} changes only
\(\sum_a\|P_{U,a}XP_{U,a}\|_1\) into
\(\|P_UXP_U\|_1\).  All factors
\(H_{uv}(E_{uv})/(\Xi_u\Xi_v)\), all payer choices, and the literal
marking map are therefore identical.
\end{proof}

\section{Regression tests}
\label{sec:V93-regressions}

\begin{proposition}[V15, V86, and V91 remain consistent]
\label{prop:V93-regression-tests}
The new obstruction and the coarse correction pass the three
mandatory regression families.
\begin{enumerate}[label=\textup{(\roman*)},leftmargin=4em]
\item In the V15 dual-shared family, the large shared representation
continues to be retained inside the legal path
\(\{12,13,24\}\).  Coarsening an outer output label does not replace
that mass by singlet product capacity.
\item In the V86 fixed-low-prefix family, every complete local heat
letter is equivariant and hence commutes with the relevant coarse
total-spin projection.  Its unpaid \(s^5\) and once-paid \(s^4\)
scales therefore acquire no additional coarse pinching loss.
\item The V91 tilted coherent packet remains paid by the explicit
path \(\{12,13,34\}\).  The estimate
\eqref{eq:V92-V91-paid} used a positive local payer cap and the lower
bound on the legal path weight, not a fine pinching-stability
inequality.
\end{enumerate}
\end{proposition}

\begin{proof}
For \textup{(i)}, use the literal marking identity in
\Cref{prop:V92-V15-V86-tests,prop:V93-routing-survives}.  The source
multiplier containing the shared mass remains inside the coarse block
by \eqref{eq:V93-raw-coarse-carrier}.  For \textup{(ii)}, apply
\Cref{lem:V93-bank-intertwiners} and
\eqref{eq:V93-unpaid-coarse-exact}, then compare with the two scales
listed in \Cref{prop:V92-V15-V86-tests}.  Item \textup{(iii)} is
exactly the proof of \Cref{thm:V92-V91-paid}; no step of that proof
uses the card refuted in \Cref{thm:V93-pinch-counterexample}.
\end{proof}

\section{The sharpened V94 gate}
\label{sec:V93-V94-gate}

\begin{programme}[Next noncircular attack]
\label{prog:V93-next}
The next version should proceed in the following order.
\begin{enumerate}[leftmargin=3em]
\item Return to the exact V85 Fourier coefficient \emph{before} an
outer-label operator norm is taken.  Identify its raw trace-class
operator \(X_U\) and use
\eqref{eq:V93-coarsen-weighted} to replace the complete
\((U,a)\)-sum by the coarse \(U\)-sum.
\item Determine the exact ordering of the coarse carrier.  It must be
recorded explicitly as a raw block \(P_UX_UP_U\) or as a symmetric
sandwich.  No cyclic reordering inside a trace norm and no
commutation with a product input may be assumed.
\item Prove a linear repeated-bank estimate for that coarse carrier.
The unpaid symmetric-sandwich branch is already
\Cref{thm:V93-unpaid-coarse-overlap}; the unresolved branch is the
coarse weighted raw trace norm.  If it fails, construct the smallest
\emph{physical} multiplicity-space witness and move farther upstream
to the unsplit Fourier coefficient.
\item Only after this analytic carrier is fixed, attach the exact
V92 marking map.  Keep paid-bank, paid-separator, and suffix-payer
cases separate.
\item Continue the V88 same-side-maximal directed-\(D_k\) branch and
the support-uniform suffix-payer sum.  Do not spend either debit using
a product of two first-moment bank estimates.
\item V94 has no compulsory companion audit.  At V95 inspect the
current SM and NS notes before choosing the next direction.  At V100
freeze this active series with suffix \texttt{\_f2} and restart at a
contextual proof-indexed V1.
\end{enumerate}
\end{programme}

\begin{firewall}[Claim boundary after V93]
\label{fire:V93-boundary}
V93 disproves arbitrary dimension-free stability of product capacity
under pinching, proves exact outer-multiplicity coarsening at the raw
trace-norm level, and closes the unpaid coarse-isotypic
contraction-sandwiched overlap.  It does not yet prove the weighted
coarse raw trace estimate required by the full V85 carrier, justify
coarsening a fine expression after insertion of a noncommuting
product density, close the general same-side-maximal directed-cross
source estimate, or sum the support-uniform suffix payer.
The remaining \(3+1\), \(2+1+1\), discrete first-connectivity, and
higher MTA families remain open.  WUFC, CPM, constructive continuum
Yang--Mills theory, Osterwalder--Schrader reconstruction, and a
strictly positive spectral gap remain open.  The full Yang--Mills
mass gap remains open.
\end{firewall}

\paragraph{Parent-version record.}
The complete V92 mathematical body was retained before this chapter
was appended.  V92 had \(44702\) source records, \(1529019\) bytes,
and SHA256
\[
 \texttt{faa1a8f09443b7ad80d607b2f551bf011b447333c738681a63a731d8572abb29}.
\]
No companion audit is due at V93; the next required audit remains
V95.

% =============================================================================
% BEGIN APPENDED COMPACT VERSION V94 MATERIAL
% =============================================================================

\clearpage
\chapter{Exact Fourier dualization of the repeated-bank carrier}
\label{chap:V94}

\section{Context: the outer label entered in a relaxation, not in V85}
\label{sec:V94-context}

The first item of the V93 mandate asked to remove an outer-label norm
from the \emph{exact V85 coefficient}.  A return to the source shows
that this wording was inaccurate.  V85 already works with the
complete simultaneous \((J,R,K)\)-isotypic projections
\(\Pi_{J,R,K}\), with the identity on every multiplicity space; see
\Cref{def:V85-local-complete-carriers,thm:V85-no-multiplicity-sum}.
There is no multiplicity-copy norm to undo in that calculation.

The fine label \(a\) enters earlier, between the exact Fourier identity
\eqref{eq:V1-exact-Fourier} and the sufficient positive carrier
\eqref{eq:V1-physical-carrier}.  After a physical decomposition of
the multiplicity space, the ideal-property estimate
\[
 \|X_{U,a}(I\otimes K_{U,a})\|_1
 \le \|K_{U,a}\|\,\|X_{U,a}\|_1
\]
replaces the multiplier by a scalar norm.  This produces the
\(a\)-dependent weights whose pinched capacity was left open in V92.
The V93 coarsening theorem remains correct, but applying it to V85
would attack an operation which V85 never made.

This chapter goes one step farther upstream.  Schatten duality is
applied to the output coefficient \emph{before} the multiplier is
normed.  The resulting dual contractions live on the irreducible
output factor \(H_U\), while the exact multiplier remains on the full
multiplicity space \(M_U\).  For a contractively separated repeated
bank, the dual contractions become part of the bounded separator.
This removes the fine pinching gate from the exact Fourier estimate.

\section{The unsplit dual Fourier carrier}
\label{sec:V94-dual-carrier}

Fix four input representations and write
\[
 {\cal H}_{\rm in}
 =\bigotimes_{i=1}^4H_{R_i}
 =\bigoplus_{U\in{\cal S}}W_U(H_U\otimes M_U),
\]
where \({\cal S}\) is the finite set of outputs occurring in this
fixed tensor product.  For \(X\in{\cal S}_1({\cal H}_{\rm in})\), put
\[
 X_U:=W_U^*XW_U.
\]
Let \(K_U\in{\cal B}(M_U)\) be the exact multiplier in
\eqref{eq:V1-exact-Fourier}, and let \(w_U\ge0\) be arbitrary scalar
weights.  Define
\[
 Y_U(X):=\Tr_{M_U}\!\left[X_U(I_{H_U}\otimes K_U)\right].
\tag{V94.1}\label{eq:V94-YU}
\]
For a family of contractions
\(\boldsymbol D=(D_U)_{U\in{\cal S}}\), \(D_U\in{\cal B}(H_U)\), set
\[
 {\cal T}_{w,K}(\boldsymbol D)
 :=
 \sum_{U\in{\cal S}}
 W_U\!\left[
  w_U(D_U^*\otimes K_U)
 \right]W_U^*.
\tag{V94.2}\label{eq:V94-dual-operator}
\]
Its dual product radius is
\[
 \Gamma_\otimes(w,K)
 :=
 \sup_{\|D_U\|\le1}
 \sup_{\rho=\rho_1\otimes\cdots\otimes\rho_4}
 \left|
  \Tr\!\left[\rho\,{\cal T}_{w,K}(\boldsymbol D)\right]
 \right|,
\tag{V94.3}\label{eq:V94-dual-radius}
\]
where the \(\rho_i\) are density matrices.

\begin{theorem}[Exact block-duality identity]
\label{thm:V94-exact-block-duality}
For every trace-class \(X\),
\[
 \boxed{
 \sum_{U\in{\cal S}}w_U\|Y_U(X)\|_1
 =
 \sup_{\|D_U\|\le1}
 \left|
  \Tr\!\left[X{\cal T}_{w,K}(\boldsymbol D)\right]
 \right|.}
\tag{V94.4}\label{eq:V94-exact-block-duality}
\]
No multiplicity basis, outer label, operator norm of \(K_U\), or
triangle inequality over output blocks occurs in this identity.
\end{theorem}

\begin{proof}
Trace-norm duality on \(H_U\) gives
\[
 \|Y_U(X)\|_1
 =
 \sup_{\|D_U\|\le1}
 \left|\Tr_{H_U}(D_U^*Y_U(X))\right|.
\]
Moreover, the defining property of the partial trace and cyclicity of
the full ordinary trace give
\begin{align*}
 \Tr_{H_U}(D_U^*Y_U(X))
 &=
 \Tr_{H_U\otimes M_U}\!\left[
  (D_U^*\otimes I)X_U(I\otimes K_U)
 \right]\\
 &=
 \Tr_{H_U\otimes M_U}\!\left[
  X_U(D_U^*\otimes K_U)
 \right].
\end{align*}
For each \(U\), choose a polar dual contraction for which the last
trace is the nonnegative number \(\|Y_U(X)\|_1\).  Since the output set
is finite, these contractions may be chosen independently.  Summing
then proves that the right side of
\eqref{eq:V94-exact-block-duality} is at least the left side.  The
reverse inequality follows from the triangle inequality and the
same one-block trace duality.
\end{proof}

\begin{corollary}[Exact weighted Fourier coefficient gate]
\label{cor:V94-exact-Fourier-gate}
Let \(X=A_1\otimes\cdots\otimes A_4\), and let
\(\widehat h(U)\) be given by the exact Fourier identity
\eqref{eq:V1-exact-Fourier}.  Then
\[
 \sum_{U\in{\cal S}}w_Ud_U\|\widehat h(U)\|_1
 =
 d_{\rm in}
 \sup_{\|D_U\|\le1}
 \left|
  \Tr\!\left[X{\cal T}_{w,K}(\boldsymbol D)\right]
 \right|
\tag{V94.5}\label{eq:V94-exact-Fourier-dual}
\]
and
\[
 \boxed{
 \sum_{U\in{\cal S}}w_Ud_U\|\widehat h(U)\|_1
 \le
 16d_{\rm in}\Gamma_\otimes(w,K)
 \prod_{i=1}^4\|A_i\|_1.}
\tag{V94.6}\label{eq:V94-Fourier-gate}
\]
If all \(A_i\) are positive, the factor \(16\) is replaced by \(1\).
\end{corollary}

\begin{proof}
Equation \eqref{eq:V1-exact-Fourier} and
\eqref{eq:V94-YU} give
\[
 d_U\|\widehat h(U)\|_1=d_{\rm in}\|Y_U(X)\|_1.
\]
Sum with weights and use
\Cref{thm:V94-exact-block-duality} to obtain
\eqref{eq:V94-exact-Fourier-dual}.  For positive \(A_i\), normalize
their traces and apply \eqref{eq:V94-dual-radius}.  For arbitrary
trace-class \(A_i\), apply the four-positive-parts decomposition
\Cref{lem:V1-positive-decomposition} in every row.  The product of
the four resulting mass losses is at most \(2^4=16\), exactly as in
\Cref{thm:V1-weighted-pinching-capacity}.
\end{proof}

\begin{remark}[The physical BFA weight]
\label{rem:V94-physical-weight}
For the quartic MTA row, take
\[
 w_U=\frac{\Xi(U)}{1-q_{\alpha,U}}
\]
on \(U\ne\boldsymbol1\).  The fusion-height factor is pulled to the
four inputs by the same subadditivity used in
\Cref{thm:V1-capacity-MTA}.  Thus
\eqref{eq:V94-Fourier-gate} has exactly the normalization and the
single resolvent required by the BFA coefficient sum.
\end{remark}

\section{Relation to the V1 positive carrier}
\label{sec:V94-positive-relaxation}

\begin{proposition}[The normed outer carrier is a relaxation]
\label{prop:V94-positive-relaxation}
Suppose
\[
 M_U=\bigoplus_aM_{U,a}
\]
is an orthogonal reducing decomposition for \(K_U\), write
\(K_{U,a}=K_U|_{M_{U,a}}\), and let
\[
 P_{U,a}=W_U(I_{H_U}\otimes Q_{U,a})W_U^*.
\]
Define
\[
 Q_{\rm norm}
 :=
 \sum_{U,a}w_U\|K_{U,a}\|P_{U,a}.
\tag{V94.7}\label{eq:V94-Qnorm}
\]
Then
\[
 \boxed{\Gamma_\otimes(w,K)\le\kappa_\otimes(Q_{\rm norm}).}
\tag{V94.8}\label{eq:V94-dual-below-positive}
\]
For the physical weights of
\Cref{rem:V94-physical-weight}, \(Q_{\rm norm}\) is the V1 carrier
\(\mathcal Q_{\rm phys}\).  Hence the V1 capacity criterion remains a
valid sufficient theorem, but it is not the exact carrier.
\end{proposition}

\begin{proof}
Fix a product density matrix \(\rho\) and contractions \(D_U\).  Set
\(\rho_{U,a}=P_{U,a}\rho P_{U,a}\).  Since the blocks reduce \(K_U\),
\[
 \Tr\!\left[\rho{\cal T}_{w,K}(\boldsymbol D)\right]
 =
 \sum_{U,a}w_U
 \Tr\!\left[
  \rho_{U,a}(D_U^*\otimes K_{U,a})
 \right].
\]
For a positive trace-class operator \(S\) and a bounded operator \(T\),
\[
 |\Tr(ST)|\le\|T\|\Tr S.
\]
Since \(\|D_U^*\otimes K_{U,a}\|\le\|K_{U,a}\|\), the absolute value
of the preceding sum is at most
\[
 \sum_{U,a}w_U\|K_{U,a}\|\Tr(\rho P_{U,a})
 =\Tr(\rho Q_{\rm norm}).
\]
Take the two suprema in \eqref{eq:V94-dual-radius}.
\end{proof}

\begin{assessment}[Where the V93 counterexample acts]
\label{ass:V94-counterexample-location}
The route
\[
 \Gamma_\otimes(w,K)
 \longrightarrow
 \kappa_\otimes(Q_{\rm norm})
\]
uses the triangle inequality and the ordinary norm independently on
every reducing outer block.  Once that route has been taken, one
faces the fine pinched absolute cards of V92 and the obstruction of
\Cref{thm:V93-pinch-counterexample}.  The exact identity
\eqref{eq:V94-exact-block-duality} does not take this route.  It
retains one freely chosen contraction \(D_U\) on \(H_U\) and the
entire multiplier \(K_U\) on \(M_U\), so there is no outer projection
to stabilize.
\end{assessment}

\section{The dual contractions belong to the separator}
\label{sec:V94-dual-sandwich}

For a bounded self-adjoint operator \(T\), introduce its one-use
energy
\[
 {\cal E}_\otimes(T):=\|T\|\,\kappa_\otimes(|T|).
\tag{V94.9}\label{eq:V94-one-use-energy}
\]

\begin{lemma}[Complete V61 separators are equivariant]
\label{lem:V94-separator-equivariant}
The complete-bank moment and cumulant operators, canonical replica
regroupings, and their finite products used in a V61 separator
commute with the diagonal \(G\)-action.  Consequently, if
\(A=A^*\), \(B=B^*\), and \(C\) are the two retained bank
occurrences and their assembled separator, then
\begin{align*}
 W_U^*AW_U&=I_{H_U}\otimes a_U,\\
 W_U^*BW_U&=I_{H_U}\otimes b_U,\\
 W_U^*CW_U&=I_{H_U}\otimes c_U
\end{align*}
on every output block, with \(a_U,b_U\) self-adjoint and
\[
 \sup_U\|c_U\|=\|C\|\le L_4.
\tag{V94.10}\label{eq:V94-separator-block-cap}
\]
The exact multiplicity multiplier of \(ACB\) is
\[
 K_U=a_Uc_Ub_U.
\tag{V94.11}\label{eq:V94-K-factorization}
\]
\end{lemma}

\begin{proof}
Complete Wilson moments and cumulants are intertwiners by
\Cref{lem:V93-bank-intertwiners}.  A canonical replica regrouping
only permutes tensor factors carrying the same diagonal action, hence
is also an intertwiner.  Sums, products, and tensor extensions of
intertwiners remain intertwiners.  Schur decomposition of the
commutant gives the three displayed block forms.  The norm of a
block-diagonal operator is the supremum of its block norms, and
\Cref{lem:V92-separator-size} gives \(\|C\|\le L_4\).  Multiplying the
three restrictions proves \eqref{eq:V94-K-factorization}.
\end{proof}

\begin{theorem}[Exact dual-sandwich estimate]
\label{thm:V94-dual-sandwich}
Assume the complete-block factorization
\eqref{eq:V94-K-factorization} and
\(\sup_U\|c_U\|\le L\).  Then, for the unpaid weights \(w_U=1\),
\[
 \boxed{
 \Gamma_\otimes(1,K)
 \le
 L\sqrt{
  {\cal E}_\otimes(A){\cal E}_\otimes(B)}
 \le
 \frac L2\bigl(
  {\cal E}_\otimes(A)+{\cal E}_\otimes(B)
 \bigr).}
\tag{V94.12}\label{eq:V94-dual-sandwich}
\]
More generally, let
\[
 Q_w=\sum_Uw_UP_U,\qquad A^{[w]}:=Q_wA.
\]
If \(A^{[w]}\) is bounded, then
\[
 \boxed{
 \Gamma_\otimes(w,K)
 \le
 L\sqrt{
  {\cal E}_\otimes(A^{[w]})
  {\cal E}_\otimes(B)}.}
\tag{V94.13}\label{eq:V94-left-paid-sandwich}
\]
The symmetric statement holds with \(B^{[w]}=Q_wB\).  If the weight
is retained instead in the separator, the same proof gives
\[
 \Gamma_\otimes(w,K)
 \le
 \left(\sup_Uw_U\|c_U\|\right)
 \sqrt{{\cal E}_\otimes(A){\cal E}_\otimes(B)}.
\tag{V94.14}\label{eq:V94-separator-paid-sandwich}
\]
\end{theorem}

\begin{proof}
For a dual family \(\boldsymbol D\), define
\[
 C_{\boldsymbol D}
 :=
 \sum_UW_U(D_U^*\otimes c_U)W_U^*.
\]
Then
\[
 \|C_{\boldsymbol D}\|\le L,\qquad
 {\cal T}_{1,K}(\boldsymbol D)=AC_{\boldsymbol D}B.
\]
For every product density matrix \(\rho\),
\begin{align*}
 \left|\Tr\!\left[
  \rho{\cal T}_{1,K}(\boldsymbol D)
 \right]\right|
 &\le
 \|\rho^{1/2}AC_{\boldsymbol D}B\rho^{1/2}\|_1\\
 &\le
 L\sqrt{\Tr(\rho A^2)\Tr(\rho B^2)}
\end{align*}
by \Cref{lem:V92-sandwiched-joint-card}.  Functional calculus gives
\[
 \Tr(\rho A^2)\le\|A\|\Tr(\rho|A|),
 \qquad
 \Tr(\rho B^2)\le\|B\|\Tr(\rho|B|).
\]
Take the supremum over product \(\rho\) and over
\(\boldsymbol D\).  This proves the first inequality in
\eqref{eq:V94-dual-sandwich}; the arithmetic--geometric mean gives
the second.

Because \(Q_w\) is central, it commutes with \(A\), and
\[
 {\cal T}_{w,K}(\boldsymbol D)
 =A^{[w]}C_{\boldsymbol D}B.
\]
Repeat the preceding argument to prove
\eqref{eq:V94-left-paid-sandwich}; the right-paid case is identical.
Finally absorb \(w_U\) into
\(D_U^*\otimes c_U\); its block norm is at most
\(\sup_Uw_U\|c_U\|\), proving
\eqref{eq:V94-separator-paid-sandwich}.
\end{proof}

\section{A one-bank energy compiler for overlap}
\label{sec:V94-energy-compiler}

\begin{corollary}[Linear stock from the geometric joint energy]
\label{cor:V94-energy-compiler}
Let \(\mathfrak h_{\cal B}\ge0\) be the scalar one-use stock assigned
to an identified bank \({\cal B}\).  Suppose two occurrences \(A,B\)
of that bank satisfy
\[
 {\cal E}_\otimes(A)\le C_A\mathfrak h_{\cal B},
 \qquad
 {\cal E}_\otimes(B)\le C_B\mathfrak h_{\cal B}.
\tag{V94.15}\label{eq:V94-unpaid-energy-cards}
\]
Then every exact unpaid multiplier \(K_U=a_Uc_Ub_U\) with
\(\sup_U\|c_U\|\le L\) obeys
\[
 \Gamma_\otimes(1,K)
 \le L\sqrt{C_AC_B}\,\mathfrak h_{\cal B}.
\tag{V94.16}\label{eq:V94-unpaid-energy-compiled}
\]
If the sole payer is placed in the \(A\)-occurrence and
\[
 {\cal E}_\otimes(A^{[w]})
 \le \frac{C_A^{\rm pay}}{s_\alpha^2}\mathfrak h_{\cal B},
 \qquad
 {\cal E}_\otimes(B)\le C_B\mathfrak h_{\cal B},
\tag{V94.17}\label{eq:V94-paid-energy-cards}
\]
then
\[
 \Gamma_\otimes(w,K)
 \le
 \frac{L\sqrt{C_A^{\rm pay}C_B}}{s_\alpha}
 \mathfrak h_{\cal B}.
\tag{V94.18}\label{eq:V94-paid-energy-compiled}
\]
The same conclusion holds with the payer in \(B\).
\end{corollary}

\begin{proof}
Insert \eqref{eq:V94-unpaid-energy-cards} into
\eqref{eq:V94-dual-sandwich}, and insert
\eqref{eq:V94-paid-energy-cards} into
\eqref{eq:V94-left-paid-sandwich}.  The square root is essential:
the two occurrences of the same stock produce
\(\sqrt{\mathfrak h_{\cal B}^2}=\mathfrak h_{\cal B}\), not a
duplicated factor \(\mathfrak h_{\cal B}^2\).
\end{proof}

\begin{corollary}[Norm and first moment imply the energy hypotheses]
\label{cor:V94-norm-capacity-to-energy}
The unpaid hypotheses \eqref{eq:V94-unpaid-energy-cards} follow from
\[
 \|A\|,\|B\|\le C_0,\qquad
 \kappa_\otimes(|A|),\kappa_\otimes(|B|)
 \le C_1\mathfrak h_{\cal B}.
\tag{V94.19}\label{eq:V94-unpaid-norm-capacity}
\]
The paid \(A\)-hypothesis in
\eqref{eq:V94-paid-energy-cards} follows from
\[
 \|A^{[w]}\|\le\frac{C_0^{\rm pay}}{s_\alpha},
 \qquad
 \kappa_\otimes(|A^{[w]}|)
 \le\frac{C_1^{\rm pay}}{s_\alpha}
 \mathfrak h_{\cal B}.
\tag{V94.20}\label{eq:V94-paid-norm-capacity}
\]
Thus the joint estimate needs one ordinary norm and one unpinched
first moment per occurrence; it never asks for a fine pinched
absolute first moment.
\end{corollary}

\begin{proof}
Multiply the two inequalities belonging to each occurrence and use
\eqref{eq:V94-one-use-energy}.
\end{proof}

\begin{theorem}[Exact Fourier closure of the overlap operator-order
gate]
\label{thm:V94-overlap-order-closure}
Let \({\cal B}_a\) be the one-bank witness selected in
\Cref{thm:V85-overlap-witness}.  Retain the complete V61 source
multiplier before outer resolution.  If its two occurrences and
assembled nonpayer separator have the V92 form
\[
 A_aC_aB_a,\qquad
 A_a=A_a^*,\quad B_a=B_a^*,\quad\|C_a\|\le L_4,
\tag{V94.21}\label{eq:V94-V92-complete-word}
\]
then its exact output Fourier coefficient mass is governed by
\Cref{cor:V94-exact-Fourier-gate} and
\[
 \Gamma_\otimes(1,K_a)
 \le
 L_4\sqrt{
  {\cal E}_\otimes(A_a){\cal E}_\otimes(B_a)}.
\tag{V94.22}\label{eq:V94-overlap-exact-unpaid}
\]
After V2 allocation, a payer in either bank occurrence is governed by
\eqref{eq:V94-left-paid-sandwich} or its right-hand analogue.  A
payer in the separator is governed by
\eqref{eq:V94-separator-paid-sandwich}, and a suffix payer remains the
separate one-sided tensor case of V92.

Consequently the fine pinched cards
\(\mathscr P_c(|A_a|)\) and \(\mathscr P_c(|B_a|)\) are not required
for the exact complete-block Fourier estimate.  Whenever the
unpinched one-use estimates
\eqref{eq:V94-unpaid-norm-capacity} hold, the unpaid overlap is
linear in the one-bank stock.  Whenever
\eqref{eq:V94-paid-norm-capacity} also holds at the selected payer
location, the paid overlap has exactly one \(s_\alpha^{-1}\) debit.
\end{theorem}

\begin{proof}
By \Cref{lem:V94-separator-equivariant}, the complete restrictions of
\eqref{eq:V94-V92-complete-word} satisfy
\(K_{a,U}=a_{a,U}c_{a,U}b_{a,U}\) on the full \(M_U\), with no outer
split and \(\sup_U\|c_{a,U}\|\le L_4\).  Apply
\Cref{thm:V94-dual-sandwich}.  The payer alternatives are its three
weighted conclusions together with the V2 one-payer allocation
\Cref{cor:V2-carrier-allocation}.  Finally apply
\Cref{cor:V94-energy-compiler,cor:V94-norm-capacity-to-energy}.
\end{proof}

\begin{assessment}[What V94 actually closes]
\label{ass:V94-closure-scope}
The theorem closes the \emph{operator-order} obstruction left by
\Cref{fire:V92-pinched-absolute-open}: output trace norm, complete
multiplicity, noncommuting bounded separation, and the two bank
occurrences are now kept in one exact dual carrier.  It does not
manufacture the scalar estimates
\eqref{eq:V94-unpaid-norm-capacity}--%
\eqref{eq:V94-paid-norm-capacity}.  Those must be cited from the
appropriate complete-bank first-moment theorem for each \(A,F,I\), or
directed-cross bank type.  In particular, a mere local payer cap
\(\|A^{[w]}\|\le C/s_\alpha\) is only half of
\eqref{eq:V94-paid-norm-capacity}; the paid unpinched capacity must
also carry the same physical bank stock.
\end{assessment}

\section{Compatibility with marking and the one-payer ledger}
\label{sec:V94-marking-ledger}

\begin{proposition}[No new fibre or support count]
\label{prop:V94-marking-compatible}
Assume the hypotheses of
\Cref{thm:V94-overlap-order-closure} and suppose the relevant one-use
stock \(\mathfrak h_{{\cal B}_a}\) is one of the complete coordinate
banks used in the V92 edge product.  Then the bounds
\eqref{eq:V94-unpaid-energy-compiled} and
\eqref{eq:V94-paid-energy-compiled} enter the exact bank-to-mark
identity \eqref{eq:V92-bank-to-mark} with the same coordinate marking
and the same routing fibre.  The output contractions \(D_U\) create
neither a coordinate choice nor an outer-label count.
\end{proposition}

\begin{proof}
The family \(\boldsymbol D\) is chosen only after the complete
coordinate multiplier has been assembled, and every \(D_U\) acts on
the final irreducible factor \(H_U\).  It is absorbed into
\(C_{\boldsymbol D}\), whose norm is already included in \(L_4\).
Thus it changes none of the edge banks \(E_{uv}\), response sums
\(H_{uv}(E_{uv})\), or payer-coordinate choices appearing in
\eqref{eq:V92-bank-to-mark}.  The marking proof
\Cref{prop:V93-routing-survives} therefore applies verbatim.
\end{proof}

\begin{assessment}[Corrected four-location ledger]
\label{ass:V94-four-location-ledger}
For an overlap word, the exact analytic cards are now:
\begin{enumerate}[label=\textup{(\roman*)},leftmargin=4em]
\item unpaid banks: \eqref{eq:V94-overlap-exact-unpaid};
\item payer in the left or right bank:
      \eqref{eq:V94-left-paid-sandwich} and its symmetric form;
\item payer in the joining separator:
      \eqref{eq:V94-separator-paid-sandwich}; and
\item payer in the suffix: the V2 one-sided tensor law, unchanged.
\end{enumerate}
This replaces the pinched ledger of
\Cref{prop:V92-four-payer-ledger} at the exact Fourier level.  It does
not combine any two payer locations and never introduces a second
resolvent.
\end{assessment}

\section{Regression tests and exact status}
\label{sec:V94-regressions}

\begin{proposition}[V15, V85, V86, V91, and V93 tests]
\label{prop:V94-regression-tests}
The exact dual carrier has the following compatibility properties.
\begin{enumerate}[label=\textup{(\roman*)},leftmargin=4em]
\item It preserves the V15 dual-singlet raw coefficient mass
\(d_R^{-1}\); no representation-dimension formation factor or
singlet operator-norm substitute is inserted.
\item It leaves the V85 complete spin-one channel calculation
unchanged, because that calculation already sums complete isotypic
operators before any multiplicity basis is chosen.
\item When \(C=I\) and \(w_U=1\),
\eqref{eq:V94-dual-sandwich} reduces to the V86 adjacent
repeated-bank estimate, with the dual output contraction playing the
role of a norm-one separator.
\item The V91 tilted coherent packet retains its path
\(\{12,13,34\}\) and its V92 payer estimate.  No output dualizer
changes a coordinate marking.
\item The V93 pinching counterexample is not contradicted or assumed
away: V94 avoids the fine pinching map rather than bounding its
product capacity.
\end{enumerate}
\end{proposition}

\begin{proof}
Item \textup{(i)} follows from the exact identity
\eqref{eq:V94-exact-Fourier-dual}, which is merely a dual
representation of \eqref{eq:V1-exact-Fourier}; compare
\eqref{eq:V1-raw-pinching}.  Item \textup{(ii)} is
\Cref{thm:V85-no-multiplicity-sum}.  For \textup{(iii)}, put
\(c_U=I\) in \Cref{thm:V94-dual-sandwich} and compare
\eqref{eq:V86-adjacent-joint-capacity}.  Item \textup{(iv)} follows
from
\Cref{thm:V92-V91-paid,prop:V94-marking-compatible}.  Item
\textup{(v)} is the distinction proved in
\Cref{prop:V94-positive-relaxation,ass:V94-counterexample-location}.
\end{proof}

\begin{theorem}[Exact status after V94]
\label{thm:V94-status}
Relative to V1--V93, this chapter proves:
\begin{enumerate}[label=\textup{(V94.\arabic*)},leftmargin=6.5em]
\item V85 contains no outer multiplicity-copy sum; the
      \(a\)-dependent norm enters only in the sufficient V1 positive
      carrier;
\item the complete weighted output trace-norm sum has the exact dual
      formula \eqref{eq:V94-exact-block-duality};
\item exact Fourier normalization converts that identity into the
      coefficient gate \eqref{eq:V94-Fourier-gate}, with factor one
      for positive row inputs and \(16\) for arbitrary row inputs;
\item the V1 normed positive carrier dominates the exact dual radius,
      so all previous positive-carrier theorems remain valid;
\item for a complete repeated-bank multiplier, every output
      dualizer is absorbed into the V61 separator without increasing
      its norm;
\item the exact unpaid overlap is bounded by the geometric mean of
      the two unpinched one-use energies, and a bank payer costs
      exactly one inverse-\(s_\alpha\) factor whenever its paid
      one-use norm and capacity have the two bounds in
      \eqref{eq:V94-paid-norm-capacity}; and
\item the operator-order obstruction caused by fine pinched absolute
      carriers is removed from the exact Fourier route.
\end{enumerate}
\end{theorem}

\begin{proof}
Items \textup{(V94.1)--(V94.4)} are
\Cref{sec:V94-context,thm:V94-exact-block-duality,cor:V94-exact-Fourier-gate,prop:V94-positive-relaxation}.
Items \textup{(V94.5)--(V94.7)} are
\Cref{lem:V94-separator-equivariant,thm:V94-dual-sandwich,cor:V94-energy-compiler,thm:V94-overlap-order-closure}.
\end{proof}

\section{The compulsory V95 audit and next proof gate}
\label{sec:V94-next}

\begin{programme}[V95 acquisition order]
\label{prog:V94-V95}
\begin{enumerate}[leftmargin=3em]
\item First perform the required read-only audit of the current
parallel SM and NS notes.  Record their exact filenames, sizes, and
hashes, and import only a type-correct proof-order lesson.
\item Build an explicit table for every clean V92 overlap occurrence
type \(A,F,I\).  For each left/right orientation, cite or prove the
unpaid norm and unpinched capacity in
\eqref{eq:V94-unpaid-norm-capacity}.
\item Repeat the table for a payer inside the identified bank.  The
required assertion is the paid unpinched capacity in
\eqref{eq:V94-paid-norm-capacity}; the universal local norm
\(\|A^{[w]}\|\le C/s_\alpha\) alone is insufficient.
\item Insert every verified row of that table into
\Cref{cor:V94-energy-compiler,prop:V94-marking-compatible}.  Any row
which fails must be returned to its complete balanced/unbalanced bank
decomposition rather than to a fine output pinching.
\item Treat the joining-separator payer and the support-uniform
suffix payer separately.  For the former prove
\(\sup_Uw_U\|c_U\|\) at marked-tree scale; for the latter retain the
ordinary one-link payer in the V2 one-sided tensor law.
\item Continue the V88 same-side-maximal directed-\(D_k\) branch
after, not before, the overlap table is closed.  Test every new
estimate on the V15, V86, and V91 families.
\end{enumerate}
\end{programme}

\begin{firewall}[Claim boundary after V94]
\label{fire:V94-boundary}
V94 gives an exact unsplit Fourier carrier and removes fine pinching
as the operator-order bottleneck for a complete repeated-bank word.
It does not prove the paid unpinched one-use capacity for every
\(A,F,I\) bank orientation, the marked joining-separator payer, or
the support-uniform suffix-payer sum.  It does not close the
same-side-maximal directed-cross branch, the remaining \(3+1\),
\(2+1+1\), discrete first-connectivity, or higher MTA families.
WUFC, CPM, constructive cutoff removal, reflection positivity,
Osterwalder--Schrader reconstruction, and a uniform positive
continuum spectral gap remain open.  The full Yang--Mills mass gap
remains open.
\end{firewall}

\paragraph{Parent-version record.}
The complete V93 mathematical body was retained before this chapter
was appended.  V93 had \(45248\) source records, \(1549170\) bytes,
and SHA256
\[
 \texttt{563575aefebdb77828356ccb478e3c263ca50194fc3e22415897d216996a2002}.
\]
No companion audit is due at V94.  The audit of the current SM and NS
notes is compulsory at V95.

% =============================================================================
% BEGIN APPENDED COMPACT VERSION V95 MATERIAL
% =============================================================================

\clearpage
\chapter{Audited \(A/F/I\) overlap energies and the local-separator
payer}
\label{chap:V95}

\section{Compulsory SM/NS checkpoint}
\label{sec:V95-audit}

\begin{record}[Active companion snapshots read at V95]
\label{rec:V95-companion-snapshots}
The read-only companion audit used the current unfrozen SM note and
the current NS note:
\begin{enumerate}
\item
\(\texttt{Renewal\_SM\_Einstein\_Continuum\_Research\_Notes\_V59.tex}\),
with \(25650\) source records, \(901982\) bytes, and SHA256
\[
 \texttt{0a7b66eca8a362245bdc7de1a4d31f82dcfb5b13f1d398ecdb396b3107062e63};
\]
\item
\(\texttt{Renewal\_NS\_Stability\_Research\_Notes\_V31.tex}\),
with \(23052\) source records, \(887537\) bytes, and SHA256
\[
 \texttt{f1121b62238ad5491bb7cf03635ea9934942edf573ed8faaa9ee79e1d38a6696}.
\]
\end{enumerate}
The SM files with suffixes \(\texttt{\_f1}\) and \(\texttt{\_f2}\)
are frozen archives rather than the active companion.  NS V31 is
unchanged from the V85 checkpoint, but its V31 formation and
sharpness conclusions were reread.  No companion file was modified.
\end{record}

\begin{record}[Late SM drift during the V95 calculation]
\label{rec:V95-late-SM-drift}
During this calculation the parallel SM programme advanced from V59
through V60--V62.  The audit was therefore refreshed against
\[
 \texttt{Renewal\_SM\_Einstein\_Continuum\_Research\_Notes\_V62.tex},
\]
which at the time of reading had \(26718\) source records,
\(942983\) bytes, and SHA256
\[
 \texttt{b8bb42c9d578c7f2a121225d5373ed656316276e2b6b22b05237308bbad3f1a6}.
\]
V60 keeps the semidirect diffeomorphism--gauge ghost operator
triangular until its joint stabilizer is formed.  V61 reconstructs
the Euclidean cold response from its local parent rather than Wick
rotating an already eliminated impedance, and isolates the interior
conformal-factor instability.  V62 rotates the entire longitudinal
BRST quartet, obtaining a finite-cutoff tangent supercycle while
leaving global reflection positivity and determinant-line phases
open.  The file was read only.
\end{record}

\begin{assessment}[Type-correct transfer from the checkpoint]
\label{ass:V95-audit-transfer}
SM V59 obtains gauge-parameter cancellation only after the cutoff is
chosen as a finite BRST complex which pairs each nonzero Neumann
scalar with its exact one-form partner.  It retains the residual
scalar determinant and keeps the chiral boundary phase separate.
The Yang--Mills lesson is proof order only: keep the complete
multiplier in the exact representation complex before taking outer
norms or absolute values, and do not infer that every unpaired factor
cancels.  This supports the V94 dual carrier but supplies no
Yang--Mills determinant estimate.

The V60--V62 refresh strengthens the same proof-order lesson.  A
coupled kernel must be assembled before separate zero-mode counts, a
reduced response must be continued through its local parent, and a
contour change must act on the whole algebraically coupled quartet.
For V95 this means that the two bank occurrences, the output
dualizer, and the payer location must remain in one complete
multiplier until the V94 sandwich is applied.  No gravitational
ellipticity, contour, or reflection-positivity statement is imported.

NS V31 proves that flat probability marks remove stopping-time
multiplicity, while unmarking costs an exact first moment.  Heat
formation supplies that moment only at the critical source action;
using a weaker source restores the missing weight and can become
circular.  The type-correct Yang--Mills consequence is that the local
cap \(C/s_\alpha\) cannot replace a paid one-use capacity carrying the
actual bank stock.  One must prove the paid capacity before applying
the V94 square-root compiler.  No Navier--Stokes energy or stopping
time theorem is imported.
\end{assessment}

\begin{assessment}[Internal nonduplication decision]
\label{ass:V95-nonduplication}
V81 already proves paired norm-and-capacity certificates for clean
own-prefix, foreign-prefix, and cut-internal dominance.  V95 does not
reprove their representation estimates.  It proves the missing
interface to V94: an independent uniform norm bound turns those
capacity certificates into one-use energies, and the geometric
repeated-bank estimate then charges the shared stock only once.
\end{assessment}

\section{The three clean bank types in one ledger}
\label{sec:V95-AFI-ledger}

Fix a same-side block \(C=\{k,k^\circ\}\).  On a clean disjoint
branch, let \(T\in\{A,F,I\}\) denote respectively own-prefix,
foreign-prefix, or cut-internal dominance at row \(k\).  Let
\({\cal N}_T\) and \({\cal P}_T\) be the corresponding complete
unpaid and internally once-paid augmented-block envelopes.  The
relevant stock is
\[
 \mathfrak h_T:=\frac{\widehat H_T}{\Xi_k},
\qquad
 0\le\widehat H_T\le H_{kk^\circ}.
\tag{V95.1}\label{eq:V95-hT}
\]
The source of \(\widehat H_T\) is:
\[
\begin{array}{c|c|c|c}
T&\text{dominant incidence}&\widehat H_T&
 \text{complete certificate}\\
\hline
A&A_k\ge\kappa\Xi_k&
 \text{selected own-prefix pair stock}&
 \text{\eqref{eq:V78-unpaid-binary-norm-compiler}--%
       \eqref{eq:V78-paid-binary-norm-compiler}}\\
F&F_k\ge\kappa\Xi_k&H_R&
 \text{\Cref{thm:V81-paid-foreign-block}}\\
I&I_k\ge\kappa\Xi_k&H_E+H_I&
 \text{\Cref{thm:V81-paid-internal-block}}
\end{array}
\tag{V95.2}\label{eq:V95-AFI-table}
\]
The consolidated statement is \Cref{cor:V81-AFI-paired}.

\begin{proposition}[Paired capacities for every \(A/F/I\) row]
\label{prop:V95-AFI-capacities}
For every row of \eqref{eq:V95-AFI-table}, with a constant depending
only on the fixed dominance fraction \(\kappa\),
\begin{align}
 \|{\cal N}_T\|,
 \ \kappa_\otimes({\cal N}_T)
 &\le C_\kappa\mathfrak h_T,
\tag{V95.3}\label{eq:V95-AFI-unpaid-paired}\\
 \|{\cal P}_T\|,
 \ \kappa_\otimes({\cal P}_T)
 &\le \frac{C_\kappa}{s_\alpha}\mathfrak h_T.
\tag{V95.4}\label{eq:V95-AFI-paid-paired}
\end{align}
These are alternative payer states of one complete augmented block,
not two factors in one source term.
\end{proposition}

\begin{proof}
For \(T=A\), apply the two ordinary-norm binary compilers
\eqref{eq:V78-unpaid-binary-norm-compiler}--%
\eqref{eq:V78-paid-binary-norm-compiler}; each also states the
capacity bound.  For \(T=F\), use
\eqref{eq:V81-foreign-unpaid}--\eqref{eq:V81-foreign-paid}.  For
\(T=I\), use
\eqref{eq:V81-internal-unpaid-certificate}--%
\eqref{eq:V81-internal-paid-certificate}.  The identification
\(\widehat H_T\le H_{kk^\circ}\) and all mixed spectator payer
locations are exactly
\Cref{def:V81-paired-block,cor:V81-AFI-paired}.
\end{proof}

\section{Uniform norms and the one-use energies}
\label{sec:V95-AFI-energies}

Let \(S_T=S_T^*\) be the exact signed complete-bank occurrence inside
the envelope \({\cal N}_T\), and let
\(S_T^{\rm pay}=(S_T^{\rm pay})^*\) be the same occurrence when the
unique V2 payer is assigned inside that augmented block.  Thus, on
the physical resolution used in V81,
\[
 |S_T|_{\rm phys}\preccurlyeq{\cal N}_T,
 \qquad
 |S_T^{\rm pay}|_{\rm phys}\preccurlyeq{\cal P}_T.
\tag{V95.5}\label{eq:V95-envelope-domination}
\]

\begin{lemma}[Uniform complete-word norms]
\label{lem:V95-AFI-uniform-norms}
For every clean quartic \(A/F/I\) occurrence,
\[
 \boxed{
 \|S_T\|\le C,\qquad
 \|S_T^{\rm pay}\|\le\frac C{s_\alpha}.}
\tag{V95.6}\label{eq:V95-AFI-uniform-norms}
\]
The constant is independent of support size, representations, the
dominance fraction, and the choice of left or right occurrence.
\end{lemma}

\begin{proof}
For \(T=A\), the unpaid binary covariance is a difference of two
complete moment contractions and hence has norm at most \(2\).
Its paid physical envelope has norm at most \(C/s_\alpha\) by
\Cref{cor:V78-paid-binary-minimum}.

For \(T=F\), use the exact factorization
\eqref{eq:V81-foreign-reduced-prefix}.  The foreign scalar and all
nonpayer moment spectators are contractions, while the reduced binary
covariance has norm at most \(2\).  In the three paid cases of
\Cref{thm:V81-paid-foreign-block}, the reduced binary payer is at most
\(C/s_\alpha\) by \Cref{cor:V78-paid-binary-minimum}, the complete
foreign-product payer is at most \(C/s_\alpha\) by
\eqref{eq:V81-product-paid-zero}, and a moment-spectator payer is at
most \(C/s_\alpha\) by \eqref{eq:V3-spectator-carrier-cap}.

For \(T=I\), retain the connected and replicated branches of
\Cref{thm:V81-paid-internal-block}.  Every unpaid factor is a
contraction or a binary covariance of norm at most \(2\).  On the
connected branch the sole paid binary factor is bounded by
\(C/s_\alpha\).  On the replicated branch the old-prefix payer has
the same bound, and either one-row complete-product payer is bounded
by \eqref{eq:V81-product-paid-zero}.  There are only two structural
branches and at most three payer positions in the replicated branch.

Finally, a fixed quartic V61 word contains only a universal finite
number of partition letters and spectator groups.  The foregoing
complete-bank bounds are applied after each coordinate group is
assembled, so no coordinate or support count is introduced.
\end{proof}

\begin{theorem}[Verified \(A/F/I\) energy table]
\label{thm:V95-AFI-energy-table}
For \(T=A,F,I\), the exact signed occurrences satisfy
\begin{align}
 {\cal E}_\otimes(S_T)
 &\le C_\kappa\mathfrak h_T,
\tag{V95.7}\label{eq:V95-AFI-unpaid-energy}\\
 {\cal E}_\otimes(S_T^{\rm pay})
 &\le
 \frac{C_\kappa}{s_\alpha^2}\mathfrak h_T.
\tag{V95.8}\label{eq:V95-AFI-paid-energy}
\end{align}
These are precisely the two hypotheses of the V94 overlap compiler;
they are now theorems for every row of
\eqref{eq:V95-AFI-table}.
\end{theorem}

\begin{proof}
By \eqref{eq:V95-envelope-domination} and
\Cref{prop:V95-AFI-capacities},
\begin{align*}
 \kappa_\otimes(|S_T|)
 &\le C_\kappa\mathfrak h_T,\\
 \kappa_\otimes(|S_T^{\rm pay}|)
 &\le \frac{C_\kappa}{s_\alpha}\mathfrak h_T.
\end{align*}
Multiply these inequalities by the two independent norm bounds in
\eqref{eq:V95-AFI-uniform-norms} and use the definition
\eqref{eq:V94-one-use-energy}.
\end{proof}

\begin{remark}[Why the two estimates are both necessary]
\label{rem:V95-two-estimates}
Using only \eqref{eq:V95-AFI-paid-paired} would give
\(\|{\cal P}_T\|\kappa_\otimes({\cal P}_T)
\lesssim s_\alpha^{-2}\mathfrak h_T^2\), which duplicates the bank
stock before the V94 square root.  Using only the uniform
\(C/s_\alpha\) norm would erase the bank stock.  The cross-use in
\Cref{thm:V95-AFI-energy-table}---uniform norm times the paired
capacity---gives \(s_\alpha^{-2}\mathfrak h_T\), the exact scale
needed for one paid occurrence against one unpaid occurrence.
\end{remark}

\section{All \(3\times3\) repeated-bank energy combinations}
\label{sec:V95-nine-combinations}

\begin{theorem}[Exact \(A/F/I\) overlap energies]
\label{thm:V95-nine-overlap-energies}
Let the two occurrences of the V85 witness bank have clean types
\(T_L,T_R\in\{A,F,I\}\), possibly with different certified rows, and
let their row-normalized stocks be
\(\mathfrak h_L,\mathfrak h_R\) as in \eqref{eq:V95-hT}.  Retain the
complete Fourier multiplier and let the nonpayer V61 separator have
norm at most \(L_4\).  Then
\begin{align}
 \Gamma_\otimes(1,K_{T_L,T_R})
 &\le
 C_\kappa L_4\sqrt{\mathfrak h_L\mathfrak h_R},
\tag{V95.9}\label{eq:V95-nine-unpaid}\\
 \Gamma_\otimes(w,K_{T_L^{\rm pay},T_R})
 +
 \Gamma_\otimes(w,K_{T_L,T_R^{\rm pay}})
 &\le
 \frac{C_\kappa L_4}{s_\alpha}
 \sqrt{\mathfrak h_L\mathfrak h_R}.
\tag{V95.10}\label{eq:V95-nine-bank-paid}
\end{align}
The constants are uniform over all nine ordered type pairs.  The two
terms in \eqref{eq:V95-nine-bank-paid} are alternative payer
locations; they are not simultaneously present in one source word.
\end{theorem}

\begin{proof}
Apply \Cref{thm:V94-dual-sandwich} and insert
\eqref{eq:V95-AFI-unpaid-energy} for both occurrences.  This proves
\eqref{eq:V95-nine-unpaid}.  If the left occurrence pays, insert
\eqref{eq:V95-AFI-paid-energy} on the left and
\eqref{eq:V95-AFI-unpaid-energy} on the right into
\eqref{eq:V94-left-paid-sandwich}; the square root produces exactly
\(s_\alpha^{-1}\sqrt{\mathfrak h_L\mathfrak h_R}\).  The right payer
is symmetric.  Adding the two mutually exclusive majorants changes
only a factor two.
\end{proof}

\begin{corollary}[Common normalized stock closes linearly]
\label{cor:V95-common-stock-overlap}
If the V92 routing supplies one scalar stock
\(\mathfrak h_{\cal B}\) such that
\[
 \mathfrak h_L,\mathfrak h_R
 \le C_0\mathfrak h_{\cal B},
\tag{V95.11}\label{eq:V95-common-stock}
\]
then
\begin{align}
 \Gamma_\otimes(1,K_{T_L,T_R})
 &\le C_{\kappa,0}\mathfrak h_{\cal B},
\tag{V95.12}\label{eq:V95-common-unpaid}\\
 \Gamma_\otimes(w,K_{\rm bank\text{-}pay})
 &\le
 \frac{C_{\kappa,0}}{s_\alpha}\mathfrak h_{\cal B}.
\tag{V95.13}\label{eq:V95-common-paid}
\end{align}
In particular, this applies when both occurrences use the same
certified row \(k\), because each table stock is bounded by
\[
 \mathfrak h_{\cal B}:=\frac{H_{kk^\circ}}{\Xi_k}.
\tag{V95.14}\label{eq:V95-common-row-stock}
\]
\end{corollary}

\begin{proof}
Equation \eqref{eq:V95-common-stock} implies
\(\sqrt{\mathfrak h_L\mathfrak h_R}
\le C_0\mathfrak h_{\cal B}\).  Apply
\Cref{thm:V95-nine-overlap-energies}.  The last assertion follows
from \eqref{eq:V95-hT}.
\end{proof}

\begin{firewall}[Different certified rows leave a scalar gate]
\label{fire:V95-row-mismatch}
If the earlier and current occurrences are certified against
different rows and no common upper stock of the form
\eqref{eq:V95-common-stock} has been proved, V95 yields only
\[
 \sqrt{
  \frac{\widehat H_L}{\Xi_{k_L}}\,
  \frac{\widehat H_R}{\Xi_{k_R}}}.
\tag{V95.15}\label{eq:V95-geometric-mismatch}
\]
This is an exact improvement over a product of two first moments, but
it is not automatically a literal MTA edge weight.  A strict-ascent
ordering of \(\Xi_{k_L},\Xi_{k_R}\) alone does not turn the geometric
mean into the stronger of the two row debits.  The mismatch must be
resolved using the actual shared-coordinate incidence or an
additional packet response; it is not hidden inside the overlap
classification.
\end{firewall}

\section{The payer in the local joining separator}
\label{sec:V95-separator-payer}

\begin{lemma}[Support-uniform local-separator payer]
\label{lem:V95-local-separator-payer}
Let \(C^{\rm pay}\) be a V61 separator in which the unique V2 output
numerator and resolvent are assigned to the local joining coordinate,
while every complete moment string away from that coordinate remains
unpaid.  Then
\[
 \boxed{\|C^{\rm pay}\|\le\frac{C}{s_\alpha}.}
\tag{V95.16}\label{eq:V95-local-separator-payer}
\]
The constant is independent of coordinate support, representations,
and the local partition.
\end{lemma}

\begin{proof}
At the joining coordinate there are at most \(B_4=15\) partition
letters.  A partition has at most four blocks.  After the V2
allocation, exactly one block contains the payer.  Its paid
moment--cumulant letter has norm at most \(C/s_\alpha\) by
\Cref{lem:V92-local-payer-cap}; every other local block has norm at
most \(26\) by \eqref{eq:V92-cumulant-caps}.  Thus the complete paid
joining sum has norm at most
\[
 15\cdot4\cdot26^3\,\frac C{s_\alpha},
\]
which is of the asserted form.  Complete moment strings away from the
joining coordinate are contractions, and canonical replica
regroupings are unitary.  They do not change the bound.
\end{proof}

\begin{theorem}[Local-separator payer has the same overlap scale]
\label{thm:V95-separator-paid-overlap}
Under the hypotheses of
\Cref{thm:V95-nine-overlap-energies}, if the sole payer lies in the
local joining separator, then
\[
 \boxed{
 \Gamma_\otimes(w,K_{\rm sep\text{-}pay})
 \le
 \frac{C_\kappa}{s_\alpha}
 \sqrt{\mathfrak h_L\mathfrak h_R}.}
\tag{V95.17}\label{eq:V95-separator-paid-overlap}
\]
Under the common-stock hypothesis
\eqref{eq:V95-common-stock}, the right side is at most
\(C_{\kappa,0}\mathfrak h_{\cal B}/s_\alpha\).
\end{theorem}

\begin{proof}
Use \eqref{eq:V94-separator-paid-sandwich} with
\(\sup_Uw_U\|c_U\|\) interpreted as the norm of the already allocated
paid separator.  Bound it by
\Cref{lem:V95-local-separator-payer}, and insert the two unpaid energy
bounds \eqref{eq:V95-AFI-unpaid-energy}.  The common-stock conclusion
is \eqref{eq:V95-common-stock}.
\end{proof}

\begin{remark}[The suffix is not part of this lemma]
\label{rem:V95-separator-not-suffix}
The proof of \Cref{lem:V95-local-separator-payer} applies to the
finite local joining letter and unpaid complete strings.  It does not
sum payer positions along an arbitrarily supported common suffix.
That is the distinct location isolated in
\Cref{fire:V92-suffix-payer-separate}; it remains open.
\end{remark}

\section{Routing consequences and regression tests}
\label{sec:V95-routing}

\begin{theorem}[Common-row \(A/F/I\) overlap closure]
\label{thm:V95-common-row-closure}
Consider a V92 overlap summand in which the repeated physical bank is
clean, each occurrence is of type \(A,F\), or \(I\), and all other
selected coordinate banks are disjoint from it.  If the two
occurrences are certified against one common row-normalized stock as
in \eqref{eq:V95-common-stock}, then the following three mutually
exclusive core states have the exact one-bank scale:
\begin{enumerate}[label=\textup{(\roman*)},leftmargin=4em]
\item no payer in the overlap word:
      \(C\mathfrak h_{\cal B}\);
\item the sole payer in either occurrence:
      \(C\mathfrak h_{\cal B}/s_\alpha\);
\item the sole payer in the local joining separator:
      \(C\mathfrak h_{\cal B}/s_\alpha\).
\end{enumerate}
After multiplication by the other genuine edge stocks of the
assembled \(S_{22}\) carrier, these bounds enter the literal marking
sum \eqref{eq:V92-bank-to-mark} with the V92 finite routing fibre.
No second use of the witness bank is made.
\end{theorem}

\begin{proof}
Items \textup{(i)--(ii)} are
\Cref{cor:V95-common-stock-overlap}; item \textup{(iii)} is
\Cref{thm:V95-separator-paid-overlap}.  The exact Fourier insertion
is \Cref{cor:V94-exact-Fourier-gate}.  Since the dual output
contractions are already inside the separator,
\Cref{prop:V94-marking-compatible} inserts the resulting scalar
stock into the exact coordinate marking identity.  The finite
connected/replicated, payer, and certificate choices are the fibre
count of \Cref{prop:V92-core-routing-fibre}; no output or support
label is added.
\end{proof}

\begin{proposition}[V15, V86, and V91 regressions]
\label{prop:V95-regression-tests}
\begin{enumerate}[label=\textup{(\roman*)},leftmargin=4em]
\item The V15 dual-shared family is not forced into the common-row
corollary when its two occurrences use different row normalizations.
It remains governed by the exact dual Fourier carrier and the legal
path \(\{12,13,24\}\); no singlet-capacity replacement is made.
\item In the V86 fixed-low-prefix family, the unpaid and paid
complete binary occurrences satisfy the \(A\)-row of
\eqref{eq:V95-AFI-table}.  The geometric compiler therefore adds no
power beyond the already recorded \(s^5\) unpaid and \(s^4\) paid
packet scales.
\item The V91 tilted coherent packet remains controlled at every
nonzero payer location by \eqref{eq:V92-V91-paid}.  Its two
same-side marks and crossing mark are unchanged by the V95 energy
interface.
\end{enumerate}
\end{proposition}

\begin{proof}
Item \textup{(i)} follows from the firewall
\Cref{fire:V95-row-mismatch} and the exact V15 routing test in
\Cref{prop:V92-V15-V86-tests}.  For \textup{(ii)}, combine
\Cref{thm:V95-AFI-energy-table} with the exact V86 scales cited in
\Cref{prop:V92-V15-V86-tests}; the V94 square root uses the same
stock once.  Item \textup{(iii)} is
\Cref{thm:V92-V91-paid,prop:V94-marking-compatible}.
\end{proof}

\section{Exact status and V96 direction}
\label{sec:V95-status}

\begin{theorem}[What V95 proves]
\label{thm:V95-status}
Relative to V1--V94, V95 proves:
\begin{enumerate}[label=\textup{(V95.\arabic*)},leftmargin=6.5em]
\item the compulsory SM/NS audit was completed read only, with exact
      fingerprints for the initial SM V59 and NS V31 snapshots and a
      late refresh through SM V62;
\item every clean \(A,F,I\) occurrence has a support-uniform unpaid
      norm \(O(1)\) and paid norm \(O(s_\alpha^{-1})\);
\item combining those norms with the V81 paired capacities gives the
      exact unpaid and paid one-use energies
      \[
       C\mathfrak h_T,\qquad
       Cs_\alpha^{-2}\mathfrak h_T;
      \]
\item all nine ordered \(A/F/I\) repeated-bank combinations obey the
      exact geometric stock bounds
      \eqref{eq:V95-nine-unpaid}--%
      \eqref{eq:V95-nine-bank-paid};
\item a common row-normalized stock turns those geometric bounds into
      one linear use of the witness bank, including a payer in either
      occurrence; and
\item a payer in the finite local joining separator has norm
      \(O(s_\alpha^{-1})\) and the same common-stock overlap scale.
\end{enumerate}
\end{theorem}

\begin{proof}
Item \textup{(V95.1)} is
\Cref{rec:V95-companion-snapshots,rec:V95-late-SM-drift,ass:V95-audit-transfer}.
Items \textup{(V95.2)--(V95.3)} are
\Cref{lem:V95-AFI-uniform-norms,thm:V95-AFI-energy-table}.
Items \textup{(V95.4)--(V95.5)} are
\Cref{thm:V95-nine-overlap-energies,cor:V95-common-stock-overlap}.
Item \textup{(V95.6)} is
\Cref{lem:V95-local-separator-payer,thm:V95-separator-paid-overlap}.
\end{proof}

\begin{programme}[V96 acquisition order]
\label{prog:V95-V96}
\begin{enumerate}[leftmargin=3em]
\item Return to the strict-ascent construction behind
\Cref{thm:V85-overlap-witness}.  For every repeated bank whose
earlier and current certificates use rows \(k_L\ne k_R\), retain its
actual common coordinate incidences and determine whether
\eqref{eq:V95-geometric-mismatch} is paid by one legal edge or by an
adjacent packet response.
\item Do not use the inequality
\(\Xi_{k_L}\le\Xi_{k_R}\) alone: it controls only one side of the
geometric mean.  If no incidence identity supplies the missing half
debit, construct the smallest exact two-row or three-row witness and
move upstream to the acquisition step which selected the bank.
\item Treat the common-suffix payer separately.  Form the complete
suffix product before summing payer positions and test whether the
V81 total-product gap gives a support-uniform \(C/s_\alpha\) norm
without erasing the coordinate mark required by V92.
\item Once the overlap and suffix rows are closed, resume the V88
same-side-maximal directed-\(D_k\) branch.  Preserve the V79 directed
incidence rather than replacing it by the larger untyped crossing
stock.
\item Recheck V15, V86, and V91.  The next compulsory companion audit
is V100.  At V100 freeze this active sequence with suffix
\texttt{\_f2}, then restart at a contextual proof-indexed V1 as
requested.
\end{enumerate}
\end{programme}

\begin{firewall}[Claim boundary after V95]
\label{fire:V95-boundary}
V95 closes the clean \(A/F/I\) overlap energy table, the common-row
bank-paid overlap, and the finite local-separator payer.  It does not
convert the different-row geometric mismatch
\eqref{eq:V95-geometric-mismatch} into a literal marked-tree edge,
sum an arbitrarily supported suffix payer, or close the general
same-side-maximal directed-cross branch.  The remaining \(3+1\),
\(2+1+1\), discrete first-connectivity, higher MTA, WUFC, and CPM
families remain open.  Constructive cutoff removal, reflection
positivity, Osterwalder--Schrader reconstruction, and a uniform
positive continuum spectral gap remain open.  The full Yang--Mills
mass gap remains open.
\end{firewall}

\paragraph{Parent-version record.}
The complete V94 mathematical body was retained before this chapter
was appended.  V94 had \(45961\) source records, \(1574286\) bytes,
and SHA256
\[
 \texttt{d4c894ad796376a88e1d87faa764d73d74d3d7a5fb4c9a5ae0def4481c2015b1}.
\]
The compulsory V95 companion audit is recorded in
\Cref{rec:V95-companion-snapshots}; the next audit is due at V100.

% The V95 document terminator is intentionally disabled here:
% \end{document}

% =============================================================================
% BEGIN APPENDED COMPACT VERSION V96 MATERIAL
% =============================================================================

\clearpage
\chapter{Ascent-overlap incidence and the grouped common-suffix payer}
\label{chap:V96}

\section{Context and proof order}
\label{sec:V96-context}

V95 reduced every clean repeated \(A/F/I\) bank to the exact geometric
one-use response
\[
 \sqrt{\mathfrak h_L\mathfrak h_R},
 \qquad
 \mathfrak h_\sigma
 =
 \frac{\widehat H_\sigma}{\Xi_{k_\sigma}},
 \quad \sigma\in\{L,R\}.
\]
When \(k_L=k_R\), this is bounded by one row-normalized physical stock.
The remaining ascent-overlap case has \(k_L\ne k_R\).  This chapter
returns to the incidence partition which produced the overlap witness.
It proves that the witness always carries an actual common-coordinate
edge and gives the exact finite graph table for using that edge.
It also constructs a two-coordinate balanced counterpacket showing that
strict mass ascent plus common incidence does not, by itself, choose the
spanning term needed by the marked-tree atlas.

The second task concerns the fourth payer location.  V61 correctly keeps
the common suffix as a complete Wilson moment, but its coordinatewise V2
majorant is too large to sum by ordinary norms.  The correct operation is
a two-stage allocation: allocate the payer first to the suffix as one
group, keep its total output denominator, and only then retain an optional
barycentric coordinate label.  The V81 total-product gap then gives the
support-uniform \(C/s_\alpha\) norm.

\section{What an ascent-overlap bank actually contains}
\label{sec:V96-overlap-incidence}

Let \({\cal B}\) be one of the coordinate banks acquired earlier along
the V84 strict-ascent path, at row \(p\).  Thus every coordinate of
\({\cal B}\) is active at \(p\).  At a later row \(r\ne p\), put
\[
 {\cal B}[r]
 :=
 \{e\in{\cal B}:a_{r,e}>0\},
 \qquad
 X_r({\cal B})
 :=
 \sum_{e\in{\cal B}[r]}a_{r,e}.
\]
Recall the standing nontrivial input floor \(a_{u,e}\ge a_*>0\)
whenever row \(u\) is active at \(e\).  Define
\begin{align}
 H_{pr}({\cal B})
 &:=
 \sum_{e\in{\cal B}[r]}a_{p,e}a_{r,e},
\label{eq:V96-common-incidence-stock}\\
 d_{p\to r}({\cal B})
 &:=
 \frac{H_{pr}({\cal B})}{\Xi_r}.
\label{eq:V96-current-oriented-response}
\end{align}

\begin{lemma}[Every ascent overlap contains a legal directed response]
\label{lem:V96-overlap-directed-response}
Suppose \({\cal B}\) is the bank selected by
\Cref{thm:V85-overlap-witness}.  If
\[
 X_r({\cal B})\ge\beta\Xi_r,
 \qquad
 \beta:=\frac{\delta}{m},
\]
then
\begin{equation}
 \boxed{
 H_{pr}({\cal B})
 \ge a_*X_r({\cal B}),
 \qquad
 d_{p\to r}({\cal B})\ge\beta a_*.}
\tag{V96.1}\label{eq:V96-directed-lower}
\end{equation}
The edge \(pr\) is an actual common-coordinate edge; it is not a
renaming of either row certificate.
\end{lemma}

\begin{proof}
Every \(e\in{\cal B}[r]\) is active at both \(p\) and \(r\).
Consequently \(a_{p,e}\ge a_*\), and
\[
 H_{pr}({\cal B})
 \ge
 a_*\sum_{e\in{\cal B}[r]}a_{r,e}
 =
 a_*X_r({\cal B}).
\]
Insert the overlap-witness inequality and divide by \(\Xi_r\).
No comparison between \(\Xi_p\) and \(\Xi_r\) is used.
\end{proof}

\begin{corollary}[Incidence-augmented geometric response]
\label{cor:V96-incidence-augmented-geometric}
For arbitrary nonnegative responses
\(\mathfrak h_L,\mathfrak h_R\), the same hypothesis gives
\begin{equation}
 \boxed{
 \sqrt{\mathfrak h_L\mathfrak h_R}
 \le
 \frac{d_{p\to r}({\cal B})}{2\beta a_*}
 \bigl(\mathfrak h_L+\mathfrak h_R\bigr).}
\tag{V96.2}\label{eq:V96-augmented-AMGM}
\end{equation}
If the unique payer is in either repeated occurrence or in the local
separator, the conclusion holds with the common factor
\(s_\alpha^{-1}\) on both sides.
\end{corollary}

\begin{proof}
Use
\(\sqrt{\mathfrak h_L\mathfrak h_R}
\le(\mathfrak h_L+\mathfrak h_R)/2\)
and \(1\le d_{p\to r}/(\beta a_*)\), the latter being
\eqref{eq:V96-directed-lower}.  The paid statement follows because
V95 introduces exactly one common factor \(s_\alpha^{-1}\).
\end{proof}

\begin{firewall}[A lower response is useful only with the right graph]
\label{fire:V96-lower-not-routing}
Equation \eqref{eq:V96-augmented-AMGM} produces the two terms
\[
 d_{p\to r}\mathfrak h_L,
 \qquad
 d_{p\to r}\mathfrak h_R.
\]
The actual edge \(pr\) must be combined with the certificate edge
retained in that term and with the remaining joining edge.  A term
which leaves one row isolated is not paid merely because its scalar
size is favorable.
\end{firewall}

\section{The exact four-row graph table}
\label{sec:V96-graph-table}

Assume that \(p\) and \(r\) lie on opposite sides of the \(2+2\) cut,
written
\[
 \{p,p^\circ\}\mid\{r,r^\circ\}.
\]
Put
\[
 e_L=pp^\circ,\qquad e_R=rr^\circ,\qquad e_\times=pr.
\]
Let \(g\) be the crossing edge retained by the assembled local joining
card.  The two terms of \eqref{eq:V96-augmented-AMGM} have candidate
edge sets
\[
 \tau_L=\{e_L,e_\times,g\},
 \qquad
 \tau_R=\{e_R,e_\times,g\}.
\]

\begin{proposition}[Complete graph compatibility table]
\label{prop:V96-graph-table}
Their spanning status is exactly
\[
\begin{array}{c|c|c}
g&\tau_L&\tau_R\\
\hline
pr&\text{not spanning}&\text{not spanning}\\
pr^\circ&\text{spanning}&\text{not spanning}\\
p^\circ r&\text{not spanning}&\text{spanning}\\
p^\circ r^\circ&\text{spanning}&\text{spanning}.
\end{array}
\tag{V96.3}\label{eq:V96-graph-table}
\]
Consequently the mismatch routes through
\eqref{eq:V96-augmented-AMGM} in each of the following cases:
\begin{enumerate}[label=\textup{(\roman*)},leftmargin=4em]
\item \(g=p^\circ r^\circ\);
\item \(g=pr^\circ\) and
      \(\mathfrak h_R\le C_0\mathfrak h_L\); or
\item \(g=p^\circ r\) and
      \(\mathfrak h_L\le C_0\mathfrak h_R\).
\end{enumerate}
The one-sided routing loss is at most
\((1+C_0)/(2\beta a_*)\).
\end{proposition}

\begin{proof}
The edges \(e_L,e_\times\) connect
\(\{p,p^\circ,r\}\), so \(\tau_L\) spans precisely when \(g\) is
incident to \(r^\circ\).  Similarly \(e_R,e_\times\) connect
\(\{p,r,r^\circ\}\), so \(\tau_R\) spans precisely when \(g\) is
incident to \(p^\circ\).  This proves the table.

For \(g=p^\circ r^\circ\), both summands in
\eqref{eq:V96-augmented-AMGM} have spanning edge sets.  If
\(g=pr^\circ\), only the left summand spans, while the stated
comparison gives
\(\mathfrak h_L+\mathfrak h_R\le(1+C_0)\mathfrak h_L\).
The third case is symmetric.
\end{proof}

\begin{remark}[Rows on the same side]
\label{rem:V96-same-side-rows}
If \(p\ne r\) lie in the same two-row block, then \(r=p^\circ\).
The common incidence edge is the already required same-side edge and
adds no new connectivity.  The scalar reduction
\[
 \sqrt{\mathfrak h_L\mathfrak h_R}
 \le\max\{\mathfrak h_L,\mathfrak h_R\}
\]
is usable when the remaining assembled packet already supplies the
other same-side edge and a crossing edge, but it cannot manufacture a
new graph edge.
\end{remark}

\section{A balanced two-coordinate counterpacket}
\label{sec:V96-counterpacket}

\begin{proposition}[Strict ascent does not choose the spanning response]
\label{prop:V96-two-coordinate-no-go}
Fix \(L>1\) and put \(b:=1+C_2(1/2)\).  For every sufficiently large
half-integer \(j\), set
\[
 N:=1+C_2(j),
\]
choose a half-integer \(J_j\in[N,N+\tfrac12]\), and put
\[
 M:=1+C_2(J_j).
\]
Then \(M\asymp N^2\), and there is a physical balanced four-row bank
\({\cal B}_j=\{e,f\}\) with weights
\[
\begin{array}{c|cccc}
 &p&p^\circ&r&r^\circ\\
\hline
e&N&N&b&b\\
f&b&b&M&M
\end{array}
\tag{V96.4}\label{eq:V96-two-coordinate-array}
\]
such that
\begin{align}
 \Xi_p=\Xi_{p^\circ}&=N+b,
 &
 \Xi_r=\Xi_{r^\circ}&=M+b>L(N+b),
\label{eq:V96-counter-masses}\\
 \mathfrak h_L
 &:=
 \frac{H_{pp^\circ}({\cal B}_N)}{\Xi_p}
 \asymp N,
 &
 \mathfrak h_R
 &:=
 \frac{H_{rr^\circ}({\cal B}_N)}{\Xi_r}
 \asymp M\asymp N^2,
\label{eq:V96-counter-responses}\\
 d_{p\to r}({\cal B}_N)&\asymp1,
 &
 \sqrt{\mathfrak h_L\mathfrak h_R}&\asymp N^{3/2}.
\label{eq:V96-counter-mismatch}
\end{align}
If \(g=pr^\circ\), the only spanning term in
\eqref{eq:V96-graph-table} is
\(d_{p\to r}\mathfrak h_L\asymp N\), which cannot dominate the
geometric mismatch uniformly in \(N\).
\end{proposition}

\begin{proof}
The half-integer \(J_j\) exists because consecutive half-integers are
separated by \(1/2\).  Since \(C_2(t)=t(t+1)\),
\(M\asymp N^2\).  Every table entry is therefore the exact
\(1+C_2\) mass of an \(SU(2)\) representation.

At each coordinate the two largest weights are equal, so no unique
weight dominates the sum of the other active weights by the V69
factor.  Both coordinates are balanced.  Directly,
\[
 H_{pp^\circ}=N^2+b^2,\qquad
 H_{rr^\circ}=M^2+b^2,\qquad
 H_{pr}=b(N+M).
\]
Division by the masses proves
\eqref{eq:V96-counter-responses}--%
\eqref{eq:V96-counter-mismatch}; the strict ascent holds for all
sufficiently large \(N\).

For \(g=pr^\circ\), the left set
\(\{pp^\circ,pr,pr^\circ\}\) spans, whereas
\(\{rr^\circ,pr,pr^\circ\}\) leaves \(p^\circ\) isolated.
The displayed asymptotics prove the failure.
\end{proof}

\begin{firewall}[Scope of the counterpacket]
\label{fire:V96-counterpacket-scope}
The counterpacket refutes an incidence-only completion of the V95
geometric mismatch.  It does not refute the complete V94 Fourier
sandwich, and it does not refute a larger first-connectivity packet
carrying an additional \(p^\circ\)-incident response.  The one-sided
and diagonal rows of \eqref{eq:V96-graph-table} must therefore be moved
upstream to that adjacent response.
\end{firewall}

\section{Grouped allocation of the common-suffix payer}
\label{sec:V96-suffix}

Let \(S\) be the finite set of coordinates strictly after a fixed V61
first-connectivity source.  On a complete coordinatewise output block
\(\boldsymbol U=(U_t)_{t\in S}\), let
\[
 A_t:=A_{U_t}\ge0,
 \qquad
 q_t:=q_{\alpha,U_t}\in[0,1],
\]
with \(A_t=0\) and \(q_t=1\) on the trivial output.  Put
\[
 Y_S(\boldsymbol U):=\sum_{t\in S}A_t,
 \qquad
 q_S(\boldsymbol U):=\prod_{t\in S}q_t.
\]
Let \(P_{\boldsymbol U,\boldsymbol a}\) be the complete physical
product-isotypic projection, including the identity on every
multiplicity space.

\begin{definition}[Aggregate paid common suffix]
\label{def:V96-aggregate-suffix}
Define
\begin{equation}
 {\cal D}_{S}^{\rm agg}
 :=
 \sum_{\boldsymbol U,\boldsymbol a}
 \frac{Y_S(\boldsymbol U)q_S(\boldsymbol U)}
      {1-q_S(\boldsymbol U)}
 P_{\boldsymbol U,\boldsymbol a},
\tag{V96.5}\label{eq:V96-aggregate-suffix}
\end{equation}
where the coefficient is zero when \(Y_S=0\), equivalently when
\(q_S=1\).
\end{definition}

\begin{lemma}[Two-stage group allocation]
\label{lem:V96-two-stage-allocation}
Let a resolved source word factor into a core and the common suffix,
with nonnegative output masses \(Y_C,Y_S\) and multipliers
\(q_C,q_S\).  Then, blockwise,
\begin{equation}
 \boxed{
 \frac{(Y_C+Y_S)q_Cq_S}{1-q_Cq_S}
 \le
 \frac{Y_Cq_C}{1-q_C}\,q_S
 +
 q_C\,\frac{Y_Sq_S}{1-q_S}.}
\tag{V96.6}\label{eq:V96-group-allocation}
\end{equation}
Thus the sole payer can be allocated between the core and suffix
without allocating it among individual suffix coordinates.
\end{lemma}

\begin{proof}
Since \(q_Cq_S\le q_C,q_S\),
\[
 1-q_Cq_S\ge1-q_C,
 \qquad
 1-q_Cq_S\ge1-q_S.
\]
Use the first inequality on the \(Y_C\) part of the numerator and the
second on the \(Y_S\) part.  A trivial group has zero output mass, so
the limiting convention is consistent.
\end{proof}

\begin{theorem}[Support-uniform complete-suffix payer]
\label{thm:V96-support-uniform-suffix}
Under the positive Wilson multiplier, one-link gap, and fixed-order
moment hypotheses of V52 and V81,
\begin{equation}
 \boxed{
 \|{\cal D}_{S}^{\rm agg}\|
 \le\frac{C}{s_\alpha}}
\tag{V96.7}\label{eq:V96-suffix-norm}
\end{equation}
uniformly in \(|S|\), all input representations, all complete output
labels, and all physical multiplicities.

If a V61 core has already been bounded by a positive carrier \(Q_C\),
then the alternative in which the sole payer belongs to the common
suffix satisfies
\begin{equation}
 \boxed{
 \kappa_\otimes(Q_C\otimes{\cal D}_{S}^{\rm agg})
 \le
 \frac{C}{s_\alpha}\kappa_\otimes(Q_C).}
\tag{V96.8}\label{eq:V96-suffix-one-sided}
\end{equation}
\end{theorem}

\begin{proof}
On a complete output block remove the trivial coordinates and call the
remaining set \(G\).  Its coefficient in
\eqref{eq:V96-aggregate-suffix} is
\[
 \frac{Y_Gq_G}{1-q_G}
 =
 \mathscr D_G^{\rm pay}.
\]
The V81 total-product estimate
\eqref{eq:V81-product-paid-zero} bounds it by \(C/s_\alpha\).
Because the complete physical blocks are mutually orthogonal, the
operator norm is the supremum of these coefficients, not their sum.
This proves \eqref{eq:V96-suffix-norm}.

The aggregate suffix is positive.  The second orientation of the
one-sided tensor law \eqref{eq:V2-one-sided-tensor} gives
\[
 \kappa_\otimes(Q_C\otimes{\cal D}_{S}^{\rm agg})
 \le
 \kappa_\otimes(Q_C)\|{\cal D}_{S}^{\rm agg}\|.
\]
Insert \eqref{eq:V96-suffix-norm}.
\end{proof}

\begin{proposition}[Barycentric coordinate labels are retained exactly]
\label{prop:V96-barycentric-suffix}
For \(t\in S\), define
\begin{equation}
 {\cal D}_{S,t}^{\rm bar}
 :=
 \sum_{\substack{\boldsymbol U,\boldsymbol a\\Y_S(\boldsymbol U)>0}}
 \frac{A_t}{Y_S(\boldsymbol U)}
 \frac{Y_S(\boldsymbol U)q_S(\boldsymbol U)}
      {1-q_S(\boldsymbol U)}
 P_{\boldsymbol U,\boldsymbol a}.
\tag{V96.9}\label{eq:V96-barycentric-piece}
\end{equation}
Then each summand is positive and
\begin{equation}
 \boxed{
 {\cal D}_{S}^{\rm agg}
 =
 \sum_{t\in S}{\cal D}_{S,t}^{\rm bar}.}
\tag{V96.10}\label{eq:V96-barycentric-sum}
\end{equation}
Thus a coordinate payer label may be retained without changing the
operator.  The sum must be formed before taking an ordinary norm or
product-state capacity.
\end{proposition}

\begin{proof}
On every block with \(Y_S>0\),
\[
 \sum_{t\in S}\frac{A_t}{Y_S}=1.
\]
Multiply by the aggregate coefficient and sum the orthogonal blocks.
\end{proof}

\begin{lemma}[Coordinatewise V2 denominators create support loss]
\label{lem:V96-coordinatewise-loss}
The coordinatewise majorant
\begin{equation}
 q_S\sum_{t\in S}\frac{A_t}{1-q_t}
\tag{V96.11}\label{eq:V96-coordinatewise-majorant}
\end{equation}
cannot be bounded by \(C/s_\alpha\) uniformly in \(|S|\).  Take \(N\)
equal nontrivial outputs with \(A_t=a>0\),
\(q_t=e^{-s_\alpha a}\), and \(s_\alpha a\le N^{-1}\).  Then the
aggregate coefficient is comparable to \(s_\alpha^{-1}\), whereas
\eqref{eq:V96-coordinatewise-majorant} is comparable to
\(Ns_\alpha^{-1}\).
\end{lemma}

\begin{proof}
Now \(q_S=e^{-Ns_\alpha a}\ge e^{-1}\), and
\[
 1-e^{-Ns_\alpha a}\asymp Ns_\alpha a,
 \qquad
 1-e^{-s_\alpha a}\asymp s_\alpha a.
\]
Therefore
\[
 \frac{Na\,q_S}{1-q_S}\asymp\frac1{s_\alpha},
 \qquad
 q_S\sum_{t=1}^N\frac{a}{1-e^{-s_\alpha a}}
 \asymp\frac N{s_\alpha}.
\]
\end{proof}

\begin{assessment}[Resolution of the suffix ledger]
\label{ass:V96-suffix-ledger}
The suffix payer is support-uniform, but the proof uses the grouped
operator \eqref{eq:V96-aggregate-suffix}, not a triangle inequality
over the coordinatewise carriers \(D_t^{\rm pay}\) of
\eqref{eq:V2-spectator-payer}.  The output coordinate \(t\) in
\eqref{eq:V96-barycentric-piece} is an auxiliary payer label; it is
not one of the three input-coordinate edge marks in the V92 path
identity.  The core already supplies those marks.  Hence the suffix is
kept as the single fourth payer location of
\Cref{fire:V92-suffix-payer-separate}, while
\eqref{eq:V96-barycentric-sum} preserves the finer label if needed.

Accordingly, the support-uniform suffix assertions in
\Cref{lem:V65-four-payers,prop:V78-suffix-payer} are valid for their
factorized complete common suffixes.  The sum over payer coordinates
in those proofs must be interpreted through the grouped
\Cref{lem:V96-two-stage-allocation,thm:V96-support-uniform-suffix},
not through the lossy majorant
\eqref{eq:V96-coordinatewise-majorant}.
\end{assessment}

\begin{corollary}[The fourth V95 overlap payer location is closed]
\label{cor:V96-overlap-suffix-closed}
Under the hypotheses of
\Cref{thm:V95-nine-overlap-energies}, suppose the unique payer belongs
to an arbitrarily supported complete common suffix and the two repeated
bank occurrences are unpaid.  Then
\begin{equation}
 \boxed{
 \Gamma_\otimes(w,K_{\rm suffix\text{-}pay})
 \le
 \frac{C_\kappa L_4}{s_\alpha}
 \sqrt{\mathfrak h_L\mathfrak h_R}.}
\tag{V96.12}\label{eq:V96-overlap-suffix-bound}
\end{equation}
Under the V95 common-stock hypothesis, this is at most
\(C_{\kappa,0}\mathfrak h_{\cal B}/s_\alpha\).
\end{corollary}

\begin{proof}
Apply the grouped allocation
\eqref{eq:V96-group-allocation}, retaining the suffix alternative as
one positive factor.  Its norm is at most \(C/s_\alpha\) by
\Cref{thm:V96-support-uniform-suffix}.  Absorb it into the bounded
separator in the V94 dual sandwich and insert the two unpaid energies
\eqref{eq:V95-AFI-unpaid-energy}.  This proves
\eqref{eq:V96-overlap-suffix-bound}.  The final assertion is
\eqref{eq:V95-common-stock}.
\end{proof}

\section{Regression tests and exact status}
\label{sec:V96-regressions}

\begin{proposition}[V15, V86, and V91 remain intact]
\label{prop:V96-regression-tests}
\begin{enumerate}[label=\textup{(\roman*)},leftmargin=4em]
\item The V15 dual-shared family receives an incidence edge only when
      the selected bank is active on both endpoint rows.  Its exact
      path \(\{12,13,24\}\) and dual Fourier mass are unchanged.
\item The V86 fixed-low-prefix packet keeps both exact binary-prefix
      responses.  The V96 counterpacket forbids discarding the
      adjacent prefix response merely because a strict ascent supplied
      a common edge.
\item The V91 tilted coherent packet keeps the path
      \(\{12,13,34\}\).  Grouping a suffix strictly after first
      connectivity changes none of those edge marks or the payer
      bounds in \eqref{eq:V92-V91-paid}.
\end{enumerate}
\end{proposition}

\begin{proof}
Item \textup{(i)} is the support condition in
\Cref{lem:V96-overlap-directed-response} and the Fourier preservation
of \Cref{prop:V94-regression-tests}.  Item \textup{(ii)} follows from
the V75 deterministic prefix order and
\Cref{prop:V96-two-coordinate-no-go}.  For item \textup{(iii)}, use
\Cref{lem:V75-common-suffix,thm:V96-support-uniform-suffix}.
\end{proof}

\begin{theorem}[What V96 proves]
\label{thm:V96-status}
Relative to V1--V95, this chapter proves:
\begin{enumerate}[label=\textup{(V96.\arabic*)},leftmargin=6.5em]
\item every V85 ascent-overlap witness supplies the actual directed
      common-coordinate response \eqref{eq:V96-directed-lower};
\item the geometric mismatch has the incidence-augmented bound
      \eqref{eq:V96-augmented-AMGM};
\item the graph table \eqref{eq:V96-graph-table} identifies the
      immediately routable and one-sided comparison rows;
\item the balanced two-coordinate family
      \eqref{eq:V96-two-coordinate-array} proves strict ascent and
      common incidence insufficient for the latter rows;
\item the complete common suffix has the aggregate paid operator
      \eqref{eq:V96-aggregate-suffix} with norm \(C/s_\alpha\);
\item the barycentric identity
      \eqref{eq:V96-barycentric-sum} retains a payer-coordinate label
      without a support count when summed before the norm; and
\item the suffix-payer row of the V94--V95 repeated-bank ledger obeys
      \eqref{eq:V96-overlap-suffix-bound}.
\end{enumerate}
\end{theorem}

\begin{proof}
Items \textup{(V96.1)--(V96.4)} are
\Cref{lem:V96-overlap-directed-response,
cor:V96-incidence-augmented-geometric,
prop:V96-graph-table,prop:V96-two-coordinate-no-go}.
Items \textup{(V96.5)--(V96.7)} are
\Cref{thm:V96-support-uniform-suffix,
prop:V96-barycentric-suffix,
cor:V96-overlap-suffix-closed}.
\end{proof}

\begin{programme}[V97 upstream acquisition order]
\label{prog:V96-V97}
\begin{enumerate}[leftmargin=3em]
\item Return to the exact V61 packet before the V94 two-occurrence
      sandwich.  In the one-sided rows of
      \eqref{eq:V96-graph-table}, retain the adjacent prefix response
      incident to the row which would otherwise be isolated.
\item Build a three-factor energy inequality for
      \[
       S_L\,C_1\,S_R\,C_2\,K_{\rm adj}
      \]
      which charges the repeated bank once, keeps both separators
      bounded, and assigns the missing graph edge to
      \(K_{\rm adj}\).
\item Test that inequality on
      \eqref{eq:V96-two-coordinate-array}.  It must recover the missing
      factor of order \(N^{1/2}\) in the \(g=pr^\circ\) orientation or
      exhibit an exact larger counterpacket.
\item Treat the diagonal row \(g=pr\) separately.  Since neither term
      in \eqref{eq:V96-graph-table} spans, retain both proper-prefix
      responses or return to the partition-state sum before fixing the
      joining card.
\item Only after these overlap rows close, resume the V88
      same-side-maximal directed-\(D_k\) branch.  Keep the V79 directed
      incidence and the V84 intermediate-output obstruction.
\item Recheck V15, V86, and V91.  No companion audit is due before
      V100.  At V100 freeze this sequence as
      \texttt{\_V100\_f2.tex} and restart with a contextual,
      proof-indexed active V1.
\end{enumerate}
\end{programme}

\begin{firewall}[Claim boundary after V96]
\label{fire:V96-boundary}
V96 closes the arbitrarily supported complete common-suffix payer and
identifies the directed edge carried by every ascent-overlap witness.
It does not close the one-sided or diagonal rows of
\eqref{eq:V96-graph-table}; the two-coordinate counterpacket proves the
present incidence and ascent data insufficient there.  The general
same-side-maximal directed-cross branch, the remaining \(3+1\),
\(2+1+1\), discrete first-connectivity, higher MTA, WUFC, and CPM
families remain open.  Constructive cutoff removal, reflection
positivity, Osterwalder--Schrader reconstruction, and a uniform positive
continuum spectral gap remain open.  The full Yang--Mills mass gap
remains open.
\end{firewall}

\paragraph{Parent-version record.}
The complete V95 mathematical body was retained before this chapter was
appended.  V95 had \(46556\) source records, \(1598351\) bytes, and
SHA256
\[
 \texttt{576ea390f65781807289220449f8d9874c43841ab0a36ae832173794bb3dc7f7}.
\]
No companion audit is due at V96; the next audit is V100.

% =============================================================================
% END APPENDED COMPACT VERSION V96 MATERIAL
% =============================================================================

% The V96 document terminator is intentionally disabled here:
% \end{document}

% =============================================================================
% BEGIN APPENDED COMPACT VERSION V97 MATERIAL
% =============================================================================

\clearpage
\chapter{The three-factor endpoint sandwich and distinct-adjacent
overlap closure}
\label{chap:V97}

\section{Context: use the middle occurrence only in norm}
\label{sec:V97-context}

V96 proved that every strict-ascent overlap contains a genuine common
incidence edge.  Its graph table also showed why that fact was not
enough.  In a one-sided row, the scalar term containing the larger
repeated-bank response can leave the opposite same-side row isolated.
The missing object is an adjacent proper-prefix response on that
same-side edge.

The exact V61 packet places this response on a coordinate group
separated from the two repeated-bank occurrences by complete
intertwining words.  The required analytic form is therefore
\[
 A\,C_1\,B\,C_2\,E.
\]
Here \(A\) is the repeated-bank occurrence whose edge response is
retained, \(B\) is the other occurrence of the same bank, and \(E\)
is the adjacent response.  The new point is to use \(A\) and \(E\)
as the two Hilbert--Schmidt endpoints and use \(B\) only through its
support-uniform ordinary norm.  Thus the repeated bank contributes
one first-moment stock, while the disjoint adjacent bank contributes
the second edge stock.

\section{A five-factor trace sandwich}
\label{sec:V97-five-factor}

For bounded self-adjoint \(A,B,E\) and bounded \(C_1,C_2\), define
\[
 \Omega_\otimes(A,C_1,B,C_2,E)
 :=
 \sup_{\rho=\rho_1\otimes\cdots\otimes\rho_4}
 \left\|
 \rho^{1/2}AC_1BC_2E\rho^{1/2}
 \right\|_1.
\]

\begin{theorem}[Three-factor endpoint sandwich]
\label{thm:V97-three-factor-sandwich}
If
\[
 \|C_1\|\le L_1,
 \qquad
 \|C_2\|\le L_2,
\]
then
\begin{align}
 \Omega_\otimes(A,C_1,B,C_2,E)
 &\le
 L_1L_2\|B\|
 \sup_{\rho\ {\rm product}}
 \sqrt{\Tr(\rho A^2)\Tr(\rho E^2)}
\label{eq:V97-five-factor-HS}\\
 &\le
 L_1L_2\|B\|
 \sqrt{
 {\cal E}_\otimes(A){\cal E}_\otimes(E)}.
\tag{V97.1}\label{eq:V97-three-factor-sandwich}
\end{align}
No commutation among the five displayed factors or with \(\rho\) is
assumed.
\end{theorem}

\begin{proof}
Put
\[
 X:=\rho^{1/2}AC_1,
 \qquad
 Y:=BC_2E\rho^{1/2}.
\]
Schatten H\"older gives
\[
 \|\rho^{1/2}AC_1BC_2E\rho^{1/2}\|_1
 \le\|X\|_2\|Y\|_2.
\]
Since \(C_1C_1^*\preccurlyeq L_1^2I\) and \(A=A^*\),
\begin{align*}
 \|X\|_2^2
 &=
 \Tr(C_1^*A\rho AC_1)
 =
 \Tr(\rho AC_1C_1^*A)\\
 &\le
 L_1^2\Tr(\rho A^2).
\end{align*}
Similarly,
\[
 C_2^*B^2C_2
 \preccurlyeq
 \|B\|^2L_2^2I,
\]
and hence
\begin{align*}
 \|Y\|_2^2
 &=
 \Tr(\rho E C_2^*B^2C_2E)\\
 &\le
 \|B\|^2L_2^2\Tr(\rho E^2).
\end{align*}
Taking square roots proves
\eqref{eq:V97-five-factor-HS}.  Functional calculus gives
\[
 A^2\preccurlyeq\|A\||A|,
 \qquad
 E^2\preccurlyeq\|E\||E|.
\]
For product \(\rho\), their expectations are bounded by
\({\cal E}_\otimes(A)\) and
\({\cal E}_\otimes(E)\), respectively.  Take the supremum.
\end{proof}

\begin{remark}[Relation to the V94 sandwich]
\label{rem:V97-relation-V94}
If \(B=I\) and \(C_2=I\), then
\eqref{eq:V97-three-factor-sandwich} is the trace-sandwiched estimate
behind \Cref{thm:V94-dual-sandwich}.  V97 does not iterate the V94
capacity bound.  It returns to Schatten H\"older once, places the
new adjacent response at the right endpoint, and keeps the second
repeated occurrence in the middle norm.
\end{remark}

\section{Exact dual Fourier form}
\label{sec:V97-dual-form}

\begin{theorem}[Five-factor exact Fourier sandwich]
\label{thm:V97-five-factor-Fourier}
Suppose the complete multiplier on every final output block has the
equivariant factorization
\[
 K_U=a_Uc_{1,U}b_Uc_{2,U}e_U,
\tag{V97.2}\label{eq:V97-block-five-factor}
\]
where \(a_U,b_U,e_U\) are the restrictions of self-adjoint complete
operators \(A,B,E\), and
\[
 \sup_U\|c_{1,U}\|\le L_1,
 \qquad
 \sup_U\|c_{2,U}\|\le L_2.
\]
Then the V94 exact dual radius satisfies
\begin{equation}
 \boxed{
 \Gamma_\otimes(1,K)
 \le
 L_1L_2\|B\|
 \sqrt{
 {\cal E}_\otimes(A){\cal E}_\otimes(E)}.}
\tag{V97.3}\label{eq:V97-Fourier-five-factor}
\end{equation}
The output dual contractions introduce no additional block or
coordinate sum.
\end{theorem}

\begin{proof}
For a V94 dual family \(\boldsymbol D=(D_U)\), absorb \(D_U^*\) into
the first separator and set
\[
 C_{1,\boldsymbol D}
 :=
 \sum_U
 W_U(D_U^*\otimes c_{1,U})W_U^*.
\]
Then \(\|C_{1,\boldsymbol D}\|\le L_1\).  Let \(C_2\) be the complete
intertwiner with restrictions \(c_{2,U}\).  The exact dual operator is
\[
 {\cal T}_{1,K}(\boldsymbol D)
 =
 A C_{1,\boldsymbol D} B C_2 E.
\]
For every product density \(\rho\),
\[
 |\Tr[\rho{\cal T}_{1,K}(\boldsymbol D)]|
 \le
 \|\rho^{1/2}{\cal T}_{1,K}(\boldsymbol D)\rho^{1/2}\|_1.
\]
Apply \Cref{thm:V97-three-factor-sandwich}, then take both suprema in
\eqref{eq:V94-dual-radius}.
\end{proof}

\begin{firewall}[The endpoint must be a complete occurrence]
\label{fire:V97-complete-endpoint}
The operator \(E\) is retained inside the complete multiplier before
the exact Fourier duality.  It is not first replaced by a scalar
capacity and then multiplied with the V94 overlap bound.  Such a
product of separately optimized capacities is not supplied by the V2
one-sided tensor law and would be vulnerable to its
entanglement-swapping regression.
\end{firewall}

\section{Quadratic edge-retaining energies}
\label{sec:V97-quadratic-energies}

V95 used a uniform norm together with a stock-carrying capacity to
obtain an energy linear in one repeated-bank stock.  The same
envelopes also give a second, deliberately stronger estimate in which
both the norm and capacity retain the edge stock.  Its square root is
exactly what the V97 endpoint sandwich needs.

\begin{lemma}[Paired bounds give quadratic energies]
\label{lem:V97-quadratic-AFI-energy}
For every clean \(T\in\{A,F,I\}\) occurrence of V95,
\begin{align}
 {\cal E}_\otimes(S_T)
 &\le C_\kappa\mathfrak h_T^2,
\label{eq:V97-quadratic-unpaid}\\
 {\cal E}_\otimes(S_T^{\rm pay})
 &\le
 \frac{C_\kappa}{s_\alpha^2}\mathfrak h_T^2.
\tag{V97.4}\label{eq:V97-quadratic-paid}
\end{align}
Both estimates are uniform in support, representations, and left or
right orientation.
\end{lemma}

\begin{proof}
By \eqref{eq:V95-envelope-domination},
\[
 \|S_T\|\le\|{\cal N}_T\|,
 \qquad
 \kappa_\otimes(|S_T|)
 \le\kappa_\otimes({\cal N}_T).
\]
Both right sides are at most
\(C_\kappa\mathfrak h_T\) by
\eqref{eq:V95-AFI-unpaid-paired}.  Multiply and use
\eqref{eq:V94-one-use-energy}.  The paid proof uses
\({\cal P}_T\) and
\eqref{eq:V95-AFI-paid-paired}; both factors then carry
\(s_\alpha^{-1}\mathfrak h_T\).
\end{proof}

\begin{remark}[Linear and quadratic energies coexist]
\label{rem:V97-two-energy-cards}
Equations
\eqref{eq:V95-AFI-unpaid-energy}--%
\eqref{eq:V95-AFI-paid-energy} and
\eqref{eq:V97-quadratic-unpaid}--%
\eqref{eq:V97-quadratic-paid} are different uses of the same
envelopes.  The linear card is optimal for two occurrences of one
physical bank because its square root charges that bank once.  The
quadratic card is used when the two Hilbert--Schmidt endpoints belong
to two distinct physical banks; its square root retains one stock
from each.
\end{remark}

\begin{corollary}[Literal two-edge product]
\label{cor:V97-two-edge-product}
In \Cref{thm:V97-five-factor-Fourier}, suppose \(A\) is one clean
\(A/F/I\) occurrence with stock \(\mathfrak h_A\), \(B\) is the other
occurrence of the repeated bank, and \(E=K_{\rm adj}\) is a clean
\(A/F/I\) occurrence on a coordinate bank disjoint from the repeated
bank, with stock \(\mathfrak h_{\rm adj}\).  Then
\begin{equation}
 \boxed{
 \Gamma_\otimes(1,K)
 \le
 C_\kappa L_1L_2
 \mathfrak h_A\mathfrak h_{\rm adj}.}
\tag{V97.5}\label{eq:V97-two-edge-product}
\end{equation}
Only the \(A\)-occurrence of the repeated bank contributes a
first-moment stock; \(B\) is used through its uniform norm.
\end{corollary}

\begin{proof}
Use \(\|B\|\le C\) from
\eqref{eq:V95-AFI-uniform-norms} and insert the two quadratic
energies \eqref{eq:V97-quadratic-unpaid} into
\eqref{eq:V97-Fourier-five-factor}.
\end{proof}

\section{All payer locations}
\label{sec:V97-payers}

\begin{theorem}[Once-paid three-factor ledger]
\label{thm:V97-paid-three-factor}
Under the hypotheses of
\Cref{cor:V97-two-edge-product}, every mutually exclusive location of
the unique payer has the bound
\begin{equation}
 \boxed{
 \Gamma_\otimes(w,K^{\rm pay})
 \le
 \frac{C_\kappa L_1L_2}{s_\alpha}
 \mathfrak h_A\mathfrak h_{\rm adj},}
\tag{V97.6}\label{eq:V97-paid-two-edge-product}
\end{equation}
with the evident replacement of \(L_1\) or \(L_2\) by the unpaid
separator norm when that separator pays.  The locations are:
\begin{enumerate}[label=\textup{(\roman*)},leftmargin=4em]
\item the edge-carrying repeated occurrence \(A\);
\item the middle repeated occurrence \(B\);
\item the adjacent occurrence \(E\);
\item either finite local separator \(C_1,C_2\); or
\item the complete common suffix.
\end{enumerate}
Only one inverse power of \(s_\alpha\) occurs.
\end{theorem}

\begin{proof}
If \(A\) pays, replace its unpaid quadratic energy by
\eqref{eq:V97-quadratic-paid}.  Its square root contributes
\(\mathfrak h_A/s_\alpha\).  The case in which \(E\) pays is
identical.

If \(B\) pays, use
\(\|B^{\rm pay}\|\le C/s_\alpha\) from
\eqref{eq:V95-AFI-uniform-norms}; both endpoint energies remain
unpaid.  If \(C_i\) pays, the same proof uses its paid norm
\(C/s_\alpha\), supplied by
\Cref{lem:V95-local-separator-payer}, and the other separator remains
bounded.  Finally,
\Cref{thm:V96-support-uniform-suffix} gives \(C/s_\alpha\) for the
grouped suffix and the V2 one-sided tensor law inserts that norm after
the five-factor carrier has been bounded.  The V2 allocation makes
these alternatives disjoint, so their estimates are never multiplied.
\end{proof}

\begin{firewall}[A local separator payer is not a suffix payer]
\label{fire:V97-separate-payer-types}
The local-separator rows of
\Cref{thm:V97-paid-three-factor} use the finite partition-letter bound
of V95.  The suffix row uses the total-output product of V96.  Replacing
the latter by a sum of coordinatewise denominators would reintroduce
the support loss proved in
\Cref{lem:V96-coordinatewise-loss}.
\end{firewall}

\section{Application to the V96 graph rows}
\label{sec:V97-graph-application}

Use the notation
\[
 \{p,p^\circ\}\mid\{r,r^\circ\},
 \qquad
 e_L=pp^\circ,\quad e_R=rr^\circ,\quad e_\times=pr
\]
of V96.  An \emph{adjacent certificate} for a missing same-side edge
means a complete \(A/F/I\) occurrence supported on a bank disjoint from
the repeated witness bank and from the already selected joining bank.

\begin{theorem}[Distinct-adjacent closure of the residual graph rows]
\label{thm:V97-distinct-adjacent-closure}
Assume the V96 ascent-overlap hypotheses and retain the actual common
incidence edge \(e_\times\).  Suppose the exact V61 word has a clean
distinct adjacent certificate on the missing same-side edge.  Then:
\begin{enumerate}[label=\textup{(\roman*)},leftmargin=4em]
\item if \(g=pr^\circ\), retain the right occurrence carrying \(e_R\)
      at the first endpoint, use the other repeated occurrence only in
      norm, and put the adjacent \(e_L\)-certificate at the second
      endpoint;
\item if \(g=p^\circ r\), retain the left occurrence carrying \(e_L\)
      and put the adjacent \(e_R\)-certificate at the second endpoint;
\item if \(g=pr\), either orientation retains the two same-side edges,
      while one copy of the parallel \(pr\) edge is surplus.
\end{enumerate}
In every case the analytic contribution has the literal product of the
two distinct same-side stocks from
\eqref{eq:V97-two-edge-product}, with the once-paid version
\eqref{eq:V97-paid-two-edge-product}.  Together with the genuine
common-incidence and joining stocks, the resulting multigraph is
connected and reduces to a quartic marked tree with a universal fibre.
No comparison between \(\mathfrak h_L\) and
\(\mathfrak h_R\) is required.
\end{theorem}

\begin{proof}
For \(g=pr^\circ\), the four retained edge labels are
\[
 rr^\circ,\quad pp^\circ,\quad pr,\quad pr^\circ.
\]
They form a connected graph.  Apply
\Cref{cor:V97-two-edge-product} to the two same-side endpoint
occurrences and use the middle repeated occurrence only through its
uniform norm.  The other two genuine stocks remain in the exact V92
mark expansion.  One edge is surplus, so
\Cref{lem:V35-surplus-edge} reduces the product to a finite sum of
spanning-tree markings.  The case \(g=p^\circ r\) is symmetric.

For \(g=pr\), the distinct edge labels
\[
 pp^\circ,\quad rr^\circ,\quad pr
\]
already form a spanning tree; the common-incidence and joining copies
of \(pr\) form a parallel surplus edge in the multigraph.  Delete one
normalized copy using the same surplus-edge argument.  The paid
conclusions follow from
\Cref{thm:V97-paid-three-factor}.  The V92 deterministic bank-to-mark
map and the finite choices in the displayed list give a universal
fibre.
\end{proof}

\begin{remark}[The opposite-corner row]
\label{rem:V97-opposite-corner}
For \(g=p^\circ r^\circ\), V96 already found both two-term routes
spanning.  The V97 theorem is unnecessary there, although it remains
available if a distinct adjacent certificate is already present in the
packet.
\end{remark}

\section{A deterministic fresh/self-adjacent split}
\label{sec:V97-adjacent-split}

\begin{lemma}[Every adjacent bank has a two-cell split]
\label{lem:V97-fresh-self-split}
Let \(E_{\rm adj}\) be the complete coordinate bank supporting the
adjacent same-side certificate, and let \({\cal B}\) be the repeated
witness bank.  Put
\[
 E_{\rm adj}^{\rm fr}:=E_{\rm adj}\setminus{\cal B},
 \qquad
 E_{\rm adj}^{\rm self}:=E_{\rm adj}\cap{\cal B}.
\]
Then
\begin{equation}
 H_{\rm adj}(E_{\rm adj})
 =
 H_{\rm adj}(E_{\rm adj}^{\rm fr})
 +
 H_{\rm adj}(E_{\rm adj}^{\rm self}).
\tag{V97.7}\label{eq:V97-adjacent-split}
\end{equation}
At least one summand is at least half the total stock.  In the first
case the hypotheses of
\Cref{thm:V97-distinct-adjacent-closure} hold after a loss of at most
two.  In the second case the obstruction is localized to a
three-occurrence carrier of the same physical bank.
\end{lemma}

\begin{proof}
The two coordinate sets are disjoint and partition
\(E_{\rm adj}\).  The edge stock is a sum of nonnegative coordinate
weights, proving \eqref{eq:V97-adjacent-split}.  The two-term
pigeonhole principle gives the alternative.  A certificate restricted
to the fresh part is disjoint from \({\cal B}\); on the self part, the
two old occurrences and the adjacent response all meet the same bank,
so applying two independent first-moment cards would charge it more
than once.
\end{proof}

\begin{corollary}[Exact residual after V97]
\label{cor:V97-exact-residual}
The one-sided and diagonal rows left by V96 split into:
\begin{enumerate}[label=\textup{(\alph*)},leftmargin=4em]
\item a distinct-adjacent row, closed by
      \Cref{thm:V97-distinct-adjacent-closure}; or
\item a self-adjacent row in which one identified physical bank occurs
      three times in the exact complete multiplier.
\end{enumerate}
No arbitrary family of overlap subsets remains.
\end{corollary}

\begin{proof}
Apply \Cref{lem:V97-fresh-self-split} to the missing same-side
certificate specified by the graph table
\eqref{eq:V96-graph-table}.
\end{proof}

\section{The V96 counterpacket with a genuine adjacent bank}
\label{sec:V97-counterpacket-test}

\begin{proposition}[The three-factor card passes the ascent test]
\label{prop:V97-counterpacket-test}
Extend the physical V96 bank
\eqref{eq:V96-two-coordinate-array} by one disjoint pure
\(pp^\circ\)-coordinate \(a\), placed in the proper prefix, with
weights
\[
 (a_{p,a},a_{p^\circ,a},a_{r,a},a_{r^\circ,a})
 =
 (N,N,0,0).
\]
Then the altered row masses still satisfy the strict ascent
\(\Xi_r>L\Xi_p\) for all sufficiently large \(j\), and the adjacent
response obeys
\[
 \mathfrak h_{\rm adj}
 =
 \frac{N^2}{\Xi_p}
 \asymp N.
\]
In the \(g=pr^\circ\) orientation,
\Cref{cor:V97-two-edge-product} retains the right response
\(\mathfrak h_R\asymp N^2\) and the adjacent response
\(\mathfrak h_{\rm adj}\asymp N\).  Hence its legal two-edge
majorant has order \(N^3\), which dominates the isolated V96
geometric scale \(N^{3/2}\).  The missing \(N^{1/2}\) is therefore
available on the distinct-adjacent branch.
\end{proposition}

\begin{proof}
The new coordinate adds \(N\) to \(\Xi_p,\Xi_{p^\circ}\) and nothing
to \(\Xi_r,\Xi_{r^\circ}\).  Thus
\[
 \Xi_p=2N+b,
 \qquad
 \Xi_r=M+b\asymp N^2,
\]
so the strict ascent persists.  Its pair stock is \(N^2\), giving the
displayed adjacent response.  The right response remains comparable to
\(M\asymp N^2\).  Apply
\eqref{eq:V97-two-edge-product}.  The coordinate \(a\) is disjoint
from \({\cal B}_j\), so no repeated-bank first moment is reused.
\end{proof}

\begin{firewall}[The scale test is not a lower bound]
\label{fire:V97-scale-test-scope}
Proposition~\ref{prop:V97-counterpacket-test} checks that the V97
upper bound does not reproduce the V96 deficit.  It does not assert
that the packet itself has order \(N^3\), nor does it address the
self-adjacent branch in which \(a\) belongs to the repeated bank.
\end{firewall}

\section{Regression tests and exact status}
\label{sec:V97-status}

\begin{proposition}[V15, V86, and V91 regressions]
\label{prop:V97-regression-tests}
\begin{enumerate}[label=\textup{(\roman*)},leftmargin=4em]
\item V15 is unchanged because the exact dual output contraction is
      still absorbed into a complete bounded separator, and no singlet
      capacity or dimension factor is introduced.
\item V86 is strengthened only when a distinct adjacent prefix
      response remains in the complete word.  Its exact equal-spin
      binary heat card is not replaced by an input-stock square tail.
\item V91 retains the path \(\{12,13,34\}\).  If a fourth response is
      present, the surplus-edge deletion occurs after the exact V92
      coordinate marking; it cannot rename the tilted joining card.
\end{enumerate}
\end{proposition}

\begin{proof}
Item \textup{(i)} follows from the exact V94 duality used in
\Cref{thm:V97-five-factor-Fourier}.  Item \textup{(ii)} follows from
the complete endpoint requirement
\Cref{fire:V97-complete-endpoint} and
\Cref{thm:V86-binary-heat-capacity}.  Item \textup{(iii)} is the
literal graph construction in
\Cref{thm:V97-distinct-adjacent-closure} together with
\Cref{prop:V94-marking-compatible}.
\end{proof}

\begin{theorem}[What V97 proves]
\label{thm:V97-status}
Relative to V1--V96, this chapter proves:
\begin{enumerate}[label=\textup{(V97.\arabic*)},leftmargin=6.5em]
\item the noncommutative five-factor endpoint inequality
      \eqref{eq:V97-three-factor-sandwich};
\item its exact unsplit Fourier form
      \eqref{eq:V97-Fourier-five-factor};
\item quadratic unpaid and paid \(A/F/I\) energies carrying the square
      of one edge stock;
\item the literal two-distinct-edge bounds
      \eqref{eq:V97-two-edge-product}--%
      \eqref{eq:V97-paid-two-edge-product};
\item once-paid closure of every core and grouped-suffix payer location
      in that five-factor word;
\item distinct-adjacent closure of all one-sided and diagonal V96 graph
      rows; and
\item the exact fresh/self-adjacent split
      \eqref{eq:V97-adjacent-split}, leaving only one identified
      three-occurrence bank in the unresolved overlap subcell.
\end{enumerate}
\end{theorem}

\begin{proof}
Items \textup{(V97.1)--(V97.5)} are
\Cref{thm:V97-three-factor-sandwich,
thm:V97-five-factor-Fourier,
lem:V97-quadratic-AFI-energy,
cor:V97-two-edge-product,
thm:V97-paid-three-factor}.
Items \textup{(V97.6)--(V97.7)} are
\Cref{thm:V97-distinct-adjacent-closure,
lem:V97-fresh-self-split,cor:V97-exact-residual}.
\end{proof}

\begin{programme}[V98 self-adjacent acquisition order]
\label{prog:V97-V98}
\begin{enumerate}[leftmargin=3em]
\item Form the exact three-occurrence multiplier of the self-adjacent
      bank before any positive envelope:
      \[
       A_1C_1A_2C_2A_3.
      \]
      Keep the three physical roles and all complete multiplicity
      spaces visible.
\item Determine whether one cyclic Schatten or noncommutative
      \(L^3\) estimate reduces this word to a single third-moment
      energy of the bank.  Do not multiply three first-moment
      capacities.
\item Prove the required unpaid and once-paid third-moment cards for
      \(A,F,I\), or construct a physical representation family
      showing which type fails.
\item Test the result on the V15 dual-shared family and on the
      two-coordinate bank \eqref{eq:V96-two-coordinate-array} without
      adding the fresh coordinate of V97.
\item If the third-moment route fails because the row state is modular
      rather than tracial, return to the complete partition-state sum
      before the three occurrences are separated.
\item Once the self-adjacent subcell closes, resume the
      same-side-maximal directed-\(D_k\) branch.  No companion audit is
      due before V100.
\end{enumerate}
\end{programme}

\begin{firewall}[Claim boundary after V97]
\label{fire:V97-boundary}
V97 closes the one-sided and diagonal ascent-overlap rows only when the
missing adjacent response has a clean coordinate bank distinct from
the repeated witness bank.  It does not close the self-adjacent
three-occurrence subcell, the general same-side-maximal directed-cross
branch, the remaining \(3+1\), \(2+1+1\), discrete
first-connectivity, higher MTA, WUFC, or CPM families.  Constructive
cutoff removal, reflection positivity, Osterwalder--Schrader
reconstruction, and a uniform positive continuum spectral gap remain
open.  The full Yang--Mills mass gap remains open.
\end{firewall}

\paragraph{Parent-version record.}
The complete V96 mathematical body was retained before this chapter was
appended.  V96 had \(47221\) source records, \(1620963\) bytes, and
SHA256
\[
 \texttt{e2b74949717d79b41f8f36d886d5f997ed7983bd198f9b1ada94b35d20b265c4}.
\]
No companion audit is due at V97; the next audit is V100.

% =============================================================================
% END APPENDED COMPACT VERSION V97 MATERIAL
% =============================================================================

% The V97 document terminator is intentionally disabled here:
% \end{document}

% =============================================================================
% BEGIN APPENDED COMPACT VERSION V98 MATERIAL
% =============================================================================

\clearpage
\chapter{Self-adjacent overlap closure and a directed-\(D\) middle
criterion}
\label{chap:V98}

\section{Context: an \(L^3\) estimate is unnecessary}
\label{sec:V98-context}

V97 closed every one-sided or diagonal ascent-overlap row for which
the response incident to the otherwise isolated row lies in a bank
disjoint from the repeated witness bank.  Its residual was called
self-adjacent: the same physical bank occurs in the two old locations
and in the adjacent response,
\[
 A_1C_1A_2C_2A_3.
\]
The V97 mandate proposed a cyclic noncommutative \(L^3\) estimate.
That is unnecessary.  The five-factor theorem already permits
\(A_1,A_2,A_3\) to act on the same coordinate Hilbert space.  It makes
no tensor-disjointness assumption.  Put \(A_1,A_3\) at the two
Hilbert--Schmidt endpoints, keep \(A_2\) in ordinary norm, and use the
quadratic endpoint energies.  This is one joint estimate of the
assembled three-occurrence word, not a product of three independently
optimized capacities.

After closing that subcell, this chapter resumes the directed-\(D\)
branch.  The complete V77 crossing covariance is placed in the middle
of the same endpoint sandwich.  Its unpaid norm is at most two and its
once-paid norm is \(O(s_\alpha^{-1})\).  Clean paired responses at the
two endpoints supply the two same-side edges; directed row dominance
inserts the actual crossing edge.  This avoids the false
same-side-maximal crossing square tail of V84--V85.

\section{The exact same-bank three-occurrence estimate}
\label{sec:V98-three-occurrence}

\begin{theorem}[Self-adjacent three-occurrence sandwich]
\label{thm:V98-self-adjacent-sandwich}
Let \({\cal B}\) be one physical coordinate bank.  Suppose three
self-adjoint complete occurrences of that bank have the exact
equivariant multiplier
\[
 K_U=a_{1,U}c_{1,U}a_{2,U}c_{2,U}a_{3,U},
\tag{V98.1}\label{eq:V98-three-occurrence-word}
\]
with
\[
 \sup_U\|c_{i,U}\|\le L_i,\qquad i=1,2.
\]
Assume the first and third occurrences have clean \(A/F/I\)
certificates with row-normalized edge stocks
\(\mathfrak h_1,\mathfrak h_3\), and the middle occurrence has any
clean \(A/F/I\) type.  Then
\begin{equation}
 \boxed{
 \Gamma_\otimes(1,K)
 \le
 C_\kappa L_1L_2
 \mathfrak h_1\mathfrak h_3.}
\tag{V98.2}\label{eq:V98-self-adjacent-unpaid}
\end{equation}
The bank is estimated only after all three occurrences and both
separators have been assembled.
\end{theorem}

\begin{proof}
Apply \Cref{thm:V97-five-factor-Fourier} with
\[
 A=A_1,\qquad B=A_2,\qquad E=A_3.
\]
The middle occurrence has the support-uniform norm
\(\|A_2\|\le C\) by
\eqref{eq:V95-AFI-uniform-norms}.  The endpoint energies obey
\[
 {\cal E}_\otimes(A_1)\le C_\kappa\mathfrak h_1^2,
 \qquad
 {\cal E}_\otimes(A_3)\le C_\kappa\mathfrak h_3^2
\]
by \Cref{lem:V97-quadratic-AFI-energy}.  Their geometric mean is
\(C_\kappa\mathfrak h_1\mathfrak h_3\), proving the claim.
No step factors the row density matrix across occurrences or
coordinates.
\end{proof}

\begin{theorem}[All payer locations in the self-adjacent word]
\label{thm:V98-self-adjacent-paid}
In the setting of
\Cref{thm:V98-self-adjacent-sandwich}, every mutually exclusive
location of the sole payer obeys
\begin{equation}
 \boxed{
 \Gamma_\otimes(w,K^{\rm pay})
 \le
 \frac{C_\kappa L_1L_2}{s_\alpha}
 \mathfrak h_1\mathfrak h_3.}
\tag{V98.3}\label{eq:V98-self-adjacent-paid}
\end{equation}
This includes a payer in any of the three occurrences, either finite
local separator, or the complete common suffix.
\end{theorem}

\begin{proof}
If the payer lies in \(A_1\) or \(A_3\), use the paid quadratic energy
\eqref{eq:V97-quadratic-paid} at that endpoint.  If it lies in \(A_2\),
use its uniform paid norm
\(\|A_2^{\rm pay}\|\le C/s_\alpha\) from
\eqref{eq:V95-AFI-uniform-norms}.  The local-separator rows use
\Cref{lem:V95-local-separator-payer}; the suffix row uses
\Cref{thm:V96-support-uniform-suffix}.  These are precisely the
alternatives in the proof of
\Cref{thm:V97-paid-three-factor}, which did not use disjointness of
the endpoint supports.
\end{proof}

\begin{assessment}[Why this is one joint use]
\label{ass:V98-joint-use}
The two endpoint energy estimates are applied inside the single
Schatten inequality
\eqref{eq:V97-five-factor-HS}.  The middle occurrence is never assigned
a capacity.  Thus V98 does not multiply three first-moment bounds for
\({\cal B}\).  The final product
\(\mathfrak h_1\mathfrak h_3\) represents two distinct edge roles of
one already assembled three-occurrence carrier.
\end{assessment}

\begin{proposition}[The V92 marking identity permits the same bank on
two edges]
\label{prop:V98-same-bank-marking}
Let \(e_1,e_3\) be the two edge labels attached to the first and third
occurrences in
\eqref{eq:V98-three-occurrence-word}.  In the exact bank-to-mark
identity \eqref{eq:V92-bank-to-mark}, it is permissible to take
\[
 E_{e_1}=E_{e_3}={\cal B}.
\]
The product
\[
 H_{e_1}({\cal B})H_{e_3}({\cal B})
\]
expands into ordered pairs of marks in
\({\cal B}\times{\cal B}\), including the diagonal pairs.  The
analytic permission for that algebraic repetition is supplied by
\Cref{thm:V98-self-adjacent-sandwich}; no two independent
one-occurrence capacity bounds are invoked.
\end{proposition}

\begin{proof}
The first statement is exactly
\Cref{fire:V92-mark-overlap}: edge-labelled markings recover both
coordinates even when their banks coincide.  What was missing there
was a joint repeated-bank carrier.  Equations
\eqref{eq:V98-three-occurrence-word}--%
\eqref{eq:V98-self-adjacent-unpaid} provide that carrier and bound it
before the scalar product is expanded.
\end{proof}

\section{Completion of the V96 overlap graph table}
\label{sec:V98-overlap-completion}

\begin{theorem}[Self-adjacent graph closure]
\label{thm:V98-self-adjacent-graph}
In the self-adjacent alternative of
\Cref{lem:V97-fresh-self-split}, use the adjacent occurrence as the
third endpoint of
\eqref{eq:V98-three-occurrence-word}.  Then every one-sided or
diagonal row of the V96 graph table has the same literal two-edge
product and once-paid estimate as in
\Cref{thm:V97-distinct-adjacent-closure}.  Explicitly:
\begin{enumerate}[label=\textup{(\roman*)},leftmargin=4em]
\item for \(g=pr^\circ\), the endpoints carry
      \(rr^\circ\) and \(pp^\circ\);
\item for \(g=p^\circ r\), the endpoints carry
      \(pp^\circ\) and \(rr^\circ\); and
\item for \(g=pr\), the two same-side endpoint edges together with one
      copy of \(pr\) form a spanning tree.
\end{enumerate}
The common-incidence and joining edges are handled exactly as in V97.
\end{theorem}

\begin{proof}
The graph argument is
\Cref{thm:V97-distinct-adjacent-closure}; only the analytic support
relation has changed.  Replace
\eqref{eq:V97-two-edge-product} by
\eqref{eq:V98-self-adjacent-unpaid}, and replace its paid version by
\eqref{eq:V98-self-adjacent-paid}.  The common incidence response is
bounded below by \eqref{eq:V96-directed-lower}, so it may be inserted
with the fixed factor \((\beta a_*)^{-1}\).  The V92 edge-bank
expansion remains legal by
\Cref{prop:V98-same-bank-marking}.  Surplus-edge deletion then produces
the stated trees.
\end{proof}

\begin{corollary}[Clean \(A/F/I\) ascent overlaps are closed]
\label{cor:V98-clean-AFI-overlap-closed}
Within the clean \(S_{22}\) calculus, suppose every repeated and
adjacent occurrence in the V96 ascent-overlap word has one of the
complete \(A,F,I\) certificates of V95 and the remaining selected
banks satisfy the stated exact V61 factorization.  Then all four
opposite-block rows of \eqref{eq:V96-graph-table}, all core payer
locations, and the grouped common-suffix payer satisfy the literal
quartic marked-tree bound with a universal routing fibre.
\end{corollary}

\begin{proof}
The opposite-corner row is V96.  The distinct-adjacent subcell is
\Cref{thm:V97-distinct-adjacent-closure}; the self-adjacent subcell is
\Cref{thm:V98-self-adjacent-graph}.  These two subcells are exhaustive
by \Cref{lem:V97-fresh-self-split}.  Their payer ledgers are
\Cref{thm:V97-paid-three-factor,thm:V98-self-adjacent-paid}, and the
suffix estimate is \Cref{thm:V96-support-uniform-suffix}.
\end{proof}

\begin{proposition}[The unextended V96 counterpacket now passes]
\label{prop:V98-counterpacket-test}
For the physical two-coordinate bank
\eqref{eq:V96-two-coordinate-array}, suppose its three exact source
roles are the left occurrence, the middle repeated occurrence, and
the right adjacent occurrence.  Then
\[
 \mathfrak h_L\asymp N,
 \qquad
 \mathfrak h_R\asymp N^2,
\]
and
\Cref{thm:V98-self-adjacent-sandwich} gives the legal two-edge
majorant
\[
 C\mathfrak h_L\mathfrak h_R\asymp N^3.
\]
Thus the physical counterpacket which defeated the incidence-only
spanning term in V96 does not defeat the assembled self-adjacent
carrier.
\end{proposition}

\begin{proof}
The response asymptotics are
\eqref{eq:V96-counter-responses}.  Insert them in
\eqref{eq:V98-self-adjacent-unpaid}.  The two edge marks are legal by
\Cref{prop:V98-same-bank-marking}.  This is an upper-bound regression
only; it asserts no matching lower order.
\end{proof}

\section{The assembled directed covariance as a middle factor}
\label{sec:V98-directed-middle}

Fix the cut \(C=\{1,2\}\), \(D=\{3,4\}\), and let \(B_D\) be a
complete V77 residual bank.  Let
\[
 {\mathfrak J}_{B_D}^{C|D}
 =
 \operatorname{Cov}_{B_D}(F_C^{B_D},F_D^{B_D})
\]
be its assembled signed covariance.  It is a difference of two
complete contractions, and therefore
\begin{equation}
 \|{\mathfrak J}_{B_D}^{C|D}\|\le2.
\tag{V98.4}\label{eq:V98-directed-unpaid-norm}
\end{equation}
Its physical once-paid envelope satisfies
\begin{equation}
 \|{\mathfrak P}_{B_D}^{C|D}\|
 \le\frac C{s_\alpha}
\tag{V98.5}\label{eq:V98-directed-paid-norm}
\end{equation}
by \Cref{lem:V78-universal-paid-cap}.

\begin{lemma}[A dominant directed cell supplies its own crossing edge]
\label{lem:V98-directed-edge-lower}
If a V79 directed cell obeys
\[
 D_{k\to\ell}(B_D)\ge\gamma\Xi_k,
 \qquad
 k\in C,\quad\ell\in D
\]
or the symmetric orientation, then
\begin{equation}
 \boxed{
 d_{k\to\ell}(B_D)
 :=
 \frac{H_{k\ell}(B_D)}{\Xi_k}
 \ge\gamma a_*.}
\tag{V98.6}\label{eq:V98-directed-edge-lower}
\end{equation}
\end{lemma}

\begin{proof}
Every incidence assigned to \(k\to\ell\) is active at both rows.
Therefore
\[
 H_{k\ell}(B_D)
 \ge
 a_*D_{k\to\ell}(B_D)
 \ge
 \gamma a_*\Xi_k
\]
by the same coordinatewise argument as
\Cref{lem:V78-pair-dominates}.  Divide by \(\Xi_k\).
\end{proof}

\begin{theorem}[Directed-\(D\) middle criterion]
\label{thm:V98-directed-middle}
Suppose the two same-side blocks have clean complete \(A/F/I\)
endpoint occurrences \(S_C,S_D\) with edge stocks
\(\mathfrak h_C,\mathfrak h_D\).  Suppose the complete source
multiplier factors blockwise as
\[
 S_C\,C_1\,
 {\mathfrak J}_{B_D}^{C|D}\,
 C_2\,S_D,
\tag{V98.7}\label{eq:V98-directed-middle-word}
\]
where \(\|C_i\|\le L_i\), and suppose
\eqref{eq:V98-directed-edge-lower} holds for an actual crossing pair
\(k\ell\).  Then the unpaid exact Fourier radius satisfies
\begin{equation}
 \boxed{
 \Gamma_\otimes(1,K_D)
 \le
 \frac{C_\kappa L_1L_2}{\gamma a_*}
 d_{k\to\ell}(B_D)
 \mathfrak h_C\mathfrak h_D.}
\tag{V98.8}\label{eq:V98-directed-middle-unpaid}
\end{equation}
The three displayed response stocks are carried by the literal path
with the two same-side edges and crossing edge \(k\ell\).
\end{theorem}

\begin{proof}
Apply \Cref{thm:V97-five-factor-Fourier} with
\[
 A=S_C,\qquad
 B={\mathfrak J}_{B_D}^{C|D},\qquad
 E=S_D.
\]
The middle norm is at most two by
\eqref{eq:V98-directed-unpaid-norm}; the endpoint quadratic energies
are bounded by
\(C_\kappa\mathfrak h_C^2\) and
\(C_\kappa\mathfrak h_D^2\).
Thus
\[
 \Gamma_\otimes(1,K_D)
 \le
 C_\kappa L_1L_2\mathfrak h_C\mathfrak h_D.
\]
By \eqref{eq:V98-directed-edge-lower},
\[
 1\le
 \frac{d_{k\to\ell}(B_D)}{\gamma a_*}.
\]
Insert this factor.  The endpoint stocks mark the two internal block
edges and \(H_{k\ell}(B_D)\) marks the actual crossing edge.  Apply
\eqref{eq:V92-S22-path-expansion}.
\end{proof}

\begin{theorem}[Once-paid directed-middle criterion]
\label{thm:V98-directed-middle-paid}
Every payer location in
\eqref{eq:V98-directed-middle-word}, including either endpoint, the
assembled crossing covariance, a finite local separator, or the
complete common suffix, obeys
\begin{equation}
 \boxed{
 \Gamma_\otimes(w,K_D^{\rm pay})
 \le
 \frac{C_\kappa L_1L_2}
      {s_\alpha\gamma a_*}
 d_{k\to\ell}(B_D)
 \mathfrak h_C\mathfrak h_D.}
\tag{V98.9}\label{eq:V98-directed-middle-paid}
\end{equation}
\end{theorem}

\begin{proof}
An endpoint payer uses
\eqref{eq:V97-quadratic-paid}.  If the assembled crossing covariance
pays, replace \eqref{eq:V98-directed-unpaid-norm} by
\eqref{eq:V98-directed-paid-norm}.  A finite separator payer uses
\Cref{lem:V95-local-separator-payer}, and a suffix payer uses
\Cref{thm:V96-support-uniform-suffix}.  In every case the proof of
\Cref{thm:V97-three-factor-sandwich} contributes exactly one
\(s_\alpha^{-1}\).  Finally insert
\eqref{eq:V98-directed-edge-lower} as in the unpaid proof.
\end{proof}

\begin{corollary}[Two-sided paired closure of the same-side-maximal
directed branch]
\label{cor:V98-directed-paired-closure}
Suppose the V92 residual
\[
 D_k>\frac{\delta_\eta}{5}\Xi_k
\]
occurs.  Partition it by its two actual opposite-block partners.
For some \(\ell\in k^c\),
\[
 D_{k\to\ell}>
 \frac{\delta_\eta}{10}\Xi_k.
\]
If both same-side endpoint blocks in the corresponding exact V61 word
have clean \(A/F/I\) certificates, then the same-side-maximal directed
configuration satisfies the unpaid and once-paid literal marked-tree
bounds
\eqref{eq:V98-directed-middle-unpaid}--%
\eqref{eq:V98-directed-middle-paid}, without any cross-retained square
tail.
\end{corollary}

\begin{proof}
There are only two opposite-block partners, so the displayed
pigeonhole bound holds.  Apply
\Cref{thm:V98-directed-middle,thm:V98-directed-middle-paid} with
\(\gamma=\delta_\eta/10\).  The proof uses the entire assembled
crossing covariance in the middle norm; it never replaces the
same-side maximal stock by a crossing-only numerator.
\end{proof}

\begin{firewall}[The remaining directed row is endpoint acquisition]
\label{fire:V98-directed-endpoint-open}
Corollary~\ref{cor:V98-directed-paired-closure} does not assert that
both endpoint blocks always possess clean \(A/F/I\) certificates.
If one endpoint is private, unbalanced, overlapping with the directed
bank, or itself directed without a paired response, its established
cell compiler or a new acquisition step must be used.  The analytic
same-side-maximal joining covariance is no longer the obstruction on
the two-sided paired subcell; the residual is now an explicitly typed
endpoint-acquisition problem.
\end{firewall}

\section{Regression tests and exact status}
\label{sec:V98-status}

\begin{proposition}[V15, V86, and V91 regressions]
\label{prop:V98-regression-tests}
\begin{enumerate}[label=\textup{(\roman*)},leftmargin=4em]
\item V15 permits one physical coordinate to mark two distinct tree
      edges, but only after a joint operator estimate.  V98 supplies
      such an estimate and does not replace the raw dual Fourier mass
      by singlet capacity.
\item The V86 fixed-low-prefix packet is compatible with
      \eqref{eq:V98-directed-middle-word}: its two binary prefix
      responses are the endpoint factors, while the joining covariance
      is used only through its complete norm.  No false isolated
      crossing square tail is restored.
\item The V91 path \(\{12,13,34\}\) is the three-edge specialization
      of \eqref{eq:V98-directed-middle-unpaid}.  The tilted coherent
      joining coordinate remains the actual \(13\) mark.
\end{enumerate}
\end{proposition}

\begin{proof}
Item \textup{(i)} uses
\Cref{prop:V98-same-bank-marking,thm:V94-exact-block-duality}.
Item \textup{(ii)} follows from
\Cref{thm:V86-binary-heat-capacity,
thm:V97-three-factor-sandwich}.  Item \textup{(iii)} is the literal
path assignment in
\Cref{thm:V92-V91-paid,thm:V98-directed-middle}.
\end{proof}

\begin{theorem}[What V98 proves]
\label{thm:V98-status}
Relative to V1--V97, this chapter proves:
\begin{enumerate}[label=\textup{(V98.\arabic*)},leftmargin=6.5em]
\item the V97 five-factor theorem already supplies the exact
      self-adjacent three-occurrence estimate
      \eqref{eq:V98-self-adjacent-unpaid};
\item every payer location in that word has the single-debit estimate
      \eqref{eq:V98-self-adjacent-paid};
\item coincident physical edge banks are compatible with the V92
      marking identity after this joint analytic estimate;
\item the self-adjacent subcell closes the final clean \(A/F/I\)
      one-sided and diagonal overlap rows;
\item the assembled V77 covariance has the middle-factor norms
      \eqref{eq:V98-directed-unpaid-norm}--%
      \eqref{eq:V98-directed-paid-norm};
\item a dominant V79 directed cell supplies the literal crossing
      response \eqref{eq:V98-directed-edge-lower}; and
\item every same-side-maximal directed row with clean paired
      certificates at both endpoints satisfies the unpaid and
      once-paid path bounds
      \eqref{eq:V98-directed-middle-unpaid}--%
      \eqref{eq:V98-directed-middle-paid}.
\end{enumerate}
\end{theorem}

\begin{proof}
Items \textup{(V98.1)--(V98.4)} are
\Cref{thm:V98-self-adjacent-sandwich,
thm:V98-self-adjacent-paid,
prop:V98-same-bank-marking,
thm:V98-self-adjacent-graph,
cor:V98-clean-AFI-overlap-closed}.
Items \textup{(V98.5)--(V98.7)} are
\Cref{eq:V98-directed-unpaid-norm,
lem:V78-universal-paid-cap,
lem:V98-directed-edge-lower,
thm:V98-directed-middle,
thm:V98-directed-middle-paid,
cor:V98-directed-paired-closure}.
\end{proof}

\begin{programme}[V99 endpoint-acquisition order]
\label{prog:V98-V99}
\begin{enumerate}[leftmargin=3em]
\item At a directed-\(D_k\) source failing the two-sided paired
      hypothesis, apply the exact V79 six-cell identity separately at
      each missing endpoint after removing all banks already used in
      \eqref{eq:V98-directed-middle-word}.
\item Route fresh private and unbalanced cells through their existing
      complete moment compilers; route fresh \(A,F,I\) cells through
      the V95 quadratic endpoint interface.
\item If the new endpoint cell is directed, retain its actual partner
      and construct the finite two-directed-edge graph table before
      using mass ascent.  Do not merge its edge stock into
      \(H_B(C,D)\).
\item If acquisition overlaps the middle directed bank, form the
      resulting repeated directed covariance word before taking a
      capacity.  Test whether the V98 middle norm and V94 endpoint
      energy already give a one-use joint estimate.
\item Recheck the V83 four-row patterns and the V84 spin-one
      intermediate-output counterexample.  No companion theorem is
      imported.
\item No companion audit is due at V99.  V100 must perform the audit,
      freeze this sequence as \texttt{\_V100\_f2.tex}, and restart at
      contextual proof-indexed V1.
\end{enumerate}
\end{programme}

\begin{firewall}[Claim boundary after V98]
\label{fire:V98-boundary}
V98 closes the clean self-adjacent \(A/F/I\) overlap and the
two-sided-paired subcell of the same-side-maximal directed-\(D\)
branch.  It does not prove endpoint acquisition for every directed
source, the remaining \(3+1\), \(2+1+1\), discrete
first-connectivity, higher MTA, WUFC, or CPM families.  Constructive
cutoff removal, reflection positivity, Osterwalder--Schrader
reconstruction, and a uniform positive continuum spectral gap remain
open.  The full Yang--Mills mass gap remains open.
\end{firewall}

\paragraph{Parent-version record.}
The complete V97 mathematical body was retained before this chapter was
appended.  V97 had \(47865\) source records, \(1643309\) bytes, and
SHA256
\[
 \texttt{1bb58202555a58de468eab7b5fe8eb1597ca0ed5fb4b478a0348a262cac90368}.
\]
No companion audit is due at V98; the next audit is V100.

% =============================================================================
% END APPENDED COMPACT VERSION V98 MATERIAL
% =============================================================================

% The V98 document terminator is intentionally disabled here:
% \end{document}

% =============================================================================
% BEGIN APPENDED COMPACT VERSION V99 MATERIAL
% =============================================================================

\clearpage
\chapter{Endpoint acquisition and the finite all-directed graph}
\label{chap:V99}

\section{Context: large directed responses may stay in norm}
\label{sec:V99-context}

V98 placed one assembled directed covariance between two paired
same-side endpoint responses.  If a same-side endpoint has no clean
\(A/F/I\) certificate, the V84 acquisition theorem must be applied at
the missing row.  A newly acquired directed cell supplies a second
crossing covariance rather than the missing same-side edge.

This does not require a root-specific capacity estimate for that
covariance.  Such an estimate would return to the false V84
cross-retained square tail.  Each assembled V77 covariance is a
uniformly bounded signed operator, while dominance of its directed
cell gives a fixed lower bound for the actual oriented crossing
response.  We therefore keep every directed covariance inside one
bounded complete word, estimate only a surviving paired endpoint when
one exists, and insert the directed edge responses afterward through
their lower bounds.

When no paired endpoint survives, the whole word is bounded.  Three
dominant directed responses close the source exactly when their
oriented edges form a spanning arborescence.  On four rows this is a
finite graph condition.  V99 proves the criterion and isolates the
only remaining disconnected matching patterns for the V100 audit and
series restart.

\section{Bounded middle words with one edge endpoint}
\label{sec:V99-one-endpoint}

\begin{lemma}[One paired endpoint behind a bounded word]
\label{lem:V99-one-endpoint-sandwich}
Let \(S_T=S_T^*\) be a clean \(A/F/I\) occurrence with stock
\(\mathfrak h_T\), and let \(B\) be any complete equivariant word
with \(\|B\|\le L\).  If the exact output multiplier is
\[
 K_U=b_Us_{T,U},
\]
with the V94 output dualizer absorbed into \(b_U\), then
\begin{equation}
 \boxed{
 \Gamma_\otimes(1,K)
 \le
 C_\kappa L\mathfrak h_T.}
\tag{V99.1}\label{eq:V99-one-endpoint}
\end{equation}
If the sole payer lies in \(S_T\) or in one factor of \(B\) whose
complete paid norm is \(O(s_\alpha^{-1})\), then
\begin{equation}
 \boxed{
 \Gamma_\otimes(w,K^{\rm pay})
 \le
 \frac{C_\kappa L}{s_\alpha}\mathfrak h_T.}
\tag{V99.2}\label{eq:V99-one-endpoint-paid}
\end{equation}
\end{lemma}

\begin{proof}
Apply \Cref{thm:V97-three-factor-sandwich} with
\[
 A=I,\qquad C_1=I,\qquad B=B,\qquad C_2=I,\qquad E=S_T.
\]
Since \({\cal E}_\otimes(I)=1\) and
\({\cal E}_\otimes(S_T)\le C_\kappa\mathfrak h_T^2\) by
\eqref{eq:V97-quadratic-unpaid}, the square root gives
\eqref{eq:V99-one-endpoint}.  A payer in \(S_T\) uses
\eqref{eq:V97-quadratic-paid}; a payer in \(B\) replaces \(L\) by
\(CL/s_\alpha\).
\end{proof}

\begin{definition}[Oriented directed response]
\label{def:V99-oriented-response}
For a dominant directed cell with tail \(u\), partner \(v\), and
complete bank \(G\), write
\[
 D_{u\to v}(G)\ge\gamma\Xi_u,
 \qquad
 d_{u\to v}(G):=\frac{H_{uv}(G)}{\Xi_u}.
\]
Then
\begin{equation}
 d_{u\to v}(G)\ge\gamma a_*
\tag{V99.3}\label{eq:V99-oriented-lower}
\end{equation}
by \Cref{lem:V98-directed-edge-lower}.
\end{definition}

\begin{lemma}[Finite products of assembled directed covariances]
\label{lem:V99-directed-word-norm}
Let \({\mathfrak J}_1,\ldots,{\mathfrak J}_q\), \(q\le4\), be
complete V77 covariances, separated by complete V61 words of uniformly
bounded norm.  Their assembled product \(B_D\) satisfies
\[
 \|B_D\|\le C_4.
\tag{V99.4}\label{eq:V99-directed-word-norm}
\]
If the sole payer is assigned to one of the \({\mathfrak J}_i\), to a
finite local separator, or to the complete common suffix, the
corresponding assembled paid word has norm at most
\[
 \frac{C_4}{s_\alpha}.
\tag{V99.5}\label{eq:V99-directed-word-paid-norm}
\]
\end{lemma}

\begin{proof}
Each unpaid covariance is a difference of two complete contractions,
so its norm is at most two by
\eqref{eq:V98-directed-unpaid-norm}.  There are at most four such
factors, and every unpaid V61 separator has a universal finite norm by
\Cref{lem:V92-separator-size}.  Multiply these finitely many bounds.

If one covariance pays, use
\eqref{eq:V98-directed-paid-norm}; every other covariance remains
unpaid.  A finite local payer uses
\Cref{lem:V95-local-separator-payer}, while a suffix payer uses
\Cref{thm:V96-support-uniform-suffix}.  The payer alternatives are
disjoint, so only one factor \(s_\alpha^{-1}\) occurs.
\end{proof}

\section{One missing same-side endpoint}
\label{sec:V99-one-missing}

\begin{theorem}[Two directed edges and one paired edge close]
\label{thm:V99-two-directed-one-paired}
Write the cut as
\[
 \{p,p^\circ\}\mid\{r,r^\circ\}.
\]
Suppose:
\begin{enumerate}[label=\textup{(\alph*)},leftmargin=4em]
\item the original directed source has a dominant edge
      \(p\to v\), \(v\in\{r,r^\circ\}\), with response \(d_1\);
\item acquisition at the missing row \(p^\circ\) produces a dominant
      fresh or assembled edge
      \(p^\circ\to w\), \(w\in\{r,r^\circ\}\), with response \(d_2\);
      and
\item the right block has a clean paired certificate on \(rr^\circ\),
      oriented from its certified row \(z\in\{r,r^\circ\}\), with
      stock
      \(\mathfrak h_D=H_{rr^\circ}/\Xi_z\).
\end{enumerate}
Then the exact unpaid source obeys
\begin{equation}
 \boxed{
 \Gamma_\otimes(1,K)
 \le
 \frac{C_\kappa}
      {\gamma_1\gamma_2a_*^2}
 d_1d_2\mathfrak h_D.}
\tag{V99.6}\label{eq:V99-two-D-one-paired}
\end{equation}
The three oriented edges form a spanning tree.  Every payer location
obeys the same estimate with exactly one additional
\(s_\alpha^{-1}\).
\end{theorem}

\begin{proof}
The undirected edges are
\[
 pv,\qquad p^\circ w,\qquad rr^\circ.
\]
The right edge connects \(r\) to \(r^\circ\); each left vertex is
attached to that connected pair.  Thus the three edges form a spanning
tree, even when \(v=w\).

Assemble the two directed covariances and all intervening complete
words into \(B_D\).  Apply
\Cref{lem:V99-one-endpoint-sandwich} to \(B_D S_D\), obtaining
\[
 \Gamma_\otimes(1,K)\le C_\kappa\mathfrak h_D.
\]
By \eqref{eq:V99-oriented-lower},
\[
 1\le
 \frac{d_i}{\gamma_i a_*},
 \qquad i=1,2.
\]
Insert both factors to obtain
\eqref{eq:V99-two-D-one-paired}.  The denominators
\(\Xi_p,\Xi_{p^\circ},\Xi_z\) orient the tree toward the right vertex
different from \(z\), so the product is exactly a legal oriented
marked-tree response.  The paid assertion follows from
\eqref{eq:V99-one-endpoint-paid} and
\eqref{eq:V99-directed-word-paid-norm}.
\end{proof}

\begin{corollary}[The first endpoint-acquisition row is closed]
\label{cor:V99-first-endpoint-closed}
At a V98 directed source with one clean paired side and one missing
same-side endpoint, a fresh dominant \(D\)-cell at the missing row
closes the source by
\Cref{thm:V99-two-directed-one-paired}.  Fresh \(A,F,I\) cells close
by \Cref{thm:V98-directed-middle}; private and unbalanced cells remain
in their previously proved complete moment compilers.  Thus this
acquisition step introduces no new analytic square-tail target.
\end{corollary}

\begin{proof}
The fresh-cell alternatives are the V84/V79 incidence alternatives.
The directed alternative has its tail at the missing row and hence
satisfies the graph hypotheses of
\Cref{thm:V99-two-directed-one-paired}.  The other alternatives cite
the stated compilers.
\end{proof}

\section{The fully bounded all-directed criterion}
\label{sec:V99-all-directed}

\begin{lemma}[Three distinct edges of \(K_{2,2}\) form a tree]
\label{lem:V99-K22-tree}
Any three distinct edges of the complete bipartite graph
\(K_{\{p,p^\circ\},\{r,r^\circ\}}\) form a spanning tree.  If the
three edges have distinct directed tails, their orientations form an
arborescence toward the unique vertex which is not a tail.
\end{lemma}

\begin{proof}
Deleting one edge from the four-cycle \(K_{2,2}\) leaves a connected
acyclic graph with four vertices and three edges, hence a tree.
If the tails are distinct, every vertex except one has out-degree one.
An oriented cycle would give an undirected cycle, which the tree does
not contain.  Therefore every directed path terminates at the unique
vertex of out-degree zero.
\end{proof}

\begin{theorem}[Bounded all-directed tree compiler]
\label{thm:V99-all-directed-tree}
Suppose an exact complete source word contains three dominant directed
covariances whose directed edges
\[
 u_i\to v_i,\qquad i=1,2,3,
\]
have distinct tails and distinct undirected labels.  Let
\[
 d_i=\frac{H_{u_iv_i}(G_i)}{\Xi_{u_i}},
 \qquad
 d_i\ge\gamma_i a_*.
\]
All remaining factors are complete uniformly bounded words.  Then
\begin{align}
 \Gamma_\otimes(1,K)
 &\le
 \frac{C_4}{a_*^3\gamma_1\gamma_2\gamma_3}
 d_1d_2d_3,
\label{eq:V99-all-D-unpaid}\\
 \Gamma_\otimes(w,K^{\rm pay})
 &\le
 \frac{C_4}
 {s_\alpha a_*^3\gamma_1\gamma_2\gamma_3}
 d_1d_2d_3.
\tag{V99.7}\label{eq:V99-all-D-paid}
\end{align}
The product on the right is the response product of a legal oriented
quartic marked tree.
\end{theorem}

\begin{proof}
Absorb the V94 output dualizer into the complete bounded word.  By
\Cref{lem:V99-directed-word-norm},
\[
 \Gamma_\otimes(1,K)\le\|K\|\le C_4.
\]
Multiply by
\[
 1\le
 \prod_{i=1}^3\frac{d_i}{\gamma_i a_*}.
\]
By \Cref{lem:V99-K22-tree}, the three edges form an arborescence whose
nonroot vertices are exactly the three tails.  Hence their three row
denominators are the oriented tree denominators, and the V92
bank-to-mark expansion applies.  If the payer lies anywhere in the
word, use \eqref{eq:V99-directed-word-paid-norm} before inserting the
same three lower responses.
\end{proof}

\begin{corollary}[More than three directed occurrences]
\label{cor:V99-more-directed}
If a word contains at most four dominant directed occurrences and
some three of them satisfy
\Cref{thm:V99-all-directed-tree}, those three pay the marked tree.
Every additional occurrence remains inside the bounded middle word
and costs only a universal norm factor.
\end{corollary}

\begin{proof}
Select the three compatible occurrences before applying the scalar
lower bounds.  All other complete covariances have norm at most two
and are included in \(C_4\).
\end{proof}

\section{The finite disconnected residual}
\label{sec:V99-disconnected}

\begin{proposition}[Classification of failed directed edge sets]
\label{prop:V99-disconnected-classification}
Let \({\cal G}\) be the undirected set of cross-pair labels carried by
all dominant directed occurrences acquired at distinct missing rows.
If \({\cal G}\) contains no spanning tree, then exactly one of the
following holds:
\begin{enumerate}[label=\textup{(\roman*)},leftmargin=4em]
\item some row is incident to no edge of \({\cal G}\); or
\item every row is incident, but
      \[
       {\cal G}\subseteq\{pr,p^\circ r^\circ\}
       \quad\hbox{or}\quad
       {\cal G}\subseteq\{pr^\circ,p^\circ r\}.
      \tag{V99.8}\label{eq:V99-disjoint-matching}
      \]
\end{enumerate}
Thus a fully incident failure is one of the two disjoint perfect
matchings, possibly with repeated occurrences of its two edges.
\end{proposition}

\begin{proof}
If every row is incident and the graph is disconnected, each connected
component must contain one left and one right vertex because all edges
are crossing.  Hence there are exactly two components, each a single
cross-pair label.  The only two partitions of \(K_{2,2}\) into such
disjoint pairs are those in
\eqref{eq:V99-disjoint-matching}.  Conversely, either displayed edge
set is disconnected.  If the graph is connected, it has a spanning
tree, and because \(K_{2,2}\) has only four edges, that tree consists
of three distinct labels.
\end{proof}

\begin{firewall}[Mass ascent does not connect a matching]
\label{fire:V99-ascent-not-connectivity}
The strict inequalities among the row masses order vertices, but they
do not add an edge between the two components of
\eqref{eq:V99-disjoint-matching}.  Applying
\Cref{lem:V83-ascent-terminates} to this state before acquiring a
partition or prefix response that joins the components would repeat
the termination-without-acquisition error prohibited in V83.
\end{firewall}

\section{Endpoint-acquisition reduction}
\label{sec:V99-reduction}

\begin{theorem}[Exact clean endpoint reduction before V100]
\label{thm:V99-endpoint-reduction}
Start from a V92 dominant directed-\(D_k\) source and remove every
bank already used in its assembled V98 word before applying the V84
fresh-cell theorem at a missing endpoint.  Within the clean \(SU(2)\)
\(S_{22}\) calculus, every resulting branch is one of:
\begin{enumerate}[label=\textup{(\alph*)},leftmargin=4em]
\item a private or unbalanced complete-moment branch handled by its
      established compiler;
\item a fresh \(A,F,I\) endpoint, handled by the V95 quadratic energy
      and V98 directed-middle criterion;
\item one missing endpoint followed by a fresh directed cell, handled
      by \Cref{thm:V99-two-directed-one-paired};
\item an all-directed word containing an oriented spanning tree,
      handled by \Cref{thm:V99-all-directed-tree};
\item an overlap with a previously used \(A,F,I\) bank, handled by
      V98, or with a directed bank, retained as one bounded complete
      word and handled whenever its edge set falls under
      \textup{(iii)} or \textup{(iv)};
\item one of the two disconnected matchings
      \eqref{eq:V99-disjoint-matching}, or a row not yet incident to
      an acquired edge; or
\item a nonclean local representation state outside the standing
      compiler.
\end{enumerate}
Items \textup{(f)--(g)} are the only new residuals.
\end{theorem}

\begin{proof}
The depleted incidence identity
\eqref{eq:V84-depleted-identity} and
\Cref{thm:V84-fresh-cell} give the fresh/overlap dichotomy and the six
cell types.  Route \(P,U,A,F,I\) as stated.  A fresh \(D\)-cell has an
actual directed partner by the V79 deterministic assignment.  With
one paired side it is item \textup{(c)}; without a paired side,
continue acquisition at a missing row.  If three acquired directed
edges form a compatible tree, use item \textup{(d)}.  Otherwise
\Cref{prop:V99-disconnected-classification} gives item \textup{(f)}.

An overlap with \(A,F,I\) is
\Cref{cor:V98-clean-AFI-overlap-closed}.  Repeated directed
covariances are not split: their complete product has the norm in
\Cref{lem:V99-directed-word-norm}, so only their finite edge-label set
matters.  The standing clean/nonclean split gives item \textup{(g)}.
\end{proof}

\section{Regression tests}
\label{sec:V99-regressions}

\begin{proposition}[V83 and V84 tests, with V15, V86, and V91]
\label{prop:V99-regression-tests}
\begin{enumerate}[label=\textup{(\roman*)},leftmargin=4em]
\item The two-large/two-small V83 pattern is routed only when the
      acquired directed labels and a paired edge form an actual tree.
      A same-side maximal mass is never renamed as a cross edge.
\item The V84 spin-one intermediate-output family is not assigned a
      small crossing square tail.  A directed covariance is used in
      norm only on a branch where its oriented response has the fixed
      lower bound \eqref{eq:V99-oriented-lower}.
\item V15 retains the literal path \(\{12,13,24\}\), including any
      repeated coordinate marks, after a joint complete-word bound.
\item V86 retains its exact binary-prefix endpoint response whenever
      that response is the sole paired edge in
      \eqref{eq:V99-two-D-one-paired}.
\item V91 retains its actual \(13\) joining mark; the all-directed
      compiler never changes an edge label after the graph is fixed.
\end{enumerate}
\end{proposition}

\begin{proof}
Items \textup{(i)--(ii)} are the graph and dominance hypotheses of
\Cref{thm:V99-two-directed-one-paired,
thm:V99-all-directed-tree}; compare
\Cref{prop:V83-four-row-test,thm:V84-cross-square-false}.
Items \textup{(iii)--(v)} follow from the exact V92 marking identity,
the V86 endpoint card, and
\Cref{thm:V92-V91-paid}, respectively.
\end{proof}

\section{Exact status and the compulsory V100 transition}
\label{sec:V99-status}

\begin{theorem}[What V99 proves]
\label{thm:V99-status}
Relative to V1--V98, this chapter proves:
\begin{enumerate}[label=\textup{(V99.\arabic*)},leftmargin=6.5em]
\item a bounded complete word followed by one paired endpoint has the
      literal edge estimate
      \eqref{eq:V99-one-endpoint} and its once-paid analogue;
\item finite products of assembled directed covariances have the
      support-uniform norms
      \eqref{eq:V99-directed-word-norm}--%
      \eqref{eq:V99-directed-word-paid-norm};
\item two dominant directed edges rooted at the two missing same-side
      rows, together with one paired opposite edge, form a legal
      oriented tree and obey
      \eqref{eq:V99-two-D-one-paired};
\item three dominant directed occurrences with distinct tails and
      distinct cross labels obey the all-directed tree estimate
      \eqref{eq:V99-all-D-paid};
\item every failure of cross-edge connectivity is a missing-incidence
      row or one of the two disjoint matching patterns
      \eqref{eq:V99-disjoint-matching}; and
\item the full clean directed endpoint-acquisition ledger reduces to
      the finite list in
      \Cref{thm:V99-endpoint-reduction}.
\end{enumerate}
\end{theorem}

\begin{proof}
These are
\Cref{lem:V99-one-endpoint-sandwich,
lem:V99-directed-word-norm,
thm:V99-two-directed-one-paired,
thm:V99-all-directed-tree,
prop:V99-disconnected-classification,
thm:V99-endpoint-reduction}.
\end{proof}

\begin{programme}[V100 audit, matching repair, and series restart]
\label{prog:V99-V100}
\begin{enumerate}[leftmargin=3em]
\item Perform the compulsory read-only audit of the newest active SM
      and NS notes.  Record exact filenames, sizes, and hashes, and
      transfer only type-correct proof order.
\item Return to the exact V61 partition-state sum on the two
      disconnected matchings
      \eqref{eq:V99-disjoint-matching}.  Determine whether a retained
      same-side prefix response joins the components before any
      directed covariance is normed.
\item If no such response exists, apply the depleted V79 identity at
      a row in each component after deleting all matching banks.
      Either acquire a joining cell or exhibit the smallest exact
      depleted configuration.
\item Keep the one-resolvent payer in the full matching packet and use
      the V96 grouped suffix theorem for the suffix location.
\item Recheck V15, V83, V84, V86, and V91.
\item After completing the V100 chapter, freeze the resulting file as
      \texttt{Renewal\_Yang\_Mills\_Mass\_Gap\_Research\_Notes\_V100\_f2.tex}.
      Remove the superseded active V99 and begin a new contextual,
      proof-indexed active
      \texttt{Renewal\_Yang\_Mills\_Mass\_Gap\_Research\_Notes\_V1.tex}
      which summarizes the proved results and continues from the
      remaining bottleneck.
\end{enumerate}
\end{programme}

\begin{firewall}[Claim boundary after V99]
\label{fire:V99-boundary}
V99 closes the clean one-missing-endpoint directed branch and every
all-directed word whose oriented edge set contains a spanning tree.
It does not close the two disconnected directed matchings, a row with
no acquired incidence, nonclean representation states, the remaining
\(3+1\), \(2+1+1\), discrete first-connectivity, higher MTA, WUFC, or
CPM families.  Constructive cutoff removal, reflection positivity,
Osterwalder--Schrader reconstruction, and a uniform positive continuum
spectral gap remain open.  The full Yang--Mills mass gap remains open.
\end{firewall}

\paragraph{Parent-version record.}
The complete V98 mathematical body was retained before this chapter was
appended.  V98 had \(48426\) source records, \(1663720\) bytes, and
SHA256
\[
 \texttt{4c9f6c0cce8e1c302415da618386435f6f990de206a5e4530df3ff90e14aa20c}.
\]
No companion audit is due at V99; the next iteration is the compulsory
V100 audit and series transition.

% =============================================================================
% END APPENDED COMPACT VERSION V99 MATERIAL
% =============================================================================

% The V99 document terminator is intentionally disabled here:
% \end{document}

% =============================================================================
% BEGIN APPENDED COMPACT VERSION V100 MATERIAL
% =============================================================================

\clearpage
\chapter{The compulsory audit and the first matching-breaking
covariance}
\label{chap:V100}

\section{Companion audit and proof-order decision}
\label{sec:V100-audit}

This is the compulsory five-version audit following V95.  The files
were read only before the V100 matching calculation.

\begin{record}[Companion snapshots used by V100]
\label{rec:V100-companion-snapshots}
The newest active companion notes at the checkpoint were:
\begin{enumerate}
\item
\(\texttt{Renewal\_SM\_Einstein\_Continuum\_Research\_Notes\_V71.tex}\),
with \(29485\) lines, \(1053799\) bytes, and SHA--256
\[
 \texttt{f1a0def363c740426b562cf0bb145271b5566c8a7cc7cb95034b525f82b40435};
\]
\item
\(\texttt{Renewal\_NS\_Stability\_Research\_Notes\_V31.tex}\),
with \(23052\) lines, \(887537\) bytes, and SHA--256
\[
 \texttt{f1121b62238ad5491bb7cf03635ea9934942edf573ed8faaa9ee79e1d38a6696}.
\]
\end{enumerate}
The archived companion files were not treated as active notes, and no
companion file was changed.
\end{record}

\begin{assessment}[Type-correct audit transfer]
\label{ass:V100-audit-transfer}
SM V71 computes a light kink--antikink singular value and a uniform
heavy resolvent inside one complete operator.  Its two adjoint matrix
elements are descriptions of one mixing event, not two independent
small factors.  The Yang--Mills proof-order consequence is to retain
both matching components and the first component-breaking coordinate
inside one exact partition-state packet.  Estimating the two matching
components as independently completed four-row carriers would count
one joining event twice or erase it.

NS V31 proves that a flat probability mark removes stopping-time
multiplicity, but unmarking restores an exact critical first moment.
The Yang--Mills counterpart is that the sum over possible
matching-breaking coordinates must be kept as the genuine response
stock of one complete covariance.  Replacing it by a uniform local cap
and then counting possible first breakers would lose support
uniformity.

No domain-wall spectral theorem, Feshbach estimate, Navier--Stokes
occupation estimate, or signed-triad theorem is imported.  Only the
order of exact assembly before norming is transferred.
\end{assessment}

\begin{assessment}[Internal nonduplication decision]
\label{ass:V100-nonduplication}
V61 proves first-connectivity factorization for an arbitrary proper
partition, and V77 proves that summing first joins across a fixed
two-block partition gives one composite covariance.  V100 does not
repeat either theorem.  It applies them to the two specific matching
partitions isolated only in V99, proves the missing ordinary-norm
response for their complete breaker bank, and combines that response
with the directed tail denominators.
\end{assessment}

\section{Disconnected matching words vanish without a breaker}
\label{sec:V100-matching-zero}

Fix one of the two V99 matchings
\[
 \mu=A\mid B,
\]
where either
\[
 A=\{p,r\},\quad B=\{p^\circ,r^\circ\},
\tag{V100.1}\label{eq:V100-straight-matching}
\]
or
\[
 A=\{p,r^\circ\},\quad B=\{p^\circ,r\}.
\tag{V100.2}\label{eq:V100-crossed-matching}
\]
A local partition letter is called \(\mu\)-refining when its partition
is at most \(\mu\).

\begin{lemma}[No breaker means zero connected carrier]
\label{lem:V100-no-breaker-zero}
If every nonzero local partition letter in a coordinate word is
\(\mu\)-refining, then the join of the word is at most \(\mu\) and its
four-row connected cumulant contribution is zero.
\end{lemma}

\begin{proof}
The join of partitions all bounded by \(\mu\) is itself bounded by
\(\mu\).  Since \(\mu<\widehat1\), the linked-word criterion
\Cref{lem:V61-linked-word} excludes the word from the connected
carrier.
\end{proof}

\begin{corollary}[Every nonzero matching packet has a first breaker]
\label{cor:V100-first-breaker}
Choose any deterministic coordinate order which places every
coordinate known to refine \(\mu\) before the remaining coordinates.
Every nonzero connected word has a unique first coordinate at which
the cumulative join leaves \(\mu\), and at that coordinate the two
blocks \(A,B\) are joined.
\end{corollary}

\begin{proof}
The cumulative join is monotone.  By
\Cref{lem:V100-no-breaker-zero}, a connected word must leave \(\mu\).
There is therefore a unique first index at which it does.  Because
\(\mu\) has exactly two blocks, any partition whose join with \(\mu\)
is larger than \(\mu\) joins \(A\) to \(B\), and the join is
\(\widehat1\).
\end{proof}

\section{The complete matching-breaker covariance}
\label{sec:V100-breaker}

Let \(G\) be a finite complete bank of coordinates which can join
\(A\) to \(B\).  Let \(F_A^G,F_B^G\) be the corresponding ordered
products of row Wilson factors.  Define
\begin{equation}
 {\mathfrak J}_G^\mu
 :=
 \mathbb E_G[F_A^G\otimes F_B^G]
 -
 \mathbb E_GF_A^G\otimes\mathbb E'_GF_B^G.
\tag{V100.3}\label{eq:V100-breaker-covariance}
\end{equation}
This is the V77 composite covariance for the matching partition, not
for the original \(2+2\) cut.  Put
\begin{equation}
 H_G(A,B)
 :=
 \sum_{e\in G}
 \sum_{\substack{x\in A\\y\in B}}
 a_{x,e}a_{y,e}.
\tag{V100.4}\label{eq:V100-breaker-stock}
\end{equation}

\begin{theorem}[Ordinary-norm matching-breaker minimum]
\label{thm:V100-breaker-norm}
Let \({\cal N}_G^\mu\) be the complete positive physical envelope of
\eqref{eq:V100-breaker-covariance}.  Then
\begin{equation}
 \boxed{
 \|{\cal N}_G^\mu\|,
 \ \kappa_\otimes({\cal N}_G^\mu)
 \le
 C\min\{1,s_\alpha H_G(A,B)\}.}
\tag{V100.5}\label{eq:V100-breaker-unpaid}
\end{equation}
If the unique output numerator and resolvent are assigned to this
complete breaker bank, its paid envelope \({\cal P}_G^\mu\) satisfies
\begin{equation}
 \boxed{
 \|{\cal P}_G^\mu\|,
 \ \kappa_\otimes({\cal P}_G^\mu)
 \le
 C\min\left\{H_G(A,B),\frac1{s_\alpha}\right\}.}
\tag{V100.6}\label{eq:V100-breaker-paid}
\end{equation}
Both bounds are uniform in \(|G|\), representations, and physical
multiplicities.
\end{theorem}

\begin{proof}
On every coordinate \(e\), let \(T_e\) be the common \(A|B\) moment
and \(R_e^\mu\) its independently replicated endpoint.  The exact
product telescope gives
\[
 {\mathfrak J}_G^\mu
 =
 \bigotimes_{e\in G}T_e-\bigotimes_{e\in G}R_e^\mu
 =
 \sum_{t\in G}
 \left(\bigotimes_{e<t}T_e\right)
 \otimes(T_t-R_t^\mu)
 \otimes
 \left(\bigotimes_{e>t}R_e^\mu\right).
\]
Every spectator factor is a contraction.  The proof of
\Cref{lem:V65-local-crossing} applies to the arbitrary two-block
partition \(A|B\) and gives
\[
 \bigl\||T_t-R_t^\mu|_{\rm phys}\bigr\|
 \le
 Cs_\alpha
 \sum_{\substack{x\in A\\y\in B}}a_{x,t}a_{y,t}.
\]
Sum \(t\) only after this complete local covariance has been formed.
This proves the ordinary-norm response
\[
 \|{\cal N}_G^\mu\|\le Cs_\alpha H_G(A,B).
\]
The two endpoints of \eqref{eq:V100-breaker-covariance} are complete
contractions, so the same envelope has the universal norm bound.
Taking the minimum proves the norm part of
\eqref{eq:V100-breaker-unpaid}; capacity is no larger.

For the paid response, apply the V2 scalar allocation to the product
telescope.  When coordinate \(t\) pays, the local paid estimate in
\Cref{lem:V65-local-crossing} is
\[
 C\sum_{\substack{x\in A\\y\in B}}a_{x,t}a_{y,t},
\]
and every other coordinate remains a contraction.  Summing over
\(t\) gives the ordinary-norm response
\[
 \|{\cal P}_G^\mu\|\le CH_G(A,B)
\]
with no cardinality factor: the coordinate sum is the stock in
\eqref{eq:V100-breaker-stock}.  The universal paid cap
\(C/s_\alpha\) is the complete-cut argument of
\Cref{lem:V78-universal-paid-cap}, which is independent of the names
of the two blocks.  Take the smaller bound.  Again capacity is bounded
by norm.
\end{proof}

\begin{firewall}[The coordinatewise allocation is harmless only after
the response is retained]
\label{fire:V100-breaker-count}
The paid proof sums
\[
 \sum_{t\in G}\sum_{x\in A,y\in B}a_{x,t}a_{y,t}
 =
 H_G(A,B).
\]
It does not replace each coordinate by the uniform cap
\(C/s_\alpha\).  The latter would introduce \(|G|\), exactly the
first-moment loss warned against by the NS V31 audit.
\end{firewall}

\begin{corollary}[Quadratic breaker energies]
\label{cor:V100-breaker-energy}
For the exact signed breaker occurrence \(S_G^\mu\) inside these
envelopes,
\begin{align}
 {\cal E}_\otimes(S_G^\mu)
 &\le
 C\min\{1,s_\alpha H_G(A,B)\}^2,
\label{eq:V100-breaker-energy}\\
 {\cal E}_\otimes((S_G^\mu)^{\rm pay})
 &\le
 C\min\left\{H_G(A,B),\frac1{s_\alpha}\right\}^2.
\tag{V100.7}\label{eq:V100-breaker-paid-energy}
\end{align}
\end{corollary}

\begin{proof}
Both the norm and capacity factors in
\eqref{eq:V94-one-use-energy} are bounded by the corresponding line of
\Cref{thm:V100-breaker-norm}.
\end{proof}

\section{The complete matching packet}
\label{sec:V100-packet}

Let \(B_\mu\) denote the complete word made of all matching-refining
prefix carriers, the two matching directed covariances, replica
regroupings, and bounded separators.  It is assembled before the
breaker bank is resolved.

\begin{lemma}[Matching word norm]
\label{lem:V100-matching-word-norm}
For the four-row matching packets of V99,
\[
 \|B_\mu\|\le C.
\tag{V100.8}\label{eq:V100-matching-word-norm}
\]
If the sole payer lies in one matching covariance or a finite local
separator, then
\[
 \|B_\mu^{\rm pay}\|\le\frac C{s_\alpha}.
\tag{V100.9}\label{eq:V100-matching-word-paid}
\]
\end{lemma}

\begin{proof}
There are at most four assembled directed occurrences.  Apply
\Cref{lem:V99-directed-word-norm}.  Matching-refining complete moments
and suffix-free regroupings are contractions or have the finite V92
separator norm.
\end{proof}

\begin{theorem}[Support-uniform first-breaker packet]
\label{thm:V100-first-breaker-packet}
For the exact multiplier \(B_\mu S_G^\mu\), with all V94 output
dualizers absorbed into \(B_\mu\),
\begin{align}
 \Gamma_\otimes(1,K_{\mu,G})
 &\le
 C\min\{1,s_\alpha H_G(A,B)\},
\label{eq:V100-packet-unpaid}\\
 \Gamma_\otimes(w,K_{\mu,G}^{\rm pay})
 &\le
 C\min\left\{H_G(A,B),\frac1{s_\alpha}\right\}
\tag{V100.10}\label{eq:V100-packet-paid}
\end{align}
for every payer location in the matching word, the breaker, a local
separator, or the complete common suffix.
\end{theorem}

\begin{proof}
Use \Cref{thm:V97-three-factor-sandwich} with \(A=I\), the matching
word in the middle, and \(E=S_G^\mu\).  Equations
\eqref{eq:V100-matching-word-norm} and
\eqref{eq:V100-breaker-energy} give the unpaid line.

If the breaker pays, use
\eqref{eq:V100-breaker-paid-energy}.  If the matching word pays, use
\eqref{eq:V100-matching-word-paid} with the unpaid breaker energy:
\[
 \frac1{s_\alpha}\min\{1,s_\alpha H_G\}
 =
 \min\left\{\frac1{s_\alpha},H_G\right\}.
\]
A local separator is included in the same paid norm.  For a suffix
payer, use the grouped theorem
\Cref{thm:V96-support-uniform-suffix}; it produces the identical
calculation.  This proves \eqref{eq:V100-packet-paid}.
\end{proof}

\section{Fully incident matching closure}
\label{sec:V100-fully-incident}

For a matching component, write \(m(x)\) for the other vertex in that
component.  Suppose every row \(x\) has an acquired dominant directed
occurrence along its matching edge, with
\begin{equation}
 d_x
 :=
 \frac{H_{x\,m(x)}(G_x)}{\Xi_x}
 \ge\gamma_xa_*.
\tag{V100.11}\label{eq:V100-all-tail-responses}
\end{equation}
The banks \(G_x\) may coincide; their exact occurrences remain in the
one matching word.

\begin{theorem}[A first breaker joins every fully incident matching]
\label{thm:V100-matching-closure}
Under \eqref{eq:V100-all-tail-responses}, every unpaid and once-paid
nonzero matching packet is bounded by a finite literal marked-tree sum.
More precisely,
\begin{align}
 \Gamma_\otimes(1,K_{\mu,G})
 &\le
 C
 \sum_{\substack{x\in A\\y\in B}}
 \frac{d_xd_y}{\gamma_x\gamma_ya_*^2}
 H_{xy}(G),
\label{eq:V100-matching-tree-unpaid}\\
 \Gamma_\otimes(w,K_{\mu,G}^{\rm pay})
 &\le
 C
 \sum_{\substack{x\in A\\y\in B}}
 \frac{d_xd_y}{\gamma_x\gamma_ya_*^2}
 H_{xy}(G).
\tag{V100.12}\label{eq:V100-matching-tree-paid}
\end{align}
For each summand, the three edges
\[
 x\,m(x),\qquad y\,m(y),\qquad xy
\]
form a spanning tree, and the scalar product is its exact aggregate
response.
\end{theorem}

\begin{proof}
Because \(0<s_\alpha\le s_0\), the unpaid line of
\Cref{thm:V100-first-breaker-packet} is at most
\(C H_G(A,B)\).  The paid line is also at most \(CH_G(A,B)\).
Expand
\[
 H_G(A,B)
 =
 \sum_{x\in A,y\in B}H_{xy}(G).
\]
For each term, \eqref{eq:V100-all-tail-responses} gives
\[
 H_{xy}(G)
 \le
 \frac{d_xd_y}{\gamma_x\gamma_ya_*^2}H_{xy}(G).
\]

The matching edges \(x\,m(x)\) and \(y\,m(y)\) lie in the two
different components of \(\mu\); the breaker edge \(xy\) joins those
components.  The result is a tree.  Its aggregate bank product is
\[
 \frac{
 H_{x\,m(x)}(G_x)
 H_{y\,m(y)}(G_y)
 H_{xy}(G)}
 {\Xi_x\Xi_y}
 =
 d_xd_yH_{xy}(G),
\]
which is exactly the V92 marked-tree expansion: \(x,y\) are the two
degree-two vertices, while the other two row masses cancel.  Coincident
banks are permitted algebraically by
\Cref{prop:V98-same-bank-marking} and analytically by the single
assembled matching-word norm.  There are only four choices of
\((x,y)\), so the routing fibre is universal.
\end{proof}

\begin{corollary}[The two disconnected perfect matchings are closed]
\label{cor:V100-two-matchings-closed}
Every fully incident occurrence of either matching in
\eqref{eq:V99-disjoint-matching} closes at all payer locations.  If no
matching-breaking coordinate exists, its connected contribution is
zero by \Cref{lem:V100-no-breaker-zero}; otherwise it is controlled by
\Cref{thm:V100-matching-closure}.
\end{corollary}

\begin{proof}
Apply the two alternatives stated in the corollary.
\end{proof}

\section{The remaining depleted-tail acquisition}
\label{sec:V100-depleted-tail}

\begin{definition}[Breaker-visible tail set]
\label{def:V100-visible-tails}
For a breaker bank \(G\), let
\[
 V_G
 :=
 \left\{
 x\in[4]:
 \sum_{y\notin\mu(x)}H_{xy}(G)>0
 \right\},
\]
where \(\mu(x)\) is the matching component containing \(x\).
These are exactly the rows which occur as endpoints of at least one
positive term in \eqref{eq:V100-breaker-stock}.
\end{definition}

\begin{proposition}[Only missing breaker endpoints require acquisition]
\label{prop:V100-missing-tail-reduction}
Suppose a matching packet is not covered by
\Cref{thm:V100-matching-closure}.  Then, after discarding zero
\(H_{xy}(G)\) terms, every remaining unproved term has an endpoint
\(x\in V_G\) for which no fixed-fraction matching-tail response
\eqref{eq:V100-all-tail-responses} has been acquired.
Applying the depleted identity
\eqref{eq:V84-depleted-identity} at that row, after removing all
matching and breaker banks already present in the word, gives:
\begin{enumerate}[label=\textup{(\roman*)},leftmargin=4em]
\item a fresh \(P,U,A,F,I\), or directed cell;
\item overlap with one identified old bank; or
\item a nonclean representation state.
\end{enumerate}
A directed cell pointing along the matching supplies the missing tail;
a directed cell pointing outside the matching is itself a
component-breaking edge.  Thus the residual is a row-local depleted
acquisition, not a disconnected four-row analytic carrier.
\end{proposition}

\begin{proof}
If both endpoints \(x,y\) of a positive breaker term have the matching
tail responses, its proof is the corresponding summand of
\Cref{thm:V100-matching-closure}.  Therefore an unproved term has a
missing endpoint response.  Apply
\Cref{thm:V84-fresh-cell} at that endpoint with the used coordinate
union enlarged to include every bank already in the matching packet.
This gives the listed alternatives.

The V79 deterministic directed label records an actual partner.  If it
is \(m(x)\), \Cref{lem:V98-directed-edge-lower} supplies the missing
matching response.  Any other partner lies in the other component and
breaks \(\mu\).  No conclusion is asserted for the nonclean state.
\end{proof}

\begin{firewall}[The local acquisition has not yet terminated]
\label{fire:V100-acquisition-open}
Proposition~\ref{prop:V100-missing-tail-reduction} identifies the exact
row and breaker term which need a new certificate.  It does not prove
that repeated depleted acquisition must reach a \(P,U,A,F,I\) compiler
or a dominant matching tail before returning to an overlap.  The V83
strict-ascent lemma may be used only after each new cell has been
acquired, and it does not by itself terminate a sequence of
partition-changing overlaps.
\end{firewall}

\begin{record}[Late companion drift at the V100 freeze]\label{rec:V100-late-SM-drift}
Before freezing this cycle I repeated the companion-file check.  The active
Navier--Stokes note was still
\[
 \texttt{Renewal\_NS\_Terminal\_Source\_Concentration\_and\_Singular\_Thread\_Closure\_Programme\_V31.tex},
\]
with the fingerprint recorded in
Record~\ref{rec:V100-companion-snapshots}.  The active Standard-Model note had
advanced to
\[
 \texttt{Renewal\_SM\_Einstein\_Continuum\_Research\_Notes\_V72.tex},
\]
having \(29\,764\) lines, \(1\,064\,651\) bytes, and SHA--256
\[
 \texttt{9666b615196ea97adb90518c85b264f774659f9ab67e6c1c549ac0951690e2bd}.
\]
The new SM chapter proves a first-order Combes--Thomas estimate for an
off-diagonal Dirac resolvent, exponential local two-wall to one-wall resolvent
convergence at fixed nonzero imaginary spectral distance, the corresponding
local Gaussian-correlator convergence, and stability under a bounded common
self-adjoint potential.  It explicitly leaves determinant phase,
zero-frequency control, and removal of the tangential cutoff open.
\end{record}

\begin{assessment}[What transfers, and what does not]\label{ass:V100-late-transfer}
The transferable point is a proof-type restriction.  Local off-diagonal
suppression does not by itself control a global trace or determinant.  In the
same way, the local depleted endpoint alternatives used below do not by
themselves constitute global four-row closure.  The matching component, its
first breaker, both matching-tail responses, and the exact coordinate response
stock must remain in one assembled packet.  This is precisely the content of
Theorem~\ref{thm:V100-first-breaker-packet} and
Theorem~\ref{thm:V100-matching-closure}; the late audit therefore reinforces
their direction but imports no analytic estimate and changes no theorem.
\end{assessment}

\section{Regression tests and checkpoint status}
\label{sec:V100-status}

\begin{proposition}[V15, V83, V84, V86, and V91 regressions]
\label{prop:V100-regression-tests}
\begin{enumerate}[label=\textup{(\roman*)},leftmargin=4em]
\item V15 permits coincident coordinate marks only after the complete
      matching word is bounded.  No singlet dimension estimate is
      inserted.
\item In the V83 two-large/two-small patterns, a component-breaking
      edge is the actual \(xy\) term of
      \eqref{eq:V100-breaker-stock}; the maximal same-side partner is
      never relabelled.
\item The V84 spin-one counterexample is not contradicted:
      \eqref{eq:V100-breaker-unpaid} uses the complete covariance
      response \(s_\alpha H_G(A,B)\), not the false
      root-square normalization.
\item V86 keeps its binary prefix inside the complete matching word.
      Its adjoint descriptions are not multiplied as independent
      events.
\item V91 keeps the literal path \(\{12,13,34\}\); a breaker edge is
      selected from its actual coordinate bank only after the complete
      packet estimate.
\end{enumerate}
\end{proposition}

\begin{proof}
Items \textup{(i)--(ii)} use
\Cref{prop:V98-same-bank-marking,prop:V83-four-row-test}.
Item \textup{(iii)} follows from the exact two-block telescope in
\Cref{thm:V100-breaker-norm}.  Item \textup{(iv)} is the one-block
assembly of \Cref{thm:V86-binary-heat-capacity}.  Item \textup{(v)}
uses \Cref{thm:V92-V91-paid,prop:V94-marking-compatible}.
\end{proof}

\begin{theorem}[What V100 proves]
\label{thm:V100-status}
Relative to V1--V99, this checkpoint proves:
\begin{enumerate}[label=\textup{(V100.\arabic*)},leftmargin=7em]
\item the compulsory SM V71/NS V31 read-only audit with exact
      fingerprints and a type-correct common-block/first-moment
      transfer;
\item a matching-refining word contributes zero unless it has a unique
      first component-breaking coordinate;
\item the complete bank of matching breakers has the support-uniform
      ordinary-norm bounds
      \eqref{eq:V100-breaker-unpaid}--%
      \eqref{eq:V100-breaker-paid};
\item every payer location in the first-breaker packet obeys
      \eqref{eq:V100-packet-paid};
\item both fully incident disconnected matching patterns of V99 close
      into literal marked trees by
      \eqref{eq:V100-matching-tree-paid}; and
\item the residual is reduced to a depleted acquisition at an explicit
      endpoint of an explicit positive breaker term.
\end{enumerate}
\end{theorem}

\begin{proof}
These are
\Cref{rec:V100-companion-snapshots,
ass:V100-audit-transfer,
lem:V100-no-breaker-zero,
cor:V100-first-breaker,
thm:V100-breaker-norm,
thm:V100-first-breaker-packet,
thm:V100-matching-closure,
cor:V100-two-matchings-closed,
prop:V100-missing-tail-reduction}.
\end{proof}

\begin{firewall}[Claim boundary at the second freeze]
\label{fire:V100-boundary}
V100 closes the fully incident disconnected directed matchings and
reduces every remaining matching term to one row-local depleted
acquisition.  It does not prove termination of that acquisition,
nonclean representation states, the remaining \(3+1\), \(2+1+1\),
discrete first-connectivity, higher MTA, WUFC, or CPM families.
Constructive cutoff removal, reflection positivity,
Osterwalder--Schrader reconstruction, and a uniform positive continuum
spectral gap remain open.  The full Yang--Mills mass gap remains open.
\end{firewall}

\paragraph{Parent-version and freeze record.}
The complete V99 mathematical body was retained before this chapter was
appended.  V99 had \(48948\) source records, \(1684025\) bytes, and
SHA256
\[
 \texttt{5220c41b35c32caf9d02aead3476fef0c7ad9d50c903e72a188d2d38e89b5349}.
\]
This V100 file is the terminal member of the second compact active
sequence.  After validation it is frozen, without mathematical
changes, as
\texttt{Renewal\_Yang\_Mills\_Mass\_Gap\_Research\_Notes\_V100\_f2.tex}.
The next active file restarts at V1 with a contextual proof index and
continues from
\Cref{prop:V100-missing-tail-reduction,fire:V100-acquisition-open}.

% =============================================================================
% END APPENDED COMPACT VERSION V100 MATERIAL
% =============================================================================

\end{document}
